\documentclass[notitlepage]{book}

%\renewcommand{\baselinestretch}{2} for the copyeditor version

\usepackage{etex}
%\setlength\textwidth{342pt}
%\usepackage{microtype}
\usepackage{makeidx}
\makeindex

\usepackage[usenames,dvipsnames]{xcolor}
\usepackage[bookmarksopen,bookmarksdepth=2]{hyperref}
\hypersetup{colorlinks=true,citecolor=NavyBlue,linkcolor=BrickRed,urlcolor=Green} 
\usepackage{tikz-cd}

\RequirePackage{amsmath}
\RequirePackage{amssymb}
\RequirePackage{amsxtra}
\RequirePackage{amsfonts}
\RequirePackage{latexsym}
\RequirePackage{euscript}
\RequirePackage{amscd}
\RequirePackage{amsthm}
\RequirePackage{xypic}
\usepackage{stmaryrd}
\usepackage{mathtools}

\usepackage{mathrsfs}
\xyoption{all}

\usepackage{enumitem}
\DeclareMathOperator{\Max}{Max}
\def\acong{\buildrel \text{almost }  \over \cong}
\renewcommand{\mathbb}{\mathbf}
\newcommand{\Knr}{K^{\mathrm{nr}}}
\newcommand{\CEGS}{\operatorname{CEGS}}
\newcommand{\BT}{\mathrm{BT}}
\newcommand{\Lotimes}{{\otimes}^{\mathbb L} \, }
\newcommand{\Acirc}{A^\circ}
\newcommand{\Bcirc}{B^\circ}
\newcommand{\Ccirc}{C^\circ}

\newcommand{\Gmhat}{\widehat{\mathbb G}_m}
\newcommand{\cYha}{\cY_{h}^a}

\newcommand{\AKnrhat}{\widehat{\A}_K^{\nr}}
\newcommand{\AKnrhatA}{\widehat{\A}_{K,A}^{\nr}}
\newcommand{\AKnrhatcK}{\widehat{\A}_{K,\cK}^{\nr}}

\newcommand{\lambdau}{\underline{\lambda}}
\newcommand{\bE}{\mathbf{E}}
\newcommand{\tbE}{\widetilde{\bE}}
\newcommand{\Kbasic}{K^{\mathrm{basic}}}
\newcommand{\piflat}{\pi^\flat}
\newcommand{\piflatprime}{(\pi')^\flat}

\newcommand{\Sets}{\underline{{Sets}}}
\newcommand{\cyc}{\operatorname{cyc}}
\newcommand{\BK}{\operatorname{BK}}
\newcommand{\Kcyc}{K_{\cyc}}
\newcommand{\Atilde}{\widetilde{\A}}
\newcommand{\Aplus}{\A^+}
\newcommand{\AAA}{\A}
\newcommand{\AKDelta}{\A_K^\Delta}
\newcommand{\AKDeltaplus}{\A_K^{\Delta,+}}
\newcommand{\AKplus}{\A_K^{+}}
\newcommand{\Cone}{\operatorname{Cone}}
\newcommand{\Fib}{\operatorname{Fib}}
\newcommand{\Lift}{\operatorname{Lift}}
\newcommand{\pro}{{\operatorname{pro\, -}}}
\def\im{\mathop{\mathrm{im}}\nolimits}
\newcommand{\negligible}{\mathrm{small}}
\newcommand{\OEA}{\cO_{\mathcal{E},A}}
\newcommand{\OEApiflat}{\cO_{\mathcal{E},\pi^\flat,A}}
\newcommand{\hOEA}{\widehat{\cO}_{\mathcal{E},A}}
\newcommand{\OEAu}{\cO_{\mathcal{E},\infty,A}}
\newcommand{\OEAtilde}{\widetilde{\cO}_{\mathcal{E},A}}
\newcommand{\tOEA}{\widetilde{\cO}_{\mathcal{E},A}}
\newcommand{\OEAprime}{\cO_{\mathcal{E},A'}}
\newcommand{\OEB}{\cO_{\mathcal{E},B}}
\newcommand{\Ginfty}{G_{K_{\infty}}}
\newcommand{\stimage}{\cX}
\newcommand{\ctX}{\widetilde{\cX}}
\newcommand{\cXbar}{\overline{\cX}}
\newcommand{\an}{\operatorname{an}}
\newcommand{\can}{\operatorname{can}}
\newcommand{\Frac}{\operatorname{Frac}}
\newcommand{\xbar}{\overline{x}}
\newcommand{\gMhat}{\widehat{\gM}}
\newcommand{\gMt}{\gM^{\inf}}
\newcommand{\gNt}{\gN^{\inf}}
\newcommand{\gFt}{\gF^{\inf}}
\newcommand{\Rhat}{\widehat{\cR}}
\newcommand{\OCp}{\cO_{\C_p}}
\newcommand{\BS}{\operatorname{BS}}
\newcommand{\ybar}{\overline{y}}
\newcommand{\Ghat}{\widehat{G}}
\newcommand{\Kbar}{\overline{K}}
\def\Lbar{\overline{L}}
\newcommand{\sigmasm}{\sigma_{\text{sm}}}
\newcommand{\sigmaalg}{\sigma_{\text{alg}}}
\newcommand{\Sbig}{S^{\text{big}}}
\newcommand{\bigS}{\mathfrak{S}}
\newcommand{\Tbig}{T^{\text{big}}}
\newcommand{\Ssmall}{S^{\text{small}}}
\newcommand{\Tsmall}{T^{\text{small}}}
\newcommand{\Rloc}{R^{\text{loc}}}
\newcommand{\PS}{\operatorname{PS}}
\newcommand{\cusp}{\operatorname{cusp}}
\newcommand{\ghat}{\widehat{g}}
\newcommand{\Xbar}{\overline{X}}
% better wide vectors - stolen from http://tex.stackexchange.com/questions/44017/dot-notation-for-derivative-of-a-vector
% --- Macro \xvec
% \usepackage{tikz}
% \usetikzlibrary{matrix,arrows,decorations.pathmorphing}
% \usepackage{tikz-cd}
% \makeatletter
% \newlength\xvec@height%
% \newlength\xvec@depth%
% \newlength\xvec@width%
% \newcommand{\xvec}[2][]{%
%   \ifmmode%
%     \settoheight{\xvec@height}{$#2$}%
%     \settodepth{\xvec@depth}{$#2$}%
%     \settowidth{\xvec@width}{$#2$}%
%   \else%
%     \settoheight{\xvec@height}{#2}%
%     \settodepth{\xvec@depth}{#2}%
%     \settowidth{\xvec@width}{#2}%
%   \fi%
%   \def\xvec@arg{#1}%
%   \def\xvec@dd{:}%
%   \def\xvec@d{.}%
%   \raisebox{.2ex}{\raisebox{\xvec@height}{\rlap{%
%     \kern.05em%  (Because left edge of drawing is at .05em)
%     \begin{tikzpicture}[scale=1]
%     \pgfsetroundcap
%     \draw (.05em,0)--(\xvec@width-.05em,0);
%     \draw (\xvec@width-.05em,0)--(\xvec@width-.15em, .075em);
%     \draw (\xvec@width-.05em,0)--(\xvec@width-.15em,-.075em);
%     \ifx\xvec@arg\xvec@d%
%       \fill(\xvec@width*.45,.5ex) circle (.5pt);%
%     \else\ifx\xvec@arg\xvec@dd%
%       \fill(\xvec@width*.30,.5ex) circle (.5pt);%
%       \fill(\xvec@width*.65,.5ex) circle (.5pt);%
%     \fi\fi%
%     \end{tikzpicture}%
%   }}}%
%   #2%
% }
% \makeatother

% % --- Override \vec with an invocation of \xvec.
% \let\stdvec\vec
% \renewcommand{\vec}[1]{\xvec[]{#1}}
% % --- Define \dvec and \ddvec for dotted and double-dotted vectors.
% \newcommand{\dvec}[1]{\xvec[.]{#1}}
% \newcommand{\ddvec}[1]{\xvec[:]{#1}}




\usepackage{dsfont}

\newcommand{\Rinfty}{{R_\infty}}
\newcommand{\Rbarinfty}{{R'_\infty}}


\let\emptyset\varnothing
\newcommand{\Jminreg}{J_{\min}^{\text{reg}}}
\newcommand{\Jmaxreg}{J_{\max}^{\text{reg}}}
\newcommand{\Jmin}{J_{\min}}
\newcommand{\Jmax}{J_{\max}}
\newcommand{\Jreg}{J^{\text{reg}}}
\newcommand{\JH}{\operatorname{JH}}


\providecommand{\bysame}{\leavevmode\hbox to3em{\hrulefill}\thinspace}
\def\emparl{{\em (}}
\def\emparr{{\em )}}
\newcommand{\hA}{\widehat{\textA}}
\def\A{\mathbb A}
\def\Abar{\overline{A}}
\def\C{\mathbb C}
\def\F{\mathbb F}
\def\Khat{\widehat{K}}
\def\Fhat{\widehat{F}}
\def\Lhat{\widehat{L}}
\def\N{\mathbb{N}}
\def\Q{\mathbb{Q}}
\def\R{\mathbb{R}}
\def\T{\mathbb{T}}
\def\Ttilde{\widetilde{\T}}
\def\Z{\mathbb{Z}}
\def\Fbar{\overline{\F}}
\def\Falg{\overline{F}}
\def\Fur{F^{\mathrm{ur}}}
\def\kbar{\overline{k}}
\def\Qpur{\Q_p^{\mathrm{ur}}}
\def\Qbar{\overline{\Q}}
\def\Ubar{\overline{U}}

\def\Zbar{\overline{\Z}}
\def\Zhat{\widehat{\Z}}
\def\Pr{\mathbb{P}}
\def\L#1{{}^\mathrm{L} #1}
\def\m{\mathfrak m}
\newcommand{\st}{\mathrm{st}}
\newcommand{\cycle}{Z}
\def\Uni{{\mathrm U}}
\def\gl{\mathfrak{gl}}
%\def\O{\mathcal O}
\def\MF{\mathrm{MF}}
\def\Gr{\mathrm{Gr}}
\def\Mod{\mathrm{Mod}}
\def\Gmod{\Mod_{G}}
\def\smGmod{\Mod^{\mathrm{sm}}_G}
\def\admGmod{\Mod^{\mathrm{adm}}_G}
\def\locadmGmod{\Mod^{\mathrm{l\, adm}}_G}
\def\fgaugGmod{\Mod^{\mathrm{fg\, aug}}_G}
\def\proaugGmod{\Mod^{\mathrm{pro\,aug}}_G}
\def\freesmGmod{\Mod^{\mathrm{fr\, sm}}_G}
\def\freeadmGmod{\Mod^{\mathrm{fr\, adm}}_G}
\def\freefgaugGmod{\Mod^{\mathrm{pro-fr\,fg.\, aug}}_G}
\def\freeproaugGmod{\Mod^{\mathrm{pro-fr\,aug}}_G}
\newcommand{\su}[1][j]{^{(#1)}}  % superscript (it is j by default, please leave it that way...)
\def\pcontGmod{\Mod^{\unif-\mathrm{cont}}_G}
\def\contGmod{\Mod^{\mathrm{cont}}_G}
\def\orthcontGmod{\Mod^{\mathrm{orth\,cont}}_G}
\def\orthadmGmod{\Mod^{\mathrm{orth\, adm}}_G}
\def\smHmod{\Mod^{\mathrm{sm}}_G}
\def\pcontHmod{\Mod^{\unif-\mathrm{cont}}_G}
\def\padmGmod{\Mod^{\mathrm{\unif-adm}}_G}
\def\BanG{\Mod^{\mathrm{u}}_G}
\def\admBanG{\Mod^{\mathrm{u\,adm}}_G}
\def\augGmod{\Mod^{\mathrm{aug}}_G}
\newcommand{\wh}{\widehat}
\def\Tmod{\Mod_{T}}
\def\smTmod{\Mod^{\mathrm{sm}}_T}
\def\admTmod{\Mod^{\mathrm{adm}}_T}
\def\locadmTmod{\Mod^{\mathrm{l\, adm}}_T}
\def\fgaugTmod{\Mod^{\mathrm{fg\, aug}}_T}
\def\proaugTmod{\Mod^{\mathrm{pro\,aug}}_T}
\def\freesmTmod{\Mod^{\mathrm{fr\, sm}}_T}
\def\freeadmTmod{\Mod^{\mathrm{fr\, adm}}_T}
\def\freefgaugTmod{\Mod^{\mathrm{pro-fr\,fg\, aug}}_T}
\def\freeproaugTmod{\Mod^{\mathrm{pro-fr\,aug}}_T}
\def\pcontTmod{\Mod^{\unif-\mathrm{cont}}_T}
\def\contTmod{\Mod^{\mathrm{cont}}_T}
\def\orthcontTmod{\Mod^{\mathrm{orth\,cont}}_T}
\def\orthadmTmod{\Mod^{\mathrm{orth\, adm}}_T}
\def\padmTmod{\Mod^{\mathrm{\unif-adm}}_T}
\def\BanT{\Mod^{\mathrm{u}}_T}
\def\admBanT{\Mod^{\mathrm{u\,adm}}_T}
\def\augTmod{\Mod^{\mathrm{aug}}_T}
\def\Bmod{\Mod_{B}}
\def\smBmod{\Mod^{\mathrm{sm}}_B}
\def\admBmod{\Mod^{\mathrm{adm}}_B}
\def\locadmBmod{\Mod^{\mathrm{l\, adm}}_B}
\def\fgaugBmod{\Mod^{\mathrm{fg.\, aug}}_B}
\def\proaugBmod{\Mod^{\mathrm{pro\,aug.}}_B}
\def\freesmBmod{\Mod^{\mathrm{fr\, sm}}_B}
\def\freeadmBmod{\Mod^{\mathrm{fr\, adm}}_B}
\def\freefgaugBmod{\Mod^{\mathrm{pro-fr\,fg\, aug}}_B}
\def\freeproaugBmod{\Mod^{\mathrm{pro-fr\,aug}}_B}
\def\pcontBmod{\Mod^{\unif-\mathrm{cont}}_B}
\def\contBmod{\Mod^{\mathrm{cont}}_B}
\def\orthcontBmod{\Mod^{\mathrm{orth\,cont}}_B}
\def\orthadmBmod{\Mod^{\mathrm{orth\, adm}}_B}
\def\padmBmod{\Mod^{\mathrm{\unif-adm}}_B}
\def\BanB{\Mod^{\mathrm{u}}_B}
\def\admBanB{\Mod^{\mathrm{u\,adm}}_B}
\def\augBmod{\Mod^{\mathrm{aug}}_B}
\def\Vcheck{V^{\vee}}
\def\Bbar{\overline{B}}
\def\Vbar{\overline{V}}
\def\Uni{{\mathrm U}}
\def\chibar{\overline{\chi}}
\def\pibar{\overline{\pi}}
\def\pihat{\widehat{\pi}}
\def\pitilde{\widetilde{\pi}}
\def\pilocalg{\widetilde{\pi}}
\def\pimod{\pi^{\mathrm{m}}}
\def\pimodp{\pi^{\mathrm{m},\, p}}
\def\pimodprime{\pi^{\mathrm{m} \prime}}
\def\LL{\pi}
\def\LLS{\pi_{\Sigma_0}}
\def\LLell{\pi_{\ell}}
\def\LLbar{\overline{\pi}}
\def\LLbarS{\overline{\pi}_{\Sigma_0}}
\def\LLcheck{\pitilde}
\def\LLcheckS{\pitilde_{\Sigma_0}}
\def\LLcheckell{\pitilde_{\ell}}
\def\LLbarcheck{\widetilde{\LLbar}}
\def\barepsilon{\overline{\varepsilon}}
%\def\varepsilonbar{\overline{\varepsilon}}
\def\unif{\varpi}
\def\pr{\mathrm{pr}}
\def\mult{\mathrm{m}}
\def\multgl{\mathrm{m}^{\mathrm{gl}}}
\def\Wgl{W^{\mathrm{gl}}}
\def\id{\mathrm{id}}
\def\tf{\mathrm{tf}}
\def\ctf{\mathrm{ctf}}
\def\fin{\mathrm{fin}}
\def\et{\mathrm{\acute{e}t}}
%\def\para{\mathrm{par}}
\def\tor{\mathrm{tor}}
\def\elalg{\mathrm{exp.l.alg}}
\def\alg{\mathrm{alg}}
\def\lalg{\mathrm{l.alg}}
\def\ladm{\mathrm{l.adm}}
\def\lan{\mathrm{an}}
\def\lp{\mathrm{lp}}
\def\ng{\mathrm{ng}}
\def\red{\mathrm{red}}
\def\rig{\mathrm{rig}}
\def\ab{\mathrm{ab}}
\def\nr{\mathrm{ur}}
\def\fl{\mathrm{fl}}
\def\sm{\mathrm{sm}}
\def\con{\circ}
\def\cont{\mathrm{cont}}
\def\ss{\mathrm{ss}}
\def\sph{\mathcal H^{\mathrm{sph}}}
\def\SL{\mathrm{SL}}
\def\sat{\mathrm{sat}}
\def\ur{\mathrm{ur}}
\def\fs{\mathrm{fs}}
\def\nil{\mathop{\mathrm{nil}}\nolimits}
\def\ord{\mathrm{ord}}
\def\nord{\mathrm{n.ord}}
\def\new{\mathrm{new}}
\def\old{\mathrm{old}}
\def\Ord{\mathrm{Ord}}
\def\GL{\operatorname{GL}}
\def\Ad{\mathrm{Ad}}
\def\Gal{\mathrm{Gal}}
\def\Pic{\mathrm{Pic}}
\def\Aut{\mathrm{Aut}}
\def\Bun{\mathrm{Bun}}
\def\Sym{\mathrm{Sym}}
\def\Ext{\mathrm{Ext}}
\def\Exp{\mathrm{Exp}}
\def\End{\mathrm{End}}
\def\Def{\mathrm{Def}}
\def\Defcrys{\mathrm{Def}^{\crys}}
\def\Defstar{\mathrm{Def}^*}
\def\Art{\mathop{\mathrm{Art}}\nolimits}
\def\Comp{\mathop{\mathrm{Comp}}\nolimits}
\def\Hom{\mathop{\mathrm{Hom}}\nolimits}
\def\RHom{\mathop{\mathrm{RHom}}\nolimits}
\def\Tan{\mathop{\mathrm{Tan}}\nolimits}
\def\Tor{\mathop{\mathrm{Tor}}\nolimits}
\def\Spec{\mathop{\mathrm{Spec}}\nolimits}
\def\Spf{\mathop{\mathrm{Spf}}\nolimits}
\def\Frob{\mathop{\mathrm{Frob}}\nolimits}
\def\Supp{\mathop{\mathrm{Supp}}\nolimits}
\def\Trace{\mathop{\mathrm{Trace}}\nolimits}
\def\trace{\mathop{\mathrm{trace}}\nolimits}
\def\Det{\mathop{\mathrm{Det}}\nolimits}
\def\Ind{\mathop{\mathrm{Ind}}\nolimits}
\def\cInd{\mathop{c\mathrm{-Ind}}\nolimits}
\def\Fil{\mathop{\mathrm{Fil}}\nolimits}
\def\Ref{\mathop{\mathrm{Ref}}\nolimits}
\def\Con{\mathop{\mathcal C}\nolimits}
\def\soc{\mathop{\mathrm{soc}}\nolimits}
\def\Eis{\mathop{{\mathcal E}\!\text{\em is}}\nolimits}
\def\rhobar{\overline{\rho}}
\def\rhomod{\rho^{\mathrm{m}}}
\def\rhomodprime{\rho^{\mathrm{m} \prime}}
\def\rhouniv{\rho^{\mathrm{u}}}
\def\rhoord{\rho^{\mathrm{ord}}}
\def\rhonord{\rho^{\mathrm{n.ord}}}
\def\abs{\lvert~\;~\rvert}
\def\cotimes{\operatorname{\widehat{\otimes}}}
\def\cboxtimes{\widehat{\boxtimes}}
\def\qotimes{{\buildrel \curlywedge \over \otimes}}
\def\crys{\mathrm{crys}}
\def\pcrys{\mathrm{pcrys}}
\def\dR{\mathrm{dR}}
\def\pst{\mathrm{pst}}
%\def\st{\mathrm{st}}
\def\BE{E\otimes_{\Q_p}B_{\crys}}
\def\BDRE{E\otimes_{\Q_p}B_{\dR}}
\def\W{\mathrm{W}}
\def\WD{\mathrm{WD}}
\def\St{\mathrm{St}}
\def\Sthat{\widehat{\mathrm{St}}}
\def\kernel{\mathop{\mathrm{ker}}\nolimits}
\def\image{\mathop{\mathrm{im}}\nolimits}
\def\cokernel{\mathop{\mathrm{coker}}\nolimits}
\def\m{\mathfrak{m}}
\def\ilim#1{\displaystyle \lim_{\longrightarrow \atop #1}}
\def\plim#1{\displaystyle \lim_{\longleftarrow \atop #1}}
\def\iso{\buildrel \sim \over \longrightarrow}
\def\Pbar{\overline{P}}
\def\GLQ{\GL_2(\Q)}
\def\GLQp{\GL_2(\Q_p)}
\def\SLQp{\SL_2(\Q_p)}
\def\PQ{\mathrm{P}(\Q)}
\def\PQp{\mathrm{P}(\Q_p)}
\def\Pbarrm{\overline{\mathrm{P}}}
\def\PbarQp{\Pbarrm(\Q_p)}
\def\NQ{\mathrm{N}(\Q)}
\def\NQp{\mathrm{N}(\Q_p)}
\def\Nbar{\overline{\mathrm{N}}}
\def\NbarQ{\overline{\mathrm{N}}(\Q)}
\def\NbarQp{\overline{\mathrm{N}}(\Q_p)}
\def\Trm{\mathrm{T}}
\def\TQ{\Trm(\Q)}
\def\TQp{\Trm(\Q_p)}
\def\glQp{\gl_2(\Q_p)}
\def\bQp{\mathfrak b(\Q_p)}
\def\zQp{\mathfrak z(\Q_p)}
\def\That{\widehat{T}}
\def\Thatzero{\widehat{T}_0}
\def\triv{\mathds{1}}
\def\ev{\mathrm{ev}}
\def\OT{\mathop{\mathrm{OT}}\nolimits}
\def\rmcap{\mathop{\mathrm{cap}}\nolimits}
\def\gauge{\cG}
\def\cosocgauge{\cG_{\textrm{cosoc}}}
\def\cosoc{\mathop{\mathrm{cosoc}}\nolimits}

\newcommand{\modules}{\cC^{\Ghat}}
\newcommand{\spaces}{\cR^{\Ghat}}
\newcommand{\spacesprime}{(\cR^{\Ghat})'}

\swapnumbers
%%%% TG macros
%% arrows
\newcommand{\lonto}{\longrightarrow \hspace{-1.44em} \longrightarrow}
\newcommand{\onto}{\twoheadrightarrow}
\newcommand{\ra}{\rightarrow}
\newcommand{\lra}{\longrightarrow}
\newcommand{\la}{\leftarrow}
\newcommand{\lla}{\longleftarrow}
\newcommand{\into}{\hookrightarrow}
\newcommand{\linto}{\hookleftarrow}
\newcommand{\mt}{\mapsto}
\newcommand{\lmt}{\longmapsto}
\newcommand{\ua}{\uparrow}
\newcommand{\da}{\downarrow}
%\newcommand{\iso}{\stackrel{\sim}{\ra}}
\newcommand{\liso}{\stackrel{\sim}{\lra}}
\newcommand{\To}{\longrightarrow}
\newcommand{\isoto}{\stackrel{\sim}{\To}}
\newcommand{\isofrom}{\stackrel{\sim}{\longleftarrow}}
%% Fontaine's rings and functors
\newcommand{\FBrModdd}[1][r]{\text{$\F$-$\operatorname{BrMod}_{\mathrm {dd}}^{#1}$}}
\newcommand{\OEModdd}[1][r]{\text{$\cO_E$-$\Mod_{\mathrm {dd}}^{#1}$}}
\newcommand{\textD}{\mathrm{D}}
\newcommand{\textT}{\mathrm{T}}
\newcommand{\textV}{\mathrm{V}}
\newcommand{\textA}{\mathrm{A}}
\newcommand{\hAst}{\widehat{\textA}_{\st}}
\newcommand{\BrMod}{\operatorname{BrMod}}
%\newcommand{\dR}{\mathrm{dR}}
\newcommand{\textB}{\mathrm{B}}
\newcommand{\Bcris}{\textB_{\cris}}
\newcommand{\Acris}{\textA_{\cris}}
 \newcommand{\BdR}{\textB_{\dR}}
\newcommand{\BHT}{\textB_{\HT}}
\newcommand{\Bst}{\textB_{\st}}
\newcommand{\Dcris}{\textD_{\cris}}
\newcommand{\DdR}{\textD_{\dR}}
\newcommand{\DHT}{\textD_{\HT}}
\newcommand{\Dst}{\textD_{\st}}
\newcommand{\Vst}{\textV_{\st}}
\newcommand{\Tst}{\textT_{\st}}
%%
\newcommand{\pfbegin}{{{\em Proof:}\;}}
\newcommand{\pfend}{$\Box$ \medskip}
\newcommand{\mat}[4]{\left( \begin{array}{cc} {#1} & {#2} \\ {#3} & {#4}
\end{array} \right)}
\newcommand{\galois}[1]{\Gal ( \overline{ {#1} }/ {#1})}
\newcommand{\smat}[4]{{\mbox{\scriptsize $\mat{{#1}}{{#2}}{{#3}}{{#4}}$}}}
\newlength{\ownl}
\newcommand{\stacktwo}[2]{\settowidth{\ownl}{\scriptsize{${#1}$}}
\parbox[t]{1\ownl}{{\scriptsize{${#1}$}} \\[-.9ex] {\scriptsize{${#2}$}}}}
\newcommand{\stackthree}[3]{\settowidth{\ownl}{\scriptsize{${#1}$}}
\parbox[t]{1\ownl}{{\scriptsize{${#1}$}} \\[-.9ex] {\scriptsize{${#2}$}}
\\[-.9ex] {\scriptsize{${#3}$}}}}
\newcommand{\<}{\langle}
\newcommand{\norm}{{\mbox{\bf N}}}
\newcommand{\ndiv}{{\mbox{$\not| $}}}
\newcommand{\hotimes}{\widehat{\otimes}}
%operators
\newcommand{\obs}{{\operatorname{obs}}}
\newcommand{\ad}{{\operatorname{ad}\,}}
\newcommand{\Aff}{{\operatorname{Aff}}}
\newcommand{\Ann}{{\operatorname{Ann}\,}}
%\newcommand{\Art}{{\operatorname{Art}\,}}
%\newcommand{\Aut}{{\operatorname{Aut}\,}}
\newcommand{\BC}{{\operatorname{BC}\,}}
\newcommand{\Br}{{\operatorname{Br}\,}}
\newcommand{\chara}{{\operatorname{char}\,}}
\newcommand{\cind}{{\operatorname{c-Ind}\,}}
\newcommand{\Cl}{{\operatorname{Cl}}}
\newcommand{\coker}{{\operatorname{coker}\,}}
\newcommand{\conju}{{\operatorname{conj}}}
\newcommand{\Cusp}{{\operatorname{Cusp}\,}}
\newcommand{\DC}{{\operatorname{DC}}}
%\newcommand{\Def}{{\operatorname{Def}}}
\newcommand{\diag}{{\operatorname{diag}}}
%\newcommand{\End}{{\operatorname{End}\,}}
\newcommand{\ep}{{\operatorname{ep}\,}}
\newcommand{\es}{\emptyset}
%\newcommand{\Ext}{{\operatorname{Ext}\,}}
%\newcommand{\Fil}{{\operatorname{Fil}\,}}
\newcommand{\Fitt}{{\operatorname{Fitt}\,}}
\newcommand{\Fix}{{\operatorname{Fix}}}
%\newcommand{\Frob}{{\operatorname{Frob}}}
\newcommand{\Fr}{{\operatorname{Fr}}}
%\newcommand{\Gal}{{\operatorname{Gal}\,}}
\newcommand{\gen}{{\operatorname{gen}}}
\newcommand{\gr}{{\operatorname{gr}\,}}
\newcommand{\Groth}{{\operatorname{Groth}\,}}
\newcommand{\hdg}{{\operatorname{hdg}}}
%\newcommand{\Hom}{{\operatorname{Hom}\,}}
\newcommand{\CHom}{{\CH om\,}}
\newcommand{\ICP}{{\operatorname{IPC}}}
\newcommand{\Id}{{\operatorname{Id}}}
\renewcommand{\Im}{{\operatorname{Im}\,}}
%\newcommand{\Ind}{{\operatorname{Ind}\,}}
\newcommand{\ind}{{\operatorname{ind}}}
\newcommand{\inv}{{\operatorname{inv}\,}}
\newcommand{\Irr}{{\operatorname{Irr}}}
\newcommand{\Isom}{{\operatorname{Isom}\,}}
\newcommand{\Iw}{{\operatorname{Iw}}}
\newcommand{\JL}{{\operatorname{JL}\,}}
\newcommand{\Lie}{{\operatorname{Lie}\,}}
\newcommand{\nind}{{\operatorname{n-Ind}\,}}
\newcommand{\nred}{{\operatorname{n-red}}}
\newcommand{\Out}{{\operatorname{Out}\,}}
\newcommand{\PC}{{\operatorname{PC}}}
\newcommand{\pic}{{\operatorname{PIC}}}
%\newcommand{\pr}{{\operatorname{pr}\,}}
\renewcommand{\Re}{{\operatorname{Re}\,}}
\newcommand{\rec}{{\operatorname{rec}}}
\newcommand{\redu}{{\operatorname{red}}}
\newcommand{\Red}{{\operatorname{Red}}}
\newcommand{\Rep}{{\operatorname{Rep}}}
\newcommand{\REP}{{\operatorname{GG}}}
\newcommand{\GrG}{{\operatorname{Rep}}}
\newcommand{\res}{{\operatorname{res}}}
\newcommand{\Res}{{\operatorname{Res}}}
\newcommand{\rk}{{\operatorname{rk}\,}}
\newcommand{\RS}{{\operatorname{RS}}}
\newcommand{\DR}{{\operatorname{DR}}}
%\newcommand{\WD}{{\operatorname{WD}}}
\newcommand{\Sd}{{\operatorname{Sd}}}
\newcommand{\sgn}{{\operatorname{sgn}\,}}
%\newcommand{\Spec}{{\operatorname{Spec}\,}}
%\newcommand{\Spf}{{\operatorname{Spf}\,}}
\newcommand{\Spp}{{\operatorname{Sp}\,}}
\newcommand{\Symm}{{\operatorname{Symm}\,}}
\newcommand{\tr}{{\operatorname{tr}\,}}
\newcommand{\Twist}{{\operatorname{Twist}\,}}
\newcommand{\val}{{\operatorname{val}\,}}
\newcommand{\vol}{{\operatorname{vol}\,}}
\newcommand{\WDRep}{{\operatorname{WDRep}}}
\newcommand{\wt}[1]{\widetilde{#1}}
%% Langlands groups and algebraic groups
%\newcommand{\GL}{\operatorname{GL}}
\newcommand{\GSp}{\operatorname{GSp}}
\newcommand{\Gdual}{{\hat{G}}}
\newcommand{\Gm}{{\mathbb{G}_m}}
\newcommand{\LG}{{{}^{L}G}}
\newcommand{\PGL}{\operatorname{PGL}}
\newcommand{\PSL}{\operatorname{PSL}}
%\newcommand{\SL}{\operatorname{SL}}
\newcommand{\Sp}{\operatorname{Sp}}
\newcommand{\Tdual}{{\hat{T}}}
%% super/subscripts
\newcommand{\APHT}{{\operatorname{APHT}}}
\newcommand{\AV}{{\operatorname{AV}}}
\newcommand{\FP}{{\operatorname{FP}}}
\newcommand{\Fsemis}{{\operatorname{F-ss}}}
\newcommand{\GB}{{\operatorname{GB}}}
\newcommand{\LA}{{\operatorname{LA}}}
\newcommand{\PHT}{{\operatorname{PHT}}}
\newcommand{\SC}{{\operatorname{SC}}}
\newcommand{\Tam}{{\operatorname{Tam}}}
%\newcommand{\ab}{{\operatorname{ab}}}
\newcommand{\ac}{{\operatorname{ac}}}
%\newcommand{\alg}{{\operatorname{alg}}}
\newcommand{\cris}{{\operatorname{cris}}}
\newcommand{\crord}{{\operatorname{cr-ord}}}
\newcommand{\cts}{{\operatorname{cts}}}
\newcommand{\der}{{\operatorname{der}}}
\newcommand{\disc}{{\operatorname{disc}}}
\newcommand{\dr}{{\operatorname{DR}}}
%\newcommand{\et}{{\operatorname{et}}}
\newcommand{\jl}{{\operatorname{JL}}}
\newcommand{\loc}{{\operatorname{loc}}}
%\newcommand{\new}{{\operatorname{new}}}
\newcommand{\noeth}{{\operatorname{noeth}}}
%\newcommand{\nr}{{\operatorname{nr}}}
%\newcommand{\old}{{\operatorname{old}}}
\newcommand{\opp}{{\operatorname{op}}}
\newcommand{\op}{{\operatorname{op}}}
\newcommand{\pnr}{{\operatorname{p-nr}}}
%\newcommand{\red}{{\operatorname{red}}}
\newcommand{\reg}{{\operatorname{reg}}}
%\newcommand{\rig}{{\operatorname{rig}}}
\newcommand{\semis}{{\operatorname{ss}}}
\newcommand{\ssord}{{\operatorname{ss-ord}}}
%\newcommand{\st}{{\operatorname{st}}}
%\newcommand{\tf}{{\operatorname{tf}}}
\newcommand{\topo}{{\operatorname{top}}}
%\newcommand{\tor}{{\operatorname{tor}}}
%\newcommand{\triv}{{\operatorname{triv}}}
\newcommand{\univ}{{\operatorname{univ}}}
%% letters
%\newcommand{\A}{{\mathbb{A}}}
\newcommand{\B}{{\mathbb{B}}}
%\newcommand{\C}{{\mathbb{C}}}
\newcommand{\D}{{\mathbb{D}}}
\newcommand{\E}{{\mathbb{E}}}
%\newcommand{\F}{{\mathbb{F}}}
\newcommand{\G}{{\mathbb{G}}}
\newcommand{\HH}{{\mathbb{H}}}
\newcommand{\I}{{\mathbb{I}}}
\newcommand{\J}{{\mathbb{J}}}
\newcommand{\K}{{\mathbb{K}}}
%\newcommand{\LL}{{\mathbb{L}}}
\newcommand{\M}{{\mathcal{M}}}
%\newcommand{\N}{{\mathbb{N}}}
\newcommand{\OO}{{\mathbb{O}}}
\newcommand{\PP}{{\mathbb{P}}}
%\newcommand{\Q}{{\mathbb{Q}}}
%\newcommand{\R}{{\mathbb{R}}}
\newcommand{\mathbbS}{{\mathbb{S}}}
%\newcommand{\T}{{\mathbb{T}}}
\newcommand{\U}{{\mathbb{U}}}
\newcommand{\V}{{\mathbb{V}}}
%\newcommand{\W}{{\mathbb{W}}}
\newcommand{\X}{{\mathbb{X}}}
\newcommand{\Y}{{\mathbb{Y}}}
%\newcommand{\Z}{{\mathbb{Z}}}
\newcommand{\CA}{{\mathcal{A}}}
\newcommand{\CB}{{\mathcal{B}}}
\newcommand{\CC}{{\mathcal{C}}}
\newcommand{\calD}{{\mathcal{D}}}
\newcommand{\CE}{{\mathcal{E}}}
\newcommand{\CF}{{\mathcal{F}}}
\newcommand{\CG}{{\mathcal{G}}}
\newcommand{\CH}{{\mathcal{H}}}
\newcommand{\CI}{{\mathcal{I}}}
\newcommand{\CJ}{{\mathcal{J}}}
\newcommand{\CK}{{\mathcal{K}}}
\newcommand{\CL}{{\mathcal{L}}}
\newcommand{\CM}{{\mathcal{M}}}
\newcommand{\CN}{{\mathcal{N}}}
\newcommand{\CO}{{\mathcal{O}}}
\newcommand{\CP}{{\mathcal{P}}}
\newcommand{\CQ}{{\mathcal{Q}}}
\newcommand{\CR}{{\mathcal{R}}}
\newcommand{\CS}{{\mathcal{S}}}
\newcommand{\CT}{{\mathcal{T}}}
\newcommand{\CU}{{\mathcal{U}}}
\newcommand{\CV}{{\mathcal{V}}}
\newcommand{\CW}{{\mathcal{W}}}
\newcommand{\CX}{{\mathcal{X}}}
\newcommand{\CY}{{\mathcal{Y}}}
\newcommand{\CZ}{{\mathcal{Z}}}
\newcommand{\cA}{\mathcal{A}}
\newcommand{\cB}{\mathcal{B}}
\newcommand{\cC}{\mathcal{C}}
\renewcommand{\cD}{\mathcal{D}}
\newcommand{\cE}{\mathcal{E}}
\newcommand{\cF}{\mathcal{F}}
\newcommand{\cG}{\mathcal{G}}
\renewcommand{\cH}{\mathcal{H}}
\newcommand{\calH}{\mathcal{H}}
\newcommand{\cI}{\mathcal{I}}
\newcommand{\cJ}{\mathcal{J}}
\newcommand{\cK}{\mathcal{K}}
\newcommand{\caL}{\mathcal{L}}
\newcommand{\cM}{\mathcal{M}}
\newcommand{\cN}{\mathcal{N}}
\newcommand{\cO}{\mathcal{O}}
\renewcommand{\O}{\cO}
\newcommand{\cP}{\mathcal{P}}
\newcommand{\cQ}{\mathcal{Q}}
\renewcommand{\cR}{\mathcal{R}}
\newcommand{\cS}{\mathcal{S}}
\newcommand{\cT}{\mathcal{T}}
\newcommand{\cU}{\mathcal{U}}
\newcommand{\cV}{\mathcal{V}}
\newcommand{\cW}{\mathcal{W}}
\newcommand{\cWW}{\mathcal{W}}
\newcommand{\cX}{\mathcal{X}}
\newcommand{\cY}{\mathcal{Y}}
\newcommand{\cZ}{\mathcal{Z}}
\newcommand{\gA}{{\mathfrak{A}}}
\newcommand{\gB}{{\mathfrak{B}}}
\newcommand{\gC}{{\mathfrak{C}}}
\newcommand{\gD}{{\mathfrak{D}}}
\newcommand{\gE}{{\mathfrak{E}}}
\newcommand{\gF}{{\mathfrak{F}}}
\newcommand{\gG}{{\mathfrak{G}}}
\newcommand{\gH}{{\mathfrak{H}}}
\newcommand{\gI}{{\mathfrak{I}}}
\newcommand{\gJ}{{\mathfrak{J}}}
\newcommand{\gK}{{\mathfrak{K}}}
\newcommand{\gL}{{\mathfrak{L}}}
\newcommand{\gM}{{\mathfrak{M}}}
\newcommand{\gN}{{\mathfrak{N}}}
\newcommand{\gO}{{\mathfrak{O}}}
\newcommand{\gP}{{\mathfrak{P}}}
\newcommand{\gQ}{{\mathfrak{Q}}}
\newcommand{\gR}{{\mathfrak{R}}}
\newcommand{\gS}{{\mathfrak{S}}}
\newcommand{\gT}{{\mathfrak{T}}}
\newcommand{\gU}{{\mathfrak{U}}}
\newcommand{\gV}{{\mathfrak{V}}}
\newcommand{\gW}{{\mathfrak{W}}}
\newcommand{\gX}{{\mathfrak{X}}}
\newcommand{\gY}{{\mathfrak{Y}}}
\newcommand{\gZ}{{\mathfrak{Z}}}
\newcommand{\ga}{{\mathfrak{a}}}
\newcommand{\gb}{{\mathfrak{b}}}
\newcommand{\gc}{{\mathfrak{c}}}
\newcommand{\gd}{{\mathfrak{d}}}
\newcommand{\goe}{{\mathfrak{e}}}
\newcommand{\gf}{{\mathfrak{f}}}
\newcommand{\gog}{{\mathfrak{g}}}
\newcommand{\gh}{{\mathfrak{h}}}
\newcommand{\gi}{{\mathfrak{i}}}
\newcommand{\gj}{{\mathfrak{j}}}
\newcommand{\gk}{{\mathfrak{k}}}
%\newcommand{\gl}{{\mathfrak{l}}}
\newcommand{\gm}{{\mathfrak{m}}}
\newcommand{\gn}{{\mathfrak{n}}}
%\newcommand{\go}{{\mathfrak{o}}}
\newcommand{\gp}{{\mathfrak{p}}}
\newcommand{\gq}{{\mathfrak{q}}}
\newcommand{\gor}{{\mathfrak{r}}}
\newcommand{\gs}{{\mathfrak{s}}}
\newcommand{\gt}{{\mathfrak{t}}}
\newcommand{\gu}{{\mathfrak{u}}}
\newcommand{\gv}{{\mathfrak{v}}}
\newcommand{\gw}{{\mathfrak{w}}}
\newcommand{\gx}{{\mathfrak{x}}}
\newcommand{\gy}{{\mathfrak{y}}}
\newcommand{\gz}{{\mathfrak{z}}}
\newcommand{\barA}{\overline{{A}}}
\newcommand{\barB}{\overline{{B}}}
\newcommand{\barC}{\overline{{C}}}
\newcommand{\barD}{\overline{{D}}}
\newcommand{\barE}{\overline{{E}}}
\newcommand{\barF}{\overline{{F}}}
\newcommand{\barG}{\overline{{G}}}
\newcommand{\barH}{\overline{{H}}}
\newcommand{\barI}{\overline{{I}}}
\newcommand{\barJ}{\overline{{J}}}
\newcommand{\barK}{\overline{{K}}}
\newcommand{\barL}{\overline{{L}}}
\newcommand{\barM}{\overline{{M}}}
\newcommand{\barN}{\overline{{N}}}
\newcommand{\barO}{\overline{{O}}}
\newcommand{\barP}{\overline{{P}}}
\newcommand{\barQ}{\overline{{Q}}}
\newcommand{\barR}{\overline{{R}}}
\newcommand{\barS}{\overline{{S}}}
\newcommand{\barT}{\overline{{T}}}
\newcommand{\barU}{\overline{{U}}}
\newcommand{\barV}{\overline{{V}}}
\newcommand{\barW}{\overline{{W}}}
\newcommand{\barX}{\overline{{X}}}
\newcommand{\barY}{\overline{{Y}}}
\newcommand{\barZ}{\overline{{Z}}}
\newcommand{\barAA}{\overline{{\A}}}
\newcommand{\barBB}{\overline{{\B}}}
\newcommand{\barCC}{\overline{{\C}}}
\newcommand{\barDD}{\overline{{\D}}}
\newcommand{\barEE}{\overline{{\E}}}
\newcommand{\barFF}{\overline{{\F}}}
\newcommand{\barGG}{\overline{{\G}}}
\newcommand{\barHH}{\overline{{\HH}}}
\newcommand{\barII}{\overline{{\I}}}
\newcommand{\barJJ}{\overline{{\J}}}
\newcommand{\barKK}{\overline{{\K}}}
\newcommand{\barLL}{\overline{{\LL}}}
\newcommand{\barMM}{\overline{{\M}}}
\newcommand{\barNN}{\overline{{\N}}}
\newcommand{\barOO}{\overline{{\OO}}}
\newcommand{\barPP}{\overline{{\PP}}}
\newcommand{\barQQ}{\overline{{\Q}}}
\newcommand{\barRR}{\overline{{\R}}}
\newcommand{\barSS}{\overline{{\SS}}}
\newcommand{\barTT}{\overline{{\T}}}
\newcommand{\barUU}{\overline{{\U}}}
\newcommand{\barVV}{\overline{{\V}}}
\newcommand{\barWW}{\overline{{\W}}}
\newcommand{\barXX}{\overline{{\X}}}
\newcommand{\barYY}{\overline{{\Y}}}
\newcommand{\barZZ}{\overline{{\Z}}}
\newcommand{\bara}{\overline{{a}}}
\newcommand{\barb}{\overline{{b}}}
\newcommand{\barc}{\overline{{c}}}
\newcommand{\bard}{\overline{{d}}}
\newcommand{\bare}{\overline{{e}}}
\newcommand{\barf}{\overline{{f}}}
\newcommand{\barg}{\overline{{g}}}
\newcommand{\barh}{\overline{{h}}}
\newcommand{\bari}{\overline{{i}}}
\newcommand{\barj}{\overline{{j}}}
\newcommand{\bark}{\overline{{k}}}
\newcommand{\barl}{\overline{{l}}}
\newcommand{\barm}{\overline{{m}}}
\newcommand{\barn}{\overline{{n}}}
%\newcommand{\baro}{\overline{{o}}}
\newcommand{\barp}{\overline{{p}}}
\newcommand{\barq}{\overline{{q}}}
\newcommand{\barr}{\overline{{r}}}
\newcommand{\bars}{\overline{{s}}}
\newcommand{\bart}{\overline{{t}}}
\newcommand{\baru}{\overline{{u}}}
\newcommand{\barv}{\overline{{v}}}
\newcommand{\barw}{\overline{{w}}}
\newcommand{\barx}{\overline{{x}}}
\newcommand{\bary}{\overline{{y}}}
\newcommand{\barz}{\overline{{z}}}
\newcommand{\tA}{\widetilde{{\A}}}
\newcommand{\tAA}{\widetilde{{\A}}}
\newcommand{\tB}{\widetilde{{B}}}
\newcommand{\tC}{\widetilde{{C}}}
\newcommand{\tD}{\widetilde{{D}}}
\newcommand{\tE}{\widetilde{{E}}}
\newcommand{\tF}{\widetilde{{F}}}
\newcommand{\tG}{\widetilde{{G}}}
\newcommand{\tH}{\widetilde{{H}}}
\newcommand{\tI}{\widetilde{{I}}}
\newcommand{\tJ}{\widetilde{{J}}}
\newcommand{\tK}{\widetilde{{K}}}
\newcommand{\tL}{\widetilde{{L}}}
\newcommand{\tM}{\widetilde{{M}}}
\newcommand{\tN}{\widetilde{{N}}}
\newcommand{\tO}{\widetilde{{O}}}
\newcommand{\tP}{\widetilde{{P}}}
\newcommand{\tQ}{\widetilde{{Q}}}
\newcommand{\tR}{\widetilde{{R}}}
\newcommand{\tS}{\widetilde{{S}}}
\newcommand{\tT}{\widetilde{{T}}}
\newcommand{\tU}{\widetilde{{U}}}
\newcommand{\tV}{\widetilde{{V}}}
\newcommand{\tW}{\widetilde{{W}}}
\newcommand{\tX}{\widetilde{{X}}}
\newcommand{\tY}{\widetilde{{Y}}}
\newcommand{\tZ}{\widetilde{{Z}}}
\newcommand{\ta}{\widetilde{{a}}}
\newcommand{\tb}{\widetilde{{b}}}
\newcommand{\tc}{\widetilde{{c}}}
\newcommand{\td}{\widetilde{{d}}}
\newcommand{\te}{\widetilde{{e}}}
\newcommand{\tif}{\widetilde{{f}}}
\newcommand{\tg}{\widetilde{{g}}}
\newcommand{\tih}{\widetilde{{h}}}
\newcommand{\ti}{\widetilde{{i}}}
\newcommand{\tj}{\widetilde{{j}}}
\newcommand{\tk}{\widetilde{{k}}}
\newcommand{\tl}{\widetilde{{l}}}
\newcommand{\tm}{\widetilde{{m}}}
\newcommand{\tn}{\widetilde{{n}}}
\newcommand{\tio}{\widetilde{{o}}}
\newcommand{\tp}{\widetilde{\mathfrak{p}}}
\newcommand{\tq}{\widetilde{{q}}}
\newcommand{\tir}{\widetilde{{r}}}
\newcommand{\ts}{\widetilde{{s}}}
\newcommand{\tit}{\widetilde{{t}}}
\newcommand{\tu}{\widetilde{{u}}}
\newcommand{\tv}{{\widetilde{{v}}}}
\newcommand{\tw}{{\widetilde{{w}}}}
\newcommand{\tx}{\widetilde{{x}}}
\newcommand{\ty}{\widetilde{{y}}}
\newcommand{\tz}{\widetilde{{z}}}
\newcommand{\tcA}{\widetilde{{\cA}}}
\newcommand{\tcB}{\widetilde{{\cB}}}
\newcommand{\tcC}{\widetilde{{\cC}}}
\newcommand{\tcF}{\widetilde{{\cF}}}
\newcommand{\tcG}{\widetilde{{\cG}}}
\newcommand{\tcH}{\widetilde{{\cH}}}
\newcommand{\tcI}{\widetilde{{\cI}}}
\newcommand{\tcK}{\widetilde{{\cK}}}
\newcommand{\tcX}{\widetilde{{\cX}}}
\newcommand{\baralpha   }{\overline{\alpha  }}
\newcommand{\barbeta         }{\overline{\beta}}
\newcommand{\bargamma   }{\overline{\gamma}}
\newcommand{\bardelta        }{\overline{\delta}}
%\newcommand{\barepsilon    }{\overline{\epsilon}}
\newcommand{\barvarepsilon  }{\overline{\varepsilon}}
\newcommand{\barzeta       }{\overline{\zeta}}
\newcommand{\bareta         }{\overline{\eta}}
 \newcommand{\bartheta    }{\overline{\theta}}
\newcommand{\barvartheta   }{\overline{\vartheta}}
\newcommand{\bariota     }{\overline{\iota}}
\newcommand{\barkappa     }{\overline{\kappa}}
 \newcommand{\barlambda   }{\overline{\lambda}}
 \newcommand{\barmu    }{\overline{\mu}}
 \newcommand{\barnu    }{\overline{\nu}}
 \newcommand{\barxi    }{\overline{\xi}}
 \newcommand{\barpi    }{\overline{\pi}}
 \newcommand{\barvarpi   }{\overline{\varpi}}
 \newcommand{\barrho   }{{\overline{\rho}}}
 \newcommand{\barvarrho   }{\overline{\varrho}}
 \newcommand{\barsigma   }{\overline{\sigma}}
 \newcommand{\barvarsigma   }{\overline{\varsigma}}
 \newcommand{\bartau     }{\overline{\tau}}
 \newcommand{\barupsilon   }{\overline{\upsilon}}
 \newcommand{\barphi    }{\overline{\phi}}
 \newcommand{\barvarphi   }{\overline{\varphi}}
 \newcommand{\barchi   }{\overline{\chi}}
 \newcommand{\barpsi   }{\overline{\psi}}
 \newcommand{\baromega   }{\overline{\omega}}
 \newcommand{\barGamma     }{\overline{\Gamma}}
\newcommand{\barDelta         }{\overline{\Delta}}
\newcommand{\barTheta       }{\overline{\Theta}}
\newcommand{\barLambda   }{\overline{\Lambda}}
\newcommand{\barXi   }{\overline{\Xi}}
\newcommand{\barPi   }{\overline{\Pi}}
\newcommand{\barSigma   }{\overline{\Sigma}}
\newcommand{\barUpsilon   }{\overline{\Upsilon}}
\newcommand{\barPhi    }{\overline{\Phi}}
\newcommand{\barPsi   }{\overline{\Psi}}
\newcommand{\barOmega   }{\overline{\Omega}}
\newcommand{\alphabar   }{\overline{\alpha  }}
\newcommand{\betabar         }{\overline{\beta}}
\newcommand{\gammabar   }{\overline{\gamma}}
\newcommand{\deltabar        }{\overline{\delta}}
\newcommand{\epsilonbar    }{\overline{\epsilon}}
\newcommand{\varepsilonbar  }{\overline{\varepsilon}}
\newcommand{\zetabar       }{\overline{\zeta}}
\newcommand{\etabar         }{\overline{\eta}}
 \newcommand{\thetabar    }{\overline{\theta}}
\newcommand{\varthetabar   }{\overline{\vartheta}}
\newcommand{\iotabar     }{\overline{\iota}}
\newcommand{\kappabar     }{\overline{\kappa}}
 \newcommand{\lambdabar   }{\overline{\lambda}}
 \newcommand{\mubar    }{\overline{\mu}}
 \newcommand{\nubar    }{\overline{\nu}}
 \newcommand{\xibar    }{\overline{\xi}}
% \newcommand{\pibar    }{\overline{\pi}}
 \newcommand{\varpibar   }{\overline{\varpi}}
% \newcommand{\rhobar   }{{\overline{\rho}}}
 \newcommand{\varrhobar   }{\overline{\varrho}}
 \newcommand{\sigmabar   }{\overline{\sigma}}
 \newcommand{\varsigmabar   }{\overline{\varsigma}}
 \newcommand{\taubar     }{\overline{\tau}}
 \newcommand{\upsilonbar   }{\overline{\upsilon}}
 \newcommand{\phibar    }{\overline{\phi}}
 \newcommand{\varphibar   }{\overline{\varphi}}
% \newcommand{\chibar   }{\overline{\chi}}
 \newcommand{\psibar   }{\overline{\psi}}
 \newcommand{\omegabar   }{\overline{\omega}}
 \newcommand{\Gammabar     }{\overline{\Gamma}}
\newcommand{\Deltabar         }{\overline{\Delta}}
\newcommand{\Thetabar       }{\overline{\Theta}}
\newcommand{\Lambdabar   }{\overline{\Lambda}}
\newcommand{\Xibar   }{\overline{\Xi}}
\newcommand{\Pibar   }{\overline{\Pi}}
\newcommand{\Sigmabar   }{\overline{\Sigma}}
\newcommand{\Upsilonbar   }{\overline{\Upsilon}}
\newcommand{\Phibar    }{\overline{\Phi}}
\newcommand{\Psibar   }{\overline{\Psi}}
\newcommand{\Omegabar   }{\overline{\Omega}}
\newcommand{\talpha   }{\widetilde{\alpha  }}
\newcommand{\tbeta         }{\widetilde{\beta}}
\newcommand{\tgamma   }{\widetilde{\gamma}}
\newcommand{\tdelta        }{\widetilde{\delta}}
\newcommand{\tepsilon    }{\widetilde{\epsilon}}
\newcommand{\tvarepsilon  }{\widetilde{\varepsilon}}
\newcommand{\tzeta       }{\widetilde{\zeta}}
\newcommand{\teta         }{\widetilde{\eta}}
 \newcommand{\ttheta    }{\widetilde{\theta}}
\newcommand{\tvartheta   }{\widetilde{\vartheta}}
\newcommand{\tiota     }{\widetilde{\iota}}
\newcommand{\tkappa     }{\widetilde{\kappa}}
 \newcommand{\tlambda   }{\widetilde{\lambda}}
 \newcommand{\tmu    }{\widetilde{\mu}}
 \newcommand{\tnu    }{\widetilde{\nu}}
 \newcommand{\txi    }{\widetilde{\xi}}
 \newcommand{\tpi    }{\widetilde{\pi}}
 \newcommand{\tvarpi   }{\widetilde{\varpi}}
 \newcommand{\trho   }{\widetilde{\rho}}
 \newcommand{\tvarrho   }{\widetilde{\varrho}}
 \newcommand{\tsigma   }{\widetilde{\sigma}}
 \newcommand{\tvarsigma   }{\widetilde{\varsigma}}
 \newcommand{\ttau     }{\widetilde{\tau}}
 \newcommand{\tupsilon   }{\widetilde{\upsilon}}
 \newcommand{\tphi    }{\widetilde{\phi}}
 \newcommand{\tvarphi   }{\widetilde{\varphi}}
 \newcommand{\tchi   }{\widetilde{\chi}}
 \newcommand{\tpsi   }{\widetilde{\psi}}
 \newcommand{\tomega   }{\widetilde{\omega}}
 \newcommand{\tGamma     }{\widetilde{\Gamma}}
 \newcommand{\Gammat     }{\widetilde{\Gamma}}
\newcommand{\tDelta         }{\widetilde{\Delta}}
\newcommand{\tTheta       }{\widetilde{\Theta}}
\newcommand{\tLambda   }{\widetilde{\Lambda}}
\newcommand{\tXi   }{\widetilde{\Xi}}
\newcommand{\tPi   }{\widetilde{\Pi}}
\newcommand{\tSigma   }{\widetilde{\Sigma}}
\newcommand{\tUpsilon   }{\widetilde{\Upsilon}}
\newcommand{\tPhi    }{\widetilde{\Phi}}
\newcommand{\tPsi   }{\widetilde{\Psi}}
\newcommand{\tOmega   }{\widetilde{\Omega}}
\newcommand{\alphat   }{\widetilde{\alpha  }}
\newcommand{\betat         }{\widetilde{\beta}}
\newcommand{\gammat   }{\widetilde{\gamma}}
\newcommand{\deltat        }{\widetilde{\delta}}
\newcommand{\epsilont    }{\widetilde{\epsilon}}
\newcommand{\varepsilont  }{\widetilde{\varepsilon}}
\newcommand{\zetat       }{\widetilde{\zeta}}
\newcommand{\etat         }{\widetilde{\eta}}
 \newcommand{\thetat    }{\widetilde{\theta}}
\newcommand{\varthetat   }{\widetilde{\vartheta}}
\newcommand{\iotat     }{\widetilde{\iota}}
\newcommand{\kappat     }{\widetilde{\kappa}}
 \newcommand{\lambdat   }{\widetilde{\lambda}}
 \newcommand{\mut    }{\widetilde{\mu}}
 \newcommand{\nut    }{\widetilde{\nu}}
 \newcommand{\xit    }{\widetilde{\xi}}
 \newcommand{\pit    }{\widetilde{\pi}}
 \newcommand{\varpit   }{\widetilde{\varpi}}
 \newcommand{\rhot   }{\widetilde{\rho}}
 \newcommand{\varrhot   }{\widetilde{\varrho}}
 \newcommand{\sigmat   }{\widetilde{\sigma}}
 \newcommand{\varsigmat   }{\widetilde{\varsigma}}
 \newcommand{\taut     }{\widetilde{\tau}}
 \newcommand{\upsilont   }{\widetilde{\upsilon}}
 \newcommand{\phit    }{\widetilde{\phi}}
 \newcommand{\varphit   }{\widetilde{\varphi}}
 \newcommand{\chit   }{\widetilde{\chi}}
 \newcommand{\psit   }{\widetilde{\psi}}
 \newcommand{\omegat   }{\widetilde{\omega}}
\newcommand{\Deltat         }{\widetilde{\Delta}}
\newcommand{\Thetat       }{\widetilde{\Theta}}
\newcommand{\Lambdat   }{\widetilde{\Lambda}}
\newcommand{\Xit   }{\widetilde{\Xi}}
\newcommand{\Pit   }{\widetilde{\Pi}}
\newcommand{\Sigmat   }{\widetilde{\Sigma}}
\newcommand{\Upsilont   }{\widetilde{\Upsilon}}
\newcommand{\Phit    }{\widetilde{\Phi}}
\newcommand{\Psit   }{\widetilde{\Psi}}
\newcommand{\Omegat   }{\widetilde{\Omega}}
\newcommand{\hatotimes}{{\widehat{\otimes}}}
\newcommand{\barCG}{\overline{{\CG}}}
\DeclareMathOperator{\proj}{proj}
\newcommand{\MK}{{\underline{M}_{\CO_K}}}
%\newcommand{\MF}{{\CM\CF}}
%\newcommand{\BT}{{Barsotti-Tate $\CO_K$-module}}
%\newcommand{\Qbar}{{\overline{\Q}}}
%\newcommand{\Fbar}{{\overline{\F}}}
\newcommand{\tbarr}{{\widetilde{\barr}}}
%
%\newcommand{\F}{\mathbb{F}}
%\newcommand{\kbar}{{\overline{k}}}
%
\def\RCS$#1: #2 ${\expandafter\def\csname RCS#1\endcsname{#2}}
\RCS$Revision: 85 $
\RCS$Date: 2011-12-16 18:54:17 -0600 (Fri, 16 Dec 2011) $
\newcommand{\versioninfo}{Version \RCSRevision; Last commit \RCSDate}
%
%\DeclareMathOperator{\Frob}{Frob}
%\DeclareMathOperator{\sgn}{sgn}
\DeclareMathOperator{\RE}{Re} \DeclareMathOperator{\IM}{Im}
\DeclareMathOperator{\ess}{ess} %\DeclareMathOperator{\ab}{ab}
%\DeclareMathOperator{\BC}{BC}
\newcommand{\notdiv}{\nmid}
%\newcommand{\m}{\mathfrak{m}}
 \newcommand{\h}{\mathcal{H}}
\newcommand{\s}{\mathcal{S}} %\newcommand{\J}{\mathcal{J}}
%\newcommand{\M}{\mathcal{M}} \newcommand{\rec}{\operatorname{rec}}
%\newcommand{\St}{\operatorname{St}}
%\newcommand{\W}{\mathcal{W}} \newcommand{\X}{\mathcal{X}}
%\newcommand{\cL}{\mathcal{L}}
\newcommand{\bigO}{\mathcal{O}} \newcommand{\calO}{\mathcal{O}}
%\newcommand{\G}{\mathcal{G}}
\newcommand{\BOP}{\mathbb{B}}
\newcommand{\BH}{\mathbb{B}(\mathcal{H})}
\newcommand{\KH}{\mathcal{K}(\mathcal{H})} %\newcommand{\R}{\mathbb{R}}
%\newcommand{\Z}{\mathbb{Z}} \newcommand{\A}{\mathbb{A}}
%\newcommand{\Q}{\mathbb{Q}} \newcommand{\N}{\mathbb{N}}
%\newcommand{\T}{\mathbb{T}} \newcommand{\gn}{\mathcal{G}_n}
\newcommand{\p}{\mathfrak{p}} \newcommand{\q}{\mathfrak{q}}
%\renewcommand{\gg}{\mathfrak{g}} \newcommand{\gh}{\mathfrak{h}}
%\renewcommand{\l}{\mathfrak{l}}
\newcommand{\frakl}{\mathfrak{l}}
\newcommand{\fp}{\mathfrak{p}}
\newcommand{\fq}{\mathfrak{q}} \newcommand{\fQ}{\mathfrak{Q}}
%\newcommand{\C}{\mathbb{C}}
\newcommand{\Field}{\mathbb{F}}
\newcommand{\Gn}{\mathcal{G}_n}
\newcommand{\sq}{\square}
\newcommand{\tensor}{\otimes}
\newcommand{\thbarpr}{\bar{\theta'}}
%% specific to TWK papers
\newcommand{\Favoid}{F^{(\mathrm{avoid})}}
%% fonts etc
\newcommand{\mcal}{\mathcal}
\newcommand{\bb}{\mathbb}
\newcommand{\mc}{\mathcal}
\newcommand{\mf}{\mathfrak}
%% Hecke operators and algebras
\newcommand{\Uop}{U_{\lambda,\varpi_{\wt{v}}}^{(j)}}
\newcommand{\TUniv}{\widehat{\bb{T}}^{T,\ord}(U(\mf{l}^{\infty}),\mc{O})}
\newcommand{\TT}{\bb{T}^{T}_{\lambda}(U,\mc{O})}
%\DeclareMathOperator{\ord}{ord}
\DeclareMathOperator{\lcm}{lcm}
\DeclareMathOperator{\charpol}{char}
\DeclareMathOperator{\sesi}{ss}
\DeclareMathOperator{\ssg}{ss}
\DeclareMathOperator{\sd}{sd}
\DeclareMathOperator{\GO}{GO}
\DeclareMathOperator{\FL}{FL}
\newcommand{\RPlus}{\Real^{+}} \newcommand{\Polar}{\mathcal{P}_{\s}}
\newcommand{\Poly}{\mathcal{P}(E)} \newcommand{\EssD}{\mathcal{D}}
\renewcommand{\o}[1]{\overline{#1}}
\newcommand{\rbar}{{\overline{r}}}
\newcommand{\sbar}{{\overline{s}}}
\newcommand{\Rbar}{\overline{R}}
\newcommand{\rbarwtv}{\rbar|_{G_{F_{\wt{v}}}}}
\newcommand{\rbarwtvL}{\rbar|_{G_{L_{\wt{v}}}}}
\newcommand{\rmbarwtv}{\rbar_{\mf{m}}|_{G_{F_{\wt{v}}}}}
\newcommand{\Lom}{\mathcal{L}}
\newcommand{\States}{\mathcal{T}}
%\newcommand{\abs}[1]{\left\vert#1\right\vert}
\newcommand{\set}[1]{\left\{#1\right\}}
\newcommand{\seq}[1]{\left<#1\right>}
%\newcommand{\norm}[1]{\left\Vert#1\right\Vert}
%\newcommand{\Gal}{\operatorname{Gal}}
%\newcommand{\gl}{\operatorname{\mathfrak{gl}}}
\newcommand{\HT}{\operatorname{HT}}
\newcommand{\Pj}{\mathbb{P}}
\newcommand{\FI}{\operatorname{FI}}
%\newcommand{\Gr}{\operatorname{Gr}}
%\newcommand{\Zbar}{\overline{\Z}}
\newcommand{\GQ}{\Gal(\Qbar/\Q)} \newcommand{\Qp}{{\Q_p}}
\newcommand{\Qpn}{{\Q_{p^n}}}
\newcommand{\GQp}{{G_{\Q_p}}}
\newcommand{\WQp}{{W_{\Q_p}}}
\newcommand{\IQp}{{I_{\Q_p}}}
\newcommand{\Zp}{{\Z_p}}
\newcommand{\Ql}{\Q_l} \newcommand{\Cp}{\C_p}
\newcommand{\Qpbar}{{\overline{\Q}_p}}
\newcommand{\Qpbartimes}{{\overline{\Q}_p^\times}}
\newcommand{\Zpbar}{{\overline{\Z}_p}}
\newcommand{\Qlbar}{\overline{\Q}_{l}}
\newcommand{\Fpbar}{{\overline{\F}_p}}
\newcommand{\Fpbartimes}{{\overline{\F}_p^\times}}
\newcommand{\Fnew}{\F'}
\newcommand{\Fnewnew}{\F''}
\newcommand{\Fl}{{\F_l}}
\newcommand{\Fp}{{\F_p}}
\newcommand{\Flbar}{\overline{\F}_l}
\newcommand{\gfv}{{G_{F_{\tilde{v}}}}}
\newcommand{\Jac}{\operatorname{Jac}}
\newcommand{\Inn}{\operatorname{Inn}}
\newcommand{\rank}{\operatorname{rank}}
\newcommand{\Ver}{\operatorname{Ver}}
%\newcommand{\Sym}{\operatorname{Sym}}
\newcommand{\sep}{\operatorname{sep}}
%\newcommand{\Mod}{\operatorname{Mod}}
\newcommand{\rmod}{R-\Mod_{cris,dd}^{k-1}}
\newcommand{\bigs}{\mathfrak{S}}
%labelled bullet points---TBL thinks these might be useful for x refs in the final argument
\newcommand{\labelbp}[1]
{\refstepcounter{equation}\hfill(\theequation)\label{#1}}
%% Marginal comments

\newif\iffinalrun
%\finalruntrue
\iffinalrun
\else
  % \usepackage[notref,notcite]{showkeys}
  % \usepackage[notref]{showkeys}
\fi


\iffinalrun
  \newcommand{\need}[1]{}
  \newcommand{\mar}[1]{}
\else
  \newcommand{\need}[1]{{\tiny *** #1}}
  \newcommand{\mar}[1]{\marginpar{\raggedright\tiny fixme #1}}
\fi


\newcommand{\Ainf}{\mathbf{A}_{\inf}}
\newcommand{\AAinf}[1]{\mathbf{A}_{\inf,#1}}
% macros added by DS
\newcommand{\emb}{\kappa}
\newcommand{\embb}{\overline{\emb}}
\newcommand{\uni}{\pi}
%
% Small matrix in text macro.
%
\def\smallmat#1#2#3#4{\bigl(\begin{smallmatrix}{#1}&{#2}\\{#3}&{#4}\end{smallmatrix}\bigr)}
%
%\usepackage{hyperref} % PDF hyperlinks - needs to go close
                      % to \begin{document} or it breaks
%%%% end of TG macros
\newtheorem{para}[subsection]{\bf}
\newtheorem{sublemma}{Sub-lemma}
\newtheorem{theorem}[subsection]{Theorem}
\newtheorem{thm}[subsection]{Theorem}
\newtheorem{lemma}[subsection]{Lemma}
\newtheorem{lem}[subsection]{Lemma}

\newtheorem{ithm}{Theorem}
\renewcommand{\theithm}{\Alph{ithm}}

\newtheorem{iconj}[ithm]{Conjecture}
%\renewcommand{\theiconj}{\Alph{iconj}}


\newtheorem{cor}[subsection]{Corollary}
\newtheorem{conj}[subsection]{Conjecture}

\newtheorem{prop}[subsection]{Proposition}


\newtheorem{atheorem}[section]{Theorem}
\newtheorem{athm}[section]{Theorem}
\newtheorem{alemma}[section]{Lemma}
\newtheorem{alem}[section]{Lemma}
\newtheorem{aconj}[section]{Conjecture}

\newtheorem{acor}[section]{Corollary}
\newtheorem{aprop}[section]{Proposition}

\theoremstyle{definition}
\newtheorem{df}[subsection]{Definition}
\newtheorem{defn}[subsection]{Definition}
\newtheorem{situation}[subsection]{Situation}
\newtheorem{adf}[section]{Definition}
\newtheorem{adefn}[section]{Definition}
\theoremstyle{remark}
\newtheorem{remark}[subsection]{Remark}
\newtheorem{rem}[subsection]{Remark}
\newtheorem{aremark}[section]{Remark}
\newtheorem{arem}[section]{Remark}
\newtheorem{question}[subsection]{Question}
\newtheorem{caution}[subsection]{Caution}
\newtheorem{example}[subsection]{Example}
\newtheorem{aexample}[section]{Example}
\newtheorem{property}[subsection]{Property}
\newtheorem{condition}[subsection]{}
\newtheorem{terminology}[subsection]{Terminology}
\newtheorem{assumption}[subsection]{Assumption}
\newtheorem{hyp}[subsection]{Hypothesis}

\def\numequation{\addtocounter{subsection}{1}\begin{equation}}
\def\nummultline{\addtocounter{subsection}{1}\begin{multline}}
\def\anumequation{\addtocounter{section}{1}\begin{equation}}
%\newenvironment{numequation}{\addtocounter{subsubsection}{1}\begin{equation}}{\end{equation}}
%\newenvironment{anumequation}{\addtocounter{subsection}{1}\begin{equation}}{\end{equation}}
%\newenvironment{nummultline}{\addtocounter{subsubsection}{1}\begin{multline}}{\end{multline}}
%\newenvironment{anummultline}{\addtocounter{subsection}{1}\begin{multline}}{\end{multline}}
%\newenvironment{numequation*}{\begin{equation*}}{\end{equation*}}
%\newenvironment{anumequation*}{\begin{equation*}}{\end{equation*}}
%\newenvironment{nummultline*}{\begin{multline*}}{\end{multline*}}
%\newenvironment{anummultline*}{\begin{multline*}}{\end{multline*}}
\renewcommand{\theequation}{\arabic{chapter}.\arabic{section}.\arabic{subsection}}
\newcommand{\ssinc}{\addtocounter{subsection}{1}}
%\def\RCS$#1: #2 ${\expandafter\def\csname RCS#1\endcsname{#2}}
%\RCS$Revision: 85 $
%\RCS$Date: 2011-12-16 18:54:17 -0600 (Fri, 16 Dec 2011) $
%\newcommand{\versioninfo}{Version \RCSRevision; Last commit \RCSDate}
%\hyphenation{rep-re-sent-a-tion}

\newenvironment{dedication}
  {%\clearpage           % we want a new page          %% I commented this
   \thispagestyle{empty}% no header and footer
   \vspace*{\stretch{1}}% some space at the top
   \itshape             % the text is in italics
   \raggedleft          % flush to the right margin
  }
  {\par % end the paragraph
   \vspace{\stretch{3}} % space at bottom is three times that at the top
   \clearpage           % finish off the page
  }

  \newcommand{\MR}{}



\title% [Moduli stacks of \'etale $(\varphi,\Gamma)$-modules]
{Moduli
  stacks of \'etale $(\varphi,\Gamma)$-modules and the existence of
  crystalline lifts}
\author{
Matthew Emerton
  \and
Toby Gee
}
\date{}
% \textbf{This is a rapid prototype}
  \begin{document}
  
% $p$-adic and mod $p$
% representations of Galois groups of $p$-adic local fields}

%\author[M. Emerton]{Matthew Emerton}%\email{emerton@math.uchicago.edu}
% \address{Department of Mathematics, University of Chicago,
% 5734 S.\ University Ave., Chicago, IL 60637}

%\author[T. Gee]{Toby Gee} %\email{toby.gee@imperial.ac.uk}
% \address{Department of
\frontmatter
Moduli stacks of \'etale $(\varphi,\Gamma)$-modules and the existence
of crystalline lifts

\cleardoublepage
\thispagestyle{empty}


\maketitle
%% better to use this in book class
 \chapter{Dedication}
  \begin{dedication}
In memory of Jean-Marc Fontaine.
\end{dedication}

\thispagestyle{empty}
\cleardoublepage
\thispagestyle{empty}
%.
%.
%.
%.

  % \chapter{Thanks}
%   %   Mathematics, Imperial College London}
% The first author was supported in part by the
%   NSF grants DMS-1303450, DMS-1601871, and DMS-1902307. The second author was 
%   supported in part by a Leverhulme Prize, EPSRC grant EP/L025485/1, Marie Curie Career
%   Integration Grant 303605,
%   ERC Starting Grant 306326, ERC Advanced Grant 884596, and a Royal Society Wolfson Research
%   Merit Award. This project has received funding from the European Research Council (ERC) under the European Union’s Horizon 2020 research and innovation programme (grant agreement No. 884596).
  % \chapter{Abstract}
%   We construct stacks which algebraize Mazur's %universal
%   formal deformation
%   rings of local Galois representations. More precisely, we construct
%   %$p$-adic
%   Noetherian formal algebraic stacks over $\Spf \Z_p$
%   which parameterize \'etale
%   $(\varphi,\Gamma)$-modules;
% % Galois-invariant lattices in $p$-adic local Galois representations,
%   the formal completions of these stacks at points in their special
%   fibres recover the universal deformation rings of local Galois
%   representations. We use these stacks to show that all mod~$p$
%   representations of the absolute Galois group of a $p$-adic local
%   field lift to characteristic zero, and indeed admit crystalline
%   lifts. We also discuss the relationship between the geometry of our
%   stacks and the Breuil--M\'ezard conjecture.

\setcounter{tocdepth}{1}
\tableofcontents
  
\mainmatter             %% and use this for main chapters.




% Main things to be done:
% \begin{itemize}
% \item Section~\ref{sec:
%   crystalline and semistable} is the biggest gap, in particular
% figuring out the treatment of Hodge--Tate weights.
% \item A couple of minor marginal comments in Section~\ref{section:
%     coefficient rings} about arguments with associated graded.
% \item Argument with dimensions and flat pullback to finish proof of
%   Theorem~\ref{thm: codimension H2 in regular weight crystalline
%     deformation ring}
% \item Appendix~\ref{app:p inverted} (Banach algebras) to be finished.
% \item Section~\ref{sec: the rank one case} (rank one case) to be
%   written. %Maybe this should actually be a section towards the end of
% %   the body of the paper?
% %   \item Appendix~\ref{app: formal algebraic
% %   stacks} (a summary of~\cite{Emertonformalstacks}) to be written.
% \end{itemize}


\chapter{Introduction}%\mar{TG: needs a total rewrite}
% \mar{TG: if only for my sake, we should explain why we
%     don't think $\cX$ is $p$-adic formal algebraic, and (related?) why
%   we don't expect its versal rings (the unrestricted deformation
%   rings) to be effective.}
% \mar{TG: should mention BM
  % in this introduction, and that we'll be coming back to that in a sequel.}
% \mar{TG: add a comparison to the $l\ne p$ case. TG: actually I'm not
% really inclined to do this, commenting out.}
%\mar{TG: % some things to add to this introduction: % explain that
  % % sections 4 and 5 are essentially independent.
  % Emphasise the
  % BKF/$(\varphi,G_K)$-perspective more, especially as
  % $(\varphi,\Gamma)$ is in the title. Explain what intuition we can
  % for the crystalline directions being (some of) the ones that
  % algebraize (probably from global considerations?).
%   Discuss further
%   what happens if you define a stack of Weil group
%   representations - is the point just to say that it isn't algebraic
%   and doesn't seem to have useful properties? Probably we will discuss
% this in Section 7 and can reference this here.
  % Explain the notion of underlying reduced substack
  % % better.
  % Is there a remark that we could make about this versus
  % actual group representations being like Betti/de Rham in geometric Langlands?
% ME: Added a comment about sections 4/5 below; see what you think TG:
% looks good, commenting out
% }
%\mar{TG: warning - we now have marginal comments!}
% \mar{ME: warning - we now have marginal comments! ME: Right now, this is the only
% comment!}
% \mar{TG: we should think about whether to add longer intros to some of
%   the chapters, and we should also refer the reader to the Bonn notes.}
%\mar{TG: run a spellcheck at the end.}
% \mar{TG: another thing to fix is that all instances in the text of
%   stuff like ``subsection blah'' should change. Perhaps we should ask
%   PUP if there's a preferred style? TG: we could of course simply not
%   worry about this! ME: I've gone through the paper and (hopefully) taken
%   care of this.}
In this book we construct moduli stacks of \'etale $(\varphi,\Gamma)$-modules
(projective, of some fixed rank, 
and with coefficients in $p$-adically complete rings),
and establish some of their basic properties.  We also present some first
applications of this construction to the theory of Galois representations.

\section{Motivation}\label{subsec: motivation}Mazur's theory of
deformations of Galois representations~\cite{MR1012172}
is modelled on the geometric study of
%infinitesimal deformations
infinitesimal neighbourhoods of points in 
moduli spaces
via 
formal deformation theory.
In the mid-2000s,
Kisin suggested %in around 2004
that some kind of
moduli spaces of local Galois representations should
exist; that is, there should be formal algebraic stacks
%\mar{TG: formal
%  algebraic over $\Zp$, if we want to say that we get the full
%  universal lifting rings? ME: Or even crystalline def.\ rings,
%if we're over $\Z_p$; made this change}
over~$\Zp$ whose
closed points correspond to representations
$\rhobar:G_K\to\GL_d(\Fpbar)$, and
whose versal rings at such points
%the versal rings at the points of the stack
should recover appropriate Galois deformation rings. This
expectation is borne out by the results of this book. (In fact, Kisin
was motivated by calculations of crystalline deformation rings
for~$\GL_2(\Qp)$ that had been carried out by Berger--Breuil using the
$p$-adic local Langlands correspondence, and suggested that the versal
rings should give crystalline deformation rings.  Thus his suggestion
is realized by the stacks~$\cX_d^{\crys,\underline{\lambda}}$ of
Theorem~\ref{thm:intro statement on crystalline moduli} below.)

% However, in the context of Galois representations,
% up till now there have been no 
% %there are no
% % \mar{ME: Rephrased, because we actually do make them!}
% corresponding ``global'' moduli spaces, and it turns
% out that naive attempts to define such spaces,
%e.g.\
A natural way to construct such a stack would be to  consider a
literal moduli stack of continuous representations
$\rho:G_K\to\GL_d(A)$, for $K$ a $p$-adic field
and $A$ a $p$-adically complete $\Z_p$-algebra; indeed such stacks
were constructed by Carl Wang-Erickson~\cite{MR3831282}. However,
the stacks constructed in this way are less ``global'' than one would
wish, and in particular the corresponding families of mod~$p$
representations $\rhobar:G_K\to\GL_d(\Fpbar)$ have constant semisimplification.
%\mar{ME: I don't think the base coefficient ring $\cO$ has been introduced at this point}
% \mar{TG: I rewrote this to discuss the
%   CWE stuff, see what you think.  ME: Nice, made minor changes.}

In this book, we instead consider moduli stacks of \'etale
$(\varphi,\Gamma)$-modules. These
contain Wang-Erickson's stacks as substacks, and coincide with them on
%the latter agree with Wang-Erickson's stacks on
the level of $\Fpbar$-points, but their geometry is quite different; in particular,
we see much larger families,  % the moduli stacks
% that we consider below are surprisingly ``large'', and 
exhibiting some unexpected features (for example, irreducible
representations arising as limits of reducible representations). The
relationship between the theory and constructions that we develop
here, and the usual formal deformation theory of Galois
representations, is the same as that between the theory of moduli
spaces of algebraic varieties, and the formal deformation theory of
algebraic varieties: the latter gives valuable local information about
the former, but moduli spaces, when they can be constructed, capture
global aspects of the situation %which are
inaccessible to the purely
infinitesimal tools of formal deformation theory.
% \mar{ME: I don't think this notation for crystalline stacks has
% actually been introduced anywhere yet in the paper! TG: at the moment
% in Section~\ref{sec: BM} I have stuff like
% $\cX^{\crys,\lambdau,\tau}_d$. I don't feel strongly about it at all
% but maybe we should agree on something once and for all? We could also
% try to make it consistent with the deformation ring notation\dots
% ME: Changed notation to (hopefully) match def.\ ring notation.} 
%\mar{TG: is that right? ME: I th% ink so; added something to emphasize this.}

% Older (but good, merge in):

% These spaces will be formal algebraic stacks,
% \mar{ME: Perhaps actually
% 	just Ind-algebraic stacks? TG: seems to be formal algebraic
%         stacks, in fact? ME: As long as the cotangent complex argument
% hangs together \dots }
% and points 
% in the special fibre will correspond to residual Galois representations.
% The formal completion of one of our stacks at such a point will then recover
% the usual formal deformation space of this residual representation.
% Thus the relationship between the theory and constructions
% that we develop here, and the usual formal deformation theory of Galois
% representations, is the same as that between the theory of moduli spaces
% of algebraic varieties, and the formal deformation theory of algebraic 
% varieties: the latter gives valuable local information about the former,
% but moduli spaces, when they can be constructed,
% capture global aspects of the situation which are
% inaccessible by the purely infinitesimal tools of formal deformation theory.


\section{Our main theorems}
Our goal in this book is to construct, and establish the basic
properties of, moduli stacks of \'etale $(\varphi,\Gamma)$-modules.
More precisely, if we fix a finite extension $K$ of $\Q_p$, and a
non-negative integer $d$ (the rank), then we let $\cX_d$ denote the
category fibred in groupoids over $\Spf \Z_p$ whose groupoid of
$A$-valued points, for any $p$-adically complete $\Z_p$-algebra $A$,
is equal to the groupoid of rank~ $d$ projective
\'etale $(\varphi,\Gamma)$-modules with $A$-coefficients.
(See Section~~\ref{subsec: phi Gamma
  coefficients}
below for a definition of these.) % \mar{TG: refer to later in the intro or
%   just to section~\label{subsec:
% defn of Xd}}). % \mar{ME: Explain
  % this notation here?  Or just use $W(k)\otimes_{\Z_p} A((T))$
  % instead? TG: I'd be inclined to do the latter ($p$-adically
  % completed), or (better?) even just to say ``with $A$-coefficients'' and then
  % make it more precise what this means after stating the main theorems.}
Our first main theorem is the following. (See
Corollary~\ref{cor:Xd is formal algebraic} and
Theorem~\ref{thm:reduced dimension}.)% \mar{TG: indicate that we will
%   elaborate on labelling by Serre weights later in the introduction.
% ME: Added some discussion just below.}

\begin{thm}\label{thm:intro statement of basic properties of X}
	The category fibred in groupoids $\cX_d$ is a 
	Noetherian formal algebraic
	stack.  Its underlying reduced substack~$\cX_{d,\red}$
        %\mar{TG: we should probably say that this is algebraic?}
	% \mar{ME: Currently the paper is full of both the notation
	% 	$\cX_{d,\red}$ and $(\cX_d)_\red$.  We should
	% 	probably settle on one or the other. TG: I think that
        %         we should do the former, else the stacks in places
        %         like Theorem~\ref{thm:Xdred is algebraic} start having
        %       horrible names.}
	{\em (}which is an algebraic stack{\em )}
	is of finite type
	over $\F_p$, and is equidimensional of dimension
	$[K:\Q_p] d(d-1)/2$. The irreducible components of 
	$\cX_{d,\red}$
%	the underlying reduced substack
	admit a natural labelling by Serre
        weights. % \mar{TG: say that its $\Fpbar$ points correspond to
          % mod $p$ repns?}
\end{thm}
%\mar{TG: insert a theorem here about the pst deformation stacks?
%ME: Definitely, especially since I've added a sentence refering to them
%below, in the context of MK's suggestions.}


%\begin{thm}\label{thm:intro statement on crystalline moduli}
%	If $\underline{\lambda}$ is a collection
%	of labeled Hodge--Tate weights,
%	then there is a closed substack $\cX_d^{\underline{\lambda},\crys}$
%	of $\cX_d$ which is a $p$-adic formal algebraic stack and flat 
%	over $\Z_p$, and which 
%	is characterized as being the unique closed substack
%	of $\cX_d$ which is flat over $\Z_p$ and whose groupoid 
%	of $A$-valued points,
%	for any finite dimensional $\Q_p$-algebra $A$,
%	is equivalent
%	{\em (}under the equivalence between $(\varphi,\Gamma)$-modules
%	and $G_K$-representations{\em )} with the groupoid of continuous
%	representations
%	$G_K\to \GL_d(A)$ which are {\em crystalline},
%	with labeled Hodge--Tate weights equal to $\underline{\lambda}$.
%
%	Similarly, 
%	there is a closed substack $\cX_d^{\underline{\lambda},\ss}$
%	of $\cX_d$ which is a $p$-adic formal algebraic stack and flat 
%	over $\Z_p$, and which 
%	is characterized as being the unique closed substack
%	of $\cX_d$ which is flat over $\Z_p$ and whose groupoid 
%	of $A$-valued points,
%	for any finite dimensional $\Q_p$-algebra $A$,
%	is equivalent
%	with the groupoid of continuous representations
%	$G_K\to \GL_d(A)$ which are {\em semistable},
%	with labeled Hodge--Tate weights equal to $\underline{\lambda}$.
%\end{thm}




We will elaborate on the labelling of components by Serre weights
further below.  For now, we mention that,
under the usual
correspondence between \'etale $(\varphi,\Gamma)$-modules
and Galois representations,
the groupoid of $\Fpbar$-points of~$\cX_d$,
which coincides with the groupoid of $\Fpbar$-points
of the underlying reduced substack $\cX_{d,\red}$,
is naturally equivalent to the groupoid of continuous representations
$\rhobar:G_K\to\GL_d(\Fpbar)$.  (More generally, if $A$ is any finite $\Z_p$-algebra,
then the groupoid $\cX_d(A)$ is canonically equivalent
to the groupoid of continuous representations $G_K \to \GL_d(A)$.) It is expected 
% \mar{ME: Not sure if I should say ``will be the case'' or ``should be
% 	the case'' here, since I'm not sure exactly how (or whether)
% 	Serre weights
% 	of not-necessarily tame $\rhobar$ are unambiguously defined,
% 	and, even if they are, whether this assertion is actually
%         known. TG: definitely only expected. Probably it would be
%         better to say that we expect that the weights are determined
%         by the labels in a slightly complicated way, and make some
%         reference to later sections (either in the introduction, or to
%       the body of the paper).}
 that our labelling of the irreducible components can be refined (by
 adding further labels to some of the components) to
 give a geometric description of the weight part of Serre's
 conjecture, so that $\rhobar$ corresponds to a point in a component
of $\cX_{d,\red}$ which is labeled by the Serre weight
$\underline{k}$ if and only if $\rhobar$ admits $\underline{k}$ as a
Serre weight; we discuss this expectation, and what is known
about it, in Section~\ref{subsec:
  geometric BM} below (and in more detail in Chapter~\ref{sec: BM}).

Again using the
correspondence between \'etale $(\varphi,\Gamma)$-modules
and Galois representations,
we see that the universal lifting ring of a 
representation $\rhobar$ as above will provide a versal ring
to~$\cX_d$ at the corresponding $\Fbar_p$-valued
point. Accordingly we expect that the stacks~$\cX_d$ will have
applications to the study of Galois representations and
their deformations. As a first
example of this, we prove the following result on
the existence of crystalline lifts; although the statement of this
theorem involves a fixed~$\rhobar$, we do not know how to prove it
 without using the stacks~$\cX_d$, over which~$\rhobar$ varies.
\begin{thm}[Theorem~\ref{thm: strong existence of crystalline lifts}]
	\label{thm:intro crystalline lifts}
	If $\rhobar: G_K \to \GL_d(\Fbar_p)$ is a continuous
	representation, then $\rhobar$
	has a lift $\rho^{\circ}: G_K \to \GL_d(\Zbar_p)$ for which
	the associated $p$-adic representation $\rho:G_K \to \GL_d(\Qbar_p)$
	is crystalline of regular Hodge--Tate weights. 

We can furthermore %  choose
        % these Hodge--Tate weights to be in the
        % interval~$[0,dp-1]$, % \mar{TG: not too sure about this - might be
        %   % better to do bounded gaps instead, although obviously a
        %   % result of this kind will be true.}
	% % {\em (}the bound depending only on $d$ and $K${\em )},
        % or we
        % can choose them to have arbitrarily large gaps, and in either
        % case we can also
        ensure that~$\rho^{\circ}$ is potentially diagonalizable.
% 	\mar{ME: We could consider making the bound more explicit, either
% 		in the statement of the result or in a subsequent remark.}
% 	\mar{ME: Maybe add potentially diagonalizable into the statement,
% 		or perhaps discuss it afterwards, if it seems a bit
% 		out of place in the statement at this particular
%                 moment. TG: also want well-spread, rather than
%                 bounded, as an option. ME: Currently, can get any plausible
% 	property of the HT weights, but can't get pot'l diagonalizable ---
% have to work on this.}
\end{thm}
(The notion of a potentially diagonalizable representation was
introduced in~\cite{BLGGT}, and is recalled as Definition~\ref{defn:
  pot diag} below.) In combination with
%together with
potential automorphy 
 theorems, this has the following application to the globalization of
local Galois representations. % (see Corollary~\ref{cor: existence of
%   global lift}; this result follows from Theorem~\ref{thm:intro
%   crystalline lifts} together with (potential) automorphy
% theorems).
% (By definition, an automorphic Galois representation lifts
% to characteristic zero, and the existence of local lifts is a
% necessary input for all known arguments to produce global lifts using
% automorphy lifting theorems or Galois deformation arguments.)

% \mar{ME: I'm sure there's a good reason for the choice of designation
% by $\rbar$ and $\rhobar$ here, but maybe we could switch them?  since
% $\rhobar$ is local throughout the rest of the introduction. TG:
% changed here and in the body of the text; commenting out.}
\begin{thm}[Corollary~\ref{cor: existence of
  global lift}]
  \label{thm: intro existence of global lifts}Suppose that~$p\nmid 2d$, and
  fix $\rhobar:G_K\to\GL_d(\Fpbar)$. Then there is an imaginary CM
  field~$F$ and an irreducible conjugate self dual automorphic Galois representation
  $\rbar:G_F\to\GL_d(\Fpbar)$ such that  for every $v|p$, we have
  $F_v\cong K$ and either~$\rbar|_{G_{F_v}}\cong \rhobar$ or~$\rbar|_{G_{F_{v^c}}}\cong \rhobar$.
\end{thm}
% Note that by definition an automorphic Galois representation lifts to
% characteristic zero. 


Another key result of the book is the following theorem,
describing moduli stacks of \'etale $(\varphi,\Gamma)$-modules corresponding to
{\em crystalline} and {\em semistable} Galois representations.

\begin{thm}[Theorem~\ref{thm: existence of ss stack}]\label{thm:intro statement on crystalline moduli}
	If $\underline{\lambda}$ is a collection
	of labeled Hodge--Tate weights,
	and if $\cO$ denotes the ring of integers
	in a finite extension $E$ of $\Q_p$ containing the Galois
	closure of $K$ {\em(}which will serve as the ring of coefficients{\em )},
	then there is a closed substack $\cX_d^{\crys,\underline{\lambda}}$
	of $(\cX_{d})_{\cO}$ which is a $p$-adic formal algebraic stack and is flat 
	over~$\cO$, and which 
	is characterized as being the unique closed substack
	of $(\cX_{d})_{\cO}$ which is flat over $\cO$ and whose groupoid 
	of $A$-valued points,
	%\mar{TG: is this correct? We should
         % discuss this interpretation somewhere if so. If not we could do
         % what is written later in the paper and stick to finite flat
         % $\Zp$-algebras. Maybe this theorem should really have some
         % coefficients that are big enough to have ``labeled HT
         % weights'' make sense, too?}
	for any finite flat $\cO$-algebra~$A$,
	is equivalent
	{\em (}under the equivalence between \'etale
	$(\varphi,\Gamma)$-modules and continuous
	$G_K$-representations{\em )} to the groupoid of continuous
	representations
	$G_K\to \GL_d(A)$ which become {\em crystalline}
	after extension of scalars to $A\otimes_{\cO} E$,
	and whose labeled Hodge--Tate weights are
	equal to~$\underline{\lambda}$.

	Similarly, 
	there is a closed substack $\cX_d^{\semis,\underline{\lambda}}$
	of $(\cX_{d})_{\cO}$ which is a $p$-adic formal algebraic stack and is flat 
	over~$\cO$, and which 
	is characterized as being the unique closed substack
	of $(\cX_{d})_{\cO}$ which is flat over $\cO$ and whose groupoid 
	of $A$-valued points,
	for any finite flat $\cO$-algebra~$A$,
	is equivalent
	to the groupoid of continuous representations
	$G_K\to \GL_d(A)$ which become {\em semistable}
	after extension of scalars to $A\otimes_{\cO} E$,
	and whose labeled Hodge--Tate weights are
	equal to $\underline{\lambda}$.
\end{thm}

\begin{remark}
  In fact, Theorem~\ref{thm: existence of ss stack} also proves
  the analogous result for potentially crystalline and potentially
  semistable representations of arbitrary inertial type, but for
  simplicity of exposition we restrict ourselves to the crystalline
  and semistable cases in this introduction.
%	\mar{ME: Add citation to where we do this.}
\end{remark}

A crucial distinction between the stacks $\cX_d$ and their
closed substacks
$\cX_d^{\crys,\underline{\lambda}}$
and
$\cX_d^{\semis,\underline{\lambda}}$
is that while $\cX_d$ is a formal algebraic stack lying over $\Spf \Z_p$,
it is not actually a $p$-adic formal algebraic stack (in the sense of Definition~\ref{adefn: p adic formal algebraic stack}); 
see Proposition~\ref{prop: X is not padic formal algebraic}.
On the other hand,
the stacks
$\cX_d^{\crys,\underline{\lambda}}$
and
$\cX_d^{\semis,\underline{\lambda}}$
{\em are} $p$-adic formal algebraic stacks, which implies 
that their mod $p^a$ reductions are in fact algebraic stacks.
This gives in particular a strong interplay between the structure of the
mod~ $p$ fibres of crystalline and semistable
lifting rings and the geometry of the underlying reduced substack~$\cX_{d,\red}$.  This plays an important role in determining the structure of this reduced
substack, and also in the proof of
Theorem~\ref{thm:intro crystalline lifts}.
As we explain in more detail in Section~\ref{subsec:
  geometric BM} below, % , \mar{ME: Give reference, once
	% the intro.\ is finalized.}
it also allows us to reinterpret the Breuil--M\'ezard conjecture
in terms of the interaction between the
structure of the mod $p$ fibres of the stacks
$\cX_d^{\crys,\underline{\lambda}}$
and
$\cX_d^{\semis,\underline{\lambda}}$,
and the geometry of $\cX_{d,\red}$.% \mar{TG: can we say something here
%   about why one might expect the crystalline directions to algebraise?
% Is there some motivation coming from something global? That's what I
% remember us saying, but I don't now recall what that is! TG: it's just
% that we don't expect the global ones do, and the crystalline
% directions are (at least heuristically) orthogonal to those, so it's
% the most we could hope for, in some sense. TG: I wasn't convinced
% this was worth saying; commenting out.}




\section{\texorpdfstring{$(\varphi,\Gamma)$}{(phi,Gamma)}-modules
  with coefficients}\label{subsec: phi Gamma coefficients}
There is quite a lot of evidence, for example from Colmez's work on
the $p$-adic local Langlands correspondence~\cite{MR2642409}, and work
of Kedlaya--Liu
\cite{MR3379653}, % \mar{TG: add citations to both of these}
that rather than considering families of representations of~$G_K$, it
is more natural to consider families of
\'etale~$(\varphi,\Gamma)$-modules.

The theory
of \'etale~$(\varphi,\Gamma)$-modules for~$\Zp$-representations was introduced
by Fontaine in~\cite{MR1106901}. There are various possible definitions that can
be made, with perfect, imperfect or overconvergent coefficient rings,
and different choices of~$\Gamma$; we discuss the various variants
that we use, and the relationships between them, at some length in the
body of the book.  For the purpose of this introduction we simply
let~$\A_K=W(k)((T))^\wedge$, where~$k$ is a finite extension of~$\Fp$
(depending on~$K$), and the hat denotes the $p$-adic completion. This
ring is endowed with a Frobenius~$\varphi$ and an action of a
profinite group~$\Gamma$ (an open subgroup of~$\Z_p^\times$) that
commutes with~$\varphi$; the formulae for $\varphi$ and for this action can be rather
complicated for general~$K$, although they admit a simple description
if~$K/\Qp$ is abelian.  (See Definition~\ref{defn: basic} and the
surrounding material.) % \mar{TG: reference wherever basic/absolutely
% abelian ends up.}

An \'etale $(\varphi,\Gamma)$-module is then, by
definition, a finite $\A_K$-module endowed with commuting semilinear
actions of~$\varphi$ and~$\Gamma$, with the property that the
linearized action of~$\varphi$ is an isomorphism. There is a natural
equivalence of categories between the category of \'etale
$(\varphi,\Gamma)$-modules and the category of continuous
representations of~$G_K$ on finite $\Zp$-modules.

Let~$A$ be a $p$-adically complete $\Zp$-algebra. We
let~$\A_{K,A}:=A\cotimes_{\Zp}\A_K$ (where the completed tensor
product is taken with respect to the $p$-adic topology on~$A$
and the so-called weak topology on~$\A_K$),
%, and
%for the topology  on~$\A_K$ for which~$W(k)[[T]]$ is open
%and has the $(p,T)$-adic topology),
%\mar{ME: Rephrase slightly (because
%	the topology on $\A_K$ is not an adic topology)? TG: better? ME: Yep,
%	commenting out ME: Sorry, I think this is still off if we're working
%$p$-adically rather than mod~$p^a$.}
and define an \'etale
$(\varphi,\Gamma)$-module with $A$-coefficients 
%in exactly the same way
just 
as in the case~$A=\Zp$ described above, but now using the ring~$\A_{K,A}$.
 In the case that~$A$ is
finite as a $\Zp$-module there is again an equivalence of categories
with the category of continuous representations of~$G_K$ on finite
$A$-modules, but for more general~$A$ no such equivalence
exists. Our moduli stack~$\cX_d$ is defined to be the stack
over~$\Spf\Zp$ with the property that~$\cX_d(A)$ is the groupoid
of projective \'etale $(\varphi,\Gamma)$-modules of rank~$d$ with $A$-coefficients. (That
this is indeed an \'etale stack, indeed even an \emph{fpqc} stack,
follows from results of
Drinfeld.) Using the machinery of our %earlier
paper~\cite{EGstacktheoreticimages} we are able to show that~$\cX_d$
is an Ind-algebraic stack, but to prove Theorem~\ref{thm:intro
  statement of basic properties of X} we need to go further and make a
detailed study of its special fibre and of the underlying reduced
substack. This study is guided by ideas coming from Galois
deformation theory and the weight part of Serre's conjecture,
in a manner that we now describe.
%as we now explain.

\section{Families of extensions}\label{subsec: families of extensions}% \mar{TG: discuss that generically
  % the representations are reducible, also weight part of Serre's
  % conjecture, referencing CEGS and our sequel papers - also could
  % reference the section on the picture in the GHS paper, after
  % checking that there's not too much nonsense there; maybe also BM for components piling
  % up}
As we have already explained, over a general base~$A$ there is no
longer an equivalence between $(\varphi,\Gamma)$-modules and
representations of~$G_K$. Perhaps surprisingly, from the point of view
of applications of our stacks to the study of $p$-adic Galois
representations, this is a feature rather than a bug. For example, an
examination of the known results on the reductions modulo~$p$ of 2-dimensional
crystalline representations of~$G_{\Qp}$ (see for
example~\cite[Thm.\ 5.2.1]{MR2906353}) suggests that any
moduli space of mod~$p$ representations of~$G_K$ should have the
feature that the representations are generically reducible, but can
specialise to irreducible representations. A literal moduli space of
representations of a group cannot behave in this way (essentially
because Grassmannians are proper), but it turns out that the
underlying reduced substack~$\cX_{d,\red}$ of~$\cX_d$ does have this
property. (See also Section~\ref{sec:
  comparison to CWE} and Remark~\ref{rem: rank 1 versus rank d} for 
further discussions of the relationship between our stacks of
$(\varphi,\Gamma)$-modules and stacks of representations of a Galois
or Weil--Deligne group.) % \mar{TG: there will presumably now be a
  % section to refer to.}
% \mar{ME: This tapers off \dots ; maybe something was accidentally deleted?  Couldn't find
% where, though, after looking through a bunch of commit messages. TG:
% finished the sentence, commented out.}

More precisely, the results of~\cite[Thm.\ 5.2.1]{MR2906353}, together
with the weight part of Serre's conjecture, suggest that each
irreducible component of~$\cX_{d,\red}$ should contain a dense set of
$\Fpbar$-points which are successive extensions of characters
$G_K\to\Fpbar$, with the extensions being as non-split as possible. This
turns out to be the case. The restrictions of these characters to the
inertia subgroup~$I_K$ are constant on the irreducible components, and
the discrete data of these characters, together with some further
information about peu- and tr\`es ramifi\'ee extensions, determines
the components. This discrete data can be conveniently and naturally
organised in terms of ``Serre weights'' $\underline{k}$, which are
tuples of integers which biject with the isomorphism classes of the
irreducible $\Fpbar$-representations of~$\GL_d(\cO_K)$. The
relationship between Serre weights and Galois representations is
important in the $p$-adic Langlands program, and in proving automorphy
lifting theorems, and we discuss it further in Section~\ref{subsec:
  geometric BM}. % but we
% will not discuss it in this paper; for a detailed explanation in the
% case~$d=2$, see the forthcoming paper~\cite{CEGSKisinwithdd}, and for
% a brief explanation of our expectations in general (which we will
% return to in a future paper), see~\cite[\S6]{2015arXiv150902527G}.

Having guessed that the $\Fpbar$-points of~$\cX_d$ should be arranged
in irreducible components in this way, an inductive strategy to prove
this suggests itself. It is easy to see that irreducible
representations of~$G_K$ are ``rigid'', in that there are up to twist
by unramified characters only finitely many in each dimension;
furthermore, it is at least intuitively clear that each such family of
unramified twists of a $d$-dimensional irreducible representation
should give rise to a zero-dimensional substack of~$\cX_d$ (there is a
$\Gm$ of twists, but also a $\Gm$ of automorphisms). On the other
hand, given
characters~$\chibar_1,\dots,\chibar_d:G_K\to\Fbar_p^\times$, a Galois
cohomology calculation suggests that there should be a substack
of~$\cX_{d,\red}$ % \mar{TG: define notation for underlying reduced}
of dimension~$[K:\Qp]d(d-1)/2$ given by successive
extensions of unramified twists of the~$\chibar_i$. Accordingly, one
could hope to construct the stacks % ~$\cX_d^{\underline{k}}$
corresponding to the Serre weights~$\underline{k}$ by inductively
constructing families of extensions of representations.

To confirm this expectation, we use the machinery originally developed
by Herr~\cite{MR1693457,MR1839766}, who gave an explicit complex
which is defined in terms of~$(\varphi,\Gamma)$-modules and computes Galois
cohomology. This definition goes over unchanged to the case with
coefficients, and with some effort we are able to adapt Herr's
arguments to our setting, and to prove finiteness and base change
properties (following Pottharst~\cite{MR3117501}, we in fact find it
helpful to think of the Herr complex of a $(\varphi,\Gamma)$-module
with $A$-coefficients as a perfect complex of $A$-modules). Using the Herr
complex, we can inductively construct irreducible closed
substacks~$\cX_{d,\red}^{\underline{k}}$ of~$\cX_{d,\red}$ of
dimension~$[K:\Qp]d(d-1)/2$ whose generic $\Fpbar$-points correspond
to successive extensions of characters as described above (the
restrictions of these characters to~$I_K$ being
determined by~$\underline{k}$). % \mar{TG:
  % explain labelling of Serre weights around here somewhere.}
Furthermore, by a rather involved induction, we can show that the
union of the~$\cX_{d,\red}^{\underline{k}}$, together possibly with a closed
substack of~$\cX_{d,\red}$ of dimension strictly less than~$[K:\Qp]d(d-1)/2$,
exhausts~$\cX_{d,\red}$. In particular, each ~$\cX_{d,\red}^{\underline{k}}$ is an irreducible component
of~$\cX_{d,\red}$, and any %other
irreducible component that is not one of the $\cX_{d,\red}^{\underline{k}}$
is of strictly
smaller dimension than these components.

One way to show that the~$\cX_{d,\red}^{\underline{k}}$ exhaust the
irreducible components of~$\cX_{d,\red}$ would be to show that every
representation $G_K\to\GL_d(\Fpbar)$ occurs as an $\Fbar_p$-valued point of
some~$\cX_{d,\red}^{\underline{k}}$.  We expect this to be 
difficult to show directly;
% \mar{ME: But we do show this in the end, right?  So maybe we mean to say
% 	that this seems to be very difficult to do directly? TG: yes, rewritten}
indeed, already for $d=2$ the
paper~\cite{CEGSKisinwithdd} shows that the closed points
of~$\cX_{d,\red}^{\underline{k}}$ are governed by the weight part of
Serre's conjecture, and the explicit description of this conjecture is
complicated (see e.g.\ \cite{BuzzardDiamondJarvis,MR3589336}). % \mar{TG: add DDR
  % citation?}
Furthermore it seems hard to explicitly understand the way
in which families of reducible $(\varphi,\Gamma)$-modules degenerate
to irreducible ones, or to reducible representations with different
restrictions to~$I_K$
(phenomena which are implied by the weight part of Serre's conjecture).

Instead, our approach is to show by a consideration of versal rings that~$\cX_{d,\red}$ is equidimensional of
dimension~$[K:\Qp]d(d-1)/2$; this suffices, since our inductive construction showed
that any other irreducible component would necessarily have dimension
strictly less than~$[K:\Qp]d(d-1)/2$. Our proof of this
equidimensionality relies on Theorems~\ref{thm:intro crystalline
  lifts} and~\ref{thm:intro statement on crystalline moduli}, as we
explain in Remark~\ref{rem: dimension and codimension} below.% the following subsections.\mar{TG: make sure that we
% do in fact explain this!}



% \cite[\S6]{2015arXiv150902527G} \cite{CEGSKisinwithdd}

% $\cX_{d,\red}^{\underline{k}}$

%  \mar{TG: here}

% Versal rings to guess dimensions, Herr complex to justify

% Somewhere discuss dimension, not $p$-adic formal algebraic, also
% effectivity crucial as begin with upper bounds

\section{Crystalline lifts}\label{subsec: crystalline lifts}Theorem~\ref{thm:intro crystalline
  lifts} solves a problem that has been considered by various
authors, in particular~\cite{MullerThesis,2015arXiv150601050G}.
It admits a well-known inductive approach (which is taken in~\cite{MullerThesis,2015arXiv150601050G}): one writes $\rhobar$ as 
a successive extension of irreducible representations, lifts each of these
irreducible representations 
to a crystalline representation, and then attempts to lift the various
extension classes.  The difficulty that arises in this approach (which
has proved an obstacle to obtaining general statements along the lines
of Theorem~\ref{thm:intro crystalline lifts} until now) is showing that the 
mod $p$ extension classes that appear in this description of $\rhobar$ 
can actually be lifted to crystalline extension classes in characteristic
zero.  The basic source of the difficulty is that the local Galois $H^2$
can be non-zero, and non-zero classes in $H^2$ obstruct the lifting 
of extension classes (which can be interpreted as classes lying in
$H^1$). In fact, the difficulty is not so much in obtaining
crystalline extension classes, as in lifting to any classes in
characteristic zero; indeed, 
it was not previously known that an arbitrary~$\rhobar$ had \emph{any} lift to
characteristic zero at all. (Subsequently a different proof of the
existence of such a lift has been found by B\"ockle--Iyengar--Pa{\v{s}}k{\=u}nas~\cite{BIP}.)
%\mar{ME: Mention \cite{BIP}? TG: done, commenting out}

Our proof of Theorem~\ref{thm:intro crystalline lifts} relies on the
inductive strategy described in the preceding paragraph, but we are able to
prove the following key result, which controls the obstructions that
can be presented by $H^2$, and is a consequence of
Theorems~\ref{thm:Xdred is algebraic} and~\ref{thm:reduced dimension}
(see also Remark~\ref{rem: dimension and codimension}).

\begin{prop}
	\label{prop:H2 dimension}
	The locus of points $\rhobar \in \cX_{d,\red}(\Fbar_p)$ 
	at which \[\dim H^2(G_K,\rhobar) \geq r\] is Zariski closed
	in $\cX_{d,\red}(\Fbar_p)$, and is of codimension $\geq r$.
      \end{prop}


Let~$R^{\square}_{\rhobar}$ denote the universal lifting ring
of~$\rhobar$, with universal lifting~$\rho^{\univ}$. For each regular
tuple of labeled Hodge--Tate weights $\underline{\lambda}$, we
let~$R_{\rhobar}^{\crys,\underline{\lambda}}$ denote the quotient
of~$R^{\square}_{\rhobar}$ corresponding to crystalline lifts
of~$\rhobar$ with Hodge--Tate weights~$\underline{\lambda}$ (of
course, this quotient is zero unless~$\rhobar$ admits such a
crystalline lift). Then $H^2(G_K, \rho^{\univ})$ is 
an $R^{\square}_{\rhobar}$-module, and Proposition~\ref{prop:H2 dimension}
implies the following corollary.

\begin{cor}
	\label{cor:H2 dimension}
	For any regular tuple of labeled Hodge--Tate weights $\underline{\lambda}$
	the locus of points 
	$x\in \Spec R^{\crys,\underline{\lambda}}_{\rhobar}/p$ for which
	$$\dim_{\kappa(x)}
 H^2(G_K,\rho^{\univ})	\otimes_{R^{\square}_{\rhobar}}\kappa(x)
	\geq r$$
	has codimension $\geq r$.
\end{cor}% \mar{TG: explain that it's crucial that we have effectivity
  % to prove this.}

\begin{remark}
	\label{rem:H2 dimension}
	Tate local duality, 
together with the compatibility of $H^2$ with base-change,
shows that
\begin{multline*}
	\dim_{\kappa(x)}
	 \bigl( H^2(G_K,\rho^{\univ})\otimes_{R^{\square}_{\rhobar}}\kappa(x)
\bigr)
	=
	\dim_{\kappa(x)}
	 H^2\bigr(G_K,\rho^{\univ}\otimes_{R^{\square}_{\rhobar}}\kappa(x)\bigl)
\\
	=
	\dim_{\kappa(x)}
	\Hom_{G_K}\bigl((\rho^{\univ})^{\vee}\otimes_{R^{\square}_{\rhobar}}\kappa(x),\epsilonbar\bigr)
\end{multline*}
	(where $\epsilonbar$ denotes the mod $p$ cyclotomic character, thought
	of as taking values in $\kappa(x)^{\times}$).
	Thus the statement of Corollary~\ref{cor:H2 dimension} is
	related to the way in which % the special fibre of the crystalline
	% lifting ring
        $\Spec R_{\rhobar}^{\crys,\underline{\lambda}}/p$ intersects
	the reducibility locus in $\Spec R^{\square}_{\rhobar}$.
\end{remark}

Given Corollary~\ref{cor:H2 dimension},
we prove Theorem~\ref{thm:intro crystalline
	lifts}
by working purely within the context of formal lifting rings.  
However we don't know how to
prove the corollary while staying within that context.
Indeed, as Remark~\ref{rem:H2 dimension} indicates, 
this corollary is related to the way in which the special fibre
of a potentially crystalline deformation ring intersects another natural
locus in $\Spec R_{\rhobar}^{\square}$ (namely, the reducibility locus).
Since the special fibre of a potentially crystalline lifting ring
is not directly defined in deformation-theoretic terms, such questions
are notoriously difficult to study directly.
Our proof of the corollary proceeds
differently, by %  using % (an expected globalization to~$\cX_d$ of)
% the  geometric Breuil--M\'ezard view-point
% of~\cite{emertongeerefinedBM} 
% %discussed in Subsection~\ref{subsec:intro Breuil--Mezard} above
% to
replacing a computation on the special fibre of the potentially
crystalline deformation ring by a computation on $\cX_{d,\red}$;
this latter space has a concrete description in terms of families of varying
~$\rhobar$, whose $H^2$ we are  able to compute, as a result of the
inductive construction of families of extensions described in Section~\ref{subsec: families of extensions}.

In order to deduce Corollary~\ref{cor:H2 dimension} from
Proposition~\ref{prop:H2 dimension}, it is crucial that we know that
the natural morphism $ \Spf R^{\crys,\underline{\lambda}}/p\to\cX_d$
is effective, in the sense that it arises from a morphism
$ \Spec R^{\crys,\underline{\lambda}}/p\to\cX_d$. More concretely, the
universal representation~$\rho^\univ$ gives an \'etale
$(\varphi,\Gamma)$-module over each Artinian quotient
of~$R_{\rhobar}^{\square}$. By passing to the limit over these
quotients, we obtain a
``universal formal \'etale $(\varphi,\Gamma)$-module''
over the completion of $(k\otimes_{\Zp}R_{\rhobar}^{\square}/p)((T))$
with respect to the maximal ideal~$\m$
of~$R_{\rhobar}^{\square}$. Since the special fibre of~$\cX_d$ is
formal algebraic but not algebraic (see Section~\ref{subsec: further
  questions} below), there is no corresponding $(\varphi,\Gamma)$-module
with $R_{\rhobar}^{\square}/p$-coefficients;
the  $\varphi$ and $\Gamma$ actions 
on the universal formal \'etale $(\varphi,\Gamma)$-module
 involve Laurent
tails of unbounded degree (with the coefficients of~$T^{-n}$ tending
to zero~$\m$-adically as $n\to\infty$).

The assertion that $ \Spf R^{\crys,\underline{\lambda}}/p\to\cX_d$ is
effective is equivalent to showing that the base change of the universal formal
\'etale $(\varphi,\Gamma)$-module to~$R^{\crys,\underline{\lambda}}/p$
arises from a genuine $(\varphi,\Gamma)$-module, i.e.\ from one that involves
only Laurent
tails of bounded degree. We deduce this from Theorem~\ref{thm:intro
  statement on crystalline moduli}. Indeed, the ring~$R^{\crys,\underline{\lambda}}/p$ is a versal ring for the special
fibre of the $p$-adic formal algebraic
stack~$\cX_{d}^{\crys,\lambdau}$, and (by the very definition of a
$p$-adic formal algebraic stack) this special fibre is an algebraic
stack; and the versal rings for algebraic stacks are always effective.

\begin{rem}
  \label{rem: dimension and codimension}As our citation of both
  Theorems~\ref{thm:Xdred is algebraic} and~\ref{thm:reduced
    dimension} for the proof of Proposition~\ref{prop:H2 dimension}
  may indicate, our proof of Proposition~\ref{prop:H2 dimension} is
  somewhat intricate. Indeed, in Theorem~\ref{thm:Xdred is algebraic},
  we show that~$\cX_{d,\red}$ has dimension at most~$[K:\Qp]d(d-1)/2$,
  and that the locus considered in Proposition~\ref{prop:H2 dimension}
  has dimension at most $[K:\Qp]d(d-1)/2-r$. This is in fact enough to
  deduce Corollary~\ref{cor:H2 dimension}, as ~$\Spec
  R^{\crys,\lambdau}/p$ is known to be equidimensional.

  Given Corollary~\ref{cor:H2 dimension}, we prove
  Theorem~\ref{thm:intro crystalline lifts}. In combination with the
  effective versality of the crystalline deformation rings discussed
  above we are then able to deduce the equidimensionality
  of~$\cX_{d,\red}$, and then also prove Proposition~\ref{prop:H2
    dimension} as stated.
\end{rem}

\section{Crystalline and semistable moduli stacks}\label{subsec: intro crystalline
  stacks}


%\mar{TG: make sure to discuss effective versality}

We now explain the proof of Theorem~\ref{thm:intro statement on
  crystalline moduli}; the proof is essentially identical in the
crystalline and semistable cases, so we concentrate on the crystalline
case. To prove the theorem, it is necessary to have a criterion for a
$(\varphi,\Gamma)$-module to come from a crystalline Galois
representation. In the case that~$K/\Qp$ is unramified, it is possible
to give an explicit criterion in terms of Wach
modules~\cite{MR1415732}, but no such direct description is known for
general~$K$. Instead, following Kisin's construction of the
crystalline deformation rings~$R_{\rhobar}^{\crys,\lambda}$
in~\cite{MR2373358}, we use the theory of Breuil--Kisin modules. More
precisely, Kisin shows that crystalline representations of~$G_K$ have
finite height over the (non-Galois) Kummer extension~$K_\infty/K$
obtained by adjoining a compatible system of $p$-power roots of a
uniformizer of~$K$; here being of finite height means that the
corresponding \'etale $\varphi$-modules admit certain $\varphi$-stable
lattices, called Breuil--Kisin modules.

While not every representation of finite height over~$K_\infty$ comes
from a crystalline representation, we are able to show in
Appendix~\ref{app: BKF pst} (jointly written by T.G.\ and Tong Liu)
that a representation $G_K\to\GL_d(\Zpbar)$ is crystalline if and only
if it is of finite height for \emph{every} choice of~$K_\infty$, and
if the corresponding Breuil--Kisin modules satisfy certain natural
compatibilities. (These compatibilities are best expressed in terms of
Bhatt--Scholze's prismatic site, as in~\cite{bhatt2021prismatic}, but
we do not make use of that perspective in this book. Instead, we write
down explicit conditions on the corresponding Breuil--Kisin--Fargues
modules; recall that Breuil--Kisin--Fargues modules are a variant of
Breuil--Kisin modules introduced by Fargues, see e.g.\ \cite[\S 4]{2016arXiv160203148B}.)% \mar{TG: reference BS for
% this?}
%\mar{TG: added sentence on BKF modules here.  ME: Great}

We use this description of the crystalline representations to
prove the existence of the stacks~$\cX_d^{\crys,\lambdau}$. % via the
% machinery developed in~\cite{EGstacktheoreticimages}. \mar{ME: I don't think we actually use this machinery; we just  define them as scheme-theoretic images
% in the regular sense}
The proof
that~$\cX_d^{\crys,\lambdau}$ is a $p$-adic formal algebraic stack
relies on an
% We studied the corresponding stacks of $\varphi$-modules in our
% earlier paper~\cite{EGstacktheoreticimages}, and it is relatively
% straightforward to deduce the required effectivity from these results
% (we also make use of
analogue of results of Caruso--Liu~\cite{MR2745530} on extensions of
the Galois action on Breuil--Kisin modules, which roughly speaking
says that the action of~$G_{K_\infty}$ determines the action of~$G_K$
up to a finite amount of ambiguity. More precisely, given a
Breuil--Kisin module over a $\Z/p^a$-algebra for some~$a\ge 1$, there
is a finite subextension $K_s/K$ of~$K_\infty/K$ depending only
on~$a$, $K$ and the height of the Breuil--Kisin module, such that there is a canonical action
of~$G_{K_s}$ on the corresponding Breuil--Kisin--Fargues
%\mar{ME: BKF modules have not been defined or  even mentioned before
%  now. TG: given that Breuil--Kisin modules have only just been
%  introduced this doesn't seem so bad, but I added a sentence about
%  them above. ME: OK}
module. This canonical action is constructed by Frobenius
amplification, and in the case that the Breuil--Kisin module arises from the
reduction modulo~$p^a$ of a crystalline representation of~$G_K$, the canonical
action coincides with the restriction to~$G_{K_s}$ of the $G_K$-action
on the Breuil--Kisin--Fargues module. % \mar{TG: probably we will want
  % to slightly reword this once we've fixed the mistake?}
(In \cite{MR2745530} a version
of this canonical action is used to prove ramification bounds on the
reductions modulo~$p^a$ of crystalline representations; in
Chapter~\ref{sec: the rank one case}, we use analogous arguments in
the setting of~$(\varphi,\Gamma)$-modules to relate our stacks to
stacks of Weil group representations in the rank one case.)

There is one significant technical difficulty, which is that we need
to define morphisms of stacks that correspond to the restriction of
Galois representations from~$K$ to~$K_\infty$. In order to do this we
have to compare $(\varphi,\Gamma)$-modules with $A$-coefficients
(which are defined via the cyclotomic extension
$K(\zeta_{p^\infty})/K$) to $\varphi$-modules with $A$-coefficients
defined via the extension~$K_\infty/K$. We do not know of a direct way to do this; we
proceed by proving a correspondence between $\varphi$-modules over
Laurent series rings with $\varphi$-modules over the perfections of
these Laurent series rings,
and proving the following descent result which may be of independent
interest; in the statement,
$\C$ denotes the completion of an algebraic closure of~$\Qp$.

%\mar{TG: mention prismatic site}
\begin{theorem}[Theorem~\ref{thm:descending projective
    modules}]\label{thm:descending projective modules intro version}
	  Let $A$ be a finite type $\Z/p^a$-algebra, for some $a \geq 1$. Let
  $F$ be a closed perfectoid subfield of $\C$, with tilt~$F^\flat$, a
  closed perfectoid subfield of $\C^\flat$.  Write
  $W(F^\flat)_A := W(F^\flat)\otimes_{\Z_p}A .$ 

Then the inclusion
  $W(F^\flat)_A \to W(\C^\flat)_A$ is a faithfully flat morphism of
  Noetherian rings, and the functor
  $M \mapsto W(\C^\flat)_A\otimes_{W(F^\flat)_A} M$ induces an
  equivalence between the category of finitely generated projective
  $W(F^\flat)_A$-modules and the category of finitely generated
  projective $W(\C^\flat)_A$-modules endowed with a continuous
  semi-linear $G_F$-action.
\end{theorem} The existence of the required morphism of stacks follows
easily from two applications of Theorem~\ref{thm:descending projective modules intro
  version}, applied with~$F$ equal to respectively the completion of~$K_\infty$ and
the completion of~$K(\zeta_{p^\infty})$. Furthermore, this
construction gives an alternative description of our stacks, as moduli
spaces of $W(\C^\flat)_A$-modules endowed with commuting
  semi-linear actions of~$G_{K}$ and~$\varphi$.  It seems plausible
  that this description will be useful in future work, as it connects
  naturally to the theory of Breuil--Kisin--Fargues modules (and
  indeed we use this connection in our construction of the potentially
  semistable moduli stacks). Note though that the description in terms
  of~$(\varphi,\Gamma)$-modules is important (at least in our
  approach) for establishing the basic finiteness properties of our stacks.

\section{The geometric Breuil--M\'ezard conjecture and the weight
  part of Serre's conjecture}\label{subsec:
  geometric BM}\sectionmark{The geometric Breuil--M\'ezard conjecture}% \mar{TG: just
  % say epsilon about how the effectivity result motivate existence of
  % crystalline stacks and that we'll come back to this in a sequel.}
%\mar{TG: again say that everything works with types.}
We will now briefly explain our results and conjectures relating our
stacks to the Breuil--M\'ezard conjecture and the weight part of
Serre's conjecture. Further explanation and motivation can be found
throughout Chapter~\ref{sec: BM}. Some of these results were previewed
in~\cite[\S 6]{2015arXiv150902527G}, and the earlier sections of that
paper (in particular the introduction) provide an overview of the
weight part of Serre's conjecture and its connections to the
Breuil--M\'ezard conjecture that may be helpful to the reader who is
not already familiar with them. As in the rest of this introduction,
we ignore the possibility of inertial types, and we also restrict to
crystalline representations for the purpose of exposition. Everything
in this section extends to the more general setting of potentially
semistable representations, and indeed as we explain in
Section~\ref{subsec: CEGS} when discussing the
papers~\cite{CEGSKisinwithdd} and~\cite{geekisin}, the additional
information provided by non-trivial inertial types is very important.

Let $\rhobar: G_K \to \GL_d(\F)$ be a continuous representation (for some finite
extension $\F$ of $\F_p$), and let $R_{\rhobar}^\square$ be the
corresponding universal lifting ring. The corresponding formal scheme $\Spf R_{\rhobar}^\square$ doesn't carry
a lot of evident geometry in and of itself; for example, its underlying reduced
subscheme is simply the closed point $\Spec \F$, corresponding 
to $\rhobar$ itself. % \mar{TG: should probably use either $k$ or $\F$ in
% the introduction, but not both; my guess is that $\F$ will be
% consistent with the rest of the paper? TG: changed, commenting out.}
On the other hand, $\cX_d$ has a quite non-trivial underlying reduced
substack~$\cX_{d,\red}$, which parameterizes all the $d$-dimensional
residual representations of $G_K$.  It is natural to ask whether this
underlying reduced substack has any significance in formal deformation
theory.  More precisely, we could ask for the meaning of the fibre
product $ \Spf R_{\rhobar}^\square \times_{\cX_d} \cX_{d,\red}.$

This fibre product is a reduced closed formal subscheme of $\Spf R_{\rhobar}^\square$
of dimension $d^2+[K:\Q_p]d(d-1)/2$.   It arises (via completion at
the closed point) from a closed subscheme of $\Spec R_{\rhobar}^\square$ (as 
does any closed formal subscheme of the $\Spf$ of a complete Noetherian local
ring), whose irreducible components,
when thought of as cycles on $\Spec R_{\rhobar}^\square$, are precisely the cycles
that (conjecturally) appear in the geometric Breuil--M\'ezard
conjecture of~\cite{emertongeerefinedBM}. % \mar{TG: usual $\Spf$ versus
  % $\Spec$, not sure whether we need to mention effectivity of
  % deformation rings here. ME: Added something to deal with this.}
%\mar{ME: Elaborate on/improve this discussion, add citation, etc.}
%As a specific application of this idea,
More precisely, we obtain the following qualitative version of the
geometric Breuil--M\'ezard
conjecture~\cite[Conj.~4.2.1]{emertongeerefinedBM}.

\begin{theorem}[Theorem~\ref{thm: qualitative BM}]\label{thm: intro geom BM without multiplicities}
	If $\rhobar: G_K \to \GL_d(\F)$ is a continuous
        representation, 
	% \mar{ME: Is there any technical reason we can't just take $\F
        %   = \Fbar_p$, and so avoid the annoyance of having to talk
        %   about a finite $\F$? TG: I suspect so; one approach would be
        % to follow BLGGT, I suspect, or we could see if it's easy to
        % do ``properly'' by directly defining the deformation
        % rings. TG: I had a go at doing this in the
        % notation/conventions section; see what you think. TG: need to
        % add coefficients here to account for possibly ramified types}
	% {\em (}with $\F$ being a finite extension of $\F_p${\em )},
	then there are finitely many cycles of dimension
        $d^2+[K:\Q_p]d(d-1)/2$ in $\Spec R^{\square}_{\rhobar}/p$,
        such that for any regular tuple of labeled Hodge--Tate
        weights $\underline{\lambda}$, %each of
        the special fibre
        $\Spec R^{\crys,\underline{\lambda}}_{\rhobar}/p$ is
        set-theoretically supported on the union of some number of
        these
        cycles.% \mar{TG: does this language imply the set-theoretic
  %         support is exactly this (which is too strong)? ME: Added
  % language to address this.}
\end{theorem}

The cycles in the statement of the theorem are precisely
% \mar{ME: Throughout this subsection there is the perennial issue of 
% 	$\Spec R$ vs.\ $\Spf R$.}
those arising from the fibre products $\Spf R_{\rhobar}^\square\times_{\cX_d}
\cX_{d,\red}^{\underline{k}}$, where~$\underline{k}$ runs over the
Serre weights. While Theorem~\ref{thm: intro geom BM without
  multiplicities} is a purely local statement, we do not know how to
prove it without using the stacks~$\cX_d$.

The full geometric Breuil--M\'ezard conjecture
of~\cite{emertongeerefinedBM} makes precise predictions about the
multiplicities of the cycles of the special fibres of
$\Spec R^{\crys,\underline{\lambda}}_{\rhobar}/p$; passing from cycles
to Hilbert--Samuel multiplicities then recovers the original
Breuil--M\'ezard conjecture~\cite{BreuilMezard} (or rather a natural
generalisation of it to~$\GL_d$), which we refer to as the ``numerical
Breuil--M\'ezard conjecture''. In particular, the multiplicities are
expected to be computed in terms of
quantities~$n_{\underline{k}}^\crys(\lambda)$ that are defined as
follows: one associates an irreducible algebraic
representation~$\sigma^{\crys}(\lambda)$ of~$\GL_d/K$ to~$\lambdau$,
defined to have highest weight (a certain shift of)~$\lambdau$. The
semisimplification of the reduction mod~$p$
of~$\sigma^{\crys}(\lambda)$ can be written as a direct sum of
irreducible representations of~$\GL_d(k)$,
and~$n_{\underline{k}}^\crys(\lambda)$ is defined to be the
multiplicity with which the Serre weight~$\underline{k}$ occurs.

In Chapter~\ref{sec: BM} we explain
that as we run over all~$\rhobar$, the geometric Breuil--M\'ezard
conjecture is equivalent to the following analogous conjecture for the
special fibres of our crystalline and semistable stacks. Here by a
``cycle'' in~$\cX_{d,\red}$ we mean a formal $\Z$-linear combination
of its irreducible components~$\cX_d^{\underline{k}}$.

\begin{conj}[Conjecture~\ref{conj: geometric BM}]
  \label{conj: intro version of geometric BM}There are
  cycles~$Z_{\underline{k}}$ in~$\cX_{d,\red}$ with the property that
  for each regular tuple of labeled Hodge--Tate weights~$\lambdau$,
  the underlying cycle of the special fibre
  of~$\cX_d^{\crys,\lambdau}$ is
  $\sum_{\underline{k}}n_{\underline{k}}^\crys(\lambda)\cdot
  Z_{\underline{k}}$.
\end{conj}
In fact we expect that the cycles~$Z_{\underline{k}}$ are effective,
i.e.\ that they are a linear combination of the irreducible
components~$\cX_d^{\underline{k}'}$ with non-negative integer
coefficients. Since there are infinitely many possible~$\lambdau$, the
cycles~$Z_{\underline{k}}$, if they exist, are hugely overdetermined
by Conjecture~\ref{conj: intro version of geometric BM}.

As first explained in~\cite{KisinFM}, the (numerical) Breuil--M\'ezard
conjecture has important consequences for automorphy lifting theorems;
indeed proving the conjecture is closely related to
proving automorphy lifting theorems in situations with arbitrarily
high weight or ramification at the places dividing~$p$. Conversely,
following~\cite{MR2827797}, one can use automorphy lifting theorems to
deduce cases of the Breuil--M\'ezard conjecture. Automorphy lifting
theorems involve a fixed~$\rhobar$, and in fact we can deduce
Conjecture~\ref{conj: intro version of geometric BM} from the
Breuil--M\'ezard conjecture for a finite set of suitably
generic~$\rhobar$.

In particular, we are able to combine results in the literature to
show that for~$\GL_2$ the cycles~$Z_{\underline{k}}$ in
Conjecture~\ref{conj: intro version of geometric BM} must have a
particularly simple form: we necessarily
have~$Z_{\underline{k}}=\cX_{d,\red}^{\underline{k}}$
unless~$\underline{k}$ is a so-called ``Steinberg'' weight, in which
case~ $Z_{\underline{k}}$ is the sum of~$\cX_{d,\red}^{\underline{k}}$
and one other irreducible component. (More precisely, what we show,
following~\cite{CEGSKisinwithdd,geekisin}, is that with these cycles~
$Z_{\underline{k}}$, Conjecture~\ref{conj: intro version of geometric
  BM} holds for all ``potentially Barsotti--Tate'' representations.)

  The weight part of Serre's conjecture predicts the weights in which
  particular Galois representations contribute to the mod~$p$
  cohomology of locally symmetric spaces. Following~\cite{geekisin},
  this conjecture is closely related to the Breuil--M\'ezard
  conjecture; indeed, if Conjecture~\ref{conj: intro version of
    geometric BM} holds, then the set of Serre weights associated to a
  representation $\rhobar:G_K\to\GL_d(\Fpbar)$ should be precisely the
  weights~$\underline{k}$ for which~$Z_{\underline{k}}$ is supported
  at~$\rhobar$. In other words, if we refine our labelling of the
  irreducible components of~$\cX_{d,\red}$ by labelling each component
  by the union of the weights~$\underline{k}$ for which that component
  contributes to~$Z_{\underline{k}}$, then we expect the set of Serre
  weights for~$\rhobar$ to be the union of the sets of weights
  labelling the irreducible components containing~$\rhobar$. This
  expectation holds for~$\GL_2$ by the main results
  of~\cite{CEGSKisinwithdd}.


% which, more intuitively, we can think of as the germ of $\cX_{d,\red}$
% at~$\rhobar$.


  \section{Further questions}\label{subsec: further questions}
  There are many other questions one could ask about the
  stacks~$\cX_d$, which we hope to return to in future
  papers. For example, we show in Proposition~\ref{prop: X is not
    padic formal algebraic} that~$\cX_d$ is not a $p$-adic formal
  algebraic stack. Indeed, if it were $p$-adic formal algebraic, then
  its special fibre would be an algebraic stack, whose dimension would
  be equal to the dimension of its underlying reduced substack. In
  turn, this would imply that the versal rings~$R_{\rhobar}^\square$
  would have dimension equal to the dimensions of the crystalline
  deformation rings~$R_{\rhobar}^{\crys,\lambdau}$, and this is known
  not to be true. In fact, it is a folklore conjecture, recently
  proved by B\"ockle--Iyengar--Pa{\v{s}}k{\=u}nas~\cite{BIP},  that the
  lifting rings~$R_{\rhobar}^\square$ are $\Zp$-flat local complete intersections
  of dimension~$1+d^2+[K:\Qp]d^2$, which should imply that the
  stacks~$\cX_d$ are $\Zp$-flat local complete intersections of
  dimension~$1+[K:\Qp]d^2$ (a notion that we do not attempt to make
  precise for formal algebraic stacks).% \mar{TG: not sure that there is actually a
    %definition of a formal algebraic stack having a dimension, and
    %being a complete intersection; should probably figure this out,
    %and of course check that we can't just prove it.}
  % \mar{TG: add reference to BIP, and maybe think about whether we can
  %   make flatness/lciness precise?}

It is natural to ask about the rigid analytic generic fibre of $\cX_d$; this 
should exist as a rigid analytic stack in an appropriate sense.  The
 generic fibres of the substacks $\cX_d^{\underline{k}}$ should admit 
morphisms to the stacks of Hartl and Hellmann~\cite{2013arXiv1312.6371H}
(although these morphisms won't be isomorphisms, since for any finite 
extension $E$ of~$\Q_p$, the $\cO_E$-points of 
$\cX_d^{\underline{k}}$,  which would coincide with the $E$-points
of its generic fibre, correspond to lattices in crystalline representations,
whereas the stacks of~\cite{2013arXiv1312.6371H} parameterize crystalline
or semistable representations themselves). 

We expect that the~$\cX_d$ will have a role to play in generalisations
of the $p$-adic local Langlands correspondence. For example, we expect
that when~$K=\Qp$ the $p$-adic local Langlands correspondence
for~$\GL_2(\Qp)$ can be extended to give rise to sheaves
of~$\GL_2(\Qp)$-representations on~$\cX_2$. More generally, we expect
 that there will be a $p$-adic analogue of the work of Fargues--Scholze 
on the local Langlands correspondence \cite{fargues2021geometrization} %\cite{MR3966767,fargues2016geometrization} % \mar{ME: Is there
% a  reference yet for Fargues' conj.?}
involving the stacks~$\cX_d$.%\mar{TG: maybe add some (forthcoming) references?}
%\mar{ME: Phrasing? TG: OK now?  ME: Great!}

  % \mar{TG: could think about whether we can show that $\cX_d$ is
  % equidimensional of dimension $d^2[K:\Qp]$; note that this would
  % require defining the dimension of Noetherian formal algebraic
  % stacks, and using versal rings. Note that we don't know in
  % general that the dimension of unrestricted deformation rings, but
  % possibly using our bounds on the codimension of the locus of $H^2$
  % being of any fixed rank we could get around this. There is also the
  % possibility of using globalization and big $R$ equals big $T$.}
% \mar{TG: we should discuss somewhere why it
%   isn't a $p$-adic formal algebraic stack. ME: At least added a discussion
%   about this below (without proof, though).}
% \mar{TG: p-adic LL}

\section{Previous work}
The description of local Galois representations in terms of \'etale $(\varphi,\Gamma)$-modules
is due to Fontaine~\cite{MR1106901}.  The importance of ``height''
as an aspect of the theory was already emphasized in~\cite{MR1106901}, 
and was further developed by Wach~\cite{MR1415732}, who explored 
the relationship between the finite height condition and crystallinity 
of Galois representations in the absolutely unramified context.

The use of what are now called Breuil--Kisin modules as a tool
for studying crystalline and semistable representations 
for general (i.e.\ not necessarily absolutely unramified) $p$-adic fields
(a study which, apart from its intrinsic importance, 
is crucial for treating potentially crystalline or
semistable representations, even in the absolutely unramified context)
was due originally to Breuil~\cite{Breuilunpublished}  %\cite{} \mar{ME: Not sure what the best citation is here.},
and was extensively developed by Kisin~\cite{KisinModularity,
  MR2373358}, who used them to study Galois deformation rings.
%idea of describing Galois representations in terms of what are now called
%\mar{ME: Add Fargues and BMS stuff here. TG: added below, but it could
%move here? ME: Looks good}

The algebro-geometric and  moduli-theoretic perspectives that already played
key roles in Kisin's work were further developed by
Pappas and Rapoport \cite{MR2562795}, who introduced moduli
stacks of Breuil--Kisin modules and of \'etale $\varphi$-modules;
it this work of Pappas and Rapoport, which can be very roughly thought
of as constructing moduli stacks of representations of the absolute Galois
groups of certain perfectoid fields, which is the immediate launching
point for our work in this book, as well as for our %earlier
paper~\cite{EGstacktheoreticimages}.
(We should also mention Drinfeld's work~\cite{MR2181808}, which underpins the
verification of the stack property for the constructions of~\cite{MR2562795},
as well as for those of the present book.) Our use of moduli stacks
of Breuil--Kisin--Fargues modules (in the construction of the
potentially crystalline and semistable substacks) was in part inspired
by the work of Fargues and Bhatt--Morrow--Scholze (see in particular
\cite[\S 4]{2016arXiv160203148B}), which taught us not to be afraid of~$\Ainf$.


Moduli stacks parameterizing crystalline and semistable representations 
have already been constructed by Hartl and Hellmann~\cite{2013arXiv1312.6371H};
as remarked upon above, these stacks should have a relationship
to the stacks $\cX_d^{\underline{k}}$ that we construct. See also the
related papers of Hellmann~\cite{MR3566522,MR3103130}.

As far as we are aware, the first construction of moduli stacks of
representations of~$G_K$ in which the residual
representation~$\rhobar$ can vary is the work of Carl
Wang-Erickson~\cite{MR3831282} mentioned above, which constructs and studies such
stacks in the case that~$\rhobar$ has fixed semisimplification. These
are literally
moduli stacks of representations of~$G_K$; % by reinterpreting them
% as moduli stacks of \'etale $(\varphi,\Gamma)$-modules,
they are isomorphic to certain 
%presumably isomorphic to
substacks of our stacks
%of \'etale $(\varphi,\Gamma)$-modules
$\cX_d$, as we explain in Section~\ref{sec:
  comparison to CWE}.
%\mar{ME: Mention these CWE stacks earlier? TG: now done}

% (and of the
% potentially crystalline and semistable stacks, as appropriate),
% although we have not checked the details of %attempted to prove
% this.   We also remark that
% the existence of 
% %and it is not clear
% %from our constructions that the stacks of~\cite{MR3831282} should exist.
% the stacks of~\cite{MR3831282} (and in particular, their realization
% as stacks of genuine Galois representations) does not seem to follow in any direct
% way from our constructions. \mar{TG: update this paragraph once that
%   section is finished.}

%\mar{fill in ME: Added a bunch of stuff; still haven't gotten to CWE.  It's a bit
%off a laundry-list/very potted history; feel free to edit!}


% Could mention that the existence of a semistable lift is conjectured
% in~\cite[\S5]{MR2745530}. (commenting out as don't want to discuss the
% comparison to their mod $p^n$ conjecture)
\section{An outline of the book}\label{subsec: outline}We finish this
introduction with a brief overview of the contents of this
book. The reader may also wish to refer to the introductions to each
chapter, as well as to the overview of this book provided by the notes~\cite{emerton2020moduli}.  %\mar{TG: update with extra sections we've added.}

In Chapter~\ref{section: coefficient rings} we recall various of the
coefficient rings used in the theories of $(\varphi,\Gamma)$-modules
and Breuil--Kisin modules, and introduce versions of these rings with
coefficients in a $p$-adically complete $\Zp$-algebra. We also prove
almost Galois descent results for projective modules, and deduce
Theorem~\ref{thm:descending projective modules intro version}.

In Chapter~\ref{sec: phi modules and phi gamma modules} we recall the
results of~\cite{EGstacktheoreticimages} on moduli stacks of
$\varphi$-modules, and use them to define our stacks~$\cX_d$ of
\'etale $(\varphi,\Gamma)$-modules. With some effort, we prove
that~$\cX_d$ is an Ind-algebraic stack. Chapter~\ref{sec: crystalline
  and semistable} defines various moduli stacks of Breuil--Kisin and
Breuil--Kisin--Fargues modules, and uses them % , together with the
% machinery of~\cite{EGstacktheoreticimages},
% \mar{ME: As I already commented  above, I'm not sure that we  actually
% do use this machinery in \S 4}
to construct our moduli
stacks of potentially semistable and potentially crystalline
representations, and in particular to prove Theorem~\ref{thm:intro
  statement on crystalline moduli}.

Chapter~\ref{sec: families of extensions} develops the theory of the
Herr complex, proving in particular that it is a perfect complex and
is compatible with base change. We show how to use the Herr complex to
construct families of extensions of $(\varphi,\Gamma)$-modules, and we
use these families to define the irreducible
substack~$\cX_{d,\red}^{\underline{k}}$ corresponding to a Serre
weight~$\underline{k}$. By induction on~$d$ we prove that~$\cX_d$ is a
Noetherian formal algebraic stack, and establish a version of
Proposition~\ref{prop:H2 dimension} (although as discussed in
Remark~\ref{rem: dimension and codimension}, we do not prove
Proposition~\ref{prop:H2 dimension} as stated at this point in the
argument). 

It may help the reader for us to point out
 that Chapters~\ref{sec: crystalline and semistable}
and~\ref{sec: families of extensions} are essentially independent of one another,
and are of rather different flavours.  %the former
Chapter~\ref{sec: crystalline and semistable} involves an interleaving of stack-theoretic
arguments with ideas from $p$-adic Hodge theory and the theory of Breuil--Kisin
modules, while in %the latter,
Chapter~\ref{sec: families of extensions},
once we complete our analysis of the Herr complex,
our perspective begins to shift: although at a technical level we of course continue
to work with $(\varphi,\Gamma)$-modules, we begin to think in terms 
of Galois representations and Galois cohomology, and the more foundational
 arguments of the preceding chapters recede somewhat into the background.
 
In Chapter~\ref{sec:properties} we combine the results of
Chapters~\ref{sec: crystalline and semistable} and ~\ref{sec: families
  of extensions} with a geometric argument on the local
deformation ring to prove Theorem~\ref{thm:intro crystalline
  lifts}. Having done this, we are then able to improve on the results
on~$\cX_d$ established in the earlier chapters by proving
Theorem~\ref{thm:intro statement of basic properties of X}. We also
deduce Theorem~\ref{thm: intro existence of global lifts}, as well as
determining the closed points of~$\cX_d$, and describing the
relationship of our stacks with Wang--Erickson's stacks of Galois representations.

Chapter~\ref{sec: the rank one case} gives explicit descriptions of
various of our moduli stacks in the case~$d~=~1$, relating them to
moduli stacks of Weil group representations. Chapter~\ref{sec: BM}
explains our geometric version of the Breuil--M\'ezard conjecture, and
proves some results towards it, particularly in the case~$d=2$.

Finally the appendices for the most part establish various technical
results used in the body of the book. We highlight in particular
Appendix~\ref{app: formal algebraic stacks}, which summarises the
theory of formal algebraic stacks developed
in~\cite{Emertonformalstacks}, and Appendix~\ref{app: BKF pst}, which
combines the theory of Breuil--Kisin--Fargues modules with Tong Liu's
theory of $(\varphi,\Ghat)$-modules to give a new characterisation of
integral lattices in potentially semistable representations, of which we
make crucial use %of
in Chapter~\ref{sec: crystalline and semistable}.

\section{Acknowledgements}\label{subsec: acknowledgements}We would like to thank Robin Bartlett,
Laurent Berger, Bhargav Bhatt, Xavier Caruso, Pierre Colmez, Yiwen
Ding, Florian Herzig, Ashwin
Iyengar, Mark Kisin, Tong Liu, George Pappas, Dat Pham, L\'eo Poyeton, Michael Rapoport, and
Peter Scholze for helpful correspondence and conversations.%  \mar{TG:
  % thank Dat Pham in text for correction.}
We would particularly like to
thank Tong Liu for his contributions to Appendix~\ref{app: BKF
  pst}. We would also like to thank the organisers (Johannes Ansch\"utz,
Arthur-C\'esar Le Bras, and Andreas Mihatsch) %and participants
of the 2019 Hausdorff School on ``the Emerton--Gee stack and related
topics'', as well as all the participants,
both for the encouragement to finish this book and for the
many helpful questions and corrections resulting from our lectures. 

Our mathematical debt to the late Jean-Marc Fontaine will be obvious
to the reader. This book benefited from several conversations with
him over the years~2011-2018, %last 8 years,
%\mar{ME: Is the right range of years? TG: yes, commenting out}
and from the interest he showed in our
results; in particular, his explanations to us of the relationship
between framing Galois representations and Fontaine--Laffaille modules
during his visit to Northwestern University in the spring of 2011 provided
an important clue as to the correct definitions of our stacks.

We owe special thanks to Colette and Therese for all of their patience and
support during the writing and revision of this book.

% \mar{TG: if we want to, we could also acknowledge our
%   spouses/families? TG: made a terrible attempt at this.  ME: Didn't look
% too bad; just jazzed it up a little to make it less repetitive}

The first author was supported in part by the
  NSF grants DMS-1303450, DMS-1601871, and DMS-1902307. The second author was 
  supported in part by a Leverhulme Prize, EPSRC grant EP/L025485/1, Marie Curie Career
  Integration Grant 303605,
  ERC Starting Grant 306326, ERC Advanced Grant 884596, and a Royal Society Wolfson Research
  Merit Award. This project has received funding from the European Research Council (ERC) under the European Union’s Horizon 2020 research and innovation programme (grant agreement No. 884596).
% for
% a helpful conversation about his paper~\cite{MR3322781}.

% \section{Old intro}
% % Our goal in this paper is to construct moduli spaces
% % of $p$-adic representations of absolute Galois groups of $p$-adic local fields.
% % These spaces will be formal algebraic stacks, % \mar{ME: Perhaps actually
% % % 	just Ind-algebraic stacks? TG: seems to be formal algebraic
% % %         stacks, in fact? ME: As long as the cotangent complex argument
% % % hangs together \dots }
% % and points 
% % in the special fibre will correspond to residual Galois representations.
% % The formal completion of one of our stacks at such a point will then recover
% % the usual formal deformation space of this residual representation.
% % Thus the relationship between the theory and constructions
% % that we develop here, and the usual formal deformation theory of Galois
% % representations, is the same as that between the theory of moduli spaces
% % of algebraic varieties, and the formal deformation theory of algebraic 
% % varieties: the latter gives valuable local information about the former,
% % but moduli spaces, when they can be constructed,
% % capture global aspects of the situation which are
% % inaccessible by the purely % infinitesimal tools of formal deformation theory.

% % \section{Statement of results}
% % In order to explain our results in more detail, we introduce some notation.
% % We fix a finite extension $K$ of $\Q_p$ (for some prime $p$),
% % and a positive integer $d$ (the dimension of the representations being
% % parameterized).  Our main construction is then a Noetherian
% % formal algebraic stack $\cX_d$,
% % equipped with a  morphism to $\Spf \Z_p$, with the following properties:
% \begin{enumerate}
% 	%\item $\cX_d$ is flat and topologically of finite type over $\Spf
% 	%	\Z_p$,
% 	%	\mar{ME: I've put flat here, since it seems to simplify the
% 	%		dimension statement, but it may not be the correct
% 	%		thing to do if we don't know that unrestricted local
% 	%		def.\ rings are flat in general.}
% 	%	\mar{ME: Investigate the correct finite type statement.}
% 	%	of relative dimension $[K:\Q_p] d^2$.
% \item
% 	The underlying reduced substack $\cX_{d,\red}$ is of finite type over
% 	$\F_p$, and is equidimensional of dimension $[K:\Qp]d(d-1)/2$.
% 	\mar{ME: I've rewritten this to avoid flatness assumptions.  We should add Noetherian, or maybe locally Noetherian, here, depending on what we eventually prove.}
% 	\item If $\F$ is an algebraic extension of $\F_p$ (in
%           particular, $\Fpbar$ or
% 		a finite extension of $\F_p$), % \mar{TG: should ``in
%                   % particular'' be ``i.e.''? ME: There are other possibilities, since
% 		  % $\widehat{\mathbb Z}$ has lots of proper closed
%                   % subgroups that are of infinite index. TG: duh,
%                   % thanks, commenting out.}
% 		then the groupoid of points $\cX_d(\F)$ is naturally equivalent
% 		\mar{ME: Think through the use of ``natural'' in various
% 			statements, to avoid accidentally invoking $3$-categories.}
% 		to the groupoid of continuous representations
% 		$G_K \to \GL_d(\F).$
% 	\item If $\F$ is a finite extension of $\F_p$, if\mar{TG: we
%             could do~$\Fpbar$ here if we want?}
% 		$\Spec \F \to \cX_d$ is a morphism, corresponding (by
% 		the previous point) to a continuous representation
% 		$\rhobar: G_K \to \GL_d(\F)$, and if $R^{\square}_{\rhobar}$
% 		denotes the universal lifting $W(\F)$-algebra %  \mar{TG: might need to
%                   % elaborate - I guess the lifting $W(k)$-algebra}
%                 of $\rhobar$,
% 		then there is a canonical morphism (of formal stacks
% 		over $\Spf \Z_p$) 
% 		% \mar{ME: If we don't know that these unrestricted deformation rings
% 		% 	 are flat in general, then we should either take
% 		% 	 the flat part, or allow $\cX_d$ to be non-flat.}
% 		$\Spf R^{\square}_{\rhobar} \to \cX_d$ extending 
% 		the given morphism on $\Spec \F$; this morphism is versal. 
% 		This last property can be made more precise: namely, the
% 		fibre product $\Spf R_{\rhobar}^{\square} \times_{\cX_d}
% 		\Spf R_{\rhobar}^{\square}$ is canonically isomorphic to
% 		the groupoid $\Spf R_{\rhobar}^\square \times \widehat{\GL}_d 
% 	\rightrightarrows \Spf R_{\rhobar}^\square$ given by change of basis.
% \end{enumerate}

% From point~(3) it follows formally that if $\cO$ is the ring
% of integers in any finite extension of $\Q_p$
% then the groupoid $\cX_d(\cO)$ (i.e.\ the groupoid of morphisms of
% formal algebraic
% stacks $\Spf \cO \to \cX_d$ over $\Spf \Z_p$) is canonically equivalent
% to the groupoid of continuous representations $G_K \to \GL_d(\cO)$.
% Similarly, the groupoid $\cX_d(\Zbar_p)$ is canonically equivalent
% to the groupoid of continuous representations
% $G_K \to \GL_d(\Zbar_p)$.
% More generally, if $A$ is any finite $\Z_p$-algebra,
% then the groupoid $\cX(A)$ is canonically equivalent
% to the groupoid of continous representations $G_K \to \GL_d(A)$.\mar{ME: This last assertion may not be true if we have to pass to flat parts, and if $A$
% 	is not flat over $\Z_p$; instead we would just get a certain
% 	subgroupoid.}

%As usual, we let $\cX_{\red}$ denote the reduced algebraic stack
%underlying the formal algebraic stack $\cX$.  It follows
%from the first
%point above that $\cX_{\red}$ is of finite type
%over $\F_p$, but we have the following more precise statement.
%
%\begin{enumerate}
%		\setcounter{enumi}{3}
%	\item $\cX_{\red}$ is equidimensional of dimension $d(d-1)/2$ over $\F_p$.
%\end{enumerate}

% Suppose now that $\underline{\lambda}$ is some $d$-tuple of labeled Hodge--Tate
% weights, and that $\tau$ is some inertial type.% \mar{TG: define this,
%   % and the notion of regular weight, in the notation section.}

% \begin{enumerate}
% 		\setcounter{enumi}{3}
% 	\item There is a closed substack $\cX_d^{\underline{\lambda},\tau,\crys}$
% 		of $\cX_d$ which is %finite type $p$-adic formal algebraic stack,
% 		flat over $\Spf \Z_p$, \mar{ME: There is a question of whether
% 			flatness makes sense, because this might require
% 			Noetherianness of the formal stack!}
% 		and which is
% 		characterized by the property that a $\Zbar_p$-valued point
% 		of $\cX_d$ lies in $\cX_d^{\underline{\lambda},\tau,\crys}$ if
% 		and only if the associated Galois representation
% 		$G_K \to \GL_d(\Zbar_p)$, when thought of as a representation
% 		into $\GL_d(\Qbar_p)$, is potentially crystalline of inertial
% 		type $\tau$ and with labeled Hodge--Tate weights 
% 		$\underline{\lambda}$.   Each of the underlying reduced algebraic
% 		substacks
% 		$(\cX_d^{\underline{\lambda},
% 			\tau,\crys})_{\red}$ is equidimensional of dimension an explicit function
%                       of~$\underline{\lambda}$, which is less than or equal
%                       to~$[K:\Qp]d(d-1)/2$, with equality if and only
%                       if $\underline{\lambda}$ is regular. The analogous
%                       statements hold for ``semi-stable'' in place of
%                       ``crystalline''. %  \mar{ME: Expected dimension
% 		% 	here; $[K:\Qp]d(d-1)/2$ when $\underline{\lambda}$ is regular.}
% 		% Similar statement for semi-stable in place of crystalline.
% \end{enumerate}

% \begin{remark}
% We expect that each of the formal algebraic
% stacks $\cX_d^{\underline{\lambda},\tau,\crys}$ is in fact a $p$-adic formal
% algebraic stack of finite type, but we are able to prove this only
% in the case when $K$ is unramified over $\Q_p$.
% \end{remark}

% It follows from~(1) and~(4) that if $\underline{\lambda}$ is regular,
% then $\cX^{\underline{\lambda}}_{d,\red}$ is a union of irreducible components
% of $\cX_{d,\red}$.  It follows from~(3) and~(4) that
% the pull-back of $\cX_d^{\underline{\lambda},\tau}$ under a versal morphism
% $\Spf R_{\rhobar}^{\square} \to \cX_d$ as in~(3) is equal to the usual
% potentially crystalline lifting space $\Spf R_{\rhobar}^{\underline{\lambda},
% 	\tau,\crys}.$

      
% There are some additional structures that $\cX_d$ carries.
% For example:
% \begin{enumerate}
% 		\setcounter{enumi}{3}
% 	\item
% There is a perfect complex of $\cO_{\cX_d}$-modules
% $\cC^{\bullet}$ on $\cX_d$,
% with cohomology lying in degrees $0$, $1$, and $2$, such that
% if $\Spf A \to \cX_d$ is any morphism with $A$ being a finite $\Z_p$-algebra,
% corresponding to a continuous
% Galois representation $\rho: G_K \to \GL_d(A)$,
% then the base-changed complex $A\otimes_{\cO_{\cX_d}} \cC^{\bullet}$ 
% computes the continuous Galois cohomology $H^{\bullet}(G_K, \rho)$.
% \end{enumerate}

% Furthermore, there are morphisms
% relating the various $\cX_d$ for differing values of $d$
% corresponding to multi-linear algebra constructions on Galois representations.
% For example, there are morphisms
% $$\cX_d \times_{\Z_p} \cX_{d'} \to \cX_{d + d'}$$
% and 
% $$\cX_d \times_{\Z_p} \cX_{d'} \to \cX_{d  d'}$$
% corresponding to direct sum and tensor product respectively.

% We may also use classes in $H^1$ to construct families
% of extensions.  More precisely:

% \begin{enumerate}
% 		\setcounter{enumi}{4}
% \item Suppose that
% 	$f_d: \cZ \to \cX_d$ is any morphism from a formal algebraic stack
% 	to $\cX_d$, and that $\phi: \cF \to H^1(f_d^*\cC^{\bullet})$ is any morphism
% 	of $\cO_{\cZ}$-modules
% 	whose domain is a locally free sheaf of locally finite rank over $\cZ$.
% 	If $\cV \to \cX_d$ is the geometric vector bundle associated
% 	to $\cF$ (so $\cF$ is the sheaf of sections of~$\cV$),
% 	then there is a morphism $f_{d+1}: \cV \to \cX_{d+1}$.
% 	On $A$-valued points, for a finite $\Z_p$-algebra $A$,
% 	this morphism has the following interpretation:
% 	giving a morphism $\Spf A \to \cV$ is equivalent to giving
%         a morphism $g: \Spf A \to \cZ$ together with
% 	a class $c \in A\otimes_{\cO_{\cZ}} \cF$; composing the
% 	morphism  $g$ with $f_d$ yields a morphism
% 	$\Spf A \to \cX_d$, corresponding to a continuous representation
% 	$\rho: G_K \to \GL_d(A)$; applying (the base-change to $A$ of)
% 	$\phi$ to the class $c$ yields a class in $H^1(G_K,\rho),$
% 	which in turn corresponds to a Galois representation $\rho'$
% 	which sits in an exact sequence $0 \to \rho \to \rho' \to \triv \to 0;$
% 	then $f_{d+1}$ is constructed in such a way that
% 	the composite $\Spf A \to \cV \buildrel f_{d+1} \over \longrightarrow
%        	\cX_{d+1}$ is precisely
% 	the morphism that classifies the representation
% 	$\rho'$.% \mar{TG: will we actually need something more general
%           than this when we want to prove existence of lifts in
%           general niveau? ME: In principle, but actually, if we have
%   a reasonable multilinear algebra formalism relating the various $\cX_d$'s,
% then the formula $\Ext^i(\rho,\psi) = H^i(\psi\otimes\rho^{\vee})$ should
% let us reduce to this case (which just has the technical advantage
% of not having to study any complexes beyond the Herr complex).  But you might
% be right, that it could be better to have this reduction from $\Ext^i$ to $H^i$
% to take place somewhere in the technical development, and to have
% as part of the formalism
% a complex over the product $\cX_{d_1} \times \cX_{d_2}$ which computes
% $\Ext^i$ between the universal families.}
%\end{enumerate}

% \section{The underlying reduced substack $\cX_{d,\red}$,
% 	and the geometric Breuil--M\'ezard philosophy (to be revised)}
% \label{subsec:intro Breuil--Mezard}

% \section{Crystalline lifts and the support of $H^2(\rhobar)$}
% \label{subsec:intro crystalline lifts}
% As an application of our constructions, % the geometric Breuil--M\'ezard ideas discussed
% % in the preceding subsection,
% % our construction,
% we prove the following result.

% \begin{thm}
% 	\label{thm:intro crystalline lifts}
% 	If $\rhobar: G_K \to \GL_d(\Fbar_p)$ is a continuous
% 	representation, then $\rhobar$
% 	has a lift $\rho^{\circ}: G_K \to \GL_d(\Zbar_p)$ for which
% 	the associated $p$-adic representation $\rho:G_K \to \GL_d(\Qbar_p)$
% 	is crystalline of regular Hodge--Tate weights. 

% We can furthermore choose
%         these Hodge--Tate weights to be in the
%         interval~$[0,dp-1]$, % \mar{TG: not too sure about this - might be
%           % better to do bounded gaps instead, although obviously a
%           % result of this kind will be true.}
% 	% {\em (}the bound depending only on $d$ and $K${\em )},
%         or we
%         can choose them to have arbitrarily large gaps, and in either
%         case we can also ensure that~$\rho^{\circ}$ is potentially diagonalizable.
% % 	\mar{ME: We could consider making the bound more explicit, either
% % 		in the statement of the result or in a subsequent remark.}
% % 	\mar{ME: Maybe add potentially diagonalizable into the statement,
% % 		or perhaps discuss it afterwards, if it seems a bit
% % 		out of place in the statement at this particular
% %                 moment. TG: also want well-spread, rather than
% %                 bounded, as an option. ME: Currently, can get any plausible
% % 	property of the HT weights, but can't get pot'l diagonalizable ---
% % have to work on this.}
% \end{thm}
% The notion of a potentially diagonalizable representation was defined
% in~\cite{BLGGT}, and is important for applications to automorphy
% lifting theorems; for example, Theorem~\ref{thm:intro crystalline
%   lifts} allows us to remove a hypothesis on the construction of a
% candidate $p$-adic Langlands correspondence in~\cite{Gpatch}. In
% addition to proving that the Hodge--Tate weights can be taken  to be in the
%         interval~$[0,dp-1]$, % the
% % case of lifts of bounded weight, \mar{ME: What does ``In the case of lifts of
% % 	bounded weight'' mean? TG: hopefully this is clearer.}
% % % \mar{TG: discuss this later in the paper}
% we are 
% able to prove a conjecture of~\cite{2015arXiv150902527G} on the
% existence of a crystalline lift whose Hodge--Tate weights correspond
% to a Serre weight.

% \mar{ME: Maybe explicitly state strengthend version of globalization result
% 	from appendix to Geom.\ Breuil--M\'ezard?}

% The problem of constructing crystalline lifts has been considered by various
% authors, in particular~\cite{MullerThesis,2015arXiv150601050G}.
% It admits a well-known inductive approach (which is taken in~\cite{MullerThesis,2015arXiv150601050G}): one writes $\rhobar$ as 
% a successive extension of irreducible representations, lifts each of these
% irreducible representations 
% to a crystalline representation, and then attempts to lift the various
% extension classes.  The difficulty that arises in this approach (which
% has proved an obstacle to obtaining general statements along the lines
% of Theorem~\ref{thm:intro crystalline lifts} up until now) is showing that the 
% mod $p$ extension classes that appear in this description of $\rhobar$ 
% can actually be lifted to crystalline extension classes in characteristic
% zero.  The basic source of the difficulty is that the local Galois $H^2$
% can be non-zero, and non-zero classes in $H^2$ obstruct the lifting 
% of extension classes (which can be interpreted as classes lying in
% $H^1$). (In fact, the difficulty is not so much in obtaining
% crystalline extension classes, as in lifting to any classes in
% characteristic zero; indeed, as far as we are aware, until now
% it was not known that an arbitrary~$\rhobar$ had \emph{any} lift to
% characteristic zero at all.)

% Our proof of Theorem~\ref{thm:intro crystalline lifts} relies on the
% inductive strategy described in the preceding paragraph, but we are able to
% prove the following key result, which controls the obstructions that
% can be presented by $H^2$.

% \begin{prop}
% 	\label{prop:H2 dimension}
% 	The locus of points $\rhobar \in \cX_{d,\red}(\Fbar_p)$ 
% 	at which $\dim H^2(G_K,\rhobar) \geq r$ is Zariski closed
% 	in $\cX_{d,\red}(\Fbar_p)$, of codimension $\geq r$.
% \end{prop}

% If we let $\rho^{\univ}$ denote the universal lifting of $\rhobar$
% over $R^{\square}_{\rhobar}$, then $H^2(G_K, \rho^{\univ})$ is 
% an $R^{\square}_{\rhobar}$-module, and Proposition~\ref{prop:H2 dimension}
% implies the following corollary.

% \begin{cor}
% 	\label{cor:H2 dimension}
% 	For any regular tuple of labeled Hodge--Tate weights $\underline{\lambda}$
% 	and inertial type $\tau$, the locus of points 
% 	$x\in \Spec R^{\underline{\lambda},\tau,\crys}/p$ for which
% 	$$\dim_{\kappa(x)}
% 	\kappa(x)\otimes_{R^{\square}_{\rhobar}} H^2(G_K,\rho^{\univ})
% 	\geq r$$
% 	has codimension $\geq r$.
% \end{cor}

% \begin{remark}
% 	\label{rem:H2 dimension}
% 	By Tate local duality, 
% 	$$\dim_{\kappa(x)}
% 	\kappa(x)\otimes_{R^{\square}_{\rhobar}} H^2(G_K,\rho^{\univ})
% 	=
% 	\dim_{\kappa(x)}
% 	\kappa(x)\otimes_{R^{\square}_{\rhobar}}
% 	\Hom_{G_K}\bigl((\rho^{\univ})^{\vee},\omega\bigr)$$
% 	(where $\omega$ denotes the mod $p$ cyclotomic character, thought
% 	of as taking values in $\kappa(x)^{\times}$).
% 	Thus the statement of Corollary~\ref{cor:H2 dimension} is
% 	related to the way in which the special fibre of the potentially crystalline
% 	lifting ring $R_{\rhobar}^{\underline{\lambda},\tau,\crys}$ intersects
% 	the reducibility locus in $\Spec R^{\square}_{\rhobar}$.
% \end{remark}

% Given Corollary~\ref{cor:H2 dimension},
% we prove Theorem~\ref{thm:intro crystalline
% 	lifts}
% by working purely within the context of formal lifting rings.  
% However we don't know how to
% prove the corollary while staying within that context.
% Indeed, as Remark~\ref{rem:H2 dimension} indicates, 
% this corollary is related to the way in which the special fibre
% of a potentially crystalline deformation ring intersects another natural
% locus in $\Spec R_{\rhobar}^{\square}$ (namely, the reducibility locus).
% Since the special fibre of a potentially crystalline lifting ring
% is not directly defined in deformation-theoretic terms, such questions
% are notoriously difficult to study directly.
% Our proof of the corollary proceeds
% differently, by using the  geometric Breuil--M\'ezard view-point
% %discussed in Subsection~\ref{subsec:intro Breuil--Mezard} above
% to replace a computation on the special fibre of the potentially
% crystalline deformation ring by a compution on $\cX_{d,\red}$;
% this latter space has a concrete description in terms of families of varying
% $\rhobar$, whose $H^2$ we are  able to compute.

% \section{The construction of $\cX_d$}
% We construct $\cX_d$ as the moduli stack of rank $d$ \'etale 
% $(\varphi,\Gamma)$-modules.  \mar{ME: Elaborate on this.}

\section{Notation and conventions}\label{subsec: notation and conventions}
\subsection*{$p$-adic Hodge theory}Let $K/\Qp$ be a finite
extension. If $\rho$ is a de Rham representation of $G_K$ on a
$\Qpbar$-vector space~$W$, then we will write $\WD(\rho)$ for the
corresponding Weil--Deligne representation of $W_K$ (see e.g.\
\cite[App. B]{CDT}), and if $\sigma:K \into \Qpbar$ is a continuous
embedding of fields then we will write $\HT_\sigma(\rho)$ for the
multiset of Hodge--Tate numbers of $\rho$ with respect to $\sigma$,
which by definition contains $i$ with multiplicity
$\dim_{\Qpbar} (W \otimes_{\sigma,K} \widehat{\barK}(i))^{G_K} $. Thus
for example if~$\epsilon$ denotes the $p$-adic cyclotomic character,
then $\HT_\sigma(\epsilon)=\{ -1\}$.
%\mar{ME: I think that we use $\epsilon$ at various points below, rather than
%$\varepsilon,$ for the $p$-adic cyclotomic character. And we also
%use $\varepsilon$ for the usual element of $A_{\inf}$ (which is maybe why
%we shifted to $\epsilon$ for the cyclotomic character, although not consistently).
%Should make this consistent. TG: I believe I've fixed this and went
%with $\epsilon$ for cyclo char. We also use $\epsilon$ for a small
%positive number internally to the proof of Lemma~\ref{lem:
%  amplification from char p to char 0} but I'm happy with that if you are.
%ME: Yep, I think this  is  the right  choice (and  using $\epsilon$ in 
%a couple of discrete locations  for a small number and for dual numbers
%seems fine to me).}

By a \emph{$d$-tuple of labeled Hodge--Tate
  weights~$\underline{\lambda}$}, \index{labeled Hodge--Tate
  weights} \index{$\underline{\lambda}$}
we mean a tuple of integers~$\{\lambda_{\sigma,i}\}_{\sigma:K\into\Qpbar,1\le i\le d}$
with $\lambda_{\sigma,i}\ge \lambda_{\sigma,i+1}$ for all~$\sigma$ and all $1\le
i\le d-1$. We will also refer to~$\underline{\lambda}$ as a
\index{Hodge type} \index{inertial type}
\emph{Hodge type}. By an \emph{inertial type~$\tau$} we mean a representation
$\tau:I_K\to\GL_d(\Qpbar)$ which extends to a  representation
of~$W_K$ with open kernel (so in particular, $\tau$ has finite
image). 

Then we say that~$\rho$ has Hodge type~$\underline{\lambda}$ (or
labeled Hodge--Tate weights~$\underline{\lambda}$) if for
each~$\sigma:K\into\Qpbar$ we
have~$\HT_\sigma(\rho)=\{\lambda_{\sigma,i}\}_{1\le i\le
  d}$, and we say that~$\rho$ has inertial
type~$\tau$ if~$\WD(\rho)|_{I_K}\cong\tau$.

We often somewhat abusively write that a representation
$\rho:G_K\to\GL_d(\Zp)$ is crystalline (or potentially crystalline, or
semistable, or\dots) if the corresponding
representation $\rho:G_K\to\GL_d(\Qp)$ is crystalline (or potentially
crystalline, or\dots).

\subsection*{Serre weights and Hodge--Tate weights}By a \emph{Serre
weight} $\underline{k}$ we mean a tuple of \index{Serre weight} \index{\underline{k}}
integers~$\{k_{\sigmabar,i}\}_{\sigmabar:k\into\Fpbar,1\le i\le d}$
with the properties that \begin{itemize}
\item $p-1\ge k_{\sigmabar,i}-k_{\sigmabar,i+1}\ge 0$ for each $1\le i\le
  d-1$, and
\item $p-1\ge k_{\sigmabar,d}\ge 0$, and not every~$k_{\sigmabar,d}$
  is equal to~$p-1$.
\end{itemize}
The set of Serre weights is in bijection with the set of irreducible
$\Fpbar$-representations of $\GL_d(k)$, via passage to highest
weight vectors (see for example the appendix to~\cite{MR2541127}). 

% If $\underline{k}$ is a Serre weight for which
% $k_{\sigmabar,1}-k_{\sigmabar,2}=p-1$ for all~$\sigmabar$ then we
% say that $\underline{k}$ is \emph{tr\`es ramifi\'ee}.
% %we are in the \emph{tr\`es ramifi\'ee case}.
% \mar{TG: not sure if we want this terminology yet. ME: I've moved
% this definition here from \S 3. TG: unlikely that we want to use it in
% this form (as it's only talking about one of the gaps in weights) but
% leaving it for now.} 

% We now recall from~\cite{2015arXiv150902527G} the notion of a
% crystalline lift of some Serre weight.
Each embedding~$\sigma:K\into\Qpbar$ induces an
embedding~$\sigmabar:k\into\Fpbar$; if $K/\Qp$ is ramified, then
each~$\sigmabar$ corresponds to multiple embeddings~$\sigma$.  We say
that~$\underline{\lambda}$ is a lift of  ~$\underline{k}$ % is a \mar{TG: here need to fix~$k$ versus
  % $\lambda$} \emph{Serre weight}
if for each
embedding~$\sigmabar:k\into\Fpbar$, we can choose an
embedding~$\sigma:K\into\Qpbar$ lifting~$\sigmabar$, with the
properties that:
\begin{itemize}
\item $\lambda_{\sigma,i}=k_{\sigma,i}+d-i$, and
\item if $\sigma':K\into\Qpbar$ is any other lift of~$\sigmabar$,
  then $k_{\sigma',i}=d-i$.
\end{itemize}


\subsection*{Lifting rings}Let $K/\Qp$ be a finite extension, and
let $\rhobar:G_K\to\GL_d(\Fpbar)$ be a continuous representation. Then
the image of~$\rhobar$ is contained in~$\GL_d(\F)$ for any
sufficiently large finite
extension~$\F/\F_p$. Let~$\cO$ be the ring of integers in some
finite extension~$E/\Qp$, and suppose that the residue field of~$E$
is~$\F$. % Fix an ordered basis for the underlying $\F$-vector space
%\mar{ME: (a) Isn't a basis ordered by definition? (b) If $\rhobar$
%maps to $\GL_d$  then haven't we already chosen a basis? (c) Can't
%we just say that $R^{\square, \cO}_{\rhobar}$ parametrerizes lifts
%of $\rhobar$ to $\rho: G_K \to \GL_d(A)$? TG: (a) no! (b) yes (c)
%how's this? ME: Great}
% of~$\rhobar$, and
Let~$R^{\square,\cO}_{\rhobar}$ be the universal \index{$R^{\square,\cO}_{\rhobar}$}
lifting~$\cO$-algebra of~$\rhobar$; by definition, this (pro-)represents the
functor given by lifts of~$\rhobar$ to representations $\rho: G_K \to
\GL_d(A)$, for~$A$ an Artin local $\cO$-algebra with residue field~$\F$. % \mar{TG: might be better to
  % use~$W(\Fpbar)$ coefficients (or rather finite extensions of that?)}
The precise choice of~$E$ is unimportant, in the sense that
if~$\cO'$ is the ring of integers in a finite extension~$E'/E$, then
by~\cite[Lem.\ 1.2.1]{BLGGT} we have
$R^{\square,\cO'}_{\rhobar}=R^{\square,\cO}_{\rhobar}\otimes_{\cO}\cO'$.
% so
% that \[R^{\square,\Zpbar}_{\rhobar}:=R^{\square,\cO}_{\rhobar}\otimes_{\cO}\Zpbar\]is
% well-defined, independently of the choice of~$\cO$. By definition, the
% $\Zpbar$-points of $\Spf R^{\square,\Zpbar}_{\rhobar}$ correspond
% precisely to the liftings of~$\rhobar$ to a
% representation~$G_K\to\GL_d(\Zpbar)$.\mar{TG: I didn't think this
%   through super carefully, so hopefully it's not a blunder.}

% By a \emph{$d$-tuple of labeled Hodge--Tate weights~$\underline{\lambda}$},
% we mean integers~$\{\lambda_{\sigma,i}\}_{\sigma:K\into\Qpbar,1\le i\le d}$
% with $\lambda_{\sigma,i}\ge \lambda_{\sigma,i+1}$ for all~$\sigma$ and all $1\le
% i\le d-1$. By an \emph{inertial type~$\tau$} we mean a representation
% $\tau:I_K\to\GL_d(\Qpbar)$ which extends to a  representation
% of~$W_K$ with open kernel (so in particular, $\tau$ has finite
% image).

Fix some Hodge type~$\underline{\lambda}$ and inertial
type~$\tau$. If~$\cO$ is chosen large enough that the inertial
type~$\tau$ is defined over~$E=\cO[1/p]$, and large enough that~$E$
contains the images of all embeddings~$\sigma:K\into\Qpbar$, then we
have the usual lifting $\cO$-algebras
~$R^{\crys,\underline{\lambda},\tau,\cO}_{\rhobar}$ \index{$R^{\crys,\underline{\lambda},\tau,\cO}_{\rhobar}$}
and~$R^{\semis,\underline{\lambda},\tau,\cO}_{\rhobar}$. By \index{$R^{\semis,\underline{\lambda},\tau,\cO}_{\rhobar}$}
definition, these are the unique $\cO$-flat quotients
of~$R_{\rhobar}^{\square,\cO}$ with the property that if~$B$ is a
finite flat~$E$-algebra, then an $\cO$-algebra homomorphism
$R_{\rhobar}^{\square,\cO}\to B$ factors through ~$R^{\crys,\underline{\lambda},\tau,\cO}_{\rhobar}$
(resp.\ through~$R^{\semis,\underline{\lambda},\tau,\cO}_{\rhobar}$)
if and only if the corresponding representation of~$G_K$ is
potentially crystalline (resp.\ potentially semistable) of Hodge
type~$\underline{\lambda}$ and inertial type~$\tau$. If~$\tau$ is
trivial, we will sometimes omit it from the notation.
% corresponds 
% to lifts of~$\rhobar$ which are crystalline with labeled Hodge--Tate
% weights~$\underline{\lambda}$. % and inertial
% type~$\tau$.
%By definition, these rings is $\cO$-flat, % \mar{TG:
  % check if they are really reduced by definition}
By
the main theorems of~\cite{MR2373358}, % it has regular generic fibre, 
these rings are (when they are nonzero) equidimensional of
dimension \[1+d^2+\sum_{\sigma}\#\{1\le i<j\le
  d|\lambda_{\sigma,i}>\lambda_{\sigma,j}\}.\]% \mar{TG: hopefully correct, didn't
% think about it too carefully.}
Note that this quantity is at most $1+d^2+[K:\Qp]d(d-1)/2$, with
equality if and only if~$\underline{\lambda}$ is \emph{regular}, in
the sense that $\lambda_{\sigma,i}> \lambda_{\sigma,i+1}$ for all~$\sigma$ and all $1\le
i\le d-1$. As above, we have
$R^{\crys,\underline{\lambda},\tau,\cO'}_{\rhobar}=R^{\crys,\underline{\lambda},\tau,\cO}_{\rhobar}\otimes_{\cO}\cO'$,
and similarly for~$R^{\semis,\underline{\lambda},\tau,\cO'}_{\rhobar}$. % Since~$R^{\crys,\underline{\lambda},\cO}_{\rhobar}$
% is~$\cO$-flat and has regular (and thus reduced) generic fibre, it is
% in particular reduced, and it can be characterised as the unique~$\cO$-flat and reduced quotient of~$R_{\rhobar}^{\square,\cO}$ with the property that if
% for all~$\cO'$ as above, a morphism of~$\cO$-algebras
% $R_{\rhobar}^{\square,\cO}\to\cO'$ factors
% through~$R^{\crys,\underline{\lambda},\cO}_{\rhobar}$ if and only if
% the corresponding representation $G_K\to\GL_d(E')$ is crystalline with
% labeled Hodge--Tate weights~$\underline{\lambda}$. 
% \mar{TG: I think
  % this characterisation is correct, but just flagging it in case I'm
  % confusing myself, as the original definition is in terms of not
  % necessarily reduced $E$-algebras.}
% set \[R^{\underline{\lambda},\tau,\crys,\Zpbar}_{\rhobar}:=R^{\underline{\lambda},\tau,\crys,\cO}_{\rhobar}\otimes_{\cO}\Zpbar,\] \[R^{\underline{\lambda},\tau,\st,\Zpbar}_{\rhobar}:=R^{\underline{\lambda},\tau,\st,\cO}_{\rhobar}\otimes_{\cO}\Zpbar.\]
% These are again well-defined independently of~$\cO$.\mar{TG:
%   eventually if we do want to use these we could just drop the
%   $\Zpbar$ from the notation? TG: more likely that we won't actually
%   use this, I imagine.}
By \cite[Thm.\ 3.3.8]{MR2373358} %\mar{ME: Cite something from MK here}
the localized rings 
$R^{\crys,\underline{\lambda},\tau,\cO}_{\rhobar}[1/p]$
are regular, and thus the rings
$R^{\crys,\underline{\lambda},\tau,\cO}_{\rhobar}$
(which embed into their localizations away from~$p$,
since they are $\cO$-flat) are reduced.



\subsection*{Algebra}Our conventions typically
follow~\cite{stacks-project}. In particular, if~$M$ is an abelian
topological group with a linear topology, then as
in~\cite[\href{https://stacks.math.columbia.edu/tag/07E7}{Tag
  07E7}]{stacks-project} we say that~$M$ is {\em complete} if the
natural morphism $M\to \varprojlim_i M/U_i$ is an isomorphism,
where~$\{U_i\}_{i \in I}$ is some (equivalently any) fundamental
system of neighbourhoods of~$0$ consisting of subgroups. Note that in
some other references this would be referred to as being~{\em complete
  and separated}.


If~$R$ is a ring, we write % ~$Ch(R)$ for the abelian category of cochain
% complexes of $R$-modules, 
$D(R)$ for the (unbounded) derived category of $R$-modules. 
% $D^b(R)$ for the bounded derived category of $R$-modules (that is, the
% triangulated subcategory of~$D(R)$ whose objects are
% cochain complexes with cohomology in only finitely many degrees), 
% and
% ~$D^{-}(R)$ for the bounded above derived category (whose objects are
% those cochain complexes whose cohomology vanishes in sufficiently high
% degree).
We say that a complex~$P^\bullet$ is \emph{good} if it is a bounded
complex of finite projective $R$-modules; then an object $C^\bullet$ of $D(R)$ is called a \emph{perfect
  complex} if there is a quasi-isomorphism
$P^\bullet \rightarrow C^\bullet$ where $P^\bullet$ is good. In fact, $C^\bullet$ is
perfect if and only if it is isomorphic in $D(R)$ to a good complex
$P^\bullet$: if we have another complex
$D^\bullet$ and quasi-isomorphisms $P^\bullet \rightarrow D^\bullet$,
$C^\bullet \rightarrow D^\bullet$, then there is a quasi-isomorphism
$P^\bullet \rightarrow C^\bullet$
(\cite[\href{http://stacks.math.columbia.edu/tag/064E}{Tag
  064E}]{stacks-project}). \index{good complex} \index{perfect complex}

\subsection*{Stacks}Our conventions on algebraic stacks and formal
algebraic stacks are those of~\cite{stacks-project}
and~\cite{Emertonformalstacks}. We recall some terminology and results
in Appendix~\ref{app: formal algebraic
  stacks}. Throughout the book, if~$A$ is a topological ring
and~$\cC$ is a stack we write~$\cC(A)$ for~$\cC(\Spf A)$; if~$A$ has
the discrete topology, this is equal to~$\cC(\Spec A)$. 
%\mar{ME: Add a remark that the Formal Stacks Appendix contains more 
%	conventions on stacks, etc.  Also, either here or there,
%	explain that $\cC(A)$ has the evident meaning, when $A$ is a
%	topological ring and $\cC$ is some kind of formal stack, Ind-stack,
%	or whatever. TG: probably you want to say something other than
%      this?}

% \chapter{Introduction}
% Our goal in this paper is to construct moduli spaces
% of $p$-adic representations of absolute Galois groups of $p$-adic local fields.
% These spaces will be formal algebraic stacks, % \mar{ME: Perhaps actually
% % 	just Ind-algebraic stacks? TG: seems to be formal algebraic
% %         stacks, in fact? ME: As long as the cotangent complex argument
% % hangs together \dots }
% and points 
% in the special fibre will correspond to residual Galois representations.
% The formal completion of one of our stacks at such a point will then recover
% the usual formal deformation space of this residual representation.
% Thus the relationship between the theory and constructions
% that we develop here, and the usual formal deformation theory of Galois
% representations, is the same as that between the theory of moduli spaces
% of algebraic varieties, and the formal deformation theory of algebraic 
% varieties: the latter gives valuable local information about the former,
% but moduli spaces, when they can be constructed,
% capture global aspects of the situation which are
% inaccessible by the purely local tools of formal deformation theory.

% \section{Statement of results}
% In order to explain our results in more detail, we introduce some notation.
% We fix a finite extension $K$ of $\Q_p$ (for some prime $p$),
% and a positive integer $d$ (the dimension of the representations being
% parameterized).  Our main construction is then a Noetherian
% formal algebraic stack $\cX_d$,
% equipped with a morphism to $\Spf \Z_p$, with the following properties:
% \begin{enumerate}
% 	%\item $\cX_d$ is flat and topologically of finite type over $\Spf
% 	%	\Z_p$,
% 	%	\mar{ME: I've put flat here, since it seems to simplify the
% 	%		dimension statement, but it may not be the correct
% 	%		thing to do if we don't know that unrestricted local
% 	%		def.\ rings are flat in general.}
% 	%	\mar{ME: Investigate the correct finite type statement.}
% 	%	of relative dimension $[K:\Q_p] d^2$.
% \item
% 	The underlying reduced substack $\cX_{d,\red}$ is of finite type over
% 	$\F_p$, and is equidimensional of dimension $[K:\Qp]d(d-1)/2$.
% 	\mar{ME: I've rewritten this to avoid flatness assumptions.  We should add Noetherian, or maybe locally Noetherian, here, depending on what we eventually prove.}
% 	\item If $\F$ is an algebraic extension of $\F_p$ (in
%           particular, $\Fpbar$ or
% 		a finite extension of $\F_p$), % \mar{TG: should ``in
%                   % particular'' be ``i.e.''? ME: There are other possibilities, since
% 		  % $\widehat{\mathbb Z}$ has lots of proper closed
%                   % subgroups that are of infinite index. TG: duh,
%                   % thanks, commenting out.}
% 		then the groupoid of points $\cX_d(\F)$ is naturally equivalent
% 		\mar{ME: Think through the use of ``natural'' in various
% 			statements, to avoid accidentally invoking $3$-categories.}
% 		to the groupoid of continuous representations
% 		$G_K \to \GL_d(\F).$
% 	\item If $\F$ is a finite extension of $\F_p$, if\mar{TG: we
%             could do~$\Fpbar$ here if we want?}
% 		$\Spec \F \to \cX_d$ is a morphism, corresponding (by
% 		the previous point) to a continuous representation
% 		$\rhobar: G_K \to \GL_d(\F)$, and if $R^{\square}_{\rhobar}$
% 		denotes the universal lifting $W(\F)$-algebra %  \mar{TG: might need to
%                   % elaborate - I guess the lifting $W(k)$-algebra}
%                 of $\rhobar$,
% 		then there is a natural morphism (of formal stacks
% 		over $\Spf \Z_p$) 
% 		% \mar{ME: If we don't know that these unrestricted deformation rings
% 		% 	 are flat in general, then we should either take
% 		% 	 the flat part, or allow $\cX_d$ to be non-flat.}
% 		$\Spf R^{\square}_{\rhobar} \to \cX_d$ extending 
% 		the given morphism on $\Spec \F$; this morphism is versal. 
% 		This last property can be made more precise: namely, the
% 		fibre product $\Spf R_{\rhobar}^{\square} \times_{\cX_d}
% 		\Spf R_{\rhobar}^{\square}$ is naturally isomorphic to
% 		the groupoid $\Spf R_{\rhobar}^\square \times \widehat{\GL}_d 
% 	\rightrightarrows \Spf R_{\rhobar}^\square$ given by change of basis.
% \end{enumerate}

% From point~(3) it follows formally that if $\cO$ is the ring
% of integers in any finite extension of $\Q_p$
% then the groupoid $\cX_d(\cO)$ (i.e.\ the groupoid of morphisms of
% formal algebraic
% stacks $\Spf \cO \to \cX_d$ over $\Spf \Z_p$) is naturally equivalent
% to the groupoid of continuous representations $G_K \to \GL_d(\cO)$.
% Similarly, the groupoid $\cX_d(\Zbar_p)$ is naturally equivalent
% to the groupoid of continuous representations
% $G_K \to \GL_d(\Zbar_p)$.
% More generally, if $A$ is any finite $\Z_p$-algebra,
% then the groupoid $\cX(A)$ is naturally equivalent
% to the groupoid of continous representations $G_K \to \GL_d(A)$.\mar{ME: This last assertion may not be true if we have to pass to flat parts, and if $A$
% 	is not flat over $\Z_p$; instead we would just get a certain
% 	subgroupoid.}

% %As usual, we let $\cX_{\red}$ denote the reduced algebraic stack
% %underlying the formal algebraic stack $\cX$.  It follows
% %from the first
% %point above that $\cX_{\red}$ is of finite type
% %over $\F_p$, but we have the following more precise statement.
% %
% %\begin{enumerate}
% %		\setcounter{enumi}{3}
% %	\item $\cX_{\red}$ is equidimensional of dimension $d(d-1)/2$ over $\F_p$.
% %\end{enumerate}

% Suppose now that $\underline{\lambda}$ is some $d$-tuple of labeled Hodge--Tate
% weights, and that $\tau$ is some inertial type.% \mar{TG: define this,
%   % and the notion of regular weight, in the notation section.}

% \begin{enumerate}
% 		\setcounter{enumi}{3}
% 	\item There is a closed substack $\cX_d^{\underline{\lambda},\tau,\crys}$
% 		of $\cX_d$ which is %finite type $p$-adic formal algebraic stack,
% 		flat over $\Spf \Z_p$, \mar{ME: There is a question of whether
% 			flatness makes sense, because this might require
% 			Noetherianness of the formal stack!}
% 		and which is
% 		characterized by the property that a $\Zbar_p$-valued point
% 		of $\cX_d$ lies in $\cX_d^{\underline{\lambda},\tau,\crys}$ if
% 		and only if the associated Galois representation
% 		$G_K \to \GL_d(\Zbar_p)$, when thought of as a representation
% 		into $\GL_d(\Qbar_p)$, is potentially crystalline of inertial
% 		type $\tau$ and with labeled Hodge--Tate weights 
% 		$\underline{\lambda}$.   Each of the underlying reduced algebraic
% 		substacks
% 		$(\cX_d^{\underline{\lambda},
% 			\tau,\crys})_{\red}$ is equidimensional of dimension an explicit function
%                       of~$\underline{\lambda}$, which is less than or equal
%                       to~$[K:\Qp]d(d-1)/2$, with equality if and only
%                       if $\underline{\lambda}$ is regular. The analogous
%                       statements hold for ``semi-stable'' in place of
%                       ``crystalline''. %  \mar{ME: Expected dimension
% 		% 	here; $[K:\Qp]d(d-1)/2$ when $\underline{\lambda}$ is regular.}
% 		% Similar statement for semi-stable in place of crystalline.
% \end{enumerate}

% \begin{remark}
% We expect that each of the formal algebraic
% stacks $\cX_d^{\underline{\lambda},\tau,\crys}$ is in fact a $p$-adic formal
% algebraic stack of finite type, but we are able to prove this only
% in the case when $K$ is unramified over $\Q_p$.
% \end{remark}

% It follows from~(1) and~(4) that if $\underline{\lambda}$ is regular,
% then $\cX^{\underline{\lambda}}_{d,\red}$ is a union of irreducible components
% of $\cX_{d,\red}$.  It follows from~(3) and~(4) that
% the pull-back of $\cX_d^{\underline{\lambda},\tau}$ under a versal morphism
% $\Spf R_{\rhobar}^{\square} \to \cX_d$ as in~(3) is equal to the usual
% potentially crystalline lifting space $\Spf R_{\rhobar}^{\underline{\lambda},
% 	\tau,\crys}.$

% There are some additional structures that $\cX_d$ carries.
% For example:
% \begin{enumerate}
% 		\setcounter{enumi}{4}
% 	\item
% There is a perfect complex of $\cO_{\cX_d}$-modules
% $\cC^{\bullet}$ on $\cX_d$,
% with cohomology lying in degrees $0$, $1$, and $2$, such that
% if $\Spf A \to \cX_d$ is any morphism with $A$ being a finite $\Z_p$-algebra,
% corresponding to a continuous
% Galois representation $\rho: G_K \to \GL_d(A)$,
% then the base-changed complex $A\otimes_{\cO_{\cX_d}} \cC^{\bullet}$ 
% computes the continuous Galois cohomology $H^{\bullet}(G_K, \rho)$.
% \end{enumerate}

% Furthermore, there are morphisms
% relating the various $\cX_d$ for differing values of $d$
% corresponding to multi-linear algebra constructions on Galois representations.
% For example, there are morphisms
% $$\cX_d \times_{\Z_p} \cX_{d'} \to \cX_{d + d'}$$
% and 
% $$\cX_d \times_{\Z_p} \cX_{d'} \to \cX_{d  d'}$$
% corresponding to direct sum and tensor product respectively.

% We may also use classes in $H^1$ to construct families
% of extensions.  More precisely:

% \begin{enumerate}
% 		\setcounter{enumi}{5}
% \item Suppose that
% 	$f_d: \cZ \to \cX_d$ is any morphism from a formal algebraic stack
% 	to $\cX_d$, and that $\phi: \cF \to H^1(f_d^*\cC^{\bullet})$ is any morphism
% 	of $\cO_{\cZ}$-modules
% 	whose domain is a locally free sheaf of locally finite rank over $\cZ$.
% 	If $\cV \to \cX_d$ is the geometric vector bundle associated
% 	to $\cF$ (so $\cF$ is the sheaf of sections of~$\cV$),
% 	then there is a morphism $f_{d+1}: \cV \to \cX_{d+1}$.
% 	On $A$-valued points, for a finite $\Z_p$-algebra $A$,
% 	this morphism has the following interpretation:
% 	giving a morphism $\Spf A \to \cV$ is equivalent to giving
%         a morphism $g: \Spf A \to \cZ$ together with
% 	a class $c \in A\otimes_{\cO_{\cZ}} \cF$; composing the
% 	morphism  $g$ with $f_d$ yields a morphism
% 	$\Spf A \to \cX_d$, corresponding to a continuous representation
% 	$\rho: G_K \to \GL_d(A)$; applying (the base-change to $A$ of)
% 	$\phi$ to the class $c$ yields a class in $H^1(G_K,\rho),$
% 	which in turn corresponds to a Galois representation $\rho'$
% 	which sits in an exact sequence $0 \to \rho \to \rho' \to \triv \to 0;$
% 	then $f_{d+1}$ is constructed in such a way that
% 	the composite $\Spf A \to \cV \buildrel f_{d+1} \over \longrightarrow
%        	\cX_{d+1}$ is precisely
% 	the morphism that classifies the representation
% 	$\rho'$.% \mar{TG: will we actually need something more general
% %           than this when we want to prove existence of lifts in
% %           general niveau? ME: In principle, but actually, if we have
% %   a reasonable multilinear algebra formalism relating the various $\cX_d$'s,
% % then the formula $\Ext^i(\rho,\psi) = H^i(\psi\otimes\rho^{\vee})$ should
% % let us reduce to this case (which just has the technical advantage
% % of not having to study any complexes beyond the Herr complex).  But you might
% % be right, that it could be better to have this reduction from $\Ext^i$ to $H^i$
% % to take place somewhere in the technical development, and to have
% % as part of the formalism
% % a complex over the product $\cX_{d_1} \times \cX_{d_2}$ which computes
% % $\Ext^i$ between the universal families.}
% \end{enumerate}

% \section{The underlying reduced substack $\cX_{d,\red}$,
% 	and the geometric Breuil--M\'ezard philosophy}
% \label{subsec:intro Breuil--Mezard}
% If $R_{\rhobar}$ is the universal lifting ring of a $d$-dimensional
% continuous representation $\rhobar: G_K \to \GL_d(\F)$ (for some finite
% extension $\F$ of $\F_p$), then $\Spf R_{\rhobar}$ doesn't carry
% a lot of evident geometry in and of itself; for example, its underlying reduced
% subscheme is simply the closed point $\Spec \F$, corresponding 
% to $\rhobar$ itself.% \mar{TG: should probably use either $k$ or $\F$ in
% % the introduction, but not both; my guess is that $\F$ will be
% % consistent with the rest of the paper? TG: changed, commenting out.}

% On the other hand, $\cX_d$ has a quite non-trivial underlying reduced
% substack, which parameterizes all the $d$-dimensional residual representations
% of $G_K$.  It's natural to ask whether this underlying reduced substack has
% any significance in formal deformation theory.  More precisely, we could ask
% for the meaning of the fibre product
% $ \Spf R_{\rhobar} \times_{\cX_d} \cX_{d,\red}.$

% This is a reduced closed subscheme of $\Spf R_{\rhobar}$, of
% dimension $[K:\Q_p]d(d-1)/2 + d^2$, whose irreducible components,
% when thought of cycles on $\Spf R_{\rhobar}$, are precisely the
% cycles that appear in the geometric Breuil--M\'ezard conjecture
% of~\cite{emertongeerefinedBM}.
% \mar{ME: Elaborate on/improve this discussion, add citation, etc.}
% %As a specific application of this idea,
% More specifically, we obtain the following qualitive 
% version of the geometric Breuil--M\'ezard
% conjecture~\cite[Conj.~4.2.1]{emertongeerefinedBM}.

% \begin{theorem}
% 	If $\rhobar: G_K \to \GL_d(\F)$ is a continuous representation
% 	\mar{ME: Is there any technical reason we can't just take $\F
%           = \Fbar_p$, and so avoid the annoyance of having to talk
%           about a finite $\F$? TG: I suspect so; one approach would be
%         to follow BLGGT, I suspect, or we could see if it's easy to
%         do ``properly'' by directly defining the deformation
%         rings. TG: I had a go at doing this in the
%         notation/conventions section; see what you think. TG: need to
%         add coefficients here to account for possibly ramified types}
% 	{\em (}with $\F$ being a finite extension of $\F_p${\em )},
% 	then there are finitely many cycles of codimension $[K:\Q_p]d(d+1)/2$
% 	in $\Spec R^{\square}_{\rhobar}$, such that for any regular tuple
% 	of labeled Hodge--Tate weights $\underline{\lambda}$, and any inertial
% 	type $\tau$, each of the special fibres
% 	$\Spec R^{\underline{\lambda},\tau,\crys}_{\rhobar}/p$ 
% 	and 
% 	$\Spec R^{\underline{\lambda},\tau,\ss}_{\rhobar}/p$ 
% 	is set-theoretically supported on the union of some number
% 	of these
% 	cycles.% \mar{TG: does this language imply the set-theoretic
%   %         support is exactly this (which is too strong)? ME: Added
%   % language to address this.}
% \end{theorem}

% The cycles in the statement of the theorem are precisely
% \mar{ME: Throughout this subsection there is the perennial issue of 
% 	$\Spec R$ vs.\ $\Spf R$.}
% the fibre product $\Spf R_{\rhobar}\times_{\cX_d} \cX_{d,\red}$,
% which, more intuitively, we can think of as the germ of $\cX_{d,\red}$
% at~$\rhobar$.


% \section{Crystalline lifts and the support of $H^2(\rhobar)$}
% \label{subsec:intro crystalline lifts}
% As an application of the geometric Breuil--M\'ezard ideas discussed
% in the preceding subsection, %our construction,
% we prove the following result.

% \begin{thm}
% 	\label{thm:intro crystalline lifts}
% 	If $\rhobar: G_K \to \GL_d(\Fbar_p)$ is a continuous
% 	representation, then $\rhobar$
% 	has a lift $\rho^{\circ}: G_K \to \GL_d(\Zbar_p)$ for which
% 	the associated $p$-adic representation $\rho:G_K \to \GL_d(\Qbar_p)$
% 	is crystalline of regular Hodge--Tate weights. 

% We can furthermore choose
%         these Hodge--Tate weights to be in the
%         interval~$[0,dp-1]$, % \mar{TG: not too sure about this - might be
%           % better to do bounded gaps instead, although obviously a
%           % result of this kind will be true.}
% 	% {\em (}the bound depending only on $d$ and $K${\em )},
%         or we
%         can choose them to have arbitrarily large gaps, and in either
%         case we can also ensure that~$\rho^{\circ}$ is potentially diagonalizable.
% % 	\mar{ME: We could consider making the bound more explicit, either
% % 		in the statement of the result or in a subsequent remark.}
% % 	\mar{ME: Maybe add potentially diagonalizable into the statement,
% % 		or perhaps discuss it afterwards, if it seems a bit
% % 		out of place in the statement at this particular
% %                 moment. TG: also want well-spread, rather than
% %                 bounded, as an option. ME: Currently, can get any plausible
% % 	property of the HT weights, but can't get pot'l diagonalizable ---
% % have to work on this.}
% \end{thm}
% The notion of a potentially diagonalizable representation was defined
% in~\cite{BLGGT}, and is important for applications to automorphy
% lifting theorems; for example, Theorem~\ref{thm:intro crystalline
%   lifts} allows us to remove a hypothesis on the construction of a
% candidate $p$-adic Langlands correspondence in~\cite{Gpatch}. In
% addition to proving that the Hodge--Tate weights can be taken  to be in the
%         interval~$[0,dp-1]$, % the
% % case of lifts of bounded weight, \mar{ME: What does ``In the case of lifts of
% % 	bounded weight'' mean? TG: hopefully this is clearer.}
% % % \mar{TG: discuss this later in the paper}
% we are 
% able to prove a conjecture of~\cite{2015arXiv150902527G} on the
% existence of a crystalline lift whose Hodge--Tate weights correspond
% to a Serre weight.

% The problem of constructing crystalline lifts has been considered by various
% authors, in particular~\cite{MullerThesis,2015arXiv150601050G}.
% It admits a well-known inductive approach (which is taken in~\cite{MullerThesis,2015arXiv150601050G}): one writes $\rhobar$ as 
% a successive extension of irreducible representations, lifts each of these
% irreducible representations 
% to a crystalline representation, and then attempts to lift the various
% extension classes.  The difficulty that arises in this approach (which
% has proved an obstacle to obtaining general statements along the lines
% of Theorem~\ref{thm:intro crystalline lifts} up until now) is showing that the 
% mod $p$ extension classes that appear in this description of $\rhobar$ 
% can actually be lifted to crystalline extension classes in characteristic
% zero.  The basic source of the difficulty is that the local Galois $H^2$
% can be non-zero, and non-zero classes in $H^2$ obstruct the lifting 
% of extension classes (which can be interpreted as classes lying in
% $H^1$). (In fact, the difficulty is not so much in obtaining
% crystalline extension classes, as in lifting to any classes in
% characteristic zero; indeed, as far as we are aware, until now
% it was not known that an arbitrary~$\rhobar$ had \emph{any} lift to
% characteristic zero at all.)

% Our proof of Theorem~\ref{thm:intro crystalline lifts} relies on the
% inductive strategy described in the preceding paragraph, but we are able to
% prove the following key result, which controls the obstructions that
% can be presented by $H^2$.

% \begin{prop}
% 	\label{prop:H2 dimension}
% 	The locus of points $\rhobar \in \cX_{d,\red}(\Fbar_p)$ 
% 	at which $\dim H^2(G_K,\rhobar) \geq r$ is Zariski closed
% 	in $\cX_{d,\red}(\Fbar_p)$, of codimension $\geq r$.
% \end{prop}

% If we let $\rho^{\univ}$ denote the universal lifting of $\rhobar$
% over $R^{\square}_{\rhobar}$, then $H^2(G_K, \rho^{\univ})$ is 
% an $R^{\square}_{\rhobar}$-module, and Proposition~\ref{prop:H2 dimension}
% implies the following corollary.

% \begin{cor}
% 	\label{cor:H2 dimension}
% 	For any regular tuple of labeled Hodge--Tate weights $\underline{\lambda}$
% 	and inertial type $\tau$, the locus of points 
% 	$x\in \Spec R^{\underline{\lambda},\tau,\crys}/p$ for which
% 	$$\dim_{\kappa(x)}
% 	\kappa(x)\otimes_{R^{\square}_{\rhobar}} H^2(G_K,\rho^{\univ})
% 	\geq r$$
% 	has codimension $\geq r$.
% \end{cor}

% \begin{remark}
% 	\label{rem:H2 dimension}
% 	By Tate local duality, 
% 	$$\dim_{\kappa(x)}
% 	\kappa(x)\otimes_{R^{\square}_{\rhobar}} H^2(G_K,\rho^{\univ})
% 	=
% 	\dim_{\kappa(x)}
% 	\kappa(x)\otimes_{R^{\square}_{\rhobar}}
% 	\Hom_{G_K}\bigl((\rho^{\univ})^{\vee},\omega\bigr)$$
% 	(where $\omega$ denotes the mod $p$ cyclotomic character, thought
% 	of as taking values in $\kappa(x)^{\times}$).
% 	Thus the statement of Corollary~\ref{cor:H2 dimension} is
% 	related to the way in which the special fibre of the potentially crystalline
% 	lifting ring $R_{\rhobar}^{\underline{\lambda},\tau,\crys}$ intersects
% 	the reducibility locus in $\Spec R^{\square}_{\rhobar}$.
% \end{remark}

% Given Corollary~\ref{cor:H2 dimension},
% we prove Theorem~\ref{thm:intro crystalline
% 	lifts}
% by working purely within the context of formal lifting rings.  
% However we don't know how to
% prove the corollary while staying within that context.
% Indeed, as Remark~\ref{rem:H2 dimension} indicates, 
% this corollary is related to the way in which the special fibre
% of a potentially crystalline deformation ring intersects another natural
% locus in $\Spec R_{\rhobar}^{\square}$ (namely, the reducibility locus).
% Since the special fibre of a potentially crystalline lifting ring
% is not directly defined in deformation-theoretic terms, such questions
% are notoriously difficult to study directly.
% Our proof of the corollary proceeds
% differently, by using the  geometric Breuil--M\'ezard view-point
% %discussed in Subsection~\ref{subsec:intro Breuil--Mezard} above
% to replace a computation on the special fibre of the potentially
% crystalline deformation ring by a compution on $\cX_{d,\red}$;
% this latter space has a concrete description in terms of families of varying
% $\rhobar$, whose $H^2$ we are  able to compute.

% \section{The construction of $\cX_d$}
% We construct $\cX_d$ as the moduli stack of rank $d$ \'etale 
% $(\varphi,\Gamma)$-modules.  \mar{ME: Elaborate on this.}

% \section{Notation and conventions}
% \subsection*{$p$-adic Hodge theory}Let $K/\Qp$ be a finite
% extension. If $\rho$ is a de Rham representation of $G_K$ on a
% $\Qpbar$-vector space~$W$, then we will write $\WD(\rho)$ for the
% corresponding Weil--Deligne representation of $W_K$ (see e.g.\
% \cite[App. B]{CDT}), and if $\sigma:K \into \Qpbar$ is a continuous
% embedding of fields then we will write $\HT_\sigma(\rho)$ for the
% multiset of Hodge--Tate numbers of $\rho$ with respect to $\sigma$,
% which by definition contains $i$ with multiplicity
% $\dim_{\Qpbar} (W \otimes_{\sigma,K} \widehat{\barK}(i))^{G_K} $. Thus
% for example if~$\varepsilon$ denotes the $p$-adic cyclotomic character,
% then $\HT_\sigma(\varepsilon)=\{ -1\}$.

% By a \emph{$d$-tuple of labeled Hodge--Tate weights~$\underline{\lambda}$},
% we mean integers~$\{\lambda_{\sigma,i}\}_{\sigma:K\into\Qpbar,1\le i\le d}$
% with $\lambda_{\sigma,i}\ge \lambda_{\sigma,i+1}$ for all~$\sigma$ and all $1\le
% i\le d-1$. By an \emph{inertial type~$\tau$} we mean a representation
% $\tau:I_K\to\GL_d(\Qpbar)$ which extends to a  representation
% of~$W_K$ with open kernel (so in particular, $\tau$ has finite
% image). 

% Then we say that~$\rho$ has Hodge type~$\underline{\lambda}$ (or
% labeled Hodge--Tate weights~$\underline{\lambda}$) if for
% each~$\sigma:K\into\Qpbar$ we
% have~$\HT_\sigma(\rho)=\{\lambda_{\sigma,i}\}_{1\le i\le
%   d}$, and we say that~$\rho$ has inertial
% type~$\tau$ if~$\WD(\rho)|_{I_K}\cong\tau$.

% \subsection*{Serre weights and Hodge--Tate weights}By a \emph{Serre
% weight} $\underline{k}$ we mean a tuple of
% integers~$\{k_{\sigmabar,i}\}_{\sigmabar:k\into\Fpbar,1\le i\le d}$
% with the properties that \begin{itemize}
% \item $p-1\ge k_{\sigmabar,i}-k_{\sigmabar,i+1}\ge 0$ for each $1\le i\le
%   d-1$, and
% \item $p-1\ge k_{\sigmabar,d}\ge 0$, and not every~$k_{\sigmabar,d}$
%   is equal to~$p-1$.
% \end{itemize}
% The set of Serre weights is in bijection with the set of irreducible
% $\Fpbar$-representations of~$\GL_d(k)$, via passage to highest
% weight vectors (see for example the appendix to~\cite{MR2541127}). 

% % We now recall from~\cite{2015arXiv150902527G} the notion of a
% % crystalline lift of some Serre weight.
% Each embedding~$\sigma:K\into\Qpbar$ induces an
% embedding~$\sigmabar:k\into\Fpbar$; if $K/\Qp$ is ramified, then
% each~$\sigmabar$ corresponds to multiple embeddings~$\sigma$.  We say
% that~$\underline{\lambda}$ is a lift of  ~$\underline{k}$ % is a \mar{TG: here need to fix~$k$ versus
%   % $\lambda$} \emph{Serre weight}
% if for each
% embedding~$\sigmabar:k\into\Fpbar$, we can choose an
% embedding~$\sigma:K\into\Qpbar$ lifting~$\sigmabar$, with the
% properties that:
% \begin{itemize}
% \item $\lambda_{\sigma,i}=k_{\sigma,i}+d-i$, and
% \item if $\sigma':K\into\Qpbar$ is any other lift of~$\sigmabar$,
%   then $k_{\sigma',i}=d-i$.
% \end{itemize}


% \subsection*{Lifting rings}Let $K/\Qp$ be a finite extension, and
% let $\rhobar:G_K\to\GL_d(\Fpbar)$ be a continuous representation. Then
% the image of~$\rhobar$ is contained in~$\GL_d(\F)$ for any
% sufficiently large finite
% extension~$\F/F_p$. Let~$\cO$ be the ring of integers in some
% finite extension~$E/\Qp$, and suppose that the residue field of~$E$
% is~$\F$. Then we let~$R^{\square,\cO}_{\rhobar}$ be the universal
% lifting~$\cO$-algebra of~$\rhobar$. % \mar{TG: might be better to
%   % use~$W(\Fpbar)$ coefficients (or rather finite extensions of that?)}
% The precise choice of~$E$ is unimportant, in the sense that
% if~$\cO'$ is the ring of integers in a finite extension~$E'/E$, then
% by~\cite[Lem.\ 1.2.1]{BLGGT} we have
% $R^{\square,\cO'}_{\rhobar}=R^{\square,\cO}_{\rhobar}\otimes_{\cO}\cO'$.
% % so
% % that \[R^{\square,\Zpbar}_{\rhobar}:=R^{\square,\cO}_{\rhobar}\otimes_{\cO}\Zpbar\]is
% % well-defined, independently of the choice of~$\cO$. By definition, the
% % $\Zpbar$-points of $\Spf R^{\square,\Zpbar}_{\rhobar}$ correspond
% % precisely to the liftings of~$\rhobar$ to a
% % representation~$G_K\to\GL_d(\Zpbar)$.\mar{TG: I didn't think this
% %   through super carefully, so hopefully it's not a blunder.}

% % By a \emph{$d$-tuple of labeled Hodge--Tate weights~$\underline{\lambda}$},
% % we mean integers~$\{\lambda_{\sigma,i}\}_{\sigma:K\into\Qpbar,1\le i\le d}$
% % with $\lambda_{\sigma,i}\ge \lambda_{\sigma,i+1}$ for all~$\sigma$ and all $1\le
% % i\le d-1$. By an \emph{inertial type~$\tau$} we mean a representation
% % $\tau:I_K\to\GL_d(\Qpbar)$ which extends to a  representation
% % of~$W_K$ with open kernel (so in particular, $\tau$ has finite
% % image).
% If~$\cO$ is chosen large enough that the inertial type~$\tau$ is defined
% over~$E=\cO[1/p]$, and large enough that~$E$ contains the images of all
% embeddings~$\sigma:K\into\Qpbar$, then we have the usual lifting
% $\cO$-algebra ~$R^{\underline{\lambda},\tau,\crys,\cO}_{\rhobar}$
% % and~$R^{\underline{\lambda},\tau,\st,\cO}_{\rhobar}$
% which corresponds to
% lifts of~$\rhobar$ which are potentially crystalline/semistable with
% labeled Hodge--Tate weights~$\underline{\lambda}$ and inertial
% type~$\tau$. By definition, these are $\cO$-flat, % \mar{TG:
%   % check if they are really reduced by definition}
% and by
% the main theorems of~\cite{MR2373358}, they are generically regular,
% and are equidimensional of
% dimension \[1+d^2+\sum_{\sigma}\#\{1\le i<j\le
%   n|\lambda_{\sigma,i}>\lambda_{\sigma,j}\}.\]% \mar{TG: hopefully correct, didn't
%   % think about it too carefully.}
% Note that this quantity is at most $1+d^2+[K:\Qp]d(d-1)/2$, with
% equality if and only if~$\underline{\lambda}$ is \emph{regular}, in
% the sense that $\lambda_{\sigma,i}> \lambda_{\sigma,i+1}$ for all~$\sigma$ and all $1\le
% i\le n-1$. As above, we have $R^{\underline{\lambda},\tau,\crys,\cO'}_{\rhobar}=R^{\underline{\lambda},\tau,\crys,\cO}_{\rhobar}\otimes_{\cO}\cO'$.  % Just as above, we
% % set \[R^{\underline{\lambda},\tau,\crys,\Zpbar}_{\rhobar}:=R^{\underline{\lambda},\tau,\crys,\cO}_{\rhobar}\otimes_{\cO}\Zpbar,\] \[R^{\underline{\lambda},\tau,\st,\Zpbar}_{\rhobar}:=R^{\underline{\lambda},\tau,\st,\cO}_{\rhobar}\otimes_{\cO}\Zpbar.\]
% % These are again well-defined independently of~$\cO$.\mar{TG:
% %   eventually if we do want to use these we could just drop the
% %   $\Zpbar$ from the notation? TG: more likely that we won't actually
% %   use this, I imagine.}


% \subsection*{Algebraic stacks}We follow the terminology of~\cite{stacks-project}; in
% particular, we write ``algebraic stack'' rather than ``Artin stack''. More
% precisely, an algebraic stack is a stack in groupoids in the \emph{fppf} topology,
% whose diagonal is representable by algebraic spaces, which admits a smooth
% surjection from a
% scheme. See~\cite[\href{http://stacks.math.columbia.edu/tag/026N}{Tag
%   026N}]{stacks-project} for a discussion of how this definition relates to
% others in the literature. If~$S$ is a scheme, then by ``a stack
% over~$S$'' we mean  a stack fibred in groupoids over the big
% \emph{fppf} site of~$S$.

% \subsection*{Ind-algebraic stacks}As in~\cite[Defn.\
% 4.2]{Emertonformalstacks}, an Ind-algebraic stack over a scheme~$S$ is
% a stack~$\cX$ over~$S$ which can be written as
% $\cX\cong\varinjlim_{i\in I}\cX_i$, where we are taking the 2-colimit
% in the 2-category of stacks of a 2-directed system $\{\cX_i\}_{i\in
%   I}$ of algebraic stacks over~$S$.

% \subsection*{Formal algebraic stacks}We refer the reader
% to~\cite{Emertonformalstacks} for full details of the framework of
% formal algebraic stacks. In particular, a formal algebraic
% stack over~$S$ is a stack~$\cX$ which admits a morphism $U\to\cX$
% which is representable by algebraic spaces, smooth, and surjective,
% and whose source is a formal algebraic space in the sense
% of~\cite[\href{http://stacks.math.columbia.edu/tag/0AHW}{Tag
%   0AHW}]{stacks-project}. The relationship between formal algebraic
% stacks and Ind-algebraic stacks is discussed in detail
% in~\cite[\S6]{Emertonformalstacks}. 

%  A \emph{$p$-adic formal algebraic stack} is
% a formal algebraic stack $\cX$ over~$\Spec\Zp$ which admits a morphism
% $\cX\to\Spf\Zp$ which is representable by algebraic stacks; so in
% particular, we may write~$\cX$ as an Ind-algebraic stack by writing
% $\cX\cong\varinjlim_n\cX\times_{\Spf\Zp}\Spec\Z/p^n\Z$. The
% following criterion for being a $p$-adic formal algebraic stack will
% be important for us.

% \begin{prop}
%   \label{prop: criterion for p-adic formal algebraic stack}Suppose
%   that~$\{\cX_n\}_{n\ge 1}$ is an inductive system of algebraic stacks
%   over~$\Spec\Zp$,
%   %whose transition morphisms are thickenings,\mar{TG:
%     %isn't this automatic from the following conditions? ME: Yes, corrected}
% such that each~$\cX_n$ in fact lies over~$\Spec\Z/p^n\Z$,  each
% of the natural morphisms
% $\cX_n\to\cX_{n+1}\times_{\Spec\Z/p^{n+1}\Z}\Z/p^n\Z$
% is an isomorphism, and each~$\cX_n$ has quasi-compact
% diagonal.\mar{TG: hypothesis on diagonal has been removed from
%   \cite{Emertonformalstacks} ?}
% Then the Ind-algebraic stack $\cX:=\varinjlim_n\cX_n$
%   is a $p$-adic formal algebraic stack.
% \end{prop}
% \begin{proof}
%   This is a special case of~\cite[Ex.\ 7.8]{Emertonformalstacks}.
% \end{proof}

% \subsection*{Derived categories}
% If~$R$ is a ring, we write~$Ch(R)$ for the abelian category of cochain
% complexes of $R$-modules, 
% $D(R)$ for the (unbounded) derived category of $R$-modules, 
% $D^b(R)$ for the bounded derived category of $R$-modules (that is, the
% triangulated subcategory of~$D(R)$ whose objects are
% cochain complexes with cohomology in only finitely many degrees), and
% ~$D^{-}(R)$ for the bounded above derived category (whose objects are
% those cochain complexes whose cohomology vanishes in sufficiently high
% degree).

% We say that a complex~$P^\bullet$ is \emph{good} if it is a bounded
% complex of finite projective $R$-modules; then an object $C^\bullet$ of $D(R)$ is called a \emph{perfect
%   complex} if there is a quasi-isomorphism
% $P^\bullet \rightarrow C^\bullet$ where $P^\bullet$ is good. In fact, $C^\bullet$ is
% perfect if and only if it is isomorphic in $D(R)$ to a good complex
% $P^\bullet$: if we have another complex
% $D^\bullet$ and quasi-isomorphisms $P^\bullet \rightarrow D^\bullet$,
% $C^\bullet \rightarrow D^\bullet$, then there is a quasi-isomorphism
% $P^\bullet \rightarrow C^\bullet$
% (\cite[\href{http://stacks.math.columbia.edu/tag/064E}{Tag
%   064E}]{stacks-project}).
\chapter{Rings and coefficients}\label{section: coefficient rings}In this
chapter we study various rings which will be the coefficients of the
$\varphi$-modules and $(\varphi,\Gamma)$-modules
%which we will consider in the later sections. 
that we study in the subsequent chapters.
Throughout the chapter,
we fix a finite extension $K$ of~$\Q_p$,
which we regard as a subfield of some fixed algebraic closure $\Qbar_p$
of~$\Q_p$.

We begin by recalling some rings used in integral $p$-adic Hodge
theory (in particular, in defining $(\varphi,\Gamma)$-modules and
Breuil--Kisin modules), before introducing versions of them with
coefficients in certain $p$-adically complete $\Zp$-algebras. We then
prove (almost) descent results that allow us to relate $\varphi$-modules over
various different rings,
%certain rings and their perfections,
before introducing the
$(\varphi,\Gamma)$-modules which we work with throughout the book.

\section{Rings}\label{subsec:rings}
\subsection{Perfectoid fields and their tilts}
% \mar{ME: Currently, haven't defined perfectoid or titls.  In the old
% 	material, the tilts are written down explicitly.  Maybe in
% 	the current epoch, this isn't really necessary?  Otherwise,
% 	could just recall the definitions. TG: I vote for not
%         bothering to define them, we can be trailblazers for them
%         being completely standard. Commenting out.}
As usual, we let $\C$ denote the completion of the algebraic closure \index{$\C$}
$\Qbar_p$ of $\Q_p$.
%and let $\cO_{\C}$ denote its ring of integers. 
It is a perfectoid field,
whose tilt
$\C^{\flat}$ is a complete  \index{$\C^{\flat}$}
non-archimedean valued perfect field of characteristic~$p$.
If $F$ is a perfectoid closed subfield of $\C$,
then its tilt $F^{\flat}$ is a closed, and perfect,
subfield of $\C^{\flat}$.

%As usual,
We let $\cO_{\C}$ denote the ring of integers in $\C$. \index{$\cO_{\C}$}
Its tilt $\cO_{\C}^{\flat}$ \index{$\cO_{\C}^{\flat}$} is then the ring of integers
in $\C^{\flat}$.
Similarly, if $\cO_F$ denotes the ring of integers
in $F$, then $\cO_F^\flat$ is the ring
of integers in~$F^\flat$.

%Since $\C^\flat$, or more generally $F^\flat$, is perfect,
We may form the rings of Witt vectors
$W(\C^\flat)$ and $W(F^\flat)$, \index{$W(\C^\flat)$}
and the rings of Witt vectors
$W(\cO_{\C^\flat})$ and $W(\cO_F^\flat)$;
	following the standard convention,
	we typically denote $W(\cO_{\C^\flat})$  by $\Ainf$.\index{$\Ainf$}

Each of these rings of Witt vectors
is a $p$-adically complete ring,
but we always consider them as topological rings \index{weak topology}
by endowing them with a finer topology, the so-called {\em weak
	topology}, which admits the following description: 
If $R$ is any of $\C^\flat$, $F^\flat$, $\cO_\C^\flat$,
or $\cO_F^\flat$, endowed with its natural (valuation)
topology,
then there is a canonical identification (of sets):
$$W_a(R) \iso R\times \cdots \times R \text{ ($a$ factors)},$$
and we endow $W_a(R)$ with the product topology.
We then endow $W(R) := \varprojlim_a W_a(R)$ with the inverse 
limit topology.

This topology admits the following more concrete description 
(in the general case of a perfectoid $F$; setting $F = \C$ 
recovers that particular case):
If $x$ is any element of $\cO_F^\flat$ of positive
valuation, and if $[x]$ denotes the Teichm\"uller 
lift of $x$, then we endow
$W_a(\cO_F^\flat)$ with the $[x]$-adic topology,
so that $W(\cO_F^\flat)$ is then endowed with
the $(p,[x]$)-adic topology.
The topology on $W_a(F^\flat)$ is then characterized 
by the fact that $W_a(\cO_F^\flat)$ is an open subring
--- concretely $W_a(F^\flat) = W_a(\cO_F^\flat)[\frac{1}{[x]}] 
= \varinjlim_i W_a(\cO_F^\flat)$ (the transition maps being
given by multiplication by~$[x]$, and each transition map
being an open and closed embedding),
while the topology on $W(F^\flat)$ is the
inverse limit topology --- concretely, $W(F^\flat)$
is the $p$-adic completion~$\widehat{W(\cO_F^\flat)[\frac{1}{[x]}]}$.

Apart from the case of $\C$ itself, 
there are two main examples of perfectoid $F$ that will be of importance
to us; see Example~\ref{ex:Sen} for a justification of the claim that
these fields are indeed perfectoid.
%In both examples, we begin by fixing a finite extension $K$ of~$\Q_p$.

\begin{example}[The cyclotomic case]
	\label{ex:cyclo}
	We write $K(\zeta_{p^{\infty}})$ to denote the 
	extension of~ $K$ obtained by adjoining all $p$-power
	roots of unity.  It is an infinite degree Galois extension
	of $K$, whose Galois group is naturally
	identified with an open subgroup of $\Z_p^{\times}$.
	We let $K_{\cyc}$ denote the unique subextension of \index{$K_{\cyc}$}
	$K(\zeta_{p^{\infty}})$ whose Galois group
	over $K$ is isomorphic to $\Z_p$ (so $K_{\cyc}$ is
	the ``cyclotomic $\Z_p$-extension'' of $K$).
	If we let $\Khat_{\cyc}$ denote the closure \index{$\Khat_{\cyc}$}
	of $K_{\cyc}$ in~$\C$,
	then $\Khat_{\cyc}$ is a perfectoid subfield of~$\C$.
\end{example}

\begin{example}[The Kummer case]
	\label{ex:kummer}
	If we choose a uniformizer $\pi$ of $K$,
	as well as a compatible system of $p$-power roots
	$\pi^{1/p^n}$ of $\pi$ (here, ``compatible'' has
	the obvious meaning, namely that $(\pi^{1/p^{n+1}})^p = \pi^{1/p^n}$),
	then we define $K_{\infty} = K(\pi^{1/p^{\infty}}) :=
	\bigcup_n K(\pi^{1/p^n}).$ \index{$K_{\infty}$}
	If we let $\Khat_{\infty}$ denote the closure
	of $K_{\infty}$ in~$\C$,
	then $\Khat_{\infty}$ is again a perfectoid subfield of~$\C$.
\end{example}

\begin{rem}
  \label{rem: Krasner Ax Tate Sen}Let~$L$ be an algebraic extension
  of~$\Qp$, and let~$\Lhat$ be the closure of~$L$ in~$\C$. By
  Krasner's lemma (see~\cite[Prop.\ 9.1.16]{Gabber--Ramero} for full
  details), the field~$\Lhat\otimes_L\Qpbar$ is an algebraic closure
  of~$\Lhat$, so that the absolute Galois groups of~$L$ and~$\Lhat$
  are canonically identified. The action of~$G_L$ on~$\Qpbar$ extends
  to an action on~$\C$, and by a theorem of
  Ax--Tate--Sen~\cite{MR0263786}, we have~$\C^{G_L}=\Lhat$. We will in
  particular make use of these facts in the cases ~$L=K_{\cyc}$
  and~$L=K_\infty$.
\end{rem}

% \mar{ME: Probably mention Ax-etc.\ theorem somewhere here, 
% 	to get control of Galois groups. TG: will this do? I suspect
%         it needs another sentence to say how we use it? ME: I just looked
% at the Gabber--Ramero reference, and it looks like this will apply to
% any algebraic extension of $\Q_p$, right?  (Because a direct limit
% of Henselian rings is Henselian --- see e.g.\ Lemma 04GI of the SP.)
% So maybe we can state it in that generality? TG: agreed, rewritten;
% comment out if you're happy.}


\subsection{Fields of norms}
If $L$ is
an infinite strictly arithmetically profinite (strictly APF) extension \index{strictly APF}
of $K$ in $\Qbar_p$, % \mar{ME: Have to add correct Fontaine--Wintenberger hypotheses here! ME: 
% APF seems to be the condition, but it also seems to assume that
% $L$ is Galois over $K$! ME: OK, not necessarily.  Should cite [FW],
% and maybe also [Br-Co].}
then we may form the field of norms $X_K(L)$, \index{field of norms}
as in~\cite{MR526137} and~\cite[\S 2]{MR719763}.
% \mar{ME: Maybe state extra structure on $X_K(L)$, e.g.\ is it valued in general? 
% 	complete?}
%\mar{ME: Find correct notation.}
%which is canonically identified with a subfield of~$L_{\infty}^{\flat}$.
This is a complete discretely valued field of characteristic $p$,
whose residue field is canonically identified with that of
$L$. 
(We don't define the notion of strictly APF here, but simply refer to
\cite[Def.~1.2.1]{MR719763}
for the definition.  We will only
apply these notions to Examples~\ref{ex:cyclo} and~\ref{ex:kummer},
in which case the end result of the field of norms
construction can be spelt out quite explicitly.)
The following theorem is well known.
\begin{thm}
	\label{thm:field of norms}
	Let $L$ be an infinite strictly APF extension of $K$ in $\Qbar_p$.
	\begin{enumerate}
		\item The closure $\widehat{L}$ of $L$ in $\C$
			is perfectoid.
		\item 
			There is a canonical
	embedding $X_K(L) \hookrightarrow (\widehat{L})^\flat$,
	which
	identifies the target with the completion of the perfect
	closure of the source.
\end{enumerate}
\end{thm}
\begin{proof}
	%\mar{ME: cite something for proof; e.g.\ this result is in [Br-Co].
	%ME: Actually, only there in the case of a virtually $\Z_p$-extension.}
  To show that $\widehat{L}$ is perfectoid, it suffices to show that
  it is not discretely valued, and that the absolute Frobenius on
  $\cO_L/p\cO_L$ is surjective.  % \mar{ME: There's probably some other
    % non-discretely valued hypothesis to check, which hopefully is
    % imposed by the infinite degree strictly APF assumption; otherwise
    % we just have to impose it. TG: I've seen something implicitly
    % saying this is automatic, but since I was exchanging emails with
    % Berger about other things I've just asked him. TG: Berger came
    % through, commenting out.}
  This surjectivity
  follows from \cite[Cor.~4.3.4]{MR719763}. It follows immediately
  from the description in~\cite[\S 1.4]{MR719763} of APF extensions as
  towers of elementary extensions that~$L$, and thus~$\widehat{L}$, is
  not discretely valued.  The canonical embedding of~(2) is
  constructed in \cite[\S 4.2]{MR719763}, and its claimed property is
   proved in \cite[Cor.~4.3.4]{MR719763}.
\end{proof}

\begin{example}
	\label{ex:Sen}
	A theorem of Sen \cite{MR0319949} 
	shows that
	if the Galois group of the Galois closure of $L$ 
	over $K$ is a $p$-adic Lie group, and the induced extension of
        residue fields is finite, then
	$L$ is a strictly APF extension of $K$.
	(Sen's theorem shows that the Galois
       closure of $L$ is strictly APF over $K$; see
       \cite[Ex. 1.2.2]{MR719763}. It then follows
       from 
       \cite[Prop.~1.2.3~(iii)]{MR719763}
       that $L$ itself is strictly APF over $K$.)
       As a consequence, we see that the theory of the field of norms,
       and in particular Theorem~\ref{thm:field of norms},
       applies in the
       cases of Examples~\ref{ex:cyclo} and~\ref{ex:kummer}.
       \end{example}

       \subsection{Thickening fields of norms}
       In the context of Theorem~\ref{thm:field of norms},
       given that one may thicken $(\widehat{L})^\flat$ to
       the flat $\Z/p^a$-algebra $W_a\bigl((\widehat{L})^\flat\bigr)$,
       it is natural to ask whether $X_K(L)$ admits a similar
       such thickening, such that the embedding $X_K(L) \hookrightarrow
       (\widehat{L})^\flat$
       lifts to an embedding of the
       corresponding thickenings.  One would furthermore like
       such a lifted embedding to be compatible with various auxiliary
       structures, such as the action of Frobenius on
       $W_a\bigl((\widehat{L})^\flat\bigr)$,
       or the action of the Galois group
       $\Gal(L/K)$, in the case when $L$ is Galois over $K$
       (i.e.\ one would like the image
       of this embedding to be stable under these actions).

       Since $X_K(L)$ is imperfect, there is no canonical
       thickening of $X_K(L)$ over $\Z/p^a$, and
       as far as we know, there is no simple or general answer to the 
       question of whether such thickenings and embeddings exist
       with desirable 
       extra properties, such as being Frobenius or Galois stable.
       However, in the cases of Examples~\ref{ex:cyclo} and~\ref{ex:kummer},
       such thickenings and thickened embeddings can be constructed
       directly, as we now recall.

       \subsection{The cyclotomic case}\label{subsubsec: cyclotomic}
       The extension $K(\zeta_{p^{\infty}})$ of $K$ is infinite and
       strictly APF, as well as being Galois over $K$.  If we write
       $\tGamma_K := \Gal(K(\zeta_{p^{\infty}})/K)$, then the
       cyclotomic character induces an embedding
       $\chi:\tGamma_K \hookrightarrow \Z_p^{\times}$.  Consequently,
       there is an isomorphism
       $\tGamma_K \cong \Gamma_K \times \Delta$, where
       $\Gamma_K \cong \Z_p$ and $\Delta$ is finite.  We have
       $K_{\cyc} = (K(\zeta_{p^{\infty}}))^{\Delta}$. Later in
       the book, $K$ will typically be fixed, and we will
       write~$\Gamma$ for~$\Gamma_K$.

       Suppose for a moment that $K = \Q_p$.
       If we choose a compatible system of $p^n$th roots of $1$,
       then these give rise in the usual way to an element
       $\varepsilon \in (\widehat{\Qp(\zeta_{p^{\infty}})})^{\flat}.$
              If we identify the field of norms $X_{\Q_p}(\Qp(\zeta_{p^{\infty}}))$ 
       with a subfield of $(\widehat{\Qp(\zeta_{p^{\infty}})})^\flat$ via the embedding
       $X_{\Q_p}(\Qp(\zeta_{p^{\infty}})) \hookrightarrow (\widehat{\Qp(\zeta_{p^{\infty}})})^\flat,$
       then $\varepsilon - 1 $ is a uniformizer of $X_{\Q_p}(\Qp(\zeta_{p^{\infty}}))$.  
       As usual, let $[\varepsilon]$ denote the Teichm\"uller 
       lift of $\varepsilon$ to an element of
       $W\bigl( \cO_{\widehat{\Qp(\zeta_{p^{\infty}})}}^\flat).$
       There is then a continuous embedding
       $$\Z_p[[T]] \hookrightarrow 
       W( \cO_{\widehat{\Qp(\zeta_{p^{\infty}})}}^\flat)$$
       (the source being endowed with its $(p,T)$-adic topology,
       and the target with its weak topology),
       defined via $T \mapsto [\varepsilon]- 1$.
       We denote the image of this embedding by $(\A'_{\Q_p})^+$.
       This embedding extends to an embedding
       $$\widehat{\Z_p((T))} \hookrightarrow W\bigl( (\widehat{\Qp(\zeta_{p^{\infty}})})^\flat
       \bigr)$$
       (here the source is the $p$-adic completion of the Laurent
       series ring $\Z_p((T))$),
       whose image we denote by $\A'_{\Q_p}$.

       We now return to the case of general $K$. We will compare 
       this case with the case of $\Q_p$. Since $K\Qp(\zeta_{p^\infty}) = K(\zeta_{p^\infty})$,
       the theory of the field of norms gives an identification
       $X_K(K(\zeta_{p^\infty})) = X_{\Q_p}(K(\zeta_{p^\infty}))$
       \cite[Rem.~2.1.4]{MR719763},
       and shows that this field 
       is a separable extension of $X_{\Q_p}(\Qp(\zeta_{p^\infty}))$
       \cite[Thm.~3.1.2]{MR719763},
       if we regard both as embedded in $\C^\flat$ via
       the embeddings of Theorem~\ref{thm:field of norms}. To ease
       notation, from now on we
       write~$\E'_K=X_K(K(\zeta_{p^\infty}))$,
       and~$\E_K:=(\E'_K)^\Delta$, so that ~$\E'_K/\E_K$ is a
       separable extension. We write $(\E'_K)^+$, $\E_K^+$ for the
       respective rings of integers. We write~$\varphi$ for the
       ($p$-power) Frobenius on~$E_K'$, $E_K$, $(\E'_K)^+$ and~ $\E_K^+$.
       % \mar{ME: Introduce $\E'_K/\E_K$ and $(\E'_K)^+/\E_K^+$ notation for 
       %         the fields of norms and their rings of integers. TG: is
       %       this the right thing to do?}


       We note now that $B'_{\Q_p} := \A'_{\Q_p}[1/p]$
       is a discretely valued field
       admitting $p$ as a uniformizer,
       with residue field $\E'_{\Q_p}$. 
       There is then a unique finite unramified extension of
       $\B'_{\Q_p}$ contained in the field $W(\C^\flat)[1/p]$
       with residue field $\E'_K$.
       We denote this extension by $\B'_K$; it is again
       a discretely valued field, admitting $p$ as a uniformizer,
       and we let $\A'_K$ denote its ring of integers;
       equivalently, we have $\A'_K = \B'_K \cap W(\C^\flat).$
       We see that $\A'_K$ is a discrete valuation ring,
       admitting $p$ as a uniformizer, and that
       $\A'_K/p \A'_K = \E'_K.$ There is a natural lift of the
       Frobenius~$\varphi$ from~$\E'_K$ to~$\A'_K$.
       
       
       The action of $\tGamma_K$ on $\E'_K=X_K(K(\zeta_{p^\infty}))$ induces
       an action of $\tGamma_K$ on $\A'_K$, and we write
       $\A_K := (\A'_K)^{\Delta}$. \index{$\A_K$}  This is the ring of integers in
       the discretely valued field $\B_K := \A_K[1/p] = (\B'_K)^{\Delta}$,
       and has residue field equal to $\E_K.$ The actions of
       $\tGamma_K$ on~$\A'_K$ and of~$\Gamma_K$ on~$\A_K$ commute with~$\varphi$.
       % \mar{ME: I don't think we've actually introduced $\varphi$ at this
       %         point; should we? TG: it was actually introduced at the
       %       end of the previous paragraph. I have now added another
       %       sentence about this. I think we definitely want it here, yes.}

       If we let $T'_K$ denote a lift of a uniformizer of
       $\E'_K$, and let $k'_{\infty}$ denote the residue field
       of $K(\zeta_{p^\infty})$, then there is an isomorphism
       $\widehat{W(k'_{\infty})((T'_K))} \iso \A'_K.$  (Here the source
       denotes the $p$-adic completion.)  Similarly if~$k_\infty$
       denotes the residue field of~$\Kcyc$, and~$T_K$ is a lift of a uniformizer of
       $\E_K$, then there is an isomorphism
       $\widehat{W(k_{\infty})((T_K))} \iso \A_K.$

       For a general extension~$K/\Qp$ it is hard to give an explicit
      formula for the actions of~$\varphi$ and~$\Gammat_K$ on
       $W(k_\infty')((T'_K))^{\wedge}$, but in some of our arguments
       it is useful to reduce to a special case where we can use
       explicit formulae. If~$K=K_0$ (that is, if~$K/\Qp$ is
       unramified) then we have such a description as follows: % . Fix an
       % element~$\varepsilon\in\C^\flat$ given by a compatible system
       % of nontrivial $p^n$th roots of unity (recall that by definition
       % we have $\cO_\C^\flat:=\varprojlim_{x\mapsto
       %   x^p}\cO_{\C}/p$). Then
       we have $\E'_{K_0}=k((\varepsilon-1))$,
       and ~$\A'_{K_0}=W(k((\varepsilon-1)))$, and
       we set~$T'_{K_0}=[\varepsilon]-1$, with the square brackets
       denoting the Teichm\"uller lift.

The actions of~$\varphi$ and~$\gamma\in\Gammat_{K_0}$
on~$T'_{K_0}\in\A_{K_0}'$ are given by the explicit formulae
\numequation\label{eqn: phi on T basic case with
  epsilon}\varphi(T'_{K_0})=(1+T'_{K_0})^p-1, \end{equation}
\numequation\label{eqn: gamma on T basic case with
  epsilon}\gamma(1+T'_{K_0})=(1+T'_{K_0})^{\chi(\gamma)},
\end{equation}where~$\chi:\Gammat_{K_0}\to\Z_p^\times$ denotes the
cyclotomic character. We set $(\A_{K_0}')^+=W(k)[[T'_{K_0}]]$, which is
visibly~$(\varphi,\Gammat_{K_0})$-stable. We set
$T_{K_0}=\tr_{\A_{K_0}'/\A_{K_0}}(T'_{K_0})$ and
$\A_{K_0}^+=W(k)[[T_{K_0}]]$; then we
have~$\A_{K_0}=W(k)((T_{K_0}))^{\wedge}$, and~$\A_{K_0}^+$ is
$(\varphi,\Gamma_{K_0})$-stable (if $p>2$ this is~\cite[Prop.\
A.3.2.3]{MR1106901}, and if~$p=2$ the same statements hold, as
explained in~\cite[\S1.1.2.1]{MR1693457}). After possibly
replacing~$T_{K_0}$ by ~$(T_{K_0}-\lambda)$ for some~$\lambda\in
pW(k)$, by~\cite[Prop.\ 4.2, 4.3]{MR3322781} we can and do assume that
$\varphi(T_{K_0})\in T_{K_0}\A_{K_0}^+$, and that $g(T_{K_0})\in
T_{K_0}\A_{K_0}^+$ for all~$g\in\Gamma_{K_0}$.


\begin{defn}\index{basic (field~$K$)}
  \label{defn: basic}We say that~$K$ is \emph{basic} if it is
  contained in~$K_0(\zeta_{p^\infty})$. % \mar{TG: check/add lemma about
    % this being equivalent to being abelian over $\Qp$?}
  (Although it won't play a role in what follows,
 we remark that this implies in particular that $K$ is abelian over $\Q_p$.)
\end{defn}% \mar{TG: for later purposes want $\A_L$ is free over $\A_K$
  % for finite extensions $L/K$}

If~$K$ is basic then we have
$K(\zeta_{p^\infty})=K_0(\zeta_{p^\infty})$, so
that~$\bE'_K=\bE'_{K_0}$, and we take $\A'_K=\A'_{K_0}$,
$T'_K=T'_{K_0}$, $(\A_K')^+=(\A_{K_0}')^+$, and similarly for~$\A_K$,
with the actions of~$\Gammat_K$ and~$\Gamma_K$ being the restrictions
of the actions of~$\Gammat_{K_0}$ and~$\Gamma_{K_0}$. 

If~$K$ is not basic then % ~$\A_K\subseteq W(\C^\flat)$ can be defined as
% being the \'etale extension of~$\A_{K_0}$ lifting the
% extension~$\bE_K/\bE_{K_0}$, and
as explained above, we may still choose some~$T_K'$ so that
$\A_K'=W(k_\infty')((T'_K))^{\wedge}$, and some~$T_K$ so that $\A_K=W(k_\infty)((T_K))^{\wedge}$. % . In all cases we
% have~$\A_K=(\A_K')^{\Delta_K}$, and we
% take~$T_K=\tr_{\A_K'/\A_K}(T_{K'})$, so that
% $\A_K=W(k_\infty)((T_K))^{\wedge}$.
% \mar{TG: hopefully this is correct
%   although I didn't try to think it through, can imagine there could
%   be some stupidity related to the extension $k'_\infty/k_\infty$ - of
% course should also check if that extension is actually just
% trivial. TG: I think there is definitely an issue here if~$p=2$ and
% $\Delta$ has order divisible by 2, but maybe it's sorted out in Herr?}
Having done so, we set $(\A_K')^+=W(k'_\infty)[[T'_K]]$,  $\A_K^+=W(k_\infty)[[T_K]]$, where the topology on~$W(k_\infty)[[T]]$
is as usual the~$(p,T)$-adic topology.\index{$\A_K^+$} % If~$K$ is basic then it follows
% from~\eqref{eqn: phi on T basic case with
%   epsilon} and~ \eqref{eqn: gamma on T basic
% case with epsilon} that~$(\A'_K)^+$ is $\varphi$- and
% $\Gammat_K$-stable, so $\A_K^+$ is $\varphi$- and $\Gamma_K$-stable. 
% $W(k_\infty)((T))^{\wedge}\isoto\A_K^+$, where the completion is with respect
% to the~$p$-adic topology.
% see p223 of Brinon--Conrad notes
% for this claim

\begin{rem}\label{rem: not phi stable T}
For an arbitrary~$K$, it is not necessarily possible to choose~$T_K$ so
  that~$\A_K^+$ is $\varphi$-stable and $\Gamma_K$-stable (and similarly
  for~$(\A_K')^+$). % \mar{TG: note
    % maybe this will all be written firstly for the $\tGamma$ version.}
  In fact, it is not even possible to choose~$T_K$ so
  that~$\A_K^+$ is $\varphi$-stable. Indeed, by~\cite[Prop.\
  4.2, 4.3]{MR3322781}, % \mar{TG: I don't understand the proof of this result
    % at the moment, but I'm inclined to believe the conclusion}
  if~$\A_K^+$ is $\varphi$-stable then it is also~$\Gamma_K$-stable. As
  explained in~\cite[\S1.1.2.2]{MR1693457}, it is possible to choose $T_K$
  so that $\A_K^+$ is $\varphi$ and $\Gamma_K$-stable precisely
  when %this implies that~$K$ is
  $K$ is contained in an unramified extension of~$K_0(\zeta_{p^\infty})$,
  i.e.\ when $K$ is %that 
  abelian over~$\Qp$.
  % contained in the maximal unramified extension of the cyclotomic
  % $\Zp$-extension of~$K_0$.% \mar{TG: check - is this the same as being
    % basic? TG: edited, commenting out.}
  % f the maximal unramified subfield of~$K$.
  % \mar{TG: I
%   think that's what that says! This whole argument should of course be
% checked.}
\end{rem}

       %\mar{ME: Define $\A_K$ and $\A^+_K$.}
%       \mar{ME: Got as far as defining $\A_K$.}

       \subsection{The Kummer case}\label{subsubsec: Kummer}
       
       % \mar{ME: Define $\gS$ and $\cO_{\cE}$.}  \mar{TG: old material
       %   from later in the paper temporarily dumped in}
       Set~$\gS=W(k)[[u]]$, \index{$\gS$} with a~$\varphi$-semi-linear
       endomorphism~$\varphi$ determined by~$\varphi(u)=u^p$, and
       let~$\cO_{\cE}$ \index{$\cO_{\cE}$} be the $p$-adic completion of~$\gS[1/u]$. The
       extension $K_\infty= K(\pi^{1/p^{\infty}})$ of~$K$ is infinite
       and strictly APF (but not Galois). The choice of compatible
       system of $p$-power roots of~$\pi$ gives an element
       $\pi^{1/p^\infty}\in \cO_{\widehat{K_\infty}}^\flat$, and there
       is a continuous $\varphi$-equivariant
       embedding \[\gS\into W(\cO_{\widehat{K_\infty}}^\flat)\]
       which lifts the morphism $X_K(K_{\infty}) \hookrightarrow
(\widehat{K}_{\infty})^{\flat},$
       defined by
       sending $u\mapsto [\pi^{1/p^\infty}]$.
%\mar{ME: This lifting statement isn't necessarily needed, but ties in
%with the theme of~(2.1.8), and is also slightly indirectly referred 
%to at the bottom of p.~22 below. TG: is there a missing $\flat$ on $\widehat{K}_{\infty}$? ME: Yes, fixed}
This embedding extends
       to a continuous $\varphi$-equivariant
       embedding \[\cO_{\cE}\into W((\widehat{K_\infty})^\flat).\]

% 	If we let $\Khat_{\infty}$ denote the closure
% 	of $K_{\infty}$ in $\C$,
% 	then $\Khat_{\infty}$ is again a perfectoid subfield of $\C$.
       
% We also need a variant which again follows from the results
% of~\cite{MR1106901}; the integral version of this theory was first
% studied by Breuil and Kisin, see~\cite[\S1]{KisinModularity}, and will
% be employed in our study of crystalline deformation rings. Fix a uniformiser~$\pi$ of~$K$, and fix a compatible
% system of~$p^n$th roots ~$\pi^{1/p^n}$ of~$\pi$ in~$\overline{K}$.
% Let~$K_\infty=K(\pi^{1/p^\infty})=\cup_nK(\pi^{1/p^n})$ be the
% infinite (non-Galois) extension of~$K$ generated by these $p$-power
% % roots of~$\pi$.
%  There is a
% $\varphi$-equivariant embedding $\gS\into\Ainf$, sending~$u$ to the
% Teichm\"uller lift of the element~$\pi^\flat=(\pi^{1/p^n})_n\in\cO_\C^\flat$
% determined by our fixed choice of~$\pi^{1/p^n}$. This extends to an
% embedding~$\cO_{\cE}\into W(\C^\flat)$, and~$G_{K_\infty}$ acts
% trivially on~$\cO_{\cE}$.
% \mar{TG: once Section~\ref{subsubsec: Kummer}
% is written we presumably just refer to that instead.}



      % \mar{ME: Now discuss about lifting the embedding of
%	       Theorem~\ref{thm:field of norms} modulo powers of $p$,
%	       and describe the 
%       cases of Examples~\ref{ex:cyclo} and~\ref{ex:kummer}
%       more explicitly.}

% \subsection{Old material still to be integrated}Let $K(\zeta_{p^\infty})$ be the Galois extension of $K$ obtained by
% adjoining all $p$-power roots of unity, and write
% $\Gal\bigl(K(\zeta_{p^\infty})/K\bigr)=\Gammat_K
% =\Gamma_K\times\Delta_K$, where~$\Delta_K$ is finite
% and~$\Gamma_K\cong\Zp$. Set~$\Kcyc:=K(\zeta_{p^\infty})^{\Delta_K}$,
% so that $\Gal(\Kcyc/K)=\Gamma_K$. We write~$k_\infty'$, $k_\infty$ for
% the residue fields of~$K(\zeta_{p^\infty})$, $\Kcyc$ respectively;
% these are both finite extensions of~$k$. In the body of the paper, $K$
% will typically be fixed, and we will write~$\Gamma_K$
% for~$\Gamma$. 
% Let~$\widehat{K(\zeta_{p^\infty})}$
% and~$\widehat{\Kcyc}$ denote the $p$-adic completions
% of~$K(\zeta_{p^\infty})$ and~$\Kcyc$ respectively; these are
% perfectoid subfields of~$\C$ (see for example~\cite[Ex.\ 1.3.5,
% 1.3.6]{MR3412585}).  The action of~$G_K$ on~$\overline{K}$ extends to
% an action on~$\C$, and by Krasner's lemma a theorem of Ax--Sen--Tate we
% have natural identifications~$G_{K(\zeta_{p^\infty})}=G_{\widehat{K(\zeta_{p^\infty})}}$,
% $G_{\Kcyc}=G_{\widehat{\Kcyc}}$. 
% \mar{TG: this is correct, right?
  % It at least gives the correct fixed field, and I think is what we
  % want to use later (and if it doesn't hold then I guess these Galois
  % groups are what we should be using, anyway). TG: replace $\C$ with
  % alg closure.}

% The theory of the field of norms associates $\varphi$-stable subfields
% ~$\bE_K'$ and~$\bE_K=(\bE'_K)^\Delta$ of~$\C^\flat$ to the extensions
% ~$K(\zeta_{p^\infty})/K$ and~$\Kcyc/K$ respectively.  These are
% complete discretely valued fields with residue fields~$k'_\infty$
% and~$k_\infty$ respectively, and inherit actions of~$\Gammat_K$
% and~$\Gamma_K$ respectively from the action of~$G_K$
% on~$\C^\flat$. More generally, for any finite extension~$L/K$, we have
% separable extensions~$\bE_L'/\bE_K'$ and~$\bE_L/\bE_K$ of respective
% degrees~$[K(\zeta_{p^\infty})L:K(\zeta_{p^\infty})]$
% and~$[\Kcyc L:\Kcyc]$.

% The field~$\bE_K'$, together with the actions of~$\varphi$
% and~$\Gammat_K$, can be lifted to a subring~$\A_K'\subset
% W(\C^\flat)$. We set~$\A_K=(\A_K')^{\Delta_K}$. % Via the theory of
% % the field of norms, one associates to~$K$ a subring
% % $\A_K\subseteq W(\C^\flat)$, which is $\varphi$- and $G_K$-stable, and on
% % which the action of~$G_K$ factors
% % through~$\Gamma=\Gal(\Kcyc/K)$.
% We give~$\A_K'$ and~$\A_K$ the
% subspace topologies, and again refer to this topology as the weak
% topology; so the action of~$\varphi$ on~$\A_K'$ and~$\A_K$ is continuous, as are
% the maps $\Gammat_K\times\A_K'\to\A_K'$ and $\Gamma_K\times\A_K\to\A_K$.
% We
% set~$\A_K^+:=\A_K\cap\Ainf$ % \mar{TG: I think/hope this is the right
%   % thing to do?}
% (and again endow it with the subspace topology) which by definition
% has continuous commuting actions of~$\varphi$
% and~$\Gamma$.
% There are isomorphisms of topological rings
% $W(k_\infty')((T'_K))^{\wedge}\isoto\A_K'$,
% $W(k_\infty)((T_K))^{\wedge}\isoto\A_K$ where $T'_K$, $T_K$ are
% variables and the completion is with respect to the~$p$-adic
% topology, %\mar{TG: I think the residue field here is correct?}
% but t

% There is in general no canonical choice of coordinates~$T'_K$,
% $T_K$.
% Indeed, the choice of~$\A_K'\subseteq W(\C^\flat)$ is itself not
% canonical, but once this choice has been made, it determines a choice
% of~$\A_K\subset\A_L\subseteq W(\C^\flat)$ for all finite extensions~$L/K$.

\subsection{Graded and rigid analytic techniques}In
Appendix~\ref{sec: rigid analytic perspective} we prove a number of
results using rigid analysis and graded techniques, which we will
apply in the next section in order to prove some basic facts about our
coefficient rings. In order to apply these results to the various
rings introduced above, we need to verify the hypotheses introduced in
\ref{subsec:setting}, which are as follows. There is an Artinian 
 local ring $R$, which in our present context we take to be the ring~$\Z/p^a$.
Then we work with a $\Z/p^a$-algebra $C^+$, and an element $u \in C^+$, satisfying the
following properties:
\begin{enumerate}[label=(\Alph*)]
	\item\label{item: Cplus A} $u$ is a regular element (i.e.\ a non-zero divisor) of $C^+$.
	\item\label{item: Cplus B}  $C^+/u$ is a flat (equivalently, free) $\Z/p^a$-algebra.
	\item\label{item: Cplus C}  $C^+/p$ is a rank one complete valuation ring, and
		the image of $u$ in $C^+/p$ (which
		is necessarily non-zero, by \eqref{item: Cplus A}) is of
		positive valuation (i.e.\ lies in the maximal
		ideal of $C^+/p$).
\end{enumerate}
The following lemma gives the examples of this construction that we
will use in the next section.

\begin{lem}
  \label{lem: actual examples of Cplus and C that we use}The following
  pairs~$(C^+,u)$ satisfy axioms \emph{\ref{item: Cplus A}--\ref{item:
    Cplus C}} above.
  \begin{enumerate}
  \item $C^+=W_a(\cO_F^\flat)$, where~$F$ is any perfectoid
    field, and~$u$ is any element of~$W_a(\cO_F^\flat)$
    whose image in $\cO_F^\flat$ is of positive
    valuation. %\mar{TG: is this a general enough~$u$? ME: Made $u$ maximally
% general.}
  \item $C^+=\A_K^+/p^a$, $u=T_K$.
  \item $C^+=W_a(k)[[u]]$.
  \end{enumerate}

\end{lem}
\begin{proof}%\mar{TG: finish ME: Had a go.}
	Axioms~\ref{item: Cplus A} and~\ref{item:
    Cplus C} are clear from the definitions in each case, so we need
  only verify that $C^+/u$ is flat over~$\Z/p^a$ in each case.
  This is immediate in cases~(2) and (3), so we focus on case~(1).
  Rather than checking~\ref{item: Cplus B} directly in this case,
  we instead check that $W(\cO_F^\flat)/ u$ is flat over $\Z_p$;
  base-changing to $\Z/p^a$ then gives us~\ref{item: Cplus B}.
  For this, it suffices to observe that multiplication by $u$ is injective
  on each term of the short exact sequence
  $$0 \to W(\cO_F^\flat) \buildrel p \cdot \over \longrightarrow
  W(\cO_F^\flat) \longrightarrow \cO_F^\flat \to 0 \, ;$$
  the snake lemma then shows that multiplication by $p$ is injective
  on $W(\cO_F^{\flat})/u$, as required.
  %By
  %Lemma~\ref{lem:graded flat}\mar{TG: is this right?} we can reduce to
  %the case~$a=1$, which is trivial.  
\end{proof}

%\mar{ME: Add statement about boundedness of Galois actions}
We also note the following lemma.

% \mar{ME: In the statement here, should $\cO_{\cE}$ actually
% be $\cO_{\cE}^{\ur},$ which I'm not sure has been introduced yet?}
\begin{lem}
  \label{lem: boundedness of Galois actions in our concrete
    settings}The actions of~$G_K$ on~$W_a(\C^\flat)$, % ;
  % of~$G_{K_\infty}$ on% ~$\gS/p^a$ and
  % ~$\cO_{\cE}/p^a$;
  and of~$\Gamma_K$
  on~$\A_K/p^a$ and $\tA_K/p^a$, are continuous and bounded in the sense
  of Definition~{\em \ref{adefn: bounded group action}}.
\end{lem}
\begin{proof}It suffices to show that the action of $G_K$
  on~$W_a(\C^\flat)$ is continuous and bounded, as the other cases follow by
  restricting this action to the corresponding subgroup and subring.
  By Remark~\ref{arem: bounded doesn't depend on u}, we may assume
  that ~$u=[v]$ for some~$v\in\cO_{\C}^\flat$ with positive valuation;
  since the action of~$G_K$ on $\cO_{\C}^\flat$
  preserves the valuation, we then
  have~$G_K\cdot u^{M}W_a(\C^\flat) =  u^{M}W_a(\C^\flat)$ for
  all~$M\in\Z$, so the action is continuous and bounded.
\end{proof}

\section{Coefficients}\label{subsec: coefficients}
%\mar{TG: we now introduce some of these rings multiple times in the paper.}
Our main concern throughout this book will be families of
\'etale $(\varphi,\Gamma)$-modules; an auxiliary role will also be played by
families of (various flavours of) \'etale $\varphi$-modules. Such a family
will be parameterized by an algebra of coefficients, typically 
denoted by $A$.   Since we will work throughout
in the context of formal algebraic stacks over $\Spf \Z_p$ (or closely 
related contexts), we will always assume that our coefficient ring $A$ is
a $p$-adically complete $\Zp$-algebra.
Frequently, we will work modulo some a fixed power $p^a$ of $p$,
and thus assume that $A$ is actually a $\Z/p^a$-algebra
(and sometimes we impose further conditions on $A$,
such as that of being Noetherian, or even of finite type over $\Z/p^a$). 

It is often convenient to introduce an auxiliary base ring for our
coefficients,
which we will take to be the ring of integers $\cO$ in a finite extension
$E$ of $\Q_p$; in this case, we will let $\varpi$ denote a uniformizer
of $\cO$, and $A$ will be taken to be a $p$-adically complete
(or, equivalently, a $\varpi$-adically complete) $\cO$-algebra,
or, quite frequently, an $\cO/\varpi^a$-algebra
(perhaps Noetherian, or even of finite type), for some power
$\varpi^a$ of $\varpi$.


For the moment, we put ourselves in the most general case;
that is, we assume that $A$ is a $p$-adically complete $\Z_p$-algebra,
and define versions of the various rings considered in
Section~\ref{subsec:rings} ``relative to $A$''.

Let $F$ be a perfectoid closed subfield of $\C$,
and if $a \geq 1$,
let $v$ denote an element of the maximal ideal
of $W_a(\cO_F^\flat)$
whose image
in $\cO_F^\flat$ is non-zero.
We then set
$$
W_a(\cO_F^\flat)_A =
W_a(\cO_F^\flat) \cotimes_{\Z_p}  A
:= \varprojlim_i \bigl( W_a(\cO_F^\flat)\otimes_{\Z_p}A\bigr)/v^i$$
(so that the indicated completion is the $v$-adic completion).
Note that any two choices of $v$ induce the same topology
on $W_a(\cO_F^\flat)\otimes_{\Z_p} A$, so
that 
$W_a(\cO_F^\flat) \cotimes_{\Z_p}  A$
is well-defined independent of the choice of $v$.
We then define
$$W_a(F^\flat)_A =  W_a(F^\flat) \cotimes_{\Z_p} A
:= W_a(\cO_F^\flat)_A[1/v];$$
this ring is again well-defined independently of the choice of $v$.

There are natural reduction maps
$$W_a(\cO_F^\flat) \cotimes_{\Z_p}  A \to
W_b(\cO_F^\flat) \cotimes_{\Z_p}  A$$
and
$$W_a(F^\flat) \cotimes_{\Z_p}  A \to
W_b(F^\flat) \cotimes_{\Z_p}  A,$$
if $a \geq b$,
so that we may define
$$
W(\cO_F^\flat)_A =
W(\cO_F^\flat)\cotimes_{\Z_p} A := \varprojlim_a W_a(\cO_F^\flat)_A,
$$
and similarly
$$
W(F^\flat)_A =
W(F^\flat)\cotimes_{\Z_p} A := \varprojlim_a W_a(F^\flat)_A.
$$

For certain choices of $F$, we introduce alternative notation
for the preceding constructions. %as we now explain.
In the case when $F$ is equal to $\C$ itself,
we write
$$\AAinf{A} :=  W(\cO_\C^\flat)_A.$$
In the case when $F$ is equal to $\widehat{K}_{\cyc}$,
we write 
% $$\tA^+_{K,A} := W\bigl((\cO_{\widehat{K}_{\cyc}})^\flat\bigr)_A$$
% and\mar{TG: I don't think we use the plus version, so I'm removing it}
$$\tA_{K,A} := W\bigl((\widehat{K}_{\cyc})^\flat\bigr)_A.$$
In the case when $F$ is equal to $\widehat{K}_{\infty}$,
we write
$$\tOEA := W\bigl((\widehat{K}_{\infty})^\flat\bigr)_A.$$

We next introduce various imperfect coefficient rings with $A$-coefficients:
If $T_K$ denotes a lift to $\A_K$ of a uniformizer of $\E_K$, %$X_K(K_{\cyc}),$
defining a subring~$\A_K^+$ of $\A_K$ as in Section~\ref{subsec:rings}, 
then we write
\[\A^+_{K,A}=\A^+_K\cotimes_{\Zp}A := \varprojlim_{n} \A^+_K/(p,T_K)^n
	\otimes_{\Z_p} A = \varprojlim_{m} (\varprojlim_n \A^+_K/(p^m,T_K^n)
	\otimes_{\Z_p} A), \]
and
\[\A_{K,A} =\A_K\cotimes_{\Zp}A :=
	\varprojlim_{m} \bigl(\varprojlim_n (\A^+_K/(p^m,T_K^n)
	\otimes_{\Z_p} A) [1/T_K]\bigr). \]
Similarly, if $u$ denotes the usual element of $\gS$, whose
reduction modulo $p$
is a uniformizer of $X_K(K_{\infty})$, then we write 
\[\gS_{A}=\gS\cotimes_{\Zp}A := \varprojlim_{n} \gS/(p,u)^n
	\otimes_{\Z_p} A = \varprojlim_{m} (\varprojlim_n \gS/(p^m,u^n)
	\otimes_{\Z_p} A), \]
and
\[\OEA =\cO_\cE \cotimes_{\Zp}A :=
	\varprojlim_{m} \bigl(\varprojlim_n (\gS/(p^m,u^n)
	\otimes_{\Z_p} A) [1/u]\bigr). \]
Note that by definition the ring~$\gS_{A}$ is a power
series ring over~$W(k)\otimes_AA$, and $\OEA$ is the $p$-adic
completion of the corresponding Laurent series ring;  and similarly
for the rings~ $\A^+_{K,A}$ and
~ $\A_{K,A}$.


%   The existence and continuity of the action of~$\varphi$ on~$\gS_A$,
% $\OEA$, $\A_{K,A}$ (and, in the case that~$K$ is basic,
% on~$\A^+_{K,A}$) is immediate from Lemma~\ref{lem: varphi extends to
%   A}. The same argument applies to our other rings, replacing~$\A^+$
% with~$W(\cO_F^\flat)$ for an appropriate~$F$.

% Similarly, the actions of $\Gamma_K$, $G_K$, $G_{K_\infty}$ (as
% appropriate) on our rings extend to continuous
% actions on the rings with coefficients.\mar{TG: need to give an
%   argument - is it best to just explain one of them?}
% (For the $\varphi$-actions, this follows from
% Lemma~\ref{lem: varphi extends to A}; the Galois
% actions can be handled in much the same way.)
% \mar{ME: Just want to flag that $\varphi$ is discussed here, but a proof
% 	if punted till later on. TG: I don't think this is a huge
%         issue, but actually I propose to move this material back to
%         here anyway.}

% Recall that if~$\Acirc$ is a $p$-adically
%    complete $\cO$-algebra, then~$\Acirc$ is said to be \emph{topologically of finite
%    type over~$\cO$} if it can be written as a quotient of a
%  restricted formal power series ring in finitely many variables~$\cO\langle\langle
%  X_1,\dots,X_n\rangle\rangle$; equivalently, if and only if~$\Acirc\otimes_{\cO}k$ is a finite
%    type $k$-algebra (\cite[\S0, Prop.\
%    8.4.2]{MR3752648}). In particular, if~$\Acirc$ is a
%    $\cO/\varpi^a$-algebra for some~$a\ge 1$, then~$\Acirc$ is
%    topologically of finite type over~$\cO$ if and only if it is of
%    finite type over~$\cO/\varpi^a$.
%      
%      We now record that, at least if ~$A$ is a finite type\mar{TG:
%        replace with topologically of finite type?}
%      $\Z/p^a\Z$-algebra, the various natural maps between these rings
%      are in fact (faithfully) flat injections.

%We make use of the
%      following results from~\cite{MR2774689}.
%
%      \begin{prop}
%%\mar{TG: might need a $(v,p)$ version of several
%%          things around here? This particular result is fine with an
%%          ideal instead of an element, of course.}\mar{TG: maybe use
%%          Remark 4.31 of \cite{2016arXiv160203148B}?}
%        \label{prop: FGK flat}Let~$R$ be a ring which is $x$-adically
%        complete for some $x\in R$, and let ~$S$ be an~$R$-algebra
%        which is also $x$-adically complete, and for which $S[1/x]$ is Noetherian.
%\begin{enumerate}
%\item
%Suppose that in addition $R[1/x]$ is Noetherian,
%and that for
%        each~$k~\ge~1$, the induced morphism $R/x^k\to S/x^kS$ is
%        flat.   Then $R\to S$ is flat.
%\item
%Suppose and that for each~$k\ge 1$, the induced morphism $R/x^k\to S/x^k$ is
%        faithfully flat. 
%Then~$R[1/x]$ is Noetherian, and $R\to S$ is faithfully flat.
%\end{enumerate}
%      \end{prop}
%      \begin{proof}
%The claim of~(1) is a special case of the discussion at the beginning
%        of~\cite[\S 5.2]{MR2774689}.
%The claim of~(2)
%        is a special case of~\cite[Prop.\ 5.2.1]{MR2774689}.
%      \end{proof}

%      We also note the following lemma.
%
%      \begin{lemma}
%	      \label{lem:Witt flatness}
%	      If $F \hookrightarrow F'$ is an inclusion of
%	      perfectoid fields in characteristic $p$,
%	      then the induced morphism $W_a(\cO_F)
%	      \to W_a(\cO_{F'})$ 
%	      is faithfully flat, for any~$a~\geq~1$.
%      \end{lemma}
%      \begin{proof}
%Since $\cO_F$ is a B\'ezout ring,
%being flat over $\cO_F$ is the same as being
%torsion free, and so the inclusion
%$\cO_{F} \hookrightarrow \cO_{F'}$
%is a flat morphism,
%and thus even faithfully flat, being a local morphism of local rings. 
%A standard grading argument, applying Lemma~\ref{lem:graded flat} to 
%the $p$-adic filtrations on source and target,
%then shows that the inclusion
%$W_a(\cO_{F}) \hookrightarrow W_a(\cO_{F'})$,
%is faithfully flat, for each~$a~\geq~1$.
%\end{proof}

\subsection{A digression on flatness and completion}
We record some results giving sufficient conditions 
for flatness to be preserved after passage to certain inverse limits,
beginning with the
      following results from~\cite{MR2774689}.

      \begin{prop}
%\mar{TG: might need a $(v,p)$ version of several
%          things around here? This particular result is fine with an
%          ideal instead of an element, of course.}\mar{TG: maybe use
%          Remark 4.31 of \cite{2016arXiv160203148B}?}
        \label{prop: FGK flat}Let~$R$ be a ring which is $x$-adically
        complete for some $x\in R$, and let ~$S$ be an~$R$-algebra
        which is also $x$-adically complete, and for which 
$R[1/x]$ and $S[1/x]$ are both Noetherian.
If for each~$k~\ge~1$, the induced morphism $R/x^k\to S/x^kS$ is {\em (}faithfully{\em )}
flat, then the morphism $R\to S$ is {\em (}faithfully{\em )} flat.
%$S[1/x]$ is Noetherian.
%\begin{enumerate}
%\item
%Suppose that in addition $R[1/x]$ is Noetherian,
%and that for
%        each~$k~\ge~1$, the induced morphism $R/x^k\to S/x^k$ is
%        flat.   Then $R\to S$ is flat.
%\item
%Suppose that for each~$k\ge 1$, the induced morphism $R/x^k\to S/x^k$ is
%        faithfully flat. 
%Then~$R[1/x]$ is Noetherian, and $R\to S$ is faithfully flat.
%\end{enumerate}
      \end{prop}
      \begin{proof}
The flatness claim is a special case of the discussion at the beginning
        of~\cite[\S 5.2]{MR2774689}.
The faithful flatness claim 
        is a special case of~\cite[Prop.\ 5.2.1~(2)]{MR2774689}.
It also follows from the flatness claim together with
Lemma~\ref{lem:faithful flatness} below.
%The claim of~(1) is a special case of the discussion at the beginning
%        of~\cite[\S 5.2]{MR2774689}.
%The claim of~(2)
%        is a special case of~\cite[Prop.\ 5.2.1]{MR2774689}.
      \end{proof}

\begin{remark}
Under the faithful flatness hypothesis of the preceding proposition,
it is proved in~\cite[Prop.\ 5.2.1~(1)]{MR2774689}
that Noetherianness of $R[1/x]$ is a consequence of
Noetherianness of $S[1/x]$.   This is a kind of {\em fpqc} descent
result for this property, which however we won't need in the present book.
\end{remark}

We next present
the following variation on
a result of Bhatt--Morrow--Scholze.%~\cite[Rem.\ 4.31]{2016arXiv160203148B}.

\begin{prop}
\label{prop:BMS flatness}
Let $R$ be a Noetherian ring which is $x$-adically complete, for some element
$x \in R$, let $S$ be an $x$-adically complete $R$-algebra,
and suppose that $S/x^n$ is {\em (}faithfully{\em )} flat over $R/x^n$ for every~$n~\geq~1$.
Then $S$ is {\em (}faithfully{\em )} flat over~$R$.
\end{prop}

\begin{remark}
In~\cite[Rem.\ 4.31]{2016arXiv160203148B}, the authors prove
the above proposition in the particular
case when $R$ and $S$ are flat $\Z_p$-algebras, and $x$ equals~$p$.
In this case they only need to assume that $S/p$ is flat over~$R/p$.
(Since $R$ and $S$ are both flat over~$\Z_p$,
it follows automatically that each $S/p^n$ is flat over~$R/p^n$.)  
Their proof makes use of the following key ingredients: that for a $p$-adically complete 
ring, any pseudo-coherent complex is derived complete; that for a complex 
over a $p$-torsion
free ring, the derived $p$-adic completion may be computed ``naively'';
and the Artin--Rees lemma for the $p$-adically complete and Noetherian ring~$R$.

Our argument is identical to theirs, once we confirm that these ingredients remain
available (with $p$-adic completions being replaced by $x$-adic completions). 
The first statement holds quite generally~\cite[\href{https://stacks.math.columbia.edu/tag/0A05}{Tag
    0A05}]{stacks-project}, and of course Artin--Rees holds
for any Noetherian ring~$R$.  Thus the main point is 
to verify that we may compute derived $x$-adic completions for the (generally
non-Noetherian ring) $S$ ``naively''.   
By~\cite[\href{https://stacks.math.columbia.edu/tag/0923}{Tag 0923}]{stacks-project}, 
this is possible provided that the $x$-power torsion in $S$ is bounded.  So
our task is to verify this boundedness (under our given hypotheses).
\end{remark}

We begin with the following general criterion for a ring to have bounded torsion.

\begin{lemma}
\label{lem:bounded torsion}
If $x$ is an element of the ring  $R$, then the following are equivalent:
\begin{enumerate}
\item
$R[x^m] = R[x^{\infty}]$
\item 
The morphism $R[x^i] \to R/x^m$ is injective for some~$i\geq 1$.
\item 
The morphism $R[x^i] \to R/x^m$ is injective for every~$i\geq 1$.
\end{enumerate}
\end{lemma}
\begin{proof}
 This is straightforward. Indeed, if the morphism $R[x^i] \to R/x^m$ is
 injective for some~$i\geq 1$, then it is in particular injective
 if~$i=1$. Suppose then that injectivity holds for~$i=1$.
 If $y\in R$ is such that~$x^{m+n}y=0$ for some~$n\ge 1$,
 then $x^{m+n-1}y$ is in the kernel of the morphism $R[x] \to R/x^m$,
 so $x^{m+n-1}y=0$, and by an easy induction we have $x^my=0$, as
 required.
 
 Conversely, if $t$ is in the kernel of the morphism $R[x^i] \to R/x^m$
 for some~$i$, then we can write $t=x^my$, and we have $x^it=0$; so
 $x^{m+i}y=0$. If $R[x^m] = R[x^{\infty}]$, then we have $x^my=0$, so
 $t=0$, as required.
\end{proof}

\begin{lemma}
\label{lem:torsion exact one}
If $x$ is any element of the ring~$R$, then for any $i,n\ge 1$ we have the exact sequence
$$R[x^i] \to R/x^n \buildrel x^i \over \longrightarrow R/x^{n+i} \to R/x^i \to 0.$$
\end{lemma}
\begin{proof}
This is immediately verified.
\end{proof}

\begin{lemma}
\label{lem:torsion exact two}
If $x$ is an element of the ring~$R$, and if $R[x^m] = R[x^{\infty}]$,
then 
$$0 \to R[x^i] \to R/x^n \buildrel x^i \over \longrightarrow R/x^{n+i} \to R/x^i \to 0$$
is an exact sequence, for any~$n~\geq~m$ and any $i \geq 1$.
\end{lemma}
\begin{proof}
This follows from Lemmas~\ref{lem:bounded torsion} and~\ref{lem:torsion exact
one}.
\end{proof}

\begin{proof}[Proof of Proposition~\ref{prop:BMS flatness}]
Since $R$ is Noetherian, we see that $R[x^m] = R[x^{\infty}]$ for some $m \geq 0$.
Thus, by Lemma~\ref{lem:torsion exact two},
for each $n~\geq~m$ and each $i \geq 1$,
we obtain an exact sequence
$$0 \to R[x^i] \to R/x^n \buildrel x^i \over \longrightarrow R/x^{n+i} \to R/x^i \to 0$$
Tensoring with the flat $R/x^{n+i}$-algebra $S/x^{n+i},$
we obtain an exact sequence
$$0 \to S\otimes_R R[x^i] \to S/x^n \buildrel x^i \over \longrightarrow S/x^{n+i} \to S/x^i \to 0$$
Passing to the inverse limit over $n$ (and taking into account that all the 
transition morphisms are surjective),
we obtain an exact sequence of $R/x^{n+i}$-modules
$$0 \to S\otimes_R R[x^i] \to S \buildrel x^i \over \longrightarrow S \to S/x^i \to 0$$
In particular, we find that $S\otimes_R R[x^i] \iso S[x^i]$,
for every~$i$, and so in particular
$S[x^m] = S[x^{\infty}]$, so that $S$ has bounded $x$-power torsion.

We now follow the proof of~\cite[Rem.\ 4.31]{2016arXiv160203148B}.
By (for example) \cite[\href{https://stacks.math.columbia.edu/tag/00M5}{Tag
    00M5}]{stacks-project}, it is enough to show that if~$M$ is a
  finitely generated $R$-module, then $M\otimes^{\mathbf{L}}_RS$ has
  cohomology concentrated in degree~$0$. Since~$R$ is Noetherian,
  \cite[\href{https://stacks.math.columbia.edu/tag/066E}{Tag
    066E}]{stacks-project} shows that $M$ is pseudo-coherent when
regarded as a complex
  of~$R$-modules in a single degree,
and so $M\otimes^{\mathbf{L}}_RS$ is a pseudo-coherent complex
  of~$S$-modules,
  by~\cite[\href{https://stacks.math.columbia.edu/tag/0650}{Tag
    0650}]{stacks-project}. Also,
since~$R$ and~$S$ are $x$-adically
  complete, it then follows
  from~\cite[\href{https://stacks.math.columbia.edu/tag/0A05}{Tag
    0A05}]{stacks-project} that both~$M$ and
  $M\otimes^{\mathbf{L}}_RS$ are derived $x$-adically complete.

Since $S$ has bounded $x$-power torsion,
it follows
from~\cite[\href{https://stacks.math.columbia.edu/tag/0923}{Tag 0923}]{stacks-project} 
that we may compute derived $x$-adic completions naively, i.e.\ via
$R\lim(\text{--}\otimes_S S/x^n).$  Thus
$$M\otimes^{\mathbf{L}}_R S
\iso R\lim( M\otimes^{\mathbf{L}}_R S/x^n) 
\iso R\lim\bigl( (M\otimes^{\mathbf{L}}_R R/x^n) \otimes^{\mathbf{L}}_{R/x^n} S/x^n\bigr).$$
Artin--Rees allows us to replace the pro-system
 $\{M\otimes^{\mathbf{L}}_R R/x^n\}$ with the pro-system
 $\{M/x^n\}$, % \mar{TG: I still don't understand this! So I guess the
 %   first pro-system can be computed by taking an (infinite) resolution
 %   of~$M$ by finite free $R$-modules, and then reducing mod~$x^n$. I
 %   can see that Artin--Rees will be relevant, but what is the precise
 %   statement that is being made? Are we claiming that there are some
 %   quasi-isomorphisms relating the two pro-systems? Is this just
 %   completely obvious if you think about it the right way? TG: vague
 %   idea for an argument: compute the first pro-system by instead
 %   taking the total complex of the tensor product with $0\to
 %   R\stackrel{x^n}{\to}R\to 0$, and then argue as in the proof of
 %   \cite[\href{https://stacks.math.columbia.edu/tag/0921}{Tag
 %     0921}]{stacks-project}. I haven't read that proof yet so I don't
 %   know if it will literally work in this unbounded context (I also
 %   don't know if the pro-systems are literally isomorphic in the
 %   unbounded case\dots) TG: if I'm doing this right then it does look
 %   like we would be able to show that for each fixed~$n$ and each~$-p<0$, there exists
 % ~$m\ge n$ such that the map $(M\otimes^{\mathbf{L}}_R
 % R/x^m)\to(M\otimes^{\mathbf{L}}_R R/x^n)$ induces the~$0$ map on the
 % homology groups in degrees $-p,\dots,-1$. I don't see that there
 % should be a uniformity though. Is this statement enough (not for the
 % replacement of pro-systems, but for the application)?}
and so we find that in fact
\begin{multline*}
M\otimes^{\mathbf{L}}_R S
\iso R\lim(M/x^n \otimes^{\mathbf{L}}_{R/x^n} S/x^n)
\\
\iso R\lim( M\otimes_R S/x^n) \iso \lim M\otimes_R S/x^n,
\end{multline*}
the penultimate isomorphism holding
since $S/x^n$ is flat over $R/x^n$, 
and the final isomorphism following from the fact that the transition morphisms
in the pro-system $\{ M\otimes_R S/x^n\}$ are surjective.
Thus indeed 
$M\otimes^{\mathbf{L}}_R S$ has cohomology supported in a single degree,
as required.

The claim about faithful flatness follows from the immediately
following Lemma~\ref{lem:faithful flatness}.
%\mar{ME: Add (easy) faithful flatness deduction}
\end{proof}

\begin{lemma}
\label{lem:faithful flatness}
If $R\to S$ is a flat morphism of rings, if $I$ is an ideal in $R$ for which
$R$ is $I$-adically complete, and if the morphism
$R/I \to S/I$ is faithfully flat, then the morphism $R\to S$ is faithfully flat.
\end{lemma}
\begin{proof}
Since $R$ is $I$-adically complete, we see that $I$ is contained in the
Jacobson radical of~$R$.  Thus $\Max\Spec R/I = \Max\Spec R$,
and so our assumption that $R/I \to S/I$ is faithfully flat shows
that $\Max\Spec R$
 is contained in the image of $\Spec S/I \subseteq \Spec S$ in~$\Spec R$.
Since flat morphisms satisfy going down, we find that in fact
$\Spec S \to \Spec R$ is surjective, as claimed.
\end{proof}

%The following lemma is an elaboration on~\cite[Rem.\
%4.31]{2016arXiv160203148B}.

%\begin{lem}\mar{TG: would like to relax some of these hypotheses,
%    maybe to bounded $p$-power torsion?}\mar{TG: would also like
%    faithfully flat, not just flat.}
%\mar{ME: As discussed in email, I think that I can remove the $p$-torsion free
%hypothesis from this result; or more precisely, can reduce the 
%general case to that case.  Can also get faithfully flat case from the flat
%case. (This is pretty formal.)}
%  Let $R\to S$ be a homomorphism of $p$-adically complete, $p$-torsion
%  free $\Zp$-algebras, and suppose that~$R$ is Noetherian. If
%  $R/p^nR\to S/\p^nS$ is flat for all~$n\ge 1$, then $R\to S$ is flat.
%\end{lem}
%\begin{proof}By (for example)
%  \cite[\href{https://stacks.math.columbia.edu/tag/00M5}{Tag
%    00M5}]{stacks-project}, it is enough to show that if~$M$ is a
%  finitely generated $R$-module, then $M\otimes^{\mathbf{L}}_RS$ is
%  concentrated in degree~$0$. Since~$R$ is Noetherian,
%  \cite[\href{https://stacks.math.columbia.edu/tag/066E}{Tag
%    066E}]{stacks-project} is a pseudo-coherent complex
%  of~$R$-modules, so $M\otimes^{\mathbf{L}}_RS$ is a pseudo-coherent complex
%  of~$S$-modules
%  by~\cite[\href{https://stacks.math.columbia.edu/tag/0650}{Tag
%    0650}]{stacks-project}. Since~$R$ and~$S$ are $p$-adically
%  complete, it follows
%  from~\cite[\href{https://stacks.math.columbia.edu/tag/0A05}{Tag
%    0A05}]{stacks-project} that both~$M$ and
%  $M\otimes^{\mathbf{L}}_RS$ are derived $p$-adically complete.
%  
%\mar{TG: see
%  also~\cite[\href{https://stacks.math.columbia.edu/tag/0BKH}{Tag
%    0BKH}]{stacks-project} for material on derived completion for
%  Noetherian rings}Since~$R$ is Noetherian, we have
%$M\cong R\lim(M\otimes^{\mathbf{L}}_RR/p^n)$. Since~$S$ is $p$-torsion
%free\mar{TG: could relax to bounded torsion?} we also have
%$M\otimes^{\mathbf{L}}_RS\cong R\lim((M\otimes^{\mathbf{L}}_RS)\otimes^{\mathbf{L}}_SS/p^n)$.\mar{TG:
%I don't follow what happens in BMS - they say that they make implicit
%use of Artin--Rees to replace $M\otimes^{\mathbf{L}}_RR/p^n$ with $M/p^n$?}
%\end{proof}

      We also note the following lemma.

      \begin{lemma}
	      \label{lem:Witt flatness}
      If $F \hookrightarrow F'$ is an inclusion of
      perfectoid fields in characteristic $p$,
	      then the induced morphism $W_a(\cO_F)
	      \to W_a(\cO_{F'})$ 
	      is faithfully flat, for any~$a~\geq~1$.
      \end{lemma}
      \begin{proof}
Since $\cO_F$ is a B\'ezout ring,
being flat over $\cO_F$ is the same as being
torsion free, and so the inclusion
$\cO_{F} \hookrightarrow \cO_{F'}$
is a flat morphism,
and thus even faithfully flat, being a local morphism of local rings. 
A standard grading argument, applying Lemma~\ref{lem:graded flat} to 
the $p$-adic filtrations on source and target,
then shows that the inclusion
$W_a(\cO_{F}) \hookrightarrow W_a(\cO_{F'})$,
is faithfully flat, for each~$a~\geq~1$.
\end{proof}

\subsection{Back to coefficient rings}
      We now return to our discussion 
of coefficient rings, and record that, at least if ~$A$ is a finite
type % \mar{TG:
%         replace with topologically of finite type? ME: Now all done, and least
% to first approximation.  You can remove this comment and the next couple if you're 
% happy.}
      $\Z/p^a\Z$-algebra, the various natural maps between these rings
      are in fact (faithfully) flat injections.
We also show that the various maps of coefficient rings induced by
a (faithfully) flat morphism of finite type $\Z/p^a$-algebras
are again (faithfully) flat.

% \mar{TG: probably also want this result for the topological finite
%   type case. I think everything is OK if we can prove that the targets
% are Noetherian outside~$(v,p)$. For this, can we perhaps reduce to considering
% the cases of~$v,p$ separately? i.e.\ to showing that each of
% $W(\C^\flat)_A[1/v]$ and $W(\C^\flat)_A[1/p]$ are Noetherian?}
\begin{prop}%\mar{TG: haven't tried to check what's true yet.}
  \label{prop: maps of coefficient rings are faithfully flat injections}
Suppose that~$A\to B$ is a flat homomorphism
  of finite type $\Z/p^a\Z$-algebras for
  some $a\ge 1$. Then all the maps in the following diagram are flat.
  Furthermore the vertical arrows are all injections,
while the  horizontal arrows are all faithfully flat {\em (}and
so in particular also injections{\em )}.
If $A\to B$ is furthermore faithfully flat, then the same is true of the
diagonal arrows.
% \mar{ME: Have to slightly rewrite proof to treat flat vs. faithfully flat
% set-up that I've just introduced. ME: Now done, but should check it
% all makes sense.}
  % \[\xymatrix{\gS_A \ar[r]\ar[dd]&\OEA\ar[d]\\&\widetilde{\cO}_{\cE,A}\ar[d]\\
  %   \AAinf{A}\ar[r] & W(\C^\flat)_A \\ &\tA_{K,A}\ar[u]\\
  %   \A_{K,A}^+\ar[uu]\ar[r]&\A_{K,A}\ar[u]} \]


\[\xymatrix{&\A_{K,B}^+\ar[rr]\ar[dd]&&\AAinf{B}\ar[dd]&&\gS_B
    \ar[dd]\ar[ll]\\
    \A_{K,A}^+\ar[rr]\ar[dd]\ar[ur]&&\AAinf{A}\ar[dd]\ar[ur]&&\gS_A
    \ar[dd]\ar[ll]\ar[ur]\\
&\A_{K,B}\ar[r]&\tA_{K,B}\ar[r]& W(\C^\flat)_B &
\widetilde{\cO}_{\cE,B}\ar[l]& \OEB\ar[l]\\
\A_{K,A}\ar[r]\ar[ur]&\tA_{K,A}\ar[r]\ar[ur]& W(\C^\flat)_A\ar[ur] &
\widetilde{\cO}_{\cE,A}\ar[l]\ar[ur]& \OEA\ar[l]\ar[ur]}\]
  
% \[\xymatrix{\A_{K,A}^+\ar[rr]\ar[d]&&\AAinf{A}\ar[d]&&\gS_A
%     \ar[d]\ar[ll]\\
% \A_{K,A}\ar[r]&\tA_{K,A}\ar[r]& W(\C^\flat)_A &
% \widetilde{\cO}_{\cE,A}\ar[l]& \OEA\ar[l]}\]
\end{prop}%\mar{TG: not sure if we also want to put some perfections}
\begin{proof}From left to right, the vertical maps are given by
  inverting the elements $T$, $v$, and~$u$ respectively (where as
  above, $v$ is any element of the maximal ideal of~$W_a(\cO^\flat_\C)$, whose
  image in~$\cO^\flat$ is nonzero; in particular,
  we can take $v$ equal to either $T$ or $u$).
  Since localizations are flat, 
  to prove the proposition for these maps it
  is enough to note that the sources of the vertical maps are
  respectively $T$-, $v$-, and $u$-torsion free, by
  Lemma~\ref{lem:non-zero divisor}. 

  We now turn to the horizontal maps, where we will make repeated use
  of Proposition~\ref{prop: FGK flat}. It evidently suffices to treat the maps
  with $A$-coefficients. Since
  each of $\A_{K,A}$, $\tA_{K,A}$, 
  $W(\C^\flat))_A$,
  $\widetilde{\cO}_{\cE,A}$, and $\OEA$
  is Noetherian by Proposition~\ref{prop:
    rigid analysis FGK}~(1), it follows from Proposition~\ref{prop: FGK flat} (and
  the flatness of the vertical maps) that we need only show that the maps
  $\gS_A/u^i\to
  (W_a(\cO_{\widehat{K_\infty}}^\flat)\otimes_{\Zp}A)/u^i\to
  \AAinf{A}/u^i$ and
  $\A^+_{K,A}/T^i\to
  (W_a(\cO_{\widehat{\Kcyc}}^\flat)\otimes_{\Zp}A)/T^i\to
  \AAinf{A}/T^i$ are faithfully flat for each~$i~\ge~1$. 

  For the faithful flatness of
  $\gS_A/u^i
  \to(W_a(\cO_{\widehat{K_\infty}}^\flat)\otimes_{\Zp}A)/u^i$,
  %\to\AAinf{A}/u^i$,
  it is
  enough to show that $\gS/p^a\to W_a(\cO_{\widehat{K_\infty}}^\flat)$
  is faithfully flat. By a standard grading argument (see
  Lemma~\ref{lem:graded flat}), it suffices in turn to prove that
  $k[[u]]\to \cO_{\widehat{K_\infty}}^\flat$ is faithfully flat, so in
  turn it is enough to check that~$\cO_{\widehat{K_\infty}}^\flat$ is
  $u$-torsion free, which is clear. To see that
  $(W_a(\cO_{\widehat{K_\infty}}^\flat)\otimes_{\Zp}A)/u^i\to
  \AAinf{A}/u^i$ is faithfully flat, it suffices to note that
  $W_a(\cO_{\widehat{K_\infty}}^\flat)\to W_a(\cO_{\C}^\flat)$ is
  faithfully flat,
  by Lemma~\ref{lem:Witt flatness}.
  %which we prove in Section~\ref{subsec:coefficients} below.
  The faithful flatness of the remaining horizontal maps
  is proved in exactly the same way.

Finally, the (faithful) flatness of the diagonal maps is now immediate
from another application of Proposition~\ref{prop: FGK flat}, together
with the (faithful) flatness of $A\to B$.% \mar{TG: maybe choose one case
  % and elaborate on this, perhaps the $W(\C^\flat)$ case?}
\end{proof}%\mar{TG: proof needed}% \mar{TG: also need $A\to B$

  % faithfully flat implies corresponding maps are faithfully flat.}

 Recall that if~$\Acirc$ is a $p$-adically
    complete $\cO$-algebra, then~$\Acirc$ is said to be \emph{topologically of finite
    type over~$\cO$} if it can be written as a quotient of a
  restricted formal power series ring in finitely many variables~$\cO\langle\langle
  X_1,\dots,X_n\rangle\rangle$; equivalently, if and only if~$\Acirc\otimes_{\cO}k$ is a finite
    type $k$-algebra (\cite[\S0, Prop.\
    8.4.2]{MR3752648}). In particular, if~$\Acirc$ is a
    $\cO/\varpi^a$-algebra for some~$a\ge 1$, then~$\Acirc$ is
    topologically of finite type over~$\cO$ if and only if it is of
    finite type over~$\cO/\varpi^a$.
    Since $\cO$ is finite over $\Z_p$, it is equivalent for $A$ to be
    topologically of finite type over $\cO$ or over $\Z_p$, and so for the
    moment we consider the maximally general case of a topologically finite type
    $\Z_p$-algebra. 
%\mar{ME: Now add result dealing with top.\ f.t.\ context.}

\begin{remark} We don't know if the analogue of 
  Proposition~\ref{prop: maps of coefficient rings are faithfully flat injections}
holds when $A$ and $B$ are taken to be merely $p$-adically complete 
and topologically of finite type over~$\Z_p$,
rather than of finite type over $\Z/p^a$ for some~$a \geq 1$.   Each of the various
coefficient
rings over $A$ and $B$ is (by definition) formed by first forming the
corresponding coefficient ring over each $A/p^a$ or~$B/p^a$, and then taking
an inverse limit.  Since the formation of inverse limits is left exact,
we see that the horizontal and vertical arrows in the diagram are injective,
but we don't know in general that the various arrows are flat (although
we have no reason to doubt it).  One can however establish some partial results,
using~Proposition~\ref{prop:BMS flatness}.  
We record here one such result,
which we will need later on.
\end{remark}

\begin{prop}
\label{prop:top f.t. flatness}
If $A$ is a $p$-adically complete $\Z_p$-algebra which is topologically of
finite type, 
then the natural morphism $\gS_A \to \AAinf{A}$ is faithfully flat.
\end{prop}
\begin{proof}
This follows from Propositions~\ref{prop:BMS flatness}
and~\ref{prop: maps of coefficient rings are faithfully flat injections},
once we note that $\gS_A$ is Noetherian when $A$ is $p$-adically
complete and topologically of finite type; indeed since $A$ is in
particular Noetherian, so is the power series ring $W(k)\otimes_{\Z_p} A [[u]].$
\end{proof}



We conclude this initial discussion of coefficient rings by 
explaining how
the action of~$\varphi$ on the various rings~$\gS$, $\Ainf$ and so on
extends to a continuous action on the corresponding rings~$\gS_A$,
$\AAinf{A},$ etc.,
and similarly for the various Galois actions.
For this,
it is convenient for us to briefly digress,
and to introduce the
following situation, which will also be useful for us in Chapter~\ref{sec:
phi modules and phi gamma modules}. %  \subsection{The action of Frobenius}

\begin{situation}
  \label{subsubsec:general framework}
  Fix a finite extension~$k/\Fp$ and
  write~$\Aplus:=W(k)[[T]]$. Write~$\A$ for the $p$-adic completion
  of~$\Aplus[1/T]$.

  % \begin{rem}In~\cite{EGstacktheoreticimages}, $\Aplus$ was
  %   denoted~$\gS$; however in this paper we will have to consider
  %   several different choices of~$\Aplus$, and we prefer to reserve
  %   the notation~$\gS$ for the ring occurring in the definition of
  %   Breuil--Kisin modules, as in Section~\ref{subsubsec: Kummer}.
  % \end{rem}
  % \mar{TG: update this remark once coefficient rings section is in
  % place, to make it clear which of those coefficient rings this
  % theory applies to.}



  If $A$ is a $p$-adically complete~$\Zp$-algebra, we write
  $\Aplus_A:=(W(k)\otimes_{\Zp}A)[[T]]$; we equip~$\Aplus_A$ with its
  $(p,T)$-adic topology, so that it is a topological $A$-algebra
  (where~$A$ has the $p$-adic topology). Let $\AAA_A$ be the $p$-adic
  completion of~ $\Aplus_A[1/T]$, which we regard as a topological
  $A$-algebra by declaring $\Aplus_A$ to be an open subalgebra. Note
  that the formation of~$\gS_A$, $\OEA$, $\A_{K,A}^+$ and~$\A_{K,A}$
  above are particular instances of this construction.


  % \mar{TG: just take $q=p$ here and below.}
  Let~$\varphi$ be a ring endomorphism of~$\A$ which is congruent to
  the ($p$-power) Frobenius endomorphism modulo~$p$. We say
  that~$\A^+$ is $\varphi$-stable if~$\varphi(\A^+)\subseteq \A^+$.

By~\cite[Lem.\ 5.2.2 and
  5.2.5]{EGstacktheoreticimages}, if~$\A^+$ is $\varphi$-stable, then
  $\varphi$ is faithfully flat, and induces the usual 
  Frobenius on~$W(k)$; the same arguments show that this is true
  for~$\varphi$ on~$\A$, even
  if $\A^+$ is not $\varphi$-stable.
\end{situation}

 We have the following variant
  of~\cite[Lem.\ 5.2.3]{EGstacktheoreticimages}.

  \begin{lem}
    \label{lem: phi on powers of T without assuming A plus
      stable}Suppose that we are in Situation~{\em \ref{subsubsec:general
      framework}}. Then for
    each integer $a\ge 1$ there is an integer~$C\ge 0$ such that for
    all~$n\ge 1$, we have $\varphi(T^n)\in T^{pn-C}\A^++p^a\A$. In
    particular, the action of~$\varphi$ on~$\A$ is continuous.
  \end{lem}
  \begin{proof}
    For some~$h\ge 0$ we have $\varphi(T)\in T^{-h}\A^++p^a\A$. Write
    $\varphi(T)=T^p+pY$, so that~$\varphi(T^n)=(T^p+pY)^n$. If we
    expand using the binomial theorem, then every term on the right
    hand side is either divisible by~$p^a$, or is a multiple of
    $T^{p(n-r)}Y^r$ for some $0\le r \le a-1$. It follows that we can
    take~$C=(a-1)(p+h)$. 
% \mar{ME: Hopefully this is right; I just think having optimal constants
% in a simple situation like this makes it easier for the reader
% to follow what's going on, since it matches with what they'll compute
% for themselves if they double-check. TG: agreed, commenting out.}
  \end{proof}

  \begin{lem}\label{lem: varphi extends to A}
	  If~$A$ is a $p$-adically complete $\Zp$-algebra,
    then in Situation {\em \ref{subsubsec:general framework}}, the
    endomorphism~$\varphi$ of~$\A$ extends uniquely to an $A$-linear continuous
  endomorphism of~$\AAA_A$ \emph{(}which we continue to denote
  by~$\varphi$\emph{)}. If~$\A^+$ is $\varphi$-stable, then $\A^+_A$
  is $\varphi$-stable for all~$A$.  
  \end{lem}
  \begin{proof}Since~$\varphi$ is $\Z_p$-linear
    by definition, and since the topologies on~$\A^+_A$ and~$\AAA_A$
 are defined by passage to the limit modulo~$p^a$
    as $a\to\infty$, we are immediately reduced to the case that~$A$
    is a $\Z/p^a\Z$-algebra. 
In this case the result is immediate from Lemma~\ref{lem: phi on powers of T without assuming A plus
      stable} and Lemma~\ref{lem:extending endomorphisms}. % bearing in
   % mind Lemma~\ref{lem: actual examples of Cplus and C that we use}.
%\mar{ME: This last citation doesn't seem relevant to this abstract set-up.}
  \end{proof}
  % \begin{proof}
  %   Now let~$C$ be as in Lemma~\ref{lem: phi on powers of T without
  %     assuming A plus stable}. Then for each~$n\ge C$, it follows from
  %   Lemma~\ref{lem: phi on powers of T without assuming A plus stable}
  %   that $\varphi(T^n\A^+)\subseteq T^n\A^+$, so that in particular for
  %   each~$n\ge 1$, $\varphi$ induces an endomorphism
  %   of~$T^C\A^+/T^{C+n}\A^+$, and thus an $A$-linear endomorphism
  %   of \[A\otimes_{\Zp}(T^C\A^+/T^{C+n}\A^+)=T^C(\A^+\otimes_{\Zp}A)/T^{C+n}(\A^+\otimes_{\Zp}A) \]
  %   (where the equality holds by the $\Z/p^a$-flatness of
  %   $\A^+/p^a$). Passing
  %   to the projective limit over~$n$, we obtain an $A$-linear
  %   endomorphism of~$T^C\A^+_A$, which is continuous for the $T$-adic
  %   topology;  inverting~$T$ gives the required endomorphism
  %   of~$\A_A$. If~$\A^+$ is $\varphi$-stable, then we can take~$C=0$,
  %   so in this case we obtain an endomorphism of~$\A^+_A$ itself. This
  %   completes the proof of the lemma.
  %   \end{proof}
  % It follows from Lemma~\ref{lem: phi on powers of T without assuming
  %   A plus stable} that $\varphi$ extends uniquely to a continuous
  % endomorphism of~$\AAA_A$ which we continue to denote
  % by~$\varphi$.\mar{TG: make this a lemma?}

%   \begin{lem}\label{lem: varphi extends to A}The
%     endomorphism~$\varphi$ of~$\A$ extends uniquely to an $A$-linear continuous
%   endomorphism of~$\AAA_A$ \emph{(}which we continue to denote
%   by~$\varphi$\emph{)}. If~$\A^+$ is $\varphi$-stable, then $\A^+_A$ is $\varphi$-stable for all~$A$.       % \mar{TG: prove this}
%   \end{lem}\mar{TG: trying below to make this into something more general.}
%   \begin{proof}Suppose firstly that~$A$ is a $\Z/p^a$-algebra, and
%     let~$C$ be as in Lemma~\ref{lem: phi on powers of T without
%       assuming A plus stable}. Then for each~$n\ge C$, it follows from
%     Lemma~\ref{lem: phi on powers of T without assuming A plus stable}
%     that $\varphi(T^n\A^+)\subseteq T^n\A^+$, so that in particular for
%     each~$n\ge 1$, $\varphi$ induces an endomorphism
%     of~$T^C\A^+/T^{C+n}\A^+$, and thus an $A$-linear endomorphism
%     of \[A\otimes_{\Zp}(T^C\A^+/T^{C+n}\A^+)=T^C(A\otimes_{\Zp}\A^+)/T^{C+n}(A\otimes_{\Zp}\A^+). \]Passing
%     to the projective limit over~$n$, we obtain an $A$-linear
%     endomorphism of~$T^C\A^+_A$, which is continuous for the $T$-adic
%     topology;  inverting~$T$ gives the required endomorphism
%     of~$\A_A$. If~$\A^+$ is $\varphi$-stable, then we can take~$C=0$,
%     so in this case we obtain an endomorphism of~$\A^+_A$ itself.

% Finally, the general case that~$A$ is a $p$-adically complete
% $\Zp$-algebra follows by noting that~$\A_A=\varprojlim_a\A_{A/p^aA}$ and~$\A^+_A=\varprojlim_a\A_{A/p^aA}^+$.
%   \end{proof}

  \begin{lem}\label{lem: varphi extends to A and Galois
      etc}Let~$A$ be a $p$-adically complete $\Zp$-algebra.
    \begin{enumerate}
\item The endomorphism
	\label{item: continuity of phi 2} $\varphi$
  of~$\gS$ {\em (}resp.\ $\cO_{\cE}$, $\A_{K}$,
  $W(\cO_F^\flat)$, $W(F^\flat)${\em )}
	extends uniquely to an $A$-linear continuous
  endomorphism
  of~$\gS_A$ {\em (}resp.\ $\OEA$, $\A_{K,A}$,
  $W(\cO_F^\flat)$, $W(F^\flat)_A${\em )}, which we again denote by~$\varphi$.
  If~$K$ is furthermore basic,
  then the endomorphism $\varphi$ of $\A_{K,A}$ preserves~$\A^+_{K,A}$.
\item\label{item: continuity of Galois action 3}
%\mar{ME: Should $\widetilde{\cO}_{\mathcal{E}}$ be $\widetilde{\cO}_{\cE}^{\ur}$?}
	The continuous $G_K$-action on 
  $W(\cO_{\C}^\flat)$ and $W(\C^\flat)$ {\em  (}%resp.\ the
 % continuous~$G_{K_\infty}$-action 
 %  on~$\widetilde{\cO}_{\mathcal{E}}$,
  resp.\ the continuous~$\Gamma_K$-action on~$\A_{K}$ and~$\tA_{K}${\em )}
  extends to a continuous $A$-linear action
	of~$G_K$ on
  $W(\cO_{\C}^\flat)_A$ and $W(\C^\flat)_A$ {\em (}% resp.\ of~$G_{K_\infty}$
  % on~$\tOEA$,
  resp.\ of~$\Gamma_K$ on~$\A_{K,A}$ and~$\tA_{K,A}${\em )}.
    \end{enumerate}
     % \mar{TG: prove this}
  \end{lem}%\mar{TG: written - leaving for ME to check.}
  \begin{proof}As in the proof of Lemma~\ref{lem: varphi extends to
      A}, we can immediately reduce to the case that~$A$ is a
    $\Z/p^a$-algebra. Part~\eqref{item: continuity of phi 2} then
    follows from Lemma~\ref{lem:extending endomorphisms} as in the
    proof of Lemma~\ref{lem: varphi extends to
      A},
%\mar{ME: I think that this is where the citation to
%    Lemma~\ref{lem: actual examples of Cplus and C that we use} belongs.}
while part~\eqref{item: continuity of Galois action 3}
%\mar{ME: Also cite boundedness result that will be put in above.}
    follows from Lemma~\ref{lem:extending group actions},     bearing in mind
    Lemmas~\ref{lem: actual examples of Cplus and C that we use}
and~\ref{lem: boundedness of Galois actions in our concrete
  settings}.\end{proof}

%We end this section by computing
For our final result of this section,
we compute
the $\varphi$-invariants in $W(\C^\flat)_A$.


\begin{lem}
\label{lem:phi invariants}
If $A$ is any $p$-adically complete $\Z_p$-algebra,
then for each $a \geq 1$,
the natural morphism
$A/p^a \to W_a(\C^\flat)_A$
induces an identification
$A/p^a = \bigl(W_a(\C^\flat)_A\bigr)^{\varphi =  1} $,
and {\em (}consequently{\em )}
the natural morphism
$A \to W(\C^\flat)_A$
induces an identification
$A = \bigl(W(\C^\flat)_A\bigr)^{\varphi = 1}.$ 
\end{lem}
\begin{proof}
Since 
$W(\C^\flat)_A = \varprojlim_a W_a(\C^\flat)_A$,
we see that
$$\bigl(W(\C^\flat)_A\bigr)^{\varphi = 1} =
\varprojlim \bigl(W_a(\C^\flat)_A\bigr)^{\varphi = 1},$$
and thus (as the phrasing of the lemma indicates),
the claim for $W(\C^\flat)_A$ follows from the 
claims for each $W_a(\C^\flat)_A.$
Thus for the remainder of the proof,
we may and do assume that $A$ is a $\Z_p/p^a$-module for some $a \geq 1$,
and prove that
$A = \bigl(W_a(\C^\flat)_A\bigr)^{\varphi = 1}.$
In fact,
for any $\Z/p^a$-module $M$, 
we may define
$$W_a(\C^\flat)_M := 
\Bigl( \varprojlim_n  \bigl( W_a(\cO_\C^\flat)/T^n\bigr)\otimes_{\Z/p^a} M \Bigr)
[1/T]$$
(this extends our preceding definition for $\Z/p^a$-algebras),
and
we will prove that
$M = \bigl(W_a(\C^\flat)_M\bigr)^{\varphi = 1}.$
(In fact this statement for modules, which obviously
implies the statement for algebras, is actually equivalent to it,
as one sees by applying the statement for algebras to 
$\Z_p/p^a$-algebras of the form $\Z_p /p^a \oplus M,$
where $M$ is an arbitrary $\Z_p/p^a$-module equipped with the
square zero multiplication.)

Suppose first that $a = 1$, so that $M$  is simply an $\F_p$-vector space..  
If we choose a basis for~$M$, i.e.\ an isomorphism $\F_p^{\oplus I} \iso M$,
then $(\cO_\cC^{\flat})_M = \prod'_{i \in I} \cO_\cC^\flat,$
where the prime indicates we consider the set of tuples $(x_i)_{i\in I}$
for which $\lim_i x_i = 0,$ where the limit is taken with respect to
the filter of cofinite subsets of~$I$.
Thus we have, correspondingly, that
$\cC^{\flat}_M  = \prod'_{i \in I} \cC^\flat$
(where the prime has the same meaning).
Now $(\cC^\flat)^{\varphi = 1} = \F_p$, and so
$$(\cC^{\flat}_M)^{\varphi= 1}  = \prod_{i \in I}{}' (\cC^\flat)^{\varphi = 1}
= \oplus_{i \in I} \F_p  = M.$$
This establishes the case~$a = 1$.

Now consider the case of general~$a$.  An application of
Lemma~\ref{lem:completion properties}~(1) % \mar{TG: can you say what
  % exactly you're applying Lemma~\ref{lem:completion properties} to
  % here? TG: I meant the rings! I've added what I think is correct.}
(taking $R=\Z/p^a$ and $C^+=W_a(\cO_{\C^\flat})$, and considering  the
exact sequence  of $\Z/p^a$-modules $0 \to pM\to M \to M/pM \to 0$;
see also Lemma~\ref{lem: actual examples of Cplus and C that we use})
shows that the sequence
$$0 \to  W_{a-1}(\cC^{\flat})_{p M} = W_a(\cC^\flat)_{p M} \to 
W_a(\cC^\flat)_M \to W_a(\cC^\flat)_{M/pM} = \cC^{\flat}_{M/pM} \to 0$$
is short exact.
Suppose that $x \in \bigl(W_a(\C^\flat)_M\bigr)^{\varphi = 1}.$
The case $a = 1$ already proved then shows that 
we may find $m \in M$ so that $y := x - m$  lies in
$W_{a-1}(\cC^{\flat})_{p M}.$    Of course $\varphi(y) = y$,
and so, by induction on $a$, we find that $y \in pM$.  Thus $x = m + y \in M$,
as claimed.
\end{proof}



We will also use the following variant on Lemma~\ref{lem:phi invariants}. Suppose  that ~$A$ is a complete local Noetherian $\cO$-algebra
  with finite residue field, and write $\widehat{W(\C^\flat)}_A$ for
the $\m_A$-adic completion of~$W(\C^\flat)_A$.
\begin{lem}% \mar{TG: needs to have $\m$-adic completion, and presumably
  % above too}
  \label{lem: phi invariants in W C flat with coefficients}Let~$A$ be
  a complete local Noetherian $\cO$-algebra with finite residue
  field. Then $(\widehat{W(\C^\flat)}_A)^{\varphi=1}=A$. % \mar{TG:
    % should at least say something about how this relates to the
    % (currently) preceding lemma.}
%  \mar{TG: replaced this with a proof using the new lemma ME: OK, looks good}
\end{lem}
\begin{proof}Since 
$\widehat{W(\C^\flat)}_A = \varprojlim_n W(\C^\flat)_{A/\m_A^n}$,
we have
\[\bigl(\widehat{W(\C^\flat)}_A\bigr)^{\varphi=1} = \varprojlim_n \bigl(W(\C^\flat)_{A/\m_A^n}\bigr)^{\varphi=1},\]
so the result follows from Lemma~\ref{lem:phi invariants} (applied to each~$A/\m_A^n$).
\end{proof}
% \begin{proof}We begin by noting that there is an exact
%   sequence % \mar{TG: need to prove this but I think it's true, e.g.\ on
%     % $\Cflat$ itself I think it's clear with surjectivity coming from
%     % $\Cflat$ being separably closed. I am surprised that I can't find
%     % it in the literature, though.}
%   \[0\to\Z_p\to W(\C^\flat)\stackrel{1-\varphi}{\to} W(\C^\flat)\to
%   0.\](To see that this sequence is exact, one easily reduces to
% showing that the sequence \[0\to\F_p\to\C^\flat\stackrel{1-\varphi}{\to} \C^\flat\to
%   0  \]is exact, which follows from the fact that~$\C^\flat$ is
% separably closed.) Since
% $W(\C^\flat)$ is $\Zp$-flat, if~$B$ is any $\Z_p$-algebra we may tensor with~$B$ and obtain a short
% exact sequence 
% \[0\to B\to W(\C^\flat)\otimes_{\Zp}B\stackrel{1-\varphi}{\to}
%   W(\C^\flat)\otimes_{\Zp}B\to 0.\]Applying this with~$B=A/\m_A^n$ and
% passing to the inverse limit over~$n$, we have a short exact sequence
% \[0\to A\to \widehat{W(\C^\flat)}_A\stackrel{1-\varphi}{\to}
%   \widehat{W(\C^\flat)}_A\to 0,\]as required.% \mar{TG: can probably prove this
%   % for any~$A$? In a similar way to the arguments in section \ref{subsec: perfect
%   % etale phi modules}.}
% % Passing to completions gives the
% % result.\mar{TG:  didn't yet check this claim about passing to
% %   completions. Not sure exactly what was used about~$A$ in this argument\dots}
% \end{proof}



%   Lemma~\ref{lem: varphi extends to A}
%   implies~\eqref{item: continuity of phi 2} for ~$\gS_A$, $\OEA$,
% $\A_{K,A}$, and~$\A^+_{K,A}$. The statement for $W(\cO_F^\flat)$,
% $W(F^\flat)_A$ follows by an identical argument, replacing~$\A^+$
% by~$W(\cO_F^\flat)$ and~$\A$ by~$W(F^\flat)$.

% Finally we turn to~\eqref{item: continuity of Galois action 3}. We show
%  that we have an $A$-linear continuous action of~$G_K$
% on $W(\cO_{\C}^\flat)_A$ and $W(\C^\flat)_A$; the proofs in the other
% cases are essentially identical, or alternatively can be deduced from
% the continuous Galois-equivariant embeddings of the corresponding
% coefficient rings into~$W(\C^\flat)_A$. We have a continuous action
% of~$G_K$ on~$W_a(\cO_{\C}^\flat)$ and on~$W_a(\C^\flat)$, so for each~$n\ge 0$,
% there is a finite index subgroup~$H$ of~$G_K$ and some~$m\ge 0$ such
% that for all~$h\in H$ we have $h(T^mW_a(\cO_{\C}^\flat))\subseteq
% T^nW_a(\cO_{\C}^\flat)$. It suffices to show that this implies that $h(T^mW_a(\cO_{\C}^\flat)_A)\subseteq
% T^nW_a(\cO_{\C}^\flat)_A$. 

% This follows by a very similar argument to the proof
% of Lemma~\ref{lem: varphi extends to A}:
% namely, since $W_a(\cO_{\C}^\flat)$ is flat
% over~$\Z/p^a$, if we tensor the short exact sequence \[0\to
%   T^mW_a(\cO_{\C}^\flat)\to W_a(\cO_{\C}^\flat)\to
%   W_a(\cO_{\C}^\flat)/T^mW_a(\cO_{\C}^\flat)\to 0\]with~$A$, we find
% that \numequation\label{eqn: powers of T and base
%   change}T^m(W_a(\cO_{\C}^{\flat})\otimes_{\Zp}A)=(T^mW_a(\cO_{\C}^{\flat}))\otimes_{\Zp}A.\end{equation}
% By assumption, the composite
% map \[T^mW_a(\cO_{\C}^\flat)\stackrel{h}{\to}W_a(\cO_{\C}^\flat)\to
%   W_a(\cO_{\C}^\flat)/T^nW_a(\cO_{\C}^\flat)\]vanishes. Tensoring this
% with~$A$ and using~\eqref{eqn: powers of T and base
%   change} we find that the composite \[T^m(W_a(\cO_{\C}^\flat)\otimes_{\Zp}A)\stackrel{h}{\to}W_a(\cO_{\C}^\flat)\otimes_{\Zp}A\to
%   (W_a(\cO_{\C}^\flat)\otimes_{\Zp}A)/T^n(W_a(\cO_{\C}^\flat)\otimes_{\Zp}A)\]also vanishes, as required.
%  \end{proof}
% Finally, the general case that~$A$ is a $p$-adically complete
% $\Zp$-algebra follows by noting that~$\A_A=\varprojlim_a\A_{A/p^aA}$ and~$\A^+_A=\varprojlim_a\A_{A/p^aA}^+$.



      
%\mar{ME: Rewrite these somewhat}
%\[\A_{K,A}:=\A_K\cotimes_{\Zp}A, \]
%\[\A^+_{K,A}:=\A^+_K\cotimes_{\Zp}A, \]
%\[\AAinf{A}:=\Ainf\cotimes_{\Zp}A, \]
%\[W(\C^\flat)_A:=W(\C^\flat)\cotimes_{\Zp}A, \]
%\[\gS_A:=\gS\cotimes_{\Zp}A, \]
%\[\OEA:=\cO_\cE\cotimes_{\Zp}A, \]


%\mar{ME: State theorem encoding various flatness and embedding properties}

% \section{Old material}
% In this section we recall some basic facts about the
% classical theory of~$(\varphi,\Gamma)$-modules, which mostly go back
% to the original paper~\cite{MR1106901}. Later in the paper we will
% need to relate these definitions to the theory of perfectoid fields,
% and in order to do so it is convenient for us to follow the approach
% taken in~\cite{MR3412585}, which uses the perfectoid theory to
% streamline the constructions of the various coefficient rings.  % ; unless otherwise
% % noted, everything that we use here is taken from~\cite{MR1106901}.
% % and for further details the reader may wish to
% % consult for example~\cite{colmeznotes,Brinon09cmisummer,MR3702016}.

% Throughout
% this section we write $K$ for a finite extension of $\Q_p$, with ring of
% integers $\cO_K$ and residue field $k$. We
% write~$K_0=\Frac(W(k))\subset K$. We fix an algebraic
% closure~$\overline{K}$ of~$K$ with ring of
% integers~$\cO_{\overline{K}}$ and residue field~$\overline{k}$, and
% regard all algebraic extensions of~$K$ as subfields
% of~$\overline{K}$. Write~$\C$ for the completion of~$\overline{K}$,
% with ring of integers~$\cO_\C$. Write
% \[\cO_\C^\flat:=\varprojlim_{x\mapsto x^p}\cO_{\C}/p=\varprojlim_{x\mapsto x^p}\cO_{\overline{K}}/p, \] a valuation
% ring with valuation~$v$, which can be defined as follows:
% write~$(x_n)$ for a sequence of elements of~$\cO_{\overline{K}}/p$
% with~$x_{n+1}^p=x_n$, and set~$v((x_n)):=p^nv_p(x_n)$, which is
% independent of~$n$ as soon as~$x_n\ne 0$. This valuation extends to
% the field of fractions~$\C^\flat:=\Frac(\cO_\C^\flat)$. We
% set~$\Ainf:=W(\cO_\C^\flat)$; % , and
% % write~$W(\C^\flat):=W(\Frac R)$.\mar{TG: not so sure about this
% %   notation, or even whether this is correct}
% both~$\Ainf$ and~$W(\C^\flat)$ have
% natural commuting actions of~$G_K$ and~$\varphi$.

%  Recall that~$\cO_\C^\flat$ and~$\C^\flat$ have a valuation~$v$, which can be
%     defined as follows: write~$(x_n)$ for a sequence of elements
%     of~$\cO_{\overline{K}}/p$ with~$x_{n+1}^p=x_n$, and
%     set~$v((x_n)):=p^nv_p(x_n)$, which is independent of~$n$ as soon
%     as~$x_n\ne 0$. Then for each $n\ge 1$, we give $\Ainf/p^n\Ainf$
%     (resp.\ $W(\C^\flat)/p^nW(\C^\flat)$) the
% topology inherited from the valuation topology, and we give~$\Ainf$
% (resp.\ $W(\C^\flat))$
% the inverse limit topology. We call these topologies the \emph{weak
%   topology} on~$\Ainf$ and~$W(\C^\flat)$. With respect to this topology the action
% of~$\varphi$ is continuous, as is the action of~$G_K$, in the sense
% that the map \[G_K\times W(\C^\flat)\to W(\C^\flat)\]is continuous.

% \mar{ME: There are many variations: the cyclotomic extension,
% 	the $p$-part of the cyclotomic extension, more general Galois 
% 	extensions (e.g.\ Lubin--Tate extensions, or extensions that
% 	incorporate descent data).  My sense is that
% 	we should probably work in some generality here, so that we
% 	get things set up flexibly.  But for now I will be in the cyclotomic
% 	case, but will try to write things so that they are easily 
% 	generalized. Also: should look into the Herr complex, to see
% which set-up(s) it is valid for; right now I'm a bit puzzled as to how
% it could work integrally when $K = \Q_2$ for the cyclotomic $\Gamma$,
% since then $\Gamma$ is pro-$2$, but not pro-cyclic.}
% We have various data attached to this situation, which we now
% introduce.
% \mar{TG: we need to get things correct here! I think
  % that~$k$ is perhaps not actually the residue field of~$K$ - rather,
  % it should be the residue field of~$K_\infty$, which I think can be
  % bigger. Also I don't know if the formulae for the actions of
  % ~$\varphi$, ~$\gamma$ are correct - we're now taking the
  % $\Delta$-trace, for starters, but I don't even know if we'd be OK
  % without that in the ramified case - it seems to me that people don't
  % actually write explicit formulae, at least when doing overconvergent
  % things. I think the point is that the uniformiser~$\pi_K$ comes
  % about from the field of norms in a rather non specific
  % way. Presumably this is all in \cite{MR1106901}, although I couldn't
  % immediately find anything very helpful. \cite[\S2.2]{MR3230818} may
  % or may not be helpful!}


%obtained by adjoining all $p$-power roots of unity.



% If~$K$ is basic, we make the following choice of~$\A_K'$, $\A_K$, $T'$
% and $T$. We have~$E_K'=E'_{K_0}$\mar{TG: am I making the
%   usual mistake here about residue fields of $K_\infty$ and $K$ not
%   agreeing? I don't think there should be an issue anyway}, and we take $\A_K'=\A'_{K_0}$\mar{TG: define this via
%   usual Teich lift} with the action of~$\Gammat_K$ being the
% restriction of the action of~$\Gamma_{K_0}$. \mar{TG: and then define
%   the version without primes via traces} 



% Note that~$\tGamma$ is a pro-finite group
% which may be written as $$\tGamma \cong \Delta \times \Gamma,$$
% where $\Delta$ is finite % of order dividing~$(p-1)$ (and in particular
% % prime to~$p$)
% and $\Gamma \cong \Z_p$. 
% \mar{TG: my notes say ``what
%   about $p=2$?}

%this is a finitely generated dense subgroup of $\Gamma$.

% It will frequently be convenient for us to consider the ring of
% invariants~$\AKDelta:=(\A_K)^{\Delta}$.\mar{TG: this is a macro, maybe it should have a
%   better name as I suspect it will eventually be widely used below?}
% This inherits from~$\A_K$ a weak topology, and the actions
% of~$\varphi$ and~$\Gamma$ on~$\A_K$ are continuous with respect to
% this topology. Again, there is a non-canonical 
% isomorphism of topological rings\mar{TG: again, I think the residue field here is correct?}
% $W(k)((T))^{\wedge}\isoto\AKDelta$. We write
% $\AKDeltaplus:=(\A_K^+)^{\Delta}$, which is $\varphi$- and
% $\Gamma$-stable, and we have an induced isomorphism of topological
% rings $W(k)[[T]]\isoto\AKDelta$.

%\section{Coefficients}\label{subsec: coefficients}
%\mar{TG: we now introduce some of these rings multiple times in the paper.}
%Let~$A$ be a $p$-adically complete $\Zp$-algebra. Then we
%write %\mar{TG: swap $A$}
%\[\A_{K,A}:=\A_K\cotimes_{\Zp}A, \]
%\[\A^+_{K,A}:=\A^+_K\cotimes_{\Zp}A, \]
%\[\AAinf{A}:=\Ainf\cotimes_{\Zp}A, \]
%\[W(\C^\flat)_A:=W(\C^\flat)\cotimes_{\Zp}A, \]
%\[\gS_A:=\gS\cotimes_{\Zp}A, \]
%\[\OEA:=\cO_\cE\cotimes_{\Zp}A, \]
%\[\cO_{\widehat{\cE^{\nr}},A}:=\cO_{\widehat{\cE^{\nr}}}\cotimes_{\Zp}A, \]
%\[\cO_{\widehat{\cE^{\nr}_\infty},A}:=\cO_{\widehat{\cE^{\nr}_\infty}}\cotimes_{\Zp}A, \]
%where in each case the completed tensor product is with respect to
%the $p$-adic topology on~$A$, and the adic topologies % \mar{TG: is that
%  % the right thing to say?}
%on the
%rings~$\A_K$, $\A_K^+$, $\Ainf$, $W(\C^\flat)$, $\gS$,
%$\cO_{\cE}$, $\cO_{\widehat{\cE^{\nr}}}$, $\cO_{\widehat{\cE^{\nr}_\infty}}$. The actions of~$\varphi$, $\Gamma_K$, $G_K$ (as appropriate) extend to continuous
%actions on these rings.



% \section{Almost Galois descent for finite group actions}
% \label{subsec:finite}
\section[Almost Galois descent for profinite group actions]{Almost Galois descent for profinite group actions
  \sectionmark{Almost Galois descent}}\sectionmark{Almost Galois descent}
\label{subsec:profinite}
We will be interested in descent results for profinite group actions,
and in this section we establish the key result that we will need.
Our setup is slightly elaborate, but accords with the situations that
will arise in practice. We begin with the case of finite group actions.

Suppose that $R$ is a ring, and that $I$ is an ideal in $R$ such
that $I^2 = I$. Assume further that~$I\otimes_R I$ is a flat $R$-module;
by~\cite[Prop.\ 2.1.7]{MR2004652}, this holds if~$I$ is a filtered
union of principal ideals. In particular, these assumptions hold
if~$R$ is a valuation ring with a non-discrete rank one valuation,
and~$I$ is the maximal ideal of~$R$.


The full subcategory of the category of
$R$-modules whose objects are the modules annihilated by $I$ is a Serre
subcategory, and so we can form the quotient category of {\em almost}
$R$-modules. \index{almost module}

Suppose that $S$ is an $R$-algebra
equipped with an action of a finite
group $G$ by $R$-algebra automorphisms
(i.e.\ the structural morphism
$R \to S$ is equivariant with
respect to the given $G$-action on $S$, and the trivial $G$-action
on $R$).

\begin{df}\index{almost Galois extension}
	We say that the morphism $R \to S$ makes $S$ an
	{\em almost Galois extension of $R$, with Galois group $G$},
	if the natural $G$-equivariant and $S$-linear morphism
$$S\otimes_R S \to \prod_{g \in G} S $$
(here $G$ acts on the source through its action on the second factor,
and on the target by permuting the factors, while $S$ acts on the source
through its action on the first factor and on the target through its
action on each factor) 
% \mar{ME: Surely the action on the target is via the regular action of $G$ on
% 	the index set $G$ for the product! TG: isn't that what
%         ``permuting the factors'' means? I'm happy to write that, but
%         I just want to check that you're proposing a more precise
%         statement of the same thing, rather than something totally different? ME: Actually, I just got confused b/w the description of the $G$-action and 
% of the $S$-action, so probably nothing needs to be changed.  Sorry!}
defined by $s_1\otimes s_2 \mapsto  \bigl(s_1 g(s_2) \bigr)_{g \in G}$
is an {\em almost isomorphism} (i.e.\ induces an isomorphism in the category
of almost $R$-modules).
\end{df}

\begin{remark}
	\label{rem:base-change}
	Suppose that $R\to R'$ is a morphism, and define $I' = I R'$,
	so that $I'^2 = I'$, allowing us to also define the category
	of almost $R'$-modules.  If $R \to S$ is almost Galois
	with Galois group~$G$, then evidently
	$R' \to S' := R'\otimes_R S$ is almost Galois with Galois
        group~$G$.

        Conversely, if $R\to R'$ is faithfully flat, and $R'\to S'$ is
        almost Galois with Galois group~$G$, then so is $R\to S$;
        indeed we may write $R'\otimes_R (S\otimes_RS)=(R'\otimes_R
        S)\otimes_{R'}(R'\otimes_R S)=S'\otimes_{R'}S'$.
\end{remark}

% For the remainder of this subsection, we suppose that $R \to S$  makes
% $S$ an almost Galois extension of $S$ with Galois group $G$,
% and furthermore that $R \to S$ is faithfully flat.


\begin{lemma}
	\label{lem:almost invariants}Suppose that $R \to S$  makes
$S$ an almost Galois extension of $S$ with Galois group $G$,
and furthermore that $R \to S$ is faithfully flat.
	Then the morphism $R \to S^G$ is an almost isomorphism.
\end{lemma}
\begin{proof}
	Since $R \to S$ is faithfully flat, it suffices to verify that
	the induced morphism 
	\numequation
        \label{eqn:hoped for a.i.}
        S \to (S\otimes_R S)^G
        \end{equation}
        is an almost isomorphism.
        By assumption the morphism $S\otimes_R S \to \prod_{g \in G} S$
        is an almost isomorphism.  One easily verifies that the induced 
        morphism on invariant subrings
        $$(S\otimes_RS)^G \to (\prod_{g \in G} S)^G = S$$
        is then also an almost isomorphism.
        The composite of this map with the morphism~\eqref{eqn:hoped for a.i.}
        is just the identity, and so we find that~\eqref{eqn:hoped for a.i.}
        is indeed an almost isomorphism.
	%Remark~\ref{rem:base-change} shows that $S \to S\otimes_R S$
	%is almost Galois with Galois group $G$,
	%so that this follows from
	%the fact that $S \to (\prod_{g\in G}  S)^G$ is evidently 
% \mar{ME: Doesn't this actually use the almost isomorphism
% $S \otimes_R S \to \prod_{g \in G} S$  coming from almost \'etaleness of $S$ over $R$ itself, rather than the base-changed statement? ME: I have (I hope) now
% corrected this argument. TG: looks good, commenting out.}
	%an isomorphism.
\end{proof}

% [TG: I couldn't follow this, mostly because I don't know what this
% ``induced morphism'' is. I wasted a couple of hours proving it myself
% (maybe not a total waste of time, as I've probably never done this
% before). Here is a sketch (not bothering to write ``almost''): we need
% to show that the kernel of $S\to\prod_{g\in G}S$, $x\mapsto
% (gx-x)_{g\in G}$ is just~$R$. It's enough to check this after
% applying~$S\otimes_R -$. (It's key that this tensor is on the left -
% tensoring on the right doesn't help!) Then replacing $S\otimes_R S$
% with $\prod_{h\in G}S$ we find that we want to compute the kernel
% of \[\prod_{h\in G}S\to\prod_{g,h\in G}S,\] \[(x_h)_h\mapsto
%   (x_{hg}-x_h)_{g,h}\]which is just the constant sequences~$(x)_h$,
% which are the things of the form $x\otimes 1\in S\otimes_R S$, which
% is the image of $S\otimes_R R\into S\otimes_R S$, as required.]

% [ME: Extending scalars gives a map $S \to S\otimes_R S^G$.  Faithful flatness 
% shows that the natural map $S\otimes S^G \to (S\otimes S)^G$ is an isomorphism.
% Putting these together gives the induced morphism
% $S \to (S\otimes S)^G .$
% Also, I think that your argument is a pretty elaborate way of noting that 
% the $G$-fixed point set of the product $\prod_{g \in G} S$, with $G$ permuting 
% the factors, is precisely the diagonal copy of $S$.]\mar{TG: thanks; I
% was being very stupid and had gotten confused about what the $G$-action on
% $S\otimes_RS$ is - I think I was thinking it was the diagonal
% action... commenting out.}

% \begin{lemma}\mar{TG: hopefully this can be removed. TG: it's
%     currently cited, as is the other result that we were thinking of
%     removing, but I do remember ME having a plan for eliminating these.}
% 	\label{lem:subgroups}
% 	If $H$ is a subgroup of $G$, then
% 	$S^H \to S$ is almost Galois with Galois group $H$. If $H$ is
%         normal in~$G$, then $R\to S^H$ is almost Galois with Galois group~$G/H$.
% \end{lemma}
% \begin{proof}Since $R\to S$ is faithfully flat, the natural map
%   $S\otimes_R S^H\to (S\otimes_RS)^H$ is an isomorphism, so by Remark~\ref{rem:base-change}, it is enough to check
%   that $(S\otimes_RS)^H\to S\otimes_RS$ is almost Galois with Galois
%   group~$H$. Equivalently, we need to check that $(\prod_{g\in
%     G}S)^H\to \prod_{g\in G}S$ is almost Galois with Galois
%   group~$H$, which is obvious.

% For the second part, we may again base change by~$S$, and reduce to
% checking that~$S\to(S\otimes_RS)^H$ is almost Galois with Galois
% group~$G/H$, and thus reduce to the same statement for $S\to (\prod_{g\in
%     G}S)^H$, which is again obvious.
% 	% Exercise.
% 	% \mar{ME: Probably should actually do this!}
% \end{proof}
	

\begin{lemma}
	\label{lem:almost descent}Suppose that $R \to S$  makes
$S$ an almost Galois extension of $S$ with Galois group $G$,
and furthermore that $R \to S$ is faithfully flat.
Then if $M$ is an $S$-module equipped with a semi-linear $G$-action,
 the induced morphism $S\otimes_R M^G \to M$ is an almost
isomorphism of $S$-modules.
\end{lemma} 
\begin{proof}
	The semi-linear $G$-action on $M$ can be reinterpreted as
	an isomorphism 
	$$(\prod_{g \in G} S) \otimes_S M \cong M \otimes_S (\prod_{g \in G}
	S)$$ (where the tensor product on the left hand side is twisted
        by the $G$-action)
	of $S\otimes_R S$-modules, and hence as an almost isomorphism
	$$ S\otimes_R M \acong M\otimes_R S$$
	of $S\otimes_R S$-modules.  The claim of the lemma now follows
	from faithfully flat descent in the almost category (for which
        see~\cite[\S 3.4.1]{MR2004652}; it is in order to make this citation
	that we have assumed that $I\otimes_R I$ is $R$-flat).% \mar{TG:
%           could cite GR03 3.4.1 - if I understand their running
%           assumptions, they're assuming that $I\otimes_RI$ is a flat
%           $R$-module at this point, so maybe we should be too (and
%           checking it in the application)? ME: It will hold in our application,
%   since [GR] show it follows from $I$ being flat, which holds (since
% it is equivalent to being torsion free) in the valuation ring setting;
% and then it is preserved by base-change (as [GR] show).  But I don't 
% know if we actually need it; I guess we have to investigate this. TG:
% not sure what ``need it'' means - do you mean that the assumption that
% $I\otimes_RI$ is flat may not
% be needed? ME: Exactly; although we do have it in application,
% it would be silly to make a song-and-dance about checking it 
% if it's not actually needed. TG: I looked, and I think it's being used
% pretty crucially in the argument in \cite{MR2004652}, in particular in
% Prop.\ 2.4.35. So I don't think it's going to be easy to avoid, and in
% the interests of writing this in finite time, I've gone ahead and put
% it in\dots}
\end{proof}

The following lemma explains our interest in almost Galois extensions.

\begin{lemma}
	\label{lem:almost etale to almost Galois}
	If $F \subseteq F'$ is a finite Galois extension
	of perfectoid fields, 
with Galois group $G$, then the corresponding inclusion of rings
of integers $\cO_F \subseteq \cO_{F'}$ 
realizes $\cO_{F'}$ as an almost 
Galois extension of $\cO_F$, with Galois group $G$
{\em (}the ``almost'' structure being understood with respect
to the maximal ideal of $\cO_F${\em )}.
\end{lemma}
\begin{proof}
	There are various more-or-less concrete ways to see this.
	For example, 
\cite[Prop.~5.23]{ScholzePerfectoid} shows that
the morphism $\cO_F \to \cO_{F'}$ is almost \'etale,
which by definition
\cite[Def.~4.12]{ScholzePerfectoid} means that the surjection
of $\cO_{F'}$-algebras $\cO_{F'}\otimes_{\cO_F} \cO_{F'} \to \cO_{F'}$
may be ``almost'' split, so that $\cO_{F'}$ is almost a direct factor
of the source.  Permuting such a splitting under the action
of the Galois group $G$ then allows one to show that the morphism
$\cO_{F'}\otimes_{\cO_F}\cO_{F'} \to \prod_{g \in G} \cO_{F'}$ is an almost
isomorphism.

However, a more direct (but less explicit) way to prove the lemma is to
use the first equivalence of categories in \cite[Thm.~5.2]{ScholzePerfectoid},
namely the equivalence between the category of perfectoid 
$F$-algebras and the category of perfectoid almost algebras
over~$\cO_F$.  
Under this equivalence, the morphism $\cO_{F'}\otimes_{\cO_F} \cO_{F'}
\to \prod_{g \in G} \cO_{F'}$ is taken to the morphism
${F'}\otimes_F {F'} \to \prod_{g \in G} {F'}$.  This latter morphism
{\em is} an isomorphism, since ${F'}$ is Galois over $F$ with Galois
group~$G$, and so we conclude 
that the former morphism is an almost isomorphism,
as required.
\end{proof}

In the case of perfectoid fields of characteristic~$p$,
we may extend the statement of the preceding lemma to the context
of truncated rings of Witt vectors.

\begin{lem}
	\label{lem:almost Galois Witt version}
	If $F \subseteq F'$ is a finite Galois extension
	of perfectoid fields in characteristic~$p$, 
	with Galois group~$G$,
	then for each $a \geq 1$, the inclusion
	$W_a(\cO_{F}) \hookrightarrow W_a(\cO_{F'})$
	is almost Galois,
	%\mar{ME: Have to explain the choice of ideal $I$ w.r.t.\ which
	%	we define ``almost'' . ME: Done}
	with Galois group $G$
{\em (}the ``almost'' structure being understood with respect
to the maximal ideal of $W_a(\cO_F)${\em )}.
\end{lem}
\begin{proof}
	For simplicity of notation, write $R := W_a(\cO_F)$
	and $S := W_a(\cO_{F'})$.  Each of $R$ and $S$ is flat
	over $\Z/p^a$, and $S$ is also flat over $R$,
	by Lemma~\ref{lem:Witt flatness}.
	Thus $S\otimes_R S$ is flat over $S$, hence over $R$,
	and hence over $\Z/p^a$ as well.

	To see that the natural morphism
	%\numequation
	%\label{eqn:test map}
	$S\otimes_R S \to \prod_{g \in G} S$
%\end{equation}
	is an almost isomorphism, it suffices to check the analogous
	condition after passing to associated graded rings for the
	$p$-adic filtration.   Lemma~\ref{lem:graded flat} allows
	us to rewrite this induced morphism on associated graded rings
	in the form
	$$\F_p[T]/(T^a) \otimes_{\F_p} \bigl((S/pS)\otimes_{(R/pR)}(S/pS)\bigr)
	\to \F_p[T]/(T^a) \otimes_{\F_p} \prod_{g \in G} (S/pS)$$
	(here $\F_p[T]/(T^a)$ appears as the associated graded ring
	to $\Z/p^a$ with its $p$-adic filtration),
	which may be identified with the base-change over $\F_p[T]/(T^a)$
	of the natural morphism
	$$(S/pS)\otimes_{(R/pR)}(S/pS) \to \prod_{g \in G} (S/pS).$$
	This latter morphism is an indeed an almost isomorphism,
	by Lemma~\ref{lem:almost etale to almost Galois}.
	%\mar{ME: Have to give argment --- via passing to the associated
	%	graded w.r.t.\ the $p$-adic filtration.  Also have
	%	to explain why we can use the maximal ideal of
	%	$W_a(\cO_F^\flat)$
	%	rather than that of
	%	$W_a(\cO_{F_m}^\flat)$ --- which amounts to the fact 
	%	that the former is contained in the latter.  And maybe
	%	the former generates the latter anyway?}
\end{proof}



% \begin{cor}\mar{TG: this can hopefully be removed,}
%   \label{cor: almost uniqueness of Galois extensions}If $S\to T$ is
%   such that the composite $R\to T$ is
%   faithfully flat, and the $G$-action on~$S$ extends to a $G$-action
%   on~$T$, making the composite $R\to T$ an almost Galois extension
%   with Galois group~$G$, then~$S\to T$ is an almost isomorphism.
% \end{cor}
% \begin{proof}
%   Combining Lemmas~\ref{lem:almost invariants} and~\ref{lem:almost
%     descent}, we have a composite of almost isomorphisms
%   $S=S\otimes_RR\to S\otimes_R T^G\to T$.% \mar{TG: I don't know if we
% %     want this lemma explicitly, but it seems to me that we're using it
% %   implicitly in 1.2. ME: Do we actually need $S\to T$ to be faithfully
% % flat?  I guess the minimum assumption we need to apply what's come
% % before is that $R\to T$ is faithfully flat, right? TG: we need $R\to
% % T$ faithfully flat to invoke Lemma~\ref{lem:almost invariants}, and
% % $R\to S$ faithfully flat to invoke Lemma~\ref{lem:almost
% %     descent}. I guess that's weaker, so we could just assume those -
% %   I've edited it accordingly (I thought when I first read this that
% %   you were dropping $R\to S$ faithfully flat, but I guess that's in
% %   place throughout). Commenting out.}
% \end{proof}




We now pass to the profinite setting. We continue to suppose that we are given the ring $R$ endowed with
an idempotent ideal $I$.  We suppose additionally that $v \in I$
is a regular element of $R$ (i.e.\ a non-zero divisor) and 
that $R$ is $v$-adically complete.

We also suppose given a $v$-adically complete $R$-algebra $S$, 
% \mar{TG: also separated? That seems to be used in the proof of
%   Lemma~\ref{lem:almost invariants two}, or if it isn't then we should
% say something more for why we can pass to the inverse limit in~$i$, I guess.
% ME: I think that throughout we use the convention (as in the stacks project)
% that {\em complete} means {\em complete and separated}. Should
% probably mention this somewhere! TG: mentioned in the notation
% section, commenting out.}
equipped with an action of a profinite group $G$ as $R$-algebra
automorphisms.  We assume that this action is continuous, 
in the sense that the action map $G \times S \to S$ is continuous
when $S$ is endowed with its $v$-adic topology.

We further suppose that
we may write $S = \widehat{\bigcup_n S_n}$ as the $v$-adic completion
of an increasing union of $G$-invariant $v$-adically complete
$R$-subalgebras $S_n$, with $R=S_0$; that the $G$-action
on $S_n$ factors through a finite quotient $G_n :=G/H_n$ of~$G$,
where~$H_n$ is a normal open subgroup of~$G$; and that for each~$i\ge
0$ and each  $m \leq n$,  the morphism $S_m/v^i \to S_n/v^i$ is faithfully flat,
 % and that the 
% where each $S_n$ is $u$-adically complete and
% faithfully flat over $R$, and in fact faithfully flat over $S_m$ 
% for each $m \leq n$, 
% \mar{ME: I think we should only require this faithful flatness 
% 	after passing mod $u^i$ for any $i > 0$.  This a weaker 
% 	condition which is preserved under base-change followed
% 	by $u$-adic completion (because the $u$-adic completion
% 	goes away after reducing mod $u^i$), and that's the 
% 	process we have to perform when passing from the
% 	classical setting to the setting with coefficients. TG:
%         agreed, that sounds good.}
% and where the $G$-action
% on $S_n$ factors through a finite quotient $G_n$ of $G$,
and realizes $S_n/v^i$ as an almost Galois
extension of $S_m/v^i$, having the subgroup $H_m/H_n$ of $G_n$ as Galois group.
 % , for each $i \geq 0$.  
% We may and do choose our labelling so
% that $S_0 = R$.

% We define the normal open subgroups $H_n$ of $G$ by writing $G_n = G/H_n$.
% If $n \geq m,$ then our assumptions imply
% (by
% %Remark~\ref{rem:base-change} and
% Lemma~\ref{lem:subgroups})
% that the faithfully flat morphism
% %\footnote{It {\em is} an inclusion,
% %       because it is a base-change of a faithfully flat morphism,
% %and hence is faithfully flat.}     
%  $S_m/u^i \to S_n/u^i$ is almost Galois,
% with Galois group $H_m/H_n$, for any $i \geq 0$.
% (More precisely,
% Lemma~\ref{lem:subgroups} shows that
% $(S_n/u^i)^{H_m/H_n} \to S_n/u^i$ is almost Galois with Galois group
% $H_m/H_n$,
% while 
% Lemma~\ref{lem:subgroups} and Corollary~\ref{cor: almost uniqueness of Galois extensions}
% \mar{ME: I'm a bit unsure how to get the necessary faithful flatness
% 	hypothesis for the Corollary, so this needs more thinking 
% 	about. In the application, the claimed almost Galoisness will
% follow (I think)
% just by replacing $R$ with $S_m$ from the beginning, so we shouldn't
% stress too much about finding the absolutely minimal set of assumptions
% from which we can deduce everything else we need (especially since
% the assumptions are a bit of a mess anyway).}
% together show that
% the natural map $S_m/u^i\to (S_n/u^i)^{H_m/H_n}$ is an almost isomorphism.)

\begin{lemma}
	\label{lem:almost invariants two}
	For any value of $m$, the morphism $S_m \to S^{H_m}$
	is an almost isomorphism, as is the induced morphism
	$S_m/v^i \to (S/v^i)^{H_m}$ for any $i \geq 0$.
\end{lemma}
\begin{proof}
	Our hypotheses, together with Lemma~\ref{lem:almost invariants},
        imply that the morphism
	$$S_m/v^i \to (S_n/v^i)^{H_m/H_n}$$
	is an almost isomorphism for each $n \geq m$ and each
	$i \geq 0.$  Passing to the direct limit over $n$
	gives the second claim,
	and then additionally passing to the inverse limit over $i$ 
	gives the first claim.%  TG: both of these are easy
        % checks from the definitions.)
	% Now since $S_m \to S_n$ is faithfully flat,
	% so is the induced morphism $S_m/u^i \to S_n/u^i$. 
       	% This morphism is almost Galois
	% with Galois group $H_m/H_n$,
	% by assumption,
	% and so the required result follows
	% from Lemma~\ref{lem:almost invariants}.
\end{proof}

\begin{cor}
	\label{cor:invariants}
	The injection $R[1/v] \hookrightarrow S[1/v]^G$ 
	is an isomorphism.
\end{cor}
\begin{proof}
	Note that $S[1/v]^G = (S^G)[1/v],$ since
	$v$ is $G$-invariant (being an element of $R$) 
	and localization is exact.
	The corollary thus follows from the $m = 0$ case of 
	Lemma~\ref{lem:almost invariants two}, as inverting $v$
	converts the almost isomorphism into a genuine isomorphism.
\end{proof}

% \begin{df}
% 	We say that an $S$-module $M$ is {\em iso-free of finite rank}
% 	if it is $u$-torsion free, and if we may find an $S$-submodule
% 	$N\subseteq M$ such that $u^c M \subseteq N$ for some $c \geq 0$,
% 	and such that $N$ is free of finite rank.
% \end{df}

% \begin{remark}
% 	\label{rem:iso-free implies complete}
% 	Since $S$ is assumed to be $u$-adically complete,
% 	any finite rank free $S$-module is again $u$-adically complete.
% 	It follows easily that any iso-free $S$-module is $u$-adically
% 	complete.
% \end{remark}


% \begin{thm}
% %	\mar{TG: this should be rewritten to do iso-projective,
% %    using an idempotent with $u$-power denominators cleared; this will
% %  introduce some extra powers of~$u$ to the error terms.}
% 	\label{thm:almost descent two}
% 	If $M$ is %an $u$-adically complete $S$-module that is
%         an iso-free $S$-module of finite rank,
% 	equipped with a semi-linear $G$-action that is continuous
% 	with respect to the $u$-adic topology on~$M$,
% 	then the kernel and cokernel 
% 	of the induced morphism $S \cotimes_R M^G \to M$
% 	{\em (}the source being the $u$-adically completed tensor
% 	product{\em )}
% 	are each annihilated by a %uniformly bounded
% 	power of $u$.%\mar{TG: delete ``uniformly bounded''? ME: Done, commenting out.}
% \end{thm}
% \begin{proof}
% 	Choose a finite rank free $R$-module $M_0$
% 	such that we can find an embedding $M \subseteq S\otimes_R M_0$
% 	whose image contains $u^c S\otimes_R M_0$ (for example,
%         take~$N$ as in the definition of ``iso-free'', let ~$N_0$ be
%         the $R$-submodule generated by an $S$-basis of~$N$, and take~$M_0=u^{-c}N_0$).
% 	The $G$-action on $M$ induces a $G$-action on
% 	$$M[1/u] := S[1/u]\otimes_S M
% 	= S[1/u]\otimes_R M_0,$$ and thus on
% 	$$ M[1/u]/M = (S[1/u]\otimes_R M_0)/ M.$$ 
% 	The $u$-adic topology on $M$
% 	induces the discrete topology on this module, and so 
% 	the $G$-action is continuous when this module
% 	is equipped with its discrete topology.
% 	Thus, for any $n \geq 0$, there is an open subgroup 
% 	of $G$ which fixes the $R$-submodule
% 	$u^{-n}M_0/(u^{-n} M_0\cap M)$ pointwise,
% 	and (since $u^{n+c}$ annihilates this quotient)
% 	$G\bigl(u^{-n} M_0/(u^{-n} M_0 \cap M)\bigr)$
% 	is contained in \[(u^{-(n+c)}T\otimes_R M_0)/
%   \bigl((u^{-(n+c)}T\otimes_R M_0) \cap M\bigr)\] %\mar{TG: I think
% %           this should say $T\otimes_R(u^{-(n+c)}M_0/M_0)$.  ME: Everything
% %   is happening inside $M[1/u]/M$, and the image is killed by $u^{n+c}$,
% %   since the source is,
% %   so it lies in $u^{-(n+c)}M/M,$ and hence in $(u^{-(n+c)} S\otimes_R M_0)/M$.
% %   Now we can replace $S$ by something f.t.\ over $R$, so
% %   it looks like it should be $(u^{-(n+c)}T\otimes_R M_0)/
% %   \bigl((u^{-(n+c)}T\otimes_R M_0) \cap M\bigr)$, I think. TG: going
% %   with that, commenting out.
% % }
% 	for some finitely generated $R$-subalgebra $T$ of $\bigcup_n S_n$.
%         Passing to a cofinal subsequence of $S_n,G_n,$ etc., and relabelling
% 	as necessary, we may assume that 
% 	$H_n$ acts trivially on $u^{-n} M_0/(u^{-n} M_0\cap M)$, and that $T \subseteq S_n$.

% 	If we let $M_n$ denote
% 	the $S_n$-submodule of $M[1/u]/M$
% 	generated by the $G$-translates of $u^{-n}M_0/(u^{-n}M_0 \cap M),$
% 	then we see that the $G$-action on $M[1/u]/M$ restricts to a
% 	$G_n := G/H_n$-action on $M_n$.
% 	If we let $SM_n$ denote the $S$-submodule of $M[1/u]/M$ generated
% 	by $M_n$, then we have inclusions
% 	$$ (u^{-n}S \otimes_R M_0)/M \subseteq S M_n \subseteq
% 	u^{-(n+c)} M/M \subseteq
% 	(u^{-(n+c)} S \otimes_R M_0)/M,$$% \mar{TG: I think $u^{-(n+c)}
% %           M/M$ should just be deleted here.  ME: You're probably right,
% %   since I don't think it's used below.  But it should probably appear
% % above as part of the explanation for where $T$ comes from (note that
% % it does appear in my response to your previous comment). TG: actually
% % I decided that I wanted it later, so I'm leaving it in; commenting out.}
% 	from which we deduce that
%         $$ u^c S M_n \subseteq  (u^{-n} S \otimes_R M_0)/M \subseteq SM_n.$$
% 	Passing to $H_n$-invariants, we obtain the inclusions
% 	\numequation
% 	\label{eqn:chain of inclusions}
% 	u^c (S M_n)^{H_n} \subseteq
% 	\bigl((u^{-n} S \otimes_R M_0)/M\bigr)^{H_n} \subseteq (SM_n)^{H_n}.
%         \end{equation}
	
% 	Lemma~\ref{lem:almost invariants two} shows that
% 	$S_n/u^n  \to (S/u^n)^{H_n}$ is an almost isomorphism.
% 	Thus $(u^{-n} S_n/S_n) \otimes_R M_0 \to
% 	(u^{-n} S/S \otimes_R M_0)^{H_n}$ is an almost
%         isomorphism (since~$M_0$ is a free $R$-module, and
%         ~$H_n$ acts trivially on~$u^{-n}M_0/(u^{-n} M_0\cap M)$), % \mar{TG: should the target be $(u^{-n} S/S
% %           \otimes_R M_0)^{H_n}$? Also, this seems to use that~$H_n$ is
% %         acting trivially and $S_n/u^n$ is flat over $R/u^n$? If so we
% %         should maybe say so. ME: {\em Yes} to the question about the target ---
% % 	this is now fixed; {\em yes} to the fact that we
% % use that $H_n$ is acting trivially; {\em no} to the flatness
% % as far as I can tell, because $M_0$ is free over $R$, so tensoring
% % with it is just taking a bunch of direct sums. TG: OK, added a little,
% % and adjusted above; feel free to comment out.}
% 	and so,
% 	since $(S\otimes_R M_0)/M$ is annihilated by $u^c$,
% 	if we let $M_n'$ denote the image of the natural map
% 	%\numequation
% 	%\label{eqn:natural map}
% 	$$
% 	(u^{-n} S_n \otimes_R M_0)/(u^cS_n \otimes_R M_0) \to
% 	\bigl((u^{-n} S \otimes_R M_0)/M\bigr)^{H_n},
% 	$$
%         %end{equation}
% 	then the cokernel of the inclusion
% 	$$M_n'\subseteq
% 	\bigl((u^{-n} S \otimes_R M_0)/M\bigr)^{H_n}
% 	$$
% 	is annihilated by $u^c I$.
% 	% \mar{ME: I expanded on the argument here slightly,
% 	% 	since it took me a while to reconstruct it
% 	% 	as I was reading through. TG: looks good, commenting out.}
	
% 	Now, by construction, we have inclusions
% 	$$M_n' \subseteq M_n \subseteq (SM_n)^{H_n}.$$ 
% 	Taking into account the result of the preceding paragraph,
% 	as well as the inclusions~(\ref{eqn:chain of inclusions}),
% 	we conclude that the cokernel of the inclusion
% 	$M'_n \subseteq (SM_n)^{H_n}$,
% 	and hence also the cokernel of the inclusion
% 	$M_n \subseteq (SM_n)^{H_n}$,
% 	is annihilated by $u^{2c} I$.
% 	Passing to $G_n$-invariants in the inclusion just mentioned,
% 	we find that there is an inclusion
% 	$M_n^{G_n}\subseteq (S M_n)^G$,
% 	whose cokernel is annihilated by $u^{2c}I$.
% 	Extending scalars to $S_n$, we obtain a morphism
% 	%Since $S_n/u^{n+c}$ is flat over $R/u^{n+c}$,
% 	%this last inclusion induces an inclusion 
% 	$$S_n\otimes_R M_n^{G_n} \to S_n\otimes_R (S M_n)^{G}$$
% 	(which is in fact an embedding,
% 	since $S_n/u^{n+c}$ is flat over $R/u^{n+c}$,
% 	although we don't need this here),
% 	whose cokernel is annihilated by $u^{2c}I$.
% 	On the other hand,
%         Lemma~\ref{lem:almost descent} implies that the natural morphism
% 	$S_n\otimes_R M_n^{G_n} \to M_n$ is an almost isomorphism.
% 	Putting all these results together,
% 	we see that the kernel and cokernel of the natural morphism
% 	$$S_n\otimes_R (S M_n)^{G} \to (S M_n)^{H_n}$$
% 	are each annihilated by $u^{2c} I$.

% 	Since $M_n \subseteq (S M_n)^{H_n},$ we see that the natural map
% 	\numequation
% 	\label{eqn:natural surjection}
% 	S\otimes_{S_n}  ( S M_n)^{H_n} \to S M_n
%         \end{equation}
% 	is surjective.
% 	We bound the exponent of its kernel as follows:
% 	The inclusion $M_n' \subseteq (S M_n)^{H_n}$,
%         whose cokernel we have shown is annihilated by $u^{2c}I$,
% 	induces an inclusion 
% 	$$S\otimes_{S_n} M_n' \subseteq S\otimes_{S_n} (S M_n)^{H_n},$$
% 	whose cokernel is again annihilated by $u^{2c}I,$
% 	while the defining surjection $(u^{-n} S_n/u^cS_n)\otimes_R M_0 \to M_n'$
% 	then induces a surjection
% 	$$(u^{-n} S\otimes_R M_0)/(u^c S\otimes_R M_0)
% 	\to S\otimes_{S_n} M_n' .$$
% 	Now the natural morphism
% 	$$(u^{-n}S\otimes_R M_0)/(u^c S\otimes_R M_0) \to M[1/u]/M$$
% 	has kernel annihilated by $u^c$,
% 	and so we find that the kernel of~(\ref{eqn:natural surjection})
% 	is annihilated by $u^{3c} I$.

% 	Putting together the results of the preceding two paragraphs,
% 	we find that the cokernel of the natural morphism
% 	$$S\otimes_R (S M_n)^G \to S M_n$$ is annihilated by $u^{2c} I$,
% 	while its kernel is annihilated by $u^{5 c} I$.

% 	Now $u^{-n} M /M \subseteq S M_n \subseteq u^{-(n+c)}M/M,$
% 	and so (recalling that $u \in I$)
% 	we conclude that the kernel and cokernel of the natural
% 	morphism
% 	$$S\otimes_R (u^{-n}M/M)^G \to u^{-n}M/M$$
% 	are annihilated by a power of $u$ that is uniformly bounded
% 	(i.e.\ bounded independently of $n$).  Passing to an
% 	inverse limit in $n$ (with respect to the morphisms induced
% 	by multiplication by $u$), the theorem follows,
% 	once we recall that $M = \varprojlim_{n} u^{-n}M/M$
% 	(by Remark~\ref{rem:iso-free implies complete}),
% 	%(we have assumed that $M$ is $u$-adically complete),
% 	and so $M^G = \varprojlim_{n} (u^{-n}M/M)^G.$ 
% 	(We should also note that the topology on $M^G$ induced
% 	by the $u$-adic topology on $M$ coincides with the $u$-adic
% 	topology on $M^G$; but this follows immediately from
% 	that fact that multiplication by $u$ on $M$ is injective
% 	(as $M$ is iso-free)
% 	and commutes with the $G$-action.)
% \end{proof}

% % \begin{remark}
% % 	In the preceding theorem, if $M$ is in fact free over $R$,
% % 	so that the constant $c$ appearing in the proof is equal to $0$,
% % 	then we see from the argument that the morphism
% % 	$S\cotimes_R M^G \to M$ is an almost isomorphism.
% % 	We won't have need of this fact in the sequel.
% % \end{remark}


\begin{df}
	% \mar{ME: I think this is the correct definition, but ultimately
	% 	we have to see exactly what is used in the argument below.}
	We say that an 
       	$S$-module $M$ is {\em iso-projective of finite rank} \index{iso-projective}
	if it is finitely generated and
	$v$-torsion free, and if $S[1/v]\otimes_S M$ 
	is a projective $S[1/v]$-module.
	%we may find an $S$-submodule
	%$N\subseteq M$ such that $u^c M \subseteq N$ for some $c \geq 0$,
	%and such that $N$ is free of finite rank.
\end{df}

%\mar{ME: Have to find proof of the remark; and (relatedly) may want to
%	upgrade it to a lemma. TG: had a go, feel free to improve.}
\begin{lem}
	\label{lem:iso-proj implies complete}
	Any iso-projective $S$-module of finite rank is $v$-adically complete.
\end{lem}
\begin{proof}
 If~$M$ is iso-projective, then we may write $M[1/v]=eF$ where~$F$ is a finite
 free $S[1/v]$-module and $e\in\End_{S[1/v]}F$ is an
 idempotent.
 Write~$F=F_0\otimes_SS[1/v]$, where
 $F_0$ is a finite free $S$-module.
 Then, since $F_0$ is finitely generated,
 we find that $eF_0$ is contained in $v^{-a}F_0$ for some sufficiently large
 value of $a$; the latter $S$-module is $v$-adically separated, and
 thus so is the former.  Since $eF_0$ is furthermore
 the image of the $v$-adically complete $S$-module $F_0$, it is in fact
 $v$-adically complete.  Since $eF_0[1/v] = M[1/v]$, and both $eF_0$
 and $M$ are finitely generated $S$-modules,
 we find that $v^b eF_0 \subseteq M \subseteq v^{-b} e F_0$ 
 for some sufficiently large value of $b$,
 and thus $M$ is also $v$-adically
 complete, as claimed.
%
%
%
 %Since~$e$ is linear, it is continuous, so $M[1/u]$ is a
 %closed $S[1/u]$-module of~$F$.
 %Writing~$F=F_0\otimes_SS[1/u]$ where
 %$F_0$ is a finite free $S$-module,
 %and multiplying by a power of~$u$
 %if necessary, we see that~$M$ is a closed submodule of~$F_0$, and is
 %therefore $u$-adically complete.
\end{proof}

\begin{thm}
	\label{thm:almost descent three}
	If $M$ is %an $u$-adically complete $S$-module that is
        an iso-projective $S$-module of finite rank,
	equipped with a semi-linear $G$-action that is continuous
	with respect to the $v$-adic topology on~$M$,
	then the kernel and cokernel 
	of the induced morphism $S \cotimes_R M^G \to M$
	{\em (}the source being the $v$-adically completed tensor
	product{\em )}
	are each annihilated by a %uniformly bounded
	power of $v$.%\mar{TG: delete ``uniformly bounded''? ME: Done, commenting out.}
\end{thm}
% \mar{ME: Currently this is just the proof in the iso-free case.  I will
% 	update it to the iso-proj.\ case soon. ME: Okay, it is now in a
% period of transition.}
\begin{proof}
        Write $P := S[1/v]\otimes_S M,$ so that (by assumption) $P$ is a
	finitely presented projective module over $S[1/v]$. 
	Since $M$ is $v$-adically complete
	(by Lemma~\ref{lem:iso-proj implies complete}),
	we have an isomorphism $M \iso \varprojlim_n M/v^n M$.
	This induces a corresponding isomorphism
	$M^G \iso \varprojlim_{n} (M/v^nM)^G.$ 
	(We should also note that the topology on $M^G$ induced
	by the $v$-adic topology on $M$ coincides with the $v$-adic
	topology on $M^G$, as follows immediately from
	that fact that multiplication by $v$ on $M$ is injective
	(as $M$ is iso-projective)
	and commutes with the $G$-action.)
	Thus, to prove the theorem,
	it suffices to show that the kernel and cokernel
	of each of the morphisms
	$$S\otimes_R (M/v^n M)^G \to M/v^n M$$
	is annihilated by a power of $v$ that is bounded independently
	of $n$.

	Multiplication by $v^{-n}$ induces an isomorphism
	$M/v^n M \iso v^{-n} M/M,$ and we will actually prove the 
	equivalent statement that each of the morphisms
	\numequation
	\label{eqn:M morphisms at level n}
	S\otimes_R (v^{-n} M/ M)^G \to v^{-n} M/M
\end{equation}
	has kernel and cokernel
	annihilated by a power of $v$ that is bounded independently
	of~$n$.

	The reason for formulating our argument in terms of
	the quotients $v^{-n}M/M$ is that these may be conveniently
	be regarded as submodules of the quotient $P/M$.
        Our argument will proceed by constructing,
	for each~$n$, a $G$-invariant $S$-submodule $Z_n$ of $P/M$,
	such that $v^{-n}M/M \subseteq Z_n \subseteq v^{-(n+c)} M/M$
	(where $c$ is independent of~$n$),
	and such that each of the morphisms
	\numequation
	\label{eqn:Z morphisms at level n}
	S\otimes_R Z_n^G \to Z_n
\end{equation}
	has kernel and cokernel
	annihilated by a power of $v$ that is bounded independently
	of~$n$.  This implies the corresponding statement for the
	morphisms~(\ref{eqn:M morphisms at level n}), and thus establishes
	the theorem.
	The remainder of the argument is devoted to constructing the
	submodules $Z_n$, and proving the requisite properties
	of the morphisms~(\ref{eqn:Z morphisms at level n}).


	Since $P$ is a projective $S[1/v]$-module, we may 
	choose a finite rank free module $F$ over $S[1/v]$, and an
	idempotent $e \in \End_{S[1/v]}(F)$, such that $P = eF.$
	We choose a free $R$-submodule $F_0$ of $F$ such that
	$S[1/v]\otimes_R F_0 \iso F.$  (More concretely, $F_0$ is simply
	the $R$-span of some chosen $S[1/v]$-basis of $F$.)
	The endomorphism $e$ may not preserve the $S$-submodule
	$S\otimes_R F_0$ of $F$, but if we choose $a$ sufficiently large,
	then $e' := v^a e$ will preserve $S\otimes_R F_0.$

	Since $M$ is finitely generated over $S$, we may and do assume
	that we have chosen $F_0$ in such a manner 
	that $M \subseteq v^a(S\otimes_R F_0)$.  (Simply replace $F_0$ by $v^{-b} F_0$
	for some sufficiently large value of $b$.)
	Then in fact $M \subseteq e'(S\otimes_R F_0).$   Furthermore,
	since $M$ and $e'(S\otimes_R F_0)$ both span $P$ as an $S[1/v]$-module,
	we find that
	$v^c e'(S\otimes_R F_0) \subseteq M$ for some sufficiently large 
	value of $c$.

	For any $n \geq 0$ we have $S/v^n = \bigcup_{i\geq 0} S_i/v^n.$ 
	Thus $e' \bmod v^n$, which is an endomorphism of $(S/v^n)\otimes_R F_0$,
	descends to an endomorphism of $(S_i/v^n)\otimes_R F_0$ for
        all sufficiently large~$i$.  Replacing
	$(S_n)$ by an appropriately chosen subsequence, we may and do assume that
	in fact for each~$n$, $e'$ descends to an endomorphism $e'_n$ of
	$(S_n/v^n)\otimes_R F_0$.

	Since $G$ acts continuously on $P$ and preserves $M$,
	it acts continuously on $P/M$.  The topology on $P/M$ is
	discrete, and thus any finite subset of $P/M$ is fixed
	by some open subgroup of $G$.  In particular,
	we find that
	\[v^{-n} e'F_0 /
	\bigl( v^{-n}e' F_0 \cap M\bigr)\]
	(which we regard in the natural way as a submodule of $P/M$,
	and which we note is finitely generated over $R$)
	is pointwise fixed by some open subgroup of $G$.  Since the $H_n$
	form a cofinal sequence of open subgroups of $G$, if we again
	replace $S_n$ and $H_n$ by appropriately chosen subsequences,
	we may and do assume that in fact
	$v^{-n} e'F_0/ \bigl( v^{-n} e'F_0 \cap M\bigr)$
	is pointwise fixed by $H_n$.

	Since $M \supseteq v^c e'(S\otimes_R F_0),$ we see that
	$v^{-n} e'F_0/ \bigl( v^{-n} e'F_0 \cap M\bigr)$ is annihilated
	by $v^{n+c}$, and thus that 
	$$G \Bigl(
	v^{-n} e'F_0/ \bigl( v^{-n} e'F_0 \cap M\bigr)\Bigr) 
	\subseteq v^{-(n+c)}M/M \subseteq 
	v^{- (n+c)} e'(S\otimes_R F_0)/ M.$$
	Since $G/H_n$ is finite, and since $H_n$ fixes $v^{-n} e'F_0/ \bigl( v^{-n} e'F_0 \cap M\bigr)$
 pointwise, we see that
	in fact
	$$G \Bigl(
	v^{-n} e'F_0/ \bigl( v^{-n} e'F_0 \cap M\bigr)\Bigr) 
	\subseteq
	v^{- (n+c)} e'(T\otimes_R F_0)/ M,$$
	for some finitely generated $R$-subalgebra $T$ of $S/v^{n+c} S.$
	Passing to a subsequence one more time, we may assume
	that $T$ is contained in $S_n$.



%	\mar{ME: Now have to consider the $G$-translates of 
%	$u^{-n} F_0/ \bigl( u^{-n} F_0 \cap M\bigr)$, and make another
%	relabelling to have them lie in the $S_n$-span of something.}

	We let $Y_n$, respectively $Z_n$,
       	denote the $S_n$-submodule, respectively the $S$-submodule, of $P/M$ 
	generated by the $G$-translates of 
	$v^{-n} e'F_0/ \bigl( v^{-n} e'F_0 \cap M\bigr)$.
	It remains to prove the requisite properties of $Z_n$.
	We begin by noting the inclusions
	$$v^{-n} e'(S\otimes_R F_0) / M \subseteq Z_n \subseteq 
	v^{-(n+c)}M/M \subseteq v^{-(n+c)} e' (S \otimes_R F_0)/M,$$
	which imply that
	$$v^c Z_n \subseteq v^{-n} e'(S\otimes_R F_0)/M \subseteq Z_n,$$
	and thus that 
	\numequation
	\label{eqn:chain of inclusions}
	v^c Z_n^{H_n} \subseteq
	(v^{-n} e'(S\otimes_R F_0)/M)^{H_n} \subseteq Z_n^{H_n}.
        \end{equation}

	Lemma~\ref{lem:almost invariants two} shows that
	$S_n/v^n  \to (S/v^n)^{H_n}$ is an almost isomorphism.
	Thus $(v^{-n} S_n/S_n) \otimes_R F_0 \to
	(v^{-n} S/S \otimes_R F_0)^{H_n}$ is an almost
        isomorphism, if we declare that $H_n$ acts trivially on $F_0$,
	and thus via its action on the first factor in the 
	target tensor product.
	%and $G$-semi-linearly on the tensor product $S\otimes_R F_0.$
	Since $(e')^2 = v^a e',$ we find that the cokernel
	of the inclusion
	$$
	e' \bigl( (v^{-n} S\otimes_R F_0)/ (S\otimes_R F_0) \bigr)^{H_n}
	\hookrightarrow
	\bigl(e' ( v^{-n} S\otimes_R F_0)/ (S\otimes_R F_0) \bigr)^{H_n}
	$$
	is annihilated by $v^a$, and thus that the natural morphism
	$$
	e' (v^{-n} S_n\otimes_R F_0)/ (S_n\otimes_R F_0) 
	\to 
	\bigl(e' ( v^{-n} S\otimes_R F_0)/ (S\otimes_R F_0) \bigr)^{H_n}
	$$
	has its kernel annihilated by $I$, and its cokernel 
	annihilated by $v^a I.$

	If we let $X_n$ denote the image of $e'(v^{-n} S_n\otimes_R F_0)$
	in 
	$\bigl(e' ( v^{-n} S\otimes_R F_0)/ M \bigr)^{H_n},$
	then certainly
	$$X_n \subseteq Y_n \subseteq Z_n^{H_n}.$$
	It follows from the conclusion of the preceding
	paragraph, along with the fact that $v^c e'(S_0\otimes_R F_0)
	\subseteq M,$
        that the cokernel of the inclusion
	$$X_n \hookrightarrow
	\bigl(e' ( v^{-n} S\otimes_R F_0)/ M \bigr)^{H_n}$$
	is annihilated by $v^{a+c} I,$
	while the chain of inclusions~(\ref{eqn:chain of inclusions})
	shows that the cokernel of the inclusion
	$$
	\bigl(e' ( v^{-n} S\otimes_R F_0)/ M \bigr)^{H_n}
	\subseteq Z_n^{H_n} 
	$$
	is annihilated by $v^c$.  
	Thus the cokernel of the inclusion $X_n \subseteq Z_n^{H_n}$
	is annihilated by $v^{a+2c} I,$ and hence so is the cokernel
	of the inclusion $Y_n \subseteq Z_n^{H_n}$.
	Passing to $G_n$-invariants in the inclusion just mentioned,
	we find that there is an inclusion
	$Y_n^{G_n}\subseteq Z_n^G$,
	whose cokernel is annihilated by $v^{a+2c}I$.
	Extending scalars to $S_n$, we obtain a morphism
	%Since $S_n/u^{n+c}$ is flat over $R/u^{n+c}$,
	%this last inclusion induces an inclusion 
	$$S_n\otimes_R Y_n^{G_n} \to S_n\otimes_R Z_n^{G}$$
	(which is in fact an embedding,
	since $S_n/v^{n+c}$ is flat over $R/v^{n+c}$,
	although we don't need this here),
	whose cokernel is annihilated by $v^{a+2c}I$.
	On the other hand,
        Lemma~\ref{lem:almost descent} implies that the natural morphism
	$S_n\otimes_R Y_n^{G_n} \to Y_n$ is an almost isomorphism.

	Putting all these results together,
	we find that the kernel and cokernel of the natural morphism
	$$S_n\otimes_R Z_n^{G} \to Z_n^{H_n}$$
	are each annihilated by $v^{a+2c} I$.
	Indeed, the composite 
        $S_n\otimes_R Y_n^{G_n} \to S_n\otimes_R Z_n^{G} \to
        Z_n^{H_n}$ factors through the almost isomorphism
        $S_n\otimes_R Y_n^{G_n} \to Y_n$, 
	and we have shown that the natural morphism $Y_n \to Z_n^{H_n}$
	is an injection whose cokernel is killed by $v^{a+2c}I$,
	while the cokernel of
        $S_n\otimes_R Y_n^{G_n} \to S_n\otimes_R Z_n^{G}$ is also
        killed by~$v^{a+2c}I$.


%	(This follows from the
%        fact that the composite
%        $S_n\otimes_R Y_n^{G_n} \to S_n\otimes_R Z_n^{G} \to
%        Z_n^{H_n}$ factors through the almost isomorphism
%        $S_n\otimes_R Y_n^{G_n} \to Y_n$, since we have shown that
%        $Y_n\to Z_n^{H_n}$ is an injection whose cokernel is killed by
%        $u^{a+2c}I$, and the cokernel of
%        $S_n\otimes_R Y_n^{G_n} \to S_n\otimes_R Z_n^{G}$ is also
%        killed by~$u^{a+2c}$.)

	Since $Y_n \subseteq Z_n^{H_n},$ and since $Y_n$ generates $Z_n$
	over $S$ by their very definitions, we see that the natural map
	\numequation
	\label{eqn:natural surjection}
	S\otimes_{S_n}  Z_n^{H_n} \to Z_n
        \end{equation}
	is surjective.
	We bound the exponent of its kernel as follows:
	The inclusion $X_n \subseteq Z_n^{H_n}$,
        whose cokernel we have shown above to be annihilated by $v^{a+2c}I$,
	induces an inclusion 
	$$S\otimes_{S_n} X_n \subseteq S\otimes_{S_n} Z_n^{H_n},$$
	whose cokernel is again annihilated by $v^{a+2c}I,$
	while the defining surjection
	$$e'\bigl((v^{-n} S_n/v^cS_n)\otimes_R F_0\bigr) \to
	X_n$$
	induces a surjection
	$$e'\bigl((v^{-n} S\otimes_R F_0)/(v^c S\otimes_R F_0)\bigr)
	\to S\otimes_{S_n} X_n .$$
	Now the natural morphism
	\numequation\label{eqn:finalsurjection}e'\bigl((v^{-n}S\otimes_R F_0)/(v^c S\otimes_R F_0)\bigr) \to
	M[1/v]/M\end{equation}
	has kernel annihilated by $v^c$,
	and so we find that the kernel of~(\ref{eqn:natural surjection})
	is annihilated by $v^{a+3c} I$. (Indeed, if $x$ is an element
        of the kernel of~(\ref{eqn:natural surjection}), then for any
        $i\in I$ we can lift $v^{a+2c}ix$ to an element of the kernel
        of~(\ref{eqn:finalsurjection}).)

	Putting together the results of the preceding two paragraphs,
	we find that the cokernel of the natural morphism
	$$S\otimes_R Z_n^G \to Z_n$$ is annihilated by $v^{a+2c} I$,
	while its kernel is annihilated by $v^{2a+5 c} I$.
	Recalling that $v\in I$, and noting that 
	the powers of $v$ just mentioned are independent 
	of $n$, and also that we have the inclusions
	$v^{-n} M /M \subseteq Z_n \subseteq v^{-(n+c)}M/M,$
	we see that the proof of the theorem is completed.
\end{proof}


We will now deduce a descent result for projective modules 
over $S[1/v]$.  For this, we need to make some additional hypotheses,
which we now describe.

\begin{hyp}
	\label{hyp:descent}\leavevmode
	\begin{enumerate} 
	\item $S/vS$ is countable.
	\item $S[1/v]$ is Noetherian.
	\item The morphism $R[1/v] \to S[1/v]$ is faithfully flat.
	\end{enumerate}
\end{hyp}

We write $S[1/v] = \varinjlim_n {v^{-n}} S.$
If we identify $v^{-n}S$ with $S$ via multiplication by~
$v^n$, then each of the transition maps becomes identified with
the closed embedding $vS \hookrightarrow S$.
Thus if we equip $S[1/v]$ with the inductive limit topology,
then $S[1/v]$ becomes a topological ring, which is completely
metrizable (since $S$ is $v$-adically complete) and in fact Polish
(since $S/vS$ is countable, so that $S$ and thus $S[1/v]$ is separable).
Note also that $vS$ is an open additive subgroup of $S[1/v]$
which is closed under multiplication, and consists of topologically
nilpotent elements.
Proposition~\ref{prop:module topologies} then
shows that finitely generated $S[1/v]$-modules
have a canonical topology, with respect to which all $S[1/v]$-homomorphisms
are continuous, with closed image. 

We now establish the following descent result in this situation.

%\mar{ME: Rephrase as an equivalence of categories between
%	f.g.\ proj.\ $R[1/v]$-modules,
%	and f.g.\ proj.\ $S[1/v]$-modules equipped with
%	semi-linear $G$-action. ME: Added corollary giving this.}
\begin{theorem}
	\label{thm:descent with coefficients}
	If Hypothesis~{\em \ref{hyp:descent}} holds,
	and if
	$M$ is a finitely generated projective $S[1/v]$-module,
	equipped with a continuous
        {\em (}when endowed with its canonical topology{\em )}
        semi-linear $G$-action,
	then $M^G$ is a finitely generated and projective $R[1/v]$-module,
	and the natural morphism
	\numequation
	\label{eqn:natural morphism}
	S[1/v]\otimes_{R[1/v]} M^G \to M
\end{equation}
       	is an isomorphism.
	%\mar{TG: is $A=R$, $B=S$? ME: Yes, fixed, commenting out}
\end{theorem}

% \begin{remark}
% %	We remark that localisation is exact, so that $M[1/v]^G = M^G[1/v]$.
% The consideration of stably free (rather than simply free) modules
% in the theorem is for purely technical reasons, motivated by 
% the logic of the proof.
% \end{remark}

\begin{proof}% [Proof of Theorem~\ref{thm:descent with
    % coefficients}]
  % \mar{TG: these first two paragraphs were commented
    %out but are then referred to below, so I'm uncommenting them for now.
%   ME: They're back in the proof now, so I'm removing this comment.}
		% Since $M$ is stably free, we may find some $n \geq 0$
		% such that $N:= M \oplus S[1/v]^n$ is free. We extend the semi-linear
		% $G$-action on $M$ to a semi-linear $G$-action on $N$ which
		% preserves its direct product decomposition by
		% having $G$ act semi-linearly on $S[1/v]^n$ via its given action
		% on $S[1/v]$.
		% The natural morphism~(\ref{eqn:natural morphism})
		% for $N$ is then the direct sum of the corresponding
		% morphisms for $M$ and $S[1/v]^n$, and so it suffices to
		% prove the theorem for $N$; i.e.\ we may assume that $M$
		% is in fact free. 

		% Actually, assuming that $M$ is free, we will first show
		% that the natural morphism~(\ref{eqn:natural morphism}) 
		% is surjective.  The argument
		% of the preceding paragraph immediately extends this result
		% to the stably free case.  We will then use this
		% as an ingredient in deducing that the
		% natural morphism is an isomorphism in the free case; the 
		% argument of the preceding paragraph then extends this
		% result to the stably free case. 

	Let $M'$
		denote any finitely generated $S$-submodule of $M$ that generates
		$M$ over $S[1/v]$; then the $S$-span $M''$ of $G M'$ 
		is finitely generated (note that since~$M'$ is finitely
                generated, there is an open subgroup~$H$ of~$G$ such
                that $HM'\subseteq M'$, and~$H$ has finite index in the
                profinite group~$G$), so it is an iso-projective $S$-submodule of $M$ which
		is $G$-invariant. % \mar{TG: is it obviously finitely
                %   generated? Is the point supposed to be that
                %   since~$M'$ is finitely generated, we can find an
                %   open subgroup of~$G$ which takes~$M'$ to itself,
                %   which then has finite index and so the span of $GM'$
                % is finitely generated?}
		Applying Theorem~\ref{thm:almost descent three} to $M''$,
		and noting that $$M^G = (M''[1/v])^G = (M'')^G[1/v]$$
		(because localisation is exact),
		we find
		that the natural morphism~(\ref{eqn:natural morphism})
		has dense image.  Since its target is finitely generated
		over $S[1/v]$, its image is also finitely generated
		(because $S[1/v]$ is Noetherian, by assumption),
		and thus closed in its target;
		combined with the density, we find 
		that~(\ref{eqn:natural morphism}) is surjective.
	

		Since $S[1/v]\otimes_{R[1/v]} M^G \to M$ is surjective,
		while $M$ is finitely generated over $S[1/v]$,
		we see that if $N$ is any sufficiently large finitely generated 
		$R[1/v]$-submodule of $M^G$, then 
		$S[1/v]\otimes_{R[1/v]}N \to M$ is surjective.
		Choose a finitely generated free $R[1/v]$-module $F$
		that surjects onto $N$, and let $E$ denote the kernel
		of the induced surjection
		$$ S[1/v]\otimes_{R[1/v]} F \to S[1/v]\otimes_{R[1/v]} N
		\to M,$$
		so that we have a short exact sequence
		$$0 \to E \to S[1/v]\otimes_{R[1/v]} F \to M \to 0.$$
		Passing to $G$-invariants,
	        and taking into account Corollary~\ref{cor:invariants},
		we obtain a left exact sequence
		$$0 \to E^G \to F \to M^G,$$
		%\mar{ME: This uses that $S[1/v]^G = R[1/v]$! So
		%we should actually prove this first. ME: Now done}
		which (by the choice of $F$)
		induces a short exact sequence
		$$0 \to E^G \to F \to N \to 0.$$
		Tensoring back up with $S[1/v]$ (which is
		flat over $R[1/v]$ by assumption),
		we obtain a morphism of short exact sequences
		$$\xymatrix{ 0 \ar[r] &
			S[1/v]\otimes_{R[1/v]} E^G \ar[r]\ar[d] &
			S[1/v]\otimes_{R[1/v]} F \ar[r]\ar[d] &
			S[1/v]\otimes_{R[1/v]} N \ar[r] \ar[d] & 0 \\
			0 \ar[r] & E \ar[r] & S[1/v]\otimes_{R[1/v]} F \ar[r] &
			M \ar[r] & 0 }
		$$
		Evidently the middle vertical arrow is the identity,
		and thus the natural morphism
		$S[1/v]\otimes_{R[1/v]} E^G \to E$ 
		is injective. Since~$E$ is finitely generated and projective
		(being the kernel of a surjection from a finitely generated
		free module to a projective module), it follows from
                what we have already proved that this morphism is also
                surjective.

                Since the left hand two vertical arrows are
                isomorphisms, so is the third, so that $S[1/v]\otimes_{R[1/v]} N \to M$ is an
                isomorphism.  This is true for any sufficiently large
                choice of $N$, and thus we find (using the faithful
                flatness of $S[1/v]$ over $R[1/v]$, which holds by
                assumption) that all these sufficiently large choices
                of~$N$ coincide, implying that $M^G$ is finitely
                generated and that~(\ref{eqn:natural morphism}) is an
                isomorphism.  The faithful flatness of $S[1/v]$ over
                $R[1/v]$ then implies that $M^G$ is projective over
                $R[1/v]$
                \cite[\href{http://stacks.math.columbia.edu/tag/058S}{Tag
                  058S}]{stacks-project}, as
                required.% \mar{TG: maybe reference stacks 058S ME: We could,
		% although this is also a pretty classical descent 
	        % statement.}
%
		% It remains to be shown that the morphism
		% $S[1/v]\otimes_{R[1/v]} E^G \to E$ 
		% is surjective.  This follows from what we
		% have already proved, taking into account the fact that $M$ 
		% and $S[1/v]\otimes_{R[1/v]} F$ are both free over $S[1/v]$,
		% so that $E$ is stably free over $S[1/v].$
\end{proof}

The following corollary provides a convenient reformulation
of the preceding theorem.

\begin{cor}
	\label{cor:descent with coefficients}
	If Hypothesis~{\em \ref{hyp:descent}} holds,
	then the functor \[M \mapsto S[1/v]\otimes_{R[1/v]} M\]
	induces an equivalence between the category
	of finitely generated projective $R[1/v]$-modules
	and the category of finitely generated projective
	$S[1/v]$-modules endowed with a continuous semi-linear
	action of $G$.  A quasi-inverse functor is given by
	$N \mapsto N^G$.
\end{cor}
\begin{proof}
	Theorem~\ref{thm:descent with coefficients}
	shows that if $N$ is a finitely generated and projective
	$S[1/v]$-module, endowed with a continuous 
	semi-linear $G$-action, then the natural map
	$S[1/v]\otimes_{R[1/v]} N^G \to N$ is an isomorphism.
	To complete the proof of the corollary, then, it suffices
	to show that if $M$ is a finitely generated and projective 
	$R[1/v]$-module, then the natural map
	$M\to (S[1/v]\otimes_{R[1/v]}M)^G$ is an isomorphism.
	Writing $M$ as the direct summand of a finite
	rank free module, we reduce to the case when $M$ is free,
	%and then to the case when $M$ is free of rank one,
	%i.e.\ equal to $R[1/v]$, 
	which (as was already observed in the proof of
	Theorem~\ref{thm:descent with coefficients})
	is established by Corollary~\ref{cor:invariants}.
	%that is, we have to show that
	%the inclusion
	%\numequation
	%\label{eqn:invariant check}
	%R[1/v] \hookrightarrow S[1/v]^G
%\end{equation}
%	is an isomorphism.
%	This follows from the $m = 0$ case of 
%	Lemma~\ref{lem:almost invariants two} (inverting $v$
%	converts the almost isomorphism into a genuine isomorphism).
%	It is also a formal consequence of Theorem~\ref{thm:descent
%	with coefficients}: namely, that theorem shows
%	that the inclusion~(\ref{eqn:invariant check}) becomes
%	an isomorphism after tensoring with $S[1/v]$ over $R[1/v]$;
%	since this is a base-change by a faithfully flat extension
%	(Proposition~\ref{prop:rigid facts}~(2)), it is already an isomorphism.
\end{proof}

% \chapter{Crystalline and semistable moduli stacks}\label{sec:
%   crystalline moduli}
% In this section we explain how to construct closed substacks of $\cX_d$
% \section{The classical theory}
% \label{subsec:classical}
% We suppose now that $R$ is a valuation ring, complete with respect to
% a non-discrete rank one valuation, that $I = \mathfrak m$ 
% is the maximal ideal of $R$, and that $u$ is an element
% of $R$ of positive valuation.

% We let $K$ denote the fraction field of $R$, and let $L$ denote
% the completion of a Galois extension of $K$.  More precisely,
% we suppose that we have an increasing sequence $L_n$ of finite Galois
% extensions of $K$ contained in $L$, with $G_n := \Gal(L_n/K)$,
% such that $L$ is the completion of $\bigcup_n L_n$.  We put
% $G := \varprojlim_n G_n$, and write $H_n := \mathrm{ker}(G \to G_n)$.
% The action of each $G_n$ on $L_n$ induces a
% a continuous action of $G$ on~$L$.

% We let $S_n$ denote the integral closure of $R$ in $L_n$,
% and let $S$ denote the completion of $\bigcup_n S_n$.  Thus $L$
% is the fraction field of $S$.
% Since $R$ is a B\'ezout ring, being flat over $R$ is the same as being
% torsion free, and so each $S_n$ is flat over $R$, and thus even
% faithfully flat (the morphism $R\to S$ being a local morphism of 
% local rings).  Similarly $R\to S$ is faithfully flat.

% Finally, we assume that $R\to S_n$ is almost Galois
% with Galois group $G_n$.  

% \begin{lemma}
% 	\label{lem:almost invariants three}
% 	The inclusion $K \to L^G$ is an isomorphism.
% \end{lemma}
% \begin{proof}
% 	It follows from Lemma~\ref{lem:almost invariants two}
%  that $R \to S^G$ is an almost isomorphism. The
% 	lemma follows by tensoring with $K$ over $R$.
% \end{proof}

% \begin{lemma}
% 	If $V$ is a finite-dimensional $L$-vector space equipped 
% 	with a continuous semi-linear $G$-action, then the natural morphism
% 	$L\otimes_K V^G \to V$ is an isomorphism.
% \end{lemma}
% \begin{proof}
% 	Let $d := \dim_L V$.
% 	We begin by showing that $\dim_K V^G \leq d$.  More precisely,
% 	we show that if $v_1,\ldots,v_n$ is a $K$-linearly independent
% 	sequence of vectors in $V^G$, then these vectors are $L$-linearly
% 	independent when thought of as lying in $V$.
% 	Suppose to the contrary that they become linearly dependent
% 	over $L$; suppose in addition (as we may) that $n$ is chosen
% 	to be as small as possible.  Clearly we must have $n \geq 2,$
% 	since $V^G$ embeds into $V$, and we may choose the linear
% 	dependence relation to be of the form
% 	$$l_1 v_1 + \cdots l_{n_1} v_{n-1} + v_n = 0.$$  
% 	Not all the $l_i$ can lie in $K$, and relabelling if necessary,
% 	we may assume that $l_1 \not\in K$.  By Lemma~\ref{lem:almost invariants
% 		three} we may find $g \in G$ such that $g(l_1) \neq l_1$,
% 	and thus we obtain the non-trivial dependence relation
% 	$$ \bigl(g(l_1) - l_1\bigr) v_1 + \cdots \bigl(g(l_{n-1}) -
% 	l_{n-1} \bigr) v_{n-1} = 0,$$ contradicting our assumption that
% 	$n$ was minimal.
%       	(This is the
% 	``admissibility'' argument going back to Serre and Fontaine.)

% 	Now choose a free $S$-submodule $M$ of $V$
% 	of rank $d$, which is preserved by $G$.  (The existence of
% 	such an $M$ follows by continuity of the $G$-action: if $M'$ 
% 	is any lattice, then --- using the compactness of $G$ ---
% 	we see that the $S$-span of $G M'$ is the required $M$.)
% 	Theorem~\ref{thm:almost descent two} shows that
% 	the $S$-span of $M^G$ is dense in $S$, and thus 
% 	that the $L$-span of $V^G$ is dense in $V$.
% 	Since (as we have seen in the previous paragraph)
% 	$V^G$ is finite-dimensional, its $L$-span is complete,
% 	and so we see that $L$-span of $V^G$ equals $V$.
% 	Thus the natural morphism $L\otimes_K V^G \to V$ is surjective.  
% 	Since the dimension of the source is at most the dimension of
% 	the target (again by the result of the preceding paragraph) 
% 	we see that this morphism is an isomorphism, as required.
% \end{proof}


%\chapter{Descending coefficients}
%\label{sec:coefficients}
%\section{The set-up}
%\section{Descending coefficients}

\section{An application of almost Galois descent}\label{subsec:coefficients}
%\mar{ME: Just begining to write this subsection}
Let $F$ be a closed perfectoid subfield of $\C$, % \mar{TG: recording
  % that we don't have~$\C$ right now, although I intend to rewrite
  % earlier so that we do.}
with tilt~$F^\flat$, a closed perfectoid subfield of $\C^\flat$.  
Recall that, for any $p$-adically complete $\Z_p$-algebra~$A$,
we defined $W(\cO_F^\flat)_A$ and $W(F^\flat)_A$
in Section~\ref{subsec: coefficients}.
%If $A$ is a finite type algebra over $\Z/p^a$
%(for some $a \geq 1$), 
% \mar{ME: Note that $\cO$ hasn't been properly introduced/fixed yet TG:
% since we are working with truncated Witt vectors, and it only come up
% a couple of times, it seemed easiest to just get rid of
% $\cO$. Commenting out.}
%then we write 
%$W(F^\flat)_A := W(F^\flat)\cotimes_{\Z_p}A .$
%(The completion applied to the tensor product is the evident one --- see below
%for a precise definition.)
%Our goal is to prove the following theorem.

\begin{theorem}
	\label{thm:descending projective modules}
			Let $A$ be a finite type $\Z/p^a$-algebra,
			for some $a \geq 1$.
			The inclusion
			$W(F^\flat)_A  \to W(\C^\flat)_A$
			is a faithfully flat morphism
			of Noetherian rings,
			and
	the functor $M \mapsto W(\C^\flat)_A\otimes_{W(F^\flat)_A} M$
	induces an equivalence between the category of finitely generated 
	projective $W(F^\flat)_A$-modules and the category of
	finitely generated projective $W(\C^\flat)_A$-modules endowed
	with a continuous semi-linear $G_F$-action. A quasi-inverse
        functor is given by $N\mapsto N^{G_F}$.
\end{theorem}



The proof will be an application of the almost Galois descent results
of Section~\ref{subsec:profinite}.  Thus we have to place ourselves
in the framework of that section.  To this end,
we note that $\C$ may be regarded as a completion $\widehat{\overline{F}}$ 
of the algebraic closure $\overline{F}$ of $F$.  We then write
$\overline{F} = \bigcup_{n\geq 1} F_n$ as the increasing union of a sequence
of finite Galois extensions of $F$; to match the notation
of Section~\ref{subsec:profinite}, we also write $F_0~:=~F$.
We write $G_n := \Gal(F_n/F),$ $H_n := \Gal(\overline{F}/F_n)$,
and $G := H_0 = \Gal(\overline{F}/F),$ so that $G \iso \varprojlim_n G_n$
and $G/H_n \iso G_n.$
Note that each of the finite Galois extensions $F_n$ of $F$
is again perfectoid. %and then $F_n^\flat$ is finite Galois over $F_m^\flat$,
%for each $n\leq m$,
%with a canonical isomorphism $\Gal(F_n^\flat/F_m^\flat) \cong
%\Gal(F_n/F_m) = H_m/H_n.$
%We now note the following lemma.

%As a consequence of Lemma~\ref{lem:almost etale to almost Galois}, we deduce the following result.
%
%\begin{cor}
%	\label{cor:confirming set-up}
%	For each $m \leq n$ and $a \geq 1$,
%	the inclusion
%	$W_a(\cO_{F_m}^\flat) \hookrightarrow W_a(\cO_{F_n}^\flat)$
%	is almost Galois,
%	%\mar{ME: Have to explain the choice of ideal $I$ w.r.t.\ which
%	%	we define ``almost'' . ME: Done}
%	with Galois group $H_m/H_n$.
%{\em (}The ``almost'' structure being understood with respect
%to the maximal ideal of $W_a(\cO_F^\flat)$.{\em )}
%\end{cor}
%\begin{proof}
%	\mar{ME: Have to give argment --- via passing to the associated
%		graded w.r.t.\ the $p$-adic filtration.  Also have
%		to explain why we can use the maximal ideal of
%		$W_a(\cO_F^\flat)$
%		rather than that of
%		$W_a(\cO_{F_m}^\flat)$ --- which amounts to the fact 
%		that the former is contained in the latter.  And maybe
%		the former generates the latter anyway?}
%\end{proof}


Assume now that $A$ is a finitely generated $\Z/p^a$-algebra.
Then
$W(F^\flat)_A  =  W_a(F^\flat)_A,$
and similarly
$ W(\C^\flat)_A =  W_a(\C^\flat)_A$,
so that from here on we may work with rings of $a$-truncated Witt vectors,
rather than with full rings of Witt vectors themselves.
To accord with the notation of Section~\ref{subsec:profinite},
we also denote these rings by $R$ and $S$ respectively.
Recall from Section~\ref{subsec: coefficients}
that $v$ denotes a non-zero element of the maximal ideal
of $W_a(\cO_F^\flat)$,
and that by definition
$$R := W_a(\cO_F^\flat)_A 
= \varprojlim_i \bigl( W_a(\cO_F^\flat)\otimes_{\Z/p^a}A\bigr)/v^i,$$
while
$$S := W_a(\cO_\C^\flat)_A
= \varprojlim_i \bigl(W_a(\cO_\C^\flat)\otimes_{\Z/p^a}A\bigr)/v^i;$$
so that both $R$ and $S$ are $v$-adically complete.

%Let $v$ denote a non-zero element of the maximal ideal
%of $W_a(\cO_F^\flat)$,
%and set % \mar{TG: this formula was garbled, you should check that this
%%   is what you meant! TG: OK what is below agrees with this, commenting
%% out.}
%$$R :=W_a(\cO_F^\flat) \cotimes_{\Z/p^a}  A
%= \varprojlim_i \bigl( W_a(\cO_F^\flat)\otimes_{\Z/p^a}A\bigr)/v^i$$
%(so that the indicated completion is the $v$-adic completion),
%and also 
%$$S := W_a(\cO_\C^\flat)\cotimes_{\Z/p^a} A
%= \varprojlim_i \bigl(W_a(\cO_\C^\flat)\otimes_{\Z/p^a}A\bigr)/v^i.$$
%Note that $R$ and $S$ are well-defined independent of the choice of $v$
%(any two different choices of $v$ determine the same topology).
%We then define
%$$W(F^\flat)_A :=  W_a(F^\flat) \cotimes_{\Z/p^a} A= R[1/v]$$
%and 
%$$W(\C^\flat)_A :=  W_a(\C^\flat)\cotimes_{\Z/p^a}A =
%S[1/v].$$% \mar{TG: got to here swapping order of factors in tensor
%  % products, stopping for a bit }
%Again these rings are well-defined independently of the choice of $v$,
%and they are precisely the rings appearing in the statement
%of Theorem~\ref{thm:descending projective modules}.

We let $I$ denote the ideal in $R$ generated by the maximal ideal
of $W_a(\cO_F^\flat),$ and let $v$ denote the image in $I$ of the
chosen element $v$ of that maximal ideal (no confusion should result
from this duplication of notation).
We further set
$$S_n := W_a(\cO_{F_n}^{\flat})_A
= \varprojlim_i \bigl(W_a(\cO_{F_n}^\flat)\otimes_{\Z/p^a}A \bigr)/v^i$$
for each $n~\geq~0$
(so in particular $S_0 = R$). 
By construction we have that
$S$ coincides with the $v$-adic completion of $\varinjlim_n S_n$.
As we will see below, the transition morphisms $S_m \to S_n$
are in fact faithfully flat, and thus injective, and so
in fact this direct limit is simply a union;
thus $S = \widehat{\bigcup_n S_n},$
as required for the setup of Section~\ref{subsec:profinite}.
Furthermore,
Lemma~\ref{lem: varphi extends to A and Galois etc}
ensures that the $G$-action on $S$ is continuous.

%Since each $\cO_{F_m}^\flat$ is a B\'ezout ring,
%being flat over $\cO_{F_m}^\flat$ is the same as being
%torsion free, and so each of the inclusions
%$\cO_{F_m}^{\flat} \hookrightarrow \cO_{F_n}^{\flat}$
%(for $m \leq n$)
%is a flat morphism,
%and thus even faithfully flat, being a local morphism of local rings. 
%Similarly the inclusion
%$\cO_F^\flat\hookrightarrow \cO_\C^\flat$
%is faithfully flat.
%A standard grading argument, applying Lemma~\ref{lem:graded flat} to 
%the $p$-adic filtrations on source and target,

%Lemma~\ref{lem:Witt flatness}
%shows that each of the inclusions
%$W_a(\cO_{F_m}^{\flat}) \hookrightarrow W_a(\cO_{F_n}^{\flat})$,
%as well as the inclusion
%$W_a(\cO_F^\flat)\hookrightarrow W_a(\cO_\C^\flat)$,
%is faithfully flat.

\begin{prop}
	\label{prop:confirming set-up}
	For each $i\geq 0,$ and each $m \leq n$,
	the morphism $S_m/v^i \to S_n/v^i$ is faithfully flat, and realizes
	$S_n/v^i$ as an almost Galois extension of $S_m/v^i$, with Galois group
	$H_m/H_n$.
\end{prop}
\begin{proof}
Lemma~\ref{lem:Witt flatness}
shows that each of the inclusions
$W_a(\cO_{F_m}^{\flat}) \hookrightarrow W_a(\cO_{F_n}^{\flat})$ 
%as well as the inclusion
%$W_a(\cO_F^\flat)\hookrightarrow W_a(\cO_\C^\flat)$,
is faithfully flat,
	%Since (as was noted above) the inclusion
	%$W_a(\cO_{F_m}^\flat)\hookrightarrow W_a(\cO_{F_n}^\flat)$
	%is faithfully flat,
	thus so is the morphism
	$W_a(\cO_{F_m}^\flat)\otimes_{\Z/p^a}A
	\to W_a(\cO_{F_n}^\flat)\otimes_{\Z/p^a} A$,\
	and hence so are each
	of the morphisms
	\begin{multline*}
	S_m/v^i S_m =
	\bigl(W_a(\cO_{F_m}^\flat)\otimes_{\Z/p^a}A)\bigr)
	/ v^i
	\bigl(W_a(\cO_{F_m}^\flat)\otimes_{\Z/p^a}A)\bigr)
	\\
	\to
	\bigl(W_a(\cO_{F_n}^\flat)\otimes_{\Z/p^a}A)\bigr)
	/v^i
	\bigl(W_a(\cO_{F_n}^\flat)\otimes_{\Z/p^a}A)\bigr)
	=S_n/v^i S_n .
\end{multline*}
An identical argument, taking into account
Lemma~\ref{lem:almost Galois Witt version}
%Corollary~\ref{cor:confirming set-up}
and Remark~\ref{rem:base-change},
shows that $S_m/v^i S_m \to S_n/v^i S_n$ is almost Galois
(with respect to the ideal in $S_m$ generated by
the maximal ideal of $W_a(\cO_{F_m}^{\flat})$,
and hence also with respect to $I S_m$, since the latter
ideal is contained
in the former),
with Galois group $H_m/H_n$.
\end{proof}


\begin{prop}
	\label{prop:rigid facts}\leavevmode
%	\mar{ME: Will probably just incorporate this into the theorem}
\begin{enumerate}
\item Each $S_m[1/v]$, as well as $S[1/v]$, is Noetherian.
	{\em (}Setting $m = 0$ gives in particular that $R[1/v]$ is
	Noetherian.{\em )}
\item Each of the morphisms $S_m \to S_n$ {\em (}for $m \leq n${\em )},
       as well as each morphism $S_m \to S$, is faithfully flat.
       In particular {\em (}setting $m = 0$ and then inverting~$v${\em )}
       the morphism $R[1/v]\to S[1/v]$ is faithfully flat.
\end{enumerate}
\end{prop}
\begin{proof}
% 	\mar{ME: Need to bootstrap from mod $\varpi$ to mod $\varpi^a$
% 		case. ME: Actually, this is is already done in \S 4.  Maybe
% 	just have to spell out slighty better how those results apply here.
% (Should propably do that above, since we also use them in the graded argument
% for faithful flatness.  Could introduce the notation for the ideal $I$,
% and $v$, at the same time.) ME: Still plan to elaborate on this, although
% $I$ and $v$ are now introduced above. TG: if you're happy with the
% proof of Proposition~\ref{prop: maps of coefficient rings are faithfully flat
%     injections}, we could probably point to that here? At the moment
%   we're writing out the argument twice. ME: Reorganized things a bit.
% There's still some duplication of arguments, but hopefully not to
% an unacceptable level.}
	The Noetherian claims of~(1) follow 
	%from~\cite[Thm.\ 1, \S5.2.6]{MR746961}.
	from~Proposition~\ref{prop: rigid analysis FGK}~(1).
	
	We already noted
	in Proposition~\ref{prop:confirming set-up}
	that the morphisms $S_m/v^i \to S_n/v^i$ are faithfully flat,
	and an identical argument shows that
	each morphism $S_m/v^i \to S/v^i$ is faithfully flat.
	Taking into account that statement of~(1),
	we find that the claims of~(2) follow from Proposition~\ref{prop:
		FGK flat}.
	%
%	Since (as was noted above) the inclusions
%	$W_a(\cO_F^\flat)\hookrightarrow W_a(\cO_\C^\flat)$
%	is faithfully flat,
%	so is the morphism
%	$W_a(\cO_F^\flat)\otimes_{\Z/p^a} A
%	\to  W_a(\cO_\C^\flat)\otimes_{\Z/p^a}A$,\
%	and hence so are each
%	of the morphisms
%	\begin{multline*}
%	R/v^i R =
%	\bigl(W_a(\cO_F^\flat)\otimes_{\Z/p^a} A\bigr)
%	/ v^i
%	\bigl(W_a(\cO_F^\flat)\otimes_{\Z/p^a} A\bigr)
%	\\
%	\to
%	\bigl(W_a(\cO_\C^\flat)\otimes_{\Z/p^a} A\bigr)
%%	/v^i
%	\bigl(W_a(\cO_\C^\flat)\otimes_{\Z/p^a} A\bigr)
%	=S/v^i S .
%\end{multline*}
%	Since~(1) implies that $S$ is Noetherian outside $vS$,
%	in the terminology of~\cite{MR2774689}, 
%	claim~(2) follows from~\cite[Prop.\ 5.2.1]{MR2774689}. 
\end{proof}

%\begin{theorem}
%	\label{thm:descending projective modules}
%			Let $A$ be a finite type $\cO/\varpi^a$-algebra,
%			for some $a \geq 1$.
%			The inclusion
%			$W(L^\flat)_A  \to W(\C^\flat)_A$
%			is a faithfully flat morphism
%			of Noetherian rings,
%			and
%	the functor $M \mapsto W(\C^\flat)_A\otimes_{W(L^\flat)_A} M$
%	induces an equivalence between the category of finitely generated 
%	projective $W(L^\flat)_A$-modules and the category of
%	finitely generated projective $W(\C^\flat)_A$-modules endowed
%	with a continuous semi-linear $G_L$-action.
%	A quasi-inverse functor is given by $N \mapsto N^{G_L}$.
%\end{theorem}
\begin{proof}[Proof of
	Theorem~{\em \ref{thm:descending projective modules}}]
	Proposition~\ref{prop:confirming set-up} verifies
	that the running assumptions imposed at the beginning
	of Section~\ref{subsec:profinite} are satisfied.
	Proposition~\ref{prop:rigid facts} verifies
	that Hypothesis~\ref{hyp:descent} is satisfied,
	and also establishes the Noetherian and faithful
        flatness claims	of
	Theorem~\ref{thm:descending projective modules}.
	The remainder of the theorem then follows from
	Corollary~\ref{cor:descent with coefficients}.
\end{proof}




%\section{A more specific situation}
%\label{subsec:coefficients}
%% \mar{ME: Some of this may end up being moved earlier, to the initial
%% 	discussion of rigid analytic stuff in \S 3.}
%\mar{ME: The theorem and corollary will be removed from this subsection
%	(they are proved in a more abstract setting in the
%	preceding subsection), and then it will be merged with the subsequent
%	subsection, where it will also be generalized to the mod
%	$\varpi^a$ setting.} 
%Let $V$ be a valuation ring, complete with respect to
% a non-discrete rank one valuation, with maximal ideal $\mathfrak m$,
% and let $v$ be an element
%of $V$ of positive valuation (equivalently, $v$ is a non-zero
%element of $\mathfrak m$).
%
% We let $F$ denote the fraction field of $V$, and let $L$ denote
% the completion of a Galois extension of $F$.  More precisely,
% we suppose that we have an increasing sequence
% $$F =: L_0 \subseteq \cdots  \subseteq L_n \subseteq \cdots \subseteq L$$
% of finite Galois
% extensions $L_n$ of $F$ contained in $L$, with $G_n := \Gal(L_n/F)$,
% such that $L$ is the completion of $\bigcup_n L_n$.  We put
% $G := \varprojlim_n G_n$, and write $H_n := \mathrm{ker}(G \to G_n)$.
% The action of each $G_n$ on $L_n$ induces a
% a continuous action of $G$ on~$L$.
%
% We let $V_n$ denote the integral closure of $V$ in $L_n$ (so in particular
% $V_0 = V$),
% and let $V'$ denote the completion of $\bigcup_n V_n$.  Thus $L$
% is the fraction field of $V'$.
% Since $V$ is a B\'ezout ring, being flat over $V$ is the same as being
% torsion free, and so each $V_n$ is flat over $V$, and thus even
% faithfully flat (the morphism $V\to V_n$ being a local morphism of 
% local rings).  Similarly $V\to V'$ is faithfully flat.
%
% Crucially,
% we assume that each of the embeddings $V_m \hookrightarrow
% V_n$ (for $m \leq n$) is almost Galois, with Galois group $H_m/H_n$.
%
%We suppose furthermore that $V$ contains a field $k$,
%and we let $A$ be a finite type $k$-algebra.
%We then let $R$ denote the $v$-adic completion of $V\otimes_k A$,
%let $S$ denote the $v$-adic completion of $V'\otimes_k A$,
%and let $S_n$ denote the $v$-adic completion of $V_n\otimes_k A$
%for each~$n$.
%%We let $F := V[1/v]$ denote the fraction field of $V$,
%%and $L := V'[1/v]$ denote the fraction field of $V'$.\mar{TG: this
%%discussion of notation is potentially quite confusing, I think - is
%%there a better way to do this? I think at the very least we should
%%make ``preceding subsection'' more precise, but at the moment we
%%change $R,S$ to $V,V'$ and then build a new $R,S$ out of them and
%%say that they satisfy the axioms of the original $R,S$, which is
%%pretty obscure! ME: I agree this should be reworded.  One thing is that
%%the ``preceding subsection'' being referred to is now commented out! (It 
%%was the discussion of the ``classical case''.)  So we should bring some
%%of that discussion into this section.  I'll try to do that soon.}
%
%We let $I$ denote the ideal $\mathfrak m R$ in $R$, where 
%$\mathfrak m$ is the maximal ideal of $V$. 
%We are then in the set-up of Subsection~\ref{subsec:profinite}.
%(To check the relevant almost Galoisness, observe that we have it
%for each $V_m \to V_n$, thus for $V_m/v^i \to V_n/v^i$ for any $i$,
%thus for the extension of scalars $S_m/v^i \to S_n/v^i$, which
%is what we need.)\mar{TG: we might also want to say why $S_m/v^i \to
%  S_n/v^i$ is faithfully flat, as that's part of the set-up; currently
%this is (almost) proved in the course of proving the following proposition.}
%% \mar{ME: Actually need it for $S_m/v^i \to S_n/v^i,$ for any $m \leq n,$
%% 	so we need it for $V_m \to V_n$.   Also, I couldn't see
%% 	where we actually made this hypothesis on $V \to V_n,$
%% 	so we need to put it in, I think (and more generally,
%% 	on $V_m\to V_n$; or at least add some framework that
%% 	lets us deduce this). ME: Now sorted out, if I haven't 
%%         gotten myself muddled.}
%
%The following proposition records some standard facts about this situation.
%
%\begin{prop}% \mar{TG: note to self: here and in proper etc, apply FC
%    % fix to indentation/hyperref issues}
%	\label{prop:rigid facts}\leavevmode
%\begin{enumerate}
%\item Each of $R[1/v]$ and $S[1/v]$ is Noetherian.
%\item The morphism $R\to S$, and hence also the morphism
%  $R[1/v]\to S[1/v]$, is faithfully flat.
%\item Every finitely generated $R[1/v]$-module has a canonical
%  topology with respect to which is is complete {\em (}the quotient
%  topology obtained by writing it as a quotient of $S[1/v]^n$ for some
%  $n$; the resulting topology is independent of choices{\em )}.  The
%  analogous statement also holds for finitely generated
%  $S[1/v]$-modules.
%\end{enumerate}
%
%\end{prop}
%\begin{proof}
%	The Noetherian claim of~(1) follows
%	from~\cite[Thm.\ 1, \S5.2.6]{MR746961}.
%	
%	Since $V\to V'$ is faithfully flat, so is the morphism
%	$V\otimes_k A \to V'\otimes_k A$, and hence so are each
%	of the morphisms
%	$$R/v^n R = (V \otimes_k A)/ v^n (V\otimes_k A) 
%	\to (V'\otimes_k A)/v^n (V'\otimes_k A) = S/v^n S.$$
%	Since~(1) implies that $S$ is Noetherian outside $vS$,
%	in the terminology of~\cite{MR2774689}, 
%	claim~(2) follows from~\cite[Prop.\ 5.2.1]{MR2774689}. 
%	
%	Claim~(3) holds by~\cite[Prop.\ 3, \S3.7.3]{MR746961}.  
%\end{proof}
%
%
%%		Recall that $R[1/v] \to S[1/v]$
%%		is faithfully flat~\cite[Prop.\ 5.2.1]{MR2774689}, that $R[1/v]$ and $S[1/v]$ are
%%                Noetherian~\cite[Thm.\ 1, \S5.2.6]{MR746961}
%%		and that any finitely generated $S[1/v]$-module
%%		is complete with respect to its natural topology
%%		(the quotient topology obtained by writing it as a quotient
%%		of $S[1/v]^n$ for some $n$; the resulting topology
%%		is independent of choices) by~\cite[Prop.\ 3, \S3.7.3]{MR746961}.  
%		% These are all standard
%		% facts of rigid analysis over $F$ and~$L$ (\cite[\S6]{MR746961}).
%                % \mar{TG: add
%                % some refs to BGR? ME: Sure! TG: I attempted to do
%                % this, but gave up the will to live. TG: maybe
%                % Noetherian is 5.2.6.1 and complete is \S 3.7.3 and ME
%                % found faithfully flat in a Bosch et al paper?}
%
%% \mar{TG: is this formal? It isn't obvious to me that
%%   a base change of an almost isomorphism is an almost isomorphism, but
%% presumably it's true! TG: ah OK, it follows from Remark 2.1.4(i) in
%% GR03, I guess ME: I think it's pretty formal: if the kernel and cokernel
%% are killed by $I$, then the same is true for all the tensor products
%% and Tors that appear in computing the kernel and cokernel
%% after making a base-change. Or, if you prefer, extension of scalars
%% induces a functor on almost categories, and functors preserve isomorphisms.}
%
%% \mar{ME: The Noetherian, faithful flatness, etc.\ results mentioned 
%% 	at the beginning of the proof of the theorem should ultimately
%% 	be recalled here instead --- where ``here'' might eventually
%% 	be moved to \S 3.}
%
%
%%\section{The descent result}
%%We maintain the notation of the previous subsection.
%
%We now prove our main result on descent in this context.
%
%%\mar{ME: Rephrase as an equivalence of categories between
%%	f.g.\ proj.\ $R[1/v]$-modules,
%%	and f.g.\ proj.\ $S[1/v]$-modules equipped with
%%	semi-linear $G$-action. ME: Added corollary giving this.}
%\begin{theorem}
%	\label{thm:descent with coefficients}
%	If $M$ is a finitely generated projective $S[1/v]$-module,
%	equipped with a continuous semi-linear $G$-action,
%	then $M^G$ is a finitely generated and projective $R[1/v]$-module,
%	and the natural morphism
%	\numequation
%	\label{eqn:natural morphism}
%	S[1/v]\otimes_{R[1/v]} M^G \to M
%\end{equation}
%       	is an isomorphism.
%	%\mar{TG: is $A=R$, $B=S$? ME: Yes, fixed, commenting out}
%\end{theorem}
%
%% \begin{remark}
%% %	We remark that localisation is exact, so that $M[1/v]^G = M^G[1/v]$.
%% The consideration of stably free (rather than simply free) modules
%% in the theorem is for purely technical reasons, motivated by 
%% the logic of the proof.
%% \end{remark}
%
%\begin{proof}% [Proof of Theorem~\ref{thm:descent with
%    % coefficients}]
%  % \mar{TG: these first two paragraphs were commented
%    %out but are then referred to below, so I'm uncommenting them for now.
%%   ME: They're back in the proof now, so I'm removing this comment.}
%		% Since $M$ is stably free, we may find some $n \geq 0$
%		% such that $N:= M \oplus S[1/v]^n$ is free. We extend the semi-linear
%		% $G$-action on $M$ to a semi-linear $G$-action on $N$ which
%		% preserves its direct product decomposition by
%		% having $G$ act semi-linearly on $S[1/v]^n$ via its given action
%		% on $S[1/v]$.
%		% The natural morphism~(\ref{eqn:natural morphism})
%		% for $N$ is then the direct sum of the corresponding
%		% morphisms for $M$ and $S[1/v]^n$, and so it suffices to
%		% prove the theorem for $N$; i.e.\ we may assume that $M$
%		% is in fact free. 
%
%		% Actually, assuming that $M$ is free, we will first show
%		% that the natural morphism~(\ref{eqn:natural morphism}) 
%		% is surjective.  The argument
%		% of the preceding paragraph immediately extends this result
%		% to the stably free case.  We will then use this
%		% as an ingredient in deducing that the
%		% natural morphism is an isomorphism in the free case; the 
%		% argument of the preceding paragraph then extends this
%		% result to the stably free case. 
%
%	Let $M'$
%		denote any finitely generated $S$-submodule of $M$ that generates
%		$M$ over $S[1/v]$; then the $S$-span $M''$ of $G M'$ 
%		is finitely generated (note that since~$M'$ is finitely
%                generated, there is an open subgroup~$H$ of~$G$ such
%                that $HM'\subseteq M'$, and~$H$ has finite index in the
%                profinite group~$G$), so it is an iso-projective $S$-submodule of $M$ which
%		is $G$-invariant. % \mar{TG: is it obviously finitely
%                %   generated? Is the point supposed to be that
%                %   since~$M'$ is finitely generated, we can find an
%                %   open subgroup of~$G$ which takes~$M'$ to itself,
%                %   which then has finite index and so the span of $GM'$
%                % is finitely generated?}
%		Applying Theorem~\ref{thm:almost descent three} to $M''$,
%		and noting that $$M^G = (M''[1/v])^G = (M'')^G[1/v]$$
%		(because localisation is exact),
%		we find
%		that the natural morphism~(\ref{eqn:natural morphism})
%		has dense image.  Since its target is finitely generated
%		over $S[1/v]$, its image is also finitely generated
%		(because $S[1/v]$ is Noetherian, by Proposition~\ref{prop:rigid facts}~(1)).  Thus its image
%		is complete (by Proposition~\ref{prop:rigid facts}~(3)),
%		and combined with the density, we find 
%		that~(\ref{eqn:natural morphism}) is surjective.
%	
%
%		Since $S[1/v]\otimes_{R[1/v]} M^G \to M$ is surjective,
%		while $M$ is finitely generated over $S[1/v]$,
%		we see that if $N$ is any sufficiently large finitely generated 
%		$R[1/v]$-submodule of $M^G$, then 
%		$S[1/v]\otimes_{R[1/v]}N \to M$ is surjective.
%		Choose a finitely generated free $R[1/v]$-module $F$
%		that surjects onto $N$, and let $E$ denote the kernel
%		of the induced surjection
%		$$ S[1/v]\otimes_{R[1/v]} F \to S[1/v]\otimes_{R[1/v]} N
%		\to M,$$
%		so that we have a short exact sequence
%		$$0 \to E \to S[1/v]\otimes_{R[1/v]} F \to M \to 0.$$
%		Passing to $G$-invariants we obtain a left exact sequence
%		$$0 \to E^G \to F \to M^G,$$ which (by the choice of $F$)
%		induces a short exact sequence
%		$$0 \to E^G \to F \to N \to 0.$$
%		Tensoring back up with $S[1/v]$ (which we recall is
%		flat over $R[1/v]$, by Proposition~\ref{prop:rigid facts}~(2)),
%		we obtain a morphism of short exact sequences
%		$$\xymatrix{ 0 \ar[r] &
%			S[1/v]\otimes_{R[1/v]} E^G \ar[r]\ar[d] &
%			S[1/v]\otimes_{R[1/v]} F \ar[r]\ar[d] &
%			S[1/v]\otimes_{R[1/v]} N \ar[r] \ar[d] & 0 \\
%			0 \ar[r] & E \ar[r] & S[1/v]\otimes_{R[1/v]} F \ar[r] &
%			M \ar[r] & 0 }
%		$$
%		Evidently the middle vertical arrow is the identity,
%		and thus the natural morphism
%		$S[1/v]\otimes_{R[1/v]} E^G \to E$ 
%		is injective. Since~$E$ is projective, it follows from
%                what we have already proved that this morphism is also
%                surjective, so that
%		$S[1/v]\otimes_{R[1/v]} N \to M$ is an isomorphism.
%		This is true for any sufficiently large choice of $N$,
%		and thus we find (using the faithful flatness of $S[1/v]$
%		over $R[1/v]$, which holds
%		by Proposition~\ref{prop:rigid facts}~(2))
%		that all these sufficiently large
%                choices of~$N$ coincide, implying that $M^G$ is finitely generated
%		and that~(\ref{eqn:natural morphism}) is an isomorphism.
%		The faithful flatness of $S[1/v]$ over $R[1/v]$ then
%		implies that $M^G$ is projective over
%                $R[1/v]$
%                (\cite[\href{http://stacks.math.columbia.edu/tag/058S}{Tag
%                  058S}]{stacks-project}), as required.% \mar{TG: maybe reference stacks 058S ME: We could,
%		% although this is also a pretty classical descent 
%	        % statement.}
%%
%		% It remains to be shown that the morphism
%		% $S[1/v]\otimes_{R[1/v]} E^G \to E$ 
%		% is surjective.  This follows from what we
%		% have already proved, taking into account the fact that $M$ 
%		% and $S[1/v]\otimes_{R[1/v]} F$ are both free over $S[1/v]$,
%		% so that $E$ is stably free over $S[1/v].$
%\end{proof}
%
%% \begin{theorem}
%% 	\label{thm:descent with coefficients}
%% 	If $M$ is a stably free $S[1/u]$-module of finite rank,
%% 	equipped with a continuous semi-linear $G$-action,
%% 	then $M^G$ is a finitely generated and projective $R[1/u]$-module,
%% 	and the natural morphism
%% 	\numequation
%% 	\label{eqn:natural morphism}
%% 	S[1/u]\otimes_{R[1/u]} M^G \to M
%% \end{equation}
%%        	is an isomorphism.
%% 	%\mar{TG: is $A=R$, $B=S$? ME: Yes, fixed, commenting out}
%% \end{theorem}
%
%% \begin{remark}
%% %	We remark that localisation is exact, so that $M[1/u]^G = M^G[1/u]$.
%% The consideration of stably free (rather than simply free) modules
%% in the theorem is for purely technical reasons, motivated by 
%% the logic of the proof.
%% \end{remark}
%
%% \begin{proof}[Proof of Theorem~\ref{thm:descent with
%%     coefficients}]%\mar{TG: these first two paragraphs were commented
%%     %out but are then referred to below, so I'm uncommenting them for now.
%% %   ME: They're back in the proof now, so I'm removing this comment.}
%% 		Since $M$ is stably free, we may find some $n \geq 0$
%% 		such that $N:= M \oplus S[1/u]^n$ is free. We extend the semi-linear
%% 		$G$-action on $M$ to a semi-linear $G$-action on $N$ which
%% 		preserves its direct product decomposition by
%% 		having $G$ act semi-linearly on $S[1/u]^n$ via its given action
%% 		on $S[1/u]$.
%% 		The natural morphism~(\ref{eqn:natural morphism})
%% 		for $N$ is then the direct sum of the corresponding
%% 		morphisms for $M$ and $S[1/u]^n$, and so it suffices to
%% 		prove the theorem for $N$; i.e.\ we may assume that $M$
%% 		is in fact free. 
%
%% 		Actually, assuming that $M$ is free, we will first show
%% 		that the natural morphism~(\ref{eqn:natural morphism}) 
%% 		is surjective.  The argument
%% 		of the preceding paragraph immediately extends this result
%% 		to the stably free case.  We will then use this
%% 		as an ingredient in deducing that the
%% 		natural morphism is an isomorphism in the free case; the 
%% 		argument of the preceding paragraph then extends this
%% 		result to the stably free case. 
%
%% 		The key facts we use are that $R[1/u] \to S[1/u]$
%% 		is faithfully flat, that $R[1/u]$ and $S[1/u]$ are
%%                 Noetherian 
%% 		and that any finitely generated $S[1/u]$-module
%% 		is complete with respect to its natural topology (the 
%% 		quotient topology obtained by writing it as a quotient
%% 		of $S[1/u]^n$ for some $n$; the resulting topology
%% 		is independent of choices).   These are all standard
%% 		facts of rigid analysis over $K$ and~$L$ (\cite[\S6]{MR746961}).\mar{TG: add
%%                 some refs to BGR? ME: Sure! TG: I attempted to do
%%                 this, but gave up the will to live. TG: maybe
%%                 Noetherian is 5.2.6.1 and complete is \S 3.7.3 and ME
%%                 found faithfully flat in a Bosch et al paper?}
%
%% 		Suppose then that $M$ is free.  If we let $M'$
%% 		denote any free $S$-submodule of $M$ that generates
%% 		$M$ over $S[1/u]$, then the $S$-span $M''$ of $G M'$ 
%% 		is an iso-free $S$-submodule of $M$ which
%% 		is $G$-invariant.
%% 		Applying Theorem~\ref{thm:almost descent two} to $M''$,
%% 		and noting that $$M^G = (M''[1/u])^G = (M'')^G[1/u]$$
%% 		(because localisation is exact),
%% 		we find
%% 		that the natural morphism~(\ref{eqn:natural morphism})
%% 		has dense image.  Since its target is finitely generated
%% 		over $S[1/u]$, its image is also finitely generated
%% 		(because $S[1/u]$ is Noetherian).  Thus its image
%% 		is complete, and combined with the density, we find 
%% 		that~(\ref{eqn:natural morphism}) is surjective.
%% 		As already noted, having proved our result in the case
%% 		where $M$ is free, it immediately extends to the case
%% 		where $M$ is stably free.
%
%% 		Since $S[1/u]\otimes_{R[1/u]} M^G \to M$ is surjective,
%% 		while $M$ is finitely generated over $S[1/u]$,
%% 		we see that if $N$ is any sufficiently large finitely generated 
%% 		$R[1/u]$-submodule of $M^G$, then 
%% 		$S[1/u]\otimes_{R[1/u]}N \to M$ is surjective.
%% 		Choose a finitely generated free $R[1/u]$-module $F$
%% 		that surjects onto $N$, and let $E$ denote the kernel
%% 		of the induced surjection
%% 		$$ S[1/u]\otimes_{R[1/u]} F \to S[1/u]\otimes_{R[1/u]} N
%% 		\to M,$$
%% 		so that we have a short exact sequence
%% 		$$0 \to E \to S[1/u]\otimes_{R[1/u]} F \to M \to 0.$$
%% 		Passing to $G$-invariants we obtain a left exact sequence
%% 		$$0 \to E^G \to F \to M^G,$$ which (by the choice of $F$)
%% 		induces a short exact sequence
%% 		$$0 \to E^G \to F \to N \to 0.$$
%% 		Tensoring back up with $S[1/u]$ (which we recall is
%% 		flat over $R[1/u]$),
%% 		we obtain a morphism of short exact sequences
%% 		$$\xymatrix{ 0 \ar[r] &
%% 			S[1/u]\otimes_{R[1/u]} E^G \ar[r]\ar[d] &
%% 			S[1/u]\otimes_{R[1/u]} F \ar[r]\ar[d] &
%% 			S[1/u]\otimes_{R[1/u]} N \ar[r] \ar[d] & 0 \\
%% 			0 \ar[r] & E \ar[r] & S[1/u]\otimes_{R[1/u]} F \ar[r] &
%% 			M \ar[r] & 0 }
%% 		$$
%% 		Evidently the middle vertical arrow is the identity,
%% 		and thus the natural morphism
%% 		$S[1/u]\otimes_{R[1/u]} E^G \to E$ 
%% 		is injective.  
%% 		We claim that it is also surjective.
%% 		Admitting this, we find that 
%% 		$S[1/u]\otimes_{R[1/u]} N \to M$ is an isomorphism.
%% 		This is true for any sufficiently large choice of $N$,
%% 		and thus we find (using the faithful flatness of $S[1/u]$
%% 		over $R[1/u]$) that these all sufficiently large
%%                 choices of~$N$ coincide, implying that $M^G$ is finitely generated
%% 		and that~(\ref{eqn:natural morphism}) is an isomorphism.
%% 		The faithful flatness of $S[1/u]$ over $R[1/u]$ then
%% 		implies that $M^G$ is projective over
%%                 $R[1/u]$ (\cite[\href{http://stacks.math.columbia.edu/tag/058S}{Tag 058S}]{stacks-project}).% \mar{TG: maybe reference stacks 058S ME: We could,
%% 		% although this is also a pretty classical descent 
%% 	        % statement.}
%
%% 		It remains to be shown that the morphism
%% 		$S[1/u]\otimes_{R[1/u]} E^G \to E$ 
%% 		is surjective.  This follows from what we
%% 		have already proved, taking into account the fact that $M$ 
%% 		and $S[1/u]\otimes_{R[1/u]} F$ are both free over $S[1/u]$,
%% 		so that $E$ is stably free over $S[1/u].$
%% \end{proof}
%
%The following corollary provides a convenient reformulation
%of the preceding theorem.
%
%\begin{cor}
%	\label{cor:descent with coefficients}
%	The functor $M \mapsto S[1/v]\otimes_{R[1/v]} M$
%	induces an equivalence between the category
%	of finitely generated projective $R[1/v]$-modules
%	and the category of finitely generated projective
%	$S[1/v]$-modules endowed with a continuous semi-linear
%	action of $G$.  A quasi-inverse functor is given by
%	$N \mapsto N^G$.
%\end{cor}
%\begin{proof}
%	Theorem~\ref{thm:descent with coefficients}
%	shows that if $N$ is a finitely generated and projective
%	$S[1/v]$-module, endowed with a continuous 
%	semi-linear $G$-action, then the natural
%        transformation\mar{TG: here and below it seems a bit weird to
%          fix~$N$ and then say ``natural transformation''? Should we
%          either say ``natural map'' or explicitly allow $N$ to vary?}
%	$S[1/v]\otimes_{R[1/v]} N^G \to N$ is an isomorphism.
%	To complete the proof of the corollary, then, it suffices
%	to show that if $M$ is a finitely generated and projective 
%	$R[1/v]$-module, then the natural transformation
%	$M\to (S[1/v]\otimes_{R[1/v]}M)^G$ is an isomorphism.
%	Writing $M$ as a direct summand of a finite
%	rank free module, we reduce to the case when $M$ is free,
%	and then to the case when $M$ is free of rank one,
%	i.e.\ equal to $R[1/v]$; that is, we have to show that
%	the inclusion
%	\numequation
%	\label{eqn:invariant check}
%	R[1/v] \hookrightarrow S[1/v]^G
%\end{equation}
%	is an isomorphism.
%	This follows from the $m = 0$ case of 
%	Lemma~\ref{lem:almost invariants two} (as inverting $v$
%	converts the almost isomorphism into a genuine isomorphism).
%	It is also a formal consequence of Theorem~\ref{thm:descent
%	with coefficients}: namely, that theorem shows
%	that the inclusion~(\ref{eqn:invariant check}) becomes
%	an isomorphism after tensoring with $S[1/v]$ over $R[1/v]$;
%	since this is a base-change by a faithfully flat extension
%	(by Proposition~\ref{prop:rigid facts}~(2)), it is already an isomorphism.
%\end{proof}

% \chapter{Crystalline and semistable moduli stacks}\label{sec:
%   crystalline moduli}
% In this section we explain how to construct closed substacks of $\cX_d$
% parameterizing lattices in (potentially) crystalline, or (potentially)
% semistable representations.  The construction is easiest in the 
% crystalline case when $K$ is absoutely unramified, since in this case 
% we describe these moduli stacks directly in terms of moduli stacks
% of Wach modules.  The semistable case, or the case when $K$ is ramified
% over $\Q_p$ (and hence also the potentially crystalline and
% potentially semistable cases) are more involved.



% \section{Crystalline moduli stacks when $K$ is absolutely
%   unramified}\label{subsec: crystalline moduli unramified base}

% We recall the following theorem,
% which was conjectured by Fontaine \cite{MR1106901},
% and one direction of which was announced by him in~\cite{MR1106901}
% and proved by Wach~\cite{MR1415732},
% the other direction then being proved subsequently by
% Colmez \cite{MR1711279}.

% \begin{thm}
% 	Suppose that $K$ is unramified over $\Q_p$.
% 	If $\rho: G_K \to \GL_d(E)$ is a continuous representation
% 	of $G_K$ with coefficients in a finite extension $E$ of $\Q_p$,
% 	then the associated $(\varphi,\Gamma)$-module $D(\rho)$
% 	contains a Wach module of height $\leq h$ and level $0$
% 	if and only if $\rho$ is crystalline with labeled Hodge--Tate weights
% 	lying in the interval $[0,h]$.
% \end{thm}
% \begin{proof}
% One direction of this result is stated in~\cite[(2.3.6)]{MR1106901},
% and proved in~\cite{MR1415732}.  (See Remark~(a) on p.~399.)
% As already remarked on above,
% the converse is proved in~\cite{MR1711279}.
% \end{proof}

% \begin{thm}
% 	Suppose that $K$ is unramified over $\Q_p$, 
% 	and fix a tuple $\underline{\lambda}$ of $d$-dimensional
% 	labeled Hodge--Tate weights.
% 	There is a closed substack $\cX_d^{\crys,\underline{\lambda}}$ of $\cX_d$
% 	which is characterized as the unique locally Noetherian, residually
% 	Jacobson, and analytically reduced closed substack of $\cX_d$
% 	which is flat over $\Spf \Z_p$, and whose groupoid of $\Zbar_p$-points,
% 	when regarded as a subgroupoid of $\cX_d(\Zbar_p)$, 
% 	coincides with the groupoid of
% 	$G_K$-invariant $\Zbar_p$-lattices in
% 	crystalline representations of $G_K$ with labeled Hodge--Tate
% 	weights equal to $\underline{\lambda}$.
% 	Furthermore, $\cX_d^{\crys,\underline{\lambda}}$ is a $p$-adic formal
% 	algebraic stack.
% \end{thm}
% \begin{proof}
% 	We have seen above that $\cX_{d,h,0}^a$ is a closed algebraic 
% 	substack of $\cX_d^a$, and thus $\cX_{d,h,0} := \varinjlim_a 
% 	\cX_{d,h,0}^a$ is a closed substack of $\cX_d$. In fact,
% 	it is a $p$-adic formal algebraic stack of finite type.
% 	\mar{ME: Also need that it's locally Noeth.; in fact it's probably
% 		better than this, i.e.\ loc.\ f.t.\, or maybe even f.t.\, 
% 		over $\Spf \Z_p$.  Sort this out.}
% 	To see this, note
%         that Proposition~\ref{prop:inductive description of W}~(1) shows that
% 	$\cW_{d,h,0}$ is a $p$-adic formal algebraic stack,
% 	while the morphism
% 	$\cW_{d,h,0} \to \cX_{d,h,0}$ is scheme-theoretically dominant,
% 	in the sense of \cite[Def.~2.12]{Emertonformalstacks},
% 	as well as being representable by algebraic spaces and proper
% 	(by Lemma~% \ref{lem:C to R is proper}
%         ).\mar{TG: fill in}
% 	It follows from \cite[Prop.~7.5]{Emertonformalstacks}
% 	that $\cX_{d,h,0}$ is indeed
% 	a $p$-adic formal algebraic stack of finite type.
% 	\mar{ME: Explain why hypotheses of the cited prop.\ hold.}

% 	We now pass to the underlying reduced substack $\cY$ of $\cX_{d,h}$,
% 	\mar{ME: Here I don't mean the underlying reduced algebraic stack
% 		of this formal algebraic stack,
% 		but rather the thing you get by passing from
% 		$\Spf A$ to $\Spf A_{\red}$ on the members of a smooth cover
% 		of $\cX_{d,h}$ by affine formal algebraic spaces.  I think
% 	this will have to be treated in {\tt formal-stacks.tex},
% 	along with flat parts; excellence will play an important role, 
%         since I want the underlying reduced to be actually analytically
%         reduced.}
%         and then pass to the $\Z_p$-flat part $\cZ$ of $\cY$.  Since
% 	$\cZ$ is a closed substack of $\cX_{d,h}$ by construction,
% 	it is a closed substack of $\cX_d$, and is also a $p$-adic
% 	formal algebraic stack.  Again by construction, it is locally
% 	Noetherian, residually Jacobson, and analytically unramified.

% 	\mar{ME: Complete proof, by fleshing out what's here, explain how
% 		to break up according to different $\underline{\lambda}$, 
% 		and prove claimed characterization.}

% \end{proof}



% Let 
% The output of Section~\ref{subsec:coefficients}, when we apply it, will be a
% finitely generated projective
% $(W_a(k)\otimes_{\Zp}A)((u^{1/p^{\infty}}))$-module $M$ equipped with an \'etale $\varphi$-module
% structure; that is, the structure of a $\varphi$-semi-linear map
% $\varphi_M:M\to M$ which induces an isomorphism of
% $(W_a(k)\otimes_{\Zp}A)((u^{1/p^{\infty}}))$-modules $\varphi^*M\to M$.
% % \mar{ME: There should be some extra $W_a(k)$-coefficients in here somewhere,
% % 	but I'm ignoring those out of laziness.}

% If $B$ is a finite type $A$-algebra, then we write
% $$M_B := 
% (W_a(k)\otimes_{\Zp}B)((u^{1/p^{\infty}}))
% \otimes_{(W_a(k)\otimes_{\Zp}A)((u^{1/p^{\infty}}))} M.$$

% Let $M$ be a perfect \'etale $\varphi$-module with
% $A$-coefficients. If $B$ is a finite type $A$-module, then we write
% $M_B:=\widetilde{\cO}_{\mathcal{E},B}\otimes_{\tA_A}M$, which is
% naturally a perfect \'etale $\varphi$-module with $B$-coefficients.

% \begin{prop}\label{prop: representability of phi-invariants} The functor on finite type $A$-algebras defined by
% 	$B \mapsto M_B^{\varphi = 1}$
% 	is representable by a finite type affine $A$-scheme.
% \end{prop}% \mar{TG: use the same argument to prove that if $M$ is an 
% % \'etale $\varphi$-module and $M'=M\otimes A(u^{1/p^\infty})$, then
% % $M^{\varphi=1}\to (M')^{\varphi=1}$ is an isomorphism.}
% \begin{proof}
% 	We suppose first that $M$ is free.  Then we may choose a
% 	free $(W_a(k)\otimes_{\Zp}A)[[u^{1/p^{\infty}}]]$-sublattice $\gM$ of $M$, and
%         after rescaling it by a power of~$u$, we may assume that~$\gM$
%         is~$\varphi$-stable. Choose an integer~$h\ge 1$ such that the
%         cokernel of $1\otimes\varphi:\varphi^*\gM\to\gM$ is killed
%         by~$u^h$.

% 	% and using its finite height and the fact that $\varphi$
% 	% takes $u^{1/p^{a+1}}$ to $u^{1/p^a}$, we see
% 	% that an element of $M^{\varphi}_B$ depends only only finitely many
% 	% of its coefficients, in a manner independent of $B$,
% 	% cutting out the desired scheme.
%  %        \mar{ME: Obviously more details 
% % 		are needed in this first para.,
% % 		but I think the overall strategy is correct. TG: added
% %               a sketch in the mod~$p$ case, hopefully not
% %               nonsense. TG: actually it seems OK mod $p^a$ as written,
% %             although I could well be blundering\dots}
% % \mar{TG: $p$ should be $q$ and all exponents should change for
% %   $u$-height, probably.}
%               % Here is more of a sketch: suppose that we are working
%               % mod~$p$. Scale~$\gM$ so that it is~$\varphi$-stable and
%               % has height at most~$h$.
%               We claim that for any~$\epsilon>0$, % \mar{TG: we might not
%               % need this~$\epsilon$, but it's harmless anyway}
%             we have $M_B^{\varphi=1}\subset
%               u^{-h/(p-1)-\epsilon}\gM_B$. Indeed, suppose that $m\in
%               M_B^{\varphi=1}$, and choose~$\alpha$ sufficiently large
%               such that $u^\alpha m\in \gM$. Then we may write
%               $\varphi(u^{(h+\alpha)/p}m)=u^{h+\alpha}m\in u^h\gM$, so
%               our assumption on~$h$ implies
%               that~$u^{(h+\alpha)/p}m\in\gM$. Iterating gives the
%               claim.

%               Consider for any $\epsilon>0$ the induced map \numequation\label{eqn: truncation with
%               unbounded denominators}M_B^{\varphi=1}\to
%               (u^{-h/(p-1)-\epsilon}\gM_B/u^\epsilon\gM_B)^{\varphi=1},\end{equation}
%             where the right hand side is a slightly abusive shorthand for $
%             u^{-h/(p-1)-\epsilon}\gM_B/u^\epsilon\gM_B\cap
%             (M/u^\epsilon\gM_B)^{\varphi=1}$. We claim that this is
%            % \mar{TG:
%            %  probably not great notation, as the RHS isn't $\varphi$-stable.}
%           is bijective for
%             any~$\epsilon>0$. Indeed, any element~$m$ of the kernel
%             satisfies~$m=\varphi(m)$ and $m\in u^\epsilon\gM$,
%             whence~$m=\varphi^n(m)\in u^{p^n\epsilon}\gM$ for
%             all~$n\ge 1$, so that~$m=0$; 
%             while given an
%             element~$x$ of right hand side, if $\tx$ is an arbitrary lift of it
%             to~$u^{-h/(p-1)-\epsilon}\gM_B$, then
%             $\tx+\sum_{i=0}^\infty\varphi^i(\varphi(\tx)-\tx)$ is the unique
%             $\varphi$-stable lift of~$x$.

%             Choose a basis $m_i$ for~$\gM$ and let~$X$ be
%             the matrix of~$\varphi$ in that basis.            Fix
%             $N>ph/(p-1)$ and choose~$e$ large enough that the entries of~$X$ are all contained
%             in~$(W_a(k)\otimes_{\Zp}B)[[u^{1/p^e}]]+u^N (W_a(k)\otimes_{\Zp}B)[[u^{1/p^\infty}]]$ (this will be
%             the case for some~$e$). Suppose that~$m\in
%             M_B^{\varphi=1}$; by the above, we may write $m=\sum \lambda_i m_i$,
%             with each~$\lambda_i\in
%             u^{-N/p}(W_a(k)\otimes_{\Zp}B)[[u^{1/p^\infty}]]$. Choose~$r$ sufficiently
%             large that each~$\lambda_i\in
%             u^{-N/p}(W_a(k)\otimes_{\Zp}B)[[u^{1/p^r}]]+(W_a(k)\otimes_{\Zp}B)[[u^{1/p^\infty}]]$. Then we
%             have~$\varphi(\lambda_i)\in
%             u^{-N}(W_a(k)\otimes_{\Zp}B)[[u^{1/p^{r-1}}]]+(W_a(k)\otimes_{\Zp}B)[[u^{1/p^\infty}]] $, so that
%             the equation $m=\varphi(m)$ and the choice of~$e$
%             guarantee that in fact each~$\lambda_i$ is contained in
%             $u^{-N/p}(W_a(k)\otimes_{\Zp}B)[[u^{1/p^{\min(r-1,e)}}]]+(W_a(k)\otimes_{\Zp}B)[[u^{1/p^\infty}]]$. 

% Iterating, we conclude that the~$\lambda_i$ are all contained in
% $u^{-N/p}(W_a(k)\otimes_{\Zp}B)[[u^{1/p^{e}}]]+(W_a(k)\otimes_{\Zp}B)[[u^{1/p^\infty}]]$, which is independent
% of the choice of~$m$. It follows from the bijection~(\ref{eqn: truncation with
%               unbounded denominators}) that the elements of $M_B^{\varphi=1}$ are
%             determined by the coefficients of the (finitely many)
%             powers of~$u^{-1/p^e}$ in the~$\lambda_i$, and writing out
%             the condition on these coefficients that~$\varphi(m)=m$ modulo~$u^\epsilon\gM$,
%             the result follows.

% % Then I \emph{think} that the
% %             equation $\varphi(m)=m$ implies that the~$\lambda_i$ are
% %             contained in $B((u^{1/p^e}))+
% %             B[[u^{1/p^\infty}]]$.
% 	   %  \mar{ME: I'm still trying to work through this.  My first point
% % 		    of confusion: since $\mathfrak M$ is $\phi$-stable,
% % 		    don't we know that entries of $X$ are integral,
% % 		    so maybe they are in
% % 		    $B[[u^{1/p^e}]] + u^N B[[u^{1/p^{\infty}}]],$ 
% % 			    rather than just
% % 		    $B((u^{1/p^e})) + u^N B[[u^{1/p^{\infty}}]]?$
% % 	   I don't know if that makes any difference to the rest
% %   of the argument, though.  Also, do you think maybe one should 
% % ask that $N$ be greater than $(ph)/(p-1)$, because when you
% % write out $\varphi(m) = m $ in coordinates, the entries of $X$
% % are multiplies against coefficients of the form $\varphi(\lambda_i)$,
% % rather than just $\lambda_i$. TG: the latter claim sounds quite
% % plausible, I'll check; and the former claim looks true too! I'll check
% % this more carefully in a bit. Sorry it's confusing (and presumably
% % possibly wrong). ME: Actually, I think I understand the argument now, and agree
% % with it! (Up to modifying the choice of $N$ as per my suggestion.) TG:
% % hopefully written out properly now.}
% 	    % Combined with the above bijection
%             % we see that the~$\lambda_i$ are determined by the finitely
%             % many coefficients of the non-negative powers of
%             % $u^{1/p^e}$ which are greater than $-h/(p-1)$, as
%             % required. 
              
% 	In the remainder of the proof, we explain how to reduce the
% 	case when $M$ is finitely generated projective to the case
% 	when $M$ is finite rank free. By the analogue of
%         Lemma~\ref{lem: projective phi module is summand of free} for
%         perfect \'etale $\varphi$-modules (which follows from an
%         identical argument), we may 
% 	write $M$ as a direct summand
% 	% (in the category of \'etale $\varphi$-modules over
% 	% $A((u^{1/p^{\infty}}))$) of a finite rank
% 	of a free  perfect \'etale $\varphi$-module $F$.   Then 
% 	$M^{\varphi = 1} = M \cap F^{\varphi = 1}$.  Since $M$ is a direct
% 	summand of $F$, the condition that an element of $F^{\varphi =1}$
% 	lies in $M$ is a closed condition, and so $M^{\varphi =1}$ 
% 	is described by a closed subscheme of $F^{\varphi =
%           1}$. (Indeed, this follows easily from the explicit
%         description of $F^{\varphi =
%           1}$ given above. More conceptually, we can argue as follows:
%         we can choose an embedding $F/M \subseteq
% E$ with $E$ free, and so we have an exact sequence\[0 \rightarrow M
% \rightarrow F \rightarrow E,\] inducing an exact sequence \[0 \rightarrow M^{\varphi = 1}
% 	\rightarrow F^{\varphi = 1} \rightarrow E^{\varphi  = 1}.\]
%       This construction is compatible with arbitrary change of coefficients,
% so that $M^{\varphi = 1}$ is the fibre (over the closed point $0 
% \in E^{\varphi = 1}$) of the morphism of finite type affine
% schemes $F^{\varphi = 1} \to E^{\varphi = 1}$, and hence is an
% affine scheme.)% \mar{TG: this seems reasonable, but I don't see how to
% %           make it rigorous - unless you mean to do so by reference to
% %           the ``explicit'' construction of $F^{\varphi=1}$ in the
% %           first paragraph? ME: I think I had the explicit construction in 
% %   mind, although perhaps one can instead make the following more
% %   conceptual argument:  we have
% % $M \subseteq F$, and we can similarly choose an embedding $F/M \subseteq
% % E$ with $E$ free again, and so we have a s.e.s.\ $0 \rightarrow M
% % \rightarrow F \rightarrow E$, inducing $0 \rightarrow M^{\varphi = 1}
% % 	\rightarrow F^{\varphi = 1} \rightarrow E^{\varphi  = 1}$
% % 	(and all this compatible with arbitrary change of coefficients),
% % so that $M^{\varphi = 1}$ is the fibre (over the closed point $0 
% % \in E^{\varphi = 1}$) of the morphism of finite type affine
% % schemes $F^{\varphi = 1} \to E^{\varphi = 1}$, and hence is an
% % affine scheme.}
% \end{proof}

%\section{An application}
%We now return to the setting of Section~\ref{subsec:
%phi gamma over Zp}, namely we fix
%a finite extension $K$ of $\Q_p$, as well as uniformizer $\pi$ of $K$,
%we let $K_{\cyc}$ denote the cyclotomic $\Z_p$-extension of $K$,
%and we write $K_{\infty} := K(\pi^{1/p^{\infty}}).$
%We also fix another finite extension $E$ of $\Q_p$,
%let $\cO$ denote the ring of integers of $E$, let $\varpi$ be a uniformizer
%of $E$,
%and fix an algebra $A$ of finite type over $\cO/\varpi^a$, for some $a \geq 1$.
%
%\mar{ME: Introduce, or recall, the notation $\Atilde_K$ and
%  $\OEAtilde$. TG: should be a matter of referring to Section~\ref{subsec: perfect
%  etale phi modules}.}
%We will apply the almost Galois descent theory of this section
%to relate finitely generated projective modules over $\Atilde_{K,A}$
%and $\OEAtilde$ to finitely generated projective modules with
%appropriate descent data over $W(\C^\flat)_A$.  
%In order to see that this theory is applicable, we rely on various
%results from the theory of perfectoid fields, and in particular 
%on the fact that 
%an a finite extension of such fields induces an {\em almost \'etale}
%extension of the corresponding rings of integers.  
%This result goes back to the intial work of Tate on $p$-adic Hodge
%theory, and has been generalized in subsequent work of Faltings,
%Gabber--Romero, and Scholze (among other authors).  We have
%found the paper \cite{ScholzePerfectoid} of Scholze to be a
%convenient reference for the results that we need,
%and it is this paper that we cite in the arguments that follow.
%
%\begin{theorem}
%	\label{thm:descending projective modules}
%			Let $A$ be a finite type $\cO/\varpi^a$-algebra,
%			for some $a \geq 1$.
%	\begin{enumerate}
%		\item
%			The inclusion
%			$\Atilde_{K,A} \to W(\C^\flat)_A$
%			is a faithfully flat morphism
%			of Noetherian rings,
%			and
%	the functor $M \mapsto W(\C^\flat)_A\otimes_{\Atilde_{K,A}} M$
%	induces an equivalence between the category of finitely generated 
%	projective $\Atilde_{K,A}$-modules and the category of
%	finitely generated projective $W(\C^\flat)_A$-modules endowed
%	with a continuous semi-linear $G_{K_{\cyc}}$-action.
%	A quasi-inverse functor is given by $N \mapsto N^{G_{K_{\cyc}}}$.
%	\item
%			The inclusion
%			$\OEAtilde \to W(\C^\flat)_A$
%			is a faithfully flat morphism
%			of Noetherian rings,
%			and
%	the functor $M \mapsto W(\C^\flat)_A\otimes_{\OEAtilde} M$
%	induces an equivalence between the category of finitely generated 
%	projective $\OEAtilde$-modules and the category of
%	finitely generated projective $W(\C^\flat)_A$-modules endowed
%	with a continuous semi-linear $G_{K_{\infty}}$-action.
%	A quasi-inverse functor is given by $N \mapsto N^{G_{K_{\infty}}}$.
%\end{enumerate}
%\end{theorem}
%\begin{proof}
%To begin with, suppose that $A$ is a finite type $\F := \O/\varpi$-algebra,
%so that $\Atilde_{K,A} := A\cotimes_{\F_p}
%\F_p((T^{1/p^{\infty}}))$\mar{TG: I guess this is
%  $k_\infty((T^{1/p^\infty}))$? But since we haven't used that
%  notation anywhere else in the paper, is there a reason not to just
%  call it $\tE_{K}$, and indeed doing something more like that would
%  let us write (1) and (2) simultaneously?}
%and
%$W(\C^\flat)_A = A\cotimes_{\F_p} \C^\flat.$
%Recall that there is an inclusion 
%$\F_p((T^{1/p^{\infty}})) \subset \C^\flat$
%which identifies the target with the completion of the algebraic closure
%of the source.\mar{TG: I don't think we have this stated anywhere
%  above yet but it should just be a matter of citing Fontaine. I guess
%we should make sure we say it for the Breuil--Kisin extension too.}
%%We consider as well the intermediate field $\F_p((T^{1/p^{\infty}}))$,
%%and the corresponding completed tensor product
%%	$\Atilde_{K,A} := A\cotimes_{\F_p} \F_p((T^{1/p^{\infty}})).$
%%Note that
%%$\F_p((T^{1/p^{\infty}}))$
%%is a perfectoid field, and that
%%$\C^\flat$ is the completion of a Galois extension of
%%$\F_p((T^{1/p^{\infty}}))$.  
%
%In order to apply the results of Section~\ref{subsec:coefficients},
%we introduce some more precise notation.
%Namely, we let $L$ denote the algebraic closure of
%$\F_p((T^{1/p^{\infty}}))$ in $\C^\flat$; then $L$ is
%an algebraic closure of 
%$\F_p((T^{1/p^{\infty}}))$.  \mar{TG: not quite sure what the point of
%the second half of the sentence is.}
%In fact $L$ is also
%a separable closure of
%$\F_p((T^{1/p^{\infty}}))$, since the latter field is perfect,\mar{TG:
%isn't it obviously perfect as it's defined to be the completion of the
%perfection?}
%being perfectoid. (This is proved by a successive approximation 
%argument, see e.g.\ \cite[Prop.~5.9]{ScholzePerfectoid}.)
%Of course $\C^\flat$ is then equal to the completion of $L$.
%
%Write (as we may) $L = \bigcup_n L_n$ as an increasing union of finite
%Galois extensions of 
%$\F_p((T^{1/p^{\infty}}))$,
%and write $V_n$ to denote the ring of integers in $L_n$.
%Each of the finite extensions $L_n$ is again perfectoid
%(see e.g.\ \cite[Prop.~5.23]{ScholzePerfectoid}), and so each 
%of the extensions $V_m \subseteq V_n$ (for $m \leq n$) is 
%%almost \'etale
%%(see \cite[Prop.~5.23]{ScholzePerfectoid} for one formulation of this statement,
%%in the language of almost mathematics),
%%and thus
%almost Galois, with Galois group $\Gal(L_n/L_m)$, by Lemma~\ref{lem:almost
%	etale to almost Galois} below.
%%\mar{ME: Have to convert from the language of ``almost \'etale'' in 
%%	the Faltings--Scholze sense to ``almost Galois'' in our 
%%	naive sense. 
%%	Also should probaby explain that we are just using \cite{ScholzePerfectoid} 
%%	as a convenient reference, but that these results go back to
%%	Tate, Faltings, Gabber--Ramero, etc.  ME: Added lemma, as
%%	well as a comment above.}
%Also, write $F := \F_p((T^{1/p^{\infty}}))$,  
%write
%$V := \F_p[[T^{1/p^{\infty}}]]$ (the ring of integers in $K$),
%and write 
%$V' := \cO_\C^\flat$ (the ring of integers in $\C^\flat$).
%We are now in the situation of Section~\ref{subsec:coefficients};
%note in particular that $\Atilde_{K,A}$ may be identified
%with the ring denoted there by $R[1/v]$, and that $W(\C^\flat)_A$
%may be identified with the ring denoted there by $S[1/v]$.
%Also the Galois group denoted there by $G$,
%which is the Galois group $\Gal(L/F)$,
%is equal (in our context) to $G_{K_{\cyc}}$.
%Applying Proposition~\ref{prop:rigid facts} and Corollary~\ref{cor:descent with coefficients},
%we deduce part~(1) of the present theorem.
%%we find that the functor $M \mapsto W(\C^\flat)_A \otimes_{\Atilde_{K,A}} M$
%%induces an equivalence of categories between
%%the category of finitely generated projective $\Atilde_{K,A}$-modules,
%%and the category of finitely generated projective 
%%$W(\C^\flat)_A$-modules endowed with a continuous semi-linear action of 
%%$G_{K_{\cyc}}$.
%
%Part~(2) is proved by an identical argument.
%%An identical argument shows that the functor
%%$M\mapsto W(\C^\flat)_A\otimes_{\OEAtilde} M$ 
%%induces an equivalence of categories between
%%the category of finitely generated projective $\OEAtilde$-modules,
%%and the category of finitely generated projective 
%%$W(\C^\flat)_A$-modules endowed with a continuous semi-linear action of 
%%$G_{K_{\infty}}$.
%\mar{ME: Finish argument by treating the case $a > 1$!}
%\end{proof}
%
%\begin{lemma}
%	\label{lem:almost etale to almost Galois}
%	If $F \subseteq E$ is a finite Galois extension
%	of perfectoid fields, 
%with Galois group $G$, then the corresponding inclusion of rings
%of integers $\cO_F \subseteq \cO_E$ realizes $\cO_E$ as an almost 
%Galois extension of $\cO_F$, with Galois group $G$.
%\end{lemma}
%\begin{proof}
%	There are various more-or-less concrete ways to see this.
%	For example, 
%\cite[Prop.~5.23]{ScholzePerfectoid} shows that
%the morphism $\cO_F \to \cO_E$ is almost \'etale,
%which by definition
%\cite[Def.~4.12]{ScholzePerfectoid} means that the surjection
%of $\cO_E$-algebras $\cO_E\otimes_{\cO_F} \cO_E \to \cO_E$
%may be ``almost'' split, so that $\cO_E$ is almost a direct factor
%of the source.  Permuting such a splitting under the action
%of the Galois group $G$ then allows one to show that the morphism
%$\cO_E\otimes_{\cO_F}\cO_E \to \prod_{g \in G} \cO_E$ is an almost
%isomorphism.
%
%However, a more direct (but less explicit) way to prove the lemma is to
%use the first equivalence of categories in \cite[Thm.~5.2]{ScholzePerfectoid},
%namely the equivalence between the category of perfectoid 
%$F$-algebras and the category of perfectoid almost algebras
%over~$\cO_F$.  
%Under this equivalence, the morphism $\cO_E\otimes_{\cO_F} \cO_E
%\to \prod_{g \in G} \cO_E$ is taken to the morphism
%$E\otimes_F E \to \prod_{g \in G} E$.  This latter morphism
%{\em is} an isomorphism, since $E$ is Galois over $F$ with Galois
%group~$G$, and so we conclude 
%that the former morphism is an almost isomorphism,
%as required.
%\end{proof}

\section{\'Etale \texorpdfstring{$\varphi$}{phi}-modules}% \mar{TG:
  % in this section we should allow $\varphi$ not to stabilise $\A^+$
  % when we define \'etale $\varphi$-modules - so we would maybe define
  % them before the Kisin modules?}
%\section{Moduli stacks of
%  \texorpdfstring{$\varphi$}{phi}-modules}
\label{subsec: EG stuff}
% \mar{TG:
%   now tempted not to use this notation for coefficient rings in this
%   first section, but to instead keep $\gS_A$ for Kisin and use
%   something like $\AKDelta$ for $(\varphi,\Gamma)$.}
% \mar{TG: probably this first subsection should be left alone for
  % now, as it will mostly be absorbed into~\cite[\S5]{EGstacktheoreticimages}.}
% \mar{ME: Probably it would be simplest just to have a section like this,
% 	which esablishes some notation and recalls the definition/existence
% 	of $\cC_h$ and $\cR$. TG: with more general~$\phi$, though?}
% We begin by generalising some of the results of~\cite{MR2562795}
% and~\cite{EGstacktheoreticimages} on moduli stacks of Kisin modules to
% allow more general choices of Frobenius; this will be important for
% us, as we will work with~$(\varphi,\Gamma)$-modules below, which use a
% choice of Frobenius adapted to the cyclotomic extension of~$K$, rather
% than to the non-Galois extensions that occur in the theory of Kisin
% modules. Most of the proofs of these results are straightforward
% adaptations of those of~\cite{MR2562795}
% and~\cite{EGstacktheoreticimages}, and we have tried to strike a
% balance between making this section easy to read for the non-expert,
% and unneccessarily repeating standard arguments with only very minor changes.
% \mar{TG: apparently general LT fg modules stuff is written up in
% \cite{schneiderPhiGamma}
% so we can just cite that below if we want.}
% Fix a finite extension $k/\Fp$, and write $\gS:=W(k)[[u]]$. Let~$q=p^r$ be
% some power of~$p$, where $r|[k:\Fp]$, % \mar{TG: added this assumption
%   % following Kisin--Ren}
% and let~$\varphi$ be a ring endomorphism of~$\gS$
% which is congruent to the~$q$-power Frobenius endomorphism
% modulo~$p$. % \mar{TG: any continuity condition needed?! ME: I think it's automatic, since $W(k)[[u]]$ is a local ring with maximal ideal $(p,u)$, $\varphi$ is 
%        % necessarily $(p,u)$-adically continuous, which is what we would want.}
     
% Note that~$\varphi$ induces the usual $q$-power Frobenius
% on~$W(k)$, and that~$\varphi$ is faithfully flat. In~\cite{MR2562795,EGstacktheoreticimages} the only case
% considered is that where $q=p$ and $\varphi(u)=u^p$, but as we have
% already mentioned, the case that~$q=p$ and~$\varphi(u)=(1+u)^p-1$ will
% be important for us in the rest of this paper. Since it costs us
% nothing to work with general~$q$ and~$\varphi$, and we expect there to
% be future
% applications of working in this generality (for example, using
% Lubin--Tate $(\varphi,\Gamma)$-modules), we have made the minimal
% assumptions that we can.
In this section
we briefly recall and generalize some definitions and results
from~% \cite{EGetalephimodules} 
% and
~\cite{EGstacktheoreticimages}. Let~$R$ be a $\Zp$-algebra, equipped
  with a  ring endomorphism~$\varphi$, which is congruent to the
  ($p$-power) Frobenius modulo~$p$. If~$M$ is an $R$-module, we write \[\varphi^*M:=R\otimes_{R,\varphi}M.\]  % , and 
% extending some of them from~$\Z/p^a\Z$-algebras to $p$-adically
% complete $\Zp$-algebras, in a mostly formal fashion.
% We also use this discussion as an opportunity to establish some notation
% related to these results.
\begin{defn}\index{\'etale $\varphi$-module}
  \label{defn: general phi module} An \emph{\'etale
    $\varphi$-module over~$R$} is a finite $R$-module~$M$, equipped
  with a $\varphi$-semi-linear endomorphism $\varphi_M:M\to M$, which
  has the property that the induced $R$-linear morphism
  \[\Phi_M:\varphi^*M\stackrel{1\otimes\varphi_M}{\longrightarrow}M \]
  is an isomorphism.  A morphism of
   \'etale $\varphi$-modules is a morphism of the underlying
  $R$-modules which commutes with the
  morphisms~$\Phi_{M}$. We say that~$M$ is \emph{projective} (resp.\
  \emph{free}) if it is projective of constant rank (resp.\ free of
  constant rank) as an $R$-module. 
\end{defn}
We will typically apply this definition with~$R$ taken to be one of
the coefficient rings defined in Section~\ref{subsec:
  coefficients}. Of particular interest to us will be the cases
corresponding to imperfect fields of norms, as it is these cases which
fit into the framework of~\cite[\S5]{EGstacktheoreticimages}, and we
will use the results of that paper to 
prove the basic algebraicity properties of our moduli
stacks.

\begin{defn}\label{defn: base change of phi modules}
  Let~$S$ be an $R$-algebra, and let~$\varphi_S$ be a ring
  endomorphism of~$S$, which is congruent to Frobenius modulo~$p$, and
  is compatible with~$\varphi$ on~$R$. Then if~$M$ is an \'etale
  $\varphi$-module over~$R$, the extension of scalars $S\otimes_RM$ is
  naturally an \'etale $\varphi$-module over~$S$, with
  $\varphi_{S\otimes_R M}:=\varphi_S\otimes\varphi_M$.
  % \mar{ME: Changed this from $1\otimes\varphi_M$ to $\varphi_S \otimes
  %         \varphi_M$. I think that's right, but we should check that
  %         I didn't misunderstand! TG: agreed, commenting out.}
\end{defn}


\subsection{Multilinear algebra}
We briefly recall the multilinear algebra
of projective \'etale $\varphi$-modules.
Firstly, if
$P$ is a projective \'etale $\varphi$-module over~$R$, then we give
its $R$-dual
$P^\vee := \Hom_{R}(P,R)$ the structure of an
\'etale $\varphi$-module by defining the isomorphism $\varphi^*P^\vee\to
P^\vee$ to be the inverse of the transpose of~$\Phi_P$.
Secondly, if~$M$ and~$N$
are projective \'etale $\varphi$-modules, then we endow $M\otimes_RN$ 
with the structure of an \'etale $\varphi$-module by
defining $\varphi_{M\otimes N}:=\varphi_M\otimes\varphi_N$. % \mar{TG: this should be earlier, and maybe a lemma,
%   and just for general $R$ - then we can cite it in the proof of the
%   following result.} Analogous considerations apply to projective perfect
% \'etale $\varphi$-modules.

\begin{lem}
  \label{lem: projective phi module duality}If $M,$
  $N$ are projective \'etale $\varphi$-modules over~$R$ then we have
a natural identification 
$\Hom_{R,\varphi}(M,N)=(M^\vee\otimes_{R}
N)^{\varphi=1}$.
\end{lem}
\begin{proof}We have
  $\Hom_{R}(M,N)=M^\vee\otimes_{R} N$. Given~$f\in \Hom_{R}(M,N)$,
  regarded as an element of~$P:=M^\vee\otimes_{R} N$, we have $f\in
  P^{\varphi=1}$ if and only if $\Phi_P(1\otimes f)=f$, and by the
  definition of the $\varphi$-structure on~$P$, this is equivalent
  to~$f$ intertwining~$\Phi_M$ and~$\Phi_N$, as required.
\end{proof}

 % We say that $\A^+_A$ is $\varphi$-stable
 %  if~$\varphi(\A^+_A)\subset\A^+_A$; note in particular that if~$\A^+$
 %  is $\varphi$-stable, then~$\A^+_A$ is $\varphi$-stable for all~$A$.

% As in Definition~\ref{defn: general phi module}, if $M$ is an $\AAA_A$-module  (resp.\ $\A^+_A$ is $\varphi$-stable, and $\gM$ is an $\Aplus_A$-module) then we
% write~$\varphi^*M$
% for~$\AAA_A\otimes_{\AAA_A,\varphi}M$  (resp.\ $\varphi^*\gM$ for~$\Aplus_A\otimes_{\Aplus_A,\varphi}\gM$). Since~$\varphi$ is faithfully
% flat % (indeed, the natural map~$\varphi^*\Aplus_A\to\Aplus_A$ is injective,
% % and its image is a direct summand), 
% % \mar{TG: what's the easiest way to see this? Is it actually true
% %   in general? I think we maybe only use it in the $p$-adically complete Noetherian case,
% %   where I think I could prove it, but if it's true in general it would
% % be nice to know.}
% the
% functors  $M\mapsto\varphi^*M$, $\gM\mapsto\varphi^*\gM$ are exact.


% \begin{defn}\label{defn: etale phi module}\mar{TG: define this for a
%     general coefficient ring~$R$; explain base change for having a
%     compatible morphism $R\to S$}\mar{TG: investigate whether we
%     should actually just drop this definition??}
%   An \emph{\'etale $\varphi$-module with $A$-coefficients} is an
%   \'etale $\varphi$-module over  % pair
%   % $(M,\varphi_M)$ consisting of a finitely generated
%   $\AAA_A$. % -module~$M$,
%   % and a $\varphi$-semi-linear map $\varphi_M:M\to M$ which induces an
%   % isomorphism of $\AAA_A$-modules $\Phi_M:=\varphi_M\otimes 1:\varphi^*M\to M$. A morphism of \'etale $\varphi$-modules is a morphism
%   % of the underlying $\AAA_A$-modules which commutes with the
%   % morphisms~$\Phi_M$.
%   % We say that an \'etale $\varphi$-module is projective of rank~$d$ %  \emph{\'etale  locally
% %   % free of rank~$d$} (resp.\ \emph{fppf  locally
% %   % free of rank~$d$}, resp.\ \emph{fpqc  locally
% %   % free of rank~$d$})
% % if it is a finitely generated projective $\AAA_A$-module of constant
% % rank~$d$. We say that it is projective if it is projective of rank~$d$
% % for some~$d\ge 1$.
% % , and
% % if \'etale (resp.\ \emph{fppf}, resp.\ \emph{fpqc}) locally on~$\Spf A$, it is finite free of rank~$d$
% % over~$\AAA_A$.
% % \mar{TG: haven't thought at all about whether this is a
%   % good definition in the $p$-adically completed world!}
% \end{defn}
% \begin{rem}
%   When $q=p$ and $\varphi(u)=u^p$, a~$\varphi$-module of finite
%   $E$-height is usually called a (Breuil--)Kisin module.
% \end{rem}

% \begin{rem}
%   \label{rem: comparing to terminology in other papers}We will
%   sometimes refer to a $\varphi$-module with finite $F$-height as a
%   {\em Breuil--Kisin--Wach-module}, or a \emph{BKW-module}.\mar{TG:
%     might want to see if we use this terminology in the end. TG: we
%     don't, removing it.} In the case
%   that~$A$ is a $\cO/\varpi^a$-algebra for some~$a\ge 1$, a $\varphi$-module
%   of $F$-height at most~$h$ is the same thing as a Kisin module of
%   height~$F^h$ in the terminology of~\cite[\S 5.2]{EGstacktheoreticimages}.
% \end{rem}


%Let $\AAA_A$ be the $p$-adic completion of~ $\Aplus_A[1/u]$.% \mar{TG:
  % hopefully this is the right definition.}


% \begin{lem}
%   \label{lem: E(u) versus u mod p^a}If $A$ is a $\Z/p^a\Z$-algebra,
%   then for each $h\ge 1$, $u^{e(a+h-1)}$ is divisible by~$E(u)^h$ %  \mar{TG:
%     % if I'm not blundering, actually $u^{e(h+a-1)}$ would work here,
%     % but at least for now I'm not going to change it.}
%   in~$\Aplus_A$, and $E(u)^h$ is divisible by~$u^{e(h-a+1)}$.
% \end{lem}
% \begin{proof}Write $u^e=E(u)-pX$ for $X\in\Aplus_A$, and use the binomial theorem.  
% \end{proof}

% \begin{lem}
%   \label{lem: phi(u) divisible by u}If $A$ is a $\Z/p^a\Z$-algebra,
%   then for each $M\ge 1$, $\varphi(u^M)$ is divisible by~$u^{(M-a+1)q}$ in~$\Aplus_A$.
% \end{lem}
% \begin{proof}
%   Write $\varphi(u)=u^{q}+pY$ and use the binomial theorem.
% \end{proof}

% \begin{cor}
%   \label{cor: exponent of phi}If $\gM$ is a $\varphi$-module and
%   $u^M\gM\ne 0$, then $u^{(M-a+1)q}\varphi^*\gM\ne 0$.
% \end{cor}
% \begin{proof}If $u^{(M-a+1)q}\varphi^*\gM=0$, then by Lemma~\ref{lem:
%     phi(u) divisible by u} we have
%   $\varphi^*(u^M\gM)=\varphi(u^M)\varphi^*\gM=0$. Since~$\varphi$ is faithfully
%   flat, this implies that $u^M\gM=0$, a contradiction.  
% \end{proof}

% \begin{lem}\label{lem: map from finite E height to etale phi
%     module}Let $\gM$ be a $\varphi$-module of finite $E$-height. Then the
%   $p$-adic completion $\widehat{\gM[1/u]}$ naturally has the structure
%   of an \'etale $\varphi$-module.
%   \end{lem}
%   \begin{proof}% \mar{TG: I suspect this is dodgy or wrong, or at the
%       % very least not the best way to be thinking about this.}
% The
%     morphism $\Phi_{\gM}$ induces
%     an isomorphism \[\varphi^*\gM[1/E(u)]\to\gM[1/E(u)]\] of
%     $\Aplus_A[1/E(u)]$-modules, and thus for each $a\ge 1$ we have an
%     induced
%     isomorphism
%     \[\varphi^*(\gM/p^a\gM)[1/E(u)]\to(\gM/p^a\gM)[1/E(u)]\] of
%     $\Aplus_{A/p^aA}[1/E(u)]$-modules. By Lemma~\ref{lem: E(u) versus u
%       mod p^a}, this implies that we have an induced
%     isomorphism \[\varphi^*(\gM/p^a\gM)[1/u]\to(\gM/p^a\gM)[1/u]\]
%     of $\Aplus_{A/p^aA}[1/u]$-modules. Passing to the limit gives the
%     claim.    
%   \end{proof}
% The following lemmas are adaptations of those of~\cite[\S
% 5.2]{EGstacktheoreticimages} to our more general setting.                                                  
% \begin{lem}
%   \label{lem: kernels and images of morphisms of Kisin modules}If
%   $f:\gM\to\gM'$ is a morphism of ~$\varphi$-modules, and $\gM$ has $E$-height
%   at most~$h$, then the $\Aplus_A$-modules $\ker f$ and $\im f$ naturally
%   inherit the structure of $\varphi$-modules of $E$-height at most~$h$.
% \end{lem}
% \begin{proof}Considering the commutative
%   diagram
%   \[\xymatrix{\varphi^*\gM\ar[r]^{\varphi^*f}\ar@{^{(}->}[d]_{\Phi_{\gM}}&\varphi^*\gM'\ar@{^{(}->}[d]^{\Phi_{\gM'}}\\
%     \gM\ar[r]^{f}&\gM'}\]it
 %   is easy to see that $\ker f$ and $\im f$ each inherit an action
%   of~$\varphi$. They are both $u$-torsion free (as they are submodules of
%   $\gM$, $\gM'$ respectively). Then to see that $\Phi_{\ker f}$ and
%   $\Phi_{\im f}$ are
%   isomorphisms after inverting~$u$, it suffices to show that they are
%   surjective; this, and the property that the cokernels of
%   $\Phi_{\ker f}$ and $\Phi_{\im f}$ are killed by~$E(u)^h$,  follow easily from a
%   brief chase through the diagram (or the snake lemma).
% \end{proof}


% \begin{prop}
%   \label{prop: faithfulness of truncation of homs II}Assume that $A$
%   is a Noetherian $\Z/p^a\Z$-algebra for some $a\ge 1$. Let $M$ be an
%   integer with $M\ge ((e+q)(a-1)+eh+1)/(q-1)$. If $f:\gM\to\gM'$ is a morphism
%   of $\varphi$-modules of $E$-height at most $h$ with the property that
%   $f(\gM)\subseteq u^{M}\gM'$, then
%   $f=0$.
% \end{prop}
% \begin{proof}

%  We claim that if for some
%   positive integer $M$ as in the statement we have $f(\gM)\subseteq u^{M}\gM'$, then
%   we also have $f(\gM)\subseteq u^{q(M-a+1)-e(a+h-1)}\gM'$. Since $q(M-a+1)-e(a+h-1)\ge
%   M+1$, it follows by 
%   induction that we have $f(\gM)\subseteq u^{M}\gM'$ for all~$M$, and
%   it follows from the Artin--Rees lemma that $f(\gM)=0$, as required.

% To see the claim, note that if $f(\gM)\subseteq u^{M}\gM'$ then using
% Lemmas~\ref{lem: E(u) versus u mod p^a} and~\ref{lem: phi(u) divisible by u} we
% have 
% \begin{multline*}
% u^{e(a+h-1)}f(\gM)=f(u^{e(a+h-1)}\gM)\subseteq f(E(u)^h\gM)\subseteq f((\Phi_{\gM})(\varphi^*\gM)) \\ =\Phi_{\gM'}((\varphi^* f)(\gM))\subseteq
% \varphi(u)^M\Phi_{\gM'}(\varphi^*\gM')\subseteq u^{q(M-a+1)}\gM'.
% \end{multline*}
%  Since $\gM'$ is $u$-torsion free, we conclude that $f(\gM)\subseteq u^{q(M-a+1)-e(a+h-1)}\gM' $, as required.
% \end{proof}

% \begin{lem}% \mar{TG: I'm a bit suspicious that this was easy to prove
%     % and that I didn't run into serious trouble with exponents - my
%     % memory was that when we wrote proper.tex we couldn't argue like
%     % this if $p=2$, so I suspect I'm missing something.}
%   \label{lem: Noetherian extension of scalars and bound on
%     torsion}Assume that $A$ is a Noetherian $\Z/p^a\Z$-algebra for
%   some $a\ge 1$, and that $B$ is an $A$-algebra which is finitely
% presented as an $A$-module. Let $\gM$
%   be a $\varphi$-module of $E$-height at most $h$. Then the $u$-torsion submodule
%   $(\gM\otimes_A B)[u^\infty]$ is annihilated by $u^{\lfloor
%     (aq+e(a+h-1))/(q-1)\rfloor}$, and the $u$-torsion free
%   quotient \[(\gM\otimes_A B)/(\gM\otimes_A B)[u^\infty]\]naturally
%   inherits the structure of a $\varphi$-module of $E$-height at most~$h$ with $B$-coefficients.
% \end{lem}
% \begin{proof}Write $\gM_B:=\gM\otimes_A B$. Certainly $\gM_B/\gM_B[u^\infty]$
%   is a finitely generated $u$-torsion free $\Aplus_B$-module, and the
%   cokernel of~$\varphi_{\gM_B}$ is killed by
%   $E(u)^h$ by right exactness of tensor products. To see that
%   $\Phi_{\gM_B}$ is injective it is enough to check
%   that~$\Phi_{\gM_B}[1/u]$ is an isomorphism, which is true
%   because~$\Phi_{\gM}[1/u]$ is an isomorphism by Lemma~\ref{lem: E(u)
%     versus u mod p^a}.

%   It remains to prove the bound on the exponent of the torsion. We
%   have an exact sequence
%   of~$\Aplus_B$-modules
%   \[\Tor^A_1(\gM/\varphi^*\gM,B) \to
%     \varphi^*\gM_B\to\gM_B\to(\gM/\varphi^*\gM)\otimes_AB\to 0.\]
%   Since $\gM/\varphi^*\gM$ is killed by $E(u)^h$, it is killed by
%   $u^{e(a+h-1)}$ by Lemma~\ref{lem: E(u) versus u mod p^a}; thus so is
%   $\Tor^A_1(\gM/\varphi^*\gM,B)$.

%   Since $\Aplus_B$ is Noetherian, there is some $t$ such that
%   $u^t\gM_B[u^\infty]=0$. Choose $t$ minimal (that is, let $t$ be
%   maximal such that $\gM_B$ has an element of exponent $u^t$). By Corollary~\ref{cor: exponent of phi},
%    $\varphi^*\gM_B$ has an element of exponent at least
%   $u^{(t-a)q}$. % To see this, note that by the choice of~$t$, we have
%   % $u^{t-1}\gM_B\ne 0$, so that since~$\varphi$ is faithfully flat, we
%   % have $\varphi(u)^{t-1}\varphi^*\gM_B\ne 0$. By Lemma~\ref{lem:
%   %   phi(u) divisible by u}, it follows that
%   % $u^{(t-a)q}\varphi^*\gM_B\ne 0$, as claimed.
% The image of this
% element in $\gM_B$, being $u$-power torsion, is killed by $u^t$; since
% $\Tor^A_1(\gM/\varphi^*\gM,B)$ is killed by
% $u^{e(a+h-1)}$, we see that $q(t-a)\le t +e(a+h-1)$, as required.
% \end{proof}

% We now introduce the moduli stacks of~$\varphi$-modules that we need
% to consider, and establish their main properties,
% following~\cite{MR2562795,EGstacktheoreticimages,CEGSKisinwithdd}. \mar{TG:
% for consistency with later sections, perhaps we should call these $\cC_d,\cR_d$?}

\section{Frobenius descent}\label{subsec: perfect
  etale phi modules}% \mar{TG: I suggest that we should move this
%   forward a bit, at least after the following section where~$\C^\flat$
% is introduced. We should also connect its assumptions to the new
% treatment of perfect coefficient rings in
% Section~\ref{subsec:coefficients}; is it clear that those are actually
% the perfections of our imperfect coefficient rings?}
  % \mar{TG: do we need to
  % take extra care to deal with the case that $\A^+$ isn't
  % $\varphi$-stable?}\mar{TG: explain that the basic idea for this is
  % that the isomorphism $\varphi^*M\isoto M$ lets us ``deperfectify'',
  % but the fact that we have a completion means we have to do a bit
  % more work}
% \mar{TG:
% I don't see where this finite type hypothesis is used right now, but I
% guess probably it will be used in the application?}
Suppose that~$\A$ is as in Situation~\ref{subsubsec:general framework},
%Section~\ref{subsec: EG stuff},
and that we have a continuous $\varphi$-equivariant embedding % \mar{TG:
% added ``continuous''; didn't think about whether this is automatic,
% but I guess if it is we should say so.}
% $\A^+\into \Ainf$ (and therefore
$\A\into W(\C^\flat)$. % ). \mar{TG: we don't use these assumptions
  % in the main arguments
  % below, but it seems easier to define the perfection if our
  % coefficient ring is a subring of a perfect ring? Note though that
  % currently $\Ainf$ is defined in the next section, so we might need
  % to reorder this material.}
% Since~$\varphi$ is
% bijective on~$W(\C^\flat)$, we have an increasing
% unions \[\A^+\subset\varphi^{-1}(\A^+)\subset\varphi^{-2}(\A^+)\subset\dots\subset
%   \Ainf.\] 
% \[\A\subset\varphi^{-1}(\A)\subset\varphi^{-2}(\A)\subset\dots\subset
  % W(\C^\flat).\]
Let~$A$ be a finite type $\Z/p^a$-algebra for some $a\ge
1$. Assume that the induced map $\A_A\to
  W(\C^\flat)_A$ is a faithfully flat injection. 
%   \mar{ME: This faithful flatness needs checking, and probably needs a hypothesis; e.g.\
% 	  that the cokernel of the embedding be $p$-torsion free,
% 	  so that it induces an embedding $k[[T]]\to \C^\flat$ 
% 	  after reducing mod $p$. Or maybe we could just assume it,
%   since in applications I think we'll get it from Prop.~\ref{prop: maps
%   of coefficient rings are faithfully flat injections}. TG: done and
% added to the following remark; commenting out.}
  Since~$\varphi$ is
bijective on~$W(\C^\flat)$, we have an increasing
union % \[\A^+_A\subset\varphi^{-1}(\A^+_A)\subset\varphi^{-2}(\A^+_A)\subset\dots\subset
% \AAinf{A},\]
\[\A_A\subset\varphi^{-1}(\A_A)\subset\varphi^{-2}(\A_A)\subset\dots\subset
  W(\C^\flat)_A.\]% \mar{TG: do we know that $\A_A\to
  % W(\C^\flat)_A$ is an embedding in this generality, or should we be
  % assuming that $A$ is finite type and the referencing the faithful flatness material above?}
We let~$\tA_A$ be the closure of % \mar{TG: ultimately need to link this
  % to earlier field of norms material.}
% completion (for the $T$-adic
% topology) of
$\cup_{n\ge 0}\varphi^{-n}(\A_A)$ in~$W(\C^\flat)_A$,
and we
set~$\tA^+_A:=\tA_A\cap\AAinf{A}$. % \mar{TG: I'm not totally sure that
  % this is what we want to do, but I think there is something to be
  % said for having $\varphi$ act on $\tA^+_A$.}
  % $\tA_A=\tA^+_A[1/T]$.
  Note that~$\varphi$ extends to a bijection
  on~$\tA_A$, which induces a bijection on~$\tA^+_A$.

  \begin{rem}
    \label{rem: abstract perfection section linked to our concrete
      settings} We will apply the results of this section in the
    setting introduced in Section~\ref{subsec:rings}, taking~$\A=\A_K$
    or~$\A=\gS$. In either case it follows from Proposition~\ref{prop: maps of coefficient rings are faithfully flat
    injections} that $\A_A\to
  W(\C^\flat)_A$ is a faithfully flat injection, and it follows easily from Theorem~\ref{thm:field of
      norms} that in the former case we have~$\tA_A=\tA_{K,A}$, and in
    the latter case we have~$\tA_A=\tOEA$.%  Note though that in the
    % former case we do not necessarily have~$\tA_A^+=\tA_{K,A}^+$, as
    % the latter ring need not be~$\varphi$-stable (see Remark~\ref{rem:
    %   not phi stable T}); since the ring~$\tA_A^+$ is only used in
    % order to prove the results of this section concerning \'etale
    % $\varphi$-modules over~$\tA_A$,
    % this does not cause us any difficulties.\mar{TG: I removed that
    % ring so I'm removing this comment.}
  \end{rem}

  \begin{rem}
    \label{rem: warning about A plus not phi stable implies not in A
      tilde plus}If~$\A^+_A$ is not $\varphi$-stable, then $\A^+_A$ is
    not a subring of ~$\tA^+_A$, and in particular it is not equal to
    $\A_A\cap\tA^+_A$. 
  \end{rem}

We will make use of the following variant of Lemma~\ref{lem: phi on powers of T without assuming A plus
      stable}; in view of Remark~\ref{rem: warning about A plus not phi stable implies not in A
      tilde plus} we cannot apply Lemma~\ref{lem: phi on powers of T without assuming A plus
      stable} in the setting of $\tA^+_A$-modules.

    
  \begin{lem}
    \label{lem: phi on powers of T for A tilde plus}There is an
    integer~$r\ge 0$ such that for all~$s\ge r$ we have
    $T^s\in\tA^+_A$, and an integer~$C\ge 0$ such that for
    all~$n\ge 1$, we have $\varphi(T^n)\in T^{pn-C}\tA^+_A$.
  \end{lem}
  \begin{proof}Since~$T$ is topologically nilpotent in~$\A_A^+$ it is
    topologically nilpotent in~$\tA_A$, so in particular
    $T^s\in\tA^+_A$ for all~$s$ sufficiently large. For the second
    statement, we follow the proof of Lemma~\ref{lem: phi on powers of
      T without assuming A plus stable}, and write
    $\varphi(T)=T^p+pY$ for some~$Y\in\A_A$, so that~$\varphi(T^n)=(T^p+pY)^n$. If we
    expand using the binomial theorem, and recall that~$p^a=0$, then
    we see that $\varphi(T^n)-T^{np}$ is a sum of terms of the form 
    $T^{p(n-r)}Y^r$ for some $0\le r \le a-1$. It therefore suffices
    to  choose~$C$ such that $T^{C}T^{-pr}Y^r\in \tA^+_A$ for $0\le r
    \le a-1$, which we can do by the topological nilpotence of~$T$.
  \end{proof}

  
% For any~$A$, $\varphi$ is injective on~$\A_A$ (indeed, it is easy to
% show by induction on~$n$ that the kernel of~$\varphi$ is contained
% in~$p^n\A_A$ for all~$n$). 

% and we
% write~$\tA_A^+$  to denote the ``perfection of $\A_A^+$ in the
% $T$-direction''.

% More precisely, we let\mar{TG: need some actual definition here, I
%   guess we just take $\varprojlim_\varphi$ and then complete?! Not
%   sure if we need to do this on the ones without $A$ and then put in
%   coefficients after or just do it once, I would have guessed the
%   latter was fine; perhaps the thing to do is to make sure it lines up
% with what we end up having in the almost \'etale descent. Should
% probably explain that $\varphi$ is injective.}
% $\F_p[[u^{1/p^{\infty}}]]$ denote the $u$-adic completion of the
% perfection of $\F_p[[u]]$, and then let
% $W_a\bigl(\F_p[[u^{1/p^{\infty}}]]\bigr)$ denote the corresponding
% ring of 
% mod $p^a$  Witt vectors; we write $u$ rather than $[u]$ to
% denote the Teichm\"uller lift of $u$ in
% $W_a\bigl(\F_p[[u^{1/p^{\infty}}]]\bigr)$. We then let
% $A[[u^{1/p^{\infty}}]]$
% denote the $u$-adic completion of
% $A \otimes_{\cO/\varpi^a} W_a\bigl(\F_p[[u^{1/p^{\infty}}]]\bigr)$.
% We also write
% $A((u^{1/p^{\infty}})) := A[[u^{1/p^{\infty}}]][1/u].$ \mar{TG: change
% notation below from $u$ to $T$ and from $\OEA$ and $\gS$ to $\A$ etc.}

% Recall that if~$A$ is an $\cO/\varpi^a$-algebra, then we write
% $\gS_A:=(W_a(k)\otimes_{\Zp}A)[[u]]$,
% $\OEA:=(W_a(k)\otimes_{\Zp}A)((u))$. Analogously, we write
% $\widetilde{\gS}_A:=(W_a(k)\otimes_{\Zp}A)[[u^{1/p^\infty}]]$, 
% $\widetilde{\cO}_{\mathcal{E},A}:=(W_a(k)\otimes_{\Zp}A)((u^{1/p^{\infty}}))$. 
% \mar{TG:
  % I am in no way wedded to this notation and terminology, just thought
  % it would be good to have something.}
% \begin{defn}
%   A \emph{perfect \'etale $\varphi$-module with $A$-coefficients} is a
%   pair~$(M,\varphi_{M})$ consisting of a finitely generated
%   $\tA_A$-module~$M$, and a
%   $\varphi$-semi-linear map $\varphi_{M}:M\to M$ which induces an
%   isomorphism of $\tA_A$-modules
%   $\Phi_{M}:=\varphi_{M}\otimes 1:\varphi^*M\to M$. A morphism of
%   perfect \'etale $\varphi$-modules is a morphism of the underlying
%   $\tA_A$-modules which commutes with the
%   morphisms~$\Phi_{M}$. We say that a perfect \'etale $\varphi$-module
%   is projective (resp.\ free) if it is projective (resp.\ free) as
%   an~$\tA_A$-module.
% \end{defn}

\begin{rem}
\label{rem:A-tilde as a localization of  A-tilde-plus}
As already noted in Remark~\ref{rem: warning about A plus not phi stable implies not in A tilde plus}, $\tA^+_A$  need not contain~$\A^+_A$, and so, although $\A_A
= \A_A^+[1/T]$, it doesn't make sense to write $\tA_A =  \tA^+_A[1/T]$.
On the other hand, Lemma~\ref{lem: phi on powers of T for A tilde plus} does
ensure that $T^r \in \tA^+_A$ for some $r > 0$, and then it does
make sense to write, and is true (since $T$ is
topologically nilpotent in $\tA_A$, while $\tA^+_A$ is open 
in $\tA_A$), that $\tA_A = \tA^+_A[1/T^r].$
\end{rem}

Given an \'etale $\varphi$-module $M$ over~$\A_A$, we may form its
\emph{perfection} $\widetilde{M}:=\tA_A\otimes_{\A_A} M$,
\index{perfection (of \'etale $\varphi$-module)}
which is an \'etale $\varphi$-module over ~$\tA_A$.
% In this section we %t ake~$q=p$\mar{TG: probably don't have to?}
% % } and \mar{TG: should this actually be in the
% %   more general setting so that we can also apply it to $(\varphi,\Gamma)$-modules?}
% % let~$\varphi$ be the usual Frobenius, so that~$\varphi(u)=u^p$.
% % We
% let $A$
% be a $\cO/\varpi^a$-algebra (for some $a \geq 1$). 


\subsection{Descending morphisms}


 % by
% demanding that
% $\varphi(\varphi_{P^\vee}(f)(\varphi_{P}(x)))=f(x)$ for
% any~$x\in P$, $f\in P^\vee$. Note that this is enough to uniquely
% specify $\varphi_{P^\vee}$, so it is enough to check the existence of
% such a map~$\varphi_{P^\vee}$. By Lemma~\ref{lem: projective phi module is summand of
%     free}, it is enough to prove this in the free case, in which case
%   we may choose a basis, and \mar{TG:
%   need to check this}


\begin{prop}% \mar{TG: presumably ultimately we will either have one of
    % these two results, or we will merge them.}
  \label{prop: descent of morphisms for phi modules}Let~$M, N$ be
  projective \'etale $\varphi$-modules. Then the natural map
  $\Hom_{\A_A,\varphi}(M,N)\to\Hom_{\tA_A,\varphi}(\widetilde{M},\widetilde{N})$
  is a bijection.
\end{prop}
\begin{proof}
  Write $P:=M^\vee\otimes N$, so that~$P$ is a projective
  \'etale $\varphi$-module over~$\A_A$.  Then by Lemma~\ref{lem: projective phi module duality}, % it follows from the
  % definitions\mar{TG: surely correct, didn't actually try to check carefully}
  % $\Hom_{\A_A,\varphi}(M,N)=P^{\varphi=1}$,
  % $\Hom_{\tA_A,\varphi}(\widetilde{M},N^{1/p^\infty})=(P^{1/p^\infty})^{\varphi
  % =1}$,
  we are reduced to checking that the natural morphism of $A$-modules
  $P^{\varphi=1}\to\widetilde{P}^{\varphi=1}$ is an isomorphism. It
  is certainly injective, so we need only show that it is
  surjective.% \mar{TG: I think we can then reduce to surjectivity for
    % $p=0$ and then literally the argument that's written is OK?!}

Since the formation of $\varphi$-invariants is compatible with direct
sums, it follows from~ \cite[Lem.\ 5.2.14]{EGstacktheoreticimages} that it is enough to consider the case that~$P$ is in fact
free. By Lemma~\ref{lem: phi on powers of T for A tilde plus}
(\emph{cf.}\ Remark~\ref{rem:A-tilde as a localization of  A-tilde-plus})
we can choose a $\varphi$-stable free
$\tA^+_A$-submodule~$\widetilde{\gP}$ of~$\widetilde{P}$, which generates $\widetilde{P}$ over $\A_A$
(indeed, choose any basis of~$\widetilde{P}$, multiply by a sufficiently large
power of~$T$, and let~$\widetilde{\gP}$ be the submodule
generated by this scaled basis).
%We write $\widetilde{\gP}:=\tA^+_A\otimes_{\A^+_A}\gP$.
%\mar{ME: Why not choose this integral basis on the $P$ level? ME: Answer: as
%TG explained, $\tA^+_A$ is  generally  {\em not} an $\A_A$-algebra!}

Consider an element $x\in \widetilde{P}^{\varphi=1}$. Choose an
integer $r\ge 1$ such that if~$s\ge r$ then $\varphi(T^s\tA^+_A)\subseteq T^{s+1}\tA^+_A$ (such
an~$r$ exists by Lemma~\ref{lem: phi on powers of T for A tilde plus}). We may write $x=x_1+x_2$
where $x_2\in T^r\widetilde{\gP}$ and $x_1\in
\varphi^{-n}(\A_A)\otimes_{\A_A}P$ for some $n\ge 0$. Choose~$n$ to be
minimal with this property. If $n>0$, then since
$x=\varphi(x)=\varphi(x_1)+\varphi(x_2)$, and $\varphi(x_2)\in T^r\widetilde{\gP}$,
we see that we may replace~$x_1$ with $\varphi(x_1)\in\varphi^{1-n}(\A_A)\otimes_{\A_A}P$, a contradiction.

Thus $x_1\in P$, and we have $x_1-\varphi(x_1)=\varphi(x_2)-x_2\in
T^r\widetilde{\gP}\cap P$. % =u\gP$ (the last equality holding by the
% freeness of~$\gP$ and the evident equality $u A[[u^{1/p^\infty}]]\cap
% A((u))=uA[[u]]$). 
% % \mar{ME: Justify the equality
% % $u\widetilde{\gP}\cap P=u\gP$? Might be good, esp.\ since $\gP$ need
% % not be projective. TG: in fact we now choose it to be free, so I think
% % it's all OK; commenting out.}
The sum
$x':=x_1+\sum_{i=0}^\infty\varphi^i(\varphi(x_1)-x_1)$ converges to an
element of~$T^r\widetilde{\gP}$, and the sum therefore converges
in~$P$. By definition we have $\varphi(x')=x'$. We have $x-x_1$,
$x'-x_1\in T^r\widetilde{\gP}$, so $x-x'\in T^r\widetilde{\gP}$; then
$x-x'=\varphi(x-x')\in T^{r+1}\widetilde{\gP}$, and iterating gives
$x=x'\in P$, as required.
\end{proof}

\subsection{Descending objects}
% Sketch:  since we can descend morphisms, to descend projectives it's enough
% to descend free objects (since we can then furthermore descend 
% idempotent endomorphisms of them).

% Now we can choose a Kisin module and truncate, via which we descend to
% $A((u^{1/p^n}))$ for some $n$, and then interpret the \'etale
% $\varphi$-structure as a concrete descent that lets us get down
% to $A((u)).$ 
% \mar{TG:literally thinking of an $A[[u]]$ module as an
%   $A[[u^p]]$-module via $A[[u^p]]\into A[[u]]$, which is what we tensor
% over when forming~$\varphi^*\gM$. To see what the truncation argument
% is doing, remember the case that
% $\varphi=(1+u^{1/p}+u^{1+1/p^2}+\dots)$; it's doing a conjugation that
% gets rid of those powers.}

We will show that every projective  \'etale $\varphi$-module over~$\tA_A$
arises as the perfection of an \'etale $\varphi$-module over~$\A_A$.
We first note the following analogue of~\cite[Lem.\ 5.2.14]{EGstacktheoreticimages}
for \'etale $\varphi$-modules over~$\tA_A$.

\begin{lemma}\label{lem: projective perfect phi module is summand of free}
If~$M$ is a projective  \'etale $\varphi$-module
  over~$\tA_A$, then~$M$ is a direct summand of a free 
  \'etale $\varphi$-module
  over~$\tA_A$. 
  \end{lemma}
  \begin{proof}
	  This can be proved in an identical fashion to~\cite[Lem.\ 5.2.14]{EGstacktheoreticimages}.
  \end{proof}

\begin{prop}
  \label{prop: descending projective modules from infinite level}Let
  $M$ be a projective  \'etale $\varphi$-module over~$\tA_A$. Then there is a projective \'etale
  $\varphi$-module $M_0$ over~$\A_A$ and an isomorphism
  $M\isoto \widetilde{M}_0$.
\end{prop}
\begin{proof}Suppose firstly that $M$ is  free of some rank~$d$ as an
  $\tA_A$-module. Let~$X$ denote the matrix
  of~$\varphi$ with respect to a choice of basis
  $e_1,\dots,e_d$. As in the proof of Proposition~\ref{prop: descent of morphisms for phi modules}, after possibly
  scaling the~$e_i$ by powers of~$T$, we may assume
  that~$X$ has entries in~$\tA^+_A$. Since $\Phi_{M}$
  is an isomorphism, it follows from Remark~\ref{rem:A-tilde as a localization of  A-tilde-plus})
that there is an integer $h\ge 0$ such that
$X^{-1} \in T^{-h} M_d(\tA^+_A).$
%and a matrix~$Z$ with entries in~$\tA^+_A$ such that $ZX=T^h\Id_d$.\mar{TG: note that
%    we don't assume that~$T^h\in\tA^+_A$, but we also don't appear to
%    need this, unless I'm blundering again.  ME: I guess if we take determinants,
%we do find that $T^{hd} \in \tA^+_A,$ fwiw.}

  If we change basis via a matrix~$Y\in M_d(\tA^+_A)\cap\GL_d(\tA_A)$, the new matrix for~$\varphi$ is
  $\varphi(Y)XY^{-1}$. It suffices to show that we can choose~$Y$ so
  that this matrix has entries in~$\A_A$ (as we can then let~$M_0$ be
  the $\A_A$-span of this new basis). %We will do this in two stages.

  Fix some~$H>\lceil (C+h+1)/(p-1)\rceil$, where~$C$ is as in
%\mar{ME: Double-check this is the correct/optimal $H$, just for the reader's benefit.}
  Lemma~\ref{lem: phi on powers of T for A tilde plus}, and write $X=X'+X''$ where
$X'\in T^{H+h}M_d(\tA^+_A)$ and
  $X''\in M_d(\varphi^{-n}(\A_A))$ % \mar{TG: replace with $\varphi^{-n}$
    % and then hopefully this argument works?}
% \mar{ME:  This  seems  to use that the union of the $\varphi^{-n}(\A^+_A)$ is dense
% in $\tA^+_A$.   I think this  would be ok if $\A^+_A$  were equal
% to $\A_A \cap \A_{\inf,A}$, but this doesn't hold in general (e.g.\ b/c
% this intersection is $\varphi$-stable, and $\A^+_A$ need not be. TG:
% hopefully now fixed but obviously needs checking! Thankfully it seems
% that for the most part one can just work with $\A_A$ instead of
% $\A^+_A$, in fact.}
  for some sufficiently large~$n$. (The reason for this choice of~$H$ is to be able
  to apply Lemma~\ref{lem: Frobenius conjugacy without phi stable} below.) Taking~$Y=X$, we may replace~$X,X',X''$ % \mar{TG: needs
    % $T^{H+h}\tA^+_A$ to be $\varphi$-stable.}
  by~$\varphi(X),\varphi(X'),\varphi(X'')$ respectively
(note that by Lemma~\ref{lem: phi on powers of T for A tilde plus}, we have
  $\varphi(X')\in T^{H+h}M_d(\tA^+_A)$, because 
  $p(H+h)-C\ge H+h$ by our choice of~$C$),
 which has the
  effect of replacing~$n$ by~$(n-1)$.
Iterating this procedure, we
  can assume that~$n=0$, so~$X''\in M_d(\A_A)$. By Lemma~\ref{lem:
    Frobenius conjugacy without phi stable} below, we can find~$Y$
  such that $\varphi(Y)XY^{-1}=X''$, completing the proof in the case
  that~$M$ is free.

%   Since $H>h$ and $ZX=u^h\Id$, we may write $ZX_1=u^h(1+u^{H-h}T)$ for
%   some~$T\in M_d(\widetilde{\gS}_A)$. Set~$Y_0=\Id_d$, and iteratively
%   define \[Y_{i+1}=u^{-h}(1+u^{H-h}T)^{-1}Z\varphi(Y_i)X.\] % An easy
%   % induction shows that each~$Y_i\in 1+u^{H-h}M_d(\gS_A)$.
% We claim that~$(Y_i)$ is a Cauchy sequence; granting this, it tends to
% a limit~$Y\in 1+u^{H-h}M_d(\gS_A)$ satisfying $\varphi(Y)XY^{-1}=X_1$,
% as required.

  % Choose $h\ge 0$ so that $\varphi_M(\widetilde{\gM})\supseteq u^h\widetilde{\gM}$. For some
  % $n\ge 0$ we can write $\varphi_M(\gM)\subseteq
  % \gS_A[[u^{1/p^n}]]\otimes_{\gS_A}\gM+u^h\widetilde{\gM}$. If we set
  % $\gM'=u^{p^n}\gM$, then we find that
  % $\varphi_M(\widetilde{\gM}')\supseteq u^h\widetilde{\gM}'$ and $\varphi_M(\gM')\subseteq
  % \gM'+u^h\widetilde{\gM}'$.
  % \mar{TG: fix this - I mean to iterate Frobenius.}

% \mar{TG: can we now conclude without choosing a basis by some cunning
%   summing up of lattices? It's basically Hensel's lemma from here, anyway.}

  We now return to the general case in which $M$ is only assumed finitely
  generated projective, rather than free. By Lemma~\ref{lem: projective perfect
	  phi module is summand of free}, we may write $M$ as a direct summand of a free 
  \'etale $\varphi$-module~$F$ over~$\tA_A$. By the case already proved, we may
  write $F\isoto\widetilde{F}_0$ for some free \'etale $\varphi$-module
  $F_0$. By Proposition~\ref{prop: descent of morphisms for phi
    modules}, the idempotent in~$\End(F)$ corresponding to ~$M$ comes
  from an idempotent in~$\End(F_0)$, and we may take~$M_0$ to be the
  \'etale $\varphi$-module corresponding to this idempotent.%  \mar{TG: this argument needs to be figured
    % out. TG: the point is to apply the variant of Prop~\ref{prop: representability of phi-invariants} indicated in
    % the margin above to $M^*\otimes N$.}
  \end{proof}The following lemma and its proof are based
  on~\cite[Prop.\ 2.2]{MR2562795}.
  \begin{lem}% \mar{TG: constant will depend on constant in Lemma~\ref{lem: phi on powers of T without assuming A plus stable}.}
    \label{lem: Frobenius conjugacy without phi stable}Suppose
    that~$X\in M_d(\tA^+_A)\cap\GL_d(\tA_A)$, and that
    $X^{-1}\in T^{-h} M_d(\tA^+_A)$ for some~$h\ge
    0$. Suppose that ~$X''\in M_d(\tA^+_A)\cap\GL_d(\tA_A)$ is such
    that $X^{-1}X''\in 1+T^{\lceil (C+h+1)/(p-1)\rceil}M_d(\tA^+_A)$, where~$C$
    is as in Lemma~{\em \ref{lem: phi on powers of T for A tilde plus}}. % \mar{TG: constant to be determined.}
Then there exists $Y\in M_d(\tA^+_A)\cap\GL_d(\tA_A)$ with $X''=\varphi(Y)XY^{-1}$.
  \end{lem}
  \begin{proof}Define sequences $X_i$, $h_i$ by $X_0=X$,
    $h_0=(X'')^{-1}X$, and for each $i\ge 1$,
    $X_i=\varphi(h_{i-1})X_{i-1}h_{i-1}^{-1}$,
    $h_i=(X'')^{-1}X_i$. Then if we set $y_i=h_ih_{i-1}\cdots h_0$, we
    have $X''=X_ih_i^{-1}=\varphi(y_{i-1})Xy_{i}^{-1}$ for
    each~$i$. We claim that the~$y_i$ tend to a limit~$Y$ as
    $i\to\infty$; then we have $X''=\varphi(Y)XY^{-1}$, as required.

    To see that the ~$y_i$ tend to a limit, it is enough to show that
    $h_i\to 1$ as~$i\to\infty$. To see this, suppose
    that $h_i\in 1+T^{s}M_d(\tA^+_A)$ for some $s\ge 
    (C+h+1)/(p-1)$. We
    have \[X''h_{i+1}=X_{i+1}=\varphi(h_{i})X_{i}h_{i}^{-1}=\varphi(h_{i})X'',\]so that
    $h_{i+1}=(X'')^{-1}\varphi(h_{i})X''$. Using Lemma~\ref{lem: phi on powers of T for A tilde plus} and the assumption that $(X'')^{-1}\in T^{-h} M_d(\tA^+_A)$,
    we see that $h_{i+1}\in 1+T^{ps-C-h}M_d(\tA^+_A)$. Since
    $ps-C-h\ge s+1$, we are done.
      \end{proof}

  \subsection{An equivalence of categories}We summarise the results
  of this section in the following proposition.
  \begin{prop}
	  \label{prop:perfect equivalence}% \mar{TG: compare to proof
            % of Cor~\ref{cor:descent with coefficients}?}
  Let~$A$ be a finite type $\Z/p^a$-algebra for some ~$a\ge 1$.  Then the functor $M\mapsto\widetilde{M}$ is an equivalence of
    categories from the category of projective \'etale
    $\varphi$-modules over~$\A_A$ to the category of
    projective \'etale $\varphi$-modules over~$\tA_A$.
  \end{prop}
  \begin{proof}
    The functor is essentially surjective by Proposition~\ref{prop:
      descending projective modules from infinite level}, and fully
    faithful by Proposition~\ref{prop: descent of morphisms for phi modules}.
  \end{proof}
\section{% Basic definitions ---
  \texorpdfstring{$(\varphi,\Gamma)$}{(phi,Gamma)}--modules}
\label{subsec: phi gamma over Zp}
% \mar{TG:
%   just slapping stuff in that is hopefully mostly true in an effort to
% make progress\dots}
% We now return to the setting of Section~\ref{section: coefficient rings}. 
% In particular,
% we fix a finite extension $K$ of $\Q_p$, and use the notation
% for the various coefficient rings introduced there, such
% as $\A_K$ and $\A_{K,A}$.\mar{TG: remove this if moved back.}

%\subsection{Recalling some definitions}
By definition, an {\em \'etale $(\varphi,\Gamma_K)$-module} % (resp.\ an \'etale
% $(\varphi,\Gammat)$-module)
is a finitely generated % \mar{TG: could make
% other hypotheses, not quite sure what we want in this setting without
% coefficients?
 % $\AKDelta$-module (resp.\ a finitely generated
$\A_K$-module $M$, equipped with
\begin{itemize}
\item a $\varphi$-linear morphism $\varphi:M\to M$ with the property
  that the corresponding morphism $\Phi_M:\varphi^*M\to M$ is an
  isomorphism (i.e.\ $M$ is given the structure of an \'etale
  $\varphi$-module over $\A_K$), and 
\item a continuous semi-linear action
  of~$\Gamma_K$ that commutes with~$\varphi$.% (resp.\ of~$\Gammat$).
\end{itemize}


% \begin{remark} In the literature on $(\varphi,\Gamma)$-modules, it is
%   traditional to use $\Gamma$ to denote both the torsion-free part of
%   $\Gal(K(\zeta_{p^\infty})/K)$, as we have done, {\em and} the entire group
%   $\Gal(K(\zeta_{p^\infty})/K)$.  (More precisely, in some sources it denotes
%   one of these groups, and in other sources, the other.)  We have
%   introduced the non-standard notation $\tGamma$ simply to enable us
%   to speak of both groups in the same discussion. As explained below
%   (see Lemma~\ref{lem: phi Gamma equivalent to Gammat}),
%   % Section~\ref{subsec: phi Gamma to Galois}, if~$A$ is a
%   % finite~$\Zp$-algebra,
%   the theories of $(\varphi,\Gamma)$- and
%   $(\varphi,\Gammat)$-modules are equivalent. The theory
%   of~$(\varphi,\Gamma)$-modules will be more convenient for us  in
%   setting up the basic properties of our moduli stacks, but we will
%   find it convenient to return to the (equivalent) setting
%   of~$(\varphi,\Gammat)$-modules in Section~\ref{section: effective crystalline deformation rings}. %  with $A$-coefficients are equivalent,
%   % and both are equivalent to the category of representations of~$G_K$
%   % with $A$-coefficients.
%   % We will not make use of~$\Gammat$ in the rest
%   % of the paper.
%  % \mar{ME: If it turns out that we {\em do} need
% 		% to refer to $\tGamma$ later on, then this statement will
% 		% have to be modified.}
% \end{remark}




\subsection{The relationship with Galois representations}
\label{subsubsec:Galois reps}
 There is an equivalence of categories between the category of
 continuous representations of~$G_K$ on finite $\Zp$-modules, and the
 category of \'etale $(\varphi,\Gamma_K)$-modules, which is given by 
functors $\mathbb{D}$ and $T$ that are defined as
follows. Let~$\AKnrhat$ denote the $p$-adic
completion of the ring of integers
 of the
maximal unramified extension of~$\A_K[1/p]$ in~$W(\C^\flat)[1/p]$;
this is preserved by the natural actions of~$\varphi$ and~$G_K$
on~$W(\C^\flat)[1/p]$.
Then for a $G_K$-representation $T$, we define % \mar{TG: check here and
  % below if that should actually be $G_K$ invariants?!}
% \mar{TG: formula for
%   $D$ is incorrect and should use $\AKnrhat$ instead,
% will need to introduce that at some point and deal with the comment
% below at the same time.}
\[\mathbb{D}(T) :=(\AKnrhat
  \otimes_{\Zp}T)^{G_{\Kcyc}},\]
while for an \'etale $(\varphi,\Gamma_K)$-module $M$ we define
 \[T(M) :=(\AKnrhat \otimes_{\A_K}M)^{\varphi=1}.\]% \mar{TG:
% obviously this needs to be filled out and made correct.}
The action of~$G_{\Kcyc}$ (resp.\ $\varphi$) in the definition
of $\mathbb D$ (resp.\ of $T$)
is the diagonal one (with the $G_K$-action on~$M$
being that inflated from the action of~$\Gamma_K$). There is a
completely analogous theory of~$(\varphi,\Gammat_K)$-modules, and
taking~$\Delta_K$-invariants gives an
equivalence of categories between ~$(\varphi,\Gammat_K)$-modules
and~$(\varphi,\Gamma_K)$-modules. 

The functor $T(M)$ is also defined for finitely generated \'etale $\varphi$-modules
over~$\A_K$ (i.e.\ in the absence of a  $\Gamma$-action on~$M$), 
and in this context yields an equivalence of categories between
finitely generated  \'etale $\varphi$-modules over~$\A_K$
and continuous $G_{K_{\cyc}}$-representations on finite $\Z_p$-modules.

% for any \'etale $(\varphi,\Gamma)$-module $M$, there is an
% isomorphism
% \[W(\C^\flat) \otimes_{\A_K}M\isoto W(\C^\flat)
%   \otimes_{\Zp}T(M), \]which is equivariant for the actions
% of~$\varphi$, $G_K$, where~$\varphi$ has the diagonal action on the
% left hand side and acts on~$W(\C^\flat)$ on the right hand side,
% and~$G_K$ has the diagonal action on each side, with the action on~$M$
% being that inflated from the action of~$\Gamma$.  (In the literature
% it is more common to use a subring
% $\AKnrhat\subset W(\C^\flat)$ in place
% of~$W(\C^\flat)$, but the two definitions are equivalent
% by~\cite[Prop.\ B.1.8.3]{MR1106901}.)  \mar{ME: This citation is just
%   to a structural statement about the relationships between various
%   subrings of $W(\C^\flat)$.  Is there a citation that literally
%   states that the functors are the same?  Or is is something that just
%   has to be deduced from the cited proposition? TG: I don't know of
%   such a citation; see \cite[Lem.\ 2.2.1]{MR2745530} for example,
%   although I think there are other places in the literature that cite
%   it. Maybe we could write out the details although presumably to do
%   so we would have to actually introduce $\AKnrhat$,
%   and it will probably take me an hour to figure out the notation in
%   Fontaine, so I'm leaving this at least for now. I think it will be
%   an easy consequence of this statement about subrings, anyway, as the
%   key point is going to be that you don't get more $\varphi$ or
%   $G_{\Kcyc}$ invariants when you extend scalars.}

Suppose now that $L/K$ is a finite extension, and
suppose further that $K = K_{\cyc} \cap L$ (or, equivalently,
that the natural embedding  $\Gamma_L \hookrightarrow \Gamma_K$ induced
by the inclusion $K_{\cyc}\subseteq L_{\cyc}$ is an isomorphism).  
We will recall the description of the functor $\Ind_{G_L}^{G_K}$
in terms of $(\varphi,\Gamma)$-modules.
(We will use this description in Section~\ref{subsec:ram bound}.)

By construction, there are inclusions $\A_K \subseteq \A_L \subseteq \AKnrhat$,
which are unramified embeddings of DVRs.  The composite of these embeddings
is $G_K$-equivariant (where $G_K$-acts on $\A_K$ through
its quotient~$\Gamma_K$),
while the second is $G_L$-equivariant (where $G_L$-acts on
$\A_L$ through its quotient~$\Gamma_L$).   Regarding the second of these embeddings
as an $\A_K$-linear morphism of $\A_K$-modules,
it induces an $\AKnrhat$-linear surjection
$\AKnrhat\otimes_{\A_K} \A_L \to \AKnrhat.$
The source of this surjection has a diagonal action of $G_K$
(the $G_K$-action on $A_L$ being via its quotient $\Gamma_K \iso \Gamma_L$),
while the surjection itself is $G_L$-equivariant.
Thus it induces  a $G_K$-equivariant morphism
$$\AKnrhat\otimes_{\A_K} \A_L \to
\Ind_{G_L}^{G_K} \AKnrhat,$$
which is easily seen to be an isomorphism.
Using the description of the induction as a tensor product,
we can express this as an isomorphism
\numequation
\label{eqn:induction formula}
\AKnrhat\otimes_{\A_K} \A_L \iso \Z_p[G_K]\otimes_{\Z_p[G_L]} \AKnrhat.
\end{equation}

If $M$ is an \'etale $(\varphi,\Gamma)$-module over~$\A_L$,
then we let $M'$  denote $M$ regarded as an \'etale $(\varphi,\Gamma)$-module
over the subring $\A_K$ of~$\A_L$. (Recall again that we are assuming
$\Gamma_L \iso \Gamma_K$.)
The isomorphism~\eqref{eqn:induction formula}
then yields an isomorphism 
\nummultline
\label{eqn:induction via phi Gamma mods}
T(M') := (\AKnrhat \otimes_{\A_K} M')^{\varphi = 1}  =
\bigl((\AKnrhat\otimes_{\A_K} \A_L) \otimes_{\A_L} M\bigr)^{\varphi =1}
\\
= 
(\Z_p[G_K]\otimes_{\Z_p[G_L]} \AKnrhat \otimes_{\A_L} M\bigr)^{\varphi = 1}
=\Z_p[G_K]\otimes_{\Z_p[G_L]} (\AKnrhat \otimes_{\A_L} M\bigr)^{\varphi = 1}
\\
= \Ind_{G_L}^{G_K} T(M). 
\end{multline}
In short, induction of Galois representations corresponds to
the restriction of coefficients on the $(\varphi,\Gamma)$-module side.

The preceding discussion also applies in the context of \'etale
$\varphi$-modules (in the sense of Definition~\ref{defn: general phi module}).
Our assumption that $\Gamma_L \iso \Gamma_K$ implies that
$G_{K_{\cyc}}/G_{L_{\cyc}}  \to G_K/G_L$  is a bijection,
so that \eqref{eqn:induction formula} may also be interpreted as an isomorphism
$$
\AKnrhat\otimes_{\A_K} \A_L \iso \Z_p[G_{K_{\cyc}}]\otimes_{\Z_p[G_{L_{\cyc}}]}
\AKnrhat.
$$
Then,
if $M$ is a finitely generated \'etale $\varphi$-module over~$\A_L$,
and if $M'$ denotes its restriction to an \'etale $\varphi$-module over~$\A_K$,
the preceding isomorphism induces an isomorphism
\numequation 
\label{eqn:induction via phi mods}
T(M') = \Ind_{G_{L_{\cyc}}}^{G_{K_{\cyc}}} T(M). 
\end{equation}

We also need a variant of the preceding theory
which again follows from the results
of~\cite{MR1106901}, using the Kummer extension~$K_\infty/K$
introduced in
Example~\ref{ex:kummer} and Section~\ref{subsubsec: Kummer}; the integral version of this theory was first
studied by Breuil and Kisin, see~\cite[\S1]{KisinModularity}. % , and will
% be employed in our study of crystalline deformation rings.
% Fix a uniformiser~$\pi$ of~$K$, and fix a compatible
% system of~$p^n$th roots ~$\pi^{1/p^n}$ of~$\pi$ in~$\overline{K}$.
% Let~$K_\infty=K(\pi^{1/p^\infty})=\cup_nK(\pi^{1/p^n})$ be the
% infinite (non-Galois) extension of~$K$ generated by these $p$-power
% roots of~$\pi$. Set~$\gS=W(k)[[u]]$, with a~$\varphi$-semi-linear
% endomorphism~$\varphi$ determined by~$\varphi(u)=u^p$, and
% let~$\cO_{\cE}$ be the $p$-adic completion of~$\gS[1/u]$. There is a
% $\varphi$-equivariant embedding $\gS\into\Ainf$, sending~$u$ to the
% Teichm\"uller lift of the element~$\pi^\flat=(\pi^{1/p^n})_n\in\cO_\C^\flat$
% determined by our fixed choice of~$\pi^{1/p^n}$. This extends to an
% embedding~$\cO_{\cE}\into W(\C^\flat)$, and~$G_{K_\infty}$ acts
% trivially on~$\cO_{\cE}$.\mar{TG: once Section~\ref{subsubsec: Kummer}
% is written we presumably just refer to that instead.}
%By definition, an \'etale $\varphi$-module over~$\cO_\cE$ is a
%finitely generated $\cO_\cE$-module~$M$, equipped with a
%$\varphi$-linear morphism $\varphi:M\to M$ with the property that the
%corresponding morphism $\Phi_M:\varphi^*M\to M$ is an
%isomorphism.
Namely,
there is an equivalence of categories between the
category of continuous representations of~$G_{K_\infty}$ on finite
$\Zp$-algebras, and the category of \'etale
%\mar{ME: This citation was wrong! (I just now changed the label of Def.\ 2.7.2
%to be semantically correct, so now this citation is just broken.) TG:
%hopefully this is a correct citation! But perhaps better would be if
%the first sentence of the preceding paragraph had this citation -- in
%fact I have now made that change, assuming that the reader can
%remember this reference from one paragraph to the next. Comment out if
%it's OK now\dots (or feel free to add another citation here, of course) ME: Looks
%good to me, commenting out}
$\varphi$-modules %(in the sense of Definition~\ref{defn: general phi module})
over~$\cO_{\cE}$, which is given by functors
$\mathbb{D}_\infty$, $T_\infty$ that are defined as
follows: Let~$\cO_{\widehat{\cE^{\nr}}}$ denote the $p$-adic
completion of the ring of integers
in the
maximal unramified extension of~$\Frac(\cO_{\cE})$
in~$W(\C^\flat)[1/p]$;
this is preserved by the natural actions of~$\varphi$ and~$G_{K_{\infty}}$
on~$W(\C^\flat)[1/p]$.
%This has actions of~$\varphi$
%and~$G_{K_\infty}$, and we have
We define
\[\mathbb{D}_\infty(T) :=( \cO_{\widehat{\cE^{\nr}}} \otimes_{\Zp}T)^{G_{K_\infty}},\]
for a $G_{K_\infty}$-representation $T$,
and define
\[T_\infty(M):=(  \cO_{\widehat{\cE^{\nr}}}\otimes_{\A_K}M)^{\varphi=1},\]
for an \'etale $\varphi$-module $M$.
% \mar{TG:
  % again correct the period ring, following Kisin.}
% There is an isomorphism
% \[W(\C^\flat) \otimes_{\A_K}M\isoto W(\C^\flat)
%   \otimes_{\Zp}T_\infty(M), \]which is equivariant for the actions
% of~$\varphi$, $G_{K_\infty}$, where~$\varphi$ has the diagonal action
% on the left hand side and acts on~$W(\C^\flat)$ on the right hand
% side, and~$G_{K_\infty}$ has the diagonal action on the right side,
% and acts on~$W(\C^\flat)$ on the left hand side.

% \section{Continuity}
% This se% ction is a place holder for my thoughts on continuity.

% \section{\texorpdfstring{$(\varphi,\Gamma)$}{(phi,Gamma)}--modules
%   with coefficients}\label{subsec: phi gamma with coefficients}
% \mar{TG: can we just define all coefficient rings
  % to be completed tensor products (with respect to the $p$-adic
  % topology on the coefficients, and the $T$-adic on the coefficient
  % rings)? Should reference Jon Dee if we do.}
% Fix a finite extension~$E/\Qp$ with ring of integers~$\cO$ and residue
% field~$F$.
% We put ourselves in the setting of Section~\ref{subsec: EG stuff}, taking $q=p$, writing~$T$ for~$u$, and taking
% $\varphi(T)=(1+T)^p-1$.
% Then $\gS$ is endowed with
% Note that in the cases of $\OEA$ and
% $\gS_A$, we can regard ourselves as being in the setting of
% Section~\ref{subsec: EG stuff}; the same is true
% of~$\A_{K,A}$, $\A^+_{K,A}$ if $K$ is basic. In particular, Theorem~\ref{thm: C is a $p$-adic formal algebraic stack and related
%     properties} proves the basic properties of the moduli stacks of
%   finite \'etale $\varphi$-modules over~$\A_{K,A}$ or~$\OEA$, and of
%   finite height $\varphi$-modules over~$\A^+_{K,A}$ or~$\gS_A$ (provided again that~$K$ is
%   basic in the case of~$\A^+_{K,A}$). Most of the rest of this
%   section is devoted to extending this theory to the case
%   of~$(\varphi,\Gamma)$-modules over~$\A_{K,A}$ for arbitrary~$K$. We begin with a definition.
% We have the ring $\gS := W(k)[[T]]$,
% \mar{ME: I will use the notation $T$ rather
% 	than $u$ for the variable, to have some conformity with various
% 	$(\varphi,\Gamma)$-module sources.  Perhaps $X$ is better, 
% 	because this is also used in some $(\varphi,\Gamma)$-module
% 	literature, and is a bit more generic --- but we're probably going
% 	to need it for schemes, stacks, etc.}
% which is endowed with its usual Frobenius endomorphism $\varphi(T)=(1+T)^p-1,$
% as well as
% a continuous action of $\tGamma$ by $W(k)\otimes_{\Zp}A$-linear automorphisms commuting
% with $\varphi$; concretely, an element~$g\in\tGamma$ acts
% by \[g(1+T)=(1+T)^{\varepsilon(g)}, \]where

% $\varepsilon:\tGamma\to\Z_p^\times$ is the $p$-adic cyclotomic character.  (Here, continuity is with respect to the $(p,T)$-adic
% topology on $\gS$; more precisely, one checks without difficulty the that ideal $(p,T)$ is stablized
% by $\varphi,$ and on any of the quotients $\gS/(p,T)^n$, the action
% of $\tGamma$ factors through a finite quotient.) The action of 
% $\tGamma$ on $\gS$ induces a continuous action of $\tGamma$ on each
% of $\gS_A$ and $\OEA$.




% If $A$ is a $p$-adically complete $\Z_p$-algebra,
% then we write $\gS_A := W(k)\otimes_{\Z_p}A[[T]]$;
% this may equally well be regarded as the $T$-adic completion of
% $\gS\otimes_{\Z_p} A$.   We regard $\gS_A$ as a topological $A$-algebra
% ($A$ being endowed with its $p$-adic topology) by equipping
% it with its $(p,T)$-adic topology.
% We also consider the ring
% $\widehat{\gS_A[1/T]}$ obtained as the $p$-adic completion
% of $\gS_A[1/T]$, which we regard as a topological $A$-algebra
% by declaring $\gS_A$ (equipped with its $(p,T)$-adic topology)
% to be an open subalgebra.

% he endomorphism $\varphi$ of $\gS$
% induces an endomorphism $\varphi$ of $\gS_A$, which extends
% to an endomorphism $\varphi$ of $\widehat{\gS_A[1/T]}$.
% Similarly, t


% \mar{ME: I have allowed $p$-adically complete $\Z_p$-modules in this
% 	discussion,
% 	rather than just working over some $\Z_p/p^a$, in anticipation
% 	of passing to $p$-adic limits.  Probably a little more explanation
% 	is required, though, because in {\tt proper.tex} the notion
% 	of locally free rank $d$ \'etale  $\varphi$-modules is only
% 	defined for $\Z_p/p^a$-algebras; we don't take the limit.}

\subsection{Various types of $\varphi$-modules with coefficients}
We now introduce coefficients.
% \mar{ME: It looks like this is where we start writing $\Gamma$
% 	rather than $\Gamma_K$.  Maybe do this consisently throughout
% 	\S 3? ME: Actually, this switch officially happens after
% Cor.~\ref{cor:deducing properties}, so maybe just write $\Gamma_K$
% consistently up till that point?}

\begin{df}
	\label{def:etale phi Gamma modules}
	% A {\em projective \'etale $(\varphi,\tGamma)$-module of rank~$d$}
	% with $A$-coefficients is an \'etale $\varphi$-module $M$
	% over $A$ which is projective of rank~$d$, and is equipped
	% with a semi-linear action of $\tGamma$ which is furthermore
        % continuous 
	% % \mar{ME: Should recall that $M$ does have a canonical topology;
	% % 	probaby a non-existent Drinfeld citation will be involved.}
	% when $M$ is endowed with its canonical topology
       
	%Similarly,
	% \mar{ME: What follows is wrong; for $(\varphi,\Gamma)$-modules,
	% 	we need to make the coefficients smaller
	% 	too, by taking $\Delta$-invariants --- this involves introducing more notation at the beginning of the discussion.}
Let~$A$ be a $p$-adically complete $\Zp$-algebra.	A {\em
  projective \'etale $(\varphi,\Gamma_K)$-module of rank~$d$} \index{\'etale $(\varphi,\Gamma_K)$-module}
	with $A$-coefficients is a projective \'etale $\varphi$-module \index{\'etale $(\varphi,\Gamma_K)$-module}
        $M$ of rank~$d$
	over $\A_{K,A}$ equipped
	with a semi-linear action of $\Gamma_K$, which commutes
        with~$\varphi$, and which is furthermore continuous
	when $M$ is endowed with its canonical topology (i.e.\ the
        topology of Remark~\ref{rem:topologies on M}).
\end{df}

% The key idea is to 
% use our descent results to give another description of our
% various moduli stacks,
If $A$ is a finite
type $\Z/p^a$-algebra, we can use Theorem~\ref{thm:descending projective modules} to give
a useful alternative description of \'etale $(\varphi,\Gamma_K)$-modules with $A$-coefficients in terms of
\'etale $\varphi$-modules with
coefficients in the ring $W(\C^\flat)_A$, as we now explain.
%  (where $A$ is any finite
% type $\cO/\varpi^a$-algebra), a ring which contains all
% the other coefficient rings that we have considered. 

%Recall that
%(in Subsection~\ref{subsec: phi gamma over Zp})
%we have fixed a uniformiser~$\pi$ of~$K$,
%and defined the extension $K_{\infty} := K(\pi^{1/p^{\infty}})$.
%Let $E$ be the Eisenstein polynomial corresponding to~$\pi$.

% Let~$A$ be a finite type
% $\cO/\varpi^a$-algebra for some~$a\ge 1$. We set \mar{TG: this is a macro,
%   I'm not claiming this is good notation!}
% \[ \AAinf{A}:=(\Ainf\otimes_\Zp A)^{\wedge}\] where the completion is
% with respect to the usual $u$-adic topology on~$\Ainf$.\mar{TG: not
%   sure if we should write $u$ or $T$ or both here.} Note that~$\AAinf{A}$
% has commuting $A$-linear actions of~$\varphi$ and~$G_K$.\mar{TG: also
%   define $W(\C^\flat)_A$ here}

\begin{defn}%\mar{TG: we should put the word ``\'etale'' in throughout.}
  \label{defn: GK module over Ainf}Let~$A$ be a $p$-adically complete $\Zp$-algebra. % finite type
% $\cO/\varpi^a$-algebra for some~$a\ge 1$.
An \'etale \emph{$(\varphi,G_K)$-module \index{\'etale $(\varphi,G_K)$-module}
with $A$-coefficients} (resp.\ an \'etale
\emph{$(\varphi,G_{\Kcyc})$-module \index{\'etale $(\varphi,G_{\Kcyc})$-module}
with $A$-coefficients}, resp.\ an \'etale
\emph{$(\varphi,G_{K_\infty})$-module \index{\'etale $(\varphi,G_{K_\infty})$-module}
with $A$-coefficients}) is by definition a finitely
generated~$W(\C^\flat)_A$-module $M$ equipped with an isomorphism
% $\varphi$-semi-linear
\[\varphi_M:\varphi^*M\isoto M\]
 of $W(\C^\flat)_A$-modules, % \mar{TG: probably shouldn't say that,
% we linearised it!}
 and
a~$W(\C^\flat)_A$-semi-linear action of~$G_K$ (resp.\ $G_{\Kcyc}$,
resp.\ $G_{K_\infty}$), which is continuous % \mar{TG:
% explain in what sense}
and commutes with~$\varphi_M$.
% Similarly, a \emph{$(\varphi,G_{K_\infty})$-module with
%   $A$-coefficients} is by definition a finitely
% generated~$W(\C^\flat)_A$-module $M$ equipped with a
% $\varphi$-semi-linear isomorphism \[\varphi_M:\varphi^*M\isoto M,\] and
% a~$W(\C^\flat)_A$-semi-linear action of~$G_{K_\infty}$, which is
% continuous and commutes with~$\varphi_M$.
We say that~$M$ is projective if it is projective of constant rank as
a~$W(\C^\flat)_A$-module.
\end{defn}

If~$A$ is a finite type $\Z/p^a$-algebra, then
there is a functor from the category of finite projective \'etale
$(\varphi,\Gamma_K)$-modules % \mar{TG: I've gone with~$\Gamma_Kt$ here as I
  % think that should probably be what we want to inflate from, but if
  % that's correct then obviously we need to discuss this somewhere.}
with~$A$-coefficients to the category of finite projective
\'etale $(\varphi,G_K)$-modules with~$A$-coefficients, which takes an \'etale
$(\varphi,\Gamma_K)$-module~$M$ to $W(\C^\flat)_A\otimes_{\A_{K,A}}M$, endowed
with the extension of scalars of~$\varphi$, and the diagonal action
of~$G_K$, with the action of~$G_K$ on~$M$ being the action inflated
from~$\Gamma_K$. 

Similarly, there is a functor from the category of finite projective
\'etale $\varphi$-modules over~$\OEA$ % \mar{TG: made up notation for
  % $\A_K$ with $u$ rather than $T$}
to the category of finite projective
\'etale $(\varphi,G_{K_\infty})$-modules with~$A$-coefficients, which takes an \'etale
$\varphi$-module~$M$ to $W(\C^\flat)_A\otimes_{\OEA}M$, endowed
with the extension of scalars of~$\varphi$, and the diagonal action
of~$G_{K_\infty}$, with the action of~$G_{K_\infty}$ on~$M$ being the
trivial action.

% In this
% section we will apply the results of Section~\ref{subsec:coefficients} to the theories of \'etale
% $\varphi$-modules and \'etale $(\varphi,\Gamma_K)$-modules introduced in Sections~\ref{subsec:
% phi gamma over Zp} and~\ref{subsec: phi gamma with coefficients}. 

% To this end, recall that if~$K$ is a finite extension of~$\Qp$
% and~$\A$ is an $\cO$-algebra, we defined a coefficient ring~$\A_{K,A}$
% for \'etale $(\varphi,\Gamma_K)$-modules, which in the case that~$A$ is
% Artin local with finite residue field correspond to
% representations of~$G_K$ with $A$-coefficients. After fixing a
% uniformiser~$\pi$ of~$K$, and a compatible system~$(\pi^{1/p^n})$ of
% $p$-power roots of unity in~$\overline{K}$, we
% set~$K_\infty=K(\pi^{1/p^\infty})=\cup_nK(\pi^{1/p^n})$, and we
% defined a coefficient ring~$\OEA$. % for \'etale $\varphi$-modules
% % with $A$-coefficients.
% Again, in the case that~$A$ is
% Artin local with finite residue field, these correspond to
% representations of~$G_{K_\infty}$ with $A$-coefficients.

% We also have the perfections~$\tA_{K,A}$ and~$\tOEA$ of these
% coefficient rings, as considered in Section~\ref{subsec: perfect
%   etale phi modules}, where we showed that the theories of \'etale
% $\varphi$-modules over~$\A_{K,A}$ and~$\tA_{K,A}$ are equivalent (and
% similarly over~$\OEA$ and~$\tOEA$).


\begin{prop}
  \label{prop: equivalences of categories to Ainf}Let~$A$ be a
  finite type $\Z/p^a$-algebra for some~$a\ge 1$.
  \begin{enumerate}
  \item The functor $M\mapsto W(\C^\flat)_A\otimes_{\A_{K,A}}M$ is an
    equivalence between the category of finite projective \'etale
$\varphi$-modules over~$\A_{K,A}$ and the category of
finite projective $(\varphi,G_{\Kcyc})$-modules
with $A$-coefficients.


It induces  an
    equivalence of categories between the  category of finite projective \'etale
$(\varphi,\Gamma_K)$-modules with $A$-coefficients and the category of
finite projective \'etale $(\varphi,G_K)$-modules
with $A$-coefficients. 

A quasi-inverse functor is given by the
composite of $N\mapsto N^{G_{\Kcyc}}$ and a quasi-inverse to the functor
of Proposition~{\em \ref{prop:perfect equivalence}}.
\item The functor $M\mapsto W(\C^\flat)_A\otimes_{\OEA}M$ is an
    equivalence of categories between the  category of finite projective \'etale
$\varphi$-modules over~$\OEA$ and the category of
finite projective \'etale $(\varphi,G_{K_\infty})$-modules
with $A$-coefficients. A quasi-inverse functor is given by $N\mapsto
N^{G_{K_\infty}}$ and a quasi-inverse to the functor of  Proposition~{\em \ref{prop:perfect equivalence}}.
  \end{enumerate}

\end{prop}
\begin{proof}
%To begin with, suppose that $A$ is a finite type $\F := \O/\varpi$-algebra,
%so that $\A_{K,A} = \widehat{A\otimes_{\F_p} \F_p((T))}$ 
%while
%$W(\C^\flat)_A = \widehat{A\otimes_{\F_p} \C^\flat}.$
%Furthermore, there is an inclusion $\F_p((T)) \subset \C^\flat$
%which identifies the target with the completion of the algebraic closure
%of the source.
%
%We consider as well the intermediate field $\F_p((T^{1/p^{\infty}}))$,
%and the corresponding completed tensor product
%	$\Atilde_{K,A} := \widehat{A\otimes_{\F_p} \F_p((T^{1/p^{\infty}}))}.$
%Note that
%$\F_p((T^{1/p^{\infty}}))$
%is a perfectoid field, and that
%$\C^\flat$ is the completion of a Galois extension of
%$\F_p((T^{1/p^{\infty}}))$.  
%
%In order to apply the results of Section~\ref{sec:coefficients},
%we introduce more precise notation.
%Namely, we let $L$ denote the algebraic closure of
%$\F_p((T^{1/p^{\infty}}))$ in $\C^\flat$; then $L$ is
%an algebraic closure of 
%$\F_p((T^{1/p^{\infty}}))$.  
%In fact $L$ is also
%a separable closure of
%$\F_p[((T^{1/p^{\infty}}))$, since the latter field is perfect,
%being perfectoid. (This is proved by a successive approximation 
%argument, see e.g.\ \cite[Prop.~5.9]{ScholzePerfectoid}.)\mar{ME: Note that that
%	this entry in the bibliography has to be fully fleshed out.}
%Of course $\C^\flat$ is then equal to the completion of $L$.
%
%Write (as we may) $L = \bigcup_n L_n$ as an increasing union of finite
%Galois extensions of 
%$\F_p((T^{1/p^{\infty}}))$,
%and write $\cO_n$ to denote the ring of integers in $L_n$.
%Each of the finite extensions $L_n$ is again perfectoid
%(see e.g.\ \cite[Prop.~5.23]{ScholzePerfectoid}), and so each 
%of the extensions $\cO_m \subseteq \cO_n$ (for $m \leq n$) is 
%almost Galois, with Galois group $\Gal(L_n/L_m)$.
%(See \cite[Prop.~5.23]{ScholzePerfectoid} for one formulation of this statement,
%in the language of almost mathematics.)
%\mar{ME: Have to convert from the language of ``almost \'etale'' in 
%	the Faltings--Scholze sense to ``almost Galois'' in our 
%	naive sense; maybe do this back in \S 5 or 6, or in a separate lemma.
%	Also should probaby explain that we are just using \cite{ScholzePerfectoid} 
%	as a convenient reference, but that these results go back to
%	Tate, Faltings, Gabber--Ramero, etc.}
%Also, write $K := \F_p((T^{1/p^{\infty}}))$,  
%write
%$\cO := \F_p[[T^{1/p^{\infty}}]]$ (the ring of integers in $K$),
%and write 
%$\cO' := \cO_\C^\flat$ (the ring of integers in $\C^\flat$).
%\mar{ME: Note that the notation $K$ and $\cO$, which are being used here
%	to match with the notation of \S~\ref{sec:coefficients},
%	are in conflict with our running notation.  I'm not going to worry
%	about this now, but it should ultimately be fixed, by introducing
%	different notation there and here. Also, maybe change the
%notation for the uniformizer from $u$ to something more neutral,
%like $v$?}
%We are now in the situation of Section~\ref{sec:coefficients};
%note in particular that $\Atilde_{K,A}$ may be identified
%with the ring denoted there by $R[1/u]$, and that $W(\C^\flat)_A$
%may be identified with the ring denoted there by $S[1/u]$.
%Also the Galois group denoted there by $G$,
%which is the Galois group $\Gal(L/K)$,
%is equal (in our context) to $G_{K_{\cyc}}$.
%Applying Theorem~\ref{thm:descent with coefficients},
%we find that the functor $M \mapsto W(\C^\flat)_A \otimes_{\Atilde_{K,A}} M$
%induces an equivalence of categories between
%the category of finitely generated projective $\Atilde_{K,A}$-modules,
%and the category of finitely generated projective 
%$W(\C^\flat)_A$-modules endowed with a continuous semi-linear action of 
%$G_{K_{\cyc}}$.
%
%An identical argument shows that the functor
%$M\mapsto W(\C^\flat)_A\otimes_{\OEA} M$ 
%induces an equivalence of categories between
%the category of finitely generated projective $\OEA$-modules,
%and the category of finitely generated projective 
%$W(\C^\flat)_A$-modules endowed with a continuous semi-linear action of 
%$G_{K_{\infty}}$.
%
% \mar{ME: Finish argument by adding in \'etale $\varphi$-structures,
% 	and in the cyclotomic case, by adding in $\Gamma_K$.  Note that
% 	we might need to add something about the fact that the equivalence
% 	identifies topologies on Hom-spaces in Theorem~\ref{thm:descent
% 		with coefficients}, so that we can deduce
% 	that {\em continuous} $\Gamma_K$-actions match with
% 	{\em continuous} $G_K$-actions.  Also have to consider the
% case $a > 1$!}
%
%
%
%
%  \mar{TG: hopefully just a question of citing earlier descent
%    results. TG: added something from an email I sent a couple of
%    months ago, might be nonsense!}
%\begin{verbatim}
%I'm a bit confused about what we want for
%  almost etaleness, but maybe the version of almost purity on p22 of
%  these notes:
%
%http://www-personal.umich.edu/~stevmatt/perfectoid2.pdf
%
%(Thm 7.10) does what we want?!
%\end{verbatim}
%
  This is a formal consequence of Theorem~\ref{thm:descending
    projective modules} and Proposition~\ref{prop:perfect
    equivalence}.  % \mar{ME: Theorem~\ref{thm:descending projective
    %   modules} is now written for an arbitrary perfectoid subextension
    % $F$ of $\C$, so we have to explain that $K_{\infty}$ and
    % $K_{\cyc}$ are perfectoid. TG: and closed; add this in the
    % introductory material of section 2. TG: actually they aren't, I
    % guess, but rather their completions are! I guess the Galois groups
    % are the same, though, so I'm guessing this won't cause any real
    % problems.}
  We begin with~(2). Firstly, by
  Proposition~\ref{prop:perfect equivalence} (and Remark~\ref{rem: abstract perfection section linked to our concrete settings}) the functor
  $M\mapsto \tM=\tOEA\otimes_{\OEA}M$ is an equivalence of categories
  between the category of finite projective \'etale $\varphi$-modules
  over~$\OEA$ and the category of finite projective \'etale
  $\varphi$-modules over~$\tOEA$, so it is enough to show that
  $\tM\mapsto W(\C^\flat)_A\otimes_{\tOEA}\tM$ is an equivalence of
  categories. By Example~\ref{ex:kummer} and Remark~\ref{rem: Krasner
    Ax Tate Sen}, $\widehat{K}_\infty$ is a perfectoid field, and we
  have a canonical identification of Galois groups
  $G_{\widehat{K}_\infty}=G_{K_\infty}$. The result then follows from
  the equivalence of categories given by Theorem~\ref{thm:descending
    projective modules}, as we can think of $\varphi$ as being an
  isomorphism of $\tOEA$-modules $\varphi^*\tM\isoto\tM$.

The same argument shows in~(1) that the functor
$M\mapsto W(\C^\flat)_A\otimes_{\A_{K,A}}M$ is an equivalence of
categories between the category of projective \'etale
$\varphi$-modules over~$\A_{K,A}$, and the category of
projective~\'etale $(\varphi,G_{\Kcyc})$-modules with
$A$-coefficients. Since~$\Gamma_K=G_K/G_{\Kcyc}$, this extends to the
claimed equivalence of categories, noting that by construction the
continuity of the $\Gamma_K$-action on~$M$ is equivalent to the
continuity of the $G_K$-action on~$W(\C^\flat)_A\otimes_{\A_{K,A}}M$
(since
Lemma~\ref{lem: varphi extends to A and Galois etc}
shows  that  the $G_K$-action on~$W(\C^\flat)_A$ is continuous).
% because the equivalences of Theorem~\ref{thm:descending projective
%   modules} and Proposition~\ref{prop:perfect equivalence} respect the
% topologies on the $\Hom$ sets.\mar{TG: we need to say something about
%   this but presumably it's easy once we're thinking about it the right
% way? I haven't thought about it so I'm leaving it for now. TG: don't
% need something abstract, it's just clear looking at it.}
%\mar{ME: Flesh this out!}
\end{proof}

% \begin{lemma}\label{lem: phi Gamma equivalent to Gammat}\mar{TG: I don't know that we now need this; I think we
%     at most want it for finite~$\Zp$-algebras and then it's in the
%     literature. TG: rewrite for general coefficients, put it there}
% 	The functor $M \mapsto M^{\Delta}$ induces an equivalence
% 	between the category fibred in groupoids classifying $(\varphi,
% 	\tGamma)$-modules over $p$-adically complete $\Z_p$-algebras,
% 	and the category fibred in groupoids classifying
% 	$(\varphi,\Gamma_K)$-modules over $p$-adically complete $\Z_p$-algebras.
% \end{lemma}
% \begin{proof}
% 	Hilbert's Thm.\ 90. \mar{ME: Should flesh this out.}
% \end{proof}

% \mar{TG: check if
% some of this should actually be about~$(\varphi,\tGamma)$-

% In the remainder of this subsection, we explicate the continuity
% condition for an action of $\Gamma_{\disc}$ on an \'etale
% $\varphi$-module. % We begin with the following general lemma.





\chapter{Moduli stacks of \texorpdfstring{$\varphi$}{phi}-modules and 
  \texorpdfstring{$(\varphi,\Gamma)$}{(phi,Gamma)}-modules}\chaptermark{Stacks of \texorpdfstring{$\varphi$}{phi}-modules and 
  \texorpdfstring{$(\varphi,\Gamma)$}{(phi,Gamma)}-modules}\label{sec:
phi modules and phi gamma modules}% \mar{TG: I would like to move
% sections 3.1-3.3 back into section 2, so that this section is all
% about stacks.}

In this chapter we build on the results
of~\cite{EGstacktheoreticimages} (which in turn built
on~\cite{MR2562795}) to construct our stacks of
$(\varphi,\Gamma)$-modules, and various related stacks of
$\varphi$-modules. We show in particular that our stack~$\cX_{K,d}$ of
$(\varphi,\Gamma)$-modules is Ind-algebraic.

\section{Moduli stacks of \texorpdfstring{$\varphi$}{phi}-modules}\label{subsec: moduli stacks of phi
  modules}

%\mar{TG: finite height can now wait until we actually define the stacks}
In this section we put ourselves in the context of Situation~\ref{subsubsec:general
	framework};
we also remind the reader that the notion of \'etale $\varphi$-module is
defined in Section~\ref{subsec: EG stuff}.
%and furthermore fix a
%polynomial $F\in W(k)[T]$ which is congruent to a positive power
%of~$T$ modulo~$p$ (for example, an Eisenstein polynomial). 
We furthermore fix
a finite extension~$E/\Qp$ with ring of
integers~$\cO$ and residue field~$\F$; all of our coefficient rings
from now on will be $\cO$-algebras. All of our constructions are
compatible with replacing~$E$ by a finite extension, and we will
typically not comment on this, although see Remark~\ref{rem: changing O} below.
%\mar{ME: This introduction and fixing of $\cO$, $E$, and $\F$
%	should be more prominent --- maybe at the very beginning
%	of the section, before \S 3.1 begins? Or perhaps postponed,
%since I don't think we realy need it in the first few subsections? TG:
%I postponed it to here, I'm not sure what you want to do about making
%it more prominent though. ME: Seems fine here.}
% \mar{TG: we could have
%   something a bit more general if we wanted to, as
%   in~\cite{EGetalephimodules}. Also in that paper we would have a
%   notation clash if we used~$E$, as that's a field, but I'm leaving it
% alone here for now.}

If we fix integers $a,d \geq 1,$
then we may follow~\cite[\S
5]{EGstacktheoreticimages} and define
an \emph{fpqc} stack in groupoids $\cR_d^a$ over
$\Spec \cO/\varpi^a$ as follows: For any $\cO/\varpi^a$-algebra~$A$,
we define $\cR^a_d(A)$ to be the groupoid of
 \'etale $\varphi$-modules over~$\A_A$ which are projective % {\em fpqc} locally free
 of rank~$d$.
If $A \to B$ is a morphism of $\cO$-algebras, and $M$
is an object of $\cR_d^a(A)$,
then the pull-back of $M$ to $\cR_d^a(B)$
is defined to be the tensor product~$\AAA_B\otimes_{\AAA_A} M$.

A key point is that this definition does not require $\A_A^+$ to be $\varphi$-stable,
although (as far as we know) this hypothesis {\em is} required to make any deductions about
$\cR^a_d$ beyond the fact that it is an {\em fpqc} stack (which relies just
on Drinfeld's general descent results, as described
in \cite{MR2181808} and \cite[\S 5.1]{EGstacktheoreticimages}).

Since we are ultimately interested in questions of algebraicity or of Ind-algebraicity,
from now on we regard $\cR^a_d$ as an {\em fppf} stack over~$\cO/\varpi^a$. 
By~\cite[\href{http://stacks.math.columbia.edu/tag/04WV}{Tag
  04WV}]{stacks-project}, 
we may also regard the stack $\cR_d^a$ as an {\em fppf} stack over $\cO$,
and as $a$ varies, we may form the $2$-colimit $\cR := \varinjlim_a \cR^a_d$,
which is again an {\em fppf} stack over $\cO$.
% (although since we are ulimately
%concerned with the context of Ind-algebraic stacks, we will typically regard
%it as an {\em fppf} stack;  note that we obtain the same stack whether we take 
%the inductive limit over the {\em fpqc} or {\em fppf} site, since the
%stackification in the formation of the colimit in fact only involves
%the Zariski topology of~$\cO$).
In fact $\cR_d$ lies over $\Spf \cO := \varinjlim_a \Spec \cO/\varpi^a$,
the formal spectrum of $\cO$ with respect to the $\varpi$-adic,
or equivalently $p$-adic, topology.

We now fix a polynomial $F\in W(k)[T]$ which is congruent to a positive power
of~$T$ modulo~$p$ (for example, an Eisenstein polynomial). 

\begin{defn}% \mar{TG: here and below should be $p$-adically complete
    % $\Zp$-algebra~$A$ and replace $\Spec A$ with $\Spf A$?}
  \label{defn: phi module of finite E-height}% \mar{TG: if $A$ is not
  %   $p$-power torsion, I haven't thought about whether we will want to
  % modify this definition. In keeping with remarks elsewhere, I'm
  % imagining that we do want to work with $p$-adic coefficients from
  % the beginning wherever we can?} 
Suppose that~$\A_A^+$ is $\varphi$-stable. Let $h$ be a non-negative
integer, and let~$A$ be a $p$-adically complete $\cO$-algebra. A  \emph{$\varphi$-module of
   $F$-height at most~$h$ over~$\A^+_A$} \index{finite height $\varphi$-module}
% \mar{ME: Maybe call these {\em Breuil--Kisin--Wach-modules}? TG: have
%   introduced this terminology, not sure if we want to use it elsewhere
% in the paper so I'll leave this marginal comment here for now.}
is  a pair $(\gM,\varphi_M)$
consisting of a finitely generated $T$-torsion free
$\Aplus_A$-module~$\gM$, and a % \mar{TG: seems a bit ugly to have to say
% ``$T$-torsion free''; do we ever use non-projective Kisin modules
% now?! If not then we could just define it for projective rank~$d$. TG:
% we do use them in arguments with versal rings, so commenting this out}
$\varphi$-semi-linear map $\varphi_\gM:\gM\to\gM$, with the further
properties that if we
write \[\Phi_{\gM}:=1\otimes \varphi_{\gM}:\varphi^*\gM\to\gM,\] then % \mar{TG: I suspect that either I'm missing something
%   obvious or this isn't the right definition - it doesn't seem so good
% for passing mod~$p^n$, for example I get a bit tangled up proving
% Lemma~\ref{lem: map from finite E height to etale phi module}
% below. In the mod $p^a$-case I think this is OK though; originally in
% proper.tex we had that it was an isomorphism after inverting~$u$ and
% then deduced injectivity from that.}
$\Phi_{\gM}$ is injective, and the cokernel of $\Phi_{\gM}$ is killed by $F^h$.

A \emph{$\varphi$-module of finite $F$-height over~$\A^+_A$} is a $\varphi$-module of
  $F$-height at most~$h$  for some $h\ge 0$. A
  morphism of $\varphi$-modules is a morphism of the underlying
  $\Aplus_A$-modules which commutes with the morphisms~$\Phi_\gM$.

  We say that a $\varphi$-module of finite $F$-height is projective of
  rank~$d$ if it is a finitely generated projective~$\Aplus_A$-module of
  constant rank~$d$.
%  and if Zariski locally free on~$\Spf A$, it is finite free
 % over~$\Aplus_A$ of rank~$d$. % \mar{TG: I think this is the sensible
    % notion regardless of $A$ being $p$-power torsion, but I think that
    % won't (at least obviously) be the case for $A((u))$-modules, so
    % I'm putting a marginal comment in now to be alert to these things.}
\end{defn}


If we maintain the assumption that $\A^+$ is~$\varphi$-stable,
and if we fix integers $a,d\ge 1$ and an integer $h\ge 0$,
then we may again follow~\cite[\S
5]{EGstacktheoreticimages} to define an 
\emph{fpqc} stack in groupoids $\cC_{d,h}^a$
over $\Spec \cO/\varpi^a$ as follows: For any $\cO/\varpi^a$-algebra~$A$,
we define $\cC^a_{d,h}(A)$ to be the groupoid of $\varphi$-modules of
$F$-height at most $h$ over~$\A^+_A$ which are projective
% (or equivalently, by Proposition~\ref{prop:improving fpqc for A[[u]]},
% {\em fpqc})
 of rank~$d$.
If $A \to B$ is a morphism of $\Z_p$-algebras, and $\mathfrak M$
is an object of $\cC_{d,h}^a(A)$ 
then the pull-back of $\gM$ to $\cC_{d,h}^a(B)$
is defined to be the tensor product
$\Aplus_{B}\otimes_{\Aplus_A} \gM$.

% That $\cC^a_d$ and $\cR^a_d$ are indeed \emph{fpqc} stacks follows from
% the results of~\cite{MR2181808}, and more specifically
% from~\cite[Thm.\ 5.1.12]{EGstacktheoreticimages}. 
%By~\cite[\href{http://stacks.math.columbia.edu/tag/04WV}{Tag
%  04WV}]{stacks-project}, %\mar{TG: check what this says}
Just as for the stack $\cR_d^a$,
we may and do also regard the stack $\cC_{d,h}^{a}$ as an {\em fppf}
stack over $\cO$,
%\mar{ME: I'm a bit confused by this, partly b/c $\cO$ is a ring, rather than
%	a scheme.  Do we mean over $\Spec \cO$ (in which case I don't
%	think this is true), or maybe $\Spf \cO$? TG: I think we mean $\Spf\cO$}
and we then, allowing $a$ to vary, define $\cC_{d,h}:=\varinjlim_{a}\cC_{d,h}^a$,
obtaining an {\em fppf} stack over $\cO$
which in fact lies over $\Spf \cO$.
% \mar{TG: need to check that
%   for $A$ $p$-adic complete, $\cC_d(A)$ and $\cR_d(A)$ have the expected
% descriptions. Insert this as a lemma once freeness statements cleared
% up above.}
There are  canonical morphisms $\cC_{d,h}^a\to\cR_d^a$ and $\cC_{d,h}\to\cR_d$ given % by
% inverting~$u$,\mar{TG: should now say ``
  by tensoring with $\A_A$ over~$\A_A^+$. % '',
  % and don't need to define mod~$\varpi^a$ first?} and therefore
  % there is a canonical morphism~$\cC_{d,h}\to\cR_d$.

  \begin{rem}
    \label{rem: changing O}If~$E'/E$ is a finite extension with ring
    of integers~$\cO'$, then by definition we have (with obvious
    notation) $\cC_{d,h,\cO'}=\cC_{d,h}\times_{\cO}\cO'$ (in the case
that $\A^+$ is $\varphi$-stable, so that these stacks are defined) and
    $\cR_{d,\cO'}=\cR_{d,\cO}\times_\cO\cO'$ (in general). 
  \end{rem}

The following lemma provides a concrete interpretation of the
$A$-valued points of the stacks we have defined, when $A$ is a $\varpi$-adically
complete $\cO$-algebra (rather than just an algebra over some $\cO/\varpi^a$).

\begin{lem}
  If $A$ is a $\varpi$-adically complete $\cO$-algebra, then there
is a canonical equivalence between the groupoid of morphisms
$\Spf A\to\cR_{d}$ and the groupoid of rank~$d$
 \'etale $\varphi$-modules over~$\A_A$. 
If $\A_A^+$ is furthermore $\varphi$-stable, 
then there is a canonical equivalence between the groupoid of morphisms
  $\Spf A\to\cC_{d,h}$%  naturally biject\mar{TG: how should this be
%     phrased? Should we be writing $\cC_{d,h}(\Spf A)$? ME: Should state
% an equivalence of groupoids.} with
and the groupoid of $\varphi$-modules of rank~$d$ and $F$-height at most~$h$
 over~$\A^+_A$.
\end{lem}
\begin{proof}
  This is immediate from Lemma~\ref{lem: can check projectivity of Kisin and etale phi modules modulo
    p^n}.
\end{proof}

We now apply the results of \cite[\S 5]{EGstacktheoreticimages} to deduce
various results about the stacks we have introduced.  
This requires the assumption that $\A^+$ is $\varphi$-stable.

\begin{thm}
  \label{thm: C is a $p$-adic formal algebraic stack and related
    properties}Suppose that~$\A^+$ is $\varphi$-stable, and let $a\ge 1$ be arbitrary.
  \begin{enumerate}
  \item The stack $\cC_{d,h}^a$ is an algebraic stack of finite
    presentation over $\Spec\cO/\varpi^a$, with affine diagonal.
  \item The morphism $\cC_{d,h}^a\to\cR_d^a$ is representable by
    algebraic spaces, proper, and of finite presentation.
  \item The diagonal morphism
    $\Delta:\cR_d^a\to\cR_d^a\times_{\cO/\varpi^a}\cR_d^a$ is
    representable by algebraic spaces, affine, and of finite
    presentation.
  \item  $\cR^a_d$ is a limit preserving Ind-algebraic stack, whose diagonal is
    representable by algebraic spaces, affine, and of finite
    presentation. % \mar{TG: any reason we can't just add that $\cR$
%   is also a limit preserving Ind-algebraic stack, whose diagonal is
%     representable by algebraic spaces, affine, and of finite
%     presentation? I think it should be formal from this (?) and seems
%     to be what we actually use later. ME:  Just to check --- the point of
% your comment is $\cR$ vs.\ $\cR^a$, right?  If so, yes, we should just
% add this! TG: yes, exactly. Presumably we could/should rewrite this
% whole thing to be more about $\cC$ and $\cR$, in fact, explaining how
% the statements reduce to the $\cC^a$, $\cR^a$? A bit tired to try to
% think about this right now though. TG: maybe keep this and put in a
% corollary for $\cC$, $\cR$.}
  \end{enumerate}
\end{thm}
\begin{proof}
  Part (1) is~\cite[Thm.\ 5.4.9~(1)]{EGstacktheoreticimages}, and
  parts~(2), (3) and~(4) are proved in~\cite[Thm.\
  5.4.11]{EGstacktheoreticimages}, except for the claim that~$\cR^a_d$
  is Ind-algebraic, which is~\cite[Thm.\
  5.4.20]{EGstacktheoreticimages}. %   -(4) are proved in~\cite[\S
  % 5]{EGstacktheoreticimages}.\mar{TG: make this reference more precise
  % once that's finished.}
  %\mar{TG: need to give a citation for~(4)!}
  \end{proof}

\begin{cor}
  \label{cor: basic properties of C and R p adic stacks}Suppose
  that~$\A^+$ is $\varphi$-stable. 
  \begin{enumerate}
  \item $\cC_{d,h}$ is a $p$-adic formal algebraic stack of finite
    presentation over $\Spf\cO$, with affine diagonal.
 \item  $\cR_d$ is a limit preserving Ind-algebraic stack, whose diagonal is
    representable by algebraic spaces, affine, and of finite
    presentation.
  \item The morphism $\cC_{d,h}\to\cR_d$ is representable by
    algebraic spaces, proper, and of finite presentation.
  \item The diagonal morphism
    $\Delta:\cR_d\to\cR_d\times_{\Spf\cO}\cR_d$ is
    representable by algebraic spaces, affine, and of finite
    presentation.
   \end{enumerate}
\end{cor}
\begin{proof}
  The first part is immediate from~ Theorem~\ref{thm: C is a $p$-adic
    formal algebraic stack and related properties}~(1) and
  Proposition~\ref{prop: criterion for p-adic formal algebraic
    stack}. Everything else is immediate from Theorem~\ref{thm: C is a
    $p$-adic formal algebraic stack and related properties}.
\end{proof}
% \begin{proof}  If~$A$ is a $\Z/p^a\Z$-algebra for some $a\ge 1$,
%   then we set\[L^+G(A):=GL_d(\Aplus_A),\] \[LG^{h}(A):=\{X\in
%     M_d(\Aplus_A)|X^{-1}\in E(u)^{-h}M_d(\Aplus_A)\},\]and let $g\in L^+G(A)$ act
%   on the right on $LG^{h}(A)$ by $\varphi$-conjugation as $g^{-1}\cdot
%   X\cdot \varphi(g)$. Then we may write % \mar{TG: may need to elaborate
%     % slightly about the morphism to $\Spf\Zp$ versus the $\Nilp$
%     % perspective}
%   \[\cC^a_d=[LG^{h}/_\varphi L^+G].\]

%   For each $n\ge 1$ we have the principal congruence subgroup $U_n$ of
%   $L^+G$ given by $U_n(A)=I+u^n\cdot M_d(\Aplus_A)$. As in~\cite[\S
%   3.b.2]{MR2562795}, for any integer $n(a)>a(eh+q)/(q-1)$, we have a
%   natural identification
%   \numequation\label{eqn:U_n phi conjugacy}
% [LG^{h}/_\varphi U_{n(a)}]_{\Z/p^a\Z}\cong [LG^{h}/
%     U_{n(a)}]_{\Z/p^a\Z}\end{equation}where
%   the $U_{n(a)}$-action on the right hand side is by right
%   translation. (This boils down to the fact that the proof
%   of~\cite[Prop.\ 2.2]{MR2562795} extends essentially without change to
%   our more general setting, using Lemmas~\ref{lem: E(u) versus u mod p^a}
%   and~\ref{lem: phi(u) divisible by u}.) % \mar{TG: fwiw I actually
%     % checked this.}
%   This quotient stack is represented by a finite type
%   scheme $(X^{h}_{n(a)})_{\Z/p^a\Z}$, and we find that
%   \[\cC^a_h\cong
%     [(X^{h}_{n(a)})_{\Z/p^a\Z}/_\varphi(\cG_{{n(a)}})_{\Z/p^a\Z}],\]
%   where $(\cG_{{n(a)}})_{\Z/p^a\Z}=(L^+G/U_{n(a)})_{\Z/p^a\Z}$ is a
%   smooth finite type affine group scheme over $\Z/p^a\Z$. Thus~$\cC_d^a$ is
%   indeed an algebraic stack of finite presentation
%   over~$\Spec\Z/p^a\Z$ with affine diagonal, which proves~(1).

%  Point~(4) is then immediate from~(1) and Proposition~\ref{prop:
%    criterion for p-adic formal algebraic stack} (alternatively it can be proved
%  using the presentations we have exhibited for the~$\cC_d^a$,  as in
%  the proof of~\cite[Prop.\ 3.1.2]{CEGSKisinwithdd}). % We next prove~(4). Set
%   % \mar{ME: The general argument with the cotangent complex should show 
%   %         that both $\cC$ and $\cR$ are formal alg.\ stacks over $\Spf \Z_p$,
%   %         albeit in a less explicit manner.  Once that's written we should
%   %         probably just cite it, and perhaps refer to the argument in [CEGS]
%   %         to say that in the case of $\cC,$ one has a more explicit 
%   %         version of the argument. TG: agreed; this argument is
%   %         literally copied from there anyway.}
%   % $Y_a := [(X_{n(a)}^{h})_{\Z/p^a\Z}/_{\varphi}
%   % (U_{n(1)})_{\Z/p^a\Z}].$ If $a \geq b,$ then there is a natural
%   % isomorphism $(Y_a)_{\Z/p^b \Z} \cong Y_b$, so we may form the
%   % $p$-adic Ind-algebraic stack $Y := \varinjlim_a Y_a.$ Since
%   % $Y_1 := (X_{n(1)}^{h})_{\F_p}$ is a scheme, each $Y_a$ is in fact a
%   % scheme~\cite[\href{http://stacks.math.columbia.edu/tag/0BPW}{Tag
%   %   0BPW}]{stacks-project}, and thus $Y$ is a $p$-adic formal scheme.
%   % The natural morphism $Y \to \cC$ is then representable by algebraic
%   % spaces, smooth, and surjective, and so witnesses the claim that
%   % $\cC$ is a $p$-adic formal algebraic stack. That it is of finite
%   % type over~$\Spf\Zp$ with affine diagonal follows from~(1).

%   Parts~(2) and~(3) may be proved in almost the same way as~
%   \cite[Thm.\ 2.5, Cor.\ 2.6]{MR2562795}, the one wrinkle being that
%   we need to make repeated use of Corollary~\ref{cor: exponent of
%     phi}. Doing this, we find that if we replace the bounds
%   $(m+\delta)/(p-1)$ and $(m+n)/(p-1)$ given there by the quantities
%   $(aq+m+\delta)/(q-1)$ and $(m+n+aq)/(q-1)$, then the same
%   conclusions apply.  

% The claim that the diagonal of~$\cR_d^a$ is affine
%   is not explicitly verified in~\cite{MR2562795}, but it just amounts
%   to checking that the corresponding $\Isom$ sheaves are represented
%   by affine schemes. This can be checked \emph{fpqc}-locally, and
%   examining the argument, we see that \emph{fpqc}-locally each $\Isom$
%   sheaf is represented by a closed subscheme of the product of a copy
%   of the group scheme~$\GL_d$ and a finite number of copies of the
%   group scheme~$M_d$ of $d\times d$ matrices, which is affine, as
%   required.
% \end{proof}

% We now follow~\cite[\S 5.5]{EGstacktheoreticimages} very closely. Every result
% there remains true as stated with our more general definitions, so
% rather than repeating arguments or technical statements, we simply
% explain the few additional arguments needed to justify this. Firstly,
% \cite[Lem.\ 5.5.2]{EGstacktheoreticimages} remains true with the same
% proof (using Lemma~\ref{lem: E(u) versus u mod p^a} above). Then
% \cite[Lem.\ 5.5.3]{EGstacktheoreticimages}  goes through completely
% unchanged, and \cite[Lem.\ 5.5.4]{EGstacktheoreticimages} also has an
% identical proof (indeed, the argument makes no use of the
% ~$\varphi$-module structure), and~ \cite[Cor.\
% 5.5.5]{EGstacktheoreticimages} also goes through completely unchanged.


% We now need to briefly discuss the connection between \'etale
% $\varphi$-modules and Galois representations,\mar{TG: replace this all
% with some references to~\cite{schneiderPhiGamma}, and possibly weaken
% hypotheses accordingly.}
% following~\cite{MR1106901,MR2565906}. % \mar{TG: apologies in
%   % advance if I make some obvious misstatement here. I added in the
%   % assumption that $\F_q\subseteq k$ above, following Kisin--Ren; I have
%   % no idea if it's really needed, as I have thought a minimal amount
%   % about things (as we only need $q=p$) at the moment, but this means
%   % that I can cite their paper easily.}
% In~\cite{EGstacktheoreticimages} we made use of the connection to
% representations of the absolute Galois group of a $p$-adic field, but
% for our purposes here it suffices to use the (easier to state in our
% general setting) correspondence to representations of the absolute
% Galois group of a local field of characteristic~$p$.

% Let $G_{k((u))}$ be the absolute Galois group of~$k((u))$. Writing~
% $l$ for the unique subfield of~$k$ of cardinality~$q$, there is an
% equivalence of categories between the category of continuous
% representations of~ $G_{k((u))}$ on finite $W(l)$-modules, and the
% category of $p$-power torsion \'etale $\varphi$-modules with $\Z_p$-coefficients. In the case
% $q=p$, this is~\cite[Prop.\ A.1.2.6]{MR1106901}, and in the case $q>p$
% it is~\cite[Thm.\ 1.6]{MR2565906} (strictly speaking, this only
% considers the case that $\varphi(u)=u^q$, but the proof works
% unchanged in the general case). If~$A$ is a finite Artinian
% $\Zp$-algebra, this automatically gives an
% equivalence of categories between the category of \'etale
% $\phi$-modules with $A$-coefficients, and the category of continuous
% representations of $G_{k((u))}$ on finite
% $W(l)\otimes_{\Zp}A$-modules. % \mar{TG: this makes me a little
% %   uncomfortable, but I'm just going to write the most naive things I
% %   can and hope that everything works out in the end. I'm concerned
% %   that when $q>p$ we're now ending up looking at Galois representations valued in
% % products of local rings? Am I doing something wrong?}
% We say that a
% representation of~$G_{k((u))}$ has $E$-height at most~$h$ if the
% corresponding \'etale $\varphi$-module~$M$ is of the form $M=\gM[1/u]$ for
% some $\varphi$-module of height at most~$h$.

% Now given a finite type point $x:\Spec\F\to\cF$ as in~\cite[Prop.\
% 5.5.6]{EGstacktheoreticimages}, we let $V_{\F}$ be the corresponding
% $G_{k((u))}$-representation, and having fixed a basis, we can consider
% the framed deformation functor~$D$ for~$V_{\F}$ to Artinian local
% $W_a(\F)$-algebras. Exactly as in the proof of~\cite[Prop.\
% 5.5.6]{EGstacktheoreticimages}, this is pro-representable by an object  $R^{\univ}$ of
%   $\pro\cC_\Lambda$, and gives a versal morphism $\Spf
%   R^{\univ}\to\widehat{\cF}_x$ to~$x$. Then~\cite[Lem.\
%   5.5.7]{EGstacktheoreticimages} goes through unchanged.

% We then define the subfunctor~$D^{\le h}$ of~$D$ to be the liftings
% of~$V_\F$ of
% $E$-height at most~$h$. If~$V_{\F}$ has height at most~$h$, then the
% proof of~\cite[Thm.\ 1.3]{MR2782840} goes through essentially
% unchanged (as usual, using Lemmas~\ref{lem: E(u) versus u mod p^a}
% and~\ref{lem: phi(u) divisible by u} above) to show
% that~$D^{\le h}$ is pro-representable by a quotient~$R^{\le h}$
% of~$R^\univ$. 

% As in~\cite[\S 1.3]{MR2373358}, for each Artinian quotient~$A$
% of~$R^\univ$, the morphism $\cC^h\times_{\cF}\Spec A\to\Spec A$ is
% represented by a projective morphism $\mathcal{L}^{\le h}_A\to\Spec A$. (As
% ever, the argument only needs to be changed to use the bounds given by
% our Lemmas~\ref{lem: E(u) versus u mod p^a} and~\ref{lem: phi(u)
%   divisible by u}.) Then the proof of~\cite{EGstacktheoreticimages}
% goes through unchanged, showing that we have Noetherian versal rings.

% We make the exact same constructions as those made after the proof
% of~\cite[Prop.\ 5.5.10]{EGstacktheoreticimages}, and~\cite[Lem.\
% 5.5.12]{EGstacktheoreticimages} goes through unchanged. The proof
% of~\cite[Lem.\ 5.5.13]{EGstacktheoreticimages} is now valid in our
% general setting, except that we use our Lemma~\ref{lem: Noetherian
%   extension of scalars and bound on torsion} instead of~\cite[Lem.\
% 5.2.9]{EGstacktheoreticimages}, and our Proposition~\ref{prop:
%   faithfulness of truncation of homs II} instead of~\cite[Prop.\
% 5.2.8]{EGstacktheoreticimages}; this last point means that we need to
% replace every instance of~$u^{eah}$ by~$u^{M}$ for some~$M$ as in Proposition~\ref{prop:
%   faithfulness of truncation of homs II}, but the argument
% is otherwise unchanged.  Everything from~\cite[Lem.\
% 5.5.14]{EGstacktheoreticimages} goes through literally unchanged, and
% consequently we have proved the following theorem.

% \begin{thm}% \mar{TG: just guessing at what we actually want to state
%     % here - obviously there are other statements we could add if we
%     % want to.}
%   \label{thm: generalisation of main results of proper.tex}For
%   each~$a\ge 1$, $\cR^a$ is an Ind-algebraic stack, whose diagonal is
%   representable by algebraic spaces, affine, and of finite presentation.
% \end{thm}
% \begin{proof}
%   We have proved that~\cite[Thm.\ 5.4.5]{EGstacktheoreticimages} is
%   valid in our context, so~$\cR^a$ is an Ind-algebraic stack. The
%   properties of the diagonal were proved in Theorem~\ref{thm: C is a
%     $p$-adic formal algebraic stack and related properties}
%   above.% \mar{TG: deduce that~$\cR$ is a $p$-adic formal algebraic
% %     stack - not quite sure what the cleanest way is to write this? A
% %     diagonal argument? ME: $\cR$ is not $p$-adic, since mod $p$ its
% % not algebraic, or even formal --- just Ind.}
% \end{proof}

%For later reference, we recall a technical result that plays a role 
%in the proof of Theorem~\ref{thm: C is a $p$-adic formal algebraic stack and related properties}~(4).
%
%\begin{lemma}
%\label{lem:lifting from R to C}
%If $T \to \cR_d^a$ is a morphism
%whose source is a Noetherian scheme,
%then there is a Noetherian scheme $Z$ and a scheme-theoretically dominant and
%surjective morphism
%$Z\to T$ such that the composite $Z \to T \to \cR_d^a$ can be lifted
%to a morphism $Z \to \cC_{d,h}^a$ for some sufficiently large value of~$h$.
%\end{lemma}
%\begin{proof}
%In case $T$ is affine, this is shown in the course of
%proving~\cite[Thm.~5.4.20]{EGstacktheoreticimages}.
%In general,
%since $T$ is Noetherian, is is quasi-compact, and so admits a scheme-theoretically
%dominant surjection from a Noetherian affine scheme (e.g.\ the disjoint union
%of the members of a finite cover of $T$ by affine open subschemes).  Thus we are
%reduced to the affine case.
%\end{proof}

\section{Moduli stacks of
  \texorpdfstring{$(\varphi,\Gamma)$}{(phi,Gamma)}-modules}\label{subsec:
defn of Xd}% \mar{TG: by this point we should explain notation for the
% two different $\varphi$-actions.}
In this section we begin the study of our main objects of interest,
namely the moduli stacks of \'etale $(\varphi,\Gamma)$-modules.
 As in Chapter~\ref{section: coefficient rings} 
% and Subsection~\ref{subsec: phi gamma over Zp}
 we fix a finite extension $K/\Q_p$.  As in Section~\ref{subsec: moduli
 stacks of phi modules}, we also fix a finite extension $E$ of $\Q_p$
 with ring of integers $\cO$, which will serve as our ring of coefficients.
As always, $k$ denotes the residue field of the ring of integers of $K$,
and $\F$ denotes the residue field of~$\cO$.

% With this notation in hand, we can make the following definition.

%which is defined as
%follows. %\mar{TG: move this defn to the start of the section}

\begin{df} \label{defn: Xd}
	% \mar{ME: Should expand on this definition a bit, explain
	% 	why we get an {\em fpqc} stack (via Drinfeld),
        %         etc. TG: this is probably not what we want to write,
        %         but wanted to add something to get us moving again}
	We let $\cX_{K,d}$ denote the moduli stack of projective \index{$\cX_{K,d}$}
	\'etale $(\varphi,\Gamma_K)$-modules of rank~$d$. More
        precisely, if~$A$ is a $p$-adically complete $\cO$-algebra,
        then we define $\cX_{K,d}(A)$ (i.e., the groupoid of morphisms
%\mar{TG: ``groupoid'' could be ``groupoid of morphisms''? ME: Agreed and done, commenting out}
        $\Spf A\to\cX_{K,d}$) to be the groupoid of 
        projective \'etale~$(\varphi,\Gamma_K)$-modules  of rank~$d$ with~$A$-coefficients,
in the sense of Definition~\ref{def:etale phi Gamma modules},
 with morphisms given by isomorphisms. If $A\to B$ is a morphism of complete
        $\cO$-algebras, and $M$ is an object of $\cX_{K,d}(A)$, then the
        pull-back of~$M$ to~$\cX_{K,d}(B)$ is defined to be the tensor
        product $\A_{K,B}\otimes_{\A_{K,A}}M$. % \mar{TG: probably not as naive
          % as this?}

        It follows from the results of~\cite{MR2181808},
        and more specifically from~\cite[Thm.\
        5.1.16]{EGstacktheoreticimages}, that~$\cX_{K,d}$ is
        an~\emph{fpqc} stack over~$\cO$. As in Remark~\ref{rem:
          changing O}, the definition of~$\cX_{K,d}$ behaves
        naturally with respect to change of the coefficient ring~$\cO$.\end{df}

One of the main results of this book is that
$\cX_{K,d}$ is a Noetherian formal algebraic stack.  However,
the proof of this is quite involved, and will only be fully
achieved at the conclusion of Chapter~\ref{sec: families of extensions}.
In this section and the two that follow it, we establish
the preliminary result that $\cX_{K,d}$ is an Ind-algebraic stack.

We begin by discussing the moduli stacks of \'etale $\varphi$-modules
over~$\A_{K,A}$, which  will play an auxiliary role in our study 
of~$\cX_{K,d}$. 

\begin{df} \index{$\cR_{K,d}$}
We let~$\cR_{K,d}$ denote the moduli stack of rank $d$ projective \'etale
$\varphi$-modules, defined as in Section~\ref{subsec: moduli stacks of
  phi modules}, taking $\A$ to be $\A_K$.
% (Note that strictly speaking we only
% defined this stack in the case that~$\A_K^+$ is $\varphi$-stable, but
% the definition of an \'etale $\varphi$-module does not require this,
% and we again have an \emph{fpqc} stack over~$\cO$.)
\end{df}

If~$\A_K^+$ is
$\varphi$-stable, then Corollary~\ref{cor: basic properties of C and R
  p adic stacks}  applies to~$\cR_{K,d}$; this
is in particular the case if~$K/\Qp$ is basic in the sense of
Definition~\ref{defn: basic}. Our first task is to establish the same
results for general~$K$, which we will do by reducing to the basic
case.


\begin{defn}\label{defn: Kbasic} \index{$\Kbasic$}
  If~$K/\Qp$ is any finite extension, we
  set~$\Kbasic:=K\cap K_0(\zeta_{p^\infty})$.% \mar{TG: check/explain; here or later remark about how
    % this, or rather than $\Gamma$ version, means that we can check
    % continuity by reduction to the basic case.}
\end{defn}
By definition, $\Kbasic$ is basic, and~$(\Kbasic)_0=K_0$. Note that the natural
  restriction map~$\Gammat_K\to\Gammat_{\Kbasic}$ is an
  isomorphism, and induces an isomorphism
  $\Gamma_K\to\Gamma_{\Kbasic}$. By~\cite[Thm.~3.1.2]{MR719763}, $\A_K$ is a
  free~$\A_{\Kbasic}$-module of
  rank~$[K:\Kbasic]=[K(\zeta_{p^{\infty}}):K_0(\zeta_{p^\infty})]$, % \mar{TG: I think
%     this is true because we can construct the one from the other by
%     adjoining a root of a monic polynomial (by using Hensel's lemma
%     for a generator of the separable extension of residue fields); is
%     there a more conceptual reason? ME: Is the issue you're worried about
% free vs.\ projective?  Anyway, I'll leave this for now, because this
% is the kind of thing we should address when we set up all our 
% coefficient rings earlier on.}
and the inclusion $\A_{\Kbasic}\subset\A_K$ is $\varphi$-equivariant,
so there is a natural morphism $\cR_{K,d}\to\cR_{\Kbasic,d[K:\Kbasic]}$ given
by forgetting the $\A_{K,A}$-algebra structure on an \'etale
$\varphi$-module.

\begin{remark}
\label{rem:explanation of reduction to basic case}
As noted in Remark~\ref{rem: not phi stable T},
in order to ensure that $\A_K^+$ is $\varphi$-stable,
it suffices for $K$ to be abelian over $\Q_p$.  Thus,
rather than relating the theory for $K$ to that for the field $\Kbasic$ introduced above,
we could just as well relate it to any other subfield $K'$ of $K$ which is abelian
over $\Q_p$, and for which the natural map $\Gamma_K \to \Gamma_{K'}$ is an isomorphism; 
e.g.\ the field $K' := K \cap \Q_p^{\ab}$ (where $\Q_p^{ab}$ denotes
the maximal abelian extension of $K$).  
It doesn't matter (for our purposes) which particular $K'$ we choose; 
$\Kbasic$ is simply a convenient choice.
\end{remark}



\begin{lem}\label{lem: relative representability for R over Kbasic}
  The morphism $\cR_{K,d}\to\cR_{\Kbasic,d[K:\Kbasic]}$ is representable by
  algebraic spaces, affine, and of finite presentation.% \mar{TG: not
    % sure if we need it to be affine, and not totally sure it is, but
    % that seems to be what the proof shows.}
\end{lem}
\begin{proof}We can prove the statement after pulling back via a
  morphism $\Spec A\to \cR_{\Kbasic,d[K:\Kbasic]}$, where~$A$ is an
  $\cO/\varpi^a$-algebra for some~$a\ge 1$. This morphism corresponds to
  a projective \'etale $\varphi$-module over~$\A_{\Kbasic,A}$ of
  rank~$d[K:\Kbasic]$, and we need to show that the functor on
  $A$-algebras taking~$B$ to the set of projective \'etale $\varphi$-modules over~$\A_{K,B}$ of
  rank~$d$, whose underlying \'etale $\varphi$-module
  over~$\A_{\Kbasic,B}$ coincides with~$M_B$, is representable by an affine scheme of finite presentation over~$\Spec A$.

  Note that since~$M_B$ is projective and in particular flat
  over~$\A_{\Kbasic,B}$, we have a natural inclusion
  $i:M_B\into \A_{K,B}\otimes_{\A_{\Kbasic,B}}M_B$.  The additional
  structure needed to make~$M_B$ into a projective \'etale
  $\varphi$-module over~$\A_{K,B}$ is the data of a morphism of
  \'etale $\varphi$-modules over~$\A_{\Kbasic,B}$
  \[f: \A_{K,B}\otimes_{\A_{\Kbasic,B}}M_B\to M_B \] satisfying the
  %additional
  conditions that
\begin{enumerate}%\mar{TG: note that ME originally suggested that
%    rather than this first condition, we just want~$f$ to be
%    surjective - if that's right, then this argument would have to change.
%ME replies:  ME just forgot to make condition~(1) explicit :)  Commenting out!}
\item the composite $M_B\stackrel{i}{\to}
  \A_{K,B}\otimes_{\A_{\Kbasic,B}}M_B\stackrel{f}{\to} M_B $ is the
  identity morphism, and
\item the kernel of~$f$ is $\A_{K,B}$-stable.
\end{enumerate}
(We can then define the ~$\A_{K,B}$-module structure
on~$M_B$ via the formula $\lambda\cdot m:=f(\lambda\otimes
m)$. The first condition guarantees that this action is compatible
with the existing ~$\A_{\Kbasic,B}$-module structure
on~$M_B$, and the second condition that $(\lambda_1\lambda_2)\cdot
m=\lambda_1\cdot(\lambda_2\cdot m)$.)

By~\cite[Prop.\ 5.4.8]{EGstacktheoreticimages}, the data of a morphism
of \'etale $\varphi$-modules
$f: \A_{K,B}\otimes_{\A_{\Kbasic,B}}M_B\to M_B$ is representable by an
affine scheme of finite presentation over~$\Spec A$, so it is enough
to show that conditions~(1) and~(2) are closed conditions, given by
finitely many equations. To see this, we follow the proof
of~\cite[Prop.\ 5.4.8]{EGstacktheoreticimages}. Exactly as in that
argument, we can reduce to the case that~$M_B$ is free, and after
choosing bases, any~$f$ is determined by the coefficients of finitely
many powers of~$T$ in the Laurent series expansions of the entries of
the matrix given by~$f$. Condition~(1) is then evidently given by
finitely many equations in these coefficients.

To see that the same is true of condition~(2), note that since~$M_B$ is
projective and condition~(1) implies in particular that~$f$ is
surjective, we have a splitting
$\A_{K,B}\otimes_{\A_{\Kbasic,B}}M_B=M_B\oplus\ker(f)$. The projection
onto~$\ker(f)$ is given by $(1-i\circ f)$, so the condition
that~$\ker(f)$ is $\A_{K,B}$-stable is the condition that for
any~$\lambda$ in~$\A_{K,B}$, and any~$m\in M_B$, we
have \[f(\lambda(m-i(f(m))))=0. \]This is evidently a closed
condition, and since $\A_{K,B}$ and~$M_B$ are both finitely
generated~$\A_{\Kbasic,B}$-modules, it is determined by finitely many
equations, as required.
\end{proof}

\begin{cor}\label{cor: R good properties all K}
  %For any~$K$,
  The stack $\cR_{K,d}$ is a limit preserving
  Ind-algebraic stack, whose diagonal is representable by algebraic
  spaces, affine, and of finite presentation.
\end{cor}
\begin{proof}
	This follows from
	Lemma~\ref{lem: relative representability for R over Kbasic},
  Corollary~\ref{cor: basic properties of C and R p adic stacks}
(which establishes the claimed properties for $\cR_{\Kbasic, d[K:\Kbasic]}$),
and Corollary~\ref{cor:deducing properties} below.
\end{proof}% \mar{TG: add the explicit information that we can get by
  % taking scheme theoretic images of C pulled back from Kbasic.}
%The following lemma gives a concrete description of~ $\cR_{K,d}^a$ as an Ind-algebraic stack,
%and will be useful in Chapter~\ref{sec: the rank one case}.  In order to state it,
%we introduce the notation
%$\cC_{\Kbasic,d[K:\Kbasic],h}^a$ for the stack
%over~$\Spec\cO/\varpi^a$ classifying rank $d$ projective
%$\varphi$-modules over~$\A_{\Kbasic,A}^+$ of
%$T$-height at most~$h$; by Theorem~\ref{thm: C is a $p$-adic formal algebraic stack and related
%    properties}, $\cC_{\Kbasic,d[K:\Kbasic],h}^a$ is an algebraic
%stack of finite presentation over~$\Spec\cO/\varpi^a$.
%\begin{lem}
%  \label{lem: explicit Ind description for R pulled back from K basic}
%  Fix~$a \geq 1$, and for each $h \geq 0$,
%  let $\cR_{K,d,\Kbasic,h}^a$ denote the scheme-theoretic image of the
%  base-changed morphism
%  \[\cC_{d[K:\Kbasic],h}^a\times_{\cR_{\Kbasic,d[K:\Kbasic]}}\cR_{K,d}\to\cR_{K,d},\]
%  so that $\cR_{K,d,\Kbasic,h}^a$ is a closed algebraic substack of $\cR^a_{K,d}$.
%  Then the canonical morphism
%  $\varinjlim_h\cR_{K,d,\Kbasic,h}^a \to \cR_{K,d}^a$
%  is an isomorphism.
%\end{lem}
%\begin{proof}
%  Since a free~$\A_{K,A}$-module of rank~$d$ is in particular a free
%  $\A_{\Kbasic,A}$-module of rank~$d[K:\Kbasic]$, and since the
%  morphism $\cR_{K,d}\to\cR_{\Kbasic,d[K:\Kbasic]}$ is representable
%  by algebraic spaces by Lemma~\ref{lem: relative representability for
%    R over Kbasic}, this can be proved in exactly the same way
%  as~\cite[Thm.\ 5.4.20]{EGstacktheoreticimages}.\mar{TG: ME to check this?}
%\end{proof}
 
The following series of results concerning morphisms of stacks
culminates in Corollary~\ref{cor:deducing properties},
which was used in the proof of Corollary~\ref{cor: R good properties all K}.

\begin{lem}
  \label{lem:finiteness of diagonals}
  Let $\cX\to\cY$ be a morphism of stacks
  over a base scheme~$S$, which is representable by algebraic
  spaces. 
%  stacks {\em (}or {\em algebraic}, in the language
%  of~\cite[\href{http://stacks.math.columbia.edu/tag/06CF}{Tag
%  06CF}]{stacks-project}{\em )}.  
   As usual, let $\Delta_f: \cX \to  \cX\times_{\cY} \cX$
	  denote the diagonal of~$f$.
%
%  \begin{enumerate}
% 	  \item If $f$ is locally of finite type,
% 		  then
% 		  $\Delta_f$
% 		  is locally of finite presentation.
% 	  \item
If $f$ is of finite type and quasi-separated,
 		  then  $\Delta_f$ is of finite presentation.
%  \end{enumerate}
\end{lem}
   \begin{proof}
	   This can be checked after
	   pulling back along an arbitrary morphism $T \to \cY$,
	   where $T$ is a scheme, and hence reduced to the case
	   of a morphism from an algebraic space to a scheme.  In this case,
	   the claim of the lemma is proved
  in~\cite[\href{http://stacks.math.columbia.edu/tag/084P}{Tag
  084P}]{stacks-project}.  %This proves the lemma.
   \end{proof}

  
  \begin{lem}
  \label{lem: abstract pullback diagonals lemma}
  Let $f:\cX\to\cY$ be a
  morphism of stacks over a base scheme~$S$ which
  is representable by
  algebraic spaces, has affine diagonal, and is of finite type.
  Suppose that the
  diagonal of~$\cY$ is representable by algebraic
  spaces, affine, and of finite presentation. Then the
  diagonal of~$\cX$ is also representable by algebraic
  spaces, affine, and of finite presentation.
  %\mar{TG: might well need
  %  additional hypotheses and conclusions, needs a proof. Am wondering
  %if some faithfulness hypothesis for restriction of Homs could be
  %useful, which will hold in all our examples (where the morphisms are
%forgetful maps). ME: Now proved.}
\end{lem}
\begin{proof}
 % The diagonal of~$\cX$ is representable by algebraic spaces and
 % locally of finite presentation by~\cite[Prop.\
 % 2.3.20]{EGstacktheoreticimages}.
  We may factor the diagonal of~$\cX$
  as \[\cX\to\cX\times_{\cY}\cX\to\cX\times_{S}\cX.\]Since the
  morphism $\cX\times_{\cY}\cX\to\cX\times_{S}\cX$ is pulled back from
  the diagonal $\cY\to\cY\times_S\cY$, it is enough to show that
  the relative diagonal
  $\Delta_f: \cX\to\cX\times_{\cY}\cX$ is representable by algebraic spaces,
  affine, and of finite presentation.
 %  \mar{TG: didn't get any further than this.}

  Now $\Delta_f$ is representable by algebraic spaces (since $f$ is),
  and affine (by assumption).  Since affine morphisms
  are quasi-compact, we see that $f$ is quasi-separated, as well as
  being of finite type (by assumption); thus $\Delta_f$ is 
  of finite presentation (by Lemma~\ref{lem:finiteness
	  of diagonals}).  \end{proof} 

\begin{cor}\label{cor:deducing properties}
	Let $\cX \to \cY$ be a morphism of stacks which
	is representable by algebraic spaces,
	affine, and of finite presentation.
	If $\cY$ is a limit preserving Ind-algebraic stack,
	whose diagonal is representable by algebraic spaces,
	affine, and of finite presentation,
	then the same is true of $\cX$.
\end{cor}
\begin{proof}
The claimed limit preserving property of $\cX$ follows,
by \cite[Lem.~2.3.20~(3)]{EGstacktheoreticimages},
from that of $\cY$
and the fact that 
$\cX\to \cY$
is of finite presentation,
while the claimed Ind-algebraic property of $\cX$ %$\cR_{K,d}$
follows from that of $\cY$
and the fact that 
$\cX\to \cY$
is representable by algebraic spaces. % \mar{TG: should be citing
  % something from \cite{Emertonformalstacks} for this? which should be
  % put into the appendix if it isn't there already Yes, \cite[Rem.\
  % 4.9]{Emertonformalstacks} should be added to the appendix}
Finally,
the claimed properties of the diagonal of $\cX$
follow from those of $\cY$
by Lemma~\ref{lem: abstract pullback diagonals lemma}.
(Note that by assumption the hypotheses of
Lemma~\ref{lem: abstract pullback diagonals lemma}
hold, since morphisms of finite presentation are in particular of finite type,
while affine morphisms are separated, and thus have their diagonals
being closed immersions, which in particular are again affine.)
%\mar{TG: revise when that lemma exists. ME: Done}
\end{proof}

We now give a concrete description of~$\cR_{K,d}^a$ as an Ind-algebraic stack
which will be useful in Chapter~\ref{sec: the rank one case}.  In order to state it,
we introduce the notation
$\cC_{\Kbasic,d[K:\Kbasic],h}^a$ for the stack
over~$\Spec\cO/\varpi^a$ classifying rank $d$ projective
$\varphi$-modules over~$\A_{\Kbasic,A}^+$ of
$T$-height at most~$h$; by Theorem~\ref{thm: C is a $p$-adic formal algebraic stack and related
    properties}, $\cC_{\Kbasic,d[K:\Kbasic],h}^a$ is an algebraic
stack of finite presentation over~$\Spec\cO/\varpi^a$.

We begin with a lemma which is a variant of one of the steps appearing 
in the proof of~\cite[Thm.~5.4.20]{EGstacktheoreticimages}.

\begin{lemma}
\label{lem:lifting from R to C}
If $T \to \cR_{K,d}^a$ is a morphism
whose source is a Noetherian scheme,
then there is a Noetherian scheme $Z$ and a scheme-theoretically dominant and
surjective morphism
$Z\to T$ such that the composite $Z \to T \to \cR_{K,d}^a 
\to \cR^a_{\Kbasic,d[K:\Kbasic]}$ can be lifted
to a morphism $Z \to \cC_{\Kbasic,d[K:\Kbasic],h}^a$
for some sufficiently large value of~$h$.
\end{lemma}
\begin{proof}
Since $T$ is Noetherian, is is quasi-compact, and so admits a scheme-theoretically
dominant surjection from a Noetherian affine scheme (e.g.\ the disjoint union
of the members of a finite cover of $T$ by affine open subsets).  Thus we are
reduced to the affine case.

%\mar{ME: Need to fix this part of the argument}
Since $T$ is affine, by~\cite[Prop.~5.4.7]{EGstacktheoreticimages},
we may find a scheme-theoretically dominant surjection $Z  = \Spec A\to T$
such that the \'etale $\varphi$-module $M$ over $A$ corresponding
to the composite $Z \to T \to \cR_{K,d}$ is free (of rank $d$)
over $\A_{K,A}$ (rather
than merely projective).  Thus $M$ is also free (of rank $d[K:\Kbasic]$) 
over $\A_{\Kbasic,A}$.  We may then choose a $\varphi$-stable
$\A_{\Kbasic,A}$-basis
of~$M$.  If we let $\gM$ denote the $\A_{\Kbasic,A}^+$-module
spanned by this basis, then $\gM$ is $\varphi$-invariant, and has some height~$h$.
Thus $\gM$ induces the desired morphism $Z \to \cC_{\Kbasic,d[K:\Kbasic],h}^a.$
\end{proof}

We now give the promised description of the Ind-algebraic stack structure on~$\cR_{K,d}$.

\begin{lem}
  \label{lem: explicit Ind description for R pulled back from K basic}
  Fix~$a \geq 1$, and for each $h \geq 0$,
  let $\cR_{K,d,\Kbasic,h}^a$ denote the scheme-theoretic image of the
  base-changed morphism
  \[\cC_{d[K:\Kbasic],h}^a\times_{\cR_{\Kbasic,d[K:\Kbasic]}}\cR_{K,d}\to\cR_{K,d},\]
  so that $\cR_{K,d,\Kbasic,h}^a$ is a closed algebraic substack of $\cR^a_{K,d}$.
  Then the canonical morphism
  $\varinjlim_h\cR_{K,d,\Kbasic,h}^a \to \cR_{K,d}^a$
  is an isomorphism.
\end{lem}
\begin{proof}
  We have to show that any morphism $T \to \cR_{K,d}^a$ whose source
is a scheme factors through
the inductive limit.  Since $\cR_{K,d}^a$ is limit preserving, we may assume
that $T$ is of finite type over $\cO/\varpi^a$.  Let $Z\to T$
be as in the statement of Lemma~\ref{lem:lifting from R to C}.
Then the composite
\numequation
\label{eqn:dominant composite}
Z\to T \to \cR^a_{K,d}
\end{equation}
 lifts to a morphism
$Z \to \cC_{d[K:\Kbasic],h}^a\times_{\cR_{\Kbasic,d[K:\Kbasic]}}\cR_{K,d},$
and hence the morphism~\eqref{eqn:dominant composite} 
factors through 
  $\cR_{K,d,\Kbasic,h}^a$.
Since $Z \to T$ is scheme-theoretically dominant, the original morphism
$T \to \cR_{K,d}^a$ also factors through $\cR_{K,d,\Kbasic,h}^a$,
and hence through the inductive limit.
Thus the lemma is proved.
(We remark that this argument is essentially identical to that used to 
prove~\cite[Thm.~5.4.20]{EGstacktheoreticimages}.)
%  Since a free~$\A_{K,A}$-module of rank~$d$ is in particular a free
%  $\A_{\Kbasic,A}$-module of rank~$d[K:\Kbasic]$, and since the
%  morphism $\cR_{K,d}\to\cR_{\Kbasic,d[K:\Kbasic]}$ is representable
%  by algebraic spaces by Lemma~\ref{lem: relative representability for
%    R over Kbasic}, this can be proved in exactly the same way
%  as~\cite[Thm.\ 5.4.20]{EGstacktheoreticimages}.%\mar{TG: ME to check this? ME: Looks good}
\end{proof}



%Because of the equivalence between $(\varphi,\Gamma_K)$-modules and
%Galois representations with~$\Zp$-coefficients, % recalled in Section~\ref{subsec: phi Gamma to Galois}
%and because of the traditional notation $\rho$ for Galois
%representations, we will often denote a family of rank $d$
%projective \'etale
%$(\varphi,\Gamma_K)$-modules over~$T$, i.e.\ a morphism
%$T \to \cX_d$, by $\rho_T$, or some similar notation.    We caution the
%reader that this notation is chosen purely for its psychological 
%suggestiveness; a family of $(\varphi,\Gamma_K)$-modules over a general
%base $T$ does not admit a literal interpretation in terms of Galois
%representations.   
%
%
%
%
%We can use the equivalence between Galois representations
%and $(\varphi,\Gamma_K)$-modules as a tool to deduce facts about
%the finite type points of $\cX_{K,d}$.  More precisely,
%if~$\F/\F_p$ is a finite extension, then the groupoid of points
%$x\in \cX_{K,d}(\F)$ is canonically equivalent to the groupoid of
% Galois representations $\rho: G_K \to \GL_d(\F)$. % , whose
%% structure is fairly well understood (as we recall below).
%For 
%this reason, we will often denote such a point $x$ simply
%by the corresponding Galois representation~$\rho$.
%% \mar{ME: Probably should
%% 	rewrite this and the above as in the introduction: describe
%% 	it an an actual equivalence of groupoids.}



We now turn to studying $\Gamma_K$-actions on our $\varphi$-modules.
To ease notation, write~$\Gamma=\Gamma_K$ from now on. 
We choose
a topological generator $\gamma$ of $\Gamma$, and let $\Gamma_{\disc} :=
\langle \gamma \rangle;$ so $\Gamma_{\disc} \cong \Z$.
\begin{remark}
	\label{rem:discrete}
Since $\Gamma_{\disc}$ is dense in $\Gamma$, while a projective \'etale $\varphi$-module $M$ is complete with respect
to its canonical topology, in order to endow $M$ with the structure
of an \'etale $(\varphi,\Gamma)$-module, it suffices to equip $M$ with
a continuous action of $\Gamma_{\disc}$ (where we equip  $\Gamma_{\disc}$ 
with the topology induced on it by $\Gamma$).
\end{remark}

% \mar{ME: Maybe this discussion has to be amplified slightly, but 
% 	I'd rather not introduce any special notation for this stack, because
% 	(a) it just ends up being $\cR^{\Gamma_{\disc}},$ and (b)
% 	it looks a bit crappy to have the definition of our main
% 	object be followed by a formal definition of an object
% 	of only very temporary significance.}
% To this end we note that we may
% also
In order to study the properties of $\cX_{K,d}$ we will take advantage
of Remark~\ref{rem:discrete}. Accordingly, we now consider the moduli stack of 
projective \'etale $\varphi$-modules of rank~$d$ equipped with
a semi-linear action of $\Gamma_{\disc}$. % \mar{TG: not bothering to put
% a $K$ into this notation for now}
% \mar{TG: $\Gamma_{\disc}$?}
We don't introduce particular
notation for this stack, since the following proposition
identifies it with a fixed point stack $\cR_{K,d}^{\Gamma_{\disc}}$, which
we now define.

There is a canonical action of~$\Gamma_{\disc}$ on~$\cR_d$ (that is, a canonical morphism $\gamma:\cR_d\to\cR_d$):
 if~$M$ is an object of~$\cR_d(A)$, then $\gamma(M)$ is given
by $\gamma^*M:=\A_{K,A}\otimes_{\gamma,\A_{K,A}}M$. (Note that this is naturally a
$\varphi$-module, because the action of~$\gamma$ on~$\A_{K,A}$ commutes
with~$\varphi$.)  Then we set % \mar{TG: in this context, $\Gamma_\gamma$
%   for the graph of~$\gamma$ looks pretty bad, but sticking with it for
% now.}
\[\cR_d^{\Gamma_{\disc}}:=
\cR_d\underset{\Delta,\cR_d\times\cR_d,\Gamma_\gamma}{\times}\cR_d, \]where
 $\Delta$ is the diagonal of~$\cR_d$ and $\Gamma_\gamma$ is the
graph of~$\gamma$, so that $\Gamma_\gamma(x)=(x,\gamma(x))$. 

% We recall from \cite{EGstacktheoreticimages}
% that there is an Ind-algebraic stack 
% $\cR_d$ classifying rank $d$ locally free \'etale $\varphi$-modules.\mar{ME: Have to decide at some point whether to be $p$-adic from the beginning, or
% 	to work modulo some power of $p$.  Since everything is Ind anyway,
% 	maybe it's just simpler to work $p$-adically from the start, and
% 	the various reductions mod powers of $p$ can just be built
% 	into the various arguments dealing with the Ind structure? If we
% do this we should put ourselves in the situation of working over $\Spf \Z_p$
% here.  And then we have to make sure that our
% algebraic stack statements in \S 2 allow $\Spf \Z_p$ as a base \dots .}

\begin{prop}
	\label{prop:Gamma-disc modules as fixed points; etale case}
	% There is an action of $\Gamma_{\disc}$ on the Ind-algebraic stack $\cR_d$,
	% with the property that t
        The moduli stack of projective
	\'etale $\varphi$-modules of rank~$d$ equipped with a semi-linear action
	of $\Gamma_{\disc}$ is isomorphic to the fixed point stack
	$\cR_d^{\Gamma_{\disc}}.$
\end{prop}
\begin{proof}Using the usual construction of the 2-fibre product, we
  see that $\cR_d^{\Gamma_\disc}$ consists of tuples
  $(x,y,\alpha,\beta)$, with $x,y$ being objects of~$\cR_d$, and
  $\alpha:x\isoto y$ and~$\beta:\gamma(x)\isoto y$ being
  isomorphisms. This is equivalent to the category fibred in groupoids
  given by pairs~$(x,\iota)$ consisting of an object~$x$ of~$\cR_d$
  and an isomorphism~$\iota:\gamma(x)\isoto x$. Thus an object
  of~$\cR_d^{\Gamma_\disc}(A)$ is a projective
  \'etale $\varphi$-module of  rank~$d$ with $A$-coefficients~$M$,
  together with an isomorphism of $\varphi$-modules
  $\iota:\gamma^*M\isoto M$; but this isomorphism is precisely the
  data of a semi-linear action of~$\Gamma_\disc=\langle\gamma\rangle$
  on~$M$, as required.
\end{proof}
Since~$\cR_d$ is an Ind-algebraic stack, so
is~$\cR_d^{\Gamma_\disc}$. More precisely, we have the following
lemma.

\begin{lem}\label{lem: good properties R Gamma disc}$\cR^{\Gamma_{\disc}}_d$ is a limit preserving
  Ind-algebraic stack, whose diagonal is representable by algebraic
  spaces, affine, and of finite presentation.
  \end{lem}
\begin{proof}%\mar{TG: needs completing, probably need to unwind the
    %fibre product in the definition of $\cR^{\Gamma_{\disc}}_d$?
% ME: Done}
The description of $\cR^{\Gamma_{\disc}}_d$ as a fibre product
shows that the forgetful morphism $\cR^{\Gamma_{\disc}}_d \to \cR_d$
(given by forgetting the $\Gamma_{\disc}$-action)
is representable by algebraic spaces, affine, and of finite presentation,
since these properties hold for the diagonal of $\cR_d$,
by Corollary~\ref{cor: R good properties all K}.
The claim of the lemma is now seen to follow via another 
application 
of Corollary~\ref{cor: R good properties all K},
together with Corollary~\ref{cor:deducing properties}.
\end{proof}

 %  \mar{TG: do we need a discussion about fibre
% products of Ind-algebraic stacks?}
% Corollary~\ref{cor:fixed points preserve Ind-algebraicity}
% shows that $\cR_d^{\Gamma_{\disc}}$ is an Ind-algebraic stack over
% $\Spf \Z_p$.
% \mar{TG: is it even formal algebraic? ME: I don't think so.}
Restricting the $\Gamma$-action on an \'etale $(\varphi,\Gamma)$-module
to $\Gamma_{\disc}$ (and taking into account Proposition~\ref{prop:Gamma-disc modules as fixed points; etale case}),
we obtain a morphism $\cX_{K,d} \to \cR_d^{\Gamma_{\disc}},$
which by Remark~\ref{rem:discrete} is fully faithful. 
Thus $\cX_{K,d}$ may be regarded as a substack of
$\cR_d^{\Gamma_{\disc}}$.


As already noted, 
the first step in proving that
$\cX_{K,d}$ is a Noetherian formal algebraic stack is to
show that it is an Ind-algebraic stack.
Although $\cX_{K,d}$ is a substack of the Ind-algebraic stack
$\cR_d^{\Gamma_{\disc}}$,
it is not a closed substack, but should rather be thought of as a
certain formal neighbourhood of $\cX_{d,\red}$ in $\cR_d^{\Gamma_{\disc}}$ (see Remark~\ref{rem: rank one X not closed substack}),
% \mar{ME: Ultimately, we'll probably want to refer to something more elaborate
% 	in App.~E.}
% \mar{ME: Do we have an example or argument
    %     showing that it truly is not closed in
    %     $\cR_d^{\Gamma_{\disc}}$? TG: I don't know - I think you wrote
    %   this and I was planning to ask you about this at some point! TG:
    % it would be good to sort this out; I got confused thinking about
    % the rank one case. The topological nilpotence modulo~$p$ suggests
    % that maybe it's ``ind-closed'' rather than closed? TG: remark has begun.}
and so even this statement will require additional work to prove. We
begin with the following lemma, which allows us to reduce to the case
that~$K$ is basic.
% \mar{ME: Related --- maybe show $\cX_d$ is equal to the formal neighbourhood
% 	of $\cX_{d,\red}$ in $\cR_d^{\Gamma_{\disc}}$. Actually, probably
% automatic:  the inclusion of $\cX_d$ in this formal n.h.\ is clear,
% and on the other hand, to check topological nilpotence of $\gamma - 1$ 
% over a some nilpotent thickening of a ring $A$, it is enough to check
% over $A$ itself (since then the ``matrix coefficients'' of some power
% of $\gamma-1$ lie in a nilpotent ideal!). TG: where should we do this?
% Do we have a definition somewhere of what it means for a formal
% algebraic stack to be a formal neighbourhood in this way?}


\begin{lem}
  \label{lem: cartesian diagram for X basic} We have a 2-Cartesian
  diagram
  \[ \xymatrix{\cX_{K,d}\ar[r]\ar[d] &
      \cX_{\Kbasic,d[K:\Kbasic]}\ar[d]\\
      \cR_{K,d}^{\Gamma_{\disc}}\ar[r]&\cR_{\Kbasic,d[K:\Kbasic]}^{\Gamma_{\disc}}} \]where
  the horizontal arrows are the natural maps \emph{(}forgetting the
  $\A_K$-module structure\emph{)}, and the vertical arrows are the
  monomorphisms given by restricting the $\Gamma$-action
  to~$\Gamma_{\disc}$. 
\end{lem}
\begin{proof}
Unwinding the definitions, we need to show that if~$A$ is an
$\cO/\varpi^a$-algebra for some~$a\ge 1$, and~$M$ is an \'etale
$\varphi$-module over $\A_{K,A}$ with a semi-linear action
of~$\Gamma_{\disc}$, then the action of~$\Gamma_{\disc}$ extends to a
continuous action of~$\Gamma$ if and only if the same is true of~$M$
regarded as a module $\A_{\Kbasic,A}$. Since~$\A_{K,A}$ is free of finite
rank over~$\A_{\Kbasic,A}$, this is clear (for example, because it
follows from Lemma~\ref{lem:lattice properties} that the set of
lattices in~$M$ regarded as an $\A_{K,A}$-module is cofinal in the set of
lattices in~$M$ regarded as an $\A_{\Kbasic,A}$-module). 
\end{proof}

As a consequence of Lemmas~\ref{lem: relative representability for R
  over Kbasic} and~\ref{lem: cartesian diagram for X basic}, we can
deduce some properties of~$\cX_d$ from the corresponding properties
of~$\cR_d$.

\begin{prop}
  \label{prop: properties of X pulled back from R}% \mar{TG: fill in.}
  \leavevmode
  \begin{enumerate}
  \item The morphism $\cX_{K,d}\to \cX_{\Kbasic,d[K:\Kbasic]}$ is
    representable by algebraic spaces, affine, and of finite presentation.
  \item The diagonal of $\cX_{K,d}$ is representable by algebraic
  spaces, affine, and of finite presentation.
  \end{enumerate}
\end{prop}
\begin{proof}To make the proof easier to read, we write~$\cR_K$
  for~$\cR_{K,d}$ and~$\cR_{\Kbasic}$ for~$\cR_{\Kbasic,d[K:\Kbasic]}$.
We factor the morphism
$\cR_{K}^{\Gamma_{\disc}}\to
  \cR_{\Kbasic}^{\Gamma_{\disc}}$
  as
  % \mar{ME: Should probably try to improve the typesetting of this
  %         displayed expression TG: hopefully OK now, commenting out.}
  \begin{multline*}
\cR_{K}^{\Gamma_{\disc}} :=
\cR_{K}\underset{\Delta,\cR_{K}\times\cR_{K},\Gamma_\gamma}{\times}\cR_{K}
\to
\cR_{K}\underset{\Delta,\cR_{\Kbasic}\times\cR_{\Kbasic},\Gamma_\gamma}{\times}\cR_{K}
\\
\to
\cR_{\Kbasic}\underset{\Delta,\cR_{\Kbasic}\times\cR_{\Kbasic},\Gamma_\gamma}{\times}\cR_{\Kbasic}
=:
  \cR_{\Kbasic}^{\Gamma_{\disc}}.
\end{multline*}
	By Lemma~\ref{lem: relative representability for R over
    Kbasic}, the morphism
$\cR_{K}\to \cR_{\Kbasic}$
  is representable by algebraic spaces, affine, and of finite presentation.
Thus so is the second arrow in the preceding displayed expression.
The first arrow is a base-change of a product of two copies
of the diagonal 
$$\cR_{K} \to \cR_{K} \times_{\cR_{\Kbasic}} 
\cR_{K}.$$
As in the proof of Lemma~\ref{lem: abstract pullback diagonals lemma},
we deduce from 
Lemma~\ref{lem: relative representability for R over Kbasic}
that this diagonal 
  is representable by algebraic spaces, affine, and of finite presentation.
  Thus so is the first arrow in the expression,
  and hence so is the morphism
$\cR_{K}^{\Gamma_{\disc}}\to
  \cR_{\Kbasic}^{\Gamma_{\disc}}$.

%It follows that
%the morphism $\cR_{K}^{\Gamma_{\disc}}\to
%  \cR_{\Kbasic,d[K:\Kbasic]}^{\Gamma_{\disc}}$
%  is representable by
%algebraic spaces and of finite presentation.
%\mar{TG: need to give an argument to deduce this! ME: Done}
It follows from Lemma~\ref{lem:
  cartesian diagram for X basic} that
 $\cX_{K,d}\to \cX_{\Kbasic,d[K:\Kbasic]}$ is then
 also representable by algebraic spaces, affine,
 and of finite presentation.
The claimed properties of the diagonal of~$\cX_{K,d}$ follow from
the corresponding properties of the diagonal
of~$\cR^{\Gamma_{\disc}}_{K,d}$ proved in Lemma~\ref{lem: good properties R Gamma disc},
together with the fact that $\cX_{K,d} \to \cR_{K,d}^{\Gamma_{\disc}}$
is a monomorphism.
%Lemma~\ref{lem: abstract pullback diagonals lemma}.
%\mar{TG: or some argument of this kind, once abstract results exist.}
\end{proof}
%We will show that in fact $\cX_d$ is also an Ind-algebraic stack,
%although it is not a closed substack of $\cR_d^{\Gamma_{\disc}}$.
% In order to prove it, we employ %do this, we have to use
% the moduli stacks of Wach modules,
% which we discuss in Section~\ref{subsec: Wach modules}.
% %(In fact, we will
% %show in Corollary~\ref{cor:Xd is formal algebraic}
% %that~$\cX_d$ is even a formal algebraic stack, but this requires
% %additional work.) % Before doing this, we first recall the relationship
% % between~$(\varphi,\Gamma)$-modules and Galois representations.
% As a necessary preliminary,
% %we show that $\cX_d$ is limit preserving, 
% %using the criterion of Lemma~\ref{lem:testing continuity on M mod T}~(3). 
% we prove the following lemma.

%\section{Interlude: continuous actions}
%\mar{ME: These are some notes that probably won't be in the final
%	version.  They are my thoughts related to continuity of
%	the $\Gamma$-action.}
%
%\begin{lemma}
%Let $\Gamma$ be a topological group, and $M$ a linearly topologized
%abelian group.  Then an action $\Gamma \times M \to M$
%is jointly continuous if and only if the following 
%conditions hold: (i) each $\gamma$ in $\Gamma$ induces
%a continuous operator on $M$; (ii) for each $m\in M$,
%the corresponding orbit map $o_m: \Gamma \to M$ is continuous
%at the identity $1 \in \Gamma$;
%(iii) the action map $\Gamma \times M \to \Gamma$ 
%is continuous at the point $(1,0)$ in the product.
%\end{lemma}
%\begin{proof}
%	The only if direction is clear.  The if direction is easy enough
%	to establish too; see e.g.\ Lemma~3.1.1 of my locally analytic
%	\mar{ME: dodgily formatted reference here, and below.}
%	vectors paper.
%\end{proof}
%
%Note that if the elements of $\Gamma$ (or of any open neighbourhood
%of $1$ in $\Gamma$) form an equicontinuous set of
%operators from $M$ to itself, then certainly the condition 
%of continuity at $(1,0)$ holds.  
%Conversely, if we are given a continuous action $\Gamma\times M \to
%M$, it is easy to see that any compact subset of $\Gamma$
%acts equicontinuously. (See e.g.\ Lemma~3.1.4 of Locally Analytic Vectors.)
%In particular, if $\Gamma$ is compact, then $\Gamma \times M \to M$
%is continous if and only if it is separately continuous,
%and if $\Gamma$ forms an equicontinuous set of operators on $M$.
%
%Suppose further that $\Gamma$
%is compact, and that $M$ is a barelled locally convex topological
%vector space over some local field $K$.
%If $\Gamma\times M \to M$ is separately continuous,
%then each orbit of $\Gamma$ is compact, and so bounded in $M$.
%Thus $\Gamma$ is a point-wise bounded, or equivalently weakly bounded,
%set of operators.  
%Since $M$ is barrelled, the collection
%of operators $\Gamma$ is then furthermore uniformly bounded,
%%and thus equicontinuous.   Hence in this situation separate continuity of 
%the action is equivalent to joint continuity. 
%
%If we continue to suppose that $M$ is a locally convex vector space
%(over some local field $K$),
%then separate continuity can also be formulated as the requirement
%that $\Gamma$ acts by continuous linear transformations,
%and that the induced morphism $\Gamma \to \CL_s(M,M)$ is continuous.
%(Here $\CL_s(M,M)$ denotes the space of continuous linear transformations
%from $M$ to itself, and `$s$' denotes 
%the weak topology, i.e.\ the topology of pointwise convergence.
%Note that evaluation at the various elements $m \in M$
%induces an embedding $\CL(M,M) \to \prod_{m \in M} M,$
%and the weak topology is precisely the topology induced by the product
%topology on the target of this embedding.  Thus our claim
%follows from that the universal property of the product topology.)
%Another way to phrase the argument of the preceding paragraph
%is that for a barrelled space, the equicontinuous sets of operators
%are precisely the bounded subsets of $\CL_s(M,M)$; so, if $\Gamma$ is 
%compact, so that its image in $\CL_s(M,M)$ is compact and thus
%bounded,
%then it acts equicontinuously.
%
%Continue to suppose that $\Gamma$ is compact,
%that $M$ is a locally convex vector space over some
%local field $K$ which is bornological.
%Then any bounded subset of 
%$\CL_b(M,M)$ is equicontinuous
%(Here `$b$' denotes the strong topology,
%i.e.\ the topology of convergence on bounded subsets.)
%Thus if $\Gamma$ acts by continuous linear transformations,
%and the induced map
%$\Gamma \to \CL_b(M,M)$ is continuous, then
%(since $\CL_b(M,M) \to \CL_s(M,M)$ is continuous) we 
%see that $\Gamma$ acts separately continuously, as well as equicontinuously,
%and thus acts jointly continuously.  
%I'm not sure when to expect that the map
%$\Gamma \to \CL_b(M,M)$ should be continuous, but I guess in
% some situations it necessarily is so.
%
%The reason for this discussion is that I think similar things should hold
%for operators on Tate modules.  And the continuity condition
%invoked (without proof!)
%in the proof of Lemma~\ref{lem:testing continuity on M mod T}
%is (I think) essentially the claim that an action $\Gamma \times M \to M$
%is jointly continuous if and only if the induced map
%$\Gamma \to \CL_b(M,M)$ is continuous.
%I want to sort this out, because I think it will help clarify the
%continuity condition, and having a slightly more general theory
%of continuity will (I expect)
%give a useful flexibility for the arguments I have 
%in mind about adding a summand to make projective objects become free.
%

%\section{$\cX_d$ is limit preserving}
%\section{$T$-quasi-linear endomorphisms}
%\mar{ME: This material might ultimately be merged elsewhere, but it seemed
%	easiest to begin developing it separately. ME: Note, e.g.,
%that Lemma~\ref{lem:continuity criterion} makes reference to
%continuity of $\Gamma$-actions, and we also use the concept of lattices.
%ME added: Right now this material is just used to show that $\cX_d$
%satisfies~[1], so maybe just should be placed there. Actually,
%just before \S~\ref{subsec: X is Ind algebraic}, could have a subsection
%entitled ``$\cX_d$ is limit preserving'', which contains this
%material, and then culminates in the proof that $\cX_d$ is limit preserving.}

%We will show that $\cX_d$ is limit preserving by
%using the criterion of Lemma~\ref{lem:testing continuity on M mod T}~(3). 
From now on we will typically drop~$K$ from the notation, simply writing
$\cX_d$, $\cR_d$ and so on. We conclude this section by showing that $\cX_d$ is limit preserving,
using the material
on lattices and continuity developed in Appendix~\ref{app: Tate modules and continuity}.
\begin{lem}%\mar{TG: hopefully correct; may not belong here}
  \label{lem: gamma on T}We have ~$\gamma(T)-T\in
  p\A_{K,A}+T^2\A^+_{K,A}$. If~$K$ is basic then $\gamma(T)-T\in
  (p,T)T\A^+_{K,A}$.
\end{lem}
%\mar{ME: I'm not sure we want coefficients~$A$ here, since e.g.\
%we want to use that $a^{p^s} = 1$ implies $a = 1$.  And the $\Z_p$
%case will imply the coefficient case anyway. TG: we apply the lemma
%with coefficients, so I think we should have them in the statement,
%but then begin the proof by reducing to the case without them?  ME: OK, done}
\begin{proof} Given the definitions of $\A_{K,A}$ and $\A^+_{K,A}$
in terms of $\A_K$ and $\A^+_K$, it suffices to prove the lemma
for these latter rings (i.e.\ in the case when $A =  \Z_p$).
We then reduce modulo~$p$, and
  write~$\gamma(T)=\sum_{i=n}^{\infty}a_iT^i$, with
  ~$a_i\in k_\infty$ and~$a_{n}\ne 0$. Since the action of~$\gamma$ is continuous for
  the ~$T$-adic topology, we see in particular that for~$m>0$
  sufficiently large, we have $\gamma(T)^m=\gamma(T^m)\in T
  k_\infty[[T]]$. Since~$\gamma(T)^m=a_{n}^mT^{nm}+\dots$, we see that~$n>0$. % it is easy to see that we have~$a_0=0$ 
  % .\mar{TG: is this OK? It's
  %   certainly going to be true that~$a_0=0$ in our case, but I think
  %   we have to be slightly careful because $\gamma:T\mapsto T+1$
  %   satisfies $\gamma^{p^s}\to 1$.} Since~$\gamma$ acts invertibly, we
  % have % and  ~$a_1\in k_\infty^\times$.
Since~$\gamma^{p^s}\to 1$ as
  $s\to \infty$, we see that~$n=1$ and that~$a_1^{p^s}=1$ for all sufficiently
  large~$s$, which implies that~$a_1=1$.

  If~$K$ is basic then %  as explained in Section~\ref{subsec: phi gamma
    % over Zp},
  we have chosen~$T$ so that $\gamma(T)\in T\A_{K,A}^+$, so
  that $\gamma(T)-T\in (p,T)T\A^+_{K,A}$ by the above.
\end{proof}
If~$K$ is basic, it follows from Lemma~\ref{lem: gamma on T}
that~\eqref{eqn: condition on gamma semi-linear} holds, so that by
Lemma~\ref{lem:gamma minus 1 T quasi linear} we can use the material
on $T$-quasi-linear morphisms developed in Appendix~\ref{app: Tate
  modules and continuity} to study the action of~$\gamma$.

\begin{lem}
  \label{lem:X is limit preserving}$\cX_d$ is limit preserving.
\end{lem}
\begin{proof}
By Proposition~\ref{prop: properties of X pulled back from R} it suffices to prove
this in the case that~$K$ is basic (since a morphism which is 
of finite presentation is in particular limit preserving). Since~$\cX_d\hookrightarrow\cR_d^{\Gamma_\disc}$ is fully faithful, and~$\cR_d^{\Gamma_\disc}$ is limit
  preserving by Corollary~\ref{cor: basic properties of C and R p adic stacks}, it suffices to prove that~$\cX_d$ is limit preserving
  on objects. Since~$\cR_d^{\Gamma_\disc}$ is limit preserving on
  objects, we are reduced to showing that if~$T=\varprojlim_iT_i$ is a
  limit of affine schemes, and  $T_{i_0}\to\cR_d^{\Gamma_\disc}$ is a morphism with
  the property that the composite \[T\to T_{i_0}\to\cR_d^{\Gamma_\disc}\]
    factors through~$\cX_d$, then for some~$i\ge i_0$, the composite \[T_i\to T_{i_0}\to\cR_d^{\Gamma_\disc}\]
    factors through~$\cX_d$.
    % By Corollary~\ref{cor: can check continuity modulo p} we may check
    % continuity after reduction modulo $p$,  %\mar{ME: Have to prove this reduction mod $p$ claim.},
    % and so t
    This follows from the equivalence of conditions~(1)
    and~(7) 
    of Lemma~\ref{lem:testing continuity on M mod T}, together with Lemma
    ~\ref{lem:quasi-linear limit preserving}.
%
%    To this end, write~$T_i=\Spec A_i$, $T_j=\Spec A_j$, $T=\Spec T$,
%    and let~$M$ be the \'etale $(\varphi,\Gamma_{\disc})$-module
%    with~$A_i$-coefficients corresponding to the morphism
%    $T_i\to\cR_d^{\Gamma_{\disc}}$. Let~$\gM$ be a lattice in~$M$, so that by
%    Lemma~\ref{lem:testing continuity on M mod T}, there is an
%    integer~$s\ge 0$ such that $(\gamma^{p^s}-1)(\gM_A)\subset T\gM_A$,
%    where~$\gM_A:=\gM\otimes_{\gS_{A_i}}\gS_A$. \mar{TG: again I'm not
%    sure what the best way is to finish from here, but the idea is
%    simple: we are asking that some map to $M/\gM$ vanishes, and we
%    know it after going up to~$A$, so it will be true after going to
%    some~$A_j$, because ``the Laurent tails are finite'', or perhaps
%    more precisely because of some finite presentation statement for
%    $\Hom$s of \'etale $\varphi$-modules? As always I'm a bit confused
%  about $\Gamma$-semi-linearity.}
\end{proof}

 
\section{Weak Wach modules}\label{subsec: Wach modules}
In this section we introduce the notion of a weak Wach module of
height at most~$h$ and level at most~$s$. These will play a purely technical
auxiliary role for us, and will be used only in order to show
that~$\cX_d$ is an Ind-algebraic stack.  %  i.e.\ those for which~$\Gamma$ acts on~$\gM$ itself, such
  % that furthermore  the induced action of $\Gamma$ on $\gM/T\gM$ is
  % trivial.

% \mar{ME: I'm not sure about terminology, but I will use Kisin module
% 	for objects with just a $\varphi$-structure (even though we're in the
% 	cyclotomic context), weak Wach module for things that have a
% 	$\Gamma$-action, and maybe (?) Wach module for things for which the
% 	$\Gamma$-action is trivial mod $T$. (Actually, I just looked
% at Berger--Breuil for a definition of Wach module, and now I'm not
% quite sure what do.  The main thing is that for a Wach module,
% the $\Gamma$-action should be finite even in the $p$-adic limit
% (I guess reflecting the abelian descent data on potentially crystalline
% Dieudonn\'e module), not just continuous.  So probably a more 
% elaborate discussion of the notion of Wach module is needed.) TG: I
% decided to change ``Kisin module'' to match the hopefully more
% standard terminology that I introduced in the previous
% section. Commenting out.}

%   \begin{defn}% \mar{TG: check that this definition coincides with
%    % Drinfeld's definition of a lattice (to the extent that makes sense
%    % - his may not require $\gS_A$-stability?)}
%    \label{defn: lattice}If~$M$ is a finite generated $\OEA$-module, then a
%    \emph{lattice} in~$M$ is a finitely generated $\gS_A$-submodule
%    $\gM\subseteq M$  whose $\OEA$-span
%    is~$M$.% \mar{TG: didn't think much about this definition! Probably
%      % torsion free is automatic, in fact\dots}
%  \end{defn}
%    \begin{rem}
%    \label{rem: lattice and Drinfeld}By~\cite[Thm.\
%    5.1.14]{EGstacktheoreticimages} (a theorem of Drinfeld), we may
%    think of a finitely generated~$\OEA$-module as a Tate $A$-module $M$
%    together with a topologically nilpotent endomorphism~$T$. In this
%    optic a lattice in our sense is a lattice in the Tate module~$M$
%    which is preserved by~$T$ (by definition, a lattice~$L\subset M$
%    is an open submodule with the property that for every open
%    submodule~$U\subset L$, the $A$-module $L/U$ is finitely
%    generated).\mar{TG: maybe the converse is true, that every
%      $T$-stable lattice is a lattice? I think we just need to deduce
%      that it's finitely generated over~$\gS_A$, which might follow
%      from the SGA results cited in \cite[Prop.\
%      5.1.8]{EGstacktheoreticimages}, but I'm not going to worry about
%    this right now.}
%  \end{rem}
%
%
%\begin{lem}
%	\label{lem:testing continuity on M mod T}
%        Suppose that $A$ is an $\cO/\varpi^a$-algebra for
%        some~$a\ge 1$. Let~$M$ be a projective \'etale $\varphi$-module with
%        $A$-coefficients, equipped with an action of~$\Gamma_\disc$. Then
%        the following are equivalent:
%        \begin{enumerate}
%        \item The action of~$\Gamma_{\disc}$ extends to a continuous
%          action of~$\Gamma$.
%        \item For any lattice~$\gM\subset M$, there exists~$s\ge 0$
%          such that $(\gamma^{p^s}-1)(\gM)\subset T\gM$.
%        \item For some lattice~$\gM\subset M$ and some~$s\ge 0$, we have
%          $(\gamma^{p^s}-1)(\gM)\subset T\gM$.
%        \end{enumerate}
%  \end{lem}
%  \begin{proof} Since $\Gamma\cong\Zp$, and the canonical topology
%        on~$M$ is induced by the $(p,T)$-adic topology on the
%        any lattice~$\gM$, we see that the action of~$\Gamma$ is
%        continuous if and only if for some (equivalently, for every)
%        lattice~$\gM\subset M$, and every~$n\ge 0$, there exists~$s(n)$
%        such that $(\gamma^{p^s}-1)(\gM)\subset (p,T)^n\gM$ for all
%        $s\ge s(n)$. If~$n\ge a$ then $(p,T)^n\subset (T)$, so we see
%        that $(1)\implies (2)\implies (3)$.
%
%        Suppose now that~(3) holds. We have
%        $(\gamma^{p^s}-1)(\gM)\subset T\gM\subset
%        (p,T)\gM$. Writing
%        \[\gamma^{p^{s+n+1}} - 1 = ((\gamma^{p^{s+n}}-1)+1)^p-1,\] it
%        follows from an easy induction on~$n$ that for each~$n\ge 0$,
%        we have
%        \[(\gamma^{p^{s+n}}-1)(\gM)\subset (p,T)^n\gM,\] as
%        required.
%  \end{proof}
We suppose throughout this
section that~$K$ is basic.
\begin{remark}
	\label{rem:testing continuity on M mod u}
        By Lemma~\ref{lem:testing continuity on M mod T}, if $A$ is an $\cO/\varpi^a$-algebra for some~$a\ge 1$, and
        $\gM$ is a rank~$d$ projective $\varphi$-module of
        $T$-height $\leq h$ over $A$, such that~$\gM[1/T]$ is equipped
        with a semi-linear action of $\Gamma_{\disc}$, then this
        action extends to a 
        continuous action of~$\Gamma$ if and only if for some~$s\ge 0$ we have $(\gamma^{p^s}-1)(\gM)\subseteq
        T\gM$.%  Indeed, $\Gamma\cong\Zp$, and the canonical topology
        % on~$\gM[1/T]$ is induced by the $(p,T)$-adic topology on the
        % open subspace~$\gM$.
        % is the product of a
                                %  finite group\mar{TG: no 
%   finite group now?}
% and a pro-$p$-group,
% a semi-linear $\Gamma$-action on  $\gM[1/T]$ is 
% continuous if and only if the induced $\Gamma$-action on $\gM/T\gM$ 
% is continuous when the latter is endowed with its $p$-adic
% topology. Indeed, both conditions are easily seen to be equivalent to
% asking that if~$m$ is sufficiently large, then~$\gamma^{p^m}$ acts
% trivially on~$\gM/(p,T)\gM$.
% \mar{TG: I think this is right, but
%   leaving this comment just in case of idiocy.}
% \mar{TG: think this through again and check I understand;
  % maybe add a remark about the proof.}
\end{remark}

\begin{df}
	\label{def:weak Wach modules}Suppose that~$K$ is basic.
	A {\em rank $d$ projective weak Wach module of $T$-height
    $\leq h$} \index{weak Wach module}
	with $A$-coefficients is a rank~$d$ projective $\varphi$-module $\gM$ 
	over $\A^+_{K,A}$, % \mar{TG: have to slightly clarify somewhere
          % what this means given that $\A^+_{K,A}$ is not
          % $\varphi$-stable; the definition should just be that after
          % inverting~$F$ it's an \'etale $\varphi$-module}
        which is of $T$-height $\leq h$, %  when viewed as a
        % $\varphi$-module over~$\A^+_{\Qp,A}$,\mar{TG: need to be
        %   careful here as there is now more than one $T$, maybe
        %   something like $T_{\Qp}$-height or $\pi$-height?}
        such that~$\gM[1/T]$ is equipped
	with a semi-linear action of $\Gamma$ which is furthermore
        continuous 
	% \mar{ME: Should recall that $\gM$ does have a canonical topology;
	% 	probaby a non-existent Drinfeld citation will be involved.}
	when $\gM[1/T]$ is endowed with its canonical topology  (see
        Remark~\ref{rem:topologies on M}).

	If~$s\ge 0$, then we say that~$\gM$ has
        \emph{level $\le s$} \index{level (of weak Wach module)} if $(\gamma^{p^s}-1)(\gM)~\subseteq~T\gM$.
	Remark~\ref{rem:testing continuity on M mod u} shows
	that any projective weak Wach module is of level $\leq s$
	for some $s \geq 0$.
\end{df}
A special role in the classical theory of $(\varphi,\Gamma)$-modules is played by the weak Wach modules of
  level~$0$.  However, this will not be important for us.


% \begin{remark}
% 	\label{rem:testing continuity on M mod u}
% Note that since $\Gamma\cong\Zp$, %  is the product of a finite group\mar{TG: no
% %   finite group now?}
% % and a pro-$p$-group,
% a semi-linear $\Gamma$-action on a $\varphi$-module $\gM$ is 
% continuous if and only if the induced $\Gamma$-action on $\gM/T\gM$ 
% is continuous when the latter is endowed with its $p$-adic
% topology. Indeed, both conditions are easily seen to be equivalent to
% asking that if~$m$ is sufficiently large, then~$\gamma^{p^m}$ acts
% trivially on~$\gM/(p,T)\gM$.\mar{TG: I think this is right, but
%   leaving this comment just in case of idiocy.}%  \mar{TG: think this through again and check I understand;
%   % maybe add a remark about the proof.}
% \end{remark}

% \mar{ME: A LaTeX remark; I am using {\tt $\backslash$cWW} rather than
% 	simply {\tt $\backslash$cW}, just
% 	in case we need to change the notation.}

\begin{df}
	We let $\cWW_{d,h}$ denote the moduli stack of rank $d$ \index{$\cWW_{d,h}$}
       	projective weak Wach modules of $T$-height $\leq h$. %  \mar{TG: at
          % this point, should this be $T$-height at most~$h$?}
        (That this
        is an \emph{fpqc} stack follows as in Definition~\ref{defn:
          Xd}.) We let~$\cWW_{d,h,s}$ denote the substack of rank $d$ \index{$\cWW_{d,h,s}$}
       	projective weak Wach modules of $T$-height $\leq h$ and
        level~$\le s$.  % \mar{TG: a
          % stack, should we say more about it?}
\end{df}

% \begin{rem}
%   % As we will see in Section~\ref{subsec: crystalline moduli unramified base}, a
% In the case that~$K$ is unramified over~$\Qp$, a special role in the theory of $(\varphi,\Gamma)$-modules (in the case when $K$ is
%   unramified over $\Q_p$) is played by the weak Wach modules of
%   level~$0$. ,%  i.e.\ those for which~$\Gamma$ acts on~$\gM$ itself, such
%   % that furthermore  the induced action of $\Gamma$ on $\gM/T\gM$ is
%   % trivial.
%   However, this will not be important for us.% \mar{TG: could
%     % just remove this}
%   % However, for the moment we will not need to
%   % consider this condition.
%   % \mar{ME: This paragraph seems like a bit of a
%     % non-sequitur at the moment; ultimately it probably should stay,
%     % but be amplified, with a forward reference to the construction of
%     % crystalline moduli stacks in the unramified case.  And perhaps
%     % should be set off as a remark?}
% \end{rem}
We recall from Section~\ref{subsec: moduli stacks of phi modules}
that there is a $p$-adic formal
algebraic stack
$\cC_{d,h}$ classifying rank $d$ projective $\varphi$-modules of
$T$-height at most~$h$.
We consider the fibre product
$\cR_d^{\Gamma_{\disc}}\times_{\cR_d}\cC_{d,h}$, where the
map~$\cR_d^{\Gamma_{\disc}}\to\cR_d$ is the canonical morphism given by
forgetting
the~$\Gamma_{\disc}$ action. By Proposition~\ref{prop:Gamma-disc
  modules as fixed points; etale case}, this is the moduli stack of rank $d$ projective $\varphi$-modules~$\gM$ of
$T$-height at most~$h$, equipped with a semi-linear action
of~$\Gamma_{\disc}$ on~$\gM[1/T]$. % \mar{TG: instead we now need to
  % consider the pullback of $\cC_{d[K:\Qp],h}\to\cR_{d[K:\Qp]}$ in the $\Qp$-case via
  % the morphism $\cR_d\to\cR_{d[K:\Qp]}^{\Qp}$ where I'm making up
  % notation on the right hand side. Not sure what properties of this
  % morphism we will need; somewhere in the arguments I expect to use
  % something like~\cite[Lem.\ 5.4.10]{EGstacktheoreticimages}.}
\begin{lem}
  \label{lem: R Gamma disc times C is p adic formal
    algebraic}$\cR_d^{\Gamma_{\disc}}\times_{\cR_d}\cC_{d,h}$ is a
  $p$-adic formal algebraic stack of finite presentation over~$\Spf\cO$.
\end{lem}
\begin{proof}The map $\cR_d^{\Gamma_{\disc}}\to\cR_d$ is  representable by algebraic spaces
	and
	of finite presentation, being a base-change of the diagonal
        of~$\cR_d$, which is representable by algebraic spaces
	and
	of finite presentation by Corollary~\ref{cor: basic properties of C and R p adic stacks}. % by Theorem~% \ref{thm: generalisation of
        %   % main results of proper.tex}
        % ),\mar{TG: add ref}
        % \mar{ME: Add reference for this; the point is that it is a base-change of the diagonal of 
		% $\cR_d$, which we know is representable by
		% algebraic spaces and of finite presentation.}
  Again by  Corollary~\ref{cor: basic properties of C and R p adic stacks}, $\cC_{d,h}$ is a $p$-adic formal
  algebraic stack of finite presentation, and thus
  $\cR_d^{\Gamma_{\disc}}\times_{\cR_d} \cC_{d,h}$ is a $p$-adic
  formal algebraic stack of finite presentation, as claimed.
  \end{proof}
% The functor which forgets the action of~$\varphi$ and sends $\gM$ to 
% $\gM/T\gM$ gives a morphism of $p$-adic
% formal algebraic stacks
% $\cC_{d,h} \to \Bun_d$, where $\Bun_d$ is the $p$-adic formal
% algebraic stack whose $\Spf A$-points are the
% projective $(W(k)\otimes_{\Zp}A)$-modules of rank~$d$.%  \mar{ME: Here $\Bun_d$ is over $\Spf \Z_p$, so
	% it is a $p$-adic formal algebraic stack. TG: should define this.}
% Just as in the case of~$\cR_d$ explained above, we have a morphism
% $\gamma:\cC_{d,h}\to\cC_{d,h}$, and we may consider the
% stack~$\cC_{d,h}^{\Gamma_\disc}$ defined in the same way
% as~$\cR_d^{\Gamma_\disc}$; that is, \[\cC_d^{\Gamma_{\disc}}:=
% \cC_d\underset{\Delta,\cC_d\times\cC_d,\Gamma_\gamma}{\times}\cC_d, \]where
%  $\Delta$ is the diagonal of~$\cC_d$ and $\Gamma_\gamma$ is the
% graph of~$\gamma$.\mar{TG is here.}

% We have
% the following analogue of
% Proposition~\ref{prop:Gamma-disc modules as fixed points; etale
%   case}. 

% \begin{prop}
% 	\label{prop:Gamma-disc modules as fixed points; Kisin case}\leavevmode
% 	% There is an action of $\Gamma$ on the $p$-adic formal
% 	% algebraic stack $\cC_{d,h}$,
% 	% with the following properties:
% 	\begin{enumerate}
% 		\item
% 	The moduli stack of rank $d$ locally free 
% 	$\varphi$-modules of $F$-height at most~$h$ equipped with a semi-linear action
% 	of $\Gamma_{\disc}$ is isomorphic to the fixed point stack
% 	$\cC_{d,h}^{\Gamma_{\disc}}$.
%       \item $\cC_{d,h}^{\Gamma_\disc}$ is a $p$-adic formal algebraic
%         stack over~$\Spf\Zp$.
% % \item The morphism $\cC_{d,h} \to \Bun_d$ is $\Gamma_{\disc}$-equivariant
% % 	if we equip $\Bun_d$ with the trivial
% %         $\Gamma_{\disc}$-action.
%         % \mar{TG: no, $W(k)$-action, and needs
%           % to be $\Bun_{d[K:\Qp]}$.}
%         % \mar{TG: $\Gamma$ versus $\Gamma_{\disc}$}
% \item If $h' \geq h,$ then the closed immersion $\cC_{d,h} \hookrightarrow
% 	\cC_{d,h'}$ is $\Gamma_{\disc}$-equivariant.
% \end{enumerate}
% \end{prop}
% \begin{proof}
% 	The construction of the action,
% 	and the proof of~(1), follows identical lines to the proof of
% Proposition~\ref{prop:Gamma-disc modules as fixed points; etale
%   case}. Since~$\cC_{d,h}$ is a $p$-adic formal
% algebraic stack with affine diagonal by Theorem~\ref{thm: C is a $p$-adic formal algebraic stack and related
%     properties}, $\cC_{d,h}^{\Gamma_\disc}$ is also a $p$-adic formal algebraic
%         stack by~\cite[Lem.\ 4.8]{Emertonformalstacks}.\mar{TG: at the
%         moment that is just a reference showing that this is an
%         Ind-algebraic stack. I think it should be easy to show that a fibre
%         product of $p$-adic formal algebraic stacks is $p$-adic formal
%       algebraic, but presumably \cite{Emertonformalstacks} would be a
%       better place to do that?} The final point is immediate from the definitions. % Since~$\Gamma_{\disc}$ acts
%         % trivially on~$\gS_A/T$, (3) is immediate from the definitions,
%         % as is~(4).
% % \mar{ME: Finish proof by adding remark about~(3), which is obvious
% % 	once we observe that $\Gamma$ acts trivially on $\gS/X\gS$,
% % and about~(4), which is just obvious.}
% \end{proof}

% Corollary~\ref{cor:fixed points preserve algebraicity}
% shows that $\cC_{d,h}^{\Gamma_{\disc}}$ is a $p$-adic formal 
% algebraic stack over
% $\Spf \Z_p$.
Restricting the $\Gamma$-action on a weak Wach module
to $\Gamma_{\disc}$, %(and taking into account the preceding proposition),
we obtain a morphism $\cWW_{d,h} \to \cR_d^{\Gamma_{\disc}}\times_{\cR_d}\cC_{d,h}$,
which, by the evident analogue for weak Wach modules
of Remark~\ref{rem:discrete}, is fully faithful. 
Thus $\cWW_{d,h}$ may be regarded as a substack of
$\cR_d^{\Gamma_{\disc}}\times_{\cR_d}\cC_{d,h}$. The following
proposition records the basic properties of the stacks~$\cWW_{d,h}$
and~$\cWW_{d,h,s}$.% ; the additional results in the case~$s=0$ will be
% used in Section~\ref{sec:
%   crystalline moduli}.

% \begin{remark}
% 	\label{rem:bundles}
% % Example~% \ref{ex:fixed points of trivial action}
% % shows that\mar{TG: missing reference to appendix that I commented out}
% As in the proof of Proposition~\ref{prop:Gamma-disc modules as fixed
%   points; etale case}, the fixed point stack \[\Bun_d^{\Gamma_{\disc}}:=\Bun_d\underset{\Delta,\Bun_d\times\Bun_d,\Gamma_\gamma}{\times}\Bun_d\] is the $p$-adic formal algebraic
% stack classifying projective $W(k)\otimes_{\Zp}A$-modules of rank~$d$
% endowed with a $W(k)\otimes_{\Zp}A$-linear action of $\Gamma_{\disc}$.
% \end{remark}

% \begin{df}\leavevmode
%   \begin{enumerate}
%   \item  For $s \geq 1$, we let $\cB_{d,s}$ denote the $p$-adic
%     formal algebraic stack, of finite presentation over $\Spf \Z_p$,
%     classifying projective $W(k)\otimes_{\Zp}A$-modules of rank~$d$
% endowed with a $W(k)\otimes_{\Zp}A$-linear action of the finite
%     group $\Gamma/\langle \gamma^{p^s}\rangle.$

    
%   \item We let $\cB_d$ denote the stack (over $\Spf \Z_p$) classifying
%     projective $W(k)\otimes_{\Zp}A$-modules of rank~$d$
% endowed with a $W(k)\otimes_{\Zp}A$-linear action of $\Gamma$.
%   \end{enumerate}

% \end{df}

% If $s' \geq s$ then there is an evident closed immersion
% \numequation
% \label{eqn:closed immersions for s}
% \cB_{d,s} \hookrightarrow \cB_{d,s'}. 
% \end{equation}
% There are also evident
% monomorphisms
% \numequation
% \label{eqn:monomorphisms for s}
% \cB_{d,s} \hookrightarrow \cB_d,
% \end{equation}
% which are compatible
% with the closed immersions~(\ref{eqn:closed immersions for s}).
% Taking into account Remark~\ref{rem:bundles}, together with 
% the density of $\Gamma_{\disc}$ in $\Gamma$,
% we furthermore see that there is a monomorphism
% \numequation
% \label{eqn:monomorphism for B}
% \cB_d \hookrightarrow
% \Bun_d^{\Gamma_{\disc}}.
% \end{equation}

% \begin{lemma}
% 	\label{lem:inductive description of B}\leavevmode
%         \begin{enumerate}
%         \item If $s$ is sufficiently large {\em (}depending on $d${\em
%             )} then the closed immersion~{\em (\ref{eqn:closed
%               immersions for s})} induces an isomorphism on underlying
%           reduced algebraic stacks.

          
%         \item  For $s \geq 1,$ the composite
% 	$$\xymatrix{\cB_{d,s}
%           \ar^-{\text{\em{(\ref{eqn:monomorphisms for s})}}}[r] &
%           \cB_d \ar^-{\text{\em{(\ref{eqn:monomorphism for B})}}}[r] &
%           \Bun_d^{\Gamma_{\disc}}\\ }$$ is a closed immersion of
%         finite presentation.


	
%       \item If we form the inductive limit $\varinjlim_s \cB_{d,s}$ with
%         respect to the closed immersions~{\em (\ref{eqn:closed
%             immersions for s})}, then the morphism
%         $\varinjlim_s \cB_{d,s} \to \cB_d$ induced by the morphisms~{\em
%           (\ref{eqn:monomorphisms for s})} is an isomorphism, and thus
%         realizes $\cB_d$ as a formal algebraic stack over $\Spf \Z_p$.
%       \end{enumerate}

% \end{lemma}

% The closed immersions~(\ref{eqn:closed immersions for s})
% induce closed immersions
% \numequation
% \label{eqn:Wach closed immersions for s}
% 	\cWW_{d,h,s} \hookrightarrow \cWW_{d,h,s'},
% \end{equation}
% and similarly
% the monomorphisms~(\ref{eqn:monomorphisms for s}) induce monomorphisms
% \numequation
% \label{eqn:Wach monomorphisms for s}
% \cWW_{d,h,s} \hookrightarrow \cWW_{d,h},
% \end{equation}
% which are compatible with the closed immersions~(\ref{eqn:Wach closed
% 	immersions for s}).

\begin{prop}
	\label{prop:inductive description of W}\leavevmode
        \begin{enumerate}% \mar{TG: I can't see why we assume $s\ge 1$
            % rather than $s\ge 0$, and I'm now referencing this in the
            % other paper to do FL theory where I'm assuming that $s=0$
            % is OK. TG: done, commenting out.}
        \item For $s \geq 0,$ the morphism
	$$\xymatrix{\cWW_{d,h,s}
          \ar% ^-{\text{\em{(\ref{eqn:Wach monomorphisms for s})}}}
          [r]%  &
          % \cWW_{d,h} \ar[r]
          & \cR_d^{\Gamma_{\disc}}\times_{\cR_d}\cC_{d,h}\\ }$$ is a
        closed immersion of finite presentation.
        In particular, each of the stacks $\cWW_{d,h,s}$ is a $p$-adic
          formal algebraic stack of finite presentation over
          $\Spf \cO$.
        \item If $s'\ge s$, then the canonical monomorphism $
          \cWW_{d,h,s} \hookrightarrow \cWW_{d,h,s'}$ is a closed
          immersion of finite presentation.
        
      %   \item   If $s$ is sufficiently large {\em (}depending on $d${\em
      %       )}, and $s' \geq s$, then the closed immersion
      %     $\cWW_{d,h,s} \hookrightarrow \cWW_{d,h,s'}$ is a finite
      %     order thickening; in particular, it %  ~{\em
      %       % (\ref{eqn:Wach closed immersions for s})}
      %     induces an
      %     isomorphism on underlying reduced algebraic substacks.

         	
       \item %  If we form the inductive limit $\varinjlim_s \cWW_{d,h,s}$, % via the
         % % closed immersions~{\em (\ref{eqn:Wach closed immersions for
         % %     s})},
         % then t
         The canonical morphism $\varinjlim_s \cWW_{d,h,s} \to \cWW_{d,h}$
         % induced by the monomorphisms~{\em (\ref{eqn:Wach monomorphisms
         %     for s})}
         is an isomorphism. In particular, $\cWW_{d,h}$
         is an Ind-algebraic stack.  %formal algebraic stack.
      \end{enumerate}

\end{prop}
\begin{proof} 	Since  $\cR_d^{\Gamma_{\disc}}\times_{\cR_d}
	\cC_{d,h}$ classifies
	$\varphi$-modules $\gM$ that are
	projective of rank $d$ and
	of $T$-height $\leq h$, which are endowed with a $\Gamma_{\disc}$-action
	on the underlying \'etale $\varphi$-module, in order to prove~(1),
	we must show that
        the condition that~$(\gamma^{p^s}-1)(\gM)\subseteq T\gM$ is a
        closed condition, and is determined by finitely many
        equations. It suffices to check this after pulling back via an
        arbitrary morphism 
          $\Spec A\to \cR_d^{\Gamma_{\disc}}\times_{\cR_d}
	\cC_{d,h}$,
where~$A$ is an $\cO/\varpi^a$-algebra for some~$a\ge 1$. 
	% We then see that $\cC_{d,h}^{\Gamma_{\disc}}$ is the sub-stack
	% of this $2$-fibre product classifying objects for which
	% $\gM$ is actually stable under the $\Gamma_{\disc}$-action. To
        % see that this is a closed condition,

%We now argue as in the proof of~\cite[Lem.\
%5.4.10]{EGstacktheoreticimages}. Write~$M=\gM[1/T]$. We may think of
%the action of~$\gamma^{p^s}$ on~$\gM[1/T]$ as an element of
%$\Hom_{\A_{K,A},\varphi}((\gamma^{p^s})^*M,M)$.  By~\cite[Prop.\
%5.4.8]{EGstacktheoreticimages}, this is represented by a finitely
%presented scheme over~$\Spec A$, so we need only show that the
%condition that $(\gamma^{p^s}-1)(\gM)\subseteq T\gM$ is a finitely presented closed
%condition.

This morphism gives rise to a projective $\A^+_{K,A}$ $\varphi$-module $\gM$,
such that $M:=  \gM[1/T]$ is \'etale, and is furthermore endowed with
a semi-linear $\Gamma_{\disc}$-action.
As already stated above, we must now check that the condition
$(\gamma^{p^s}-1)(\gM) \subseteq T\gM$  is a finitely presented closed condition
(in the sense that it holds after replacing $\gM$  by  $\gM_B := \A^+_{K,B}
\otimes_{\A^+_{K,A}} \gM$ for an $A$-algebra $B$ if and only if
$\Spec B \to \Spec A$ factors through a certain finitely presented closed
subscheme of~$\Spec A$).

Note firstly that for~$n$ sufficiently large, we have
$(\gamma^{p^s}-1)(T^n\gM)\subseteq T\gM$. Indeed,
writing \[\gamma^{p^s}-1=((\gamma-1)+1)^{p^s}-1\]and expanding via the
binomial theorem, it suffices to show that for ~$n$ sufficiently large, we have
$(\gamma-1)^m(T^n\gM)\subseteq T\gM$ for all $0\le m\le p^s$; and this
is immediate from Lemma~\ref{lem: how f behaves on lattices}.

We next % note that by \cite[
% Lem.~5.2.14]{EGstacktheoreticimages}, we may find a projective \'etale
% $\varphi$-module $N$ such that $F:=M \oplus N$ is free.
choose a finitely generated projective $\A_{K,A}^+$-module ~$\gN$ such
that ~$\gF:=\gM\oplus\gN$ is free, and write~$N:=\gN[1/T]$,
$F:=\gF[1/T]$. Extend the morphism~$\gamma^{p^s}-1:M\to M$ to the
morphism~$f:F\to F$ given by
$(\gamma^{p^s}-1,0):M\oplus N\to M\oplus N$, so that the condition
that $(\gamma^{p^s}-1)(\gM)\subseteq T\gM$ is equivalent to asking that
$f(\gF)\subseteq T\gF$.
% \mar{ME: The argument in this paragraph seems slightly garbled to me; e.g.\ it's not explained why $T^n \gF$ is in the kernel for large $n$, and I don't really understand the role of the splitting.  But it is very similar to some arguments in App.~D, and hopefully can be rephrased to make sense (assuming I'm not 
% 	just confused.}
By our choice of~$n$ above,  the morphism $\gF\to F/T\gF$
induced by~$f$ factors through a morphism $\gF/T^n\gF\to F/T\gF$;
since~$\gF$ is finitely generated, after enlarging~$n$ if necessary,
it factors through
$\gF/T^n\gF\to T^{-n}\gF/T\gF$. %  Since
% the morphism $\gF\to\gF/T^n\gF$ splits (because $\gF/T^n\gF$ is
% projective), the condition that~$f(\gF)\subseteq T\gF$  is equivalent to
% the vanishing of the morphism $\gF/T^n\gF\to T^{-n}\gF/T\gF$ of finite
% free~$A$-modules. This
The vanishing of this morphism is obviously a closed condition; indeed
it is given by the vanishing of finitely many matrix entries, so is
closed and of finite presentation, as required.
% \mar{ME: I think we also want to note that it's given by the vanishing 
% 	of finitely many matrix entries, so that it's closed
% 	and of finite presentation, right? TG: yes, good point.}
        
% \mar{TG: still stuck on this, annoyingly! What we had here before just
%   seems to be bogus to me - it was an argument about matrix entries
%    that didn't seem to make sense. What I'd like to do is to
%   argue that we want to consider the condition that the map
%   $(\gamma^{p^s}-1):\gM\to M/T\gM$ vanishes, and then to argue as in
%   the proof of~\cite[Lem.\
%         5.4.10]{EGstacktheoreticimages}: namely this map factors
%         through a map of finite rank projective $A$-modules, and then
%         we're surely in good shape? Perhaps we can in fact just do
%         that? I keep being thrown off by ~$\gamma$ being only
%         $A$-linear and not $\A_{K,A}$-linear, but maybe this is a red herring.}
        %By~\cite[Lem.\
        % 5.2.14]{EGstacktheoreticimages} there is a finite rank projective
        % \'etale $\varphi$-module~$P$ such that~$F:=\gM[1/u]\oplus P$ is a
        % free \'etale $\varphi$-module of finite rank. We can think of
        % the action of~$\gamma\in\Gamma_{\disc}$ as an element
        % of~$\End_{\OEA}(F)$ by letting it act by~$0$ on~$P$. Choosing
        % a basis for~$F$, 
	% % If we work \'etale locally and so choose a basis for $\gM$,
	% % then
        % the condition that~$(\gamma^{p^s}-1)(\gM)\subset T\gM$ % stability of $\gM$ under the $\Gamma_{\disc}$-action
	% is equivalent to the vanishing of certain entries in the
        % corresponding 
	% matrix describing this action of~$(\gamma^{p^s}-1)$;
        % in
        % particular, % \mar{TG: need to rewrite this proof a bit as in
        %   % proper.tex, using that we can write projective as a summand
        %   % of free}
	% it is a closed condition of finite presentation, as required.
	% (Since the basis of $F$ consists of finitely many elements,
	% and since $\gamma^{p^s}-1$ takes each of these generators
	% to an an element of $F[1/u]$, there are indeed only finitely
	% many equations that have to be imposed in order to ensure
	% that $(\gamma^{p^s}-1)(\gM) \subset T\gM$.)
	%\mar{TG: this is hopefully not too
        %  far from a correct argument\dots ME: Is finite presentation clear?
%ME: OK, looks good, commenting out.}
      Since~$\cR_d^{\Gamma_{\disc}}\times_{\cR_d}\cC_{d,h}$ is a
      $p$-adic formal algebraic stack of finite presentation over
      $\Spf \cO$, it follows that so is each~$
      \cWW_{d,h,s}$. That~$\cWW_{d,h,s} \hookrightarrow \cWW_{d,h,s'}$
      is a closed immersion of finite presentation follows immediately
      from a consideration of the composite
      \[\cWW_{d,h,s} \hookrightarrow \cWW_{d,h,s'}\hookrightarrow
        \cR_d^{\Gamma_{\disc}}\times_{\cR_d}\cC_{d,h},\]as we have
      just shown that both the composite and the second morphism are
      closed immersions of finite presentation. Thus we have
      proved~(1) and~(2).

For~(3), we need to show that  every morphism $\Spec A\to\cWW_{d,h}$,
where~$A$ is an $\cO/\varpi^a$-algebra for some~$a\ge 1$, factors
through~$\cWW_{d,h,s}$  for~$s$ sufficiently large. As
we already noted in Definition~\ref{def:weak Wach modules}, this follows from
Remark~\ref{rem:testing continuity on M mod u}. Since
each~$\cWW_{d,h,s}$ is an Ind-algebraic stack, so is~$\cWW_{d,h}$. % The
% thickenings~$\cWW_{d,h,s} \hookrightarrow \cWW_{d,h,s'}$ are of finite order
% (since we are comparing the conditions that~$\gamma^{p^s}=1$ and that
% ~$\gamma^{p^{s'}}=1$), so~$\cB_d$ is a formal algebraic stack
% over~$\Spf\Zp$ by~\cite[Lem.\ 6.3]{Emertonformalstacks}.
\end{proof}
% For~(2), note that since the composite is a morphism between formal
% algebraic stacks of finite presentation over~$\Spf\Zp$, it is a
% morphism of finite presentation. To see that it is a closed immersion,
% note that it is given by imposing the equation~$\gamma^{p^s}=1$;
% locally this is given by equations in the corresponding matrix
% entries, as required.


% Since each of~$\cB_{d,s}$, $\Bun_d^{\Gamma_{\disc}}$
%   and~$\cC_{d,h}^{\Gamma_{\disc}}$ is a $p$-adic formal algebraic
%   stack of finite presentation over~$\Spf\Zp$, so is the fibre
%   product~$\cWW_{d,h,s}=\cB_{d,s}\times_{\Bun_d^{\Gamma_{\disc}}}\cC_{d,h}^{\Gamma_{\disc}}$. Parts~(2)-(4)
%   then follow from Lemma~\ref{lem:inductive description of B} by
%   pulling back via the diagram~(\ref{eqn:base-change description of
%     W}).


% To see that  $\cWW_{d,h,s} \hookrightarrow \cWW_{d,h,s'}$ %  ~{\em
%             % (\ref{eqn:Wach closed immersions for s})}
%           induces an
%           isomorphism on underlying reduced algebraic substacks, it is
%           enough to show that there is\mar{TG: currently stuck here
%             and concerned that it might not even be true!}
%   some~$s$ such that if~$s'\ge s$, $A$ is a
%   reduced $\Fp$-algebra, and~$\gM$ is a weak Wach module of
%   $T$-height~$\le h$ and level~$\le s'$, then it in fact has
%   level~$\le s$. Since~$A$ is an~$\Fp$-algebra, we have~$\gamma^{p^{s'}}-1=(\gamma-1)^{p^{s'}}$.

%   $W(k)\otimes_\Zp A$-module with an action of~$\gamma$
%   satisfying~$\gamma^{p^{s'}}=1$, then ~$\gamma^{p^s}=1$. We claim that
%   it is in fact enough to take~$p^s\ge d$. To see this, write
%   $B=W(k)\otimes_\Zp A=k\otimes_{\Fp}A$, which is reduced. % Working
%   % locally on~$\Spec B$, we may suppose that~$M$ is free.
%   Since~$B$ is
%   reduced, it embeds into the product of the algebraic closures of its
%   residue fields at prime ideals, so we may assume that~$B$ is an
%   algebraically closed field of characteristic~$p$. Then since
%   $(\gamma-1)^{p^{s'}}=\gamma^{p^{s'}}-1=0$, we deduce
%   that~$(\gamma-1)^d=0$, so that $\gamma^{p^s}-1=(\gamma-1)^{p^s}=0$,
%   as required.
% \mar{TG: here and at several points below I won't worry too much about
% justifying things like finite presentation for $p$-adic formal stacks;
% hopefully eventually this will be a bunch of citations to
% \cite{Emertonformalstacks}.}



% \end{proof}

% It follows from
% Proposition~\ref{prop:Gamma-disc modules as fixed points; Kisin
%   case}~(2) % ,
% % together with, \mar{ME: Have to add citation to discussion of equivariant 
% % 	morphisms of groupoids, once that's written TG: probably not
% %         needed now?}
% that there is morphism $\cC_{d,h}^{\Gamma_{\disc}} \to \Bun_d^{\Gamma_{\disc}}.$
% Furthermore, Remark~\ref{rem:testing continuity on M mod u}
% shows that we have the following $2$-commutative diagram:
% \numequation
% \label{eqn:base-change description of W}
% \xymatrix{\cWW_{d,h} \ar[r] \ar[d] & \cC_{d,h}^{\Gamma_{\disc}} \ar[d] \\
% \cB_d \ar^-{\text{(\ref{eqn:monomorphism for B})}}[r] & \Bun_d^{\Gamma_{\disc}} }
% \end{equation}

% \begin{df} We write $\cWW_{d,h,s} := \cB_{d,s} \times_{\cB_d} \cWW_{d,h},$
% 	and refer to $\cWW_{d,h,s}$ as the moduli stack of {\em rank $d$
% 		locally free Wach modules
% 		of height $\leq h$ and level $\leq s$.}
% 	%\mar{ME: I just made up some terminology here.}
% \end{df}


%\mar{ME: Once we sort out the definition of Wach module, there'll be
%	a similar discussion.  But now the analogue of $\cB$ will
%	be some {\em algebraic stack} classifying bundles on which
%	the $\Gamma$-action factors through some fixed finite quotient,
%	and so the corresponding moduli of Wach modules will be 
%	algebraic (not just formal) as well.  Actually, since these
%algebraic stacks are exactly the things whose inductive limit is
%the formal stack $\cB_d$, maybe I should just rewrite things slightly
%to introduce them as part of the above discussion. ME: Okay, I did this.}

% As noted in 
% Proposition~\ref{prop:Gamma-disc modules as fixed points; Kisin case}~(3),





\section{$\cX_{\lowercase{d}}$ is an Ind-algebraic stack}
\label{subsec: X is Ind algebraic}
% \mar{TG: this section is
% bogus, and needs to be rethought, to be more like \cite[Thm.\
%   5.3.15]{EGstacktheoreticimages} - we no longer have the kind of
%   \'etale local statements that were once in proper.tex, so we need
%   something a bit more elaborate. ME: Here and below, I think Thm.\ 5.3.15
% should be Thm.\ 5.4.15, or maybe even Thm.\ 5.4.16.}

% Inverting $T$ and $p$-adically completing induces a morphism
% \mar{ME: Maybe flesh out this rather abbreviated description
% of the morphism.}
% The morphism $\cC_{d,h} \to \cR_d$ of Section~\ref{subsec: EG stuff} is $\Gamma_{\disc}$-equivariant,
% and so induces a morphism
% \numequation
% \label{eqn:C to R}
% \cC_{d,h}^{\Gamma_{\disc}} \to \cR_d^{\Gamma_{\disc}}.
% \end{equation}
% This restricts to a morphism
% \numequation
% \label{eqn:W to X}
% \cWW_{d,h} \to \cX_d.
% \end{equation}
% Since an action of $\Gamma_{\disc}$ on a $\varphi$-module
% $\gM$ is continuous if and only if it is continuous on the 
% underlying \'etale $\varphi$-module $\widehat{\gM[1/T]},$
% we see that the $2$-commutative square
% \numequation
% \label{eqn:C to R square}
% \xymatrix{\cWW_{d,h} \ar[r]\ar_-{\text{(\ref{eqn:W to X})}}[d] &
% 	\cC_{d,h}^{\Gamma_{\disc}} \ar^-{\text{(\ref{eqn:C to R})}}[d] \\
% \cX_d \ar[r] & \cR_d^{\Gamma_{\disc}} }
% \end{equation}
% is $2$-Cartesian.

%\begin{lem}
%  \label{lem:X is limit preserving}$\cX_d$ is limit preserving.
%\end{lem}
%\begin{proof}
%  Since~$\cX_d\subset\cR_d^{\Gamma_\disc}$ is fully faithful, and~$\cR_d^{\Gamma_\disc}$ is limit
%  preserving by Theorem~\ref{thm: C is a $p$-adic formal algebraic stack and related
%    properties}, it suffices to prove that~$\cX_d$ is limit preserving
%  on objects. Since~$\cR_d^{\Gamma_\disc}$ is limit preserving on
%  objects, we are reduced to showing that if~$T=\varinjlim_iT_i$ is a
%  limit of affine schemes, and  $T_i\to\cR_d^{\Gamma_\disc}$ is a morphism with
%  the property that the composite \[T\to T_i\to\cR_d^{\Gamma_\disc}\]
%    factors through~$\cX_d$, then for some~$j\ge i$, the composite \[T_j\to T_i\to\cR_d^{\Gamma_\disc}\]
%    factors through~$\cX_d$.
%    We may check continuity after reduction modulo $p$ \mar{ME: Have to prove this reduction mod $p$ claim.},
%    and so this follows from Lemma~\ref{lem:continuity criterion}
%    and~\ref{lem:quasi-linear limit preserving}.
%%
%%    To this end, write~$T_i=\Spec A_i$, $T_j=\Spec A_j$, $T=\Spec T$,
%%    and let~$M$ be the \'etale $(\varphi,\Gamma_{\disc})$-module
%%    with~$A_i$-coefficients corresponding to the morphism
%%    $T_i\to\cR_d^{\Gamma_{\disc}}$. Let~$\gM$ be a lattice in~$M$, so that by
%%    Lemma~\ref{lem:testing continuity on M mod T}, there is an
%%    integer~$s\ge 0$ such that $(\gamma^{p^s}-1)(\gM_A)\subset T\gM_A$,
%%    where~$\gM_A:=\gM\otimes_{\gS_{A_i}}\gS_A$. \mar{TG: again I'm not
%%    sure what the best way is to finish from here, but the idea is
%%    simple: we are asking that some map to $M/\gM$ vanishes, and we
%%    know it after going up to~$A$, so it will be true after going to
%%    some~$A_j$, because ``the Laurent tails are finite'', or perhaps
%%    more precisely because of some finite presentation statement for
%%    $\Hom$s of \'etale $\varphi$-modules? As always I'm a bit confused
%%  about $\Gamma$-semi-linearity.}
%\end{proof}

%\mar{TG: still need to revise the rest of this section}


Continue to assume that~$K$ is basic. By definition, we have a $2$-Cartesian diagram
\numequation
\label{eqn:C to R square}
\xymatrix{\cWW_{d,h} \ar[r]\ar[d] &
	\cR_d^{\Gamma_{\disc}}\times_{\cR_d}\cC_{d,h} \ar[d] \\
\cX_d \ar[r] & \cR_d^{\Gamma_{\disc}} }
\end{equation}
If $h' \geq h$ then
the closed immersion $\cC_{d,h} \hookrightarrow \cC_{d,h'}$ is
compatible with the morphisms from each of its source and target
to $\cR_d$, 
%$\Gamma_{\disc}$-equivariant, %  and so % (by )\mar{ME: Cite non-existent discussion
% 	% of equivariant closed immersions here TG: probably don't need
%         % anything now?}
% induces a closed immersion
% $\cC_{d,h}^{\Gamma_{\disc}} \hookrightarrow \cC_{d,h'}^{\Gamma_{\disc}}.$
% This is compatible with the morphisms of source and target 
% to $\cR_d$, and
and so % pulling back over $\cR_d^{\Gamma_{\disc}}$,
we obtain a closed immersion 
\numequation
\label{eqn:closed immersion of Wachs}
\cWW_{d,h} \hookrightarrow \cWW_{d,h'}.
\end{equation}


By construction,
the morphisms $\cWW_{d,h}\to\cX_d$  are compatible,
as $h$ varies,
with the closed immersions~(\ref{eqn:closed immersion of Wachs}). 
Thus we also obtain a morphism
\numequation
\label{eqn:Ind W to X}
\varinjlim_h \cWW_{d,h} \to \cX_d.
\end{equation}
Roughly speaking, we will prove that~$\cX_d$ is an Ind-algebraic stack
by showing that it is the ``scheme-theoretic image'' of the morphism
$\varinjlim_h \cWW_{d,h} \to \cR_d^{\Gamma_{\disc}}$ induced
by~\eqref{eqn:Ind W to X}. More precisely, choose $s \geq 0,$ and consider the composite
        \numequation
        \label{eqn:W to R composite}
        \cWW_{d,h,s} \to \cWW_{d,h} \to \cX_d \to \cR_d^{\Gamma_{\disc}}.
        \end{equation}
        This admits the alternative factorization
        $$\cWW_{d,h,s} \to \cWW_{d,h} \to 	\cR_d^{\Gamma_{\disc}}\times_{\cR_d}\cC_{d,h} 
        \to \cR_d^{\Gamma_{\disc}}.$$
        Proposition~\ref{prop:inductive description of W} shows
        that the composite of the first two arrows is a closed 
        embedding of finite presentation,
        while  Corollary~\ref{cor: basic properties of C and R p adic stacks} % Lemma~\ref{lem:C to R is proper}
        % Theorem~\ref{thm: C is a $p$-adic formal algebraic stack and related
    % properties}\mar{TG: might need to augment this reference, as it's
    % not about fixed point stacks as stated.}
        shows that the third arrow is representable by algebraic spaces, proper, 
        and of finite presentation. 
        Thus~(\ref{eqn:W to R composite}) is representable by
        algebraic spaces, proper, and of
        finite presentation.   

	Fix an integer $a \geq 1,$ and write $\cWW^a_{d,h,s} := \cWW_{d,h,s}
	\times_{\Spf \cO} \Spec \cO/\varpi^a$.   
	Proposition~\ref{prop:inductive description of W}
	shows that $\cWW_{d,h,s}$ is a $p$-adic formal
	algebraic stack of finite presentation over $\Spf \cO$,
	and so $\cWW^a_{d,h,s}$ is an algebraic stack, 
	and a closed substack of $\cWW_{d,h,s}$.

\begin{df}  We let $\cX^a_{d,h,s}$ denote the \index{$\cX^a_{d,h,s}$}
scheme-theoretic image of the composite 
	\numequation
	\label{eqn:another composite}
	\cWW^a_{d,h,s} \hookrightarrow \cWW_{d,h,s}
	\buildrel \text{(\ref{eqn:W to R composite})} \over
	\longrightarrow \cR_d^{\Gamma_{\disc}},
        \end{equation}
defined via the formalism of scheme-theoretic images for morphisms of Ind-algebraic
stacks developed in Appendix~\ref{app: formal algebraic stacks}.
\end{df}

More concretely,
	since $\cR_d^{\Gamma_{\disc}}$ is an Ind-algebraic stack, % inductive limit
	% of algebraic stacks under closed immersions, 
        constructed as the $2$-colimit of a directed system of algebraic
        stacks whose transition morphisms are closed immersions,
	the morphism~\eqref{eqn:another composite},
	which is representable by algebraic spaces, proper, and of finite presentation,
	factors through a closed algebraic substack $\cZ$ of
	$\cR_d^{\Gamma_{\disc}}.$   We then define $\cX^a_{d,h,s}$
	to be the 
	% \mar{TG: is ``it'' $\cWW^a_{d,h,s}$? ME: Yes, clarified.}
        scheme-theoretic image of
	$\cWW^a_{d,h,s}$ in $\cZ$.   Note that
	$\cX^a_{d,h,s}$ % \mar{TG: $\cZ$ should be $\cX^a_{d,h,s}$? ME: Yes, fixed.}
        is a closed algebraic substack of $\cR_d^{\Gamma_{\disc}}$,
        and is independent of the choice of~$\cZ$.

\begin{remark}
As already observed above, the morphism~\eqref{eqn:another composite} 
factors through $\cX_d$. However, since at this point we don't know that $\cX_d$
is Ind-algebraic, we can't directly define a scheme-theoretic image of 
$\cW^a_{d,h,s}$ in $\cX_d$.  It might be possible to do this using
the formalism of \cite{EGstacktheoreticimages}; since~\eqref{eqn:another composite} 
is proper, this scheme-theoretic image would then coincide with $\cX^a_{d,h,s}$.
We don't do this here; but we do prove somewhat more directly,
in Lemma~\ref{lem: scheme theoretic images factor through X} below,
that $\cX^a_{d,h,s}$ is a substack of~$\cX_d$.
\end{remark}

        As in Definition~\ref{defn: lattice}, a \emph{lattice} $\gM$ \index{lattice}
        in a projective \'etale $\varphi$-module~$M$ is a finitely
        generated ~$\A_{K,A}^+$-submodule of~$M$ whose $\A_{K,A}$-span
        is~$M$. Note that~$\gM$ is not assumed to be projective.
	\begin{lem}
		\label{lem:Artinian points of scheme theoretic images}
		Suppose that $M$ is a projective \'etale $\varphi$-module of rank $d$
		over a finite type Artinian $\cO/\varpi^a$-algebra $A$,
		and that~$M$ is endowed with an action of $\Gamma_{\disc}$,
		such that the corresponding morphism $\Spec A \to
		\cR_d^{\Gamma_{\disc}}$ 
		factors through $\cX^a_{d,h,s}$.
		Then $M$ contains a $\varphi$-invariant
		lattice % {\em (}in the sense of Definition~{\em \ref{defn: lattice})}
                $\gM$ of $T$-height $\leq h$,
		such that $(\gamma^{p^s}-1) (\gM) \subseteq
		T\gM.$
	\end{lem}
	\begin{proof}Since an Artinian ring is a direct product of
          Artinian local rings, it suffices to treat the case that~$A$
          is local.
          % By~\cite[Lem.\ 3.2.29, Defn.\
          % 3.2.5]{EGstacktheoreticimages}, the assumption that~$\Spec
          % A\to\cR_d^{\Gamma_{\disc}}$ factors through~$\cX_{d,h,s}^a$
          % implies that there is an Artinian local
          % $\cO/\varpi^a$-algebra~$B$, such that $\Spec
          % A\to\cR_d^{\Gamma_{\disc}}$ factors through a closed
          % immersion~$\Spec A\into\Spec B$, and the morphism
          % $\cW_{d,h,s}\times_{\cR_d^{\Gamma_{\disc}}}\Spec B\to \Spec
          % B$ is scheme-theoretically dominant. Replacing~$A$ by~$B$,
          % we can and do assume from now on that the morphism
          % $\cW_{d,h,s}\times_{\cR_d^{\Gamma_{\disc}}}\Spec A\to \Spec
          % A$ is scheme-theoretically dominant. (Note that the existence
          % of a lattice of the required form for the \'etale
          % $\varphi$-module $M_B$ corresponding to $\Spec
          % B\to\cR_d^{\Gamma_{\disc}}$ implies the existence of one
          % for~$M$, simply by taking the image of the lattice 
	  % under the natural surjection $M_B \to \A_{K,A}\otimes_{\A_{K,B}}
	  % M_B \cong M.$)
          Let the residue field of~$A$ be~$\F'$, a finite
          extension of~$\F$, and write~$\cO'$ for the ring of integers
          in the compositum of~$E$ and~$W(\F')[1/p]$, so that~$\cO'$
          has residue field~$\F'$; note that $A$ is naturally on $\cO'$-algebra. 

          The projective
          \'etale $\varphi$-module $M$ is in fact free of rank~$d$,
          since $\A_{K,A}$ is again a local ring, and we fix an (ordered)
          basis for $M$ as an $\A_{K,A}$-module.
          Write~$M_{\F'} := \A_{K,\F'} \otimes_{\A_{K,A}} M$; the $\A_{K,A}$-basis
          of $M$ gives rise to a corresponding $\A_{K,\F'}$-basis of~$M_{\F'}$.
          %Let~$M_{\F'}$ be the projective (equivalently, free)
          %\'etale $\varphi$-module of rank~$d$ over~$\A_{K,\F'}$
          %(endowed with an action of~$\Gamma_{\disc}$) that
          %corresponds to the composite
          %morphism $\Spec \F'\to\Spec A\to\cR_d^{\Gamma_{\disc}}$.
          Let~$\cC_{\cO'}$ be the category of Artinian local
          $\cO'/\varpi^a$-algebras for which the structure map induces
          an isomorphism on residue fields.  %Fix a basis of~$M_{\F'}$; then
By a \emph{lifting} of~$M_{\F'}$ to an object ~$\Lambda$
         of~$\cC_{\cO'}$, we mean a triple consisting of an \'etale
 $\varphi$-module~$M_\Lambda$ which is free of rank~$d$, a choice of (ordered)
 $\A_{K,\Lambda}$-basis of~$M_\Lambda$, and an isomorphism $M_\Lambda\otimes_\Lambda
 \F'\cong M_{\F'}$ of \'etale $\varphi$-modules which takes the chosen basis
 of~$M_\Lambda$ to the fixed basis of~$M_{\F'}$. Let~$D$ be the functor
          $\cC_{\cO'}\to\Sets$ taking~$\Lambda$ to the set of isomorphism
          classes of liftings of~$M_{\F'}$ to an \'etale $\varphi$-module~$M_\Lambda$
          with $\Lambda$-coefficients,
	  endowed with an action of~$\Gamma_{\disc}$ lifting
	  that on~$M_{\F'}$.
          Note that $A$  is an object of $\cC_{\cO'}$, and that
our originally chosen basis for $M$ is classified by a 
          continuous morphism $R  \to A.$

          The functor~$D$ is pro-representable by an object~$R$
          of~$\pro\cC_{\cO'}$, by the same argument as in the
          proof % \mar{TG: note that we should fix reference numbers once
          % the dust settles in proper.tex.}
          of~\cite[Prop.\ 5.3.6]{EGstacktheoreticimages}: namely, 
          by Grothendieck's representability theorem, it suffices to
          prove the compatibility of~$D$ with fibre products
          in~$\cC_{\cO'}$, which is obvious. The universal lifting
          $M_R$ gives a morphism $\Spf R\to\cR_d^{\Gamma_{\disc}}$. 
          The composite  $\Spec A \to \Spf R \to \cR_d^{\Gamma_{\disc}}$ 
          is of course just the morphism that classifies~$M$.

We let~$D'$ denote the
          subfunctor of~$D$ consisting of those lifts for which there
          is a lattice  ~$\gM_\Lambda$ %  with
          % $\Lambda$-coefficients,
          of $F$-height~$\le h$ and level~$\le s$,
          with~$\gM_\Lambda[1/T]=M$. (More precisely, we require that
          $T^h\gM_\Lambda\subseteq (1\otimes\varphi_{M_\Lambda})(\varphi^*\gM_\Lambda)$  and
          $(\gamma^{p^s}-1)(\gM_\Lambda)\subseteq T\gM_\Lambda$.) We claim that the functor~$D'$ is
          pro-representable by a quotient~$S$ of~$R$. To see this, it
          is again enough to show that~$D'$ preserves fibre products,
          and in turn it is enough to show that the property of an
          \'etale $\varphi$-module~$M$ with an action
          of~$\Gamma_{\disc}$ admitting a ~$\varphi$-stable
          lattice~$\gM$ of $T$-height~$\le h$, and such that
          $(\gamma^{p^s}-1)(\gM)\subseteq T\gM$, is stable under
          taking direct sums and subquotients. This is obvious for
          direct sums, and the case of subquotients follows as
          in~\cite[Lem.\ 5.3.10]{EGstacktheoreticimages} (which is the
          same result without the conditions
          on~$\Gamma_{\disc}$). More precisely, once checks that if we
          have a short exact sequence $0\to M'\to M\to M''\to 0$ of
          \'etale $\varphi$-modules with $\Lambda$-coefficients and
          actions of~$\Gamma_{\disc}$, and $\gM$ is a lattice of the
          appropriate kind in~$M$, then the kernel and image of the
          map $\gM\to M''$ give the appropriate lattices~$\gM'$,
          $\gM''$ in~$M'$, $M''$ respectively. (The properties that
          $(\gamma^{p^s}-1)(\gM')\subseteq T\gM'$ and
          $(\gamma^{p^s}-1)(\gM'')\subseteq T\gM''$ follow from a
          short diagram chase, using that the composite
          $\gM\stackrel{(\gamma^{p^s}-1)}{\to}M\to M/T\gM$ vanishes by
          hypothesis.)% \mar{TG: I had imagined that as in the revisions
          % to proper.tex, we would want to take some quotient to cut
          % down to something Noetherian, but actually I don't see a
          % need to right now. Leaving this comment here as a reminder
          % to be careful, though!}

          % By Lemma~\ref{alem: scheme theoretic image of versal is
          %   versal}, \mar{ME: What does this lemma have to do with anything here?}
 By the definition of~$D'$, the statement of the lemma is equivalent to
          the statement that the %tautological
map $R\to A$ factors
          through~$S$. Let $X=\cW_{d,h,s}\times_{\cR_d^{\Gamma_{\disc}}}\Spf
          R$, a formal algebraic space, and let~$\Spf T$ be the
          scheme-theoretic image of the morphism $X\to\Spf
          R$. Since the morphism~$\Spec A\to\cR_d^{\Gamma_{\disc}}$
          factors through $\cX^a_{d,h,s}$ by hypothesis, the morphism
          $\Spec A\to\Spf R$ factors through~$\Spf T$, so it is enough
to show that~$\Spf T$ is a closed formal subscheme of~$\Spf S$.
          %\mar{TG: is this actually clear now? I think
    %         this is implicitly using that the versal ring for the
    %         scheme-theoretic image is the scheme-theoretic image of
    %         the versal rings. I wondered if we might have to go back
    %         to something Noetherian, but actually it looks like we can
    %       probably just cite \cite[Lem.\
    %       3.2.16]{EGstacktheoreticimages}; need to make sure that this
    %     is legitimate; though, as e.g.\ probably in that setting the
    %     source of $\cX\to\cF$ is always assumed to be algebraic, not
    %     formal algebraic. Probably something needs to be added to the
    %     discussion around Lemma~\ref{lem: criterion for Artin to map to scheme theoretic
    % image}.}
By Lemma~\ref{lem: criterion for Artin to map to scheme theoretic
    image}, 
% 	  \mar{ME: Should we import 5.4.15 into this paper? The idea
% 	  would be to also include the step of reduction-to-scheme-theoretic
%   dominance that is used at the beginning of this proof and also in
% the proof of 4.8.8}
          % taking~$X=\cW_{d,h,s}\times_{\cR_d^{\Gamma_{\disc}}}\Spf
          % R$. To see this, we need to
it in turn suffices to          show that if~$A'$ is a finite
          type Artinian local~$R$-algebra for which $R\to A'$ factors
          through a discrete quotient of~$R$, and for which the
          morphism $\cW_{d,h,s}\times_{\cR_d^{\Gamma_{\disc}}}\Spec
          A'\to\Spec A'$ admits a section, then~$R\to A'$ factors
          through~$S$.

          % As in the proof of~\cite[Prop.\
          % 5.4.14]{EGstacktheoreticimages}, we can (b
          By replacing~$R$
          with~$R\otimes_{W(\F')}\kappa(A')$ we can reduce to the case
          that~$A'$ has residue field~$\F'$ (\emph{cf.}\ \cite[Cor.\ 5.3.18]{EGstacktheoreticimages}), so that $R\to A'$ factors
          through~$S$ if and only if the \'etale $\varphi$-module
          corresponding to $\Spec A'\to \cR_{d}^{\Gamma_{\disc}}$
          admits a $\varphi$-stable lattice~$\gM_{A'}$ of
          $T$-height~$\le h$ for
          which~$(\gamma^{p^s}-1)(\gM_{A'})\subseteq T\gM_{A'}$. But
          a section to the morphism $\cW_{d,h,s}\times_{\cR_d^{\Gamma_{\disc}}}\Spec
          A'\to\Spec A'$ gives us a morphism $\Spec A'\to\cW_{d,h,s}$, and the
          corresponding weak Wach module provides the required lattice.
          \end{proof}
          
		% \mar{ME: Sketch has to be fleshed out! TG: I will do
                %   this (and the next one) soon.}
	% 	Proof sketch:  We define a subring of the versal
	% 	ring to $\cR_d^{\Gamma_{\disc}}$ classifying 
	% 	\'etale $\varphi$-modules with a $\Gamma_{\disc}$-action
	% 	which admit a height $\leq h$ Kisin module $\gM$ such
	% 	that $(\gamma^{p^s}-1)(\gM) \subseteq T\gM.$  This will
	% 	be Noetherian, and using \cite[Lem.~5.4.15]{EGstacktheoreticimages} (the ``axiomatic lemma'') we will see that any Artinian point
	% 	of $\cX^a_{d,h,s}$ factors through this ring, and so
	% 	admits a lattice of the required type.
	% \end{proof}

        \begin{lem}
  \label{lem: scheme theoretic images factor through X}% \mar{TG: I
    % don't yet know how to prove this!}
  Each~$\cX^a_{d,h,s}$ is a
  closed substack of~$\cX_d$.
\end{lem}
\begin{proof}
  Since $\cX^a_{d,h,s}$ is a closed substack of $\cR_d^{\Gamma_{\disc}}$,
  if it is a substack of $\cX_d$, it will in fact
  be a closed substack.
  Thus it suffices to show that $\cX^a_{d,h,s}$ is indeed a substack
  of~$\cX_d$.
  Since~$\cX^a_{d,h,s}$ is limit preserving, it is enough to check
  that if~$A$ is a finite type $\cO/\varpi^a$-algebra, then for any
  morphism $\Spec A\to\cX^a_{d,h,s}$, the composite morphism
  $\Spec A\to\cX^a_{d,h,s}\to\cR_d^{\Gamma_{\disc}}$ factors
  through~$\cX_d$.  Equivalently,
  if $M$ denotes the \'etale $\varphi$-module over~$A$, 
  endowed with a $\Gamma_{\disc}$-action,
  associated to the given point $\Spec A \to \cR_d^{\Gamma_{\disc}},$
  then we must show that the $\Gamma_{\disc}$-action
  on $M$ is continuous. %\mar{TG: the rest of this argument is new,
%    please check carefully! ME: It looked good to me.}

  By Lemma~\ref{lem:testing continuity on M mod T}, to see that the
  $\Gamma_{\disc}$-action on $M$ is continuous, it suffices to produce
  a lattice $\gM \subseteq M$ such that
  $(\gamma^{p^s}-1)(\gM)\subseteq T\gM$. We will produce such a
  lattice by reduction to the Artinian case, as follows. Let
  $\{A_i\}_{i \in I}$ be the directed system of Artinian quotients of
  $A$.  Since $A$ is of finite type over $\cO/\varpi^a$, and so
  Noetherian, the natural map $A \to B := \prod_i A_i$ is injective.
  Lemma~\ref{lem:Artinian points of scheme theoretic images} shows
  that each base-changed module $M_{A_i}$ admits a $\varphi$-invariant
  lattice $\gM_i$ of $T$-height~$\le h$ satisfying the condition
  $(\gamma^{p^s}- 1)(\gM_i) \subseteq T \gM_i.$

  Write~$\gM_B:=\prod_i\gM_i\subset M_B$, and set
  $\gM:=M\cap\gM_B\subset M_B$. Since  $(\gamma^{p^s}- 1)(\gM_B)
  \subseteq T \gM_B$, we have  $(\gamma^{p^s}- 1)(\gM) \subseteq T
  \gM$, so it only remains to check that~$\gM$ is a lattice. To see
  this, let~$\gM'$ be a finite height lattice in~$M$ (i.e.\ a lattice
  in $M$ which is~$\varphi$-stable, and for which the cokernel
  of~$\Phi$ is killed by a power of~$T$; such a lattice exists
  by~\cite[Lem.\ 5.2.15]{EGstacktheoreticimages}). By
  Lemma~\ref{lem:lattice properties}~(3), it suffices to prove that
  there is an integer~$l\ge 0$ such that $T^l\gM'\subseteq
  \gM\subseteq T^{-l}\gM$. Accordingly, if for each~$i$ we
  let~$\gM'_i$ denoted the image of~$\gM_i\otimes_AA_i$ in~$M_{A_i}$,
  we need to show that we have $T^l\gM'_i\subseteq
  \gM_i\subseteq T^{-l}\gM_i'$.

  The existence of such an~$l$ is established in the course of the
  proof of~\cite[Lem.\ 5.3.14]{EGstacktheoreticimages}, although this
  is not explicitly recorded there. For the convenience of the reader,
  we recall the argument in our present setting. Increasing~$h$ if
  necessary, we may assume that~$\gM'$ has $T$-height at most~$h$. It
  follows that~$\gM_i'$ and~$\gM_i$ are each $\varphi$-stable lattices
  of height at most~$h$ in~$M_{A_i}$. We claim that we may
  take~$l=\lfloor (h+ap)/(p-1)\rfloor$.

  To see this, let~$j$ be minimal such that $T^j\gM_i\subseteq\gM_i'$;
  then we have
  \[\Phi_{\gM_i}(\varphi^*\gM_i)\subseteq \gM_i\subseteq
    T^{-j}\gM'_i\subseteq T^{-h-j}\Phi_{\gM'_i}(\varphi^*\gM'_i),\] so
  that $T^{h+j}\varphi^*\gM_i\subseteq\varphi^*\gM'_i$. It follows
  from~\cite[Lem.\ 5.3.13]{EGstacktheoreticimages} that $h+j>(j-a)p$,
  so that $j<(h+ap)/(p-1)$, and~$j\le l$ by definition.  Thus
  $T^l\gM_i\subseteq\gM_i'$, as claimed. Reversing the roles
  of~$\gM_i$ and~$\gM'_i$, we have $T^l\gM'_i\subseteq\gM_i$, and we
  are done.
\end{proof}    
  % For each~$A_i$, we
  % let~$S(A_i)$ denote the set of such lattices. By~\cite[Lem.\
  % 5.3.14]{EGstacktheoreticimages}, the inverse system~$\{S(A_i)\}$
  % satisfies the Mittag-Leffler property, so we \mar{TG: this argument
  %   is incomplete! It seems that it shouldn't be too hard to fix,
  %   though. We need a consequence of the proof of ~\cite[Lem.\
  % 5.3.14]{EGstacktheoreticimages}, which is that lattices of bounded
  % height are a bounded distance apart, independently of the
  % coefficients (i.e.\ depending only on the exponent of the
  % coefficients). So then what we do is choose any lattice in~$M$ (not
  % necc.\ $\varphi$-stable or $\gamma$-stable), and compare this fixed
  % lattice to the $\gM_i$, and so see that the product of the $\gM_i$
  % is a lattice over $B[[T]]=(\prod_iA_i)[[T]]$.}
  % % see that
  % % \mar{ME: The following argument should probably be fleshed out.}
  % % Arguing as in \cite[Lem.\ 5.3.12]{EGstacktheoreticimages},
  % % i.e.\ using Mittag--Leffler and the theory of minimum/maximum
  % % height $\leq h$ lattices,
  % find that $M_B$ admits a lattice $\gM_B$
  % (which is $\varphi$-invariant, and of height $\leq h$, although we don't need this)
  % such that $(\gamma^{p^s}-1)(\gM_B)\subseteq T\gM_B$.
  % Intersecting $\gM_B$ with $M$, inside $M_B$,
  % we obtain (by Lemma~\ref{lem:intersecting with lattice})
  % a lattice $\gM \subseteq M$ 
  % such that $(\gamma^{p^s}-1)(\gM)\subseteq T\gM$.
  % Lemma~\ref{lem:testing continuity on M mod T}
  % then shows that the $\Gamma_{\disc}$-action
  % on $M$ is continuous, as required.

  

  % Let $\{A_i\}_{i \in I}$ be the directed system of Artinian
  % quotients of $A$.   Since $A$ is of finite type over $\cO/\varpi^a$,
  % and so Noetherian, the natural map
  % $A \to B := \varprojlim_i A_i$
  % is injective.
  % Lemma~\ref{lem:Artinian points of scheme theoretic images}
  % shows that each base-changed module $M_{A_i}$ admits
  % a $\varphi$-invariant lattice $\gM_i$ of $T$-height~$\le h$ satisfying the condition
  % $(\gamma^{p^s}- 1)(\gM_i) \subseteq T \gM_i.$ For each~$A_i$, we
  % let~$S(A_i)$ denote the set of such lattices. By~\cite[Lem.\
  % 5.3.14]{EGstacktheoreticimages}, the inverse system~$\{S(A_i)\}$
  % satisfies the Mittag-Leffler property, so we \mar{TG: this argument
  %   is incomplete! It seems that it shouldn't be too hard to fix,
  %   though. We need a consequence of the proof of ~\cite[Lem.\
  % 5.3.14]{EGstacktheoreticimages}, which is that lattices of bounded
  % height are a bounded distance apart, independently of the
  % coefficients (i.e.\ depending only on the exponent of the
  % coefficients). So then what we do is choose any lattice in~$M$ (not
  % necc.\ $\varphi$-stable or $\gamma$-stable), and compare this fixed
  % lattice to the $\gM_i$, and so see that the product of the $\gM_i$
  % is a lattice over $B[[T]]=(\prod_iA_i)[[T]]$.}
  % % see that
  % % \mar{ME: The following argument should probably be fleshed out.}
  % % Arguing as in \cite[Lem.\ 5.3.12]{EGstacktheoreticimages},
  % % i.e.\ using Mittag--Leffler and the theory of minimum/maximum
  % % height $\leq h$ lattices,
  % find that $M_B$ admits a lattice $\gM_B$
  % (which is $\varphi$-invariant, and of height $\leq h$, although we don't need this)
  % such that $(\gamma^{p^s}-1)(\gM_B)\subseteq T\gM_B$.
  % Intersecting $\gM_B$ with $M$, inside $M_B$,
  % we obtain (by Lemma~\ref{lem:intersecting with lattice})
  % a lattice $\gM \subseteq M$ 
  % such that $(\gamma^{p^s}-1)(\gM)\subseteq T\gM$.
  % Lemma~\ref{lem:testing continuity on M mod T}
  % then shows that the $\Gamma_{\disc}$-action
  % on $M$ is continuous, as required.
%
%
%
%
%
%  \mar{ME: I think everything that follows this comment will be removed.}
%  As in the proof of~\cite[Lem.\
%  3.2.25]{EGstacktheoreticimages}, it
%  suffices in turn to prove this for every local Artinian quotient
%  of~$A$, so we may assume that~$A$ is Artin local. Since
%  $\cW_{d,h,s}^a\to\cR_d^{\Gamma_{\disc}}$ is a proper morphism whose
%  target is limit preserving, it follows from~\cite[Lem.\
%  3.2.29]{EGstacktheoreticimages}\mar{TG: this is a bit of a
%    convoluted way of arguing; if it's correct, I don't know if it
%    would be easier to just extract whatever is needed from the
%    arguments in~\cite[\S 3.2]{EGstacktheoreticimages}, as we're only
%    working with algebraic stacks right here?} that the 
%  definition of~$\cX^a_{d,h,s}$ coincides with the notion of
%  scheme-theoretic image of proper morphisms defined in~\cite[Defn.\
%  3.2.6]{EGstacktheoreticimages},  and there is therefore a closed
%  immersion $\Spec A\to \Spec B$, where~$B$ is Artin local, and a
%  factorisation \[\Spec A\to \Spec B\to\cR_d^{\Gamma_{\disc}},\]so that the pullback
%  $(\cW_{d,h,s}^a)_B\to\Spec B$ is scheme-theoretically dominant.
%
%  %Since~$B$ is Artin local, so is~$B((T))$, and the \'etale
%  %$\varphi$-module~$M_B$ corresponding to the morphism
%  %$\Spec B\to\cR_d^{\Gamma_{\disc}} $ is therefore automatically
%  %free.
%  Let~$\F'$ be the residue field of~$B$, and write~$M_{\F'}$ for
%  the base change of~$M_B$ to~$\F'$; since  $(\cW_{d,h,s}^a)_B\to\Spec B$ is scheme-theoretically
%  dominant, we see in particular that the fibre of~$\cW_{d,h,s}^a$
%  over~$\Spec\F'$ is non-empty, so that there is a projective weak
%  Wach module~$\gM_{F'}$ of height at most~$h$ and level at most~$s$,
%  with~$\gM_{\F'}[1/T]=M_{\F'}$. In particular, it follows from
%  Lemma~\ref{lem:testing continuity on M mod T} that the action
%  of~$\Gamma_\disc$ on~$M_{\F'}$ extends to a continuous action
%  of~$\Gamma$ (in other words, the morphism $\Spec
%  \F'\to\cR_d^{\Gamma_{\disc}}$ factors through~$\cX_d$).
%
%  Now let~$\gM_B$ be any lattice in~$M_B$. Applying
%  Lemma~\ref{lem:testing continuity on M mod T}, we see that there
%  exists an~$s'$ such
%  that \[(\gamma^{p^{s'}}-1)(\gM_B\otimes_B\F')\subset T
%    (\gM_B\otimes_B\F'),\] so that \[(\gamma^{p^{s'}}-1)(\gM_B)\subset
%    (T,\m_B)\gM_B,\] where~$\m_B$ is the maximal ideal
%  of~$B$. Since~$p\in\m_B$, it follows exactly as in the proof of
%  Lemma~\ref{lem:testing continuity on M mod T} that we have \[(\gamma^{p^{s'+n}}-1)(\gM_B)\subset
%    (T,\m_B)^n\gM_B\]for each~$n\ge 0$. Choosing~$n$ such
%  that~$\m_B^n=0$, we have $(\gamma^{p^{s'+n}}-1)(\gM_B)\subset
%  T\gM_B$, and by one more application of Lemma~\ref{lem:testing
%    continuity on M mod T}, it follows that the morphism $\Spec
%  B\to\cR_d^{\Gamma_{\disc}} $ factors through~$\cX_d$. In
%  particular, the morphism $\Spec
%  A\to\cR_d^{\Gamma_{\disc}} $ factors through~$\Spec\cX_d$, as required.

%   It follows that after possibly increasing~$h$, we may lift it
%   to a morphism
%   $\Spec B\to \cR_d^{\Gamma_{\disc}}\times_{\cR_d}\cC_{d,h}$. We
%   therefore have a commutative diagram \numequation
%   \xymatrix{(\cWW_{d,h,s})_{B}\ar[dr]\ar[r] & \cWW_{d,h,s}
%     \ar[r]\ar[d] &
%     \cR_d^{\Gamma_{\disc}}\times_{\cR_d}\cC_{d,h} \ar[d] \\
%     &\Spec B \ar[r]\ar[ur] & \cR_d^{\Gamma_{\disc}} }\end{equation}
% where the upper right horizontal arrow is a closed immersion, and the
% left hand diagonal arrow is scheme-theoretically dominant.\mar{TG:
%   this is as far as I got!}


% \begin{lem}
%   \label{lem: deducing finite level from closed point}
% \end{lem}
        
	% Since~(\ref{eqn:another composite}) factors through $\cX_d$,
	% we see that $\cX^a_{d,h,s}$ is in fact a closed algebraic substack
	% of $\cX_d$.\mar{ME: Cite a result from {\tt proper.tex} for
        %   this. TG: this isn't completely clear to me! \cite[Thm.\
        %   3.2.34]{EGstacktheoreticimages} might apply if we know that
        %   $\cX_d\into\cR_d^{\Gamma_{\disc}}$ was of finite type - do
        %   we know this? But that wouldn't tell us that
        %   $\cX^a_{d,h,s}$ is closed in $\cX_d$. Hopefully this is all easy\dots}
	
	% Noting that $\varinjlim_h \cWW_{d,h} = \varinjlim_{h,s,a} \cWW_{d,h,s}^a$
	% (by Proposition~\ref{prop:inductive description of W}), 
	% we deduce from Lemma~\ref{lem:surjectivity}
	% that the morphism $\varinjlim_{h,s,a} \cX^a_{d,h,s} \to \cX_d$
	% is an isomorphism.   This isomorphism realizes $\cX_d$ as
	% an Ind-algebraic stack.\mar{ME: Add more details; closed immersions,
	% 	finite presentation, etc.}



% \begin{lemma}
% 	\label{lem:C to R is proper}
% 	Each of the morphisms~{\em (\ref{eqn:C to R})}
% 	and~{\em (\ref{eqn:W to X})} is representable by algebraic spaces,
% 	proper, and of finite presentation.
% %	\mar{ME: Think about whether this is actually true.}
% \end{lemma}
% \begin{proof}
% 	Since~(\ref{eqn:C to R square}) is $2$-Cartesian,  
% 	it suffices to prove the lemma for the morphism~(\ref{eqn:C to R}).
% 	To this end,
% 	we consider the $2$-commutative diagram
% 	$$\xymatrix{
% 		\cC_{d,h}^{\Gamma_{\disc}} \ar[r] \ar[rd] &
% 		\cR_d^{\Gamma_{\disc}}\times_{\cR_d} \cC_{d,h} \ar[r] \ar[d]
% 	       	& \cC_{d,h} \ar[d] \\
% 	        & \cR_d^{\Gamma_{\disc}} \ar[r] & \cR_d}
% 	$$
% 	in which the square is $2$-Cartesian.  Since the right-hand vertical
% 	arrow is representable by algebraic spaces, proper, and of finite
%         presentation (by Theorem~\ref{thm: C is a $p$-adic formal algebraic stack and related
%     properties}), % \mar{TG: is representability, rather than
%           % Ind-representability, obvious? Does it follow from
%           % properness, maybe?}
%         % \mar{ME: Here I am using facts from {\tt proper.tex} that should be recalled above.}
%         so is the left-hand vertical arrow.
% 	Thus to complete the proof of the lemma, it suffices to show that
% 	upper left horizontal arrow is a closed immersion of finite
% 	presentation. 

% 	Since the lower horizontal arrow is representable by algebraic spaces
% 	and
% 	of finite presentation (being a base-change of the diagonal
%         of~$\cR_d$, which is representable by algebraic spaces
% 	and
% 	of finite presentation (again by Theorem~\ref{thm: C is a $p$-adic formal algebraic stack and related
%     properties}) % by Theorem~% \ref{thm: generalisation of
%         %   % main results of proper.tex}
%         % ),\mar{TG: add ref}
%         % \mar{ME: Add reference for this; the point is that it is a base-change of the diagonal of 
% 		% $\cR_d$, which we know is representable by
% 		% algebraic spaces and of finite presentation.}
% 	so is the upper horizontal arrow.  Furthermore, $\cC_{d,h}$ is
% 	a $p$-adic formal algebraic stack of finite presentation,
% 	and thus $\cR_d^{\Gamma_{\disc}}\times_{\cR_d} \cC_{d,h}$ is
% 	a $p$-adic formal algebraic stack of finite presentation.
% 	Now again using the fact that $\cC_{d,h}$ is a $p$-adic
% 	formal algebraic stack of finite presentation, we see
% 	that the same is true of $\cC_{d,h}^{\Gamma_{\disc}}.$ 
% 	% \mar{ME: Maybe note this earlier, as part of some general
% 	% 	discussion of fixed point stacks. TG: I added it as a
%         %         statement earlier now, but won't reference it yet in
%         %         case we move things around.}
% 	Thus the upper left horizontal arrow in the diagram
% 	is representable by algebraic stacks and 
% 	of finite presentation (being a morphism of $p$-adic
% 	formal algebraic stacks, both of which are of finite
%         presentation).\mar{TG: again, probably there should be a
%           citation of something in formal stacks at this point}
%         It remains to show that it is a closed immersion.

% 	The $2$-fibre product  $\cR_d^{\Gamma_{\disc}}\times_{\cR_d}
% 	\cC_{d,h}$ classifies projective rank $d$ $\varphi$-modules
% 	$\gM$ of $F$-height $\leq h$ endowed with a $\Gamma_{\disc}$-action
% 	on the underlying \'etale $\varphi$-module.  
% 	We then see that $\cC_{d,h}^{\Gamma_{\disc}}$ is the sub-stack
% 	of this $2$-fibre product classifying objects for which
% 	$\gM$ is actually stable under the $\Gamma_{\disc}$-action. To
%         see that this is a closed condition, we argue as in the proof
%         of~\cite[Prop.\ 5.3.8]{EGstacktheoreticimages}. By~\cite[Lem.\
%         6.1.9]{EGetalephimodules} there is a finite rank projective
%         \'etale $\varphi$-module~$P$ such that~$F:=\gM[1/u]\oplus P$ is a
%         free \'etale $\varphi$-module of finite rank. We can think of
%         the action of~$\gamma\in\Gamma_{\disc}$ as an element
%         of~$\End_{\OEA}(F)$ by letting it act by~$0$ on~$P$. Choosing
%         a basis for~$F$, 
% 	% If we work \'etale locally and so choose a basis for $\gM$,
% 	% then
%         the stability of $\gM$ under the $\Gamma_{\disc}$-action
% 	is equivalent to the vanishing of certain entries in the
% 	matrix describing this action of~$\gamma$; in
%         particular, % \mar{TG: need to rewrite this proof a bit as in
%           % proper.tex, using that we can write projective as a summand
%           % of free}
% 	it is a closed condition.\mar{TG: hopefully not a nonsense
%           argument at the end here - I think at worst this should be
%           close to what we need.} % \mar{ME: It remains to be seen that it is
% % 		of finite presentation, which is not obvious from this description, in which we ask infinitely many matrix entries to vanish.  Have to
% % 		show that source and target are locally of finite presentation,
% % 		so that we get that this map is automatically locally
% % 		of finite presentation, which is enough to conclude.  This
% % 	is one of many finiteness statements that has to be dealt with
% % throughout! ME: Now done, modulo various references.}
% \end{proof}

% \textbf{This is just some rough thoughts that I don't want to forget.}
% The hardest thing to prove seems to be that if as below we
% let~$\cX_{d,h,s}^a$ be the scheme-theoretic images of
% the~$\cW_{d,h,s}^a$ in~$\cR_d^{\Gamma_{\disc}}$, then these are actually substacks
% of~$\cX_d$. While we don't at this point know that~$\cX$ is limit
% preserving, we do know that~$\cX_{d,h,s}^a$ is, so I think that as in
% the proof of~\cite[Lem.\ 3.2.25]{EGstacktheoreticimages} it's
% reasonable to reduce to the case of showing that if we have~$\Spec
% A\to\cX_{d,h,s}^a$ with~$A$ Artinian, then the composite $\Spec
% A\to\cX_{d,h,s}^a\to\cR_d$ factors through~$\cX_d$. We can now go back
% to the definition of the scheme-theoretic image, again as in
% the proof of~\cite[Lem.\ 3.2.25]{EGstacktheoreticimages}, and maybe
% this will help\dots

% \begin{lem}
% 	\label{lem:surjectivity}\mar{TG: this isn't what is being
%           proved by the argument, need to think about this! Also we
%           don't currently know that $\cX$ is limit preserving at this
%           point in the argument. % this is wrong, but isn't
%   %         really what we're using below anyway - do something more
%   %         like \cite[Thm.\
%   % 5.3.15]{EGstacktheoreticimages}
% }
% 	The morphism~{\em (\ref{eqn:Ind W to X})} is surjective as a morphism
% 	of stacks on the Zariski site,
% 	i.e.\ any morphism to $\cX_d$ lifts Zariski locally on the source
% 	to a morphism to $\varinjlim_h \cWW_{d,h}.$
% \end{lem}
% \begin{proof} From the 2-Cartesian
%   diagram~(\ref{eqn:C to R square}), we see that it is enough to show
%   that every morphism to~$\cR_d^{\Gamma_{\disc}}$ lifts Zariski
%   locally on the source to
%   some~$\cR_d^{\Gamma_{\disc}}\times_{\cR_d}\cC_{d,h}$. Accordingly,
%   it is enough to show that every morphism to~$\cR_d$  lifts Zariski
%   locally on the source to
%   some~$\cC_{d,h}$. In the course of the proof of~\cite[Thm.\
%   5.4.16]{EGstacktheoreticimages}, it is shown that\mar{TG: what is
%     shown is that if we start with something Noetherian, then after
%     replacing it with a scheme-theoretic surjection, we can lift to
%     some $\cC_h$. The theorem itself proves that $\cR_d$ is the union
%     of the $\cR_{d,h}$, and maybe that's what we should be using here?}
  % , The analogous property holds for
  % the morphisms $\cC_{d,h}\to\cR_d$. Accordingly, we find ourselves in
  % the following situation: we have a projective \'etale
  % $\varphi$-module $M$ of rank~$d$ with $A$-coefficients and a
  % semi-linear action of~$\Gamma_{\disc}$, and a projective
  % $\varphi$-module $\gM$ of finite $T$-height with $\gM[1/u]=M$, but
  % we do not know that~$\gM$ is $\Gamma_{\disc}$-stable.
  % \mar{TG:
    % getting a bit stuck at this point! Are we supposed to use that the
    % morphism to $\cR^{\Gamma_{\disc}}$ actually comes from~$\cX_d$?
    % Then Lemma~\ref{lem: existence of phi gamma lattice} below shows
    % that there is a $(\varphi,\Gamma)$-stable lattice - but we don't
    % know that it's projective. Could we then try to do something like
    % the argument of the proof of~\cite[Prop 5.3.7]{EGstacktheoreticimages} to
    % massage this into something projective?!}
%\end{proof}

We now prove our first key structural result for $\cX_d$, in the case when $K$ is basic.

\begin{prop}
	\label{prop:X is an Ind-stack basic case}
If~$K$ is basic, then the canonical morphism
	$\varinjlim \cX^a_{d,h,s} \to \cX_d$ is an isomorphism.
	Thus $\cX_d$ is an Ind-algebraic stack, and may in fact
	be written as the inductive limit of
        algebraic stacks of finite presentation, with the transition maps
        being closed immersions.
	% \mar{ME: Extract appropriate finiteness statements from the
	% 	argument and add them.  Basically: $\cX_d$ will
	% 	be locally of finite presentation/limit preserving,
	% 	i.e.\ will satisfy~[1]
	% 	(just b/c it's an inductive limit of things
	% 	satisfying~[1]).
	% 	Maybe we can also get that it's quasi-separated, by
	% 	some general limit argument? 
	% 	Quasi-compactness will be less obvious; e.g.\ the
	% 	fact that $\cX_{d,\red}$ is quasi-compact is a
	% 	genuine theorem we prove below. (But add a remark
	%         here noting that we'll prove it below.)}
\end{prop}%\mar{TG: add a sentence to reduce to the basic case.}
\begin{proof}Note firstly that if $a'\ge a$, $h'\ge h$ and $s'\ge s$
  then the canonical morphism~$\cX^a_{d,h,s}\to\cX^{a'}_{d,h',s'}$ is a
  closed immersion by construction. By Lemma~\ref{lem: scheme
    theoretic images factor through X}, each~$\cX^a_{d,h,s}$ is a
  closed substack of~$\cX_d$, so it remains to show that any morphism
  $T\to\cX_d$ whose source is an affine scheme
  factors through some~$\cX^a_{d,h,s}$,
  or equivalently, that the closed immersion
  \numequation
  \label{eqn:closed immersion into test scheme}
  \cX^a_{d,h,s}\times_{\cX_d} T \to T
  \end{equation}
  is an isomorphism,
  for some choice of $h$ and $s$.

  Since~$\cX_d$ is
  limit preserving, by Lemma~\ref{lem:X is limit preserving}, we can
  reduce to the case where $T = \Spec A$ for a Noetherian
  $\cO/\varpi^a$-algebra~$A$.
  If $M$ denotes the \'etale $(\varphi,\Gamma)$-module corresponding
  to the morphism $\Spec A \to \cX_d$,
  then an application of \cite[Prop.~5.4.7]{EGstacktheoreticimages}
  shows that
  we may find a scheme-theoretically dominant morphism 
  $\Spec B \to \Spec A$ such that $M_B$ is free of rank $d$.
  If we show that the composite $\Spec B \to \Spec A \to \cX_d$
  factors through $\cX^a_{d,h,s}$ for some $h$ and $s$,
  then we see that the morphism $\Spec B \to \Spec A$
  factors through the closed subscheme
  $\cX^a_{d,h,s} \times_{\cX_d} \Spec A$ of $\Spec A.$ 
  Since $\Spec B \to \Spec A$ is scheme-theoretically dominant,
  this implies that~(\ref{eqn:closed immersion into test scheme})
  is indeed an isomorphism, as required.

  Since $M_B$ is free,
  we may choose a $\varphi$-invariant free lattice $\gM \subseteq M_B$,
  of height $\leq h$ for some sufficiently large value of $h$.
  Since the $\Gamma_{\disc}$-action on~$M$, and hence
  on~$M_B$, is continuous by assumption,
  Lemma~\ref{lem:testing continuity on M mod T}
  shows that $(\gamma^{p^s}-1)(\gM)\subseteq T\gM$ 
  for some sufficiently large value of $s$.
  Then $\gM$ gives rise to a $B$-valued point
  of $\cW^a_{d,h,s}$, whose image in $\cR^{\Gamma_{\disc}}_d$
  is equal to the \'etale $\varphi$-module $M_B$.
  Thus the morphism $B \to \cX_d$
  corresponding to $M_B$
  does indeed factor through $\cX^a_{d,h,s}$.
%
%
%
%
%
%  \mar{ME: Everything that follows will probably be removed.}
%  as in the proof of Lemma~\ref{lem: scheme theoretic images factor
%    through X} reduce to the case that~$T=\Spec A$ is an Artin local
%  $\cO/\varpi^a$-algebra.
%
%  Let~$M$ be the corresponding \'etale $(\varphi,\Gamma)$-module with
%  $A$-coefficients.  Since~$A((T))$ is local, $M$ is free, so we may
%  choose a lattice~$\gM\subset M$ which is a free rank $d$
%  $\varphi$-module of finite $T$-height (by choosing any basis for~$M$
%  and scaling it by an appropriate power of~$T$). By
%  Remark~\ref{rem:testing continuity on M mod u}, $\gM$ is a weak Wach
%  module of height at most~$h$ and level at most~$s$ for some
%  sufficiently large~$h,s$. The morphism $\Spec A\to\cX_d$ therefore
%  factors through~$\cW^a_{d,h,s}$ and thus through~$\cX^a_{d,h,s}$, as
%  required.
\end{proof}  

Finally, we drop our assumption that $K$ is basic.

\begin{prop}
	\label{prop:X is an Ind-stack}
Let~$K$ be an arbitrary finite extension of~$\Qp$. % Then the canonical morphism
	% $\varinjlim \cX^a_{d,h,s} \to \cX_d$ is an isomorphism.
	Then $\cX_d$ is an Ind-algebraic stack, and may in fact
	be written as the inductive limit of
        algebraic stacks of finite presentation over~$\Spec\cO$, with the transition maps
        being closed immersions. Furthermore the diagonal of~$\cX_d$
        is representable by algebraic spaces, affine, and of finite presentation.
	% \mar{ME: Extract appropriate finiteness statements from the
	% 	argument and add them.  Basically: $\cX_d$ will
	% 	be locally of finite presentation/limit preserving,
	% 	i.e.\ will satisfy~[1]
	% 	(just b/c it's an inductive limit of things
	% 	satisfying~[1]).
	% 	Maybe we can also get that it's quasi-separated, by
	% 	some general limit argument? 
	% 	Quasi-compactness will be less obvious; e.g.\ the
	% 	fact that $\cX_{d,\red}$ is quasi-compact is a
	% 	genuine theorem we prove below. (But add a remark
	%         here noting that we'll prove it below.)}
\end{prop}
\begin{proof}
  This is immediate from Propositions~\ref{prop: properties of X
    pulled back from R} and~\ref{prop:X is an Ind-stack basic case},
  together with Corollary~\ref{cor:deducing properties}.
  %together with Proposition~\ref{prop: properties of X pulled back from R}.
\end{proof}
%         Choose $s \geq 0,$ and consider the composite
%         \numequation
%         \label{eqn:W to R composite}
%         \cWW_{d,h,s} \to \cWW_{d,h} \to \cX_d \to \cR_d^{\Gamma_{\disc}}.
%         \end{equation}
%         This admits the alternative factorization
%         $$\cWW_{d,h,s} \to \cWW_{d,h} \to 	\cR_d^{\Gamma_{\disc}}\times_{\cR_d}\cC_{d,h} 
%         \to \cR_d^{\Gamma_{\disc}}.$$
%         Proposition~\ref{prop:inductive description of W} shows
%         that the composite of the first two arrows is a closed 
%         embedding of finite presentation,
%         while Theorem~\ref{thm: C is a $p$-adic formal algebraic stack and related
%     properties}  % Lemma~\ref{lem:C to R is proper}
%         % Theorem~\ref{thm: C is a $p$-adic formal algebraic stack and related
%     % properties}\mar{TG: might need to augment this reference, as it's
%     % not about fixed point stacks as stated.}
%         shows that the third arrow is representable by algebraic spaces, proper, 
%         and of finite presentation. 
%         Thus~(\ref{eqn:W to R composite}) is representable by
%         algebraic spaces, proper, and of
%         finite presentation.   

% 	Fix an integer $a \geq 1,$ and write $\cWW^a_{d,h,s} := \cWW_{d,h,s}
% 	\times_{\Spf \cO} \Spec \cO/\varpi^a$.   
% 	Proposition~\ref{prop:inductive description of W}
% 	shows that $\cWW_{d,h,s}$ is a $p$-adic formal
% 	algebraic stack of finite presentation over $\Spf \cO$,
% 	and so $\cWW^a_{d,h,s}$ is an algebraic stack, 
% 	and a closed substack of $\cWW_{d,h,s}$.
% 	Since $\cR_d^{\Gamma_{\disc}}$ is an Ind-algebraic stack, % inductive limit
% 	% of algebraic stacks under closed immersions, 
% 	the composite 
% 	\numequation
% 	\label{eqn:another composite}
% 	\cWW^a_{d,h,s} \hookrightarrow \cWW_{d,h,s}
% 	\buildrel \text{(\ref{eqn:W to R composite})} \over
% 	\longrightarrow \cR_d^{\Gamma_{\disc}},
%         \end{equation}
% 	which is representable by algebraic spaces, proper, and of finite presentation,
% 	factors through a closed algebraic substack $\cZ$ of
% 	$\cR_d^{\Gamma_{\disc}}.$   Let $\cX^a_{d,h,s}$
% 	denote the 
% 	% \mar{TG: is ``it'' $\cWW^a_{d,h,s}$? ME: Yes, clarified.}
%         scheme-theoretic image of
% 	$\cWW^a_{d,h,s}$ in $\cZ$; then
% 	$\cX^a_{d,h,s}$ % \mar{TG: $\cZ$ should be $\cX^a_{d,h,s}$? ME: Yes, fixed.}
%         is a closed algebraic substack of $\cR_d^{\Gamma_{\disc}}$.
% 	Since~(\ref{eqn:another composite}) factors through $\cX_d$,
% 	we see that $\cX^a_{d,h,s}$ is in fact a closed algebraic substack
% 	of $\cX_d$.\mar{ME: Cite a result from {\tt proper.tex} for
%           this. TG: this isn't completely clear to me! \cite[Thm.\
%           3.2.34]{EGstacktheoreticimages} might apply if we know that
%           $\cX_d\into\cR_d^{\Gamma_{\disc}}$ was of finite type - do
%           we know this? But that wouldn't tell us that
%           $\cX^a_{d,h,s}$ is closed in $\cX_d$. Hopefully this is all easy\dots}
	
% 	Noting that $\varinjlim_h \cWW_{d,h} = \varinjlim_{h,s,a} \cWW_{d,h,s}^a$
% 	(by Proposition~\ref{prop:inductive description of W}), 
% 	we deduce from Lemma~\ref{lem:surjectivity}
% 	that the morphism $\varinjlim_{h,s,a} \cX^a_{d,h,s} \to \cX_d$
% 	is an isomorphism.   This isomorphism realizes $\cX_d$ as
% 	an Ind-algebraic stack.\mar{ME: Add more details; closed immersions,
% 		finite presentation, etc.}
% \end{proof}
\section{Canonical actions and weak Wach modules}\label{subsec:
  canonical weak Wach}% \mar{TG: might need to weaken these introductory
% comments.}
We now explain an alternative perspective on some of the above
results, which gives more information about the moduli stacks of weak
Wach modules. % , at the expense of using more material from $p$-adic
% Hodge theory as an input (in particular, we use our results on almost
% \'etale descent, in the shape of Proposition~\ref{prop: equivalences
%   of categories to Ainf}).
The results of this section are only used
in Chapter~\ref{sec: the rank one case}, where we use them to
establish a concrete description of~$\cX_d$ in the case~$d=1$.
%\mar{TG: this section is likely not in final form.}
For each~$s\ge 0$, we write~$K_{\cyc,s}$
for the unique subfield of~$K_{\cyc}$ which is cyclic over~$K$ of
degree~$p^s$. The following lemma can be proved in exactly the same
way as Lemma~\ref{lem: Caruso Liu Galois action on Kisin} below.
\begin{lem}\label{lem: Caruso Liu Galois action on weak Wach}
Assume that~$K$ is basic.  For any fixed $a,h$,  there is a constant~$N(a,h)$ such that if $N\ge N(a,h)$, there is a
  positive integer~$s(a,h,N)$ with the property that for any finite type
  $\cO/\varpi^a$-algebra~$A$, any finite
  projective $\varphi$-module~$\gM$ over~$\A_{K,A}^+$  of $T$-height at most~$h$, and 
  any~$s\ge s(a,h,N)$, there is a unique continuous action of~ $G_{K_{\cyc,s}}$ on
  $\gMt:=\AAinf{A}\otimes_{\A_{K,A}^+}\gM$ which commutes with~$\varphi$
  and is semi-linear with respect to the natural action of~$G_{K_{\cyc,s}}$
  on~$\AAinf{A},$ with the additional property that for
  all~$g\in G_{K_{\cyc,s}}$ we have $(g-1)(\gM)\subset
  T^N\gMt$. % \mar{TG: should probably
    % explain how we are thinking of~$\gM$ as living inside $\AAinf{A}\otimes_{\gS_A}\gM$}
\end{lem}
% We can reinterpret this result in the language
% of~$(\varphi,\Gamma)$-modules as follows.
It is possible to
prove the following result purely in the world of $(\varphi,\Gamma)$-modules
(by following the proof of Lemma~\ref{lem: Caruso Liu Galois action on
  Kisin}, interpreting the existence of a semilinear action of~$\Gamma_{K_{\cyc,s}}$ in
terms of linear maps between twists of $\varphi$-modules, satisfying
certain compatibilities),
but we have found it more straightforward to argue with Proposition~\ref{prop: equivalences of
  categories to Ainf}.
% \mar{TG: following corollary should maybe say that in particular we
%   have a $(\varphi,\Gamma_{K_{\cyc,s}})$-module? TG: done, leaving
%   this comment so that you can check that I didn't do something stupid.}
 \begin{cor}\label{cor: Caruso Liu Galois action phi Gamma discrete version}
Assume that~$K$ is basic.  For any fixed $a,h$,  there is a constant~$N(a,h)$ such that if $N\ge N(a,h)$, there is a
  positive integer~$s(a,h,N)$ with the property that for any finite type
  $\cO/\varpi^a$-algebra~$A$, any finite
  projective $\varphi$-module~$\gM$ over~$\A_{K,A}^+$  of $T$-height at most~$h$, and 
  any~$s\ge s(a,h,N)$, there is a  unique semi-linear  action of~
  $\langle \gamma^{p^s}\rangle$ on
  $\gM$ which commutes with~$\varphi$,
  and  with the additional property that $(\gamma^{p^s}-1)(\gM)\subseteq
  T^N\gM$.

  In particular, this action gives $\gM[1/T]$ the structure of a projective \'etale $(\varphi,\Gamma_{K_{\cyc,s}})$-module.
  % \mar{TG: need to check that this implies the same for all
    % powers of~$\gamma^{p^s}$; probably follows from some quasi-linearity.}
  % \mar{TG: should probably
    % explain how we are thinking of~$\gM$ as living inside $\AAinf{A}\otimes_{\gS_A}\gM$}
\end{cor}%\mar{TG: proof will be added next}
\begin{proof}By Lemma~\ref{lem: Caruso Liu Galois action on weak
    Wach}, there is a a unique continuous semi-linear action of~ $G_{K_{\cyc,s}}$ on
  $\gMt:=\AAinf{A}\otimes_{\A_{K,A}^+}\gM$ which commutes with~$\varphi$,
   with the additional property that for
  all~$g\in G_{K_{\cyc,s}}$ we have $(g-1)(\gM)\subset
  T^N\gMt$. Write~$M:=\A_{K,A}\otimes_{\A_{K,A}^+}\gM$, $\widetilde{M}:=W(\C^\flat)_A\otimes_{\AAinf{A}}\gMt$. By Proposition~\ref{prop: equivalences
  of categories to Ainf}, the~$G_{K_{\cyc,s}}$-action on $\widetilde{M}$
endows~$M$ with the structure of a 
projective $(\varphi,\Gamma_{K_{\cyc,s}})$-module with $A$-coefficients,
and $\widetilde{M}$ with its~$G_{K_{\cyc,s}}$-action is recovered as
$$\widetilde{M}=W(\C^\flat)_A\otimes_{\A_{K,A}}M.$$ % \mar{TG: usual
  % $K_{\cyc,s}$ notational disaster.}

  
% We claim that $M=\A_{K,A}\otimes_{\A_{K,A}^+}\gM$.\mar{TG: uniqueness
%   of descent - should be similar to arguments we have made below} 

Since~$\Gamma_{K_{\cyc,s}}$ is topologically generated by~$\gamma^{p^s}$, and  since
$(g-1)(\gM)\subset
  T^N\gMt$ for $g \in G_{K_{\cyc,s}}$, we have $(\gamma^{p^s}-1)(\gM)\subset
  T^N\gMt$. We also have~$(\gamma^{p^s}-1)(\gM)\subseteq M$, and
  since~$M\cap\gMt=\gM$ (as is easily checked, by reducing to the case
  that~$\gM$ is free, and then to the case that it is free of rank
  one), % \mar{TG: need to justify this}
  we have $(\gamma^{p^s}-1)(\gM)\subset
  T^N\gM$, as required. For the uniqueness, note that if we
  have two such actions, then taking their difference gives a nonzero
  $\varphi$-linear morphism $(\gamma^{p^s})^*\gM\to T^N\gM$, which is
  impossible for~$N$ sufficiently large (see e.g.\ the first part of
  the proof of Lemma~\ref{lem: Frobenius amplification lemma over
    Ainf} below).

Finally, the claim that this gives $\gM[1/T]$ the structure of a
$(\varphi,\Gamma_{K_{\cyc,s}})$-module is immediate from Remark~\ref{rem:testing continuity on M mod u}.
  % \mar{TG: do we need to argue something for the converse,
    % i.e.\ from the ~$\gamma^{p^s}$ condition to the ~$G_{K_{\cyc,s}}$ version?}
\end{proof}

% For our applications 
% in Section~\ref{sec: the rank one case},
% we will need a version of the preceding corollary in the not-necessarily-basic
% case.  The difficulty in that case, as always, is that it is difficult
% to work meaningfully over $\A^+_{K,A}$, since this in general this
% ring is not preserved by~$\varphi$ or $\Gamma_K$.  So, once again,
% we will treat the general case by an appropriate reduction to the
% basic case.    

%Consider, then, a not-necessarily-basic~$K$. 
%Since $\A_K$ and $\A_{\Kbasic}$ are rings of integers
%in the discretely valued fields $\B_K$ and $\B_{\Kbasic}$ respectively,
%with the former being an unramified extension of the latter (by construction),
%we see that $\A_K$ is finite \'etale over $\A_{\Kbasic}$, 
%and in particular, the multiplication map
%$$\mu: \A_K \otimes_{\A_{\Kbasic}} \A_K \to \A_K$$
%splits as a homomorphism of rings: there is an idempotent
%$e \in \A_K \otimes_{\A_{\Kbasic}} \A_K$ such that $\mu$
%restricts to an isomorphism 
%$$e(\A_K \otimes_{\A_{\Kbasic}} \A_K) \iso \A_K.$$
%Since $\mu$ is $(\varphi,\Gamma_K)$-equivariant (here we endow the
%source with the
%diagonal actions of $\varphi$ and$\Gamma_K$, the latter action
%being well-defined, since $\Gamma_K = \Gamma_{\Kbasic}$),
%the idempotent $e$ is $\Gamma_K$-fixed. 

% Suppose now that $M$ is a rank $d$ projective
% \'etale $\varphi$-module
% over $\A_{K,A}$, for some finite type $\cO/\varpi^a$-algebra~$A$.
% \mar{ME: Can retrench to the case of a finite algebra if that helps.}
% %(We restrict to the finite case, since that is all we will need
% %in applications, and it makes the  manipulation of coefficient rings
% %easier, since we then have simply $\A_{K,A} = \A_K \otimes_{\Z/p^a} A$
% %and $\A_{\Kbasic,A} = \A_{\Kbasic} \otimes_{\Z/p^a} A$; no completions
% %are needed.)
% We let $M'$ denote $M$ regarded as an \'etale $\varphi$-module
% over $\A_{\Kbasic,A}$-module
% (now projective of rank $d[K:\Kbasic]$).
%We can then form 
%$$\A_{K,A} \otimes_{\A_{\Kbasic,A}} M'
%= \A_K\otimes_{\A_{\Kbasic}} M' = (\A_K \otimes_{\A_{\Kbasic}} \A_K)
%\otimes_{\A_K} M',$$
%and
%$e(\A_{K,A} \otimes_{\A_{\Kbasic,A}} M')$ is then a  copy of $M$
%contained in
%$\A_{K,A} \otimes_{\A_{\Kbasic,A}} M'$.

% Related to this %the preceding
% discussion, we make the following definition.

% \begin{defn}
%   Let~$A$ be a finite type $\cO/\varpi^a$-algebra for
%   some~$a\ge 1$, and let~$M$ be a projective \'etale $\varphi$-module
%   over~$\A_{K,A}$.  As above, let $M'$ denote the
% restriction of $M$ to~$\A_{\Kbasic,A}$.
% We say that~$M$ \emph{has
%     $T_{\Kbasic}$-height at most~$h$} if there is a
%   projective $\varphi$-module~$\gM$ over~$\A_{\Kbasic,A}^+$ of
% $T_{\Kbasic}$-height at
%   most~$h$ such that $M'=\gM[1/T_{\Kbasic}]$.
% \end{defn}

% Suppose, then, that $M$
% %now that $M$ is a rank $d$ projective \'etale $\varphi$ module
% %over $\A_{K,A}$, for some finite type $\cO/\varpi^a$-algebra, which
% is of $T_{\Kbasic}$-height at most~$h$.
% Let $N(a,h)$ and $s(a,h,N)$ (for some $N \geq N(a,h)$)
% be as in Corollary~\ref{cor:
% Caruso Liu Galois action phi Gamma discrete version}
% (applied to $\Kbasic$),
% so that $M'$ is endowed with a canonical
% $\Gamma_{\Kbasic_{\cyc,s}} = \Gamma_{K_{\cyc,s}}$-action.  

% \begin{lemma}
% \label{lem:non-basic canonical action semi-linearity}
% In the context of the preceding discussion,
% for some $s' := s'(a,h,N) \geq s(a,h,n)$,
% the canonical $\Gamma_{K_{\cyc,s'}}$-action
% on $M' =  M$ is semi-linear with respect to
% the $\Gamma_{K_{\cyc,s'}}$-action on $\A_{K,A}$
% and the $\A_{K,A}$-module structure on~$M$.
% \end{lemma}
% \begin{proof}
% \mar{ME: Proof needed.  Seems trickier than I'd thought,
% b/c $\A_{K,A}$ is not even a $\varphi$-module (let alone
% of finite height), so we can't just use off-the-shelf uniqueness
% of canonical actions.}
% \end{proof}

%\mar{ME: Add remark about compatibility with $\varpi$-adic graded pieces.}

We continue to assume that $K$ is basic.
%We now return to the case that $K$ is basic.
For any~$s~\geq~1$,
we write $\Gamma_{s,\disc}$ to denote the subgroup
of $\Gamma_{\disc}$ generated by $\gamma^{p^s}$.
We write $\cR_d^{\Gamma_{s,\disc}}$ in obvious analogy to the notation
$\cR_d^{\Gamma_{\disc}}$ introduced above.
The result of 
Corollary~\ref{cor: Caruso Liu Galois action phi Gamma discrete version}
may be interpreted as constructing a morphism
\numequation
\label{eqn:CL morphism cyclo case}
\cC^a_{d,h} \to \cR_d^{\Gamma_{s,\disc}}
\end{equation}
for sufficiently large values of $s$ (depending on $a$ and $h$),
lying over the morphism $\cC^a_{d,h} \to \cR_d.$ Of course there is also a morphism
$\cR_d^{\Gamma_{\disc}} \to \cR_d^{\Gamma_{s,\disc}}$
given by restricting the action of $\Gamma_{\disc}$ to $\Gamma_{s,\disc}$.

Since the morphism $\cR_d^{\Gamma_{s,\disc}} \to \cR_d$ given by
forgetting the $\Gamma_{s,\disc}$-action is representable by algebraic 
spaces and separated (indeed, even affine --- it is a base-change
of the morphism $\cR_d \to \cR_d \times_{\cO} \cR_d$ giving
the graph of the action of~$\gamma^{p^s}$, and this latter
morphism is representable by algebraic spaces and affine, since
the diagonal morphism of $\cR_d$ is so),
the diagonal morphism 
$$\cR_d^{\Gamma_{s,\disc}} \to
\cR_d^{\Gamma_{s,\disc}} \times_{\cR_d} \cR_d^{\Gamma_{s,\disc}}
$$
is a closed immersion.
We may thus define a closed substack $\cZ_{d,h,s}^a$ of 
$\cR_d^{\Gamma_\disc}\times_{\cR_d} \cC^a_{d,h}$
via the following $2$-Cartesian diagram:
$$\xymatrix{
\cZ^a_{d,h,s} \ar[r] \ar[d] &  \cR_d^{\Gamma_\disc}\times_{\cR_d} \cC^a_{d,h} \ar[d] \\
\cR_d^{\Gamma_{s,\disc}} \ar[r] & \cR_d^{\Gamma_{s,\disc}} \times_{\cR_d}
\cR_d^{\Gamma_{s,\disc}} }
$$
in which the lower horizontal arrow is the diagonal, and the right-hand vertical
arrow is the product of the restriction morphism and the morphism~\eqref{eqn:CL 
morphism cyclo case}.
In less formal language, the stack $\cZ^a_{d,h,s}$ parameterizes
projective rank~$d$ $\varphi$-modules $\gM$  over $\A_{K,A}^+$ of $T$-height at most~$h$,
for which $\gM[1/T]$ is endowed with a $\Gamma_{\disc}$-action extending
the canonical action of $\Gamma_{s,\disc}$ given by 
Corollary~\ref{cor: Caruso Liu Galois action phi Gamma discrete version}.


\begin{prop}\label{prop:ind structure on X via canonical
    actions}Assume that~$K$ is basic.
Then each $\cZ^a_{d,h,s}$ is contained {\em (}as a closed algebraic substack{\em )} in
$\cX_d\times_{\cR_d} \cC_{d,h},$
and the natural morphism
$\varinjlim_{a,s} \cZ_{d,h,s}^a \to \cX_d\times_{\cR_d} \cC_{d,h}$
is an isomorphism.
\end{prop}
\begin{proof}That $\cZ^a_{d,h,s}$ is a substack of
$\cX_d\times_{\cR_d} \cC_{d,h}$ is immediate from
Remark~\ref{rem:testing continuity on M mod u}; indeed, if we are given $\gM$ over $\A_{K,A}^+$ of $T$-height at most~$h$,
for which $\gM[1/T]$ is endowed with a $\Gamma_{\disc}$-action extending
the canonical action of $\Gamma_{s,\disc}$ given by 
Corollary~\ref{cor: Caruso Liu Galois action phi Gamma discrete
  version}, then we have~$(\gamma^{p^s}-1)(\gM)\subseteq
T^N\gM\subseteq T\gM$. Since $\cZ_{d,h,s}^a$ is a closed substack of 
$\cR_d^{\Gamma_\disc}\times_{\cR_d} \cC^a_{d,h}$, it is a closed
substack of~$\cX_d\times_{\cR_d} \cC_{d,h}$.

To see that the natural morphism
$\varinjlim_{a,s} \cZ_{d,h,s}^a \to \cX_d\times_{\cR_d} \cC_{d,h}$ is an isomorphism,
we need to show that it is surjective, and so we need to show that for any
finite type $\cO/\varpi^a$-algebra~$A$, and any projective $\varphi$-module $\gM$ over $\A_{K,A}^+$ of $T$-height at most~$h$,
for which $\gM[1/T]$ is endowed with a continuous
$\Gamma_{\disc}$-action, there is some~$s\ge 1$ such that 
the restriction of this action to~$\gamma_{s,\disc}$ agrees with the canonical action of $\Gamma_{s,\disc}$ given by 
Corollary~\ref{cor: Caruso Liu Galois action phi Gamma discrete
  version}.

By Remark~\ref{rem:testing continuity on M mod u}, there is
some~$s'\ge 1$ such that $(\gamma^{p^{s'}}-1)(\gM)\subseteq
T\gM$. Let~$N=N(a,h)$ be as in the statement of Lemma~\ref{lem: Caruso
  Liu Galois action on weak Wach}. By the equivalence of conditions
(3) and~(4) of Lemma~\ref{lem:testing continuity on M mod T}, there is
some~$s\ge s'$ such that $(\gamma^{p^s}-1)(\gM)\subseteq
T^N\gM$, as required.
%\mar{ME: Needs proof TG: how about this? ME: Looks good!}
\end{proof}

For arbitrary~$K$ (not necessarily basic) we can deduce the following description
of~$\cX_d^a$ as an Ind-algebraic stack, in the style of Lemma~\ref{lem:
  explicit Ind description for R pulled back from K basic}.

\begin{lem}
  \label{lem: explicit Ind description for X pulled back from K basic}
Fix $a \geq 1$.
For each $h$, and each sufficiently large {\em (}depending on $a$ and $h${\em )}
value of~$s$,
%  $\cX_{K,d}^a=\varinjlim_{h,s}\cX_{K,d,\Kbasic,h,s}^a$
let $\cX_{K,d,\Kbasic,h,s}^a$ denote the scheme-theoretic image of the
  base-changed
  morphism
  \[\cZ_{d[K:\Kbasic],h,s}^a\times_{\cX_{\Kbasic,d[K:\Kbasic]}}\cX^a_{K,d}\to\cX^a_{K,d},\]
so that
$\cX_{K,d,\Kbasic,h,s}^a$ is a closed algebraic substack of $\cX_{K,d}$.
Then the natural morphism
  $\varinjlim_{h,s}\cX_{K,d,\Kbasic,h,s}^a \to \cX_{K,d}^a$
is an isomorphism.
\end{lem}
\begin{proof}
%  Since the morphism $\cX_{K,d}\to\cX_{\Kbasic,d[K:\Kbasic]}$ is
%  representable by algebraic spaces by Proposition~\ref{prop: properties of X pulled back from R},
This is proved in the
  same way as Lemma~\ref{lem: explicit Ind description for R pulled
    back from K basic}, bearing in mind Proposition~\ref{prop:ind structure on X
    via canonical actions}. 
%\mar{TG: hopefully this is OK; I didn't make much of an effort to check it though.}
\end{proof}




%We now explain how the preceding results can be used to give another proof of the 
%properties of~$\cX$ established in
%Proposition~\ref{prop:X is an Ind-stack}. For the moment,
%we continue to assume that $K$ is basic. 
%
%\begin{df} For any $h \geq 0,$ define $\cR_{d,h}$ 
%to be the scheme-theoretic image of the morphism $\cC_{d,h} \to \cR_d.$
%\end{df}
%
%Each $\cR_{d,h}$ is a closed substack of $\cR_d$,
%and is furthermore an algebraic stack of finite presentation over $\cO/\varpi^a$,
%and the proof of \cite[Thm.~5.4.20]{EGstacktheoreticimages},
%which shows that $\cR_d$ is Ind-algebraic,
%in fact shows that there is an isomorphism $\varinjlim_h \cR_{d,h} \iso \cR_d.$
%
%We let $\cR_{d,h}^{\Gamma_{\disc}}$ denote the preimage of $\cR_{d,h}$ under the forgetful 
%morphism $\cR_{d}^{\Gamma_{\disc}} \to \cR^{d}$ (given by
%forgetting the $\Gamma_{\disc}$-action). 
%This latter morphism is representable by algebraic spaces, affine,
%and of finite presentation
%(since by definition it is a base-change of the graph morphism
%$\Gamma_{\gamma}: \cR_d \to \cR_d \times_{\cO/\varpi^a} \cR_d$,
%the graph of the automorphism of $\cR_d$ given by the action of~$\gamma$, 
%which inherits these properties from the diagonal of $\cR_d$),
%and thus $\cR_{d,h}^{\Gamma_{\disc}}$ is a closed algebraic substack
%of~$\cR_d^{\Gamma_{\disc}}$ which is of finite presentation over $\cR_{d,h}$ (and
%thus over $\cO/\varpi^a$), the forgetful morphism
%$\cR_{d,h}^{\Gamma_{\disc}} \to \cR_{d,h}$ is representable by algebraic spaces
%and affine, and there is an isomorphism
%$\varinjlim_h \cR_{d,h}^{\Gamma_{\disc}} \iso \cR_d^{\Gamma_{\disc}}.$
%
%We can now give another description of the Ind-algebraic stack
%structure of $\cX_d$.  Namely, since $\varinjlim_h \cR_{d,h}^{\Gamma_{\disc}}
%\iso \cR_d^{\Gamma_{\disc}}$, with the transition morphisms being closed immersions,
%we have an induced isomorphism
%$$\varinjlim_h (\cR_{d,h}^{\Gamma_{\disc}} \times_{\cR_d^{\Gamma_{\disc}}} \cX_d)
%\iso \cX_d,$$
%again with the transition morphisms being closed immersions.
%The following result then shows that this isomorphism realizes 
%$\cX_d$ as an Ind-algebraic stack with the properties stated in
%Proposition~\ref{prop:X is an Ind-stack}.
%
%\begin{prop}
%\label{prop:algebraicity result}
%For each $h \geq 1$, the fibre product
%$\cR_{d,h}^{\Gamma_{\disc}} \times_{\cR_d^{\Gamma_{\disc}}} \cX_d$
%is an algebraic stack of finite presentation over $\cO/\varpi^a$.
%\end{prop}
%
%The proof of this proposition will use the canonical action of
%Corollary~\ref{cor: Caruso Liu Galois action phi Gamma discrete version}.
%To explain this,
%%Fix $a \geq 1$ and $h \geq 0$, as well as
%fix some $N \geq N(a,h)$ and $s \geq s(a,h,N)$,
%in the notation of that corollary.
%%Corollary~\ref{cor: Caruso Liu Galois action phi Gamma discrete version}.
%Let $\cR_d^{\Gamma_{s,\disc}}$ and $\cR_{d,h}^{\Gamma_{s,\disc}}$
%have the evident meaning (the analogue of $\cR_d^{\Gamma_{\disc}}$
%and $\cR_{d,h}^{\Gamma_{\disc}}$
%that one obtains by replacing $\Gamma_{\disc}$ with its subgroup $\Gamma_{s,\disc}$
%generated by $\gamma^{p^s}$).
%We may interpret the canonical action of
%Corollary~\ref{cor: Caruso Liu Galois action phi Gamma discrete version}
%as giving
%a morphism of stacks $\cC_{d,h} \to \cR_{d,h}^{\Gamma_{s,\disc}}$.
%
%%\begin{prop}
%%\label{prop:canonical action morphism in the cyclotomic setting}
%%There is a morphism of algebraic stacks
%%$\cR_{d,h} \to \cR_{d,h}^{\Gamma_{\disc_s}}$,
%%uniquely determined by the fact that its composite with
%%the canonical morphism $\cC_{d,h} \to \cR_{d,h}$ coincides
%%with the construction of
%%Corollary~{\em \ref{cor: Caruso Liu Galois action phi Gamma discrete version}}.
%%\end{prop}
%%\begin{proof}
%%The construction of
%%Corollary~{\em \ref{cor: Caruso Liu Galois action phi Gamma discrete version}}
%%yields a morphism
%%\numequation
%%\label{eqn:canonical action on C}
%%\cC_{d,h} \to \cR_{d,h}^{\Gamma_{s,\disc}},
%%\end{equation}
%%which we wish to factor through the scheme-theoretically dominant morphism
%%$\cC_{d,h} \to \cR_{d,h}$.  
%%
%%Since the morphism
%%$\cR_{d,h}^{\Gamma_{s,\disc}} \to \cR_{d,h}$
%%given by forgetting
%%the $\Gamma_{s,\disc}$-action is representable by algebraic spaces and separated (since
%%it is even affine), the diagonal morphism
%%$$\cR_{d,h}^{\Gamma_{s,\disc}} \to 
%%\cR_{d,h}^{\Gamma_{s,\disc}} \times_{\cR_{d,h}} \cR_{d,h}^{\Gamma_{s,\disc}}$$
%%is a closed immersion.  The scheme-theoretic dominance of the morphism
%%$\cC_{d,h} \to \cR_{d,h}$ then implies that the sought-after factorization
%%of~\eqref{eqn:canonical action on C} is unique, if it exists.
%%
%%%The morphism   
%%%$$\cC_{d,h}^{\Gamma_{s,\disc}}\times_{\cR_{d,h}}\cR_{d,h}^{\Gamma_{s,\disc}}
%%%\to \cC_{d,h}$$
%%%is again representable by algebraic spaces and separated (being the base-change of
%%%such a morphism), 
%%%and so the graph of~\eqref{eqn:canonical action on C} is a closed
%%%substack of 
%%%$\cC_{d,h}^{\Gamma_{s,\disc}}\times_{\cR_{d,h}}\cR_{d,h}^{\Gamma_{s,\disc}}$
%%%which projects isomorphically onto the first factor.
%%%Consider now the scheme-theoretic image of this graph in
%%%$$\cR_{d,h}^{\Gamma_{s,\disc}}\times_{\cR_{d,h}}\cR_{d,h}^{\Gamma_{s,\disc}}
%%%\iso \cR_{d,h}^{\Gamma_{s,\disc}};$$
%%%we wish to show that this projects isomorphically to $\cR_{d,h}$ under
%%%the projection (i.e.\ the forgetful morphism).
%%
%%Consider the scheme-theoretic image $\cZ$ of~\eqref{eqn:canonical action on C};
%%this is a closed subscheme of 
%%$\cR_{d,h}^{\Gamma_{s,\disc}}.$   If we show that $\cZ$ maps isomorphically
%%to $\cR_{d,h}$ under the forgetful morphism, then the inverse of this isomorphism
%%provides the desired morphism.
%%
%%
%%
%%
%%
%%\end{proof}
%\mar{needed}
%\begin{proof}[Proof of Proposition~\ref{prop:algebraicity result}]
%%The closed substack 
%%$\cR_{d,h} \times_{\cR^{\Gamma_{\disc}}} \cX_d$
%%of $\cX_d$ admits the alternative description 
%%$$\cR_{d,h}^{\Gamma_{\disc}}
%%\times_{p_1,\cR^{\Gamma_{s,\disc}},p_2} \cR_{d,h}^{\Gamma_{\disc}},$$
%%where $p_1$ is the restriction morphism (restrict the $\Gamma_{\disc}$ action
%%to a $\Gamma_{s,\disc}$-action), while $p_2$ is the morphism
%%given by passing to the canonical $\Gamma_{s,\disc}$-action
%%on the underlying \'etale $\varphi$-module (given by the morphism
%%of Proposition~\ref{prop:canonical action morphism in the cyclotomic setting}).
%%\mar{ME: Probably should elaborate.}
%\end{proof}
%
%\mar{TG: then we deduce Proposition~\ref{prop:X is an Ind-stack}
%  again, with exactly the same one sentence argument to go from the
%  basic to the general case.}

\section{The connection with Galois representations}
\label{subsec:Galois reps}
In Section~\ref{subsubsec:Galois reps}  we  recalled
the  relationship between  \'etale $(\varphi,\Gamma)$-modules
without coefficients (which is to say, with coefficients in~$\Z_p$)
and $p$-adic representations of~$G_K$,
and the similar relationship between \'etale $\varphi$-modules
over $\cO_{\mathcal E}$ and $p$-adic representations of~$G_{K_{\infty}}$.
In this section, we revisit those topics in the context of
\'etale $(\varphi,\Gamma)$-modules (resp.\ \'etale $\varphi$-modules)
and $G_K$-representations (resp.\ $G_{K_{\infty}}$-modules) with coefficients.


\subsection{Galois representations with coefficients}
\label{subsec:Galois reps with coeffs}
 For a general $p$-adically complete
$\Zp$-algebra~$A$, a projective \'etale $(\varphi,\Gamma)$-module with
$A$-coefficients need not correspond to a family
of~$G_K$-representations. It will be useful, though, to have a version of
such a correspondence in the case that~$A$ is complete local Noetherian
with finite residue field. In this case, following ~\cite{MR1805474}
we let~$\widehat{\A}_{K,A}$ denote the $\m_A$-adic completion
of~${\A}_{K,A}$, and we define a \emph{formal \'etale
  $(\varphi,\Gamma)$-module} with $A$-coefficients \index{formal \'etale
  $(\varphi,\Gamma)$-module}
% \mar{ME: Added ``over $A$'' here; 
%is that correct terminology? TG: no, corrected ME: Also, explain that $\cX(A)$ classifies
%these. ME: Sorted}
to be an \'etale $(\varphi,\Gamma)$-module over
$\widehat{\A}_{K,A}$ in the obvious sense.

\begin{remark}
	If $A$ is a complete local Noetherian $\cO$-algebra,
	with finite residue field,
	then the groupoid of formal \'etale $(\varphi,\Gamma)$-modules
	is equivalent to the groupoid $\cX_d(A)$ (which we remind the reader
	refers to the groupoid of morphisms $\Spf A \to \cX_d$; here
	$\Spf A$ is taken with respect to the $\mathfrak m_A$-adic
	topology on $A$).  Indeed, by definition, the 
	latter groupoid may be identified with the $2$-limit
	$\varprojlim_i \cX_d(A/\mathfrak m_A^i)$, which is easily
	seen to be equivalent 
	to the groupoid of formal \'etale $(\varphi,\Gamma)$-modules,
	via an application of~\cite[Prop.\
        0.7.2.10(ii)]{MR3075000}. % \mar{TG: should be inverse limit? ME:
% I agree, and have made the change, and replaced ``$2$-colimit''
% by  ``$2$-limit'', but just want to double-check this terminology/concept,
% since I've not consciously used it before. TG: google seems to confirm
% this is fine, commenting out.} 
\end{remark}



We let $\AKnrhatA$ denote 
  ~$\AKnrhat\cotimes_{\Zp} A$, where the completed
  tensor product is with respect to the usual topology
  on~$\AKnrhat$, and the  $\m_A$-adic topology
  on~$A$.
  %\mar{ME: Should we also introduce the notion of a {\em formal 
%		  \'etale $\varphi$-module}? TG: it doesn't look like
%                we ever use it?}
%  We define
%  $\widehat{\cO}_{\widehat{\cE^{\nr}},A}:=\cO_{\widehat{\cE^{\nr}}}\cotimes_{\Zp}A$
%  in the same way.
%   \mar{ME: The notation
% $\AKnrhatA$ and
%   $\widehat{\cO}_{\widehat{\cE^{\nr}},A}$ seems to inconsistent
% with the notation introduced in \S 2.}
%
%
% \mar{ME: Not dealing with these last two for now, because perhaps
% 	we don't really need them? TG: deal with them here, I guess
%         just recalling Dee's definitions is what we should do.}
% \[\cO_{\widehat{\cE^{\nr}},A}:=\AKnrhat\cotimes_{\Zp}A, \]
% \[\cO_{\widehat{\cE^{\nr}_\infty},A}:=\cO_{\widehat{\cE^{\nr}_\infty}}\cotimes_{\Zp}A, \]
%where in each case the completed tensor product is with respect to
%the $p$-adic topology on~$A$, and the adic topologies % \mar{TG: is that
  % the right thing to say?}
%on the
%rings~$\A_K$, $\A_K^+$, $\Ainf$, $W(\C^\flat)$, $\gS$,
%$\cO_{\cE}$, $\AKnrhat$, $\cO_{\widehat{\cE^{\nr}_\infty}}$.  It is shown in~\cite{MR1805474}
The functors~$T_A$, $\mathbf{D}_A$ defined by  \[\mathbb{D}_A(T)=(\AKnrhatA 
  \otimes_{A}T)^{G_{\Kcyc}},\]  \[T_A(M)=(  \AKnrhatA \otimes_{\widehat{\A}_{K,A}}M)^{\varphi=1}\]% \mar{TG:
% obviously this needs to be filled out and made correct.}
then give equivalences of categories between the category of finite projective formal \'etale
$(\varphi,\Gamma)$-modules with $A$-coefficients, and the category of
finite free $A$-modules with a continuous action of~$G_K$.
% Again, there is a~$G_K$, $\varphi$-equivariant
% isomorphism
% \[W(\C^\flat)_A \otimes_{\A_{K,A}}M\isoto W(\C^\flat)_A
%   \otimes_{A}T_A(M). \] 
(In fact the results of~\cite{MR1805474} do
not require the $(\varphi,\Gamma)$-modules to be projective, but we
will for the most part only consider projective modules in this book.)

Continuing to assume that~$A$ is complete local Noetherian with finite
residue field, we let~$\hOEA$ denote the $\m_A$-adic completion
of~$\OEA$,
  and we define
  $\widehat{\cO}_{\widehat{\cE^{\nr}},A}:=\cO_{\widehat{\cE^{\nr}}}\cotimes_{\Zp}A$,
  where the completed
  tensor product is with respect to the usual topology
  on~$\cO_{\widehat{\cE^{\nr}}}$, and the  $\m_A$-adic topology
  on~$A$.
Then the analogous statements to those of the previous
paragraph, relating representations
of~$G_{K_\infty}$ on finite free $A$-modules and \'etale
$\varphi$-modules over~$\hOEA$\footnote{One might refer to
an \'etale $\varphi$-module over~$\hOEA$ as a
``formal \'etale $\varphi$-module over $A$'', but since there
are several different species of \'etale $\varphi$-modules
under consideration throughout the book, we avoid using
this potentially ambiguous terminology.} via the functors
\[\mathbb{D}_{\infty,A}(T)=(\widehat{\cO}_{\widehat{\cE^{\nr}},A} 
  \otimes_{A}T)^{G_{K_\infty}},\]  \[T_{\infty,A}(M)=(  \widehat{\cO}_{\widehat{\cE^{\nr}},A} \otimes_{\hOEA}M)^{\varphi=1},\] can be proved in exactly the same way as in~\cite{MR1805474}
(by passage to the limit over~$A/\m_A^n$).

We will also occasionally
apply these statements in the case~$A=\Fpbar$.
% or~$A=\Zpbar$ (in which case we do not need to make an $\m_A$-adic
% completion, as our coefficient rings are already $p$-adically complete
% by definition).
Their validity in this case 
% \mar{ME: Note that $\Zpbar$  is {\em  not} $p$-adically  complete!
% Which might also explain why the limit-preserving  argument that
% follows only made sense to me for $\Zpbar/p^a$ for some fixed~$a$,
% and not for $\Zpbar$ itself. TG: removed it here and elsewhere. I have
% left it alone in places where we are purely talking about Galois
% representations, e.g.\ in stating the existence of crystalline lifts,
% but obviously we could change it there too if you like.}
follows from their validity in the case when~$A$ is
a finite extension of~$\Fp$, % or the ring of integers in a
% finite extension of~$\Qp$, 
and the fact that both Galois representations
and  projective \'etale $(\varphi,\Gamma)$-modules over~$\Fpbar$ %and~$\Zpbar$
arise as base changes from such finite contexts; for
Galois representations, this  %is a well-known consequence of % the Baire
% category theorem and
follows from the compactness of~$G_K$ and~$G_{K_{\infty}}$,
and for
\'etale $(\varphi,\Gamma)$-modules, it is an immediate consequence of
Lemma~\ref{lem:X is limit preserving}.
%
%
%\mar{TG:  is this the correct generality for the
%   connection to Galois representations? Mark proves it for finite
%   cardinality~$A$. Probably it should also be
%   moved to the section where $(\varphi,\Gamma)$-modules are actually
%   introduced? ME: Should also perhaps beef this up to $\Zbar_p$-algebras,
% so e.g.\ the case of $\Fbar_p$ is also included. TG: is the idea to 
% take a direct limit to cover these?}\mar{TG: add a forward citation to
% finite presentation on the $(\varphi,\Gamma)$-side, and discuss
% together with Baire category on the Galois side.}


% (see~\cite{MR1106901} % \mar{TG: some Fontaine
% % ref}
% for the case of~$\Zp$-coefficients, and (the proof of)~\cite[Lem.\
% 1.2.7]{KisinModularity} for the case that~$A$ has finite cardinality;
% the general case follows from this by writing
% $A=\varprojlim_nA/p^nA$). Let $\Rep_{G_K}(A)_d$ % \mar{TG: terrible notation,
% % and may well not be the category we actually want. ME: Slightly tweaked the
% % notation.}
% denote the category of
% continuous representations of~$G_K$ on free $A$-modules of
% rank~$d$. Then there is a $\Zp$-algebra $\AKnrhat$, with commuting
% actions of~$G_{K_\infty}$ and~$\varphi$, and an equivalence of categories between
% $\Rep_{G_K}(A)_d$ and the category of free \'etale
% $(\varphi,\Gamma)$-modules with $A$-coefficients, given by the
% functors $\mathbb{D}$, $T$ which are defined by \[\mathbb{D}(V)=(
%   \otimes_{\Zp}V)^{G_{K_\infty}},\] \[T(M)=(\AKnrhat \otimes_{\OEA}M)^{\varphi=1}.\]% \mar{TG
%:
% obviously this needs to be filled out and made correct.}

% \mar{TG: need to add stuff in the Kisin setting, too}
       



We can use the equivalence between Galois representations
and $(\varphi,\Gamma)$-modules as a tool to deduce facts about
the finite type points of $\cX_{d}$.  % \mar{ME: Is this the only place
	% we write $\cX_{K,d}$ rather than just $\cX_d$?}
      More precisely,
if~$\F'/\F$ is a finite extension, then the groupoid of points 
% \mar{ME: Shouldn't this be $\F'/\F$, since we have $\cO$ as our
% base coefficients, not $\Z_p$? TG: yes, fixed, commenting out.}
$x\in \cX_{d}(\F')$ is canonically equivalent to the groupoid of
 Galois representations $\rhobar: G_K \to \GL_d(\F')$. % , whose
% structure is fairly well understood (as we recall below).
For 
this reason, we will often denote such a point $x$ simply
by the corresponding Galois representation~$\rhobar$.
% \mar{ME: Probably should
% 	rewrite this and the above as in the introduction: describe
% 	it an an actual equivalence of groupoids.}


Suppose now that~$\F'/\F$ is a finite extension, and
that~$x:\Spec\F'\to\cX_d$ is a finite type point, % \mar{TG:
% maybe should have $\F'$, $\cO'$ here, and then annoyingly~$\cO''$ below?}
with  corresponding Galois representation 
$\rhobar:G_K\to\GL_d(\F')$. Let $\widehat{(\cX_d)}_x$ be the category
of Definition~\ref{adefn: deformation category} (with~$\cF$ there
being $\cX_d$). Let ~$\cO'=\cO\otimes_{W(\F)}W(\F')$ be the ring of integers in
the finite extension~$E'=W(\F')E$. % with residue field~$\F'$. 
% \mar{ME: Somewhere we should probably note that the whole theory is
% 	compatible with change of coefficients, e.g.\ from $\cO$
% 	to $\cO'$. TG: e.g.\ that
%         $\cX_{d,\cO'}=\cX_d\otimes_{\cO}\cO'$, which is
%         tautological. TG: here and in the Hodge theory section,
%         probably $\cO'$ should just be~$\cO W(\F')$?}
\begin{prop}
	\label{prop:versal rings}
	%\mar{ME: Made this statement into a prop.\ for later reference;
	%	maybe proof can be tidied up/elaborated on if necessary,
	%	now that it's been highlighted a bit more than it was.}
There is a morphism
$\Spf R_{\rhobar}^{\square,\cO'}
\to\cX_d$
which is versal at the point $x$ corresponding to~$\rhobar$,
an isomorphism 
$\Spf R_{\rhobar}^{\square,\cO'}
\times_{\cX_d}
\Spf R_{\rhobar}^{\square,\cO'}
\iso
\widehat{\GL}_{d,R_{\rhobar}^{\square,\cO'},Z(\rhobar)},$
where
$\widehat{\GL}_{d,R_{\rhobar}^{\square,\cO'},Z(\rhobar)}$
denotes the completion of $(\GL_{d})_{R_{\rhobar}^{\square,\cO'}}$
along the closed subgroup of $(\GL_{d})_{\F'}$
given by the centraliser of~$\rhobar$,
and an isomorphism 
$\Spf R_{\rhobar}^{\square,\cO'}
\times_{\widehat{(\cX_d)}_x}
\Spf R_{\rhobar}^{\square,\cO'}
\iso
\widehat{\GL}_{d,R_{\rhobar}^{\square,\cO',1}},$
where
%\mar{ME: There is duplication of notation for the two different completions
%of $\GL_{d,R_{\rhobar}^{\square,\cO'}},$, but I won't worry about it here.
%Also, this is notation for a base-change of schemes that follows a different
%convention to our $(\cX_d)_{\Fbar_p}$ notation.  Perhaps we should settle
%on a consistent notation for base-changes? TG: hopefully better now? ME: Looks good,
%commenting out}
$\widehat{\GL}_{d,R_{\rhobar}^{\square,\cO'},1}$
denotes the completion of $(\GL_{d})_{R_{\rhobar}^{\square,\cO'}}$
along the identity of $(\GL_{d})_{\F'}$.
% \mar{ME: I'm not sure  
% $\widehat{(\cX_d)}_x$ has been defined at all, but assuming it means
% the completion of $\cX_d$ at~$x$,  it will  surely be a substack of $\cX_d$,
% so  the  fibre product over it and over $\cX_d$ should be the same,
% in contrast  to  what is claimed here. ME: OK, I found a definition
% in App.~A (in the discussion of versal rings),
% which is not at all what I'd imagined it would be,
% and with that definition (which should definitely be recalled, or at
% least referred to, here!) I think the statement is OK. TG: agreed, but
% leaving this comment here for now because we might want to clarify
% having two statements, and I presume we might actually apply the one
% we were sceptical of (just because it got introduced later, and
% presumably it was introduced for a reason\dots) TG: note that I
% recalled the definition immediately above, and we do use both fibre
% products, so I suggest that we simply remove this comment, unless you
% want to add more explanation.}
\end{prop}
\begin{proof}
The existence of the
morphism follows from the theorem of Dee
recalled above, and the descriptions of the two fibre products are
clear from its very definition.
% \mar{ME: I agree with the description of the first fibre product, but not the 
% second --- see my preceding marginal comment.}
%that the corresponding morphism $\Spf
%R_{\rhobar}^{\square,\cO'}\to\cX_d$ is a torsor for the formal
%completion of $\GL_d$ along the closed
%subgroup given by the centraliser of~$\rhobar$. 
To see that this morphism is versal, it suffices to show that 
if $\rho: G_K \to \GL_d(A)$ is a representation with $A$ a
finite Artinian $\cO$-algebra, and if $\rho_B: G_K \to \GL_d(B)$
is a second representation, with $B$ a finite Artinian $\cO$-algebra
admitting a surjection onto $A$, such that the base change
$\rho_A$ of $\rho_B$ to $A$ is isomorphic to $\rho$
(more concretely, so that there exists $M \in \GL_d(A)$ with
$\rho = M \rho_A M^{-1}$),
then we may find $\rho': G_K \to \GL_d(B)$ which lifts $\rho$,
and is isomorphic to $\rho_B$.  But this is clear: the natural
morphism $\GL_d(B) \to \GL_d(A)$ is surjective, 
and so if $M'$ is any lift of $M$ to an element of $\GL_d(B)$,
then we may set $\rho' = M' \rho_B (M')^{-1}$.
% \mar{ME: Why ``in particular''?  I agree that the morphism is versal,
% 	but why/how are we deducing it from the torsor statement? TG:
%         needs to be fixed in proper.tex p121 as well and maybe even in
%       CEGS? ME: Added proof}
\end{proof}
% \begin{rem}\label{rem: versal Galois is iso to versal ring for phi Gamma}
%    In fact
% \mar{ME: Is the assertion of this remark really true?  Doesn't framing
% a $(\varphi,\Gamma)$-module lead to the kind of huge pro-Artinian def.\ rings
% we (crazily) considered in {\tt proper.tex}? TG: agreed; I am just
% going to comment out this remark.}
% $R_{\rhobar}^{\square,\cO'}$ is isomorphic to the universal lifting ring for 
% % \mar{ME: $\cO'$ is being used here, before its actually defined. TG:
% %   rearranged to hopefully sort this out?}
% $(\varphi,\Gamma)$-modules lifting the $(\varphi,\Gamma)$-module
% corresponding to $\Spec\F'\to\cX_d$, an isomorphism being given by fixing
% a choice of bases for the universal Galois lifting and universal
% $(\varphi,\Gamma)$-lifting; but there is no canonical such
% isomorphism.
% \end{rem}
% , and it 
% \mar{ME: There's probably some stupid issue here where choosing a basis
% 	for a Galois rep'n is not the same as choosing a basis for a $(\varphi,
% 	\Gamma)$-modules, so they are not literally the same ring. TG:
%      Good point! I guess the real problem here is that we never defined ``lifting
%     ring'', so we never say anything about bases, and we don't do that
%   on the $(\varphi,\Gamma)$-side either, although we did in
%   \cite{EGstacktheoreticimages}. Leaving this for now but I guess we
%   need to expand on this, probably actually defining the lifting rings
% would be a good start? I guess the correct way to set this up is that
% the two are related by fixing bases on each side
% (Galois/$(\varphi,\Gamma)$) for the universal objects? TG: OK wrote
% something, feel free to improve. TG: plan is to somewhere have the
% deformation ring with framings for both $\rhobar$ and $M$. In fact it
% probably belongs here. But in fact perhaps since versal is just a
% lifting property it can be directly verified on the Galois side
% without ever needing to frame twice.}

Because of the equivalence between $(\varphi,\Gamma)$-modules and
Galois representations with~$\Zp$-coefficients, % recalled in Section~\ref{subsec: phi Gamma to Galois}
and because of the traditional notation $\rho$ for Galois
representations, we will often denote a family of rank $d$
projective \'etale
$(\varphi,\Gamma)$-modules over~$T$, i.e.\ a morphism
$T \to \cX_d$, by $\rho_T$, or some similar notation.    We caution the
reader that this notation is chosen purely for its psychological 
suggestiveness; a family of $(\varphi,\Gamma)$-modules over a general
base $T$ does not admit a literal interpretation in terms of Galois
representations.
%\mar{ME: Here and in 3.7 below we have $(G_K,\varphi)$ (and sometimes
%$(G_{K_\infty},\varphi)$) rather
%than $(\varphi, G_K)$ (or $(\varphi, G_{K_{\infty}})$).  Is there a
%reason for that? TG: surely there is no reason; do you mean that by
%analogy with $(\varphi,\Gamma)$ the $\varphi$ should be first, or do
%we permute the two variables (which would be worse)? TG: pretty sure
%this is now fixed, e.g.\ I have searched for all occurrences of
%``phi)'' in the source, and they're now just Herr complex stuff. Feel
%free to comment out. ME: Thanks!}
\subsection{Galois representations associated to
  \texorpdfstring{$(\varphi,G_K)$}{(phi,Galois)}-modules}\label{subsubsec:
Galois reps for GK phi}
At times
it will be convenient to consider the Galois representations
associated to \'etale $(\varphi,G_K)$-modules. Let~$A$ be a finite
$\cO/\varpi^a$-algebra for some~$a\ge 1$, and let ~$M$ be a finite
projective \'etale $(\varphi,G_K)$-module (resp.\
\'etale $(\varphi,G_{K_\infty})$-module) with $A$-coefficients. Then
we may apply the equivalence of 
Proposition~\ref{prop: equivalences of categories to Ainf}
to obtain an \'etale $(\varphi,\Gamma)$-module (resp.\
an \'etale $\varphi$-module) from $M$;
applying
the functor~$T_A$ (resp.\
$T_{\infty,A}$) of Section~\ref{subsec:Galois reps with coeffs} to this 
latter object yields a 
$G_K$-representation (resp.\
$G_{K_\infty}$-representation)
on a finite free~$A$-module,
which we denote simply by~$T_A(M)$
(resp.\ ~$T_{\infty,A}(M)$).

By passage to the limit over~$a$, we can extend these
functors to the case where~$A$ is a finite $\cO$-algebra, and in
particular to the case that~$A=\cO'$ is the ring of integers in a
finite extension~$E'/E$. As in  Section~\ref{subsec:Galois reps with coeffs},
we can and do
then further
extend these functors to the case~$A=\Fpbar$.

While we have defined the functors~$T_A(M)$ (resp.\ $T_{\infty,A}(M)$) 
via our equivalences of categories,
we note that both also admit a more direct description: % \mar{TG:
  % I made no effort to think about this carefully, just writing off the
  % top of my head, so we should check this! I think it should be OK at
  % least if~$A$ is finite.}
namely $T_A(M)=M^{\varphi=1}$
(resp.\ $T_{\infty,A}(M)=M^{\varphi=1}$). Similarly, the composites of
the functors $\mathbb{D}_A$ and $\mathbb{D}_{\infty,A}$ with the equivalences of 
Proposition~\ref{prop: equivalences of categories to Ainf} admit the
following simple description: if $T_A$ is a finite free $A$-module
with a continuous action of $G_K$ (resp.\ $G_{K_\infty}$), then the
corresponding \'etale $(\varphi,G_K)$-module (resp.\
$(\varphi,G_{K_\infty})$-module) is given by
$W(\C^\flat)_A\otimes_AT_A$ (with~$\varphi$ acting on the first
factor in the tensor product, and $G_K$ (resp.\ $G_{K_\infty}$) acting diagonally).


% \mar{TG: please check that I'm not writing nonsense! ME: It looks good.  Do we
% actually have a proof? (More constructively:  is there a quick way to deduce/explain
% why in terms of our almost Galois descent results?}


\subsection{Galois representations over certain
  fields}\label{subsubsec: Galois repns into certain fields}
There is one more context in which we will need to consider the correspondence 
between Galois representations and \'etale $(\varphi,\Gamma)$-modules.
Namely, suppose that $A$ is a complete Noetherian local $\Z_p$-algebra
which is a domain,
and which has finite residue field.
Let $\cK$ denote the fraction field of~$A$, and as usual, let $\mathfrak m$
denote the maximal ideal of~$A$.
We say that a representation $\rho: G_K \to \GL_d(\cK)$ is {\em continuous}
if there exists a finitely generated $A$-submodule $L$ of $\cK^d$ which spans
$\cK^d$ over $\cK$ (a ``lattice''), such that $G_K$ preserves $L$ and
acts continuously (when $L$ is given its $\mathfrak m$-adic topology).

We write $\widehat{\A}_{K,\cK}:= \widehat{\A}_{K,A}\otimes_A \cK$,
and also write
$\AKnrhatcK := \AKnrhatA\otimes_{A}\cK.$
We have the obvious notion of a projective $(\varphi,\Gamma)$-module with
coefficients in $\widehat{\A}_{K,\cK}$.  We say that such a $(\varphi,\Gamma)$-module
$D$ is {\em \'etale} if there exists a (not necessarily projective) \'etale
$(\varphi,\Gamma)$-module $D_A$ over $\widehat{\A}_{K,A}$ contained in $D$ such
that the evident morphism
$\widehat{\A}_{K,\cK} \otimes_{\widehat{\A}_{K,A}} D_A \to D$ is an isomorphism.

Then, analogously to the case of $A$ itself, we have
functors~$T_{\cK}$, $\mathbf{D}_{\cK}$ defined by  \[\mathbb{D}_{\cK}(T)=(\AKnrhatcK
  \otimes_{\cK}T)^{G_{\Kcyc}},\]  \[T_{\cK}(M)=(  \AKnrhatcK \otimes_{\widehat{\A}_{K,\cK}}M)^{\varphi=1}\]% \mar{TG:
% obviously this needs to be filled out and made correct.}
which give equivalences of categories between the category of finite projective formal \'etale
$(\varphi,\Gamma)$-modules with $\cK$-coefficients, %\mar{ME: Have to define these in the
%$\cK$-context --- \'etale will again be defined via existence of a lattice; now an
%$\widehat{\A}_{K,A}$-lattice!}
and the category of finite dimensional $\cK$-vector spaces with a continuous action of~$G_K$.
(This is proved by passing to lattices on each side, and using 
the results of~\cite{MR1805474}; note that the lattices 
involved need not be projective in general, and so we apply those results 
in their full generality.)

This formalism is most often applied in the literature in the case when $A = \Z_p$,
so that $\cK = \Q_p$.  However, we will not consider 
\'etale $(\varphi,\Gamma)$-modules with $\Q_p$-coefficients in this book. Rather
we will apply the preceding formalism only once, namely in our analysis
of the closed $\Fbar_p$-points of $\cX_d$,
which we make in Section~\ref{subsec:closed points};
and in this application,
we will take $A$ to be a complete local domain of characteristic~$p$.




\section{\texorpdfstring{$(\varphi,G_K)$}{(phi,Galois}-modules
and restriction}
\label{subsec:restriction}
% \mar{TG: all of this material can be moved back into
% the previous section, basically as soon as the relevant objects are defined.}
Our goal in this section is to define certain morphisms of stacks
which are the analogues,
for families of $(\varphi,\Gamma)$-modules,
of the restriction functors on Galois representations
$\rho \mapsto \rho_{| G_{K_{\infty}}}$
and $\rho \mapsto \rho_{|G_L}$ (for any finite extension $L$ of~$K$).

%\section{The morphism \texorpdfstring{$\cX_d\to
%    \cR_{\BK,d}$}{X->R}}

\begin{defn}\index{$\cR_{\BK,d}$}
	% \mar{ME: Probably should cite \cite{EGstacktheoreticimages}
	% 	for details/background on this}
We let~$\cR_{\BK,d}$ denote the moduli stack
of rank~$d$ projective \'etale $\varphi$-modules over~$\OEA$; that is,
the stack~$\cR_d$ constructed in Section~\ref{subsec: moduli stacks of phi
  modules} in the case when~$\A_A=\OEA$.
\end{defn}

Recall that we also have the stack~$\cX_d$ of \'etale $(\varphi,\Gamma)$-modules
of Definition~\ref{defn: Xd}. We now construct a %can now show the existence of a
morphism~$\cX_d\to\cR_{\BK,d}$ % \mar{TG: put in the definition of
  % the target; it's just $u$-\'etale $\varphi$-modules.}
which corresponds to the restriction of Galois representations from~$G_K$
to~$G_{K_\infty}$. 
\begin{prop}% \mar{TG: do we actually need this level of generality
    % in~$A$?}
  \label{prop: natural morphism X to R}There is a canonical
  morphism~$\cX_{d}\to\cR_{\BK,d}$. If~$A$ is a complete local Noetherian
  $\cO$-algebra with finite residue field, or equals  $\Fpbar$, then the
  morphism~$\cX_d(A)\to\cR_{\BK,d}(A)$ is given by
  restriction of the corresponding representation of~$G_K$
  to~$G_{K_\infty}$.% \mar{TG: I'm not sure in what generality we use
    % this claim, presumably it will come up when we
    % cite~\cite{MR2745530}, but it's surely true in this generality.}
\end{prop}
\begin{proof}% \mar{TG: have to justify the claim about the Galois
    % representations, which will just come from the construction.}
  If~$A$ is a finite type $\cO/\varpi^a$-algebra, for some~$a\ge
  1$, then we define a morphism of groupoids $\cX_{d}(A) \to\cR_{\BK,d}(A)$
  via the equivalences of categories of Proposition~\ref{prop:
    equivalences of categories to Ainf}; that is, given a finite
  projective \'etale $(\varphi,\Gamma)$-module with $A$-coefficients,
  we form the corresponding finite projective \'etale $(\varphi,G_K)$-module,
  which yields a finite projective \'etale $(\varphi,G_{K_\infty})$-module by
  restricting the~$G_K$-action to~$G_{K_\infty}$, and thus gives an \'etale
  $\varphi$-module over~$\OEA$. Since both~$\cX_d$ and~$\cR_{\BK,d}$
  are limit preserving (by  Corollary~\ref{cor: basic properties of C and R p adic stacks} and Lemma~\ref{lem:X is limit preserving}), it follows from~\cite[Lem.\ 2.5.4, Lem.\
  2.5.5~(1)]{EGstacktheoreticimages} that this construction determines
  a morphism ~$\cX_{d}\to\cR_{\BK,d}$.

%\mar{ME: This argument is written in a strange way --- e.g. ``Recall
%  that we have {\em a} functor \dots'' rather than ``Recall that {\em
%    the functor} \dots''. TG: rewrote slightly, but not sure quite
%  what you have in mind. ME: Modified it a bit.} 
  Suppose now that ~$A$ is a complete local Noetherian $\cO$-algebra
  with finite residue field. % \mar{TG: need to deal with formal \'etale
    % $(\phi,\Gamma)$-modules}
  %Recall that the
  %functor~$T_{A}$ takes formal \'etale $(\varphi,\Gamma)$-modules with
  %$A$-coefficients to representations of~$G_K$ with $A$-coefficients,
  %and similarly the functor $T_{\infty,A}$ takes  \'etale
  %$\varphi$-modules over ~$\hOEA$
  %to~$G_{K_\infty}$-representations.
  Let $M$ be a formal
  \'etale $(\varphi,\Gamma)$-module corresponding to a morphism
  $\Spf A\to\cX_{d}$, and~$M_\infty$ be the \'etale $\varphi$-module
  over~$\hOEA$ corresponding to the composite
  $\Spf A\to \cX_d\to\cR_{\BK,d}$.
  The relationship between $M$ and $M_{\infty}$ is
  expressed as an isomorphism
  % \mar{ME: I think that $\A_{K,A}$ in the following 
  % discussion should probably be $\widehat{\A_{K,A}}$ instead. TG:
  % correct, fixed, commenting out.}
  \numequation
  \label{eqn:module isomorphism}
  \widehat{W(\C^\flat)}_A \otimes_{\widehat{\A}_{K,A}} M
  \iso 
    \widehat{W(\C^\flat)}_A\otimes_{\hOEA}M_\infty.
  \end{equation}
  
  Recall that the Galois representation associated to $M$ is 
  defined via $T_A(M) := (\AKnrhatA\otimes_{\widehat{\A}_{K,A}}M)^{\varphi = 1},$
  and that the evident
  %We have a natural $(\varphi,G_{K})$-equivariant isomorphism
  $(\varphi,G_{K})$-equivariant $\AKnrhatA$-linear morphism
  \[\AKnrhatA\otimes_{A}T_A(M)\to
    \AKnrhatA\otimes_{\widehat{\A}_{K,A}}M \]
  is then an isomorphism (see for
  example~\cite[Prop.\ 2.1.26]{MR1805474}).   % and thus a
  Thus there is an induced natural
  $(\varphi,G_{K})$-equivariant
  isomorphism
  \[\widehat{W(\C^\flat)}_A\otimes_{A}T_A(M)\isoto
    \widehat{W(\C^\flat)}_A\otimes_{\widehat{\A}_{K,A}}M, \] where we recall
  that we write $\widehat{W(\C^\flat)}_A$ for
the $\m_A$-adic completion of~$W(\C^\flat)_A$.  Similarly, we obtain a natural
  $(\varphi,G_{K_\infty})$-equivariant isomorphism
  \[ \widehat{W(\C^\flat)}_A\otimes_{A}T_{\infty,A}(M_\infty)\isoto
    \widehat{W(\C^\flat)}_A\otimes_{\hOEA}M_\infty.\]
%It follows that we have
Combining these two isomorphisms with~\eqref{eqn:module isomorphism},
we obtain a natural $(\varphi,G_{K_\infty})$-equivariant isomorphism
\[\widehat{W(\C^\flat)}_A\otimes_{A}T_A(M)\isoto
  \widehat{W(\C^\flat)}_A\otimes_{A}T_{\infty,A}(M_\infty). \]
Each of $T_A(M)$ and $T_{\infty,A}(M_{\infty})$ 
is a free $A$-module with trivial~$\varphi$-action,
and so if we pass to $\varphi$-invariants in this isomorphism,
and take into account
Lemma~\ref{lem: phi invariants in W C flat with coefficients},
we obtain an isomorphism
$T_A(M)|_{G_{K_\infty}}\isoto T_{\infty,A}(M_\infty)$, 
which is what we wanted to show.
%We claim that taking~$\varphi$-invariants induces an isomorphism
%$T_A(M)|_{G_{K_\infty}}\isoto T_{\infty,A}(M_\infty)$. To see this,
%note that since~$T_A(M)$ and~$T_{\infty,A}(M_\infty)$ are free
%$A$-modules with a trivial $\varphi$-action, we are reduced to showing
%that $(\widehat{W(\C^\flat)}_A)^{\varphi=1}=A$, which is Lemma~\ref{lem: phi
%  invariants in W C flat with coefficients}.
Finally, if $A=\Fpbar$, % or~$\Zpbar$,
the result follows by taking direct limits.
\end{proof}% The following lemma presumably holds for general
% coefficient rings~$A$, but since we can give a simple proof in the
% only case that we need, we have not investigated this.

%\mar{ME: I don't think this lemma should ultimately reside here --- probably
%it should be back in the initial discussion of $(\varphi,\Gamma)$-modules and
%Galois reps.}

%Recall the following standard result about invariants under absolute Frobenius.
%
%\begin{lem}
%\label{lem:Frob invariants}
%If $X$  is an $\F_p$-scheme, if $(\F_p)_X$ denotes
%the sheaf of locally constant $\F_p$-valued functions on~$X$ {\em (}which is 
%a subsheaf of~$\cO_X$ --- if $f:X\to \Spec \F_p$ is the structure morphism,
%then $(\F_p)_X = f^{-1} \cO_{\Spec \F_p}${\em )},
%and if $F: X \to X$ denotes the absolute 
%Frobenius, the $(\F_p)_X = \cO_X^{F = 1}$.
%\end{lem}
%\begin{proof}
%Evidently the embedding of
%$(\F_p)_X$ into $\cO_X$ has image landing  in $\cO_X^{F=1}$, 
%and so induces an embedding 
%$ (\F_p)_X \hookrightarrow \cO_X^{F=1}.$
%%
%%Note that the closed immersion $X_{\red} \hookrightarrow X$
%%induces an isomorphism $(\F_p)_{X_{\red} \iso (\F_p)_X$ (because
%%$(F_p)_X$ depends only on the underlying topological space of~$X$)
%%and also induces an embedding 
%%$\cO_{X_{\red}}^{F=1} \iso \cO_X^{F=1}$
%%(because any section $a$ of the kernel over an open subset $U$
%%of $X$ is locally nilpotent on~$U$, and also satisfies $a = a^p$;
%%and hence vanishes identically on~$U$).
%%
%To show that this is an isomorphism,
%we may work locally, and so assume $X = \Spec A$ for some $\F_p$-algebra~$A$.
%Let $ a\in A$ satisfy $a^p = a$,
%and consider the morphism $\F_p[x] \to A$ defined by $x \mapsto a$.
%Since the kernel of this morphism is a factor of the polynomial
%$x^p - x = \prod_{i \in \F_p} (x - i),$
%its image $B$ is a product of copies of $\F_p$,
%and so $\Spec B$ is  a disjoint union of copies of~$\Spec \F_p$.
%Since by construction $a$ is pulled back from an element of $\Gamma(\Spec B, \O_{\Spec B})$
%along the morphism $\Spec A \to \Spec B$, we see that indeed $a$ is a locally
%constant $\F_p$-valued function on~$\Spec A$.
%\end{proof}


% \begin{lem}\mar{TG: n.b.\ I think we're back to only needing this in
%     the CNL case where there is a proof commented out in the file.}
%   \label{lem: phi invariants in W C flat with coefficients}Let~$A$ be
%   a $p$-adically complete $\cO$-algebra. Then
%   $(W(\C^\flat)_A)^{\varphi=1}=A$.\mar{TG: sketching an attempt at
%     a proof, didn't try to make it rigorous.}
% \end{lem}
% \begin{proof}We first show that
%   $(W(\C^\flat)_A)^{\varphi=1}=(\AAinf{A})^{\varphi=1}$.\mar{TG: going
%   to write an argument for $\cO/\varpi^a$-algebras, didn't think yet
%   about passing to the limit; also being a big vague about ``u''.}
% Suppose that $x\in (W(\C^\flat)_A)^{\varphi=1}$; then for
% some~$\epsilon>0$ we have $x\in
% u^{-\epsilon}\AAinf{A}$. Since~$varphi$ is a bijection on both
% $W(\C^\flat)_A$ and $\AAinf{A}$, it follows that $x=\varphi^{-1}(x)\in u^{-\epsilon/p}\AAinf{A}
% $. Iterating, we conclude that $x\in
% \cap_{\epsilon>0}u^{-\epsilon}\AAinf{A}$, so that
% $x\in\AAinf{A}$.\mar{TG: this presumably needs an argument but is
%   hopefully true?!}

% We claim that there exists some~$a\in A$ such that $x-a\in
% I\AAinf{A}$. \mar{TG: $I$ is some notation I'm making up for the stuff
%   with positive $u$-adic valuation. Basic idea is that if it has
%   positive valuation we get a contradiction, although for general $A$
%   it's not so clear that there is a well-defined valuation?}

% Since~$\C^\flat$ is separably closed, we have a short exact sequence \[0\to\F_p\to\C^\flat\stackrel{1-\varphi}{\to} \C^\flat\to
%   0, \] from which it follows easily that there is a short exact
% sequence \[0\to\Z_p\to \Ainf \stackrel{1-\varphi}{\to} \Ainf \to
%   0.\]  Since
% $\Ainf$ is $\Zp$-flat, if~$B$ is any $\Z_p$-algebra we may tensor with~$B$ and obtain a short
% exact sequence 
% \[0\to B\to \Ainf\otimes_{\Zp}B\stackrel{1-\varphi}{\to}
%   \Ainf\otimes_{\Zp}B\to 0.\]\mar{TG: now that I'm writing this it's
%   less clear how to argue, so I'm abandoning it for now!}
% \end{proof}

% \begin{prop}\label{prop: X to R separated ind representable}
%   \mar{TG: not sure quite what we want here but this is what a label on an
%     arrow says}The morphism~$\cX_d\to\cR_{\BK,d}$ is separated and ind-representable.
% \end{prop}
% \begin{proof}
% 	\mar{ME: Sketch} 
% 	The product $\cX_d\times_{\cR_{\BK}} \cX_d$ can be described explicitly
% 	as follows: its $A$-valued points are pairs $(M_1,M_2)$ of 
% 	$(\varphi,G_K)$-modules with a $G_{K_{\infty}}$-equivariant isomorphism
% 	between them.  The diagonal is just $M \mapsto (M,M),$ with
% 	the isomorphism being the identity.
	
% 	For ind-representability:  The preceding explicit description
%         shows that the diagonal is a monomorphism,
% 	and the source and target are Ind-algebraic.  I think this should be 
% 	enough to give Ind-representability. 
% 	(Just as a morphism of algebraic stacks is representable
% 	iff its diagonal is a monomorphism.)\mar{TG: fwiw I don't
%           think we have defined Ind-representability, and it's not
%           clear to me where we've used it.}

% 	For separatedness: the diagonal is already seen to a monomorphism,
% 	so we have to show that it is a proper monomorphism, i.e.\ a closed
% 	immersion.  For this, we have to check that if we 
% 	are given a pair of $(\varphi,G_K)$-modules,
% 	with a $G_{K_{\infty}}$-equivariant isomorphism,
% 	then the locus where this isomorphism becomes $G_K$-equivariant
% 	is closed.  This should just be a matter of observing that 
%         equivariance is determined by one
% 	equation (in $W(\C^\flat)$, so really an infinite
% 	collection of equations) for each element of $G_K$.\mar{TG: I
%           was wondering about this too - do we have to try to say
%           something like we can regard~$W(\C^\flat)$ as an
%           infinitely generated free (?) A-module?}
% \end{proof}

% \mar{ME: This was already proved in the preceding result, so probably
% 	should just be folded into that.}
\begin{prop}
  \label{prop: diagonal of X to R}The diagonal morphism~$\Delta:
  \cX_d\to\cX_d\times_{\cR_{\BK,d}}\cX_d$
  induced by the morphism of
  Proposition~{\em \ref{prop: natural morphism X to R}}
  is a closed immersion.
\end{prop}
\begin{proof}	%\mar{ME: Sketch} 
	The product $\cX_d\times_{\cR_{\BK}} \cX_d$ can be described explicitly
	as follows: its $A$-valued points are pairs $(M_1,M_2)$ of \'etale
	$(\varphi,G_K)$-modules with a $G_{K_{\infty}}$-equivariant isomorphism
	between them.  The diagonal is defined by $M \mapsto (M,M),$ with
	the isomorphism being the identity.
%	
%  The diagonal is evidently a monomorphism.
	% so we have to show that it is a proper monomorphism, i.e.\
To see that this defines a closed
	immersion,  we have to check that if we 
	are given a pair~$(M_1,M_2)$ of $(\varphi,G_K)$-modules with $A$-coefficients,
	together with a $G_{K_{\infty}}$-equivariant isomorphism $f:M_1\isoto M_2$,
	then the locus where~$f$ becomes $G_K$-equivariant
	is closed.   That this is so follows from Corollary~\ref{cor:equivariance locus}.  
	(That corollary shows that, for each $g \in G_K$,
       	there is an ideal $J_g \subseteq A$
	which cuts out the condition for the given
       isomorphism to commute with~$g$.   The ideal $J := \sum_g J_g$
       % \mar{TG: Dat Pham pointed out that we had $J := \bigcap_g J_g$
       %   before, which is wrong! Assuming that you agree, I will try
       %   to find/fix the other instances of this mistake in the
       %   paper. TG: others are now marked by marginal comments.}
       then cuts out the condition for the given isomorphism to be
       $G_K$-equivariant.)
%	\mar{TG: ME suggests thinking about $(M_1\otimes
%          M_2^*)^{G_{K_\infty}, \varphi=1}$. Maybe we can see that
%          being fixed by $G_K$ in particular means fixed by
%          $G_{\Kcyc}$ and so you get down to the original
%          $(\varphi,\Gamma)$-module. Is that a closed condition?}
%
%        This should just be a matter of observing that 
%        equivariance is determined by one
%	equation (in $W(\C^\flat)$, so really an infinite
%	collection of equations) for each element of $G_K$.\mar{TG: I
%          was wondering about this too - do we have to try to say
%          something like we can regard~$W(\C^\flat)$ as an
%          infinitely generated free (?) A-module? TG: still confused
%          about this! We could reduce to the free case, and then maybe
%        we could do something like asking that for all~$g$ we have
%        $g\circ g^*f=f\circ g$ in
%        $\Hom_{W(\C^\flat)_A,\varphi}(g^*M_1,M_2)$, but this seems
%        like a bit of a mess. Alternatively we could consider
%        $\Hom_{G_K,\varphi}(M_1,M_2)\subseteq
%        \Hom_{G_{K_\infty},\varphi}(M_1,M_2)$ and just try to show that this is
%      closed?}
\end{proof}
%  For each finite extension~$L/K$, we can
% carry out the constructions of the earlier sections with~$K$ replaced
% by~$L$, and we write~$\cX_{L,d}$ for the corresponding moduli
% stack of \'etale $(\varphi,\Gamma)$-modules.

%We can also define a restriction map for finite extensions, as in the
%following lemma.
We now construct the restriction maps corresponding to finite extensions
of~$K$.

\begin{lem}
  \label{lem: restriction from X to Xs}Let~$L/K$ be a finite
  extension. There is a canonical
  morphism~$\cX_{K,d}\to\cX_{L,d}$, which is representable by
  algebraic spaces, affine, 
  %\mar{ME: The proof seems to show that it's even affine.}
  and of finite
  presentation. If~$A$ is a complete local Noetherian
  $\cO$-algebra with finite residue field,
  % \mar{ME: Need to constrain the residue field in the complete
  %         local case, I guess? In the preceding result, too. TG: which
  %       previous result? I don't see another missing one. TG: pretty
  %       sure this is fine, commenting out.}
  or equals~$\Fpbar$, then the
  morphism~$\cX_{K,d}(A)\to\cX_{L,d}(A)$ is given by
  restriction of the corresponding representation of~$G_K$
  to~$G_{L}$.
  %If~$A$ is a % finite
  % $\cO$-algebra and $M$ is an \'etale $(\varphi,\Gamma)$-module
  % corresponding to a morphism $\Spec A\to\cX_{d}$, and~$M_s$ is
  % the \'etale $\varphi$-module over~$\OEA$ corresponding to the
  % composite $\Spec A\to \cX_d\to\cX_{K_s,d}$, then we have a
  % natural isomorphism\mar{TG:
  %   making up notation for passage to Galois
  %   representations} \[T(M)|_{G_{K_s}}\cong T(M_s). \]
\end{lem}
% \mar{ME: Since at other points in the development of the theory
% 	we'll need restriction to arbitrary finite extensions,
% 	this lemma
% 	should handle the case of arbitrary finite extensions $L/K$,
% 	not just $K_s$.}
\begin{proof}% \mar{TG: previously ME had sketched an argument that
    % didn't involve reduction to the Galois case, but instead went
    % to~$L_\infty$ - I couldn't understand this argument, so I
    % commented it out.}
  % \mar{TG: hopefully it's easy to
  % describe the morphism $\cX_d\to\cX_{K_s,d}$ purely on the level of
  % $(\varphi,\Gamma)$-modules; didn't think about this yet. Then have
  % to check representable and finite type.} 
%\mar{ME: Proof sketch, written for general $L$ in place of $K_s$}
As in the case of Proposition~\ref{prop: natural morphism X to R},
the morphism is defined first in the case of $\cO$-algebras that
are of finite type over $\cO/\varpi^a$ for some $a$ by
  applying the equivalences of categories of Proposition~\ref{prop:
    equivalences of categories to Ainf}:
we send an \'etale $(\varphi,G_K)$-module~$M$ 
to an \'etale $(\varphi,G_L)$-module~$M_L$ via restricting the continuous
$G_K$-action to a continuous $G_L$-action.
Since $\cX_{K,d}$ and $\cX_{L,d}$ 
  are limit preserving (by Lemma~\ref{lem:X is limit preserving}),
  it follows from~\cite[Lem.\ 2.5.4, Lem.\ 2.5.5~(1)]{EGstacktheoreticimages}
  that this construction determines a morphism $\cX_{K,d} \to \cX_{L,d}$.

Before verifying the various claimed properties of this morphism,
we confirm that it has the claimed effect on associated Galois representations.
To this end, we note that
if~$A$ is a complete local Noetherian
  $\cO$-algebra, then, as in the proof of Proposition~\ref{prop:
    natural morphism X to R}, it follows from the definition that we
  have a %\mar{TG: again probably need completion}
  natural $(\varphi,G_{L})$-equivariant isomorphism
  \[\widehat{W(\C^\flat)}_A\otimes_{A}T_A(M)\isoto
    \widehat{W(\C^\flat)}_A\otimes_{A}T_{A}(M_L), \] and by Lemma~\ref{lem: phi
    invariants in W C flat with coefficients},
  taking~$\varphi$-invariants induces an isomorphism
  $T_A(M)|_{G_L}\isoto T_{A}(M_L)$. We deduce the same statement if
    ~$A=\Fpbar$ % or~$A=\Zpbar$
    by taking direct limits.  % \mar{ME: I
      % suppose we should check this does the claimed thing on Galois
      % rep.\ in the complete local case. TG: done, commenting out.}
%As in the proof of Proposition~\ref{prop: diagonal of X to R},
%it follows from Corollary~\ref{cor:equivariance locus}
%that the diagonal morphism $\cX_{K,d} \to \cX_{K,d}\times_{\cX_{L,d}} \cX_{K,d}$
%is a closed immersion, which shows that $\cX_{K,d} \to \cX_{L,d}$
%is Ind-representable and separated.  
%
% But we can do better: Let $\{g_i\}$ be a set of finitely 
% many coset representatives for $G_L$ in $G_K$.  Then to
% endow an \'etale $(\varphi,G_L)$-module $M$ with a $G_K$-structure
% amounts to giving an isomorphism of
% \'etale $(\varphi,G_{L_{\infty}})$-modules
% \mar{ME: This could be the cyclotomic $L_{\infty}$ 
% 	or the Breuil--Kising $L_{\infty}$,
% 	or any other appropriate $L_{\infty}$.}
%To go further,

We now verify the claimed properties of the restriction morphism $\cX_{K,d}
\to \cX_{L,d}$.
Suppose firstly that~$L/K$ is Galois, and let $\{g_i\}$
be a set of (finitely many) coset representatives for $G_L$ in
$G_K$. For each~$i$, we may give~$g_i^*M$ the structure of a
$(\varphi,G_L)$-module by letting each $h\in G_L$ act as~$g_i^{-1}h g_i$
acts on~$M$; then to extend the action of~$G_L$ to an action of~$G_K$
is to give an isomorphism of~$(\varphi,G_L)$-modules $g_i^*M \to M$
for each index~$i$, satisfying a slew of compatibilities. (Since $G_L$
is open in $G_K$, the continuity of the $G_K$-action is automatic,
given that the $G_L$-action that it is extending is continuous.)

We let $\cY_i$ denote the stack classifying objects $M_A$ 
of~$\cX_{L,d}(A)$,
endowed with an isomorphism 
$g_i^*M_A\isoto M_A$.  If we regard $g_i^*$  as an automorphism
of~$\cX_{L,d}(A)$,
then we  may form  its  graph~$\Gamma_i$,
and we then have an isomorphism of stacks
$$\cY_i \iso
\cX_{L,d}(A)\times_{\Delta,
\cX_{L,d}(A) \times
\cX_{L,d}(A),\Gamma_i}
\cX_{L,d}(A).$$
(Here $\Delta$ denotes the diagonal of~$\cX_{L,d}$;
\emph{cf.}\ the definition of $\cR_d^{\Gamma_{\disc}}$ in 
Section~\ref{subsec: defn of Xd}
above, together 
with Proposition~\ref{prop:Gamma-disc modules as fixed points; etale case}.)
The projection onto the second factor $\cY_i \to 
\cX_{L,d}(A)$,
which corresponds to
forgetting the isomorphism, is a base-change of the
diagonal~$\Delta,$
and so is representable by algebraic spaces, affine and of finite presentation by Proposition~\ref{prop: properties of X pulled back from R}.

We can rephrase our interpretation of objects 
of~$\cX_{K,d}(A)$ %(for some~$A$)
as objects of $\cX_{L,d}(A)$
endowed with isomorphisms
$g_i^*M_A\isoto M_A$ satisfying certain compatibilities
as the existence of  a closed immersion
%\numequation\label{eqn: semistable in inertia stack product
$$
\cX_{K,d}\into
\cY_1 \times_{\cX_{L,d}}\times \cdots \times_{\cX_{L,d}}
\cY_n
$$  %\end{equation}
(the point  being that the compatibilities
arise as
base-changes of the double diagonal, and so impose closed conditions).
Since the morphisms ~$\cY_i\to\cX_{L,d}$ are
representable by algebraic spaces, affine and of finite presentation,
it follows that the same is true of the morphism
$\cX_{K,d}\to\cX_{L,d}$, as claimed.


% Since the~$\Isom$ spaces for~$\cX_{L,d}$ are finite type affine group
% schemes, by Proposition~\ref{prop:X is an Ind-stack}, and the
% aforementioned slew of compatibilities evidently give closed
% conditions, we see that $\cX_{K,d}\to \cX_{L,d}$ is representable by
% algebraic spaces, and is in fact affine and of finite presentation, as
% claimed. 
% \mar{TG: is there a
  % reason not to change this (and the statement of the result) to finite presentation? ME: No, and I've done this now.}
%\mar{TG: need to check that $\Isom$ is really enough and we don't need
%a graph argument or similar. ME: I think we need to argue as in
%4.5.26,  unfortunately. TG: have basically copied that argument to here.
% ME: Looks good, thanks!}



Finally, let~$L/K$ be a general finite extension. Let~$M/K$ be the
normal closure of~$L/K$, so that we have
morphisms \[\cX_{K,d}\to\cX_{L,d}\to\cX_{M,d}.\]
By what we have already proved, both the second arrow and the composite
of the two arrows are representable by algebraic
spaces and affine, and of finite presentation.
An affine morphism is separated, and thus has affine and finite type diagonal.
One sees (e.g.\ by
Lemma~\ref{lem:finiteness of diagonals})
that the diagonal of an affine morphism of finite type (and so, in particular,
the diagonal of an affine morphism of finite presentation)
is furthermore of finite presentation.
A standard graph argument,
in which we factor the first morphism as
$$\cX_{K,d} \hookrightarrow \cX_{K,d} \times_{\cX_{M,d}} \cX_{L,d} 
\to \cX_{L,d}$$
(the first arrow in this factorization
being the graph of the first morphism,
which is a base change of the diagonal
$\cX_{L,d} \to \cX_{L,d} \times_{\cX_{M,d}} \cX_{L,d} $,
and the second arrow being the projection, which is a base change
of the composite $\cX_{K,d} \to \cX_{L,d} \to \cX_{M,d}$),
then shows that the first morphism is also
representable by algebraic spaces and affine, and of finite presentation.% \mar{TG: this ``standard
  % graph argument'' is exactly what I failed to find, so I'd like to
  % see it written, if you don't mind - I ended up wanting to know
  % that the diagonal $\cX_{K,d}\to\cX_{K,d}\times_{\cX_{L,d}}\cX_{K,d}$
  % was finite type. ME: Done.}
%
%We may factor the
%morphism $\cX_{K,d}\to\cX_{L,d}$ as the composite \[\cX_{K,d}\to
%  \cX_{K,d}\times_{\cX_{M,d}}\cX_{L,d}\to\cX_{L,d}.\] The first
%morphism is a base change of the diagonal
%$\cX_{K,d}\to\cX_{K,d}\times_{\cX_{L,d}}\cX_{K,d}$, which we already
%noted is a closed immersion, and in particular representable and
%affine. The second morphism is a base change
%of~$\cX_{K,d}\to\cX_{M,d}$, which is representable and
%affine because~$M/K$ is Galois. It follows that the composite
%$\cX_{K,d}\to\cX_{L,d}$ is representable and affine.
%Finally, to see that~$\cX_{K,d}\to\cX_{L,d}$ is of finite
%type\mar{TG: hopefully easy? If they were algebraic stacks, not just
%  Ind-algebraic, it would be \cite[\href{https://stacks.math.columbia.edu/tag/06Q6}{Tag 06Q6}]{stacks-project}.}
\end{proof}

% \begin{lem}\mar{TG: needs a proof, might not have correct
%     hypotheses. Hopefully something like this is true! TG: OK it's
%     fine, we already have a closed immersion internal to the lemma
%     above, so we can make this work. Will fill in after lunch.}
%   \label{lem: cancellation affine}Suppose that $\cX\to\cY\to\cZ$ are
%   morphisms of Ind-algebraic stacks over a scheme~$S$, such that each
%   of $\cX\to\cZ$, $\cY\to\cZ$, $\Delta_{\cZ}:\cZ\to\cZ\times_S\cZ$ are
%   representable, affine, and of finite type. Then $\cX\to\cY$ is also
%   representable, affine, and of finite type.
% \end{lem}
% \begin{proof}  We factor $\cX\to\cY$ as the
%   composite \[\cX\to\cX\times_{\cZ}\cY\to\cY\]where the first morphism
%   is base change of $\cY\to\cY\times_{\cZ}\cY$, and the second is the
%   base change of $\cX\to\cZ$. It is therefore sufficient to show that
%   $\cY\to\cY\times_{\cZ}\cY$ is   representable, affine, and of finite
%   type. 
% \end{proof}


\section{Tensor products and duality} \label{subsec: tensor product
  of phi gamma and duality}
% \mar{ME: Maybe this can be moved to the end of \S 2,
% 	because it interupts the flow a bit. TG: agreed, could also
% %  discuss restriction and induction here too.}
% \mar{TG: I don't think that at the moment we need
%   induction/restriction, but if we do we could also discuss them here,
% referencing \cite{2018arXiv180508103D}.}
If~$M$ is a projective \'etale $(\varphi,\Gamma)$-module with $A$-coefficients, we define the dual
$(\varphi,\Gamma)$-module by \[M^\vee:=\Hom_{\A_{K,A}}(M,\A_{K,A}); \] if $A$
is a finite $\cO$-module, then there is a natural isomorphism
$T(M^\vee)\cong T(M)^\vee$. Similarly, for
  each~$M$, we let~$M(1):=M\otimes_{\A_{K,A}}\A_{K,A}(1)$ denote the Tate
  twist, where $\A_{K,A}(1)$ denotes the free $(\varphi,\Gamma)$-module of
  rank~$1$ with a generator~$v$ on which $\Gamma$ acts via the
  cyclotomic character and~$\varphi(v)=v$.  % for
  % some generator~$v$.
  % \mar{TG: didn't try to check normalisation here}
  (Note that if~$A=\cO$ then $T(\A_{K,A}(1))=\cO(1)$, the usual
  Tate twist on the Galois side.)  We also define the Cartier dual
$M^*:=M^\vee(1)$, which again in the case that $A$ is a finite
$\cO$-algebra is compatible with the usual notion for Galois
modules. If $T\to\cX_d$ corresponds to a family $\rho_T$,
then we adopt the natural convention of writing ~$\rho_T(1)$, $\rho_T^\vee$ and~$\rho_T^*$ for the
corresponding families.

Given two projective \'etale $(\varphi,\Gamma)$-modules
$M_1$, $M_2$ with $A$-coefficients of respective ranks~$d_1,d_2$, we may form the tensor
product $(\varphi,\Gamma)$-module $M_1\otimes M_2$ (given by the
tensor product on  underlying
$\A_{K,A}$-modules, and the tensor products of the actions of each of~$\varphi$
and $\Gamma$). If~$A$ is a finite $\cO$-module, then we have a natural
isomorphism \[T(M_1\otimes M_2)\cong T(M_1)\otimes T(M_2).\] The
tensor product induces a natural morphism \[\cX_{d_1}\times_{\cO}
  \cX_{d_2}\to\cX_{d_1d_2};\] in particular, twisting by (the
$(\varphi,\Gamma)$-module corresponding to) any
character~$G_K\to\cO^\times$ gives an automorphism of each~$\cX_d$. Given morphisms $S\to\cX_{d_1}$,
$T\to\cX_{d_2}$, denoted % (as in Section~\ref{subsec: phi Gamma to
  % Galois})
by $\rho_S$ and $\rho'_T$, we write $\rho_{S}\boxtimes
\rho'_T$ for the family corresponding to the composite % \mar{TG: should
  % have a base for all these products}
\[S\times_{\cO}
  T\to\cX_{d_1}\times_{\cO}\cX_{d_2}\to\cX_{d_1d_2}.\] If $S=T$ then we
write $\rho_T\otimes\rho'_T$ for the composite  \[T\to T\times_{\cO}
  T\to\cX_{d_1}\times_{\cO}\cX_{d_2}\to\cX_{d_1d_2},\]where the first map is
the diagonal embedding.

\chapter{Crystalline and semistable moduli stacks}\label{sec:
  crystalline and semistable}
In this chapter we build on the results of Appendix~\ref{app: BKF pst} to
define the potentially semistable and potentially crystalline
substacks of~$\cX_{K,d}$. Using the results of the earlier chapters,
and Kisin's results on potentially semistable lifting
rings~\cite{MR2373358}, we show that these stacks are $p$-adic formal
algebraic stacks, and compute the dimensions of their special fibres.
\section{Notation} Recall from Section~\ref{subsubsec: Kummer}
that the embedding~$\gS_A\into\AAinf{A}$ depends on the choice of a
uniformizer~$\pi$ of~$K$, and of a compatible
family~$\pi^{1/p^\infty}$ of $p$-power roots of unity
in~$\overline{K}$. In this chapter we will want to consider all
possible choices of~$\pi$ and~$\pi^{1/p^\infty}$, and consequently we
introduce some notation to do so. We write~$\pi^\flat\in\cO_\C^\flat$ for
the element determined by~$\pi^{1/p^\infty}$. For each~$s\ge 0$ we
write~$K_{\pi^\flat,s}$ for $K(\pi^{1/p^s})$,
and~$K_{\pi^\flat,\infty}$ for~$\cup_sK_{\pi^\flat,s}$. We
write~$\gS_{\pi^\flat,A}$ for the image of~$\gS_A$ in ~$\AAinf{A}$ via
the homomorphism determined by~$u\mapsto[\pi^\flat]$,
and~$\OEApiflat$ for the image of~$\OEA$ in~$W(\C^\flat)_A$
under the corresponding homomorphism. When~$\piflat$ is fixed, we will often write~$u$
for~$[\pi^\flat]$ when discussing~$\gS_{\piflat,A}$.  We write~$E_\pi$ for the
Eisenstein polynomial corresponding to~$\pi$. We also have a natural
map~$\theta:\AAinf{A}\to\cO_{\C,A}$ (where the target is the tensor
product  $\cO_{\C}\cotimes A$, which is 
the quotient of the source by the principal ideal % \mar{TG: is this latter ring
  % defined? Is it better to say we're just quotienting out by a
  % principal ideal?}
generated by the kernel of the usual map $\theta:\Ainf\to\cO_{\C}$).
%\mar{ME: Is this just $\cO_{\C}\cotimes A$? TG: yes, edited to say this.}

\begin{rem}\label{rem: we don't twist our embeddings by phi}In
  contrast to many papers in the literature (e.g.\
  \cite{2016arXiv160203148B,liulattice2}) we do not twist the
  embedding $\gS\to\Ainf$ by~$\varphi$. While such a twist is
  important in applications involving comparisons to Fontaine's period
  rings, it is more convenient for us to use the embedding we have
  given here. % \mar{TG: of course we could use that embedding if we
    % preferred.}
  Since~$\varphi$ is an automorphism of~$\Ainf$, it is
  in any case straightforward to pass back and forth between these two
  conventions. 
\end{rem}


%\mar{ME: This subsection still needs editing/rearranging.}

% Let $E$ be the Eisenstein polynomial corresponding to our fixed
% uniformiser~$\pi$ of~$K$.
\section[Breuil--Kisin modules and Breuil--Kisin--Fargues
  modules]{Breuil--Kisin modules and Breuil--Kisin--Fargues
  modules \sectionmark{Breuil--Kisin--Fargues
  modules}}\sectionmark{Breuil--Kisin--Fargues
  modules}\label{subsec: BK and BKF modules}
\begin{defn}%\mar{TG: add Fargues}
Fix a choice of~$\pi^\flat$.	Let $A$ be a $p$-adically complete
$\cO$-algebra which is topologically of finite type. % finite type $\cO/\varpi^a$-algebra
	% for some $a \geq 1$.
We define a {\em projective Breuil--Kisin module} \index{Breuil--Kisin module} 
(resp.\ a {\em projective Breuil--Kisin--Fargues module}) \index{Breuil--Kisin--Fargues module}  {\em of height at
  most~$h$ with $A$-coefficients} to be a finitely generated projective $\gS_{\pi^\flat,A}$-module
  (resp.\ $\AAinf{A}$-module)
  $\gM$, equipped with a~$\varphi$-semi-linear morphism
  $\varphi:\gM\to\gM$, with the property that the corresponding
  morphism % of $\AAinf{A}$-modules
  $\Phi_{\gM}:\varphi^*\gM\to\gM$ is
  injective, with cokernel killed by~$E_\pi^h$. % \mar{TG: understand how
    % this compares to usual BKF definition}  
If~$\gM$ is a Breuil--Kisin module   then
$\AAinf{A}\otimes_{\gS_{\pi^\flat,A}}\gM$ is a Breuil--Kisin--Fargues module.
\end{defn}
% See Remark~\ref{rem: all BK BKF are effective} for a comparison of
% this definition to the notion of a Breuil--Kisin--Fargues module as
% defined in~\cite{2016arXiv160203148B}.
\begin{rem}
  \label{rem: all BK BKF are effective}Note that our definition of a
  Breuil--Kisin--Fargues module is less general than the definition
  made in~\cite{2016arXiv160203148B}, in that we require~$\varphi$ to
  take~$\gM$ to itself; this corresponds to only considering Galois
  representations with non-negative Hodge--Tate weights. This
  definition is convenient for us, as it allows us to make direct
  reference to the literature on Breuil--Kisin modules. The
  restriction to non-negative Hodge--Tate weights is harmless in our
  main results, as we can reduce to this case by twisting by a large
  enough power of the cyclotomic character (the interpretation of
  which on Breuil--Kisin--Fargues modules is explained in~\cite[Ex.\
  4.24]{2016arXiv160203148B}).
\end{rem}

\begin{defn}
  \label{defn: BKF module with GK action coefficients}Let $A$ be a $p$-adically complete
$\cO$-algebra which is topologically of finite type.
A Breuil--Kisin--Fargues $G_K$-module with $A$-coefficients is \index{Breuil--Kisin--Fargues module!$G_K$-module}
  a Breuil--Kisin--Fargues module with $A$-coefficients, equipped with a continuous
  semilinear action of~$G_K$ which commutes with~$\varphi$.
\end{defn}
Note that if~$\gMt$ is a Breuil--Kisin--Fargues $G_K$-module with
$A$-coefficients, then $W(\C^\flat)_A\otimes_{\AAinf{A}}\gMt$ is
naturally an \'etale $(\varphi,G_K)$-module in the sense of
Definition~\ref{defn: GK module over Ainf}.  The following definition
is motivated by Theorem~\ref{athm: admits all descents if and only if
  semistable} and Corollary~\ref{acor: admits all descents if and only
  if potentially semistable}.
% \mar{TG: that definition
%   currently requires $A$ to be finite type over $\cO/\varpi^a$, but we
% should presumably extend it?}
%\mar{TG: to here with indexing defns}
\begin{defn}%\mar{TG: need to update this to match the appendix}
  \label{defn: descending BKF to BK coefficients}
% \mar{ME: Does this definition implictly assume that $A$ is topologically of
% finite type over $\cO$?  E.g.\ are we implicitly assuming that
% $\gS_A$ is a subring of $\AAinf{A}$?  Or at least, does the definition 
% only seem natural when this holds?}
Let $A$ be a $p$-adically complete
$\cO$-algebra which is topologically of finite type over~$\cO$ (and recall
then that $\AAinf{A}$ is faithfully flat over $\gS_{\pi^\flat,A}$,
by Proposition~\ref{prop:top f.t. flatness}).
We say that a
%let $\gMt$ be a
  Breuil--Kisin--Fargues $G_K$-module with $A$-coefficients
% Then we say that~$\gMt$
$\gMt$
\emph{descends for~$\pi^\flat$} or \emph{descends
    to~$\gS_{\piflat,A}$} if there is a
%\mar{ME: May have to modify this definition slightly if we can't prove a 
%general faithful flatness statement in the top.\ f.t.\ context.  The rephrasing would be
%that we assume given a projective BK module $\gM_{\pi^\flat}$ and an isomorphism
%$\AAinf{A} \otimes_{\gS_{\pi^\flat,A}} \gM_{\pi^\flat} \iso \gMt$
%which embeds $\gM_{\pi^\flat}$ into $(\gMt)^{G_{K_{\pi^\flat,\infty}}}$.}
  Breuil--Kisin module~$\gM_{\pi^\flat}$ with
  $\gM_{\pi^\flat}\subseteq(\gMt)^{G_{K_{\pi^\flat,\infty}}}$ and
for which the natural map
  $\AAinf{A}\otimes_{\gS_{\pi^\flat,A}}\gM_{\pi^\flat} \to \gMt$ is an isomorphism.

We say
  that~$\gMt$ \emph{admits all descents} if it descends for every
  \index{Breuil--Kisin--Fargues module!admitting all descents}
  choice of~$\pi^\flat$ (for every choice of~$\pi$), and if
  furthermore 

  % Then
  % we say that~$\gMt$ \emph{admits all descents}  if the following
  % conditions hold.
  % \begin{enumerate}
  % \item\label{item: existence of descent Zp version} For every choice of~$\pi$ and~$\piflat$, there is a
  %   Breuil--Kisin module~$\gM_{\pi^\flat}$ of height at most~$h$ with
  %   $\gM_{\pi^\flat}\subset(\gMt)^{G_{K_{\pi^\flat,\infty}}}$ and
  %   $\gMt=\AAinf{A}\otimes_{\gS_{\pi^\flat,A}}\gM_{\pi^\flat}$.
  \begin{enumerate}
  \item\label{item: M mod u descends} The $W(k)\otimes_{\Zp}A$-submodule $\gM_{\piflat}/[\piflat]\gM_{\piflat}$ of
  $W(\overline{k})_A\otimes_{\AAinf{A}}\gMt$ is independent of the
    choice of~$\pi$ and~$\piflat$.
  \item\label{item: M mod E descends} The $\cO_K\otimes_{\Zp}A$-submodule $\varphi^*\gM_{\piflat}/E_{\piflat}\varphi^*\gM_{\piflat}$
    of  $\cO_{\C,A}\otimes_{\AAinf{A},\theta}\varphi^*\gMt$  is independent of the
    choice of~$\pi$ and~$\piflat$.%\mar{TG: have to define $\theta$ above.}
% \mar{ME: TG suggests that we add a compatibity with canonical actions here;
% more precisely, for all $a$ and $N$, we should have compatibility with
% the canonical action for some $s$ suff.\ large depending on $a$ and $N$.}
  \end{enumerate}
  If~$\gMt$ admits all descents, then we say that it is
  \emph{crystalline} if for each~$\piflat$, and each~$g\in G_K$, we
  have \numequation\label{eqn: condition on BK GK for
    crystalline}(g-1)(\gM_{\piflat})\subseteq \varphi
  ^{-1}(\mu)[\piflat]\gMt \end{equation} (where ~$\mu=[\varepsilon]-1$
for some compatible choice of roots of unity
$\varepsilon=(1,\zeta_p,\zeta_{p^{2}},\dots)\in\cO_\C^\flat$).

  
 If~$L/K$ is a finite Galois extension, then we say that~$\gMt$
 \emph{admits all descents over~$L$} if the corresponding
 Breuil--Kisin--Fargues $G_L$-module (obtained by restricting
 the~$G_K$-action to~$G_L$) admits all descents.
\end{defn}

\begin{rem}
  \label{rem: descents are automatically projective}There is
  considerable redundancy and rigidity in Definition~\ref{defn:
    descending BKF to BK coefficients}. Note in particular that if % ~$A$
  % is a finite type $\cO/\varpi^a$ algebra and
  % ~$\gMt$ is a Breuil--Kisin--Fargues module with $A$-coefficients,
  % and
  for some~$\piflat$ there exists a $\gS_{\piflat,A}$-module
  $\gM_{\piflat}$ with
  $\gMt=\AAinf{A}\otimes_{\gS_{\pi^\flat,A}}\gM_{\pi^\flat}$, then
  since~$\gS_{\piflat,A}\to\AAinf{A}$ is faithfully flat,
%(by Proposition~\ref{prop: maps of coefficient rings are faithfully flat injections}),
 the $\gS_{\piflat,A}$-module $\gM_{\piflat}$ is automatically finite projective
  (by~\cite[\href{https://stacks.math.columbia.edu/tag/058S}{Tag
    058S}]{stacks-project}).
\end{rem}
We will find the following basic lemmas useful.
\begin{lem}
  \label{lem: descent for GK action depends only on pi}Fix a choice
  of~$\pi$. Suppose that $\gMt$ is a
Breuil--Kisin--Fargues $G_K$-module with~$A$ coefficients, where~$A$
is a $p$-adically complete
$\cO$-algebra which is topologically of finite type over~$\cO$. Then~$\gMt$ descends for some
choice of~$\pi^\flat$ \emph{(}for our fixed~$\pi$\emph{)} if and only if it descends
for all such choices.
\end{lem}
\begin{proof}% \mar{TG: I am currently confused by this ``accordingly'',
  % but hopefully the result is actually true?! I don't think it would
  % be a disaster if it wasn't, we don't make that much use of it\dots
  % TG: leaving for a moment but I think it's fine, I think it's just
  % transitivity of the Galois group on the roots of an irreducible
  % polynomial, but I'm going to leave it for now.}
  % If~$(\pi^\flat)'$ is another choice of~$\pi^\flat$, then
  % $(\pi^\flat)'/\pi^\flat$ is given by a choice of primitive $p^n$th
  % roots of unity. Accordingly
  We can choose an element~$g\in G_K$ with
  $g(\pi^\flat)=(\pi^\flat)'$, so
  that~$\gS_{(\pi^\flat)',A}=g(\gS_{\pi^\flat,A})$.  Then if~$\gM_\pi$
  is a descent to~$\gS_{\pi,A}$, $g(\gM_\pi)$ is a descent
  to~$\gS_{(\pi^\flat)',A}$.
\end{proof}
% In view of Lemma~\ref{lem: descent for GK action depends only on pi},
% we will sometimes say that $\gMt$ descends for~$\pi$ if it
% descends for~$\pi^\flat$.\mar{TG: do we actually ever say this?}

%The following lemma is reassuring.
%\mar{ME: This should bootstrap to top.\ f.t.\ algebras, by taking a $\varpi$-adic limit.
%But maybe will also apply to them directly, if we get a general faithful flatness
%statement in the top.\ f.t.\ context. ME: I think this comment and the preceding one 
%are dealt with now; if you agree,  feel free to remove these comments. (ME: Actually
%I removed the preceding one; leaving this one for now just until we check
%the implict $p$-adic completion mentioned in the subsequent comment
%doesn't cause trouble.) TG: as written it didn't, but now it
%does. Does this look OK?}
\begin{lem}
  \label{lem: uniqueness of descent for G K infty Ainf}Let $A$ be a $p$-adically complete
$\cO$-algebra which is topologically of finite type over~$\cO$. 
%   finite type 
% $\cO/\varpi^a$-algebra for some~$a\ge 1$.
  Suppose that $\gMt$ is a
  Breuil--Kisin--Fargues  $G_K$-module. Then if $\gMt$ descends for~$\pi^\flat$, the Breuil--Kisin
  module~$\gM_{\pi^\flat}$ is uniquely determined.
\end{lem}% \mar{TG: tensor products should have rings on the left, fix
  % at some point.}
\begin{proof}
Assume to begin with that $A$ is an $\cO/\varpi^a$-algebra of finite type,
for some $a \geq 1$.
  By Proposition~\ref{prop: equivalences of categories to Ainf},
%for each~$a\ge 1$
the $\varphi$-module % \mar{TG: I guess at this point we're implicitly
    % invoking Corollary~\ref{cor:descent with coefficients} for the
    % quasi-inverse functor. Maybe that should be incorporated into the
    % statements of Theorem~\ref{thm:descending projective modules} and
    % Proposition~\ref{prop: equivalences of categories to Ainf}?}
  $(W(\C^\flat)_{A}\otimes_{\AAinf{A}}\gMt)^{G_{K_{\pi^\flat,\infty}}}$ uniquely descends to an
  \'etale $\varphi$-module~$M$ over~$\cO_{\mathcal{E},\pi^\flat,A}$.
%It follows\
%%mar{TG:
%%    hopefully! I didn't think this through carefully\dots TG: actually
%%  it would seem to use that formation of $G_{K_{\pi^\flat,\infty}}$
%%  invariants commutes with going mod $\varpi^a$, which isn't obvious
%%  to me. But do we really need to use~\ref{prop: equivalences of
%%    categories to Ainf} at all? Can we not just directly use the
%%  faithful flatness of $\gS\to\Ainf$? ME: I was bit confused about the role of \ref{prop: equivalences of categories to Ainf} actually.
%%But I'm also a bit confused as to why $\gM_1 + \gM_2$
%%is also a descent. TG: I'm confused too. I think it should be the case
%%that \ref{prop: equivalences of categories to Ainf} can be used to
%%show that the descent $\gM_{\piflat}\otimes\OEA$ is unique. Then
%%$\gM_1,\gM_2$ are both lattices in here, so I think their span should
%%(??) be another finitely generated module whose base extension
%%to~$\AAinf{A}$ is $\gMt$? But I am basically in a state of confusion
%%about how this argument is supposed to go.}
%that  $(\gMt\otimes_{\AAinf{A}}W(\C^\flat)_{A})^{G_{K_{\pi^\flat,\infty}}}$ uniquely descends to an
%  \'etale $\varphi$-module $M=\varprojlim_a M_a$ over~$\cO_{\mathcal{E},\pi^\flat,A}$.
Thus
if~$\gM_1$, $\gM_2$ are two descents of~$\gMt$
  to~$\gS_{\piflat,A}$, then $\cO_{\mathcal{E},\pi^\flat,A}\otimes_{\gS_{\piflat,A}}\gM_1$
and  $\cO_{\mathcal{E},\pi^\flat,A}\otimes_{\gS_{\piflat,A}}\gM_2$ coincide; they are both
equal to~$M$.
%\mar{ME: There should be $p$-adic completions here, right? TG: how
%  about this? ME: Since inverting $u$ (which is what required $p$-adic completions)
%is gone for now, I'm removing this comment.}

This allows us to show that 
  $\gM_1+\gM_2$ is also a descent of~$\gMt$.  For this, we have to show that
the natural map $ \AAinf{A}\otimes_{\gS_{\piflat,A}}(\gM_1  + \gM_2) \to \gMt$
is an isomorphism. 
(Note that it will then follow automatically that $\gM_1+\gM_2$
is projective over $\gS_{\piflat,A}$,
  by Remark~\ref{rem: descents are automatically projective}.)
That it is a surjection is clear, since it is already a surjection
when restricted to each summand.  To see that it is an injection, we may check
that its composite with
the injection $\gMt \hookrightarrow  W(\C^\flat)_A\otimes_{\AAinf{A}}\gMt$ is
injective.  But this composite may be factored as the composite of the embedding
$$\AAinf{A}\otimes_{\gS_{\piflat,A}}(\gM_1 + \gM_2) \hookrightarrow
 \AAinf{A} \otimes_{\gS_{\piflat,A}}M $$
(obtained by tensoring the embedding $\gM_1 + \gM_2 \hookrightarrow M$ with
the extension $\AAinf{A}$ of $\gS_{\piflat, A}$, which is flat by
Proposition~\ref{prop:top f.t. flatness}),
and the isomorphism %\mar{TG: ME to check if it's an iso, add an argument}
\begin{multline*}
 \AAinf{A} \otimes_{\gS_{\piflat,A}}M\cong
 \bigl( \AAinf{A}\otimes_{\gS_{\pi^{\flat},A}} \cO_{\mathcal{E},\pi^{\flat},A} 
  \bigr)
\otimes_{\cO_{\mathcal{E},\pi^\flat,A}} M
\\
\cong  W(\C^{\flat})_A  \otimes_{\cO_{\mathcal{E},\piflat,A}}M
\cong W(\C^\flat)_A \otimes_{\AAinf{A}}\gMt.
\end{multline*}
%(note that it is automatically
%  projective by Remark~\ref{rem: descents are automatically projective}),
%  and
  % since~$\gS_{\piflat,A}\to\AAinf{A}$ is faithfully flat,
  % $\gM_1+\gM_2$ is finite projective
  % by~\cite[\href{https://stacks.math.columbia.edu/tag/058S}{Tag
  %   058S}]{stacks-project}.\mar{TG: move this projectivity up to where
  %   we define things, noting that we don't need to check projectivity
  %   if we have a descent} R
  %and
Since $\gM_1 + \gM_2$ is also a descent,
  %The result of the preceding paragraph allows us to
  we may replace~$\gM_2$ by~$\gM_1+\gM_2$, and 
  therefore assume that
  $\gM_1\subseteq\gM_2$.
  Since~$\AAinf{A}\otimes_{\gS_{\piflat,A}}(\gM_2/\gM_1)=0$, and
  as ~$\gS_{\piflat,A}\to\AAinf{A}$ is faithfully flat
  (again by Proposition~\ref{prop:top f.t. flatness}),
  we see that~$\gM_2/\gM_1=0$, as required.

Suppose now that $A$ is $p$-adically complete and topologically of finite type.
Let $\gM_1$ and $\gM_2$ be two descents of $\gMt$.  Since $\gS_A \to \AAinf{A}$
is faithfully flat
(again by Proposition~\ref{prop:top f.t. flatness}),
we find that $\gM_1/\varpi^a$ and $\gM_1/\varpi^a$ both embed into
$\gMt/\varpi^a$, for any~$a~\geq~1$, and both provide descents for $\piflat$ of $\gMt/\varpi^a$.
Applying the uniqueness result we've already proved (our coefficients now being
$A/\varpi^a$), we find that $\gM_1/\varpi^a$ and $\gM_2/\varpi^a$ coincide
as submodules of $\gMt/\varpi^a$.  Passing to the inverse limit over~$a$,
we find that $\gM_1$ and $\gM_2$ coincide as submodules of~$\gMt$,
proving the desired uniqueness in general.
\end{proof}


% \begin{proof}The case $A=\Zpbar$ follows from the case that $A$ is a
%   finite flat $\cO$-algebra by taking direct limits, so we assume that~$A$
%   is a finite flat $\cO$-algebra for the rest of the proof.  Suppose firstly
%   that~$T_A(M)$ is semistable with Hodge--Tate weights
%   in~$[-h,0]$.\mar{TG: this direction is fine for~$p=2$} Then the
%   existence of~$\gMt$ is immediate from~\cite[Thm.\ 2.3.1]{MR2558890}
%   and~\cite[Thm.\ 2.2.1]{1302.1888}, bearing in mind Remark~\ref{rem: we don't
%     twist our embeddings by phi}.

%   Conversely, suppose that we are given~$\gMt$. Let~$s$ be the
%   greatest integer such that~$K$ contains a primitive $p^s$th root of
%   unity. Then by~\cite[Thm.\ 3]{MR3127808}, for each~$\pi$,
%   $T_A(M)|_{G_{K_{\pi^\flat,s}}}$ is semistable (necessarily with
%   Hodge--Tate weights in~$[-h,0]$, by the results
%   of~\cite{KisinCrys}). We claim that this implies that~$T_A(M)$
%   itself is semistable. It is enough to check
%   that~$T_A(M)|_{G_{\Knr}}$ is semistable, where~$\Knr/K$ is the
%   maximal unramified extension.

%   We claim that it in turn suffices  to show that as we run
%   over all~$\pi$, we have \[\cap_\pi \Knr(\pi^{1/p^s})=\Knr\](note that by
%   the definition of~$s$, $K_{\pi^\flat,s}=K(\pi^{1/p^s})$ is
%   independent of the choice of~$\pi^{1/p^s}$,
%   so~$\Knr(\pi^{1/p^s})=\Knr\cdot K_{\pi^\flat,s}$). To see this, note
%   that for any~$L/K$, $T_A(M)|_{G_L}$ is semistable if and only if the
%   restriction to the inertia group~$I_L$ of the Weil--Deligne
%   representation associated to~$T_A(M)$ is trivial, so it is enough to
%   show that the inertia groups~$I_{\Knr(\pi^{1/p^s})}$ generate the
%   inertia group~$I_{\Knr}$. By the definition of~$\Knr$, these inertia
%   groups are equal to the corresponding absolute Galois groups. If
%   the $G_{\Knr(\pi^{1/p^s})}$ are all contained in a proper subgroup
%   of~$G_{\Knr}$, then the fixed field of this subgroup would be
%   contained in each~$\Knr(\pi^{1/p^s})$, as required.

%   Suppose then that this intersection strictly contains~$\Knr$; then
%   the degree~$p$ extension~$\Knr(\pi^{1/p})/\Knr$ must be independent
%   of~$\pi$, so by Kummer theory, all choices of~$\pi$ must differ by
%   elements of~$((\Knr)^\times)^p$. This implies that all units
%   in~$K^\times$ are $p$th powers of elements of~$(\Knr)^\times$, a
%   contradiction (for example, consider the unit~$1+\pi$).
% \end{proof}



%\begin{remark}
%	% \mar{ME: Note that projective Breuil--Kisin modules over $\gS_A$ haven't 
%	% 	actually been defined.}
%  If~$\gM$ is a projective
%  Breuil--Kisin module over~$\gS_A$, of height at most~$h$ (for some
%  $h \geq 0$),
%  then
%$\gMt:=\AAinf{A}\otimes_{\gS_A}\gM$
%  is a projective
%  Breuil--Kisin module over~$\AAinf{A}$, of height at most~$h$.
%Since~$\gM$ is a
%flat~$\gS_A$-module, we may regard~$\gM$ as a
%$\gS_A$-submodule of~$\gMt$.
%\end{remark}

\section{Canonical extensions of $G_{K_{\infty}}$-actions}
\label{subsec:canonical actions}
% \mar{ME: This subsection might have to  be moved up, so that
% its results can be incorporated into the definition of ``BKF module
% admitting all descents''.}
% \mar{ME: There are no finiteness assumptions being imposed
% on the coefficients $A$ appearing
% in this section, but we should check carefully that we're not implicitly using
% something like the fact that $\gS_A \to \AAinf{A}$ is an embedding with various
% properties, etc., which only holds when $A$ is of finite type over $\cO/\varpi^a$.}


In this section we will consider a fixed choice of~$\pi^\flat$, which
we will accordingly drop from the notation. For each~$s\ge 1$,
write~$K_s=K(\pi^{1/p^s})$. We now use a variant of the arguments of
~\cite{MR2745530}, which show that the trivial action
of~$G_{K_\infty}$ on a Breuil--Kisin module with
$\cO/\varpi^a$-coefficients can be extended to some~$G_{K_s}$. We
begin with some preliminary lemmas, the first of which is an analogue
for Breuil--Kisin modules of~\cite[Lem.\
5.2.14]{EGstacktheoreticimages}, and is proved in a similar way.

\begin{lem}
  \label{lem: adding projective Kisin to get free}Let $A$ be a $p$-adically complete
$\cO$-algebra which is topologically of finite type, and let~$\gM$ be a finite projective Breuil--Kisin module with $A$-coefficients of height at
  most~$h$. Then~$\gM$ is a direct summand of a finite free Breuil--Kisin
  module with $A$-coefficients of height at
  most~$h$.
\end{lem}
\begin{proof}
  Since the map $\Phi_{\gM}:\varphi^*\gM\to\gM$ is an injection with
  cokernel killed by~$E^h$, there is a
  map~$\Psi_{\gM}:\gM\to\varphi^*\gM$ with $\Psi_{\gM}\circ\Phi_{\gM}$
  and $\Phi_{\gM}\circ\Psi_{\gM}$ both being given by multiplication
  by~$E^h$. Let~$\gP$ be a finite projective $\gS_A$-module
  such that~$\gF:=\gM\oplus\gP$ is a finite free~$\gS_A$
  module. Set~$\gQ:=\gP\oplus\gM\oplus\gP$, a finite
  projective~$\gS_A$-module.

  Note that~$\gM\oplus\gQ\cong\gF\oplus\gF$ is a finite
  free~$\gS_A$-module, so it suffices to show that~$\gQ$ can
  be endowed with the structure of a Breuil--Kisin module of height at
  most~$h$. Since~$\gF$ is free, we can choose an
  isomorphism~$\Phi_{\gF}:\varphi^*\gF\isoto\gF$, and we then
  define~$\Phi_{\gQ}:\varphi^*\gQ\to\gQ$ as the composite
  \begin{align*}
    \varphi^*\gQ
    &=\varphi^*\gP\oplus\varphi^*\gM\oplus\varphi^*\gP\\&\isoto
    \varphi^*\gM\oplus\varphi^*\gP\oplus\varphi^*\gP\\&=\varphi^*\gF\oplus\varphi^*\gP\\&\stackrel{\Phi_{\gF}}{\to}\gF\oplus\varphi^*\gP\\&\isoto
    \gP\oplus\gM\oplus\varphi^*\gP\\&\stackrel{\Psi_{\gM}}{\to} \gP\oplus\varphi^*\gM\oplus\varphi^*\gP\\&= \gP\oplus\varphi^*\gF\\&\stackrel{\Phi_{\gF}}{\to}\gP\oplus\gF=\gQ.
  \end{align*}Every morphism in this composite other than~$\Psi_{\gM}$
  is an isomorphism. Since
  $\Psi_{\gM}\circ\Phi_{\gM}=\Phi_{\gM}\circ\Psi_{\gM}=E^h$, we see
  that~$\Psi_{\gM}$ is an injection with cokernel killed
  by~$E^h$, so the same is true of~$\Phi_{\gQ}$, as required.
\end{proof}


The following lemma is proved by a standard Frobenius amplification
argument, which is in particular almost identical to the proof
of~\cite[Lem.\ 4.1.9]{CEGSKisinwithdd}.
\begin{lem}
  \label{lem: Frobenius amplification lemma over Ainf}Let~$\gMt$, $\gNt$
  be projective Breuil--Kisin--Fargues modules  of height at
  most~$h$, where $A$ is a finite type
  $\cO/\varpi^a$-algebra. Suppose that~$N\ge e(a+h)/(p-1)$. Let \[f:\gMt\to
    \gNt\] be an $\AAinf{A}$-linear map,
  and suppose that for all~$m\in\gMt$,
  $(f\circ\Phi_{\gMt}-\Phi_{\gNt}\circ\varphi^*f)(m)\in
  u^{pN}\gNt$.%  \emph{(}where we
  % write~$\Phi_{\gMt},\Phi_{\gNt}$ for their $\AAinf{A}$-linear
  % extensions\emph{)}.
  % \mar{TG:
    % sort out precise value for $N$}

  Then there is a unique $\AAinf{A}$-linear, $\varphi$-linear morphism
  % \mar{ME: Probably the correct adjective for $f$ 
  %         is just $\varphi$-linear, rather than semi-linear, right?}
  $f':\gMt\to
    \gNt$ with the property that for all~$m\in\gMt$,
  $(f'-f)(m)\in u^N\gNt$; in fact, we have $(f'-f)(m)\in u^{N+1}\gNt$.
\end{lem}
\begin{proof}We firstly prove uniqueness. Suppose that~$f''$ also
  satisfies the properties that~$f'$ does, and write~$g=f'-f''$, so that by
  assumption we have
  that~$\im g\subseteq u^N \gNt$ and
  $g\circ\Phi_{\gMt}=\Phi_{\gNt}\circ\varphi^*g$. We claim that this implies
  that~$\im g\subseteq u^{N+1} \gNt$; if
  this is the case, then by induction on~$N$ we have
  $\im g\subseteq u^N \gNt$ for all~$N$,
  and so~$g=0$ (note that as~$\gNt$ is a finite
  projective~$\AAinf{A}$-module,
  %$\gNt$
  it is $u$-adically complete and separated).

  To prove the claim, note that
  since~$\im g\subseteq u^N \gNt$ we
  have~$\im \varphi^*g\subseteq u^{pN}
  \gNt$. By~\cite[Lem.\
  5.2.6]{EGstacktheoreticimages} (and its proof), and the assumption
  that~$\gMt$ has height at most~$h$, we
  have~$\im \Phi_{\gMt}\supseteq u^{e(a+h-1)}\gMt$.  Since
  $g\circ\Phi_{\gMt}=\Phi_{\gNt}\circ\varphi^*g$, it follows that
  \[u^{e(a+h-1)}\im g\subseteq\im g\circ\Phi_{\gMt}= \im \Phi_{\gNt}\circ\varphi^*g
    \subseteq u^{pN} \gNt\]and since
  $\gNt$ is $u$-torsion free, we have
  $\im g\subseteq u^{pN-e(a+h-1)}
  \gNt$. Our assumption on~$N$ implies
  that $pN-e(a+h-1)\ge N+1$, as required.

  We now prove the existence of~$f'$. For any~$\AAinf{A}$-linear
  map $h: \gMt\to
  \gNt$ we
  set~$\delta(h):=\Phi_{\gNt}\circ\varphi^*h-h\circ\Phi_{\gMt}:\varphi^*\gMt\to\gNt$. We
  claim that for
  any~$s:\varphi^*\gMt\to
  u^{pN}\gNt$, we can find~$t:\gMt\to
  u^{N+1}\gNt$ with~$\delta(t)=s$. Given this claim, the existence of~$f'$ is immediate,
taking~$s=\delta(f)$ and~$f':=f-t$, so that~$\delta(f')=0$.


To prove the claim, we first prove the weaker claim that for
any~$s$ as above we can
find~$h:\gMt\to
u^{N+1}\gNt$
with~$\im(\delta(h)-s)\subseteq u^{p(N+1)}\gNt$. % \mar{TG: getting confused
  % about precise numbers here!}
Admitting this second claim, we prove
the first claim by successive approximation: we set~$t_0=h$, so
that~$\im(\delta(t_0)-s)\subseteq u^{p(N+1)}\gNt$. Applying the second
claim again with~$N$ replaced by~$N+1$, and $s$ replaced
by~$s-\delta(t_0)$, we
find~$h:\gMt\to
u^{N+2}\gNt$
with~$\im(\delta(h)-(s-\delta(t_0)))\subseteq u^{p(N+2)}\gNt$. % \mar{TG:
  % confused again, sort out exponents!}.
Setting~$t_1=t_0+h$, and
proceeding inductively, we obtain a Cauchy sequence converging (in
the~$u$-adically complete finite $\AAinf{A}$-module
$\Hom_{\AAinf{A}}(\gMt,\gNt)$)
to the required~$t$.

Finally, we prove this second
claim.
If~$h:\gMt\to
u^{N+1}\gNt$ then 
% \mar{ME: I agree that mapping into $u^{N+1}\gNt$ is even better than
% 	mapping into $u^N\gNt$, but doesn't this mean that all the
% 	exponents above can be improved by $1$? TG: possibly; but I
%         got pretty confused about exponents at various points when
%         writing it, so I don't want to just change one here. If you've
%       checked it all and can see what they should all be then feel fee
%     to improve! Otherwise I would be happy enough to just change the
%     $N+1$ to $N$ here, to be honest, as the actual bounds are unimportant.}
$\im\varphi^*h\subseteq u^{p(N+1)}\gNt$, so it suffices to show that we can
find~$h$ such that $h\circ\Phi_{\gMt}=-s$; but this is immediate,
because the cokernel of~$\Phi_{\gMt}$ is killed by~$u^{e(a+h-1)}$, and
$pN-e(a+h-1)\ge N+1$.% \mar{TG: should be easy to get exponents
  % straight; come back to this.}
\end{proof}

% \begin{remark}
% 	% \mar{ME: Note that projective Breuil--Kisin modules over $\gS_A$ haven't 
% 	% 	actually been defined.}
%   If~$\gM$ is a projective
%   Breuil--Kisin module with $A$-coefficients, of height at most~$h$ (for some
%   $h \geq 0$),
%   then
% $\gMt:=\AAinf{A}\otimes_{\gS_A}\gM$
%   is a projective
%   Breuil--Kisin--Fargues module with $A$-coefficients, of height at most~$h$.
% Since~$\gM$ is a
% flat~$\gS_A$-module, we may regard~$\gM$ as a
% $\gS_A$-submodule of~$\gMt$.
% \end{remark}

The following lemma is proved by a reinterpretation of some of the arguments
of~\cite[\S2]{MR2745530}.  
%If~$\gM$ is a projective Kisin module
%over~$\gS_A$, we write
%$\gMt:=\AAinf{A}\otimes_{\gS_A}\gM$. Since~$\gM$ is a
%flat~$\gS_A$-module, we may regard~$\gM$ as a
%$\gS_A$-submodule of~$\gMt$.

\begin{lem}\label{lem: Caruso Liu Galois action on Kisin}
  For any fixed $a,h$, and any $N\ge e(a+h)/(p-1)$, there is a
  positive integer~$s(a,h,N)$ with the property that for any finite type
  $\cO/\varpi^a$-algebra~$A$, any
  projective Breuil--Kisin
  module~$\gM$  of height at most~$h$, and 
  any~$s\ge s(a,h,N)$, there is a unique continuous action of~ $G_{K_s}$ on
  $\gMt:=\AAinf{A}\otimes_{\gS_{A}}\gM$ which commutes with~$\varphi$
  and is semi-linear with respect to the natural action of~$G_{K_s}$
  on~$\AAinf{A},$ with the additional property that for
  all~$g\in G_{K_s}$ we have $(g-1)(\gM)\subset
  u^N\gMt$. % \mar{TG: should probably
    % explain how we are thinking of~$\gM$ as living inside $\AAinf{A}\otimes_{\gS_A}\gM$}
\end{lem}
\begin{proof}By
  definition, to give a semi-linear action of~$G_{K_s}$  on~$\gMt$ is
  to give for each~$g\in G_{K_s}$ a morphism of Breuil--Kisin--Fargues
  modules \[\beta_g:g^*\gMt\to\gMt\] with the property that for
  all~$g,h\in G_{K_s}$, we have $\beta_{gh}=\beta_h\circ
  h^*\beta_g$. % \mar{TG: I think this is the right way round, but it
    % would be sensible to check for nonsense.}
  Now, the requirement that
  $(g-1)(\gM)\subseteq u^N\gMt$ implies that if we have two such
  morphisms~$\beta_g,\beta'_g$ then
  $(\beta_g-\beta'_g)(g^*\gMt)\subseteq u^N\gMt$, so
  that~$\beta'_g=\beta_g$ by the uniqueness assertion of
  Lemma~\ref{lem: Frobenius amplification lemma over Ainf}.

  We now use the existence part of Lemma~\ref{lem: Frobenius
    amplification lemma over Ainf} to construct the required
  action. Since the homomorphism $\gS_A\into\AAinf{A}$ takes~$u$ to the
  Teichm\"uller lift of~$(\pi^{1/p^n})_n$, % \mar{TG: have we actually
%     discussed this above? I suspect that we use the injectivity of
%     various coefficient rings into $W(\C^\flat)_A$ throughout the
%     paper - do we prove it? ME: The faithful flatness result
% in \S 3 on Galois descent gives injectivy for perfect rings.  Hopefully
% there is an argument that shows that the non-perfect rings embed into
% (rather than just map to) the perfect ones.}
we can
  and do choose~$s(a,h,N)$ sufficiently large that for all~$s\ge
  s(a,h,N)$ and~$g\in
  G_{K_s}$, we have~$g(u)-u\in u^{pN}\AAinf{A}$; thus for
  all~$\lambda\in \gS_A$, we have $g(\lambda)-\lambda\in
  u^{pN}\AAinf{A}$. % \mar{TG: need to choose~$N$
    % and of course need to justify that~$s$ exists, but I think this
    % should be just the continuity of the action of $G_K$ on~$\Ainf$.}

We claim that for each~$g\in G_{K_s}$, we may define
  an~$\AAinf{A}$-linear map  \[\alpha_g:g^*\gMt\to
    \gMt\] with the following two properties:
\begin{itemize}
	\item
$(\alpha_g\circ\Phi_{g^*\gMt}-\Phi_{\gMt}\circ\varphi^*\alpha_g)(\varphi^*g^*\gMt)\subseteq u^{pN}\gMt.$
\item
For each $m \in \gM,$ we have $\alpha_g(1\otimes m) - m \in u^N \gMt.$
\end{itemize}
It suffices to construct such maps in the case when~$\gM$ is free;
indeed, by Lemma~\ref{lem: adding projective Kisin to get free},
%we can and do assume from now on that~$\gM$ is free; indeed, in general
  we may write~$\gF=\gM\oplus\gP$ where~$\gF,\gP$ are respectively free and
  projective Breuil--Kisin modules of height at most~$h$, and given a
  morphism~$\alpha_g:g^*\gF^{\inf}\to\gF^{\inf}$ satisfying our
  requirements (where we write $\gF^{\inf}=\AAinf{A}\otimes_{\gS_{A}}\gF$),
  we may obtain the desired
  morphism~$\alpha_g:g^*\gMt\to\gMt$ by projection onto the
  corresponding direct summands.
  % \mar{ME: Do we also have to see that this construction preserves the
  %         cocyle property?  Or is the cocycle property a consequence
  %         of the uniquess?  If so, we should note this above! TG: see
  %         next marginal comment.}

If then~$\gF$ is a free Breuil--Kisin module of height at most~$h$, we choose a
basis~$f_1,\dots,f_d$ for~$\gF$, and then, for each~$g\in G_{K_s}$, we define
  an~$\AAinf{A}$-linear map  \[\alpha_g:g^*\gFt\to
    \gFt\] (which depends on our choice of basis) by
  %
  % (Indeed, in general we we may
  % write~$\gF=\gM\oplus\gP$ where~$\gF,\gP$ are respectively free and
  % projective BK modules of height at most~$h$. If we prove that there
  % is a canonical action of~$G_{K_s}$ on
  % $\AAinf{A}\otimes_{\gS_A}\gF[1/u]$, then applying the
  % projector to $\AAinf{A}\otimes_{\gS_A}\gM[1/u]$, we obtain
  % an action of~$G_{K_s}$ on
  % $\AAinf{A}\otimes_{\gS_A}\gM[1/u]$. To see that this action
  % is independent of the choice of~$\gF$ and thus canonical,)\mar{TG:
  %   have to figure out in what sense the action on free modules is
  %   canonical, I guess; hopefully this isn't hard. It's presumably
  %   related to checking that the construction that we make below in
  %   the free case is independent of the choice of basis? In fact I
  %   think it should just characterised by demanding that
  %   for all $g\in G_{K_s}$, we have
  %   $(g-1)(\gM)\subset u^N\AAinf{A}\otimes_{\gS_A}\gM$,
  %   because if we have two different actions that both satisfy this,
  %   then applying their difference will yield a $\varphi$-equivariant
  %   map $g^*\AAinf{A}\otimes_{\gS_A}\gM\to
  %   u^N\AAinf{A}\otimes_{\gS_A}\gM$ which will vanish for the
  %   usual reasons.}
  %
  %
  % Choose a basis~$m_1,\dots,m_d$ for~$\gM$ as an~$\AKplus$-module, and
  % for each element~$g\in G_{K_s}$,
  % define the $\AAinf{A}$-linear map % \mar{TG: should probably justify
  %   % that the LHS is actually a Kisin module of height at most~$h$.}
  % \[\alpha_g:g^*\AAinf{A}\otimes_{\gS_A}\gM\to
  %   \AAinf{A}\otimes_{\gS_A}\gM\]
  %
  \[\alpha_g(\sum_i\lambda_i\otimes f_i)=\sum_i\lambda_if_i.\] We now check
  that~$\alpha_g$ has the claimed properties.
  If~$\Phi_{\gF}(1\otimes f_i)=\sum_j\theta_{i,j}f_j$ (so
  that~$\theta_{i,j}\in\gS_A$), then 
  for any $\sum_i\lambda_i\otimes
    f_i\in\varphi^*g^*\gFt, $ we compute that % \mar{TG: Dat Pham said
      % that on the RHS there should be $\varphi(\lambda_i)$ rather
      % than~$\lambda_i$, but I don't see why -- the calculation still
      % looks correct to me. What do you think? ME: Looks ok to me. TG:
      % commenting out.}
    \[(\alpha_g\circ\Phi_{g^*\gFt}-\Phi_{\gFt}\circ\varphi^*\alpha_g)(\sum_i\lambda_i\otimes
    f_i)=\sum_{i,j}\lambda_i(g(\theta_{i,j})-\theta_{i,j})f_j\in
    u^{pN}\gFt,\] 
  while for any $\sum_i \lambda_i f_i \in \gF$
  (note then that each $\lambda_i \in \gS_A$),
  we compute that %\mar{TG: too wide}
\begin{multline*}
    \alpha_g\bigl(1\otimes (\sum_i \lambda_i f_i)\bigr)
    - \sum_i \lambda_i f_i
\\
    = 
				 \alpha_g(\sum_ig(\lambda_i)\otimes
                                 g_i)-\sum_i\lambda_ig_i = 
				 \sum_i(g(\lambda_i)-\lambda_i)g_i
				 \in u^N \gFt;
				 \end{multline*}
  in both computations we have used the fact that, by our choice of~$s$,
  if~$\lambda\in\gS_A$, 
  then~$g(\lambda)-\lambda\in u^{pN}\gFt \subseteq u^N\gFt$.

 Returning to the general case of a projective Breuil--Kisin
 module~$\gM$ of height at most~$h$, it follows from % our
  % assumption on~$s$, and
  Lemma~\ref{lem: Frobenius amplification lemma
    over Ainf} that there is a unique % $\AAinf{A}$-linear,
  % $\varphi$-semi-linear 
morphism of Breuil--Kisin--Fargues modules
  \[\beta_g:g^*\gMt\to
 \gMt\]
  with~$\im(\beta_g-\alpha_g)\subseteq
  u^N\AAinf{A}\otimes_{\gS_A}\gM$. Since we have
  $\alpha_{gh}=\alpha_h\circ h^*\alpha_g$, it follows from this
  uniqueness that~$\beta_{gh}=\beta_h\circ h^*\beta_g$, % \mar{TG: does
    % this answer the previous marginal comment?}
  so we have
  constructed an action of~$G_{K_s}$ on~$\gMt$. It remains to check
  that~$(g-1)(\gM)\subseteq u^N\gMt$; for this, note that,
  for any $m \in \gM$, we have
    $$(g-1)m = \beta_g(1\otimes m)  - m = (\beta_g - \alpha_g)(1\otimes m)
    + \alpha_g(1\otimes m ) - m \in u^N \gMt;
    $$
here we have used the fact $\im(\beta_g-\alpha_g)\subseteq
  u^N\AAinf{A}\otimes_{\gS_A}\gM$, together with the second of the
  properties satisfied by the maps $\alpha_g$.
 % \mar{TG: this
    % looks wrong after all, because as written we don't have ~$\lambda_i\in\AKplus$!}
% \mar{TG: as ever not
    % worrying about the exponents here.} 
%
  % It remains to check that this construction does indeed
  % give~$\gM[1/u]$ the structure of a
  % $(\varphi,G_{K_s})$-module. Firstly, note that the uniqueness
  % of~$\beta_g$ implies that $\beta_{gh}=\beta_h\circ
  % h^*\beta_g$\mar{TG: check} so that the~$\beta_g$ do indeed define an
  % action of~$G_{K_s}$ on $\AAinf{A}\otimes_{\gS_A}\gM$ and
  % thus on $\AAinf{A}\otimes_{\gS_A}\gM[1/u]$.\mar{TG: have to
  % also justify continuity but once we've sorted that out elsewhere it
  % will presumably be formal.}  
\end{proof}


\section{Breuil--Kisin--Fargues $G_K$-modules and canonical
  actions}\label{subsec: BKF action canonical}
% \mar{ME: Similar warning to that for the previous subsection: are there
% implicit finiteness assumptions on $A$ being made?}
%\mar{TG: this section, and in particular the statement and then the
%  notation in the proof of the proposition, are new, so should be
%  carefully checked! ME: Removing this comment, since this has now
%been read by both of us}

In an earlier draft of this book, we claimed to show that for
Breuil--Kisin--Fargues $G_K$-modules admitting descents for
all~$\piflat$, the actions of the~$G_{K_{\piflat,s}}\subseteq G_K$
(for sufficiently large values of~$s$)
are necessarily the canonical actions considered in
Section~\ref{subsec:canonical actions}. We then used this claim to
establish finiteness properties of the stacks of
Breuil--Kisin--Fargues $G_K$-modules that we consider later in this
chapter. We are grateful to Dat Pham for pointing out a mistake in our
argument, and we do not know whether this is in fact
automatic. However, the following result of Caruso--Liu (which
motivated our consideration of canonical actions in the first place)
shows that the Breuil--Kisin--Fargues $G_K$-modules coming from
potentially semistable $G_K$-representations do necessarily have this
property; so we can (and do) impose this as an extra condition in the
definition of our stacks.

% \mar{ME: Right now this citation is circular, since we are in
% 	Section~\ref{subsec: BKF action canonical}!  Not sure
% what the final status of this subsection and citation is
% envisaged to be,
% but maybe citation should be to Subsection~\ref{subsec:canonical
% actions}? TG: yes, fixed, commenting out; this was as a result of
% earlier rearrangements.}

% \mar{TG: in the end I couldn't prove this lemma! The essential issue
%   seems to be that we have basically no a priori control of what
%   $g(M_{\piflat})$ is for $g\in G_{K_s}$; I currently think that all arguments that have
%   been written more or less presume that it agrees with the same thing
% where~$g$ is constructed from the canonical action.}
% \begin{lem}\mar{TG: this is provisional, and intended to fill a gap in
%   the following proposition.  Originally I only proposed to prove this lemma
%   where~$\pi^{1/p^s}$ was the same for the two extensions, but I don't
% seem to need it\dots}Let $A$ be a finite type
% $\cO/\varpi^a$-algebra, and let $\gMt$ be a Breuil--Kisin--Fargues
% $G_K$-module of height at most~$h$ with $A$-coefficients, which admits
% all descents. Consider two choices ~$\pi^\flat,(\pi^\flat)'$
% \emph{(}possibly with~$\pi\ne \pi'$\emph{)}, with the property that
% for some~$s\ge s(a,h,N)$ as in Lemma~{\em \ref{lem: Caruso Liu Galois
%     action on Kisin}}, we have~$K_{\pi^\flat,s}=K_{(\pi^\flat)',s}$.

% Then the two canonical actions of~$G_{K_{\pi^\flat,s}}$ on~$\gMt$
% obtained from Lemma~{\em
%   \ref{lem: Caruso Liu Galois action on Kisin}} (one
% from~$\gM_{\piflat}$, and the other from~$\gM_{(\piflat)'}$) agree.  
% \end{lem}
% \begin{proof}Fix some~$g\in
%   G_{K_{\pi^\flat,s}}=G_{K_{(\pi^\flat)',s}}$. Write~$\alpha_g,\beta_g$
%   for the morphisms $g^*\gMt\to\gMt$ constructed in the proof of Lemma~{\em \ref{lem: Caruso Liu Galois
%     action on Kisin}} for~$\pi^{\flat}$, and $\alpha'_g,\beta'_g$ for
% the corresponding morphisms for~$(\pi^\flat)'$. We need to show that
% $\beta_g=\beta'_g$; by construction, it suffices to show that
% $\im(\alpha'_g-\alpha_g)\subseteq
% [\pi^\flat]^N\gMt=[(\pi^\flat)']^N\gMt$. 

% As in the proof of Lemma~\ref{lem: Caruso Liu Galois action on Kisin},
% let~$f_1,\dots ,f_d$  and $f'_1,\dots, f'_d$ be bases
% for~$\gM_{[\pi^\flat]}$ and~$\gM_{[(\pi^\flat)']}$
% respectively. \mar{TG: for now I am ignoring the issue that as written
% the reduction to the free case in Lemma~\ref{lem: Caruso Liu Galois
%   action on Kisin} depends on the choice of~$\pi^\flat$; I don't think
% this can possibly be an issue.} 
% \end{proof}
% \mar{TG: I don't see how to fix this proof! Currently I think we
%   should replace it with a definition. ME: I agree.  TG
% suggests we should incorporate  its statement (quantified over all
% $a$ and $N$)  into the definition of BKF module admitting all descents.}
% \begin{prop}
%   \label{prop: BKF has canonical actions}Let $A$ be a finite type
%   $\cO/\varpi^a$-algebra, and let $\gMt$ be a Breuil--Kisin--Fargues
%   $G_K$-module of height at most~$h$ with $A$-coefficients, which
%   admits all descents. Then for each choice of~$\pi^\flat$, and
%   each~$s\ge s(a,h,N)$ as in Lemma~{\em \ref{lem: Caruso Liu Galois action
%     on Kisin}}, the action of~$G_{K_{\pi^\flat,s}}$ agrees with the
%   canonical action obtained from~$\gM_{\piflat}$ in Lemma~{\em \ref{lem:
%     Caruso Liu Galois action on Kisin}}.  % \mar{TG: am assuming
%     % throughout that there is no primitive $s(a,h,N)$th root of unity
%     % in~$K$, which obviously we can insist on, but I suspect is
%     % probably automatic from the construction of that integer anyway.}
% \end{prop}
% \begin{proof}
%   We note firstly that for each~$\pi^\flat$ this holds for~$s\gg 0$
%   (possibly depending on~$\pi^\flat$). To see this, let~$m_1,\dots
%   m_t$ be a set of generators for~$\gM_{\piflat}$ as a $\gS_{\piflat,A}$-module. By the continuity of
%   the action of~$G_K$ on~$\gMt$, there is an open subgroup~$G_L\subseteq G_K$ such
%     that~$g\in G_L$ implies $(g-1)m_i\in [\piflat]^N\gMt$ for each~$i$. Since
%     this property also holds for each~$g\in G_{K_\infty}$, it holds
%     for each element of the group generated by~$G_L$
%     and~$G_{K_\infty}$, and in particular for some~$G_{K_s}$. After
%     possibly increasing~$s$, we may suppose that we have
%     $(g-1)(\gS_{\piflat,A})\subseteq [\piflat]^N\gS_{\piflat,A}$ for each $g\in G_{K_s}$. It
%     follows that  $(g-1)(\gM_{\piflat})\subseteq [\piflat]^N\gMt$ for all~$g\in G_{K_s}$,
%     so the action is canonical by Lemma~\ref{lem: Caruso Liu Galois action on Kisin}.
% %    take
%   % any~$h\in G_{K_s}$, and note that~$h^{p^t}\to 1$\mar{TG: not true as
%   % written, mean something more like it's tending to the trivial map on~$\gM$}
%   % as~$t\to\infty$. Since~$G_{K_{\pi^\flat,\infty}}$ acts trivially
%   % on~$\gM_{\piflat}$, \mar{TG: maybe aren't even explicitly using that?!} we certainly have~$(h^t-1)(\gM)\subset
%   % u^N\gM^{\inf}$ for $t$ sufficiently large.\mar{TG: need to now
%   %   explain why this is enough; at least if~$p>2$ it should just be a
%   %   matter of explaining $\tau$, if I'm not blundering. Is it
%   %   obvious/true that $G_{K_s}^p=G_{K_{s+1}}$ for $s$ sufficiently
%   %   large? I guess it would follow from that and topological finite
%   %   generation, although presumably this is overkill!}

%   We claim that in fact we can choose~$s\gg 0$ depending only
%   on~$\pi$, and not on~$\pi^\flat$. Indeed, this follows as in the
%   proof of Lemma~\ref{lem: descent for GK action depends only on pi}:
% any other choice of~$\pi^\flat$ is of the form~$g(\pi^\flat)$ for
% some~$g\in G_K$, and if $(h-1)(\gM_{\pi^\flat})\subseteq [\piflat]^N\gMt$ for
% all $h\in G_{K_{\pi^\flat,s}}$, then since $G_{K_{g(\pi^\flat),s}}=g
% G_{K_{\pi^\flat,s}}g^{-1}$ and
% $(ghg^{-1}-1)(g(\gM_{\pi^\flat}))\subseteq g([\piflat])^N\gMt$, we see that the
% corresponding condition holds for $\gM_{g(\piflat)}=g(\gM_{\piflat})$.

% Now fix some~$\pi$ and~$\pi^{1/p^s}$, and write
% $K_{\piflat,s}=K(\pi^{1/p^s})$. Assume that  $s\ge s(a,h,N)$, that~$s$ is
% large enough that~$K$ does not contain a primitive~$p^{s+1}$st root of
% unity, and that the
% action of each~$G_{K_{\pi^\flat,s+1}}$ is the canonical one. We claim that 
% the action of each~$G_{K_{\pi^\flat},s}$ is also canonical. It is enough
% to show that the various groups~$G_{K_{\pi^\flat,s+1}}$ (as~$\piflat$
% varies over the various choices with our fixed choice of~$\pi^{1/p^s}$)
% generate~$G_{K_{\pi^\flat},s}$. Since each~$G_{K_{\pi^\flat,s+1}}$ is
% an index~$p$ subgroup of~$G_{K_{\pi^\flat},s}$, and the
% extensions~$K_{\pi^\flat,s+1}/K_{\pi^\flat,s}$ are not all equal
% (because ~$K_{\pi^\flat,s}$ does not contain a primitive~$p^{s+1}$st
% root of unity), the claim follows.

% We have now shown that for all~$\pi$, the action
% of~$G_{K_{\pi^\flat},s}$ is canonical, provided that~$s\ge s(a,h,N)$
% and~$K$ does not contain a primitive~$p^{s+1}$st root of unity. To
% complete the proof, it is enough to show that if~$s\ge s(a,h,N)$, $K$
% contains a primitive~$p^{s+1}$st root of unity, and the action of each
% $G_{K_{\pi^\flat},s+1}$ is canonical, then the action of
% each~$G_{K_{\pi^\flat},s}$ is also canonical. In this case the
% extensions~$K_{\pi^\flat,s}/K$ and~$K_{\pi^\flat,s+1}/K$ depend only
% on~$\pi$, so we temporarily write them as~$K_{\pi,s}$,
% $K_{\pi,s+1}$. Let~$t\ge 0$ be the maximal integer such that~$1+\pi$
% is a~$p^t$th power in~$K_{\pi,s}$. % \mar{TG: does $t=0$, in fact? This
%   % wasn't clear to me, but I could be missing something obvious!}.
% Then
% we have~$K_{\pi(1+\pi)^{p^{s-t}},s}=K_{\pi,s}$, so arguing as in the
% previous paragraph, it is enough to show
% that~$K_{\pi(1+\pi)^{p^{s-t}},s+1}\ne K_{\pi,s+1}$. By Kummer theory,
% this is equivalent to showing that~$(1+\pi)^{1/p^{t+1}}$ is not a
% $p$th power in~$K_{\pi,s}$, which is true by the definition of~$t$.
% \end{proof}

% %\mar{ME: Can we move this result (4.2.6) to the end of \S 4.2, or the start
% %	of \S 4.3?  As far as I can tell, the results in \S 4.2 
% %	that comes after it (4.2.7, 4.2.8, and the subsequent
% %	discussion) don't use 4.2.6, but rather just 4.2.5,
% %       	so I think the logical flow would 
% %	be clearer if they followed directly after 4.2.5. ME: Made this
% % move.}
% \mar{TG: material that follows is uncommented out now. Might want to
%   look at the various commented out citations of it!}
The following lemma is essentially the content
of~\cite[\S3]{MR2745530} (although they are concerned with getting
precise bounds on the size of~$s$, which is unimportant for
us). Recall that, by Theorem~\ref{athm: admits all descents if and only
  if semistable},
%for each integer~$h\ge 0$,
to any $\cO$-lattice in a
semistable~$E$-representation of~$G_K$ with Hodge--Tate weights
in~$[0,h]$
(for some integer~$h\ge 0$),
there corresponds a
Breuil--Kisin--Fargues $G_K$-module~$\gMt$ of height at most~$h$ which
admits all descents.

% We
% put ourselves in the following situation:  let~$B$ be a finite flat
% $\cO$-algebra, and suppose that~$M_B$ is an \'etale
% $(\varphi,\Gamma)$-module with $B$-coefficients, with the property
% that the $E$-representation of~$G_K$ obtained from~$M_B$ is semistable with
% Hodge--Tate weights contained in the interval~$[-h,0]$. For each~$a\ge
% 1$, write~$M_{B,\infty}^a$ for the \'etale $\varphi$-module
% corresponding to the composite $\Spec (B/\varpi^a) \to\cX_d^a\to\cR_{\BK,d}^a$.
\begin{prop}
 \label{prop: Caruso Liu canonical action semistable}For any
 fixed~$K,a,h$ and~$N\ge e(a+h)/(p-1)$ there is a
 constant~$s'(K,a,h,N)\ge s(a,h,N)$ {\em (}where~$s(a,h,N)$ is as in
 Lemma~{\em \ref{lem: Caruso Liu Galois action on Kisin})} with the
 following property: for any Breuil--Kisin--Fargues $G_K$-module~$\gMt$
 corresponding as above to a semistable
$G_K$ representation with Hodge--Tate weights in~$[0,h]$, and for any
choice of~$\piflat$ with corresponding descent~$\gM_{\piflat}$, 
 if~$s\ge s'(K,a,h,N)$, then %  there is a Breuil--Kisin module~$\gM$
 % \mar{ME: Projective, height $\leq h$?}
 % which is a lattice in ~$M_{B,\infty}^a$, for which
 the restriction to~$G_{K_{\piflat,s}}$ of the action of~$G_K$ 
 on $\gMt\otimes_{\cO}\cO/\varpi^a$ % ~\[W(\C^\flat)_{\cO/\varpi^a}\otimes_{\gS_{\cO/\varpi^a}}\gM_{\piflat}\subset
 %   W(\C^\flat)_{\cO/\varpi^a}\otimes_{\A_{K,\cO}}\gMt\]\emph{(}given by the
 % restriction of the natural action of~$G_K$ on the
 % $(\varphi,G_K)$-module
 % $W(\C^\flat)_{\cO/\varpi^a}\otimes_{\A_{}}M_B$\emph{)}
 agrees with the
 action obtained from~$\gM_{\piflat}$ by Lemma~{\em \ref{lem: Caruso Liu Galois action on Kisin}}.
\end{prop}
\begin{proof}
 This is essentially the content of~\cite[Prop.\
 3.3.1]{MR2745530}. % More precisely, the existence of an appropriate
%\mar{ME:  I don't understand the second sentence. Aren't all our
%  BK-modules already given? TG: agreed. What happened was I commented
%  back in some much older argument that was written before Appendix F,
%but now that we have the discussion above this is not what we
%want. I've rewritten slightly. ME: Looks good} 
 % Kisin module is immediate from~\cite[Thm.\ 3.1.3]{MR2745530}.
 By the
 statement of Lemma~\ref{lem: Caruso Liu Galois action on Kisin}, it is enough to check that if~$s$ is sufficiently
 large, then for all~$g\in G_{K_{\piflat,s}}$ we have
 $(g-1)(\gM_{\piflat}\otimes_{\cO}\cO/\varpi^a)\subset u^N\Ainf\otimes_\gS\gM_{\piflat}\otimes_{\cO}\cO/\varpi^a$, where the
 action of~$g$ is via
%that given by
the given action of~$G_K$ on~$\gMt$. % restricting the natural action on~
 % $W(\C^\flat)_{\cO/\varpi^a}\otimes_{\A_{K}}M_B$ (note
 % that~$B/\varpi^a$ is a finite $\cO/\varpi^a$-module, so there is no
 % need to take a completed tensor product here).
 This is immediate
 from~\cite[Lem.\ 3.2.1]{MR2745530} (taking~$\mathfrak{L}=\gM_{\piflat}$, and
 noting that in their notation, $\widehat{\mathfrak{L}}$ is an
 extension of scalars of~$\gL$ to a subring
 of~$\AAinf{\cO/\varpi^a}$; this extension of scalars involves a
 twist by~$\varphi$, but since~$\varphi$ is a continuous automorphism
 of~$\Ainf$, this is harmless).
\end{proof}

\section[Stacks of semistable and crystalline Breuil--Kisin---Fargues
modules]{Stacks of semistable and crystalline Breuil--Kisin---Fargues
  modules \sectionmark{Stacks of semistable and crystalline BKF modules}}\sectionmark{Stacks of semistable and crystalline BKF modules}\label{sec:
  semistable stack}% \mar{TG: need to start fixing from here to agree
  % with the definitions above.}
We now define moduli stacks parameterizing the Breuil--Kisin--Fargues $G_K$-modules
admitting all descents that were introduced in Definition~\ref{defn:
  descending BKF to BK coefficients}. As discussed in
Section~\ref{subsec: BKF action canonical}, we also impose a condition
that the~$G_K$-actions agree with the canonical actions introduced in
Section~\ref{subsec:canonical actions}. To this end, we make the
following definition.
\begin{defn}
  \label{defn: canonical actions modulo anything}For each ~$K,a,h$
  and~$N\ge e(a+h)/(p-1)$, we fix a constant $s'(K,a,h,N)$ as in the
  statement of Proposition~\ref{prop: Caruso Liu canonical action
    semistable}. Let~$A$ be a $p$-adically complete $\cO$-algebra
  which is topologically of finite type, and let $\gMt$ be a
  projective Breuil--Kisin--Fargues $G_K$-module with $A$-coefficients
  which admits all descents. Then we say that the $G_K$-action on
  \index{canonical $G_K$-action}
  $\gMt$ is {\em canonical} if for any~$a\ge 1$, any~$\piflat$, and
  any~$s\ge s'(K,a,h,N)$, the restriction to~$G_{K_{\piflat,s}}$ of
  the action of~$G_K$ on $\gMt\otimes_{\cO}\cO/\varpi^a$ agrees with
  the action obtained from~$\gM_{\piflat}$ by Lemma~{\ref{lem:
      Caruso Liu Galois action on Kisin}}.
\end{defn}
\begin{rem}
  \label{rem: why canonical actions}The point of this definition is
  that it is immediate from
  Proposition~\ref{prop: Caruso Liu canonical action semistable} that
  the Breuil--Kisin--Fargues modules coming from lattices in
  semistable $G_K$-representations have canonical $G_K$-actions.
\end{rem}

We now define our moduli stacks of Breuil--Kisin--Fargues $G_K$-modules.
\begin{defn}\label{defn:stacks of semistable cris BKF modules}For any~$h\ge 0$ we let~$\cC_{d,\semis,h}^a$ % and~$\cC_{d,\cris,h}$
  denote the
 limit preserving category of groupoids over
  $\Spec\cO/\varpi^a$ determined by decreeing, for any finite
  type $\cO/\varpi^a$-algebra~$A$, that
  $\cC_{d,\semis,h}^a(A)$ is the
  groupoid of rank~$d$ projective Breuil--Kisin--Fargues $G_K$-modules with
  $A$-coefficients, which are of height at most~$h$, which admit
  all
  descents, and whose $G_K$-action is canonical.
% \mar{ME: While this definition may not change in its phrasing, 
% the  definition of ``BKF module admitting all descents'' will probably
% change.}

  We let $\cC_{d,\crys,h}^a$ denote the limit preserving subcategory of
  groupoids of~$\cC_{d,\semis,h}^a$ % determined by decreeing that $\cC_{d,\cris,h}^a(\Spec A)$ is the subgroupoid of
  % $\cC_{d,\semis,h}^a(\Spec A)$
  consisting of those
  Breuil--Kisin--Fargues $G_K$-modules~$\gMt$ which are furthermore crystalline. % furthermore have
  % the property that for each~$\piflat$ and each~$g\in G_K$, we have $(g-1)(\gM_{\piflat})\subset
  % \mu[\piflat]\gMt$.\mar{TG: make sure $\mu$ defined.}
  % .\mar{TG: fill
    % in whatever is needed.}
  \index{$\cC_{d,\semis,h}$} \index{$\cC_{d,\crys,h}$}
  We let~$\cC_{d,\semis,h}:=\varinjlim_a\cC_{d,\semis,h}^a$,
  and let~$\cC_{d,\crys,h}:=\varinjlim_a\cC_{d,\crys,h}^a$.% \mar{TG: why do we
  % do it this way -- can't we just define the $p$-adic one to start
  % with? Is it harder to give some justification of the procedure of
  % passing to general objects by taking limits if we do that?}
\end{defn}% \mar{TG: this definition might change if the $p$-adic HT
  % stuff does, but it shouldn't be too serious, and I propose to leave
  % it alone for the begin.}
\begin{rem}
  This definition uniquely determines limit preserving categories
  fibred in
  groupoids over $\cO$ by~\cite[Lem.\
  2.5.4]{EGstacktheoreticimages}.
  We will see shortly that the inductive limits~$\cC_{d,\semis,h}$ and~$\cC_{d,\crys,h}$ are in fact $p$-adic formal algebraic stacks
  (Theorem~\ref{thm: semi stable stack is formal algebraic}).
\end{rem}

If $A$ is a $p$-adically complete $\cO$-algebra,
then as usual we write $\cC_{d,\semis,h}(A)$ 
(resp. $\cC_{d,\crys,h}(A)$) to denote the groupoid
of morphisms from $\Spf A$ (defined using the $p$-adic topology on $A$)
to $\cC_{d,\semis,h}$ (resp.\ $\cC_{d,\crys,h}$).
More concretely, we have
$$\cC_{d,\semis,h}(A) = \varprojlim_a \cC_{d,\semis,h}^a(A/\varpi^a),$$
and similarly for $\cC_{d,\crys,h}(A)$. % \mar{TG: should be inverse limit? ME:
% I agree, and made this change.}
We will typically apply this convention in the case when $A$ is furthermore
topologically of finite type over $\cO$.  In this case each
of the quotients $A/\varpi^a$ is of finite type over $\cO/\varpi^a$
(by the definition of topologically of finite type),
and so the objects of $\cC_{d,\semis,h}^a(A/\varpi^a)$
admit a concrete interpretation as Breuil--Kisin--Fargues $G_K$-modules over~$A$.
The objects of $\cC_{d,\semis,h}(A)$ and of $\cC_{d,\crys,h}(A)$
then similarly admit a concrete interpretation.

% \mar{ME: Surely this lemma needs $A$ to be topologically of finite type
% over~$\cO$, because $\cC^a$ is only explicitly defined on finite type $\cO/\varpi^a$-algebras, and then is extended to general algebras by the limit-preserving construction.
% So I added this hypothesis.}
\begin{lem}
  \label{lem: finite O algebra points of semistable Kisin stack}Let
  $A$ be a $p$-adically complete
  $\cO$-algebra which is topologically of finite type.
Then $\cC_{d,\semis,h}(A)$ %\mar{TG: $\Spf$ versus
  %Spec issue? Also there are the same issues in Lemma~\ref{lem: finite
  %  O algebra points of potentially semistable Kisin stack} below.}
  is the groupoid of rank $d$ projective
  Breuil--Kisin--Fargues $G_K$-modules with $A$-coefficients, which
  are of height at most~$h$, which admit
  all
  descents, and whose $G_K$-action is canonical. In addition,
  $\cC_{d,\crys,h}(A)$ is the subgroupoid consisting of those
  Breuil--Kisin--Fargues $G_K$-modules~$\gMt$ which are furthermore
  crystalline.
\end{lem}
\begin{proof}
  % If~$A$ is an~$\cO/\varpi^a$-algebra for some~$a\ge 1$ then this is
  % immediate from the definition.\mar{TG: hopefully in general it's
  % then pretty much immediate, but didn't think this through yet.}
  %\mar{TG: is there a stackification issue to discuss? ME: Don't think so}
  This follows from Lemma~\ref{lem: can check projectivity of Kisin and etale phi modules modulo p^n}.
\end{proof}

The following lemma is crucial: it shows that Breuil--Kisin--Fargues modules which
admit all descents give rise to~$(\varphi,\Gamma)$-modules.

\begin{lemma}
\label{lem:C to X map}
There is a natural morphism $\cC_{d,\semis,h}\to\cX_d$, which
on $\Spf A$-points {\em (}for $p$-adically complete $\cO$-algebras $A$ that are topologically
of finite type{\em )} is given by extending scalars to~$W(\C^\flat)_A$.
\end{lemma}
\begin{proof}
As indicated in the statement of the lemma, this morphism is defined, for finite type 
$\cO/\varpi^a$-algebras, via
$\gMt \mapsto W(\C^\flat)_A \otimes_{\AAinf{A}} \gMt,$
with the target object being regarded as an $A$-valued point of $\cX_d$
via the equivalence of Proposition~\ref{prop: equivalences of categories to Ainf}.
Since both $\cC_{d,\semis,h}$ and $\cX_d$ are limit preserving
(the former by its very construction, and the latter by 
Lemma~\ref{lem:X is limit preserving}), it follows from~\cite[Lem.\ 2.5.4, Lem.\
  2.5.5~(1)]{EGstacktheoreticimages} that this construction determines
a morphism $\cC_{d,\semis,h}\to\cX_d$.
\end{proof}

\begin{remark}
\label{rem:restricting to finite type algebras}
The proof of Lemma~\ref{lem:C to X map} illustrates
a general principle (which was also applied in the proofs
of the results in Section~\ref{subsec:restriction}),
namely that to construct a morphism between limit preserving
stacks over~$\Spec \cO/\varpi^a$,
or over~$\Spf \cO$,
it suffices to define the intended morphism on finite
type $\cO/\varpi^a$-algebras (with $a$ either fixed or varying,
depending on which case we are in).

Similarly, if $P$ is any property of such a morphism that is preserved under
base-change, and that is tested by pulling back over $A$-valued points,
then to test for the property~$P$, it suffices to consider such pull-backs
in the case when $A$ is of finite type over~$\cO/\varpi^a$.

We will apply these principles consistently throughout the remainder of this section.
\end{remark}

For any~$h\ge 0$ and any choice of~$\pi^\flat$, we write~$\cC_{\piflat,d,h}$ for the moduli stack of
rank~$d$ projective Breuil--Kisin modules over~$\gS_{\pi^\flat,A}$ of height at
most~$h$, and~$\cR_{\piflat,d}$ for the corresponding stack of rank $d$
projective \'etale
$\varphi$-modules.
We also write $\cC^a_{\piflat,d,h}$ and $\cR^a_{\piflat,d}$ for their
base-change over $\cO/\varpi^a$, for any $a \geq 1$.

Recall that,
by Proposition~\ref{prop: natural morphism X to R}, we have a natural morphism
$\cX_{K,d}\to\cR_{\piflat,d}$, which for each~$s\ge 1$ can be factored,
via Lemma~\ref{lem: restriction from X to Xs},
as\[\cX_{K,d}\to\cX_{K_{\piflat,s},d}\to\cR_{\piflat,d}. \]
\begin{prop}
  \label{prop: morphism C to Xs}For any fixed $a,h$, and any~$N$ and~$s(a,h,N)$ as in
  Lemma~{\em \ref{lem: Caruso Liu Galois action on Kisin}}, then for
  any~$s\ge s(a,h,N)$ there is a canonical morphism
  $\cC^a_{\piflat,d,h}\to\cX^a_{K_{\piflat,s},d}$ obtained from the canonical
  action of Lemma~{\em \ref{lem: Caruso Liu Galois action on Kisin}}. This
  morphism fits into a commutative
  triangle% \mar{TG: probably eventually add something about how this
    % works with crystalline representations, too.}
  \[\xymatrix{& \cC^a_{\piflat,d,h}\ar[dl]\ar[d]\\ \cX^a_{K_{\piflat,s},d}\ar[r]&\cR^a_{\piflat,d}}\]
%  $\cC^a_{d,\infty,h}\to\cR^a_{d,\infty}$ factors through the morphism $\cX^a_{K_s,d}\to\cR^a_{d,\infty}$.
\end{prop}%\mar{TG: need to be rewritten}
\begin{proof}% \mar{TG: needed, should basically be Frobenius
    % amplification}
As discussed in Remark~\ref{rem:restricting to finite type algebras},
%  As in the proof of Proposition~\ref{prop: natural morphism X to R},
% \mar{ME: I don't understand why Prop.~\ref{prop: natural morphism X to R} was being
% cited here at all.  I think it was just because we're using the general principle that
% we can define morphisms by restricting to finite-type algebras, and then
% extending via the limit preserving property, and so I rewrote things to indicate this
% more directly.}
it is enough to show that if~$A$ is a finite type
%\mar{ME: finite type algebra, maybe?}
  $\cO/\varpi^a$-algebra, and~$\gM$ is a finite projective \'etale Breuil--Kisin
  module with $A$-coefficients, then we can canonically % \mar{TG: need
    % to think about why the extension will actually be canonical!}
  extend the action of~$G_{K_{\pi^\flat,\infty}}$ on
  $W(\C^\flat)_A\otimes_{\gS_{\piflat,A}}\gM$ to an action
  of~$G_{K_{\pi^\flat,s}}$, if~$s\ge s(a,h,N)$. Such an extension
is provided by Lemma~\ref{lem: Caruso Liu Galois action on Kisin}.
\end{proof}

\begin{lem}\label{lem: C to Xs is proper}
  The morphism $\cC^a_{\piflat,d,h}\to\cX^a_{K_{\piflat,s},d}$ of
  Proposition~{\em \ref{prop: morphism C to Xs}} is representable by
  algebraic spaces, proper, and of finite presentation.% \mar{TG: do
  % we explicitly use this properness? It should help scheme-theoretic
  % images to be better behaved, at the very least?}
\end{lem}
\begin{proof} This follows by a standard graph argument.  Namely,
	we write this morphism as the
  composite \[\cC^a_{\piflat,d,h}\to\cC^a_{\piflat,d,h}\times_{\cR^a_{\piflat,d}}\cX^a_{K_{\piflat,s},d}
    \to\cX^a_{K_{\piflat,s},d}. \] The first arrow is a closed immersion, 
  being
  the base change of the diagonal morphism $\cX_d\to\cX_d\times_{\cR_{\piflat,d}}\cX_d$, which is a closed immersion
%\mar{ME: What about f.p.? ME: Now sorted out.}
  by Proposition~\ref{prop: diagonal of X to R}. The second arrow
  is representable by  algebraic spaces  and proper, %and of finite presentation
because it is the base change of the morphism
  $\cC^a_{\piflat,d,h}\to{\cR^a_{\piflat,d}}$ of Theorem~\ref{thm: C is a $p$-adic formal algebraic stack and related
    properties}.
The  composite morphism  is thus representable by algebraic spaces
and  proper.  Such a morphism is in particular of finite type
and separated (and so also quasi-separated), and hence Lemma~\ref{lem:finite presentation from closed immersion}
below allows us to conclude that it is also of finite presentation
(once we take into account Lemma~\ref{lem:X is limit preserving},
 which shows that its target is limit preserving).
\end{proof}

\begin{lem}
\label{lem:finite presentation from closed immersion}
Let $\cZ \hookrightarrow \cX$ be a morphism of % closed immersion of %algebraic
stacks over a locally
Noetherian base scheme~$S$ 
which is representable by algebraic spaces, of finite type,
and quasi-separated.
If $\cX$ is limit preserving, %locally of finite presentation over~$S$,
then this morphism is furthermore of finite presentation.
%In particular, $\cZ$ is also locally of finite presentation over~$S$,
%while if $\cX$ is furthermore of finite presentation over~$S$,
%then the same is true of~$\cZ$.
\end{lem}
\begin{proof}
It is enough to show, for any morphism $T \to \cX$ whose source
is an affine $S$-scheme $T$, that the pulled back morphism
(of algebraic spaces)
$T\times_{\cX} \cZ \to T$ is of finite presentation.  Since $\cX$ is
%locally of finite presentation over~$S$, it is
limit preserving,
 the morphism $T \to \cX$ factors through a morphism
of $S$-schemes
$T \to T'$, where $T'$ is affine and of finite presentation
over~$S$.  Thus, replacing $T$ by~$T'$, we may assume
that $T$ is of finite presentation over~$S$, and consequently Noetherian.

The base-changed morphism $T\times_{\cX} \cZ \to T$ is thus a
finite type and quasi-separated morphism  from
an algebraic space to a Noetherian scheme, and hence is of finite presentation
by~\cite[\href{https://stacks.math.columbia.edu/tag/06G4}{Tag 06G4}]{stacks-project}.
This proves the lemma.
\end{proof}

\begin{df}
	\label{def:2-Cartesian}
The morphisms of Lemma~\ref{lem: restriction from X to Xs} and
Proposition~\ref{prop: morphism C to Xs} allow us to define a
stack~$\cC^a_{\piflat,s,d,h}$ by the requirement that it fits into a  \index{$\cC^a_{\piflat,s,d,h}$}
2-Cartesian diagram \numequation\label{eqn: 2 cartesian canonical C}\xymatrix{\cC^a_{\piflat,s,d,h} \ar[r]\ar[d] &  \cC^a_{\piflat,d,h}\ar[d]\\
  \cX^a_{K,d}\ar[r] & \cX^a_{K_{\piflat,s},d}}\end{equation}
The stack~$\cC^a_{\piflat,s,d,h}$ classifies rank $d$
projective Breuil--Kisin modules $\gM$ of height at most $h$
over $\gS_{\piflat,A}$
equipped with an extension of the canonical $G_{K_{\piflat,s}}$-action
on $W(\C^\flat)_A\otimes_{\gS_{\piflat,A}} \gM$ to an action of $G_K$
(making it an \'etale $(\varphi,G_K)$-module).

Since the lower horizontal arrow in~\eqref{eqn: 2 cartesian canonical C}
%the diagram defining $\cC^a_{s,d,h}$
%this diagram
is representable by algebraic spaces and of 
finite presentation, by Lemma~\ref{lem: restriction from X to Xs},
and since~$\cC^a_{\piflat,d,h}$ is an
algebraic stack of finite presentation over $\cO/\varpi^a$,
by Theorem~\ref{thm: C is a $p$-adic formal algebraic stack and related
    properties}, we see that~$\cC^a_{\piflat,s,d,h}$   is also an algebraic
stack of finite presentation over $\cO/\varpi^a$.
Since the right hand vertical arrow in this diagram is representable by
algebraic spaces, proper, and of
finite presentation,
by Lemma~\ref{lem: C to Xs is proper}, so is the left
hand vertical arrow.%  Since~$\cX^a_{K,d}$ is an Ind-algebraic stack by Proposition~\ref{prop:X is an Ind-stack}, we can
% form the scheme-theoretic image~$\cZ^a_{\piflat,s,d,h}$ of the morphism $\cC^a_{\piflat,s,d,h}\to \cX^a_{K,d}$, which is a
% closed algebraic substack of~$\cX^a_{K,d}$, and so, by the following lemma
% (together with Proposition~\ref{prop:X is an Ind-stack}),
% is in particular of 
% finite presentation over $\cO/\varpi^a$.
% \mar{ME: I think this last claim (and the implied argument for it --- that a closed
% immersion into something of finite presentation is of finite presentation) is correct,
% but maybe needs a tiny bit of justification. This might also help with the finite 
% presentation claims in what follows.}
% \mar{ME: Actually, I couldn't see where the scheme-theoretic image $\cZ^a$
% is actually used.  All mentions of it that I could find by searching the \TeX \ 
% file seem to be commented out. So I haven't actually added ``the following lemma'' yet,
% since maybe it's not needed, at least not at this point? TG:
% correct! This is something vestigal from the earlier arguments which
% only constructed effective versal rings.}
\end{df}

\begin{lem}
\label{lem:auxiliary object has affine diagonal}
The diagonal of~$\cC^a_{\piflat,s,d,h}$ is affine and of finite
presentation.
\end{lem}
\begin{proof}
As we already noted, 
the morphism $\cC^a_{\piflat,s,d,h} \to \cX_{K,d}$ is representable by algebraic
spaces and proper (and so in particular is both separated and of finite type);
the claim of the lemma thus follows from Proposition~\ref{prop: properties
of X pulled back from R}, together with Lemma~\ref{lem:controlling diagonal} below.
%
%The diagonal morphism
%$$\cC^a_{\piflat,s,d,h} \to \cC^a_{\piflat,s,d,h}\times_{\cO/\varpi^a}\cC^a_{\piflat,s,d,h}$$
%may be factored as the composite of the relative diagonal
%\numequation
%\label{eqn:first factor}
%\cC^a_{\piflat,s,d,h} \to \cC^a_{\piflat,s,d,h}\times_{\cX_{K,d}^a}\cC^a_{\piflat,s,d,h}
%\end{equation}
%and the morphism
%\numequation
%\label{eqn:second factor}
%\cC^a_{\piflat,s,d,h}\times_{\cX_{K,d}^a}\cC^a_{\piflat,s,d,h}
%\to \cC^a_{\piflat,s,d,h}\times_{\cO/\varpi^a}\cC^a_{\piflat,s,d,h}
%\end{equation}
%which is a base change of the diagonal morphism
%$$\cX_{K,d}^a \to \cX_{K,d}^a \times_{\cO/\varpi^a} \cX_{K,d}^a$$
%As we already noted, 
%the morphism $\cC^a_{\piflat,s,d,h} \to \cX_{K,d}$ is representable by algebraic
%spaces and proper (and so in particular is both separated and of finite type);
%thus~\eqref{eqn:first factor}
%is a closed immersion (and in particular
%affine), and is also
%of finite presentation, by Lemma~\ref{lem:finiteness of diagonals}.
%The morphism~\eqref{eqn:second factor}
%is also affine and of finite presentation,
%by Proposition~\ref{prop: properties of X pulled back from R}.
%Thus their composite is affine and of finite presentation, as claimed.
\end{proof}

\begin{lem}
\label{lem:controlling diagonal}
Suppose that $f:\cX \to \cY$ is a morphism of stacks over a base scheme $S$
which is representable
by algebraic spaces, separated, and of finite type.  Suppose also that
the diagonal $\Delta_{\cY}: \cY \to \cY\times_S \cY$ is affine and of finite
presentation.  Then the diagonal $\Delta_{\cX}: \cX \to \cX \times_S \cX$
is affine and of finite presentation.
\end{lem}
\begin{proof}
The diagonal morphism 
$$\cX \to \cX\times_S \cX$$
may be factored as the composite of the relative diagonal
\numequation
\label{eqn:first factor}
\cX \to \cX \times_{\cY} \cX
\end{equation}
and the morphism
\numequation
\label{eqn:second factor}
\cX\times_{\cY} \cX \to \cX\times_S \cX,
\end{equation}
which is a base change of the diagonal morphism~$\Delta_{\cY}$. 
Our assumption on $\Delta_{\cY}$, then, implies that
the morphism~\eqref{eqn:second factor} 
is affine and of finite presentation.
Since $\cX \to \cY$ is representable by algebraic spaces and separated,
its diagonal~\eqref{eqn:first factor} is a closed immersion, and thus affine. 
Since it is furthermore
of finite type, the morphism~\eqref{eqn:first factor} is also 
of finite presentation, by Lemma~\ref{lem:finiteness of diagonals}.
Thus $\Delta_{\cX}$ is the composite of morphisms that are affine and
of finite presentation, and the lemma follows.
\end{proof}

%\begin{lem}
%Let $\cZ \hookrightarrow \cX$ be a closed immersion of stacks over a locally
%Noetherian base scheme~$S$, with $\cZ$ being algebraic,
%and with $\cX$ being an inductive limit of  
%        algebraic stacks of finite presentation over~$S$, with the transition maps
%        being closed immersions.   Then $\cZ$ is of finite presentation over~$S$.
%\end{lem}
%\begin{proof}
%\end{proof}


%%\mar{ME: We apply this lemma in the formal alg.\ stack context,
%%in Theorem~\ref{thm: semi stable stack is formal algebraic},
%%and again in Proposition~\ref{prop: HT weights are a closed condition},
%%so maybe should state something in that level of generality.}
%\begin{lem}
%\label{lem:finite presentation from closed immersion}
%Let $\cZ \hookrightarrow \cX$ be a closed immersion of %algebraic
%stacks over a locally
%Noetherian base scheme~$S$. 
%If $\cX$ is limit preserving, %locally of finite presentation over~$S$,
%then this closed immersion is necessarily of finite presentation.
%%In particular, $\cZ$ is also locally of finite presentation over~$S$,
%%while if $\cX$ is furthermore of finite presentation over~$S$,
%%then the same is true of~$\cZ$.
%\end{lem}
%\begin{proof}
%It is enough to show, for any morphism $T \to \cX$ whose source
%is an affine $S$-scheme $T$, that the pulled back closed immersion
%$T\times_{\cX} \cZ \to T$ is of finite presentation.  Since $\cX$ is
%%locally of finite presentation over~$S$, it is
%limit preserving,
% the morphism $T \to \cX$ factors through a morphism
%of $S$-schemes
%$T \to T'$, where $T'$ is affine and of finite presentation
%over~$S$.  Thus, replacing $T$ by~$T'$, we may assume
%that $T$ is of finite presentation over~$S$, and consequently Noetherian.
%
%The base-changed morphism $T\times_{\cX} \cZ \to T$ is thus a
%closed immersion into a Noetherian scheme, and hence is of finite presentation.
%This proves the lemma.
%\end{proof}

\begin{prop}
  \label{prop: BKF closed in BK}For each~$\piflat$ and each~$s\ge s'(K,a,h,N)$, there are  natural
  closed immersions $\cC^a_{d,\crys,h}\to\cC^a_{d,\semis,h}\to
  \cC^a_{\piflat,s,d,h}$. In particular, $\cC^a_{d,\semis,h}$
  and~$\cC^a_{d,\crys,h}$ are
 algebraic stacks of finite presentation over~$\cO/\varpi^a$, and have
 affine diagonals. % \mar{TG: needs proving, didn't check it yet!
  % This should be the key statement for having control of the
  % semistable stack.}\
%\mar{ME: Add affine diagonal statement}
\end{prop}%\mar{TG: remember to update this if the definitions change.}
\begin{proof} %As observed above,
We have already observed that $\cC^a_{\piflat,s,d,h}$ is of
finite presentation over~$\cO/\varpi^a$, and has affine diagonal by
Lemma~\ref{lem:auxiliary object has affine diagonal}. Once we
construct the claimed closed immersions, 
the claims of finite presentation will follow from
Lemma~\ref{lem:finite presentation from closed immersion}; and since closed
immersions are in particular monomorphisms,
the diagonals of  $\cC_{d,\semis,h}$ and  $\cC_{d,\crys,h}$  are then obtained via base change
from the diagonal of~ $\cC^a_{\piflat,s,d,h}$, and are in particular affine.
We thus focus on constructing these closed immersions.
Throughout the proof, we take into
Remark~\ref{rem:restricting to finite type algebras},
which allows us to restrict our attention to the points
of the various stacks in question that are defined
over finite type $\cO/\varpi^a$-algebras.

Lemma~\ref{lem:C to X map}
  constructs a morphism $\cC^a_{d,\semis,h}\to\cX^a_{K,d}$, given
  by extending scalars to~$W(\C^\flat)$, and it follows from
  Lemma~\ref{lem: uniqueness of descent for G K infty Ainf} that there
  is a natural morphism $\cC^a_{d,\semis,h}\to\cC^a_{\piflat,d,h}$, defined via
  $\gMt \mapsto \gM_{\piflat}$. The composite morphisms
  $\cC^a_{d,\semis,h} \to \cC^a_{\piflat,d,h} \to
  \cX^a_{K_{\piflat,s},d}$ and
  $\cC^a_{d,\semis,h} \to \cX^a_{K,d}\to \cX^a_{K_{\piflat,s},d}$
  coincide by definition (see Definitions~\ref{defn: canonical actions
    modulo anything} and~\ref{defn:stacks of semistable cris BKF modules}). %  Proposition~\ref{prop: BKF has canonical actions}.
% \mar{ME: We will replace this citation of  the  bogus proposition
% by an appeal to the  (to-be-modified) definition of BKF module admitting
% all descents.}
  Thus these morphisms induce a morphism
\numequation
\label{eqn:induced map}
\cC^a_{d,\semis,h}\to \cC^a_{\piflat,s,d,h}.
\end{equation}

  To see that~\eqref{eqn:induced map}
%the induced map $\cC^a_{d,\semis,h}\to \cC^a_{\piflat,s,d,h}$
is a monomorphism, it is enough to note
  that if~$A$ is any finite type $\cO/\varpi^a$-algebra~$A$,
  and~$\gMt$ is any Breuil--Kisin--Fargues module over $A$ which
  admits all descents and whose $G_K$-action is canonical,
 corresponding to an $A$-valued point of~$\cC^a_{d,\semis,h},$
  then~$\gMt=\AAinf{A}\otimes_{\gS_{\piflat,A}}\gM_{\piflat}$ is
  determined by~$\gM_{\piflat}$, and the $G_K$-action on~$\gMt$ is
  determined by the $G_K$-action on
  $W(\C^\flat)_A\otimes_{\AAinf{A}}\gMt$.

% \mar{ME: If we add conditions to the  definition of ``BKF modules admitting
% all descents'', we'll have to confirm here that those are closed
% conditions.  This should similar to verifying that ``crys'' is closed
% in ``ss''.}
We now show that~\eqref{eqn:induced map}
%$\cC^a_{d,\semis,h}\to \cC^a_{\piflat,s,d,h}$
is a closed immersion. It is enough to do
  this after pulling back to some finite type
  $\cO/\varpi^a$-algebra~$A$, where we need to show that the conditions that
$\gMt=\AAinf{A}\otimes_{\gS_{\piflat,A}}\gM_{\piflat}$ is $G_K$-stable
and admits all descents, and that the $G_K$-action is canonical, are closed conditions. We begin with the first
of these. It is enough to show that for each~$g\in G_K$, the condition
that $g(\gMt)\subset\gMt$ is closed. This condition is equivalent to
the vanishing of the composite morphism \[g^*\gMt\to
  \gMt[1/u]\to\gMt[1/u]/\gMt,\]and this is a closed condition by
Lemma~\ref{lem: inclusion lattices closed condition}. Let~$\cC^a_{d,G_K}$ denote the closed substack of
$\cC^a_{\piflat,s,d,h}$ for which~$\gMt$ is $G_K$-stable. % \mar{TG: not quite sure how to finish
  % from here! Not sure if we want to use something like the arguments
  % from Appendix~\ref{sec: rigid analytic perspective}? If we had
  % smaller coefficients it would be the proof of~\cite[Lem.\ 5.4.10]{EGstacktheoreticimages}.}

 We next show that for each choice of uniformizer~$\pi'$, and
each~$(\pi')^\flat$, the condition that~$\gMt$ admits a descent
to~$\gS_{A,(\pi')^\flat}$ is a closed condition. To this end, note
that by Proposition~\ref{prop: equivalences of categories to Ainf},
$\gMt[1/u]$ descends uniquely to~$\gS_{(\pi')^\flat,A}[1/u]$, so we
have a morphism $\cC^a_{d,G_K}\to\cR^a_{(\pi')^\flat,d}$. Then the
fibre product
$\cC^a_{d,G_K}\times_{\cR^a_{(\pi')^\flat,d}}\cC^a_{(\pi')^\flat,d,h}$
is the moduli space of pairs~$(\gM,\gM')$ where~$\gM := \gM_{\piflat}$ is as above
(an object classified by~$\cC^a_{d,G,K}$),
and~$\gM'$ is a Breuil--Kisin module for~$\gS_{(\pi')^\flat,A}$ of
height at most~$h$ which
satisfies \[W(\C^\flat)_A\otimes_{\gS_{(\pi')^\flat,A}}\gM'=W(\C^\flat)_A\otimes_{\gS_{\pi^\flat,A}}\gM.\]
Consider the substack~$\cC^a_{d,G_K,(\pi')^\flat}$ of this fibre
product satisfying the condition that
\numequation\label{eqn: pi pi prime
  agree}\AAinf{A}\otimes_{\gS_{(\pi')^\flat,A}}\gM'=\gMt.\end{equation} Note that this is exactly the condition that~$\gMt$ admits
a descent for~$(\pi')^\flat$, and that~$\gM'$ is this descent (which
is uniquely determined by Lemma~\ref{lem: uniqueness of descent for G K infty Ainf}),
so the
projection $\cC^a_{d,G_K,(\pi')^\flat}\to \cC^a_{d,G_K}$ is a
monomorphism, and we need to show that it is a closed immersion. Since the projection
$\cC^a_{d,G_K}\times_{\cR^a_{(\pi')^\flat,d}}\cC^a_{(\pi')^\flat,d,h}\to
\cC^a_{d,G_K}$ is proper (being a base change of the proper morphism
$\cC^a_{(\pi')^\flat,d,h}\to\cR^a_{(\pi')^\flat,d}$), and since proper
monomorphisms are closed immersions, it is enough to show that
$\cC^a_{d,G_K,(\pi')^\flat}$ is a closed substack of
$\cC^a_{d,G_K}\times_{\cR^a_{(\pi')^\flat,d}}\cC^a_{(\pi')^\flat,d,h}$;
that is, we must show that~\eqref{eqn: pi pi prime
  agree} is a closed condition. This again follows from
Lemma~\ref{lem: inclusion lattices closed condition}.

We now consider the closed substack  of $\cC^a_{\piflat,s,d,h}$ for
which~$\gMt$ is $G_K$-stable and admits a descent for
each~$(\pi')^\flat$. We need to show that for each~$(\pi')^\flat$, the
further conditions that  \[\gM_{\piflat}/[\piflat]\gM_{\piflat}=\gM_{\piflatprime}/[\piflatprime]\gM_{\piflatprime}\]and  \[\varphi^*\gM_{\piflat}/E_{\piflat}\varphi^*\gM_{\piflat}=\varphi^*\gM_{\piflatprime}/E_{\piflatprime}\varphi^*\gM_{\piflatprime}\]
are closed conditions.

The arguments in both cases are very similar (and in turn are similar
to the proof of Lemma~\ref{lem: inclusion lattices closed condition}),
so we only give the argument in the second case, leaving the first to
the reader. Each side is a projective $\cO_K\otimes_{\Zp}A$-submodule
of the projective $\cO_{\C}\otimes_{\Zp}A$-module
$\cO_{\C}\otimes_{\AAinf{A},\theta}\varphi^*\gMt$, and indeed spans this module
after extension of scalars to~$\cO_{\C}$. By symmetry, it is enough to
show that for each
element~$m\in \varphi^*\gM_{\piflatprime}/E_{\piflatprime}\varphi^*\gM_{\piflatprime}$,
the condition that $m\in \varphi^*\gM_{\piflat}/E_{\piflat}\varphi^*\gM_{\piflat}$ is
closed.

Write $P:=\varphi^*\gM_{\piflat}/E_{\piflat}\varphi^*\gM_{\piflat}$, and choose
a finite projective $\cO_K\otimes_{\Zp}A$-module~$Q$ such
that~$F:=P\oplus Q$ is free. We can think of~$m$ as an element of
$\cO_{\C}\otimes_{\cO_K}P$ and thus as an element of
$\cO_{\C}\otimes_{\cO_K}F$, and it is enough to check that the
condition that~$m\in F$ is closed. Choosing a basis for~$F$, we reduce
to the case that~$F$ is one-dimensional, and thus to showing that the
condition that an element of~$\cO_{\C}\otimes_{\Zp}A$ lies
in~$\cO_{K}\otimes_{\Zp}A$ is closed. Since~$\cO_{\C}$ is a torsion-free
and thus flat~$\cO_K$-module, $\cO_{\C}/\varpi^a$ is flat and thus
free as a module for the Artinian ring $\cO_K/\varpi^a$. Thus
$\cO_{\C}\otimes_{\Zp}A$ is a free~$\cO_{K}\otimes_{\Zp}A$-module, and
the result follows by choosing a basis.

We now need to show that the condition that the~$G_K$-action is
canonical is a closed condition. Explicitly, by Lemma~\ref{lem: Caruso
  Liu Galois action on Kisin}, we need to show that for each~$b\le a$,
each $N\ge e(b+h)/(p-1)$, each~$s\ge s'(K,b,h,N)$, 
each~$\piflat$,
and each~$g\in G_{K_{\piflat,s}}$,
the condition that 
\[(g-1)(\gM_{\piflat}\otimes_{\cO}\cO/\varpi^b)\subset
  u^N(\gMt\otimes_{\cO}\cO/\varpi^b)\] is closed. This
follows from Lemma~\ref{lem: element of lattice closed condition}.

Finally, it remains to show that the monomorphism
~$\cC^a_{d,\crys,h}\to \cC^a_{d,\semis,h}$ is a closed
immersion; that is, for each~$g\in G_K$ and each~$\piflat$ we need to
show that the condition that \[(g-1)(\gM_{\piflat})\subseteq \varphi
  ^{-1}(\mu)[\piflat]\gMt \] is a closed condition. It is enough to show that
for each~$m\in\gM_{\piflat}$, the condition that
$(g-1)(m)\in \varphi^{-1}(\mu)[\piflat]\gMt$ is closed. Since $ \varphi^{-1}(\mu)[\piflat]\gMt$ is a
finite projective $\AAinf{A}$-module, this follows from another
application of Lemma~\ref{lem: element of lattice closed condition}.%  \mar{TG: this should be fine but I'm not going to try to do
  % it until we have the definitions settled.}
% \mar{TG: needs to be finished - again it's basically the same question as the
%   previous marginal comment, of showing that the condition that one
%   lattice contains another is a closed condition.}
% \mar{TG: can we show closed
%     immersion by formulating $\piflat$ descents as a faithfully flat
%     descent problem? Hopefully being $G_K$-stable is just easy and is
%     done first. ME points out that the descent should just
%     be~$M_\pi\cap\gMt$ (but of course that won't always be a descent).}
\end{proof}


% The following definition is again motivated by the results of
% Appendix~\ref{app: BKF pst}.\mar{TG: put in a more precise citation
%   eventually.}

% \begin{defn}
%   \label{defn: BKF module over L with GK action}Let~$L/K$ be a finite
%   extension. A Breuil--Kisin--Fargues $(G_K,G_L)$-module of height at
%   most~$h$ with $A$-coefficients is a Breuil--Kisin--Fargues
%   $G_L$-module of height at most~$h$ with $A$-coefficients~$\gMt$,
%   equipped with an extension of the~$G_L$-action on~$\gMt$ to a
%   continuous semilinear action of~$G_K$ which commutes with~$\varphi$.

%   We say that~$\gMt$ admits all descents if the underlying
%   Breuil--Kisin--Fargues $G_L$-module admits all descents.
% \end{defn}
% Note that if~$\gMt$ is a Breuil--Kisin--Fargues $(G_K,G_L)$-module
% which admits all descents, then the $W(k)\otimes_{\Zp}A$-module $\gM$
% and the $\cO_K\otimes_{\Zp}A$-module \mar{TG: fill in: I think the
%   assumption that these things are independent of $\piflat$
%   automatically implies
%   that $\Gal(L/K)$ acts, because $G_K$ permutes the~$\piflat$ and $G_L$
%   acts trivially? Then next need to show that the moduli stack of
%   these is closed in the appropriate fibre product.}


\begin{thm}\label{thm: semi stable stack is formal algebraic}% \mar{TG:
    % more finite type finite presentation. ME: Added to statement and proof.}
  $\cC_{d,\semis,h}$ is a $p$-adic formal algebraic stack
of finite presentation,
  as is its closed substack~$\cC_{d,\crys,h}$. In addition,
both of these stacks have affine diagonal,
  and each of the 
  morphisms
$\cC_{d,\semis,h}\to\cX_{K,d}$
 and
$\cC_{d,\crys,h}\to\cX_{K,d}$
is representable by algebraic
  spaces, proper, and of finite presentation. 
% \mar{ME: Add affine diagaonal. ME: Added to statement; still have to
%   add to proof TG: hopefully more or less immediate from Prop~\ref{prop: BKF closed in BK}?
% ME: Yes, added a sentence to this effect. Remove this comment if you're happy with what
% I wrote.}
\end{thm}%\mar{TG: needs proof}
\begin{proof}Since ~$\cC_{d,\crys,h}$ is a closed substack
  of~$\cC_{d,\semis,h}$, it suffices to prove the statements of the lemma
  for~$\cC_{d,\semis,h}$.
(Here we use the fact that a closed immersion is a monomorphism
to see that the diagonal of $\cC_{d,\crys,h}$ is obtained via base change
from the diagonal of $\cC_{d,\semis,h}$, while we use Lemma~\ref{lem:finite
presentation from closed immersion} to transfer finite presentation
properties from $\cC_{d,\semis,h}$ to $\cC_{d,\crys,h}$).

 By Proposition~\ref{prop: criterion for
    p-adic formal algebraic stack}, to see that~$\cC_{d,\semis,h}$ is
  a $p$-adic formal algebraic stack of finite presentation, it is enough to show
  that each~$\cC_{d,\semis,h}^a$ is an algebraic stack of finite
presentation over $\cO/\varpi^a$,
  which was proved in Proposition~\ref{prop: BKF closed in BK}.
  Since each $\cC_{d,\semis,h}^a$ has affine diagonal, by the same
proposition, we see as well that $\cC_{d,\semis,h}$ has affine diagonal.

  That $\cC_{d,\semis,h}\to\cX_{K,d}$ is representable by algebraic
  spaces, proper, and of finite presentation also follows from
  Proposition~\ref{prop: BKF closed in BK}. Indeed, for each~$a$ the
  morphism $\cC_{d,\semis,h}^a\to\cX_{K,d}$ factors as
  $\cC_{d,\semis,h}^a\to\cC^a_{\piflat,s,d,h}\to \cX_{K,d}^a$, where the
  first morphism is a closed immersion of finite type algebraic
  stacks, so it is enough to prove the same properties for the
  morphism~$\cC^a_{\piflat,s,d,h}\to \cX_{K,d}^a$; as explained in
  Definition~\ref{def:2-Cartesian}, this follows from Lemma~\ref{lem:
    C to Xs is proper}.
\end{proof}


\subsection{Potentially crystalline and potentially semistable
	stacks}
We now consider the corresponding potentially semistable and
potentially crystalline versions of these moduli stacks of
Breuil--Kisin--Fargues modules. % \mar{TG: do this next; just a
  % question of the lattice being preserved by~$\Gal(L/K)$ being a
  % closed condition, which is hopefully immediate from stuff in the
  % appendix that we already used above.}

  \begin{df}\label{defn: semistable crys stacks L over K}
For a Galois extension~$L/K$,
%and % \mar{TG: for now defining these as $p$-adic stacks rather than
  % mod~$\varpi^a$}
define stacks~$\cC_{d,\semis,h}^{L/K}$ and \index{$\cC_{d,\semis,h}^{L/K}$}
\index{$\cC_{d,\crys,h}^{L/K}$}
$\cC_{d,\crys,h}^{L/K}$ as follows: for each~$a\ge 1$ we
let~$\cC_{d,\semis,h}^{L/K,a}$ denote the limit preserving category of
groupoids over $\Spec\cO/\varpi^a$ determined by decreeing,
for any finite type $\cO/\varpi^a$-algebra~$A$,
that $\cC_{d,\semis,h}^{L/K,a}(A)$ is the groupoid of
Breuil--Kisin--Fargues $G_K$-modules with $A$-coefficients, which are
of height at most~$h$, which admit all descents over~$L$, and whose
$G_L$-actions are canonical. 

We let
$\cC_{d,\crys,h}^{L/K,a}$ denote the limit preserving subcategory of
groupoids of~$\cC_{d,\semis,h}^a$ consisting of those
Breuil--Kisin--Fargues $G_K$-modules~$\gMt$ for which the action
of~$G_L$ is crystalline.

We let~$\cC_{d,\semis,h}^{L/K}:=\varinjlim_a\cC_{d,\semis,h}^{L/K,a}$,
and let~$\cC_{d,\crys,h}^{L/K}:=\varinjlim_a\cC_{d,\crys,h}^{L/K,a}$.% \mar{TG: again,
  % I guess it would be cleaner if we could just define a $p$-adic
  % version immediately?}
\end{df}

% If $A$ is an $p$-adically complete $\cO$-algebra,
% then we write 
% $\cC_{d,\semis,h}^{L/K}(A)$
% (resp.\
% $\cC_{d,\crys,h}^{L/K}(A)$)
% to denote the groupoid of morphisms
% from $\Spf A$ (defined with respect to the $p$-adic topology on $A$)
% to 
% $\cC_{d,\semis,h}^{L/K}$
% (resp.\
% $\cC_{d,\crys,h}^{L/K}$).
The following lemma is proved in exactly the same way as Lemma~\ref{lem: finite O algebra points of  semistable Kisin stack}.

\begin{lem}
  \label{lem: finite O algebra points of potentially semistable Kisin
    stack}Let $A$ be a $p$-adically complete
  $\cO$-algebra which is topologically of finite type over~$\cO$. Then $\cC_{d,\semis,h}^{L/K}(A)$ is the groupoid of
  Breuil--Kisin--Fargues $G_K$-modules with $A$-coefficients, which
  are of height at most~$h$, and which admit all descents over~$L$, and whose
$G_L$-actions are canonical; and
  $\cC_{d,\crys,h}^{L/K}(A)$ is the subgroupoid of those
  Breuil--Kisin--Fargues $G_K$-modules~$\gMt$ which are furthermore
  crystalline.
\end{lem}
% \begin{proof}
% \mar{ Should be the same as the proof of Lemma~\ref{lem: finite O
%     algebra points of potentially semistable Kisin stack} once that exists.}
% \end{proof}

\begin{prop}% \mar{TG:
    % more finite type finite presentation}
  \label{prop: pst and pcrys stacks are padic}For any Galois
  extension~$L/K$, any $d$, and any~$h$, 
  ~$\cC_{d,\semis,h}^{L/K}$ is a 
  $p$-adic formal algebraic stack of finite presentation, and
  ~$\cC_{d,\crys,h}^{L/K}$ is a closed substack, which is thus
again a $p$-adic formal algebraic stack of finite presentation.
Both $\cC_{d,\semis,h}^{L/K}$ and $\cC_{d,\crys,h}^{L/K}$ have affine diagonal,
and the natural morphisms
  $\cC_{d,\semis,h}^{L/K}\to\cX_{K,d}$
  and $\cC_{d,\crys,h}^{L/K}\to\cX_{K,d}$
are representable by algebraic
  spaces, proper, and of finite presentation.
% \mar{ME: And affine diagonal! ME: Added statement about this, and about finite presentation;
% added proofs too!}
\end{prop}
\begin{proof}
  We have a natural
  morphism~$\cC_{d,\semis,h}^{L/K}\to\cC_{d,\semis,h}^{L/L}$ (given by
  restricting the action of~$G_K$ to~$G_L$; the target is of course
  just the stack denoted~$\cC_{d,\semis,h}$ above, but for~$L$ instead
  of~$K$), and a natural morphism
  $\cC_{d,\semis,h}^{L/K}\to\cX_{K,d}$, and thus a natural morphism
  \numequation\label{eqn: mapping pst BKF to fibre
    product}\cC_{d,\semis,h}^{L/K}\to \cC_{d,\semis,h}^{L/L}\times_{\cX_{L,d}}\cX_{K,d}   \end{equation}

 We claim that~\eqref{eqn: mapping pst BKF to fibre
    product} is a closed immersion. Given this, the proposition
  follows. Indeed, by Theorem~\ref{thm: semi stable stack is formal algebraic}, $\cC_{d,\semis,h}^{L/L}$ is a
  $p$-adic formal algebraic stack of finite presentation, and the morphism
  $\cC_{d,\semis,h}^{L/L}\to\cX_{L,d}$ is representable by algebraic
  spaces, proper, and of finite presentation. Furthermore,
  $\cX_{K,d}\to\cX_{L,d}$ is representable by algebraic spaces %, affine,
  % \mar{ME: Just to check: we don't use affineness of the morphism
  %         $\cX_{K,d} \to \cX_{L,d}$, do we?  In which case,
  %         maybe we can omit it from the list of recalled properties,
  %         just to avoid unnecessary confusion? TG: agreed, done,
  %         commenting out.}
  and of
  finite presentation by Lemma~\ref{lem: restriction from X to Xs}. That~$\cC_{d,\crys,h}^{L/K}$ is a closed substack of
  $\cC_{d,\semis,h}^{L/K}$ follows from Proposition~\ref{prop: BKF
    closed in BK}.
  The claims on the diagonals are proved by appeal to Lemma~\ref{lem:controlling
diagonal} (\emph{cf.}\ the proof of Lemma~\ref{lem:auxiliary object has affine diagonal}).

  It remains to prove the claim. It suffices to prove this after
  pulling back via a morphism $\Spec A\to\cX_{L,d}$ for some finite
  type $\cO/\varpi^a$-algebra~$A$. Unwinding the definitions, we have
  to show that if~$\gMt$ is a Breuil--Kisin--Fargues $G_L$-module with
  $A$-coefficients, with a compatible action of~$G_K$ on
  $W(\C^{\flat})_A\otimes_{\AAinf{A}}\gMt$, then the condition
  that~$\gMt$ is $G_K$-stable is a closed condition. It suffices to
  show that for each~$g\in G_K$, the condition that
  $g(\gMt)\subseteq\gMt$ is a closed condition; as in the proof of
  Proposition~\ref{prop: BKF closed in BK}, this follows from
  Lemma~\ref{lem: inclusion lattices closed condition}, applied to the composite morphism \[g^*\gMt\to
  \gMt[1/u]\to\gMt[1/u]/\gMt.\qedhere\]
\end{proof}

% By the result (or rather the method)
% of~\cite{MR2745530} we have
% \[\xymatrix{& (alg,ft)\ar[dl]\ar[r]&\cC_a^{\le
%       h}\ar[dl]^{proper}\ar[d]^{proper}\\ \cX_a\ar[r]^{f.t., rep}&\cX_{s,a}\ar[r]^{sep,ind-rep}&\cR_a} \]
% \begin{itemize}
% 	\item Using a reinterpretation of Caruso--Liu in terms of $(\varphi,
% 		G_K)$-modules, we show that (for some $s$ depending on
% 		$a$ and $h$) the right-side arrow of the parallelogram exists.
% 		We want to show that this map is proper (given
% 		that the vertical arrow on the extreme right is known
% 		to be proper)
% \item To check the properness of the previous point,
% 	the key thing is the properties of the lower right
%   arrow; namely, its diagonal should be a closed immersion (equivalently,
%  a proper monomorphism). 
%  This is basically by hand, using the fact that
% % : separability is saying that the
% %  diagonal is proper, and that's not so hard, because in fact it's a
% %  closed immersion:
% endomorphisms of the
%   $G_K$-representation are the same as endomorphisms of the
%   $G_{K_\infty}$-representation which commute with~$G_K$. 
% \item Once we have the desired properness,
% we can pullback over $\cX_a$ and obtain a finite-type alg.\ stack
% with a proper morphism to $\cX_a$.  (Here we use the fact 
% that $\cX_a \to \cX_{s,a}$ is finite type and representable,
% which is easy to see when thought of as restricting $(\varphi,\Gamma)$-modules
% to $(\varphi,\Gamma_s)$-modules.)
% Its scheme-theoretic image is then a closed finite-type alg.\ substack
% of $\cX_a$.  We now check, for any crystalline versal ring
% $R_{\crys}^{\le h}$,
% that the morphism $\Spf R^{\leq h}_{\crys} \to \cX_a$ factors through 
% this scheme-theoretic image, and so is effective.
% \end{itemize}



% \begin{itemize}
% \item Discuss compatibility with Galois side.
% \item Should probably cite 1-ification material from
%   \cite{EGstacktheoreticimages}.
% \item Want to show it's separated and ind-representable.
% \end{itemize}

We end this section with the following lemma, which will be used in
the proof of Proposition~\ref{prop: versal rings for pst stacks}.

\begin{lem}
  \label{lem: algebraising C}Let~$R$ be a complete local
  Noetherian~$\cO$-algebra with residue field~$\F$, % \mar{TG: what
    % should the hypotheses be? I think this is probably what Mark is
    % using although he doesn't say Noetherian.}
  together with a morphism $\Spf
  R\to\cX_{K,d}$, and
  let~$\widehat{C}:=\cC_{d,\semis,h}^{L/K}\times_{\cX_{K,d}}\Spf R$. Then
  there is a projective morphism of schemes $C\to\Spec R$ whose
  $\m_R$-adic completion is isomorphic to~$\widehat{C}$.
\end{lem}
\begin{proof}
Note that since $\cC_{d,\semis,h}^{L/K}\to{\cX_{K,d}}$ 
  is proper and representable by algebraic spaces, 
  we see that $\widehat{C}$ is a formal algebraic space,
  and  the morphism $\widehat{C}\to\Spf R$ is proper and representable
  by algebraic spaces. In particular, $\widehat{C}$ is a proper
  $\m_R$-adic formal algebraic space over~$\Spf R$.
%R_{\rhobar}^{\square,\cO'}$.

  To show that we can algebraize~$\widehat{C}$, we will use an
  argument of Kisin, see \cite[Prop.\ 2.1.10]{KisinModularity}, which
  proves a similar statement for moduli of Breuil--Kisin modules. In
  order to use this, we show that we can
  realise~$\cC_{d,\semis,h}^{L/K}$ as a closed substack of a product
  of moduli stacks of Breuil--Kisin modules.

Let~$\pi,\pi'$ be two choices of uniformizers of~$L$, chosen such
that~$\pi/\pi'$ is not a $p$th power, and choose compatible
systems of $p$-power roots~$\piflat,(\pi')^\flat$. Note that the
closure of the subgroup of~$G_L$ generated by~$G_{L_{\piflat,\infty}}$
and~$G_{L_{(\pi')^\flat,\infty}}$ is just~$G_L$, because its fixed
field is contained in $L_{\piflat,\infty}\cap
L_{(\pi')^\flat,\infty}=L$, so that a continuous action of~$G_L$ is
determined by its restrictions to~$G_{L_{\piflat,\infty}}$
and~$G_{L_{(\pi')^\flat,\infty}}$.

Suppose firstly that~$L=K$. We have a natural morphism \numequation\label{eqn: L equals K mapping to two
  uniformizers}\cC_{d,\semis,h}\to\cC_{\piflat,d,h}\times_{\Spf \cO}
\cC_{(\pi')^\flat,d,h},\end{equation} which we claim is a  monomorphism. To see this, it is enough to note
that~$\gMt=\AAinf{A}\otimes_{\gS_{\piflat,A}}\gM_{\piflat}$ is
determined by~$\gM_{\piflat}$, while the $G_K$-action on~$\gMt$ is
determined by the actions of ~$G_{K_{\piflat,\infty}}$
and~$G_{K_{(\pi')^\flat,\infty}}$, which are in turn determined by the
conditions that they act trivially on~$\gM_{\piflat}$
and~$\gM_{(\pi')^\flat}$ respectively.


We may factor \eqref{eqn: L equals K mapping to two uniformizers} as
the composite
\[\cC_{d,\semis,h}\to(\cC_{\piflat,d,h}\times_{\Spf \cO}
  \cC_{(\pi')^\flat,d,h})\times_{(\cR_{\piflat,d}\times_{\Spf \cO}\cR_{(\pi')^\flat,d})}\cX_{{K,d}}\to\cC_{\piflat,d,h}\times_{\Spf \cO}
  \cC_{(\pi')^\flat,d,h}.\] Since the composite is a monomorphism, the
first morphism is a monomorphism.
This first morphism is also proper
(since the morphism $\cC_{d,\semis,h} \to \cX_{K,d}$ obtained by
composing it with the projection to $\cX_{K,d}$ is proper,
by Theorem~\ref{thm: semi stable stack is formal algebraic}, while 
this projection is separated, being a base-change of the product of
the morphisms $\cC_{(\pi)^\flat,d,h}\to\cR_{\piflat,d}$ and
  $\cC_{(\pi')^\flat,d,h}\to\cR_{(\pi')^\flat,d}$, each
of which is proper, and so in
  particular separated, by Theorem~\ref{thm: C is a $p$-adic formal algebraic stack and related
    properties}) % \mar{ME: Maybe need
% to say slightly more here?  The basic idea is that if $f\circ g$ is proper than
% $g$ is, under mild hypotheses on $f$, maybe separatedness? TG: I guess
% that's exactly the hypothesis we need, so if we unwind it should come
% down to some maps $\cC\to \cR$ being separated, hopefully? (I have to
% run now though). TG: how about this?}
%(by Corollary~\ref{cor: basic properties of C and R p adic stacks}),
and so it is in fact a closed immersion. 
% where the second morphism is separated,
% being the base change of the
% morphism~$\cX_d\to\cR_{\piflat,d}\times\cR_{(\pi')^\flat,d}$, whose
% diagonal is a closed immersion by Proposition~\ref{prop: diagonal of X
%   to R}. Since the composite is a closed immersion, it follows that
% the first morphism is also a closed immersion (for example, because
% closed immersions are proper monomorphisms). 
% We claim that~\eqref{eqn: L
% equals K mapping to two uniformisers} is a closed immersion.  That~\eqref{eqn: L equals K
%   mapping to two uniformisers} is furthermore a closed immersion can
% then be proved exactly as in the proof of Proposition~\ref{prop: BKF
%   closed in BK}.\mar{TG: not obviously correct! But it's no problem,
%   we get that the first map below is a monomorphism, and it's also
%   proper because  $\cC_{d,\semis,h}\to\cX_d$ is proper and $\cC\to\cR$ is
%   separated (even proper), and so it's a closed immersion, and that's
%   all we need.}
Thus~$\widehat{C}$ is a
closed algebraic subspace of
\[(\cC_{\piflat,d,h}\times_{\cR_{\piflat,d}}\Spf
  R)\times_{\Spf  R}(\cC_{(\pi')^\flat,d,h}\times_{\cR_{(\pi')^\flat,d}}\Spf
  R).\] By the Grothendieck Existence theorem for algebraic
spaces~\cite[Thm.\ V.6.3]{MR0302647} (and the symmetry
between~$\piflat$ and~$(\pi')^\flat$), we are therefore reduced to
showing that the morphism
$\cC_{\piflat,d,h}\times_{\cR_{\piflat,d}}\Spf R\to \Spf R$ can be
algebraized to a projective morphism. In the case~$h=1$, this
is~\cite[Prop.\ 2.1.10]{KisinModularity} (bearing in mind the main
result of~\cite{MR1320381}, as in the proof of~\cite[Prop.\
2.1.7]{KisinModularity}),
and the case of general~$h$ can be proved in
exactly the same way, as explained in the proof of~\cite[Prop.\
1.3]{MR2827797}; the key point is that
$\cC_{\piflat,d,h}\times_{\cR_{\piflat,d}}\Spf R$ inherits a natural very
ample (formal) line bundle from the affine Grassmannian.

We now consider the case of a general finite Galois
extension~$L/K$, where we argue as in the proof of Proposition~\ref{lem: restriction from X to Xs}. Let~$\{g_i\}_{i=1,\ldots,n}$ be a set of coset representatives
for~$G_L$ in~$G_K$. % \mar{TG: here and above, specify whether left or
  % right!}
By definition, to give an object of~$\cC_{d,\semis,h}^{L/K}(A)$ (for some~$A$) is
the same as giving an object of~$\cC_{d,\semis,h}^{L/L}(A)$, together
with an extension of the action of~$G_L$ on the underlying Breuil--Kisin--Fargues
module~$\gMt_A$ to an action of~$G_K$. Now, for each~$i$ we can
give~$g_i^*\gMt_A$ the structure of a Breuil--Kisin--Fargues module by
letting~$h\in G_L$ act as~$g_i^{-1}hg_i$ acts on~$\gMt_A$; it follows
from the definitions that $g_i^*\gMt_A$ is also an object
of~$\cC_{d,\semis,h}^{L/L}(A)$. Then to give an extension of the
action of~$G_L$ on~$\gMt_A$ to an action of~$G_K$ is to give for
each~$i$ an isomorphism $g_i^*\gMt_A\iso\gMt_A$ of objects
of~$\cC_{d,\semis,h}^{L/L}(A)$, satisfying a slew of compatibilities.

We let $\cY_i$ denote the stack classifying objects $\gMt_A$ 
of~$\cC_{d,\semis,h}^{L/L}(A)$, 
endowed with an isomorphism 
$g_i^*\gMt_A\to\gMt_A$.  If we regard $g_i^*$  as an automorphism
of~$\cC_{d,\semis,h}^{L/L}(A)$,
then we  may form  its  graph~$\Gamma_i$,
and we then have an isomorphism of stacks
$$\cY_i \iso
\cC_{d,\semis,h}^{L/L}(A)\times_{\Delta,
\cC_{d,\semis,h}^{L/L}(A) \times
\cC_{d,\semis,h}^{L/L}(A),\Gamma_i}
\cC_{d,\semis,h}^{L/L}(A).$$
(Here $\Delta$ denotes the diagonal;
\emph{cf.}\ the definition of $\cR_d^{\Gamma_{\disc}}$ in 
Section~\ref{subsec: defn of Xd}
above, together 
with Proposition~\ref{prop:Gamma-disc modules as fixed points; etale case}.)
The projection onto the second factor $\cY_i \to 
\cC_{d,\semis,h}^{L/L}(A)$,
which corresponds to
forgetting the isomorphism, is a base-change of the
diagonal~$\Delta,$
and so is affine,
since~$\Delta$ is affine by Theorem~\ref{thm: semi stable stack is formal algebraic}.
In particular,
$\cY_i\times_{\cX_{L,d}}\Spf R$ admits a natural ample formal line
bundle, pulled back from the natural very ample formal line bundle on
$\cC_{d,\semis,h}^{L/L}\times_{\cX_{L,d}}\Spf R$, whose existence we
established above.

We can rephrase our interpretation of objects 
of~$\cC_{d,\semis,h}^{L/K}(A)$ %(for some~$A$)
as objects of $\cC_{d,\semis,h}^{L/L}(A)$
endowed with isomorphisms
$g_i^*\gMt_A\to\gMt_A$ satisfying compatibilities
as the existence of  a closed immersion
%\numequation\label{eqn: semistable in inertia stack product
$$
\cC_{d,\semis,h}^{L/K}\into
\cY_1 \times_{\cC_{d,\semis,h}^{L/L}}\times \cdots \times_{\cC_{d,\semis,h}^{L/L}}
\cY_n
$$  %\end{equation}
(the compatibilities that cut out 
$\cC_{d,\semis,h}^{L/K}$ arise as
base-changes of the double diagonal, and so impose closed conditions).
%
%These various compatibilities are closed conditions, so that
%if~$\cI^{L/L}$ denotes the inertia stack of~$\cC_{d,\semis,h}^{L/L}$,
%we see that there is a closed immersion
%\numequation\label{eqn: semistable in inertia stack product}\cC_{d,\semis,h}^{L/K}\into
%\cI^{L/L} \times_{\cC_{d,\semis,h}^{L/L}}\times \cdots \times_{\cC_{d,\semis,h}^{L/L}}
%\cI^{L/L}\end{equation}
%(where the right-hand contains one copy of $\cI^{L/L}$ for each coset
%representative~$g_i$). Now, we
%have
%\[\cI^{L/L}=\cC_{d,\semis,h}^{L/L}\times_{\cC_{d,\semis,h}^{L/L}\times
%    \cC_{d,\semis,h}^{L/L}}\cC_{d,\semis,h}^{L/L},\]and
%since~$\cC_{d,\semis,h}^{L/L}$ has affine diagonal by Theorem~\ref{thm: semi stable stack is formal algebraic}, % \mar{TG: this
%  % doesn't seem to be stated anywhere right now, so we should put it in
%  % above; it will follow from Corollary~\ref{cor: basic properties of C
%  %   and R p adic stacks}~(1).}
%$\cI^{L/L}$ is affine
%over~$\cC_{d,\semis,h}^{L/L}$, and therefore
%$\cI^{L/L}\times_{\cX_{L,d}}\Spf R$ admits a natural ample formal line
%bundle, pulled back from the natural very ample formal line bundle on
%$\cC_{d,\semis,h}^{L/L}\times_{\cX_{L,d}}\Spf R$, whose existence we
%established above.
%\[\cC_{d,\semis,h}^{L/K}\into
%  \cC_{d,\semis,h}^{L/L}\times\prod_{i}\cI^{L/L}. \]
%It follows from~\eqref{eqn: semistable in inertia stack product} that
It follows that $\cC_{d,\semis,h}^{L/K}\times_{\cX_{L,d}}\Spf R$ inherits a natural ample formal line
%\mar{ME: Unfortunately, we need a base-change over $\cX_{K,d}$ here!
%So maybe we need to also describe $\cX_{K,d}$ as being closed
%inside some analogous $n$-fold fibre product over $\cX_{L,d}$ --- which
%is more-or-less what happens in the proof of Lemma~\ref{lem:
%  restriction from X to Xs}. TG: proposal is just to apply that lemma
%and deduce the fibre product over ~$K$ from that over~$L$. TG: had a
%go at doing this. ME: Looks good.  Added briefest comment on why
%the imposed conditions are closed, and am now removing this comment!}
bundle. As noted in the proof of Lemma~\ref{lem: restriction from X to
  Xs}, since the morphism $\cX_{K,d}\to\cX_{L,d}$ is affine, it  has affine
diagonal, so that the natural morphism
\[\cC_{d,\semis,h}^{L/K}\times_{\cX_{K,d}}\Spf
  R\to\cC_{d,\semis,h}^{L/K}\times_{\cX_{L,d}}\Spf R \]is affine. It
follows that that $\cC_{d,\semis,h}^{L/K}\times_{\cX_{K,d}}\Spf R$ in
turn inherits a natural ample formal line bundle (see e.g.\
\cite[\href{https://stacks.math.columbia.edu/tag/0892}{Tag
  0892}]{stacks-project}), and the  result then follows from another application of \cite[Thm.\
V.6.3]{MR0302647}.% again, we are reduced to showing that
% the morphisms \[\cI^{L/L}\times_{\cX_{L,d}}\Spf R\to\Spf R\] can
% be
%\mar{TG: no, we rather have the line bundles on the RHS and pull
%  them back to the LHS TG: OK now? ME: Looks good, commenting out}
%algebraized. 
\end{proof}

\section{Inertial types}
\label{subsec:inertial types}% \mar{TG: rewrite this to just be over $E$
  % from the beginning - there are issues of mixed types!}
  In this section we examine how to extract the inertial type of the
  Galois representation associated to a Breuil--Kisin--Fargues
  $G_K$-module; in Section~\ref{subsec: HT
  weights} we study the same problem for Hodge--Tate weights, and in
Section~\ref{subsec: moduli stacks of pst GK repns} we use these
results to define our moduli stacks of potentially semistable and
potentially crystalline representations of fixed inertial and Hodge
type.

In contrast to the rest of the book, in this and the following
sections we will need to consider Kisin modules with coefficients
with~$p$ inverted. To this end, we will write ~$\Acirc$ and~$\Bcirc$
for $p$-adically complete flat $\cO$-algebras which are topologically
of finite type over~$\cO$, and
write~$A=\Acirc[1/p]$, $B=\Bcirc[1/p]$. We hope that this change of
notation will not cause any confusion. We will freely use that
if~$\Acirc$ is topologically of finite type over~$\cO$, then
$A:=\Acirc[1/p]$ is Noetherian and Jacobson (\cite[\S0, Prop.\ 9.3.2,
9.3.10]{MR3752648}), and the residue fields of the maximal ideals
of~$A$ are finite extensions of~$K$ (\cite[\S0, Cor.\
9.3.7]{MR3752648}).
  % \begin{rem}\label{rem: affinoid algebra facts}
  %   We will freely use the following standard facts
  %   about affinoid algebras, for which~\cite{MR3752648} is a
  %   convenient reference. Recall that if~$\Acirc$ is a $p$-adically
  %   complete $\cO$-algebra, then~$\Acirc$ is topologically of finite
  %   type over~$\cO$ if and only if~$\Acirc\otimes_{\cO}k$ is a finite
  %   type $k$-algebra (\cite[\S0, Prop.\
  %   8.4.2]{MR3752648}). If~$\Acirc$ is topologically of finite type
  %   over~$\cO$, then $A:=\Acirc[1/p]$ is Noetherian and Jacobson
  %   (\cite[\S0, Prop.\ 9.3.2, 9.3.10]{MR3752648}), and the residue
  %   fields of the maximal ideals of~$A$ are finite extensions of~$K$
  %   (\cite[\S0, Cor.\ 9.3.7]{MR3752648}).
  % \end{rem}

Let~$L/K$ be a finite Galois
extension with inertia group~$I_{L/K}$, and %  and let $\Rep_{I_{L/K},d}$ be
% the algebraic stack over~$\Spec\cO$ classifying $d$-dimensional
% representations of~$I_{L/K}$. More precisely, if~$X$ is an
% $\cO$-scheme, then $\Rep_{I_{L/K},d}(X)$ is the groupoid of locally
% free sheaves of $\cO_X$-modules of rank~$d$, equipped with an
% $\cO_X$-linear action of~$I_{L/K}$.
%
% We 
suppose now that~$E$ is large enough that it contains the images of
all embeddings $L\into\Qpbar$, that all irreducible
$E$-representations of~$I_{L/K}$ are absolutely irreducible, and that
every irreducible~$\Qpbar$-representation of~$I_{L/K}$ is defined
over~$E$. Write~$l$ for the residue field of~$L$, and
write~$L_0=W(l)[1/p]$. % Then $\Rep_{I_{L/K},d}\times_{\Spec\cO}\Spec E$ admits the
% following simple description (see also~\cite[Lem.\
% 3.2.2]{CEGSKisinwithdd} for an analogous statement
% for~$\Rep_{I_{L/K},d}$ itself in the case that~$I_{L/K}$ has order
% prime to~$p$, with an almost identical proof).

% \begin{lem}
%   \label{lem: generic fibre of inertial type stack}$\Rep_{I_{L/K},d}\times_{\Spec\cO}\Spec E$
%   is the disjoint union of irreducible components
%   $\Rep_{I_{L/K},d}(\tau)_E$, where the~$\tau$ run over the isomorphism
%   classes of semisimple $E$-representations
%   of~$I_{L/K}$. Writing~$G_\tau:=\Aut_{E[I_{L/K}]}(\tau)$ \emph{(}a
%   reductive, hence smooth, group scheme over~$E$\emph{)}, we have
%   $\Rep_{I_{L/K},d}(\tau)_E=[\Spec E/G_\tau]$, so that in particular a
%   morphism $X\to\Rep_{I_{L/K},d}\times_{\Spec\cO}\Spec E$ factors
%   through~$\Rep_{I_{L/K},d}(\tau)_E$ if and only if the corresponding
%   locally free $I_{L/K}$-sheaf on~$X$ is Zariski locally isomorphic to~$\tau\otimes_E\cO_X$.
% \end{lem}
% \begin{proof}Since~$E$ has characteristic zero, the
%   representation~$P:= \oplus_\sigma\sigma$ is a projective generator
%   of the category of $E[I_{L/K}]$-modules, where~$\sigma$ runs over a set of
%   representatives for the isomorphism classes of irreducible
%   $E$-representations of~$I_{L/K}$.  Our assumption that each $\sigma$ is
%   absolutely irreducible furthermore shows that
%   $\End_{I_{L/K}}(\sigma) = E$ for each~ $\sigma$, so that
%   $\End_{I_{L/K}}(P) = \prod_{\sigma} E$.

% Standard Morita theory then shows that the functor
% $M \mapsto \Hom_{I_{L/K}}(P,M)$ induces an equivalence between the category
% of $E[I_{L/K}]$-modules and the category of $\prod_{\sigma} E$-modules.
% Of course, a $\prod_{\sigma}E$-module is just given by a tuple
% $(N_{\sigma})_{\sigma}$ of $E$-vector spaces,
% and in this optic, the functor $\Hom_{I_{L/K}}(P,\text{--})$ 
% can be written as $M \mapsto \bigl(\Hom_{I_{L/K}}(\sigma,M)\bigr)_{\sigma}$,
% with a quasi-inverse functor being given by $(N_{\sigma}) \mapsto
% \bigoplus_{\sigma} \sigma\otimes_{E} N_{\sigma}.$
% It is easily seen (just using the fact that $\Hom_{I_{L/K}}(P,\text{--})$
% induces an equivalence of categories)
% that $M$ is a finitely generated projective $A$-module,
% for some $E$-algebra $A$,
% if and only if each $\Hom_{I_{L/K}}(\sigma,M)$
% is a finitely generated projective $A$-module.

% %The first two parts of the lemma are then straightforward consequences
% %of this discussion.
% %of Morita theory, together with the observation
% %that~$\End_G(\tau^\circ)=\cO$ (by Schur's lemma). Indeed we see that
% %for any~$\cO[G]$-module~$M$ the natural
% %map \[\oplus_\tau\Hom_G(\tau^\circ,M)\otimes_{\cO}\tau^\circ\to M\]is
% %an isomorphism.\mar{TG: slightly lost on what I should be writing
% %  here}

% The preceding discussion shows that
% giving a rank $d$ representation of $I_{L/K}$ over an $E$-algebra $A$
% amounts to giving a tuple $(N_{\sigma})_{\sigma}$ of projective
% $A$-modules, of ranks~$n_{\sigma}$, such that $\sum_{\sigma} n_{\sigma}
% \dim \sigma = d.$
% For each such tuple of ranks $(n_\sigma)$, we obtain a corresponding
% moduli stack
% $\Rep_{(n_{\sigma})}(I_{L/K})$ classifying rank $d$ representations of $I_{L/K}$
% which decompose in this manner,
% and $\Rep_d(I_{L/K})$ is isomorphic to the disjoint union
% of the various stacks~$\Rep_{(n_{\sigma})}(I_{L/K}).$

% If we write $\tau = \oplus_{\sigma} \sigma^{n_{\sigma}}$, 
% then we may relabel $\Rep_{(n_{\sigma})}(I_{L/K})$ as~$\Rep_{I_{L/K},d}(\tau)_E$.
% By construction, there is an isomorphism  
% $$\Rep_{I_{L/K},d}(\tau)_E= \Rep_{(n_{\sigma})}(I_{L/K}) \iso \prod_{\sigma}
% [\Spec E/\GL_{n_{\sigma}}].$$
% Noting that $G_{\tau} := \Aut(\tau) = 
% \prod_\sigma\GL_{n_\sigma/E}$, 
% we are done.
% \end{proof}% \mar{TG: probably want more details, but iirc ME plans to
%   rewrite the proof in \cite{CEGSKisinwithdd}, so it doesn't make much
% sense for me to copy that out!}
% \begin{defn}
%   \label{defn: tau inertia rep stack integral via closure}We
%   let~$\Rep_{I_{L/K},d}(\tau)$ be the closure of
%   $\Rep_{I_{L/K},d}(\tau)_E$ in~$\Rep_{I_{L/K},d}(\tau)$.
% \end{defn}% \mar{TG: quite possibly we won't actually need to use this,
%   given that the Hodge--Tate weights are only being defined after
%   inverting~$p$, this could be too? ME: I think it's better to keep it
% this way.  Since we don't have a completely developed formalism for passing
% from $\cO$ to $E$ in formal-stack-world, it seems better to do it this way,
% where this passage happens on the algebraic stack $\Rep_{I_{L/K},d},$
% where it makes perfect sense within the current state of
% technology. TG: sounds good, commenting out.}

Let~$\gM_{\Acirc}$ be a Breuil--Kisin--Fargues $G_K$-module with~$\Acirc$-coefficients which
  admits all descents over~$L$, and write
  $\overline{\gM}_{\Acirc}:=\gM_{\Acirc,\piflat}/[\piflat]\gM_{\Acirc,\piflat}$ for the
  module considered in Definition~\ref{defn: descending BKF to BK
    coefficients}~\eqref{item: M mod u
    descends} (for some choice of~$\piflat$, with~$\pi$ a uniformizer
  of~$L$; note that by the definition of $\gM_{\Acirc}$ admitting
all descents over~$L$, the quotient $\overline{\gM}_{\Acirc}$
is actually well-defined as a $W(l)\otimes_{\Z_p}A^{\circ}$-submodule
of $W(\overline{k})_A\otimes_{\AAinf{A}} \gM_{\Acirc}$,
independently of the choice $\pi$ or~$\piflat$).
Then~$\overline{\gM}_{\Acirc}$ has a natural $W(l)\otimes_{\Zp}A$-semilinear
  action of~$\Gal(L/K)$, which is defined as follows: if~$g\in \Gal(L/K)$,
  then~$g(\gM_{\Acirc,\piflat})=\gM_{\Acirc,g(\piflat)}$ (see the proof of
  Lemma~\ref{lem: descent for GK action depends only on pi}), so the
  morphism $g:\gM_{\Acirc,\piflat}\to g(\gM_{\Acirc,\piflat})=\gM_{\Acirc,g(\piflat)}$
  induces a morphism \[g: \gM_{\Acirc,\piflat}/[\piflat]\gM_{\Acirc,\piflat} \to
    \gM_{\Acirc,g(\piflat)}/[g(\piflat)]\gM_{\Acirc,g(\piflat)},\]and the source and
  target are both
  canonically identified with~$\overline{\gM}_{\Acirc}$.

% \mar{TG: I don't think what's written here is what we mean. Don't we
%   mean to just take the action that we just wrote on
%   $\overline{\gM}_{\Acirc}$ and just invert~$p$?}
  % Write~$K_0=W(k)[1/p]$.  
  This action of~$\Gal(L/K)$ on~$\overline{\gM}_{\Acirc}$ induces an
  $L_0\otimes_{\Qp}A$-linear action of $I_{L/K}$ on the projective
  $L_0\otimes_{\Qp}A$-module
  $\overline{\gM}_{\Acirc}\otimes_{\Acirc}A$. Fix a choice of
  embedding $\sigma:L_0\into E$, and
  let~$e_\sigma\in L_0\otimes_{\Qp} E$ be the corresponding
  idempotent. Then
  $e_\sigma(\overline{\gM}_{\Acirc}\otimes_{\Acirc}A$) is a projective
  $A$-module of rank~$d$, with an $A$-linear action
  of~$I_{L/K}$. (Indeed, since~$\gM_{\Acirc}$ is a finite
  projective~$\gS_A$-module of rank~$d$, $\overline{\gM}_{\Acirc}$ is
  a finite projective $W(k)\otimes_{\Zp}\Acirc$-module of rank~$d$,
  and $\overline{\gM}_{\Acirc}\otimes_{\Acirc}A$ is a finite
  projective $L_0\otimes_{\Qp}A$-module of rank~$d$. It follows that
  $e_\sigma(\overline{\gM}_{\Acirc}\otimes_{\Acirc}A)$ is a finite
  projective $A$-module of rank~$d$.)

  Note that up to canonical isomorphism, this module does
  not depend on the choice of~$e_\sigma$: indeed the induced action
  of~$\varphi$ commutes with~$I_{L/K}$ and induces isomorphisms
  between the 
  $e_\sigma(\overline{\gM}_{\Acirc}\otimes_{\Acirc}A)$ (with~$\sigma$
  varying), because the cokernel of~$\varphi$ is killed
  by~$E_{\piflat}$, which is a unit in our setting (in
  which~$[\piflat]=0$ and~$p$ is a unit).

  \begin{defn}
    \label{defn: WD representation associated to BKF
      module}Let~$\gMt_{\Acirc}$ be as above. Then we write \index{$\WD(\gMt_{\Acirc})$}
    \[\WD(\gMt_{\Acirc}):=e_\sigma(\overline{\gM}_{\Acirc}\otimes_{\Acirc}A),\]
    a projective~$A$-module of rank~$d$ with an $A$-linear action of $I_{L/K}$.
  \end{defn}
% \mar{TG: following paragraph added to be cited in 4.7.12. ME: Looks good,
% except that shouldn't $V$ and $T$ be applied to $M$ rather than~$\gM$? (Also,
% is it  actually cited in 4.7.12?) TG: unsure about how to cite it so
% this comment was left in. I don't think there was an~$M$, but I've
% introduced one, and have added an implicit citation.}

\begin{remark}
\label{rem:inertial types point}
  The point of this definition is that if~$\Acirc$ is a finite flat
  $\cO$-algebra, then it computes inertial types in the following
  sense:  writing~$A:=\Acirc[1/p]$, and~$M_{\Acirc}:=W(\C^\flat)_{\Acirc}\otimes_{\AAinf{\Acirc}}\gMt_{\Acirc}$, we have a potentially semistable representation of~$G_K$ on
  a free $A$-module given by
  $V_A(M_{\Acirc}):=T_{\Acirc}(M_{\Acirc})\otimes_{\Acirc}A$.
As explained in Section~\ref{subsec: Hodge and inertial types}, the
  inertial type $D_{\pst}(V_A(M_{\Acirc}))|_{I_K}$ is then given
  by~$\WD(\gMt_{\Acirc})$ with its action of~$I_{L/K}$.
\end{remark}

  \begin{prop}
  \label{prop: good behaviour of inertial type}Let~$\Acirc$ be a $p$-adically complete flat
$\cO$-algebra which is topologically of finite type over~$\cO$, and write~$A=\Acirc[1/p]$. Let~$\gMt_{\Acirc}$
and~$\gM_{\Acirc}$ be as
above, and fix~$\sigma:L_0\into E$. Then~$\WD(\gMt_{\Acirc})$ % \mar{TG: have to
  % explain what this is! Maybe $\varphi$ should be something like~$\sigma$}
is a finite
  projective $A$-module of rank~$d$ with an action of~$I_K$, whose formation is
  compatible with base changes $\Acirc\to \Bcirc$ of $p$-adically complete flat
$\cO$-algebras which are topologically of finite type over~$\cO$.%\mar{TG: again some finiteness condition}
\end{prop}
\begin{proof}Writing~$u$ for~$[\piflat]$, the compatibility of
  formation with base change reduces to the observation
  that the natural map $\gS_A\otimes_AB\to\gS_B$ induces an isomorphism \[\gS_B/u\gS_B \cong (\gS_A/u\gS_A)\otimes_AB. \qedhere\]
\end{proof}
Let~$\tau$ be a $d$-dimensional $E$-representation of~$I_{L/K}$.
\begin{defn}%\mar{TG: definition should just be Zariski local on~$A$.}
  \label{defn: having inertial type tau}In the setting of
  Proposition~\ref{prop: good behaviour of inertial type}, we say that
  $\gMt_{\Acirc}$ % \mar{TG: or perhaps~$\gMt_{\Acirc}$
    % itself has type~$\tau$?}
  has inertial type~$\tau$ \index{inertial type}
  if Zariski locally on~$\Spec A$, $\WD(\gMt_{\Acirc})$ is isomorphic to the base change of~$\tau$
  to~$A$.
\end{defn}

\begin{cor}
  \label{cor:inertial type decomposition}In the setting of
  Proposition~{\em \ref{prop: good behaviour of inertial type}}, we  can
  decompose~$\Spec A$ as the disjoint union of open and closed
  subschemes~$\Spec A^\tau$, where~$\Spec A^\tau$ is the locus over
  which $\gMt_{\Acirc}$ has inertial type~$\tau$. Furthermore, the
  formation of this decomposition is compatible with base changes $\Acirc\to \Bcirc$ of $p$-adically complete flat
$\cO$-algebras which are topologically of finite type over~$\cO$.
\end{cor}
\begin{proof}
  By Proposition~\ref{prop: good behaviour of inertial type} (more
  precisely, by the compatibility with base change), we can define
  $\Spec A^\tau$ to be the locus over which $\gMt_{\Acirc}$
  has inertial type~$\tau$. That~$\Spec A$ is actually the
  disjoint union of the $\Spec A^\tau$ follows easily from our
  assumptions on~$E$, and standard facts about the representation
  theory of finite groups in characteristic zero. For lack of a
  convenient reference, we sketch a proof as follows.

  Since~$E$ has characteristic zero, the
  representation~$P:= \oplus_rr$ is a projective generator
  of the category of $E[I_{L/K}]$-modules, where~$r$ runs over a set of
  representatives for the isomorphism classes of irreducible
  $E$-representations of~$I_{L/K}$.  Our assumption that~$E$ is large
  enough that each $r$ is
  absolutely irreducible furthermore shows that
  $\End_{I_{L/K}}(r) = E$ for each~ $r$, so that
  $\End_{I_{L/K}}(P) = \prod_{r} E$.

Standard Morita theory then shows that the functor
$M \mapsto \Hom_{I_{L/K}}(P,M)$ induces an equivalence between the category
of $E[I_{L/K}]$-modules and the category of $\prod_{r} E$-modules.
Of course, a $\prod_{r}E$-module is just given by a tuple
$(N_{r})_{r}$ of $E$-vector spaces,
and in this optic, the functor $\Hom_{I_{L/K}}(P,\text{--})$ 
can be written as $M \mapsto \bigl(\Hom_{I_{L/K}}(r,M)\bigr)_{r}$,
with a quasi-inverse functor being given by $(N_{r}) \mapsto
\bigoplus_{r} r\otimes_{E} N_{r}.$
It is easily seen (just using the fact that $\Hom_{I_{L/K}}(P,\text{--})$
induces an equivalence of categories)
that $M$ is a finitely generated projective $A$-module,
for some $E$-algebra $A$,
if and only if each $\Hom_{I_{L/K}}(r,M)$
is a finitely generated projective $A$-module. Writing the various
~$\tau$ in the form~$\oplus_{r} r^{n_{r}}$, we are done.
\end{proof}
%   Accordingly, we have morphisms of stacks over $\cO$
% $$\cC_{d,\crys,h}^{L/K}\to \cC_{d,\ss,h}^{L/K} \to \Rep_{I_{L/K},d}$$
% taking~$\gMt$ to $\overline{\gM}$. We make the following
% definition of moduli stacks of Breuil--Kisin--Fargues modules of fixed
% inertial type.
% \mar{ME: Want to define this via $I_{L/K}$-action on $\gM/\pi^{\flat}\gM$,
% 	but I couldn't see from the definition that we actually had a
%         $G_K$-action after descending the BKF module, rather than just
%         a $G_L$-action. TG: see email}

% \begin{defn}\mar{TG: now want to treat as we do the Kisin ones, by
%     taking flat parts first. Also have to put in some $\varphi$-H90
%     stuff to explain that it's well-defined, maybe taking
%     $\varphi$-invariants is the best way  to do this.}
%   \label{defn: fixed inertial type Kisin stacks}We set
%   $\cC_{d,\ss,h}^{L/K,\tau}:=\cC_{d,\ss,h}^{L/K}\times_{\Rep_{I_{L/K},d}}\Rep_{I_{L/K},d}(\tau)$,
%   and
%   $\cC_{d,\crys,h}^{L/K,\tau}:=\cC_{d,\crys,h}^{L/K}\times_{\Rep_{I_{L/K},d}}\Rep_{I_{L/K},d}(\tau)$. By
%   Proposition~\ref{prop: pst and pcrys stacks are padic}, both
%   $\cC_{d,\ss,h}^{L/K,\tau}$ and $\cC_{d,\crys,h}^{L/K,\tau}$ are
%   $p$-adic formal algebraic stacks, and the morphisms
%   $\cC_{d,\ss,h}^{L/K,\tau}\to\cX_d$ and
%   $\cC_{d,\crys,h}^{L/K,\tau}\to\cX_d$ are representable by algebraic
%   spaces, proper, and of finite presentation.
% \end{defn}

% \mar{ME: Next, define substacks of
% 	$\cC_{d,\ss,h}^{L/K}$
% 	and
% 	$\cC_{d,\crys,h}^{L/K}$
% 	of fixed type $\tau$, via pull-back.  Have to choose notation, too!}


%\section{The flat part with Hodge and inertial types}\label{subsec: HT and crystalline in flat part}% \mar{TG:
  % approach to crystalline should be to work out how to interpret
  % $\gM/u\gM$ as Dieudonne module over the stack}
% \begin{itemize}
% \item Define flat part.
% \item Prove that flat part commutes with scheme-theoretic image (if we
%   can).
% \item Impose HT weights following Kisin.
% \end{itemize}
%\mar{ME: Change title}
\section{Hodge--Tate weights}\label{subsec: HT
  weights}Let~$\Acirc$ be a $p$-adically complete flat
$\cO$-algebra which is topologically of finite type over~$\cO$, and let % \mar{TG: not sure what the right hypotheses are
% - topologically of finite type? TG: this works and I don't see
% anything more general, so I've gone with it}
~$\gMt_{\Acirc}$ be a Breuil--Kisin--Fargues $G_K$-module of
height at most~$h$ with $\Acirc$-coefficients, which admits all
descents. % \mar{TG: might in theory want to do this over~$L$ but it
  % shouldn't matter, because of course it's OK to fix Hodge--Tate
  % weights after a base change.}
% Suppose that~$A$ is an $E$-algebra.\mar{TG: is that actually
  % allowed in our setup?! I think we have baked in an assumption of
  % being $p$-adically complete. Not going to worry about that for now -
  % maybe we actually want to work integrally but the invert~$p$, for
  % example.}
We now explain how to interpret the condition that~$\gMt_{\Acirc}$
has a fixed Hodge type, following~\cite[Lem.\ 2.6.1, Cor.\
2.6.2]{MR2373358} (but bearing in mind the corrections to these
results explained in~\cite[A.4]{KisinFM}). We would like to thank Mark
Kisin for a helpful conversation about these results, and for some
suggestions regarding the proof of Proposition~\ref{prop: good behaviour of Kisin
    filtration}.

  
Fix some choice of~$\piflat$, and write~$\gM_{\Acirc}$
for~$\gM_{\piflat,\Acirc}$, and~$u$ for~$[\piflat]$. For
each~$0\le i\le h$ we define
$\Fil^i\varphi^*\gM_{\Acirc}=\Phi_{\gM_{\Acirc}}^{-1}(E(u)^i\gM_{\Acirc})$,
and we set $\Fil^i\varphi^*\gM_{\Acirc}=\varphi^*\gM_{\Acirc}$
for~$i<0$.

\begin{remark}
\label{rem:Hodge filtration point}
The point of this definition is that it captures the Hodge
filtration. Indeed, if~$\Acirc$ is a finite flat $\cO$-algebra,
then writing~$A=\Acirc[1/p]$, and~$M_{\Acirc}:=W(\C^\flat)_{\Acirc}\otimes_{\AAinf{\Acirc}}\gMt_{\Acirc}$, we have a representation of~$G_K$ on
  a free $A$-module given by
  $V_A(M_{\Acirc}):=T_{\Acirc}(M_{\Acirc})\otimes_{\Acirc}A$. As
explained in Section~\ref{subsec: Hodge and inertial types} (or see
the proof of~\cite[Cor.\ 2.6.2]{MR2373358}), there is a natural
identification of~$\DdR(V_A(M_{\Acirc}))$
with~$(\varphi^*\gM_{\Acirc}/E(u)\varphi^*\gM)\otimes_{\Acirc}A$, under
which $\Fil^i\DdR(V_A(M_{\Acirc}))$ is identified with
\[(\Fil^i\varphi^*\gM_{\Acirc}/E(u)\Fil^{i-1}\varphi^*\gM_{\Acirc})\otimes_{\Acirc}A;\]
in particular, this latter module is a finite projective
$K\otimes_{\Qp}A$-module.
%$\cO_K\otimes_{\Zp}A$-module.
\end{remark}

The following proposition shows that the
final conclusion of the preceding remark
holds for more general choices of~$\Acirc$; the proof
uses that particular case (i.e.\ the case that $\Acirc$ is a finite flat $\cO$-algebra)
as an input.


% \begin{rem}\mar{TG: haven't thought through how this all works with
%     projectivity etc}
%   \label{rem: how to read off the Hodge type}Suppose that~$A$ is flat
%   over~$\cO$.\mar{TG: probably want this assumption?} Let~$\gMt$ be a
%   Breuil--Kisin--Fargues $G_K$-module which admits all descents
%   over~$L$. Write~$\gM'$ for the module
%   $\varphi^*\gM_{\piflat}/E_{\piflat}\varphi^*\gM_{\piflat}$
%   considered in Definition~\ref{defn: descending BKF to BK
%     coefficients}~\eqref{item: M mod u descends} (for some choice
%   of~$\piflat$, with~$\pi$ a uniformiser of~$L$).  We have a filtration on the
%   $L\otimes_{\Qp}A[1/p]$-module $D_L:=\overline{\gM'}[1/p]$, which is
%   defined as follows.  We write
%   $\Phi_{\gM_{\piflat}}:\varphi^*\gM_{\piflat}\to\gM_{\piflat}$ and
%   $f_\pi: \varphi^*\gM_{\piflat}[\frac 1 p] \twoheadrightarrow D_L $.
%   For each~$ i\geq 0$ we define
%   $\Fil^i\varphi^*\gM_{\piflat}=\Phi_{\gM_{\piflat}}^{-1}(E_{\piflat}^i\gM_{\piflat})$
% %We set $\Fil^i\varphi^*\gM_{\piflat}=\varphi^*\gM_{\piflat}$
% %for~$i<0$.  
% %Then for each $i$ we set
% and $\Fil^i D_L = f_\pi (\Fil ^i \varphi ^* \gM_{\piflat})$. 

% By Hilbert 90, $D_L$ and
% its filtration descend to a filtration on a
% $K\otimes_{\Qp}E$-module~$D_K$.\mar{TG: I didn't
%   check this! Does it follow from our assumptions to here? Also, is
%   the independence of~$\piflat$ clear?! In the appendix we use
%   $p$-adic Hodge theory to see this in the case of $\Qp$-coefficients.}  Then for each $\sigma:K\into E$, and
% each~$i\in [-h,0]$, the multiplicity of~$i$ in the multiset
% $\HT_{\sigma}(T(M))$ is the dimension of the $(-i)$-th graded piece of
% $e_{\sigma}D_K$, where $e_\sigma\in (K\otimes_{\Qp}E)$ is the
% idempotent corresponding to~$\sigma$.



\begin{prop}% \mar{TG: just guessing at the best things we could hope to
    % prove, final result may well need modification. For example, I'm
    % completely ignoring finite type hypotheses for now, so what is
    % written below is not a complete proof in the claimed generality.}
  \label{prop: good behaviour of Kisin
    filtration}Let~$\Acirc$ be a $p$-adically complete flat
$\cO$-algebra which is topologically of finite type over~$\cO$, and write~$A=\Acirc[1/p]$. Let~$\gMt_{\Acirc}$
and~$\gM_{\Acirc}$ be as above.  For each~$0\le i\le h$,
  \[(\Fil^i\varphi^*\gM_{\Acirc}/E(u)\Fil^{i-1}\varphi^*\gM_{\Acirc})\otimes_{\Acirc}A\] is a finite
  projective $K\otimes_{\Qp}A$-module, whose formation is
  compatible with base changes $\Acirc\to \Bcirc$ of $p$-adically complete flat
$\cO$-algebras which are topologically of finite type over~$\cO$.%\mar{TG: again some finiteness condition}
\end{prop}
\begin{proof}For any morphism of $\cO$-algebras $\Acirc\to \Bcirc$, we write~$\gM_{\Bcirc}:=\gS_{\Bcirc}\otimes_{\gS_{\Acirc}}\gM_{\Acirc}$. We begin by showing (following the proof of \cite[Lem.\
  2.6.1]{MR2373358}) that for any $p$-adically complete $\cO$-algebra
  $\Acirc$ which is topologically of finite type over~$\cO$, both $\gM_{\Acirc}/\im\Phi_{\gM_{\Acirc}}$ and $\varphi^*\gM_{\Acirc}/\Fil^h\varphi^*\gM_{\Acirc}$
  are finite projective $\cO_K\otimes_{\Zp}\Acirc$-modules, whose formation
  is compatible with arbitrary base changes $\Acirc\to \Bcirc$
  (with~$\Bcirc$ also topologically of finite type). % Note firstly that for any
  % integers $m\le n$, $E(u)^m\gM_{\Acirc}/E(u)^n\gM_{\Acirc}$ is a finite projective
  % $W(k)_{\Acirc}$-module, and that
  % $E(u)^m\gM_{\Bcirc}/E(u)^n\gM_{\Bcirc}=B\otimes_{\Acirc}
  % E(u)^m\gM_{\Acirc}/E(u)^n\gM_{\Acirc}$.\mar{TG: do we need to say anything to
  %   justify this? Here and below I've gotten a bit confused about
  %   where we want to be talking about $\cO_K\otimes_{\Zp}A$-modules
  %   and where~$W(k)_{\Acirc}$-modules so for now I've just written $W(k)_{\Acirc}$ everywhere.}

  Indeed, that~$\gM_{\Acirc}/\im\Phi_{\gM_{\Acirc}}$ is a finite projective
 $\cO_K\otimes_{\Zp}\Acirc$-module whose formation is compatible with base change
  follows as in the proof of~\cite[Lem.\ 2.6.1]{MR2373358} from the
  fact that~$\Phi_{\gM}$ remains injective after any base change
  (because~$E(u)$ remains a non-zero-divisor after any base change),
  and the result for $\varphi^*\gM_{\Acirc}/\Fil^h\varphi^*\gM_{\Acirc}$ follows
  from this and the short exact
  sequence \[0\to\varphi^*\gM_{\Acirc}/\Fil^h\varphi^*\gM_{\Acirc}\to \gM_{\Acirc}/E(u)^h\gM_{\Acirc}\to    \gM_{\Acirc}/\im\Phi_{\gM_{\Acirc}}\to 0.\]


 We claim that for each~$0\le i\le h$, 
  $(\Fil^h\varphi^*\gM_{\Acirc}/E(u)^i\Fil^{h-i}\varphi^*\gM_{\Acirc})\otimes_{\Acirc}A$ is a finite
  projective $K\otimes_{\Qp}A$-module whose formation is compatible with base
  change. Admitting the claim, the proposition follows from the short
  exact sequence% \mar{TG: sort this out; numbers are wrong in the
    % middle/right, probably an $i+1$ in the middle?}
%\mar{TG: still too wide}
  \nummultline\label{eqn: more Kisin
    fun}
  %\begin{split}
0\longrightarrow\Fil^i\varphi^*\gM_{\Acirc}/E(u)\Fil^{i-1}\varphi^*\gM_{\Acirc}
\\
\stackrel{E(u)^{h-i}}{\longrightarrow}
    \Fil^h\varphi^*\gM_{\Acirc}/E(u)^{h-i+1}\Fil^{i-1}\varphi^*\gM_{\Acirc}\\ \longrightarrow
    \Fil^h\varphi^*\gM_{\Acirc}/E(u)^{h-i}\Fil^{i}\varphi^*\gM_{\Acirc}\longrightarrow 0
%  \end{split}
\end{multline}(because after tensoring with~$A$, the second and third
terms are projective and compatible with base change, so that the
sequence splits, and the first term is also projective and compatible
with base change).

To prove the claim, we firstly consider for each~$0\le i\le h$ the finite $\cO_K\otimes_{\Zp}\Acirc$-module
  \[\varphi^*\gM_{\Acirc}/(E(u)^i\varphi^*\gM_{\Acirc}+\Fil^h\varphi^*\gM_{\Acirc}).\] Since this is
  the cokernel of the morphism
  \[E(u)^i\varphi^*\gM_{\Acirc}% /E(u)^h\varphi^*\gM_{\Acirc}
    \to\varphi^*\gM_{\Acirc}/\Fil^h\varphi^*\gM_{\Acirc}, \]we see that its formation is
  compatible with base change. We have a short exact sequence
of finite type $\Acirc$-modules
  \numequation\label{eqn: first Kisin nightmare}
  \begin{split}
    0\to \Fil^h\varphi^*\gM_{\Acirc}/E(u)^i\Fil^{h-i}\varphi^*\gM_{\Acirc}\to
    \varphi^*\gM_{\Acirc}/E(u)^i\varphi^*\gM_{\Acirc} \\ \to
    \varphi^*\gM_{\Acirc}/(E(u)^i\varphi^*\gM_{\Acirc}+\Fil^h\varphi^*\gM_{\Acirc})\to 0,
  \end{split}
\end{equation}in which the second and third terms are compatible with
base change, and we need to prove that after inverting~$p$, the first
term is projective and is of formation compatible with base change. To
this end, note firstly that as discussed above, if~$\Acirc$ is a
finite flat $\cO$-algebra, then
each~\[(\Fil^i\varphi^*\gM_{\Acirc}/E(u)\Fil^{i-1}\varphi^*\gM_{\Acirc})\otimes_{\Acirc}A\]
is a finite projective $K\otimes_{\Qp}A$-module, and it follows
from \eqref{eqn: more Kisin fun} and an easy induction that the same
is true
of~$(\Fil^h\varphi^*\gM_{\Acirc}/E(u)^i\Fil^{h-i}\varphi^*\gM_{\Acirc})\otimes_{\Acirc}A$. 
  
% Suppose now that~$B$ is a finite  $\Qp$-algebra. As in the proof
% of~\cite[Cor.\ 2.6.2]{MR2373358}, for each~$i$ the
% $\cO_K\otimes_{\Zp}A$-module
% $\Fil^i\varphi^*\gM_{\Bcirc}/E(u)\Fil^{i-1}\varphi^*\gM_{\Bcirc}$ is finite
% projective,\mar{TG: can probably chuck in a reference to something in
%   this paper once the appendix is complete?} so it follows by induction on~$i$ using the analogous
% short exact sequence to~\eqref{eqn: more Kisin
%     fun} for~$B$ that
%   $\Fil^h\varphi^*\gM_{\Bcirc}/E(u)^i\Fil^{h-i}\varphi^*\gM_{\Bcirc}$ is a finite
%   projective $W(k)_{\Bcirc}$-module.

To simplify notation, we write~\eqref{eqn: first Kisin nightmare} as
\[0\to K(\Acirc)\to M_1\to M_2\to 0,\] and for any $\Acirc$-algebra $C$,
we write $K(C)$ for the kernel of the surjection
$M_1\otimes_{\Acirc}C\to M_2\otimes_{\Acirc}C$. In particular
for~$\Bcirc$ as in the statement of the proposition, if as usual we write 
$B := \Bcirc[1/p]$, then we have
\[K(B)=(\Fil^h\varphi^*\gM_{\Bcirc}/E(u)^i\Fil^{h-i}\varphi^*\gM_{\Bcirc})\otimes_{\Bcirc}B;\]
so we need to show that~$K(A)$ is projective, and
that~$K(B)=K(A)\otimes_{A}
B$. % By the previous paragraph, if~$B$ is a finite $E$-algebra,
  % then ~$K(B)$ is a finite projective $W(k)_{\Bcirc}$-module.
  
  Let~$\m$ be a maximal ideal of~$A$, so that $A/\m^i$ is a finite
  $E$-algebra for each~$i$. % \mar{TG: needs a finiteness
    % hypothesis on~$A$! Affinoid algebra?} 
  Then since~$\widehat{A}_\m$ is a flat~$A$-algebra,
since the Artin--Rees Lemma shows that tensoring with
$\widehat{A}_\m$ coincides with  $\m$-adic completion for finite type modules,
and since the formation of kernels commutes with limits,  we see that
  \[K(A)\otimes_{\Acirc}\widehat{A}_\m=K(\widehat{A}_\m)=\varprojlim_iK(A/\m^i).\]
%(Note that we're using the Artin--Rees Lemma to know that tensoring with
%$\widehat{A}_\m$ coincides with  $\m$-adic completion for finite type modules.)
  Since ~$A/\m^i$ is a finite~$E$-algebra, % \mar{TG: do we need to
    % choose a lattice in it which is a finite flat $\Acirc$-algebra?}
  we see that~$K(A/\m^i)$ is a finite projective
  $(K\otimes_{\Qp}A)/\m^i$-module of rank bounded independently
  of~$i$. It follows that $K(A)\otimes_{\Acirc}\widehat{A}_\m$ is a finite
  projective~$\cO_K\otimes_\Zp{\widehat{A}_\m}$-module. Since~$A_\m$
  is Noetherian, % \mar{TG: for finiteness reasons}
  $\widehat{A}_\m$ is a
  faithfully flat~$A_\m$-algebra, so that $K(A)\otimes_{\Acirc}{A}_\m$ is a finite
  projective~ $\cO_K\otimes_\Zp A_\m$-module. Since this holds for
  all~$\m$, we see that~$K(A)$ is a finite
  projective~$\cO_K\otimes_\Zp A$-module, as claimed.% \mar{TG: do we want to
  % use something like: if~$A_\m$ is Noetherian then $A_\m\to
  % \widehat{A}_\m$ is faithfully flat, so $K(A)\otimes_{\Acirc}{A}_\m$
  % is finite projective, so~$K(A)$ is finite projective?}

  It remains to show that we have $K(B)=K(A)\otimes_{A} B$. We have a
  natural surjective morphism of finite projective
  $K\otimes_{\Qp}B$-modules \numequation\label{eqn: base change of
    K map}K(A)\otimes_A B\to K(B),\end{equation}which we need to show
is an isomorphism. Note firstly that if~$B=A/\m$ for some maximal
ideal~$\m$ of~$A$, then this follows from the previous paragraph. Now
suppose that~$B$ is general. Since the kernel of~\eqref{eqn: base
  change of K map} is in particular a finite projective $B$-module, is
enough to prove that for any maximal ideal~$\m_B$ of~$B$, \eqref{eqn:
  base change of K map} becomes an isomorphism after tensoring
with~$B/\m_B$. % \mar{TG: because the kernel is finite projective, so we
  % can check modulo maximal ideals.}
Suppose that~$\m_B$ lies over a
maximal ideal~$\m$ of~$A$, so that $B/\m_B$ is a finite field
extension of~$A/\m$; % \mar{TG: using some finiteness}
% that is, it is
% enough
we need to show that the induced surjection
\[K(A/\m)\otimes_{A/\m} B/\m_B\to K(B/\m_B)\]is an
isomorphism. % Since~,\mar{TG: hopefully true for some finiteness
  % reason?} and e
Each side is determined by the de Rham filtration on
the corresponding $G_K$-representation, which is compatible with the
extension of scalars, so we are done.  % \mar{TG: now MK has some claim
  % that one can check compatibility with base change by thinking about
  % ranks after going to finite $E$-algebras, which I didn't think
  % about. TG: maybe something like this?!}
\end{proof}
In what follows, we will work in the relative setting of an
extension~$L/K$.  % We now set up this context more precisely,
% and explain how to extend the previous results to it.
%\mar{ME: Have to edit what follows in accordance with this introductory paragraph.}
To this end,
we let~$\Acirc$ be a $p$-adically complete flat $\cO$-algebra, with $A := \Acirc[1/p]$,
let~$L/K$
be a finite Galois extension, and let~$\gMt_{\Acirc}$ be a
Breuil--Kisin--Fargues $G_K$-module of rank $d$ and height at most~$h$ with
$\Acirc$-coefficients, which admits all descents
over~$L$. % \mar{TG: might in theory want to do this over~$L$ but it
  % shouldn't matter, because of course it's OK to fix Hodge--Tate
  % weights after a base change.}
% Suppose that~$A$ is an $E$-algebra.\mar{TG: is that actually
  % allowed in our setup?! I think we have baked in an assumption of
  % being $p$-adically complete. Not going to worry about that for now -
  % maybe we actually want to work integrally but the invert~$p$, for
  % example.}
Fix some choice of~$\piflat$ a uniformizer of~$L$, write~$\gM_{\Acirc}$
for~$\gM_{\piflat,\Acirc}$, and~$u$ for~$[\piflat]$.
Applying Proposition~\ref{prop: good behaviour of Kisin filtration},
with $L$ in place of $K$,
we obtain a projective $L\otimes_{\Q_p} A$-module 
	$(\varphi^*\gM_{\Acirc}/E(u)\varphi^*\gM_{\Acirc})
	\otimes_{A^{\circ}} A,$  
	which is filtered by projective submodules.
	This filtered module has a natural action of $\Gal(L/K)$,
	which is semi-linear with respect to the action
	of $\Gal(L/K)$ on $L\otimes_{\Q_p} A$ induced by its action
	on the first factor.
	Since $L/K$ is a Galois extension, the tensor product
	$L\otimes_{\Q_p} A$ is an \'etale $\Gal(L/K)$-extension 
	of $K\otimes_{\Q_p} A$, and so \'etale descent allows
	us to descend 
	$(\varphi^*\gM_{\Acirc}/E(u)\varphi^*\gM_{\Acirc})
	\otimes_{A^{\circ}} A$  
	to a filtered module over $K\otimes_{\Q_p} A$;
	concretely, this descent is achieved by taking $\Gal(L/K)$-invariants.
This leads to the following definition. %\mar{TG: perhaps the following
  %two definitions, and the paragraphs following them, should be
  %slightly rearranged? ME: Did something.}

\begin{df}
	\label{def:DdR L/K case}
	In the preceding situation, we write \index{$D_{\dR}(\gMt_{\Acirc})$}
	$$D_{\dR}(\gMt_{\Acirc}) := 
	\bigl((\varphi^*\gM_{\Acirc}/E(u)\varphi^*\gM_{\Acirc})
	\otimes_{A^{\circ}} A\bigr)^{\Gal(L/K)},$$  
		and more generally, for each $i \geq 0,$
		we write
	$$\Fil^i
	D_{\dR}(\gMt_{A^{\circ}}) :=
	\bigl((\Fil^i\varphi^*\gM_{\Acirc}/E(u)\Fil^{i-1}\varphi^*\gM_{\Acirc})\otimes_{\Acirc}A\bigr)^{\Gal(L/K)}$$
(and for $i < 0$, we write 
	$\Fil^i D_{\dR}(\gMt_{A^{\circ}}) :=
	D_{\dR}(\gMt_{A^{\circ}})$).
The property of being a finite rank projective module is preserved
under \'etale descent, and so we find that $D_{\dR}(\gMt_{\Acirc})$
is a rank $d$ projective $K\otimes_{\Q_p} A$-module, filtered by projective
submodules. %with projective associated graded modules. %  Thus %, just as above,
% we obtain a tuple $\bigl(e_{\sigma} D_{\dR}(\gMt_{\Acirc})\bigr)_{\sigma:
% 	K\hookrightarrow L}$ of filtered vector bundles over $\Spec A$.% ,
% and we have the following analogue of Corollary~\ref{cor:Hodge decomposition}.



%\subsection{Inertial types}
%\label{subsubsec:inertial types}Let~$L/K$ be a finite Galois
%extension with inertia group~$I_{L/K}$, and let $\Rep_{I_{L/K},d}$ be
%the algebraic stack over~$\Spec\cO$ classifying $d$-dimensional
%representations of~$I_{L/K}$. More precisely, if~$X$ is an
%$\cO$-scheme, then $\Rep_{I_{L/K},d}(X)$ is the groupoid of locally
%free sheaves of $\cO_X$-modules of rank~$d$, equipped with an
%$\cO_X$-linear action of~$I_{L/K}$.
%
%Suppose that~$E$ is large enough that all irreducible
%$E$-representations of~$I_{L/K}$ are absolutely irreducible, and that
%every irreducible~$\Qpbar$-representation of~$I_{L/K}$ is defined
%over~$E$. Then $\Rep_{I_{L/K},d}\times_{\Spec\cO}\Spec E$ admits the
%following simple description (see also~\cite[Lem.\
%3.2.2]{CEGSKisinwithdd} for an analogous statement
%for~$\Rep_{I_{L/K},d}$ itself in the case that~$I_{L/K}$ has order
%prime to~$p$).
%
%\begin{lem}
%  \label{lem: generic fibre of inertial type stack}$\Rep_{I_{L/K},d}\times_{\Spec\cO}\Spec E$
%  is the disjoint union of irreducible components
%  $\Rep_{I_{L/K},d}(\tau)_E$, where the~$\tau$ run over the isomorphism
%  classes of semisimple $E$-representations
%  of~$I_{L/K}$. Writing~$G_\tau:=\Aut_{E[I_{L/K}]}(\tau)$ \emph{(}a
%  reductive, hence smooth, group scheme over~$E$\emph{)}, we have
%  $\Rep_{I_{L/K},d}(\tau)_E=[\Spec E/G_\tau]$, so that in particular a
%  morphism $X\to\Rep_{I_{L/K},d}\times_{\Spec\cO}\Spec E$ factors
%  through~$\Rep_{I_{L/K},d}(\tau)_E$ if and only if the corresponding
%  locally free $I_{L/K}$-sheaf on~$X$ is Zariski locally isomorphic to~$\tau\otimes_E\cO_X$.
%\end{lem}
%\begin{proof}
%  Since~$E$ has characteristic zero, everything except the reductivity
%  of~$G_\tau$ follows easily from Maschke's theorem. Arguing as in the
%  proof of Lemma~\ref{lem:automorphisms are smooth}, the reductivity
%  follows easily from Schur's lemma; indeed
%  if~$\tau=\oplus\tau_i^{\oplus m_i}$ with the~$\tau_i$ irreducible
%  and pairwise non-isomorphic, then we find
%  that~$G_\tau=\prod_i\GL_{n_i/E}$.
%\end{proof}\mar{TG: probably want more details, but iirc ME plans to
%  rewrite the proof in \cite{CEGSKisinwithdd}, so it doesn't make much
%sense for me to copy that out!}
%\begin{defn}
%  \label{defn: tau inertia rep stack integral via closure}We
%  let~$\Rep_{I_{L/K},d}(\tau)$ be the closure of
%  $\Rep_{I_{L/K},d}(\tau)_E$ in~$\Rep_{I_{L/K},d}(\tau)$.
%\end{defn}\mar{TG: quite possibly we won't actually need to use this,
%  given that the Hodge--Tate weights are only being defined after
%  inverting~$p$, this could be too?}
%

%\begin{df}
%	\label{def:DdR}
	% In the preceding situation, % context of
% 	% Proposition~\ref{prop: good behaviour of Kisin filtration},
% 	we write
% 	$D_{\dR}(\gMt_{A^{\circ}})$
% 	to denote the
% 	projective $K\otimes_{\Q_p} A$-module
% 	$\varphi^*\gM_{\Acirc}/E(u)\varphi^*\gM_{\Acirc}
% 	\otimes_{A^{\circ}} A.$  
% 	For each $i \geq 0,$
% 	we write
% 	$\Fil^i
% 	D_{\dR}(\gMt_{A^{\circ}}) :=
% (\Fil^i\varphi^*\gM_{\Acirc}/E(u)\Fil^{i-1}\varphi^*\gM_{\Acirc})\otimes_{\Acirc}A$
% (and for $i < 0$, we write 
% 	$\Fil^i D_{\dR}(\gMt_{A^{\circ}}) :=
% 	D_{\dR}(\gMt_{A^{\circ}})$).

	Since $A$ is an $E$-algebra,
	we have the product decomposition  $K\otimes_{\Q_p} A
	\iso \prod_{\sigma: K \hookrightarrow E} A,$
	and so, if we write $e_{\sigma}$ for the idempotent
	corresponding to the factor labeled by $\sigma$
	in this decomposition,
	we find that 
	$$D_{\dR}(\gMt_{A^{\circ}}) = 
	\prod_{\sigma: K \hookrightarrow E} e_{\sigma} D_{\dR}(\gMt_{A^{\circ}}),$$
	where each
	$e_{\sigma}D_{\dR}(\gMt_{A^{\circ}})$
	is a projective $A$-module of rank~$d$.
	For each $i$, we write
	$$\Fil^i 
	e_{\sigma}D_{\dR}(\gMt_{A^{\circ}})
	=
	e_{\sigma}\Fil^i D_{\dR}(\gMt_{A^{\circ}}).$$
	Each $\Fil^i
	e_{\sigma}D_{\dR}(\gMt_{A^{\circ}})$ is again a projective $A$-module.
\end{df}

%If $\gMt_{A^{\circ}}$ is of rank $d$,
%then each  
%	$e_{\sigma} D_{\dR}(\gMt_{A^{\circ}})$ is a projective $A$-module
%	of rank $d$. 
The base-change property proved in
	Proposition~\ref{prop: good behaviour of Kisin filtration}
shows
%that each filtered piece
%$\Fil^i
%e_{\sigma} D_{\dR}(\gMt_{A^{\circ}})$
%is a projective $A$-submodule of 
%$e_{\sigma} D_{\dR}(\gMt_{A^{\circ}})$,
%and the base-change invariance proved in that proposition shows
that the various quotients
$\Fil^i
e_{\sigma} D_{\dR}(\gMt_{A^{\circ}})/
\Fil^{i+1}
e_{\sigma} D_{\dR}(\gMt_{A^{\circ}})$
are again projective $A$-modules.  Phrased more geometrically,
then, we see that each $e_{\sigma} D_{\dR}(\gMt_{A^{\circ}})$
gives rise to a vector bundle over $\Spec A$ which is endowed 
with a filtration by subbundles.
We may thus decompose $\Spec A$ into a disjoint union of
open and closed subschemes over which the ranks of the
various subbundles $\Fil^i e_{\sigma} D_{\dR}(\gMt_{\Acirc})$
(or equivalently, the ranks of the various constituents
\[\Fil^i
e_{\sigma} D_{\dR}(\gMt_{A^{\circ}})/
\Fil^{i+1}
e_{\sigma} D_{\dR}(\gMt_{A^{\circ}})\] of the associated graded bundle)
are constant.  

To encode this rank data, and the corresponding decomposition
of~$\Spec A$, it is traditional to use the terminology of
Hodge types, which we now recall.

\begin{df}
	\label{def:Hodge type}\index{Hodge type}
A Hodge type~$\lambdau$ of rank $d$ is by definition a set of tuples of
integers $\{\lambda_{\sigma,j}\}_{\sigma:K\into\Qpbar,1\le j\le d}$
with $\lambda_{\sigma,j}\ge \lambda_{\sigma,j+1}$ for all~$\sigma$ and
all $1\le j\le d-1$. % If~$\lambdau$ is a Hodge type, we
% write~$\lambda|_L$ for the Hodge type for~$L$ determined by
% $(\lambda|_L)_{\sigma,i}=\lambda_{\sigma|_K,i}$ for each
% $\sigma:L\into\Qpbar$.

If $\underline{D} := (D_{\sigma})_{\sigma:K \hookrightarrow E}$
is a collection
of rank $d$ vector bundles over $\Spec A$, labeled (as indicated) by the embeddings
$\sigma:K \hookrightarrow E$,
then we say that $\underline{D}$ has Hodge type $\lambdau$ if
  $\Fil^i D_{\sigma}$
  has constant rank equal 
  to $\#\{j\mid\lambda_{\sigma,j}\ge i\}.$
% \mar{ME: I think the subscript $\sigma|K,j$ should just be  $\sigma,j$, but 
% am possibly confused. TG: agreed, fixed, commenting out. fwiw I
% checked the repository and as I suspected, it dates
% back to some original version where we weren't descending back down
% from ~$L/K$}
\end{df}

% \begin{cor}\label{cor:Hodge decomposition}
% 	In the context of 
% 	Proposition~{\em \ref{prop: good behaviour of Kisin filtration}},
% 	if $\lambdau$ is a Hodge type of rank~$d$,
% 	then we let 
% 	$\Spec A^{\lambdau}$ 
% 	denote the open and closed subset of $\Spec A$
% 	over which the tuple $\bigl(e_{\sigma} D_{\dR}(\gMt_{A^{\circ}})\bigr)$
% 	of filtered vector bundles is of Hodge type $\lambdau$.
% 	We have a corresponding decomposition
% 	\numequation
% 	\label{eqn:Hodge type decomposition}
% 	\Spec A = \coprod_{\lambdau}
% 	\Spec A^{\lambdau}, 
% \end{equation}
% 	labeled by the set of Hodge types $\lambdau$ of rank $d$.
% \end{cor}
% \begin{proof}
% 	This follows from the preceding discussion.
% \end{proof}

% \begin{remark}
%   \label{rem: Hodge filtration not depending on choice of piflat}
%  {\em A priori},
% the tuple of filtered vector bundles
% $\bigl(e_{\sigma} D_{\dR}(\gMt_{\Acirc})\bigr)_{\sigma: K \hookrightarrow E}$
% on~$\Spec A$,
% and hence the decomposition~\eqref{eqn:Hodge type decomposition}
% of~$\Spec A$,
% depends on the descent $\gM_{\piflat, \Acirc}$, 
% and hence on the choice of~$\piflat$.
% However, 
%   Theorem~\ref{thm: admits all descents if and only if semistable}
%   implies that in fact that the decomposition~\eqref{eqn:Hodge type
% 	  decomposition} is independent of the choice
%   of~$\piflat$; see
%   Remark~\ref{rem: Hodge filtration not depending on choice of piflat explained}
%   for a more detailed explanation of this.
%   \end{remark}


\begin{cor}\label{cor:Hodge decomposition L/K case}
	In the preceding context,
	if $\lambdau$ is a Hodge type of rank~$d$,
	then we let 
	$\Spec A^{\lambdau}$ 
	denote the open and closed subscheme of $\Spec A$
	over which the tuple $\bigl(e_{\sigma} D_{\dR}(\gMt_{A^{\circ}})\bigr)$
	of filtered vector bundles is of Hodge type $\lambdau$.
	We have a corresponding decomposition
	\numequation
	\label{eqn:Hodge type decomposition L/K case}
	\Spec A = \coprod_{\lambdau}
	\Spec A^{\lambdau}, 
\end{equation}
labeled by the set of Hodge types $\lambdau$ of rank $d$.  This
decomposition is compatible with base changes $\Acirc\to \Bcirc$ of
$p$-adically complete flat $\cO$-algebras which are topologically of
finite type over~$\cO$.
\end{cor}
\begin{proof}
	This follows from the preceding discussion.
\end{proof}

\begin{remark}
  \label{rem: Hodge filtration not depending on choice of piflat}
   {\em A priori},
the tuple of filtered vector bundles
$\bigl(e_{\sigma} D_{\dR}(\gMt_{\Acirc})\bigr)_{\sigma: K \hookrightarrow E}$
on~$\Spec A$,
and hence the decomposition~\eqref{eqn:Hodge type decomposition L/K case}
of~$\Spec A$,
depends on the descent $\gM_{\piflat, \Acirc}$, 
and hence on the choice of~$\piflat$.
However, 
  Theorem~\ref{thm: admits all descents if and only if semistable}
below
  implies in fact that the decomposition~\eqref{eqn:Hodge type
	  decomposition L/K case} is independent of the choice
  of~$\piflat$; see
  Remark~\ref{rem: Hodge filtration not depending on choice of piflat explained}
  for a more detailed explanation of this.
\end{remark}

As in Section~\ref{subsec:inertial types}, suppose now that~$E$ is large enough that it contains the images of
all embeddings $L\into\Qpbar$, that all irreducible
$E$-representations of~$I_{L/K}$ are absolutely irreducible, and that
every irreducible~$\Qpbar$-representation of~$I_{L/K}$ is defined
over~$E$. We can then immediately combine Corollaries~\ref{cor:inertial type decomposition} and~\ref{cor:Hodge
	  decomposition L/K case}, obtaining the following result.
%\mar{ME: Have to enlarge $E$ to deal with types. TG: how about this? ME: Looks good}

        \begin{cor}
          \label{cor: Hodge and inertial decomposed together}In the
          preceding situation, 	we have a  decomposition
	\numequation
	\label{eqn:Hodge inertia type decomposition L/K case}
	\Spec A = \coprod_{\lambdau,\tau}
	\Spec A^{\lambdau,\tau}, 
\end{equation}
labeled by the set of Hodge types $\lambdau$ of rank $d$, and the set
of inertial types~$\tau$ of~$I_{L/K}$.  This
decomposition is compatible with base changes $\Acirc\to \Bcirc$ of
$p$-adically complete flat $\cO$-algebras which are topologically of
finite type over~$\cO$. The Breuil--Kisin--Fargues
$G_K$-module~$\gMt_{\Acirc}$
is of Hodge type~$\lambdau$ and inertial type~$\tau$ if and only if $A^{\lambdau,\tau}=A$.
        \end{cor}
%  \mar{ME: Check/fix reference once intertial types material is complete}
 
%\begin{remark}
%  \label{rem: Hodge filtration not depending on choice of piflat}
%   {\em A priori},
%the tuple of filtered vector bundles
%$\bigl(e_{\sigma} D_{\dR}(\gMt_{\Acirc})\bigr)_{\sigma: K \hookrightarrow E}$
%on~$\Spec A$,
%and hence the decomposition~\eqref{eqn:Hodge type decomposition L/K case}
%of~$\Spec A$,
%depends on the descent $\gM_{\piflat, \Acirc}$, 
%and hence on the choice of~$\piflat$.
%However, 
%  Theorem~\ref{thm: admits all descents if and only if semistable}
%  implies in fact that the decomposition~\eqref{eqn:Hodge type
%	  decomposition L/K case} is independent of the choice
%  of~$\piflat$; see
%  Remark~\ref{rem: Hodge filtration not depending on choice of piflat explained}
%  for a more detailed explanation of this.
%  % Just as in the case $L = K$ considered above,
%%
%% 
%% The decomposition~\eqref{eqn:Hodge type decomposition L/K case}
%% of~$\Spec A$
%% depends {\em a priori} on the descent $\gM_{\piflat, \Acirc}$, 
%% and hence on the choice of~$\piflat$.
%% However, 
%%   Theorem~\ref{thm: admits all descents if and only if semistable}
%%   implies that in fact that this decomposition
%% 	  is independent of the choice
%%   of~$\piflat$; see
%%   Remark~\ref{rem: Hodge filtration not depending on choice of piflat explained}
%%   for a more detailed explanation of this.
%  \end{remark}


%\begin{defn}
%  \label{defn: Hodge type of family of Kisin modules}Let~$\Acirc$ be a
%  $p$-adically complete flat $\cO$-algebra which is topologically of
%  finite type over~$\cO$, and let~$\gMt_{\Acirc}$
%  and~$\gM_{\Acirc}$ be as above. % \mar{TG: need finiteness assumptions
%    % on~$\Acirc$.}
%  We say that $\gMt_{\Acirc}$ has \emph{Hodge
%    type~$\lambdau$} if for each~$0\le i\le h$, and
%  each~$\sigma:L\into E$, the finite projective $A$-module
%  \[e_\sigma\left((\Fil^i\varphi^*\gM_{\Acirc}/E(u)\Fil^{i-1}\varphi^*\gM_{\Acirc})\otimes_{\Acirc}A\right)\]
%  has constant rank \[\#\{j\mid-\lambda_{\sigma|K,j}\ge i\},\] where~$e_\sigma\in
%  (\cO_L\otimes_{\Zp}A)$ is the idempotent corresponding to~$\sigma$.
%\end{defn}

%The point of Definition~\ref{defn: Hodge type of family of Kisin modules} is the following theorem.

The following key theorem relates the constructions of this section, 
% \mar{ME: As far as I can tell, this result doesn't actually use any of the
% stacky results from this subsection or the preceding two.  Maybe we can
% move it up?  It could be somewhere in \S 4.5. TG: I think the reason
% it's here is that in the final paragraph of the statement we refer to
% the definitions of the Hodge and inertial types, which have only just
% happened. I'm not opposed to moving it, but we might want to avoid
% forward references (and I don't want to move all of that material, at
% least not without a careful look at it\dots) ME: Ah, I'd missed
% that point.  Why is there no discussion of Hodge/inertial types
% in the proof, though?  Maybe there should be {\em some}
% justification of that final claim in  the statement of the theorem?
% TG: good point! Have added something.}
and the previous one, to the $p$-adic Hodge theory of
$G_K$-representations. If~$B$ is a finite $E$-algebra, then we say
that a sub-$\cO$-algebra $\Bcirc\subseteq B$ is an \emph{order of~$B$}
if \index{order}
$\Bcirc[1/p]=B$, and~ $\Bcirc$ is a finite $\cO$-algebra.

% If~$A$ is a finite flat~$\cO$-algebra or $A=\Zpbar$, then we may
% associate a $G_K$-representation $T_A(M)$ to~$M$ as in
% Section~\ref{subsec: Galois reps for GK phi}. 
% Definition~\ref{defn:
%   descending BKF to BK coefficients} is motivated by the following
% result.

% \begin{thm}%\mar{TG: probably move this thm earlier again?}
%   \label{thm: admits all descents if and only if semistable}% Suppose
%   % that~$p>2$. 
%   % \mar{TG: hopefully we can also deduce this from the
%   %   literature if~$p=2$, but I'm not worrying about that for now. Note
%   % that one direction is known even if~$p=2$, and at worst if we don't
%   % know it for~$p=2$, we'll just have to do a bit of rewriting to make
%   % sure we still prove crystalline lifts in this case. Also, I suspect
%   % the characterisation of the crystalline substack will turn out only
%   % to be known for~$p=2$.}
% Suppose
%   that~$A$ is a  finite flat $\cO$-algebra,
%   or that~$A=\Zpbar$. Let~$M$ be a $(G_K,\varphi)$-module with
%   $A$-coefficients, and write~$V_A(M)=T_A(M)[1/p]$. Let~$L/K$ be a finite Galois extension. % \mar{TG: currently we only define these in the
% %     $\varpi^a$-torsion case. Perhaps we could add something to the end
% %     of Section~\ref{subsec:
% % phi gamma over Zp} which allows us to work in this case too? Not going
% % to worry about it for now.}
% Then~$V_A(M)|_{G_L}$ % \mar{TG: probably also not
% % defined! But will be via equivalence to $(\varphi,\Gamma)$-modules, if
% % nothing else - should probably explain the direct functor though?}
% is
% semistable with Hodge--Tate weights in~$[-h,0]$ if and only if there
% is a Breuil--Kisin--Fargues $G_K$-module $\gMt$ with~$A$-coefficients,
% which is of height at most~$h$, which
% admits all descents over~$L$, and which satisfies
% $M=W(\C^\flat)_A\otimes_{\AAinf{A}}\gM$.

% Furthermore~$V_A(M)|_{G_L}$ is crystalline if and only if~$\gMt$ is
% crystalline as a Breuil--Kisin--Fargues $G_L$-module with~$A$
% coefficients.

% In either case, the inertial type of~$V_A(M)$ is given by the action
% of~$I_{L/K}$ on~$\overline{\gM}[1/p]$ described in
%   Section~{\em \ref{subsec:inertial types}}, % \mar{TG: reference to where
%   % this ends up}
% and the Hodge type is determined by Definition~{\em \ref{defn: Hodge type of family of Kisin modules}}.% \mar{TG:
%   % reference to this too.}
% \end{thm}
% \begin{proof}
%   The case $A=\Zpbar$ follows from the case that $A$ is a   finite
%   flat $\cO$-algebra by taking direct limits. In turn, that case
%   easily reduces to the case~$A=\cO$, which is immediate from
%   Corollaries~\ref{acor: admits all descents if and only if potentially
%     semistable} and ~\ref{acor: reading off inertial type and HT weights from BKF}.% \mar{TG: add appendix
%     % ref. Might want to have pst version with HT weights etc. TG: add
%     % HT weights and types once it's settled in the appendix.}
% \end{proof}\mar{TG: this is probably not true but maybe will be true
%   after inverting~$p$? ME: I put in a variant of this result which
%   is hopefully the correct version.}

\begin{thm}
  \label{thm: admits all descents if and only if semistable}
Suppose
  that~$\Acirc$ is a  finite flat $\cO$-algebra,
  %or that~$\Acirc=\Zpbar$.
 let~$M$ be a projective \'etale $(\varphi,G_K)$-module with
  $\Acirc$-coefficients, and write~$V_A(M)=T_{\Acirc}(M)[1/p]$. Let~$L/K$ be a finite Galois extension. % \mar{TG: currently we only define these in the
%     $\varpi^a$-torsion case. Perhaps we could add something to the end
%     of Section~\ref{subsec:
% phi gamma over Zp} which allows us to work in this case too? Not going
% to worry about it for now.}
Then~$V_A(M)|_{G_L}$ % \mar{TG: probably also not
% defined! But will be via equivalence to $(\varphi,\Gamma)$-modules, if
% nothing else - should probably explain the direct functor though?}
is
semistable with Hodge--Tate weights in~$[0,h]$ if and only if there
is an order $(\Acirc)'$ of $A := \Acirc[1/p]$ that contains $\Acirc,$
and a Breuil--Kisin--Fargues $G_K$-module $\gMt_{(\Acirc)'}$ with~$(\Acirc)'$-coefficients,
which is of height at most~$h$, which
admits all descents over~$L$, whose $G_L$-action is canonical, and which satisfies
$M_{(\Acirc)'}=W(\C^\flat)_{(\Acirc)'}\otimes_{\AAinf{(\Acirc)'}}\gMt_{(\Acirc)'}$.

Furthermore~$V_A(M)|_{G_L}$ is crystalline if and only if% , for an
% appropriate choice of $(\Acirc)'$, we may choose
~$\gMt_{(\Acirc)'}$ is
crystalline as a Breuil--Kisin--Fargues $G_L$-module with~$(\Acirc)'$
coefficients.

In either case, % the action
% of~$I_{L/K}$ on~$\overline{\gM}[1/p]$ described in
%   ,
  % \mar{TG: reference to where
  % this ends up}
the Hodge type is determined by
applying Definition~{\em \ref{def:Hodge type}} to the tuple
$\bigl(e_{\sigma} D_{\dR}(\gMt_{(\Acirc)'})\bigr)_{\sigma:K\hookrightarrow E}$
arising from Definition~{\em \ref{def:DdR L/K case}}, and the inertial type of~$V_A(M)$ is given by
Definition~{\em \ref{defn: having inertial type tau}}.
%Definition~{\em \ref{defn: Hodge type of family of Kisin modules}}.% \mar{TG:
  % reference to this too.}
\end{thm}
% \mar{ME: Have to determine if this is really true, and find a proof.
% The {\em if} direction should be immediate, since $(\Acirc)'[1/p] = \Acirc[1/p](= A)$ but the {\em only if} direction involves some kind of ``descent'' from
% $A$ to an appropriately chosen order $(\Acirc)'$.}
\begin{proof}
%By an easy limit argument, the case that~$\Acirc=\Zpbar$
 % follows from the case that~$\Acirc$ is a finite flat $\cO$-algebra,
 % which we assume is the case from now on. Then $A$
  Suppose firstly that~$(\Acirc)'$ and~$\gMt_{(\Acirc)'}$ exist. Then
  if we simply forget the $(\Acirc)'$-coefficients and
  consider~$\gMt_{(\Acirc)'}$ as a Breuil--Kisin--Fargues $G_K$-module
  with~$\Zp$-coefficients, the theorem follows from
  Corollaries~\ref{acor: admits all descents if and only if
    potentially semistable} and~\ref{acor: reading off inertial type
    and HT weights from BKF}.  Conversely, suppose
  that~$V_A(M)|_{G_L}$ is semistable with Hodge--Tate weights
  in~$[0,h]$. If we show that~$(\Acirc)'$ and~$\gMt_{(\Acirc)'}$
  exist, then (as discussed
  in Remarks~\ref{rem:inertial types point} and~\ref{rem:Hodge filtration point})
%at the beginning of
%  Section~\ref{subsec:inertial types}, and the beginning of the current section),
%\mar{ME: Note to self: investigate why here we cite F.24, while
%both remarks cite F.23.}
it again follows from Corollary~\ref{acor: reading off inertial
    type and HT weights from BKF} that the inertial and Hodge types
  are given by the claimed recipes; so it suffices to prove this
  existence. We begin with a preliminary reduction.
  
The ring $A$ is a product of Artinian local $E$-algebras,
say $A  = \prod_i A_i$; then if $\Acirc_i$ denotes
the image of $\Acirc$ in~$A_i$, 
the product $\prod_i \Acirc_i$ is an order of $A$ containing~$\Acirc$.
Thus it is no loss of generality to replace $\Acirc$ by this product,
and hence, by working one factor at a time, to assume that $A$ is local.
 % replacing $A$ by one of these factors (and $\Acirc$ by the image
 % of $\Acirc$ in this factor)
 % we may thus assume that $A$ is local. %  Since~$\cO$ is complete
  % local Noetherian, $\Acirc$ is a product of complete local flat
  % $\cO$-algebras, so we can and do assume from now on that~$\Acirc$ is
  % local, in which case~$A$ is an Artin local $E$-algebra.\mar{TG: no,
  %   it's a product of these, check that we can reduce to this; should
  %   be immediate as an order in each factor gives an order in the
  %   product, but write this and check that we don't have the same
  %   issue elsewhere.}
  The residue field~$E'$ of~$A$ is then a finite extension of~$E$, and $A$ is naturally
  an~$E'$-algebra. The compositum $\cO_{E'}\Acirc$ is an order
in $A$ containing~$\Acirc$ which is furthermore an $\cO_{E'}$-algebra.
Replacing $\Acirc$ by this compositum, we may thus assume
that $\Acirc$ is an $\cO_{E'}$-algebra.
Then, relabelling~$E'$ as~$E$ if necessary, we can and do
  assume that~$E'=E$ . % \mar{TG: any issue about $\cO$ changing when we
    % do this? ME: Dealt with this, hopefully in a reasonable way. TG:
    % commenting out.}
  % \mar{ME: Presumably can base-change/relabel so that $E = E'$? TG: indeed.}
Thus, reformulating the preceding discussion slightly,
we have $A_{\red} = E$ and $\Acirc_{\red} = \cO_E$.

  
  % \mar{TG: may want to reduce to~$A$ being local, actually?}
% \mar{TG: we should investigate whether we could just do all this in
%   the~$\Ainf$ world. TG: if we can't, we could remark that we make an
%   initial $\piflat$ descent to use~$\gS$ commutative algebra.}

  By Corollary~\ref{acor: admits all
    descents if and only if potentially semistable} there is a unique
  Breuil--Kisin--Fargues $G_K$-module~$\gMt$ with $\Zp$-coefficients
  which is of height at most~$h$ and admits all descents over~$L$, and
  satisfies \[M=W(\C^\flat)\otimes_{\Ainf}\gMt.\] For each~$\piflat$
  we have by definition a descent~$\gM_{\piflat}$ of~$\gMt$. Fix some
  choice of~$\piflat$, write~$u$ for~$[\piflat]$, and write~$\gS$
  for~$\gS_{\piflat}$ and~$\gM$ for~$\gM_{\piflat}$.

  We now follow the proof of~\cite[Prop.\ 1.6.4]{MR2373358}. In
  particular, we work for the most part with the Breuil--Kisin
  module~$\gM$, rather than the Breuil--Kisin--Fargues
  module~$\gMt$. It is possible that by using~\cite[Prop.\
  4.13]{2016arXiv160203148B}, we could make our arguments with~$\gMt$
  itself, but since this does not seem likely to significantly
  simplify the proof, and would make it harder for the reader to
  compare to the arguments of~\cite{MR2373358}, we have not attempted
  to do this.

 Note
  firstly that~$\gM\subset\gMt\subset M$ are stable under the action
  of~$\Acirc$ on~$M$. Indeed, since~$\Acirc$ is local, it is enough to
  check stability under the action of~$(\Acirc)^\times$. In the case
  of~$\gMt$ it is immediate from the unicity of~$\gMt$ that we
  have~$\gMt=a\gMt$ for any~$a\in(\Acirc)^\times$, and similarly
  for~$\gM$ it follows from Lemma~\ref{lem: uniqueness of descent for
    G K infty Ainf}.

  In particular~$\gM$ is naturally a
  finite~$\gS_{\Acirc}$-module, which is projective
  as an $\gS$-module.  While~$\gM$ need not be a projective
  $\gS_{\Acirc}$-module, $\cO_{\cE,\Acirc}\otimes_{\gS_{\Acirc}}\gM$ is a projective~$\cO_{\cE,\Acirc}$-module,
  because after the faithfully flat base extension
  $\cO_{\cE,\Acirc}\into W(\C^\flat)_{\Acirc}$ it is identified with~$M$.
It follows
%from~\ref{lem:MK projectivity lemma}
from~\cite[Lem.\ 1.6.1]{MR2373358}
that $\gM[1/p]$ is a projective
  $\gS_{\Acirc}[1/p]$-module, necessarily of rank~$d$, %\mar{TG: it's
    %equivalently free? I think $\gS_{\Acirc}[1/p]$ is local? We use
    %this below, I think. No, PID?}
and that $\gM[1/u]$ is a projective $\gS_{\Acirc}[1/u]$-module,
again of rank~$d$.

  % \mar{TG: $E'$ is making a reappearance! ME: Now quashed!}
  In fact, it will be useful to note that $\gM[1/p]$ is actually
  free of rank~$d$.  To see this, it suffices to prove it
  after base-changing to $$(\gS_{\Acirc})[1/p]_{\red}
  = (\gS_{\Acirc_{\red}})[1/p] = \gS_{\cO}[1/p] 
  %Now $\Acirc_{\red}[1/p]$ is
  %a finite extension $E'$ of $E$, so that $\Acirc_{\red}$ is an order
  %in~$E'$, and in particular is of finite index 
  %in~$\cO_{E'}$. 
%Thus we have isomorphisms
% $$(\gS_{\Acirc_{\red}})[1/p] \cong \gS_{\cO_{E'}}[1/p]
  \cong \prod_{\sigma:W(k) \hookrightarrow \cO} \cO[[u]][1/p].$$
  Since $\cO[[u]][1/p]$ is a PID (by the Weierstrass preparation
  theorem), 
%  \mar{ME: Give reference if we can find one }
  any finitely generated projective module over
  $(\gS_{\Acirc_{\red}})[1/p]$ of constant rank is necessarily free.
For later use,  we note that since
$(\cO_{\cE,\Acirc})_{\red}=
\cO\otimes_{\Z_p} \cO_{\cE} = \prod_{\sigma: W(k) \hookrightarrow \cO}
\cO_{\cE}$,
and $\cO_{\cE}$ is a local ring (indeed a discrete valuation ring,
with uniformizer~$p$), 
any projective module of constant rank over
$\cO_{\cE,\Acirc}$ is also necessarily free.
  

  Now let~$\gM'_{\cO}$ denote the image of~$\gM$ under the
  projection~$\gM[1/u]\to\gM[1/u]\otimes_AE$ (the map~$A\to E$ being
  the projection from~$A$ to its residue field), and
  let~$\gM'_{\cO}\subset\gM_{\cO}$ be the canonical inclusion (with
  finite cokernel) of~$\gM'_{\cO}$ into the corresponding finite
  projective~$\gS_{\cO}$-module. (The existence of~$\gM_{\cO}$
  follows from the structure theory of $\gS$-modules, see~\cite[Prop.\
  4.3]{2016arXiv160203148B}. Concretely, we have
  $\gM_{\cO}=\gM'_{\cO}[1/u]\cap\gM'_{\cO}[1/p]$.) The module~$\gM_{\cO}$ is again a
  Breuil--Kisin module of height at most~$h$.

  Choose an $\gS_{\cO}$-basis for~$\gM_{\cO}$, and lift it to an
  $\gS_{\Acirc}[1/p]$-basis of $\gM[1/p]$,
  %Choose~$(\Acirc)'$ such
  %that the matrix of~$\Phi_{\gM[1/p]}$ with respect to this basis has
  %entries in~$(\Acirc)'$; indeed, we can choose~$(\Acirc)'$ to be the
  and let~$(\Acirc)''$ be the
  $\Acirc$-subalgebra of~$A$ generated by these matrix coefficients.
  These coefficients have bounded powers of $p$  in their
  denominators (when we think of $A$ as equalling $\Acirc[1/p]$)
  and lie in $\cO$ after passing to the quotient of $A$ by its
  maximal ideal (equivalently its nilradical); in other words
  these entries lie in $\Acirc + p^{-N} \nil(\Acirc)$ for some $N \geq 0$,
  and this  $\Acirc$-submodule of $A$  generates a finite $\Acirc$-subalgebra
  of~$A$.
  %$(\Acirc)' \subseteq  A
  %which is a finitely generated subalgebra because the images of these
  %coefficients modulo the (nilpotent) maximal ideal of~$A$ all lie
  %in~$\cO$.
Thus $(\Acirc)''$  is an  order in~$A$. 
Let~$\gM_{(\Acirc)''}$ be the $\gS_{(\Acirc)''}$-submodule
  of~$\gM[1/p]$ spanned by this basis. Then~$\gM_{(\Acirc)''}$ is
  $\varphi$-stable by construction.
  Again by construction we have that
 $$\gM_{(\Acirc)''}[1/p] 
  = (\Acirc)''\otimes_{\Acirc} \gM[1/p],$$  
  and that
$$\cO\otimes_{(\Acirc)''} \cO_{\cE,(\Acirc)''} \otimes_{\gS_{(\Acirc)''}}
\gM_{(\Acirc)''}
=  \cO_{\cE}\otimes_{\gS} \gM_{\cO}
=  \cO_{\cE}\otimes_{\gS} \gM'_{\cO}$$
(where the first tensor product is with respect to the surjection
$(\Acirc)'' \to  (\Acirc)''_{\red} = \cO$).
Thus $\cO_{\cE,(\Acirc)''} \otimes_{\gS_{(\Acirc)''}} \gM_{(\Acirc)''}$ and 
  $\cO_{\cE,(\Acirc)''}\otimes_{\gS_{\Acirc}} \gM$
  are two rank $d$ projective $\cO_{\cE,(\Acirc)''}$-modules which
  coincide after inverting $p$, and also after reducing modulo
 $\nil\bigl( (\Acirc)'').$ In fact, they are both free of rank~$d$, by
our above observation that a projective $\cO_{\cE,\Acirc}$-module of
constant rank is necessarily free
(applied now  with $(\Acirc)''$  in  place of $\Acirc$).
If we choose bases for each of them which coincide 
modulo $\nil\bigl((\Acirc)''\bigr)$, then we may regard
each of these bases as a basis of $\cO_{\cE,(\Acirc)''}\otimes_{\gS_{\Acirc}}
\gM[1/p]$
over $\cO_{\cE,(\Acirc)''}[1/p]$, and so they differ by a
change-of-basis matrix lying in $1 + \nil\bigl((\Acirc)'') 
M_d\bigl(\cO_{\cE,(\Acirc)''}[1/p]\bigr).$
Just as we argued above in the construction of~$(\Acirc)''$,
we see that we may enlarge $(\Acirc)''$ to an
order $(\Acirc)'$ such that the entries of this change-of-basis matrix
lie in~$(\Acirc)'$.  Consequently, if we set
$\gM_{(\Acirc)'}   := \gS_{(\Acirc)')} \otimes_{\gS_{(\Acirc)''}} 
\gM_{(\Acirc)''}$, then we find that
$$\cO_{\cE,(\Acirc)'} \otimes_{\gS_{\Acirc}} \gM
= \cO_{\cE,(\Acirc)'} \otimes_{\gS_{(\Acirc)''}} \gM_{(\Acirc)'}.$$
Extending scalars further, we find that
\begin{multline*}
M_{(\Acirc)'} :=  W(\C^\flat)_{(\Acirc)'}\otimes_{W(\C^\flat)_{\Acirc}}
M \\
 = W(\C^\flat)_{(\Acirc)'} \otimes_{\gS_{\Acirc}} \gM
= W(\C^\flat)_{(\Acirc)''}\otimes_{\gS_{(\Acirc)'}} \gM_{(\Acirc)'}.
\end{multline*}


%and satisfies \mar{TG: this looks
%    wrong to me! I think $W(\C^\flat)_{(\Acirc)'}$ should be $\cO_{\cE,(\Acirc)'}$? 
%ME: Hopefully fixed now.}
%  $W(\C^\flat)_{(\Acirc)'}\otimes_{\gS_{(\Acirc)'}}\gM_{(\Acirc)'}=\cO_{\cE}\otimes_{\gS}\gM$.

  We claim that $\gM_{(\Acirc)'}$ is a Breuil--Kisin module of height at most~$h$,
  i.e.\ that the cokernel of~$\Phi_{\gM_{(\Acirc)'}}$ is killed
  by~$E(u)^h$. This cokernel is, as an $\gS$-module, a successive
  extension of copies of the cokernel of~$\Phi_{\gM_{\cO}}$ (the
  injectivity of~$\Phi$ implies that the formation of the cokernel
  of~$\Phi$ is exact), % \mar{TG: find
    % something in MK's papers to justify exactness of the formation of cokernels}
  and is in
  particular $p$-torsion free (by \cite[Lem.\ 1.2.2~(2)]{KisinModularity}). % \mar{TG: again justify following MK}
  After inverting~$p$ the cokernel is a
  base change of the cokernel of~$\Phi_{\gM}$, and is therefore killed
  by~$E(u)^h$, as required.
  
  We now
  set~$\gMt_{(\Acirc)'}:=\AAinf{(\Acirc)'}\otimes_{\gS_{(\Acirc)'}}\gM_{(\Acirc)'}$,
  a Breuil--Kisin--Fargues $G_K$-module of height at most~$h$ with
  $(\Acirc)'$-coefficients which by construction satisfies
  $M_{(\Acirc)'}=W(\C^\flat)_{(\Acirc)'}\otimes_{\AAinf{(\Acirc)'}}\gMt_{(\Acirc)'}$. It
  remains to prove that $\gMt_{(\Acirc)'}$ admits all descents
  over~$L$ and that the $G_L$-action is canonical. Applying Corollary~\ref{acor: admits all descents if and
    only if potentially semistable} again to~$M_{(\Acirc)'}$ regarded
  as an \'etale $(\varphi,G_K)$-module with $\Zp$-coefficients, we see that
  there is a Breuil--Kisin--Fargues $G_K$-module~$(\gMt)'$ of height
  at most~$h$ with $\Zp$-coefficients which admits all descents
  over~$L$ and satisfies
  $M_{(\Acirc)'}=W(\C^\flat)\otimes_{\Ainf}(\gMt)'$. Write~$(\gM)'$
  for the descent of~$(\gMt)'$, for our particular choice of~$\piflat$;
  then by the uniqueness of Breuil--Kisin modules
  with~$\Zp$-coefficients (\cite[Prop.\ 2.1.12]{KisinCrys}) we have
  $(\gM)'=\gM_{(\Acirc)'}$, so that in
  fact~$(\gMt)'=\gMt_{(\Acirc)'}$. Thus~$\gMt_{(\Acirc)'}$ admits all
  descents over~$L$ to Breuil--Kisin modules with~$\Zp$-coefficients,
  and by another application of Lemma~\ref{lem: uniqueness of descent
    for G K infty Ainf}, these Breuil--Kisin modules
  are~$(\Acirc)'$-stable. Finally, the $G_L$-action is canonical by Proposition~\ref{prop: Caruso Liu canonical action semistable}, and we are done.
\end{proof}

\begin{rem}
  \label{rem: Hodge filtration not depending on choice of piflat explained}It
  follows from
  Theorem~\ref{thm: admits all descents if and only if semistable}
  that the decomposition %s~\eqref{eqn:Hodge type decomposition}
  ~\eqref{eqn:Hodge type decomposition L/K case} is
 % Definition~\ref{defn: Hodge type of family of Kisin modules} is
  independent of the choice of~$\piflat$. Indeed,
  the condition that
  $\Fil^i e_{\sigma} D_{\dR}(\gMt_{\Acirc})$
 % $e_\sigma\left((\Fil^i\varphi^*\gM_{\Acirc}/E(u)\Fil^{i-1}\varphi^*\gM_{\Acirc})\otimes_{\Acirc}A\right)$
  has given constant rank can be checked after base changing
  to % \mar{TG: as ever writing something modulo finiteness, and didn't
    % think through modulo maximal ideal versus powers of it}
  all~$A/\m$, for~$\m$ a maximal ideal of~$A$, so it follows from
  Proposition~\ref{prop: good behaviour of Kisin filtration} and
  Theorem~\ref{thm: admits all descents if and only if semistable}
  that the direct factor $A^{\lambdau}$ of $A$ is characterized as follows:
  %that~$\gMt_{\Acirc}$ has Hodge type~$\lambdau$ if and only if the
  %following condition holds:
  for any $\Acirc$-algebra $\Bcirc$ which
  is finite and flat as an~$\cO$-algebra,
  the canonical morphism $A \to \Bcirc[1/p]$ factors through $A^{\lambdau}$
  if and only if the $G_K$-representation
  $V_B(\gMt_{\Bcirc})$ has Hodge type~$\lambdau$.
\end{rem}

%The following result is minor modification (in its statement)
%of~\cite[Lem.\ 1.6.1]{MR2373358}, which follows directly
%from the proof of that result.
%
%\begin{lemma} 
%\label{lem:MK projectivity lemma}
%Let $\gM$ be a Breuil--Kisin module over $\gS_{\Acirc}$,
%for some finite flat $\cO$-algebra~$\Acirc$,
%of height at most $h$ {\em (}for some choice of~$h${\em )},
%such that $\gM$  is projective over~$\gS$
%and
%$\cO_{\cE,\Acirc}\otimes_{\gS_{\Acirc}}\gM$ is projective
%over~$\cO_{\cE,\Acirc}$.
%Then $\gM[1/p]$ is  projective over~$\gS_{\Acirc}[1/p]$,
%and $\gM[1/u]$ is projective over~$\gS_{\Acirc}[1/u].$
%\end{lemma}
%\begin{proof}
%\end{proof}

%\begin{cor}
%	\label{cor:Hodge decomposition L/K case}
%\end{cor}
%\begin{proof}
%\mar{ME: This material has to actually be reworked into a corollary, which
%we then have to state.}
%  By Proposition~\ref{prop: good behaviour of Kisin
%    filtration}, for each~$0\le i\le h$, and each~$\sigma:K\into E$,
%  \[
%    e_\sigma((\Fil^i\varphi^*\gM_{\Bcirc}/E(u)\Fil^{i-1}\varphi^*\gM_{\Bcirc})\otimes_{\Bcirc}B)\]
%  is a finite projective $B := \Bcirc[1/p]$-module whose formation is compatible with
%  base change.
%  Thus we may decompose $\Spec B$ as the disjoint union of closed subschemes
%  $\Spec B^{\lambdau}$,
%  where, for each $\lambdau \leq h$,
%  we let $\Spec B^{\lambdau}$ denote the locus over 
%  which ~$\gMt_{\Bcirc}$ has Hodge type~$\lambdau$.
%  \mar{ME: Make this a bit more explicit: I think we are just talking
%	  about a filtration on the projective module $\varphi^*\gM_B/E(u)
%	  \varphi^*\gM_B$.}
%  Passing to rings, we may thus factor $B$ as a product
%  $B = \prod_{\lambdau} B^{\lambdau}$.
%  \end{proof}



\section[{Moduli stacks of potentially semistable
  representations}]{Moduli stacks of potentially semistable
  representations \sectionmark{Stacks of potentially semistable
  representations}}\sectionmark{Stacks of potentially semistable
  representations}\label{subsec: moduli stacks of pst GK repns}
% \subsection{Flat parts and scheme-theoretic images}\label{subsubsec: scheme
%   theoretic image}
% We apply this formalism as follows.
% Fix  an inertial type~$\tau$.
We are now in
a position to define the main objects of interest in this chapter,
which are the closed substacks of~$\cX_d$ classifying representations
which are potentially crystalline or potentially semistable of
% inertial type~$\tau$ and
fixed Hodge and inertial types. 
% \mar{TG: make
%   sure these are now defined again in the notation section, as well as
%   the versal rings.}



In this section ~$L/K$ will denote a finite Galois
extension with inertia group~$I_{L/K}$; we will always assume
(without specifically recalling this) that~$E$ is large enough that all irreducible
$E$-representations of~$I_{L/K}$ are absolutely irreducible, and that
every irreducible~$\Qpbar$-representation of~$I_{L/K}$ is defined
over~$E$. % Suppose that~$h$ is a sufficiently large integer
% with the property that~$\lambdau$ is bounded by~$h$.
For any such~$L/K$ we have the
stacks   $\cC_{d,\ss,h}^{L/K}$ and $\cC_{d,\crys,h}^{L/K}$
that we defined in Definition~\ref{defn: semistable crys stacks L over K}, % \mar{TG: where defined now?}
and we write $\cC_{d,\ss,h}^{L/K,\fl}$ and
$\cC_{d,\crys,h}^{L/K,\fl}$ respectively for their flat parts (in
the sense recalled in Appendix~\ref{app: formal algebraic
  stacks}).

% Recall that a Hodge type~$\lambdau$ is by definition a set of tuples
% integers $\{\lambda_{\sigma,i}\}_{\sigma:K\into\Qpbar,1\le i\le d}$
% with $\lambda_{\sigma,i}\ge \lambda_{\sigma,i+1}$ for all~$\sigma$ and
% all $1\le i\le d-1$. 
\begin{df} \index{Hodge type!effective}
We say that a Hodge
type~$\underline{\lambda}$ is \emph{effective} % \index{Hodge type|effective}
if $\lambda_{\sigma,i}
\geq 0$ for each $\sigma$ and $i$,
and that $\lambdau$ is \emph{bounded by~$h$} \index{Hodge type!bounded
by~$h$} if
$\lambda_{\sigma,i}\in [0,h]$ for each~$\sigma$ and $i$.
% \mar{TG: or perhaps the
  % opposite sign, I think it would be wise to check our conventions
  % throughout at some point.} 
\end{df}
% If~$\lambdau$ is a Hodge type, we write~$\lambda|_L$ for the Hodge
% type for~$L$ determined by
% $(\lambda|_L)_{\sigma,i}=\lambda_{\sigma|_K,i}$ for each
% $\sigma:L\into\Qpbar$.

\begin{prop}
  \label{prop: HT weights are a closed condition}
  Let~$L/K$ be a finite
  Galois extension. Then the stacks
  $\cC_{d,\ss,h}^{L/K,\fl}$ and $\cC_{d,\crys,h}^{L/K,\fl}$
  are scheme-theoretic unions of closed substacks
  $\cC_{d,\ss,h}^{L/K,\fl,\lambdau,\tau}$ and
  $\cC_{d,\crys,h}^{L/K,\fl,\lambdau,\tau}$, where~$\lambdau$ runs
  over all effective Hodge types that are bounded by~$h$, and~$\tau$ runs over
  all $d$-dimensional $E$-representations of~$I_{L/K}$. These latter closed substacks are
  uniquely characterised by the following property: if~$\Acirc$ is a
  finite flat~$\cO$-algebra, then an $\Acirc$-point of
  $\cC_{d,\ss,h}^{L/K,\fl}$ \emph{(}resp.\
  $\cC_{d,\crys,h}^{L/K,\fl}$\emph{)} is a point of
  $\cC_{d,\ss,h}^{L/K,\fl,\lambdau,\tau}$ \emph{(}resp.\
  $\cC_{d,\crys,h}^{L/K,\fl,\lambdau,\tau}$\emph{)} if and only if
  the corresponding Breuil--Kisin--Fargues module~$\gMt_{\Acirc}$ has
  Hodge type~$\lambdau$ in the sense given by
applying Definition~{\em \ref{def:Hodge type}} to the tuple
$\bigl(e_{\sigma} D_{\dR}(\gMt_{\Acirc})\bigr)_{\sigma:K\hookrightarrow E}$
arising from Definition~{\em \ref{def:DdR L/K case}},
 % Definition~{\em \ref{defn: Hodge type of family of Kisin modules}},
  and inertial type~$\tau$ in the
sense of Definition~{\em \ref{defn: having inertial type tau}};
or, more succinctly {\em (}and writing $A := \Acirc[1/p]$, as usual{\em )},
if and only if $A^{\lambdau,\tau} = A$. %\mar{TG: need this to
                                %actually be defined somewhere}

Finally, each stack $\cC_{d,\ss,h}^{L/K,\fl,\lambdau,\tau}$ or
$\cC_{d,\crys,h}^{L/K,\fl,\lambdau,\tau}$ is a
%finite type\mar{TG: more finite type finite presentation}
$p$-adic
formal algebraic stack of finite presentation which is flat over~$\Spf\cO$
and whose diagonal is affine, and the natural morphisms to $\cX_{d}$ are
representable by algebraic spaces, proper, and of finite
presentation. % \mar{TG: add something about all these properties to the proof: add
  % finite presentation, just need to check lfp, might be
  % automatic? ME: Should be all done now}
% \mar{TG: affine diagonal?}
\end{prop}%\mar{TG: I am making a mess of this, but I suspect that
%  getting this to be rigorous is more a matter of writing it correctly
%than of proving anything new, so I'm leaving it as is. ME: Should 
%cite more abstract result that is proved first (or maybe in the appendix),
%and then apply Prop.~\ref{prop: good behaviour of Kisin filtration}
%to get this. (Projectivity statement of that result will give us the
%relevant decomposition of $\Spf B^{\circ}$, and then the base-change
%statement will show that it descends.)}
\begin{proof}
  % \mar{TG: need to prove this!}
  We give the proof
  for~$\cC_{d,\ss,h}^{L/K,\fl}$, the crystalline case being
  formally identical.
  Since~$\cC_{d,\ss,h}^{L/K,\fl}$ is a flat
  $p$-adic formal algebraic stack of finite presentation
  over~$\Spf\cO$, we can choose a smooth surjection
  \numequation
  \label{eqn:B to C}
  \Spf\Bcirc\to
  \cC_{d,\ss,h}^{L/K,\fl}
  \end{equation}
  where~$\Bcirc$ is topologically of
  finite type over~$\cO$.
  %\mar{ME: Sketching proof for Hodge types, ignoring inertial types for now}

As usual, we write $B := \Bcirc[1/p]$.
By Corollary~\ref{cor: Hodge and inertial decomposed together} %   Combining Corollaries~\ref{cor:inertial type decomposition} and~\ref{cor:Hodge
	  % decomposition L/K case},
%  \mar{ME: Check/fix reference once intertial types material is complete}
  we may write $\Spec B$ as a disjoint union $\Spec B = \coprod_{\lambdau,
	  \tau} B^{\lambdau,\tau},$
  and so correspondingly factor $B$ as a product
  $B = \prod_{\lambdau,\tau} B^{\lambdau,\tau}.$
  If we let $B^{\lambdau,\tau,\circ}$ denote the image of $\Bcirc$
  in $B^{\lambdau,\tau}$, then we obtain an induced injection
  $B^{\circ} \hookrightarrow \prod_{\lambdau,\tau} B^{\lambdau,\tau,\circ}$,
  which induces a scheme-theoretically dominant morphism
  \numequation
  \label{eqn:dominant B}
  \coprod_{\lambdau,\tau} \Spf B^{\lambdau,\tau,\circ} \to \Spf \Bcirc.
  \end{equation}

  Write $R = \Spf \Bcirc \times_{\cC_{d,\ss,h}^{L/K,\fl}} \Spf \Bcirc,$
	  so that 
	  $[\Spf \Bcirc/R] \iso \cC_{d,\ss,h}^{L/K,\fl}.$
	  Then we claim that each $\Spf B^{\lambdau,\tau,\circ}$ is $R$-invariant.
	  Granting this,
	  %\mar{ME: Have to check $R$-invariance; this uses 
	  %the base-change statement from Prop.~\ref{prop: good behaviour
	  %of Kisin filtration}.}
	  if we write $R^{\lambdau,\tau}$ for the
	  restriction of $R$ to $\Spf B^{\lambdau,\tau,\circ}$,
	  and then define
	  $\cC_{d,\ss,h}^{L/K,\fl,\lambdau,\tau}
	  := [\Spf B^{\lambdau,\tau,\circ}/R^{\lambdau,\tau}],$
	  it follows from the discussion of closed substacks
	  in Appendix~\ref{app: formal
		 algebraic stacks} % \mar{TG: we should find a way to
                 % make this a more precise citation ME: Added something}
               that 
	  $\cC_{d,\ss,h}^{L/K,\fl,\lambdau,\tau}$
	  embeds as a closed sub-formal algebraic stack of 
	  $\cC_{d,\ss,h}^{L/K,\fl}.$
	  Furthermore, the induced morphism
	  \numequation
	  \label{eqn:dominant C}
	  \coprod_{\lambdau,\tau} 
	  \cC_{d,\ss,h}^{L/K,\fl,\lambdau,\tau}
	  \to 
	  \cC_{d,\ss,h}^{L/K,\fl}
  \end{equation}
	  induces the morphism~\eqref{eqn:dominant B} after
	  pull-back via the morphism~\eqref{eqn:B to C},
	  which is representable by algebraic spaces,
	  smooth, and surjective.  Since this latter morphism is
	  scheme-theoretically dominant, so is~\eqref{eqn:dominant C}.

	  We now verify that $\Spf B^{\lambdau,\tau,\circ}$ is 
	  $R$-invariant, i.e.\ that
	  $R_0 := \Spf B^{\lambdau,\tau,\circ} \times_{\cC_{d,\ss,h}^{L/K,\fl}} \Spf \Bcirc$
	  and
	  $R_1 := \Spf \Bcirc \times_{\cC_{d,\ss,h}^{L/K,\fl}} \Spf B^{\lambdau,\tau,\circ}$
	  coincide as closed sub-formal algebraic spaces of 
	  $R = \Spf \Bcirc \times_{\cC_{d,\ss,h}^{L/K,\fl}} \Spf \Bcirc$
	  (and thus that both coincide with
	  $R^{\lambdau,\tau}
	  := \Spf B^{\lambdau,\tau,\circ} \times_{\cC_{d,\ss,h}^{L/K,\fl}}
	  \Spf B^{\lambdau,\tau,\circ}$).

	  Since $\Spf \Bcirc \to \cC_{d,\ss,h}^{L/K,\fl}$
	  is representable by algebraic spaces and smooth,
	  it is in particular representable by algebraic spaces,
	  flat, and locally of finite type, and thus so are each 
	  of the projections
	  %$$\Spf B^{\lambdau,\circ} \times_{\cC_{d,\ss,h}^{L/K,\fl}} \Spf \Bcirc \to \Spf B^{\lambdau,\circ}$$
	  $R_0 \to \Spf B^{\lambdau,\tau,\circ}$
	  and 
	  $R_1 \to \Spf B^{\lambdau,\tau,\circ}.$
	 % $$\Spf \Bcirc \times_{\cC_{d,\ss,h}^{L/K,\fl}} \Spf B^{\lambdau,\circ} \to \Spf B^{\lambdau,\circ}.$$
	  It follows from Lemmas~\ref{lem: alem closed formal algebraic scheme
    adic*} and~\ref{alem: Noetherian flat affine formal algebraic
    spaces} that any affine formal algebraic space which is \'etale
	  over either of
	 % $\Spf B^{\lambdau,\circ} \times_{\cC_{d,\ss,h}^{L/K,\fl}} \Spf \Bcirc$
	  $R_0$ or $R_1$
	  %$\Spf \Bcirc \times_{\cC_{d,\ss,h}^{L/K,\fl}} \Spf B^{\lambdau,\circ}$
	  is of the form $\Spf \Acirc$, where $\Acirc$ is $\varpi$-adically
	  complete, topologically
	  of finite type, and flat over~$\cO$. 
	  % (Here we are applying \cite[Lem.~8.18]{Emertonformalstacks}
	%   to deduce flatness on the level of rings from flatness
	%   of a morphism of affine formal algebraic spaces,
	%   and~\cite[\href{https://stacks.math.columbia.edu/tag/0ANT}{Tag 0ANT}]{stacks-project}  
	%   to deduce being topologically of finite type on the level
	%  of rings from a morphism of affine
	% formal algebraic spaces being of finite type.)
	Thus, to show that $R_0$ and $R_1$ coincide as subspaces
	of $R$, it suffices to show that if $\Acirc$ is any
	$\varpi$-adically complete, topologically of finite type,
	and flat $\cO$-algebra endowed with a morphism 
	$\Spf \Acirc \to R$, then this morphism factors through
	$R_0$ if and only if it factors through $R_1$.
	Unwinding the definitions of $R,$ $R_0$, and $R_1$
	as fibre products, this amounts to showing that
	if we are given a pair of morphism
	$\Spf \Acirc \rightrightarrows \Spf \Bcirc$
	which induce the same morphism to $\cC_{d,\ss,h}^{L/K,\fl}$,
	then one factors through $\Spf B^{\lambdau,\tau,\circ}$
	if and only if the other does.

	The pair of morphisms
	$\Spf \Acirc \rightrightarrows \Spf \Bcirc$
	correspond to a pair of morphisms
	$f_0,f_1: \Bcirc \rightrightarrows \Acirc$.
	The fact that these morphisms induce the same morphism to
	$\cC_{d,\ss,h}^{L/K,\fl}$ may be rephrased as saying 
	that $f_0^*\gMt$ and $f_1^*\gMt$ are isomorphic
% \mar{ME: BKF modules, right?  (Not just BK modules.) TG: fixed,
%   commenting out.}
	Breuil--Kisin--Fargues modules over $\Acirc$.  % In particular,
	% writing $A := \Acirc[1/p]$, these two pull-backs 
	% give rise to the same Breuil--Kisin module $\gMt_A$ 
	% over $A$.
        Now, applying the base change statements
	from Propositions~\ref{prop: good behaviour of inertial type}
	and~\ref{prop: good behaviour of Kisin filtration},
	we find that both
	% \mar{ME: Writing $D_{\dR}(\gMt)$ to denote $\varphi^*\gMt/E(u)
	% 	\varphi^*\gMt$. ME: Also writing $\WD(\gMt)$
	% to denote $E\otimes_{E\otimes K_0}\gMt/u\gMt$}
	$f_0^*\WD(\gMt_{\Bcirc})$
	and
	$f_1^*\WD(\gMt_{\Bcirc})$
	are isomorphic, as $I_K$-representations,
	to $\WD(\gMt_{\Acirc})$,
        and that both
	$f_0^*D_{\dR}(\gMt_{\Bcirc})$
	and
	$f_1^*D_{\dR}(\gMt_{\Bcirc})$
	are isomorphic, as filtered $K\otimes_{\Q_p} B$-modules,
	to $D_{\dR}(\gMt_{\Acirc})$.  Thus either, and hence both,
	of
	$f_0^*\WD(\gMt_{\Bcirc})$
	and
	$f_1^*\WD(\gMt_{\Bcirc})$
	are of inertial type $\tau$
	if and only if
	$\WD(\gMt_{\Acirc})$ is of inertial type~$\tau$,
	while either, and hence both,
	of 
	$f_0^*D_{\dR}(\gMt_{\Bcirc})$
	and
	$f_1^*D_{\dR}(\gMt_{\Bcirc})$
	are of Hodge type $\lambdau$ if and only if 
	$D_{\dR}(\gMt_{\Acirc})$ is of Hodge type~$\lambdau$.  
	Consequently $f_0$ factors through
	$\Spf B^{\lambdau, \tau,\circ}$
	if and only $f_1$ does.  This shows that
	$\Spf B^{\lambdau, \tau,\circ}$ is indeed $R$-invariant.
        
% Suppose now that ~$\Acirc$ is a
%   finite flat~$\cO$-algebra, and that we have a morphism
%   $\Spec\Acirc\to\cC_{d,\ss,h}^{L/K,\fl}$.
	
We now show that~ $\cC_{d,\ss,h}^{L/K,\fl,\lambdau,\tau}:= [\Spf
B^{\lambdau,\tau,\circ}/R^{\lambdau,\tau}]$ satisfies the required
property, i.e.\ that for any finite flat $\cO$-algebra~$\Acirc$,
a morphism $\Spf A^{\circ} \to \cC_{d,\ss,h}^{L/K,\fl}$
factors through $\cC_{d,\ss,h}^{L/K,\fl,\lambdau,\tau}$
if and only if % the associated Breuil--Kisin--Fargues module
% $\gMt_{\Acirc}$ has Hodge type $\lambdau$ and inertial type $\tau$, 
% Writing $A := \Acirc[1/p]$ as usual,
% this latter condition is equivalent to asking that
$A^{\lambdau,\tau}  = A$.
	% \mar{ME: Combining notation
	% 	from earlier corollaries --- introduce this notation
	% 	in a more formal fashion earlier on.}
By construction, such a factorization occurs if and only if the induced
morphism
$\Spf \Acirc \times_{\cC_{d,\ss,h}^{L/K,\fl}} \Spf B^{\circ}
	\to \Spf B^{\circ}$ 
	factors through
	$\Spf B^{\lambdau,\tau,\circ}$.

	Since the surjection $\Spf B^{\circ} \to
	\cC_{d,\ss,h}^{L/K,\fl}$ is representable by algebraic
	spaces, smooth, and surjective,
	the fibre product 
$\Spf \Acirc \times_{\cC_{d,\ss,h}^{L/K,\fl}} \Spf B^{\circ}$
is a formal algebraic space,
	and the surjection
$\Spf \Acirc \times_{\cC_{d,\ss,h}^{L/K,\fl}} \Spf B^{\circ}
	\to \Spf \Acirc$
	is representable by algebraic spaces, smooth, and surjective.
	We may thus find an open cover of 
$\Spf \Acirc \times_{\cC_{d,\ss,h}^{L/K,\fl}} \Spf B^{\circ}$
by affine formal algebraic spaces  $\Spf C^{\circ}$,
with $C^{\circ}$ being a faithfully flat $\varpi$-adically complete
$\Acirc$-algebra which is topologically of
finite type. 
	  (We are again applying Lemmas~\ref{lem: alem closed formal algebraic scheme
    adic*} and~\ref{alem: Noetherian flat affine formal algebraic
    spaces}.)
 % \cite[Lem.~8.18]{Emertonformalstacks}
	%   to deduce faithful flatness on the level of rings from smoothness 
	%   and surjectivity,
	%   and hence faithful flatness,
	%   of a morphism of affine formal algebraic spaces,
	%   and~\cite[\href{https://stacks.math.columbia.edu/tag/0ANT}{Tag 0ANT}]{stacks-project}  
	%   to deduce being topologically of finite type on the level
	%  of rings from a morphism of affine
	% formal algebraic spaces being of finite type.)
	If we write $C := C^{\circ}[1/p]$,
then the base-change property of Corollary~\ref{cor: Hodge and
  inertial decomposed together} % ~\ref{cor:inertial type decomposition}
	% and~\ref{cor:Hodge decomposition L/K case}
        shows that
	the hypothesized factorization occurs if and only
	if $C = C^{\lambdau,\tau}$, 
	% \mar{ME: Combining notation
	% 	from earlier corollaries --- introduce this earlier.}
	and also that $C^{\lambdau,\tau} = A^{\lambdau,\tau}\otimes_A C,$
	so that (since $C$ is faithfully flat over $A$) 
	$C^{\lambdau,\tau} = C$ if and only if $A^{\lambdau,\tau} = A$.
	In conclusion, we have shown that 
$\Spf A^{\circ} \to \cC_{d,\ss,h}^{L/K,\fl}$
factors through $\cC_{d,\ss,h}^{L/K,\fl,\lambdau,\tau}$
if and only if $A^{\lambdau,\tau} = A$, as required.

Finally, the uniqueness of~$\cC_{d,\ss,h}^{L/K,\fl,\lambdau,\tau}$ is
immediate from Proposition~\ref{prop: uniqueness of a flat stack with
  local ring points} below; the properties of being of finite
presentation,%^\mar{TG: finite presentation?}
of having affine diagonal,
and
of the morphism $\cC_{d,\ss,h}^{L/K,\fl,\lambdau,\tau}\to\cX_d$ being representable by algebraic spaces, proper, and of finite
presentation, are immediate from Proposition~\ref{prop: pst and pcrys
  stacks are padic} and Lemma~\ref{lem:finite presentation from closed immersion}.% \ref{alem: closed immersion is
  % lfp}.\mar{TG: check hypotheses of latter once it's filled in.}
\end{proof}

	% \mar{ME: Have to check that $\cC_{d,\ss,h}^{L/K,\fl,\lambdau}$ has required property}\mar{TG: also need to check
        %           uniqueness}









	  

	  






  %For each Hodge type~$\lambdau$ which is bounded
  %by~$h$, and each inertial type $\tau$,
  %we can define an $\cO$-flat subcategory fibred in groupoids
  %$\cC_{d,\ss,h}^{L/K,\fl,\lambdau,\tau}$ of
  %~$\cC_{d,\ss,h}^{L/K,\fl}$ by decreeing that for each
  %$p$-adically complete flat $\cO$-algebra~$\Acirc$ which is
  %topologically of finite type over~$\cO$, a $\Spf\Acirc$-valued point
  %of ~$\cC_{d,\ss,h}^{L/K,\fl}$ is a point
  %of~$\cC_{d,\ss,h}^{L/K,\fl,\lambdau,\tau}$ if and only
  %if~$\gMt_{\Acirc}$ has Hodge type~$\lambda$ and inertial type~$\tau$.
  %%\mar{TG: probably not correct, might need some stackification? ME: I think it's ok}}
%
%  %Since~$\cC_{d,\ss,h}^{L/K,\fl}$ is a flat
%  %$p$-adic formal algebraic stack of finite presentation
%  %over~$\Spf\cO$, we can choose a smooth surjection~$\Spf\Bcirc\to
%  %\cC_{d,\ss,h}^{L/K,\fl}$, where~$\Bcirc$ is topologically of
%  %finite type over~$\cO$.
%By Proposition~\ref{prop: good behaviour of Kisin
%    filtration}, for each~$0\le i\le h$, and each~$\sigma:K\into E$,
%   \[
%    e_\sigma((\Fil^i\varphi^*\gM_{\Bcirc}/E(u)\Fil^{i-1}\varphi^*\gM_{\Bcirc})\otimes_{\Bcirc}B)\]
%  is a finite projective $B := \Bcirc[1/p]$-module whose formation is compatible with
%  base change. In particular, there is a closed sub-formal scheme
%  of~$\Spf\Bcirc$ given by the locus for which ~$\gMt_{\Bcirc}$ has Hodge type~$\lambda$. \mar{TG: need to then say something to demonstrate why
%    this descends to a closed substack of~$\cC_{d,\ss,h}^{L/K,\fl}$,
%    and need to do inertial type.}
%
%  Finally, to check that~$\cC_{d,\ss,h}^{L/K,\fl}$ is the union
%  of the~$\cC_{d,\ss,h}^{L/K,\fl,\lambdau,\tau}$, we may pull back
%  to~$\Spf\Bcirc$. Since~$B$ %$B=\Bcirc[1/p]$
%  is Jacobson, it then suffices
%  to note that if~$\Acirc$ is a $\Bcirc$-algebra which is finite and
%  flat as an $\cO$-algebra, then ~$\gMt_{\Acirc}$ necessarily has
%  Hodge type~$\lambdau$ for some~$\lambdau$ bounded by~$h$, by
%  Theorem~\ref{thm: admits all descents if and only if
%    semistable}.\mar{TG: not true; I guess it is true if~$\Acirc$ is
%    local though, so maybe we should be doing something like arguing
%    with each~$B/\m_B^i$?}
%\end{proof}

\begin{prop}\label{prop: uniqueness of a flat stack with local ring points}
	Suppose that $\cY$ and $\cZ$ are two Noetherian
	%locally of finite type $p$-adic
	formal algebraic stacks, both lying over $\Spf \cO$ and both flat
	over $\cO$, and both embedded as closed substacks
	of a stack $\cX$ over $\Spec \cO$ % \mar{TG: should be $\Spf \cO$?}
        {\em (}in the usual sense
       that the inclusions of substacks $\cY, \cZ \hookrightarrow \cX$
       are representable by algebraic spaces, and induce closed immersions
when pulled back over any scheme-valued point of $\cX${\em )}.      
%whose diagonal is representable by algebraic spaces.  
Suppose also that each of $\cY_{\red}$ and $\cZ_{\red}$ are locally
of finite type over $\F$,
	and suppose further that, for any finite flat $\cO$-algebra~$\Acirc$,
	a morphism $\Spf \Acirc \to \cX$ factors through ~$\cY$ if
	and only if it factors through~ $\cZ$.  
	Then $\cY$ and $\cZ$ coincide.
	%\mar{ME: Put a uniqueness result here, in suff.\ generality for all applications}
\end{prop}
\begin{proof}%Since the diagonal of~$\cX$ is representable by algebraic spaces,
	The $\cO$-flat part $\cW := (\cY\times_{\cX} \cZ)_{\fl}$
        of the $2$-fibre product~$\cY\times_{\cX}\cZ$ embeds
	as a closed substack of each of $\cY$ and~$\cZ$,
	and it suffices to show that each of those embeddings
        are isomorphisms.
	The stack $\cW$
	inherits all of the hypotheses shared by $\cY$ and~$\cZ$,
	including the fact that, for any finite flat $\cO$-algebra~$\Acirc$,
	the $\Acirc$-valued points of $\cW$ coincide with 
	the $\Acirc$-valued points of each of $\cY$ and~$\cZ$.
        Thus, replacing the pair $\cY$, $\cZ$ by the pairs
        $\cW, \cY$ and $\cW, \cZ$ in turn, we may reduce to the case
  %we may replace~$\cY$ b ~$\cY\times_{\cX}\cZ$, and assumed
  when~$\cY$ is actually a closed substack of~$\cZ$.
  %\mar{ME: A problem: this fibre product is not obviously flat, and
%	  presumably is not in general.  Have to figure out how to deal
%	  with this. ME: This turned out to be easier than I anticipated!}

  Since~$\cZ$ is Noetherian, lies over $\Spf \cO$, and is flat over $\cO$,
  %locally of finite type $p$-adic formal stack,\mar{TG: this isn't an
  %  explicit assumption, so probably we need to elaborate a bit here?}
  we can find a smooth
  surjection $U\to\cZ$, whose source is a disjoint union of affine
  formal algebraic spaces $\Spf\Ccirc$, where~$\Ccirc$ is a flat
  $I$-adically complete Noetherian $\cO$-algebra, for some ideal
  $I$ which contains $\varpi$ and defines the topology on $\Ccirc$.

 % $\varpi$-adically complete $\cO$-algebra which is also topologically of finite type.
  The assumption on $\cZ_{\red}$ further implies that $\Ccirc/I$ 
  is a finite type $\F$-algebra.
  Then~$\Spf\Ccirc\times_{\cZ}\cY$ is a closed sub-formal
  algebraic space of~$\Spf\Ccirc$, and (since it is also~$\cO$-flat)
  it follows from Lemmas~\ref{lem: alem closed formal algebraic scheme
    adic*} and~\ref{alem: Noetherian flat affine formal algebraic
    spaces} that it is of the form~$\Spf\Bcirc$, where~$\Bcirc$ is an $\cO$-flat
  quotient of~$\Ccirc$, 
  % \mar{ME: In other contexts, we have given a bunch of citations
  %         to justify that closed formal subschemes translate to
  %         quotient rings in the Noetherian context; probably should
  %         give those again, or perhaps isolate the relevant statement
  %         and put it in the appendix. Note: tags 0ANQ and 0APT are the 
  % relevant citations.}
  %and is topologically of
  %finite type and flat over~$\cO$.
  again endowed with the $I$-adic topology.
  Since the property of being an isomorphism can be checked {\em fppf}
  locally on the target, we may replace the closed 
  embedding $\cY \hookrightarrow \cZ$ 
  by $\Spf \Bcirc \hookrightarrow \Spf \Ccirc.$
  In summary, we have a surjection $\Ccirc \to \Bcirc$ of flat,
  $I$-adically complete %, and topologically of finite type
  $\cO$-algebras, with $\Ccirc/I$ (and hence also $\Bcirc/I$) being 
  of finite type over~$\F$,
  and with the further property that any morphism $\Ccirc \to \Acirc$
  of $\cO$-algebras in which $\Acirc$ is finite flat over $\cO$ factors
  through $\Bcirc$; we then have to prove that this surjection is an
  isomorphism.

  Write~$B=\Bcirc[1/p]$, $C=\Ccirc[1/p]$;
  since~$\Bcirc$ and~$\Ccirc$ are flat over~$\cO$, it is enough to
  check that the surjection $C\to B$ is also injective. Now, we can
  embed~$C$ into the product of its localizations at all maximal
  ideals, and since~$C$ is Noetherian, it in fact embeds into the
  product of the completions of these localizations. It therefore
  embeds into the product of all of its local Artinian quotients~$A$. We
  claim that any such quotient is in fact obtained from a morphism
  ~$\Ccirc\to\Acirc$,
  where~$\Acirc$ is a finite flat~$\cO$-algebra; since any such
  morphism~$\Ccirc\to\Acirc$ factors through~$\Bcirc$ by assumption,
  it follows that any such~$A$ is a quotient of~$B$, so that 
  the surjection $C \to B$ is indeed injective.

It remains to prove the claim. Let~$\Acirc$ be the image of the
composite $\Ccirc~\to~C\to~A$; we need to show that~$\Acirc$ is a
finite $\cO$-algebra.  This follows from Lemma~\ref{lem:order in the court}
below.\end{proof}
%\mar{TG: ME to fill in, the point is to show
%  boundedness, and this should be shown on the Pontryagin dual side?} The residue field~$E'$ of~$A$ is a
%finite extension of~$E$, so the image of~$\Acirc$ in this residue
%field is a finite $\cO$-algebra.  %  That~$\Acirc$ is also finite follows
% from the topological Nakayama lemma (recalling that the maximal ideal~$\m_A$ of~$A$
% is nilpotent).
% \mar{ME: Not sure that the field case is easier
% 	than the general Artinian case; think about it. ME: Also
%  have to treat more general Noetherian case introduced above.}

%  It is enough to show
% that~$\Bcirc$ embeds into the product of all such~$\Acirc$.\mar{TG: if
% we inverted~$p$ we'd be Noetherian and could embed into the product of
% the completions and then the product of their (local) Artinian quotients, for
% example. Don't see why we can't basically do that - since~$\Bcirc$ is
% flat, isn't it enough to check after inverting~$p$??}
% \mar{ME: Add proof! TG: maybe done? ME: Yes, modulo marginal comments}
%\end{proof}
 
% Choose ~$h\ge
% 0$ so that~$\underline{\lambda}$ is bounded by~$h$ (this requires that all of the Hodge--Tate weights
% of~$\underline{\lambda}$ are non-negative; we make this assumption for
% now, but we will later remove it by a simple twisting
% argument).\mar{TG: make sure this twisting argument actually exists!}


% Let~$\cS=\cC_{d,\semis,h}$, and let~$\cC$ denote
% the category of locally Noetherian formal algebraic stacks
% over~$\cC_{d,\semis,h}$, so that an object of~$\Aff_\cC$ is a morphism
% $\Spf B\to \cC_{d,\semis,h}$, where~$\Spf B$ is a Noetherian affine
% formal algebraic space. In particular, $B$ is complete with respect to
% an ideal containing~$pB$. We write~$B_\fl$ for the quotient of~$B$ by
% its ideal of~$p$-power torsion elements. Then we set
% $F(\Spf B)=\Spf
% B_\fl\times_{\cC_{d,\semis,h}}\cC_{d,\semis,h}$.\mar{TG: pst notation
%   and so on. Note that this definition looks a bit stupid, but I'm
%   guessing the point is going to be that the HT weight locus (and
%   perhaps the type locus too, depending on how we define it) is not
%   closed unless we're in the flat setting.}
% \mar{TG: fix $h$ and~$L$ but say ultimately independent of
%   choice. Similarly deal with twisting somewhere. Put in a lemma for this!}

% To see that~$F$ satisfies the required properties, we must show
% that~$F(\Spf B)\to \Spf B$ is a closed immersion, and that if~$\Spf
% C\to\Spf B$ is representable by algebraic spaces and smooth, then
% $F(\Spf C)=\Spf C\times_{\Spf B}F(\Spf B)$. 



% (Include finite flat $\Zp$-algebra points.)

% \begin{itemize}
% \item Define pst and pcris just by fibre products.\mar{TG: probably
%     now taken care of above -- we have the stacks of BKF modules.}
% \item Define locus with a particular type - could do this on the flat
%   part, maybe, as there we probably have access to~$\gM/u\gM$?
%   Although maybe that suggests that we do it for~$\cC$ rather
%   than~$\cX$? Although even if we do that it's not immediately clear
%   to me that $\Gal(L/K)$ automatically acts on~$\gM/u\gM$, e.g.\ I
%   don't know if it's automatically the case that
%   $\Gal(L/K)$ is forced to stabilise~$\gMt$?
% \end{itemize}



% We will make use of the following basic result about
% scheme-theoretic images.
% % \begin{lem}\mar{TG: rather than this, just cite \cite[Prop
% %     10.5]{Emertonformalstacks} --- or rather put a version of that
% %     into the appendix, and then cite it.}
% %   \label{lem: scheme theoretic image inherits flatness and
% %     finiteness}Let~$\cC\to\cX$ be a proper morphism of Noetherian
% %   formal algebraic stacks over~$\Spf\cO$, and let~$\cZ$ be its
% %   scheme-theoretic image. Assume that $\cC$ is a $p$-adic formal
% %   algebraic stack which is of finite type and flat over~$\Spf
% %   \cO$. Then~$\cZ$ is also a $p$-adic formal algebraic which is of
% %   finite type flat over~$\Spf\cO$, and the morphism $\cC\to\cZ$ is
% %   proper and scheme-theoretically surjective.
% % \end{lem}
% % \begin{proof}
% %  
% %\end{proof}
% \begin{lem}\label{lem: lifting through proper surjection after a
%     finite extension}
%   Let ~$\cX\to\cZ$ be a morphism 
% of Noetherian $p$-adic formal
% algebraic stacks which are flat over~$\Spf \cO$
% which is representable by algebraic spaces, proper, and
% scheme-theoretically dominant.
% % \mar{TG: check that we have defined
%   % scheme-theoretic dominance in Appendix~\ref{app: formal algebraic
%   % stacks} }
% %Let~$\cO'$ be the ring
% %of integers in some finite extension~$E'/\Qp$. Then any morphism
% %$\Spf\cO'\to\cZ$ lifts, after possibly replacing~$E'$ with a further
% %finite extension, to a morphism $\Spf\cO'\to\cX$.\mar{TG: since we
% %  aren't assuming reduced, should this maybe be replaced by something
% %  like $\Acirc\subset(\Acirc)'$ with $(\Acirc)'$ a finite
% %  flat~$\cO$-algebra? TG: or ideally the statement with orders?}
% Given a morphism $\Spf A^{\circ} \to \cZ$,  with $\Acirc$
% being a finite flat $\cO$-algebra,
% we may find an order $(\Acirc)'$ in $\Acirc[1/p]$ containing $\Acirc$ 
% such that the composite $\Spf (\Acirc)' \to \Spf \Acirc
% \to \cZ$ lifts to a morphism $\Spf (\Acirc)' \to \cX$.
% \end{lem}
% \begin{proof}
%   \mar{TG: ME to prove! ME: Have to find correct statement}
% \end{proof}
% \section{Versal deformation rings}\label{subsec: versal deformation
% rings}
% We now return to the
% potentially semistable deformation rings of~\cite{MR2373358}, which we recalled
% in Section~\ref{subsec: notation and conventions}.

% \begin{df}
% We say that a Hodge
% type~$\underline{\lambda}$ is \emph{bounded by~$h$} if
% for each~$\sigma$, $i$ we have $\lambda_{\sigma,i}\in [-h,0]$.% \mar{TG: or perhaps the
%   % opposite sign, not worrying about conventions} 
% \end{df}



%\mar{ME: Can we move this result (4.2.6) to the end of \S 4.2, or the start
%	of \S 4.3?  As far as I can tell, the results in \S 4.2 
%	that comes after it (4.2.7, 4.2.8, and the subsequent
%	discussion) don't use 4.2.6, but rather just 4.2.5,
%       	so I think the logical flow would 
%	be clearer if they followed directly after 4.2.5. ME: Made this
%       move.}
% The following result is essentially the content
% of~\cite[\S3]{MR2745530} (although they are concerned with getting
% precise bounds on the size of~$s$, which is unimportant for us). We
% put ourselves in the following situation: let~$B$ be a finite flat
% $\cO$-algebra, and suppose that~$M_B$ is an \'etale
% $(\varphi,\Gamma)$-module with $B$-coefficients, with the property
% that the $E$-representation of~$G_K$ obtained from~$M_B$ is semistable with
% Hodge--Tate weights bounded by~$h$.
% %contained in the interval~$[-h,0]$.
% For each~$a\ge 1$, write~$M_{B,\infty}^a$ for the \'etale $\varphi$-module
% corresponding to the composite $\Spec (B/\varpi^a) \to\cX_d^a\to\cR_{\piflat,d}^a$.
% \begin{prop}\mar{TG: assuming descent action is always canonical, we
%     can remove this result}
%   \label{prop: Caruso Liu canonical action semistable}For any
%   fixed~$a,h$ and~$N\ge e(a+h)/(p-1)$ there is a
%   constant~$s'(a,h,N)\ge s(a,h,N)$ {\em (}where~$s(a,h,N)$ is as in
%   Lemma~{\em \ref{lem: Caruso Liu Galois action on Kisin})} with the
%   following property: if~$B$ and~$M_B$ are as above, and
%   if~$s\ge s'(a,h,N)$, then there is a projective Breuil--Kisin
%   module~$\gM$ of height at most~$h$
% %  \mar{ME: Projective, height $\leq h$?}
%   which is a lattice in ~$M_{B,\infty}^a$, for which the action
%   of~$G_{K_s}$
%   on~\[W(\C^\flat)_{B/\varpi^a}\otimes_{\gS_{B/\varpi^a}}\gM\subset
%     W(\C^\flat)_{B/\varpi^a}\otimes_{\A_{K,B}}M_B\]\emph{(}given by the
%   restriction of the natural action of~$G_K$ on the
%   $(\varphi,G_K)$-module
%   $W(\C^\flat)_{B/\varpi^a}\otimes_{\A_{K,B}}M_B$\emph{)} agrees with the
%   action coming from Lemma~{\em \ref{lem: Caruso Liu Galois action on
%     Kisin}}.
% \end{prop}
% \begin{proof}\mar{TG: I suspect this might be a bit dodgy! I think in
%     the theory of $(\varphi,\widehat{G})$-modules one has a
%     $\varphi$-twist to the embedding of~$\gM$, so it's possible that
%     whenever we talk about Breuil--Kisin--Fargues modules, we should
%     be doing $\otimes_{\gS_A,\varphi}\AAinf{A}$? Probably we should
%     compare to whatever \cite{2016arXiv160203148B} do, too.}
%   This is essentially the content of~\cite[Prop.\
%   3.3.1]{MR2745530}. More precisely, the existence of an appropriate
%   Kisin module is immediate from~\cite[Thm.\ 3.1.3]{MR2745530}. By the
%   statement of Lemma~\ref{lem: Caruso Liu Galois action on Kisin}, we
%   see that it is then enough to check that if~$s$ is sufficiently
%   large, then for all~$g\in G_{K_s}$ we have
%   $(g-1)(\gM)\subset u^N\Ainf\otimes_\gS\gM_{B,\infty}^a$, where the
%   action of~$g$ is that given by restricting the natural action on
%   $W(\C^\flat)_{\cO/\varpi^a}\otimes_{\A_{K}}M_B$ (note
%   that~$B/\varpi^a$ is a finite $\cO/\varpi^a$-module, so there is no
%   need to take a completed tensor product here). This is immediate
%   from~\cite[Lem.\ 3.2.1]{MR2745530} (taking~$\mathfrak{L}=\gM$, and
%   noting that in their notation, $\widehat{\mathfrak{L}}$ is an
%   extension of scalars of~$\gL$ to a subring
%   of~$\AAinf{\cO/\varpi^a}$; this extension of scalars involves a
%   twist by~$\varphi$, but since~$\varphi$ is a continuous automorphism
%   of~$\Ainf$, this is harmless).
% \end{proof}



%Suppose now that~$\F'/\F$ is a finite extension, and
%that~$\Spec\F'\to\cX_d$ is a finite type point, % \mar{TG:
%% maybe should have $\F'$, $\cO'$ here, and then annoyingly~$\cO''$ below?}
%with  corresponding Galois representation 
%$\rhobar:G_K\to\GL_d(\F')$. Let ~$\cO'$ be the ring of integers in
%some finite extension~$E'/E$ with residue field~$\F'$.  Then there is a versal morphism $\Spf
%R_{\rhobar}^{\square,\cO'}\to\cX_d$; indeed, by the theorem of Dee
%recalled in Section~\ref{subsec: phi gamma with coefficients},
%$R_{\rhobar}^{\square,\cO'}$ is isomorphic to the universal lifting ring for 
%% \mar{ME: $\cO'$ is being used here, before its actually defined. TG:
%%   rearranged to hopefully sort this out?}
%$(\varphi,\Gamma)$-modules lifting the $(\varphi,\Gamma)$-module
%corresponding to $\Spec\F'\to\cX_d$ (an isomorphism being given by any
%choice of ordered basis for the universal Galois lifting and universal
%$(\varphi,\Gamma)$-lifting), and it is clear from the definitions that the
%corresponding morphism $\Spf
%R_{\rhobar}^{\square,\cO'}\to\cX_d$ is a $\widehat{\GL}_d$-torsor, and
%in particular is versal.
%\mar{ME: There's probably some stupid issue here where choosing a basis
%	for a Galois rep'n is not the same as choosing a basis for a $(\varphi,
%	\Gamma)$-modules, so they are not literally the same ring. TG:
%     Good point! I guess the real problem here is that we never defined ``lifting
%    ring'', so we never say anything about bases, and we don't do that
%  on the $(\varphi,\Gamma)$-side either, although we did in
%  \cite{EGstacktheoreticimages}. Leaving this for now but I guess we
%  need to expand on this, probably actually defining the lifting rings
%would be a good start? I guess the correct way to set this up is that
%the two are related by fixing bases on each side
%(Galois/$(\varphi,\Gamma)$) for the universal objects? TG: OK wrote
%something, feel free to improve.}

 % Then the corresponding lifting ring
% $R^{\crys,\underline{\lambda},\cO'}_{\rhobar}$  is nonzero, and the
% versal morphism
% \mar{TG: we haven't set this up yet, but it is
  % presumably easy to show that this does give the right versal
  % rings?!}

The following lemma is no doubt standard, but we recall a proof
for the sake of completeness.

\begin{lemma}
\label{lem:order in the court}
If $R \to A$ is a morphism of $\cO$-algebras, with $R$ being $p$-adically complete
and Noetherian {\em (}e.g.\ a complete local $\cO$-algebra
with finite residue field{\em )}, and $A$ being 
a finite-dimensional $E$-algebra, then the image of $R$ in $A$
is finite over~$\cO$.
\end{lemma}
\begin{proof}
Since $R$ is Noetherian 
and $p$-adically complete, the Artin--Rees lemma shows
%\mar{TG: I would like this to have either more explanation or a reference to
%  Matsumura (I did take a look there but I didn't immediately see what
%  to cite\dots) ME: Prop.\ 10.13  of Atiyah--MacDonald looks like an ok
%citation to me.}
that the same is true of its image in~$A$ (see e.g.\ \cite[Prop.\ 10.13]{MR0242802}).  Thus,
by replacing $R$ by its image in $A$ and $A$ itself by the $E$-span
of this image, we are reduced to checking the following statement:
if $M$ is a $p$-adically complete and separated torsion free
$\cO$-module, and if $M[1/p]$ is finite-dimensional,
then $M$ is finite over $\cO$.  There are many ways to see this,
of course; here is one. 
% \mar{ME: What follows is perhaps overly detailed ---feel free to shorten it,
% or just comment out and replace by the claim that this is standard!
% TG: I'm happy with it, commenting out.}

Since $V:= M[1/p]$ is finite dimensional, we may find a finite spanning
set for this vector space contained in $M;$ if $L$ denotes its
$\cO$-span, then $L$ is a lattice in~$V$.  Since $L$ is compact
with respect to its $p$-adic topology,
the $p$-adic topology on $M$ induces the $p$-adic topology on $L$,
and in particular $L$ is closed in
the $p$-adic topology on~$M$. Thus $M/L$ is a $p$-adically separated submodule
of $V/L$.   Note that $V/L$ consists entirely of $p$-power torsion elements,
and thus that the same is true of $M/L$.

This implies that $\bigcap_n p^n\bigl((M/L)[p^{n+1}]\bigr) \subseteq \bigcap_n p^n(M/L) = 0.$
On the other hand,
$p^n\bigl((M/L)[p^{n+1}]\bigr) \subseteq  (M/L)[p] \subseteq (V/L)[p] \iso L/pL,$
and so $\{p^n\bigl((M/L)[p^{n+1}]\bigr)$ is a decreasing sequence of finite sets,
which thus eventually stabilizes.  We conclude that $p^n\bigl((M/L)[p^{n+1}]\bigr) = 0$
for some sufficiently large value of~$n$;
equivalently, $(M/L)[p^n] = (M/L)[p^{n+1}]$, and thus in fact
$(M/L)[p^n] = (M/L)[p^{\infty}] =  M/L.$
Consequently $L \subseteq M \subseteq p^{-n}L,$
showing that $M$ itself is finite over~$\cO$, as claimed.
\end{proof}

%\mar{ME: For a while now, we've been careful to write
%$\cX_{K,d}$ to make the field clear. (In the few cases where we hadn't, 
%I edited it to do this, for the sake of clarity and consistency. 
%Below  we revert to $\cX_d$ --- is
%this deliberate? TG: I don't believe so, so I've changed it throughout
%the rest of the section. ME: Thanks, commenting out  now}
\begin{defn}
  \label{defn: our actual pcris pst stacks}
  Let~$\tau$ be an inertial type,
  and let  $\lambdau$ be a Hodge type.
  To begin with, assume in addition 
  %that~$\lambda_{\sigma,i}\ge 0$ for all~$\sigma,i$.
  that~$\lambdau$ is effective, and
  choose~$h\ge 0$ such
  that~$\lambdau$ is bounded by~$h$, as well as a finite Galois
  extension ~$L/K$ %a finite Galois extension such that
  for which the kernel of~$\tau$ contains~$I_{L}$. 
  We then define~$\cX_{K,d}^{\crys,\lambda,\tau}$ \index{$\cX_{K,d}^{\crys,\lambda,\tau}$}
  to be the scheme-theoretic image of
  the morphism $\cC_{d,\crys,h}^{L/K,\fl,\lambdau,\tau}\to\cX_{K,d}$,
  and~$\cX_{K,d}^{\semis,\lambdau,\tau}$ to be the scheme-theoretic
  image of \index{$\cX_{K,d}^{\semis,\lambda,\tau}$}
  the morphism $\cC_{d,\semis,h}^{L/K,\fl,\lambdau,\tau}\to\cX_{K,d}$.
  (The characterization of these closed substacks of $\cX_{K,d}$
  given in Theorem~\ref{thm: existence of ss stack}
  will show that they are well-defined substacks of $\cX_{K,d}$,
  independently of the auxiliary choices of $h$ and $L/K$ that were used
  in their definition.)

%\mar{ME: Note that we are reverting here to $A$ being a $p$-adically
%complete ring, and abandoning our $\Acirc$  notation. TG: since we go
%back to that notation soon after, I've changed this. ME: Great, commenting out}
  In general (i.e.\ if $\lambdau$ is not effective), then we choose an
  integer~$h'$ so
  that~$\lambda'_{\sigma,i}:=\lambda_{\sigma,i}+h'\ge 0$ for
  all~$\sigma$ and all~$i$, and define
  $\cX_{K,d}^{\semis,\lambdau,\tau}$ to be the unique substack
  of~$\cX_{K,d}$ with the property that if~$B$ is a
  $\cO/\varpi^a$-algebra for some~$a\ge 1$, then a morphism $\rho:\Spec B\to \cX_{K,d}$
  factors through~ $\cX_{K,d}^{\semis,\lambdau,\tau}$ if and only if the
  morphism $\rho\otimes\epsilon^{h'}:\Spec B\to \cX_{K,d}$ factors
  through ~ $\cX_{K,d}^{\semis,\lambdau',\tau}$. (Here we are using the
  notation of Section~\ref{subsec: tensor product
  of phi gamma and duality}; more precisely, we are twisting by the
\'etale $(\varphi,\Gamma)$-module corresponding to~$\epsilon^{h'}$.)
  We define
  $\cX_{K,d}^{\crys,\lambdau,\tau}$ in the same way. The point of this
  definition is that if~$A^\circ$ is a finite flat $\cO$-algebra, then a
  representation $\rho:G_K\to\GL_d(\Acirc[1/p])$ is potentially
  semistable (resp.\ crystalline) of Hodge type~$\lambdau$ and
  inertial type~$\tau$ if and only if~$\rho\otimes\epsilon^{h'}$ is potentially
  semistable (resp.\ crystalline) of Hodge type~$\lambdau'$ and
  inertial type~$\tau$; again, Theorem~\ref{thm: existence of ss stack} below
  shows that these stacks are defined independently of the choices
  made in this definition, and in particular independently of the
  choice of~$h'$.
  % \mar{ME: Finish this
  %         by explaining the twist a little more carefully, with back-reference
  %         to \S~3 (which means material should probably be added there).
  %         Then remove what follows.}\mar{TG: I'm struggling to find a
  %         good way to put this in here, without (say) introducing a
  %         notion of a non-effective BKF module. My current inclination
  %         is not to do this, but to make a brief remark saying that it
  %         might be more natural; and then to prove the main results
  %         below by just twisting on the level of \'etale
  %         $\varphi$-modules, which should be pretty harmless - I think
  %         we can just remark that this gives an isomorphism of
  %         crystalline deformation rings etc. What do you think? TG:
  %         actually we can just twist here on the level of \'etale
  %         $\varphi$-modules!! TG: done, commenting out.}


  % By twisting by a sufficiently large power of the cyclotomic
  % character (or rather by the corresponding ``Tate twist''
  % construction for Breuil--Kisin--Fargues modules, \cite[Ex.\
  % 4.2.4]{2016arXiv160203148B}), we can and do assume
  % that~$\lambda_{\sigma,i}\ge 0$ for all~$\sigma,i$.
 % Choose~$h\ge 0$ such
 % that~$\lambdau$ is bounded by~$h$, and choose~$L/K$ a finite Galois extension such that the kernel
 % of~$\tau$ is contained in~$I_{L/K}$. 

  %Then we let~$\cX_d^{\crys,\lambda,\tau}$ be the scheme-theoretic image of
  %the morphism $\cC_{d,\crys,h}^{L/K,\fl,\lambdau,\tau}\to\cX_d$, and let
  %~$\cX_{d}^{\semis,\lambdau,\tau}$ be the scheme-theoretic image of
  %the morphism $\cC_{d,\semis,h}^{L/K,\fl,\lambdau,\tau}\to\cX_d$.
\end{defn}
\begin{rem}
  \label{rem: could use non effective BKF modules but no}Rather than
  introducing the twist by~$\epsilon^{h'}$, it would perhaps be
  more natural to work throughout with Breuil--Kisin and
  Breuil--Kisin--Fargues modules which are not necessarily
  $\varphi$-stable, as in~\cite{2016arXiv160203148B}. However, it
  would still frequently be useful to make the corresponding Tate
  twists (see for example the proof of ~\cite[Lem.\
  4.26]{2016arXiv160203148B}), and in particular we would need to make
  such twists in order to use the results
  of~\cite{MR2562795,EGstacktheoreticimages}, so we do not see any
  advantage in doing so.
\end{rem}

% \begin{rem}
%   \label{rem: Tate twisting BKF modules}It will be useful to
% allow the Hodge--Tate weights of our representations to be
% negative. With this in mind, for any integers~$h,h'\ge 0$, we define 
% \end{rem}
The rest of this section is devoted to proving some fundamental
properties of the stacks ~$\cX_{K,d}^{\crys,\lambda,\tau}$ and
~$\cX_{K,d}^{\semis,\lambdau,\tau}$.
%in particular, we will show that
%they are independent of the choice of~$h$ and~$L$.
%
%\begin{thm}\label{thm: existence of ss stack}Let~$\tau$ be an inertial type, and let  $\lambdau$ be a
%  Hodge type. Then there are uniquely determined closed
%  substacks~$\cX_{d}^{\crys,\lambdau,\tau}$ and
%  ~$\cX_{d}^{\semis,\lambdau,\tau}$ of~$\cX_d$ with the following properties:
%  \begin{enumerate}
%  \item $\cX_{d}^{\crys,\lambdau,\tau}$ and
%  ~$\cX_{d}^{\semis,\lambdau,\tau}$ are $p$-adic formal stacks which
%  are of finite type and flat over~$\Spf\cO$.
%\item\label{item: the finite flat points of the crystalline stack} If~$\Acirc$ is a finite flat~$\cO$-algebra, then
%  $\cX_d^{\semis,\lambda,\tau}(\Acirc)$ \emph{(}resp.\
%  $\cX_d^{\crys,\lambda,\tau}(\Acirc)$\emph{)} is precisely the
%  subgroupoid of~$\cX_d(\Acirc)$ consisting of
%  $G_K$-representations which are potentially semistable
%  \emph{(}resp.\ potentially crystalline\emph{)} of Hodge
%  type~$\lambdau$ and inertia type~$\tau$.
%  \end{enumerate}
We begin with an analysis of versal
rings. To this end, fix a point $\Spec \F'\to\cX_{K,d}$  %(\F')$ 
for some finite extension $\F'$ of $\F$, giving rise to a continuous representation $\rhobar: G_K \to \GL_d(\F')$.
% and also fix a $p$-adic Hodge type~$\underline{\lambda}$ and an
% inertial type~$\tau$.
Let $\cO'=W(\F')\otimes_{W(\F)}\cO$, the ring of integers in a finite extension $E'$ of $E$
having residue field $\F'$. Then the versal morphism $\Spf
R_{\rhobar}^{\square,\cO'}\to\cX_{K,d}$ of Proposition~\ref{prop:versal rings}
induces  
morphisms $\Spf R^{\crys,\underline{\lambda},\tau,\cO'}_{\rhobar}\to\cX_{K,d}$
and $\Spf R^{\semis,\underline{\lambda},\tau,\cO'}_{\rhobar}\to\cX_{K,d}$
(where these deformation rings are as in Section~\ref{subsec: notation and conventions}).


\begin{prop}
  \label{prop: versal rings for pst stacks}The morphisms  $\Spf R^{\crys,\underline{\lambda},\tau,\cO'}_{\rhobar}\to\cX_{K,d}$
and $\Spf R^{\semis,\underline{\lambda},\tau,\cO'}_{\rhobar}\to\cX_{K,d}$ factor
through versal morphisms $\Spf
R^{\crys,\underline{\lambda},\tau,\cO'}_{\rhobar}\to\cX_{K,d}^{\crys,\underline{\lambda},\tau}$
and  $\Spf
R^{\semis,\underline{\lambda},\tau,\cO'}_{\rhobar}\to\cX_{K,d}^{\semis,\underline{\lambda},\tau}$ respectively.
\end{prop}% \mar{TG: proof needs to be finished.}
% \mar{TG: have to prove this! Here's a guess as to how to
%   argue: we know that versal rings for scheme-theoretic images of
%   proper morphisms are given by scheme-theoretic images to the versal
%   ring. Hopefully it's easy to show that if you are flat over $\Zp$
%   then so are your versal rings. Then we should just have to show that the
%   two quotients of the universal deformation ring have the same points
% over finite flat $E$-algebras. This must more or less be what Mark
% does! Showing that the things that factor through the pst deformation
% rings factor through the sch-th image is easy. Conversely I think the
% key might be (as proved in Kisin) that after you pull back to an
% $E$-algebra, the morphism $\cC\to\cR$ becomes a closed immersion
% (presumably because it's a proper monomorphism - the latter coming
% from uniqueness of Kisin modules)? TG: wrote something, not making use
% of the unicity of Kisin modules, so I don't know whether I'm
% blundering\dots Or if indeed Lemma~\ref{lem: lifting through proper surjection after a
%     finite extension} is closely related to uniqueness of Kisin modules.}
\begin{proof}By definition, we can and do assume that~$\lambdau$ is
  effective. We give the proof for
  $\Spf R^{\crys,\underline{\lambda},\tau,\cO'}_{\rhobar}$, the
  argument in the semistable case being identical. % We use the notation
  % and setup of the proof of Theorem~\ref{thm: existence of ss
  %   stack}.\mar{TG: that's now a forward reference.}
  We begin by introducing 
  the fibre product $\widehat{C} :=
  \cC_{d,\crys,h}^{L/K,\fl,\lambdau,\tau}\times_{\cX_{K,d}}\Spf R_{\rhobar}^{\square,\cO'}$.
  Since 
  $\cC_{d,\crys,h}^{L/K,\fl,\lambdau,\tau}\to {\cX_{K,d}}$
  is proper and representable by algebraic spaces, 
  we see that $\widehat{C}$ is a formal algebraic space,
  and that the morphism
  \numequation
  \label{eqn:C to versal R}
  \widehat{C} \to
  \Spf R_{\rhobar}^{\square,\cO'}
  \end{equation}
  is proper and representable by algebraic spaces. The latter
  condition can be reexpressed by saying that~\eqref{eqn:C to versal R}
  is an {\em adic} morphism; thus we see that
  $\widehat{C}$ is a proper
  $\mathfrak m_{R_{\rhobar}^{\square,\cO'}}$-adic
  formal algebraic space over $\Spf R_{\rhobar}^{\square,\cO'}$.

  If we let $\Spf R$ denote the scheme-theoretic image of the
  morphism~\eqref{eqn:C to versal R}, 
  then
  by  Propositions~\ref{prop:X is an Ind-stack}, ~\ref{prop: HT
    weights are a closed condition} and Lemma~\ref{alem: scheme theoretic
      image of versal is versal}, % \mar{TG: we 
%     should perhaps verify that the hypotheses of this lemma hold TG: I
%   think the issue is just what the hypotheses of that lemma actually
%   are - right now there is a qcqs assumption on source and target and
%   not much more, and right now I think the things that we know
%   about~$\cX_d$ are just what is in Proposition~\ref{prop:X is an
%     Ind-stack}. TG: so I guess the question is maybe: do we have qcqs
%   for $\cX_d$? If not, how do we relax that lemma? I think almost by its construction
% as an Ind-stack, we'll know that $\cX_d$ is a colimit of qcqs alg.\ stacks, 
% and I think that's what the lemma requires --- see comment I added
% there. TG: the statement currently there does what we want here,
% commenting out.}
the versal morphism
  $\Spf R_{\rhobar}^{\square,\cO'}\to\cX_{K,d}$ induces a versal morphism
  $\Spf R\to\cX_{K,d}^{\crys,\underline{\lambda},\tau}$,
 % where~$\Spf R$ is
 % the scheme-theoretic image of the
 % morphism
 % \[\cC_{d,\crys,h}^{L/K,\fl}\times_{\cX_d}\Spf
 %   R_{\rhobar}^{\square,\cO'}
 %   \to \Spf
 %   R_{\rhobar}^{\square,\cO'}, \]this scheme-theoretic image being
 % defined via taking the scheme-theoretic images on Artinian quotients
 % (see ~\cite[Defn.\ 3.2.15]{EGstacktheoreticimages}).
  and so the assertion of the proposition may be rephrased as the claim that
  $R$ and $R^{\crys,\lambdau,\tau,\cO'}_{\rhobar}$
  coincide as quotients of $R_{\rhobar}^{\square,\cO'}.$

  Since
  $\cC_{d,\crys,h}^{L/K,\fl}$ is flat over~$\Spf\cO$, it
  follows from Lemma~\ref{lem:versal flatness} %~\cite[\href{https://stacks.math.columbia.edu/tag/0DR2}{Tag 0DR2}]{stacks-project} % (for example using Lemmas~\ref{lem: alem closed formal algebraic scheme
    % adic*} and~\ref{alem: Noetherian flat affine formal algebraic
    % spaces})
    %        \mar{TG: needs correct justification - using some flatness
    %          of versality and will also need to justify things being
    %          Noetherian.}
%  \mar{TG: need formal stacks version of~\cite[\href{https://stacks.math.columbia.edu/tag/0DR2}{Tag 0DR2}]{stacks-project}! ME: Done~}
  that $R$ is an~$\cO$-flat quotient of~$
  R_{\rhobar}^{\square,\cO'}$. Since
  $R^{\crys,\underline{\lambda},\tau,\cO'}_{\rhobar}$ is also an
  $\cO$-flat quotient of~$R_{\rhobar}^{\square,\cO'}$, we are reduced
  by Proposition~\ref{prop: uniqueness of a flat stack with local ring
            points}
  to showing that if~$\Acirc$ is any finite flat $\cO'$-algebra, then an
  $\cO'$-homomorphism $R_{\rhobar}^{\square,\cO'}\to\Acirc$ factors
  through~$R$ if and only if it factors
  through~$R^{\crys,\underline{\lambda},\tau,\cO'}_{\rhobar}$.
  % \mar{ME: I agree that this reduction is true, but probably also needs
  %         justification. TG: is it maybe now a special case of
  %         Prop~\ref{prop: uniqueness of a flat stack with local ring
  %           points}? At the very least the proof will give us this, I think.}

  To this end, note that if the morphism factors
  through~$R^{\crys,\underline{\lambda},\tau,\cO'}_{\rhobar}$, then by
  Theorem~\ref{thm: admits all descents if and only if semistable}
  there is an order~$(\Acirc)'\supseteq \Acirc$ in~$\Acirc[1/p]$ such that the induced
  morphism $\Spf(\Acirc)'\to \Spf R_{\rhobar}^{\square,\cO'}$ lifts to~$\widehat{C}$.
  %$\cC_{d,\crys,h}^{L/K,\fl}\times_{\cX_d}\Spf R_{\rhobar}^{\square,\cO'}$.
  Consequently the morphism
  $\Spf(\Acirc)'\to \Spf R_{\rhobar}^{\square,\cO'}$ factors
  through~$\Spf R$, and since $\Spf(\Acirc)'\to\Spf\Acirc$ is
  scheme-theoretically dominant, the morphism $\Spf\Acirc\to \Spf
  R_{\rhobar}^{\square,\cO'}$ also factors through~$\Spf R$.% \mar{TG:
    % anything more needed here?!}

    % \mar{ME: Proof sketch here}
  It remains to prove the converse, namely that if 
  $\Spf \Acirc \to \Spf R_{\rhobar}^{\square,\cO'}$ factors through
  $\Spf R$, then it factors 
  through~$R^{\crys,\underline{\lambda},\tau,\cO'}_{\rhobar}$.
  %We can and do assume that~$\Acirc$ is local.
  If we 
  write $S := R\otimes_{R_{\rhobar}^{\square,\cO'}} R^{\crys,
		  \underline{\lambda},\tau,\cO'}_{\rhobar}$
	  (a quotient of $R$),
	  then equivalently,
	  we wish to show that 
	  any morphism $\Spf \Acirc \to \Spf R$
	  in fact factors through $\Spf S$.  

  %The idea is to use~\cite[Lem. 5.4.15]{EGstacktheoreticimages},
  %but to apply this, we need to be considering Artinian rings.
  %So we should invert $p$, and there is a question of how to handle
  %this technically.  (I don't know if we need to use some Kisin-style
  %projectivity statement to pass from a proper formal scheme to
  %something algebraic and projective.)
  %\mar{ME: Our object is closed in the fixed $\pi^\flat$ object.  
%	  We algebraize the pull-back of the latter via Kisin,
%	  and then use Grothendieck Existence for algebraic spaces
%	  (citing Knutson if necessary) to get that our pulled-back
%	  object is projective too.}

  To do this, we will apply Lemma~\ref{lem: criterion for Artin to map to scheme theoretic
    image}, % \mar{ME: Add citation to whatever version of
  % 5.4.15 we add to this paper}
  not in the context of $\cO'$-algebras, but rather in the context of
  $E'$-algebras.  That is, we will invert $p$, and work with the rings
  $R[1/p]$ and $S[1/p]$ (or rather on certain completions of these
  rings).  We would also like to invert $p$ on the object
  $\widehat{C}$, but since the latter is a formal algebraic space,
  doing so directly would lead us into considerations of rigid
  analytic geometry that we prefer to avoid.  Thus, we begin by
  observing that by Lemma~\ref{lem: algebraising C} (together
  with~\cite[Thm.\ V.6.3]{MR0302647}, to pass from
  $\cC_{d,\semis,h}^{L/K}$ to its closed
  substack~$\cC_{d,\crys,h}^{L/K,\fl}$), $\widehat{C}$ can be promoted
  to a projective scheme over $\Spec R$.  That is, there is a
  projective morphism of schemes
  $C \to \Spec R_{\rhobar}^{\square,\cO'}$, with scheme-theoretic
  image equal to $\Spec R,$ whose
  $\mathfrak m_{R_{\rhobar}^{\square,\cO'}}$-adic completion is
  isomorphic to $\widehat{C}$.  Write $C[1/p] := C\otimes_{\cO'} E'.$
  Since scheme-theoretic dominance is preserved by flat base-change,
  the morphism $C[1/p] \to \Spec R[1/p]$ is proper and
  scheme-theoretically dominant.
  
 % But morally we write $X$ to denote the generic fibre (i.e.\ invert $p$)
 % of 
 % $\cC_{d,\crys,h}^{L/K,\fl}\times_{\cX_d}\Spf R_{\rhobar}^{\square,\cO'}.$ 
 % We then have a scheme-theoretically dominant morphism
 % $X \to \Spec R[1/p].$  
  Returning to our argument,
  we want to show, for any finite flat $\cO'$-algebra~$\Acirc$,
  that a morphism $R\to \Acirc$ necessarily factors through~$S$,
  or, equivalently,
  that the induced morphism
  $R[1/p] \to A := \Acirc[1/p]$ 
  factors through~$S[1/p]$.
  Since $A$ is a product of Artinian local $E$-algebras,
  we can check this assertion factor-by-factor;
  replacing $A$ by one of these factors (and $\Acirc$ by the image
  of $\Acirc$ in this factor)
  we may thus assume that $A$ is local.  
  We then let $\widehat{R[1/p]}$ denote the completion of
  $R$ at the kernel of its map to $A$, 
  and we let $\widehat{S[1/p]}$ denote the corresponding 
  completion of $S[1/p]$;
  by Artin--Rees, we also have that $\widehat{S[1/p]} = \widehat{R[1/p]} 
  \otimes_{R[1/p]} S[1/p]$, and so the map $R[1/p] \to A$ factors through
  $S[1/p]$
  if and only if the induced morphism $\widehat{R[1/p]} \to A$
  factors through $\widehat{S[1/p]}$.

  % \mar{ME: This para.\ has to be rewritten slightly to 
  %         cite version of 5.4.15 that's in the present paper.}
  % We may replace $A$ by some Artinian quotient of $\widehat{R[1/p]}$
  % for which $C_A \to \Spec A$ is scheme-theoretically dominant.
  % (Cf.\ the proof of
  % Lemma~\ref{lem:Artinian points of scheme theoretic images}.)
  To show the desired factorization, 
  we see from Lemma~\ref{lem: criterion for Artin to map to scheme theoretic
    image} % \mar{TG: there was a conflict of notation here where~$A$ was
  % being used for an arbitrary Artin algebra, rather than our fixed
  % $A$; I've tried to ungarble the next paragraph, but it should be
  % gone over again.}
  that it suffices to show that if $\widehat{R[1/p]} \to B$
  is a morphism to a finite Artinian $E$-algebra
  for which $C_B \to \Spec B$ admits a section,
  then this morphism factors through~$\widehat{S[1/p]}$.

Write~$B^\circ$ to denote the image of the composite
$R_{\rhobar}^{\square,\cO'}\to \widehat{R[1/p]} \to B$; by Lemma~\ref{lem:order in
the court}, we see that $B^{\circ}$ is
an order in $B$.  Let $Z\hookrightarrow C_{B^{\circ}}$ denote the scheme-theoretic
  closure of the section $\Spec B\to C_B$  in $C_{B^{\circ}}$. 
  %\mar{ME: Need to (re)introduce $A^{\circ}$.}
  This is proper over $\Spec \Bcirc$, flat over $\cO$,
  and irreducible of dimension one.
  Thus $Z$ is finite over $\Spec \Bcirc$, and hence $Z = \Spec (\Bcirc)'$
  for some order $(\Bcirc)'$ in $B$ containing~ $\Bcirc$.
  Thinking of $Z$ as a section of $C$ over $(\Bcirc)'$,
  we find that the induced morphism $\Spf (\Bcirc)' \to \cX_{K,d}$
  lifts to a morphism $\Spf (\Bcirc)' \to 
  \cC_{d,\crys,h}^{L/K,\fl,\lambdau, \tau}$. 
  Theorem~\ref{thm: admits all descents if and only if
    semistable} % \mar{TG: plus usual remark about extension of scalars
    % of deformation rings? TG: I decided to suppress this, commenting out.}
  then implies that the morphism $R_{\rhobar}^{\square, \cO'} 
  \to \Bcirc$ factors through $R_{\rhobar}^{\crys,\lambdau,\tau,\cO'}$,
  and thus the morphism $\widehat{R[1/p]} \to B$ does indeed
  factor through $\widehat{S[1/p]}$, as required.
%
%  But this follows (at least morally) from
%  Theorem~\ref{thm: admits all descents if and only if semistable}.
%
%
%   \mar{ME: Old material begins here}
% Conversely, if the morphism $\Spf\Acirc\to \Spf
%   R_{\rhobar}^{\square,\cO'}$  factors through~$\Spf S$, then there is
%   some~$(\Acirc)'$ as above such that the induced morphism $\Spf(\Acirc)'\to \Spf
%   R_{\rhobar}^{\square,\cO'}$ 
%   by Lemma~\ref{lem: lifting through proper surjection after a finite
%     extension}\mar{TG: again, idealised version of it} lifts to
%   $\cC_{d,\crys,h}^{L/K,\fl}\times_{\cX_d}\Spf
%   R_{\rhobar}^{\square,\cO'}$, and we conclude by another application of
%   Theorem~\ref{thm: admits all descents if and only if
%     semistable}.
  % \mar{TG: need to revisit this argument in light of
    % there being an $(\Acirc)'$}
\end{proof}


\begin{thm}\label{thm: existence of ss stack}Let~$\tau$ be an inertial type, and let  $\lambdau$ be a
  Hodge type. Then the closed substacks~$\cX_{K,d}^{\crys,\lambdau,\tau}$ and
  ~$\cX_{K,d}^{\semis,\lambdau,\tau}$ of~$\cX_{K,d}$ are  $p$-adic formal
  algebraic stacks % \mar{TG: make sure this is proved somewhere}
  which
  are of finite type and flat over~$\Spf\cO$, and  are
  uniquely % \mar{TG: tweak this slightly to get the uniqueness statement
  % correctly ME: Did something}
  determined as $\cO$-flat
  closed substacks of~$\cX_{K,d}$ by the following property:
%  $\cX_{d}^{\crys,\lambdau,\tau}$ and
%   ~$\cX_{d}^{\semis,\lambdau,\tau}$ are
% \item\label{item: the finite flat points of the crystalline stack}
  if~$\Acirc$ is a finite flat~$\cO$-algebra, then
  $\cX_{K,d}^{\semis,\lambda,\tau}(\Acirc)$ \emph{(}resp.\
  $\cX_{K,d}^{\crys,\lambda,\tau}(\Acirc)$\emph{)} is precisely the
  subgroupoid of~$\cX_{K,d}(\Acirc)$ consisting of
  $G_K$-representations which become potentially semistable
  \emph{(}resp.\ potentially crystalline\emph{)} of Hodge
  type~$\lambdau$ and inertia type~$\tau$ after inverting~$p$.
% \mar{ME: At some point,  double check that we've  explained
% that ``crys./ss'' means ``crys./ss'' after  inverting $p$. TG: I put
% this into the notation section at the very beginning, but of course we
% can add it here too if you like.}
\end{thm}
%\mar{ME: I think we should switch the order of this result and the next one;
%	then~(2) will just follow from the next theorem. ME: Done now}
\begin{proof}
We can and do assume that~$\lambdau$ is effective. By Propositions~\ref{prop:X is an Ind-stack}, ~\ref{prop: HT weights are a closed condition} and~\ref{aprop: scheme theoretic image is p adic finite
    type II}, $\cX_{K,d}^{\crys,\lambdau,\tau}$ and
  ~$\cX_{K,d}^{\semis,\lambdau,\tau}$ of~$\cX_{K,d}$ are  $p$-adic formal
  algebraic stacks  which
  are of finite type and flat over~$\Spf\cO$.  % \mar{TG: need to check
    % hypotheses once that result is proved. TG: it's immediate,
    % commenting out.}
  % , and at the moment it
    % doesn't include flatness.}% \mar{TG: is
%  there actually a slight
    % problem, in that as things are set up in the appendix at the
    % moment we might want~$\cX_d$ to be formal algebraic rather than
    % just Ind-algebraic? At worst we could presumably reorder the
    % sections, but maybe we don't really need that assumption there?}

  We now verify the claimed description of the $\Acirc$-valued points
  of $\cX_{K,d}^{\semis,\lambda,\tau}(\Acirc)$ and
  $\cX_{K,d}^{\crys,\lambda,\tau}(\Acirc)$, for finite flat $\cO$-algebras
  $\Acirc$. Since the argument is identical in either case, we give it
  in the semistable case. Any such algebra~$\Acirc$ is a product of
  finitely many finite flat local $\cO$-algebras, and so we immediately
  reduce to verifying the claim in the case when $\Acirc$ is
  furthermore local.  Since $\Acirc$ is then a complete local finite
  flat $\cO$-algebra, the claimed description follows from
  Proposition~\ref{prop: versal rings for pst stacks}. Indeed, the
  residue field~$\F'$ is a finite extension of~$\F$, and if as above
  we set~$\cO'=W(\F')\otimes_{W(\F)}\cO$, then the morphism
  $\Spf\Acirc\to\cX_{K,d}$ factors
  through~$\cX_{K,d}^{\semis,\underline{\lambda},\tau}$ if and only if it
  factors through the versal morphism
  $\Spf
  R^{\semis,\underline{\lambda},\tau,\cO'}_{\rhobar}\to\cX_{K,d}^{\semis,\underline{\lambda},\tau}$
  of Proposition~\ref{prop: versal rings for pst stacks}. In turn,
  this happens if and only if the corresponding $G_K$-representation
  is potentially semistable of Hodge type~$\lambdau$ and inertial
  type~$\tau$, by the defining property
  of~$R^{\semis,\underline{\lambda},\tau,\cO'}_{\rhobar}$. %\mar{ME:
    %Maybe elaborate? TG: how about this? ME: Looks fine, commenting out.}


It remains to show the claimed uniqueness statement. We prove in Corollary~\ref{cor:Xd is formal algebraic} below
	that $\cX_{K,d}$ is a Noetherian formal algebraic stack,
	and in Theorem~\ref{thm:Xdred is algebraic} below
	that $\cX_{d,\red}$ is of finite presentation over $\F$.
	The reader can easily confirm that
	Chapter~\ref{sec: families of extensions}
	contains no citation to the present chapter,
	and thus that those results are independent of the present ones;
	in particular, it is safe to invoke them here.
        These results imply that any closed substack $\cY$ of $\cX_{K,d}$
	is Noetherian (so that it makes sense to speak of $\cY$ being
	flat over $\cO$), and that $\cY_{\red}$ is of finite type
	over $\cF$.  The claimed uniqueness then follows from
	Proposition~\ref{prop: uniqueness of a flat stack with local ring
    points}.
%
%	The uniqueness statements are immediate from
%  Proposition~\ref{prop: uniqueness of a flat stack with local ring
%    points},
%so we need only prove existence.
%  We give the proof in the crystalline case, the proof in the
%  semistable case being formally identical. Note that as in
%  the proof of Proposition~\ref{prop: HT weights are a closed
%    condition}, the unicity of $\cX_{d}^{\crys,\lambdau,\tau}$ is
%  immediate from Proposition~\ref{prop: uniqueness of a flat stack
%    with local ring points}. %  condition~\eqref{item: the finite flat points of the
%  %   crystalline stack}, together with the requirement that it is
%  % finite type and flat over~$\Spf\cO$. 
%
%%   By twisting by a sufficiently large power of the cyclotomic
%%   character (or rather by the corresponding ``Tate twist''
%%   construction for Breuil--Kisin--Fargues modules, \cite[Ex.\
%%   4.2.4]{2016arXiv160203148B}), we can and do assume
%%   that~$\lambda_{\sigma,i}\ge 0$ for all~$\sigma,i$. Choose~$h\ge 0$ such
%%   that~$\lambdau$ is bounded by~$h$, and choose~$L/K$ a finite Galois extension such that the kernel
%%   of~$\tau$ is contained in~$I_{L/K}$. Let~$\cX_d^{\crys,\lambda,\tau}$
%%   be the scheme-theoretic image of the morphism
%%   $\cC_{d,\crys,h}^{L/K,\fl,\lambdau,\tau}\to\cX_d$. By Propositions~\ref{prop: pst
%%     and pcrys stacks are padic}, \ref{prop: HT weights are a closed condition} and~\ref{aprop: scheme theoretic image is p adic finite
%%     type II},\mar{TG: this latter result has not yet been proved, and
%%     we should probably put references in for its hypotheses - also
%%     might want to add flatness to its statement.} this is $\cO$-flat and of finite type over~$\Spf\cO$,
%%   so we need only check that ~\eqref{item: the finite flat points of the
%%     crystalline stack} holds.
%
%% \mar{ME: I think this material should just be commented out.}
%  In one direction, note that if $\Spf\Acirc\to\cX_d$ factors through
%  ~$\cX_d^{\crys,\lambda,\tau}$, then by Lemma~\ref{lem: lifting through proper surjection after a
%    finite extension}\mar{TG: or rather by the imagined improvement to
%  it indicated in the margin above - perhaps this should also be in
%  terms of orders?} there is a finite
%flat~$\Acirc$-algebra $(\Acirc)'$ such that the composite morphism
%$\Spf(\Acirc)'\to\Spf\Acirc\to\cX_d$ factors
%through~$\cC_{d,\crys,h}^{L/K,\fl,\lambdau,\tau}$. It follows from
%Theorem~\ref{thm: admits all descents if and only if semistable} that
%the corresponding $G_K$-representation is potentially crystalline of Hodge
%type~$\lambdau$ and inertial type~$\tau$, and so the same holds for
%the $G_K$-representation corresponding to $\Spf\Acirc\to\cX_d$.
%
%Conversely, if $\Spf\Acirc\to\cX_d$ corresponds to a
%$G_K$-representation which is potentially crystalline of Hodge
%type~$\lambdau$ and inertial type~$\tau$, then by Theorem~\ref{thm:
%  admits all descents if and only if semistable} there is an
%order~$(\Acirc)'$ of~$A:=\Acirc[1/p]$ which contains~$\Acirc$ and has
%the property that  we
%can lift the morphism $\Spf(\Acirc)'\to\cX_d$
%to~$\cC_{d,\crys,h}^{L/K,\fl,\lambda,\tau}$. In particular the
%morphism $\Spf(\Acirc)'\to\cX_d$ factors through the
%scheme-theoretic image ~$\cX_d^{\crys,\lambda,\tau}$, and since
%$\Spf(\Acirc)'\to\Spf\Acirc$ is scheme-theoretically dominant, the
%morphism $\Spf\Acirc\to\cX_d$ also factors through
%~$\cX_d^{\crys,\lambda,\tau}$, as required. \mar{TG: insert some justification?}
\end{proof}

% \begin{defn}Let $\cX_{d,\semis,h}$ be the scheme-theoretic image of
%   $\cC_{d,\semis,h}\to\cX_d$.% \mar{TG: appropriate citations so that
%     % this makes sense.}  
% \end{defn}

%The remaining results in this section rely on an analysis of versal
%rings. To this end, fix a point $\Spec \F'\to\cX_d(\F')$ 
%for some finite extension $\F'$ of $\F$, giving rise to a continuous representation $\rhobar: G_K \to \GL_d(\F')$.
%% and also fix a $p$-adic Hodge type~$\underline{\lambda}$ and an
%% inertial type~$\tau$.
%Let $\cO'=W(\F')\otimes_{ W(\F)}\cO$, the ring of integers in a finite extension $E'$ of $E$
%having residue field $\F'$. Then the versal morphism $\Spf
%R_{\rhobar}^{\square,\cO'}\to\cX_d$ of Proposition~\ref{prop:versal rings}
%induces  
%morphisms $\Spf R^{\crys,\underline{\lambda},\tau,\cO'}_{\rhobar}\to\cX_d$
%and $\Spf R^{\semis,\underline{\lambda},\tau,\cO'}_{\rhobar}\to\cX_d$
%(where these deformation rings are as in Section~\ref{subsec: notation and conventions}).
%
%\begin{prop}
%  \label{prop: versal rings for pst stacks}The morphisms  $\Spf R^{\crys,\underline{\lambda},\tau,\cO'}_{\rhobar}\to\cX_d$
%and $\Spf R^{\semis,\underline{\lambda},\tau,\cO'}_{\rhobar}\to\cX_d$ factor
%through versal morphisms $\Spf
%R^{\crys,\underline{\lambda},\tau,\cO'}_{\rhobar}\to\cX_d^{\crys,\underline{\lambda},\tau}$
%and  $\Spf
%R^{\semis,\underline{\lambda},\tau,\cO'}_{\rhobar}\to\cX_d^{\semis,\underline{\lambda},\tau}$ respectively.
%\end{prop}% \mar{TG: proof needs to be finished.}
%% \mar{TG: have to prove this! Here's a guess as to how to
%%   argue: we know that versal rings for scheme-theoretic images of
%%   proper morphisms are given by scheme-theoretic images to the versal
%%   ring. Hopefully it's easy to show that if you are flat over $\Zp$
%%   then so are your versal rings. Then we should just have to show that the
%%   two quotients of the universal deformation ring have the same points
%% over finite flat $E$-algebras. This must more or less be what Mark
%% does! Showing that the things that factor through the pst deformation
%% rings factor through the sch-th image is easy. Conversely I think the
%% key might be (as proved in Kisin) that after you pull back to an
%% $E$-algebra, the morphism $\cC\to\cR$ becomes a closed immersion
%% (presumably because it's a proper monomorphism - the latter coming
%% from uniqueness of Kisin modules)? TG: wrote something, not making use
%% of the unicity of Kisin modules, so I don't know whether I'm
%% blundering\dots Or if indeed Lemma~\ref{lem: lifting through proper surjection after a
%%     finite extension} is closely related to uniqueness of Kisin modules.}
%\begin{proof}We give the proof for
%  $\Spf R^{\crys,\underline{\lambda},\tau,\cO'}_{\rhobar}$, the
%  argument in the semistable case being identical. We use the notation
%  and setup of the proof of Theorem~\ref{thm: existence of ss
%    stack}. By~\cite[Lem.\ 3.2.16]{EGstacktheoreticimages}, the
%  morphism
%  $\Spf
%  R_{\rhobar}^{\square,\cO'}\to\cX_d^{\crys,\underline{\lambda},\tau}$
%  factors through a versal morphism
%  $\Spf S\to\cX_d^{\crys,\underline{\lambda},\tau}$, where~$\Spf S$ is
%  the scheme-theoretic image of the
%  morphism
%  \[\cC_{d,\crys,h}^{L/K,\fl}\times_{\cX_d}\Spf
%    R_{\rhobar}^{\square,\cO'} \to \Spf
%    R_{\rhobar}^{\square,\cO'}, \]this scheme-theoretic image being
%  defined via taking the scheme-theoretic images on Artinian quotients
%  (see ~\cite[Defn.\ 3.2.15]{EGstacktheoreticimages}). Since
%  $\cC_{d,\crys,h}^{L/K,\fl}$ is flat over~$\Spf\cO$, it
%  follows (for example using~ \cite[Lem.~8.18]{Emertonformalstacks})
%	   \mar{TG: needs justification?}
%  that $S$ is an~$\cO$-flat quotient of~$
%  R_{\rhobar}^{\square,\cO'}$. Since
%  $R^{\crys,\underline{\lambda},\tau,\cO'}_{\rhobar}$ is also an
%  $\cO$-flat quotient of~$R_{\rhobar}^{\square,\cO'}$, we are reduced
%  to showing that if~$\Acirc$ is a finite flat $\cO'$-algebra, then an
%  $\cO'$-homomorphism $R_{\rhobar}^{\square,\cO'}\to\Acirc$ factors
%  through~$S$ if and only if it factors
%  through~$R^{\crys,\underline{\lambda},\tau,\cO'}_{\rhobar}$.
%  \mar{ME: I agree that this reduction is true, but probably also needs
%	  justification.}
%
%  To this end, note that if the morphism factors
%  through~$R^{\crys,\underline{\lambda},\tau,\cO'}_{\rhobar}$, then by
%  Theorem~\ref{thm: admits all descents if and only if semistable}
%  there is an order~$(\Acirc)'$ in~$\Acirc[1/p]$ such that the induced
%  morphism $\Spf(\Acirc)'\to \Spf R_{\rhobar}^{\square,\cO'}$ lifts to
%  $\cC_{d,\crys,h}^{L/K,\fl}\times_{\cX_d}\Spf
%  R_{\rhobar}^{\square,\cO'}$. Consequently the morphism
%  $\Spf(\Acirc)'\to \Spf R_{\rhobar}^{\square,\cO'}$ factors
%  through~$\Spf S$, and since $\Spf(\Acirc)'\to\Spf\Acirc$ is
%  scheme-theoretically dominant, the morphism $\Spf\Acirc\to \Spf
%  R_{\rhobar}^{\square,\cO'}$ also factors through~$\Spf S$.
%
%  \mar{ME: Proof sketch here}
%  It remains to prove the converse, namely that if 
%  $\Spf \Acirc \to \Spf R_{\rhobar}^{\square,\cO'}$ factors through
%  $\Spf S$, then it factors 
%  through~$R^{\crys,\underline{\lambda},\tau,\cO'}_{\rhobar}$.
%  The idea is to use~\cite[Lem. 5.4.15]{EGstacktheoreticimages},
%  but to apply this, we need to be considering Artinian rings.
%  So we should invert $p$, and there is a question of how to handle
%  this technically.  (I don't know if we need to use some Kisin-style
%  projectivity statement to pass from a proper formal scheme to
%  something algebraic and projective.)
%
%  But morally we write $X$ to denote the generic fibre (i.e.\ invert $p$)
%  of 
%  $\cC_{d,\crys,h}^{L/K,\fl}\times_{\cX_d}\Spf R_{\rhobar}^{\square,\cO'}.$ 
%  We then have a scheme-theoretically dominant morphism
%  $X \to \Spec S[1/p].$   We want to show that the morphism
%  $R_{\rhobar}^{\square,\cO'}[1/p] \to A := \Acirc[1/p]$, which
%  factors through $S[1/p]$, factors 
%  through~$R^{\crys,\underline{\lambda},\tau,\cO'}_{\rhobar}[1/p]$.
%  We first replace $A$ by some Artinian quotient of $S[1/p]$
%  for which $X_A \to \Spec A$ is scheme-theoretically dominant.
%  (Cf.\ the proof of
%  Lemma~\ref{lem:Artinian points of scheme theoretic images}.)
%  To show the desired factorization, 
%  we see from~\cite[Lem. 5.4.15]{EGstacktheoreticimages}
%  that it suffices to show that if $A$ is an Artinian quotient
%  of $S[1/p]$ for which $X_A \to A$ admits a section,
%  then $S[1/p]\to A$ factors 
%  through~$R^{\crys,\underline{\lambda},\tau,\cO'}_{\rhobar}[1/p]$.
%  But this follows (at least morally) from
%  Theorem~\ref{thm: admits all descents if and only if semistable}.
%
%
%  \mar{ME: Old material begins here}
%Conversely, if the morphism $\Spf\Acirc\to \Spf
%  R_{\rhobar}^{\square,\cO'}$  factors through~$\Spf S$, then there is
%  some~$(\Acirc)'$ as above such that the induced morphism $\Spf(\Acirc)'\to \Spf
%  R_{\rhobar}^{\square,\cO'}$ 
%  by Lemma~\ref{lem: lifting through proper surjection after a finite
%    extension}\mar{TG: again, idealised version of it} lifts to
%  $\cC_{d,\crys,h}^{L/K,\fl}\times_{\cX_d}\Spf
%  R_{\rhobar}^{\square,\cO'}$, and we conclude by another application of
%  Theorem~\ref{thm: admits all descents if and only if
%    semistable}.% \mar{TG: need to revisit this argument in light of
%    % there being an $(\Acirc)'$}
%\end{proof}

We will make use of the following corollary in Chapter~\ref{sec:properties}.
\begin{cor}
  \label{cor: crystalline deformation rings are effective
    versal}
% \mar{ME: Rephrased this to avoid assuming existence of a lift,
% 	b/c I don't think that's relevant/necessary.}
For any Hodge type $\underline{\lambda}$, and
%Suppose that~$\rhobar$ admits a crystalline lift $\rho:G_K\to\GL_d(\cO')$ of Hodge
 % type ~$\underline{\lambda}$ for some Hodge
 % type~$\underline{\lambda}$. Then
  for any~$a\ge 1$,
  % \mar{ME: Shouldn't we introduce an auxiliary $\cO'$ here, as we did
  %         above? TG: done, commenting out.}
  the corresponding morphism
  $\Spf R^{\crys,\underline{\lambda},\cO'}_{\rhobar}/\varpi^a \to
  \cX_{K,d}^a$ is effective, i.e.\ is induced by a morphism $\Spec R^{\crys,\underline{\lambda},\cO'}_{\rhobar}/\varpi^a \to
  \cX_{K,d}^a$.
\end{cor}% \mar{TG: add proof; should be immediate as we now have
  % algebraic stacks}
\begin{proof}By Proposition~\ref{prop: versal rings for pst stacks},
  the morphism $\Spf R^{\crys,\underline{\lambda},\cO'}_{\rhobar}/\varpi^a \to
  \cX_{K,d}^a$ factors through a versal morphism $\Spf R^{\crys,\underline{\lambda},\cO'}_{\rhobar}/\varpi^a \to
  \cX_{K,d}^{\crys,\lambda}\times_{\Spf\cO}\Spec\cO/\varpi^a$. Since~$\cX_{K,d}^{\crys,\lambda}$
  is a $p$-adic formal algebraic stack (by Theorem~\ref{thm: existence of ss stack})
  the base-change
  $\cX_{K,d}^{\crys,\lambda}\times_{\Spf\cO}\Spec\cO/\varpi^a$ is an algebraic
  stack, and so the result follows from~\cite[\href{https://stacks.math.columbia.edu/tag/07X8}{Tag 07X8}]{stacks-project}.  
\end{proof}

Finally, we can compute the dimensions of our potentially crystalline
and semistable stacks. It is presumably possible to develop the
dimension theory of $p$-adic formal algebraic stacks in some
generality, but we do not need to do so, as we can instead work with
the special fibres, which are algebraic stacks.
\begin{thm}%\mar{TG: put in pst pcris versions, say that it is uniquely
%   determined by this property.}
  \label{thm: dimension of ss stack}  % it has regular generic fibre, 
The algebraic stacks~$\overline{\cX}_{K,d}^{\crys,\lambdau,\tau}:=\cX_{K,d}^{\crys,\lambdau,\tau}\times_{\Spf\cO}\Spec \F$ and
  ~$\overline{\cX}_{K,d}^{\semis,\lambdau,\tau}:=\cX_{K,d}^{\semis,\lambdau,\tau}\times_{\Spf\cO}\Spec \F$ are equidimensional of dimension \[\sum_{\sigma}\#\{1\le i<j\le
  d|\lambda_{\sigma,i}>\lambda_{\sigma,j}\}.\]%
In particular, if~$\lambdau$ is regular, then the algebraic stacks~$\overline{\cX}_{K,d}^{\crys,\lambdau,\tau}$ and
  ~$\overline{\cX}_{K,d}^{\semis,\lambdau,\tau}$ are equidimensional of
  dimension~$[K:\Qp]d(d-1)/2$.   % $\cX_{d,\semis,h}(\Spf\Zpbar)$
  % \emph{(}resp.\ \emph{)} is naturally equivalent to the
  % groupoid of $\Zpbar$-lattices in potentially semistable \emph{(}resp.\
  % potentially crystalline\emph{)} $d$-dimensional
  % $\Qpbar$-representations of~ $G_K$ of Hodge type~$\lambdau$ and
  % inertial type~$\tau$. 
\end{thm}% \mar{TG: not sure if we actually have dimension theory of
%   $p$-adic formal algebraic stacks set up anywhere? I presume that
%   once we do this should be a legitimate argument, though. TG: in fact
% maybe we will just say that they are flat and we can make assertions
% about the dimensions of the special fibres. TG: change notation to
% special fibres, have forgotten what that is.}
\begin{proof}Once again, we give the argument in the crystalline case,
  the semistable case being identical. Write
  $d_{\lambdau}:=\sum_{\sigma}\#\{1\le i<j\le
  d|\lambda_{\sigma,i}>\lambda_{\sigma,j}\}$.
  % Since~$\cX_{d}^{\crys,\lambdau,\tau}$ is a $p$-adic formal algebraic
  % stack which is flat and of finite type over~$\Spf\cO$,
  % it suffices to show that t
 The
  algebraic stack
  $\cX_{K,d}^{\crys,\lambdau,\tau}\times_{\Spf\cO}\Spec \F$
  is of finite type over ~$\F$, and we need to show that it is
  equidimensional of dimension~$d_{\lambdau}$.
  Let~$x:\Spec \F'\to\cX_{K,d}^{\crys,\lambdau,\tau}$ be a finite type
  point corresponding to a Galois representation~$\rhobar$,
with corresponding versal morphism
  $\Spf
  R^{\crys,\underline{\lambda},\cO'}_{\rhobar}/\varpi\to\overline{\cX}_{K,d}^{\crys,\lambdau,\tau}
  $. % \mar{TG: shouldn't this rather be versal to the special fibre $\cX_{K,d}^{\crys,\lambdau,\tau}\times_{\Spf\cO}\Spec \F$, and
  % then we should make a few corresponding changes below? Maybe we
  % should have some abbreviated notation for this special fibre if so,
  % e.g.\ $(\cX_{K,d}^{\crys,\lambdau,\tau})_{\F}$ or $\overline{\cX}_{K,d}^{\crys,\lambdau,\tau}$}
%This versal morphism is easily seen % \mar{TG: maybe even something
%    % we can reference above about torsors once versal ring stuff is fixed?}
%  to be a torsor for the formal
%\mar{ME: Not a torsor!}
We have the fibre product 
% \mar{ME: Fibre product blunder here, has to be fixed TG: I think it's
%   less this and more that we should add slightly more, and I think I
%   know what to do.}
  $$\Spf R^{\crys,\underline{\lambda},\cO'}_{\rhobar}/\varpi
  \times_{\widehat{(\overline{\cX}_{K,d}^{\crys,\lambdau,\tau})}_x} 
  \Spf R^{\crys,\underline{\lambda},\cO'}_{\rhobar}/\varpi
\iso \widehat{\GL}_{d,R^{\crys,\underline{\lambda},\cO'}_{\rhobar}/\varpi,1},$$
where $\widehat{\GL}_{d,R^{\crys,\underline{\lambda},\cO'}_{\rhobar}/\varpi,1}$
denotes the completion of 
$({\GL}_{d})_{R^{\crys,\underline{\lambda},\cO'}_{\rhobar}/\varpi}$
along the identity element in its special fibre.
  By~\cite[Lem.\ 2.40]{EGcomponents}, % \mar{TG: I thought that we
%     needed to keep this citation here because this wasn't in the
%     Stacks project, but isn't it actually
%     \cite[\href{https://stacks.math.columbia.edu/tag/0DSB}{Tag
%       0DSB}]{stacks-project} ? If so we should probably cite this both
%   here and in the other use of this lemma, and then there would only
%   be one citation of \cite{EGcomponents} remaining, which I suspect we
% might also be able to get rid of.}
it is therefore enough to recall that since
  $ R^{\crys,\underline{\lambda},\cO'}_{\rhobar}$ is equidimensional
  of dimension~$d_{\lambdau}+1$, $ R^{\crys,\underline{\lambda},\cO'}_{\rhobar}/\varpi$ is
  equidimensional of dimension~$d_{\lambdau}$ (see~\cite[Lem.\ 2.1]{MR3248725}).
\end{proof}
% \mar{TG: needs proof - although actually I might be inclined to wait
%   and instead do this together with types etc later on. Also put in
%   finite flat $E$-algebra case?}


% \mar{ME: A technical/expositional point:  the ring is still defined even if it
% 	is zero; it just gives us the empty scheme mapping to $\cX_d^a$.
% 	So we really need to assume the existence of a crystalline lift
% 	at this point? We could just introduce the auxiliary extensions
% $E'/E$ and $\cO/\cO'$. TG: agreed, commenting out.}
% \begin{prop}
%   \label{prop: crystalline deformation ring factors through the scheme
%     theoretic image}Suppose that~$\underline{\lambda}$ is bounded by~$h$. Fix~$a\ge 1$, and fix~$s$ sufficiently large in
%   the sense of Lemma~{\em \ref{lem: Caruso Liu Galois action on Kisin}}
%   and 
%   Propositions~%{\em \ref{prop: Caruso Liu canonical action semistable}}
%   and~{\em \ref{prop: morphism C to Xs}}. Then for every % discrete
%   % \mar{ME: *Every* Artinian quotient is discrete, right?}
%   Artinian quotient~$A^\crys$ of $
%   R^{\crys,\underline{\lambda},\cO'}_{\rhobar}/\varpi^a$, the
%   composite \[\Spec A^\crys\hookrightarrow \Spf R^{\crys,\underline{\lambda},\cO'}_{\rhobar}/\varpi^a \to
%   \cX_d^a\]factors through a morphism $\Spec A^\crys\to \cC^a_{\piflat,s,d,h}$ 
% {\em (}where $\cC^a_{\piflat,s,d,h}$ denotes the $2$-Cartesian product
% of Definition~{\em \ref{def:2-Cartesian})}.
% % the morphism
%   % $\Spf R^{\crys,\underline{\lambda},\cO'}_{\rhobar}/\varpi^a \to
%   % \cX_d^a$ factors through~$\cZ^a_{d,K_s,h}$.
% \end{prop}
% \begin{proof}We make use of the results
%   of~\cite[\S4]{MR3449204}. Define a quotient~$R^{\underline{\lambda},\semis,\cO'}_{\rhobar}$
%   of~$R^{\square,\cO'}_{\rhobar}$ as follows: if~$A$ is an Artinian
%   local $\cO'$-algebra with residue field~$\F'$, then a morphism
%   $R^{\square,\cO'}_{\rhobar}\to A$ factors
%   through~$R^{\underline{\lambda},\semis,\cO'}_{\rhobar}$ if and only
%   if the corresponding representation $\rho_A:G_K\to\GL_d(A)$ satisfies:
%   \begin{itemize}
%   \item there exists a surjection of~$\cO'$-algebras $B\to A$,
%     where~$B$ is finite flat over $\cO'$, and a
%     representation~$\rho_B:G_K\to\GL_d$ lifting~$\rho_A$, such
%     that~$\rho_B\otimes_{\Zp}\Qp$ is semistable of Hodge type~$\underline{\lambda}$.
%   \end{itemize}


%   The existence of this quotient is straightforward: in the case
%   that~$\rhobar$ admits a universal deformation ring, it
%   is~\cite[Prop.\ 4.23]{MR3449204}, and for general~$\rhobar$ an
%   identical argument works for lifting
%   % \mar{ME: Just for consistencty in terminology:  ``framed deformation ring''
%   %         = ``lifting ring''?}
%   rings. By~\cite[Thm.\ 4.25]{MR3449204} (or rather the corresponding
%   statement for lifting rings, which again has an identical
%   proof), 
% %  \mar{ME: Have we explained that $E'$ equals frac.\ field of $\cO'$?}
%   if $E''/E'$ is a finite extension with ring of
%   integers~$\cO''$, then a homomorphism of $\cO'$-algebras
%   $R_{\rhobar}^{\square,\cO'}\to\cO''$ factors through
%   $R^{\underline{\lambda},\semis,\cO'}_{\rhobar}$ if and only if the
%   corresponding Galois representation is semistable of Hodge
%   type~$\underline{\lambda}$.

%   Recall from Section~\ref{subsec: notation and
%     conventions} that $R^{\crys,\underline{\lambda},\cO'}_{\rhobar}$ is
%   characterised as the unique quotient of~$R_{\rhobar}^{\square,\cO'}$
%   which is~$\cO'$-flat and reduced and has the property that for all
%   finite extensions~$E''/E'$ with ring of integers~$\cO''$, a morphism of
%   $\cO'$-algebras $R_{\rhobar}^{\square,\cO}\to\cO''$ factors through
%   $R^{\crys,\underline{\lambda},\cO'}_{\rhobar}$ if and only if the
%   corresponding Galois representation is crystalline of Hodge
%   type~$\underline{\lambda}$.

%   It is immediate that $R^{\crys,\underline{\lambda},\cO'}_{\rhobar}$
%   is a quotient of
%   $R^{\underline{\lambda},\semis,\cO'}_{\rhobar}$. Let~$A$ be any
%   Artinian quotient of
%   $R^{\crys,\underline{\lambda},\cO'}_{\rhobar}/\varpi^a$; we must
%   show that the morphism $\Spec A\to\cX^a_d$ lifts to~$\cC^a_{\piflat,s,d,h}$. Since~$A$ is also a quotient of
%   $R^{\underline{\lambda},\semis,\cO'}_{\rhobar}$, there is a
%   surjection of~$\cO'$-algebras $B\to A$ satisfying the property
%   above. Let~$M_{B/\varpi^a}$ be the \'etale $(\varphi,\Gamma)$-module
%   corresponding to the representation~$\rho_B\otimes_BB/\varpi^a$;
%   then by Proposition~% \ref{prop: Caruso Liu canonical action
%     % semistable}
%   , the morphism $\Spec B/\varpi^a\to\cX_{K,d}$ lifts
%   to~$\cC^a_{\piflat,s,d,h}$. Since $\Spec A\to\cX_{K,d}$ factors
%   through the closed immersion~$\Spec A\to\Spec B/\varpi^a$,
%   it follows that it lifts to~$\cC^a_{\piflat,s,d,h}$, as required.
%   %since~$A$ was arbitrary, we are done.
% \end{proof}


%   In particular, if we consider all such
%   $x:R^{\crys,\underline{\lambda},\cO'}_{\rhobar}\to\cO''$ (where in a
%   slight abuse of notation, $\cO''$
%   depends on~$x$), we see that there is an injection
%   $R^{\crys,\underline{\lambda},\cO'}_{\rhobar}\to\prod_x\cO''$, and
%   thus (since $R^{\crys,\underline{\lambda},\cO'}_{\rhobar}$ is
%   $\cO'$-flat) an injection
%   $R^{\crys,\underline{\lambda},\cO'}_{\rhobar}/\varpi^a\to\prod_x\cO''/\varpi^a$.
  
%   We are therefore reduced\mar{TG: ME says this needs some justification.} to showing that for each
%   such~$R^{\crys,\underline{\lambda},\cO'}_{\rhobar}\to\cO''$, the
%   corresponding morphism $\Spec \cO''/\varpi^a\to\cX_d^a$ factors
%   through~$\cZ^a_{d,K_s,h}$. To see this, it is enough to show that we
%   can lift this morphism to a
%   morphism~$\Spec\cO''/\varpi^a\to\cC^a_{d,K_s,h}$.

% To this end, note firstly that the composite
% morphism~$\Spf\cO''\to\cX_d\to\cR_{\BK,d}$ lifts uniquely
% to~$\cC_{\BK,d,h}$ as a consequence of the main theorems
% of~\cite{KisinCrys} (that is, the corresponding \'etale
% $\varphi$-module over~$\cO_{\cE,\cO''}$ contains a unique Breuil--Kisin module~$\gM$ over
% $\gS_{\cO''}$ of height at most~$h$). We need to show that the
% action of~$G_{K_s}$ on
% $W(\C^\flat)_{\cO''/\varpi^a}\otimes_{\gS_{\cO''}}\gM[1/u]$ which comes from
% its incarnation as  the $(G_K,\varphi)$-module corresponding to the
% morphism $\Spec \cO''/\varpi^a\to\cX_d$ satisfies the
% property of Lemma~\ref{lem: Caruso Liu Galois action on Kisin}; that
% is, we must show that this action preserves
% $\AAinf{\cO''/\varpi^a}\otimes_{\gS_A}\gM$, and that for
% all~$g\in G_{K_s}$, this action satisfies
% $(g-1)(\gM)\subseteq u^N\AAinf{\cO''/\varpi^a}\otimes_{\gS_A}\gM$.\mar{TG:
% should now be able to complete by citing \cite{MR2558890} and
% \cite{MR2745530}, the former for the structure of phi Ghat and the
% latter for this congruence.}
% \end{proof}

% \begin{thm}
%   \label{thm: potentially crystalline and semistable stacks are padic
%     formal algebraic flat of expected dimension}
% \end{thm}
% \mar{ME: By twisting into the effective
% 	situation, we can obtain the same result for any 
% 	crystalline deformation ring, so probably we should just state the 
% 	result in that level of generality?}
% \begin{proof}Twisting by~$\varepsilon^{p^n(p-1)}$ for some 
%   %       \mar{ME: Somewhere we should note that twisting by a character
%   %       	induces an automorphism of $\cX_d$! TG: did this back
%   %               in Section~\ref{subsec: tensor product
%   % of phi gamma and duality}, commenting out.}
%   sufficiently large~$n$ (chosen so that in particular
%   $\varepsilon^{p^n(p-1)}$ is trivial mod~$\varpi^a$), we may assume
%   that 
% for each~$\sigma$, $i$ we have $\lambda_{\sigma,i}\le 0$. Choose~$h\ge
% 0$ so that~$\underline{\lambda}$ is bounded by~$h$. Fix~$s$ sufficiently large in
%   the sense of Lemma~\ref{lem: Caruso Liu Galois action on Kisin}
%   % Proposition~\ref{prop: Caruso Liu canonical action semistable},
%   and
%   Proposition~\ref{prop: morphism C to Xs}. 



% Let~$\Spf R^\cC$ be the scheme-theoretic image (in the sense
% of~\cite[Defn.\ 3.2.15]{EGstacktheoreticimages}) of the pulled-back
% proper morphism  \[(\cC^a_{\piflat,s,d,h})_{\Spf
%     R_{\rhobar}^{\square,\cO'}/\varpi^a}\to \Spf
%   R_{\rhobar}^{\square,\cO'}/\varpi^a.\]Since~$\cC^a_{\piflat,s,d,h}\to\cX_d^a$ is
% proper, it follows from~\cite[Lem.\ 3.2.16]{EGstacktheoreticimages}
% that the composite $\Spf R^\cC\into  \Spf
%   R_{\rhobar}^{\square,\cO'}/\varpi^a\to\cX_d^a$ factors
%   through~$\cZ^a_{\piflat,s,d,h}$ (the scheme-theoretic image of 
%   $\cC^a_{\piflat,s,d,h}$; see Definition~\ref{def:2-Cartesian}) and the morphism $\Spf
%   R^\cC\to\cZ^a_{\piflat,s,d,h}$ is versal. Since~$\cZ^a_{\piflat,s,d,h}$ is an
%   algebraic stack, its versal rings are effective. It follows that the
%   morphism $\Spf R^\cC\to\cX_d^a$ is induced by a  morphism $\Spec R^\cC\to\cX_d^a$.
  
% It is therefore enough to show that the morphism
% $\Spf R^{\crys,\underline{\lambda},\cO'}_{\rhobar}/\varpi^a\hookrightarrow
% \Spf R_{\rhobar}^{\square,\cO'}/\varpi^a$ factors through
% $\Spf R^\cC$. In turn, it suffices to show that if~$A$ is an % discrete
% % \mar{ME: All Artinian quotients are discrete.}
% Artinian quotient of
% $R^{\crys,\underline{\lambda},\cO'}_{\rhobar}/\varpi^a$, the morphism
% $\Spec A\to\Spf R_{\rhobar}^{\square,\cO'}/\varpi^a$ factors through~$\Spf R^\cC$. To see this, note
% that by Proposition~\ref{prop: crystalline deformation ring factors through the scheme
%     theoretic image}, the morphism
%   $\Spec A\to \cX_d^a$ factors through 
%   $\cC^a_{\piflat,s,d,h}\to\cX_d^a$; so the morphism $\Spec A\to\Spf
%   R_{\rhobar}^{\square,\cO'}/\varpi^a$ factors
%   through~$(\cC^a_{\piflat,s,d,h})_{R_{\rhobar}^{\square,\cO'}/\varpi^a}$,
%   and therefore through the scheme-theoretic image~$\Spf R^\cC$, as required. 
% \end{proof}


% \chapter{Temporary notes}\label{sec:temp}I think that we can detect
% Hodge--Tate weights as in Kisin. The key thing seems to be to define
% $\Fil^i\varphi^*\gM=\Phi_{\gM}^{-1}(E(u)^i\gM)$, and then look at
% jumps
% in \[\Fil^i\varphi^*\gM/(E(u)\varphi^*\gM\cap\Fil^i\varphi^*\gM).\] I
% think we can show that this last thing is always (maybe assuming~$A$
% Noetherian or $p$-power torsion) a finite projective
% $\cO_K\otimes_{W(k)}A$-module.

% \begin{question}
%   I don't know if this is obviously independent of the choice of
%   descent (I don't see why it should be abstractly) although
%   presumably it is in the de Rham case; so in general we'll probably
%   also want to impose the condition for all descents. Also, I don't
%   think it can be imposed just using $\gMt$ because that only knows
%   about the base change of this to $\cO_{\C}$ which on the face of it
%   isn't enough information? (Similar confusion exists for inertial
%   types, although there we are only base changing from $W(k)$ to
%   $W(\overline{k})$ which maybe doesn't lose the information?
%   Hopefully Tong can clear this all up in any case.)
% \end{question}

% \begin{question}
%   I don't think we want to impose the condition that this has the
%   required rank integrally, because after all we know that the
%   reductions of crystalline representations can have inertial weights
%   that differ from the Hodge--Tate weights. On the other hand as there
%   is an intersection in the definition I don't think it's compatible
%   with all base changes, and I'm guessing that instead
%   what we want to do is to use the ``flat part'' machinery to pass to
%   the closed substack where this holds for flat things. Is that right?
% \end{question}

% \begin{question}
%   For scheme-theoretic images for formal algebraic stacks, are we just
%   saying that they are Ind-algebraic and so there's an obvious
%   definition and that works?
% \end{question}

% \begin{question}
%   Do we ever want to figure out what the ``flat parts'' construction
%   is doing, or do we just say that the versal rings will be quotients
%   of the universal ones, and by flatness they're flat and thus
%   determined by $\Zpbar$-points (or rather probably by possibly
%   reduced finite $\Zp$-algebra points).
% \end{question}

% \begin{question}
%   How do we see that we have the right thing on $\Zpbar$-points? Is
%   this going to use properness/some kind of formal functions argument?
% \end{question}

% \section{Restriction of
%   \texorpdfstring{$(\varphi,\Gamma)$}{(phi,Gamma)}-modules}\label{section:
%   restriction of phi gamma}

\chapter{Families of extensions}\label{sec: families of extensions}
In this chapter we extend the theory of the Herr complex 
to the context of $(\varphi,\Gamma)$-modules with coefficients,
%using the Herr complex.
and use it to
develop a theory of families of extensions of
$(\varphi,\Gamma)$-modules. %, using the Herr complex.
We then %use this to
inductively construct families which cover the underlying reduced
stack $\cX_{K,d,\red}$, and employ obstruction theory to deduce
that~$\cX_{K,d}$ is a Noetherian formal algebraic stack.
% \mar{ME: Here we will discuss the Herr complex, how to convert it into
% 	a perfect complex over $\cX_d$, and then the theory of universal
% 	extensions, families, etc.}

\section{The Herr complex}\label{subsec: Herr complex}% \mar{TG: I
%   glanced at the foundations of perfect complexes for stacks and it's
%   safe to say that I haven't understood them at all, but I guess we
%   should make some effort
%   to. \cite[\href{http://stacks.math.columbia.edu/tag/09RG}{Tag
%     09RG}]{stacks-project} has something about limits of algebraic
%   spaces, the stack analogue of which might be helpful. It's a bit
%   unclear to me right now how we will end up with a perfect complex on
% $\cX_d$ in the most naive sense, but I guess if nothing else then a
% lot of people in Chicago probably understand everything, so I think
% it's not a good use of my time right now to think about this!}
In this chapter we consider the Herr complex; our approach is informed
by the papers~\cite{MR1693457,MR1839766,MR1645070,MR2416996,MR3117501,MR3230818}, and in
particular we follow~\cite{MR3117501,MR3230818} in considering it as an object
of the derived category. Our main technical result is to show that it
is a perfect complex, which we will do by showing that it satisfies
the following well-known criterion.% \mar{TG: maybe throughout this
  % section $\Zp$ should be $\cO$?}

\begin{lem}
  \label{lem: criterion for perfectness of complex} Let~$A$ be a
  Noetherian commutative ring, and let~$C^\bullet$ be an object
  of~$D(A)$. Then the following two conditions are equivalent:
  \begin{enumerate}
  \item There is a quasi-isomorphism $F^\bullet\to C^\bullet$
    where~$F^\bullet$ is a complex of flat $A$-modules, concentrated
    in a finite number of degrees;  and the cohomology groups of~$C^\bullet$ are all finite
    $A$-modules.
  \item $C^\bullet$ is perfect.
  \end{enumerate}
\end{lem}
\begin{proof}
  This follows
  from~\cite[\href{http://stacks.math.columbia.edu/tag/0658}{Tag
    0658}]{stacks-project},
  \cite[\href{http://stacks.math.columbia.edu/tag/0654}{Tag
    0654}]{stacks-project} and
  \cite[\href{http://stacks.math.columbia.edu/tag/066E}{Tag
    066E}]{stacks-project}.
\end{proof}
 % \mar{TG: I guess we need
% to recall the connection of $(\varphi,\Gamma)$-modules to Galois
% representations first.}
% \mar{TG: I'm still not quite clear what the formalism is that we're
%   working in here; I guess we literally write down a complex of
%   $A$-modules for any~$A$, and we show that it's well-behaved in the
%   Noetherian case. I don't know whether we want to invoke a theory of
%   derived categories of sheaves for stacks such as the one
%   in~\cite[\href{http://stacks.math.columbia.edu/tag/07B5}{Tag
%     07B5}]{stacks-project}, but I won't attempt to do that for now, at
% least.}
% We follow the notation and terminology of~\cite{stacks-project} for
% perfect complexes, as we now recall. If~$A$ is a commutative ring,
% then we write~$D(A)$ for the derived category of the abelian category
% of $A$-modules, and we say that an object~$K$ of~$D(A)$ is a
% \emph{perfect complex} if it is quasi-isomorphic to a bounded complex
% of finite projective $A$-modules.\mar{TG: write more precise
%   definition at the start of the paper, and say for stacks that we
%   just ask for all pullbacks.}



Let $A$ be a $p$-adically complete $\cO$-algebra, and let $M$ be an
$A$-module with commuting $A$-linear endomorphisms $\varphi,\Gamma$
(in our main applications, $M$ will be a
projective \'etale $(\varphi,\Gamma)$-module
with~$A$-coefficients, but it will be convenient in our arguments to be
able to consider more general possibilities, such as subquotients of
base changes of such modules). % \mar{TG: right now I'm going to use the
  % $\Zp$-version of the theory; there should be some discussion above
  % about how the two relate.}
Then the~\emph{Herr complex} of~$M$ is by \index{Herr complex}
definition the complex of~$A$-modules $\cC^\bullet(M)$ in degrees
$0,1,2$ given by% \mar{TG: I imagine we want to consider it as an object
  % of the bounded derived category, but leaving that aside for
  % now}
\[\xymatrix{M\ar[rr]^{(\varphi-1,\gamma-1)}&&
      M\oplus M\ar[rr]^{(\gamma-1)\oplus(1-\varphi)}&&M.}\]
% \[0\to M\stackrel{(\varphi-1,\gamma-1)}{\to} M\oplus
%   M\stackrel{(\gamma-1)\oplus(1-\phi)}{\to}M\to 0.\] 
% \mar{TG: didn't
%   think about the signs here.}
% Similarly, if~$M'$ is a $(\varphi,\Gamma)$-stable subquotient of~$M$,
% we define %  locally free
% % weak Wach module\mar{TG: not sure if ``weak'' is the correct thing
% %   here; actually would be better to make a definition that we could
% %   apply to subs or quotients of~$M$, perhaps by defining a larger
% %   category of~$(\varphi,\Gamma)$-modules?} then we define
% the Herr
% complex~$\cC^\bullet(M')$ in exactly the same way. By definition,
% the Herr complex is an object of~$D^b(A)$ and in particular of~$D^-(A)$.

 We have the
following useful interpretation of the cohomology groups of the Herr complex.
\begin{lem}
  \label{lem: Yoneda for H0 H1 of Herr complex}Let $M_1$, $M_2$ be
  projective \'etale $(\varphi,\Gamma)$-modules with
  $A$-coefficients. Then for $i=0,1$ there are natural
  isomorphisms \[H^i\bigl(\cC^\bullet(M_1^\vee\otimes M_2)\bigr)\cong\Ext^i_{\A_{K,A},\varphi,\Gamma}(M_1,M_2). \]
\end{lem}
\begin{proof}
  This is straightforward; in
  particular the case $i=0$ is immediate. When $i=1$, any extension
  $M$ of~$M_1$ by~$M_2$ splits on the level of the underlying projective
  $\A_{K,A}$-modules, and such a splitting is unique up to an element of
  $\Hom_{\A_{K,A}}(M_1,M_2)=M_1^\vee\otimes M_2$. Given such a
  splitting of~$M$, we obtain two elements~$f,g$ of $M_1^\vee\otimes M_2$ by
  writing \[\varphi_M=
    \begin{pmatrix}
      \varphi_{M_2}&\varphi_{M_2}\circ f\\0&\varphi_{ M_1}
    \end{pmatrix},\ \gamma_M=
    \begin{pmatrix}
      \gamma_{M_2}&\gamma_{M_2}\circ g\\0&\gamma_{M_1}
    \end{pmatrix}.
\]The condition that~$\varphi_M$, $\gamma_M$ commute shows that
$(f,g)$ determines a class in $H^1\bigl(\cC^\bullet(M_1^\vee\otimes M_2)\bigr)$
(this class is easily seen to be well-defined, by our above remark
about the ambiguity in the choice of splitting). We leave the
verification that this gives the claimed bijection to the sufficiently
enthusiastic reader.
\end{proof}

In order to show that the Herr complex is a perfect complex, we will
need to develop a little of the theory of the~$\psi$-operator
on~$(\varphi,\Gamma)$-modules. We will only need this in the case
that~$A$ is an~$\F$-algebra, and we mostly work in this context from
now on. Note that for any~$K$, if $A$ is an~$\F$-algebra, then
$\A_{K,A}^+$ is $(\varphi,\Gamma)$-stable, and~$(\A_{K,A}')^+$ is
$(\varphi,\Gammat)$-stable; indeed we have~$\varphi(T)=T^p$
and~$\gamma(T)-T\in T^2\A^+_{K,A}$ by Lemma~\ref{lem: gamma on T}, and
the stability under the action of~$\Gammat$ follows by an identical
argument to the proof of Lemma~\ref{lem: gamma on T} (that is, it
follows from the continuity of the action of~$\Gammat$).
% since restricting to this case
% significantly simplifies our arguments, we for the most part restrict
% our consideration to~$\psi$ to this setting. Recall firstly that
% if~$A=\Zp$, we have the operator $\psi:\A_{K,\Zp}\to\A_{K,\Zp}$
% of~\cite[\S 3.1.1]{MR1693457}, which is continuous, $\Zp$-linear, commutes
% with~$\Gamma$, and satisfies $\psi\circ\varphi=\id$. \mar{TG: remark
%   here that char $p$ implies $A^+$ $\Gamma$-stable}

We claim that for any~$K$, ~$1,\varepsilon,\dots,\varepsilon^{p-1}$ is
a basis for~$\bE'_{K}$ as a $\varphi(\bE'_{K})$-vector
space. % (Here~$\varepsilon$ and $\bE'_K=\A'_{K,\Fp}$ are as in
% Section~\ref{subsec: phi gamma over Zp}.)
To see this, note that since the
extension $\bE'_K/\varphi(\bE'_K)$ is inseparable of degree~$p$, while
$\bE'_K/\bE'_{\Qp}$ is a separable extension, it is enough to show the
result for~$\bE'_{\Qp}$. In this case we have
$\bE'_{\Qp}=\Fp((\varepsilon-1))$, so $\varphi(\bE'_{\Qp})=\Fp((\varepsilon^p-1))$
and the claim is clear. Then for any~$x\in\bE'_K$ we can write
$x=\sum_{i=0}^{p-1}\varepsilon^i\varphi(x_i)$, and we define
$\psi:\bE'_K\to\bE'_K $ by % it follows from the
% discussion in~\cite[I.5]{MR1645070} that\mar{TG: I think this is OK,
%   but it makes me a bit nervous in the ramified case, just because
%   people don't seem to argue this way\dots}
$\psi(x)=x_0$. By definition~$\psi$ is continuous, $\Fp$-linear, commutes
with~$\Gammat$, and satisfies $\psi\circ\varphi=\id$. Since~$\psi$
commutes with~$\Delta$, we have an induced map $\psi:\bE_K\to\bE_K$,
which again is continuous, $\Fp$-linear, commutes
with~$\Gamma$, and satisfies $\psi\circ\varphi=\id$.

% The induced
% operator % \mar{TG: this is going to be my notation for the
%   % coefficient ring with the full cyclotomic extension. Have to decide
%   % whether to be defining things here first and observing that they
%   % descend to the invariants. Quite probably we should actually move
%   % this discussion of~$\psi$ back to the section where we recall stuff
%   % from Fontaine? Also it would be more traditional to use $\mathbf{E}$
%   % rather than~$\A_{\Fp}$ to denote these.}
% $\psi:\bE'_K\to\bE'_K$ admits an explicit description as
% follows.Recall firstly that
% if~$A=\Zp$, we have the operator $\psi:\A_{K,\Zp}\to\A_{K,\Zp}$
% of~\cite[\S 3.1.1]{MR1693457}, which is continuous, $\Zp$-linear, commutes
% with~$\Gamma$, and satisfies $\psi\circ\varphi=\id$.

If~$A$ is an $\Fp$-algebra, then since~$\psi$ is continuous we can
extend scalars from~$\Fp$ to~$A$ and complete to obtain continuous
$A$-linear maps $\psi:\A'_{K,A}\to\A'_{K,A}$ and $\psi:\A_{K,A}\to\A_{K,A}$. Again, these are continuous
and commute with~$\Gammat$ and~$\Gamma$, and satisfy $\psi\circ\varphi=\id$. We
have \[\A'_{K,A}=\oplus_{i=0}^{p-1}\varepsilon^i\varphi(\A'_{K,A}),\]and
% by definition and the discussion of the previous paragraph,
if $x\in\A'_{K,A}$ and we
write $x=\sum_{i=0}^{p-1}\varepsilon^i\varphi(x_i)$, then
$\psi(x)=x_0$. 

% \begin{lem}\mar{TG: redo this as an inline explanation? Should be
%     including continuity in the case $A=\Zp$ so that the extension to
%     the completed tensor product actually makes sense.}
%   \label{lem: existence of psi on rings}Let $A$ be an $\F_p$-algebra. %  $p$-adically
%   % complete $\cO$-algebra.
%   Then there is a continuous $A$-linear
%   operator~$\psi:\A_{K,A}\to\A_{K,A}$ 
%   % which is continuous with respect to the canonical topology {\em (}that is,
%   % the topology induced by the $(p,T)$-adic topology on~$\gS_A${\em )}, 
%   which
%   commutes with~$\Gamma$, and which
%   satisfies~$\psi\circ\varphi=\id$. % Furthermore
%   % $\psi(\A^+_{K,A})\subset\A^+_{K,A}$.
%   % \mar{TG: need to prove this last statement}
% \end{lem}
% \begin{proof}
%   It suffices to prove this in the case $A=\Zp$, as we then
%   define~$\psi$ in the general case by extension of scalars. % (as there is then a
%   % unique $A$-linear continuous extension of
%   % $\id\otimes\psi:A\otimes_{\cO}\cO_{\cE,\cO}\to
%   % A\otimes_{\cO}\cO_{\cE,\cO}$ to an endomorphism of~$\A_{K,A}$).
%   In this case the result is proved in~\cite[\S 3.1]{MR1693457}), the
%   key point being that the trace map of an inseparable extension of
%   characteristic~$p$ fields of degree~$p$ vanishes.
%   % one shows (by reducing modulo~$\varpi$) that~$\cO_{\cE,\cO}$ is free
% %   over~$\varphi(\cO_{\cE,\cO})$, with a basis given by $(1+T)^i$,
% %   $0\le i\le p-1$, and then~$\psi$ is given by projecting onto the
% %   basis vector with~$i=0$ and applying~$\varphi^{-1}$ (so concretely we
% %   have $\psi(\sum_{i=0}^{p-1}\varphi(a_i)(1+T)^i)=a_0$).
% %
% % Since~$\varphi$ commutes with~$\Gamma$, so does~$\psi$. To see
% % that~$\psi(\gS)\subseteq \gS$, it is enough to show that if
% % $\sum_{i=0}^{p-1}\varphi(a_i)(1+T)^i\in \gS$, then each
% % $a_i\in\gS$. If not, by clearing denominators we can assume that each
% % $a_i\in\gS$, and at least one $a_i\notin p\gS$, but that
% % $\sum_{i=0}^{p-1}\varphi(a_i)(1+T)^i\in p\gS$. Reducing mod~$p$, we
% % have $\sum_{i=0}^{p-1}a_i^p(1+T)^i=0$, and equating coefficients
% % of~$T^i$ we easily see that each $a_i=0$, a contradiction.
% % \mar{TG: hopefully not nonsense; actually
%     % everyone works with the maximal unramified extension and then
%     % takes Galois invariants, so maybe I'm missing something?}
% \end{proof}

  
\begin{prop}
  \label{prop: existence of psi on phi Gamma modules}Let $A$ be an
  $\F$-algebra, 
  % $p$-adically complete $\cO$-algebra,
  and let $M$ be a projective
  \'etale $(\varphi,\Gamma)$-module with $A$-coefficients. Then there
  is a continuous and open $A$-linear surjection~$\psi:M\to M$ such that $\psi$
  commutes with~$\Gamma$, and we
  have \[\psi(\varphi(a)m)=a\psi(m), \] \[\psi(a\varphi(m))=\psi(a)m\]
  for any $a\in \A_{K,A}$, $m\in M$.
\end{prop}
\begin{proof}We define $\psi:M\to M$ to be the
  composite
  \[M\stackrel{\Phi_M^{-1}}{\to} \A_{K,A}\otimes_{\A_{K,A},\varphi}M
    \stackrel{\psi\otimes 1}{\to}\A_{K,A}\otimes_{\A_{K,A}}M=M.  \]
% where the endomorphism $\psi$ of $\A_{K,A}$ is constructed in
%   Lemma~\ref{lem: existence of psi on rings}.
  The relations
  $\psi(\varphi(a)m)=a\psi(m)$, $\psi(a\varphi(m))=\psi(a)m$ are
  immediate from the definitions.
That ~$\psi$ is continuous follows
  from the continuity of $\Phi_M^{-1}$ and  the continuity of~ $\psi$
  on~$\A_{K,A}$, 
  % of 
  % Lemma~\ref{lem: existence of psi on rings},
  while the
   surjectivity of~$\psi$ is immediate from the relation $\psi(\varphi(m))=m$.
  % Since the ideals $(p,T)^n$ are cofinal with the
  % ideal~$(p^n,T^n)$, 
  % \mar{TG: may well rewrite everything to be
    % $(p,T)$-adic in the end, in fact. TG: it isn't clear to me that
    % there's a good way to do this; commenting out.}
 That~$\psi$ is open follows from Lemma~\ref{lem: bounds on psi} below.
%   (which also yields another proof of continuity).
  % \mar{ME: Have to check openness, which is less obvious than the
  %         other claims, I guess. TG: should perhaps give more details
  %         once lemma below is written. TG: actually I think continuity
  %       is trickier than openness\dots}
 \end{proof}

 % \begin{lem}
 %   \label{lem: phi versus p including negative powers}Let~$A$ be an
 %   $\cO/\varpi^a$-algebra, and suppose that either~$n\ge 0$,
 %   or~$n\le 1-a$. Then we have
 %   $\varphi(T^n)\in T^{(n+1-a)p}\A_{K,A}^+$, and
 %   $T^{np}\in\varphi(T^{n+1-a)})\A_{K,A}^+$.
 % \end{lem}
 % \begin{proof}
 %   The case~$n\ge a-1$ is~\cite[Lem.\
 %   5.2.3]{EGstacktheoreticimages}. If~$n\le 0$, then the result
 %   follows from the corresponding statement with~$n$ replaced
 %   with~$a-1-n\ge a-1$; for example, we have $T^{p(a-1-n)}\in
 %   varphi(T^{-n})\A_{K,A}^+$, which is equivalent to $\varphi(T^n)\in
 %   T^{(n+1-a)p}\A_{K,A}^+$, as required.
 % \end{proof}

  \begin{lem}\label{lem: bounds on psi}Let~$A$ be an
    $\F$-algebra, let~$M$ be a projective \'etale
    $(\varphi,\Gamma)$-module with $A$-coefficients, and let~$\gM$ be
    a $\varphi$-stable lattice in~$M$. % \mar{TG: have to put in lemma
      % saying this exists TG: that exists, have to decide whether or
      % not to make~$A$ finite type here.}
    Then there is an integer~$h\ge 0$ such that % if either~$n\ge 0$
    % or~$n\le 1-a$, we
    for any integer~$n$, we
    have \[\psi(T^{h+np}\gM)\subseteq T^n\gM\subseteq \psi(T^{np}\gM). \]
    
  \end{lem}%\mar{TG: haven't 100 percent thought this through yet\dots}
  \begin{proof}
%Choose~$h$ such that \[T^h\gM\subseteq\Phi_{M}(\varphi^*\gM)\subseteq\gM.\]
Since
    $\psi\circ\varphi=\id$, for any~$n$ we
    have \[\psi(T^{np}\varphi(\gM))=T^n\gM.\]
   % By Lemma~\ref{lem: phi versus p including negative powers} t
    This implies that \[T^n\gM=\psi(T^{np}\varphi(\gM))\subseteq
      \psi(T^{np}\gM). \]
For the other direction,
choose~$h$ such that \[T^h\gM\subseteq\Phi_{M}(\varphi^*\gM)\subseteq\gM.\]
Then we have
 \[\psi(T^h\gM)\subseteq\psi(\Phi_{M}(\varphi^*\gM))=% \psi(\A^+_{K,A})\psi(\varphi(\gM))=\psi(\A^+_{K,A})\gM\subseteq
    \gM.\]
   % By another application of Lemma~\ref{lem: phi versus p including
   %   negative powers}, it
  It follows that \[\psi(T^{h+np}\gM)=T^n\psi(T^h\gM)\subseteq
    T^n\gM, \]as required.
  \end{proof}



  \begin{lem}%\mar{TG: probably this should move}
    \label{lem: existence of phi gamma lattice}Let~$A$ be a Noetherian
    $\F$-algebra, and let~$M$ be a projective \'etale
    $(\varphi,\Gamma)$-module with $A$-coefficients. Then~$M$ contains
    a  $(\varphi,\Gamma)$-stable lattice.
  \end{lem}
  \begin{proof}By~\cite[Lem.\ 5.2.15]{EGstacktheoreticimages}, $M$ contains
    a $\varphi$-stable lattice~$\gM$. Let~$\Gamma\gM$ be the
    $\A^+_{K,A}$-submodule of~$M$ generated by the elements~$\gamma m$
    with~$\gamma\in\Gamma$, $m\in\gM$. We claim that~$\Gamma \gM$ is a
    lattice; note that since~$\Gamma$ and~$\varphi$ commute,~$\Gamma \gM$ is
    $\varphi$-stable, and it is $\Gamma$-stable because~$\A^+_{K,A}$
    is $\Gamma$-stable. To see that it is a lattice, it is
    enough (since $\A^+_{K,A}$ is Noetherian) to show that it is contained
    in a lattice (as it certainly spans~$M$). In particular, it is
    enough to show that $\Gamma\gM\subseteq T^{-n}\gM$ for some $n\ge
    0$. This follows easily from the compactness of~$\Gamma$; for
    example, if $e_1,\dots,e_m$ are generators of the finitely
    generated $\A^+_{K,A}$-module $\gM$, then we have a continuous map
    $j:\Gamma\to M^m$, $\gamma\mapsto (\gamma e_1,\dots,\gamma e_d)$,
    and  $\Gamma=\cup_nj^{-1}(T^{-n}\gM^m) $ is an open cover
    of~$\Gamma$. Since this has a finite subcover we are done.
    % \mar{TG: proof to be supplied by ME, using compactness
    %   of~$\Gamma$. TG: hopefully going along the right lines; this
    %   last statement
    %   should be almost immediate from compactness, I would think, but
    %   didn't check it carefully yet. TG: filled in, commenting out.}
  \end{proof}

    \begin{lem}%\mar{TG: probably want a variant of some kind}
    \label{lem: gamma on T unit}Let~$A$ be an~$\F$-algebra. Then for any~$i\in \Z$ we have
    $\gamma(T^{i})\in T^i(\A_{K,A}^+)^\times$. % , and $\gamma(T^i)-T^i\in
    % T^{i+1}\A_{K,A}^+$. The analogous statements hold
    % for~$(\A_{K,A}')^+$. 
  \end{lem}
  \begin{proof}
    This follows from Lemma~\ref{lem: gamma on T}. Indeed, we can
    write~$\gamma(T)=T+\lambda$ with $\lambda\in T^2\A_{K,A}^+$, so for
    any~$i\ge 1$, we have $\gamma(T^i)=(T+\lambda)^i$, and
    $\gamma(T^i)-T^i\in T^{i+1}\A_{K,A}^+$, whence $\gamma(T^{i})\in
    T^i(\A_{K,A}^+)^\times$. Since~$\gamma(T^i)\gamma(T^{-i})=1$, the
    result then also holds for~$i\le 0$.
  \end{proof}
\begin{cor}
    \label{cor: existence of psi gamma lattice}Let~$A$ be a Noetherian
    $\F$-algebra, and let~$M$ be a projective \'etale
    $(\varphi,\Gamma)$-module with $A$-coefficients. Then~$M$ contains
    a lattice~$\gM'$ such that for all~$m\ge 0$, $T^{-m}\gM'$ is $(\psi,\Gamma)$-stable.
  \end{cor}
  \begin{proof}
    By Lemma~\ref{lem: existence of phi gamma lattice} there is a
    $(\varphi,\Gamma)$-stable lattice~$\gM$, and by Lemma~\ref{lem:
      bounds on psi}, the lattice $T^{-n}\gM$ is
    $\psi$-stable for all $n\gg 0$. It is also $\gamma$-stable for
    all~$n\ge 0$ by Lemma~\ref{lem: gamma on T unit}, so we may take $\gM'=T^{-n}\gM$ for any sufficiently
    large value of~$n$.
  \end{proof}
  % \begin{lem}\mar{TG: might not want this in the end}
  %   \label{lem: ker psi is a direct summand}Let~$A$ be an~$\F$-algebra, and let~$M$ be a projective \'etale
  %   $(\psi,\Gamma)$-module with $A$-coefficients. Then for any
  %   integer~$r\ge 1$, we have decompositions of
  %   $\varphi^r(\A_{K,A})$-modules\mar{TG: check correct ring to be a
  %     module over} $M=\ker(\psi^{r})\oplus\varphi^{r}(M)$,
  %   and $\ker(\psi^{r+1})=\ker(\psi^{r})\oplus\varphi^{r}(\ker\psi)$.

  %   % ~$\ker(\psi)$ is a
  %   % direct summand of~$M$ as an~$A$-module; more precisely, we have~$M=\ker(\psi)\oplus\varphi(M)$.
  % \end{lem}
  % \begin{proof}
  %   For any~$m\in M$ we can write
  %   $m=(m-\varphi^r(\psi^r(m)))+\varphi^r(\psi^r(m))$, and since
  %   $\psi^r(m-\varphi^r(\psi^r(m)))=0$, we
  %   have~$M=\ker(\psi^r)+\varphi^r(M)$. If~$m\in\ker(\psi^r)\cap\varphi^r(M)$, we
  %   can write~$m=\varphi^r(m')$; but then $0=\psi^r(m)=m'$, so~$m=0$,
  %   as required. Finally, if~$m\in\ker(\psi^{r+1})$
  %   then~$\psi^r(m)\in\ker(\psi)$, so that
  %   $\varphi^r(\psi^r(m))\in\varphi^r(\ker(\psi))$; conversely, we
  %   have $\varphi^r(\ker(\psi))\subseteq\ker(\psi^{r+1})$, as required.
  % \end{proof}

  If~$M$ is an \'etale $(\varphi,\Gamma)$-module with $A$-coefficients, then we
  write \[M':=\A'_{K,A}\otimes_{\A_{K,A}}M,\] an \'etale
  $(\varphi,\Gammat)$-module with $A$-coefficients.
  \begin{lem}%\mar{TG: might not want this in the end}
    \label{lem: ker psi is a direct summand}Let~$A$ be an~$\F$-algebra, and let~$M$ be a projective \'etale
    $(\varphi,\Gamma)$-module with $A$-coefficients. Then we have a
    decomposition of
    $\varphi(\A'_{K,A})$-modules $M'=\ker(\psi)\oplus\varphi(M')$,
    and we may write $\ker(\psi)=\oplus_{i=1}^{p-1}\varepsilon^i\varphi(M')$.
    % ~$\ker(\psi)$ is a
    % direct summand of~$M$ as an~$A$-module; more precisely, we have~$M=\ker(\psi)\oplus\varphi(M)$.
  \end{lem}% \mar{TG: need to fix this - the second part is over
    % $\A'$. Probably just introduce some notation $\widetilde{M}$ for
    % tensoring up to this, but probably don't need to worry about $\Gammat$ action?}
  \begin{proof}
    Since $\psi\circ \phi$ is the identity on~$M'$, the endomorphism $\phi\circ \psi$
    of $M'$ is idempotent, and so $M'$ decomposes as the direct sum of its
    kernel and image.  Since $\phi$ is injective, while $\psi$ is surjective,
    we may express this decomposition as $M' = \ker(\psi) \oplus \varphi(M')$,
    as claimed.
    %For any~$m\in M'$ we can write
    %$m=(m-\varphi(\psi(m)))+\varphi(\psi(m))$, and since
    %$\psi(m-\varphi(\psi(m)))=0$, we
    %have~$M'=\ker(\psi)+\varphi(M')$. If~$m\in\ker(\psi)\cap\varphi(M')$, we
    %can write~$m=\varphi(m')$; but then $0=\psi(m)=m'$, so~$m=0$,
    %as required.

    As~$M'$ is \'etale, we have $\A'_{K,A}\otimes_{\varphi(\A'_{K,A})}
\varphi(M') \iso M'$. Since
    $\A'_{K,A}=\oplus_{i=0}^{p-1}\varepsilon^i\varphi(\A'_{K,A})$, we
    find that
    $M'=\oplus_{i=0}^{p-1}\varepsilon^i\varphi(\A'_{K,A})\varphi(M')=\oplus_{i=0}^{p-1}\varepsilon^i\varphi(M')$. Now,
    for~$1\le i\le p-1$ we have
    $\psi(\varepsilon^i\varphi(M'))=\psi(\varepsilon^i)M'=0$, so
    $\oplus_{i=1}^{p-1}\varepsilon^i\varphi(M')\subseteq\ker(\psi)$. Since
    we have
    $\oplus_{i=1}^{p-1}\varepsilon^i\varphi(M')\oplus\varphi(M')=\ker(\psi)\oplus\varphi(M')$,
    it follows that this inclusion is an equality, as required.
  \end{proof}

  % \begin{lem}
  %   \label{lem: reducing powers of gamma}Let~$A$ be an $\F$-algebra, % $p$-adically
  %   % complete $\cO$-algebra,
  %   and let~$M$ be a projective \'etale
  %   $(\psi,\Gamma)$-module with $A$-coefficients. Suppose that for
  %   some~$s\ge 1$, $(1-\gamma^{p^s})$ acts invertibly
  %   on~$\ker(\psi^{s+1})$. Then $(1-\gamma^{p^{s-1}})$ acts invertibly
  %   on~$\ker(\psi^{s})$
  % \end{lem}
  % \begin{proof}
  %   We may
  %   write~$(1-\gamma^{p^s})=(1-\gamma^{p^{s-1}})(1+\gamma^{p^{s-1}}+\dots+\gamma^{p^{s-1}(p-1)})$,
  %   so certainly~$(1-\gamma^{p^{s-1}})$ acts invertibly
  %   on~$\ker(\psi^{s+1})$. By Lemma~\ref{lem: ker psi is a direct
  %     summand}, we may write
  %   $\ker(\psi^{s+1})=\ker(\psi^{s})\oplus\varphi^{s}(\ker\psi)$.
  %   Since~$\gamma$ commutes with both~$\varphi$ and~$\psi$, each
  %   summand is stable under the action of~$(1-\gamma^{p^s})$, and the
  %   result follows.
  % \end{proof}
  The proof of the following result follows the approach
  of~\cite{MR1645070}, although one difference in our situation is
  that because we are working in characteristic~$p$, we do not need to
  reduce to the case that~$K/\Qp$ is unramified. % \mar{TG: tbh I'm back
    % to being a bit confused as to why this reduction is really needed
    % in char 0! I guess it's always possible that it's actually only
    % being used for convenience?}
  % On the other hand,
  % we need a minor additional argument relating to the difference
  % between $(\varphi,\Gamma)$-modules for the cyclotomic
  % $\Zp$-extension and the full $\Z_p^\times$-extension.
  \begin{prop}% \mar{TG: we need to figure out what the proof should be
    %   here now that we aren't using explicit formulae.}\mar{TG: can we
    % use Lemma~\ref{cor: existence of psi gamma lattice} to reduce to
    % the Artinian case, by considering all Artinian quotients of~$A$? I think it should show that the situation is
    % compatible with base change, and in the Artinian case we can just
    % cite Herr.}\ref
  \label{prop: inverting 1-gamma}Let~$A$ be a Noetherian % \mar{TG:
  %   rewriting this; hopefully have reduced it to a statement about the
  % coefficient rings with $\Zp$-coefficients, which can hopefully be
  % extracted from cO.}
  $\F$-algebra. % \mar{TG: maybe here and
    % below should be $\cO/\varpi^a$?} 
  % \mar{TG: should say we're
    % following Colmez's notes}
  % $p$-adically
  % complete $\Zp$-algebra which is topologically of finite type
  % over~$\Zp$ \emph{(}and so in particular Noetherian\emph{)}.
  % \mar{TG: didn't think about how $p$-adically
    % complete interacts with the reduction to Artinian; on the other
    % hand I think we only use this with $\Z/p^a\Z$-algebras at the
    % moment, so we could just write it in that case if needed.}
  If~$M$ is a projective \'etale
  $(\varphi,\Gamma)$-module with $A$-coefficients, then $(1-\gamma)$
  is bijective on~$\ker(\psi)$.
\end{prop}%\mar{TG: not yet proved}
\begin{proof}% We firstly reduce to the case that~$A$ is an
  % $\Fp$-algebra. Indeed, \mar{TG: irritatingly my argument for this
  % seems to be broken!}
  % \mar{TG: firstly need to explain the straightforward reduction to
  %   case of $\tA$, which is just that $\psi$ and $\gamma$ both commute
  %   with~$\Delta$.}\mar{TG: planning to change $\tA$ to $\A'$, not
  %   least to avoid conflict with notation for perfection of
  %   coefficients}
We begin with some preliminaries on the action of~$\gamma$ on~$\varepsilon\in(\A'_{K,A})^+$. We have $\gamma(\varepsilon)=\varepsilon^{\chi(\gamma)}$, and we 
write $\chi(\gamma)=1+p^Nu$ where $N\ge 1$ and~$u\in\Z_p^\times$. Then
for each $n\ge 1$ we can write $\chi(\gamma^{p^{n-1}})=1+p^{n+N-1}u_n$
with $u_n\in\Z_p^\times$, and if~$0\le r\le n+N-1$ and~$i\in\Z$ then we have
\numequation\label{eqn: moving epsilon into phi}
\gamma^{p^{n-1}}(\varepsilon^i)=\varepsilon^i\varphi^r(\varepsilon^{ip^{n+N-1-r}u_n}). \end{equation}Note
that since~$u_n$ is a $p$-adic unit, $(\varepsilon^{u_n}-1)/(\varepsilon-1)$ is a
unit in~$(\A'_{K,A})^+$, and we
can
write \numequation\label{eqn: gamma-1 on epsilon minus 1 equation}(\gamma^{p^{n-1}}-1)(\varepsilon-1)=\varepsilon(\varepsilon^{u_n}-1)^{p^{n+N-1}}\in
  (\varepsilon-1)^{p^{n+N-1}}((\A'_{K,A})^+)^\times. \end{equation}It follows that for any~$r\in\Z$, we
have \numequation\label{eqn: gamma minus 1 on epsilon minus 1}(\gamma^{p^{n-1}}-1)((\varepsilon-1)^r)\in
  (\varepsilon-1)^{r+p^{n+N-1}-1}(\A'_{K,A})^+. \end{equation} Indeed
if~$r\ge 0$ then we may
write \[(\gamma^{p^{n-1}}-1)((\varepsilon-1)^r)=(\gamma^{p^{n-1}}-1)((\varepsilon-1))\cdot\sum_{j=0}^{r-1}(\gamma^{p^{n-1}}(\varepsilon-1)^{r-1-j})\cdot(\varepsilon-1)^j, \]so
\eqref{eqn: gamma minus 1 on epsilon minus 1} follows from~\eqref{eqn:
  moving epsilon into phi} and the fact that
$\gamma^{p^{n-1}}(\varepsilon-1)\in(\varepsilon-1)(\A'_{K,A})^+$,
which in turn follows from~\eqref{eqn: gamma-1 on epsilon minus 1
  equation}. The case~$r\le 0$ then follows from the result for~$-r$ by
writing \[(\gamma^{p^{n-1}}-1)((\varepsilon-1)^r)=-\frac{(\gamma^{p^{n-1}}-1)((\varepsilon-1)^{-r})}{(\gamma^{p^{n-1}}(\varepsilon-1)^{-r})(\varepsilon-1)^{-r}},\]and
noting that $(\gamma^{p^{n-1}}(\varepsilon-1))/(\varepsilon-1)=(\varepsilon^{1+p^{n+N-1}u_n}-1)/(\varepsilon-1)$ is a unit.


   Since both~$\psi$ and~$\gamma$
  commute with the action of~$\Delta$, in order to prove the
  proposition it suffices to show
  that~$(1-\gamma)$ is bijective on the kernel of $\psi:M'\to
  M'$.  By Lemma~\ref{lem: ker psi is a direct summand}, this kernel equals
$\ker(\psi)=\oplus_{i=1}^{p-1}\varepsilon^i\varphi(M')$. % \mar{TG: justify this!}
It follows from~\eqref{eqn: moving epsilon into phi} that $\gamma$
preserves~$\varepsilon^i\varphi(M')$ for any integer~$i$, so it is
enough to prove that if $(i,p)=1$  that~$(1-\gamma)$ acts
invertibly on $\varepsilon^i\varphi(M')$.


In fact, we will prove the %\mar{TG: rework this reduction slightly}
stronger statement that
for each~$n\ge 1$, each $0\le t\le N-1$, and each~$(i,p)=1$, the operator $(1-\gamma^{p^{n-1}})$ acts
invertibly on~$\varepsilon^i\varphi^{n+t}(M')$ (note that
$(1-\gamma^{p^{n-1}})$ acts on this space by~\eqref{eqn: moving epsilon into phi}).


Writing \numequation\label{eqn: yet another eqn in psi invertible gamma proof}\varepsilon^i\varphi^{n+t-1}(M')=\varepsilon^i\varphi^{n+t-1}(\oplus_{j=0}^{p-1}\varepsilon^j\varphi(M'))=\oplus_{j=0}^{p-1}\varepsilon^{i+p^{n+t-1}j}\varphi^{n+t}(M'),\end{equation}we
see that if~$t>0$ then the statement for some pair~$(n,t)$ (and
all~$i$) implies the statement for~$(n,t-1)$. Writing
$(1-\gamma^{p^{n-1}})=(1-\gamma^{p^{n-2}})(1+\gamma^{p^{n-2}}+\dots+\gamma^{p^{n-2}(p-1)})$,
and again using~\eqref{eqn: yet another eqn in psi invertible gamma proof},
we see also that the statement for~$(n,t)$ implies the statement
for~$(n-1,t)$. We can therefore assume that~$n$ is arbitrarily large
and that~$t=N-1$.  

Let~$\gM$ be a
$(\psi,\Gamma)$-stable lattice in~$M$ (which exists by
Corollary~\ref{cor: existence of psi gamma lattice}), let~$\gM'=(\A'_{K,A})^+\otimes_{\A_{K,A}}\gM$, and choose~$n$
large enough that % $(\gamma^{p^{n-1}}-1)((\A'_{K,A})^+)\subset
% T^2(\A'_{K,A})^+$
$p^{n+N-1}\ge 3$
and \numequation\label{eqn: squeezing by 2 on M}(\gamma^{p^{n-1}}-1)(\gM')\subseteq (\varepsilon-1)^2\gM'.\end{equation}It follows that \numequation\label{eqn:gamma -1 squeezing}(\gamma^{p^{n-1}}-1)((\varepsilon-1)^r\gM')\subseteq (\varepsilon-1)^{r+2}\gM'\end{equation}
for all $r\in \Z$, because if $m\in \gM'$
then \[(\gamma^{p^{n-1}}-1)((\varepsilon-1)^rm)=(\gamma^{p^{n-1}}-1)((\varepsilon-1)^r)\cdot\gamma^{p^{n-1}}(m)+(\varepsilon-1)^r(\gamma^{p^{n-1}}-1)(m),\]so
that ~\eqref{eqn:gamma -1 squeezing} follows from ~\eqref{eqn: gamma
  minus 1 on epsilon minus 1} and~\eqref{eqn: squeezing by 2 on M},
together with our assumption that~$p^{n+N-1}\ge 3$.


%\mar{TG: need to think about
   % what should replace this claim for~$\gamma$ with $K$ being
%   arbitrary, and how the rest of the argument goes}
% $\chi(\gamma^{p^{n-1}})=1+p^nu$ with $u\in\Z_p^\times$,
% then \[(\gamma^{p^{n-1}}-1)(T)/T=(\varepsilon^{1+p^nu}-1)/T\]is a unit
% in~$\A^+_{K,A}$.
% \mar{TG: here and below, once $(\varphi,\Gamma)$ stuff is
  % sorted out, should check that we aren't using any explicit formulae
  % that need to be modified.}
% \mar{TG: ugh! OK that's
% not good enough except mod $p$, presumably just have to go deeper
% depending on~$a$. Come back and fix this. Also maybe it's even
% supposed to be for all~$m$, not just all~$m\ge 0$?}

By~\eqref{eqn: moving epsilon into phi}, we
have \[(\gamma^{p^{n-1}}-1)(\varepsilon^i\varphi^{n+N-1}(x))=\varepsilon^{i}\varphi^{n+N-1}(\varepsilon^{iu_n}\gamma^{p^{n-1}}(x)-x), \]so
it is enough to check that the map $f:M'\to M'$ given
by \[f(x)=\varepsilon^{iu_n}\gamma^{p^{n-1}}(x)-x \]is
invertible.

Let $\alpha=\varepsilon^{iu_n}-1$, so that~$\alpha/(\varepsilon-1)$ is a unit
in~$(\A'_{K,A})^+$. Then for any $r\in\Z$ and~$x\in (\varepsilon-1)^r\gM'$,  % \mar{TG: check versus
  % the above}
it follows from~(\ref{eqn:gamma -1 squeezing})
that
\[\left(\frac{1}{\alpha}f-1\right)(x)=\frac{\varepsilon^{iu_n}}{\alpha}(\gamma^{p^{n-1}}-1)(x)\in
  (\varepsilon-1)^{r+1}\gM'. \] In particular the
sum
\[g(x):=\sum_{j=0}^\infty\left(1-\frac{1}{\alpha}f\right)^j(x) \]
converges. Since $f,g$ are additive by definition, we have \[\left(1-\frac{1}{\alpha}f\right)\circ g(x)=g\circ
\left(1-\frac{1}{\alpha}f\right)(x)=g(x)-x, \] so that the function $g:M\to M$ is a
left and right inverse to~$\frac{1}{\alpha}f$. Thus~$f$ is
invertible, as required. % since by definition we
% have \[\left(1-\frac{1}{\alpha}f\right)(g(x))=g(x)-x/\alpha \]
% and
% \[g(f(x))=\sum_{j=0}^\infty\left(1-\frac{1}{\alpha}f\right)^j\left(\frac{1}{\alpha}f(x)\right)=x, \]
% so that $f(g(x))=x=g(f(x))$, as required.
\end{proof}

% \mar{ME: The statement says {\em bijective} (which we definitely need below),
% 	but doesn't the argument just prove {\em injective}?}
% \begin{proof}
%   \mar{TG: needed. ME observed that we should be able to reduce to the
%   field case, by arguing as in \cite[Prop.\ 4.1.16]{CEGSKisinwithdd},
%   i.e.\ going to Noetherian, and then embedding in product of Artinian
% quotients of completions. TG: wrote something, hopefully not too far
% from being correct. TG: it seems unclear whether we can just reduce to
% the Artinian case, so instead we might want to read Colmez and see if
% we can generalise to a Tate module setting.}We need to show that if~$m\in M$ satisfies
% $(1-\gamma)(m)=\psi(m)=0$, then~$m=0$. The ring~$A$ injects into the product of
% its localisations at all maximal ideals, and since~$A$ is Noetherian,
% it follows from the Artin--Rees lemma that it injects into the product
% of its completions at all maximal ideals. This product is equal to the
% projective limit of the Artinian quotients~$A/I$ of~$A$, so we see
% that it is enough to check that~$m$ vanishes modulo each
% such~$I$. Since~$M$ is projective, we may check this by
% replacing~$A$ by~$A/I$, and we may therefore assume that~$A$ is
% Artinian. By an obvious d\'evissage, we may assume further that~$A$ is
% a local~$\Fp$-algebra, and thus of finite length over its residue
% field~$\kappa(A)$. We can therefore reduce to the case that~$A$ is a finite
% field, which is standard~\cite[Thm.\ 3.8]{MR1693457}.% \mar{TG: do we need to do some truncation to
% %   justify a further reduction to a finite field?! ME: Maybe just assume
% % that $A$ is actually top.\ of finite type over $\mathbb Z_p,$ so
% % that its Artinian truncations automatically have finite residue
% % field. TG: thanks, commenting out.}
% \end{proof}

Suppose that~$A$ is an $\F$-algebra. Using the~$\psi$ operator, we can give an alternative description of
the Herr complex, which will be important in establishing that it is a
perfect complex. Let $\cC_\psi^\bullet(M)$ be the complex in degrees
$0,1,2$ given by \[\xymatrix{M\ar[rr]^{(\psi -1,\gamma-1)}&&
      M\oplus M\ar[rr]^{(\gamma-1)\oplus(1 - \psi)}&&M. }\]

  We have a morphism of complexes
  $\cC^\bullet(M)\to\cC^\bullet_\psi(M)$ given by
  \numequation
  \label{eqn:Herr quasi-iso}
  \xymatrix{M\ar[d]^{1}\ar[rr]^{(\varphi-1,\gamma-1)}&&
      M\oplus M\ar[d]^{(-\psi,1)}\ar[rr]^{(\gamma-1)\oplus(1-\varphi)}&&M\ar[d]^{-\psi} \\ M\ar[rr]^{(\psi-1,\gamma-1)}&&
      M\oplus M\ar[rr]^{(\gamma-1)\oplus(1-\psi)}&&M} 
\end{equation}
(That this is a morphism of complexes follows from the facts that
  $\psi\circ\varphi=\id$, and that~$\psi$ commutes with~$\gamma$.)

\begin{prop}
  \label{prop: Herr complex in terms of psi}
  Let~$A$ be a Noetherian % $p$-adically
  % complete $\cO$
  $\F$-algebra, % which is topologically of finite type
  % over~$\Zp$,
  and let $M$ be a projective
  \'etale $(\varphi,\Gamma)$-module with $A$-coefficients. Then
  %there is a quasi-isomorphism
  the morphism of complexes
  $\cC^\bullet(M)\to\cC^\bullet_\psi(M)$
  defined by~{\em (\ref{eqn:Herr quasi-iso})}
  is a quasi-isomorphism.
\end{prop}
\begin{proof}
%  Since $\psi\circ\phi=\id$, and~$\psi$ commutes with~$\gamma$,
%  we have a morphism of complexes $\cC^\bullet(M)\to\cC^\bullet_\psi(M)$ given
%  by
%  \[\xymatrix{0\ar[r]&M\ar[d]^{\id}\ar[rr]^{(\varphi-1,\gamma-1)}&&
%      M\oplus M\ar[d]^{(-\psi,1)}\ar[rr]^{(\gamma-1)\oplus(1-\varphi)}&&M\ar[d]^{-\psi}\ar[r]&
%    0 \\ 0\ar[r]&M\ar[rr]^{(\psi-1,\gamma-1)}&&
%      M\oplus M\ar[rr]^{(\gamma-1)\oplus(1-\psi)}&&M\ar[r]&
%    0 } \]
The cokernel of~(\ref{eqn:Herr quasi-iso})  %this morphism
is the zero complex, while its %and the
kernel is the complex \[\xymatrix{0\ar[r]&
     \ker(\psi)\ar[r]^{(\gamma-1)}&\ker(\psi), }\]which is acyclic by Proposition~\ref{prop: inverting
  1-gamma}. 
% \mar{ME: Doesn't that Prop.\ just show that there is no kernel?  To show there
% 	is no cokernel, don't we need the stronger statement that $1 -\gamma$
% 	is a bijection on $\ker(\psi)$? (Which is surely true, and is known
% in the case of classical coefficients. Or am I confused? TG: thanks,
% fixed above.}
\end{proof}

In fact, we require a slightly stronger statement than the quasi-isomorphism
of the preceding lemma.  Namely, we need a statement that takes into
account topologies, which will make use of the following lemma.  % We first % prove and
% recall
% a result from topology. 
%some results from topology.
%some results about quotient topologies
%and direct summands of Tate modules (which may well be known,
%but for which we don't have a reference). As usual, by a morphism of
%topological spaces, we mean a continuous map.

\begin{lemma}\label{lem:open maps}
	If $f:X \to Y$ is a continuous open map of topological spaces, 
	and if $Y'\subseteq Y$ is the inclusion of a subspace
	into $Y$ {\em (}i.e.\ $Y'$ is a subset of $Y$
	endowed with the induced topology{\em )},
	then the base-changed map $X' := f^{-1}(Y') \to Y'$ is again open,
	when $X'$ is endowed with the topology induced by that of $X$.
\end{lemma}
\begin{proof}
	Let $U'$ be an open subset of $X'$; then we may write
	$U' = U \cap X'$, for some open subset $U$ of $X$.
	Thus $$f(U') = f(U \cap X') = f\bigl(U\cap f^{-1}(Y')\bigr)
	= f(U) \cap Y'.$$
	Since $f$ is open, we conclude that $f(U')$ is an open
	subset of $Y'$, as required.
\end{proof}

% Recall that a topology on an abelian group is called {\em linear}
% if $0$ admits a neighbourhood basis consisting of open subgroups.
% If $M$ is a linearly topologized abelian group, 
% and $N$ is a subquotient of $M$,
% then $N$ admits a canonical linear topology, obtained
% either by regarding it as a quotient of a subgroup of $M$
% (equipped with the quotient topology on the quotient arising from
% the induced topology on the subgroup) or as a subgroup 
% of a quotient of $M$ (equipped with the induced topology arising
% from the quotient topology on the given quotient).
% \mar{TG: left this
  % paragraph alone for now, but imagine we might not want it.}
%\mar{ME: Add prefaratory remark here, reminding reader what a linear topology
%	on an abelian group is (it's one for which $0$ has a n.h.\ basis
%	of open subgroups), and that a subquotient of a linear topological
%	group has a canonical topology (obtained by thinking of it 
%	as a sub of a quotient or quotient of a sub, in either order --- this
%	is more-or-less equivalent to the fact that the lattice of 
%	subgroups is a modular lattice, I think). ME: Done}

%  \begin{lemma}%\mar{TG: where does the proof use that the topology is linear?
% % 	ME: I guess it doesn't, really.  It's just that if the topology
% % is not linear, I think we have to be more careful specifying the topology
% % on the $H^i$; it the quotient topology obtained by regarding $H^i$ as a
% % quotient of $Z^i$.  Maybe I should rewrite things slightly to clarify
% % this? TG: thanks - I think we should probably be quite explicit about all the topological
% % things, since they seem to be new. I'm happy to try to add some stuff
% % next week. }
% 	\label{lem:open maps of complexes}
% 	Let $f^{\bullet}: C^{\bullet} \to D^{\bullet}$ be a morphism
% 	of bounded above
% 	complexes of linearly
% 	topologized abelian groups which is open {\em (}in the 
% 	sense of being an open morphism term-by-term{\em )} and a
% 	quasi-isomorphism {\em (}in the usual algebraic sense, i.e.\ 
% 	it induces an isomorphism on cohomology groups, thought
% 	of simply as abstract abelian groups, ignoring their topologies{\em )}.
% 	Then, if we equip the cohomology groups with their canonical
% 	subquotient topologies, 
% 	the induced morphisms
% 	$$H^{\bullet}(f^{\bullet}): H^{\bullet}(C^{\bullet})
% 	\to H^{\bullet}(D^{\bullet})$$
% 	are isomorphisms of topological abelian groups {\em (}i.e.\
% 	are homeomorphisms, as well as isomorphisms of groups{\em )}.
% \end{lemma}
% \begin{rem}%\mar{TG: similarly I left this remark}
%   We are assuming that the topologies are linear in order to have
%   canonical subquotient topologies on the cohomology groups. More
%   generally (as the proof shows), the lemma holds for any topological
%   abelian groups, with the topology on the cohomology groups being the
%   quotient topology induced from the cocycle groups (which are
%   themselves given the subspace topology).
% \end{rem}
% \begin{proof}[Proof of Lemma~\ref{lem:open maps of complexes}]
% 	Let $d$ be such that $C^{d'} = D^{d'} = 0$ if $d' > d$.
% 	Then we have a commutative diagram
% 	$$ \xymatrix{
% 		0 \ar[r] & Z^d(C^{\bullet}) \ar[r]\ar[d] & 
% 		 C^d \ar[r]\ar[d] & H^d(C^{\bullet}) \ar[r]\ar[d] & 0 \\
% 		0 \ar[r] & Z^d(D^{\bullet}) \ar[r] & 
% 		 D^d \ar[r] & H^d(D^{\bullet}) \ar[r] & 0 }
% 	 $$
% 	 By assumption the right-most vertical arrow is an isomorphism
% 	 of groups.
% 	 Since the middle
% 	 arrow % (which is the right-hand arrow of this square)
%          is open,
% 	 the right-most vertical arrow is also open, when each $H^d$
% 	 is equipped with its quotient topology.
% 	 %\mar{ME: Need to ensure that these quotient maps are open!}
% 	 (Here we use the fact that a quotient map of topological groups
% 	 is necessarily open.)
% 	 Thus this arrow is a homeomorphism, and so the lemma 
% 	 is proved for cohomology in the top-most degree.
         
% 	 Again using that the right-most vertical arrow is
% 	 an isomorphism, we see that
% 	 the left-hand square is Cartesian.   The openness
% 	 of the middle arrow, together with Lemma~\ref{lem:open maps},
% 	 shows that the left-most arrow is open.   Thus we 
% 	 see that the truncated morphism
% 	 $$\tau_{< d} f^{\bullet}: \tau_{<d} C^{\bullet} \to \tau_{<d}D^{\bullet}$$
% 	 is again open.  The lemma now follows by descending induction on $d$.
% \end{proof}

% \mar{ME: This is now in Appendix~\ref{app:topological groups}
% 	as Corollary~\ref{cor:open polish}, so assuming 
% 	things stay that way, it should ultimately be removed from here.}
% We also recall the following result (which we will apply in
% conjunction with Lemma~\ref{lem:polish}).
% \begin{lemma}
% 	\label{lem:open polish}
% 	A continuous surjective homomorphism
%        	of Polish topological groups is necessarily
% 	open.\mar{ME: I think we need to assume that the source is
% 		actually complete, and not just completely metrizable ---
% 		unless this somehow holds automatically? If so,
% 	maybe replace ``Polish'' by ``complete and second countable''?
% ME: These concerns/confusions are now dealt with in Appendix~\ref{app:topological groups}.}
% \end{lemma}
% \begin{proof}
%   % \mar{ME: Cite something. TG: still couldn't find a precise
%   %   statement; you should check that I haven't written anything
%   %   stupid, as I learned this as I wrote it.}
%   This is well known, and
%   originally due to Banach~\cite{Banach1931}. For a more modern
%   reference, we note that \cite[Ex.\ IX.5.26]{MR0205211} shows that if
%   $f:G\to G'$ is a continuous homomorphism, where $G$ is complete and
%   has a countable base, and~$G'$ is non-meagre, then $f$ is strict;
%   this implies the statement of the lemma, as we now explain. 

%   By~\cite[Prop.\ III.2.24]{MR0205210}, being strict is equivalent to
%   being open.
%   \mar{ME: Maybe we mean strict {\em and surjective} is 
% 	  equivalent to being open?}
%   A Polish space is separable by definition, and so has a
%   countable base by~\cite[Prop.\ IX.2.12]{MR0205211}. In addition, a
%   Polish space is completely metrisable, and thus Baire by the Baire
%   category theorem~\cite[Thm.\ XI.5.1]{MR0205211}, and in particular
%   non-meagre (by the definition of a Baire space), as required.
% \end{proof}

We now have the following strengthening of
Proposition~\ref{prop: Herr complex in terms of psi},
which takes into account the topologies on the complexes.

\begin{prop}
	\label{prop:topological quasi-iso}
  Let~$A$ be a countable Noetherian  $\F$-algebra
  %for which $A/\varpi$ is countable,
  %which is topologically of finite type over~$\Zp$,
  and let $M$ be a projective
  \'etale $(\varphi,\Gamma)$-module with $A$-coefficients. Then the morphism
  of complexes
  $\cC^\bullet(M)\to\cC^\bullet_\psi(M)$
  defined by~{\em (\ref{eqn:Herr quasi-iso})}
  induces topological isomorphisms on each of the associated cohomology modules.
\end{prop}
\begin{proof}
  Proposition~\ref{prop: Herr complex in terms of psi} shows
  that~(\ref{eqn:Herr quasi-iso}) is a quasi-isomorphism. By
  definition, the induced map on $H^0$ is
  the identity map, and is therefore a homeomorphism. Since~$\psi$
  is continuous and open, we see that the maps $\cC^i(M)\to\cC^i_\psi(M)$ are
  continuous and open for each~$i$; in particular, the
 isomorphism on~$H^2$ is induced from the continuous open map
 $-\psi:\cC^2(M)\to\cC^2_\psi(M)$, and is therefore a homeomorphism.
 (Here and below
 we use the standard fact that a quotient morphism of topological
 groups is necessarily open.)

 The case of~$H^1$ requires a little more work. Since the maps
 $\cC^0(M)\to\cC^0_\psi(M)$ and $H^1\bigl(\cC^{\bullet}(M)\bigr) \to 
 H^1\bigl(\cC_{\psi}^{\bullet}(M)\bigr)$ are isomorphisms,
 we see that the continuous morphism
 $\cC^1(M)\to\cC^1_\psi(M)$ induces a continuous isomorphism of the
 modules of cocycles $Z^1(M)\to Z^1_\psi(M)$, and it suffices to show
 that this map is open (since this will imply that the induced
 bijection on $H^1$ is both continuous and open, and thus is an isomorphism
 of topological groups).

 If we let $\tZ^1(M)$ denote the preimage of~$Z^1_\psi(M)$ in $\cC^1(M)$,
and if we let $K\subseteq \cC^1(M)=M\oplus M$ denote $\ker(\psi)\oplus 0$,
 then we have inclusions $Z^1(M), K \subseteq \tZ^1(M)$,
 and hence a continuous morphism
 \numequation
 \label{eqn:continuous bijection}
 Z^1(M) \oplus K \to \tZ^1(M),
 \end{equation}
 which is in fact a bijection.
 Each of $Z^1(M)$, $K$, and $\tZ^1(M)$ is closed
 subgroup of $\cC^1(M) = M\oplus M$.  Lemma~\ref{lem:polish}
 shows that this latter topological group is Polish,
 and thus so is any of its closed subgroups.  
 Corollary~\ref{cor:open polish} then shows that~(\ref{eqn:continuous bijection})
 is in fact a homeomorphism,
while Lemma~\ref{lem:open maps} shows that the morphism $\tZ^1(M) \to 
Z^1_\psi(M)$ is open.  Consequently, we find that the morphism
$Z^1(M) \to Z^1_{\psi}(M)$ is open, as required.
% and so, by Lemma~\ref{lem:open maps of complexes},
% it suffices to show that the morphism
% $\cC^{\bullet}(M)\to\cC^{\bullet}_{\psi}(M)$
% induces an open mapping on each term.
% From the definition of~(\ref{eqn:Herr quasi-iso}),
% this follows from
%   Proposition~\ref{prop: existence of psi on phi Gamma modules},
% which shows that $\psi:M \to M$ is open.
% \mar{TG: make sure this states
% both continuous and open.}
\end{proof}%\mar{TG: got this far.}

\begin{lem}
  \label{lem: discrete quotients of projective}Let~$A$ be a Noetherian
  $\F$-algebra,
  and let~$X$ be an
  $A$-module   subquotient of a finitely generated projective
  $A[[T]]$-module, endowed with its natural subquotient
  topology. If the topology on~$X$ is discrete, then $X$ is finitely
  generated as an $A$-module.
\end{lem}
\begin{proof}Write $X=Y/Z$, where~$Y$ is an $A$-submodule of a
  finitely generated projective $A[[T]]$-module~$\gM$. Since the
  topology on~$X$ is discrete, $Z$ is open in~$Y$, and thus
 %we can write $Z=U\cap Y$, where $U$ is an open neighbourhood of zero in $\gM$.
  $Z\supseteq U\cap Y$, for some open neighbourhood $U$ of zero in $\gM$.
  Such a neighbourhood $U$ contains
  $T^n\gM$ for some $n\gg 0$, and so we find that~$X$ is a subquotient of the
  finitely generated $A$-module $\gM/T^n\gM$, as required.
  % \mar{TG: probably not precisely correct, but hopefully the right idea?
  % ME: Feel free to leave this, since I think I know how to do it,
  % and it won't take me long once I get to it. ME: Just read this,
  % and it looks absolutely correct to me (and is exactly what I planned
  % to write!)}
\end{proof}

% \mar{ME: Rewrite the proof of the following theorem so that, via base-change,
% 	we deduce perfection for arbitrary rings. TG: now that we've
%         had to add an extra hypothesis here, I guess we won't do this?
%       TG: removing this comment}
\begin{thm}
  \label{thm: Herr complex is perfect}\leavevmode Let~$A$ be a Noetherian % \mar{TG:
    % probably don't need}
    $\cO/\varpi^a$-algebra such that $A/\varpi$ is
    countable, % \mar{TG: eventually just
    % $p$-adic complete?}
    and let~$M$ be a projective \'etale $(\varphi,\Gamma)$-module with
    $A$-coefficients. 
  \begin{enumerate}
  \item  The Herr complex~$\cC^\bullet(M)$ is a
    perfect complex of $A$-modules, concentrated in
    degrees~$[0,2]$.% \mar{TG: probably have some notation
    % for this, and also add in duality etc}

  
  \item   If either (i) $B$ is a finite $A$-algebra;
or (ii) $B$ is a finite type $A$-algebra and $A$ itself 
is of finite type over~$\cO/\varpi^a$;
    then there is a natural isomorphism in the derived
    category
    \[\cC^\bullet(M)\otimes^{\mathbb{L}}_AB\isoto\cC^\bullet(M\otimes_{\A_{K,A}}\A_{K,B}). \]
    In particular, there is a natural isomorphism
    \[H^2(\cC^\bullet(M))\otimes_AB\isoto
      H^2(\cC^\bullet(M\otimes_{\A_{K,A}}\A_{K,B})). \]
  \end{enumerate}

\end{thm}
\begin{proof}Since~$A$ is Noetherian, $\A_{K,A}$ is a flat~$A$-algebra
  (being a localisation of the power series ring~$\A^+_{K,A}$), so
  ~$\cC^\bullet(M)$ is a complex of flat $A$-modules.  By
  \cite[\href{https://stacks.math.columbia.edu/tag/07LU}{Tag
    07LU}]{stacks-project} (for part~(1), taking~$R=\cO/\varpi^a$), and \cite[Prop.\
  2.2.2]{pilloniHidacomplexes} (for part~(2), again taking~$R=\cO/\varpi^a$), it suffices to prove
  the result in the case that~$A$ is an $\F$-algebra, which we assume
  from now on. We begin with~(1). By Lemma~\ref{lem: criterion for
    perfectness of complex}, we need only check that the cohomology
  groups of~$\cC^\bullet(M)$ are finitely generated
  $A$-modules. % \mar{TG: should probably have a reference to
    % somewhere in one of the documents for this? I guess it's obvious
    % as it's a localisation of the power series ring, which is a
    % countable product of copies of~$A$.} 
In order to do this, we will make two truncation arguments. % The first
% will show that these cohomology groups are $T$-power torsion
% $\gS_A$-modules, while the second (which uses the quasi-isomorphic
% complex $\cC_\psi^\bullet(M)$) shows that they are finitely presented
% $\gS_A$-modules; combining these two facts gives the result.
  % By~\cite[\href{http://stacks.math.columbia.edu/tag/068T}{Tag
%     068T}]{stacks-project} we can work \emph{fppf} locally,
% \mar{ME: Doesn't this localization require~(2) ? TG: yes; but this is
%   a holdover from some old version before we understood projective
%   versus locally free, so I guess we can hopefully argue as in your
%   next comment.}
% and therefore we can in particular assume that~$M$ is a free
%   $\OEA$-module.
  % \mar{TG: this is so that we can choose a lattice;
    % perhaps it will be better handled by more Drinfeld stuff.}

Firstly, by Lemma~\ref{lem: existence of phi gamma lattice}, we can choose a
 $(\varphi,\Gamma)$-stable lattice $\gM\subseteq M$. % (that is, a free~$\gS_A$-submodule with
%   $\gM\otimes_{\gS_A}\OEA=M$; such a lattice exists by the
%   \mar{ME: Should elaborate on the existence of $\gM$, I think.
%   It's not obvious to me right now that we can get a free lattice
%   that's $\Gamma$-invariant.  On the other hand, I think we {\em can}
%   find a torsion-free Kisin submodule that's $\Gamma$-invariant,
%   and I think the argument will work just as well with this.  (And 
%   this would also obviate the need for the initial {\em fppf}
%   localization, I think.)}
%   continuity of the
%   $\Gamma$-action). For each $n\ge 0$, the sublattice
%   $T^n\gM$ is $\Gamma$-stable (because~$\gamma(T)$ is divisible
%   by~$T$), so that after replacing~$\gM$ with $T^n\gM$ for $n\gg 0$,
%   we can assume that~$\gM$ is also $\varphi$-stable. Furthermore, for
%    each $n\ge 0$,  $T^n\gM$
% is $(\varphi,\Gamma)$-stable.
  % \mar{TG: introduce this terminology elsewhere.}
%  $(\varphi,\Gamma)$-stable lattice % \mar{TG: do this; $\gM$ is in,
    % and we use~$\psi$ for $M/\gM$.}
%
%
% We begin by showing that $\cC^\bullet(\gM)$ is perfect.
% Note that for each $n\ge 0$, the sub-$\gS_A$-module $T^n\gM$
% is $(\varphi,\Gamma)$-stable
%   (because~$\varphi(T)$ and~$\gamma(T)$ are both divisible
%   by~$T$)
   % \mar{TG: something here is making me nervous, but I think
    % this is OK?}
We claim that for every~$n\ge 1$ the Herr complex
$\cC^\bullet(T^n\gM)$ is acyclic; consequently,  the natural morphism of
complexes $\cC^\bullet(M)\to\cC^\bullet(M/T^n\gM)$ is a
quasi-isomorphism.%  In particular, the cohomology groups of~$\cC^\bullet(M)$ are
% $T$-power torsion.
% quasi-isomorphism., and since $\gM/T^n\gM$ is a finite rank projective
% $A$-module, $\cC^\bullet(\gM)$ is  perfect.

To see the claim, it suffices to show that $(1-\varphi):T^n\gM\to T^n\gM$
  is an isomorphism; indeed, the
  exactness of~$\cC^\bullet(T^n\gM)$ is a formal consequence of this. We begin by checking injectivity. If $m\in T^n\gM$ and
$(1-\varphi)(m)=0$, then we have
$m=\varphi(m)\in T^{pn}\gM$ % by
% Lemma~\ref{lem: phi versus p including negative powers}, \mar{TG: here
% and below simplify because we will have $a=1$}
% Lemma~\mar{TG: ref to lemma in etale phi modules}
% \ref{lem: phi(u) divisible by u}
so that in particular % by our
% assumption on~$n$,
we have $m\in T^{n+1}\gM$. By induction, we see
that $m\in T^n\gM$ for all~$n$, so that~$m=0$, as required.

We now prove surjectivity. If $m\in T^n\gM$, then we have seen in
the previous paragraph that~$\varphi(m)\in T^{n+1}\gM$,
$\varphi^2(m)\in T^{n+2}\gM$, and so on. Since $\gM$ is $T$-adically
complete, we can set $x=\sum_{i\ge 0}\varphi^i(m)\in T^n\gM$; then
$(1-\varphi)(x)=m$, as required.  

 % We have an
 %  exact sequence of
 %  complexes \[0\to\cC^\bullet(\gM)\to\cC^\bullet(M)\to\cC^\bullet(M/\gM)\to
 %    0.\] so it is enough to prove that each of~$\cC^\bullet(\gM)$
 %  and~$\cC^\bullet(M/\gM)$ is perfect.

We now turn to our other truncation argument. Let $\gM'$ be as in
Lemma~\ref{cor: existence of psi gamma lattice}, so that for each $n'\le 0$, % \mar{TG: maybe give some
%   justification of this? TG: can now cite lemma above, but this whole
%   proof should be gone through in any case} 
% We now choose a $(\psi,\Gamma)$-stable lattice $\gM'\subseteq M$ (for
% example, we could take $\gM'=T^m\gM$ with $m\ll 0$). For each $n'\le 0$,
$T^{n'}\gM'$ is a $(\psi,\Gamma)$-stable lattice in~$M$. % \mar{TG: I think this is true? Certainly OK if $n'\ll 0$.}
In particular, for each~$n'\le 0$, $\cC_\psi^\bullet(T^{n'}\gM)$ is a
subcomplex of $\cC_\psi^\bullet(M)$. We claim that if
$n'\le -2$ then $\cC_\psi^\bullet(M/T^{n'}\gM')$ is acyclic,
so that the natural morphism of complexes
$\cC_\psi^\bullet(T^{n'}\gM')\to\cC_\psi^\bullet(M)$ is a
quasi-isomorphism. % It follows from Proposition~\ref{prop: Herr complex
%   in terms of psi} that the cohomology groups of~$\cC^\bullet(M)$ are
% finitely generated $\gS_A$-modules. We have already seen that they are
% $T$-power torsion, so we see that for all $N\gg 0$, the cohomology
% groups of~$\cC^\bullet(M)$ are finitely generated $\gS_A/T^M$-modules,
% and thus finitely generated $A$-modules, as required.

% We now need to show that~$\cC^\bullet(M/\gM)$ is perfect. By
% Proposition~\ref{prop: Herr complex in terms of psi}, it suffices to
% show that~$\cC^\bullet_\psi(M/\gM)$ is perfect.\mar{TG: actually this
%   is all going a bit wrong. It's presumably roughly the correct
%   approach, but I'm getting confused about what we have to do to get
%   $\psi$-stability etc.} We can do this by an
% argument dual to the case of~$\cC(\gM)$. Note firstly that for each
% $n\ge 0$, $T^{-n}\gM\subseteq M$ is $(\psi,\Gamma)$-stable.\mar{TG: is
%   it obvious that $T\psi(T)$ is integral? I presume it's true\dots} It suffices to show
% that for $n\ge pa/(p-1)$, \[(1-\psi):M/T^{-n}\gM\to M/T^{-n}\gM\] is an
% isomorphism. 

% It remains to prove the claim. 
As above, it is enough to show that~$(1-\psi)$ is bijective on
$M/T^{n'}\gM'$. Note firstly that if $r\le n'$
then $\psi(T^r\gM')\subseteq T^{r+1}\gM'$. Indeed, %by Lemma~\ref{lem: phi versus p including negative powers},  % Lemma~\mar{TG: ref
% to lemma in etale phi modules}
% \ref{lem:
  % phi(u) divisible by u}
for each integer $r\le 0$ we
have \[\psi(T^r\gM')\subseteq\psi(T^{p\lfloor r/p\rfloor}\gM')% subseteq
  % \psi(\varphi(T^{\lfloor r/p\rfloor+1-a})\gM')
  =T^{\lfloor r/p\rfloor}\psi(\gM')\subseteq
  T^{\lfloor r/p\rfloor}\gM',\] and the claim follows since
if~$r\le -2$ then 
$\lfloor r/p\rfloor\ge r+1$.

We now show that~$(1-\psi)$ is injective on $M/T^{-n}\gM'$. Suppose that $m\in M$ with
$(1-\psi)(m)\in T^{n'}\gM'$. Then $m\in T^r\gM'$ for some $r\ll 0$. If
$r\ge n'$ then we are done, and if not then since $\psi(m)\in
T^{r+1}\gM'$ and $(1-\psi)(m)\in T^{n'}\gM\subseteq T^{r+1}\gM$, we
have $m\in T^{r+1}\gM$. It follows by induction that $m\in
T^{n'}\gM'$, as required.

For surjectivity, take $m\in M$, and choose $r\ll 0$ such that $m\in
T^r\gM'$. If $r< n'$ then $\psi(m)\in T^{r+1}\gM'$, so by induction we
see that there is some $s\ge 0$ such that $\psi^s(m)\in
T^{n'}\gM'$. If we set $x=\sum_{i=0}^{s-1}\psi^i(m)$, then
$(1-\psi)(x)=m-\psi^{s}(m)\in m+ T^{n'}\gM$, as required.

 We now consider the quasi-isomorphisms
%\mar{TG: seem to have relabeled
%   these lattices without saying so?!}
\[\cC^\bullet_\psi(T^{n'}\gM')\to\cC^\bullet_\psi(M)\leftarrow
  \cC^\bullet(M)\to\cC^\bullet(M/T^n\gM).  \] 

Proposition~\ref{prop:topological quasi-iso}
shows that the middle quasi-isomorphism
induces a topological isomorphism on cohomology modules,
and thus altogether we obtain morphisms
$$H^i\bigl( \cC^{\bullet}_{\psi}(T^{n'}\gM')\bigr) \to H^i\bigl(
\cC^{\bullet}(M/T^n\gM)\bigr)$$
which are isomorphisms of $A$-modules, and continuous
with respect to the natural topologies on the source
and target.
Since the target is endowed with the discrete topology,
we find that the source is also endowed with the discrete
topology. By Lemma~\ref{lem: discrete quotients of projective}, it
follows that %we  % On the other hand, the source has
% the property that any discrete quotient is finitely generated
% over $A$. \mar{ME: Explain/prove this, perhaps in a preparatory lemma.}
%find that 
these cohomology modules are indeed finitely generated
over $A$, and hence so are the cohomology modules of
$\cC^{\bullet}(M)$, as required.
%\textbf{Notes on where we are} We have (ignoring the powers of~$T$)
%quasi-isomorphisms
%\[\cC^\bullet_\psi(\gM')\to\cC^\bullet_\psi(M)\leftarrow
%  \cC^\bullet(M)\to\cC^\bullet(M/\gM).  \] If the middle arrow went
%the other way, I think we'd be OK, because the images of the maps of
%complexes would be termwise finitely generated $A$-modules. I had
%hoped to use that we could lift maps from complexes of projectives to
%get some maps in other directions, but in fact only the final term is
%of that form, so that's not much help!
%
%An idea: the map $\cC^\bullet_\psi(M)\leftarrow
%  \cC^\bullet(M)$ is actually quite close to being an isomorphism of
%  complexes; it's termwise surjective, and the kernel only has
%  $\ker(\psi)$, which I think is finitely generated over $A[[u]]$. So
%  maybe we still basically have something of the form that I
%  want??\mar{TG: needs fixing}
% \mar{TG: I would like to prove base change following the proof
%   of~\cite[Thm.\ 4.4.3]{MR3230818}. Unfortunately, I don't understand
%   the last quasi-isomorphism in the second paragraph of the proof of
%   that result. TG: seems to be easy, in fact, because under the
%   hypotheses that we are now working with, the residue fields of~$B$
%   will be finite over those of~$A$. So maybe this is a good way to
%   argue?}

We now turn to~(2), where we follow the proof of~\cite[Thm.\
4.4.3]{MR3230818}. Since both~(i) and~(ii) in particular require~$B$
to be a finite type $A$-algebra, we see
that~$B$ is Noetherian, and~$B$ is countable, so that in particular
$\cC^\bullet(M\otimes_{\A_{K,A}}\A_{K,B})$ is a perfect complex of
$B$-modules by part~(1) (replacing $A$ by~$B$, and~$M$ by $M\otimes_{\A_{K,A}}\A_{K,B}$). Since~$\cC^\bullet(M)$ is a perfect complex
of $A$-modules, $\cC^\bullet(M)\otimes^{\mathbb{L}}_AB$ is also a perfect
complex of $B$-modules, and the natural map $\A_{K,A}\otimes_AB\to\A_{K,B}$
induces a morphism \numequation\label{eqn: base change for Herr complex}\cC^\bullet(M)\otimes^{\mathbb{L}}_A B\to
  \cC^\bullet(M\otimes_{\A_{K,A}}\A_{K,B}).\end{equation}

In case~(i),
when~$B$ is in fact finite as an $A$-module, the natural map
$\A_{K,A}\otimes_AB\to\A_{K,B}$ is an isomorphism (indeed, we have
$A[[T]]\otimes_AB=B[[T]]$, because $A[[T]]\otimes_AB$ is finitely
generated over~$A[[T]]$ and thus $T$-adically complete and separated
by the Artin--Rees lemma), so~(\ref{eqn: base change for Herr
  complex}) is certainly an isomorphism in this case.

We now turn to  %the general
case~(ii),
when~$B$ is only assumed to be a finite type $A$-algebra,
but~$A$, and hence also $B$, is furthermore of finite type over~$\cO/\varpi^a$.
This implies in particular that $A$ is Jacobson,
so that if~$\m$ is any
maximal ideal of~$B$, then~$B/\m$ is finite as an $A$-module.
%\mar{ME: Here we seem to be implicitly asserting that a f.g.\ $A$-algebra
%which is also a field is finite over~$A$, which requires  $A$  to be
%Jacobson.  So I think for part~(2) we should be assuming that $A$
%is f.t.\  over some $\cO/\varpi^a$, to get this Jacobsonness.  Of course,
%we then have to check all our citations of part (2) \dots . Actually,
%maybe we can have two version, a more relaxed one for $B/A$ finite,
%and then the version where $B$ and $A$  are both f.t.\ over $\cO/\varpi^a$.}
From this, and the already-proved case~(i),
we conclude that (\ref{eqn: base change for Herr
  complex}) is an isomorphism if we replace~$B$ by~$B/\m$. 

It follows that we have a chain of
quasi-isomorphisms \[(\cC^\bullet(M)\otimes^{\mathbb{L}}_A
  B)\otimes^{\mathbb{L}}_B B/\m \isoto \cC^\bullet(M)\otimes^{\mathbb{L}}_A B/\m \isoto \cC^\bullet(M\otimes_A B/\m)\]\[\isoto
  \cC^\bullet(M\otimes_{\A_{K,A}}\A_{K,B})\otimes^{\mathbb{L}}_B B/\m, \]
where the final quasi-isomorphism comes from replacing~$A$ by~$B$, $B$
by~$B/\m$, and~$M$ by $M\otimes_{\A_{K,A}}\A_{K,B}$ in~(\ref{eqn: base change
  for Herr complex}). (That this is indeed a quasi-isomorphism follows
by case~(i), %by the discussion above,
since~$B/\m$ is a finite $B$-module, while we have
$(M\otimes_{\A_{K,A}}\A_{K,B})\otimes_{\A_{K,B}}\A_{K,B/\m}
=M\otimes_{\A_{K,A}}\A_{K,B/\m}=M\otimes_{\A_{K,A}}(\A_{K,A}\otimes_AB/\m)=M\otimes_AB/\m$).  
%$M\otimes_{\A_{K,A}}\A_{K,B}\otimes_B
%B/\m=M\otimes_{\A_{K,A}}\A_{K,B/\m}=M\otimes_{\A_{K,A}}(\A_{K,A}\otimes_AB/\m)=M\otimes_AB/\m$).  

Thus~ (\ref{eqn: base change for Herr complex}) becomes a
quasi-isomorphism after applying~$\otimes^{\mathbb{L}}_B B/\m$, for
any maximal ideal~$\m$ of~$B$. It follows from~\cite[Lem.\
4.1.5]{MR3230818} that (\ref{eqn: base change for Herr complex}) is a
quasi-isomorphism, as required. Finally, the compatibility of the formation of~$H^2$ with base change
is an immediate consequence of this isomorphism in the derived
category, together with the vanishing of all of the higher degree
cohomology groups.
\end{proof}

% Since the modules in the original complex
% $\cC^\bullet(M)$ are flat $A$-modules, the derived tensor product
% $\cC^\bullet(M)\otimes^{\mathbb{L}}_AB$ can be computed by the ordinary
% tensor product $\cC^\bullet(M)\otimes_AB=\cC^\bullet(M\otimes_A B)$, and we have a natural continuous
% morphism of %\mar{TG: starting a new argument} 
% \mar{ME: Explain topology on $M\otimes_A B$, once we work out what it
%   is! TG: is it just induced by $T$-adic norm? TG: wrote this, not
%   sure if we are actually going to use continuity in the end though?}
% complexes % \numequation\label{eqn: completion of tensor product A to
%   % B}
% \[\cC^\bullet(M\otimes_A B)\to\cC^\bullet(M\otimes_{\OEA}\OEB)\]%end{equation}% \mar{TG:
% % strictly speaking $\cC^\bullet(M\otimes_AB)$ isn't defined, but
% % suspect we might want to mention it?}
% induced by the ring homomorphism $\OEA\otimes_AB\to\OEB$
% (here~$M\otimes_A B$ has the topology for which the image
% of~$\gM\otimes_AB$ is open and has the $T$-adic topology). We need to
% show that this morphism induces an isomorphism on cohomology
% groups. 

% % By our assumption on~$B$ and the above, the cohomology groups of the
% % target of this morphism have the discrete topology, so it is enough to
% % check that the map has dense image.\mar{TG: why is this? Should also
% %   probably explain somewhere that the maps on cohomology groups are
% %   indeed continuous\dots}



% % Since $M$ is a projective $\OEA$-module and $\OEA\otimes_AB\to\OEB$ has dense image, the morphism of
% % complexes~(\ref{eqn: completion of tensor product A to B}) has
% % termwise dense image.
% Write  $M_B:= M\otimes_{\OEA}\OEB$, and let~$\gM_B$
% be the image of $\gM\otimes_{\gS_A}\gS_B$ in $M_B$. % , and let $\gM^B$ be
% % the image of $\gM\otimes_A B$ in $M\otimes_AB$. 
% % , so that $\gM_B$ is the
% % closure of $\gM^B$ in~$M_B$.
% Since~$B$ is a finite type $A$-algebra, we see that~$B$ is Noetherian,
% and~$B/p$ is countable. By the proof of part~(1), applied to~$M_B$ in place of~$M$, we have a
% natural quasi-isomorphism  
% % \mar{ME: How is this proved?  Note that $\gM^B$ is not $T$-adically complete,
% % 	so the proof that $1-\varphi$ is surjective doesn't apply
% % 	to $\cC^{\bullet}(T^n\gM^B)$. TG: need to think about this -
% %         it's quite possible that there's a mistake here...TG: maybe
% %         assume that $B$ is finite type over~$A$, and then
% %         $\gS_A\otimes_AB$ is finite type over~$\gS_A$, so Noetherian,
% %         so $T$-adically separated works OK, and we should be able to
% %         deduce that $1-\varphi$ is an isomorphism, because it is for $T^n\gM$.}
%  \[\cC^\bullet(M_B)\isoto
% \cC^\bullet(M_B/T^n\gM_B)=\cC^\bullet((M/T^n\gM)\otimes_{\gS_A}{\gS_B}).\]%Now, since $M_B/T^n\gM_B$ has the discrete
% % topology, the induced morphism $(M/T^n\gM)\otimes_AB=(M\otimes_AB)/T^n\gM^B\to M_B/T^n\gM_B=(M/T^n\gM)\otimes$ is
% % surjective (as it has dense image),\mar{TG: I think this might actually
% %   an isomorphism, but I get confused - if it isn't, just need to
% %   replace $T^n\gM^B$ by the kernel of  $M\otimes_AB\to M_B/T^n\gM_B$,
% %   and either make a topological argument or maybe something like
% %Artin--Rees??}


% We saw in the proof of part~(1) that $(1-\varphi):T^n\gM\to T^n\gM$ is
% an isomorphism. It follows that $(1-\varphi):T^n\gM\otimes_AB\to T^n\gM\otimes_AB$ is
% an isomorphism, and from the right exact sequence \[T^n\gM\otimes_AB\to
%   M\otimes_AB\to(M/T^n\gM)\otimes_AB\to 0 \]\mar{TG: I now get
%   completely stuck! TG: ah, we think that in fact $(1-\varphi)$ is
%   an isomorphism on~$T^n\gM^B$, where~$\gM^B$ is the image of $\gM\otimes_AB\to
%   M\otimes_AB$. The point is that surjectivity follows from
%   surjectivity on~$\gM\otimes_AB$ (which follows from being an
%   isomorphism on $\gM\otimes_AB$), and then the argument for
%   injectivity in part~(1) works, provided that we know that $\gM^B$ is
% $T$-adically separated. For this it is probably enough by Artin--Rees
% to check that $A[[T]]\otimes_AB$ is $T$-adically separated, and we're
% hoping that can't be too hard, e.g.\ the case where $B$ is a
% polynomial ring looks easy as we should just inject into~$B[[T]]$, and
% maybe in general we use Noether normalisation? But maybe it will be
% easier to just use the KPX argument above in the end?}

% % \[\cC^\bullet(M\otimes_AB)\isoto
% % \cC^\bullet((M\otimes_AB)/T^n\gM^B)=\cC^\bullet((M/T^n\gM)\otimes_AB),\]

% It is therefore enough to note that the natural
% map \[(M/T^n\gM)\otimes_AB=(M/T^n\gM)\otimes_{\gS_A}(\gS_A\otimes_AB)\to
%   (M/T^n\gM)\otimes_{\gS_A}\gS_B\] is an isomorphism, because every
% element of
% $(M/T^n\gM)$ is $T$-power torsion, and $\gS_A\otimes_AB$ is $T$-adically dense in~$\gS_B$.

% Let $\gM^B$ be the inverse image of~$\gM_B$ in $M\otimes_AB$. We claim
% that~$\gM^B$ is a $(\varphi,\Gamma)$ lattice in $M\otimes_AB$;   % \mar{TG: again, not
%   % actually defined for this object.}
% it then follows as above that
% $\cC^\bullet(M\otimes_AB)\to\cC^\bullet(M\otimes_AB/T^n\gM^B)$ is a
% quasi-isomorphism, and since
% $\cC^\bullet(M\otimes_AB/T^n\gM^B)\to\cC^\bullet(M_B/T^n\gM_B)$ is an
% isomorphism of complexes, we will be done.

% Since $\gM_B$ is open in~$M_B$, $\gM^B$ is open in $M\otimes_AB$, and
% therefore $\gM^B\otimes_{\gS_B}\OEB=M\otimes_AB$. It therefore
% suffices to show that $\gM^B$ is finitely generated as an
% $\gS_B$-module, and since $\gS_B$ is Noetherian, it suffices in turn
% to show that $\gM^B$ is contained in a lattice.
% \mar{TG: haven't
  % thought this through, but hopefully not hard to show by comparing
  % everything to a lattice $\gM$ in $M$?}



% To do this, we return to the beginning of the proof of part~(1), and
% to the quasi-isomorphism $\cC^\bullet(M)\isoto \cC^\bullet(M/T^n\gM)$. Since
% $\gM\otimes_{\gS_A}\gS_B$ is a $(\varphi,\Gamma)$-stable lattice
% in  \mar{TG: warning: should check that this is true if we don't have
%   free lattices, but just torsion-free, and if it isn't true, think
%   about what we need to say instead. TG: presumably need to either check that
%   $\gM\otimes_{\gS_A}\gS_B$ is $T$-torsion free, or to replace it with
% its image in $M_B$ and see if the next claim is OK; probably it is?
% TG: still needs some thought!}
% $M\otimes_{\OEA}\OEB$, the same argument shows that the natural map
% $\cC^\bullet(M\otimes_{\OEA}\OEB)\to
% \cC^\bullet((M\otimes_{\OEA}\OEB)/T^n(\gM\otimes_{\gS_A}\gS_B))$ is a
% quasi-isomorphism. It is therefore enough to show that there is a natural
% quasi-isomorphism \[\cC^\bullet(M/T^n\gM)\otimes^{\mathbb{L}}_AB\isoto
%   \cC^\bullet((M\otimes_{\OEA}\OEB)/T^n(\gM\otimes_{\gS_A}\gS_B)).\]Now, since
% $M/T^n\gM$ is a flat $A$-module, the derived tensor product is
% computed by the ordinary tensor product $\cC^\bullet(M/T^n\gM)\otimes_AB$. Since
% we have a natural isomorphism \[(A((T))/A[[T]])\otimes_AB\isoto B((T))/B[[T]], \] we
% have a natural isomorphism \[ (M/T^n\gM)\otimes_AB\isoto
%   (M\otimes_{\OEA}\OEB)/T^n(\gM\otimes_{\gS_A}\gS_B),\] which
% induces a termwise isomorphism of the two complexes, as required.

% Finally, for surjectivity, it is enough to check that for each $m\in
% M$, $\psi^i(m)\to 0$ in $M/T^{-n}\gM$, as we can then write
% $m=(1-\psi)(\sum_{i\ge 0}\psi^i(m))$. \mar{TG: I don't think I'm
%   considering the topological issues very well, but hopefully this is
%   essentially correct.}
% Considering the long exact sequence of homology groups associated to
% the short exact sequence of
% complexes \[0\to\cC^\bullet(T^n\gM)\to\cC^\bullet(\gM)\to\cC^\bullet(\gM/T^n\gM)\to
%   0,\]
% we see that $\cC^\bullet(\gM)\to\cC^\bullet(\gM/T^n\gM)$ is a 
% quasi-isomorphism. Since $\cC^\bullet(\gM/T^n\gM)$ is concentrated in
% degrees $[0,2]$, and all its terms are finite rank projective
% $A$-modules, we are done.

% By Proposition~\ref{prop: Herr complex in terms of psi}, we have a
% sequence of
% quasi-isomorphisms \[\cC_\psi^\bullet(T^{n'}\gM')\to\cC_\psi^\bullet(M)\to
%   \cC^\bullet(M)\to\cC^\bullet(M/T^n\gM).\]\mar{TG: this argument is
%   bogus - the map in the middle doesn't exist! So I don't think we can
% argue in this way. Instead I'm guessing we should argue that we can
% truncate in each direction, showing (?) that the cohomology groups are
% both projective and dual-projective and thus finite (?).} 

% These morphisms are
% continuous in all degrees, and it follows from~\cite[Lem.\
% 3.1]{MR2181808} that each term of the image
% of~$\cC_\psi^\bullet(T^{n'}\gM')$ in~$\cC^\bullet(M/T^n\gM)$ is
% finitely generated as an~$A$-module. In particular, the cohomology groups of
% the Herr complex are finitely generated, so
% by~\cite[\href{http://stacks.math.columbia.edu/tag/066E}{Tag
%   066E}]{stacks-project} the Herr complex is pseudo-coherent. It then
% follows from~\cite[\href{http://stacks.math.columbia.edu/tag/0658}{Tag
%   0658}]{stacks-project} that it is a perfect complex concentrated in
% degrees~$[0,2]$, as required.


\begin{cor}% \mar{TG: not sure if we ultimately use this result or
    % not.}
  \label{cor: Herr complex is perfect for p-adic A}Let $A$ be a
  $p$-adically complete Noetherian $\cO$-algebra such that~$A/\varpi$ is countable, and let $M$ be a projective
  \'etale $(\varphi,\Gamma)$-module. Then the Herr complex
  $\cC^\bullet(M)$ is a perfect complex concentrated in degrees $[0,2]$.
\end{cor}%\mar{TG: might need a Noetherian hypothesis?}
\begin{proof}This follows
  from~\cite[\href{https://stacks.math.columbia.edu/tag/0CQG}{Tag 0CQG}]{stacks-project} and Theorem~\ref{thm: Herr complex is perfect}.
(We apply the base-change statement of part~(2) of the theorem
to the surjective (and hence finite) morphisms
$A/\varpi^{a+1} \to A/\varpi^a$.)
%  \mar{TG: proof needed}
%\mar{ME: This seems to implicitly use the base-change statement for  Herr complexes,
%but only for the surjections $A/\varpi^{a+1} \to A/\varpi^a$ (which
%are finite morphisms).}
\end{proof}

We also note the following rather technical corollary
Theorem~\ref{thm: Herr complex is perfect},
which we will use in our discussion of families of extensions below.

\begin{cor}
\label{cor:Herr complex length}
If $A$  is of finite type over $\cO/\varpi^a$ for some $a  \geq 1$,
and $M$ is a projective \'etale  $(\varphi,\Gamma)$-module,
then the Herr complex $\cC^{\bullet}(M)$ can be represented
by a complex $C^0 \to C^1 \to C^2$ of finite projective $A$-modules
in degrees~$[0,2]$.
\end{cor}
\begin{proof}
Theorem~\ref{thm: Herr complex is perfect}
shows that $\cC^{\bullet}(M)$  is perfect,
and so 
by~\cite[\href{https://stacks.math.columbia.edu/tag/0658}{Tag 0658}]{stacks-project}, 
the present corollary will follow  provided that we show that
$\cC^{\bullet}(M)$ has tor-amplitude in~[0,2].
Since any $A$-module is a filtered direct limit of finite $A$-modules,
it suffices to  show that
$\cC^{\bullet}(M)\Lotimes_A M$
has cohomological amplitude lying in~$[0,2]$
for any finite $A$-module~$M$.   
If we let $B  = A  \oplus M$, thought of as an $A$-algebra by  declaring
$M$ to be  a square-zero  ideal,
then it is  equivalent to show that
$\cC^{\bullet}(M)\Lotimes_A B$
has cohomological amplitude lying in~$[0,2]$.
But Theorem~\ref{thm: Herr complex  is perfect}
shows  that this latter complex is isomorphic (in the derived category of~$B$-modules)
to
$\cC^{\bullet}(M\otimes_{\A_{K,A}}  \A_{K,B}),$
whose cohomological amplitude  lies in~[0,2] by construction.
\end{proof}



We now explain the comparison between the Herr complex and Galois
cohomology, and the relationship to Tate local duality. These results
are essentially due to Herr and are proved
in~\cite{MR1693457,MR1839766} but we follow~\cite{MR3230818} and
formulate them as statements in the derived category.



%\mar{TG: now want cup product, Tate duality map}

Exactly as in~\cite[Defn.\ 2.3.10]{MR3230818}, if $M_1$, $M_2$ are
projective \'etale $(\varphi,\Gamma)$-modules with $A$-coefficients,
then there is a cup
product \[\cC^\bullet(M_1)\otimes_A\cC^\bullet(M_2)\to\cC^\bullet(M_1\otimes_{\A_{K,A}}M_2).\]
%\mar{TG: put in details.}

More precisely, the cup product arises from the following
generalities. If we have two complexes of $A$-modules~$C^\bullet$
and~$D^\bullet$, then the tensor product $C^\bullet\otimes D^\bullet$
has differential given by \[d(x\otimes y)=dx\otimes y+(-1)^i x\otimes
  y\] if $x\in C^i$, $y\in D^j$. If $f^\bullet:C^\bullet\to
C^\bullet$, then we write $\Fib(f|C^\bullet):=\Cone(f)[-1]$, which by definition
has $\Fib(f)^i=C^i\oplus C^{i-1}$ and
$d^i((x,y))=(d^i(x),-d^{i-1}(y)-f^i(x))$.

As a special case of~\cite[Lem.\ 2.3.9]{MR3230818}, if
$f_1:C^\bullet_1\to C^\bullet_1$, $f_2:C^\bullet_2\to C^\bullet_2$,
then we have a natural
morphism \numequation\label{eqn: cup product of fibre complex}\Fib(1-f_1|C_1^\bullet)\otimes \Fib(1-f_2|C_2^\bullet)
  \to
  \Fib(1-f_1\otimes f_2|C_1^\bullet\otimes C^\bullet_2) \end{equation}%\mar{TG: here}

Then by definition we have
$\cC^\bullet(M)=\Fib\bigl(1-\gamma|\Fib(1-\varphi|M)\bigr)$, so that by multiple
applications of~(\ref{eqn: cup product of fibre complex}) we have morphisms
\begin{align*}
  \cC^\bullet(M_1)\otimes_A\cC^\bullet(M_2)&=\Fib\bigl(1-\gamma|\Fib(1-\varphi|M_1)\bigr)\otimes_A\Fib\bigl(1-\gamma|\Fib(1-\varphi|M_2)\bigr)\\
  &\to \Fib\Bigl(1-\gamma|\bigl(\Fib(1-\varphi|M_1)\otimes_A\Fib(1-\varphi|M_2)\bigr)\Bigr)\\& \to
  \Fib\Bigl(1-\gamma|\Fib\bigl(1-\varphi|(M_1\otimes_A M_2)\bigr)\Bigr)\\&\onto  \Fib\Bigl(1-\gamma|\Fib\bigl(1-\varphi|(M_1\otimes_{\A_{K,A}} M_2)\bigr)\Bigr)\\
&=\cC^\bullet(M_1\otimes M_2)
\end{align*}
whose composite defines the cup product.

% \begin{lem}
%   \label{lem: H2 commutes with base change}
% \end{lem}

\begin{lem}\label{lem: H2 of Tate twist}
  If~$A$ is a finite type $\cO/\varpi^a$-algebra, then there is an isomorphism
%\mar{ME: Looking at the citation of this just below,
%I think we want $A$ to be finite type, not  just finite.  But maybe
%also over some $\cO/\varpi^a$ rather than $\cO$, since our base-change
%results etc.\  are stated in that context.}
  \[H^2(\cC^\bullet(\A_{K,A}(1)))\cong A,\] compatible with base change.
\end{lem}
\begin{proof}
  Since by Theorem~\ref{thm: Herr complex is perfect}~(2) the formation of~$H^2(\cC^\bullet(M))$ is
  compatible with base change, the result follows from the
  case~$A=\cO/\varpi^a$, which is immediate from Theorem~\ref{thm: Herr complex
    and Galois reps} below (that is, from the natural isomorphism
  $H^2\bigl(G_K,(\cO/\varpi^a)(1)\bigr)~\cong~\cO/\varpi^a$). 
\end{proof}

If $M$ is a projective \'etale $(\varphi,\Gamma)$-module %of rank $d$
over a finite type $\cO/\varpi^a$-algebra,
then we  define the Tate duality
pairing between the Herr complexes of~$M$ and of its Cartier dual~$M^*$ as the following composite of the cup product, truncation, and the
isomorphism of Lemma~\ref{lem: H2 of Tate twist}: \[\cC^\bullet(M)\times\cC^\bullet(M^*)\to
  \cC^\bullet(\A_{K,A}(1)))\to H^2(\cC^\bullet(\A_{K,A}(1)))[-2]\cong A[-2].\]

%  we have the following
% \mar{TG: now Tate duality and Euler char - 2.3.7 in \cite{MR3230818}
%   but actually they refer to elsewhere anyway.}
% If~$V$ is a finitely
%   generated $\Zp$-module with a continuous action of~$G_K$, then we
%   have a natural\mar{TG: ?} quasi-isomorphism \[ \]
  % write $\mathbb{D}(V)$ for the corresponding
%   $(\varphi,\Gamma)$-module, and $H^i(\mathbb{D}(V))$ for the
%   cohomology of the Herr complex of~$\mathbb{D}(V)$.\mar{TG: actually,
%   rather than comparing cohomology groups, we should presumably
%   compare in the derived category if we can? ME: Makes sense, but
% perhaps then also cite/mention Herr's original paper for the isomorphism
% on the level of cohomologies?}

\begin{prop}
  \label{prop: Herr complex Tate duality and Euler char}Let $A$ be a
  finite type $\cO/\varpi^a$-algebra, and let $M$ be a projective
  \'etale $(\varphi,\Gamma)$-module. % of rank~$d$. %\mar{TG: maybe just mod $p^a$?}
  \begin{enumerate}
%		  \mar{ME: What does {\em Euler char.} mean when the
%			  cohomology may not be locally free? TG: I
%                          think it was supposed to be a field, so I've
%                        adjusted this. ME: OK, commenting out}
  \item The Tate duality pairing induces a quasi-isomorphism \[\cC^\bullet(M)\isoto\RHom_A(\cC^\bullet(M^*),A))[-2]. \]
  \item If $A$ is a finite extension of~$\F$, then the Euler characteristic~$\chi_A(\cC^\bullet(M))$ is equal to
    $-[K:\Qp]d$.

  \end{enumerate}%\mar{TG: define $M^*$.}
\end{prop}
\begin{proof}For part~(1), by~\cite[Lem.\ 4.1.5]{MR3230818}, it is enough to treat
  the case that~$A$ is a field, in which case~$A$ is a finite extension
  of~$\F$. Then both parts  follow from Theorem~\ref{thm: Herr complex
    and Galois reps} below (that is, from the corresponding statements for
  Galois representations).
%
% Sketch: I think we can probably follow KPX, at least if
%   we're happy to work with finite type $\Z/p^a\Z$-algebras. The point
%   is that I think we can then define the Tate pairing etc just using
%   the field case and extension of scalars, and then checking that we
%   have a quasi-isomorphism reduces to the finite field case. If we
%   need something more general then we might have to go into the actual
%   arguments!
\end{proof}%\mar{TG to write all this properly}

Finally we recall the relationship between the Herr complex and Galois
cohomology. As in~\cite[\S2]{MR3117501} it is possible to upgrade the
following theorem to an isomorphism in the derived category, but as we
do not need this we do not give the details here. If~$A$ is a complete local Noetherian
  $\cO$-algebra with finite residue field, and~$M$ is a formal
  projective \'etale $(\varphi,\Gamma)$-module with $A$-coefficients,
  then we can define the Herr complex~$\cC^\bullet(M)$ exactly as for
  \'etale $(\varphi,\Gamma)$-modules, so that by definition we have
  $\cC^\bullet(M)=\varprojlim_n\cC^\bullet(M_{A/\m_A^n})$. By \cite[\href{https://stacks.math.columbia.edu/tag/0CQG}{Tag 0CQG}]{stacks-project}, 
  $\cC^\bullet(M)$ is a perfect complex, concentrated in degrees~$[0,2]$.
\begin{thm}% \mar{TG: should probably change to $A$ is CNL as above in
    % Galois discussion, and check if it's in Dee.}
%\mar{TG: need to sort out the Herr complex in the CNL case!}
  \label{thm: Herr complex and Galois reps}If~$A$ is a complete local Noetherian
  $\cO$-algebra with finite residue field, and $T$ is a finitely
  generated projective $A$-module with a continuous action
  of~$G_K$, % \mar{TG: not 100\% sure this is the correct hypothesis but
    % must be close!}
  then there
  are isomorphisms of $A$-modules \[H^i(G_K,T) \isoto H^i(\cC^\bullet(\mathbb{D}_A(T)))
    \] % in the derived category of perfect
  % complexes over~$A$,
  % are isomorphisms \[H^i(\mathbb{D}(V))\cong H^i(G_K,V).\],
  which are
  functorial in~$T$ and compatible with cup products and duality.% , and
  % base change.
\end{thm}
\begin{proof}This is~\cite[Prop.\ 3.1.1]{MR1805474}, which is deduced from the results
  of~\cite{MR1839766,MR1693457} by passage to the limit. % There is a short exact sequence\mar{TG: guessing about this}\[ 0\to\Zp\to
%     \AKnrhat\stackrel{\varphi-1}{\to}\AKnrhat\to 0. \]
%
% It is enough to prove this in the case that~$A$ is a
%   finite $\Z/p^a\Z$-algebra.\mar{TG: justify!} We follow~\mar{TG: pottharst?}
  %
% \mar{TG: I guess we should probably be explaining the relationship to $\Gamma$-cohomology?}
\end{proof}
\begin{rem}
  \label{rem: writing Herr complex as Galois cohomology}In accordance
  with our general convention of writing $\rho_T$ for a family
  $T\to\cX_d$, if~$T$ is an affine scheme of finite type
  over~$\cO/\varpi^a$, then we write $H^2(G_K,\rho_T)$ for the
  pull-back to~$T$ of the cohomology group $H^2(\cC^\bullet(M))$
  on~$\cX_d$, where~$M$ is the \'etale $(\varphi,\Gamma)$-module
  corresponding to the morphism $T\to\cX_d$. By Theorem~\ref{thm: Herr
    complex is perfect}, $H^2(G_K,\rho_T)$ is a coherent sheaf, whose
  formation is compatible with arbitrary finite type base-change, and
  so in particular by Theorem~\ref{thm: Herr complex and Galois reps}
  its specialisations at $\Fpbar$-points coincide with the usual
  Galois cohomology groups. This compatibility with base-change also
  allows us to extend the definition of $H^2(G_K,\rho_T)$ to arbitrary
  (not necessarily affine) schemes~$T$ of finite type over~$\cO/\varpi^a$.
\end{rem}
\begin{lem}\label{lem: Tate duality for H2 to H0 in families}
Let~$T$ be a scheme of finite type over~$\cO/\varpi^a$, % \mar{TG: at some
  % point we have to sort out what our coefficients are, e.g.\ to allow
  % $\Fpbar$.}
and let~$\rho_T$, $\rho'_T$ be families of Galois representations
  over~$T$. Then sections of~$H^2(G_K,\rho_T\otimes(\rho'_{T})^\vee(1))$
  are in natural bijection with homomorphisms $\rho_T\to\rho'_{T}$.
\end{lem}
\begin{rem}
  \label{rem: clarifying what a morphism of families means}Here by a
  homomorphism $\rho_T\to\rho'_T$ we mean a homomorphism of the
  corresponding $(\varphi,\Gamma)$-modules.
\end{rem}
\begin{proof}
  [Proof of Lemma~\ref{lem: Tate duality for H2 to H0 in
    families}]% We may suppose that~$T$ is affine, say~$T=\Spec
  % A$.\mar{TG: I guess we should justify this!} Let $M$, $M'$ be the
  % projective \'etale $(\varphi,\Gamma)$-modules with $A$-coefficients
  % corresponding to $\rho_T$, $\rho'_T$ respectively.
  % By
This follows from   Proposition~\ref{prop: Herr complex Tate duality
  and Euler char} and Lemma~\ref{lem: Yoneda for H0 H1 of Herr complex}.
  % sections of $H^2(G_K,\rho_T\otimes(\rho'_{T})^\vee(1))$ are in
  % natural bijection with elements of $H^0(\cC^\bullet(M^\vee\otimes
  % M'))$, which by definition correspond to homomorphisms $M\to M'$, as required.
\end{proof}
\subsection{Obstruction theory}\label{subsubsec: obstruction
  theory}% \mar{TG: this is where we will prove that we have a nice
  % obstruction theory}
We now show that the Herr complex provides~$\cX_d$ with a nice
obstruction theory in the sense of Definition~\ref{df:obstruction}.% recalled in Appendix~\ref{app: formal algebraic
  % stacks}.
% of~\cite[\S
% 10]{Emertonformalstacks}.

Let $A$ be a $p$-adically complete $\cO$-algebra, and let $x:\Spec A\to
\cX_d$ be a morphism, corresponding to a projective $(\varphi,\Gamma)$-module
$M$. We wish to consider the problem of deforming~$M$ to a square zero
thickening of~$A$. Specifically, if  \[0\to I\to A'\to A\to 0\] is a
square zero extension, then we let $\Lift(x,A')$ be the set of
isomorphism classes of projective $(\varphi,\Gamma)$-modules $M'$ with
$A'$-coefficients which have the property that $M'\otimes_{A'}A\cong M$. 

For any such thickening~$A'$, we define a corresponding obstruction
class as follows. The underlying $\A_{K,A}$-module of~$M$ has a unique (up
to isomorphism) lifting to a projective $\A_{K,A'}$-module~$\tM$, and
we may lift $\varphi,\gamma$ to semi-linear endomorphisms~$\varphit,\gammat$
of~$\tM$. (Indeed, $\A_{K,A'}$ is a square zero thickening of~$\A_{K,A}$,
and a finite projective module~$P$ over any ring~$R$ deforms
uniquely through any square zero extension $R'\to R$, as its
deformations are controlled by
\[H^1(\Spec R,\ker(R'\to R)\otimes\End_R(P)),\] which vanishes (as
$\Spec R$ is affine). To see that $\varphi,\gamma$ lift, think of them
as $\A_{K,A}$-linear maps $\varphi^*M\to M$, $\gamma^*M\to M$.)

However, there is no guarantee that we can find lifts
~$\varphit,\gammat$ which commute. To measure the obstruction to the
existence of such lifts, let $\ad M=\Hom_{\A_{K,A}}(M,M)$
be the adjoint of~$M$. This naturally has the structure of
a~$(\varphi,\Gamma)$-module; indeed, we have a natural identification $\varphi^*\ad
M=\Hom_{\A_{K,A}}(\varphi^*M,\varphi^*M)$, and we define $\Phi_{\ad M}:\varphi^*\ad
M\to\ad M$ by \[(\Phi_{\ad M}(f))(x)=\Phi_M(f(\Phi_M^{-1}(x))). \]We
define the action of~$\Gamma$ in the analogous way. Then we let $o_x(A')$ be the image in
$H^2(\cC^\bullet(\ad M))\otimes_A I=H^2(\cC^\bullet(\ad M\otimes_A I))$ of
\[\varphit\gammat\varphit^{-1}\gammat^{-1}-1\in \ad M\otimes_A
I=\cC^2(\ad M\otimes_AI).\]

\begin{lem}
  \label{lem: obstruction class well defined and controls existence of
  lifts}% Suppose that~$A'$ is \mar{TG: why are we assuming~$A'$
  % Noetherian? TG: I can see no reason for this, removing and
  % commenting out.}
  % Noetherian. Then t
  The cohomology class $o_x(A')$ is well-defined independently
  of the choice of~$\varphit,\gammat$, and vanishes if and only if
  $\Lift(x,A')\ne 0$.
\end{lem}
\begin{proof}Any other liftings $\varphit',\gammat'$ are  obtained from
  our given liftings $\varphit,\gammat$ by setting
  $\varphit'=(1+X)\varphit$, $\gammat'=(1-Y)\gammat$, for some $X,Y\in\ad
  M\otimes_A I$. A simple computation shows
  that \[\varphit'\gammat'(\varphit')^{-1}(\gammat')^{-1}-\varphit\gammat\varphit^{-1}\gammat^{-1}=(1-\gamma)X+(1-\varphi)Y, \]which
  shows that the cohomology class $o_x(A')$ is well-defined, and that it vanishes if and
  only if we can choose $\varphit',\gammat'$ so that
  $\varphit'\gammat'(\varphit')^{-1}(\gammat')^{-1}=1$, which is in
  turn equivalent to $\Lift(x,A')\ne 0$, as required.  
\end{proof}

% We sometimes want to consider a particular class of thickenings which
% is defined as follows.
If~$F$ is an $A$-module, we let~$A[F]:=A\oplus
F$ be the $A$-algebra with multiplication given
by \[(a,m)(a',m'):=(aa',am'+a'm).\] This is a square zero thickening
of~$A$, and $\Lift(x,A[F])\ne 0$, because we have the trivial lifting
given by $M\otimes_AA[F]$. 



\begin{lem}
  \label{lem: H1 of Herr complex of adjoint controls lifts}
 Suppose that~$F$ is a finitely generated $A$-module. Then
    there is a natural isomorphism of $A$-modules $\Lift(x,A[F])\isoto
    H^1(\cC^\bullet(\ad M)\otimes_A^{\mathbb{L}}F)$. % \mar{TG: need to
      % explain the $A$-module structure on $\Lift(x,A[F])$. Presumably
      % it is just scaling $\varphit-\varphi$, $\gammat-\gamma$, but I
      % don't know if there is a short way to put it?}
  % \item Suppose that $0\to I\to A'\to A\to 0$ is a square zero
  %   thickening of~$A$, and that~$A'$ is Noetherian. Then $o_x(A')\ne 0$
  %   if and only if $\Lift(x,A')\ne 0$.
\end{lem}
\begin{proof}We begin by constructing the isomorphism on the level of sets. Note that $\cC^\bullet(\ad M)\otimes^{\mathbb{L}}F$ is
  computed by $\cC^\bullet(\ad M\otimes_AF)$. Liftings of~$M$
  to~$A[F]$ are determined by the corresponding liftings
  $\varphit,\gammat$ of $\varphi,\gamma$, and we obtain a class in
  $H^1(\cC^\bullet(\ad M\otimes_AF))$ by taking the image of
  of \[(\varphit\varphi^{-1}-1,\gammat\gamma^{-1}-1)\in\ad M\otimes_A F\oplus\ad M\otimes_A F =\cC^1(\ad
    M\otimes_AF).\] An elementary calculation shows that
  $(\varphit\varphi^{-1}-1,\gammat\gamma^{-1}-1)$ is in the kernel of $\cC^1(\ad
    M\otimes_AF)\to\cC^2(\ad M\otimes_AF)$ if and only if
    $\varphit\gammat=\gammat\varphit$, so it only remains to check
    that the lifting given by $\varphit,\gammat$ is trivial if and
    only if the corresponding cohomology class vanishes. 
  
    To see this, note that the endomorphisms of the trivial lifting
    are of the form $1+X$, for some $X\in\ad M\otimes F$. The corresponding
    $\varphit,\gammat$ are given by $\varphit=(1+X)\varphi(1+X)^{-1}$,
    $\gammat=(1+X)\gamma(1+X)^{-1}$, which is equivalent to
    $(\varphit\varphi^{-1}-1,\gammat\gamma^{-1}-1)=((1-\varphi)
    X,(1-\gamma) X)$, which by definition is equivalent to
    $(\varphit\varphi^{-1}-1,\gammat\gamma^{-1}-1)$ being a
    coboundary, as required.

To compare the $A$-module structures, recall that by definition the
$A$-module structure on $\Lift(x,A[F])$ is defined as follows
(see~\cite[\href{http://stacks.math.columbia.edu/tag/07Y9}{Tag
  07Y9}]{stacks-project}). If $r\in A$, then we have a homomorphism
$f_r:A[F]\to A[F]$ given by $f_r(a,f)=(a,rf)$, and given a lifting
$\tM$ of~$M$, we let~$r\tM$ be the base change of~$\tM$
via~$f_r$. Explicitly, this means that we replace~$(\varphit-\varphi)$
and~$(\gammat-\gamma)$ by ~$r(\varphit-\varphi)$
and~$r(\gammat-\gamma)$.

Similarly, the addition map
$\Lift(x,A[F])\times\Lift(x,A[F])\to \Lift(x,A[F])$ comes from the
obvious identification
$\Lift(x,A[F])\times\Lift(x,A[F])=\Lift(x,A[F\times F])$ together with
base change via the homomorphism $A[F\times F]\to A[F]$ given by
$(a,f_1,f_2)\mapsto (a,f_1+f_2)$. With obvious notation, this amounts
to setting $\varphit_1\boxplus\varphit_2:=\varphit_1+\varphit_2-\varphi$, 
$\gammat_1\boxplus\gammat_2:=\gammat_1+\gammat_2-\gamma$.

On the other hand, by definition the $A$-module structure on
$H^1(\cC^\bullet(\ad M)\otimes^{\mathbb{L}}F)$ is given by the obvious
$A$-module structure on the pairs
$(\varphit\varphi^{-1}-1,\gammat\gamma^{-1}-1)$. This is obviously the
same as the $A$-module structure that we have just explicated on~$\Lift(x,A[F])$.
% :=r\varphit+(1-r)\varphi$ and
% $r\gammat:=r\gammat+(1-r)\gamma$, and setting
% $\varphit_1\boxplus\varphit_2:=\varphit_1+\varphit_2-\varphi$, 
% $\gammat_1\boxplus\gammat_2:=\gammat_1+\gammat_2-\gamma$.\mar{TG:
%   obviously if this is correct the notation should be explained}
%
% To compare these, we need to identify the pullback of~$M\otimes_AA[F]$
% via~$f_r$ with $M\otimes_A A[F]$, and compute the corresponding change
% to~$\varphit,\gammat$ (and similarly for the addition map). It is
% obviously enough to compare the action on elements of the form $fm$,
% $f\in F$, $m\in M$, and then the result is clear.\mar{TG: I got quite
%   confused when I actually tried to do this, but hopefully there is
%   enough information here for us to sort out the correct thing.}
\end{proof}

\begin{prop}
  \label{prop: Xd has a good obstruction theory}$\cX_d$ admits a nice
  obstruction theory in the sense of Definition~{\em \ref{df:obstruction}}.
\end{prop}
\begin{proof}
  Since~$\cX_d$ is limit preserving by Lemma~\ref{lem:X is limit
    preserving}, % \mar{TG: limit preserving will
  %   hopefully be proved in Prop~\ref{prop:X is an Ind-stack} when that
  % is finished.}
it follows from Theorem~\ref{thm: Herr complex is perfect} and % ~\cite[Rem.\
  % 11.9]{Emertonformalstacks} and
  Lemmas~\ref{lem: obstruction class
    well defined and controls existence of lifts} and~\ref{lem: H1 of
    Herr complex of adjoint controls lifts} that the Herr
  complex~$\cC^\bullet(\ad M)$ provides the required obstruction
  theory.%  (with the auxiliary module $C_{(x,A)}$ in~\cite[Defn.\
  % 11.6]{Emertonformalstacks} being zero).
\end{proof}
% \begin{prop}
%   \label{prop: Herr complex is nice obstruction}Let $A$ be a
%   Noetherian $\Z/p^a\Z$-algebra, and let~$M$ be a locally free
%   $(\varphi,\Gamma)$-module of rank~$d$ with % \mar{TG: note that this
%     % notation is used above but conflicts with the notation in \cite{Emertonformalstacks}.}
%   $A$-coefficients. Let \[0\to I\to A'\to A\to 0\] be a square zero
%   extension, where $I$ is a finite type $A$-module.
% \end{prop}\mar{TG: to be continued}
% \begin{proof}
%   \mar{TG: needed}
% \end{proof}

% Quick statements: we use $\cC^\bullet(\ad M)$. For the obstruction
% theory, we want to use $\varphit\gammat\varphit^{-1}\gammat^{-1}$. For
% the tangent space, we consider $\varphit\varphi^{-1}-1$ and $\gammat\gamma^{-1}-1$.
%\mar{TG: didn't type up 5.2 edits yet}
\section{Residual gerbes and isotrivial families}
\label{subsec:isotrivial}
In this section we briefly discuss the notion 
of isotrivial families of $(\varphi,\Gamma)$-modules
over reduced $\F$-schemes; i.e.\ of families
which are pointwise constant.   
The language of residual gerbes (see Appendix~\ref{app:residual gerbes}) provides
a convenient framework for doing this.
%
%We also slightly modify our notational conventions. \mar{ME: Not sure what
%conventions are being referred to here.}  Namely, we allow
In the following discussion we allow $\Fnew$ to denote
any algebraic extension of $\F$. %; in particular, this means that~$\F'$ contains $k$, the residue field of the
%ring of integers $\cO_K$ of~$K$.  (The reason for allowing this
%slightly more general choice of coefficients
%is that we want to consider finite type points of~$\cX_d$, which will
%be defined over finite extensions of~$\F$; it is also sometimes convenient to be able to
%work over $\Fbar_p$.)% , rather than over the given finite extension $\F$ of $\Fbar_p$
%% that was fixed in the discussion of $\cX_d$ up until now.)
%%\mar{ME: Have to replace $\Fbar_p$ by $\F$ more-or-less
%%  throughout. TG: I wasn't convinced that changing the definition
%%  of~$\F$ in this way was a good idea, so at least for now I've gone
%%  with a macro. Also it seemed like a good idea for~$\Fnew$ to be an~$\F$-algebra.}


We have seen that $\cX_d$ is a quasi-separated Ind-algebraic stack,
which by Proposition~\ref{prop:X is an Ind-stack} may be written as the inductive limit of algebraic stacks 
with closed immersions (and so, in particular, monomorphisms) as
transition morphisms.  Thus the
discussion at the end of Appendix~\ref{app:residual gerbes} applies,
and in particular, for each point $x \in |\cX_d|,$
the residual gerbe $\cZ_x$ at $x$ in  $\cX_d$ exists.
Furthermore, if $x$ is a finite type point, then the canonical
monomorphism $\cZ_x \hookrightarrow \cX_d$ is an immersion.

%Since $\cX_d$ is quasi-separated and locally Noetherian,
%the results of Appendix~\ref{app:residual gerbes} apply.  In particular,
%every point of $\cX_d$ admits a residual gerbe.\mar{TG: note that at
%  the moment these properties are only established rather later in the
%paper, namely in Corollary~\ref{cor:Xd is formal algebraic}, so I
%don't know if we should move this material to the end of section 3?
%I'm not quite sure if/where it's actually used - I thought it might be
%used in section 3, but the word ``gerbe'' doesn't appear as far as I
%can see\dots}

If we consider an $\Fnew$-valued point 
$x: \Spec \Fnew \to \cX_{d}$
(by abuse of notation we use $x$ to denote both this point,
and its image $x \in |\cX_d|$, which is a finite type point of~$\cX_d$),
then we may base-change the residual gerbe to $\cX_{d}$ at $x$
(which is a gerbe over~$\F$) %\mar{TG: confused about this; since
 % $\cX_d$ is defined over~$\F$, I would have expected this to be a gerbe over~$\F$?}
over $\Fnew$ via the composite $\Spec \Fnew \buildrel x \over
\longrightarrow \cX_{d} \to \Spec \F$; equivalently,
we may regard $x$ as a point of the base-change $(\cX_{d})_{\Fnew}$, 
and then consider the residual gerbe of $(\cX_{d})_{\Fnew}$ at
this point.
This residual gerbe is then of the form $[\Spec \Fnew/G]$,
for a finite type affine group scheme $G$ over $\Fnew$.
(By~\cite[\href{http://stacks.math.columbia.edu/tag/06QG}{Tag
  06QG}]{stacks-project} the group $G$ may be described as the fibre
product $x \times_{\cX_{d}} x$,
and hence its claimed properties follow from the
fact that by Proposition~\ref{prop:X is an Ind-stack}, the diagonal of $\cX_{d}$ is affine and of finite
type.) % \mar{TG: should refer to somewhere in the paper for these
%   properties of~$\cX_{d}$. TG: I guess these properties of the
%   diagonal could be added to Proposition~\ref{prop:X is an Ind-stack}?
% I guess the point is that they can be deduced
% from~$\cR_d^{\Gamma_{\disc}}$ and thus from~$\cR_d$?}
%If $\rhobar:G_K \to \GL_d(\Fbar_p)$ is the Galois representation 
%\mar{ME: May be the first mention of Galois reps.\ after the intro.}
%associated to $x$, then $G = \Aut_{G_K}(\rhobar).$
%
In fact, we have the following more precise result regarding $G$.

\begin{lemma}
	\label{lem:automorphisms are smooth}
	If $x$
%	$\rhobar: G_K \to \GL_d(\Fbar_p)$ is
%	a Galois representation,
%	thought of as
        is an $\Fnew$-valued point of $(\cX_d)_{\Fnew}$,
	then $\Aut(x)$
	is an irreducible smooth closed algebraic subgroup
	of~$\GL_{d/\Fnew}$.
\end{lemma}
\begin{proof}As already noted, it follows from Proposition~\ref{prop:X
    is an Ind-stack}  that~$\Aut(x)$ is a finite type affine group
  scheme over~$\F'$.
	Let $D$ denote the \'etale $(\varphi,\Gamma)$-module over $\Fnew$
	corresponding to $x$,
	and let $R:= \End_{\A_{K,\F'},\varphi,\Gamma}(D)$.
	% \mar{ME: Maybe explain the notation in the subscript of
	% 	$\End$, here and below.}
	Then, if $\rhobar:G_K \to \GL_d(\Fnew)$
	denotes the Galois representation corresponding to $x$,
        we see that also $R = \End_{G_K}(\rhobar)$,
	and hence that $R$ is an $\Fnew$-subalgebra
	of $M_d(\Fnew)$. Furthermore, if $A$ is any finite
        type $\Fnew$-algebra,
	and if $D_A$ denotes the base-change of $D$ over~$A$,
	then the natural morphism
	$R\otimes_{\F'} A \to \End_{\A_{K,\F'},\varphi,\Gamma}(D_A)$
%	\mar{ME: Have to prove this!}
	is an isomorphism by Theorem~\ref{thm: Herr complex is
          perfect}~(2) and Lemma~\ref{lem: Yoneda for H0 H1 of Herr
          complex} (note that~$A$ is automatically flat over~$\F'$).
	Thus $$\Aut(x)(A) :=
	\Aut_{\A_{K,A},\varphi,\Gamma}(D_A) = (R\otimes_{\Fnew} A)^{\times}
	= (R\otimes_{\Fnew} A) \bigcap \GL_d(A),$$
	so that $\Aut(x)$, as a scheme over $\Fnew$,
	is precisely the open subscheme $R\cap \GL_d$, where
	we think of $R$ as an affine subspace of the affine space
	$M_d$.  In particular, we see that $\Aut(x)$ is an open
	subscheme of an affine space, and thus smooth and irreducible.
	% \mar{ME: I guess I'm using here implicitly the fact that
	% 	the formation of the ring of endomorphisms is compatible
	% 	with arbitrary base-change.  This should be some
	% 	faithfully flat descent result, combined with the fact
	% 	that every $\Fbar_p$-algebra is flat. ME: Dealt with,
	% to some extent.}
\end{proof}


% If~$M$, $N$ are \'etale $(\varphi,\Gamma)$-module with $A$-coefficients,
% then we write $\Hom_{A,(\varphi,\Gamma)}(M,N)$ for the $A$-linear
% homomorphisms~$M\to N$ which commute with the $\varphi$- and
% $\Gamma$-actions. Similarly, we write $\End_{A,(\varphi,\Gamma)}(M)$
% for the  $A$-linear
% endomorphisms of~$M$ which commute with the $\varphi$- and
% $\Gamma$-actions.
% \begin{lem}\mar{TG: needs to be proved and moved}
%   \label{lem: base change property for endomorphisms of phi Gamma}Let
%   $A$ be a $p$-adically complete $\cO$-algebra, and let~$B$ be a
%   $p$-adically complete flat $A$-algebra. Let~$M,N$ be projective \'etale
%   $(\varphi,\Gamma)$-modules with $A$-coefficients. Then the natural
%   map 	$\Hom_{A,(\varphi,\Gamma)}(M,N)\otimes_{A} B
%   \to \Hom_{B,(\varphi,\Gamma)}(M_B,N_B)$ is an isomorphism.
% \end{lem}
% \begin{proof}As in the proof of Proposition~\ref{prop: descent of
%     morphisms for phi modules}, if we write~$P:=M^\vee\otimes N$ ,
%   then we are reduced to showing
%   that~$P^{\varphi=1,\Gamma=1}\otimes_AB=(P_B)^{\varphi=1,\Gamma=1}$.\mar{TG:
%   not clear to me how to reduce to the free case, or indeed what the
%   argument is from here! I don't know if it might be easier to work
%   with $(\varphi,G_K)$-modules?}  
% \end{proof}

Suppose that $S$ is a reduced $\Fnew$-scheme of finite type,
and that $S \to \cX_d$ is a morphism such that every closed point of
$S$ maps to some fixed $\F'$-valued point $x \in |\cX_d|$.  We can think of
this morphism as classifying a family of Galois representations
$\rhobar_S$ over $S$ whose fibre at each closed point is isomorphic
to the fixed representation $\rhobar: G_K \to \GL_d(\Fnew)$
classified by $x$.  If $G := \Aut_{G_K}(\rhobar)$,
then Lemma~\ref{lem:isotrivial} shows
that the morphism $S \to \cX$ factors through the residual gerbe
$[\Spec \Fnew/G],$ and thus corresponds to an {\em \'etale}
locally trivial $G$-bundle $E$ over $S$.
({\em A priori}, the $G$-bundle $E$ over $S$ is {\em fppf} locally
trivial.  However, since $G$ is smooth, by Lemma~\ref{lem:automorphisms
	are smooth}, we see that $E$ is in fact smooth locally trivial.
By taking an \'etale slice of a smooth cover
over which $E$ trivializes, we see that $E$ is in fact \'etale
locally trivial.)
We can then describe the family $\rhobar_S$ concretely
as a twist by $E$ of the constant family 
$S \times \rhobar,$ i.e.\ $\rhobar_S = E\times_G \rhobar$.
%We may find an \'etale cover $S' \to S$ such that
%the pull-back $E' := E\times_S S'$ is trivial.
%(Choose a smooth surjection $U \to [\Spec \F_p/G]$ with source
%a scheme, and then
%take $S'$ to be an \'etale slice of the smooth morphism
%$U\times_{[\Spec \F_p/G]} S \to S$.)
%\mar{ME: Possibly $G$ is actually smooth itself ---
%	in which case we should note this (if true, it's because
%	it's a general property of automorphisms of representations; 
%	and incidentally, since we're over a perfect field, it suffices
%	to check reduced) ---  and then $S'$ is just an \'etale
%        slice of the smooth (because $G$ is) morphism $E\to S$.}
%	The pull-back $\rhobar_{S'}$ of $\rhobar_S$ to $S'$ 
%	is then isomorphic to the constant family $S' \times \rhobar$.


%\chapter{Almost Galois descent}\label{sec: almost Galois}


\section{Twisting families}\label{subsec: twists of
  families}We now begin our study of the dimensions of certain families of
$(\varphi,\Gamma)$-representations, and the behaviours of these
dimensions under the operations of twisting by families of
1-dimensional representations, and forming extensions of families.%  We
% will abusively write such a family over a scheme~$T$ as~$\rhobar_T$;
% we caution the reader that unless in general such a family of
% $(\varphi,\Gamma)$-modules does not correspond to a Galois
% representation, but since the specialisation of the family at any
% $\Fpbar$-point \emph{does} correspond to a Galois representation, this
% abuse is hopefully not too serious, and is convenient for (for
% example) studying irreducible components of the stacks~$\cX_{d,\red}$
% in terms of the weight part of Serre's conjecture.

% \mar{ME: Put in explanation that we now denote things as $\rho$ rather
%   than $D$. TG: wrote something, which could certainly be improved upon.}
We will need in particular to be able to twist families of representations by
unramified characters. Given an $\Fp$-algebra~$A$ and an element~$a\in
A^\times$, we have a $(\varphi,\Gamma)$-module $M_a$ whose underlying
$\A_{K,A}$-module is free of rank~$1$, generated by some~$v\in M_a$ for which
$\varphi(v)=av$ and $\gamma(v)=v$. % \mar{TG: have made no attempt to get
% the normalisation right for the following statement!}
If~$A=\F$
then the corresponding representation of~$G_K$ is the unramified
character~$\ur_a$ taking a geometric Frobenius to~$a$. The universal instance of this
construction comes by taking $a = x \in A = \F[x,x^{-1}]$;
the corresponding $(\varphi,\Gamma)$-module $M_x$ is then classified by 
% and  corresponds to a  morphism~$\Gm\to\cX_{1}$, which induces
a morphism which we denote ~$\ur_x:\Gm := \Spec \F[x,x^{\pm 1}] \to \cX_1$,
%\cX_{1,\red}$  (where $x$ denotes the variable on 
%  $\Gm$, so that $\Gm := \Spec \F[x,x^{-1}]$). % Since twisting $G_K$-representations by
          % unramified characters corresponds to multiplying the action
          % of~$\varphi$ on $(\varphi,\Gamma)$-modules by scalars, and
          % accordingly we also
which evidently factors through~$\cX_{1,\red}.$
         

Given any morphism $T \to \cX_{d,\red},$ with $T$ a reduced finite
type $\F$-scheme, corresponding to a family $\rhobar_{T}$ of
$G_K$-representations over $T$, we may consider the 
% composite \[\times T\to\cX_{1,\red}\times\cX_{d,\red}\to\cX_{d,\red},\]where\mar{TG:
%   I would have thought that these $\times$ should be $\times_{\Fp}$,
%   but it could be $\times_{\Fpbar}$?}
% the first factor in the first morphism is~$\ur_x$, and the second morphism corresponds to the tensor product of
% $(\varphi,\Gamma)$-modules.\mar{TG: probably this morphism will be
%   discussed elsewhere? It already appears in the introduction, for
%   example.}
% The corresponding family $\rhobar_{T\times\Gm}$
% corresponds to the twists of~$\rhobar_T$ by unramified characters, and
family ~$\rhobar_T \boxtimes\ur_x$ over $T\times\Gm$, as in Section~\ref{subsec: tensor product
  of phi gamma and duality}. We refer to this operation on
$(\varphi,\Gamma)$-modules as \index{unramified twisting} \emph{unramified
          twisting}.  %\mar{TG: again, $\Fp$?}% \mar{TG: also
  % need definition for twist by~$(1)$.}
%over % $T\times_{\Spec \Fbar_p} \Gm$
 


\begin{df}
	\label{def:twistable} \index{twistable} \index{essentially twistable}
We say that~$\rhobar_T$ is \emph{twistable} if whenever
$\rhobar_{t}\cong\rhobar_{t'}\otimes\ur_a$ where $t,t'\in T(\Fpbar)$
and~$a\in\Fpbartimes$, then~$a=1$. We say that it is
\emph{essentially twistable} if for each~$t\in T(\Fpbar)$,
 the set of~$a\in\Fpbartimes$ for which there exists~$t'\in
T(\Fpbar)$  with $\rhobar_{t}\cong\rhobar_{t'}\otimes\ur_a$ is finite.
\end{df}

\begin{lem}\label{lem: dimension of Gm twist}
  If the dimension of the scheme-theoretic image of~$T$
  in~$\cX_{d,\red}$ is~$e$, then the dimension of the scheme-theoretic
  image of~$T\times\Gm$ in~$\cX_{d,\red}$ is at most~$e+1$. If~$T$ 
   contains a dense open subscheme~$U$ such
  that~$\rhobar_U$ is essentially twistable, then equality holds.
\end{lem}
\begin{rem}
  \label{rem: twistable is enough}Since a twistable
  representation is essentially twistable, we see that if~$T$ 
   contains a dense open subscheme~$U$ such
  that~$\rhobar_U$ is twistable, then equality holds in Lemma~\ref{lem: dimension of Gm twist}.
\end{rem}
\begin{proof}[Proof of Lemma~\ref{lem: dimension of Gm twist}]
  % As in the proof of Proposition~\ref{prop:dimension control}, w
  We may assume that~$T$ is irreducible. Write~$f$ for the morphism
  $T\to\cX_{d,\red}$ and~$g$ for the morphism
  $T\times\Gm\to\cX_{d,\red}$.  By~\cite[\href{https://stacks.math.columbia.edu/tag/0DS4}{Tag 0DS4}]{stacks-project}, we
  may, after possibly replacing~$T$ by a nonempty open subscheme,
  assume that for each $\Fpbar$-point $t$ of~$T$ we have
  $\dim T_{f(t)}=\dim T-e$.

  Let~$v=(t,\lambda)$ be an $\Fpbar$-point of~$T\times\Gm$. Then
  $(T\times \Gm)_{g(v)}$ contains~$T_{f(t)}\times\{\lambda\}$, % \mar{TG:
  % hopefully not nonsense here}
so that
  $\dim (T\times \Gm)_{g(v)}\ge \dim T_{f(t)}=\dim T-e=\dim (T\times
  \Gm)-(e+1)$, so the first claim follows from another application
  of~\cite[\href{https://stacks.math.columbia.edu/tag/0DS4}{Tag 0DS4}]{stacks-project}.

  For the last part, we may replace~$T$ by~$U$, and then the
  hypothesis that~$\rhobar_T$ is essentially twistable means that we can identify
  $(T\times \Gm)_{g(v)}$ with a finite union of fibres $T_{f(t')}$
  (indexed by the finitely many~$a$ for which there is a ~$t'$ and an isomorphism $\rhobar_{t}\cong\rhobar_{t'}\otimes\ur_a$), so that equality holds in
  the above inequality, and the result again follows
  from~\cite[\href{https://stacks.math.columbia.edu/tag/0DS4}{Tag 0DS4}]{stacks-project}.
\end{proof}

Recall that for each embedding~$\sigmabar:k\into\F$, we have a
corresponding fundamental character
$\omega_{\sigmabar}:I_K\to\F^\times$, 
% \mar{ME: Based on the formulas below, I think that $\omega_{\sigmabar}$ 
% 	is the right notation for the fundamental character, but TG
% 	should check. TG: yes.}
which
corresponds via local class field theory (normalised to take
uniformizers to geometric Frobeneii) to the composite of~$\sigmabar$
and the natural map $\cO_K^\times\to k^\times$. It is well-known,
and easy to check,
that
$\prod_{\sigmabar}\omega_{\sigmabar}^{e(K/\Qp)}=\epsilonbar|_{I_K}$. If~$\Fnew$
is an algebraic extension of~$\F$, then we
can and do identify the embeddings $k\into\F$ and $k\into\Fnew$.

Let~$\underline{n}=(n_{\sigmabar})_{\sigmabar:k\into\F}$ be a tuple of integers $0\le
n_{\sigmabar}\le p-1$. % , with not all~$n_{\sigmabar}=p-1$.
The characters~$\omega_{\sigmabar}$ can all be extended to~$G_K$, and the restriction
to~$I_K$ of any character $G_K\to\Fpbartimes$ is equal
to~$\prod_{\sigmabar}\omega_{\sigmabar}^{-n_{\sigmabar}}$ for 
some~$\underline{n}$ (which is unique unless the character is unramified). We write~$\psi_{\underline{n}}:G_K\to\F^\times$ for a
fixed choice of an extension
of~$\prod_{\sigmabar}\omega_{\sigmabar}^{-n_{\sigmabar}}$
to~$G_K$. If~$n_{\sigmabar}=0$ for all~$\sigmabar$, or
$n_{\sigmabar}=p-1$ for all~$\sigmabar$ (these are exactly
the cases corresponding to unramified characters), then we fix the choice
$\psi_{\underline{n}}=1$. We can and do make  our choice so that
if~$\underline{n}$, $\underline{n}'$ are such that
$(\psi_{\underline{n}}\psi_{\underline{n}'}^{-1})|_{I_K}=\epsilonbar|_{I_K}$,
then in fact $\psi_{\underline{n}}\psi_{\underline{n}'}^{-1}=\epsilonbar$.

Somewhat abusively, we will also
write~$\psi_{\underline{n}}$ for the constant family of
$(\varphi,\Gamma)$-modules $\mathbb{D}(\psi_{\underline{n}})$ over any
$\Fpbar$-scheme~$T$.

 \begin{rem}\label{rem: irreducible representations finite up
  to twist} The irreducible representations
          $\alphabar: G_K \to \GL_d(\Fpbar)$ are easily classified;
          indeed, since the wild inertia subgroup must act trivially,
          they are tamely ramified, so that the restriction
          of~$\alphabar$ to~$I_K$ is diagonalizable. Considering the
          action of Frobenius on tame inertia, it is then easy to see
          that~$\alphabar$ is induced from a character of the
          unramified extension of~$K$ of degree~$d$. It follows that
          each irreducible~$\alphabar$ is absolutely
          irreducible.% \mar{TG: this is false; here and below we should
            % go back to~$\Fpbar$, and then descend the irreducible
            % components at the end, because the Galois action is trivial.}
          
          As we have just seen, there are only finitely many
          such characters up to unramified twist, so
          that %\mar{ME: Recall this}
          in particular there are only finitely many such~$\alphabar$
          up to unramified twist. It follows that there are up to unramified twist
          only finitely many irreducible $(\varphi,\Gamma)$-modules
          over~$\Fpbar$ of any fixed rank. 
          % As explained in Section~\ref{subsec: twists of
  % families},
        \end{rem}


% \begin{df}
% 	\label{def:maximally nonsplit}
% % \mar{ME: Once we replace $\Fbar_p$ by $\F$ here, we should add a lemma proving that maximally
% % non-split is a base-change-invariant notion.}
% We say that a
% representation~$\rhobar:G_K\to\GL_d(\Fpbar)$ is \emph{maximally
%   nonsplit of niveau~$1$} if it has a unique filtration by~$G_K$-stable
% $\Fpbar$-subspaces such that all of the graded pieces are one-dimensional
% representations of~$G_K$. We say that the family~$\rhobar_T$ is
% \emph{maximally nonsplit of niveau~$1$} if~$\rhobar_t$ is maximally
% nonsplit of niveau~$1$ for
% all~$t\in T(\Fpbar)$, and that~$\rhobar_T$ is
% \emph{generically maximally nonsplit of niveau~$1$} if there is a dense open subscheme~$U$
% of~$T$ such that~$\rhobar_U$ is maximally nonsplit of niveau~$1$.
% \end{df}
% The following lemma, which we will employ without comment in the
% sequel, shows that the notion of being maximally nonsplit of
% niveau~$1$ is compatible with extension of scalars.
% \begin{lem}
%   \label{lem: maximally nonsplit extension of scalars}If
%   $\Fnewnew/\Fnew$ is an algebraic extension, then a representation
%   $\rhobar:G_K\to\GL_d(\Fnew)$ is maximally nonsplit of niveau~$1$ if
%   and only if
%   $\rhobar':=\rhobar\otimes_{\Fnew}\Fnewnew:G_K\to\GL_d(\Fnewnew)$ is
%   maximally nonsplit of niveau~$1$.
% \end{lem}\mar{TG: this proof could well be nonsense, please check!}
% \begin{proof}We observed in Remark~\ref{rem: irreducible representations finite up
%   to twist} that the irreducible $\Fnew$-representations of~$G_K$ are
% absolutely irreducible. In particular, we see that~$\rhobar$ admits a
% filtration whose graded pieces are one-dimensional if and only
% if~$\rhobar'$ admits such a filtration. By induction on~$d$, it
% suffices to show that~$\rhobar$ admits a unique one-dimensional
% $G_K$-stable $\Fnew$-subspace if and only if $\rhobar'$ admits a unique one-dimensional
% $G_K$-stable $\Fnewnew$-subspace; this follows from the compatibility
% of the formation of~$H^0(G_K,-)$ with extension of scalars from~$\Fnew$ to~$\Fnewnew$.
% \end{proof}


% We now wish to prove a structure theorem for families which are
% generically maximally nonsplit of niveau~$1$. 


 % In particular if $d=1$, then the set of Serre weights~$\underline{k}$
 % is in bijection with the set of~$\alphabar_i$ of dimension~$1$,
 % because there is a unique~$\alphabar_i:G_K\to\Fpbartimes$ extending
 %  $\prod_{\sigmabar:k\into\Fpbar}\omega_{\sigmabar}^{-k_{\sigmabar,1}}$. Write~$\psi_{\underline{k}}$
 %  for this character.

% \begin{prop}\label{prop: structure of maximally
%     nonsplit}If~$T$ is an integral finite type $\Fpbar$-scheme, and~$\rhobar_T$ is generically maximally
%   nonsplit of niveau~$1$, then there exist:
%   \begin{itemize}% \mar{TG: this needs to be fixed! rather a finite
% %       dominant map? Call the source a ``diagonalizing scheme''? Also,
% %       remark that we can't just fix it by shrinking. TG: perhaps the
% %       counterexample is something like $T=\Gm$ and the family is $
% %       \begin{pmatrix}
% %         0&1\\ x&0
% %       \end{pmatrix}
% % $? Confused about the maximally nonsplit aspect of this though.}\mar{TG: so maybe we can work over $\F$ after all!}
%   \item
%     a nonempty \emph{(}so dense\emph{)} open subscheme~$U$ of~$T$;
% %  \item a decomposition $d=d_1+\dots+d_n$, 
%   \item%  irreducible
%     % representations~$\alphabar_i:G_K\to\Fpbartimes$,
% tuples $\underline{n}_i$,    $1\le i\le d$;
%   \item morphisms~$\lambda_i:U\to\Gm$, $1\le i\le d$;
%   \item morphisms~$x_i:U\to\cX_{i}$, $0\le i\le d$, corresponding to
%     families of representations~$\rho_i$ over~$U$;
%   \end{itemize}
% such that $x_d$ is the restriction of~$\rhobar_T$ to~$U$, and for
% each~$0\le i\le d-1$, we have short exact
% sequences of families over~$U$ \[0\to\rho_i\to\rho_{i+1}\to
%   \ur_{\lambda_{i+1}}\otimes\psi_{\underline{n}_{i+1}}\to 0.\] % \mar{TG: this is pretty
%   % abusive at this point; should we be rewriting this in terms of the
%   % corresponding $(\varphi,\Gamma)$-modules?} 

% % if for each $1\le i\le n$ we write
% % $\delta_i:U\to\cX_{d_i,\red}$ for the morphism given by
% % twisting~$\alphabar_i$ by~$\lambda_i$, then
% In particular, for each~$t\in U$, we
% have \numequation\label{eqn:maximally nonsplit structure}\rhobar_t\cong \begin{pmatrix}
%       \ur_{\lambda_1(t)}\otimes\psi_{\underline{n}_1} &*&\dots &*\\
%       0& \ur_{\lambda_2(t)}\otimes\psi_{\underline{n}_2}&\dots &*\\
%       \vdots&& \ddots &\vdots\\
%       0&\dots&0& \ur_{\lambda_d(t)}\otimes\psi_{\underline{n}_d}\\
%     \end{pmatrix}.  \end{equation}% \mar{TG: probably not the precise statement we want, but
% %   hopefully close. ME: We'll surely have to cut down $T$ to a dense
% %   open subset,right? TG: why? The family is assumed maximally nonsplit
% % rather than generically maximally nonsplit.}
  
% \end{prop}

% \begin{proof}[Proof of Proposition~\ref{prop: structure of maximally
%     nonsplit}]% \mar{TG: should the statement actually be stronger? At
%     % the moment it's only asserting the pointwise existence of the filtration}
% %  \mar{TG: needed} 
%   By definition, we can replace by $T$ by a dense open subscheme,
% and assume that that~$\rhobar_T$ is maximally nonsplit of
%   niveau~$1$ i.e.\ that~$\rhobar_t$ is maximally nonsplit of
%   niveau~$1$ for each~$t\in T(\Fpbar)$. For each $\underline{n}$, we may form the family
%   $$\rho_T \boxtimes (\ur_x^{-1} \otimes \psi_{\underline{n}}^{-1})(1)$$ 
%   over % \mar{TG: here and below we move back and forward between~$\F$,
%     % $\Fpbar$ and sometimes even~$\Fp$. Which do we actually want/need to use?}
%   $T\times_{\Fpbar} \Gm$ (where as above $x$ denotes the variable on 
%   $\Gm$, so that $\Gm := \Spec \Fbar_p[x,x^{-1}]$),
%   and then consider (in the notation of Remark~\ref{rem: writing Herr complex as Galois cohomology})
%   \[H_{\underline{n}} := H^2\bigl(G_K, \rho_T\boxtimes (\ur_x^{-1} \otimes \psi_{\underline{n}}^{-1})(1)\bigr),\]
%   a coherent sheaf on $T\times_{\F} \Gm$.% \mar{TG: be sure that $\otimes$
%     % and $\boxtimes$ both explained above.}

%   Since the formation of $H^2(G_K,\text{--})$ is compatible
%   with  arbitrary finite type base-change, we see from Tate local duality that for each finite type point
%   $(t,x_0) \in T\times_{\F} \Gm$, the corresponding fibre $(H_{\underline{n}})_{(t,x_0)}$
%  % that $H_t$, the coherent sheaf on $\Gm$
%  % obtained by pulling back $H$ via the closed embedding
%  % $\Gm = t\times \Gm \hookrightarrow T\times_{\F} \Gm$,
%   admits the description
%   $$(H_{\underline{n}})_{(t,x_0)} \cong 
%   H^2\bigl(G_K, \rho_t\otimes (\ur_{x_0}^{-1} \otimes \psi_{\underline{n}}^{-1})(1)\bigr)
% 	  \cong \Hom_{G_K}(\rho_t, \ur_{x_0}\otimes \psi_{\underline{n}}).
% 	  $$ 
% 	  The assumption that $\rho_t$ is maximally non-split
% 	  ensures that for each choice of $t$,
% 	  this fibre is non-zero 
% 	  for exactly one choice of $\underline{n}$ and one choice of $x_0$,
% 	  and is then precisely one-dimensional.

% 	  Thus, if $\Supp(H_{\underline{n}}) \hookrightarrow T\times_{\Fpbar} \Gm$ denotes
% 	 the closed subscheme on which $H_{\underline{n}}$ is supported,
% 	 then the morphism 
% 	  $\coprod_{\underline{n}} \Supp(H_{\underline{n}}) \to T$
% 	  given by the various projections onto $T$ induces a bijection
% 	  on closed points.

% 	  Since a quasi-finite morphism 
% 	  becomes finite after restricting to a dense open subset of the
% 	  target~\cite[\href{http://stacks.math.columbia.edu/tag/03I1}{Tag 03I1}]{stacks-project},
% 	  % \mar{ME: This citation is overkill in the context of morphisms
% 	  %         b/w finite type schemes over a field, but I guess
% 	  %         that's fine.}
% 	  and since $T$ is irreducible, we see that
% 	  if we replace $T$ by a suitably chosen non-empty open subset,
% 	  then we may assume that $\Supp(H_{\underline{n}}) \to T$ is finite and 
% 	  bijective on closed points
% for one choice of $\underline{n}$, and that $H_{\underline{n}} = 0$ for
% 	  all the other possible choices of $\underline{n}$.   We make
% 	  such a replacement, and now write $\psi_{\underline{n}_d}$ rather than
% 	  $\psi_{\underline{n}}$ and $H$ rather than $H_{\underline{n}}$.

% 	  The morphism $\Supp(H) \to T$ induces a morphism $\Supp(H)_{\red}
% 	  \to T$,
% 	  which is then a bijection between reduced finite type $\Fpbar$-schemes,
% 	  and thus restricts to an isomorphism over a dense open subset
% 	  of the target.  Thus, replacing $T$ by this open subset,
% 	  we may assume that $\Supp(H)_{\red} \to T$ is an isomorphism,
% 	  and thus identify $\Supp(H)_{\red}$ with $T$.
% 	  Let $\lambda: T \to \Gm$ denote the composite of this
% 	  identification with the projection $\Supp(H)_{\red} \to \Gm$.
% 	  Again using the fact that the formation of $H^2$ commutes
% 	  with base-change, we find that the pull-back $H_T$
% 	  of $H$ to $T \cong \Supp(H)_{\red}$ may be identified with
% 	  $$H^2\bigl( \rho_T \boxtimes (\ur_{\lambda}^{-1}\otimes \psi_{\underline{n}_d}^{-1})(1)\bigr).$$
% 	  The computation already made of the fibres of $H$ shows that
% 	  the fibres of $H_T$ are all one-dimensional.  Since $T$
% 	  is reduced, we conclude that $H_T$ is locally free of rank one;
% 	  replacing $T$ by a non-empty open subset once more, we may
% 	  in fact assume that $H_T$ is free of rank one, and
% 	  so choose a nowhere zero section of~ $H_T$, which by Lemma~\ref{lem: Tate duality for H2 to H0 in families} we 
% 	  may and do regard as a surjection % \mar{TG: should we say
%             % anything for why nowhere zero implies this is a surjection?}
% 	  % \mar{ME: Have to explain this, perhaps in an earlier discussion
% 	  %         of $H^2$ and its interpretation, rather than right here.}
% 	  $$ \rho_T \onto \ur_{\lambda}\otimes
%           \psi_{\underline{n}_d}.$$ The kernel of this surjection is a
%           rank~$(d-1)$ family of projective \'etale
%           $(\varphi,\Gamma)$-modules, so the result follows by
%           induction on~$d$.
% %\mar{TG: now finish by induction, I guess}
% \end{proof}
% \mar{TG: at the end, we will add a version of this
%   for~$(\cX_{d,\red})_\Fpbar$ rather than for a scheme, and then the
%   eigenvalue morphism will make sense without having a particular
%   affine chart in mind.}
% We now deduce a more intrinsic version of the previous result, without
% reference to a parameterising scheme~$T$.

It is convenient to extend the notation~$\rho_T$ to algebraic stacks.
To this end, if~$\cT$ is a reduced algebraic stack of finite type
over~$\Fpbar$, we will denote
% by~$\rhobar_{\cT}$ 
a morphism
$\cT\to(\cX_{d,\red})_{\Fpbar}$ 
by~$\rhobar_{\cT}$, and will sometimes abusively refer
to $\rho_{\cT}$ as a family of $G_K$-representations parameterized by~$\cT$.

\begin{df}
	\label{def:maximally nonsplit}We say that a
representation~$\rhobar:G_K\to\GL_d(\Fpbar)$ is \emph{maximally
  nonsplit of niveau~$1$} \index{maximally
  nonsplit of niveau~$1$} if it has a unique filtration by~$G_K$-stable
$\Fpbar$-subspaces such that all of the graded pieces are one-dimensional
representations of~$G_K$. We say that the family~$\rhobar_{\cT}$ is
\emph{maximally nonsplit of niveau~$1$} if~$\rhobar_t$ is maximally
nonsplit of niveau~$1$ for
all~$t\in T(\Fpbar)$, and that~$\rhobar_T$ is
\emph{generically maximally nonsplit of niveau~$1$} if there is a
dense open substack~$\cU$ \index{generically maximally nonsplit of niveau~$1$} 
of~$\cT$ such that~$\rhobar_{\cU}$ is maximally nonsplit of niveau~$1$.

 In particular, we say that an algebraic substack~$\cT$ of~$(\cX_{d,\red})_\Fpbar$ is
\emph{maximally nonsplit of niveau~$1$} if~$\rhobar_t$ is maximally
nonsplit of niveau~$1$ for
all~$t\in \cT(\Fpbar)$, and that~$\cT$ is
\emph{generically maximally nonsplit of niveau~$1$} if there is a dense open substack~$\cU$
of~$\cT$ such that~$\cU$ is maximally nonsplit of niveau~$1$.
\end{df}
\begin{rem}
  \label{rem: everything is maximally nonsplit niveau 1}We will
  eventually see in Theorem~\ref{thm:reduced dimension} below that 
~$(\cX_{d,\red})_{\Fpbar}$ itself is an algebraic stack of finite type
over~$\Fpbar$, and is generically maximally nonsplit of
  niveau 1.
\end{rem}

Our next goal is to prove a structure theorem (Proposition~\ref{prop: structure of maximally
    nonsplit}) for families which are
generically maximally nonsplit of niveau~$1$. We begin with the
following technical lemma.


\begin{lemma}
\label{lem:Fitting fibres}
Let  $T$ be a reduced scheme of finite type over an algebraically closed field~$k$,
let $M$ be a coherent sheaf on $T\times_k \Gm$, and for each $t \in T(k)$,
let $M_t$ denote the restriction of $M$ to
$\Gm \iso t\times_k \Gm \hookrightarrow T \times_k \Gm$.
Suppose that for each point $t \in T(k)$, $M_t$ is a length one skyscraper sheaf supported at a  single point
of~$\Gm$.  Then there is a dense open subscheme $U$ of
%\mar{TG: here and in the proof of Prop \ref{prop: structure of maximally
%    nonsplit}, better to write ``subscheme'' than ``subset''? ME: Yes, good idea, fixed.}
$T$ such that, over  $U$, the composite $\Supp(M) \hookrightarrow   T\times_k \Gm
\to  T$  pulls back to  an isomorphism,
while the pullback of $M$ over $U\times_k \Gm$ is locally free of rank one over
its support.
\end{lemma}
\begin{proof}
Let $\Fitt(M)\subseteq \cO_{T\times \Gm}$  denote the Fitting ideal sheaf of $M$,
and write $Z := \Spec \cO_{T\times \Gm}/\Fitt(M);$ recall that  $\Supp(M)$ 
is a closed subscheme of~$Z$, supported on  the same underlying  closed subset
of~$T\times_k \Gm.$
For each $t  \in T(k)$, we find that the fibre $Z_t$  of $Z$  over $t$
is equal to
$\Spec  \cO_{\Gm}/\Fitt(M_t)$ (recall that the formation
of Fitting  ideals is compatible  with  base-change), which, by assumption,
is a single reduced point of $\Gm$.
Thus the composite $Z \to T\times_k \Gm \to T$  is a  morphism with
reduced singleton fibres,
and hence we may find a dense open subscheme $U$ of
$T$ such that this morphism restricts to an isomorphism over~$U$.
In particular, the pull-back of $Z$ over $U$ is reduced (since $U$ is),
and thus coincides with the pull-back of $\Supp (M)$  over $U$.
This shows that the morphism $\Supp (M) \to T$  pulls back to an isomorphism
over~$U$.

If we use the  inverse of the isomorphism just constructed
to identify $U$ with an open subscheme of $\Supp(M)$,
then our assumption on the nature of $M_t$  for  $t  \in T(k)$ implies
that the fibre of $M$ over each $k$-point of $U$ is one-dimensional.  Since
$U$  is reduced, we find that the restriction of $M$ to $U$ is  furthermore locally
free of rank one, as claimed.
%\mar{ME: Add in argument to get local freeness statement.}
\end{proof}

\begin{prop}\label{prop: structure of maximally
    nonsplit}If~$\cT$ is a reduced finite type algebraic stack over~$\Fpbar$, and~$\rhobar_{\cT}$ is generically maximally
  nonsplit of niveau~$1$, then there exist:
  \begin{itemize}
  \item
    a dense %nonempty \emph{(}so dense\emph{)}
open substack~$\cU$ of~$\cT$;
%  \item a decomposition $d=d_1+\dots+d_n$, 
  \item%  irreducible
    % representations~$\alphabar_i:G_K\to\Fpbartimes$,
tuples $\underline{n}_i$,    $1\le i\le d$;
  \item morphisms~$\lambda_i:\cU\to\Gm$, $1\le i\le d$;
  \item morphisms~$x_i:\cU\to(\cX_{i,\red})_{\Fpbar}$, $0\le i\le d$, corresponding to
    families of $i$-dimensional representations~$\rho_i$ over~$\cU$;
  \end{itemize}
such that $x_d$ is the restriction of~$\rhobar_{\cT}$ to~$\cU$, and for
each~$0\le i\le d-1$, we have short exact
sequences of families over~$\cU$ \[0\to\rho_i\to\rho_{i+1}\to
  \ur_{\lambda_{i+1}}\otimes\psi_{\underline{n}_{i+1}}\to 0.\] 
In particular, for each~$t\in \cU$, we
have \numequation\label{eqn:maximally nonsplit structure}\rhobar_t\cong \begin{pmatrix}
      \ur_{\lambda_1(t)}\otimes\psi_{\underline{n}_1} &*&\dots &*\\
      0& \ur_{\lambda_2(t)}\otimes\psi_{\underline{n}_2}&\dots &*\\
      \vdots&& \ddots &\vdots\\
      0&\dots&0& \ur_{\lambda_d(t)}\otimes\psi_{\underline{n}_d}\\
    \end{pmatrix}.  \end{equation}%
  
\end{prop}
\begin{rem}
  \label{rem: maximally nonsplit doesn't have a universal family}% Even
  % if~$\rhobar_T$ is maximally nonsplit (not just generically maximally
  % nonsplit),
  We cannot necessarily attain~(\ref{eqn:maximally nonsplit
    structure}) over all of~$\cT$. For example, there are families of two-dimensional representations of~$G_K$ which
  are generically maximally nonsplit of niveau one, but which also specialise to irreducible
  representations. (The existence of such families follows
  from Theorem~\ref{thm:reduced dimension} below.)% \mar{TG: actually I
% don't think it does; rather, we need to cite a theorem yet to be
% written at the end, which shows equidimensionality, so that we know
% that irreducible representations show up in the irreducible components
% we're writing down.}
    % \mar{TG: citation of either our final results, or CEGS, or both.}
\end{rem}

\begin{proof}[Proof of Proposition~\ref{prop: structure of maximally
    nonsplit}]% \mar{TG: should the statement actually be stronger? At
    % the moment it's only asserting the pointwise existence of the filtration}
%  \mar{TG: needed} 
We begin by assuming that~$\cT$ is a scheme, say~$\cT=T$.  %By definition,
Without loss of generality, we can replace by $T$ by a dense open subscheme,
and thus assume that~$\rhobar_{T}$ is maximally nonsplit of
  niveau~$1$, i.e.\ that~$\rhobar_t$ is maximally nonsplit of
  niveau~$1$ for each~$t\in T(\Fpbar)$, and also assume
that $T$ is a disjoint union of finitely many integral open and closed subschemes.  Replacing
$T$ by each of these subschemes in turn, we may assume
that $T$ itself is integral. 

For each $\underline{n}$, we may form the family
  $$\rho_T \boxtimes (\ur_x^{-1} \otimes \psi_{\underline{n}}^{-1})(1)$$ 
  over % \mar{TG: here and below we move back and forward between~$\F$,
    % $\Fpbar$ and sometimes even~$\uFp$. Which do we actually want/need to use?}
  $T\times_{\Fpbar} \Gm$ (where as above $x$ denotes the variable on 
  $\Gm$, so that $\Gm := \Spec \Fbar_p[x,x^{-1}]$),
  and then consider (in the notation of Remark~\ref{rem: writing Herr complex as Galois cohomology})
  \[H_{\underline{n}} := H^2\bigl(G_K, \rho_T\boxtimes (\ur_x^{-1} \otimes \psi_{\underline{n}}^{-1})(1)\bigr),\]
  a coherent sheaf on $T\times_{\Fbar_p} \Gm$. % \mar{TG: be sure that $\otimes$
    % and $\boxtimes$ both explained above.}
  Since the formation of $H^2(G_K,\text{--})$ is compatible
  with  arbitrary finite type base-change, we see that  for any point
$ t \in T(\Fbar_p)$, the pull-back of
$H_{\underline{n}}$ over $t \times_{\Fbar_p} \Gm \cong \Gm$
admits the description
\[(H_{\underline{n}})_t\cong H^2\bigl(G_K, \rho_t\otimes (\ur_{x}^{-1} \otimes \psi_{\underline{n}}^{-1})(1)\bigr),\]
and that the fibre $(H_{\underline{n}})_{(t,x_0)}$ of this pull-back
at a point $x_0 \in \Gm(\Fbar_p)$,
admits the description
  $$(H_{\underline{n}})_{(t,x_0)} \cong 
  H^2\bigl(G_K, \rho_t\otimes (\ur_{x_0}^{-1} \otimes \psi_{\underline{n}}^{-1})(1)\bigr)
	  \cong \Hom_{G_K}(\rho_t, \ur_{x_0}\otimes \psi_{\underline{n}})
	  $$ 
(the second isomorphism following
from Tate local duality).
%that for each finite type point
%  $(t,x_0) \in T\times_{\F} \Gm$, the corresponding fibre $(H_{\underline{n}})_{(t,x_0)}$
% % that $H_t$, the coherent sheaf on $\Gm$
% % obtained by pulling back $H$ via the closed embedding
% % $\Gm = t\times \Gm \hookrightarrow T\times_{\F} \Gm$,
%  admits the description
%  $$(H_{\underline{n}})_{(t,x_0)} \cong 
%  H^2\bigl(G_K, \rho_t\otimes (\ur_{x_0}^{-1} \otimes \psi_{\underline{n}}^{-1})(1)\bigr)
%	  \cong \Hom_{G_K}(\rho_t, \ur_{x_0}\otimes \psi_{\underline{n}}).
%	  $$ 
	  The assumption that $\rho_t$ is maximally non-split
%	 ensures that for each choice of $t$,
	  shows that this fibre is non-zero 
	  for exactly one choice of $\underline{n}$ and one choice of $x_0$,
	  and is then precisely one-dimensional.
In particular, we see, for this distinguished choice 
of~$\underline{n}$, that $(H_{\underline{n}})_t$ is set-theoretically supported at~$x_0$.
%(and that $(H_{\underline{n}})_t$ vanishes for all other  choices of~$\underline{n}$).
%\mar{ME: Need a slightly different/stronger statement to apply Lemma~\ref{lem:Fitting fibres} below.}
          Now let $\xi_0:\Spec  \Fbar_p[\epsilon]/(\epsilon^2) \hookrightarrow
\Gm$ be the (unique up to scalar) non-trivial tangent vector to $x_0$, 
and let
$(H_{\underline{n}})_{(t,\xi_0)}$
denote the pull-back of $(H_{\underline{n}})_t$
over $\xi_0$.  A similar computation using Tate local duality and the maximal non-splitness
of $\rho_t$ shows that 
$(H_{\underline{n}})_{(t,\xi_0)}$
is also one-dimensional; thus $(H_{\underline{n}})_t$  is in fact a  skyscraper sheaf of
length one supported at~$x_0$.
For all other choices of~$\underline{n}$,
we see that $(H_{\underline{n}})_t$ vanishes (since all its fibres do).

          Applying Lemma~\ref{lem:Fitting fibres} to
          $M  := \bigoplus_{\underline{n}} H_{\underline{n}},$ 
          and replacing $T$ by an appropriately chosen dense open subscheme,
          we find that $\Supp(M)$ maps isomorphically to~$T$,
          and that $M$ is locally free of rank one over its support.
          Taking into account the definition of $M$ as a direct sum,
          and the irreducibility of~$T$,
          we find that in fact $M = H_{\underline{n}}$  for exactly one choice
          $\underline{n}_d$ of~$\underline{n}$,
          and that $H_{\underline{n}} = 0$ for all other
          possible choices. 
%
%	  Thus, if $\Supp(H_{\underline{n}}) \hookrightarrow T\times_{\Fpbar} \Gm$ denotes
%	 the closed subscheme on which $H_{\underline{n}}$ is supported,
%	 then the morphism 
%	  $\coprod_{\underline{n}} \Supp(H_{\underline{n}}) \to T$
%	  given by the various projections onto $T$ induces a bijection
%	  on closed points.
%
%	  Since a quasi-finite morphism 
%	  becomes finite after restricting to a dense open subset of the
%	  target~\cite[\href{http://stacks.math.columbia.edu/tag/03I1}{Tag 03I1}]{stacks-project},
%	  % \mar{ME: This citation is overkill in the context of morphisms
%	  %         b/w finite type schemes over a field, but I guess
%	  %         that's fine.}
%	  and since $T$ is irreducible, we see that
%	  if we replace $T$ by a suitably chosen non-empty open subset,
%	  then we may assume that $\Supp(H_{\underline{n}}) \to T$ is finite and 
%	  bijective on closed points
%for one choice of $\underline{n}$, and that $H_{\underline{n}} = 0$ for
%	  all the other possible choices of $\underline{n}$.   We make
%	  such a replacement, and now write $\psi_{\underline{n}_d}$ rather than
%	  $\psi_{\underline{n}}$ and $H$ rather than $H_{\underline{n}}$.
%\mar{TG: I think Yiwen pointed out that we need to be a bit more
%  careful about inseparable extensions? Idea is: the map is radicial,
%  use uniqueness of~$x_0$ (the unramified twist in~$\Gm$). The
%  intuitive idea is that if~$t$ is now the generic point of~$T$ there
%  should still be a unique~$x_0$ (and not some $p$th root phenomenon)\dots}
%	  The morphism $\Supp(H) \to T$ induces a morphism $\Supp(H)_{\red}
%	  \to T$,
%	  which is then a bijection between reduced finite type $\Fpbar$-schemes,
%	  and thus restricts to an isomorphism over a dense open subset
%	  of the target.  Thus, replacing $T$ by this open subset,
%	  we may assume that $\Supp(H)_{\red} \to T$ is an isomorphism,
%	  and thus identify $\Supp(H)_{\red}$ with $T$.
	  Let $\lambda: T \to \Gm$ denote the composite of the inverse
          isomorphism $T \iso \Supp(M)$
	  with the projection $\Supp(M) \to \Gm$.
	  Again using the fact that the formation of $H^2$ commutes
	  with base-change, we find that the pull-back $M_T$
	  of $M$ to $T \cong \Supp(M)$ may be identified with
	  $$H^2\bigl( \rho_T \boxtimes (\ur_{\lambda}^{-1}\otimes \psi_{\underline{n}_d}^{-1})(1)\bigr).$$
%	  The computation already made of the fibres of $H$ shows that
%	  the fibres of $H_T$ are all one-dimensional.  Since $T$
%	  is reduced, we conclude that
          We have already remarked that $M_T$ is locally free of rank one;
	  replacing $T$ by a non-empty open subscheme once more, we may
	  in fact assume that $M_T$ is free of rank one, and
	  so choose a nowhere zero section of~ $M_T$, which by Lemma~\ref{lem: Tate duality for H2 to H0 in families} we 
	  may and do regard as a surjection % \mar{TG: should we say
            % anything for why nowhere zero implies this is a surjection?}
	  % \mar{ME: Have to explain this, perhaps in an earlier discussion
	  %         of $H^2$ and its interpretation, rather than right here.}
	  $$ \rho_T \onto \ur_{\lambda}\otimes
          \psi_{\underline{n}_d}.$$ The kernel of this surjection is a
          rank~$(d-1)$ family of projective \'etale
          $(\varphi,\Gamma)$-modules, so the result for~$T$ follows by
          induction on~$d$.

%\mar{ME: Will add stuff here}
          We now return to the general case that~$\cT$ is a reduced
          finite type algebraic stack over~$\Fpbar$. Since~$\cX_d$ is
          Ind-algebraic by Proposition~\ref{prop:X is an Ind-stack},
          we can form the scheme-theoretic image~$\cT'$ of the
          morphism $\cT\to(\cX_{d,\red})_{\Fpbar}$, which is again a reduced finite
          type algebraic stack over~$\Fpbar$. The
          family~$\rhobar_{\cT'}$ is again generically maximally
          nonsplit of niveau~$1$ (since by the stacky version of
          Chevalley's theorem (see~\cite[App.\ D]{MR2818725}),
          the image of $\cT$ in $\cT'$ is constructible, and so contains
          a dense open subset of $\cT'$), and since the morphism $\cT\to\cT'$
          is scheme-theoretically dominant, any dense open substack
          of~$\cT'$ pulls back to a dense open substack of~$\cT$. It
          therefore suffices to prove the result
          for~$\cT'$. Replacing~$\cT$ by~$\cT'$, then, we can and do
          assume that~$\cT$ is a closed substack of~$(\cX_{d,\red})_{\Fpbar}$.

          Let $T\to\cT$ be a smooth cover by a (necessarily) reduced
          and finite type $\Fpbar$-scheme (we can ensure that $T$ 
is finite type over $\Fbar_p$, since $\cT$ is so), and let~$\rhobar_T$ be the
          induced family. We wish to deduce the result for~$\cT$ from
          the result for~$T$, which we have already proved.  We thus
replace $T$ by a dense open subscheme (and correspondingly replace $\cT$ by
a dense open substack, namely the image of this open subscheme of $T$),
          so that the family $\rhobar_T$ is maximally non-split,
and so that the morphisms $\lambda_i:T\to\Gm$
          and ~$x_i:T\to(\cX_{i,\red})_{\Fpbar}$ exist.  We will then
show that these morphisms factor through the smooth
          surjection~$T\to\cT$. This amounts to showing that we can
          factor the morphism  of groupoids
\numequation
\label{eqn:groupoid to product}
T\times_{\cT}T\to T\times T
\end{equation}
through a morphism of groupoids $T\times_{\cT} T \to T\times_{\Gm}T$ or
$T\times_{\cT} T \to T\times_{(\cX_{i,\red})_{\Fpbar}}T$ respectively.
%\mar{TG: reword?  Those are groupoids rather than morphisms thereof\dots}

Since $\cT$ is a substack of $\cX$, the natural morphism induces an isomorphism
$T\times_{\cT} T \iso T\times_{\cX} T$.  Thus for any test scheme $T'$,
a $T'$-valued point of $T\times_{\cT} T$ consists of a pair of morphisms
$f_0,f_1: T' \rightrightarrows T$, and an isomorphism of families
\numequation
\label{eqn:family isomorphism}
f_0^*\rho_{T} \iso f_1^*\rho_T
\end{equation}
(where $f_i^* \rho_T$ denotes the pull-back of $\rho_T$ 
to $T'$ via~$f_i$).    


In the first case, in which we consider one of the morphisms $\lambda_i$,
we need to show (employing the notation just introduced)
 that the morphism $f_0 \times f_1: T'\times T' \to T\times T$ necessarily factors through
$T \times_{\G_m} T$ (this latter fibre product formed using the morphisms $\lambda_i$).
More concretely, this amounts to checking that $\lambda_i \circ f_0 = \lambda_i \circ f_1$.
Since the source of~\eqref{eqn:groupoid to product} is reduced and of finite type
over $\Fbar_p$, it suffices to check this for reduced test schemes
$T'$ of finite type over~$\Fbar_p$, and thus to check that these morphisms coincide
at the $\Fbar_p$-valued points $t'$ of $T'$.   This follows from
the assumed isomorphism~\eqref{eqn:family isomorphism}, and the uniqueness 
of the filtration at each $\Fbar_p$-valued point of the family. 



%since $T\times_{\cT}T$ is a reduced algebraic space of
%          finite type over~$\Spec\Fpbar$, while $T\times_{\Gm}T$ is a
%          subscheme of~$T\times T$, we are reduced to checking this
%          factorisation on finite type points, where it follows
%          immediately from the definition of a maximally nonsplit
%          representation (i.e.\ from the uniqueness of the
%          filtration).

          In the second case,
in which we consider one of the morphisms $x_i$,
the morphism
          $T\times_{(\cX_{i,\red})_{\Fpbar}}T\to T\times T$ is no longer a monomorphism,
and so we have to actually construct a morphism of groupoids
          $T\times_{\cT}T\to T\times_{(\cX_{i,\red})_{\Fpbar}}T$
(rather than simply verify a factorization).
%
%
%
%By our assumption
%          that~$\cT$ is a substack of~$(\cX_{d,\red})_{\Fpbar}$, we have
%          $T\times_{\cT}T=T\times_{(\cX_{d,\red})_{\Fpbar}}T$, so that a morphism
%          $S\to T\times_{\cT}T$ is the data of two families of
%          representations~$\rhobar_S$, $\rhobar_S'$ together with an
%          isomorphism $\rhobar_S\cong\rhobar_S'$.%  \mar{TG: wtf? only if it's a substack of
%            % $\cX_d$?!}
%
So, given $f_0, f_1: T' \to T,$ and an isomorphism~\eqref{eqn:family isomorphism},
we have to construct a corresponding isomorphism
\numequation
\label{eqn:desired isomorphism}
f_0^*(\rho_T)_i \iso f_1^*(\rho_T)_i,
\end{equation}
where $(\rho_T)_i$ denotes the family of $i$-dimensional subrepresentations
of $\rho_T$ corresponding to the morphism~$x_i$.
We will define the desired isomorphism~\eqref{eqn:desired isomorphism}
simply to be the restriction of
the isomorphism~\eqref{eqn:family isomorphism}.
For this definition to make sense, we have to show that~\eqref{eqn:family isomorphism}
does in fact restrict to an isomorphism~\eqref{eqn:desired isomorphism}.
In fact, it suffices to check that~\eqref{eqn:family isomorphism}
induces an embedding 
$f_0^*(\rho_T)_i \hookrightarrow f_1^*(\rho_T)_i;$
reversing the roles of $f_0$ and $f_1$ then allows us to promote this
to an isomorphism.

Again, it suffices to check this on reduced test schemes $T'$ of finite type
over~$\Fbar_p$, and then it suffices to show that the composite
$$f_0^*(\rho_T)_i \hookrightarrow f_0^*\rho_T \buildrel
\text{\eqref{eqn:family isomorphism}} \over \longrightarrow
f_1^*\rho_T \to f_1^*\rho_T/ f_1^*(\rho_T)_i$$ 
vanishes.  This can be checked on $\Fbar_p$-valued points,
where it again follows from the uniqueness of the filtrations that define
the~$(\rho_T)_i$.
%
%eGiven such an isomorphism $\rhobar_S\cong\rhobar_S'$ of families of
%r epresentations, to construct the composite morphism
%$S\to T\times_{(\cX_{d,\red})_{\Fpbar}}T\to T\times_{(\cX_{i,\red})_{\Fpbar}}T$   we need to show that we have an
%induced isomorphism~$\rho_i\cong\rho'_i$ of the corresponding
%subrepresentations. This amounts to the statement that the composite
%morphism of $(\varphi,\Gamma)$-modules $\rho_i\to\rhobar_S\cong\rhobar_S'\to\rhobar_S'/\rho'_i$ vanishes. We can assume that~$S$ is reduced (and finite
 %type over~$\Fpbar$), so this vanishing can be checked on finite type
%points, where it is again immediate from the uniqueness of the
%filtration; it follows in the same way that the morphism
%$T\times_{\cT}T\to T\times T$ follows through this morphism, as required.
\end{proof}
%\mar{TG: hopefully basically OK, see what you think.}

          % it suffices to check these properties on finite type points,
          % where they follow immediately from the definition of a
          % maximally nonsplit representation (i.e.\ from the uniqueness
          % of the filtration).
% \mar{TG: this needs to be done, and then rearrange to
%   remove the earlier version, and fix the labels.}
%\mar{TG: now finish by induction, I guess}



\section{Dimensions of families of extensions}\label{subsec:
  dimensions of families of extensions}
Suppose that we have a morphism $T \to (\cX_{d,\red})_{\Fpbar},$
with $T$ a reduced finite type $\Fpbar$-scheme,
% \mar{ME: Ultimately we will want $T$ to be an alg.\ stack, not just
% 	a scheme, but I think it might be easier to do the scheme case first.
% ME: Actually, maybe the scheme case is fine; the stacky case takes care
% of itself just by passing to a smooth cover by a scheme.}
corresponding to a family $\rhobar_{T}$ of $G_K$-representations over 
$T$.  
Fix a representation $\alphabar: G_K \to \GL_a(\Fpbar)$.

\begin{lemma}\label{lem: Ext2 locally free} If $\Ext^2_{G_K}(\alphabar,\rhobar_t)$ is of constant rank $r$ 
	for all $t \in T(\Fbar_p)$, then $H^2(G_K,\rhobar_T\otimes \alphabar^{\vee})$ is
	locally free of rank $r$ as an $\cO_T$-module.
\end{lemma}
\begin{proof}
	Since the formation of $H^2$
	is compatible with arbitrary finite type base-change (by Theorem~\ref{thm: Herr complex is perfect}~(2)),
	we see that $H^2(G_K,\rhobar_T\otimes\alphabar^{\vee})$
	is a coherent sheaf on $T$ of constant
	fibre rank $r$.  Since $T$ is reduced, the lemma follows.
\end{proof}


By Corollary~\ref{cor:Herr complex length},
%By Theorem~\ref{thm: Herr complex is perfect}
if $T$ is affine, 
then we can and do choose a good complex (that is, a bounded
complex of finite rank locally free $\cO_T$-modules) 
% \mar{ME: Should we explain how we know we can find a complex in these
% specific degrees? ME: Added an explanation. (Caveat: as written,
% it is only valid  in the affine case, so maybe  restrict to that
% case?) TG: will this ever cause a problem when we do things like
% passing to an open subscheme? I guess it's OK if we've fixed this once
% and for all but we might need to be a bit careful to check we never do
% something stupid? ME: I think I dealt with this in a minimally
% invasive manner. TG: looks good, commenting out.} 
$$C^0_T \to C^1_T \to C^2_T$$ 
% \mar{ME: I'm not sure if perfect complex is the correct word; I just mean
% 	a complex of finite rank locally free sheaves. TG: introduced
%         terminology of ``good'', although only for rings so far, I guess.}
computing $H^{\bullet}(G_K,\rhobar_T\otimes\alphabar^{\vee})$. Suppose that we are in the context of the preceding lemma,
i.e.\ that $\Ext^2_{G_K}(\alphabar,\rhobar_T)$ is locally free of some rank $r$. It
follows that the truncated complex
$$C^0_T \to Z^1_T$$ is again good (here $Z^1_T:=\ker(C^1_T\to
C^2_T)$). As in Remark~\ref{rem: writing Herr complex as Galois
  cohomology}, we write \[H^1(G_K,\rhobar_T\otimes\alphabar^\vee)\] for
the cohomology group of this complex in degree~$1$; its formation is
compatible with arbitrary finite type base-change, and its specialisations to
finite type points of~$T$ agree with the usual Galois cohomology.

In particular, if we choose another integer $n$,
we may use the surjection $Z^1_T \to
H^1(G_K,\rhobar_T\otimes\alphabar^{\vee})$ (together with Lemma~\ref{lem: Yoneda for H0 H1 of Herr complex})
to construct a universal family of extensions % \mar{TG: perhaps we should explain
  % this slightly more, like we do in \cite{CEGSKisinwithdd}?}
\numequation\label{eqn: universal family of extensions}0 \to \rhobar_T \to \cE \to \alphabar^{\oplus n} \to 0\end{equation}
parameterized by the vector bundle $V$ corresponding to the finite rank
locally free sheaf $(Z_T^1)^{\oplus n}$,
giving rise to a morphism
\numequation
\label{eqn:induced morphism}
V \to (\cX_{d+ a n,\red})_{\Fpbar}.
\end{equation}(See~\cite[\S 4.2]{CEGSKisinwithdd} for a similar construction.)
 % \mar{TG: I guess we need to explain why $H^1$ is
  % classifying extensions of $(\varphi,\Gamma)$-modules?}
%\mar{TG: should this and the next instance be $\cX_{d+n}$? ME: Yes, thanks, now fixed}

Write $G_{\alphabar} := \Aut_{G_K}(\alphabar),$
thought of as an affine algebraic group over $\Fpbar$.



% \mar{ME: Maybe better to work with $\cX_{d,\red,\Fbar_p}$ here, to avoid issues
% later on about irreducibility vs.\ abs.\ irreducibility?  Should just be a matter
% of changing notation, but not changing any arguments. TG: done,
% commenting out.}
\begin{prop}Maintain the notation and assumptions above; so we assume in particular
  that $\Ext^2_{G_K}(\alphabar,\rhobar_t)$ is of constant rank $r$ 
	for all $t \in T(\Fbar_p)$.
	\label{prop:dimension control}\leavevmode
        \begin{enumerate}
        \item If $e$ denotes the dimension of the scheme-theoretic
          image % \mar{TG: should probably have an explanation that this
          % exists, and is an algebraic stack. ME: I rewrote the
          % discussion to emphasize that the morphisms in question
          % factor through underlying reduced algebraic substacks,
          % which makes this clearer.}
          of $T$ in $(\cX_{d,\red})_{\Fpbar},$ then the dimension of the
          scheme-theoretic image,
          % {\em (}
          with respect to the morphism~{\em (\ref{eqn:induced
              morphism})}, %)}
          of $V$ in $(\cX_{d+an,\red})_{\Fpbar}$ is bounded above by
          \[e + n([K:\Q_p] ad +r) -n^2 \dim G_{\alphabar}.\]

       
        \item  Suppose that:
          \begin{enumerate}
          \item $n=1$;
          \item  $\alphabar$ is one-dimensional;
            
          \item $\rhobar_T$ is generically maximally nonsplit of
            niveau~$1$;
          \item\label{item: weird condition over Qp} if~$K=\Qp$, then
after replacing $T$ by a dense open subscheme,
so that the tuples $\underline{n}_i$ exist,
and so that the morphisms $\lambda_i$ and $x_i$ as in
%in the notation of
the statement of Proposition~{\em \ref{prop: structure of maximally nonsplit}}
are defined on~$T$, 
%	  we further assume that after possibly shrinking~$U$, for all~$t\in U(\Fbar_p)$
at least one of the following conditions holds at each~$t~\in~T(\Fbar_p)$:
  \begin{enumerate}
  \item
    $\ur_{\lambda_{d-1}(t)}\otimes\psi_{\underline{n}_{d-1}}\ne\alphabar(1)$.
  \item
    $\ur_{\lambda_{d}(t)}\otimes\psi_{\underline{n}_{d}}=\alphabar$.
  \item $\ur_{\lambda_{d}(t)}\otimes\psi_{\underline{n}_{d}}=\alphabar(1)$.
  \end{enumerate}
% \numequation\label{item: weird condition over Qp}\ur_{\lambda_{d-1}(t)}\otimes\psi_{\underline{n}_{d-1}}\ne\alphabar(1)\end{equation}
        \end{enumerate}
Then equality
          holds in the inequality of~(1). Furthermore the family of
          extensions~$\cE_V$ corresponding to~$V$ is generically
          maximally nonsplit of niveau~$1$.
%\mar{TG: not sure if we can prove this, but will have a go.}
\item Suppose that we are in the setting of~(2), and that 
  after replacing $T$ by a dense open subscheme, the following conditions hold:
\begin{enumerate}
\item
For each
  $t~\in~T(\Fbar_p)$ the representation $\rhobar_t$ is maximally nonsplit,
and furthermore the unique quotient character of~$\rhobar_t$
  {\em(}which exists by the assumption that~$\rhobar_T$ is generically
  maximally nonsplit of niveau~$1${\em)} is equal
  to~$\alphabar(1)$.
\item
Furthermore, if we write $$0 \to \rbar_t \to
  \rhobar_t \to \alphabar(1) \to 0$$ for the corresponding filtration of $\rhobar_t$,
and write~$\gamma_t$ for the unique quotient character of~$\rbar_t$ %of~$\rhobar_t$,
then for each $t\in T(\Fpbar)$ either
  \begin{enumerate}
  \item we have~$\gamma_t\ne \alphabar(1)$, or
\item the cyclotomic character~$\epsilonbar$ is trivial, we have
  $\gamma_t=\alphabar(1)=\alphabar$, and the extension of~$\alphabar$
  by~$\gamma_t=\alphabar(1)$ induced by~$\rhobar_t$ is tr\`es ramifi\'ee.
  \end{enumerate}
\end{enumerate}
  Then, after replacing~$V$ by a dense open subscheme, 
% \mar{ME: Maybe ``after repacing~$V$ by a dense open subset''?  It is not $T$ that
% parameterizes extensions, but~$V$. TG: agreed, changed, commenting out.}
we have that for all
  $t\in V(\Fpbar)$ the extension of $\alphabar$ by~$\alphabar(1)$
  induced by the extension of~$\alphabar$ by $\rhobar_t$ is tr\`es
  ramifi\'ee.
\end{enumerate}
\end{prop}
\begin{rem}
  \label{rem: see below for weird condition}See Remark~\ref{rem: further explanation of weird condition} below for a
  further discussion of condition~(\ref{item: weird condition over Qp}).
\end{rem}

\begin{proof}[Proof of Proposition~\ref{prop:dimension control}]
	Writing $T$ as the union of its irreducible components,
	we obtain a corresponding decomposition of $V$ into the
	union of its irreducible components, and we can prove the
	theorem one component at a time.  Thus we may assume
	that $T$ (and hence $V$) is irreducible.
	% \mar{ME: We may not need this reduction to the irred.\ case, but 
	% 	just in case \dots }
        Replacing $T$ by a non-empty affine open subscheme, we may furthermore
        assume that it is affine, so that we may find a good complex $C^{\bullet}_T$
        supported in degrees $[0,2]$ that is
        isomorphic in the derived category to the Herr complex, as above.

	Let $f: T \to (\cX_{d,\red})_{\Fpbar}$ be the given morphism, and let
	$g: V \to (\cX_{d+ a n,\red})_{\Fpbar}$ denote the
	morphism~(\ref{eqn:induced morphism}).
	%in the statement of the proposition.
        In order to relate the scheme-theoretic
	image of $g$ to the scheme-theoretic image of $f$,
	we will also consider (in a slightly indirect manner)
	the scheme-theoretic image of the morphism
	$V \to (\cX_{d,\red})_{\Fpbar} \times(\cX_{d+ a n,\red})_{\Fpbar}$
	(the first factor being the composite
	of $f$ and the projection from $V$ to $T$, and the second factor
	being $g$).  If $(\cX_{d,\red})_{\Fpbar}$ were a scheme or an algebraic space,
	then the dimension of this latter scheme-theoretic image
	would be greater than or equal to the dimension of the former;
	but since $(\cX_{d,\red})_{\Fpbar}$ is a stack, we have to be slightly careful
	in comparing the dimensions of these two images.
	% \mar{ME: In the end, it may make sense to remove
	% 	this cautionary note, in so far as the actual 
	% 	computation below is fairly routine, and it serves
	% 	more as an explanation of the mistake in the previous
	% 	version.  But I'll leave it
	% 	in for now at least for our benefit.}


	%Rather than proving the bound of the lemma, we will prove
	%something stronger, namely we will prove that this bound
	%holds for the scheme-theoretic image of $V$ in
	%the product $\cX_d\times \cX_{d+n}$ (where the morphism onto
	%first factor is the composite of the projection $V \to T$
	%with $f$).

	Let $v:\Spec \Fbar_p \to V$ be an $\Fbar_p$-point, and let $t$ 
	denote the composite $\Spec \Fbar_p \to V \to T$, so that
	$t$ is an $\Fbar_p$-point of $T$.  We also write $f(t)$ for
	the composite $f\circ t$, and $g(v)$ for the composite $g\circ v$.
	We write $T_{f(t)}$ for the fibre product
	$T\times_{(\cX_{d,\red})_{\Fpbar}} \Spec \Fbar_p$ (with respect
        to the morphisms $f: T \to (\cX_{d,\red})_{\Fpbar}$
	and $f(t): \Spec \Fbar_p \to (\cX_{d,\red})_{\Fpbar}$), and write $V_{g(v)}$ 
	and $V_{(f(t),g(v))}$ for the evident analogous fibre
	products. 

	We begin with~(1). In order to prove the required bound on the dimension of the
	scheme-theoretic image of $g$, it suffices (e.g.\
	by~\cite[\href{https://stacks.math.columbia.edu/tag/0DS4}{Tag 0DS4}]{stacks-project}) to show
	%In order to prove the proposition, it suffices to show 
	that %  \mar{TG: this claim needs some justification? Not quite
  %         in the appendix to CEGS at the moment, I think? ME: I improved
  % the appendix so that it now gives this claim.}
\[\dim V_{g(v)} \buildrel ? \over \geq
	\dim V - e - n( [K:\Q_p] ad +r) + n^2 \dim G_{\alphabar}\]
	for $v$ lying in some non-empty open subscheme of $V$.
	As already alluded to above, what we will actually be able to
	do is to estimate the dimension of the fibre $V_{(f(t),g(v))},$
	and so our first job is to compare the dimension of this fibre
	with that of $V_{g(v)}$.  

	Let $\rhobar_{f(t)}$ denote the Galois
	representation corresponding to $f(t)$, and
	let $G_t  := \Aut(\rhobar_{f(t)}).$  Then,
        by the discussion of Section~\ref{subsec:isotrivial},
        the morphism
	$t: \Spec \Fbar_p \to (\cX_{d,\red})_{\Fpbar}$ induces an immersion
	$$[\Spec \Fbar_p/G_t] \hookrightarrow (\cX_{d,\red})_{\Fpbar},$$
	which in turn induces a monomorphism
\numequation
\label{eqn:residual gerbe}
	[\Spec \Fbar_p/G_t]\times_{(\cX_{d,\red})_{\Fpbar}} V_{g(v)} \hookrightarrow V_{g(v)}.
\end{equation}
	Since $\Spec \Fbar_p$ is a $G_t$-torsor over $[\Spec \Fbar_p/G_t],$
	we see that $V_{(f(t),g(v))}$ is a $G_t$-torsor over
	$[\Spec \Fbar_p/G_t]\times_{(\cX_{d,\red})_{\Fpbar}} V_{g(v)},$
	and hence that
		\numequation\label{eqn: first dimension inequality}\dim V_{(f(t),g(v))}  \leq \dim G_t + \dim V_{g(v)}.\end{equation}
	Thus it suffices to prove the inequality
	$$\dim V_{(f(t),g(v))} \buildrel ? \over \geq
	\dim V - e - n( [K:\Q_p] ad +r) + n^2\dim G_{\alphabar} + \dim G_t.$$
	For $t$ lying in some non-empty open subscheme of $T$,
	we have that $e = \dim T - \dim T_{f(t)},$
	and so, replacing $T$ by this non-empty open subscheme 
	and $V$ by its preimage,
	we have to show that
	$$\dim V_{(f(t),g(v))} \buildrel ? \over \geq
	\dim V - \dim T + \dim T_{f(t)} - n( [K:\Q_p] ad +r) + n^2\dim G_{\alphabar} + \dim G_t$$
	for $v$ lying in some non-empty open subscheme of $V$.
	In fact we will show that
	this inequality holds for every $\Fbar_p$-point
	$v$ of $V$.

	The right hand side of our putative inequality can be rewritten
	as
	$$n \rk Z_T^1 + \dim T_{f(t)} - n( [K:\Q_p] ad + r) +
	n^2\dim G_{\alphabar} + \dim G_t,$$
	while the local Euler characteristic formula
	shows that
	$$H^0(G_K,\rhobar_{f(t)}\otimes\alphabar^{\vee}) - H^1(G_K,\rhobar_{f(t)}\otimes\alphabar^{\vee}) + r = 
	- [K:\Q_p] ad.$$
	Thus we may rewrite our desired inequality as 
	\begin{multline*}
	\dim V_{(f(t),g(v))}
	\buildrel ? \over \geq
	n\bigl(\rk Z_T^1 - \dim H^1(G_K, \rhobar_{f(t)}\otimes \alphabar^{\vee})\bigr) \\ +
	n \dim H^0(G_K, \rhobar_{f(t)}\otimes \alphabar^{\vee}) + \dim T_{f(t)} + n^2 \dim G_{\alphabar} + \dim G_t.
\end{multline*}

	There is a canonical isomorphism
	$V_{(f(t),g(v))} \cong T_{f(t)}\times_T V_{g(v)},$
	and so, writing $S := (T_{f(t)})_{\red}$,
	we find that $\dim V_{(f(t),g(v))} = \dim S\times_T V_{g(v)}.$
	If we let $\rhobar_S$ denote
	the pull-back of the family $\rhobar_T$ over $S$,
	then the Galois representations attached to all the
	closed points of $S$ are isomorphic to $\rhobar_{f(t)}$.
	The discussion of Section~\ref{subsec:isotrivial} % \mar{TG:
  %         forward reference to gerbe material ME: This material is
  % now moved back to \S 2, and is hopefully correct.}
	shows we may find an \'etale cover $S' \to S$ such
	that the pulled-back family $\rhobar_{S'}$ is constant.  We may replace
	$S$ by $S'$ without changing the dimension
	we are trying to estimate, and thus we may furthermore
	assume that $\rhobar_S$ is the trivial family with fibres
	$\rhobar_{f(t)}$.

	If we let $C^0_S \to Z^1_S$
	denote the pull-back of the good complex $C^0_T \to Z^1_T$
	to $S$, then this complex is also the pull-back to
	$S$ of the good complex
	$C^0_t \to Z^1_t$ corresponding to the representation $\rhobar_t$.
	%Over $S$,
	%the Galois representations attached to all the
	%closed points are isomorphic to $\rhobar_{f(t)}$,
	%and so the cokernel $H^1(G_K,\rhobar_S)$
	%of $C^0_S \to Z^1_S$ has constant fibre dimension
	%(namely $\dim H^1(G_K, \rhobar_{f(t)})$) at each closed point.
	%Thus this cokernel is locally free. 
	If we let $W$ denote the vector space 
	$H^1(G_K,\rhobar_t\otimes\alphabar^{\vee})^{\oplus n}$, 
	thought of as an affine space over $\Fbar_p$,
	then there is a projection $S \times_T V \to 
	S\times_{\Fbar_p} W$, whose kernel is a (trivial) vector bundle
	which we denote by~$V'$.
        
	
	The affine space
	$W$ parameterizes a universal family of extensions
        $0 \to \rhobar_{f(t)} \to \rhobar_W \to \alphabar^{\oplus n} \to 0$,
	giving rise to a morphism $h: W \to (\cX_{d+ a n,\red})_{\Fpbar},$
        and the composite morphism
	$S\times_T V \to V \buildrel g \over \longrightarrow 
	(\cX_{d+ a n,\red})_{\Fpbar}$ admits an alternative 
	factorization as
       	$$S \times_T V \to S\times_{\Fbar_p} W \to W \buildrel h \over 
	\longrightarrow (\cX_{d+ a n,\red})_{\Fpbar}.$$
	Thus 
	$$S \times_T V_{g(v)} = S\times_T V \times_W W_{h(w)}$$
	(where $w$ denotes the image of $v$ in $W$ under the projection),
	and hence
	\begin{multline*} \dim S \times_T V_{g(v)} 
	= \rk V' + \dim S + \dim W_{h(w)} \\  =
	n\bigl(\rk Z_T^1 - \dim H^1(G_K,\rhobar_{f(t)}\otimes \alphabar^{\vee})\bigr)
	+ \dim S + \dim W_{h(w)},\end{multline*}
	so that our desired inequality becomes equivalent to the inequality
	\numequation\label{eqn: second dimension inequality}\dim W_{h(w)} 
	\buildrel ? \over \geq
	n \dim H^0(G_K, \rhobar_{f(t)}\otimes \alphabar^{\vee}) + n^2\dim G_{\alphabar} + \dim G_t.\end{equation}
        There is an action of $G_t \times \Aut_{G_K}(\alphabar^{\oplus n})$ on
	$$W = H^1(G_K,\rhobar_{f(t)}\otimes \alphabar^{\vee})^{\oplus n} \cong \Ext^1_{G_K}(\alphabar^{\oplus n},
	\rhobar_t),$$
	which lifts to an action on the family $\rhobar_W.$
	There is furthermore a (unipotent)
	action of $H^0(G_K,\rhobar_t\otimes \alphabar^{\vee})^{\oplus n}$
	on $\rhobar_W$ (lying over the trivial action on $W$). 
	Altogether, we find that the morphism $h$
	factors through $[W/ \bigl( H^0(G_K,\rhobar_t \otimes \alphabar^{\vee})^{\oplus n} 
	\rtimes (G_t \times \Aut_{G_K}(\alphabar^{\oplus n})) \bigr) ]$,
	implying that 
	$$\dim W_{h(w)} \geq n\dim H^0(G_K,\rhobar_t\otimes \alphabar^{\vee}) + n^2\dim G_{\alphabar} + \dim G_t,$$
	as required.

        We now turn to~(2), so that we assume that~$n=1$,
        that~$\alphabar$ is one-dimensional, and
        that~$\rhobar_T$ is generically maximally nonsplit of niveau~$1$. By
        Proposition~\ref{prop: structure of maximally nonsplit}, after
        possibly replacing~$T$ with a non-empty open subscheme, we can and do
        assume that~$\rhobar_T$ is maximally nonsplit of niveau one, that we can
        write~$\rhobar_T$ as an extension of
        families \numequation\label{eqn: filtration on rhobarT}0\to\rbar_T\to\rhobar_T\to\betabar_T\to 0\end{equation}where the
        family~$\betabar_T$ arises from twisting a
        character~$\betabar:G_K\to\Fpbartimes$ via a morphism
        $T\to\Gm$, and that if~$K=\Qp$, then~(\ref{item: weird
          condition over Qp}) holds for all finite type points~$t$
        of~$T$. (Indeed, $\betabar_T$ is the family corresponding to
        the character~$\ur_{\lambda_d}\otimes\psi_{\underline{n}_d}$
        in the notation of Proposition~\ref{prop: structure of maximally
    nonsplit}.)% \mar{TG: should make sure that our argument is OK in
        % the case $\rbar_T=0$.}

We want to show that~$\cE_V$ is generically maximally nonsplit of niveau~$1$.  Assume that this is the case; % then~$\cE_V$ is certainly
% generically maximally nonsplit of niveau~$1$, as claimed,
we now show that we then have equality in the bound of part~(1). After
replacing~$V$ with a non-empty open subscheme, we can and do suppose that~
$\cE_V$ is maximally nonsplit of niveau~$1$, so 
that in particular, for every finite type point~$v$ of~$V$, the extension
\[0\to\rhobar_{f(t)}\to\cE_{g(v)}\to\alphabar\to 0\]is maximally
nonsplit. From the proof of part~(1), we see that it
is enough to show that (after shrinking
~$V$ as we have) % after replacing~$V$ by a nonempty (so dense)
% subset,
equality holds in both~(\ref{eqn: first dimension inequality})
and~(\ref{eqn: second dimension inequality}).% \mar{TG: have to finish
%   this. Need to delete a scheme-theoretic image; should shrink the base
% so that we have constant rank, so actual vector bundles? Should
% actually note that it's not surjective on fibres! TG: I don't know
% what this remark means, so I'm leaving it here for now. Probably it is
% related to the argument below not being completely justified\dots}

% and
% that any automorphism of~$\cE_{g(v)}$ induces automorphisms
% of~$\rhobar_t$ and~$\alphabar$. Note in particular that
% since~$\rhobar_{f(t)}$ is maximally nonsplit of niveau~$1$, $\cE_{g(v)}$ will then
% be maximally nonsplit of niveau~$1$, so that~$\cE_V$ will certainly be generically maximally
% nonsplit of niveau~$1$, as claimed. We will see below that this condition also
% guarantees that we have equality in the bound of~(1).

% Suppose then that for some~$v$ we have an automorphism~$\theta$
% of~$\cE_{g(v)}$. There is an induced
% morphism~$\rhobar_{f(t)}\to\alphabar$, and if this morphism vanishes,
% then~$\theta$ induces automorphisms of~$\rhobar_{f(t)}$
% and~$\alphabar$, as required. Suppose instead that the induced
% morphism~$\rhobar_{f(t)}\to\alphabar$ does not vanish; then
% since~$\alphabar$ is irreducible, we see from~(\ref{eqn: filtration on
%   rhobarT}) that this morphism induces an isomorphism
% $\beta_t\isoto\alphabar$, and that the extension class corresponding to~$\cE_{g(v)}$
% is in the image of the natural map $H^1(G_K,\rbar_{f(t)})\to
% H^1(G_K,\rhobar_{f(t)})$. 


% We now turn to showing that  equality
% holds in the bound of part~(1). 

We begin by considering~(\ref{eqn: first dimension
  inequality}). For equality to hold here, it is enough to check
that the immersion~\eqref{eqn:residual gerbe}
induces an isomorphism on underlying reduced substacks
$$([\Spec \Fbar_p/G_t]\times_{(\cX_{d,\red})_{\Fpbar}} V_{g(v)})_{\red}
\iso (V_{g(v)})_{\red}.$$
We can check this on the level of $\Fbar_p$-valued points,
%\mar{ME: Do we literally prove this, or just a weaker statement on
%the level of points (which would still be enough)?}
%For this, we need
for which it suffices
to show that the
representation~$\cE_{g(v)}$ has a unique $d$-dimensional
subrepresentation, namely~$\rhobar_{f(t)}$; but this is immediate from
$\cE_{g(v)}$ being maximally non-split of niveau~$1$.

Similarly, to show that equality holds in ~(\ref{eqn: second dimension
  inequality}), it is enough to show that the morphism
\[[W/ \bigl( H^0(G_K,\rhobar_ {f(t)}\otimes \alphabar^{\vee}) \rtimes \bigl(G_t
  \times \Aut_{G_K}(\alphabar)\bigr) \bigr) ]\to(\cX_{d+a,\red})_{\Fbar_p}\]is a
monomorphism. To see this, note that  it follows from the definition of ``maximally nonsplit'' that any
automorphism of $\cE_{g(v)}$ induces automorphisms of~$\rhobar_{f(t)}$
and~$\alphabar$, so that we
have \[W\times_{(\cX_{d+a,\red})_{\Fbar_p}}W = W \times_{\Fbar_p}
\Bigl( H^0(G_K,\rhobar_t \otimes \alphabar^{\vee}) \rtimes \bigl(G_t
  \times \Aut_{G_K}(\alphabar)\bigr)\Bigr).\] Hence it is enough to show that the
morphism $[W/W\times_{(\cX_{d+a,\red})_{\Fbar_p}}W]\to(\cX_{d+a,\red})_{\Fbar_p}$ is a
monomorphism; this follows from Lemma~\ref{lem:stacks as quotients}.%  Lemma~\ref{lem:monomorphism
  % criterion}.\mar{TG: need to say more?}% , or equivalently
%that the diagonal
% morphism \numequation\label{eqn: diagonal to check is an iso}[W/W\times_{\cX_{d+a,\red}}W]\to
%   [W/W\times_{\cX_{d+a,\red}}W]\times_{\cX_{d+a,\red}}[W/W\times_{\cX_{d+a,\red}}W] \end{equation}
% is an isomorphism. Since  $W\to [W/W\times_{\cX_{d+a,\red}}W]$ is
% faithfully flat (because $W\times_{\cX_{d+a,\red}}W$ is a flat
% groupoid over~$W$),\mar{TG: justify this} we can check this after pulling
% back via the morphism \[W  \times_{\cX_{d+a,\red}} W  \to   [ W / W
%   \times_{\cX_{d+a,\red}} W ]  \times _{\cX_{d+a,\red}}   [ W / W
%   \times_{\cX_{d+a,\red}} W ],  \]whereupon~(\ref{eqn: diagonal to check is an iso}) becomes the
% identity morphism \[W\times_{\cX_{d+a,\red}}W\to
%   W\times_{\cX_{d+a,\red}}W\]which is certainly an isomorphism, as required.
 % for each finite type point~$t\in T$,
 %        the representation~$\rhobar_t$ is maximally
 %        nonsplit of the form~(\ref{eqn:maximally nonsplit
 %          structure}).\mar{TG: probably rewrite as an extension of an
 %          irreducible $\Gm$ family by a family~$\rbar_T$.}
	%\mar{ME: Have to flesh this out.}
	%Sketch of proof: The relative dimension of $V$ over $T$
	%is $n \rk Z^1_T$; while the dimension in the statement
        %of the proposition is $e + n(\rk Z^1_T - \rk C^0_T) - n^2.$
	%The ``$-n^2$'' terms comes from automorphisms of $\triv^{\oplus n}$.
	%Now the image of an extension classified by a point in $V$
	%depends only on its image in $H^1$, so this lets us cut down
	%from relative dimension $n \rk Z^1_T$ to relative dimension
	%$n \rk H^1$.  But elements of $H^0$ give
	%(unipotent) automorphisms of the extension, and so this lets us
	%cut down to $n(\rk H^1 - \rk H^0)$, which is the same as
	%$n(\rk Z^1_T - \rk C^0_T).$ This gives the
        %proposition.\mar{TG: this seems slightly backwards - how do
        %  you identify $\rk Z^1_T - \rk C^0_T$ with $[K:\Q_p] d +r$
        %  without going via $H^1-H^0$ and Euler char? ME: Probably; I'm going
	%to try to write the argument carefully soon, and then I'll
        %sort this out.}

In order to prove~(2), it remains to prove that~$\cE_V$ is generically maximally nonsplit of niveau~$1$. We need to show that after possibly
shrinking~$V$, for each $v\in V$ the image in
$\Ext^1_{G_K}(\alphabar,\betabar_{f(t)})$ of the element
of~$\Ext^1_{G_K}(\alphabar,\rhobar_{f(t)})$ corresponding to~$\cE_{g(v)}$ is
nonzero. To this end,
%\mar{ME: I looked in the proof of this Prop., and it didn't seem
%very related.}
note that,
%^as in the proof of Proposition~\ref{prop: structure of maximally nonsplit},
after possibly shrinking~$T$ (and remembering that $T$ is reduced),
we may assume that
  $H^2(G_K,\betabar_T\otimes\alphabar^\vee)$ is locally free of some
constant  rank. We can therefore repeat the construction that we
carried out after Lemma~\ref{lem: Ext2 locally free}, and find a good
complex \[C^0_T(\betabar)\to Z^1_T(\betabar)\] % whose formation is compatible
% with arbitrary base change, and
whose cohomology in degree~$1$, which we denote by~$H^1_T(\betabar)$,  is compatible with base change and
computes $H^1(G_K,\betabar_t\otimes\alphabar^\vee)=\Ext^1_{G_K}(\alphabar,\betabar_t)$ at finite type
points~$t$ of~$T$. By~\cite[\href{http://stacks.math.columbia.edu/tag/064E}{Tag
  064E}]{stacks-project} we have a morphism of complexes
\[(C^0_T\to Z^1_T)\to(C^0_T(\betabar)\to Z^1_T(\betabar)),\] compatible
with the natural maps on the cohomology groups
induced by the corresponding morphism of Herr complexes.
%induced by those for
%\mar{ME: I found this phrasing a bit hard to follow.}
%the Herr complexes.%  \mar{TG: this last paragraph below
%  % is probably not really correct, and this always confuses me; but
%  % hopefully something close to this is true?}

The kernel of~$Z^1_T\to H^1_T(\betabar)$ is a coherent sheaf, so after
possibly shrinking the reduced scheme~$T$, we can suppose that it is a
vector subbundle of~$Z^1_T$. By definition, we see that if we delete
from~$V$ the corresponding subbundle
 then the
required condition holds. It therefore suffices to show that
we are deleting a {\em proper} subbundle of~$V$.  %  there is no~$t$
%\mar{ME: Since we  are deleting a subbundle, surely this hold either
%for  all $t$ or none of them?}
%for which the fibre of this
%subbundle over~$t$ is the whole fibre ~$V_t$.
If this were not
the case, then (considering the fibre of $V$ and the subbundle under
consideration over some closed point~$t$), we would have an exact
sequence \[0\to\Ext^1_{G_K}(\alphabar,\betabar_t)\to\Ext^2_{G_K}(\alphabar,\rbar_t)\to\Ext^2_{G_K}(\alphabar,\rhobar_t)\to\Ext^2_{G_K}(\alphabar,\betabar_t)\to
  0; \]equivalently, by Tate local duality we would have an exact
sequence \[0\to\Hom_{G_K}(\betabar_t,\alphabar(1))\to\Hom_{G_K}(\rhobar_t,\alphabar(1))\to\Hom_{G_K}(\rbar_t,\alphabar(1))\to\]\[\Ext^1_{G_K}(\betabar_t,\alphabar(1))\to
  0. \]Assume
for the sake of contradiction that this is the case. Since $\rhobar_t$ is maximally nonsplit, the map
$\Hom_{G_K}(\betabar_t,\alphabar(1))\to\Hom_{G_K}(\rhobar_t,\alphabar(1))$ is an
isomorphism; so it suffices to show that we cannot have an isomorphism
$\Hom_{G_K}(\rbar_t,\alphabar(1))\to\Ext^1_{G_K}(\betabar_t,\alphabar(1))$. 

Since~$\rbar_t$ is maximally nonsplit, $\Hom_{G_K}(\rbar_t,\alphabar(1))$ is
at most 1-dimensional with equality if and only if $\ur_{\lambda_{d-1}(t)}\otimes\psi_{\underline{n}_{d-1}}=\alphabar(1)$; while by the local Euler characteristic
formula, ~$\Ext^1_{G_K}(\betabar_t,\alphabar(1))$ has
dimension at least $[K:\Qp]$, with equality if and only
if~$\betabar_t\ne\alphabar$ and~$\betabar_t\ne\alphabar(1)$. %  , and dimension strictly greater than~$[K:\Qp]$
% if~$\betabar=\alphabar$ or~$\alphabar(1)$.
It follows that we must have~$K=\Qp$, $\ur_{\lambda_{d-1}(t)}\otimes\psi_{\underline{n}_{d-1}}=\alphabar(1)$, ~$\betabar_t\ne\alphabar$ and~$\betabar_t\ne\alphabar(1)$; % , and $\betabar\ne\alphabar,
% \alphabar(1)$;
but this case was excluded by assumption~(\ref{item: weird condition
  over Qp}), and we have our a contradiction.

Finally, suppose that we are in the situation of~(3). We wish to argue
in the same way as part~(2), but rather than deleting the split locus,
we need to delete the peu ramifi\'ee locus.
%\mar{ME: How do we know the p.r.\ locus is closed? ME: Put in argument}
To be more precise, note that the surjection $\rhobar_T \to \alphabar(1)_T$
(where $\alphabar(1)_T$ denotes the constant rank one family given
by spreading out $\alphabar(1)$  over~$T$)
induces a composite morphism $Z^1_T \to H^1(G_K, \rhobar_T\otimes \alphabar^{\vee})
\to H^1(G_K, \epsilonbar_T)$
(where $\epsilonbar_T$ denotes
the spreading out of the  mod $p$ cyclotomic character  $\epsilonbar$
over~$T$),  
which in turn induces 
a morphism $V \to Y$ of total spaces of vector bundles
over $T$, with $Y$ denoting the total space of the trivial bundle
over~$T$ having fibre~$H^1(G_K, \epsilonbar)$.   Now $Y$ contains
a trivial subbundle $Y_{\mathrm{p.r.}}$ of codimension~$1$, parameterizing 
the peu ramifi\'ee classes.    The preimage of $Y_{\mathrm{p.r.}}$ is
a closed subscheme $V_{\mathrm{p.r.}}$ of~$V$,
and we wish to show that $V_{\mathrm{p.r.}}$ is a proper closed subscheme
of~$V$.

For this, it is enough to show that for each~$t$ (perhaps after replacing
$T$ by a non-empty open subset),
there is some extension of~$\alphabar$ by~$\rhobar_{f(t)}$ with the
property that the induced extension of~$\alphabar$
by~$\betabar_t=\alphabar(1)$ is tr\`es ramifi\'ee. We have an exact sequence
\[\Ext^1_{G_K}(\alphabar,\rhobar_t)\to
  \Ext^1_{G_K}(\alphabar,\betabar_t)\to\Ext^2_{G_K}(\alphabar,\rbar_t),\]so
we are done if $\Ext^2_{G_K}(\alphabar,\rbar_t)=0$. Since~$\rbar_t$ is
maximally nonsplit, this means that we are done unless~$\rbar_t$
admits~$\alphabar(1)$ as its unique quotient character. 

By hypothesis, this implies that~$\epsilonbar$ is trivial, 
and we can assume that the extension
of~$\betabar_t=\alphabar(1)=\alphabar$ by~$\alphabar(1)=\alphabar$
induced by~$\rhobar_t$ is tr\`es ramifi\'ee. Write~$c_t$ for the
corresponding class in~$H^1(G_K,\epsilonbar)$. As in
Remark~\ref{rem: further explanation of weird condition} below, it is
enough to show that there is a tr\`es ramifi\'ee extension
of~$\alphabar$ by~$\alphabar(1)$, corresponding to a
class~$d\in H^1(G_K,\epsilonbar)$ such that the cup product of~$c_t$
and~$d$ vanishes. If this is not the case, then~$c_t$ must be an
unramified class in $H^1(G_K,1)=H^1(G_K,\epsilonbar)$; but the
extension of~$K$ cut out by~$c_t$ considered as a class in
$H^1(G_K,1)$ is an extension given by adjoining a $p$th root of a
uniformizer, so is in particular a ramified extension, and we are done. 
% \mar{TG: at the moment it looks to me as if we could
%   actually have a problem in this case? This seems very odd though, so
% I'm probably making a mistake.}\mar{TG: apparently the idea is that
% everything is over some $\Gm\times\Gm$, and we can delete the bad locus.}
\end{proof}
\begin{rem}
  \label{rem: further explanation of weird condition}The curious
  looking condition Proposition~\ref{prop:dimension control}~(\ref{item: weird condition over Qp})  is not
  merely an artefact of our arguments. Indeed, if we fix
  characters~$\alphabar,\betabar:G_{\Qp}\to\Fpbartimes$, with
  $\betabar\ne\alphabar$, $\betabar\ne\alphabar(1)$, and let~\[\rbar=
  \begin{pmatrix}
    \alphabar(1)&*\\0&\betabar
  \end{pmatrix}\]be a non-split extension, then any extension of~$\alphabar$ by~$\rbar$ induces
the trivial extension of~$\alphabar$ by~$\betabar$. To see this, note that if we fix
an extension of~$\alphabar$ by~$\betabar$, then the condition that we
can form an extension  of~$\alphabar$ by~$\rbar$ realising this
extension  of~$\alphabar$ by~$\betabar$ is that the cup product of the corresponding classes
in~$H^1(G_{\Qp},\alphabar\betabar^{-1}(1))$
and~$H^1(G_{\Qp},\betabar\alphabar^{-1})$ vanishes. But by Tate
local duality, this cup product is a perfect pairing
of~$1$-dimensional vector spaces, so one of the two classes must
vanish, as required.\end{rem}
% It is therefore sufficient to prove that there is a proper closed
% subscheme of~$V$ which contains all those~$v$ for which
%  that the extension class corresponding to~$\cE_{g(v)}$
% is  in the image of this map. We construct this subscheme as the
% support of\mar{TG: what?}


% In our arguments we will also need the following result, which gives a
% criterion for the upper bound of Proposition~\ref{prop:dimension
%   control} to be attained.
% We say that a
% representation~$\rhobar:G_K\to\GL_d(\Fpbar)$ is \emph{maximally
%   nonsplit} if it has a unique filtration by~$G_K$-stable
% $\Fpbar$-subspaces, all of whose graded pieces are irreducible
% representations of~$G_K$. We say that the family~$\rhobar_T$ is
% maximally nonsplit if~$\rhobar_t$ is maximally nonsplit for
% all~$t\in T(\Fpbar)$. Take~$n=1$ in the construction before
% Proposition~\ref{prop:dimension control}, so that we have a morphism
% $V\to\cX_{d,a,\red}$, and the $\Fpbar$-points of the scheme-theoretic
% image of~$V$ are exactly the
% extensions \[0\to\rhobar_t\to\cE_t\to\alphabar\to 0\]where~$t\in T(\Fpbar)$.\mar{TG:
%   hopefully correct - we should maybe state this more clearly above if
%   so? ME: The statement about scheme-theoretic images isn't quite
% true, since scheme-theoretic image involves taking a closure.}

% Assume that~$\alphabar$ is irreducible. If~$\rhobar_T$ is maximally
% nonsplit then by Tate local duality, for each~$t\in T$ the
% space~$\Ext^2_{G_K}(\alphabar,\rhobar_t)\cong\Hom_{G_K}(\rhobar_t(1),\alphabar)$
% is at most $1$-dimensional. We say that~$\alphabar$ and~$\rhobar_T$
% are \emph{unrelated} if $\Hom(\rhobar_t(1),\alphabar)=0$ for
% all~$t\in T$, and we say that they are \emph{highly correlated} if
% $\Hom_{G_K}(\rhobar_t(1),\alphabar)\ne 0$ for all~$t\in T$. 

% By Lemma~\ref{lem: Ext2 locally free}, if~$\alphabar$ and~$\rhobar_T$
% are unrelated then $\Ext^2_{G_K}(\alphabar,\rhobar_T)$ vanishes, while
% if they are highly correlated, it is locally free of rank~$1$. In the
% former case we set~$r=0$, and in the latter we set~$r=1$. 

% \begin{prop}
%   \label{prop: building families of maximal dimension}Suppose that~$T$
%   is reduced and irreducible, that~$\alphabar$ is irreducible, and
%   that~$\rhobar_T$ is maximally nonsplit. Assume that ~$\alphabar$
%   and~$\rhobar_T$ are either unrelated or highly correlated. If
%   $e$ denotes the dimension of the scheme-theoretic
%   image % \mar{TG: should probably have an explanation that this
% %           exists, and is an algebraic stack. ME: I rewrote the discussion
% %   to emphasize that the morphisms in question factor through underlying
% % reduced algebraic substacks, which makes this clearer.}
% 	of $T$ in $(\cX_{d,\red})_{\Fpbar},$ then the dimension of the
%         scheme-theoretic image of $V$ in $\cX_{d+a n,\red}$ is exactly
%         $e + [K:\Q_p] ad+r -1$.\mar{TG: add statement about this
%           image being irreducible and having a dense open of maximally
%         nonsplit.}
%       \end{prop}\mar{TG: this might not quite be true, of course! The
%         idea is that in the induction below we will keep track both of
%         the case that~$\alphabar$ is fixed, and of the case that it
%         isn't; the former spaces (from which one gets the latter) will
%         have dimension one less, but will be inductively used in
%         building tres ramifiee extensions.}
%       \begin{proof}Examining the proof of
%         Proposition~\ref{prop:dimension control}, we see that it is
%         enough to check that there is a non-empty set of points
%         $v\in V$ such that equality holds in~(\ref{eqn: first
%           dimension inequality}), and a non-empty set of points
%         $w\in W$ such that equality holds in~(\ref{eqn: second
%           dimension inequality}).\mar{TG: now start following ME emails.}
% \end{proof}


% \begin{lemma}
% 	\label{lem:monomorphism criterion}
% 	If $X \to \cY$ is a morphism from an algebraic space~$X$
% 	to a stack~$\cY$ whose diagonal is representable by algebraic spaces
% 	{\em (}e.g.\ an algebraic stack{\em )},
% 	and if $R:= X\times_{\cY} X$, which is a groupoid in algebraic
% 	spaces over $X$, %is a flat groupoid over $X$,
% 	has the property that the projections $R\rightrightarrows X$
% 	are flat and locally of finite presentation,
% 	then the induced morphism
% 	$[X/R] \to \cY$ is a monomorphism.
% \end{lemma}
% \begin{proof}
% 	This follows from~\cite[Lem.~3.19]{Emertonformalstacks}.
% 	%If we write $\cX: = [X/R]$, then we have to show
% 	%that the morphism $\cX \to \cX\times_{\cY} \cX$ 
% 	%is an isomorphism.  The natural  morphism $X \to \cX$
% 	%is flat and surjective,
% 	%and so the natural morpism $R = X\times_{\cY} X
% 	%\to \cX\times_{\cY} \cX$ is also flat and surjective,
% 	%and so it suffices to check that the induced morphism
% 	%$\cX \times_{\cX\times_{\cY} \cX} R \to R$
% 	%is an isomorphism.   This follows by a direct
% 	%comparison of the source and target, remembering
% 	%that $\cX$ is defined to be the quotient stack $[X/R]$.
% \end{proof}

\section{$\cX_{\lowercase{d}}$ is a formal algebraic stack}\label{subsec: formal
  algebraic stack}We now use the theory of families of extensions to
show that~$\cX_d$ is a formal algebraic stack. The key ingredient is
Theorem~\ref{thm:Xdred is algebraic}, which we prove by induction
on~$d$. We begin by setting up some useful terminology, which is
motivated by the generalisations of the weight part of Serre's
conjecture formulated in~\cite{2015arXiv150902527G}. Recall that the
notion of a \emph{Serre weight} was defined in Section~\ref{subsec:
  notation and conventions}. Since~$\F$ is assumed to contain~$k$, we
can and do identify embeddings $k\into\Fpbar$ and $k\into\F$.

% Recall that for each embedding~$\sigmabar:k\into\Fpbar$, we have a
% corresponding fundamental character
% $\omega_{\sigmabar}:I_K\to\Fpbartimes$, 
% % \mar{ME: Based on the formulas below, I think that $\omega_{\sigmabar}$ 
% % 	is the right notation for the fundamental character, but TG
% % 	should check. TG: yes.}
% which
% corresponds via local class field theory (normalised to take
% uniformisers to geometric Frobeneii) to the composite of~$\sigmabar$
% and the natural map $\cO_K^\times\to k^\times$. It is easy to check
% that
% $\prod_{\sigmabar}\omega_{\sigmabar}^{e(K/\Qp)}=\varepsilonbar|_{I_K}$.

\begin{defn}\label{defn: max nonsplit weight k}
% We say that a representation $\rhobar:G_K\to\GL_d(\Fpbar)$ is of
% niveau~$1$ if it can be written as a successive extension of
% $1$-dimensional representations.
% We say that a representation of
% niveau~$1$ is maximally nonsplit if it can be written as such an
% extension in a unique way.
  If~$\underline{k}$ is a Serre weight, and~$\rhobar:G_K \to \GL_d(\Fbar_p)$ is
  maximally nonsplit
  of niveau~$1$, then we say that~$\rhobar$
  is of weight~$\underline{k}$ if we can \index{maximally nonsplit of
    niveau 1!of weight~$\underline{k}$}
  write % \mar{TG: note to self - this might become a bit cleaner if we
    % reversed the numbering of the~$\chi_i$}
  \[\rhobar\cong
    \begin{pmatrix}
      \chi_1 &*&\dots &*\\
      0&\chi_2&\dots &*\\
      \vdots&& \ddots &\vdots\\
      0&\dots&0&\chi_d\\
    \end{pmatrix}
  \] where
  \begin{itemize}
  \item
    $\chi_i|_{I_K}=\epsilonbar^{1-i}\prod_{\sigmabar:k\into\F}\omega_{\sigmabar}^{-k_{\sigmabar,d+1-i}}$, 
  % \item for each~$1\le i\le d-1$,
  %   $(\chi_{i+1}\chi_{i}^{-1})|_{I_K}=\varepsilonbar^{-1}\prod_{\sigmabar:k\into\Fpbar}\omega_{\sigmabar}^{k_{\sigmabar,d+1-i}-k_{\sigmabar,d-i}}$,
    and
  \item if $(\chi_{i+1}\chi_{i}^{-1})|_{I_K}=\epsilonbar^{-1}$,
    then %  $k_{\sigmabar,d+1-i}-k_{\sigmabar,d-i}=0$ for all~$\sigmabar$ if
    % $\chi_{i+1}\chi_{i}^{-1}\ne\varepsilonbar^{-1}$, and
    $k_{\sigmabar,d-i}-k_{\sigmabar,d+1-i}=p-1$ for all $\sigmabar$ if
    and only if
    $\chi_{i+1}\chi_{i}^{-1}=\epsilonbar^{-1}$ and the element
    of~$\Ext^1_{G_K}(\chi_i,\chi_{i+1})=H^1(G_K,\epsilonbar)$
    determined by~$\rhobar$ is
    tr\`es ramifi\'ee  (and otherwise
    $k_{\sigmabar,d-i}-k_{\sigmabar,d+1-i}=0$ for all~$\sigmabar$).
  \end{itemize}
\end{defn}% \mar{TG: I think this should be updated with the correct
%   definition, i.e.\ the Steinberg condition should also include tr\`es
% ramifi\'ee. Then we will have to modify the proofs in the rest of this
% section a bit, but I think it should be fine, because this is an open
% condition. TG: if we do this then maybe it is worth pointing out
% explicitly that the openness of the eigenvalue morphism means that
% e.g.\ on the component for the trivial weight, there is a dense open
% where none of the ratios of characters is cyclotomic.}

% It is easy to see that every~$\rhobar$ of niveau~$1$ is of weight~$\underline{k}$
% for at least one~$\underline{k}$, and that if~$\rhobar$ is furthermore
% maximally nonsplit then~$\underline{k}$ is unique.

Note that each maximally nonsplit representation of niveau one is of
weight~$\underline{k}$ for a unique~$\underline{k}$. For each~$\underline{k}$,
there exists a~$\rhobar$ which is maximally nonsplit of niveau~$1$ and
weight~$\underline{k}$; % , and is not of weight~$\underline{k}'$ for
% any~$\underline{k}'\ne\underline{k}$ (namely, one can take any
% $\rhobar$ of niveau~$1$ and weight~$\underline{k}$ which is maximally
% nonsplit; 
the existence of such a~$\rhobar$ follows by an easy
induction, and is in any case an immediate consequence of Theorem~\ref{thm:Xdred
  is algebraic} below.
% \mar{ME: Should we give a citation or explanation for why such
% 	a $\rhobar$ exists?  Or are we taking it as
%         given/known/obvious? TG: probably we should explain more,
%         although I think it's pretty clear - I guess it's just that
%         the $H^1$ is always positive-dimensional. TG: on further
%         thought, presumably if we are able to prove the more precise
%         version of the below with equality for the dimensions of the
%         components, then we get this as a consequence.}

\begin{defn}
  \label{defn: shifted weight}We say that a weight~$\underline{k}'$ is
  a \emph{shift} of a weight~$\underline{k}$ \index{shift (of a Serre weight)}
  if~$k_{\underline{\sigma},d}'=k_{\underline{\sigma},d}$ for
  all~$\sigmabar$, and if for each~$1\le i\le d-1$, we either have
  $k'_{\sigmabar,i}-k'_{\sigmabar,i+1}=k_{\sigmabar,i}-k_{\sigmabar,i+1}$
  for all~$\sigmabar$, or we have
  $k'_{\sigmabar,i}-k'_{\sigmabar,i+1}=p-1$ and
  $k_{\sigmabar,i}-k_{\sigmabar,i+1}=0$ for all~$\sigmabar$.
(Note in particular that any weight is a shift of itself.)
\end{defn}

\begin{defn}
  \label{defn: abuse of terminology crystalline lattice}We say that a
  representation~$\rho:G_K\to\GL_d(\cO)$ or~$\rho:G_K\to\GL_d(\Zpbar)$
  is \emph{crystalline}, or~\emph{crystalline of weight~$\lambda$},
  if~$\rho[1/p]$ is crystalline, resp.\ crystalline of
  weight~$\lambda$. We say that~$\rho$ is a \emph{crystalline lift} of
  $\rhobar:G_K\to\GL_d(\F)$ (resp.\ of $\rhobar:G_K\to\GL_d(\Fpbar)$).
\end{defn}

%Note in particular that any weight is a shift of itself.

The
motivation for Definitions~\ref{defn: max nonsplit weight k}
and~\ref{defn: shifted weight} is the following lemma.
%\mar{ME: Maybe say something about field of rationality,
%since (based on the proof) it  seems to be uncontrolled.} ME: Also,
%\mar{ME:
%Since we make a song-and-dance about crystalline lifts elsewhere,
%point out more explicitly that this hinges on maximal non-splitness. In fact
%maybe cite this lemma later in the crys. lift discussion, pointing out that
%it is only for maximally non-split $\rho$.}
\begin{lem}
  \label{lem: maxnonsplit weight k implies crystalline lifts to
    shifts}If~$\rhobar$ is maximally nonsplit of niveau~$1$ and
  weight~$\underline{k}$, if~$\underline{k}'$ is a shift
  of~$\underline{k}$, and~$\underline{\lambda}$ is a lift
  of~$\underline{k}'$, then~$\rhobar$ has a crystalline lift {\em (}defined
over $\Zbar_p${\em )} of 
% \mar{ME: Here and below, should explaing that ``crystalline'' for an integral
% rep'n means that the corresponding $p$-adic rep'n is crystalline.  (We use this
% abuse of terminology a lot here and in \S \ref{sec:properties}
% in discussing ``crystalline lifts''.)}
  weight~$\underline{\lambda}$. Furthermore, this lift can be chosen
  to be ordinary in the sense of Definition~{\em \ref{defn: ordinary}} below.
\end{lem}
\begin{proof}Write~$\rhobar$ as in Definition~\ref{defn: max nonsplit
    weight k}, and assume (by replacing~$E$ by a finite unramified extension if
  necessary) that both~$\rhobar$ and the filtration in Definition~\ref{defn: max nonsplit
    weight k} are defined over~$\F$.  Note firstly that by~\cite[Lem.\
  5.1.6]{2015arXiv150902527G} (which uses the opposite conventions for
  the sign of the Hodge--Tate weights to those of this book),
  together with the observation made above that
%\mar{ME: What role does this play? TG: the relationship between the
%  $\chi_i$ and the $k_i$ has a cyclo-shift, and the relationship
%  between $k$ and $\lambda$ a $\rho$-shift, and this is what makes it work.
%ME: OK, I'm happy with this.}
  $\prod_{\sigmabar}\omega_{\sigmabar}^{e(K/\Qp)}=\epsilonbar|_{I_K}$,
  we can find crystalline characters~$\psi_i:G_K\to\cO^\times$ such
  that~$\psibar_i=\chi_i$ and~$\psi_i$ has Hodge--Tate weights given
  by~$\lambda_{\sigma,d+1-i}$.

  It is enough to prove that~$\rhobar$ has a crystalline lift which is
  a successive extension of characters~$\psi_i$ of the form \[    \begin{pmatrix}
      \psi'_1 &*&\dots &*\\
      0&\psi'_2&\dots &*\\
      \vdots&& \ddots &\vdots\\
      0&\dots&0&\psi'_d\\
    \end{pmatrix},
  \]where each~$\psi'_i$ is an unramified twist of~$\psi_i$; note that
  any such representation is by definition ordinary in the sense of
  Definition~\ref{defn: ordinary}.  We prove this by induction on~$d$,
  the case $d=1$ being trivial. In the general case, write $\rhobar$
  as an extension \[0\to\rbar\to\rhobar\to\chibar_d\to 0,\]and suppose
  inductively that~$\rbar$ has a lift of the required form. Write
  ~$r:G_K\to\GL_{d-1}(\cO)$ for this lift; then we inductively seek a
  lift of~$\rhobar$ the form \numequation\label{eqn: rho lift is an
    extension}0\to r\to \rho\to \psi'_d \to 0, \end{equation}
where~$\psi'_d$ is an unramified twist of~$\psi_d$. We insist also
that~\numequation\label{eqn: ratio of crystalline chars isn't
  cyclotomic}\psi'_{d-1}(\psi'_d)^{-1}\ne \epsilon;\end{equation}
note that this could only fail
if~$k'_{\sigmabar,d}=k'_{\sigmabar,d-1}$ for all~$\sigmabar$, and in
this case we are only ruling out a single unramified twist.  Then
since~$\rhobar$ is maximally nonsplit, we have
$\Hom_{G_K}(r,\psi_d\epsilon)=0$, and it follows as in the proof of
Lemma~\ref{lem: extensions are crystalline} that such a~$\rho$, if it exists,
is automatically
crystalline.

It is therefore enough to show that there is a lift of the
form~\eqref{eqn: rho lift is an extension} satisfying~\eqref{eqn:
  ratio of crystalline chars isn't cyclotomic}.  As in the proof
of~\cite[Thm.\ 2.1.8]{2015arXiv150601050G}, we have an exact sequence
% \mar{ME: Why $\Hom(X,Y)$ notation rather than the $X^{\vee}\otimes Y$
% notation used above?}
\nummultline\label{eqn: lifting ses H1
  H2}H^1\bigl(G_K, (\psi'_d)^{-1}\otimes r\bigr) \to %\Hom_{\cO}(\psi'_d,r))\to
H^1(G_K, \chi_d^{-1} \otimes \rbar) %\Hom_{\cO}(\chi_d,\rbar))
\stackrel{\delta}{\to}
H^2\bigl(G_K, (\psi'_d)^{-1} \otimes r\bigr) %\Hom_{\cO}(\psi'_d,r)),
\end{multline}
Thus,
if~$c\in H^1(G_K, \chi_d^{-1} \otimes \rbar)$ %\Hom_{\cO}(\chi_d,\rbar))$
is the class determined
by~$\rhobar$, it is enough to show that we can choose the unramified
twist~$\psi'_d$ so that~$\delta(c)=0$. Now,
$H^2\bigl(G_K, (\psi'_d)^{-1}\otimes r\bigr)$ %\Hom_{\cO}(\psi'_d,r))$
is dual to
$H^0\bigl(G_K, r^{\vee} \otimes (\psi'_d \epsilon) \otimes E/\cO\bigr),$
%\Hom_{\cO}(r,\psi'_d\epsilon)\otimes E/\cO)$,
and
since~$\rbar$ is maximally nonsplit, we see that this latter group is equal
to $H^0\bigl(G_K,  (\psi'_{d-1})^{-1}\psi'_d\epsilon \otimes E/\cO\bigr).$ %\Hom_{\cO}(\psi'_{d-1},\psi'_d\epsilon)\otimes E/\cO)$.
In other
words, to show that~$\delta(c)=0$ in~\eqref{eqn: lifting ses H1 H2}, it
is enough to show the corresponding statement where~$r$ is replaced
by~$\psi'_{d-1}$; so we need only show that we can choose ~$\psi'_{d}$
so that the extension of~$\chi_d$ by~$\chi_{d-1}$ lifts to a
crystalline extension of~$\psi'_d$ by~$\psi'_{d-1}$. This is an
immediate consequence of the weight part of Serre's conjecture
for~$\GL_2$, as proved in~\cite{MR3324938}; see in
particular~\cite[Thm.\ 6.1.18]{MR3324938} and the references therein
(and note that these results also hold for~$p=2$,
by~\cite{2017arXiv171109035W}).
\end{proof}

\begin{remark}
We emphasize that the preceding lemma applies only to $\rhobar$ that
are maximally split of niveau~$1$.  
In Theorem~\ref{thm: strong existence of crystalline lifts}
below we prove an analogous result for arbitrary~$\rhobar$.
\end{remark}

%  This is the motivation for our
  % terminology, but we cannot establish this result at this point in
  % the development of our theory.
% \mar{TG: prove this, and probably also prove that there are
%   also lifts for all the shifts of weights, and that all of these
%   lifts (that we are constructing here) are ordinary. Also need to
%   justify that they are crystalline!}

% \begin{thm}\mar{TG: rewrite in terms of dimension}
% 	\label{thm:dimension of H2}
% 	If we fix an irreducible representation $\alphabar: G_K \to
% 	\GL_a(\Fbar_p)$,
% 	then the locus of $\rhobar$ in $(\cX_{d,\red})_{\Fpbar}(\Fbar_p)$
% 	for which $\dim \Hom_{G_K}(\rhobar,\alphabar) \geq r$ 
% 	is of codimension $\geq \lceil {r\bigl((a^2+1)r-a\bigr)}/{2}
%         \rceil$.\mar{TG: perhaps we should be more precise about what
%           this means, e.g.\ it could be defined and checked on
%           pullbacks to schemes?}
% 	% \mar{ME: I originally thought we could get by with $c_r = r$,
% 	% 	but this seems to be too gross of an underestimate,
% 	% 	and doesn't fit well with the induction.  I think we
% 	% 	need to get closer to the truth for the induction  In the course of doing this, we will
% 	% 	also prove a result on the existence of crystalline lifts
% 	% 	which is of independent interest.
% 	% 	to work, i.e.\ we have to make a good estimate for $c_r$.
% 	% ME: OK, hopefully what I've written works. TG: it seems OK to
%         % me, barring the fact of it not being as big as you claim if $r=1$\dots
%         % ME: I think I had an $r$ vs.\ $r+1$ confusion; see my response below.}
% \end{thm}




% Using the theory of families of extensions, we are able to obtain
% finer control over the structure of $\cX_d$,
% beginning with the following result.
% \mar{TG: assuming that everything works out, do the stack version
%   instead of the scheme version here; this takes care of the comment
%   about the eigenvalue morphism below. It's just a matter of making
%   the letters calligraphic and replacing ``subscheme'' with ``substack''.} 

Suppose that~$\rhobar_{\cT}$ is
generically maximally nonsplit of niveau~$1$, where~$\cT$ is some
integral finite type $\Fpbar$-stack. We say that~$\rhobar_{\cT}$ is
generically of weight~$\underline{k}$ if there is a dense open
substack~$\cU$ of~$\cT$ such that every $\Fpbar$-point of~$\rhobar_{\cU}$ is
maximally nonsplit of niveau~$1$, and is of weight~$\underline{k}$.
\begin{lem}
\label{lem:weighing families}
  If~$\rhobar_{\cT}$ is generically maximally nonsplit of
niveau~$1$, then~$\rhobar_{\cT}$ is generically of weight~$\underline{k}$ for a unique
Serre weight~$\underline{k}$.
\end{lem}
\begin{proof}
  Consider a dense open substack~$\cU\subseteq \cT$,  tuples~$\underline{n}_i$ and 
  morphisms~$\lambda_i:U\to\Gm$ as in Proposition~\ref{prop: structure of maximally
    nonsplit}. We inductively determine the~$k_{\sigmabar,i}$ as
  follows. We let~$k_{\sigmabar,d}$ be the unique choice with~$p-1\ge
  k_{\sigmabar,d}\ge 0$ not all equal to~$p-1$ and
  satisfying \[\psi_{\underline{n}_1}|_{I_K}=\prod_{\sigmabar:k\into\Fpbar}\omega_{\sigmabar}^{-k_{\sigmabar,d}}.\]
  Inductively, if $k_{\sigmabar,i+1}$ has been determined, then we
  demand that~$k_{\sigmabar,i}$ satisfies~$p-1\ge
  k_{\sigmabar,i}-k_{\sigmabar,i+1}\ge 0$  and
  \[(\psi_{\underline{n}_{d-i}}\psi_{\underline{n}_{d+1-i}}^{-1})|_{I_K}=\epsilonbar\prod_{\sigmabar:k\into\Fpbar}\omega_{\sigmabar}^{k_{\sigmabar,i}-k_{\sigmabar,i+1}}.\]This
  determines~$k_{\sigmabar,i}$, except in the case
  that~$(\psi_{\underline{n}_{d-i}}\psi_{\underline{n}_{d+1-i}}^{-1})|_{I_K}=\epsilonbar$.
  In this case, we consider the
  character~$\lambda_{d+1-i}\lambda_{d-i}^{-1}:U\to\Gm$. Since~$\cU$ is
  integral, this character is either constant or has open image. After
  possibly shrinking~$\cU$, we may therefore assume that either
  ~$\lambda_{d+1-i}(t)\ne\lambda_{d-i}(t)$ for all~$t\in \cU$, in which
  case we set $k_{\sigmabar,d+1-i}-k_{\sigmabar,d-i}=0$ for
  all~$\sigmabar$, or that the character is trivial.

  If the character is trivial, we may consider the locus in~$\cU$ over which
% \mar{ME: Again, how do we know it is  closed? ME: Added explanation}
  the extension between~$\chi_{i+1}$ and~$\chi_i$ is peu
  ramifi\'ee. The argument used in in the proof of part~(3) of
  Proposition~\ref{prop:dimension control}
  shows that this is a constructible subset of~$\cU$; since~$\cU$ is integral,
  after shrinking~$\cU$ we can thus suppose that it
  is either equal to~$\cU$, or
  %after shrinking~$\cU$ we can assume
  that it is empty. In the former
  case we set $k_{\sigmabar,d+1-i}-k_{\sigmabar,d-i}=0$ for
  all~$\sigmabar$, and in the latter we take
  $k_{\sigmabar,d+1-i}-k_{\sigmabar,d-i}=p-1$ for
  all~$\sigmabar$. % \mar{TG: I think this is backwards - also need
    % to change Steinberg definition anyway as indicated above.}
  It follows from the construction that every $\Fpbar$-point
  of~$\rhobar_{\cU}$ is maximally nonsplit of niveau~$1$ and
  weight~$\underline{k}$, as required.
  \end{proof}

 It will be convenient to use the following variant of the notation
 established in Proposition~\ref{prop: structure of maximally
    nonsplit}. % \mar{TG: of course in the end we might be best off just
    % having one set of notation, but while things are in flux I don't
    % want to mess around with it. TG: actually I'm inclined to leave
    % things as they are!}
  If~$\rhobar_{\cT}$ is generically maximally nonsplit of
  niveau one and weight~$\underline{k}$, then we
  set~$\nu_i:=\lambda_{d+1-i}$, and write~$\omega_{\underline{k},i}$
  for the character~$\epsilonbar^{d-i}\psi_{\underline{n}_{d+1-i}}$,
  so that \[\omega_{\underline{k},i}|_{I_K}=\prod_{\sigmabar:k\into\Fpbar}\omega_{\sigmabar}^{-k_{\sigmabar,i}},\] and~\eqref{eqn:maximally nonsplit structure} can be rewritten as
   \numequation\label{eqn:weight k version of maximally nonsplit structure}\rhobar_t\cong \begin{pmatrix}
      \ur_{\nu_d(t)}\otimes\omega_{\underline{k},d} &*&\dots &*\\
      0& \ur_{\nu_{d-1}(t)}\otimes\epsilon^{-1}\omega_{\underline{k},d-1}&\dots &*\\
      \vdots&& \ddots &\vdots\\
      0&\dots&0& \ur_{\nu_1(t)}\otimes\epsilon^{1-d}\omega_{\underline{k},1}\\
    \end{pmatrix}.  \end{equation}
We have the \emph{eigenvalue morphism} % \mar{TG: needs a better name! These
  % aren't slopes at all but I can't think of anything better right now.}
~$\underline{\nu}:\cU\to(\Gm)^d$ given
by~$(\nu_1,\dots,\nu_d)$. \index{eigenvalue morphism}

\begin{df}
If $\underline{k}$ is a Serre weight, we let
$(\Gm)^d_{\underline{k}}$ denote the closed subgroup scheme of $(\Gm)^d$ parameterizing 
tuples~$(x_1,\dots,x_d)$ for which~$x_i=x_{i+1}$
whenever~$k_{\sigmabar,i}-k_{\sigmabar,i+1}=p-1$ for all~$\sigmabar$. \index{$(\Gm)^d_{\underline{k}}$}
\end{df}

Note that $(\Gm)_{\underline{k}}^d$ is closed under the simultaneous multiplication
action of $\Gm$ on $(\Gm)_{\underline{k}}^d$ (i.e.\ the action $a \cdot (x_1,
\ldots,x_d) := (a x_1,\ldots, a x_d)$).

It follows from the definitions that
%after possibly shrinking~$\cU$,
the eigenvalue morphism $\underline{\nu}$
is valued in~$(\Gm)^d_{\underline{k}}.$
%the subgroup~$(\Gm)^d_{\underline{k}}$ consisting of
%tuples~$(x_1,\dots,x_d)$ for which~$x_i=x_{i+1}$
%whenever~$k_{\sigmabar,i}-k_{\sigmabar,i+1}=p-1$ for all~$\sigmabar$. 

% We now record a lemma which provides an addendum to the proof of
% Proposition~\ref{prop:dimension control}
% in the tr\`es ramifie\'e context.

% \begin{lemma}
% \label{lem:tres ramifiee}
% Suppose that we are in the context of
% Proposition~{\em \ref{prop:dimension control}~(2)},
% and that $T$ is furthermore irreducible.
% Let $\underline{k}$ denote the Serre weight of the family $\rhobar_T$
% {\em (}which is well-defined, by Lemma~{\em \ref{lem:weighing families})},
% and let $\underline{\nu}: U \to (\Gm)^d_{\underline{k}}$
% denote the eigenvalue morphism
% {\em (}where $U$ denotes a non-empty open subset of $T$ on which the eigenvalue 
% morphism is defined{\em)}.
% Suppose that the family $\rhobar_T$ sits in a short exact sequence
% $$0 \to \rbar_T \to \rhobar_T \to \alphabar(1) \to 0,$$
% and finally,
% suppose that the morphism $\underline{\nu}': \Gm \times U \to (\Gm)_{\underline{k}}^d$
% defined via $(a,u) \mapsto a\cdot \underline{\nu}(u)$
% has dense image.
% Then we may find a non-empty open {\em(}so dense{\em)} subset $U'$ of $T$
% such that for each $t \in U'(\Fbar_p)$,
% there is an extension of $\alphabar$  by $\rhobar_t$
% which induces a tr\`es ramifi\'ee extension
% of $\alphabar$ by~$\alphabar(1).$\mar{TG: I don't understand the
%   statement of this lemma! Or at least I don't understand the
%   application of it below - don't we want something that says that the
%   universal family that we construct contains tr\`es ramifi\'ee
%   extensions? Whereas this just says that it isn't automatically false?}
% \end{lemma}
% \begin{proof}
% We first of all replace $T$ by $U$, so that the eigenvalue morphism $\underline{\nu}$
% is defined on~$T$.  Note that our assumption on $\rhobar_T$ (that it admits
% the constant character $\alphabar(1)$ as a quotient) shows that 
% the image of $\nu$ is contained in a hyperplane section $Z$
% of $(\Gm)^d_{\underline{k}}$
% on which $x_d$ is contant; the assumption that $\underline{\nu}'$ 
% has dense image then amounts to assuming that the image of $\underline{\nu}$
% is dense in $Z$.   Since this image is also constructible, it then contains 
% an open subscheme $V$ of $Z$.
% %\mar{TG: I don't
% %  understand this - how is it contained in~$Z$ but also containing an open? ME: Fixed typo} 
% Replacing $T$ by
% the preimage $\underline{\nu}^{-1}(V)$, we may then assume that
% $\underline{\nu}$ surjects onto $V$.

% \mar{ME: What follows is a cut-and-paste --- has to be cleaned up}
% We now return to the end of the
% proof of Proposition~\ref{prop:dimension control}; we see (now switching to the
% notation of that proof)
% \mar{ME: This argument mixes material from the proof of Prop.~\ref{prop:dimension control}
% with the material of the current proof --- e.g.\ by its invokation of $\cX_{d-1}$ ---
% but the notation of the two arguments is in horrible conflict, e.g.\ $\rhobar$ in one
% is $\rbar$ in the other.  Can you think of an easy way to improve this, e.g.\
% by adding a statement to~\ref{prop:dimension control} that can then just
% be cited here.}
% that it is enough to show that for each~$t$,
% there is some extension of~$\alphabar$ by~$\rhobar_{f(t)}$ with the
% property that the induced extension of~$\alphabar$
% by~$\betabar_t=\alphabar(1)$ is tr\`es ramifi\'ee. We have an exact sequence
% \[\Ext^1_{G_K}(\alphabar,\rhobar_t)\to
%   \Ext^1_{G_K}(\alphabar,\betabar_t)\to\Ext^2_{G_K}(\alphabar,\rbar_t),\]so
% we are done if $\Ext^2_{G_K}(\alphabar,\rbar_t)=0$. Since~$\rbar_t$ is
% maximally nonsplit, this means that we are done unless~$\rbar_t$
% admits~$\alphabar(1)$ as its unique quotient character. By our
% inductive assumption that the image of the eigenvalue morphism is
% open, this is only possible if~$\varepsilonbar$ is trivial, and we
% have~$k_{\sigmabar_2}-k_{\sigmabar_3}=p-1$ for all~$\sigmabar$.

% In
% this case, since by induction we know
% that~$\cX_{d-1,\red,\Fpbar}^{\underline{k}_{d-1}}$ is generically maximally nonsplit of niveau~$1$ and
% weight~$\underline{k}_{d-1}$, we can assume that the extension
% of~$\betabar_t=\alphabar(1)=\alphabar$ by~$\alphabar(1)=\alphabar$
% induced by~$\rhobar_t$ is tr\`es ramifi\'ee. Write~$c_t$ for the
% corresponding class in~$H^1(G_K,\varepsilonbar)$. As in Remark~\ref{rem: further explanation of weird condition}, it is
% enough to show that there is a tr\`es ramifi\'ee extension
% of~$\alphabar$ by~$\alphabar(1)$, corresponding to a
% class~$d\in H^1(G_K,\varepsilonbar)$ such that the cup product of~$c_t$
% and~$d$ vanishes. If this is not the case, then~$c_t$ must be an
% unramified class in $H^1(G_K,1)=H^1(G_K,\varepsilonbar)$; but the
% extension of~$K$ cut out by~$c_t$ considered as a class in
% $H^1(G_K,1)$ is an extension given by adjoining a $p$th root of a
% uniformizer, so is in particular a ramified extension. 
% \end{proof}


% In the following theorem, we say that a closed substack~$\cY$ of~$(\cX_{d,\red})_{\Fpbar}$
% is \emph{preserved under twist} if for every $\Fpbar$-point $t$
% of~$\cY$, and every $a\in\Fpbartimes$, there is an $\Fpbar$-point $t'$
% of~$\cY$ with $\rhobar_{t'}\cong\rhobar_t\otimes\ur_a$. (This notion
% is in some sense the opposite of being twistable.)\mar{TG: check
%   if we use this notion; also, might prefer to just actually define
%   the (universal?) unramified twist of a closed substack above (I don't think we do
%   right now), and then say that the substack is equal to its
%   unramifie,d twist.}

We have seen in Proposition~\ref{prop:X is an Ind-stack}
that~$\cX_{d}$ is an Ind-algebraic stack, which is in fact the $2$-colimit 
of algebraic stacks with respect to transition morphisms that
are closed immersions.   We let~$\cX_{d,\red}$ be the
underlying reduced substack of~$\cX_d$; it is then a closed substack of $\cX_d$
with the same structure:  namely, it is the $2$-colimit of closed algebraic stacks (which
can even be assumed to be reduced) with respect to transition morphisms that are closed
immersions. % , in the sense of~\cite[Defn.\ 3.27]{Emertonformalstacks}.
\begin{thm} 
	\label{thm:Xdred is algebraic}\leavevmode
        \begin{enumerate}
        \item 	The Ind-algebraic stack $\cX_{d,\red}$
is an algebraic stack, of finite presentation over $\F$.
\item  We can write~$(\cX_{d,\red})_{\Fpbar}$ 
% \mar{ME: Any reason not to write $\cX_{d,\red,\Fbar_p}$ instead, to fit
% with the notation for the irred.\ comps.?}
as a union of closed algebraic substacks of finite
presentation
over~$\Fpbar$ 
% \mar{TG: stupidity: can we descend the small guy, or should
  % this result just be over~$\Fpbar$, and then the descent can happen
  % at the very end? Or we could descend the weight~$\underline{k}$
  % stacks now but not the small one.}
\[(\cX_{d,\red})_{\Fpbar}=\cX_{d,\red,\Fpbar}^{\negligible}\cup\bigcup_{\underline{k}}\cX_{d,\red,\Fpbar}^{\underline{k}},\]
where:
\begin{itemize}
\item $\cX_{d,\red,\Fpbar}^{\negligible}$ is empty if~$d=1$, and otherwise is
  non-empty of dimension
  strictly less than~$[K:\Qp]d(d-1)/2$. %, and is preserved under twist,
\item each $\cX_{d,\red,\Fpbar}^{\underline{k}}$ is a closed 
irreducible substack of dimension 
% \mar{ME: Added {\em absolutely} to the irreducibility statement here, should address
% it in the proof too. TG: changed this to ``geometric'', which I think
% is what you mean? Wasn't sure how you meant to prove it (I think it
% just follows from the compatibility of our constructions with
% extension of scalars?), but I added a
% sentence below.}
  ~$[K:\Qp]d(d-1)/2$, and is generically maximally nonsplit of niveau
  one and weight~$\underline{k}$.
  % \mar{ME: Emphasize that this description of the generic points
  %         of $\cX_{d,\red}^{\underline{k}}$ shows that they are mutually
  %         distinct as $\underline{k}$ varies. TG: see Remark~\ref{rem:
  %           different k give different components}
  %        immediately below?
%          We could say it here too, though.}
  The corresponding eigenvalue morphism % \mar{TG: does this quite make
  %   sense as written? The eigenvalue morphism involves a choice
  %   of~$U$, I think the statement should be easily seen to be
  %   independent of this. TG: maybe can say ``for any diagonalizing
  %   scheme''? TG: is there a way to either set up some equivalence
  %   relation on the possible choices of~$U$, or just to actually
  %   define this morphism on an open in the stack itself? Also if we
  %   rewrite, need to change other instances of ``eigenvalue morphism''
  % in the next page or so.}
is dominant {\em (}i.e.\ has dense image in~$(\Gm)^d_{\underline{k}}${\em )}.%\mar{TG: I think ME
   % pointed out that this should say dense image, here and below? Not
   % sure so not changing it.}
 %, and is preserved under twist, % \mar{TG: actually given that
    % Proposition~\ref{prop:dimension control} is an inequality,
    % presumably we only get an upper bound for the dimension here, too?
    % Assuming this is the case, we should be careful later on to show
    % that each of these is actually an irreducible component, and not
    % something smaller, I guess. Probably we could also refine the
    % arguments here and get that equality is holding? I suspect that we
    % may need to do that if we want to actually enumerate the
    % irreducible components rather than just get an upper bound for
    % them? I guess we would like to have a criterion for when equality
    % holds in Proposition~\ref{prop:dimension control}} 
  % \mar{TG: should add in the statement that the
  %   representations are maximally nonsplit, and that the $H^2$ are
  %   generically free of rank $0$ or $1$, the latter being exactly when
  %   the jumps in the Serre weight are all~$p-1$.}
% \item there is a nonempty open {\em (}so dense{\em )} % \mar{TG: check open is proved}
%   substack
%   of~$\cX_{d,\red}^{\underline{k}}$ whose $\Fpbar$-points are all of
%   niveau~$1$ and weight~$\underline{k}$, and are maximally
%   nonsplit.
  % \mar{TG: probably going to end up also saying that all
    % (maximally nonsplit?)
    % points of niveau~$1$ and this weight are contained in this?}
  %o9\mar{TG: might end up being precisely these?!}
% \item for each irreducible
%   representation~$\alphabar:G_K\to\GL_a(\Fpbar)$, there is a nonempty
%   open (so dense) substack of~$\cX_{d,\red}^{\underline{k}}$  whose
%   $\Fpbar$-points all correspond to representations which do not have
%   any Jordan--H\"older factor isomorphic to~$\alphabar$.
\end{itemize}


\item  If we fix an irreducible representation
$\alphabar: G_K \to \GL_a(\Fbar_p)$ {\em (}for some $a \geq 1${\em )},
then the locus of $\rhobar$ in
$\cX_{d,\red}(\Fbar_p)$ for which
$\dim \Hom_{G_K}(\rhobar,\alphabar) \geq r$ {\em (}for any $r \geq 1${\em )}
is {\em (}either empty,
or{\em )} of dimension at most
\[[K:\Qp]d(d-1)/2- \lceil {r\bigl((a^2+1)r-a\bigr)}/{2} \rceil.\] % \mar{TG: perhaps we
  % should be more precise about what this means, e.g.\ it could be
  % defined and checked on pullbacks to schemes?}
Furthermore, the locus
of $\rhobar$ in
$\cX_{d,\red,\Fpbar}^{\negligible}(\Fbar_p)$ for which
$\dim \Hom_{G_K}(\rhobar,\alphabar) \geq r$ is of dimension strictly
less than this.%\mar{TG: think we need this for the induction.}

\item  % If we fix an irreducible representation
% $\alphabar: G_K \to \GL_a(\Fbar_p)$, then t
If we fix an irreducible representation
$\alphabar: G_K \to \GL_a(\Fbar_p)$ {\em (}for some $a \geq 1${\em )},
then the locus of $\rhobar$ in
$\cX_{d,\red}(\Fbar_p)$ for which
$\dim \Ext^2_{G_K}(\alphabar,\rhobar) \geq r$ is of dimension at most
\[[K:\Qp]d(d-1)/2-r.\]%  \mar{TG: also add $H^2$
% statement explicitly}
        \end{enumerate}
\end{thm}
\begin{rem}
  Note that in parts~(3) and~(4), the locus of points in question
  corresponds to a closed substack of~$(\cX_{d,\red})_\Fpbar$, by
  upper-semicontinuity of fibre dimension. 
\end{rem}

\begin{rem}\label{rem: different k give different components}
  Since the~$\cX_{d,\red,\Fpbar}^{\underline{k}}$ are irreducible, have
  dimension equal to that of~$(\cX_{d,\red})_\Fpbar$, and 
  have pairwise disjoint open substacks (corresponding to maximally
  nonsplit representations of niveau~$1$ and weight~$\underline{k}$),
  they are in fact distinct irreducible components of~$(\cX_{d,\red})_\Fpbar$.
\end{rem}

\begin{remark}
	We can, and will, be much more precise about the structure of $\cX_{d,\red}$.
	% Namely, we are able to show that it is equidimensional, of
	% dimension $[K:\Q_p]d(d-1)/2$, and also to
	% enumerate its irreducible components.
	% In Subsection~\ref{subsec: dimensions of families} 
	% we begin our preparations for this, by studying more carefully
	% the dimensions of the various families of extensions constructed
	% in the proof of Proposition~\ref{prop:Xdred is algebraic}.  
	Namely, in Chapter~\ref{sec:properties} we combine
        Theorem~\ref{thm:Xdred is algebraic} with additional arguments
        to show that $\cX_{d,\red}$ is equidimensional of dimension
        $[K:\Q_p]d(d-1)/2$; % so that  we may
        % take~$\cX_{d,\red}^{\negligible}=0$; 
        accordingly, the
        irreducible components of~$(\cX_{d,\red})_\Fpbar$ are
        precisely the
        ~$\cX_{d,\red,\Fpbar}^{\underline{k}}$, and in particular
        are in bijection with the Serre
        weights~$\underline{k}$. We also show that these irreducible
        components are all defined over~$\F$. % to compute its dimension, and to enumerate
        % its components.
        % \mar{ME: Have to remember
% 		to add this last result to
%                 Section~\ref{sec:properties}. TG: do we really
%                 enumerate its components? Maybe we do, but this seems
%                 a little painful even conjecturally, keeping track of
%                 the tres ram extensions.}
      \end{remark}
      \begin{rem}\label{rem: open eigenvalue morphism implies ratio
          not cyclo}
        Note that since the eigenvalue morphism has dense image, it
        follows that~$\cX_{d,\red}^{\underline{k}}$ contains a dense
        open substack with the property that $\nu_i(t)\ne
        \nu_{i+1}(t)$ unless $k_{\sigmabar,i}-k_{\sigmabar,i+1}=p-1$
        for all~$\sigmabar$.
      \end{rem}
	\begin{remark}
          As will be evident from the proof, the upper bound of
          Theorem~\ref{thm:Xdred is algebraic}~(3) is quite crude when
          $[K:\Q_p] > 1$, although it is reasonably sharp in the case
          $K = \Q_p$.  However, it suffices for our purposes, and
          indeed we will only use the (even cruder) consequence Theorem~\ref{thm:Xdred is
            algebraic}~(4) in the rest of the book. % for
          % which it is sufficient to have an upper bound of
          % $[K:\Qp]d(d-1)/2-r$ on the dimension as in
          % , an upper bound which
          % is certainly satisfied by
          % $[K:\Qp]d(d-1)/2-\lceil
          % r\bigl((a^2+1)r-a\bigr)/2\rceil$.
          % \mar{TG: this is only true
%                 if $r>1$! Hopefully not an issue? Earlier you were saying~$r$\dots ME: Yes, sorry, we need $r+1$ when we work over $\Z_p$, and (equivalently)
% 	$r$ when we work on the special fibre.  I think this same
% confusion (on my part) is underlying your comment above; I'll work through
% the argument and try to sort this out.}
	\end{remark}


\begin{proof}[Proof of Theorem~{\em \ref{thm:Xdred is algebraic}}] 
 Recall that a closed immersion of
  reduced algebraic stacks that are locally of finite type over $\F_p$
  which is surjective on finite-type points is necessarily an
  isomorphism. As recalled above, $\cX_{d,\red}$ is an inductive limit of such
  stacks (indeed, by Lemma~\ref{alem: underlying reduced of Ind algebraic}, we have~$\cX_{d,\red}=\varinjlim\cX^a_{d,h,s,\red}$,
where the~$\cX^a_{d,h,s}$ are as in Section~\ref{subsec: X is Ind
  algebraic}), and so 
if we produce closed algebraic substacks
  $\cX_{d,\red,\Fpbar}^{\negligible}$ and $\cX_{d,\red,\Fpbar}^{\underline{k}}$
 of $\cX_d$,
  the union of whose $\Fpbar$-points exhausts those
  of~$\cX_{d,\red}$, then $\cX_{d,\red,\Fpbar}$ will in fact
be an algebraic stack which is the union of its closed substacks
 $\cX_{d,\red,\Fpbar}^{\negligible}$ and $\cX_{d,\red,\Fpbar}^{\underline{k}}$. % \mar{TG: also need the irreducibility of~
    % $\cX_{d,\red}^{\underline{k}}$, which hopefully follows from the
    % construction.}
  % In fact, it suffices to construct
  % the closed substacks~$\cX_{d,\red}^{\underline{k}}$,
  % and then to show that the remaining
  % $\Fpbar$-points of $\cX_{d,\red}$ are contained in a finite union of
  % finite type algebraic substacks of dimension less than
  % $[K:\Qp]d(d-1)/2$.
Thus~(1) is an immediate consequence of~(2) (where the ``union'' statement
in~(2) is now to be understood on the level of $\Fbar_p$-points).

%  Since~(1) is an immediate consequence of~(2) (given the known Ind-algebraic
%stack structure of $\cX_{\red}$), while
Claim~(4) follows
  from~(3) (with~$\alphabar$ replaced by
  $\alphabar\otimes\epsilonbar$) % \mar{TG: check sign}
  by Tate local duality and the easily verified inequality
  \[\lceil {r\bigl((a^2+1)r-a\bigr)}/{2} \rceil\ge r.\]  Thus it is enough to
  prove~(2) and~(3), which we do simultaneously by induction
  on~$d$. 

As recalled in Remark~\ref{rem: irreducible representations finite up
  to twist}, there are up
to twist by unramified characters only finitely many irreducible
$\Fpbar$-representations of~$G_K$ of any fixed dimension. % ; furthermore,
% these representations are all absolutely irreducible, and up to twist
% by unramified characters they exhaust the irreducible
% $\Fpbar$-representations of~$G_K$ of this fixed dimension.
Accordingly,
we let $\{\alphabar_i\}$ be a finite set of irreducible continuous
representations $\alphabar_i: G_K \to \GL_{d_i}(\Fpbar)$, such that
any irreducible continuous representation of $G_K$ over $\Fbar_p$ of
dimension at most $d$ arises as an unramified twist of exactly one of
the~$\alphabar_i$. We let the 1-dimensional representations in this set be the
characters  ~$\psi_{\underline{n}}$ defined in Section~\ref{subsec: twists of
  families}. % \mar{TG: reference Galois representations discussion
% above for explanation of this.}

Each~$\alphabar_i$ % (in our chosen set of
  % representatives for the classes of irreducible representations up to
  % unramified twist) of dimension~$d$
  corresponds to a finite type point
  of~$\cX_{d_i,\red}$, whose associated residual gerbe
  is a substack of~$\cX_{d_i,\red}$ of dimension~$-1$:
  the morphism $\Spec\Fpbar\to\cX_{d_i,\red}$ corresponding
  to $\alphabar_i$ factors through
  a monomorphism $[\Spec\Fpbar/\Gm]\to\cX_{d_i,\red}$.
  %\mar{ME: Maybe this monomorphism claim is a residual gerbe fact(?), in which
  %	  case we can actually apply the gerbe discussion here!}
  It follows from
  Lemma~\ref{lem: dimension of Gm twist} that for each~$\alphabar_i$
  there
  is an irreducible closed zero-dimensional
  algebraic substack of~$\cX_{d_i,\red}$ of finite presentation
  over~$\Fpbar$ %and dimension~$0$
 % (and in the case~$d_i=1$, of dimension exactly~$0$)
 whose $\Fpbar$-points are exactly the unramified
  twists of~$\alphabar_i$.
 % (The dimension of this stack is easily seen to be~$0$ even
 % if~$d_i>1$, but we do not need to know this for the arguments that follow.) % \mar{TG: could presumably prove
    % exactly~$0$ here if we wanted to.}

 In particular, if $d=1$, then % the set of Serre weights~$\underline{k}$
 % is in bijection with the set of~$\alphabar_i$ of dimension~$1$,
 % because there is a unique~$\alphabar_i:G_K\to\Fpbartimes$ extending
 %  $\prod_{\sigmabar:k\into\Fpbar}\omega_{\sigmabar}^{-k_{\sigmabar,1}}$. Write
% we have the characters  ~$\psi_{\underline{k}}$ defined in section~\ref{subsec: twists of
%   families}, 
% %  for this character,
%   % of~$G_K$ 
%   and
  we let $\cX_{1,\red}^{\underline{k}}$ be the
  zero-dimensional stack constructed in the previous paragraph, whose $\Fpbar$-points
  are the unramified twists of~$\psi_{\underline{k}}$; this satisfies the
  required properties by definition, so~(2) holds when~$d=1$. For~(3),
  note that if~$r>0$ then
  we must have~$a=1$, and then the locus where
  $\Hom_{G_K}(\rhobar,\alphabar)$ is non-zero (equivalently, $1$-dimensional) is exactly the closed
  substack of dimension~$-1$ corresponding to~$\alphabar$, so the
  required bound holds. % \mar{TG: need to make sure that the rest of
    % the argument takes account of the exceptional case Proposition~\ref{prop:dimension control}~(\ref{item: weird condition over Qp}).}

%   Assume that we have proved that~(2) and~(3) hold in
%   dimension~$d$.  Let~$\underline{k}$ be a Serre
%   weight, let~$\underline{k}_{d-1}$ be the Serre weight in
%   dimension~$(d-1)$ obtained by deleting the first entry
%   in~$\underline{k}$, % let~$\underline{k}_d$ be the Serre weight in
%   % dimension~$1$ obtained by taking the last entry in~$\underline{k}$,
%   and let~$\underline{k}'$ be the Serre weight in dimension~$1$ given
%   by the first entry. % ~$\{k_{\sigmabar,d}-k_{\sigmabar_{d-1}}\}$, unless all of these
%   % differences are~$p-1$, in which case we set~$\underline{k}'=0$.
%   %If $k_{\sigmabar,1}-k_{\sigmabar,2}=p-1$ for all~$\sigmabar$ then we
%   %say that we are in the \emph{tr\`es ramifi\'ee case}.\mar{TG: not
%   %  sure if we want this terminology yet.} 
% Then we define the closed
%   algebraic substack~$\cX_{d,\red}^{\underline{k},\textrm{fixed}}$
%   of~$\cX_{d,\red}^{\underline{k}}$ to be the substack given by the
%   support of the sheaf
%   $H^2(G_K,\varepsilonbar^{d}\psi^{-1}_{\underline{k}'}\rhobar_t)$, so
%   that by Tate local duality, its closed points are exactly those for
%   which~$\rhobar_t$ admits~$\epsilonbar^{1-d}\psi_{\underline{k}'}$ as a
%   quotient.
%   \mar{ME: If we change the field of definition from $\Fbar_p$ to $\F$,
% 	  then ``closed point'' should be changed to ``$\Fbar_p$-point''.}
%   Since by hypothesis~$\cX_{d,\red}^{\underline{k}}$ is
%   stable under twisting, and admits a dense open substack consisting
%   of maximally nonsplit representations of niveau~$1$ and
%   weight~$\underline{k}$ (all of which are unramified twists of points
%   of $\cX_{d,\red}^{\underline{k},\textrm{fixed}}$, by definition), we
%   see that the universal unramified twist\mar{TG: introduce this
%     terminology earlier, or something similar}
%   of~$\cX_{d,\red}^{\underline{k},\textrm{fixed}}$ is equal
%   to~$\cX_{d,\red}^{\underline{k}}$. In particular, it follows
%   thatt~$\cX_{d,\red}^{\underline{k},\textrm{fixed}}$ is non-empty,
%   and it follows from Lemma~\ref{lem: dimension of Gm twist}
%   that~$\cX_{d,\red}^{\underline{k},\textrm{fixed}}$ has dimension at
%   least~$[K:\Qp]d(d-1)/2-1$. On the other hand, it follows from~(4)
%   that~$\cX_{d,\red}^{\underline{k},\textrm{fixed}}$ has dimension at
%   most~$[K:\Qp]d(d-1)/2-1$, so in fact it has dimension
%   exactly~$[K:\Qp]d(d-1)/2-1$.

%   We claim that~$\cX_{d,\red}^{\underline{k},\textrm{fixed}}$ has only
%   one irreducible component of dimension exactly~$[K:\Qp]d(d-1)/2-1$
%   (although we do not rule out that it might also have irreducible components of smaller
%   dimension).\mar{TG: I don't think this is proved right now! So I
%     guess this argument at the moment is probably good enough to get
%     everything except that we might currently have some weird extra
%     irreducible components (of the expected dimension).} To see this,
%   let~$\cY_1$ and~$\cY_2$ be two such irreducible components. By~(4),
%   there is a dense open substack of each such irreducible component,
%   over which
%   $\Hom_{G_K}(\rhobar_t,\epsilonbar^{1-d}\psi_{\underline{k}'})$ is
%   exactly $1$-dimensional. We note
%   that~$\cX_{d,\red}^{\underline{k},\textrm{fixed}}$ is essentially
%   twistable (because any fixed representation has only finitely many
%   twists which admit a fixed character as a quotient), 
%   % \mar{TG: need to introduce this notation and use it above
%     % - the idea is that if all closed points are isomorphic to at most
%     % finitely many of their twists, we still get the dimension bound -
%     % actually maybe we don't need this, as we're about to observe the
%     % existence of a dense open with a unique quotient character}
%   so
%   that given any such component, its universal unramified twist has
%   dimension~$[K:\Qp]d(d-1)/2$, and is therefore equal
%   to~$\cX_{d,\red}^{\underline{k}}$.  Accordingly, if there are two such
%   components~$\cY_1$ and~$\cY_2$, then (by removing their intersection
%   from each), we can find find dense open subsets~$\cU_1$, $\cU_2$
%   of~$\cY_1$, $\cY_2$ respectively, with the properties that there are
%   no closed points of~$\cX_{d,\red}^{\underline{k}}$ which are
%   simultaneously a twist of a point of~$\cU_1$ and~$\cU_2$. Since the
%   points which are unramified twists of points of~$\cU_1$ (resp.\
%   of~$\cU_2$) are dense in~$\cX_{d,\red}^{\underline{k}}$ (which is
%   equal to the universal unramified twist of~$\cY_1$, resp.\
%   of~$\cY_2$), this contradicts the irreducibility
%   of~$\cX_{d,\red}^{\underline{k}}$.\mar{TG: this
%     isn't literally quite correct, as it's not clear that the points
%     which are twists of~$\cU_i$ are actually open, and I don't know if
%   it's actually true. Plausibly one wants instead to show that there
%   is only one irreducible component of maximal dimension which meets
%   the maximally nonsplit locus, and work with that?}

  We now begin the inductive proof of~(2) and~(3) for~$d>1$. In fact,
  it will be helpful %\mar{TG: I just put this back in!}
  to simultaneously prove another statement~(2'), which is as follows:
  for each~$\underline{k}$, there is a closed irreducible algebraic
  substack~$\cX_{d,\red,\Fpbar}^{\underline{k},\textrm{fixed}}$
  of~$(\cX_{d,\red})_\Fpbar$ of finite presentation over~$\Fpbar$ and
  dimension~$[K:\Qp]d(d-1)/2-1$, which is generically maximally nonsplit of niveau~$1$ and
  weight~$\underline{k}$, and furthermore has the property that the
  corresponding character~$\nu_1$ is trivial. %  quotient character~$\chi_d$ is equal to 
  % ~$\epsilonbar^{1-d}\psi_{\underline{k}'}$, % that we fixed in the previous
  % % paragraph,
  % where~$\underline{k}'=\{k_{\sigmabar,1}\}$. (Note that by the definition of being maximally nonsplit
  % of weight~$\underline{k}$, $\chi_d$ is necessarily equal
  % to an unramified twist of~$\epsilonbar^{1-d}\psi_{\underline{k}'}$.) 
  % is determined by the equation
  % \[(\chi_{d}\chi_{d-1}^{-1})|_{I_K}=\varepsilonbar^{-1}\prod_{\sigmabar:k\into\Fpbar}\omega_{\sigmabar}^{k'_{\sigmabar,1}}. \]
  The discussion of the previous paragraph also proves~(2')
  when~$d=1$. The point of 
  hypothesis~(2') is that we can use Proposition~\ref{prop:dimension
    control}~(2) to
  construct~$\cX_{d,\red,\Fpbar}^{\underline{k},\textrm{fixed}}$ by using~(2)
  and~(2') in dimension $(d-1)$, and then
  construct~$\cX_{d,\red,\Fpbar}^{\underline{k}}$ from it using
  Lemma~\ref{lem: dimension of Gm twist}.

 %Recall that a closed immersion of
 % reduced algebraic stacks that are locally of finite type over $\F_p$
 % which is surjective on finite-type points is necessarily an
 % isomorphism. Since~$\cX_{d,\red}$ is an inductive limit of such
 % stacks (indeed, by Lemma~\ref{alem: underlying reduced of Ind algebraic}, we have~$\cX_{d,\red}=\varinjlim\cX^a_{d,h,s,\red}$,
%where the~$\cX^a_{d,h,s}$ are as in Section~\ref{subsec: X is Ind
%  algebraic}), we
% see that % \mar{TG: we only know that the source is finite type,
%    % not the target, but maybe we established that the target is
%    % locally of finite type in the previous section? TG: probably need
%    % to cite somewhere saying that $\cX_d$ is limit-preserving, and
%    % maybe some Ind-stack stuff from \cite{Emertonformalstacks}.} 
%% to prove~(2),
%if we produce closed algebraic substacks
%  $\cX_{d,\red,\Fpbar}^{\negligible}$ and $\cX_{d,\red,\Fpbar}^{\underline{k}}$
% of $\cX_d$,
% %satisfying the claims in
% %whose $\Fpbar$-points and dimensions are as in
% % the statement of the
% %  theorem, and
%  the union of whose $\Fpbar$-points exhausts those
%  of~$\cX_{d,\red}$, then $\cX_{d,\red}$ will be the union of  $\cX_{d,\red,\Fpbar}^{\negligible}$ and $\cX_{d,\red,\Fpbar}^{\underline{k}}$. % \mar{TG: also need the irreducibility of~
%    % $\cX_{d,\red}^{\underline{k}}$, which hopefully follows from the
%    % construction.}
%  % In fact, it suffices to construct
%  % the closed substacks~$\cX_{d,\red}^{\underline{k}}$,
%  % and then to show that the remaining
%  % $\Fpbar$-points of $\cX_{d,\red}$ are contained in a finite union of
%  % finite type algebraic substacks of dimension less than
%  % $[K:\Qp]d(d-1)/2$.

 

  % Before proving the inductive step, we introduce a hypothesis~(2')
  % which will be a useful tool in the inductive argument; namely,
  % \mar{TG: now need to
  %   actually put this into the induction. Also, we're going to need to
  % say that the characters in~$\cX_{d,\red}^{\underline{k}}$ are all
  % ranging freely, except for the tres ram extensions; need to figure
  % out where/how best to do that. }

  We now prove the inductive step, so we assume that~(2), (2') and~(3)
  hold in dimension less than~$d$. We begin by constructing
  the closed substacks~$\cX_{d,\red}^{\underline{k},\textrm{fixed}}$
  and~$\cX_{d,\red,\Fpbar}^{\underline{k}}$.  Let~$\underline{k}_{d-1}$ be the Serre weight in
  dimension~$(d-1)$ obtained by deleting the first entry
  in~$\underline{k}$. % let~$\underline{k}_d$ be the Serre weight in
  % dimension~$1$ obtained by taking the last entry in~$\underline{k}$,
  % and let~$\underline{k}'$ be the Serre weight in dimension~$1$ given
  % by the first entry.
  We
  set~$\alphabar:=\epsilonbar^{1-d}\omega_{\underline{k},1}$. If~$k_{\sigmabar,1}-k_{\sigmabar,2}=p-1$
  for all~$\sigmabar$ then we say that we are in the \index{tr\`es
    ramifi\'ee} 
  \emph{tr\`es ramifi\'ee case}, and we let~$\cT$ denote
  $\cX_{d-1,\red,\Fpbar}^{\underline{k}_{d-1},\textrm{fixed}}$; otherwise,
  we let~$\cT$ denote $\cX_{d-1,\red,\Fpbar}^{\underline{k}_{d-1}}$. In the
  latter case, it follows from our inductive
  hypothesis (more precisely, from the assumption that~(2) and (3) hold in dimension~$(d-1)$), together with Tate
  local duality, that after replacing~$\cT$ with an open substack we
  can and do assume that for each $\Fpbar$-point~$t$ of~$\cT$, we have
  $\Ext^2_{G_K}(\alphabar,\rbar_t)=0$, where~$\rbar_t$
  is the $(d-1)$-dimensional representation corresponding to~$t$. If
  we are in the tr\`es ramifi\'ee case then it follows similarly % from the inductive
  % hypothesis that~(2') holds in dimension~$(d-1)$
  that after replacing~$\cT$ with an open substack we
  can and do assume that for each $\Fpbar$-point~$t$ of~$\cT$,
  % \mar{TG: should also say that~$\rbar_T$ is maximally nonsplit, and
  %   explain why this~$\Ext^2$ is the one that we want to know about in
  % order to apply Prop \ref{prop:dimension control}.}
  $\Ext^2_{G_K}(\alphabar,\rbar_t)$ is $1$-dimensional. In either case
  we can in addition assume that~$\rbar_t$ is maximally nonsplit of
  niveau one and weight~$\underline{k}_{d-1}$.% \mar{TG: also need to
    % make sure tres ram in Steinberg case.}

  We let~$T$ be an irreducible scheme which smoothly covers~$\cT$, and
  we let~$\cX_{d,\red,\Fpbar}^{\underline{k},\textrm{fixed}}$ be the
  irreducible closed
  substack of $(\cX_{d,\red})_{\Fbar_p}$ constructed as the scheme-theoretic
image of $V$ in the notation of Proposition~\ref{prop:dimension
    control}. % , applied with~$\alphabar=\varepsilonbar^{1-d}\psi_{\underline{k}'}$.
  Part~(2) of that proposition, together with
  the inductive hypothesis, implies
  that~$\cX_{d,\red,\Fpbar}^{\underline{k},\textrm{fixed}}$ has the
  claimed % \mar{TG: need to explain why this is OK in the special case
    % of Prop \ref{prop:dimension control} for~$K=\Qp$.}
  dimension. (Note in particular that if~$K=\Qp$, then condition~\eqref{item: weird
    condition over Qp} of
  Proposition~\ref{prop:dimension control} holds. Indeed, if we are in the tr\`es
  ramifi\'ee case then condition~ \eqref{item: weird
    condition over Qp}(iii) holds, and otherwise the inductive hypothesis that
  the image of the eigenvalue morphism is dense in $(\Gm)^{d-1}_{\underline{k}_{d-1}}$
%(and hence contains a dense open subset of $(\Gm)^d_{\underline{k}}$, since this image is also constructible)
implies that condition \eqref{item: weird
    condition over Qp}(i) holds after shrinking~$T$.) We then let $\cX_{d,\red,\Fpbar}^{\underline{k}}$ be the
  substack obtained from~$\cX_{d,\red}^{\underline{k},\textrm{fixed}}$
  by twisting by unramified characters, which has the claimed
  dimension by Lemma~\ref{lem: dimension of Gm twist} (and the fact
  that by Proposition~\ref{prop:dimension control}~(2), there is a
  dense open substack of $\cX_{d,\red}^{\underline{k},\textrm{fixed}}$
  whose $\Fpbar$-points correspond to representations which are of
  niveau~$1$ and are maximally nonsplit, so that in particular the
  corresponding family is twistable). % \mar{TG: now need to finish
  % justifying the properties.}
  
By construction, we see that the
stacks~$\cX_{d,\red,\Fpbar}^{\underline{k},\textrm{fixed}}$
and~$\cX_{d,\red,\Fpbar}^{\underline{k}}$ satisfy the properties required
of % \mar{TG: should also explain slightly more why there is a dense open subset of
  % maximally nonsplit of correct weight.}
them, except possibly that in the tr\`es ramifi\'ee case, we need to check
that ~$\cX_{d,\red,\Fpbar}^{\underline{k},\textrm{fixed}}$ (and consequently
~$\cX_{d,\red,\Fpbar}^{\underline{k}}$) is generically maximally nonsplit of
niveau~$1$ and weight~$\underline{k}$. More precisely, in this case we need to check
the condition that the final extension is not just nonsplit, but is
(generically) tr\`es ramifi\'ee. This follows from
Proposition~\ref{prop:dimension control}~(3). Indeed by our
inductive assumption 
that~$\cX_{d-1,\red,\Fpbar}^{\underline{k}_{d-1}}$ is generically maximally nonsplit of niveau~$1$ and
weight~$\underline{k}_{d-1}$, and the image of the eigenvalue morphism is
dense,  we see that we are in case~(3)(b)(i) unless ~$\epsilonbar$ is trivial and
furthermore we
have~$k_{\sigmabar_2}-k_{\sigmabar_3}=p-1$ for all~$\sigmabar$, in
which case we satisfy~(3)(b)(ii).

% In
% this case, since by induction we know
% that~$\cX_{d-1,\red,\Fpbar}^{\underline{k}_{d-1}}$ is generically maximally nonsplit of niveau~$1$ and
% weight~$\underline{k}_{d-1}$, we can assume that the extension
% of~$\betabar_t=\alphabar(1)=\alphabar$ by~$\alphabar(1)=\alphabar$
% induced by~$\rhobar_t$ is tr\`es ramifi\'ee. 

% But this case follows from Lemma~\ref{lem:tres ramifiee} (taking
% the scheme $T$ there to be some integral scheme of finite type
% over $\Fbar_p$ which
% maps with dense image to $\cX_{d-1,\red,\Fbar_p}^{\underline{k}_{d-1},\textrm{fixed}}$,
% and the family $\rhobar_T$ to be the pull-back of the universal family;
% note that the hypotheses of the lemma follow by construction together 
% with our inductive hypotheses).

% (To see that~$\cX_{d,\red}^{\underline{k},\textrm{fixed}}$
% and~$\cX_{d,\red}^{\underline{k}}$ are geometrically irreducible, and
% not merely irreducible, it suffices to note that all of our
% constructions are obviously compatible with extension of scalars
% of~$E$ (and thus~$\F$) to any finite extension.)
% \mar{TG: will this do
  % for geometrically irreducible?} 
%%%
%To see this, we return to the
%end of the
%proof of Proposition~\ref{prop:dimension control}; we see (now switching to the
%notation of that proof)
%\mar{ME: This argument mixes material from the proof of Prop.~\ref{prop:dimension control}
%with the material of the current proof --- e.g.\ by its invokation of $\cX_{d-1}$ ---
%but the notation of the two arguments is in horrible conflict, e.g.\ $\rhobar$ in one
%is $\rbar$ in the other.  Can you think of an easy way to improve this, e.g.\
%by adding a statement to~\ref{prop:dimension control} that can then just
%be cited here.}
%that it is enough to show that for each~$t$,
%there is some extension of~$\alphabar$ by~$\rhobar_{f(t)}$ with the
%property that the induced extension of~$\alphabar$
%by~$\betabar_t=\alphabar(1)$ is tr\`es ramifi\'ee. We have an exact sequence
%\[\Ext^1_{G_K}(\alphabar,\rhobar_t)\to
%  \Ext^1_{G_K}(\alphabar,\betabar_t)\to\Ext^2_{G_K}(\alphabar,\rbar_t),\]so
%we are done if $\Ext^2_{G_K}(\alphabar,\rbar_t)=0$. Since~$\rbar_t$ is
%maximally nonsplit, this means that we are done unless~$\rbar_t$
%admits~$\alphabar(1)$ as its unique quotient character. By our
%inductive assumption that the image of the eigenvalue morphism is
%open, this is only possible if~$\varepsilonbar$ is trivial, and we
%have~$k_{\sigmabar_2}-k_{\sigmabar_3}=p-1$ for all~$\sigmabar$.
%
%In
%this case, since by induction we know
%that~$\cX_{d-1,\red,\Fpbar}^{\underline{k}_{d-1}}$ is generically maximally nonsplit of niveau~$1$ and
%weight~$\underline{k}_{d-1}$, we can assume that the extension
%of~$\betabar_t=\alphabar(1)=\alphabar$ by~$\alphabar(1)=\alphabar$
%induced by~$\rhobar_t$ is tr\`es ramifi\'ee. Write~$c_t$ for the
%corresponding class in~$H^1(G_K,\varepsilonbar)$. As in Remark~\ref{rem: further explanation of weird condition}, it is
%enough to show that there is a tr\`es ramifi\'ee extension
%of~$\alphabar$ by~$\alphabar(1)$, corresponding to a
%class~$d\in H^1(G_K,\varepsilonbar)$ such that the cup product of~$c_t$
%and~$d$ vanishes. If this is not the case, then~$c_t$ must be an
%unramified class in $H^1(G_K,1)=H^1(G_K,\varepsilonbar)$; but the
%extension of~$K$ cut out by~$c_t$ considered as a class in
%$H^1(G_K,1)$ is an extension given by adjoining a $p$th root of a
%uniformizer, so is in particular a ramified extension.
%So $\cX_{d,\red,\Fpbar}^{\underline{k},\textrm{fixed}}$ is generically maximally
%nonsplit of niveau~$1$ and weight~$\underline{k}$, as required.






% We say that a
%   set of representations $\rhobar:G_K\to\GL_d(\Fpbar)$ is \emph{negligible}
%   if it is contained in a closed finite type algebraic substack of dimension
%   less than $[K:\Qp]d(d-1)/2$; so for example it follows from the
%   above discussion that the set of irreducible representations is
%   negligible. We will now use the inductive hypothesis to show that in
%   fact the set of all representations which do not correspond to
%   points of any~$\cX_{d,\red}^{\underline{k}}$ is negligible, thereby
%   completing the proof of~(2).
  % \mar{TG: prove something stronger,
    % namely that removing this gives the claimed open substacks of the $\cX_{d,\red}^{\underline{k}}$.}

  % Choose for each~$a<d$ a (finite) set of irreducible
  % representations~$\alphabar:G_K\to\GL_a(\Fpbar)$, one for each
  % equivalence class of twists by unramified characters, extending the
  % choice we already made if~$a=1$.
 To complete the proof of~(2), we need only
construct~$\cX_{d,\red,\Fpbar}^{\negligible}$. By
  Proposition~\ref{prop:dimension control}, Tate local duality, upper
  semi-continuity of the fibre dimension, and the inductive
  hypothesis, we see that for each~$1\le i\le d$, each~$\alphabar_i$
  (of dimension~$a_i$, say),  and each~$s\ge 0$
  % \mar{TG: do we need to spell this out more? Maybe this should be
  %   pulled out as a lemma about the dimension of the space of
  %   extensions of an irreducible by a family?}
  there is a finitely presented closed algebraic substack $\cX_{s,\alphabar_i,\Fpbar}$
  of~$(\cX_{d,\red})_\Fpbar$, whose $\Fpbar$-points contain all the
  representations of the form
  $0\to\rhobar_{d-{a_i}}\to\rhobar\to\alphabar_i\to 0$ for which
  $\dim_{\Fpbar}\Ext^2_{G_K}(\alphabar_i,\rhobar_{d-{a_i}})=s$, and whose
  dimension is at most 
% \mar{ME: Should we include the ceiling in the notation? TG: I don't
%   think so, but we can discuss -- actually the claim about $s=1$ when
%   $a=1$ is confusing! TG: OK, I think we should have done this, so
%   hopefully this is better. It currently seems to me that the
%   statement shouldn't change, but that what was written before was
%   false if~$a=1$, and we really want the ceiling. But this should be
%   checked!}
\[[K:\Qp](d-{a_i})(d-{a_i}-1)/2-\lceil
  s(({a_i}^2+1)s-{a_i})/2\rceil+[K:\Qp]{a_i}(d-{a_i})+s-1; \] furthermore, the locus
where~$\rhobar_{d-{a_i}}$ is an $\Fpbar$-point
of~$\cX_{d-{a_i},\red,\Fpbar}^{\negligible}$ is of dimension strictly less
than this. These stacks are only nonzero for finitely many values
of~$s$. For fixed~${a_i}$, we see that as a function of~$s$, this quantity
is maximised by~$s=0$, as well as by~$s=1$ when ${a_i}=1$. (To see this, we have to maximise
the quantity $s-\lceil s(({a_i}^2+1)s-{a_i})/2\rceil$. Suppose firstly that
${a_i}>1$. Then if~$s=0$ we have $0$, while if~$s>0$ we
have
$s-\lceil s(({a_i}^2+1)s-{a_i})/2\rceil\le s-s(({a_i}^2+1)s-{a_i})/2\le
s-s({a_i}^2+1-{a_i})/2\le s-3s/2<0$. Meanwhile if~${a_i}=1$, then for~$s=0$ we
have $0$, for~$s=1$ we have $1-\lceil 1/2\rceil=1$,
while for ~$s>1$ we have
$s-\lceil s(2s-1)/2\rceil\le s-s(2s-1)/2 \le s-3s/2<0$.) It follows that
  as a function of~${a_i}$ the bound is maximised at ${a_i}=1$ and $s=0$
  or~$1$, when it is equal to $[K:\Qp]d(d-1)/2-1$, and it is otherwise
  strictly smaller.

  By Lemma~\ref{lem: dimension of Gm twist}, it follows that the locus
  in~$(\cX_{d,\red})_\Fpbar$ of representations of the form
  $0\to\rhobar_{d-{a_i}}\to\rhobar\to\alphabar'\to 0$, with~$\alphabar'$
  an unramified twist of~$\alphabar_i$ for which
  $\dim_{\Fpbar}\Ext^2_{G_K}(\alphabar',\rhobar_{d-{a_i}})=s$, is of
  dimension at most~$[K:\Qp]d(d-1)/2$, with equality holding only if
  ${a_i}=1$ and $s=0$ or~$1$. Furthermore, the locus of those
  representations for which ~$\rhobar_{d-{a_i}}$ is an $\Fpbar$-point
  of~$\cX_{d-{a_i},\red,\Fpbar}^{\negligible}$ is of dimension strictly less
  than~$[K:\Qp]d(d-1)/2$. Putting this together, we see that~(2) holds
  in dimension~$d$ if we take~$\cX_{d,\red,\Fpbar}^{\negligible}$
to be the union of the twists by unramified characters of the
substacks~$\cX_{s,\alphabar_i}$ for which $\dim\alphabar_i>1$ or~$s>1$, together
with the union of the twists by unramified characters of the
substacks of the~$\cX_{s,\alphabar_i,\Fpbar}$ for which~$\dim\alphabar_i=1$, $s=0$
or~$1$, and~$\rhobar_{d-1}$ is an $\Fpbar$-point
  of~$\cX_{d-1,\red,\Fpbar}^{\negligible}$.%   .\mar{TG:
  % presumably also need extensions of characters by things in~$\cX_{d-1,\red}^{\negligible}$?}
% \mar{TG: my feeling is that we need to write slightly more, to justify
% that when we constructed things by scheme-theoretic closures, we have
% the things of niveau 1 that aren't maximally nonsplit (and maybe we
% also need to justify that we have all the maximally nonsplit ones, too?).}
% It follows in particular that the only representations which could
% fail to be negligible are\mar{TG: I think it should be trivial to
%   finish from here, hopefully, although I'm mildly confused as to
%   whether it is literally trivial or actually needs an argument!}

  % It follows in particular that which are not of niveau~$1$ are
  % negligible. More precisely, we see that those representations which
  % are not an extension of a character by a maximally nonsplit
  % niveau~$1$ representation are negligible. % \mar{TG: in fact we should
  %   % have managed to show that the ones that are generically of }
  % To see that the 
  % representations of this form which are not themselves maximally nonsplit of
  % niveau~$1$ are also negligible, % \mar{TG: I don't think we can quite
  %   % do this with the machinery we have set up right now, annoyingly -
  %   % I think we need to either have something that lets us bound
  %   % representations which are extensions of~$\chi_1\oplus\chi_2$, or
  %   % maybe we make an argument via considering the projection map
  %   % to~$\cX_2$. Leaving this for now, as I'm sure it can in any case
  %   % be handled by the kinds of arguments we are making here.}

  % Let~$\rhobar$ be maximally nonsplit of niveau~$1$. We can
  % write~$\rhobar$ as an extension of a character~$\chi_d$ by a
  % representation~$\rhobar_{d-1}$, which in turn admits a unique
  % $1$-dimensional quotient~$\chi_{d-1}$. Then by Tate local duality,
  % $\Ext^2_{G_K}(\chi_d,\rhobar_{d-1})$ had dimension at most~$1$, with
  % equality holding if and only if
  % $\chi_d\chi_{d-1}^{-1}=\epsilonbar^{-1}$.
  % \mar{TG: now just need to
    % justify irreducibility of the corresponding families of
    % extensions; perhaps again this should be pulled out as a separate lemma?}

Finally we prove~(3) in dimension~$d$. %  by induction on~$d$. We assume that~(3)
% holds in dimension less than~$d$.
% As already explained,
        % we assume that Theorem~\ref{thm:reduced dimension} holds for dimensions
	% less than or equal to $d$,
	% and that the present theorem holds for dimensions less than $d$.
	In the case $r = 0$, there is nothing to prove.
	If $\dim \Hom_{G_K}(\rhobar,\alphabar) \geq r \geq 1,$
	then we may place $\rhobar$ in a short exact sequence
	$$0 \to \thetabar \to \rhobar \to \alphabar^{\oplus r} \to 0,$$
	where $\thetabar$ is of dimension $d- ra < d$.
        We may apply  part~(2) %Theorem~\ref{thm:reduced dimension}
		so as to find that $\cX_{d-ar,\red,\Fpbar}$ 
has dimension at most $[K:\Q_p](d-ar)(d-ar-1)/2$.
		Let $\cU_s$ be the locally closed substack %  \mar{TG:
                  % need to justify that this is locally closed}
		of $\cX_{d-ar,\red,\Fpbar}$ over which
		$\dim H^2(G_K,\thetabar\otimes \alphabar^{\vee}) =
                s$; % \mar{TG: presumably the idea is to stratify by~$s$,
        %         and note that the argument below works for all~$s$? ME: Yes,
	% that's right. TG: should reference elsewhere in the paper for
        % this being locally closed; also the $\cW$ notation currently conflicts. ME: Changed notation to $\cU_s$ to eliminate conflict.}
		by the inductive hypothesis, this locus has dimension at most
		$[K:\Qp](d-ar)(d-ar-1)/2-s\bigl((a^2+1)s-a\bigr)/2$, %  in $\cX_{d-a
                  % r,\red},$\mar{TG: here}
		and over this locus
		we may construct a universal family of extensions 
		$$0 \to \thetabar \to \rhobar_{\cU_s}
		\to \alphabar^{\oplus r} \to 0.$$
		The locus of $\rhobar$ we are interested in is
		contained in the scheme-theoretic
		image of this family in $(\cX_{d,\red})_\Fpbar$,
		and Proposition~\ref{prop:dimension control}
		shows that this scheme-theoretic image 
		has dimension bounded above by
		%relative dimension $r([K:\Q_p](d-r) + s)- r^2$
		%\mar{ME: Have to explain this above.  The point is that
		%	over this locus, $H^2$ has constant rank $s$ and so
		%	is locally free, and then it should be true
		%	that $H^1$ is generically free of 
		%	rank $r([K:\Q_p](d-r) + s)$, which
		%	lets us compute the dimension of the universal
		%	extension.   Maybe have to think more carefully
		%	about what happens when $H^0 \geq 0$ ,
		%	especially for $K$ where cyclo is trivial,
		%	but I think it works out.}
		%over $\cW_s$, and thus of total dimension at most
                % $$[K:\Q_p](d-ar)(d-ar-1)/2 -s((a^2+1)s-a)/2 +r([K:\Qp]a(d-ar)+s) -r^2 .$$
                % \mar{TG: check this with the
                  % new notation}
                % Using the fact that $\cX_{d,\red}$ is
                % equidimensional of dimension $[K:\Q_p]d (d-1)/2,$ we
                % find that the image of this family in $\cX_{d,\red}$
                % has codimension at least
		% \begin{multline*} 
                \[  [K:\Q_p](d-ar)(d-ar-1)/2 -s((a^2+1)s-a)/2 +r([K:\Qp]a(d-ar)+s) -r^2\]
\begin{align*}
	= &[K:\Qp]d(d-1)/2- (r(a^2+1)r-a)/2-(r-s)^2/2-(ar(ar-1))/2\\
\le &[K:\Qp]d(d-1)/2- (r(a^2+1)r-a)/2.
 \end{align*}

% \dfrac{[K:\Q_p]ar(ar-1) + s\bigl((a^2+1)s-1\bigr)}{2} + r^2 - rs 
% 		\\
% 		\geq \dfrac{r\bigl((a^2+1)r - a\bigr) + (r-s)^2 + as(as-1)}{2} \geq
% 		\dfrac{r\bigl((a^2+1)r-a\bigr)}{2}.
                % \end{multline*}
		% Since this codimension is an integer it is in fact
		% at least $\lceil r\bigl((a^2+1)r-a\bigr)/2 \rceil.$
                Since this conclusion holds for each of the finitely
                many values of~$s$ (and since the dimension is an
                integer, allowing us to take the floor of this upper
                bound), we are done.
	% 	\mar{ME: The displayed equality is just a computation;
	% 		the $r^2$ comes from the $GL_r$ action on
	% 		$\chibar^{\oplus r}$ as automorphisms.  
	% 		The lower bound $c_r$ for the codimension
	% 		is then supposed to be worked out/chosen
	% 		so that the claimed inequality is true! ME: OK,
	% 	I made choice of $c_r$ and wrote down somehopefully
	% correct computations. TG: they seem correct to me.}
\end{proof}



%We next deduce the following corollary.

\begin{cor}\label{cor:Xd is formal algebraic}
	$\cX_d$ is a  Noetherian formal algebraic
        stack.% \mar{TG: can probably upgrade this to Noetherian, i.e.\
          % qcqs and locally Noetherian; note that this notion doesn't
          % seem to be defined in \cite{Emertonformalstacks} right now,
          % though. TG: yes it is! Commenting out.}
\end{cor}
\begin{proof}
	Theorem~\ref{thm:Xdred is algebraic} shows that
	$\cX_{d,\red}$ is an algebraic stack that 
	is of finite presentation over $\F$ (and hence
	quasi-compact and quasi-separated).  
	Together with Proposition~\ref{prop:X is an Ind-stack},
	which shows that $\cX_d$ is isomorphic to the inductive limit
	of a sequence of finitely presented algebraic stacks with respect to transition
	morphisms that are closed immersions, 
	% \mar{ME: Make sure that that proposition is formulated
	% 	so that this summary  is valid.}
	this implies that $\cX_d$
      is indeed a (locally countably indexed) formal algebraic stack,
      by Proposition~\ref{prop: formal stack 6.6}. By
      Remark~\ref{rem:Ind-loc. f.t. criterion}, $\cX_d$ is Ind-locally
of finite type over $\cO$,
%By Lemma~\ref{lem:X
%        is limit preserving}, $\cX_d$ is limit preserving.
%(Of course, this also follows from it begin Ind-locally of finite type.)
%Since~$\cX_d$ is locally Ind-finite type over~$\F$,
and it then follows from
      Proposition~\ref{prop: Xd has a good obstruction theory}
      and Theorem~\ref{thm:noetherian criterion} that~$\cX_d$ is
      locally Noetherian, and hence Noetherian (as we have already
      seen that it is quasi-compact and quasi-separated).
%\mar{TG: just
%        flagging that there is a locally Ind-finite type hypothesis in the
%        statement of Theorem~\ref{thm:noetherian criterion} which we
%        don't seem to be addressing here.}
      \end{proof}

%      \mar{ME: Maybe add Noetherianness here, by invoking Herr complex obstruction theory.}
\chapter{Crystalline lifts and the finer structure of $\cX_{\lowercase{d,\red}}$}\chaptermark{Crystalline lifts and $\cX_{\lowercase{d,\red}}$}
\label{sec:properties}
%\chapter{Crystalline deformation rings}\label{sec: crystalline deformation rings}
%\mar{ME: Should we merge this section with the following one? ME: Done}
Up to this point, we have shown that $\cX_d$ is
a Noetherian formal algebraic stack. 
In this chapter,
we will additionally show that $(\cX_d)_\red$ is
equidimensional of dimension
$[K:\Q_p]d(d-1)/2$, and enumerate its irreducible components. 
In the course of doing this, we will
also prove a result on the existence of crystalline lifts
which is of independent interest. We furthermore determine the closed points
of~$(\cX_d)_\red$, and describe the maximal substack of~$\cX_d$ over which the
universal $(\varphi,\Gamma)$-module gives rise to a continuous $G_K$-representation.% \mar{TG: should also mention results
  % on closed points, and the substack of actual Galois representations.}
% \mar{TG: I commented out two
  % subsections here, at least for now.}

% \section{Extending the $G_{K_{\infty}}$-action}
% %\mar{ME: This subsection still needs editing/rearranging.}

% Let $E$ be the Eisenstein polynomial corresponding to our fixed
% uniformiser~$\pi$ of~$K$.

% \begin{defn}%\mar{TG: add Fargues}
% 	Let $A$ be a finite type $\cO/\varpi^a$-algebra
% 	for some $a \geq 1$.
% We define a projective Breuil--Kisin module 
% (resp.\ a Breuil--Kisin--Fargues module) of height at
%   most~$h$ with $A$-coefficients to be a finitely generated projective $\gS_A$-module
%   (resp.\ $\AAinf{A}$-module)
%   $\gM$, equipped with a~$\varphi$-semi-linear morphism
%   $\varphi:\gM\to\gM$, with the property that the corresponding
%   morphism % of $\AAinf{A}$-modules
%   $\Phi_{\gM}:\varphi^*\gM\to\gM$ is
%   injective, with cokernel killed by~$E^h$.
% \end{defn}

% %\begin{remark}
% %	% \mar{ME: Note that projective Breuil--Kisin modules over $\gS_A$ haven't 
% %	% 	actually been defined.}
% %  If~$\gM$ is a projective
% %  Breuil--Kisin module over~$\gS_A$, of height at most~$h$ (for some
% %  $h \geq 0$),
% %  then
% %$\gMt:=\AAinf{A}\otimes_{\gS_A}\gM$
% %  is a projective
% %  Breuil--Kisin module over~$\AAinf{A}$, of height at most~$h$.
% %Since~$\gM$ is a
% %flat~$\gS_A$-module, we may regard~$\gM$ as a
% %$\gS_A$-submodule of~$\gMt$.
% %\end{remark}

% For any~$h\ge 1$, we write~$\cC_{\BK,d,h}$ for the moduli stack of
% rank~$d$ Breuil--Kisin modules of height at most~$h$.

% For each~$s\ge 1$, write~$K_s=K(\pi^{1/p^s})$. We now use a variant of the arguments of ~\cite{MR2745530}, which show
% that the action of~$G_{K_\infty}$ on an \'etale $\varphi$-module
% over~$\OEA$ can be extended to some~$G_{K_s}$. We begin with some
% preliminary lemmas, the first of which is an analogue for Breuil--Kisin modules
% of~\cite[Lem.\ 5.2.14]{EGstacktheoreticimages}, and is proved in a
% similar way.

% \begin{lem}
%   \label{lem: adding projective Kisin to get free}Let~$A$ be an
%   $\cO$-algebra, and let~$\gM$ be a finite projective Breuil--Kisin module with $A$-coefficients of height at
%   most~$h$. Then~$\gM$ is a direct summand of a finite free Breuil--Kisin
%   module over with $A$-coefficients of height at
%   most~$h$.
% \end{lem}
% \begin{proof}
%   Since the map $\Phi_{\gM}:\varphi^*\gM\to\gM$ is an injection with
%   cokernel killed by~$E^h$, there is a
%   map~$\Psi_{\gM}:\gM\to\varphi^*\gM$ with $\Psi_{\gM}\circ\Phi_{\gM}$
%   and $\Phi_{\gM}\circ\Psi_{\gM}$ both being given by multiplication
%   by~$E^h$. Let~$\gP$ be a finite projective $\gS_A$-module
%   such that~$\gF:=\gM\oplus\gP$ is a finite free~$\gS_A$
%   module. Set~$\gQ:=\gP\oplus\gM\oplus\gP$, a finite
%   projective~$\gS_A$-module.

%   Note that~$\gM\oplus\gQ\cong\gF\oplus\gF$ is a finite
%   free~$\gS_A$-module, so it suffices to show that~$\gQ$ can
%   be endowed with the structure of a Breuil--Kisin module of height at
%   most~$h$. Since~$\gF$ is free, we can choose an
%   isomorphism~$\Phi_{\gF}:\varphi^*\gF\isoto\gF$, and we then
%   define~$\Phi_{\gQ}:\varphi^*\gQ\to\gQ$ as the composite
%   \begin{align*}
%     \varphi^*\gQ
%     &=\varphi^*\gP\oplus\varphi^*\gM\oplus\varphi^*\gP\\&\isoto
%     \varphi^*\gM\oplus\varphi^*\gP\oplus\varphi^*\gP\\&=\varphi^*\gF\oplus\varphi^*\gM\\&\stackrel{\Phi_{\gF}}{\to}\gF\oplus\varphi^*\gP\\&\isoto
%     \gP\oplus\gM\oplus\varphi^*\gP\\&\stackrel{\Psi_{\gM}}{\to} \gP\oplus\varphi^*\gM\oplus\varphi^*\gP\\&=\gP\oplus\varphi^*\gF\\&\stackrel{\Phi_{\gF}}{\to}\gP\oplus\gF=\gQ.
%   \end{align*}Every morphism in this composite other than~$\Psi_{\gM}$
%   is an isomorphism. Since
%   $\Psi_{\gM}\circ\Phi_{\gM}=\Phi_{\gM}\circ\Psi_{\gM}=E^h$, we see
%   that~$\Psi_{\gM}$ is an injection with cokernel killed
%   by~$E^h$, so the same is true of~$\Phi_{\gQ}$, as required.
% \end{proof}


% The following lemma is proved by a standard Frobenius amplification
% argument, which is in particular almost identical to the proof
% of~\cite[Lem.\ 4.1.9]{CEGSKisinwithdd}.
% \begin{lem}
%   \label{lem: Frobenius amplification lemma over Ainf}Let~$\gM$, $\gN$
%   be projective Breuil--Kisin--Fargues modules  of height at
%   most~$h$, where $A$ is an
%   $\cO/\varpi^a$-algebra. Suppose that~$N\ge e(a+h)/(p-1)$. Let \[f:\gM\to
%     \gN\] be an $\AAinf{A}$-linear map,
%   and suppose that for all~$m\in\gM$,
%   $(f\circ\Phi_{\gM}-\Phi_{\gN}\circ\varphi^*f)(m)\in
%   u^{pN}\gN$.%  \emph{(}where we
%   % write~$\Phi_{\gM},\Phi_{\gN}$ for their $\AAinf{A}$-linear
%   % extensions\emph{)}.
%   % \mar{TG:
%     % sort out precise value for $N$}

%   Then there is a unique $\AAinf{A}$-linear, $\varphi$-linear morphism
%   % \mar{ME: Probably the correct adjective for $f$ 
%   %         is just $\varphi$-linear, rather than semi-linear, right?}
%   $f':\gM\to
%     \gN$ with the property that for all~$m\in\gM$,
%   $(f'-f)(m)\in u^N\gN$; in fact, we have $(f'-f)(m)\in u^{N+1}\gN$.
% \end{lem}
% \begin{proof}We firstly prove uniqueness. Suppose that~$f''$ also
%   satisfies the properties that~$f'$ does, and write~$g=f'-f''$, so that by
%   assumption we have
%   that~$\im g\subseteq u^N \gN$ and
%   $g\circ\Phi_{\gM}=\Phi_{\gN}\circ\varphi^*g$. We claim that this implies
%   that~$\im g\subseteq u^{N+1} \gN$; if
%   this is the case, then by induction on~$N$ we have
%   $\im g\subseteq u^N \gN$ for all~$N$,
%   and so~$g=0$ (note that since~$\gN$ is a finite
%   projective~$\gS_A$-module,
%   $\gN$ is $u$-adically complete and separated).

%   To prove the claim, note that
%   since~$\im g\subseteq u^N \gN$ we
%   have~$\im \varphi^*g\subseteq u^{pN}
%   \gN$. By~\cite[Lem.\
%   5.2.6]{EGstacktheoreticimages} (and its proof), and the assumption
%   that~$\gM$ has height at most~$h$, we
%   have~$\im \Phi_{\gM}\supseteq u^{e(a+h-1)}\gM$.  Since
%   $g\circ\Phi_{\gM}=\Phi_{\gN}\circ\varphi^*g$, it follows that
%   \[u^{e(a+h-1)}\im g\subseteq\im g\circ\Phi_{\gM}= \im \Phi_{\gN}\circ\varphi^*g
%     \subseteq u^{pN} \gN\]and since
%   $\gN$ is $u$-torsion free, we have
%   $\im g\subseteq u^{pN-e(a+h-1)}
%   \gN$. Our assumption on~$N$ implies
%   that $pN-e(a+h-1)\ge N+1$, as required.

%   We now prove the existence of~$f'$. For any~$\gS_A$-linear
%   map $h: \gM\to
%   \gN$ we
%   set~$\delta(h):=\Phi_{\gN}\circ\varphi^*h-h\circ\Phi_{\gM}:\varphi^*\gM\to\gN$. We
%   claim that for
%   any~$s:\varphi^*\gM\to
%   u^{pN}\gN$, we can find~$t:\gM\to
%   u^{N+1}\gN$ with~$\delta(t)=s$. Given this claim, the existence of~$f'$ is immediate,
% taking~$s=\delta(f)$ and~$f':=f-t$, so that~$\delta(f')=0$.


% To prove the claim, we first prove the weaker claim that for
% any~$s$ as above we can
% find~$h:\gM\to
% u^{N}\gN$
% with~$\im(\delta(h)-s)\subseteq u^{p(N+1)}$. % \mar{TG: getting confused
%   % about precise numbers here!}
% Admitting this second claim, we prove
% the first claim by successive approximation: we set~$t_0=h$, so
% that~$\im(\delta(t_0)-s)\subseteq u^{p(N+1)}$. Applying the second
% claim again with~$N$ replaced by~$N+1$, and $s$ replaced
% by~$s-\delta(t_0)$, we
% find~$h:\gM\to
% u^{N+2}\gN$
% with~$\im(\delta(h)-(s-\delta(t_0)))\subseteq u^{p(N+2)}$. % \mar{TG:
%   % confused again, sort out exponents!}.
% Setting~$t_1=t_0+h$, and
% proceeding inductively, we obtain a Cauchy sequence converging (in
% the~$u$-adically complete finite $\AAinf{A}$-module
% $\Hom_{\AAinf{A}}(\gM,\gN)$)
% to the required~$t$.

% Finally, we prove this second
% claim.
% If~$h:\gM\to
% u^{N+1}\gN$ then 
% % \mar{ME: I agree that mapping into $u^{N+1}\gN$ is even better than
% % 	mapping into $u^N\gN$, but doesn't this mean that all the
% % 	exponents above can be improved by $1$? TG: possibly; but I
% %         got pretty confused about exponents at various points when
% %         writing it, so I don't want to just change one here. If you've
% %       checked it all and can see what they should all be then feel fee
% %     to improve! Otherwise I would be happy enough to just change the
% %     $N+1$ to $N$ here, to be honest, as the actual bounds are unimportant.}
% $\im\varphi^*h\subseteq u^{p(N+1)}$, so it suffices to show that we can
% find~$h$ such that $h\circ\Phi_{\gM}=-s$; but this is immediate,
% because the cokernel of~$\Phi_{\gM}$ is killed by~$u^{e(a+h-1)}$, and
% $pN-e(a+h-1)\ge N+1$.% \mar{TG: should be easy to get exponents
%   % straight; come back to this.}
% \end{proof}

% \begin{remark}
% 	% \mar{ME: Note that projective Breuil--Kisin modules over $\gS_A$ haven't 
% 	% 	actually been defined.}
%   If~$\gM$ is a projective
%   Breuil--Kisin module with $A$-coefficients, of height at most~$h$ (for some
%   $h \geq 0$),
%   then
% $\gMt:=\AAinf{A}\otimes_{\gS_A}\gM$
%   is a projective
%   Breuil--Kisin--Fargues module with $A$-coefficients, of height at most~$h$.
% Since~$\gM$ is a
% flat~$\gS_A$-module, we may regard~$\gM$ as a
% $\gS_A$-submodule of~$\gMt$.
% \end{remark}

% The following lemma is proved by a reinterpretation of some of the arguments
% of~\cite[\S2]{MR2745530}.  
% %If~$\gM$ is a projective Kisin module
% %over~$\gS_A$, we write
% %$\gMt:=\AAinf{A}\otimes_{\gS_A}\gM$. Since~$\gM$ is a
% %flat~$\gS_A$-module, we may regard~$\gM$ as a
% %$\gS_A$-submodule of~$\gMt$.

% \begin{lem}\label{lem: Caruso Liu Galois action on Kisin}
%   For any fixed $a,h$, and any $N\ge e(a+h)/(p-1)$, there is a
%   positive integer~$s(a,h,N)$ with the property that for any
%   $\cO/\varpi^a$-algebra~$A$, any
%   projective Breuil--Kisin
%   module  of height at most~$h$, and 
%   any~$s\ge s(a,h,N)$, there is a unique continuous action of~ $G_{K_s}$ on
%   $\gMt$ which commutes with~$\varphi$
%   and is semi-linear with respect to the natural action of~$G_{K_s}$
%   on~$\AAinf{A},$ with the additional property that for
%   all~$g\in G_{K_s}$ we have $(g-1)(\gM)\subset
%   u^N\gMt$. % \mar{TG: should probably
%     % explain how we are thinking of~$\gM$ as living inside $\AAinf{A}\otimes_{\gS_A}\gM$}
% \end{lem}
% \begin{proof}By
%   definition, to give a semi-linear action of~$G_{K_s}$  on~$\gMt$ is
%   to give for each~$g\in G_{K_s}$ a morphism of Breuil--Kisin--Fargues
%   modules \[\beta_g:g^*\gMt\to\gMt\] with the property that for
%   all~$g,h\in G_{K_s}$, we have $\beta_{gh}=\beta_h\circ
%   h^*\beta_g$. % \mar{TG: I think this is the right way round, but it
%     % would be sensible to check for nonsense.}
%   Now, the requirement that
%   $(g-1)(\gM)\subseteq u^N\gMt$ implies that if we have two such
%   morphisms~$\beta_g,\beta'_g$ then
%   $(\beta_g-\beta'_g)(g^*\gMt)\subseteq u^N\gMt$, so
%   that~$\beta'_g=\beta_g$ by the uniqueness assertion of
%   Lemma~\ref{lem: Frobenius amplification lemma over Ainf}.

%   We now use the existence part of Lemma~\ref{lem: Frobenius
%     amplification lemma over Ainf} to construct the required
%   action. Since the injection $\gS_A\into\AAinf{A}$ takes~$u$ to the
%   Teichm\"uller lift of~$(\pi^{1/p^n})_n$, % \mar{TG: have we actually
% %     discussed this above? I suspect that we use the injectivity of
% %     various coefficient rings into $W(\C^\flat)_A$ throughout the
% %     paper - do we prove it? ME: The faithful flatness result
% % in \S 3 on Galois descent gives injectivy for perfect rings.  Hopefully
% % there is an argument that shows that the non-perfect rings embed into
% % (rather than just map to) the perfect ones.}
% we can
%   and do choose~$s(a,h,N)$ sufficiently large that for all~$s\ge
%   s(a,h,N)$ and~$g\in
%   G_{K_s}$, we have~$g(u)-u\in u^{pN}\AAinf{A}$; thus for
%   all~$\lambda\in \gS_A$, we have $g(\lambda)-\lambda\in
%   u^{pN}\AAinf{A}$. % \mar{TG: need to choose~$N$
%     % and of course need to justify that~$s$ exists, but I think this
%     % should be just the continuity of the action of $G_K$ on~$\Ainf$.}

% We claim that for each~$g\in G_{K_s}$, we may define
%   an~$\AAinf{A}$-linear map  \[\alpha_g:g^*\gMt\to
%     \gMt\] with the following two properties:
% \begin{itemize}
% 	\item
% $(\alpha_g\circ\Phi_{g^*\gMt}-\Phi_{\gMt}\circ\varphi^*\alpha_g)(\varphi^*g^*\gMt)\subseteq u^{pN}\gMt.$
% \item
% For each $m \in \gM,$ we have $\alpha_g(1\otimes m) - m \in u^N \gMt.$
% \end{itemize}
% It suffices to construct such maps in the case when~$\gM$ is free;
% indeed, by Lemma~\ref{lem: adding projective Kisin to get free},
% %we can and do assume from now on that~$\gM$ is free; indeed, in general
%   we may write~$\gF=\gM\oplus\gP$ where~$\gF,\gP$ are respectively free and
%   projective Breuil--Kisin modules of height at most~$h$, and given a
%   morphism~$\alpha_g:g^*\gF^{\inf}\to\gF^{\inf}$ satisfying our requirements,
%   we may obtain the desired
%   morphism~$\alpha_g:g^*\gMt\to\gMt$ by projection onto the
%   corresponding direct summands.
%   % \mar{ME: Do we also have to see that this construction preserves the
%   %         cocyle property?  Or is the cocycle property a consequence
%   %         of the uniquess?  If so, we should note this above! TG: see
%   %         next marginal comment.}

% If then~$\gF$ is a free Breuil--Kisin module of height at most~$h$, we choose a
% basis~$f_1,\dots,f_d$ for~$\gF$, and then, for each~$g\in G_{K_s}$, we define
%   an~$\AAinf{A}$-linear map  \[\alpha_g:g^*\gFt\to
%     \gFt\] (which depends on our choice of basis) by
%   %
%   % (Indeed, in general we we may
%   % write~$\gF=\gM\oplus\gP$ where~$\gF,\gP$ are respectively free and
%   % projective BK modules of height at most~$h$. If we prove that there
%   % is a canonical action of~$G_{K_s}$ on
%   % $\AAinf{A}\otimes_{\gS_A}\gF[1/u]$, then applying the
%   % projector to $\AAinf{A}\otimes_{\gS_A}\gM[1/u]$, we obtain
%   % an action of~$G_{K_s}$ on
%   % $\AAinf{A}\otimes_{\gS_A}\gM[1/u]$. To see that this action
%   % is independent of the choice of~$\gF$ and thus canonical,)\mar{TG:
%   %   have to figure out in what sense the action on free modules is
%   %   canonical, I guess; hopefully this isn't hard. It's presumably
%   %   related to checking that the construction that we make below in
%   %   the free case is independent of the choice of basis? In fact I
%   %   think it should just characterised by demanding that
%   %   for all $g\in G_{K_s}$, we have
%   %   $(g-1)(\gM)\subset u^N\AAinf{A}\otimes_{\gS_A}\gM$,
%   %   because if we have two different actions that both satisfy this,
%   %   then applying their difference will yield a $\varphi$-equivariant
%   %   map $g^*\AAinf{A}\otimes_{\gS_A}\gM\to
%   %   u^N\AAinf{A}\otimes_{\gS_A}\gM$ which will vanish for the
%   %   usual reasons.}
%   %
%   %
%   % Choose a basis~$m_1,\dots,m_d$ for~$\gM$ as an~$\AKplus$-module, and
%   % for each element~$g\in G_{K_s}$,
%   % define the $\AAinf{A}$-linear map % \mar{TG: should probably justify
%   %   % that the LHS is actually a Kisin module of height at most~$h$.}
%   % \[\alpha_g:g^*\AAinf{A}\otimes_{\gS_A}\gM\to
%   %   \AAinf{A}\otimes_{\gS_A}\gM\]
%   %
%   \[\alpha_g(\sum_i\lambda_i\otimes f_i)=\lambda_if_i.\] We now check
%   that~$\alpha_g$ has the claimed properties.
%   If~$\Phi_{\gF}(1\otimes f_i)=\sum_j\theta_{i,j}f_j$ (so
%   that~$\theta_{i,j}\in\gS_A$), then 
%   for any $\sum_i\lambda_i\otimes
%     f_i\in\varphi^*g^*\gFt, $ we compute that
%     \[(\alpha_g\circ\Phi_{g^*\gFt}-\Phi_{\gFt}\circ\varphi^*\alpha_g)(\sum_i\lambda_i\otimes
%     f_i)=\sum_{i,j}\lambda_i(g(\theta_{i,j})-\theta_{i,j})f_j\in
%     u^{pN}\gFt,\] 
%   while for any $\sum_i \lambda_i f_i \in \gF$
%   (note then that each $\lambda_i \in \gS_A$),
%   we compute that
%     $$\alpha_g\bigl(1\otimes (\sum_i \lambda_i f_i)\bigr)
%     - \sum_i \lambda_i f_i
%     = 
% 				 \alpha_g(\sum_ig(\lambda_i)\otimes
%                                  g_i)-\sum_i\lambda_ig_i = 
% 				 \sum_i(g(\lambda_i)-\lambda_i)g_i
% 				 \in u^N \gFt;
% 				 $$
%   in both computations we have used the fact that, by our choice of~$s$,
%   if~$\lambda\in\gS_A$, 
%   then~$g(\lambda)-\lambda\in u^{pN}\gFt \subset u^N\gFt$.

%  Returning to the general case of a projective Breuil--Kisin
%  module~$\gM$ of height at most~$h$, it follows from % our
%   % assumption on~$s$, and
%   Lemma~\ref{lem: Frobenius amplification lemma
%     over Ainf} that there is a unique % $\AAinf{A}$-linear,
%   % $\varphi$-semi-linear 
% morphism of Breuil--Kisin--Fargues modules
%   \[\beta_g:g^*\gMt\to
%  \gMt\]
%   with~$\im(\beta_g-\alpha_g)\subseteq
%   u^N\AAinf{A}\otimes_{\gS_A}\gM$. Since we have
%   $\alpha_{gh}=\alpha_h\circ h^*\alpha_g$, it follows from this
%   uniqueness that~$\beta_{gh}=\beta_h\circ h^*\beta_g$, % \mar{TG: does
%     % this answer the previous marginal comment?}
%   so we have
%   constructed an action of~$G_{K_s}$ on~$\gMt$. It remains to check
%   that~$(g-1)(\gM)\subseteq u^N\gMt$; for this, note that,
%   for any $m \in \gM$, we have
%     $$(g-1)m = \beta_g(1\otimes m)  - m = (\beta_g - \alpha_g)(1\otimes m)
%     + \alpha_g(1\otimes m ) - m \in u^N \gMt;
%     $$
% here we have used the fact $\im(\beta_g-\alpha_g)\subseteq
%   u^N\AAinf{A}\otimes_{\gS_A}\gM$, together with the second of the
%   properties satisfied by the maps $\alpha_g$.
%  % \mar{TG: this
%     % looks wrong after all, because as written we don't have ~$\lambda_i\in\AKplus$!}
% % \mar{TG: as ever not
%     % worrying about the exponents here.} 
% %
%   % It remains to check that this construction does indeed
%   % give~$\gM[1/u]$ the structure of a
%   % $(\varphi,G_{K_s})$-module. Firstly, note that the uniqueness
%   % of~$\beta_g$ implies that $\beta_{gh}=\beta_h\circ
%   % h^*\beta_g$\mar{TG: check} so that the~$\beta_g$ do indeed define an
%   % action of~$G_{K_s}$ on $\AAinf{A}\otimes_{\gS_A}\gM$ and
%   % thus on $\AAinf{A}\otimes_{\gS_A}\gM[1/u]$.\mar{TG: have to
%   % also justify continuity but once we've sorted that out elsewhere it
%   % will presumably be formal.}  
% \end{proof}

% %\mar{ME: Can we move this result (4.2.6) to the end of \S 4.2, or the start
% %	of \S 4.3?  As far as I can tell, the results in \S 4.2 
% %	that comes after it (4.2.7, 4.2.8, and the subsequent
% %	discussion) don't use 4.2.6, but rather just 4.2.5,
% %       	so I think the logical flow would 
% %	be clearer if they followed directly after 4.2.5. ME: Made this
% %        move.}
% %The following result is essentially the content
% %of~\cite[\S3]{MR2745530} (although they are concerned with getting
% %precise bounds on the size of~$s$, which is unimportant for us). We
% %put ourselves in the following situation: let~$B$ be a finite flat
% %$\cO$-algebra, and suppose that~$M_B$ is an \'etale
% %$(\varphi,\Gamma)$-module with $B$-coefficients, with the property
% %that the $E$-representation of~$G_K$ obtained from~$M_B$ is semistable with
% %Hodge--Tate weights contained in the interval~$[-h,0]$. For each~$a\ge
% %1$, write~$M_{B,\infty}^a$ for the \'etale $\varphi$-module
% %corresponding to the composite $\Spec (B/\varpi^a) \to\cX_d^a\to\cR_{\BK,d}^a$.
% %\begin{prop}
% %  \label{prop: Caruso Liu canonical action semistable}For any
% %  fixed~$a,h$ and~$N\ge e(a+h)/(p-1)$ there is a
% %  constant~$s'(a,h,N)\ge s(a,h,N)$ {\em (}where~$s(a,h,N)$ is as in
% %  Lemma~{\em \ref{lem: Caruso Liu Galois action on Kisin})} with the
% %  following property: if~$B$ and~$M_B$ are as above, and
% %  if~$s\ge s'(a,h,N)$, then there is a Breuil--Kisin module~$\gM$
% %  \mar{ME: Projective, height $\leq h$?}
% %  which is a lattice in ~$M_{B,\infty}^a$, for which the action
% %  of~$G_{K_s}$
% %  on~\[W(\C^\flat)_{B/\varpi^a}\otimes_{\gS_{B/\varpi^a}}\gM\subset
% %    W(\C^\flat)_{B/\varpi^a}\otimes_{\A_{K,B}}M_B\]\emph{(}given by the
% %  restriction of the natural action of~$G_K$ on the
% %  $(\varphi,G_K)$-module
% %  $W(\C^\flat)_{B/\varpi^a}\otimes_{\A_{K,B}}M_B$\emph{)} agrees with the
% %  action coming from Lemma~{\em \ref{lem: Caruso Liu Galois action on
% %    Kisin}}.
% %\end{prop}
% %\begin{proof}
% %  This is essentially the content of~\cite[Prop.\
% %  3.3.1]{MR2745530}. More precisely, the existence of an appropriate
% %  Kisin module is immediate from~\cite[Thm.\ 3.1.3]{MR2745530}. By the
% %  statement of Lemma~\ref{lem: Caruso Liu Galois action on Kisin}, we
% %  see that it is then enough to check that if~$s$ is sufficiently
% %  large, then for all~$g\in G_{K_s}$ we have
% %  $(g-1)(\gM)\subset u^N\Ainf\otimes_\gS\gM_{B,\infty}^a$, where the
% %  action of~$g$ is that given by restricting the natural action on
% %  $W(\C^\flat)_{\cO/\varpi^a}\otimes_{\A_{K}}M_B$ (note
% %  that~$B/\varpi^a$ is a finite $\cO/\varpi^a$-module, so there is no
% %  need to take a completed tensor product here). This is immediate
% %  from~\cite[Lem.\ 3.2.1]{MR2745530} (taking~$\mathfrak{L}=\gM$, and
% %  noting that in their notation, $\widehat{\mathfrak{L}}$ is an
% %  extension of scalars of~$\gL$ to a subring
% %  of~$\AAinf{\cO/\varpi^a}$; this extension of scalars involves a
% %  twist by~$\varphi$, but since~$\varphi$ is a continuous automorphism
% %  of~$\Ainf$, this is harmless).
% %\end{proof}
  
% It follows from
% the definitions % \mar{TG: hopefully doesn't need any work?}
% that for each~$s\ge 1$, the morphism
% $\cX_{K,d}\to\cR_{\BK,d}$ % \mar{TG: subscripts are permuting themselves!}
% factors as\[\cX_{K,d}\to\cX_{K_s,d}\to\cR_{\BK,d}. \]
% \begin{prop}
%   \label{prop: morphism C to Xs}For any fixed $a,h$, there is
%   some~$s=s(a,h)$ such that there is a morphism
%   $\cC^a_{\BK,d,h}\to\cX^a_{K_s,d}$ fitting into a commutative
%   triangle% \mar{TG: probably eventually add something about how this
%     % works with crystalline representations, too.}
%   \[\xymatrix{& \cC^a_{\BK,d,h}\ar[dl]\ar[d]\\ \cX^a_{K_s,d}\ar[r]&\cR^a_{\BK,d}}\]
% %  $\cC^a_{d,\infty,h}\to\cR^a_{d,\infty}$ factors through the morphism $\cX^a_{K_s,d}\to\cR^a_{d,\infty}$.
% \end{prop}%\mar{TG: need to be rewritten}
% \begin{proof}% \mar{TG: needed, should basically be Frobenius
%     % amplification}
%   As in the proof of Proposition~\ref{prop: natural morphism X to
%     R}, it is enough to show that if~$A$ is a finite type
% %\mar{ME: finite type algebra, maybe?}
%   $\cO/\varpi^a$-algebra, and~$\gM$ is a finite projective \'etale Breuil--Kisin
%   module with $A$-coefficients, then we can canonically % \mar{TG: need
%     % to think about why the extension will actually be canonical!}
%   extend the action of~$G_{K_\infty}$ on
%   $W(\C^\flat)_A\otimes_{\gS_A}\gM$ to an action
%   of~$G_{K_s}$, if~$s$ is sufficiently large. This is immediate from Lemma~\ref{lem: Caruso Liu Galois action on Kisin}.
% \end{proof}

% \begin{lem}\label{lem: C to Xs is proper}
%   The morphism $\cC^a_{\BK,d,h}\to\cX^a_{K_s,d}$ of
%   Proposition~{\em \ref{prop: morphism C to Xs}} is proper.% \mar{TG: do
%   % we explicitly use this properness? It should help scheme-theoretic
%   % images to be better behaved, at the very least?}
% \end{lem}
% \begin{proof} This follows by a standard graph argument.  Namely,
% 	we write this morphism as the
%   composite \[\cC^a_{\BK,d,h}\to\cC^a_{\BK,d,h}\times_{\cR^a_{\BK,d}}\cX^a_{K_s,d}
%     \to\cX^a_{K_s,d}. \] The first arrow is proper, because it is
%   the base change of the diagonal morphism $\cX_d\to\cX_d\times_{\cR_{\BK,d}}\cX_d$, which is a closed immersion
%   by Proposition~\ref{prop: diagonal of X to R}. The second arrow
%   is proper because it is the base change of the proper morphism
%   $\cC^a_{\BK,d,h}\to{\cR^a_{\BK,d}}$ of Theorem~\ref{thm: C is a $p$-adic formal algebraic stack and related
%     properties}.
% \end{proof}

% \begin{df}
% 	\label{def:2-Cartesian}
% The morphisms of Lemma~\ref{lem: restriction from X to Xs} and
% Proposition~\ref{prop: morphism C to Xs} allow us to define a
% stack~$\cC^a_{s,d,h}$ by the requirement that it fits into a 
% 2-Cartesian diagram \[\xymatrix{\cC^a_{s,d,h} \ar[r]\ar[d] &  \cC^a_{\BK,d,h}\ar[d]\\
%   \cX^a_{K,d}\ar[r] & \cX^a_{s,d}}\]
% Since the lower horizontal arrow in
% %the diagram defining $\cC^a_{s,d,h}$
% this diagram
% is representable by algebraic spaces and of 
% finite type, by Lemma~\ref{lem: restriction from X to Xs},
% and since~$\cC^a_{\BK,d,h}$ is an
% algebraic stack of finite type over $\Z/p^a$,
% by Theorem~\ref{thm: C is a $p$-adic formal algebraic stack and related
%     properties}, we see that~$\cC^a_{s,d,h}$   is also an algebraic
% stack of finite type over $\Z/p^a$.
% Since the right hand vertical arrow in this diagram is proper,
% by Lemma~\ref{lem: C to Xs is proper}, so is the left
% hand vertical arrow. Since~$\cX^a_{K,d}$ is an Ind-algebraic stack by Proposition~\ref{prop:X is an Ind-stack}, we can
% form the scheme-theoretic image~$\cZ^a_{s,d,h}$ of the morphism $\cC^a_{s,d,h}\to \cX^a_d
% $, which is a
% closed algebraic substack of~$\cX^a_d$, and is in particular of 
% finite type over $\Z/p^a$.
% \end{df}

% % By the result (or rather the method)
% % of~\cite{MR2745530} we have
% % \[\xymatrix{& (alg,ft)\ar[dl]\ar[r]&\cC_a^{\le
% %       h}\ar[dl]^{proper}\ar[d]^{proper}\\ \cX_a\ar[r]^{f.t., rep}&\cX_{s,a}\ar[r]^{sep,ind-rep}&\cR_a} \]
% % \begin{itemize}
% % 	\item Using a reinterpretation of Caruso--Liu in terms of $(\varphi,
% % 		G_K)$-modules, we show that (for some $s$ depending on
% % 		$a$ and $h$) the right-side arrow of the parallelogram exists.
% % 		We want to show that this map is proper (given
% % 		that the vertical arrow on the extreme right is known
% % 		to be proper)
% % \item To check the properness of the previous point,
% % 	the key thing is the properties of the lower right
% %   arrow; namely, its diagonal should be a closed immersion (equivalently,
% %  a proper monomorphism). 
% %  This is basically by hand, using the fact that
% % % : separability is saying that the
% % %  diagonal is proper, and that's not so hard, because in fact it's a
% % %  closed immersion:
% % endomorphisms of the
% %   $G_K$-representation are the same as endomorphisms of the
% %   $G_{K_\infty}$-representation which commute with~$G_K$. 
% % \item Once we have the desired properness,
% % we can pullback over $\cX_a$ and obtain a finite-type alg.\ stack
% % with a proper morphism to $\cX_a$.  (Here we use the fact 
% % that $\cX_a \to \cX_{s,a}$ is finite type and representable,
% % which is easy to see when thought of as restricting $(\varphi,\Gamma)$-modules
% % to $(\varphi,\Gamma_s)$-modules.)
% % Its scheme-theoretic image is then a closed finite-type alg.\ substack
% % of $\cX_a$.  We now check, for any crystalline versal ring
% % $R_{\crys}^{\le h}$,
% % that the morphism $\Spf R^{\leq h}_{\crys} \to \cX_a$ factors through 
% % this scheme-theoretic image, and so is effective.
% % \end{itemize}



% % \begin{itemize}
% % \item Discuss compatibility with Galois side.
% % \item Should probably cite 1-ification material from
% %   \cite{EGstacktheoreticimages}.
% % \item Want to show it's separated and ind-representable.
% % \end{itemize}

% \section{Crystalline deformation rings}
% We now return to the
% crystalline deformation rings of~\cite{MR2373358}, which we recalled
% in Section~\ref{subsec: notation and conventions}.

% \begin{df}
% We say that a Hodge
% type~$\underline{\lambda}$ is \emph{bounded by~$h$} if
% for each~$\sigma$, $i$ we have $\lambda_{\sigma,i}\in [-h,0]$.% \mar{TG: or perhaps the
%   % opposite sign, not worrying about conventions} 
%   \end{df}

% %\mar{ME: Can we move this result (4.2.6) to the end of \S 4.2, or the start
% %	of \S 4.3?  As far as I can tell, the results in \S 4.2 
% %	that comes after it (4.2.7, 4.2.8, and the subsequent
% %	discussion) don't use 4.2.6, but rather just 4.2.5,
% %       	so I think the logical flow would 
% %	be clearer if they followed directly after 4.2.5. ME: Made this
% %       move.}
% The following result is essentially the content
% of~\cite[\S3]{MR2745530} (although they are concerned with getting
% precise bounds on the size of~$s$, which is unimportant for us). We
% put ourselves in the following situation: let~$B$ be a finite flat
% $\cO$-algebra, and suppose that~$M_B$ is an \'etale
% $(\varphi,\Gamma)$-module with $B$-coefficients, with the property
% that the $E$-representation of~$G_K$ obtained from~$M_B$ is semistable with
% Hodge--Tate weights bounded by~$h$.
% %contained in the interval~$[-h,0]$.
% For each~$a\ge 1$, write~$M_{B,\infty}^a$ for the \'etale $\varphi$-module
% corresponding to the composite $\Spec (B/\varpi^a) \to\cX_d^a\to\cR_{\BK,d}^a$.
% \begin{prop}
%   \label{prop: Caruso Liu canonical action semistable}For any
%   fixed~$a,h$ and~$N\ge e(a+h)/(p-1)$ there is a
%   constant~$s'(a,h,N)\ge s(a,h,N)$ {\em (}where~$s(a,h,N)$ is as in
%   Lemma~{\em \ref{lem: Caruso Liu Galois action on Kisin})} with the
%   following property: if~$B$ and~$M_B$ are as above, and
%   if~$s\ge s'(a,h,N)$, then there is a projective Breuil--Kisin
%   module~$\gM$ of height at most~$h$
% %  \mar{ME: Projective, height $\leq h$?}
%   which is a lattice in ~$M_{B,\infty}^a$, for which the action
%   of~$G_{K_s}$
%   on~\[W(\C^\flat)_{B/\varpi^a}\otimes_{\gS_{B/\varpi^a}}\gM\subset
%     W(\C^\flat)_{B/\varpi^a}\otimes_{\A_{K,B}}M_B\]\emph{(}given by the
%   restriction of the natural action of~$G_K$ on the
%   $(\varphi,G_K)$-module
%   $W(\C^\flat)_{B/\varpi^a}\otimes_{\A_{K,B}}M_B$\emph{)} agrees with the
%   action coming from Lemma~{\em \ref{lem: Caruso Liu Galois action on
%     Kisin}}.
% \end{prop}
% \begin{proof}\mar{TG: I suspect this might be a bit dodgy! I think in
%     the theory of $(\varphi,\widehat{G})$-modules one has a
%     $\varphi$-twist to the embedding of~$\gM$, so it's possible that
%     whenever we talk about Breuil--Kisin--Fargues modules, we should
%     be doing $\otimes_{\gS_A,\varphi}\AAinf{A}$? Probably we should
%     compare to whatever \cite{2016arXiv160203148B} do, too.}
%   This is essentially the content of~\cite[Prop.\
%   3.3.1]{MR2745530}. More precisely, the existence of an appropriate
%   Kisin module is immediate from~\cite[Thm.\ 3.1.3]{MR2745530}. By the
%   statement of Lemma~\ref{lem: Caruso Liu Galois action on Kisin}, we
%   see that it is then enough to check that if~$s$ is sufficiently
%   large, then for all~$g\in G_{K_s}$ we have
%   $(g-1)(\gM)\subset u^N\Ainf\otimes_\gS\gM_{B,\infty}^a$, where the
%   action of~$g$ is that given by restricting the natural action on
%   $W(\C^\flat)_{\cO/\varpi^a}\otimes_{\A_{K}}M_B$ (note
%   that~$B/\varpi^a$ is a finite $\cO/\varpi^a$-module, so there is no
%   need to take a completed tensor product here). This is immediate
%   from~\cite[Lem.\ 3.2.1]{MR2745530} (taking~$\mathfrak{L}=\gM$, and
%   noting that in their notation, $\widehat{\mathfrak{L}}$ is an
%   extension of scalars of~$\gL$ to a subring
%   of~$\AAinf{\cO/\varpi^a}$; this extension of scalars involves a
%   twist by~$\varphi$, but since~$\varphi$ is a continuous automorphism
%   of~$\Ainf$, this is harmless).
% \end{proof}



% %Suppose now that~$\F'/\F$ is a finite extension, and
% %that~$\Spec\F'\to\cX_d$ is a finite type point, % \mar{TG:
% %% maybe should have $\F'$, $\cO'$ here, and then annoyingly~$\cO''$ below?}
% %with  corresponding Galois representation 
% %$\rhobar:G_K\to\GL_d(\F')$. Let ~$\cO'$ be the ring of integers in
% %some finite extension~$E'/E$ with residue field~$\F'$.  Then there is a versal morphism $\Spf
% %R_{\rhobar}^{\square,\cO'}\to\cX_d$; indeed, by the theorem of Dee
% %recalled in Section~\ref{subsec: phi gamma with coefficients},
% %$R_{\rhobar}^{\square,\cO'}$ is isomorphic to the universal lifting ring for 
% %% \mar{ME: $\cO'$ is being used here, before its actually defined. TG:
% %%   rearranged to hopefully sort this out?}
% %$(\varphi,\Gamma)$-modules lifting the $(\varphi,\Gamma)$-module
% %corresponding to $\Spec\F'\to\cX_d$ (an isomorphism being given by any
% %choice of ordered basis for the universal Galois lifting and universal
% %$(\varphi,\Gamma)$-lifting), and it is clear from the definitions that the
% %corresponding morphism $\Spf
% %R_{\rhobar}^{\square,\cO'}\to\cX_d$ is a $\widehat{\GL}_d$-torsor, and
% %in particular is versal.
% %\mar{ME: There's probably some stupid issue here where choosing a basis
% %	for a Galois rep'n is not the same as choosing a basis for a $(\varphi,
% %	\Gamma)$-modules, so they are not literally the same ring. TG:
% %     Good point! I guess the real problem here is that we never defined ``lifting
% %    ring'', so we never say anything about bases, and we don't do that
% %  on the $(\varphi,\Gamma)$-side either, although we did in
% %  \cite{EGstacktheoreticimages}. Leaving this for now but I guess we
% %  need to expand on this, probably actually defining the lifting rings
% %would be a good start? I guess the correct way to set this up is that
% %the two are related by fixing bases on each side
% %(Galois/$(\varphi,\Gamma)$) for the universal objects? TG: OK wrote
% %something, feel free to improve.}

%  % Then the corresponding lifting ring
% % $R^{\crys,\underline{\lambda},\cO'}_{\rhobar}$  is nonzero, and the
% % versal morphism
% % \mar{TG: we haven't set this up yet, but it is
%   % presumably easy to show that this does give the right versal
%   % rings?!}
% Fix a continuous representation $\rhobar: G_K \to \GL_d(\F')$,
% for some finite extension $\F'$ of $\F$, giving rise to a point
% $x \in \cX_d(\F')$,  
% and also fix a $p$-adic Hodge type~$\underline{\lambda}$.
% Let $\cO'$ be the ring of integers in a finite extension $E'$ of $E$
% having residue field $\F'$.  Then the versal morphism $\Spf R_{\rhobar}^{\square,\cO'}\to\cX_d$
% (described in Subsection~\ref{subsec:Galois reps})
% induces a
% morphism $\Spf R^{\crys,\underline{\lambda},\cO'}_{\rhobar}\to\cX_d$, and
% thus a morphism $\Spf
% R^{\crys,\underline{\lambda},\cO'}_{\rhobar}/\varpi^a\to\cX_d^a$ for
% each~$a\ge 1$. 
% % \mar{ME: A technical/expositional point:  the ring is still defined even if it
% % 	is zero; it just gives us the empty scheme mapping to $\cX_d^a$.
% % 	So we really need to assume the existence of a crystalline lift
% % 	at this point? We could just introduce the auxiliary extensions
% % $E'/E$ and $\cO/\cO'$. TG: agreed, commenting out.}
% \begin{prop}
%   \label{prop: crystalline deformation ring factors through the scheme
%     theoretic image}Suppose that~$\underline{\lambda}$ is bounded by~$h$. Fix~$a\ge 1$, and fix~$s$ sufficiently large in
%   the sense of Lemma~{\em \ref{lem: Caruso Liu Galois action on Kisin}}
%   and 
%   Propositions~{\em \ref{prop: Caruso Liu canonical action semistable}}
%   and~{\em \ref{prop: morphism C to Xs}}. Then for every % discrete
%   % \mar{ME: *Every* Artinian quotient is discrete, right?}
%   Artinian quotient~$A^\crys$ of $
%   R^{\crys,\underline{\lambda},\cO'}_{\rhobar}/\varpi^a$, the
%   composite \[\Spec A^\crys\hookrightarrow \Spf R^{\crys,\underline{\lambda},\cO'}_{\rhobar}/\varpi^a \to
%   \cX_d^a\]factors through a morphism $\Spec A^\crys\to \cC^a_{s,d,h}$ 
% {\em (}where $\cC^a_{s,d,h}$ denotes the $2$-Cartesian product
% of Definition~{\em \ref{def:2-Cartesian})}.
% % the morphism
%   % $\Spf R^{\crys,\underline{\lambda},\cO'}_{\rhobar}/\varpi^a \to
%   % \cX_d^a$ factors through~$\cZ^a_{d,K_s,h}$.
% \end{prop}
% \begin{proof}We make use of the results
%   of~\cite[\S4]{MR3449204}. Define a quotient~$R^{\underline{\lambda},\semis,\cO'}_{\rhobar}$
%   of~$R^{\square,\cO'}_{\rhobar}$ as follows: if~$A$ is an Artinian
%   local $\cO'$-algebra with residue field~$\F'$, then a morphism
%   $R^{\square,\cO'}_{\rhobar}\to A$ factors
%   through~$R^{\underline{\lambda},\semis,\cO'}_{\rhobar}$ if and only
%   if the corresponding representation $\rho_A:G_K\to\GL_d(A)$ satisfies:
%   \begin{itemize}
%   \item there exists a surjection of~$\cO'$-algebras $B\to A$,
%     where~$B$ is finite flat over $\cO'$, and a
%     representation~$\rho_B:G_K\to\GL_d$ lifting~$\rho_A$, such
%     that~$\rho_B\otimes_{\Zp}\Qp$ is semistable of Hodge type~$\underline{\lambda}$.
%   \end{itemize}


%   The existence of this quotient is straightforward: in the case
%   that~$\rhobar$ admits a universal deformation ring, it
%   is~\cite[Prop.\ 4.23]{MR3449204}, and for general~$\rhobar$ an
%   identical argument works for lifting
%   % \mar{ME: Just for consistencty in terminology:  ``framed deformation ring''
%   %         = ``lifting ring''?}
%   rings. By~\cite[Thm.\ 4.25]{MR3449204} (or rather the corresponding
%   statement for lifting rings, which again has an identical
%   proof), 
% %  \mar{ME: Have we explained that $E'$ equals frac.\ field of $\cO'$?}
%   if $E''/E'$ is a finite extension with ring of
%   integers~$\cO''$, then a homomorphism of $\cO'$-algebras
%   $R_{\rhobar}^{\square,\cO'}\to\cO''$ factors through
%   $R^{\underline{\lambda},\semis,\cO'}_{\rhobar}$ if and only if the
%   corresponding Galois representation is semistable of Hodge
%   type~$\underline{\lambda}$.

%   Recall from Section~\ref{subsec: notation and
%     conventions} that $R^{\crys,\underline{\lambda},\cO'}_{\rhobar}$ is
%   characterised as the unique quotient of~$R_{\rhobar}^{\square,\cO'}$
%   which is~$\cO'$-flat and reduced and has the property that for all
%   finite extensions~$E''/E'$ with ring of integers~$\cO''$, a morphism of
%   $\cO'$-algebras $R_{\rhobar}^{\square,\cO}\to\cO''$ factors through
%   $R^{\crys,\underline{\lambda},\cO'}_{\rhobar}$ if and only if the
%   corresponding Galois representation is crystalline of Hodge
%   type~$\underline{\lambda}$.

%   It is immediate that $R^{\crys,\underline{\lambda},\cO'}_{\rhobar}$
%   is a quotient of
%   $R^{\underline{\lambda},\semis,\cO'}_{\rhobar}$. Let~$A$ be any
%   Artinian quotient of
%   $R^{\crys,\underline{\lambda},\cO'}_{\rhobar}/\varpi^a$; we must
%   show that the morphism $\Spec A\to\cX^a_d$ lifts to~$\cC^a_{s,d,h}$. Since~$A$ is also a quotient of
%   $R^{\underline{\lambda},\semis,\cO'}_{\rhobar}$, there is a
%   surjection of~$\cO'$-algebras $B\to A$ satisfying the property
%   above. Let~$M_{B/\varpi^a}$ be the \'etale $(\varphi,\Gamma)$-module
%   corresponding to the representation~$\rho_B\otimes_BB/\varpi^a$;
%   then by Proposition~\ref{prop: Caruso Liu canonical action
%     semistable}, the morphism $\Spec B/\varpi^a\to\cX_{K,d}$ lifts
%   to~$\cC^a_{s,d,h}$. Since $\Spec A\to\cX_{K,d}$ factors
%   through the closed immersion~$\Spec A\to\Spec B/\varpi^a$,
%   it follows that it lifts to~$\cC^a_{s,d,h}$, as required.
%   %since~$A$ was arbitrary, we are done.
% \end{proof}


% %   In particular, if we consider all such
% %   $x:R^{\crys,\underline{\lambda},\cO'}_{\rhobar}\to\cO''$ (where in a
% %   slight abuse of notation, $\cO''$
% %   depends on~$x$), we see that there is an injection
% %   $R^{\crys,\underline{\lambda},\cO'}_{\rhobar}\to\prod_x\cO''$, and
% %   thus (since $R^{\crys,\underline{\lambda},\cO'}_{\rhobar}$ is
% %   $\cO'$-flat) an injection
% %   $R^{\crys,\underline{\lambda},\cO'}_{\rhobar}/\varpi^a\to\prod_x\cO''/\varpi^a$.
  
% %   We are therefore reduced\mar{TG: ME says this needs some justification.} to showing that for each
% %   such~$R^{\crys,\underline{\lambda},\cO'}_{\rhobar}\to\cO''$, the
% %   corresponding morphism $\Spec \cO''/\varpi^a\to\cX_d^a$ factors
% %   through~$\cZ^a_{d,K_s,h}$. To see this, it is enough to show that we
% %   can lift this morphism to a
% %   morphism~$\Spec\cO''/\varpi^a\to\cC^a_{d,K_s,h}$.

% % To this end, note firstly that the composite
% % morphism~$\Spf\cO''\to\cX_d\to\cR_{\BK,d}$ lifts uniquely
% % to~$\cC_{\BK,d,h}$ as a consequence of the main theorems
% % of~\cite{KisinCrys} (that is, the corresponding \'etale
% % $\varphi$-module over~$\cO_{\cE,\cO''}$ contains a unique Breuil--Kisin module~$\gM$ over
% % $\gS_{\cO''}$ of height at most~$h$). We need to show that the
% % action of~$G_{K_s}$ on
% % $W(\C^\flat)_{\cO''/\varpi^a}\otimes_{\gS_{\cO''}}\gM[1/u]$ which comes from
% % its incarnation as  the $(G_K,\varphi)$-module corresponding to the
% % morphism $\Spec \cO''/\varpi^a\to\cX_d$ satisfies the
% % property of Lemma~\ref{lem: Caruso Liu Galois action on Kisin}; that
% % is, we must show that this action preserves
% % $\AAinf{\cO''/\varpi^a}\otimes_{\gS_A}\gM$, and that for
% % all~$g\in G_{K_s}$, this action satisfies
% % $(g-1)(\gM)\subseteq u^N\AAinf{\cO''/\varpi^a}\otimes_{\gS_A}\gM$.\mar{TG:
% % should now be able to complete by citing \cite{MR2558890} and
% % \cite{MR2745530}, the former for the structure of phi Ghat and the
% % latter for this congruence.}
% % \end{proof}
% \begin{cor}
%   \label{cor: crystalline deformation rings are effective
%     versal}
% % \mar{ME: Rephrased this to avoid assuming existence of a lift,
% % 	b/c I don't think that's relevant/necessary.}
% For any Hodge type $\underline{\lambda}$, and
% %Suppose that~$\rhobar$ admits a crystalline lift $\rho:G_K\to\GL_d(\cO')$ of Hodge
%  % type ~$\underline{\lambda}$ for some Hodge
%  % type~$\underline{\lambda}$. Then
%   for any~$a\ge 1$,
%   % \mar{ME: Shouldn't we introduce an auxiliary $\cO'$ here, as we did
%   %         above? TG: done, commenting out.}
%   the corresponding morphism
%   $\Spf R^{\crys,\underline{\lambda},\cO'}_{\rhobar}/\varpi^a \to
%   \cX_d^a$ is effective, i.e.\ is induced by a morphism $\Spec R^{\crys,\underline{\lambda},\cO'}_{\rhobar}/\varpi^a \to
%   \cX_d^a$.
% \end{cor}
% % \mar{ME: By twisting into the effective
% % 	situation, we can obtain the same result for any 
% % 	crystalline deformation ring, so probably we should just state the 
% % 	result in that level of generality?}
% \begin{proof}Twisting by~$\varepsilon^{p^n(p-1)}$ for some 
%   %       \mar{ME: Somewhere we should note that twisting by a character
%   %       	induces an automorphism of $\cX_d$! TG: did this back
%   %               in Section~\ref{subsec: tensor product
%   % of phi gamma and duality}, commenting out.}
%   sufficiently large~$n$ (chosen so that in particular
%   $\varepsilon^{p^n(p-1)}$ is trivial mod~$\varpi^a$), we may assume
%   that 
% for each~$\sigma$, $i$ we have $\lambda_{\sigma,i}\le 0$. Choose~$h\ge
% 0$ so that~$\underline{\lambda}$ is bounded by~$h$. Fix~$s$ sufficiently large in
%   the sense of Lemma~\ref{lem: Caruso Liu Galois action on Kisin},
%   Proposition~\ref{prop: Caruso Liu canonical action semistable}, and
%   Proposition~\ref{prop: morphism C to Xs}. 



% Let~$\Spf R^\cC$ be the scheme-theoretic image (in the sense
% of~\cite[Defn.\ 3.2.15]{EGstacktheoreticimages}) of the pulled-back
% proper morphism  \[(\cC^a_{s,d,h})_{\Spf
%     R_{\rhobar}^{\square,\cO'}/\varpi^a}\to \Spf
%   R_{\rhobar}^{\square,\cO'}/\varpi^a.\]Since~$\cC^a_{s,d,h}\to\cX_d^a$ is
% proper, it follows from~\cite[Lem.\ 3.2.16]{EGstacktheoreticimages}
% that the composite $\Spf R^\cC\into  \Spf
%   R_{\rhobar}^{\square,\cO'}/\varpi^a\to\cX_d^a$ factors
%   through~$\cZ^a_{s,d,h}$ (the scheme-theoretic image of 
%   $\cC^a_{s,d,h}$; see Definition~\ref{def:2-Cartesian}) and the morphism $\Spf
%   R^\cC\to\cZ^a_{s,d,h}$ is versal. Since~$\cZ^a_{s,d,h}$ is an
%   algebraic stack, its versal rings are effective. It follows that the
%   morphism $\Spf R^\cC\to\cX_d^a$ is induced by a  morphism $\Spec R^\cC\to\cX_d^a$.
  
% It is therefore enough to show that the morphism
% $\Spf R^{\crys,\underline{\lambda},\cO'}_{\rhobar}/\varpi^a\hookrightarrow
% \Spf R_{\rhobar}^{\square,\cO'}/\varpi^a$ factors through
% $\Spf R^\cC$. In turn, it suffices to show that if~$A$ is an % discrete
% % \mar{ME: All Artinian quotients are discrete.}
% Artinian quotient of
% $R^{\crys,\underline{\lambda},\cO'}_{\rhobar}/\varpi^a$, the morphism
% $\Spec A\to\Spf R_{\rhobar}^{\square,\cO'}/\varpi^a$ factors through~$\Spf R^\cC$. To see this, note
% that by Proposition~\ref{prop: crystalline deformation ring factors through the scheme
%     theoretic image}, the morphism
%   $\Spec A\to \cX_d^a$ factors through 
%   $\cC^a_{s,d,h}\to\cX_d^a$; so the morphism $\Spec A\to\Spf
%   R_{\rhobar}^{\square,\cO'}/\varpi^a$ factors
%   through~$(\cC^a_{s,d,h})_{R_{\rhobar}^{\square,\cO'}/\varpi^a}$,
%   and therefore through the scheme-theoretic image~$\Spf R^\cC$, as required. 
% \end{proof}


\section{The fibre dimension of $H^2$ on crystalline deformation rings}
%\mar{ME: Provisional title}
%\mar{TG: before the following theorem need to actually explain the
%  Herr complex over the crystalline deformation ring and so on}
% Note that we can work over $\Spec R$ rather than just
%		$\Spf R$ for the following reasons: firstly, since $R$ is
%		a complete Noetherian local ring,
%		coherent sheaves, perfect complexes, etc.\ on $\Spec R$
%		and $\Spf R$ are the same.\mar{TG: I suspect this
%                  should be explained somewhere outside of this proof,
%                and similarly the claim that things are pulled back
%                TG: just dumped in some material about this moved back
%              from the final section}
%              % Secondly,
%	      %   since $\cX_d^{\crys,\underline{\lambda}}$
%	      %   is a $p$-adic formal algebraic stack, its mod $p$ fibre is algebraic,
%	      %   so that its versal rings will be effective; thus we will have
%	      %   a versal morphism
%By Theorem~% \ref{thm: crystalline deformation rings are effective
%  % versal}
%, we have a versal morphism		$\Spec R/p \to \Spec \F_p\times_{\Spf \Z_p}\cX_d^{\crys,\underline{\lambda}},$ 
%		and so we can in particular apply the property
%                Theorem~\ref{thm:Xdred is algebraic}~(4) on  the support of $\Ext^2$ over $\cX_{d,\red}$
%		is inherited by $\Ext^2$ over $\Spec R$. \mar{TG: I don't
%                  understand what's being used/assumed here to get
%                  that the statements about the codimension of the support of $H^2$ pull
%                  back to $\Spec R$, and in particular over~$X$. ME: Added ``versal'' to the discussion;
%	  I guess smoothness/flatness is the key point. TG: presumably
%        again this will be discussed more fully - and in fact replaced
%      with something more legitimate.  ME: This first paragraph will
%be more-or-less entirely replaced by a citation to



%\mar{TG: say that we mean Theorem~\ref{thm:Xdred is algebraic}~(4)}
Let $\rhobar: G_K \to \GL_d(\F')$ be a continuous representation,
for some finite extension $\F'$ of $\F$, and let $\cO'$ be the ring 
of integers in a finite extension $E'/E$, with residue field~ $\F'$.
Let $\underline{\lambda}$ be  a Hodge type,
and consider the lifting ring $R :=
R_{\rhobar}^{\crys,\underline{\lambda},\cO'}$.
Let $M$ denote the universal \'etale % \mar{TG: might want to change
  % this terminology, and discuss the $\m$-adic completion aspect
  % further in the earlier sections.}
$(\varphi,\Gamma)$-module
over $R$ (note that there is a universal \'etale
$(\varphi,\Gamma)$-module over~$R$, and not merely a universal
\emph{formal} \'etale $(\varphi,\Gamma)$-module,  as a consequence of Corollary~\ref{cor: crystalline deformation rings are effective
    versal}). Let $\alpha^{\circ}: G_K \to \GL_a(\cO')$ be a representation
	with $\alphabar: G_K \to \GL_a(\F')$ being absolutely irreducible,
and let~$N$ be the base change to~$R$ of the \'etale $(\varphi,\Gamma)$-module
over~$\cO'$ which corresponds to~$\alpha^{\circ}$. 

Let $\cC^{\bullet}(N^\vee\otimes M)$ denote the Herr complex
associated to $N^\vee\otimes M$, as in Section~\ref{subsec: Herr complex}; this
is a perfect complex of~$R$-modules, with cohomology concentrated in degrees~$[0,2]$,
as a consequence of
~\cite[\href{https://stacks.math.columbia.edu/tag/0CQG}{Tag
  0CQG}]{stacks-project} and Theorem~\ref{thm: Herr complex is
  perfect} (applied to the quotients~$R/\m_R^a$).
Write $\Ext^2 := H^2\bigl(\cC^{\bullet}(N^\vee\otimes M)\bigr)$.
For each $r \geq 1,$
write $$X_r := \{x \in \Spec R \, | \, \dim \kappa(x)\otimes_R \Ext^2 \geq r \}.$$
The theorem on upper semi-continuity of fibre dimension for coherent sheaves
shows that $X_r$ is a closed subset of $\Spec R$.


\begin{thm}%\mar{TG: needs to be about codim not dim}\mar{TG: should be
  %moved into section~\ref{sec: families of extensions}}
  \label{thm: codimension H2 in regular weight crystalline deformation ring}
  If $\Ext^2$
  is $\varpi$-power torsion,
  then for each~$r\ge 1$, $X_r$ is of codimension $\geq r+1$ in $\Spec R$.
\end{thm}% \mar{TG: I think this needs to be rewritten to be about
  % $\Ext^2$ for general~$\alpha$ TG: attempted to do this.}
\begin{proof}
%	\mar{ME: Write complete proof!}  
	The assumption that $\Ext^2$ is $\varpi$-power torsion
	implies that $\Ext^2$ is (set-theoretically) supported
	on the closed subscheme $\Spec (R/\varpi)_{\red}$ of $\Spec R$.
        Write $\overline{M}$ and $\overline{N}$ to denote the
        pull-backs of $M$ and~$N$ respectively,
	over $(R/\varpi)_{\red}$,
	write $\overline{\Ext}^2 := \Ext^2\bigl(\cC^{\bullet}(\overline{N}^\vee\otimes\overline{M})\bigr),$
	and define
        $$\overline{X}_r := \{x \in \Spec (R/\varpi)_{\red} \, | \,
	\dim \kappa(x)\otimes_{(R/\varpi)_{\red}} \overline{\Ext}^2 \geq r \}.$$
	Since~$\Ext^2$ is compatible with base change, and is supported
        on~$\Spec (R/\varpi)_{\red}$, while~$R$ is~$\cO$-flat,  it suffices to show that
	% \mar{ME: Maybe elaborate on why this suffices?  Use that
	% $H^2$ is compatible with base-change!}
the codimension of	$\overline{X}_r$  in
	$\Spec (R/\varpi)_{\red}$ is at least~$r$.
	Since $\Spec (R/\varpi)_{\red}$
	is equidimensional of dimension $[K:\Q_p]d(d-1)/2 + d^2,$
	it suffices in turn to show that $\overline{X}_r$
	is of dimension at most $([K:\Q_p]d(d-1)/2) + d^2 - r.$

	%By the discussion of Section~\ref{subsec:Galois reps},
	By Proposition~\ref{prop:versal rings},
	we have a versal morphism
	$\Spf R^{\square,\cO'}_{\rhobar}
	\to \cX_d$.
	Consider the fibre product
	$\Spf R^{\square,\cO'}_{\rhobar} \times_{\cX_d} \cX_{d,\red}.$
	This is a closed formal algebraic subspace of 
	$\Spf R^{\square,\cO'}_{\rhobar}$,
	and so by Lemma~\ref{lem: alem closed formal algebraic scheme
    adic*} 
is of the form $\Spf S$ for some quotient $S$ of 
	$R^{\square,\cO'}_{\rhobar}$. 
	Its description as a fibre product
	shows that the morphism $\Spf S \to \cX_{d,\red}$ is versal.
	Since $\cX_{d,\red}$ is an algebraic stack,
	this morphism is effective, i.e.\ arises from a morphism
	\numequation
	\label{eqn:S effective}
	\Spec S \to \cX_{d,\red}
\end{equation}
       	(\cite[\href{https://stacks.math.columbia.edu/tag/0DR1}{Tag 0DR1}]{stacks-project}), 
which is furthermore 
flat, by~\cite[\href{https://stacks.math.columbia.edu/tag/0DR2}{Tag 0DR2}]{stacks-project}.
	Let $x_0 = \Spec \F'$, thought of as the closed point of $\Spec S$,
	and endowed with a morphism $x_0 \to \cX_{d,\red}$
	corresponding to the representation~$\rhobar$.
	Since $\Spf S \to \cX_{d,\red}$
	is versal,
	it follows from the Artin Approximation Theorem that
	we may find a finite type $\F$-scheme $U$, equipped with a point
	$u_0$ with residue field $\F'$, such that the complete local
	ring of $U$ at $u_0$ is isomorphic to $S$,
	and 
	such that~\eqref{eqn:S effective} may be promoted 
	to  a smooth morphism
	\numequation
	\label{eqn:U promotion}
	U\to \cX_{d,\red},
\end{equation}
in the sense that on $u_0 (= \Spec \F')$,
the morphism~\eqref{eqn:U promotion}
induces the given morphism $x_0 (= \Spec \F') \to \cX_{d,\red}$,
and on $\Spec \widehat{\cO}_{U,u_0} \cong \Spec S$, 
the morphism~\eqref{eqn:U promotion}
induces the morphism~\eqref{eqn:S effective}.
	(See~\cite[\href{https://stacks.math.columbia.edu/tag/0DR0}{Tag 0DR0}]{stacks-project}.)


%	the morphism~\eqref{eqn:S effective} is flat by~\cite[\href{https://stacks.math.columbia.edu/tag/0DR2}{Tag 0DR2}]{stacks-project}.
%	% \mar{ME: Not sure if we need this effectivity statement,
%	% 	or the flatness, but if we do, give references.}

	%\mar{ME: The argument in this paragraph maybe needs some elaboration,
	%	and also some citations.}
%	Let $x_0 = \Spec \F'$, thought of as the closed point of $\Spec S$,
%	and endowed with a morphism $x_0 \to \cX_{d,\red}$
%	corresponding to the representation $\rhobar$.
	It follows from Proposition~\ref{prop:versal rings},
	by pulling back over $\cX_{d,\red}$,
	that
	$\Spf S \times_{\cX_{d,\red}} \Spf S$
	is isomorphic to
        $\widehat{\GL}_{d,S},$
	a certain completion
	of~$(\GL_{d})_S$,  
	and thus that
	$\Spf S \times_{\cX_{d,\red}} x_0$
	is isomorphic to
        $\widehat{\GL}_{d,\F'},$
	a certain completion
	of~$(\GL_{d})_{\F'}$.  
	This latter fibre product may be identified
	with the completion of
	$\Spec S \times_{\cX_{d,\red}} x_0$ along
	its closed subspace $x_0\times_{\cX_{d,\red}} x_0,$
	and so we deduce that the morphism~\eqref{eqn:S effective}
	has relative dimension $d^2$ at the point $x_0$ in its domain.
	Thus the morphism~\eqref{eqn:U promotion} 
	has relative dimension $d^2$ at the point $u_0$ in its domain.
	Since this latter morphism is also smooth, its relative dimension
	is locally constant on its
	domain~\cite[\href{https://stacks.math.columbia.edu/tag/0DRQ}{Tag
    0DRQ}]{stacks-project},
       	and thus (since $S$
	is a local ring, so that $\Spec S$ is connected),
	we see that~\eqref{eqn:S effective}
	in fact has relative dimension~$d^2$.

	%and of relative dimension~$d^2$.
	%\mar{ME: This relative dimension statement
	%	is coming from the corresponding fact for the versal rings
	%	of $\cX_d$.  Have to explain this, wihtout going to far
	%	down a dimension-theory rabbit-hole in the formal situation.
	%ME: Idea is to deduce this from Prop.~\ref{prop:versal rings}.}

	Let $\widetilde{M}$ denote the universal \'etale
	$(\varphi,\Gamma)$-module over $S$, and~$\widetilde{N}$ denote
        the base change of~$N$ to~$S$.
	Write $\widetilde{\Ext}^2 := H^2\bigl(\cC^{\bullet}(\widetilde{N}^\vee\otimes\widetilde{M})\bigr),$
	and 
write $$Y_r := \{x \in \Spec S \, | \, \dim \kappa(x)\otimes_S \widetilde{\Ext}^2
\geq r \}.$$ We claim that~$Y_r$ has dimension at most
$([K:\Q_p]d(d-1)/2) + d^2 - r. $ In order to see this, we note
that~$(Y_r)_{\Fpbar}$ is the pull-back to~$\Spec S$ via the flat morphism~\eqref{eqn:S effective}
%$\Spec S \to \cX_{d,\red}$
of the locus considered in Theorem~\ref{thm:Xdred
  is algebraic}~(4) %(with~$\alphabar$ the trivial representation), 
%\mar{TG: probably we should base change $Y_r$ to $\Fpbar$ to say this?}
which has dimension at most~$[K:\Q_p]d(d-1)/2 - r. $
The required bound on the dimension of $Y_r$ follows, for example, from
  \cite[\href{https://stacks.math.columbia.edu/tag/02RE}{Tag
    02RE}]{stacks-project}.

%\mar{ME: Now bound dimension of $Y_r$ from above. TG: hopefully easy
%  to finish from here?! Maybe we are supposed to use something like
%  \cite[\href{https://stacks.math.columbia.edu/tag/02RE}{Tag
%    02RE}]{stacks-project}, or maybe it's easier? ME: Looks good,
% added this to argument and removed this comment}

	The composite
	$\Spf R/\varpi \to 
	\Spf R^{\square,\cO'}_{\rhobar}
	\to \cX_d$
	is effective
        by Corollary~\ref{cor: crystalline deformation rings are effective
    versal},
        i.e.\ arises from a morphism
	$\Spec R/\varpi \to \cX_d.$
	Thus the composite
	$\Spf (R/\varpi)_{\red} \to
	\Spf R/\varpi \to \cX_d$ 
	arises from a morphism
	$\Spec (R/\varpi)_{\red} \to \cX_d,$
	and hence factors through
	$\cX_{d,\red}$.  
	Consequently, we find that the surjection
	$R^{\square,\cO'}_{\rhobar} \to (R/\varpi)_{\red} $
	factors through $S$.
	The closed immersion $\Spec (R/\varpi)_{\red} \hookrightarrow \Spec S$
	then induces a closed immersion $\overline{X}_r \hookrightarrow Y_r$
	for each $r \geq 0$.  The upper bound on the dimension 
	of~ $Y_r$ computed in the preceding paragraph then provides
	the desired upper bound on the dimension of $\overline{X}_r$,
	completing the proof of the theorem.
%
%        
%	
%
%
%	Sketch:  We have $\Spf R^{\square,\cO'}_{\rhobar} \to \cX_d$ versal
%	inducing $\Spf S \to \cX_{d,\red}$ versal.
%	We also have
%	$\Spec (R/\varpi)_{\red} \to \cX_{d,\red},$ by versality;
%	thus the surjection $R^{\square,\cO'}_{\rhobar} \to (R/\varpi)_{\red}$
%	factors through~$S$.  Since $S \to \cX_{d,\red}$ has
%	$\widehat{GL}_d$ fibres, we get upper bound on dimension
%	of~$Y_d$, the $S$-analogue of~$X_d$,  and hence
%	upper bound on dimension of $X_d$ itself.  Using equidimensionality 
%	of~$R$, get lower bound on codimension of~$X_d$.
\end{proof}



%\chapter{Crystalline lifts and the finer structure of $\cX_{d,\red}$.}
%%\chapter{Properties of $\cX_d$}
%\label{sec:properties}
%Up to this point, we have seen that $\cX_d$ is
%%an Ind-algebraic stack;
%%more precisely, that it can be written as an inductive limit 
%%of algebraic stacks over $\Spf \Zp$ of finite presentation, with
%the transition morphisms being closed immersions.
%%In this section, we will establish some finer properties
%%of $\cX_d$: namely, we will show that it is
%a formal algebraic stack. 
%In this section,
%%and
%we will show that $\cX_{d, \red}$ is
%%of finite presentation over~$\F_p$, and is furthermore
%equidimensional of dimension
%$[K:\Q_p]d(d-1)/2$, and enumerate its irreducible components. 
%In the course of doing this, we will
%also prove a result on the existence of crystalline lifts
%which is of independent interest.

\section{Two geometric lemmas}
In this section we establish two lemmas which will be used in the proof
of Theorem~\ref{thm:crystalline lifts} below. 
We begin with the following simple lemma. % \mar{TG: can we provide any
  % motivation/intuition for these lemmas?}

\begin{lemma}
	\label{lem:torsion support}
	Suppose that $0 \to \cF \to \cG \to \cH \to 0$ is 
	a short exact sequence of coherent sheaves on a 
	reduced Noetherian scheme $X$, with $\cF$ and $\cG$ being locally free,
	and $\cH$ being torsion, in the sense that its support
	does not contain any irreducible component of $X$.
	Then there is an effective Cartier divisor~$D$ which contains
	the scheme-theoretic support of $\cH$, and whose 
	set-theoretic support coincides with the set-theoretic support
	of $\cH$,
	with the property that if $f: T \to X$
	is any morphism which meets $D$ properly,
	in the sense that the pull-back $f^*\cO_X \to f^*\cO_X(D)$ 
	of the canonical morphism $\cO_X \to \cO_X(D)$ is injective,
	then $\mathbb L_i f^*\cH = 0$ if $i  > 0$.
\end{lemma}
\begin{proof}
	Replacing $X$ by each of its finitely many connected components
	in turn, we see that it is no loss of generality to assume
	that $X$ is connected, and we do so; this ensures that
	each locally free sheaf on $X$ is of constant rank.
	If  $\eta$ is any generic point of $X$, then $\cH_{\eta} = 0$
	by assumption, and so $\cF_{\eta} \to \cG_{\eta}$ is an isomorphism.
	In particular, $\cF$ and $\cG$ are of the same rank,
        say~$r$, %\mar{TG: the rank is only constant on connected
          %components, so perhaps we should either assume connectedness
        %(it's irreducible in the application) or make a remark that we
      %reduce to that case? ME: Remark added.}
	and the induced morphism
	\numequation
	\label{eqn:wedge map}
	\wedge^r \cF \to \wedge^r \cG
\end{equation}
	is an injection of invertible sheaves. (Indeed, it follows
        from~\cite[III, \S7, Prop.\ 3]{MR1727844} that a map of
        locally free sheaves of rank~$r$ is injective if and only if
        the map on~$\wedge^r$ is injective.)  %\mar{TG: probably worth
          %checking I'm not being an idiot; the reference is to the
          %statement that over any commutative ring, an endomorphism of
        %a free module is injective iff its determinant is not a zero divisor. ME: That seems like an equivalent statement to what's written here, so I think
%this is fine.}
	If we let $D$ denote the support of the cokernel of~(\ref{eqn:wedge
	       map}),	
        then $D$ is an effective Cartier divisor,
	and twisting by $\wedge^r \cF^{\vee}$
	identifies~(\ref{eqn:wedge map}) with the canonical section
	$\cO_X \hookrightarrow \cO_X(D)$.   The scheme-theoretic
	support of $\cH$ is contained in $D$, and the set-theoretic
	supports of $\cH$ and of $D$ coincide (this is a special
	case of the general fact that the Fitting ideal is
	contained in the annihilator, and has the same radical as it).

	Now suppose that $f:T \to X$
	meets $D$ properly in the sense described in the statement
	of the theorem.  Then we see that
	$$\wedge^r f^*\cF \to \wedge^r f^*\cG$$
	is injective, and thus that
	$f^*\cF \to f^*\cG$ is also injective. 
	Consequently
	$$0 \to f^*\cF \to f^*\cG \to f^*\cH \to 0$$ is
	short exact, and so $\mathbb L_i f^* \cH = 0$ for $i > 0,$
	as claimed.
\end{proof}


We now introduce some notation related to the second of the lemmas
(Lemma~\ref{lem:sections} below).
%To this end, we 

\begin{hyp}\label{hyp: codimension hypothesis}
  Let $X$ be a Noetherian scheme, and suppose that we have a short
  exact sequence of coherent sheaves \numequation
  \label{eqn:exact sequence}
  0 \to \cF \to \cG \to \cH \to 0,
\end{equation}
such that
\begin{itemize}
\item $\cG$ is locally free; %(so that $\cA$ is torsion-free);
\item $\cH$ %is torsion, with
  has the property that, for each $r\geq 1,$ the locus
  \[X_r:= \{x \in X \, | \, \dim \kappa(x) \otimes_{\cO_X} \cH \geq
    r\}\] is of codimension $\geq r+1$ in $X$.
\end{itemize}
\end{hyp}
In the setting of Hypothesis~\ref{hyp: codimension hypothesis}, if $\pi:\tX \to X$ is a morphism of Noetherian
schemes, then
we write $\tcG := \pi^* \cG$ and $\tcH := \pi^*\cH$,
and we let $\tcF$ denote the image of the pulled-back morphism $\pi^* \cF \to \pi^* \cG = \tcG$,
%(so that $\tcA$ is the maximal
%torsion-free quotient of $\pi^*\cA$), and 
%we define $\tcC$ to be the quotient $\tcB/\tcA$,
so that we again have a short exact sequence
\numequation
\label{eqn:tilde ses}
0 \to \tcF \to \tcG \to \tcH \to 0.
\end{equation}

\begin{remark} If $\cE$ is a locally free coherent sheaf on a 
	Noetherian scheme~ $X$, and $s:\cO_X \to \cE$ is a morphism,
	then we can think of $\cE$ as being the sheaf of sections
	of a vector bundle over $X$, and think of $s$ as being 
	a section of this vector bundle.  We can then speak of the zero locus
	of $s$; it is the closed subscheme of $X$ which locally,
        if we choose an isomorphism $\cE \cong \cO_X^r$ and write
$s = (a_1,\ldots,a_r),$ is cut out by the ideal sheaf $(a_1,\ldots,a_r)$;
more globally, it is cut out by the ideal sheaf which is the
image of the composite
$$\cE^{\vee}=\cO_X\otimes_{\cO_X}\cE^{\vee} \buildrel s\otimes \id \over \longrightarrow 
\cE\otimes_{\cO_X}\cE^{\vee} \to \cO_X,$$
where the second arrow is the canonical pairing.

We note that $s$ is nowhere-vanishing, i.e.\ the zero locus of $s$ is empty,
if and only if  $s$ is a split injection (so that
$s$ gives rise to a morphism of vector bundles, rather than merely
of sheaves).
We also note that
if $\cE$ has rank $r$, then the Hauptidealsatz shows that
the zero locus of a section of $\cE$ has codimension at most $r$ around
each of its points. % \mar{FH:"s is non-vanishing iff coker(s) is locally free" 
% isn't technically correct: s could be zero... but what you mean is e.g. 
% coker is locally free of rank 1 less (or that it's a split injection).}
\end{remark}

\begin{lem}%\mar{TG: it might be good to recall the hypotheses here}
	\label{lem:sections}
	Suppose that we are in the setting of Hypothesis~{\em \ref{hyp:
          codimension hypothesis}}, that the $\pi: \tX \to X$ is a surjective
morphism whose domain is a Noetherian scheme,
and that $\tcF$ is locally
	free.  Suppose further that $\cO_X \to \cF$ is a morphism with
	the property that the pulled-back morphism $\cO_{\tX} \to \tcF$
%	has locally free cokernel. 
is a nowhere-vanishing section.
	Then the composite 
	$\cO_X \to \cF \to \cG$ is nowhere-vanishing. %has locally free cokernel.
\end{lem}
\begin{proof} The statement may be checked by working
	locally at each of the points of~$X$; more precisely, we may replace
$X$ by $\Spec \cO_{X,x}$ and $\tX$ by its pull-back over $\Spec \cO_{X,x}$.  
We have the stratification
	of $X$ by the closed subsets $X_r$ (where in the
case $r = 0$, we declare $X_0:= X$),
and if $x \in X$, then we
	set $r(x) := \dim \kappa(x)\otimes_{\cO_X} \cH$, or equivalently,
	the maximal value of $r$ for which $x \in X_r$; we then prove the
	theorem by induction on $r(x)$.  In fact, we assume that 
        the statement is true  after localizing at points $x$ for which $r(x) < r$ for {\em all} exact sequences~(\ref{eqn:exact
		sequence}) which satisfy the hypotheses of the lemma, and prove it for our given
	exact sequence after localizing at an $x$ for which $r(x) =
        r$. %(Note that the inductive hypothesis is vacuous in the case $d = 0$.)

In the case $r = 0$, the sheaf $\cH$ is zero, %in a neighbourhood of $x$,
so that $\cF \iso \cG$. % (in this neighbourhood, which we may assume is all of $X$,  since were are working  locally around~$x$). 
Thus the induced morphism $\cO_{\tX} \to \tcG$
gives a nowhere-vanishing section.  Since $\pi$ is surjective, this ensures
that the induced morphism $\cO_X \to \cG$ is also nowhere-vanishing, as required.

We now consider the case $r \geq 1$.
	Since we have replaced $X$ by $\Spec \cO_{X,x}$,
we may assume that $\cG$ is free.
%	say $\cG = \cO_X^n$.
By Nakayama's lemma and our assumption
        that~$r(x)=r$, we may choose a surjection
%\mar{FH says this - deal with this verbatim!}
%\begin{verbatim}
%to get (6.2.8) modding out by the maximal torsion object does not seem to 
%be the right thing. here's how i understand it:
%
%we have two exact sequences with cokernel H and a connecting arrow. now 
%pull the whole diagram back via pi^*. note that L_1 pi^* of the middle 
%term is zero in both rows, as the modules are locally free, hence flat.
%
%that means that \tilde F = pi^* F/L_1pi^*H and similarly \tilde K = pi^* 
%K/L_1pi^*H, where \tilde K is defined as earlier in the section, namely 
%image of pi^*K in pi^*(O_X^d).
%
%now just mod out L_1pi^*H on both sides of the equation just before 
%(6.2.8) to get (6.2.8).
% this way \tilde K is defined correctly so that induction applies below.
%\end{verbatim}
%
%
	\numequation
	\label{eqn:surjection}
        \cO_X^r \to \cH,
\end{equation}
	which we may then lift to a morphism $\cO_X^r \to \cG$.
        If we let $\cK$ denote the kernel of~(\ref{eqn:surjection}), 
then we obtain a morphism of short exact sequences
\numequation
\label{eqn:K to F map}
\xymatrix{
0 \ar[r] & \cK \ar[r] \ar[d] & \cO_X^r \ar[r]\ar[d] & \cH \ar[r]\ar@{=}[d] & 0
\\ 
0 \ar[r] & \cF  \ar[r] & \cG \ar[r] & \cH \ar[r] & 0
}
\end{equation}
Pulling back this diagram back via $\pi$, and letting~$\tcK$ denote
the image of~$\pi^*\cK$ in~$\cO_{\tX}^r$,
we  obtain a corresponding morphism of short exact sequences
\numequation
\label{eqn:tK to tF map}
\xymatrix{
0 \ar[r] & \tcK \ar[r] \ar[d] & \cO_{\tX}^r \ar[r]\ar[d] & \tcH \ar[r]\ar@{=}[d] & 0
\\ 
0 \ar[r] & \tcF  \ar[r] & \tcG \ar[r] & \tcH \ar[r] & 0
}
\end{equation}
The morphism of short exact sequences~\eqref{eqn:K to F map}
induces a short exact sequence
\numequation
\label{eqn:K to F plus O to G}
	0 \to \cK \to \cF \oplus \cO_X^r \to \cG \to 0,
\end{equation}
	which is furthermore split (since $\cG$ is free and we are over an
	affine scheme), so that we may write
\numequation
\label{eqn:direct sum iso}
	 \cF \oplus \cO_X^r \cong \cK \oplus \cG.
\end{equation}
Correspondingly, the short exact sequence~\eqref{eqn:tK to tF map}
induces a short exact sequence
	$$0 \to \tcK \to \tcF \oplus \cO_{\tX}^r \to \tcG \to 0,$$
which can be thought of as being obtained from~\eqref{eqn:K to F plus O to G}
by pulling back via $\pi$ and then taking the quotient by the (naturally identified)
copies of $\mathbb L_1 \pi^*\cH$ sitting inside $\pi^* \cK$ and $\pi^* \cF$.   
Since~\eqref{eqn:K to F plus O to G} is split, so is this latter
short exact sequence, and so we also obtain an isomorphism
	\numequation
	\label{eqn:pulled-back iso}
	\tcF \oplus \cO_{\tX}^r \cong \tcK \oplus \tcG
\end{equation}
which can be thought of as being obtained via by pulling back the
isomorphism~\eqref{eqn:direct sum iso} along~$\pi,$
and then taking the quotient of each side by the appropriate
copy of $\mathbb L_1 \pi^* \cH$.
Since all the other summands appearing in the isomorphism~\eqref{eqn:pulled-back
iso} are locally free, we see that the same is true of~$\tcK$.

	%(where $\tcK$ denotes the maximal torsion free quotient
	%of $\pi^* \cK$). In particular, $\tcK$ is locally free.
        
	Suppose, by way of obtaining a contradiction, that 
	the composite $\cO_X \to \cF \to \cG$
	is not nowhere-vanishing;
	%does not have locally free cokernel;
	then we see that its image lies in $\mathfrak m_x \cG.$
        Thus the pulled back morphism $\cO_{\tX} \to \tcG$ 
        vanishes at each point $\tx$ lying over $x$ (and 
there is at least one such point, since $\pi$ is surjective by assumption).
	Since the pulled-back morphism $\cO_{\tX} \to \tcF$
	{\em is} nowhere-vanishing,
	%{\em does} have locally free cokernel by assumption,
	so is the induced morphism
        $\cO_{\tX}\to \tcF \oplus \cO_{\tX}^r.$
	A consideration of~(\ref{eqn:pulled-back iso}) then shows
	that the induced morphism $\cO_{\tX} \to \tcK$,
	which is pulled-back from the induced morphism $\cO_X \to \cK,$
	must be nowhere-vanishing.
       	%must have locally free cokernel.

	Thus, if we consider the exact sequence
	$$0 \to \cK \to \cO_X^r \to \cH \to 0,$$
	it satisfies Hypothesis~\ref{hyp: codimension hypothesis}, % the same hypotheses as~(\ref{eqn:exact sequence}),
	and we are given a morphism $\cO_X\to \cK$ whose pull-back to $\tX$
	is a nowhere-vanishing section of $\tcK$.
	%embeds $\cO_{\tX}$ into $\tcK$ with locally free cokernel.
        Applying the inductive hypothesis to this situation, we find 
	that the composite $\cO_{X} \to \cO_{X}^r$ is a section
	of a rank $r$ locally free sheaf
	%vector bundle
	whose zero locus is contained
        in~$X_r$ (because if we localise at a point not in~$X_r$, then
        the section is nowhere vanishing by the inductive hypothesis), and  is therefore of codimension
	at least $r+1$ by hypothesis.  This zero locus contains $x$ (since this
	section factors through $\cK$, and $\cH \cong \cO_X^r/\cK$
	has fibre dimension exactly $r$ at $x$, so the map
        $\cO_X^r\otimes_{\cO_X}\kappa(x)\to\cH\otimes_{\cO_X}\kappa(x)$
        is an isomorphism), and
	hence is non-empty. %, and thus is of codimension exactly $d+1$.
	On the other hand, as we noted above,
	the zero locus of a section of a rank $r$ locally free sheaf,
	if it is non-empty, has codimension at most~$r$. 
	This contradiction completes the proof of the lemma.
\end{proof}


\section{Crystalline lifts} % and the dimension of $\cX_{d,\red}$}
% \mar{TG: note to self - niveau-one.tex is probably a good thing to
%   read to get this argument in my head.}

% \mar{ME: We eventually need a version in which $\chibar$ is
% 	generalized to higher dimensional irreps. ME: Now added this;
% hopefully it's correct.}
		
Given a~$d$-tuple of labeled Hodge--Tate weights~$\underline{\lambda}$, and
a~$d'$-tuple of labeled Hodge--Tate weights~$\underline{\lambda}'$, we say
that~$\underline{\lambda}'$ is \emph{slightly greater than~$\underline{\lambda}$} (and that $\underline{\lambda}$ is
\emph{slightly less than~$\underline{\lambda}'$}) if for each
$\sigma:K\into\Qpbar$ we have $\lambda'_{\sigma,d}\ge \lambda_{\sigma,1}+1$, and
the inequality is strict for at least one~$\sigma$. 
% \mar{ME: Don't want to change anything right now, b/c I'm a bit tired and punchy,
% but it seems that the roles of $\lambda$ and $\lambda'$ are reversed in this
% def'n, compared to in the statement of the Thm.\ just below, whereas it might be 
% better if they matched. TG: swapped, commenting out}

This is not standard terminology, but it will be convenient for us; it
is motivated by the following well-known result. Here, as throughout
the book, we
slightly abuse terminology and refer to lattices in crystalline
representations as themselves being crystalline. % as in Definition~\ref{defn: abuse of terminology crystalline lattice},
% \mar{ME: As in \S \ref{subsec: formal algebraic stack}, we should explain/recall
% that ``crystalline lift'' means ``integral lattice in a crystalline lift''.} 
\begin{lem}
  \label{lem: extensions are crystalline}Let
  $0\to\rho_1^\circ\to\rho^\circ\to\rho_2^\circ\to 0$ be an extension of
  $\Zpbar$-valued representations of~$G_K$, with~$\rho_1^\circ$
  and~$\rho_2^\circ$ being crystalline of Hodge--Tate
  weights~$\underline{\lambda}_1$, $\underline{\lambda}_2$
  respectively. If~$\underline{\lambda}_1$ is slightly less
  than~$\underline{\lambda}_2$, then~$\rho^{\circ}$ is 
  crystalline.% \mar{TG: warning that HT weights might be the wrong way
    % around. Also as I've written Theorem~\ref{thm:crystalline lifts}
    % right now, we want the pst version of this too. Don't need pst,
    % do need to change sign and change to ``slightly''.}
\end{lem}
\begin{proof}This follows easily from the formulae in~\cite[Prop.\
  1.24]{MR1263527}. Indeed, if we let~$\rho_1$, $\rho_2$ be
  the~$\Qpbar$-representations corresponding to~$\rho_1^\circ$,
  $\rho_2^\circ$, and we set $V:=\rho_1\otimes\rho_2^\vee$, then we
  need to show that $h^1_f(V)=h^1(V)$. Since the Hodge--Tate weights
  of~$V$ are all negative, and for at least one
  embedding~$\sigma:K\into\Qpbar$ the $\sigma$-labeled Hodge--Tate
  weights are all less than~$-1$, we have
  $h^1(V)-h^1_f(V)=h^2(V)=h^0(V^\vee(1))=0$, as required.
\end{proof}

%We now prove our main theorems on crystalline lifts, beginning with
% \mar{ME: Change phrasing? just b/c the subsequent results are
% in a new subsection. TG: how about this? I incorporated the remark
% that had been just below the theorem}
% The following result will allow us to construct crystalline lifts by
% induction on the dimension.
Arguing inductively,
the following theorem will allow us to construct 
crystalline lifts of any given $\rhobar$. % with more or less 
% arbitrary Hodge--Tate weights (compatible with the inertial weights 
% of $\rhobar$), and
In particular it implies Theorem~\ref{thm:intro crystalline lifts} from
the introduction, in the more refined form of Theorem~\ref{thm: strong
  existence of crystalline lifts} below.
\begin{theorem}
	\label{thm:crystalline lifts}
	Suppose given a representation $\rhobar_d: G_K \to \GL_d(\Fbar_p)$
	that admits a lift
	$\rho_d^{\circ} : G_K \to \GL_d(\Zbar_p)$ which 
%	is a lift of $\rhobar_d$ such that
        % for which the associated
	% $p$-adic representation $\rho_d: G_K \to \GL_d(\Qbar_p)$
	is crystalline with labeled Hodge--Tate weights
        $\underline{\lambda}$. 
% \mar{ME: Add condition on HT wts.\ that ensure $H^2
%           = 0$. TG: I think that since we changed from the trivial
%           representation, we can't just write a condition like
%           this. Instead I've baked in crystalline conditions for both
%           factors, and an inequality on the weights. Maybe we should
%           do something a bit more general, but I'm inclined to stick
%           with at least potentially semistable with weights satisfying
%         this condition?} 
	% \mar{ME: Ultimately need to generalize from extension by $\triv$
	% 	to an extension by an irreducible rep.\ of arb.\ dim'n. 
	% ME: I've now done this (hopefully without blundering).}
	Let $0 \to \rhobar_d \to \rhobar_{d+a} \to \alphabar \to 0$
	be any extension of $G_K$-representations over $\Fbar_p$,
	with $\alphabar: G_K \to \GL_a(\Fbar_p)$ irreducible,
	and let $\alpha^{\circ}: G_K \to \GL_a(\Zbar_p)$ be any
        crystalline lifting
	of $\alphabar$ with labeled Hodge--Tate
        weights~$\underline{\lambda}'$, which we assume to be slightly greater than~$\underline{\lambda}$.


	Then we may find a lifting of the given extension to an extension
	$$0 \to \theta^{\circ}_d \to \theta_{d+a}^{\circ} \to \alpha^{\circ} \to 0$$
	of $G_K$-representations over $\Zbar_p$, where
        $\theta^{\circ}_d: G_K \to \GL_d(\Zbar_p)$ again has the
        property that the associated $p$-adic representation
        $\theta_d: G_K \to \GL_d(\Qbar_p)$ is crystalline with
        labeled Hodge--Tate weights $\underline{\lambda}$.
        Furthermore, $\theta_{d+a}^{\circ}$ is crystalline, and we may
        choose $\theta_d^{\circ}$ to lie on the same irreducible
        component of
        $\Spec R_{\rhobar_d}^{\crys,\underline{\lambda}}$ that
        $\rho^{\circ}_d$ does.
	% \mar{ME: We would like to actually have some relationship b/w properties
	% 	of $\theta_d^{\circ}$ and $\rho^{\circ}_d$, e.g.\ the former can
	% 	be taken to be niveau $1$ if the latter is.   But I don't know
	% 	how to do this yet.}
\end{theorem}

% \begin{remark}
% 	\label{rem:crystalline lifts}
% % 	In the context of the preceding theorem, if $\alpha^{\circ}$ is 
% % 	chosen so that the associated $p$-adic representation 
% % 	$\alpha: G_K \to \GL_a(\Qbar_p)$ is crystalline,\mar{TG: I
% %           just baked this condition in\dots In fact if this is OK, the
% %         whole remark should be adjusted. Do we ever want the lifts to
% %         be anything more general than de Rham? }
% % 	and if the tuple of weights
% % 	$\underline{\lambda}$ satisfies \mar{ME: write down some positivity condition
% % 	here related to the HT weights of $\alpha$},
% % then the $p$-adic representation $\theta_{d+a}:
% % G_K \to \GL_{d+a}(\Qbar_p)$ associated to $\theta_{d+a}^{\circ}$ is necessarily
% % crystalline.  Thus, 
% Arguing inductively,
% Theorem~\ref{thm:crystalline lifts} allows us to construct 
% crystalline lifts of any given $\rhobar$. % with more or less 
% % arbitrary Hodge--Tate weights (compatible with the inertial weights 
% % of $\rhobar$), and
% In particular it implies Theorem~\ref{thm:intro crystalline lifts} from
% the introduction, in the more refined form of Theorem~\ref{thm: strong
%   existence of crystalline lifts} below.
% \end{remark}


% We now present the proofs of these results.
% The logic of induction is that we assume that
% Theorem~\ref{thm:reduced dimension} holds in the rank $\leq d$ case,
% and that Theorem~\ref{thm:dimension of H2} holds for ranks $< d$. 
% We then deduce
% Theorems~\ref{thm:dimension of H2} and~\ref{thm:crystalline lifts}
% in the rank $d$ case.   From this we then deduce Theorem~\ref{thm:reduced
% 	dimension} in the rank $d+1$ case, completing the inductive
%       step.\mar{TG: probably worth saying that we start by assuming
%         the trivial $d=0$ case?}





		%We now use Theorem~\ref{thm:dimension of H2} to
		%prove the rank $d+1$ case of Theorem~\ref{thm:reduced
		%	dimension}.
		%For this,  we first prove a more precise version
		%of the statement about
		%crystalline lifts given in Theorem~\ref{thm:intro crystalline
		%	lifts}.  We use the rank $d$ case 
		%of Theorem~\ref{thm:dimension of H2} as input.

\begin{proof}[Proof of Theorem~{\em \ref{thm:crystalline lifts}}]
%\mar{ME: Deal with field of definition stuff, i.e.\ choose $E$ appropriately.}
We may and do choose our field of coefficients $E$ to be large enough
that the various Galois representations are defined over $\cO$ or $\F$
as the case may be.
	We write $R := R_{\rhobar}^{\crys,\underline{\lambda}}$, 
	and we let $X$ denote the component of the crystalline lifting scheme 
	%$$X := X_{\rhobar}^{\crys,\underline{\lambda}} =
	%\Spec R.$$
	$\Spec R_{\rhobar}^{\crys,\underline{\lambda}}$ on which $\rho^{\circ}_d$
	lies. As recalled in Section~\ref{subsec: notation and
          conventions}, $X$ is reduced. % \mar{TG: is it recalled there?
        % I guess the point is $\cO$-flat and regular generic fibre? ME: I added
        % something there.}
	% \mar{ME: The following explanation needs to be smoothed out.  The
	% 	main reason I want to work on $\Spec R$ rather than $\Spf R$
	% 	is to not have to think about the geometric lemma over formal
	% 	schemes. }
%        Fix the lift $\alpha^{\circ}$ of $\alpha$.
	Over $X$ we have a good complex $C^{\bullet}$ supported in degrees~$[0,2]$ computing
	$\Ext^{\bullet}_{G_K}(\alpha^{\circ},\rho^{\circ})$ for the
        universal deformation~$\rho^{\circ}$; indeed, by
a straightforward variant of Corollary~\ref{cor:Herr complex length},
%Theorem~\ref{thm: Herr complex and Galois reps},
%\mar{ME: Maybe cite Cor.~5.1.25 here?  But also figure out
%which citations are valid, just b/c the def.\ ring $R$ is a complete
%local ring, and so is not top.\ of f.t.\ over $\cO$ when we think
%of it as a $p$-adically completee ring.}
we can choose a good complex
        supported in degrees~$[0,2]$ and quasi-isomorphic to the Herr
        complex for~$\rho^\circ\otimes(\alpha^{\circ})^\vee$.  Our assumption
	on the Hodge--Tate weights % tuple of weights $\underline{\lambda}$
        implies that $\Ext^2$ is
	supported (set-theoretically) on the special fibre $\overline{X}$
	of $X$. Indeed, the formation of~$\Ext^2$ is compatible with
        base change, and at any closed point~$x$ of the generic fibre
        of~$R$, we have
        $\Ext^{2}_{G_K}(\alpha^{\circ},\rho^{\circ}_x)=\Hom_{G_K}(\rho^{\circ}_x,\alpha^{\circ}(1))$
        by Tate local duality, and this space vanishes by the
        assumption that~$\underline{\lambda}'$ is slightly greater than~$\underline{\lambda}$.

	As usual, we let $B^2$ denote the image of $C^1$ in $C^2$; 
	it is a subsheaf of the locally free sheaf $C^2$, and so is 
%         \mar{TG: should probably add an explanation somewhere of
%           why~$X$ is reduced? Should be coming from~$R$ being flat
%           over~$\Zp$ and having regular generic fibre, I think.
%   ME: Flatness says that $R$ embeds into $R[1/p]$, and these rings
% are reduced, since $R$ is generically regular, and so in particular
% generically reduced. This is standard, I guess, but we can recall
% it back when we first introduce crystalline def.\ rings.}
	torsion free.\footnote{Since $X$ is reduced, {\em torsion free}
		is a reasonable notion for a finitely
		generated module, equivalent to the associated
		primes lying among the minimal primes.}
	By~ \cite[\href{http://stacks.math.columbia.edu/tag/0815}{Tag
          0815}]{stacks-project}, we may find a blow up $\pi:\tX \to X$, whose centre lies
	(set-theoretically) in the special fibre of $X$, % \mar{TG: do we
%         have a reference for this? I presume it relates to
%         \cite[\href{http://stacks.math.columbia.edu/tag/080X}{Tag
%           080X}]{stacks-project}. TG: probably
%         \cite[\href{http://stacks.math.columbia.edu/tag/0815}{Tag
%           0815}]{stacks-project} but I need to think this stuff through.
% ME: That looks like the right reference, taking $X = S = \Spec R$
% and taking $U$ to be the complement of the locus where $p = 0$.}
	such that the torsion-free quotient of $\pi^*B^2$ becomes
	locally free (recall that~$\Ext^2$ is supported on the special
        fibre of~$X$); in other words, if we let $\tC^{\bullet}$
	denote the pull-back by $\pi$ of $C^{\bullet}$, 
	the corresponding $\cO_{\tX}$-module of $2$-coboundaries
	$\tB^2$ (that is, the image of~$\tC^1$ in~$\tC^2$)
	is locally free. Thus, if we let $\tZ^1$ denote the $\cO_{\tX}$-module of
	$1$-cocycles for $\tC^{\bullet}$, then $\tZ^1$ is also locally
	free; so in particular, the complex $\tC^0 \to \tZ^1$ is a
        good complex. %\mar{TG: we don't seem to explicitly use this fact. TG: I
        %think this is still the case - just for my understanding, are
        %we actually using it somewhere? ME: No, I don't think we do,
        %so I'm commenting it out, along with this comment. TG:
        %amusingly, this is exactly what FH wanted us to add, it turns
        %out! }

	The given extension $\rhobar_{d+a}$ is classified by a
	class in $\Ext^1_{G_K}(\alphabar,\rhobar_d)$, which is to say, a class
	in $H^1$ of the complex $\kappa\otimes C^{\bullet}$.
	(Here we write $\kappa$ to denote $\F$ thought of as the
	residue field of $R$.)
        Lifting this class to an element of $\kappa \otimes C^1,$
	and then to an element of $C^1$, we find ourselves in the 
	following situation:
	we have a morphism $c:R \to C^1$, whose image under the coboundary
	lies in $\mathfrak m_R C^2$ (reflecting the fact that we have a cochain
	which becomes a cocycle at the closed point). % \mar{TG: see
          % FLorian emails}

	Consider the composite $b: R \buildrel c \over \longrightarrow
	C^1 \to B^2,$ 
	which pulls back to a section $\tb: \cO_{\tX} \to \tB^2.$
	It follows from Lemma~\ref{lem:sections} above (applied with
        $\cF=B^2$, $\cG=C^2$, and $\cH=\Ext^2$),
	together with Theorem~\ref{thm: codimension H2 in regular
          weight crystalline deformation ring}, %  Theorem~\ref{thm:Xdred is algebraic}~(4),
%         \mar{TG:  replace with a reference to an explanation above
%           of why this holds on~$\Spec R$ - note that the result cited
%           is a statement about dimension rather than codimension, and
%           we don't know that ~$\cX_d$ is equidimensional at this
%           point, so we may have to do more work! ME: Actually,
%   will cite
% Theorem~\ref{thm: codimension H2 in regular weight crystalline deformation ring}, which will already be formulated with the correct codimension statement.}
	% \mar{ME: Have to explain that  Theorem~\ref{thm:dimension of H2} actually
	% 	does control $H^2$, because of Tate local duality.}
	that the lifted section~$\tb$ must have non-empty zero locus,
	which (being closed) must contain a point $\tx$ lying
	over the closed point ~$x\in X$.   The section $c$ itself
	pulls back to a section $\tc: \cO_{\tX} \to \tC^1$, whose value at the 
	point $\tx$ lies in the fibre of $\tZ^1$.   In other
	words, the fibre of $\tc$ at $\tx$ is a one-cocycle
	in the complex $\kappa(\tx)\otimes [\tC^0 \to \tZ^1]$, giving rise to
	a class
	$\overline{e} \in H^1\bigl(\kappa(\tx)\otimes [\tC^0 \to \tZ^1]\bigr)$
	lifting the original class in $\Ext^1_{G_K}(\alphabar,\rhobar_d)$ that classifies~
	$\rhobar_{d+a}$.

	Since $\tX$ is flat over $\Z_p$ (because~$X$ is), % \mar{TG: presumably this is
%           obvious from the definition as a blowup, but since I can't
%           ever remember how that works I'm leaving this as a note to myself.
%   ME: Blowing up something flat over $\Z_p$ gives something flat over
% $\Z_p$, as follows e.g.\ from an explicit consideration of the relevant Proj.}
we may find a morphism
	$\tif:\Spec \Zbar_p \to \tX$
	passing through the point $\tx \in \tX(\Fbar_p)$ (in the sense
        that the closed point of $\Spec\Zbar_p$ maps to~$\tx$). 
        % \mar{TG: does this just mean that the close point maps
        %   to~$\tx$? Perhaps we should say so if so.  ME: Yep,
	%   and I guess we use this below.}
	The morphism $\tif$ determines (and, by the valuative
	criterion of properness, is determined by) a morphism 
	$f: \Spec \Zbar_p \to X$,
	lifting the closed point $x \in X$, % \mar{TG: do you mean
          % lifting the closed point of~$X$? ME: Yes, fixed.}
	which in turn corresponds to a representation
	$\theta_d^{\circ}: G_K \to \GL_d(\Zbar_p),$ as in the statement
	of the theorem.
	Now $H^2(\tC^{\bullet})$ is the cokernel of an inclusion of
	locally free sheaves (the inclusion $\tB^2 \hookrightarrow
        \tC^2$), and it is torsion (because it is pulled back
        from~$H^2(C^{\bullet})$, which is set-theoretically supported in the special fibre)
	and so Lemma~\ref{lem:torsion support} above shows that
	there is an effective Cartier divisor $D$ contained (set-theoretically)
	in the
        special % \mar{TG: is ``contained in'' a statement about
  %         set-theoretic support (in which case this is fine)? ME: Yep,
  % added clarification.}
	fibre of $\tX$ with the property that for any morphism
	to $\tX$ that meets $D$ properly, the higher derived
	pull-backs of $H^2(\tC^{\bullet})$ under this morphism
	vanish.  Since the domain of $\tif$ is $p$-torsion free,
	it meets the special fibre of $\tX$ properly,
	and in particular meets $D$ properly; 
	%	in particular
	%is supported on a Cartier divisor on $\tX$,
  %       \mar{TG: I don't
  %         follow this - can you elaborate? ME: I added a lemma which hopefully
  % makes things clearer.}
        %which furthermore
	        %lies (set-theoretically) on the special fibre of $\tX$ (and thus is,
	        %set-theoretically, a union of irreducible components
	        %of this special fibre).
	thus we infer that
       	%$$\Tor_i^R\bigl(\tif_*\cO_{\Spf \Zbar_p}, H^2(\tC^{\bullet})\bigr) = 0$$
	$$\mathbb L_i\tif^* H^2(\tC^{\bullet}) = 0$$
	for $i > 0.$ From this vanishing we deduce 
	that % \mar{TG: I don't follow this either. ME: I may
% 	have had the wrong statement there; I've rewritten it for the
% higher derived pull-backs by $\tif$, which is what I actually wanted.}
	\begin{multline*}
		\Ext^1_{G_K}(\alpha^{\circ},\theta_d^{\circ}) = H^1(f^*C^{\bullet}) =
	H^1(\tif^*\tC^{\bullet}) = \tif^*H^1(\tC^{\bullet})
	\\
	= \tif^*H^1\bigl( [ \tC^0 \to \tZ^1]) = H^1\bigl(\tif^*[\tC^0 \to
	\tZ^1]\bigr)
	.\end{multline*}
	(The first of these identifications is the general base-change
	property of the perfect complex $C^{\bullet}$, % \mar{TG: ref for
        % this to relevant Herr complex statement}
      the second
	follows from the fact that $f = \pi\circ \tif$,  the third
	follows from the vanishing of $\mathbb L_i\tif^* H^2(\tC^{\bullet}) $
        for $i>0$ and the base-change spectral sequence
%	\mar{TG: can we
%        give a reference for this? ME: Probably,
%but I guess it's also a standard thing. TG suggests:
%\cite[\href{https://stacks.math.columbia.edu/tag/0662}{Tag 0662}]{stacks-project}}
\[E_2^{p,q} := L_{-p}\tif^*H^q(\tC^\bullet)\implies H^{p+q}(\tif^*\tC^\bullet) \] % computing $H^i(f^*\tC^\bullet)$,
        % $\Tor$-vanishing that we just
        % observed
        % \mar{TG: I could also do with this third point being expanded.
	% ME: The third point uses the remarked-upon Tor vanishing together
        % with the base-change spectral sequence for $H^i$ (which, if I 
        % write it for rings rather than schemes,
        % is the second quadrant sequence
        % $E_2^{p,q} := \Tor_{-p}^A\bigl(B,H^q(G_K, \rho_A) \bigr)
        % \implies H^{p+q} (G_K, \rho_B).$}
(see e.g.~\cite[\href{https://stacks.math.columbia.edu/tag/0662}{Tag 0662}]{stacks-project}),
%	\mar{TG: can we
	the fourth holds by definition of $H^1$, and the fifth follows
	from right-exactness of~$\tif^*$.)
	%Thus $H^1*(G_K,\theta_d^{\circ})$ \mar{TG: I think this
        %  sentence fragment should just be removed? ME: I think so, and so
        % I removed it.}

	Since $\tif$ identifies the closed point of $\Spec \Zbar_p$
	with the point $\tx,$ we may find a class
	$e \in H^1\bigl(\tif^*[\tC^0 \to \tZ^1]\bigr)$ lifting the
	class $\overline{e}$, which is then identified with an element
	of $\Ext^1_{G_K}(\alpha^{\circ},\theta_d^{\circ}).$   This element classifies
	an extension $0 \to \theta_d^{\circ} \to \theta_{d+a}^{\circ} 
	\to \alpha^{\circ} \to 0$ which by construction lifts the
        given extension $0 \to \rhobar^{\circ}_d \to \rhobar_{d+a}^{\circ} \to \alphabar^{\circ} \to 0$, and so by Lemma~\ref{lem: extensions are crystalline} satisfies the conclusions of the theorem.
\end{proof}

%Finally, we can prove the rank $d+1$ case of Theorem~\ref{thm:reduced dimension}.
%We now prove the following result by induction on $d$.

       % Let $\rhobar: G_K \to \GL_{d}(\Fbar_p)$ be a continuous representation,
       % corresponding to a point of $\cX_{d+1,\red}(\Fbar_p)$.
       % If $\rhobar$ is irreducible, then we see directly that it
       % admits a crystalline lift whose labeled Hodge--Tate weights
       % \mar{ME: Maybe this argument should be tightened or rephrased,
       %         now that we are discussing these sorts of issues about 
       %         lifts much more carefully further below.}
       % are uniformly bounded (in terms of $d$).
       % On the other hand, if $\rhobar$ is reducible,
       % then this follows from Theorem~\ref{thm:crystalline lifts} 
       % % \mar{ME: Need the (currently unwritten) version of
       % %         that theorem in which $\triv$ is replaced by
       % %         an arbitrary irrep.}
       % (taking into account Remark~\ref{rem:crystalline lifts}), 
       % together with the present theorem in ranks $\leq d$.


% \section{Crystalline lifts and the dimension of $\cX_{d,\red}$ -
%   old version}
% The following results are proved together, by induction on $d$.

% \begin{thm}
% 	\label{thm:reduced dimension}
% 	There is a finite set of tuples $\underline{\lambda}$ of regular labeled Hodge--Tate
% 	weights such that $\cX_{d,\red}$ is the union of corresponding
% 	underlying reduced substacks $\cX_d^{\crys,\underline{\lambda}}.$
% 	In particular,
% 	each irreducible component of $\cX_{d,\red}$ is of dimension
% 	$[K:\Q_p] d(d-1)/2$.
% \end{thm}

% % \mar{ME: We eventually need a version in which $\chibar$ is
% % 	generalized to higher dimensional irreps. ME: Now added this;
% % hopefully it's correct.}
% \begin{thm}
% 	\label{thm:dimension of H2}
% 	If we fix an irreducible representation $\alphabar: G_K \to
% 	\GL_a(\Fbar_p)$,
% 	then the locus of $\rhobar$ in $\cX_{d,\red}(\Fbar_p)$
% 	for which $\dim \Hom_{G_K}(\rhobar,\alphabar) \geq r$ 
% 	is of codimension $\geq \lceil {r\bigl((a^2+1)r-a\bigr)}/{2}
%         \rceil$.\mar{TG: perhaps we should be more precise about what
%           this means, e.g.\ it could be defined and checked on
%           pullbacks to schemes?}
% 	% \mar{ME: I originally thought we could get by with $c_r = r$,
% 	% 	but this seems to be too gross of an underestimate,
% 	% 	and doesn't fit well with the induction.  I think we
% 	% 	need to get closer to the truth for the induction  In the course of doing this, we will
% 	% 	also prove a result on the existence of crystalline lifts
% 	% 	which is of independent interest.
% 	% 	to work, i.e.\ we have to make a good estimate for $c_r$.
% 	% ME: OK, hopefully what I've written works. TG: it seems OK to
%         % me, barring the fact of it not being as big as you claim if $r=1$\dots
%         % ME: I think I had an $r$ vs.\ $r+1$ confusion; see my response below.}
% \end{thm}

% 	\begin{remark}
% 		As will be evident from the proof, the lower bound
% 		of Theorem~\ref{thm:dimension of H2} is quite crude
% 		when $[K:\Q_p] > 1$, although it is reasonably sharp
% 		in the case $K = \Q_p$.
% 		However, it suffices for our purposes, for which it
% 		is sufficient to have a lower bound of $r$ on the codimension,
% 		a lower bound which is certainly satisfied
% 		by $\lceil r\bigl((a^2+1)r-a\bigr)/2\rceil$.% \mar{TG: this is only true
% %                 if $r>1$! Hopefully not an issue? Earlier you were saying~$r$\dots ME: Yes, sorry, we need $r+1$ when we work over $\Z_p$, and (equivalently)
% % 	$r$ when we work on the special fibre.  I think this same
% % confusion (on my part) is underlying your comment above; I'll work through
% % the argument and try to sort this out.}
% 	\end{remark}
		
% Given a~$d$-tuple of labeled Hodge--Tate weights~$\underline{\lambda}$, and
% a~$d'$-tuple of labeled Hodge--Tate weights~$\underline{\lambda}'$, we say
% that~$\lambda'$ is \emph{slightly less than~$\lambda$} (and that $\lambda$ is
% \emph{slightly greater than~$\lambda'$} if for each
% $\sigma:K\into\Qpbar$ we have $\lambda_{\sigma,d}\ge \lambda'_{\sigma,1}+1$, and
% the inequality is strict for at least one~$\sigma$. (This is not
% standard terminology, but it will be convenient for us.)

% \begin{theorem}
% 	\label{thm:crystalline lifts}
% 	Suppose given a representation $\rhobar_d: G_K \to \GL_d(\Fbar_p)$
% 	that admits a lift
% 	$\rho_d^{\circ} : G_K \to \GL_d(\Zbar_p)$
% %	is a lift of $\rhobar_d$ such that
%         for which the associated
% 	$p$-adic representation $\rho_d: G_K \to \GL_d(\Qbar_p)$
% 	is crystalline with labeled Hodge--Tate weights $\underline{\lambda}$.
% \mar{ME: Add condition on HT wts.\ that ensure $H^2
%           = 0$. TG: I think that since we changed from the trivial
%           representation, we can't just write a condition like
%           this. Instead I've baked in crystalline conditions for both
%           factors, and an inequality on the weights. Maybe we should
%           do something a bit more general, but I'm inclined to stick
%           with at least potentially semistable with weights satisfying
%         this condition?} 
% 	% \mar{ME: Ultimately need to generalize from extension by $\triv$
% 	% 	to an extension by an irreducible rep.\ of arb.\ dim'n. 
% 	% ME: I've now done this (hopefully without blundering).}
% 	Let $0 \to \rhobar_d \to \rhobar_{d+a} \to \alphabar \to 0$
% 	be any extension of $G_K$-representations over $\Fbar_p$,
% 	with $\alphabar: G_K \to \GL_a(\Fbar_p)$ irreducible,
% 	and let $\alpha^{\circ}: G_K \to \GL_a(\Zbar_p)$ be any
%         potentially semistable lifting
% 	of $\alphabar$ with labeled Hodge--Tate
%         weights~$\underline{\lambda}'$, which we assume to be slightly greater than~$\underline{\lambda}$.


% 	Then we may find a lifting of the given extension to an extension
% 	$$0 \to \theta^{\circ}_d \to \theta_{d+a}^{\circ} \to \alpha^{\circ} \to 0$$
% 	of $G_K$-representations over $\Zbar_p$,
% 	where $\theta^{\circ}_d: G_K \to \GL_d(\Zbar_p)$
% 	again has the property that the associated
% 	$p$-adic representation $\theta_d: G_K \to \GL_d(\Qbar_p)$
% 	is crystalline with labeled Hodge--Tate weights $\underline{\lambda}$.
% 	Furthermore, $\theta_{d+a}^{\circ}$ is potentially semistable
%         {\em (}and crystalline if~$\alpha^\circ$ is crystalline{\em )}, and we may choose $\theta_d^{\circ}$ to lie on
% 	the same irreducible component of $\Spec R_{\rhobar_d}^{\underline{\lambda},
% 		\crys}$ that $\rho^{\circ}_d$ does.
% 	% \mar{ME: We would like to actually have some relationship b/w properties
% 	% 	of $\theta_d^{\circ}$ and $\rho^{\circ}_d$, e.g.\ the former can
% 	% 	be taken to be niveau $1$ if the latter is.   But I don't know
% 	% 	how to do this yet.}
% \end{theorem}

% \begin{remark}
% 	\label{rem:crystalline lifts}
% % 	In the context of the preceding theorem, if $\alpha^{\circ}$ is 
% % 	chosen so that the associated $p$-adic representation 
% % 	$\alpha: G_K \to \GL_a(\Qbar_p)$ is crystalline,\mar{TG: I
% %           just baked this condition in\dots In fact if this is OK, the
% %         whole remark should be adjusted. Do we ever want the lifts to
% %         be anything more general than de Rham? }
% % 	and if the tuple of weights
% % 	$\underline{\lambda}$ satisfies \mar{ME: write down some positivity condition
% % 	here related to the HT weights of $\alpha$},
% % then the $p$-adic representation $\theta_{d+a}:
% % G_K \to \GL_{d+a}(\Qbar_p)$ associated to $\theta_{d+a}^{\circ}$ is necessarily
% % crystalline.  Thus, 
% Arguing inductively,
% Theorem~\ref{thm:crystalline lifts} allows us to construct 
% crystalline lifts of any given $\rhobar$. % with more or less 
% % arbitrary Hodge--Tate weights (compatible with the inertial weights 
% % of $\rhobar$), and
% In particular it implies Theorem~\ref{thm:intro crystalline lifts} from
% the introduction, in the more refined form of Theorem~\ref{thm: strong
%   existence of crystalline lifts} below.
% \end{remark}


% We now present the proofs of these results.
% The logic of induction is that we assume that
% Theorem~\ref{thm:reduced dimension} holds in the rank $\leq d$ case,
% and that Theorem~\ref{thm:dimension of H2} holds for ranks $< d$. 
% We then deduce
% Theorems~\ref{thm:dimension of H2} and~\ref{thm:crystalline lifts}
% in the rank $d$ case.   From this we then deduce Theorem~\ref{thm:reduced
% 	dimension} in the rank $d+1$ case, completing the inductive
%       step.\mar{TG: probably worth saying that we start by assuming
%         the trivial $d=0$ case?}


% \begin{proof}[Proof of Theorem~\ref{thm:dimension of H2} in the rank $d$ case]
% 	As already explained,
%         we assume that Theorem~\ref{thm:reduced dimension} holds for dimensions
% 	less than or equal to $d$,
% 	and that the present theorem holds for dimensions less than $d$.

% 	In the case $r = 0$, there is nothing to prove.
% 	If $\dim \Hom_{G_K}(\rhobar,\alphabar) \geq r \geq 1,$
% 	then we may place $\rhobar$ in a short exact sequence
% 	$$0 \to \thetabar \to \rhobar \to \alphabar^{\oplus r} \to 0,$$
% 	where $\thetabar$ is of dimension $d- ra < d$.
%         We may apply Theorem~\ref{thm:reduced dimension}
% 		so as to find that $\cX_{d-ar,\red}$ is equidimensional
% 		of dimension $[K:\Q_p](d-ar)(d-ar-1)/2$.
% 		Let $\cU_s$ be the locally closed substack
% 		of $\cX_{d-ar,\red}$ over which
% 		$\dim H^2(G_K,\thetabar\otimes \alphabar^{\vee}) =
%                 s$;\mar{TG: presumably the idea is to stratify by~$s$,
%                 and note that the argument below works for all~$s$? ME: Yes,
% 	that's right. TG: should reference elsewhere in the paper for
%         this being locally closed; also the $\cW$ notation currently conflicts. ME: Changed notation to $\cU_s$ to eliminate conflict.}
% 		by induction, this locus has codimension at least
% 		$s\bigl((a^2+1)s-a\bigr)/2$ in $\cX_{d-a r,\red},$
% 		and over this locus
% 		we may construct a universal family of extensions 
% 		$$0 \to \thetabar \to \rhobar_{\cU_s}
% 		\to \alphabar^{\oplus r} \to 0.$$
% 		The locus of $\rhobar$ we are interested in is
% 		contained in the scheme-theoretic
% 		image of this family in $\cX_{d,\red}$,
% 		and Proposition~\ref{prop:dimension control}
% 		shows that this scheme-theoretic image 
% 		has dimension bounded by
% 		%relative dimension $r([K:\Q_p](d-r) + s)- r^2$
% 		%\mar{ME: Have to explain this above.  The point is that
% 		%	over this locus, $H^2$ has constant rank $s$ and so
% 		%	is locally free, and then it should be true
% 		%	that $H^1$ is generically free of 
% 		%	rank $r([K:\Q_p](d-r) + s)$, which
% 		%	lets us compute the dimension of the universal
% 		%	extension.   Maybe have to think more carefully
% 		%	about what happens when $H^0 \geq 0$ ,
% 		%	especially for $K$ where cyclo is trivial,
% 		%	but I think it works out.}
% 		%over $\cW_s$, and thus of total dimension at most
%                 $$\dfrac{[K:\Q_p]\bigl(d(d-1) - ar(ar-1)\bigr)
% 			- s\bigl((a^2+1)s-a\bigr)}{2} -r^2 + rs .$$
% 		Using the fact that $\cX_{d,\red}$ is equidimensional of
% 		dimension $[K:\Q_p]d (d-1)/2,$ we find that the image
% 		of this family in $\cX_{d,\red}$ has codimension 
% 		at least 
% 		\begin{multline*}
% 		\dfrac{[K:\Q_p]ar(ar-1) + s\bigl((a^2+1)s-1\bigr)}{2} + r^2 - rs 
% 		\\
% 		\geq \dfrac{r\bigl((a^2+1)r - a\bigr) + (r-s)^2 + as(as-1)}{2} \geq
% 		\dfrac{r\bigl((a^2+1)r-a\bigr)}{2}.
%                 \end{multline*}
% 		Since this codimension is an integer it is in fact
% 		at least $\lceil r\bigl((a^2+1)r-a\bigr)/2 \rceil.$ Since this
%                 conclusion holds for each of the finitely many values
%                 of~$s$, we are done.
% 	% 	\mar{ME: The displayed equality is just a computation;
% 	% 		the $r^2$ comes from the $GL_r$ action on
% 	% 		$\chibar^{\oplus r}$ as automorphisms.  
% 	% 		The lower bound $c_r$ for the codimension
% 	% 		is then supposed to be worked out/chosen
% 	% 		so that the claimed inequality is true! ME: OK,
% 	% 	I made choice of $c_r$ and wrote down somehopefully
% 	% correct computations. TG: they seem correct to me.}
% 	\end{proof}


% 		%We now use Theorem~\ref{thm:dimension of H2} to
% 		%prove the rank $d+1$ case of Theorem~\ref{thm:reduced
% 		%	dimension}.
% 		%For this,  we first prove a more precise version
% 		%of the statement about
% 		%crystalline lifts given in Theorem~\ref{thm:intro crystalline
% 		%	lifts}.  We use the rank $d$ case 
% 		%of Theorem~\ref{thm:dimension of H2} as input.

% \begin{proof}[Proof of Theorem~\ref{thm:crystalline lifts} in the rank $d$ case]
% 	We write $R := R_{\rhobar}^{\crys,\underline{\lambda}}$, 
% 	and we let $X$ denote the component of the crystalline lifting scheme 
% 	%$$X := X_{\rhobar}^{\crys,\underline{\lambda}} =
% 	%\Spec R.$$
% 	$\Spec R_{\rhobar}^{\crys,\underline{\lambda}}$ on which $\rho^{\circ}_d$
% 	lies.
% 	\mar{ME: The following explanation needs to be smoothed out.  The
% 		main reason I want to work on $\Spec R$ rather than $\Spf R$
% 		is to not have to think about the geometric lemma over formal
% 		schemes. } Note that we can work over $\Spec R$ rather than just
% 		$\Spf R$ for the following reasons: firstly, since $R$ is
% 		a complete Noetherian local ring,
% 		coherent sheaves, perfect complexes, etc.\ on $\Spec R$
% 		and $\Spf R$ are the same. Secondly,
% 		since $\cX_d^{\crys,\underline{\lambda}}$
% 		is a $p$-adic formal algebraic stack, its mod $p$ fibre is algebraic,
% 		so that its versal rings will be effective; thus we will have
% 		a versal morphism
% 		$\Spec R/p \to \Spec \F_p\times_{\Spf \Z_p}\cX_d^{\crys,\underline{\lambda}},$ 
% 		and so the results on the support of $H^2$ over $\cX_{d,\red}$
% 		are inherited by $H^2$ over $\Spec R$.\mar{TG: I don't
%                   understand what's being used/assumed here to get
%                   that the statements about the codimension of the support of $H^2$ pull
%                   back to $\Spec R$, and in particular over~$X$. ME: Added ``versal'' to the discussion;
% 	  I guess smoothness/flatness is the key point.}

%         Fix the lift $\alpha^{\circ}$ of $\alpha$.
% 	Over $X$ we have the perfect complex $C^{\bullet}$ computing
% 	$\Ext^{\bullet}_{G_K}(\alpha^{\circ},\text{--})$ for the universal deformation.  Our assumption
% 	on the tuple of weights $\underline{\lambda}$ implies that $\Ext^2$ is
% 	supported (set-theoretically) on the special fibre $\overline{X}$
% 	of $X$.

% 	As usual, we let $B^2$ denote the image of $C^1$ in $C^2$; 
% 	it is a subsheaf of the locally free sheaf $C^2$, and so is 
% 	torsion free.\footnote{Since $X$ is reduced, {\em torsion free}
% 		is a reasonable notion for a finitely
% 		generated module, equivalent to the associated
% 		primes lying among the minimal primes.}
% 	We may find a blow up $\pi:\tX \to X$, whose centre lies
% 	(set-theoretically) in the special fibre of $X$,\mar{TG: do we
%         have a reference for this? I presume it relates to \cite[\href{http://stacks.math.columbia.edu/tag/080X}{Tag 080X}]{stacks-project}.}
% 	such that the torsion-free quotient of $\pi^*B^2$ becomes
% 	locally free (recall that~$\Ext^2$ is supported on the special
%         fibre of~$X$); in other words, if we let $\tC^{\bullet}$
% 	denote the pull-back by $\pi$ of $C^{\bullet}$, 
% 	the corresponding $\cO_{\tX}$-module of $2$-coboundaries
% 	$\tB^2$
% 	is locally free.
% 	Thus, if we let $\tZ^1$ denote the $\cO_{\tX}$-module of
% 	$1$-cocyles for $\tC^{\bullet}$, then $\tZ^1$ is also locally
% 	free.\mar{TG: we don't seem to explicitly use this fact.}

% 	The given extension $\rhobar_{d+a}$ is classified by a
% 	class in $\Ext^1_{G_K}(\alphabar,\rhobar_d)$, which is to say, a class
% 	in $H^1$ of the complex $\kappa\otimes C^{\bullet}$.
% 	(Here we write $\kappa$ to denote $\Fbar_p$ thought of as the
% 	residue field of $R$.)
%         Lifting this class to an element of $\kappa \otimes C^1,$
% 	and then to an element of $C^1$, we find ourselves in the 
% 	following situation:
% 	we have a morphism $c:R \to C^1$, whose image under the coboundary
% 	lies in $\mathfrak m_R C^2$ (reflecting the fact that we have a cochain
% 	which becomes a cocycle at the closed point).

% 	Consider the composite $b: R \buildrel c \over \longrightarrow
% 	C^1 \to B^2,$ 
% 	which lifts to a section $\tb: \cO_{\tX} \to \tB^2.$
% 	It follows from Lemma~\ref{lem:sections} below (applied with
%         $\cA=B^2$, $\cB=C^2$, and $\cC=\Ext^2$),
% 	together with Theorem~\ref{thm:dimension of H2},
% 	\mar{ME: Have to explain that Theorem~\ref{thm:dimension of H2} actually
% 		does control $H^2$, because of Tate local duality.}
% 	that the lifted section must have non-empty zero locus,
% 	which (being closed) must contain a point $\tx$ lying
% 	over the closed point of~$X$.   The section $c$ itself
% 	lifts to a section $\tc: \cO_{\tX} \to \tC^1$, whose value at the 
% 	point $\tx$ lies in the fibre of $\tZ^1$.   In other
% 	words, the fibre of $\tc$ at $\tx$ is a one-cocycle
% 	in the complex $\kappa(\tx)\otimes [\tC^0 \to \tZ^1]$, giving rise to
% 	a class
% 	$\overline{e} \in H^1\bigl(\kappa(\tx)\otimes [\tC^0 \to \tZ^1]\bigr)$
% 	lifting the original class in $\Ext^1_{G_K}(\alphabar,\rhobar_d)$ that classifies
% 	$\rhobar_{d+1}$.

% 	Since $\tX$ is flat over $\Z_p,$ we may find a morphism
% 	$\tif:\Spf \Zbar_p \to \tX$
% 	passing through the point $\tx \in \tX(\Fbar_p)$.	
% 	The morphism $\tif$ determines (and, by the valuative
% 	criterion of properness, is determined by) a morphism 
% 	$f: \Spf \Zbar_p \to X$,
% 	lifting the closed point $x \in X$, % \mar{TG: do you mean
%           % lifting the closed point of~$X$? ME: Yes, fixed.}
% 	which in turn corresponds to a representation
% 	$\theta_d^{\circ}: G_K \to \GL_d(\Zbar_p),$ as in the statement
% 	of the theorem.
% 	Now $H^2(\tC^{\bullet})$ is the cokernel of an inclusion of
% 	locally free sheaves (the inclusion $\tB^2 \hookrightarrow \tC^2$),
% 	and so Lemma~\ref{lem:torsion support} below shows that
% 	there is an effective Cartier divisor $D$ contained in the special
% 	fibre of $\tX$ with the property that for any morphism
% 	to $\tX$ that meets $D$ properly, the higher derived
% 	pull-backs of $H^2(\tC^{\bullet})$ under this morphism
% 	vanish.  Since the domain of $\tif$ is $p$-torsion free,
% 	it meets the special fibre of $\tX$ properly,
% 	and in particular meets $D$ properly; 
% 	%	in particular
% 	%is supported on a Cartier divisor on $\tX$,
%   %       \mar{TG: I don't
%   %         follow this - can you elaborate? ME: I added a lemma which hopefully
%   % makes things clearer.}
%         %which furthermore
% 	        %lies (set-theoretically) on the special fibre of $\tX$ (and thus is,
% 	        %set-theoretically, a union of irreducible components
% 	        %of this special fibre).
% 	thus we infer that
%        	%$$\Tor_i^R\bigl(\tif_*\cO_{\Spf \Zbar_p}, H^2(\tC^{\bullet})\bigr) = 0$$
% 	$$\mathbb L_i\tif^* H^2(\tC^{\bullet}) = 0$$
% 	for $i > 0.$ From this vanishing we deduce 
% 	that % \mar{TG: I don't follow this either. ME: I may
% % 	have had the wrong statement there; I've rewritten it for the
% % higher derived pull-backs by $\tif$, which is what I actually wanted.}
% 	\begin{multline*}
% 	H^1(G_K,\theta_d^{\circ}) = H^1(f^*C^{\bullet}) =
% 	H^1(\tif^*\tC^{\bullet}) = \tif^*H^1(\tC^{\bullet})
% 	\\
% 	= \tif^*H^1\bigl( [ \tC^0 \to \tZ^1]) = H^1\bigl(\tif^*[C^0 \to
% 	\tZ^1]\bigr)
% 	.\end{multline*}
% 	(The first of these identifications is the general base-change
% 	property of the perfect complex $C^{\bullet}$, the second
% 	follows from the fact that $f = \pi\circ \tif$,  the third
% 	follows from the vanishing of $L_i\tif^* H^2(\tC^{\bullet}) $
%         for $i>0$ and the base-change spectral sequence computing $H^i(f^*\tC^\bullet)$, % $\Tor$-vanishing that we just
%         % observed
%         % \mar{TG: I could also do with this third point being expanded.
% 	% ME: The third point uses the remarked-upon Tor vanishing together
%         % with the base-change spectral sequence for $H^i$ (which, if I 
%         % write it for rings rather than schemes,
%         % is the second quadrant sequence
%         % $E_2^{p,q} := \Tor_{-p}^A\bigl(B,H^q(G_K, \rho_A) \bigr)
%         % \implies H^{p+q} (G_K, \rho_B).$}
% 	the fourth holds by definition of $H^1$, and the fifth follows
% 	from right-exactness of $\tif^*$.)
% 	%Thus $H^1*(G_K,\theta_d^{\circ})$ \mar{TG: I think this
%         %  sentence fragment should just be removed? ME: I think so, and so
%         % I removed it.}

% 	Since $\tif$ identifies the closed point of $\Spf \Zbar_p$
% 	with the point $\tx,$ we may find a class
% 	$e \in H^1\bigl(\tif^*[\tC^0 \to \tZ^1]\bigr)$ lifting the
% 	class $\overline{e}$, which is then identified with an element
% 	of $\Ext^1_{G_K}(\alpha^{\circ},\theta_d^{\circ}).$   This element classifies
% 	an extension $0 \to \theta_d^{\circ} \to \theta_{d+a}^{\circ} 
% 	\to \alpha^{\circ} \to 0$ which satisfies the conclusions of the theorem.
% \end{proof}

% %Finally, we can prove the rank $d+1$ case of Theorem~\ref{thm:reduced dimension}.

% \begin{proof}[Proof of Theorem~\ref{thm:reduced dimension} in the rank $d+1$ case]
%        We know that $\cX_{d+1,\red}$ is the inductive limit
%        of reduced algebraic stacks that are locally of 
%        finite type over $\Spec \F_p$, 
%        with the transition morphisms being closed immersions.
%        For any tuple $\underline{\lambda}$ of labeled Hodge--Tate weights,
%        we know that $\cX^{\crys,\underline{\lambda}}_{d+1,\red}$ is 
%        an algebraic stack over $\F_p$, locally of finite type
%        \mar{ME: It may be that we already know that the various
%        algbebraic stacks appearing here are  actually loc.\ f.p.,
%        or even f.p.\,
%        over $\F_p$, but I think we only need loc.\ f.p.\ for the argument,
% 	       so that's what I'll write for now.}
%        and equidimensional of dimension
%        $[K:\Q_p](d+1)d/2$, which is also a closed substack of
%        $\cX_{d+1,\red}$.    

%        A closed immersion of reduced algebraic stacks that are locally
%        of finite type over $\F_p$ which is surjective on finite-type 
%        points is necessarily an isomorphism, and so we see that
%        it suffices to show that $\cX_{d+1,\red}(\Fbar_p)$ is contained
%        in the union of finitely many of its subsets
%        $\cX_{d+1,\red}^{\crys,\underline{\lambda}}(\Fbar_p).$  
%        Let $\rhobar: G_K \to \GL_{d+1}(\Fbar_p)$ be a continuous representation,
%        corresponding to a point of $\cX_{d+1,\red}(\Fbar_p)$.
%        If $\rhobar$ is irreducible, then we see directly that it
%        admits a crystalline lift whose labeled Hodge--Tate weights
%        \mar{ME: Maybe this argument should be tightened or rephrased,
% 	       now that we are discussing these sorts of issues about 
% 	       lifts much more carefully further below.}
%        are uniformly bounded (in terms of $d$).
%        On the other hand, if $\rhobar$ is reducible,
%        then this follows from Theorem~\ref{thm:crystalline lifts} 
%        % \mar{ME: Need the (currently unwritten) version of
%        %         that theorem in which $\triv$ is replaced by
%        %         an arbitrary irrep.}
%        (taking into account Remark~\ref{rem:crystalline lifts}), 
%        together with the present theorem in ranks $\leq d$.
% \end{proof}



\section{Potentially diagonalizable crystalline
  lifts}\label{subsec: precise crystalline lifts}We now use
Theorem~\ref{thm:crystalline lifts} to prove Theorem~\ref{thm:intro
  crystalline lifts}.  There are at least two differences between the
proof of Theorem~\ref{thm:intro crystalline lifts} and previous work
on the problem (in particular the results
of~\cite{MullerThesis,2015arXiv150601050G}.) One is that, rather than
working only with lifts which are extensions of inductions of characters,
we utilise extensions of more general potentially diagonalizable
representations. The other difference
is that our argument exploits our control
of the support of~$H^2$. In previous work, the only tool used to deal
with classes in~$H^2$ was twisting the various irreducible
representations by unramified characters; at points where~$H^2$ has
dimension greater than~$1$, this seems to be insufficient.


In fact, we prove various refinements of this result, allowing us to
produce potentially diagonalizable lifts; such lifts are important for
automorphy lifting theorems. We can also control the Hodge--Tate
weights of these potentially diagonalizable lifts; for example, we can
insist that the gaps between the weights are arbitrarily large, which
is useful in applications to automorphy lifting. Alternatively, we can
produce lifts whose Hodge--Tate weights lift some Serre weight,
proving a conjecture which is important for the formulation of general
Serre weight conjectures (see the discussion after~\cite[Rem.\
5.1.8]{2015arXiv150902527G}).

We begin by recalling some definitions and some basic lemmas about
extensions of crystalline representations.  Let~$\underline{\lambda}$ be a
regular $d$-tuple of Hodge--Tate weights.

% \mar{ME: Am I correct that everything that follows is in the service of
% 	proving Thm.~\ref{thm: strong existence of crystalline lifts}, 
% 	and that for this Thm.~\ref{thm:crystalline lifts} will be
% 	used as input exactly as it's stated, but without needing
% 	any further geometric input?  If so, then maybe this should
% 	be in its own section, just to separate it out a bit from
% 	the previous theorems, which, thematically, are more about
% 	determining the basic structural properties of $\cX_d$? TG:
%         hopefully done this, commenting out.}
% \mar{ME: Add preamble explaining what we're doing now.}




% \begin{rem}\label{rem:it's not always exts of inds, and
%     the multiple directions in H^2}

%\end{rem}


\begin{defn}\label{defn: ordinary} \index{ordinary}
  A crystalline representation of weight~$\underline{\lambda}$
  is~\emph{ordinary} if it has a $G_K$-invariant decreasing % \mar{TG:or
    % increasing? hopefully I have this the right way around}
  filtration
  whose associated graded pieces are all one-dimensional, such that
  the $\sigma$-labeled Hodge--Tate weight of the~$i$th graded piece
  is~$\lambda_{\sigma,i}$. In other words, we can write the
  representation in the form \[\begin{pmatrix}
      \chi_1 &*&\dots &*\\
      0& \chi_2&\dots &*\\
      \vdots&& \ddots &\vdots\\
      0&\dots&0& \chi_d\\
    \end{pmatrix} \]where~$\chi_i$ is a crystalline
  character whose $\sigma$-labeled Hodge--Tate weight
  is~$\lambda_{\sigma,d+1-i}$. 
\end{defn}

If~$\rho_1^\circ,\rho_2^\circ:G_K\to\GL_d(\Zpbar)$ are two crystalline
representations, then we write $\rho_1^\circ\sim\rho_2^\circ$, and say
that $\rho_1^\circ$ \emph{connects to} $\rho_2^\circ$, if and
only if the following conditions hold: $\rho_1^\circ$ and~$\rho_2^\circ$ have the same labeled
Hodge--Tate weights, have isomorphic reductions modulo~$\m_{\Zpbar}$, and
determine points on the same irreducible component of the
%\mar{ME: deformation rings $\to$ lifting  rings?}
corresponding crystalline lifting rings. The following definition
was originally made in~\cite[\S 1.4]{BLGGT}.
\begin{defn}\label{defn: pot diag}\index{potentially diagonalizable}
  A representation~$\rho^\circ:G_K\to\GL_d(\Zpbar)$ is
  \emph{potentially diagonalizable} if there is a finite
  extension~$K'/K$ and crystalline
  characters~$\chi_1,\dots,\chi_d:G_{K'}\to\overline{\Z}_p^\times$
  such that~$\rho^\circ|_{G_{K'}}\sim\chi_1\oplus\cdots\oplus\chi_d$.
\end{defn}

\begin{lem}
  \label{lem: extensions of PD are PD}Let
  $0\to\rho_1^\circ\to\rho^\circ\to\rho_2^\circ\to 0$ be an extension of
  $\Zpbar$-valued representations of~$G_K$. If $\rho^\circ$ is
  crystalline, and~$\rho_1^\circ$ and~$\rho_2^\circ$ are potentially
  diagonalizable, then~$\rho^\circ$ is also potentially
  diagonalizable.
\end{lem}
\begin{proof}We may choose $K'/K$ such that $\rho_1^\circ|_{G_{K'}}$
  and $\rho_2^\circ|_{G_{K'}}$ both connect to direct sums of crystalline
  characters, and such that $\rhobar|_{G_{K'}}$ is trivial. It then
  follows from points~(5) and~(7) of the list before~\cite[Lem.\
  1.4.1]{BLGGT} that~$\rho^\circ|_{G_{K'}}$ connects to the direct sum
  of the union of the sets of crystalline characters for $\rho_1^\circ|_{G_{K'}}$
  and $\rho_2^\circ|_{G_{K'}}$, as required.
\end{proof}
Note that it follows in particular from Lemma~\ref{lem: extensions of
  PD are PD} that ordinary representations are potentially diagonalizable.

% We now recall from~\cite{2015arXiv150902527G} the notion of a
% crystalline lift of some Serre weight. Each
% embedding~$\sigma:K\into\Qpbar$ induces an
% embedding~$\sigmabar:k\into\Fpbar$; if $K/\Qp$ is ramified, then
% each~$\sigmabar$ corresponds to multiple embeddings~$\sigma$.  We say
% that~$\underline{k}$ is a\mar{TG: here need to fix~$k$ versus $\lambda$}
% \emph{Serre weight} if for each embedding~$\sigmabar:k\into\Fpbar$, we
% can choose an embedding~$\sigma:K\into\Qpbar$ lifting~$\sigmabar$,
% with the properties that:
% \begin{itemize}
% \item $p\ge k_{\sigma,i}-k_{\sigma,i+1}\ge 1$ for each $1\le i\le
%   n-1$, and
% \item if $\sigma':K\into\Qpbar$ is any other lift of~$\sigmabar$,
%   then $k_{\sigma',i}=n-i$ for each $1\le i\le n$.
% \end{itemize}
% (We caution the reader that this terminology is not standard, but it
% is convenient for us here.)

\begin{thm}
  \label{thm: strong existence of crystalline lifts}Let $K/\Qp$ be a
  finite extension, and let~$\rhobar:G_K\to\GL_d(\Fpbar)$ be a
  continuous representation. Then~$\rhobar$ admits a lift to a
  crystalline representation~$\rho^\circ:G_K\to\GL_d(\Zpbar)$ of some
  regular labeled Hodge--Tate weights~$\underline{\lambda}$. Furthermore:
  \begin{enumerate}
  \item $\rho^\circ$ can be taken to be potentially diagonalizable.
  \item If every Jordan--H\"older factor of~$\rhobar$ is
    one-dimensional, then~$\rho^\circ$ can be taken to be
    ordinary.%\mar{TG: add defn}
  \item $\rho^\circ$ can be taken to be potentially diagonalizable,
    and~$\underline{\lambda}$ can be taken to be a lift of a Serre weight.
  \item $\rho^\circ$ can be taken to be potentially diagonalizable,
    and~$\underline{\lambda}$ can be taken to have arbitrarily spread-out
    Hodge--Tate weights: that is, for any $C>0$, we can
    choose~$\rho^\circ$ such that for each~$\sigma:K\into\Qpbar$, we
    have $\lambda_{\sigma,i}-\lambda_{\sigma,i+1}\ge C$ for each $1\le i\le n-1$.
  \end{enumerate}

\end{thm}
\begin{proof}We prove all of these results by induction on~$d$, using
  Theorem~\ref{thm:crystalline lifts}. We begin by proving~(1), and
  then explain how to refine the proof to give each of (2)-(4).

Suppose firstly that~$\rhobar$ is irreducible (this is the base case
of the induction). Then~$\rhobar$ is of
the form~$\Ind_{G_{K'}}^{G_K}\psibar$ for some
character~$\psibar:G_{K'}\to\Fpbartimes$, where $K'/K$ is unramified
of degree~$d$. We can choose (for example by~\cite[Lem.\
7.1.1]{2015arXiv150902527G}) a crystalline
lift~$\psi:G_{K'}\to\overline{\Z}_p^\times$ of~$\psibar$ such
that~$\rho^\circ:=\Ind_{G_{K'}}^{G_K}\psi$ has regular Hodge--Tate
weights. Then~$\rho^\circ$ is crystalline (because $K'/K$ is
unramified), and it is potentially diagonalizable,
because~$\rho^\circ|_{G_K'}$ is a direct sum of crystalline
characters.

For the inductive step, we may therefore suppose that we can
write~$\rhobar$ as an extension
$0\to\rhobar_1\to\rhobar\to\rhobar_2\to 0$, with~$\rhobar_2$
irreducible, and we may assume that we have regular potentially
diagonalizable lifts~$\rho_1^\circ$, $\rho_2^\circ$ of~$\rhobar_1$,
$\rhobar_2$ respectively. Twisting~$\rho_1^\circ$ by a crystalline
character with sufficiently large Hodge--Tate weights and trivial
reduction modulo~$p$ (which exists by~\cite[Lem.\
7.1.1]{2015arXiv150902527G}), we may suppose that the Hodge--Tate
weights of~$\rho_1^\circ$ are slightly less than those
of~$\rho_2^\circ$. By Theorem~\ref{thm:crystalline lifts}, there is a
lift~$\theta_1^{\circ}$ of~$\rhobar_1$
with~$\theta_1^\circ\sim\rho_1^\circ$, and a lift~$\rho^\circ$
of~$\rhobar$ which is an
extension \[0\to\theta_1^\circ\to\rho^\circ\to\rho_2^\circ\to 0.\]
Since~$\rho_1^\circ$ is potentially diagonalizable, so
is~$\theta_1^\circ$, and it follows from Lemmas~\ref{lem: extensions
  are crystalline} and~\ref{lem: extensions of PD are PD}
that~$\rho^\circ$ is crystalline and potentially diagonalizable. This
completes the proof of~(1).

We claim that if every Jordan--H\"older factor of~$\rhobar$ is
    one-dimensional, then the lift that we produced in proving~(1) is
    actually automatically ordinary. Examining the proof, we see that
    it is enough to show that if~$\rho_1^\circ$ is ordinary
    and~$\theta_1^\circ\sim\rho_1^\circ$ then $\theta_1^\circ$ is also
    ordinary; this is immediate from~\cite[Lem.\ 3.3.3(2)]{ger}.

    To show~(3) and~(4) we have to check that we can choose the
    Hodge--Tate weights of the lifts of the irreducible pieces
    of~$\rhobar$ appropriately. This requires substantially more
    effort in the case of~(3), but that effort has already been made
    in the proof of~\cite[Thm.\ B.1.1]{2015arXiv150902527G} (which
    proves the case that~$\rhobar$ is semisimple). Indeed, examining
    that proof, we see that it produces lifts satisfying all of the
    conditions we need; we need only note that the ``slightly less''
    condition is automatic, except in the case that for each~$\sigma$
    we have (in the notation of the proof of ~\cite[Thm.\
    B.1.1]{2015arXiv150902527G}) $h_\sigma+x_\sigma=H_\sigma+1$, in
    which case we can take $h_\sigma+x_\sigma=H_\sigma+p$ for
    each~$\sigma$ instead.

Finally, to prove~(4) it is enough to check the case that~$\rhobar$ is
irreducible (because we can choose the Hodge--Tate weights of the
character that we twisted~$\rho_1^\circ$ by to be arbitrarily large). For this, note that in our application of~\cite[Lem.\
7.1.1]{2015arXiv150902527G}, we can change any labeled Hodge--Tate by
any multiple of~$(p^d-1)$, so we can certainly arrange that the gaps
between labeled Hodge--Tate weights are as large as we please.
\end{proof}

\begin{rem}
  \label{rem: bound on weights}To complete the proof of
  Theorem~\ref{thm:intro crystalline lifts}, it is enough to note that
  by definition the lifts produced in Theorem~\ref{thm: strong
    existence of crystalline lifts}~(3) have all their Hodge--Tate
  weights in the
  interval~$[0,dp-1]$. % Indeed, by the definition of a Serre weight,
  % for each~$\sigma$ the $\sigma$-labeled Hodge--Tate weights are
  % contained in an interval of length at most~$(d-1)p$, so it is enough
  % to note that by the proof of~\cite[Prop.\
  % B.1.2]{2015arXiv150902527G}, we can ensure that each~$\lambda_{\sigma,d}$
  % is in the interval~$[0,p]$.\mar{TG: this is again unclear to me - my
  %   notes make a remark about twisting, and say ``just fix ``$[0,pd-1]$''\dots}

    % \mar{TG: I'm not entirely sure this is
    % true. I wonder if it would be better if we just made a statement
    % about the $\sigma$-labeled weights being contained in an interval
    % of length~$(d-1)p$? That would be more natural from the
    % perspective of discussing Serre weights, anyway.}

  % It is plausible that in the case that~$K/\Qp$ is unramified, this
  % result should be best possible, in the sense that there exist
  % representations~$\rhobar$ (given by successive tr\`es ramifi\'ee
  % extensions) for which there are no crystalline lifts with
  % $\sigma$-labeled Hodge--Tate weights contained in an interval of
  % length strictly less than~$(d-1)p$ for each~$\sigma$, but the
  % evidence is very limited beyond the case~$d=2$. In the case
  % that~$K/\Qp$ is ramified it is plausible that the required interval
  % should be shorter; it is possible that the Breuil--M\'ezard
  % conjecture predicts that this is the case, but we have not
  % investigated this.\mar{TG: we should probably check that we still
  %   agree with this remark.}
  % This is because we
% anticipate\mar{TG: probably there will be a discussion of this
%   elsewhere? Maybe we're trying to prove existence of labeling of
%   components, for example. ME: I agree; let's leave this comment for now
% to remind ourselves that this will have to be revisited at some point.} that the irreducible components
% of~$\cX_{d,\red}$ are labeled by Serre weights, in such a way that a
% generic point of a component correspond to a single Serre
% weight. Choosing the weight so that $k_{\sigma,i}-k_{\sigma,i+1}=p$
% for all~$1\le i\le d-1$ would give the claim.
% (As a concrete example,
% consider the case of $K=\Qp$, $d=2$, and~$\rhobar$ is a ramified
% extension of the trivial character by itself.)
\end{rem}
\begin{rem}
  \label{rem: get all lifts of a Serre weight}The proof of
  Theorem~\ref{thm: strong existence of crystalline lifts}~(3)
  actually proves the slightly stronger statement that there is a
  Serre weight such that if~$\underline{\lambda}$ is any lift of that
  Serre weight, then~$\rhobar$ admits a potentially diagonalizable
  crystalline lift of weight~$\underline{\lambda}$. (Note that in the
  case that~$K/\Qp$ is ramified, there may be many such choices of a
  lift of a given Serre weight, because by definition,
  the choice of a lift of $\underline{\lambda}$ depends on a choice of
  a lift of each embedding~$k\into\Fpbar$ to an embedding
  $K\into\Qpbar$. However, it is noted in the
  first paragraph of the proof of~\cite[Thm.\
  B.1.1]{2015arXiv150902527G} that in the case that~$\rhobar$ is
  semisimple, there is a crystalline lift of
  weight~$\underline{\lambda}$ for any choice of
  lift~$\underline{\lambda}$, so the same is true in our construction.)  % \mar{ME: The
    % expression ``the independence of the existence of such lifts'' is
    % a bit unclear.  I have a vague sense of what it means, but maybe
    % it can be clarified? TG: yeah that wasn't good, is this better?}
\end{rem}


%\section{$\cX_d$ is a formal algebraic stack}
%
%We now prove the following result.
%
%\begin{thm}\label{thm: Xd is formal algebraic}
%	$\cX_d$ is a formal algebraic stack.
%\end{thm}
%\begin{proof}
%	Theorem~\ref{thm:reduced dimension} shows that
%	$\cX_{d,\red}$ is an algebraic stack that 
%	is of finite presentation over $\F_p$ (and hence
%	quasi-compact and quasi-separated).  \mar{ME: Here we
%		are using that the $\cX_d^{\underline{k},\crys}$ 
%		are f.p., not just l.f.p. If we don't already know
%		this, we can probably deduce it from the fact
%		that $\cX_{d,\red}$ can be covered by a finite
%		union of explicit families of extensions; but since
%		I'm not quite sure yet of the logical order of these
%		various finiteness statements, I won't do anything more about
%		them for now than adding these marginal comments!}
%	Together with
%	Proposition~\ref{prop:X is an Ind-stack},
%	which shows that $\cX_d$ is isomorphic to the inductive limit
%	of a sequence of algebraic stacks with respect to transition
%	morphisms that are closed immersions,
%	\mar{ME: Make sure that that proposition is formulated
%		so that this summary  is valid.}
%	this implies that $\cX_d$
%      is indeed a formal algebraic stack,
%      by \cite[Cor.~2.5]{Emertonformalstacks}.
%      \end{proof}
%
%%\begin{lemma}
%%        Suppose that $\cX$ is an Ind-algebraic stack 
%%	that can be written as the inductive limit
%%	$\cX \cong \varinjlim \cY_n$ of an inductive sequence
%%	$(\cY_n)_{n \geq 1}$ in which the $\cY_n$ are
%%	algebraic stacks, and the transition morphisms
%%	are closed immersions.
%%	If $\cX_{\red}$ is a quasi-compact and quasi-separated
%%       	algebraic stack, then $\cX$ 
%%	is a formal algebraic stack.\mar{TG: should this be moved to
%%          the end of the introduction, or maybe even to \cite{Emertonformalstacks}? ME: I moved it there.}
%%\end{lemma}
%%\begin{proof}
%%	It suffices to show that
%%	the induced closed immersions $\cY_{n,\red} \to \cY_{n+1,\red}$
%%	are isomorphisms, for $n$ sufficiently large.
%%	By assumption,
%%	$\cX_{\red}$ is a quasi-compact and
%%	quasi-separated algebraic stack,
%%	and thus a consideration of the induced isomorphism
%%	$$\cX_{d,\red} \cong \varinjlim_n \cY_{n,\red}$$
%%	shows that indeed the inductive system $(\cY_{n,\red})_{n \geq 1}$
%%	stablizes when $n$ becomes sufficiently large.\mar{TG: I think
%%        this should be spelled out a bit more, e.g.\ some reference to
%%      see that qcqs is doing something. ME: Yes, there is a remark in 
%%      \S 4 of {\tt proper.tex} that discusses this that we can cite.}
%%\end{proof}
As a consequence of Theorem~\ref{thm: strong existence of crystalline
  lifts}, we can remove a hypothesis made in~\cite[App.\
A]{emertongeerefinedBM}, proving in particular the  following result
on the existence of globalizations of local Galois representations.
\begin{cor}
  \label{cor: existence of global lift}Suppose that~$p\nmid 2d$. Let $K/\Qp$ be a finite extension, and let
  $\rhobar:G_K\to\GL_n(\Fpbar)$ be a continuous representation. Then
  there is an imaginary CM field $F$ and a continuous irreducible
  representation $\rbar:G_{F}\to\GL_n(\Fpbar)$ such that
  \begin{itemize}
  \item each place $v|p$ of~$F$ satisfies $F_v\cong K$,
  \item for each place $v|p$ of $F$, either $\rbar|_{G_{F_v}}\cong
    \rhobar$ or $\rbar|_{G_{F_{v^c}}}\cong   \rhobar$, and
   \item $\rbar$ is automorphic, in the sense that it may be lifted
     to a representation $r:G_{F}\to\GL_n(\Qpbar)$ coming from a
     regular algebraic conjugate self dual cuspidal automorphic
     representation of~$\GL_n/F$.
  \end{itemize}
\end{cor}
\begin{proof}
  This is immediate from~\cite[Cor.\ A.7]{emertongeerefinedBM}
  (since~\cite[Conj.\ A.3]{emertongeerefinedBM} is a special case of
  Theorem~\ref{thm: strong existence of crystalline lifts}).
\end{proof}
\section{The irreducible components of $\cX_{\lowercase{d,\red}}$}\label{subsec:
  irred components}
%\mar{ME: I forget, we do write $\cX_{d,\red}$ or $(\cX_d)_{\red}$?  We have both
%in the paper right now. TG: pretty sure that the former is preferred,
%but the latter is not an outrage. In particular in commit 3309 I
%changed all instances of the latter that I could find to the former --
%where are the ones I've missed? ME: I thought I saw some, but couldn't find them
%again, so will just remove this comment}
We now complete the analysis of the irreducible
components of~$\cX_{d,\red}$ that we began in Section~\ref{subsec: formal
  algebraic stack}. Recall that Theorem~\ref{thm:Xdred is algebraic}
shows that~$\cX_{d,\red}$ is an algebraic stack of finite presentation
over~$\F$, and has dimension~$[K:\Qp]d(d-1)/2$. Furthermore, for each
Serre weight~$\underline{k}$, there is a corresponding irreducible
component~$\cX_{d,\red,\Fpbar}^{\underline{k}}$
of~$(\cX_{d,\red})_{\Fpbar}$, and the components for
different weights~$\underline{k}$ are distinct (see Remark~\ref{rem: different k give different components}).
%\mar{ME: Show that they are naturally labeled by Serre weights.}

We are now finally in a position to prove the following
result.% \mar{TG: assuming maximally nonsplit works out for stacks, add
  % a sentence here saying that we are maximally nonsplit}
\begin{thm}
	\label{thm:reduced dimension}  $\cX_{d,\red}$ is
        equidimensional of dimension
	$[K:\Q_p] d(d-1)/2$, and the irreducible components of~$(\cX_{d,\red})_\Fpbar$ are
        precisely the various closed substacks~$\cX_{d,\red,\Fpbar}^{\underline{k}}$
	of~Theorem~{\em \ref{thm:Xdred is algebraic}}; in particular,
        $(\cX_{d,\red})_\Fpbar$ is maximally nonsplit of niveau~$1$.  Furthermore
        each~$\cX_{d,\red,\Fpbar}^{\underline{k}}$ can be defined
        over~$\F$, i.e.\ is the base change
        of an irreducible component~$\cX_{d,\red}^{\underline{k}}$ of~$\cX_{d,\red}$.
	% There is a finite set of tuples $\underline{\lambda}$ of regular labeled Hodge--Tate
	% weights such that $\cX_{d,\red}$ is the union of corresponding
	% underlying reduced substacks $\cX_d^{\crys,\underline{\lambda}}.$
	% In particular,
	% each irreducible component of $\cX_{d,\red}$ is of dimension
	% $[K:\Q_p] d(d-1)/2$.
\end{thm}
\begin{proof}%[Proof of Theorem~\ref{thm:reduced dimension} in the rank $d+1$ case]
     By Theorem~\ref{thm:Xdred is algebraic}, it is enough to prove
     that each irreducible component of 
     $(\cX_{d,\red})_\Fpbar$ is of dimension of at least
     $[K:\Q_p] d(d-1)/2$. Indeed, this shows that the irreducible
     components of~$(\cX_{d,\red})_\Fpbar$ are
        precisely the~$\cX_{d,\red,\Fpbar}^{\underline{k}}$; to see
        that these may all be defined over~$\F$, we need to show that
        the action of~$\Gal(\Fpbar/\F)$ on the irreducible components
        of~$(\cX_{d,\red})_\Fpbar$ is trivial. This follows
        immediately by considering its action on the maximally
        nonsplit representations of niveau~$1$ (since the action
        of~$\Gal(\Fpbar/\F)$ preserves the property of being maximally
        nonsplit of niveau~$1$ and weight~$\underline{k}$).

In order to see that each  irreducible component of 
     $(\cX_{d,\red})_\Fpbar$ is of dimension of at least
     $[K:\Q_p] d(d-1)/2$, it suffices to show that every finite type
     point~$x$     of $(\cX_{d,\red})_\Fpbar$ is contained in an irreducible
     substack of $(\cX_{d,\red})_\Fpbar$ of dimension at least
     $[K:\Q_p] d(d-1)/2$; so it suffices in turn to show that each~$x$
     is contained in an equidimensional substack of dimension
     $[K:\Q_p] d(d-1)/2$. But by Theorem~\ref{thm: strong existence of
       crystalline lifts}, there is a regular Hodge
     type~$\underline{\lambda}$ such that $x$ is contained in
     $\cX_d^{\crys,\underline{\lambda}}\times_{\Spf \cO}\Spec\F$
     (for~$\cO$ sufficiently large), and this stack is equidimensional
     of dimension $[K:\Q_p] d(d-1)/2$ by Theorem~\ref{thm: dimension
       of ss stack}, as required.\end{proof}
% \mar{TG: one more $\Fpbar$ issue here, have to decide whether~$\cU$ is
% over~$\F$ or~$\Fpbar$. TG: with this and the one above, just base
% change deformation rings (by completed tensor product)
% to~$W(\Fpbar)$. TG: is this OK?!}
%         \mar{TG: it looks like we could now make this shorter/clearer
%           by just using the special fibre of
%           some~$\cX^{\crys,\lambda}_d$, rather than passing to
%           versal rings.}
%      Given an irreducible component of~ $(\cX_{d,\red})_\Fpbar$, choose a non-empty
%      open substack $\cU$ 
%      of $\cX_{d}\times_{W(\F')}W(\Fpbar)$ such that $\cU_{\red}$ is an open substack of
%      this component, disjoint from all the others (so in particular
%      $\cU_{\red}$ is irreducible).
%      It suffices (e.g.\ by \cite[\href{https://stacks.math.columbia.edu/tag/0DRX}{Tag 0DRX}]{stacks-project})
%      to show that $\dim \, \cU_{\red} \geq 
%      [K:\Q_p] d(d-1)/2$.
%      %\mar{ME: Citation for this.}

%      Let~$x$ be a finite type point of~$\cU_{\red}$, corresponding
%      %\mar{TG: rewrite as finite type point with some $\F$ and $\cO$?}
%      to a Galois representation $\rhobar: G_K \to \GL_d(\F')$ for some
%      finite extension~$\F'/\F$.
%      By Theorem~\ref{thm: strong existence of crystalline lifts},
%      after possibly enlarging~$\F'$, we may find a finite
%      extension~$E'/\Qp$ with ring of integers~$\cO'$ and residue field~$\F'$, and a
%      %and Theorem~% \ref{thm: crystalline deformation rings are effective versal}
%       % \mar{TG: and cite something
%        % for versal ring}
%      tuple of regular labeled Hodge--Tate
%      weights~$\underline{\lambda}$
%      %there is a versal morphism ~$\Spf
%      %R_{\rhobar}^{\overline{\lambda},\crys}\to\cX_{d,\red}$ at~$x$.
% such that     the corresponding crystalline lifting ring
%      $R_{\rhobar}^{\overline{\lambda},\crys,\cO'}$ 
%      is non-zero.
%      Furthermore, the natural morphism %\mar{TG: $\cU$ versus $\cX$ in
%        %this paragraph. ME: Dealt with}
%      $\Spec R_{\rhobar}^{\overline{\lambda},\crys,\cO'}\cotimes_{W(\F')}W(\Fpbar) /\varpi
%      \to \cU\times_{W(\Fpbar)}\Fpbar$ is effective, %\mar{TG: should be mod $\varpi$?}
%      by Corollary~\ref{cor: crystalline deformation rings are effective versal}.
%      If we let $R_{\rhobar}^{\square,\cO'}$ denote the corresponding universal lifting ring of~$\rhobar$,
%      then
%      %by the discussion
%      %of Subsection~\ref{subsec:Galois reps},
%      %%following Proposition~\ref{prop: Caruso
%      %%  Liu canonical action semistable})\mar{TG: this discussion will
%      %%  probably move. ME: Dealt with this now}
%      %the canonical
%      Proposition~\ref{prop:versal rings} yields a
%      morphism $\Spf R_{\rhobar}^{\square,\cO'} \to \cX_d$
%      that is versal to~$x$.
%      This morphism factors through the open substack $\cU$ (since $x \in
%      \cU$), and is of course (after base changing to~$W(\Fpbar)$) also versal to $\cU$ at~$x$.
%      % \mar{ME: Citation
% 	     % for this}
%    Writing $R_{\rhobar}^{\square,\cO',\Fpbar}$ for $R_{\rhobar}^{\square,\cO'}\cotimes_{W(\F')}W(\Fpbar)$, Proposition~\ref{prop:versal rings} provides an isomorphism
%     \begin{multline*}
%      \Spf R_{\rhobar}^{\square,\cO',\Fpbar} \times_{\widehat{\cU}_x}
% \Spf R_{\rhobar}^{\square,\cO',\Fpbar}
% \\
%      \iso
%      \Spf R_{\rhobar}^{\square,\cO',\Fpbar}
% \times_{\widehat{\bigr(\cX_d\times_{W(\F')}W(\Fpbar)\bigr)}_x} \Spf R_{\rhobar}^{\square,\cO',\Fpbar}
% \\
%      \iso \widehat{\GL}_{d,R_{\rhobar}^{\square,\cO',\Fpbar},1}.
% \end{multline*}
% %\mar{TG:  $S$ is maybe $\CO$?}
%      %this morphism is a~$\widehat{\GL}_d$-torsor, and it
%      %follows from~\cite[Lem.\ 2.40]{EGcomponents}\mar{TG: we should
%      %  probably cite some other things too, e.g.\ \cite[Lem.\
%      %  2.25]{EGcomponents}, to show that it's enough to compute the
%      %  dimension at closed points.
%      %ME: Will rewrite proof a little
% %to inclued a reference to Lemma~\ref{lem:bounding dimensions} below.}
% It follows from Lemma~\ref{lem:bounding dimensions} below
% that~$\dim \, \cU_{\red} \geq [K:\Q_p]
%      d(d-1)/2$, as required.
%     % Since this holds for all $\Fpbar$-points~$x$, we are done.  %  an
%   %    algebraic stack of finite presentation over~$\Fp$, and its
%   %    irreducible components all have dimension at most
%   %    $[K:\Qp]d(d-1)/2$.
%   % % We know that $\cX_{d+1,\red}$ is the inductive limit
%   % %      of reduced algebraic stacks that are locally of 
%   % %      finite type over $\Spec \F_p$, 
%   % %      with the transition morphisms being closed immersions.
%   %      For any tuple $\underline{\lambda}$ of labeled Hodge--Tate weights,
%   %      we know that $\cX^{\crys,\underline{\lambda}}_{d,\red}$ is 
%   %      an algebraic stack over $\F_p$, locally of finite type
%   %      \mar{ME: It may be that we already know that the various
%   %      algbebraic stacks appearing here are  actually loc.\ f.p.,
%   %      or even f.p.\,
%   %      over $\F_p$, but I think we only need loc.\ f.p.\ for the argument,
%   %              so that's what I'll write for now.}
%   %      and equidimensional of dimension
%   %      $[K:\Q_p]d(d-1)/2$, which is also a closed substack of
%   %      $\cX_{d,\red}$.    
% %
%   %      A closed immersion of reduced algebraic stacks that are locally
%   %      of finite type over $\F_p$ which is surjective on finite-type 
%   %      points is necessarily an isomorphism, and so we see that
%   %      it suffices to show that $\cX_{d,\red}(\Fbar_p)$ is contained
%   %      in the union of finitely many of its subsets
%   %      $\cX_{d,\red}^{\crys,\underline{\lambda}}(\Fbar_p).$   It is
%   %      therefore enough to show that if $\rhobar: G_K \to
%   %      \GL_{d}(\Fbar_p)$ is a continuous representation, then it
%   %      admits a crystalline lift whose labeled Hodge--Tate weights
%   %      are uniformly bounded (in terms of $d$). This follows easily by
%   %      induction on~$d$ from Theorem~\ref{thm:crystalline lifts}, and
%   %      in particular is immediate from Theorem~\ref{thm: strong
%   %        existence of crystalline lifts}~(3) below.
% \end{proof}

%\mar{ME: The following lemma seems related to the argument in
%  the proof of
%  Theorem~\ref{thm: codimension H2 in regular weight crystalline deformation ring}.  Maybe I should investigate this! ME: Don't quite understand
%all the hypotheses in this lemma.  I wonder if it's formulation changed
%over time? Actually, the Stacks Project just seems to have a more
%general statement than our write-up; have to figure out how/why! Also,
%relative dimension doesn't make sense --- strictly speaking --- in the
%formal scheme world. TG: I think we can just scrap most of the
%hypotheses, Johan improved a lot of things at the end of our
%dimensions paper when he put them into the stacks project. I guess
%hopefully we can reformulate the relative dimensions in terms of a
%presentation for the relation, as that's what we actually use (this
%lemma is only used once, immediately above)? TG: this is the same
%issue that came up in CEGS, so we should definitely stick with these hypotheses.}

% In order to state our next lemma,
% we first introduce some notation.
% Namely, we let $S$ be a Jacobson, universally
% 	catenary, locally Noetherian scheme, all of whose 
% 	local rings %at finite type points
% 	are $G$-rings,
% 	and with the further property that each %locally closed
% 	irreducible component %subset
% 	of $S$ is of finite dimension
% 	and equicodimensional.
% (This somewhat technical list of assumptions is what is required
% for in order to apply certain of the dimension theory results 
% of~\cite{EGcomponents}.)
% We let $\cX$ be a formal algebraic stack over~$S$,
% 	whose underlying reduced algebraic stack $\cX_{\red}$ is locally 
% 	of finite type over $S$, and irreducible. 
% We suppose that $x: \Spec k \to \cX$ is a morphism from a finite type
% $\cO_S$-field to $\cX$ (or, equivalently, to $\cX_{\red}$).
% Finally, we suppose that $A$ is an object of $\cC_{\Lambda}$ 
% (see the discussion of versality in Appendix~\ref{app: formal algebraic stacks}
% for an explanation of this notation)
% endowed with a morphism $\Spf A \to \cX$ that realizes it as a versal ring
% to $\cX$ at~$x$.
% By definition, this means that the induced morphism $\Spf A 
% \to \widehat{\cX}_x$ is smooth in the sense 
% of~\cite[\href{https://stacks.math.columbia.edu/tag/06HG}{Tag 06HG}]{stacks-project}.
% Thus, if we form the fibre product $R := \Spf A \times_{\widehat{\cX}_x} \Spf A$,
% then $R$ is a groupoid sheaf over $\Spf A$,
% and by~\cite[\href{https://stacks.math.columbia.edu/tag/06L1}{Tag 06L1}]{stacks-project} there is an isomorphism $[\Spf A/R] \iso \widehat{\cX}_x$.  % \mar{ME:
% % Give citation to SP for this.}


% \begin{lemma}
% 	\label{lem:bounding dimensions} 
% %	Let $S$ be a Jacobson, universally
% %	catenary, locally Noetherian scheme, all of whose 
% %	local rings %at finite type points
% %	are $G$-rings,
% %	and with the further property that each %locally closed
% %	irreducible component %subset
% %	of $S$ is of finite dimension
% %	and equicodimensional.
% %	Let $\cX$ be a %locally countably indexed and locally Ind-finite type
% %	formal algebraic stack over~$S$,
% %	whose underlying reduced algebraic stack $\cX_{\red}$ is locally 
% %	of finite type over $S$, and irreducible. 
% %	% topological space $|\cX|$ is irreducible.
% %	Suppose that $x: \Spec k \to \cX_{\red}$ is a morphism
% %	whose source is a field which is of finite type as an $\cO_S$-algebra,
% %	that $A$ is a complete Noetherian local ring %pro-Artinian topological ring
% %	equipped with a morphism
% %	$\Spf A \to \cX$ which realizes $A$ as a versal
% %	ring to $\cX$ at~$x$,
% %and that $B$ is a quotient of $A$
% %	%by a closed ideal
% %	for which the induced morphism $\Spf B \to \cX$ 
% %	is effective.
% %	Suppose furthermore that $\Spec B$
% %	contains an irreducible component of dimension $\geq d$, 
% %	and that each of the two formally smooth projections 
% %	$\Spf A \times_{\cX} \Spf A
% %	\to \Spf A$ is of relative dimension $e$.
% %	Then $\dim_x(\cX_{\red}) \geq d-e$.
% In the preceding situation,
% suppose that $B$ is a quotient of $A$ for which
% the induced morphism $\Spf B \to \cX$
% is effective,
% and such that $\Spec B$ contains an irreducible component 
% of dimension $\geq d$.
% Suppose further that $R$ is Noetherianly pro-representable \emph{(}in the
% sense of~\cite[Defn.\ 2.2.19]{EGstacktheoreticimages};
% e.g.\ if the diagonal of $\cX$ is locally of finite type,
% this holds by \cite[Lem.~2.7.2]{EGstacktheoreticimages}{\em )}, 
% %\mar{ME: Check whether or not Noetherian pro-representability is
% %  automatic TG: there is 2.7.2 in \cite{EGstacktheoreticimages}, and
% %  we are assuming that~$\cX_{\red}$ is locally of finite type, but I
% %  don't see that it implies that~$\cX$ has diagonal which is locally
% %  of finite type? ME: Incorporated this into statement}
% say $R = \Spf C$, and that $\dim C - \dim A = e$.
% Then $\dim_x(\cX_{\red}) \geq d-e$.
% \end{lemma}
% \begin{proof}
% 	Base-changing the closed immersion $\cX_{\red} \hookrightarrow
% 	\cX$ over the morphism $\Spf A \to~\cX$,
% 	we obtain a closed immersion
% 	$\cX_{\red} \times_{\cX} \Spf A \hookrightarrow \Spf A.$
% 	Since $\Spf A$ is locally countably indexed
% 	we see %\mar{ME: Give citations to [SP]}
% 	that $\cX_{\red} \times_{\cX} \Spf A = \Spf D$
% 	for some quotient $D$ of $A$ (by Lemma~\ref{lem: alem closed formal algebraic scheme
%     adic*}).
% 	%by a closed ideal.
% 	The construction of $\Spf D$ as a fibre product shows that
% 	the morphism $\Spf D \to \cX_{\red}$ realizes $D$
% 	as a versal ring to $\cX_{\red}$ at $x$.
% 	%Since $\cX_{\red}$ is an algebraic stack,
% 	%this morphism is effective.

% 	Since $\Spec B_{\red} \to \Spec B$ is a closed immersion which
% 	induces an isomorphism on underlying topological spaces,
% 	it it no loss of generality to replace $B$ by $B_{\red}$,
% 	and therefore to suppose that $B$ is reduced.   The assumption
% 	that the morphism $\Spf B \to \cX$ is effective means (by definition)
% 	that this morphism is induced by a morphism $\Spec B \to \cX$.
% 	Since $B$ is reduced, this morphism factors as a morphism
% 	$\Spec B \to \cX_{\red}$. 
%         The surjection $A\to B$ thus factors through $D$,
% 	and so $\Spec D$ contains an irreducible component of dimension
% 	$\geq d$.
% 	%It is no loss of generality, then, 
% 	%to replace $B$ by $D$ and thus to assume that 
% 	%$\Spec B \iso \cX_{\red} \times_{\cX} \Spf A,$ 
% 	%and that $B$ is versal to $\cX_{\red}$ at $x$.
  
%         Since each projection
% $\Spf D \times_{(\widehat{\cX_{\red}})_x} \Spf D \to
% 	\Spf D$ is obtained
% 	as the base-change via $\Spf D \to \Spf A$ 
% 	of the corresponding projection $\Spf A \times_{\widehat{\cX}_x} \Spf A 
% 	\to \Spf A,$ we find that these projections are formally smooth,
% 	of relative dimension $e$.  (More precisely,
% $\Spf D \times_{(\widehat{\cX_{\red}})_x} \Spf D = \Spf C\otimes_A D,$
% and $ \dim C\otimes_A D - \dim D = \dim C - \dim A = e$.)
%         Applying~\cite[Lem.\ 2.40]{EGcomponents}, the lemma follows.
% 	(In the notation of that result, we have $A$ equal to our ring $D,$
% 	so $\dim A = \dim D \geq d,$ while $B$ in that result is equal
%         to our present $C\otimes_A D$, so that we have
% 	$\dim B = \dim A + e.$)
% %	Since $\cX_{\red} \hookrightarrow \to \cX$ is a closed immersion
% %	(and so in particular a monomorphism), we find that
% %	$\Spec D \times_{\cX_{\red}} \Spec D \iso \Spec D \times_{\cX} \Spec D.$
% \end{proof}
We end this section with the following result, showing that in
contrast to its substacks~$\cX_{d}^{\crys,\lambdau,\tau}$
and~$\cX_d^{\semis,\lambdau,\tau}$, the formal algebraic stack~$\cX_d$ is
not a $p$-adic formal algebraic stack.

\begin{prop}\label{prop: X is not padic formal algebraic}
   $\cX_d$ is not a $p$-adic formal algebraic stack.
\end{prop}
\begin{proof}
  Assume that~$\cX_d$ is a $p$-adic formal algebraic stack, so that
  its special fibre~$\overline{\cX}_d:=\cX_d\times_\cO\F$ is an
  algebraic stack, which is furthermore of finite type over~$\F$ (since~$\cX_d$ is a
  Noetherian formal algebraic stack, by Corollary~\ref{cor:Xd is
    formal algebraic}, and~$\cX_{d,\red}$ is of finite type
  over~$\F$, by Theorem~\ref{thm:Xdred is algebraic}). Since the underlying reduced substack
  of~$\overline{\cX}_d$ is~$\cX_{d,\red}$, which is equidimensional of
  dimension~$[K:\Qp]d(d-1)/2$, we see that~$\overline{\cX}_d$ also has
  dimension~$[K:\Qp]d(d-1)/2$.

  Consider a finite type point~$x:\Spec\F'\to\cX_{d,\red}$,
  corresponding to a representation $\rhobar:G_K\to\GL_d(\F')$.
  By Proposition~\ref{prop:versal rings} there
  is a corresponding versal morphism
  $\Spf R_{\rhobar}^\square/\varpi\to\overline{\cX}_d$, and
  applying~\cite[Lem.\ 2.40]{EGcomponents} as in the proof of Theorem~\ref{thm:
    dimension of ss stack},
% \mar{ME:  Check this, b/c Thm.~4.8.14 is where our fibre product blunder
% has to be fixed. TG: I think this should be fine, especially as
% Proposition~\ref{prop:versal rings} has the description at the
% completion, which is what we would apply. TG: I think it's fine now
% and propose commenting this out. Note we have $\overline{\cX}$ here already!}
we conclude
  that~$R_{\rhobar}^\square/\varpi$ must have
  dimension~$d^2+[K:\Qp]d(d-1)/2$. However, it is known that there are
  representations~$\rhobar$ for which~$R_{\rhobar}^\square/\varpi$ is
  formally smooth of dimension $d^2+[K:\Qp]d^2$ (see for
  example~\cite[Lem.\ 3.3.1]{Allen2019}; indeed, a Galois cohomology
  calculation shows that the dimension is always at
  least~$d^2+[K:\Qp]d^2$, which is all that we need). Thus we must
  have~$d^2=d(d-1)/2$, a contradiction. % \mar{TG: in fact the dimension
    % is always at least this?}
  %\mar{TG: should double check
 %   that I have these dimensions correct; I guess the only possible
 %   issue would be~$d=1$, which will be explicitly checked in the
 %   following section anyway. ME: Looks fine to me}
\end{proof}

\begin{remark}
\label{rem:non-effective}
A similar argument shows that the versal morphisms $\Spf R_{\rhobar}^{\square,\cO'}
\to \cX_d$ and $\Spf R_{\rhobar}^{\square,\cO'}/\varpi \to \cX_d \otimes_{\cO} \F$
of Proposition~\ref{prop:versal rings} are {\em not} effective.
%\mar{ME: Should we explain in more detail?  At least  we should convince ourselves that a similar argument proving this
%does exist! TG: I guess the idea is that if the former is effective,
%so is the latter, and then we would also get an effective versal (?
%didn't check, hopefully true) morphism
%$\Spec (R_{\rhobar}^{\square,\cO'}/\varpi)_{\red} \to (\cX_d
%\otimes_{\cO} \F)_{\red}$; and now we're in good shape, since the
%target is now an algebraic stack, and then we would get a
%contradiction? Is that right? ME: That looks good to me!}
\end{remark}

% \section{The relationship with Galois
%   representations}\label{subsec: phi Gamma to Galois}\mar{TG: should
%   probably merge with the following section}

% The irreducible representations $\alphabar: G_K \to \GL_d(\Fbar_p)$
% are easily classified; indeed, since the wild inertia subgroup must
% act trivially, they are tamely ramified, so that the restriction
% of~$\alphabar$ to~$I_K$ is diagonalisable. Considering the action of
% Frobenius on tame inertia, it is then easy to see that~$\alphabar$ is
% induced from a character of the unramified extension of~$K$ of
% degree~$d$. \mar{TG: move these two paragraphs}

% By local class field theory, there are only finitely many characters
% up to unramified twist, so that %\mar{ME: Recall this}
% in particular there are only finitely many such~$\alphabar$ up to
% unramified twist. Via the equivalence explained above, twisting
% $G_K$-representations by unramified characters corresponds to
% multiplying the action of~$\varphi$ on $(\varphi,\Gamma)$-modules by
% scalars, and accordingly we also refer to this operation as unramified
% twisting. It follows that there are up to unramified twist only
% finitely many irreducible $(\varphi,\Gamma)$-modules over~$\Fpbar$ of
% any fixed rank.
% % we may choose a finite set $\{\alphabar\}$ of such 
% % irreducible representations of dimension $d$ such that any
% % irreducible $d$-dimensional representation may be obtained
% % as an unramified twist of precisely one of the finitely many fixed
% % $\alphabar$.

\section{Closed points}
\label{subsec:closed points}
Recall that if $\cY$ is  an algebraic stack, then $|\cY|$ denotes its underlying topological
space.  The points of $|\cY|$, which we also refer to as the points of~$\cY$,
 are the equivalence classes of morphisms $\Spec k \to \cY$,
with $k$ a field; two morphisms $\Spec k \to \cY$ and $\Spec l \to \cY$ are deemed to be
equivalent if we may find a field $\Omega$ and
embeddings $k, l \hookrightarrow \Omega$ such that
the pull-backs to $\Spec \Omega$ of the two given morphisms coincide.
A point of $\cY$ is called {\em finite type} if it is representable by
\index{finite type point}
a morphism $\Spec k \to \cY$ which is locally of finite type.   If $\cY$ is locally of
finite type over a base-scheme~$S$, then a point of $\cY$ is of finite type
if and only if its image in $S$ is a finite type point of $S$ in the usual
sense~\cite[\href{https://stacks.math.columbia.edu/tag/01T9}{Tag 01T9}]{stacks-project}.

We say that a point of $\cY$ is {\em closed} if the corresponding point of $|\cY|$ is closed
in the topology on~$|\cY|$.  Closed points are necessarily of finite type.
However, if $\cX$ is not quasi-DM, then finite type 
points of $\cX$ need not be closed, even if $\cX$ is Jacobson (in
contrast to the situation for
schemes~\cite[\href{https://stacks.math.columbia.edu/tag/01TB}{Tag 01TB}]{stacks-project}).

If $\cY$ is locally of finite type over an algebraically closed field~$k$,
then the map $\cY(k) \to |\cY|$ (taking a $k$-valued point of $\cY$ to the
point it represents) identifies $\cY(k)$ with the set of finite type points 
of~$\cY$.  (This is a standard consequence of the Nullstellensatz, applied
to a scheme chart of~$\cY$.)

Throughout this section,
we fix a value of $d \geq 1$, and %as usual,
write $\cX$ rather than
$\cX_d$. In order to describe the finite type and closed points
of~$(\cX_{\red})_{\Fbar_p}$, we introduce the following terminology.

\begin{defn}\label{defn: partial ss}
  Let $\rhobar, \thetabar:G_K\to\GL_d(\Fpbar)$ be continuous
  \index{partial semi-simplification}
  representations. Then we say that~$\thetabar$ is a \emph{partial
    semi-simplification} of~$\rhobar$ if  there are short exact
  sequences \[0\to \rbar_i\to\rhobar_i\to\rbar'_i\to 0\]
  for~$i=1,\dots$ such that $\rhobar_1 = \rhobar,$
  $\rhobar_{i+1}=\rbar_i\oplus\rbar_i'$,
  %$\rhobar_1=\rhobar$,
  and $\rhobar_n=\thetabar$ for~$n$ sufficiently large.
\end{defn}
Note in particular that the semi-simplification of~$\rhobar$ is a
partial semi-simplification of ~$\rhobar$. We also consider the
following related notion.

\begin{defn}\label{defn: virtual partial ss}Let $\rhobar, \thetabar:G_K\to\GL_d(\Fpbar)$ be continuous
  representations. Then we say that~$\thetabar$ is a \emph{virtual partial
    semi-simplification} \index{virtual partial
    semi-simplification} of~$\rhobar$ if there is a short exact
  sequence of the form \[0\to\rbar\to\rbar\oplus
    \rhobar\to\thetabar\to 0.\]  
\end{defn}
It follows from an evident induction on the integer~$n$ in Definition~\ref{defn: partial ss} that if ~$\thetabar$ is a partial
semi-simplification of~$\rhobar$, then it is in particular a virtual
partial semisimplification; however, the converse does not hold in
general. (See the introduction to~\cite{MR1757882}.)
%Recall also that $\F$ denotes the residue field of our chosen coefficient ring~$\cO$.

Our promised description of the finite type and closed points
of~$(\cX_{\red})_{\Fbar_p}$ is given by the following theorem.
\begin{theorem}
\label{thm:closed points}
% \mar{ME: I think there's a command that fixes the spacing in this context, 
% but I forget what it is!}
\leavevmode
\begin{enumerate}
\item
The morphism $\cX(\Fbar_p) = \cX_{\red}(\Fbar_p) \to |(\cX_{\red})_{\Fbar_p}|$ that sends each
morphism $\Spec \Fbar_p \to \cX$ to the point
of $(\cX_{\red})_{\Fbar_p}$  that it represents is injective, and its image
consists precisely of the finite type points.  Thus the finite type points
of $(\cX_{\red})_{\Fbar_p}$ are in natural bijection with the isomorphism
classes of continuous representations $\rhobar: G_K \to \GL_d(\Fbar_p)$.
\item A finite type point
of $|(\cX_{\red})_{\Fbar_p}|$
is closed if and only if the associated Galois representation $\rhobar$ is semi-simple.
\item If $x \in |(\cX_{\red})_{\Fbar_p}|$ is a finite type point, corresponding
to the Galois representation~$\rhobar$, then the closure $\overline{\{x\}}$
contains a unique closed point, whose corresponding Galois
representation is the semi-simplification $\rhobar^{\ss}$ of~$\rhobar$.
More generally, if $x$ and $y$ are two finite type points 
of~$|(\cX_{\red})_{\Fbar_p}|$, corresponding to Galois representations
$\rhobar$ and $\thetabar$ respectively,
then $y$ lies in $\overline{\{x\}}$ if and only if $\thetabar$ is a
virtual partial
semi-simplification of~$\rhobar$. % \mar{ME: Should probably explain what ``partial semi-simplification'' means.}
\end{enumerate}
\end{theorem}
\begin{proof}% \mar{TG: the argument that was written here before
    % (without the ``virtual'') ignored $H^2$ and seemed to implicitly
    % assume that~$H^1$ is compatible with base change. I don't think
    % the argument now written works in the same way and I don't see an
    % issue right now, but I'm leaving a comment here in case I'm
    % missing something!}
Since $\Fbar_p$-valued points of $\cX$ coincide with $\Fbar_p$-valued points of $\cX_{\red}$,
and so also with $\Fbar_p$-valued points of $(\cX_{\red})_{\Fbar_p}$, the claim of~(1)
is a particular case of the more general consequence of the Nullstellensatz that we
recalled above.

Suppose now that we have a short exact sequence  \numequation\label{eqn:
  virtual ses}0\to\rbar\to\rbar\oplus
    \rhobar\to\thetabar\to 0.\end{equation} We now follow the proof
  of~\cite[Prop.\ 3.4]{MR868301}
 to show that the point~$y$ corresponding
  to~$\thetabar$ lies in the closure of the point~$x$ corresponding
  to~$\rhobar$. Denote the morphism $\rbar\to\rbar\oplus\rhobar$ in~\eqref{eqn:
  virtual ses} by~$(f,g)$. Then for~$t$ in a sufficiently small open
neighbourhood~$U$ of $0\in\A^1$, the morphism
$(f+t\mathbf{1}_{\rbar},g):\rbar\to \rbar\oplus\rhobar$ is an
injection with projective cokernel (to see this one can for example
consider a splitting of\eqref{eqn:
  virtual ses} on the level of vector spaces); so we have a  morphism
(where we as usual abuse notation by writing families of
$(\varphi,\Gamma)$-modules as families of Galois representations)
$U\to\cX$ given by $t\mapsto \rhobar_t:=
(\rbar\oplus\rhobar)/\im(f+t\mathbf{1}_{\rbar},g)$. After possibly
shrinking~$U$ further we may assume that~$f+t\mathbf{1}_{\rbar}$ is an
automorphism of~$\rbar$ for all $\Fpbar$-points $t\ne 0$ of~$U$, so
that $\rhobar_t\cong\rhobar$; while for~$t=0$ we
have~$\rhobar_t\cong\thetabar$. Thus~$\thetabar$ is indeed
in the closure of~$\rhobar$, as claimed. This proves the ``if'' direction of~(3), % \mar{TG: maybe needs some kind
  % of induction on the length? ME: Yes, that's right}
and the ``only if'' direction of~(2).


% the $G_K$-representation $\rhobar$ sits in a short exact sequence
% \[0 \to \alpha \to \rhobar \to \beta\to 0.\]  Let  $e \in \Ext^1_{G_K}(\beta,\alpha)$
% be the class classifying this extension, and let $L$ denote the line spanned
% by $e$ over~$\Fbar_p$.  Then, under the embedding
% $$ L^{\vee} \otimes L \hookrightarrow
% L^{\vee}\otimes \Ext^1_{G_K}(\beta , \alpha) \iso
% \Ext^1_{G_K}(\beta ,  L^{\vee}\otimes \alpha),$$
% % \mar{TG: confused - isn't $L$ a line? Is this tensor product over
% %   something other than~$\Fpbar$? The RHS seems too small to me ME: Fixed}
% the canonical ``trace'' element in the source maps to an extension
% $$0 \to L^{\vee} \otimes \alpha \to \xi \to \beta \to 0.$$
% Identifying $L^{\vee}$ with the space of homogeneous linear functions on 
% the affine line $L$, and passing to $(\varphi,\Gamma)$-modules (but without
% changing our notation),
% the extension $\xi$ gives rise to an extension
% of $(\varphi,\Gamma)$-modules
% $$0 \to \alpha_L \to \widetilde{\xi} \to \beta_L \to 0$$ over $L$
% which interpolates the family of extensions classified by the points of~$L$.
% In particular, all the non-zero points of $L$ give rise to extensions which
% are isomorphic (as $(\varphi,\Gamma)$-modules) to  $\rho$,
% while the fibre of the family over $0 \in L$ is the direct  sum $\alpha \oplus \beta$.
% Thus this family is classified by a morphism
% $L \to (\cX_{\red})_{\Fbar_p}$ which factors through a morphism
% $L/\Gm \to (\cX_{\red})_{\Fbar_p}$; under this morphism the open point of $L/\Gm$ maps 
% to  $\rhobar$, while the  closed point maps to $\alpha \oplus \beta$, and so the
% latter  point  of $\cX_{\red}$ is in the closure of the former.
% This proves the ``if'' direction of~(3), % \mar{TG: maybe needs some kind
%   % of induction on the length? ME: Yes, that's right}
% and the ``only if'' direction of~(2).

The ``if'' direction of~(2) follows from the ``only if'' direction of~(3),
and so it remains to prove this latter statement.   To this end,
we fix  $\rhobar,\thetabar: G_K \to \GL_d(\Fbar_p)$, and assume that the point $y$
corresponding to $\thetabar$ lies in the closure of the point $x$ corresponding to~$\rhobar$.
In order to avoid discussing deformation theory over $\Fbar_p$, we extend 
our coefficients $\cO$ if necessary
so that $\rhobar$ and $\thetabar$ are both defined over $\F$.
We may and do also assume that $x$ and $y$ are distinct; i.e.\ that 
$\rhobar$ and $\thetabar$ are not isomorphic (as representations defined over $\Fbar_p$,
or equivalently, as representations defined over~$\F$).
Let $x_0$ and $y_0$ denote the images of $x$ and $y$ respectively in~$\cX_{\red}$;
then $x_0$ and $y_0$ are again distinct.
Let $\cZ$ denote the closure of $\{x_0\}$, thought of as a reduced closed
substack of $\cX_{\red}$; then $y_0$ is a point of~$\cZ$. 
%Similarly, let $\cW$ denote the closure of $\{y_0\}$.  Then $\cU := \cZ\setminus \cW$
%is an open substack of $\cZ$ which contains $\{x_0\}$,
%and thus is dense in~$\cZ$.
%\mar{This uses that $|\cX_{\red}|$
%is $T_0$, which follows from it being sober, which follows from $\cX_{\red}|$
%being quasi-separated.}

Let $R$ be the framed deformation ring of $\thetabar$ over~$\cO$; then $R$ is a versal
ring to $\cX$ at~$y$. 
Let $\Spf S = \Spf R \times_{\cX}\cZ$; then $S$
is a versal ring to $\cZ$ at $y_0$.   Since $\cZ$ is algebraic,
the versal ring $S$ is effective~\cite[\href{https://stacks.math.columbia.edu/tag/07X8}{Tag 07X8}]{stacks-project},  
and the corresponding morphism $\Spec S \to \cZ$
is furthermore flat~\cite[\href{https://stacks.math.columbia.edu/tag/0DR2}{Tag
  0DR2}]{stacks-project}. 
The morphism $\Spec \F \to \cZ$ corresponding to $x_0$
is quasi-compact and scheme-theoretically dominant;
thus the base-changed morphism
\numequation
\label{eqn:dominant}
\Spec S \times_{\cZ, x_0} \Spec \F \to \Spec S
\end{equation}
 is again scheme-theoretically dominant.
% \mar{ME: This uses preservation by flat
% base-change of scheme-theoretic images under quasi-compact
% morphisms. TG: are you suggesting that we add something to the text
% about this? I'm happy either way.} 

Let $\eta$ be a point in the image of~\eqref{eqn:dominant}, and let
$\Spec T$ be the Zariski closure of $\eta$ in $\Spec S$.  Then $T$ is a quotient
of $R$ which is an integral domain.  If $\cK$ denotes the fraction field of~$T$,
then the morphism $\Spec \cK \to \cX_{\red}$ factors through the morphism
$\Spec \F \to \cX_{\red}$ corresponding to $x_0$, and so the Galois representation
$G_K \to \GL_d(\cK)$ classified by the point $\eta$ of $\Spec T$ is isomorphic
to $\rhobar$ (see Section~\ref{subsubsec: Galois repns into certain
  fields} for the definition of this Galois representation).  % \mar{ME: This final conclusion requires some manipulations moving
% b/w $(\varphi,\Gamma)$-modules and Galois representations that are probably not
% literally justified in the text --- e.g.\ because we have to work with coefficients
% in the field $\cK$.}
Thus, if $L$ denotes the splitting field of~$\rhobar$,
then we see that the deformations of $\thetabar$ classified by $T$ all factor
through the finite group $\Gal(L/K)$. It follows immediately from~\cite[Thm.\ 1]{MR1757882}  % The general theory of 
% deformations of representations of finite groups over $\F$ then shows
that $\thetabar$ is a virtual partial semi-simplification
of~$\rhobar$, as required.
%\mar{ME: Haven't looked into CWE to find reference for this.}
%\mar{ME: To be finished \dots}
\end{proof}

\section{The substack of $G_K$-representations}\label{sec:
  comparison to CWE}% \mar{TG: another new section, which should be
  % carefully checked. We should probably also run it past CWE?}
We now briefly explain the relationship between our
stacks and the stacks of $G_K$-representations constructed by Wang-Erickson
in~\cite{MR3831282}. % The results in this section are
% essentially due to Wang--Erickson, as will be clear from the
% proofs. \mar{TG: that might change depending on the proofs!}
Write~$\cX_d^{\Gal}$ for the formal algebraic
stack characterised by the following property: if~$A$ is an
$\cO$-algebra in which $p$ is nilpotent, then $\cX_{d}^{\Gal}(A)$ is the
groupoid of continuous morphisms \[\rho:G_{K}\to\GL_d(A)\](where~$A$
has the discrete topology, and~$G_{K}$ its natural profinite
topology). That this \emph{is} a formal algebraic stack follows
from~\cite[Thm.\ 3.8, Rem.\ 3.9]{MR3831282}. Equivalently, we can
think of $\cX_{d}^{\Gal}(A)$ as the groupoid
of rank~$d$ projective $A$-modules~$T_A$ with a continuous action of~$G_K$
(where each~$T_A$ has the discrete topology).%\mar{TG: added final sentence to
%join up with the approach below. ME: Looks good to me} % \mar{TG: I'm confused by
  % this definition and citation! Maybe it's OK if $A$ is finite type,
  % but I'm not even sure of this.}



For any finite type $\cO/\varpi^a$-algebra~$A$, and
any~$T_A\in\cX_d^{\Gal}(A)$, we
set \[\mathbb{D}_A(T_A):=W(\C^\flat)_A\otimes_AT_A,\] which naturally has
the structure of a rank~ $d$ projective $(\varphi,G_K)$-module with
$A$-coefficients (with~$\varphi$ acting on the first factor in the
tensor product, and ~$G_K$ acting diagonally). (In the case that~$A$
is actually a finite $\cO/\varpi^a$-algebra, this agrees with the
construction given in Section~\ref{subsubsec:
Galois reps for GK phi}.)% \mar{TG: assuming
%   this isn't all nonsense, want to refer back to something which
%   should be added to Section~\ref{subsubsec:
% Galois reps for GK phi} giving this
%   description there.}

By Proposition~\ref{prop: equivalences of categories to Ainf}, for
each finite type $\cO/\varpi^a$-algebra~$A$ the assignment
$T_A\mapsto \mathbb{D}_A(T_A)$ gives a
functor~$\cX_d^{\Gal}(A)\to\cX_d(A)$.  Since both~$\cX_d^{\Gal}$
and~$\cX_d$ are limit preserving (for $\cX_d^{\Gal}$ this follows
easily from the definition, since any~$\rho$ as above factors through
a finite quotient of~$G_K$), this defines a
morphism of stacks
\numequation
\label{eqn:Galois stacks to ours}
\cX_d^{\Gal}\to\cX_d.
\end{equation}
We now show that this
morphism is in fact a monomorphism.

\begin{thm}
  \label{thm: map from Galois stack to ours}
  %There is a natural monomorphism
  %$\cX_d^{\Gal}\to\cX_d$, which  for  any finite
  %type~$\cO/\varpi^a$-algebra $A$, is given by the functor
  %$\cX_d^{\Gal}(A)\to\cX_d(A)$ taking $T_A\to \mathbb{D}_A(T_A)$.
  The morphism~{\em \eqref{eqn:Galois stacks to ours}} is a monomorphism,
  and is furthermore versal at finite type points in the sense of Definition~{\em \ref{defn:formally smooth points}} 
below.
\end{thm}
\begin{proof}
  The versality statement follows from the fact that for any
  finite Artinian $\Z_p$-algebra~$A$, the $A$-valued points
  of $\cX_d^{\Gal}$ and $\cX_d$ coincide.  (Over such a ring~$A$,
  \'etale $(\varphi,\Gamma)$-modules do arise from Galois representations.)

  To prove the monomorphism claim, we need to show that for  any finite
  type~$\cO/\varpi^a$-algebra $A$, the functor
  $\cX_d^{\Gal}(A)\to\cX_d(A)$ is fully faithful; that is, we need to
  show that given $T_1,T_2\in \cX_d^{\Gal}(A)$, we
  have \[\Hom_{G_K,A}(T_1,T_2)\stackrel{?}{=}\Hom_{\varphi,G_K}(\mathbb{D}_A(T_1),\mathbb{D}_A(T_2)).\]
  Writing $T:=T_1^\vee\otimes_AT_2$, it suffices to show
  that \[T^{G_K}\stackrel{?}{=}\mathbb{D}_A(T)^{\varphi=1,G_K},\]or even
  that  \[T\stackrel{?}{=}\mathbb{D}_A(T)^{\varphi=1}.\] Since~$T$ is
  a finite projective $A$-module, we
  have  \[\mathbb{D}_A(T)^{\varphi=1}:=(W(\C^\flat)_A\otimes_AT)^{\varphi=1}=(W(\C^\flat)_A)^{\varphi=1}\otimes_AT,\]and
  the result follows from Lemma~\ref{lem:phi invariants}. % \mar{TG: note that my earlier version with
    % $(W(\C^\flat)_A)^{\varphi=1,G_L}=A$ would also suffice for this
    % argument, by applying it with~$G_L$ acting trivially on~$T$, and
    % then further passing to $G_K$-invariants.}
\end{proof}



% \mar{TG: hopefully not nonsense, I
%   think the $A$-action just comes along for the ride? ME: Should at least note
% that we naturally get an $\A_{K,A}$-module. (Now added)}
% % functors~$T_A$, $\mathbf{D}_A$ defined by
% \[\mathbb{D}_A(V):=\A_{K,A}\otimes_{\A_K\otimes_{\Zp}A}(\AKnrhat
%   \otimes_{\Zp}V)^{G_{\Kcyc}},\] which has the structure of a
% $(\varphi,\Gamma)$-module over~$\A_{K,A}$.\mar{TG: note that I changed
% the definition; hopefully not nonsense!} % $\A_{K} \otimes_{\Z_p} A =: \A_{K,A}$.
% It  is not
% obvious that it is a finite projective % formal \mar{TG: confused as to
%   why we have ``formal'' here. ME: I don't think we should; where did
%   it come from? TG: probably I wrote it in another document when I was
% bashing stuff out, so I've deleted it here. Commented out.}
% \'etale
% $(\varphi,\Gamma)$-module with $A$-coefficients; this is the content
% of the following theorem, which is essentially due to Wang--Erickson.   % \[T_A(M)=(  
%\AKnrhat \otimes_{\A_{K}}M)^{\varphi=1}\]
% \mar{TG:
% obviously this needs to be filled out and made correct.}
% giving equivalences of categories between the category of finite projective formal \'etale
% $(\varphi,\Gamma)$-modules with $A$-coefficients, and the category of
% finite free $A$-modules with a continuous action of~$G_K$.

% It is presumably possible to prove the following lemma directly, and
% to prove the stronger statement that if $A$ is a finite type
%   $\cO/\varpi^a$-algebra then $W(\C^\flat)_A^{\varphi=1}=A$, but
% we prefer to give a proof using our description of the stack of rank
% $1$ \'etale $\varphi$-modules in Chapter~\ref{sec: the rank one case}
% (which does not make use of the results of this section).
% \begin{lem}\mar{TG: since this is true for all~$K$ I guess one could
%     try to deduce the statement about $\varphi$-invariants from it,
%     but I have not attempted to do so. It should definitely be checked
%   for nonsense\dots}
%   \label{lem: phi invariants in W C flat}If~$A$ is a finite type
%   $\cO/\varpi^a$-algebra, then $W(\C^\flat)_A^{\varphi=1,G_K}=A$.
% \end{lem}
% \begin{proof}Let~$M$ be the rank $1$ $(\varphi,Gamma)$-module over
%   $\A_{K}$ which corresponds to the trivial $G_K$-representation, and
%   let~$M_A:=\A_{K,A}\otimes_{\A_K}M$. Set~$\tM_A:=W(\C^\flat)_A\otimes_{\A_{K,A}}M_A$;
%   this is a $(\varphi,G_K)$-module with $A$-coefficients, and by
%   construction we
%   have \[\Hom_{G_K,\varphi}(\tM_A,\tM_A)=W(\C^\flat)_A^{\varphi=1,G_K}.\]
%   By Proposition~\ref{prop: equivalences of categories to Ainf}, we
%   also
%   have  \[\Hom_{G_K,\varphi}(\tM_A,\tM_A)=\Hom_{\A_{K,A},\varphi,\Gamma}(M_A,M_A).\]
%   By Proposition~\ref{prop:rank 1 X}, we in turn
%   have \[\Hom_{\A_{K,A},\varphi,\Gamma}(M_A,M_A)=\Hom_{G_K}(A,A)\] (where
%   $A$ has the trivial action of~$G_K$), and this latter group is
%   obviously equal to~$A$.
% \end{proof}
% \begin{thm}
%   \label{thm: map from Galois stack to ours}There is a natural monomorphism
%   $\cX_d^{\Gal}\to\cX_d$, which  for  any finite
%   type~$\cO/\varpi^a$-algebra $A$, is given by the functor
%   $\cX_d^{\Gal}(A)\to\cX_d(A)$ taking $V\to \mathbb{D}_A(V)$.
% \end{thm}
% \begin{proof}
%   The theorem amounts to the statement that for any finite type
%   $\cO/\varpi^a$-algebra $A$, and any object~$V\in\cX_d^{\Gal}(A)$,
%   $\mathbb{D}_A(V)$ is an object of~$\cX_d(A)$; and $\mathbb{D}_A$
%   defines an equivalence of categories from $\cX_d^{\Gal}(A)$ to its
%   essential image in $\cX_d(A)$. These statements follow from
%   arguments extremely similar to those used to prove~\cite[Defn./Lem.\
%   4.2, Prop.\ 4.4, (4.10)]{MR3831282}, which prove closely related
%   statements for~$G_{K_\infty}$-representations (albeit for
%   $\varphi$-modules over $(\gS\otimes_{\Zp}A)[1/u]$, rather than over
%   the completed ring~$\OEA$ that we use in this book).

%   We briefly sketch the proof. We define a functor~$V_A$ from the essential
%   image of  $\mathbb{D}_A$ to $\cX_d^{\Gal}(A)$ by
% \mar{ME: Shouldn't the tensor be over $\A_K$?}
% \[V_A(M):= (\AKnrhat
%   \otimes_{\Zp}M)^{\varphi=1},\] which inherits a $G_K$-action from
% the diagonal action on~$\AKnrhat$ and~$M$ (in the latter case, via
% inflation of the action of~$\Gamma_K$). It is enough to show that the composite
%   $V_A\circ \mathbb{D}_A$ is naturally isomorphic to the identity functor.

%   If~$A=\cO/\varpi^a$, this holds, and indeed (as we recalled in
%   Section~\ref{subsec: phi gamma over Zp}) the functors~$\mathbb{D}_A$
%   and~$V_A$ extend with the same definitions to an equivalence of
%   categories between continuous representations of~$G_K$ on finite
%   $A$-modules and \'etale $(\varphi,\Gamma)$-modules with
%   $A$-coefficients.

%   For general~$A$, since the action of~$G_{K}$ % \mar{TG: even~$G_K$?}
%   on~$V$ factors through a finite quotient, $V$ can be written as an
%   injective direct limit $\varinjlim_iV_i$, where the~$V_i$ run
%   through the (literally) finite $G_K$-stable
%   $\cO/\varpi^a$-submodules of~$V$. Exactly as in the proof
%   of~\cite[Defn./Lem.\ 4.2]{MR3831282}, the functors $\mathbb{D}_A$
%   and~$V_A$ commute with injective direct limits
% \mar{ME: Just looking at the citation, I think these should maybe 
% be $\mathbb{D}_{\cO/\varpi^a}$  and $V_{\cO/\varpi^a}$ here. I'm
% then also not sure I believe the following argument as written (although
% I'm sure that the result is true, and think 
% that a variant argument will work).}
%   (the only additional ingredient which is required in our setting is to note that
%   $\A_{K,A}$ is faithfully flat over ${\A_K\otimes_{\Zp}A}$; this is
%   true because it is a localisation of the statement that $\A_{K,A}^+$
%   is faithfully flat over ${\A_K^+\otimes_{\Zp}A}$, which is true
%   because the former ring is a completion of the latter, and the
%   latter is Noetherian).  From this it is easy to deduce the statement
%   for general~$A$ %from the % that $V_A$ provides a one-sided inverse
%   % (up to natural isomorphism) to~$\mathbb{D}_A$,
%   by passage to limits
%   from the case that ~$A=\cO/\varpi^a$, as
%   required. % , by~\cite[Cor.\ 3.10]{MR3831282}.
% \end{proof}%\mar{TG: there are also issues of $A((u))$ versus $\cO((u))\otimes A$?}

% \begin{rem}
%   \label{rem: fixed determinant map Galois stack}We can define a fixed
%   determinant stack
%   $\cX_d^{\Gal,\chi}$ in exactly the same way as we
%   defined~$\cX_d^\chi$, and it follows from the definition that we
%   have an induced monomorphism $\cX_d^{\Gal,\chi}\to\cX_d^\chi$.
% \end{rem}

\begin{rem}%\mar{TG: this remark is intended to replace the earlier
    %claim that we're just following CWE, since we now aren't.}
  \label{rem: comparison to what CWE did}In~\cite[\S4]{MR3831282},
  Wang-Erickson explains how to associate \'etale $\varphi$-modules to
  $G_{K_\infty}$-representations with open kernel. A similar argument
  would allow us to associate projective \'etale
  $(\varphi,\Gamma)$-modules to objects of~$\cX_d^{\Gal}(A)$. However,
  if we directly followed the strategy of~\cite{MR3831282},
  the resulting \'etale $(\varphi,\Gamma)$-modules
  would have as coefficients the rings $\A_{K}\otimes_{\Zp}A$,
  rather than the rings~$\A_{K,A}$ that we use throughout this book.
  Thus this approach would not obviously yield Theorem~\ref{thm: map from
    Galois stack to ours}, which is why we've preferred to make a more direct
  argument in terms of \'etale~$(\varphi,G_K)$-modules.
\end{rem}


\begin{rem}
  \label{rem: CWE stack is exactly the substack where there is a
    Galois representation}It follows immediately from
  Theorem~\ref{thm: map from Galois stack to ours} that~$\cX_d^{\Gal}$
  is largest substack of~$\cX_d$ over which the
  universal $(\varphi,\Gamma)$-module can be realised as a
  $G_K$-representation.
  
In fact, it should be possible
to strengthen Theorem~\ref{thm: map from Galois stack to ours}.
For example, \cite[Thm.\ 3.8]{MR3831282}
  shows that we can write \numequation\label{eqn: Galois breaks up
    over
    pseudorepns}\cX_d^{\Gal}=\coprod_{D}\cX_{d,D}^{\Gal},\end{equation}where~$D$
runs over the isomorphism classes of $d$-dimensional semisimple
$\Fpbar$-representations of~$G_K$, and~$\cX_{d,D}^{\Gal}(A)$ is the
groupoid of those~$\rho$ the semisimplification of whose reductions
modulo~$\varpi$ is~$D$. In view of Theorem~\ref{thm:closed points},
we expect that the monomorphism 
$(\cX_{d,D}^{\Gal})_{\red} \to \cX_{d,\red}$ (induced by the monomorphism of
Theorem~\ref{thm: map from Galois stack to ours})
will actually be a closed immersion, and hence (given the versality statement of
Theorem~\ref{thm: map from Galois stack to ours}) that
$\cX_{d,D}^{\Gal}$ will be identified with the formal completion
of $\cX_d$ along the image of this closed immersion.
%the decomposition~\eqref{eqn: Galois breaks up over pseudorepns} suggests
% \mar{ME: Softened ``means'' to ``suggests'', since the next  sentence
% indicates that we are not being precise here. TG: agreed, commenting
% out.}
%that we can regard $\cX_d^{\Gal}$ as a disjoint union of certain
In fact, this image will contain a unique closed point (corresponding
to the semisimple Galois representation~$D$),
and we imagine that $\cX_{d,D}^{\Gal}$ could also be identified
with the {\em coherent completion}
(in the sense ~\cite[Defn.\ 2.1]{MR4088350})
of $\cX_d$ a this closed point.
However, we don't pursue these ideas further here.
%completions of~$\cX_d$ at the closed points of~$\cX_{d,\red}$. We
%anticipate that this can be made precise by using the language of
%coherent completeness~\cite[Defn.\ 2.1]{MR4088350}, but we do not pursue this here.
\end{rem}
\begin{rem}
  \label{rem: could also consider WD}As well as considering
  $G_K$-representations, it is also possible to consider a larger
  substack of~$\cX_d$ over which our $(\varphi,\Gamma)$-modules can be
  realised as ``Weil--Deligne''-representations. (Here the notion of
  Weil--Deligne representations is not quite the usual one, but rather
  is given by representations of certain discretizations of~$G_K$, as described
  in~\cite[\S 1.2]{emerton2020moduli} and the references therein.)

  In the case~$K=\Qp$ and~$d=2$, this locus is discussed in~\cite[\S
  1]{emerton2020moduli}; roughly speaking, the difference between it
  and $\cX_d^{\Gal}$ is that it contains families of semisimple
  representations given by direct sums of unramified twists of fixed
  irreducible representations (the key point being that the universal
  unramified character does not correspond to a $G_K$-representation,
  but does correspond to a representation of the Weil group~$W_K$,
  see Remark~\ref{rem: rank 1 versus rank d} below).
\end{rem}
% \begin{rem}\label{rem: existence of pseudodeformation space is
%     mysterious}As is shown in~\cite{MR3831282}, the stack~$\cX_d^{\Gal}$ has a well-defined ``coarse
%   moduli space'', given by the moduli stack of
%   pseudorepresentations. More precisely, we have the
%   decomposition~\eqref{eqn: Galois breaks up over pseudorepns}, and
%   the coarse moduli space of $\cX_{d,D}^{\Gal}$ is the quotient
%   stack\mar{TG: should be a quotient of $\Spf R_D$? Not actually
%     finding this in CWE right now.}

%   It is unclear whether to expect that there is a reasonable notion of
%   a coarse moduli space for the stacks~$\cX_d$. In the case $d=2$,
%   $K=\Qp$ such a space does exist, as shown in~\cite{DottoEG}, but its
%   underlying reduced scheme is not affine, in contrast to the case of
%   ~$\cX_d^{\Gal}$.

%   % Say: existence of some kind of trace map on our stack (to a
%   % ``pseudodeformation space'') is harder, for $\GL_2(\Qp)$ it will be
%   % in~\cite{DottoEG}. \mar{TG: finish this}
  
% \end{rem}

% \chapter{Fixing determinants}\label{sec: fixing determinants}\mar{TG:
%   this might or might not be its own section in the end, but it seems
%   easiest to leave it as one for now.}
% \mar{TG: maybe remove this again!}

% \section{Fixing the determinant}\label{subsec: fixed determinant
%   stacks}

% There is a natural morphism $\wedge^d:\cX_d\to\cX_1$, sending a rank~$d$
% projective \'etale $(\varphi,\Gamma)$-module $D$ to~$\wedge^dD$. 

% We define~$\cY_d$ as the fibre product 

% \[ \begin{tikzcd}
%       \cY_d\ar{d}\ar{r}&\cX_d\times\cX_1\ar{d}{\wedge^d\times\pr_2}\\
%   \cX_1\ar{r}{\Delta}&\cX_1\times\cX_1
% \end{tikzcd}\]
   


% Let~$\chi:G_K\to\cO^\times$ be a character, 
% determining a morphism
% $\Spf\cO\to\cX_1$. Write~$\chibar:G_K\to\F^\times$ for the reduction
% of~$\chi$ modulo$~\varpi$. We
% set \[\cX_d^{\chi}:=\cY_d\times_{\cX_1}\Spf\cO,\]so that $\cX_d^\chi(A)$
% is the groupoid of pairs $(D,\theta)$ where~$D$ is a rank~$d$
% projective \'etale $(\varphi,\Gamma)$-module with $A$-coefficients,
% and $\theta$ is an identification of~$\wedge^dD$ with~$\chi$.


% If~$\chi$ is a crystalline character and~$\underline{\lambda}$ is a
% Hodge type, we define \[\cX_d^{\crys,\lambdau,\chi}:=\cX_d^{\crys,\lambdau}\times_{\cX_d}\cX_d. \] By
% construction, $\cX_d^{\crys,\lambdau,\chi}$  is a closed substack
% of~$\cX_d^\chi$, and by~\cite[Thm.\ 4.8.12]{emertongeepicture}, it is
% a $p$-adic formal algebraic stack. It is nonzero if and only
% if~$\lambdau$ is compatible with ~$\chi$ in the following sense.

% \begin{defn}
%   \label{defn: Hodge type compatible with chi }We say that ~$\lambdau$
%   is \emph{compatible with} ~$\chi$ if $\chi$ is crystalline, and for
%   each $\sigma:K\into\Qpbar$ we have $\HT_\sigma(\chi)=\sum_{i=1}^d\lambda_{\sigma,i}$.
% \end{defn}

% % \mar{TG: need a notion of crystalline version of fixed determinant
% %   stack which we can define as a fibre product, should be easy to
% %   prove equidimensionality of it by the same argument with versal
% %   rings, thankfully.}
% \begin{cor}\label{cor: fixed determinant crystalline
%     stacks}If~$\lambdau$ is a regular Hodge type which is compatible
%   with~$\chi$, then $\cX_d^{\crys,\lambdau,\chi}\times_{\Spf
%     \cO}\Spec\F$ is equidimensional of dimension~$[K:\Qp]d(d-1)/2$. 
% \end{cor}
% \begin{proof}The proof of ~\cite[Thm.\
%   4.8.14]{emertongeepicture} goes over immediately, replacing the
% crystalline   deformation rings with their fixed determinant variants
% (whose dimensions are known by the proof of~\cite[Thm.\
% 3.3.8]{MR2373358}; for a precise statement see for example~\cite[Thm.\
% A]{zbMATH07042062}), and the reductive group~$\GL_d$ by~$\SL_d$.  
% \end{proof}

% \begin{defn}
%   \label{defn: compatible Serre weight}We say that a Serre
%   weight~$\underline{k}$ is \emph{compatible with~$\chi$}
%   if
%   \[\chibar|_{I_K}=\varepsilonbar^{-d(d-1)/2}\prod_{\sigmabar:k\into\Fpbar}\omega_{\sigmabar}^{-\sum_{i=1}^dk_{\sigmabar,i}}.\]
% \end{defn}

% We now recall the \emph{eigenvalue morphism}.   If~$\rhobar_{\cT}$ is generically maximally nonsplit of
%   niveau one and weight~$\underline{k}$, then we
%   have
%    \numequation\label{eqn:weight k version of maximally nonsplit structure}\rhobar_t\cong \begin{pmatrix}
%       \ur_{\nu_d(t)}\otimes\omega_{\underline{k},d} &*&\dots &*\\
%       0& \ur_{\nu_{d-1}(t)}\otimes\epsilon^{-1}\omega_{\underline{k},d-1}&\dots &*\\
%       \vdots&& \ddots &\vdots\\
%       0&\dots&0& \ur_{\nu_1(t)}\otimes\epsilon^{1-d}\omega_{\underline{k},1}\\
%     \end{pmatrix},  \end{equation}
% and the eigenvalue morphism % \mar{TG: needs a better name! These
%   % aren't slopes at all but I can't think of anything better right now.}
% ~$\underline{\nu}:\cU\to(\Gm)^d_{\underline{k}}$ is given
% by~$(\nu_1,\dots,\nu_d)$. (Here
% $(\Gm)^d_{\underline{k}}$ denotes the closed subgroup scheme of $(\Gm)^d$ parameterizing 
% tuples~$(x_1,\dots,x_d)$ for which~$x_i=x_{i+1}$
% whenever~$k_{\sigmabar,i}-k_{\sigmabar,i+1}=p-1$ for all~$\sigmabar$.)


% % Note that $(\Gm)_{\underline{k}}^d$ is closed under the simultaneous multiplication
% % action of $\Gm$ on $(\Gm)_{\underline{k}}^d$ (i.e.\ the action $a \cdot (x_1,
% % \ldots,x_d) := (a x_1,\ldots, a x_d)$).

% The following is the analogue for~$\cX_d^\chi$ of our main results for~$\cX_d$.

% \begin{thm}
%   \label{thm: fixed determinant stack in general}$\cX_d^\chi$ is a Noetherian
%   formal algebraic stack. Its underlying reduced substack $\cX_{d,\red}^\chi$
%   is of finite type over~$\Fp$, and is equidimensional of
%   dimension~$[K:\Qp]d(d-1)/2$. The irreducible components of~$\cX_{d,\red}$ admit a
%   natural bijection with the Serre weights which are compatible
%   with~$\chi$.\mar{TG: is this actually true?? It looks like
%     unramified quadratic twists are an issue for $\Sym^{p-1}$
%     components, for example??}
% \end{thm}
% \begin{proof}We follow the proofs of~\cite[Thm.\ 5.5.11,
%   6.5.1]{emertongeepicture} closely. In particular, the argument of
%   the proof of~\cite[Thm.\ 6.5.1]{emertongeepicture} goes through
%   almost unchanged to show that every irreducible component
%   of~$\cX_{d,\red}^\chi$ has dimension at least~$[K:\Qp]d(d-1)/2$, % and
%   % that every irreducible component of~$\cY_{d,\red}$ has dimension at
%   % least~$[K:\Qp]d(d-1)/2+1$,
%   replacing the appeal to~\cite[Thm.\
%   4.8.14]{emertongeepicture} with one to Corollary~\ref{cor: fixed
%     determinant crystalline stacks}. (The one possibly subtle point is
%   that~ $\cX_{d,\red}^\chi$ only depends on~$\chibar$, so in making
%   the comparison to $\cX_d^{\crys,\lambdau,\chi}$ we are free to
%   replace~$\chi$ by a crystalline character which is compatible with~$\lambdau$.)

%   Let~$\cZ$ be an irreducible component of~$\cX_{d,\red}^\chi$.  We now
%   consider the morphism \numequation\label{eqn: twisting from fixed
%     determinant to unfixed}f:\cX_{d,\red}^\chi\times\Gm\to
%   \cX_{d,\red}\end{equation} given by forgetting~$\theta$ and taking
%   unramified twists, as in~\cite[\S 5.3]{emertongeepicture}. Note that
%   any family of $(\varphi,\Gamma)$-modules with determinant~$\chi$ is
%   essentially twistable in the sense of~\cite[\S
%   5.3.1]{emertongeepicture} (because any twist preserving the
%   determinant would have to be by a $d$th root of unity). Let~$\cW$ be
%   the scheme-theoretic image of~$\cZ\times\Gm$ in
%   $\cX_{d,\red}^\chi\times\Gm$.

%   It follows from~\cite[Lem.\ 5.3.2]{emertongeepicture}, together with
%   the lower bound proved above for the dimension of~$\cZ$ that the
%   scheme-theoretic image of~$f(\cW)$ under~\ref{eqn:
%     twisting from fixed determinant to unfixed} has dimension at least
%   $[K:\Qp]d(d-1)/2$. (Here we made implicit use of
%   ~\cite[\href{https://stacks.math.columbia.edu/tag/0DS4}{Tag
%     0DS4}]{stacks-project}, together with the fact that the forgetful
%   morphism $\cX_{d,\red}^\chi\to\cX_{d,\red}$ has fibres of dimension
%   (at most) $1$.) % \mar{TG: we need to put in something about the
%   %   fibres of $\cX_{d,\red}^\chi\to\cX_{d,\red}$ having dimension 1 to justify this?
%   % ~\cite[\href{https://stacks.math.columbia.edu/tag/0DS4}{Tag 0DS4}]{stacks-project}} 
% Since~$\cX_{d,\red}$ is
%   equidimensional of this dimension, we see that~$\cZ$ must have
%   dimension exactly $[K:\Qp]d(d-1)/2$, and that~$f(\cW)$ is
%   an irreducible component of~$\cX_{d,\red}$, so
%   that~$f(\cW)=\cX_{d,\red}^{\underline{k}}$ for some Serre
%   weight~$\underline{k}$.

% It follows from Definition~\ref{defn: compatible Serre weight}
% and~\cite[Defn.\ 5.5.1]{emertongeepicture} (together with the defining
% property of $\cX_{d,\red}^{\underline{k}}$) that~$\underline{k}$ is
% compatible with~$\chi$. It therefore remains to show that the
% association of~$\underline{k}$ to~$\cZ$ is injective and surjective.


% Now, \mar{TG: this is the claim that is not currently proved
%   in~\cite{emertongeepicture}, but it will be easy to add; since I am
%   making statements that hold at generic points, it might be necessary
% to do some further shrinking. The idea about the irreducibility
% statement is just to prove it by induction, using that $\Ext^1$
% groups are themselves irreducible, being vector spaces.} there is a
% dense open substack~$\cU$ of~$\cX_{d,\red}^{\underline{k}}$ which is
% maximally nonsplit of niveau one and weight~$\underline{k}$, and for
% which the eigenvalue morphism \[\nu:\cU\to (\Gm)^d_{\underline{k}}\]
% has the following property: there is a dense open subset~$U\subset
% (\Gm)^d_{\underline{k}}$ such that for any $\Fpbar$-point $x\in U$,
% the fibre~$\cU_x$ is irreducible. Furthermore, the determinant map
% $\det:(\Gm)^d_{\underline{k}}\to\Gm$ taking $(x_1,\dots,x_d)\to
% x_1\cdots x_d$ restricts to a surjection $U\to\Gm$.\mar{TG: maybe just
%   want to say surjective on $\Fpbar$-points?}

% Let~$\ur_{\alpha}$ be the character such that
% \[\chibar=\ur_{\alpha}\varepsilonbar^{-d(d-1)/2}\prod_{\sigmabar:k\into\Fpbar}\omega_{\sigmabar}^{-\sum_{i=1}^dk_{\sigmabar,i}},\]
% and let~$\cU_{\alpha}$ be defined as the fibre product 
% \[ \begin{tikzcd}
%       \cU_\alpha\ar{d}\ar{r}&\cU\times\cX_1\ar{d}{\det\times\pr_2}\\
%   \Spec\Fpbar\ar{r}{\Delta_\alpha}&\cX_1\times\cX_1.
% \end{tikzcd}\]where~$\Delta_\alpha$ is the composite of the
% diagonal~$\Delta$ and the map~$\Spec\Fpbar\to\cX_1$ corresponding
% to~$\ur_\alpha$. By construction, $\cU_\alpha$ is a substack
% of~$\cX_{d,\red}^\chi$.\mar{TG: this definition is not currently
%   correct, $\det$ is valued in~$\Gm$. Leaving for now as easy to fix
%   if this is the right way to be arguing.}

% We claim that~$\cU_{\alpha}$ is
% irreducible.\mar{TG: hopefully this follows from the above in
%   some way??} Now, $\cU_{\alpha}$ is non-empty by the surjectivity
% of $\det:U\to\Gm$, and the closure of~$\cU_{\alpha}$ in
% $\cX_{d,\red}^\chi$ has dimension~$[K:\Qp]d(d-1)/2$ (for example by
% another appeal to ~\cite[\href{https://stacks.math.columbia.edu/tag/0DS4}{Tag
%     0DS4}]{stacks-project}), so it is an irreducible component~$\cZ'$
%   of~$\cX_{d,\red}^\chi$. It remains to prove that~$\cZ=\cZ'$. To this
%   end, suppose to the contrary that~$\cZ\ne\cZ'$. \mar{TG: I still
%     seem to be failing to finish from here! The following is what we
%     had before.}


% % Both
% % statements follow from the following claim: if $\underline{k}$ is
% % compatible with~$\chi$, then there is an irreducible
% % substack~$\cX'$ of~$\cX_{d,\red}^{\underline{k}}$ of
% % dimension~$[K:\Qp]d(d-1)/2$, all of whose finite type points have
% % determinant~$\chi$, which is maximally nonsplit of niveau one and
% % weight~$\underline{k}$, and having the property that the restriction
% % of the eigenvalue morphism to~$\cX'$ dominates the subscheme
% % of~$(\Gm)_{\underline{k}}^d$ of elements whose product is determined by~$\chi$. \mar{TG: is this obviously enough? The idea is that
% %   for injectivity, any two such components would have to have generic
% %   points in common.}

% Look at the locus of $(\rhobar,\alpha)$ where
% $\alpha^d=\chi\det\rhobar$ (where this equality is a structure; and
% maybe we want to be working on the part where the last character is
% fixed, or maybe we don't), and
% map this to $\alpha^{-1}\circ\det\otimes\rhobar$. This will be a
% finite morphism back again. To see that the composite in each
% direction is the identity it will be helpful to think about
% $(\rhobar,1)$.\mar{TG: fill it out!}

% This claim follows from the inductive proof of~\cite[Thm.\ 5.5.11,
%   6.5.1]{emertongeepicture}.\mar{TG: cop out! Although I think this
%     does indeed follow easily from the proof. Can we do
%     better/clearer?}\mar{TG: generic fibre of eigenvalue morphism is
%     irreducible seems to be the key statement, and this will follow by
%   induction - we're pulling back from a map  $(\Gm)^{d-1}\to(\Gm)^d$
%   sending $(\alpha_1,\dots,\alpha_{d-1})\to
%   (\alpha_1\pi,\dots,\alpha_{d-1}\pi,1)$, or something like this.}
% \end{proof}

%   % We deduce the equidimensionality and the labelling of the
%   % irreducible components by Serre weights from the analogous results
%   % for~$\cX_d$, which follow from the proof
%   % of~\cite[Thm.\ 5.5.11]{emertongeepicture}.\mar{TG: have to figure out precisely
%   %   how this goes! Maybe we are trying to show that something is a
%   %   $\Gm$-torsor? But I \emph{think} that the following should be the
%   %   key point.} More precisely, the key point is the following: by
%   % induction on~$d$ (the case~$d=1$ being trivial), one sees that the
%   % determinant of the eigenvalue morphism~$\underline{\nu}$ is surjective.

%   % \mar{TG: I think there should be a better statement about it being
%   %   homogeneous; e.g.\ if you think about the case~$d=3$, we've ended
%   %   up with the unramified parts of the three characters being given
%   %   by $(abc,bc,c)$ as $(a,b,c)$ varies over~$(\Gm)^3$, so the
%   %   determinant is $ab^2c^3$ and in particular the factor of~$a$ is
%   %   somehow key?}
  



% % \begin{thm}
% %   \label{thm: fixed determinant stack for GL2 Qp}$\cX$ is a Noetherian
% %   formal algebraic stack. Its underlying reduced substack $\cX_{\red}$
% %   is of finite type over~$\Fp$, and is equidimensional of
% %   dimension~$[K:\Qp]d(d-1)/2$. The irreducible components of~$\cX_{\red}$ admit a
% %   natural bijection with the Serre weights which are compatible
% %   with~$\chi$.
% % \end{thm}
% % \begin{proof}
  
% % \end{proof}


\chapter{The rank one case}\label{sec: the rank one case}% \mar{TG: this
% needs to be done.}
In this chapter we describe some of our key constructions explicitly
in the case where $d = 1$.
More precisely, following the notation of
Chapter~\ref{sec: phi modules and phi gamma modules},
we give explicit descriptions of the stacks $\cR_d$, $\cR_d^{\Gamma_{\disc}}$,
and $\cX_d$, all in the case when $d = 1$ (and $K$ is arbitrary).

\section{Preliminaries}
While it may be possible to find an explicit description 
of the stacks we are interested in
directly from their definitions, this is not how 
we proceed.  Rather, we use the relationship between \'etale $\varphi$-modules
(resp.\ \'etale $(\varphi,\Gamma)$-modules) and Galois representations
to construct certain affine formal algebraic spaces $U$,
$V$, and $W$,
along with morphisms $U \to \cR_1$, $V \to \cR^{\Gamma_{\disc}}_1$,
and $W \to \cX_1$,
each satisfying the conditions of Lemma~\ref{lem:presentations} below.
An application of this lemma then yields a description
of each of $\cR_1,$ $\cR^{\Gamma_{\disc}}_1,$ and $\cX_1$.

%Recall that if $\cF$ is a category fibred in groupoids over a locally Noetherian
%scheme~$S$, if $k$ is a finite type $\cO_S$-field (i.e.\ if there is
%a morphism $\Spec k \to \Spec \Lambda \to S$, with the first morphism being of finite type
%and the second being an open immersion --- in what follows we fix such
%a choice of $\Lambda$ for each such~$k$,
%although the notions are independent of this choice),
%and if $x:\Spec k \to \cF$ is a morphism,
%then we may define the category $\widehat{\cF}_x$, which is a category
%cofibred in groupoids over the category $\cC_{\Lambda}$ of Artinian local $\Lambda$-algebra
%thickenings of $k$ \cite[\S 2.2]{EGstacktheoreticimages}:  an object of $\widehat{\cF}_x$ consists of an object
%$A$ of $\cC_{\Lambda}$ together with a morphism $\Spec A \to \cF$ which prolongs
%the given morphism $x: \Spec k \to \cF$, and a morphism between two such objects
%lying over $A$ is a morphism in $\cF$ which induced the identity when restricted to $x$.
%
%%We now make a definition.
%%\mar{TG: making something up; perhaps there should be a lemma then a definition.}

\begin{defn}\index{versal at finite type points}
  \label{defn:formally smooth points}Let~$S$ be a locally Noetherian
  scheme, let $\cX$ and $\cY$ be stacks over~$S$ and let $f:\cX\to\cY$ be a
  morphism.   We say that $f$ is {\em versal
at finite type points} if for each morphism $x:\Spec k \to \cX$,
with $k$ a finite type $\cO_S$-field,
%inducing a morphism $y:\Spec k \to Y$ by
%composing with~$f$, the induced morphism
%$\widehat{X}_x \to \widehat{\cY}_y$
%of categories cofibred in groupoids over $\cC_{\Lambda}$ 
%is smooth,
%in the sense
%of~\cite[\href{https://stacks.math.columbia.edu/tag/06HG}{Tag 06HG}]{stacks-project}.
%%Let~$x\in X$ be a finite type
%%  point. Assume that~$f$ is representable by algebraic spaces. Then we
%%  say that~$f$ is \emph{versal at~$x$} if for each
%%  morphism~$Y\to\cY$ whose source is a scheme, and each finite type
%%  point~$x'$ of the algebraic space ~$X':=X\times_{\cY}Y$ lying
%%  over~$x$, the morphism $X'\to Y'$ is versal at~$x'$.
the morphism $f$ is versal at~$x$, in the sense of Definition~\ref{def:versal}.
\end{defn}

\begin{remark}
As the discussion of versality in Appendix~\ref{app: formal algebraic stacks}
should make clear,
this definition is related to the various properties considered in
\cite[Def.~2.4.4]{EGstacktheoreticimages}
and the surrounding discussion.  As noted in 
\cite[Rem.~2.4.6]{EGstacktheoreticimages}, it is somewhat complicated to define
the notion of versality at literal points of a stack, since the points of a stack
are (by definition) equivalence classes of morphisms from the spectrum of a field
to the stack, and it is not immediately clear in general whether
or not  the versality condition would be
independent of the choice of equivalence class representative.  
In the preceding definition we obviate this point, by directly requiring 
versality at all (finite type) representatives of all (finite type) points.
\end{remark}

\begin{lemma}
\label{lem:base-changing versality}
If $\cX \to \cY$ is a morphism of stacks over a locally Noetherian base $S$ which 
is versal at finite type points, %and which is representable by algebraic stacks,
and if $\cZ \to \cY$ is a morphism of stacks, then the induced morphism
$\cX \times_{\cY} \cZ \to \cZ$ is again versal at finite type points.
\end{lemma}
\begin{proof}
%\mar{ME: fill in}
Let $\tx:\Spec k \to \widetilde{\cX} := \cX\times_{\cY} \cZ$ with $\Spec k$ of finite type over $S$;  from the definition of the $2$-fibre product $\widetilde{\cX}$,
we see that giving the morphism $\tx$ amounts to giving the 
pair of morphisms $x: \Spec k \to \cX$ and $z: \Spec k \to \cZ$ (the result
of composing $\tx$ with each of the projections),
as well as an isomorphism $y \iso y'$, where $y$ and $y'$ are the $k$-valued
points of $\cX$ obtained respectively by composing~ $x$ with the morphism $\cX \to \cY$
and by composing $z$ with the morphism $\cZ \to \cY$.
It is then a straightforward diagram chase,
working from the various definitions, 
%\mar{ME: Recall this $\widehat{\phantom{\cX}}$ notation from app.\
%  A. TG: done, feel free to comment out if this is OK}
to deduce the versal property of the morphism $\widehat{(\widetilde{\cX})}_{\widetilde{x}}
\to \widehat{\cZ}_z$ (where we use the notation of Definition~\ref{adefn: deformation category}) from the versal property of the morphism 
$\widehat{\cX}_x \to \widehat{\cY}_y.$
\end{proof}

The following result is standard, but we indicate the proof for the sake of
completeness.

\begin{lemma}
\label{lem:versality --- the algebraic case}
If $f: \cX \to \cY$ is a morphism of algebraic stacks, each of finite type
over a locally Noetherian scheme $S$,
then $f$ is versal at finite type points if and only if $f$ is smooth.
\end{lemma}
\begin{proof}
The ``if'' direction follows from an application of the infinitesimal lifting
property of smooth morphisms; see e.g.\ the implication ``$\text{(4)}
\implies  \text{(1)}$''
of
\cite[Lem.~2.4.7~(4)]{EGstacktheoreticimages}.
For the converse, choose a smooth
surjective morphism $U \to \cX$ with $U$ a scheme (necessarily
locally of finite type over~$S$); then by the direction
already proved, this morphism is also versal at finite type points,
and hence so is the composite morphism $U \to \cX \to \cY$.  To show that $f$
is smooth, it suffices to show the same for this composite.  
It follows from the implication ``(1) $\implies$ (4)'' of
\cite[Lem.~2.4.7~(4)]{EGstacktheoreticimages} that $f$ is smooth
in a neighbourhood of each finite type point of $U$; but since these points
are dense in $U$, the desired result follows.
(We have cited \cite{EGstacktheoreticimages} in this argument
purely for our own convenience; of course
the present lemma is just a version of Grothendieck's result relating smoothness
and formal smoothness.)
\end{proof}

We now strengthen the previous result to cover the case where $\cX$ and $\cY$ are
not necessarily assumed to be algebraic, but merely the map between them is
assumed to be representable by algebraic stacks.

\begin{lemma}
\label{lem:versality --- the representable case}
If $f: \cX \to \cY$ is a morphism of limit preserving
stacks over a locally Noetherian base $S$,
and~$f$ is representable by algebraic stacks, 
then $f$ is versal at finite type points if and only if $f$ is smooth.
\end{lemma}
\begin{proof}
%\mar{ME: fill in; note that this is very similar to the lemma that we ultimately
%want to prove!}
If $f$ is smooth, then by Lemma~\ref{lem:characterizing properties}~(2)
it satisfies the infinitesimal lifting property,
%\mar{TG: can we be a bit more specific about what exactly
%  this infinitesimal lifting property is?! Are we literally
%  formulating it on the level of~$\cX$ and~$\cY$? If so then is it
%  clear that it follows from the map being smooth (or are there
%  various base changes to consider?)}
which immediately implies that it is versal at finite type points.
Conversely, suppose now that $f$ is versal at finite type points.

It follows from~\cite[Cor.~2.1.8]{EGstacktheoreticimages} 
that  $f$
is limit preserving on objects,
and since it is also representable by algebraic stacks,
it is locally of finite presentation,
by Lemma~\ref{lem:characterizing properties}~(1).
%~(this is an evident
%variant of
%\cite[\href{https://stacks.math.columbia.edu/tag/06CX}{Tag 06CX}]{stacks-project},
%which treats the case of morphisms representable by algebraic spaces).
We now have to verify that for any test morphism $Z \to \cY$ whose source is a scheme, the induced morphism 
\numequation
\label{eqn:morphism induced by f}
\cX\times_{\cY} Z \to Z ,
\end{equation}
which is again locally of finite presentation,
is in fact smooth. Of course it suffices to do this for affine schemes, and then,
since $\cY$ is limit preserving, for affine schemes that are of finite type
over $S$.  Since $f$ is representable by algebraic stacks,
we find that the source of~\eqref{eqn:morphism induced by f} 
is an algebraic stack.  Furthermore, since this morphism is locally of finite
presentation, its source is again locally of finite type over~$S$. 
Finally, Lemma~\ref{lem:base-changing 
versality} shows that~\eqref{eqn:morphism induced by f} 
is versal at finite type points.
The claimed result then follows from Lemma~\ref{lem:versality --- the algebraic case}.
\end{proof}

We may apply the previous lemma to obtain a criterion 
for representing a stack as the quotient of a sheaf by a smooth groupoid.
We first make another definition.

\begin{defn}
  \label{defn:surjective on f.t. points}
We say that a morphism $\cX\to \cY$ 
of stacks over a locally Noetherian scheme~$S$
%We say that $f$
is {\em surjective
on finite type points} \index{surjective
on finite type points} if for each morphism $x:\Spec k \to \cY$,
with $k$ a finite type $\cO_S$-field, we may find a field extension
$l$ of $k$ and a morphism $\Spec l \to \cX$ which makes the diagram
$$\xymatrix{ \Spec l \ar[r]\ar[d] & \cX \ar[d] \\
\Spec k \ar[r] & \cY}
$$
commutative.
\end{defn}

%\mar{ME: This lemma also needs a surjectivity hypothesis of some sort}
\begin{lemma}
\label{lem:presentations}% \mar{TG: perhaps refer to
  % \cite{EGstacktheoreticimages} for definition of formally smooth?}
Suppose that $\cZ$ is a stack over a locally Noetherian scheme~$S$,
and that $f:U \to \cZ$ is a morphism from a sheaf to $\cZ$ which
is representable by algebraic spaces\footnote{Since the source
of the morphism is a sheaf,
this is equivalent to being representable by algebraic stacks.}
%stack is representable by algebraic spaces if and only if it is representable
%by algebraic stacks (this is easily verified,
%or see e.g.\ \cite[Lem.~3.5]{Emertonformalstacks}), and so this is the
%natural hypothesis to impose in the present context.} 
and surjective on finite type points {\em(}in the sense of Definition~{\em \ref{defn:surjective
on f.t. points})}.
Suppose furthermore that both $U$ and $\cZ$ are limit preserving,
and that~$f$ is 
versal at  finite type points.
Then the morphism $f$ is in fact smooth and surjective,
and, if we let $R := U\times_{\cZ} U$,
then the induced morphism $[U/R] \to \cZ$ is an isomorphism.
\end{lemma}
\begin{proof}
It follows from Lemma~\ref{lem:versality --- the representable case}
that the morphism $f$ is smooth.
We claim that it is furthermore surjective; by definition,
this means that if $T \to \cZ$ is any test morphism from an affine scheme,
we must show that the morphism of algebraic spaces
$g:T\times_{\cZ} U \to T$ induced by $f$ is surjective.
Since $\cZ$ is limit preserving, we may assume that $T$ is of finite type
over~$S$; also, since $g$ is a base-change of the smooth morphism $f$,
it is itself smooth, and so in particular has open image.  Our assumption that
$f$ is surjective on finite type points implies furthermore that the image
of $g$ contains all the finite type points of $T$; thus $g$ is indeed surjective.

Since smooth morphisms are in particular flat, the second assertion
of the lemma follows from Lemma~\ref{lem:stacks as quotients}.
\end{proof}

\begin{remark}
If we omit
the assumption in Lemma~\ref{lem:presentations} that the morphism $U \to \cZ$ be
representable by algebraic spaces, then the morphism $[U/R] \to \cZ$ need 
not be an isomorphism.

An example illustrating this is given by taking $\cZ$ to be a scheme~$Z$,
locally of finite type 
% \mar{ME: Here I am assuming that $S$ is locally Noetherian, 
% so that for schemes locally finite type is equivalent to locally finite presentation,
% open subschemes are then also locally finite type, etc.}
over a locally Noetherian base scheme~$S$, choosing a closed subscheme $Y$
of $Z$ which is {\em not} also open, and defining $U$ to be the formal
scheme obtained as the disjoint union $U := (Z\setminus Y) \coprod \widehat{Z}$,
where $\widehat{Z}$ denotes the completion of $Z$ along~$Y$.
We let $U \to Z$ be the obvious morphism; this is then a surjective monomorphism,
which is not an isomorphism, although it is versal at
every finite type point, and (since it is a monomorphism)
we have that $R := U\times_Z U$ coincides with the diagonal
copy of $U$ inside $U \times_S U$ (so that $[U/R] = U$).
\end{remark}

We will find the following lemmas useful in verifying the hypotheses of
Lemma~\ref{lem:presentations} in our application.
%(more precisely, the hypothesis of representability by algebraic spaces).

\begin{lem}
\label{lem:Artinian test}
Let $f:\cX  \to \cY$ and $g:\cY \to \cZ$ be morphisms of stacks over a locally
Noetherian scheme~$S$. Assume that~$f$ is representable by algebraic stacks
and locally of finite type, and that both $g$ and the morphism $\cZ \to S$
are limit preserving on objects.
Then $f$ is an isomorphism if and only if
%it induces an isomorphism
the induced morphism
\numequation
\label{eqn:A pullback}
\Spec A \times_{\cZ} \cX \to \Spec A \times_{\cZ} \cY
\end{equation}
is an isomorphism
for every morphism $\Spec A \to \cZ$ with 
%on $A$-valued points for every
$A$ a locally of finite type Artinian local $\cO_S$-algebra. %~$A$.
\end{lem}
\begin{proof}
The ``only if'' direction is evident,
and so we focus on the ``if'' direction. We begin by reducing to the
case $\cZ=S$.
We first recall a basic fact about fibre products:
if $T \to \cY$ is any morphism from a scheme to the stack~$\cY$,
then the base-change $T\times_{\cY} \cX$
may be described as an iterated fibre product
$T \times_{(T\times_{\cZ} \cY)} (T\times_{\cZ} \cX)$
(where we regard $T$ as lying over $\cZ$ via the composite
$T \to \cY \to \cZ$, and as lying over $T \times_{\cZ} \cY$
via the graph of the given morphism $T \to \cY$).

In order to show that $f$ is an isomorphism,
we have to show that for any morphism $T \to \cY$,
the base-changed morphism $T\times_{\cY} \cX \to T$
is an isomorphism.  By the fact about fibre products just recalled
(and remembering that the base-change of an isomorphism 
is an isomorphism),
it suffices to show that for any morphism $T \to \cZ$ whose source is
a scheme, the base-changed morphism
\numequation
\label{eqn:base-changed morphism}
T\times_{\cZ} \cX \to T\times_{\cZ} \cY
\end{equation}
is an isomorphism.  Since $\cZ \to S$ is limit preserving on objects, we may and do assume
that $T$ is locally of finite type over~$S$.
We will then replace both $\cZ$ and $S$ by~$T$ (mapping identically to itself),
and the morphism $f$ by~\eqref{eqn:base-changed
morphism}.
Note that the projection
$T\times_{\cZ} \cY \to T$ is again limit preserving on objects (being
a base-change of~$g$, which has this property). 
We now have to reinterpret our original hypothesis on $A$-valued
points of $\cZ$ (for locally of finite type Artinian local $\cO_S$-algebras~$A$)
in this new context.
To this end, note that if $\Spec A\to T$
is locally of finite type with $A$ being Artinian local,
%\numequation
%\label{eqn:test morphism}
% \Spec A \to T \times_{\cZ} \cY
%\end{equation}
% \mar{TG: I'm a bit lost as to what we are doing here; note
%   that this equation isn't actually referenced anywhere\dots ME: Hopefully corrected/clarified.}
%is a morphism with $A$ being Artinian local,
%and with the induced morphism $\Spec A \to T$ being
%locally of finite type,
then $\Spec A$ is also 
locally of finite type over~$S$. 
The fibre product of~\eqref{eqn:base-changed morphism} 
with $\Spec A$ over $T$ may then be identified with
the morphism~\eqref{eqn:A pullback},
%$\Spec A \times_{\cZ} \cX \to \Spec A\times_{\cZ} \cY$ considered in the statement of the lemma,
and is thus an isomorphism (by assumption).
Putting all this together, we see that we have
reduced to the case when $\cZ~=~S$;
so we return to our original notation, but assume in addition 
that $\cZ = S$ (mapping to itself via the identity).
%$\cX \to \cY$ of stacks over a locally Noetherian scheme~$S$,
%which is representable by algebraic stacks and locally of finite type,
%with $\cY \to S$ being limit preserving on objects,
%and with the base-change 
%$\Spec A \times_S \cX \to \Spec A \times_S \cY$ being an isomorphism.

%Now the fibre product fact recalled above (applied to $T = \Spec A$) 
%shows that if $\Spec A \to \cY$ is a morphism,
%with $A$ being an Artinian local finite type $\cO_S$-algebra,
%then $\Spec A \times_{\cY} \cX \to \Spec A$ is an isomorphism.
%
% \mar{TG: $T$ is in danger of being overloaded at this point? ME: Hopefully the
% preceding comment (that I added) about resetting the notation eliminates this danger.}
What we have to show is that if $T \to \cY$ is any morphism,
then
\numequation
\label{eqn:T pullback}
T\times_{\cY} \cX \to T
\end{equation}
is an isomorphism.  Since $\cY \to S$
is limit preserving, we may assume that $T$ is locally of finite type over~$S$.
Any finite type Artinian local $T$-algebra is then also finite type over $S$,
and so, replacing both $\cY$ and $S$ by $T$
(and again applying the fibre product fact recalled above, with $T$ replaced by $\Spec A$) 
we find that we may make another reduction, to the case when $\cY = S$.
The representability assumption on $f$ then implies
that $\cX$ is an algebraic stack.
In this case the lemma is standard, but we recall a proof for completeness.

%Thus, replacing $S$ by $T$,
%and $f$ by its base-change,
%we may assume that we're in the following situation: we have a morphism
%$f: \cX \to \cY$ of stacks over $S$, representable by algebraic stacks
%and locally of finite type, with the morphism $\cY \to S$ being limit
%preserving on objects, 
%Our assumption then implies
%that if $A$ is a finite type Artinian local $\cO_S$-algebra,
%%(and so also a finite type Artinian local $\cO_S$-algebra),
%the base-change of~$f$ %of ~\eqref{eqn:base-changed morphism} 
%over $\Spec A$ is an isomorphism.  


%Thus, replacing both $\cY$ and $S$ by $T$,
% we reduce to the case when $\cY = S$, 
%and $\cX$ is an algebraic stack that is locally of finite type over $S$.
%In this case the lemma is standard, but we recall a proof for completeness.

Our hypothesis that~$f$ induces an isomorphism on $A$-valued points implies that the morphism $f$
is versal at finite type points, and so by Lemma~\ref{lem:versality --- the algebraic case},
it is smooth.  It also implies that $f$ contains all finite
type points in its image; since the image of $f$ is constructible, we see that $f$
is in fact surjective.  This same argument implies that the diagonal
morphism $\cX \to \cX\times_S \cX$ 
is surjective, and thus that $f$ is universally injective.  
Since $f$ is smooth, it is flat and locally of finite type.  Since it 
is also universally injective, it is an open immersion.  Since it is surjective,
it is an isomorphism, as claimed.
\end{proof}

\begin{lem}
  \label{lem: checking closed subsheaf inclusion on versal rings}
Let~$\cY$ be a stack over a locally Noetherian scheme~$S$, 
  let $\cZ$ be a closed substack of~$\cY$,
%Assume that~$\cX$ is limit preserving. % Suppose that~$X''$ is a
  % scheme.
  % \mar{TG: hypotheses to be figured out, proof to be given.}
and let $f:\cX \to \cY$ be a morphism of stacks. 
Suppose that both~ $f$ and the morphism $\cY \to S$ are limit preserving on objects;
then~$f$ factors through~ $\cZ$ if and only 
  %Then $\cX'$ is a substack of~$X''$ % \emph{(}and is in particular a scheme\emph{)}
  if and only if for any finite local Artinian $\cO_S$-algebra~$A$, and any
  morphism $\Spec A\to \cY$ over~$S$, the morphism
$\Spec A\times_{\cY} \cX \to \cY$ factors through $\cZ$. %$\Spec A\times \cZ$.
\end{lem}% \mar{TG: is this statement/proof OK? Do we actually need it
  % to be closed?!}
\begin{proof}
%\mar{ME: Modify proof to fit new statement. TG: made some corrections,
%I think the statement is still garbled though!}
The ``only if'' direction is trivial. For the converse,
consider the base-change via $f$ %closed immersion
%\mar{TG: this isn't a base change of~$f$, is it?! So why is it called~$f$? ME: Typo, fixed}
of the closed immersion $\cZ \hookrightarrow \cY$;
this is a closed immersion
$g: \cX\times_{\cY} \cZ \hookrightarrow \cX$ of stacks over~$\cY$,
which we must show is an isomorphism.
The hypothesis implies that for
any morphism $\Spec A \to \cY$ whose source is a finite type local Artinian 
$\cO_S$-algebra,  the pull-back of~ $g$ \[\Spec
  A\times_\cY(\cX\times_{\cY}\cZ)\to\Spec A\times_{\cY}\cX \] is an isomorphism; % \mar{TG: this tails off; I'm not sure what it should
%   say, to be honest, as I'm a bit unclear whether you want to be base
%   changing the map~$f$ in the statement of the lemma, in which case I
%   don't know what you want to write; or the map here, in which case
%   since the target of that map is~$\cX$, not~$\cY$, I'm confused about
% the base change. ME: Tried to clarify; but maybe I am invoking Lemma~\ref{lem:Artinian test}
% in a stronger form than it is written (since I am using $A$-valued points of a base-stack $\cY$,
% rather than of just the base scheme~$S$), and also, have to think about where the
% limit preserving hypothesis (which is need for Lemma~\ref{lem:Artinian test} should
% actually be imposed in this lemma.}
it then follows from Lemma~\ref{lem:Artinian test}
that $g$ is an isomorphism. %\mar{TG: I guess this is the~$f$ in the
 % proof, rather than in the statement? ME: Yep}
%  set~$\cY=\cX'\times_{\cX} \cX''$; then~$\cY\to \cX'$ is a closed immersion, and we
%  need to show that it is an isomorphism. Since~$\cX$ is limit preserving,
%  and since the morphism $\cX' \hookrightarrow \cX$ is representable
%  by algebraic spaces and of finite type (being a closed
%  immersion), and so also locally of finite presentation (since $S$ is locally
%  Noetherian), we find that $\cX'$ is limit preserving too.
%  The present lemma thus follows from Lemma~\ref{lem:Artinian test}.
%  %it is enough to check this after pulling back by an
%  %arbitrary morphism $\Spec A\to \cX'$ whose source is the spectrum of a
%  %finite type $\cO_S$-algebra. Denoting this pulled-back closed
%  %immersion by $\Spec B\into\Spec A$, it follows from the hypothesis
%  %that the kernel of the surjection $A\onto B$ has vanishing image in
%  %any Artinian quotient of $A$. 
%%
%%  Now, $A$ injects into the product of its localizations at maximal
%%  ideals, and since~$A$ is Noetherian, it injects into the product of
%%  the completions of these localizations. This product can be
%%  identified with the projective limit of the finite length quotients
%%  of~$A$, so~$A$ injects into the product of its local Artinian
%%  quotients.  Thus we find that $A \to B$ is an isomorphism, and we are done.
\end{proof}

The next lemma lets us compute fibre products over stacks in certain situations.
We recall that if $\cY$ is a stack (over some base scheme~$S$),
then the inertia stack $\cI_{\cY}$ is the stack over $S$ whose $T$-valued points
(for any $S$-scheme $T$) consist of pairs $(y,\alpha)$, where $y$ is a $T$-valued
point of $\cY$, and $\alpha$ is an automorphism of~$y$. 
It is a group object in the category of stacks lying over $\cY$ 
via morphisms that are representable by sheaves.
% \mar{ME: Trying to explain that we consider the category of morphisms
% $\cX \to \cY$ which are representable by sheaves (possibly there's
% another, more foundational/categorical way to express this condition). I'm
% a bit worried that without this condition (i.e.\ if one
% just considers the category of aribtrary morphisms of stacks
% $\cX \to \cY$), one lands in a painful $2$-category
% context.}
%(and so can also be thought of as a groupoid over $\cY$
%in which the source and target morphisms coincice).
We also remind the
reader that there is a canonical isomorphism
$$\cI_{\cY} \iso \cY\times_{\Delta_{\cY},\cY\times_S \cY, \Delta_{\cY}} \cY,$$
where the subscripts $\Delta_{\cY}$ indicate that both copies of $\cY$ are regarded 
as lying over the product $\cY\times_S \cY$ via the diagonal morphism $\Delta_{\cY}:
\cY \to\cY\times_S \cY$, and where, under this identification,
 the forgetful morphism $\cI_{\cY} \to \cY$ may be identified with projection onto
the first copy of~$\cY$.

\begin{lemma}
\label{lem:fibre product and inertia}
Let $X \to \cY$ be a morphism from a sheaf to a stack {\em(}both over some base-scheme~$S${\em )}, and suppose 
that the natural morphism $X\times_{\cY} X \to X\times_S X$ factors through the diagonal
copy of~$X$ lying in the target {\em (}so that the groupoid $X\times_{\cY} X$ is in fact
a group object in the category of sheaves over~$X${\em )}.  
Then there is a canonical isomorphism
$X\times_{\cY} X \iso X\times_{\cY} \cI_{\cY}$
of group objects over~$X$.
\end{lemma}
\begin{proof}
By assumption, we have a commutative diagram
$$\xymatrix{X\times_{\cY} X \ar[d] \ar[r] & X \ar^-{\Delta_X}[r]\ar[dr] & X\times_S X \ar[d] \\
\cY \ar^-{\Delta_{\cY}}[rr] & & \cY\times_S \cY}
$$
in which the outer rectangle is $2$-Cartesian.  
Since $X$ is a sheaf, the diagonal $\Delta_X: X \to X \times_S X$ is a monomorphism,
and so one immediately checks that the left-hand trapezoid is  also 
$2$-Cartesian.  Thus we obtain the required isomorphism
\[X\times_{\cY} X \iso X \times_{\cY\times_S \cY} \cY \iso X \times_{\cY}
(\cY\times_{\cY\times_S\cY} \cY) \iso X\times_{\cY} \cI_{\cY}. \qedhere\]
\end{proof}

\section{Moduli stacks in the rank one case}
We now give concrete descriptions of our various stacks in the rank one case.
Before doing this, we prove a result that describes the stack structure
on families of rank one objects.  (It encodes the fact that the
automorphisms of a rank one object are simply the scalars.)

%\mar{ME: This lemma is nonsense as stated; trying to figure out correct formulation}
\begin{lemma}
\label{lem:rank one automorphisms}
If $\cZ$ denotes either of the stacks $\cR_1$ or $\cX_1$,
then there is a canonical isomorphism
\numequation
\label{eqn:inertia stack of R1}
\cZ \times_{\cO} \widehat{\mathbb G}_m \iso \cI_{\cZ}
\end{equation}
of group objects over~$\cZ$  (where $\widehat{\G}_m$ is the $\varpi$-adic
completion of $\Gm$ over~$\cO$).
\end{lemma}
\begin{proof}
For definiteness we give the proof in the case of $\cR_1$; the proof in the $\cX_1$
case is identical.
We begin by defining the morphism~\eqref{eqn:inertia stack of R1}:
if $A$ is any $\varpi$-adically complete $\cO$-algebra,
and $M$ is an \'etale $\varphi$-module over $\A_{K,A},$
then~\eqref{eqn:inertia stack of R1} is defined on $A$-valued points
of the source lying over~$M$ by mapping an element $a \in A^{\times}$
to the automorphism of $M$ given by multiplication by~$a$.

Since $\cR_1$ has affine diagonal which is of finite presentation,
the morphism~\eqref{eqn:inertia stack of R1}
is a morphism between finite type group objects over $\cR_1$,
each of whose structure morphisms is representable by algebraic spaces,
indeed affine,
and of finite presentation.    Furthermore, $\cR_1$ itself is limit
preserving over~$\cO$. Thus 
to show that~\eqref{eqn:inertia stack of R1} is an isomorphism,
it suffices (by Lemma~\ref{lem:Artinian test})
to show that it induces an isomorphism on $A$-valued points
for $A$ an Artinian local $\cO$-algebra of finite type.
Returning to the notation of the preceding paragraph (but now assuming
that $A$ is Artinian local),
we have to show that any automorphism of $M$ is 
given by multiplication by an element of~$A^{\times}$.
Since $M$ is of rank~$1$,
we find that
$$\Hom_{\A_{K,A},\varphi}(M,M) \cong (M^{\vee}\otimes_{\A_{K,A}}M)^{\varphi=1} \cong (\A_{K,A})^{\varphi=1}
\cong A$$
(the final isomorphism following from~Lemma~\ref{lem:phi invariants}, since $\A_{K,A} \subseteq
W(\C^{\flat})_A$ by Proposition~\ref{prop: maps of coefficient rings
  are faithfully flat injections}). 
% \mar{ME: Maybe connect this citation to the required statement more clearly,
% by observing that $\A_{K,A} \subseteq W(\C^{\flat})_A.$ TG: done,
% commenting out.}
The lemma follows.
%\mar{ME: Have to give proof}
\end{proof}


%\subsection{Local class field theory}
\subsection{Local Galois theory}
We briefly recall the local Galois theory that is relevant to our
computation. If~$L/K$ is an algebraic extension, we write~$I_L^\ab$
for the image of~$I_L$ in~$G_L^{\ab}$ (note that this is not the
abelianization of~$I_L$).
We have short exact sequences
\numequation
\label{eqn:K ramification s.e.s.}
1 \to I_K \to G_K \to \Frob_K^{\Zhat} \to 1 
\end{equation}
and
\numequation
\label{eqn:K-cyc ramification s.e.s}
1 \to I_{K_{\cyc}} \to G_{K_{\cyc}} \to \Frob_{K}^{f \Zhat} \to 1,
\end{equation}
where $f := [k_{\infty}: k].$
%Since $I_{K_{\cyc}}$ is equal to the absolute Galois group
%of the compositum of $K_{\cyc}$ with the maximal unramified extension
%of $K$, we see that $G_K/I_{K_{\cyc}}$ is abelian. 

These induce corresponding short exact sequences
\numequation
\label{eqn:abelian K ramification s.e.s.}
1 \to I_K^{\ab} \to G_K^{\ab} \to \Frob_K^{\Zhat} \to 1 
\end{equation}
and
\numequation
\label{eqn:abelian K-cyc ramification s.e.s}
1 \to I_{\Kcyc}^{\ab} \to G_{K_{\cyc}}^{\ab} \to \Frob_K^{f \Zhat} \to 1.
\end{equation}
% where $I_K^\ab$ (resp.\ $I^\ab_{\Kcyc}$) denotes the image of $I_K$
% (resp.\ of $I_{K_{\cyc}}$).
%and of course local class field theory yields isomorphisms\mar{TG: do
%  we need to make some slight adjustments if $p=2$?}
%$I \cong (1+\mathfrak m_K) \times k^{\times}$
%and $I_{\cyc} \cong (1 + \mathfrak m_{K_{\cyc}} ) \times k_{\infty}^{\times}.$
%
% There is also an induced short exact sequence\mar{TG: presumably this
%   will be either labeled and referenced or removed.}
% $$1 \to G_{K_{\cyc}}^{\ab}/[G_{K_{\cyc}}^{\ab},\Gamma] 
% \to G_K^{\ab} \to \Gamma \to 1.$$

If $G$ is any of the various profinite Galois groups appearing 
in the preceding discussion, we let $\cO[[G]]$ denote the corresponding
completed group ring over~$\cO$.  This is a pro-Artinian ring, which
gives rise to the affine formal algebraic space $\Spf \cO[[G]]$.  We endow
this space with the trivial action of $\widehat{\G}_m$ (the $\varpi$-adic
completion of $\Gm$ over~$\cO$).

\subsection{Galois lifting rings for characters}If $\F'/\F$ is a
finite extension, and $\rhobar:G_{\Kcyc}\to(\F')^\times$ is a
continuous character, then we can write
~$\cO'=\cO\otimes_{W(\F)}W(\F')$, and consider the universal lifting
$\cO$-algebra $R_{\rhobar}^{\square,\cO'}$. If $x:\Spf\F'\to\cR_1$ is
the corresponding finite type point, then we have a versal morphism
$\Spf R_{\rhobar}^{\square,\cO'}\to \cR_1$; indeed it follows exactly
as in the proof of Proposition~\ref{prop:versal rings} that we have an
isomorphism \numequation\label{eqn: deformation ring character
  torsor}\Spf R_{\rhobar}^{\square,\cO'} \times_{\cR_1} \Spf
R_{\rhobar}^{\square,\cO'} \iso
(\Gmhat)_{R_{\rhobar}^{\square,\cO'}},\end{equation} where
$(\Gmhat)_{R_{\rhobar}^{\square,\cO'}}$ denotes the completion of
$(\Gm)_{R_{\rhobar}^{\square,\cO'}}$ along $(\Gm)_{\F'}$.

\subsection{Descriptions of the rank one stacks}
We can now establish our explicit description of~$\cR_1$.

%We will find
%the following definition useful.
%
%\begin{defn}\label{defn: Kbasic T height at most h}
%  Let~$A$ be a finite local Artinian $\cO/\varpi^a$-algebra for
%  some~$a\ge 1$, and let~$M$ be a rank one \'etale $\varphi$-module
%  over~$\A_{K,A}$. As in Section~\ref{subsec: canonical weak Wach}, we
%  may regard~$M$ as a rank~$[K:\Kbasic]$ \'etale $\varphi$-module~$M'$
%  over~$\A_{\Kbasic,A}$, and we say that~$M$ \emph{has
%    $(\Kbasic,T)$-height at most~$h$} if there is a
%  $\varphi$-module~$\gM$ over~$\A_{\Kbasic,A}^+$ of $T$-height at
%  most~$h$ such that $M'=\gM[1/T]$.
%\end{defn}

\begin{prop}
	\label{prop:rank 1 R}
       	There is an isomorphism
	$$
\Bigl[ \Bigl ( \Spf \cO[[I^{\ab}_{\Kcyc}]] \times \widehat{\mathbb G}_m
	\Bigr) / \widehat{\mathbb G}_m \Bigr]
\iso
	\cR_{1}$$
{\em (}where, in the formation of the quotient stack, the $\widehat{\mathbb
G}_m$-action is taken to be trivial{\em )}.
\end{prop}
\begin{proof}
We begin by constructing a morphism
	\numequation\label{eqn: covering R1}\Spf \cO[[I^{\ab}_{\Kcyc}]] \times \widehat{\mathbb G}_m
	\to \cR_1.\end{equation} For this, we choose a lift of
        $\sigma_{\cyc} \in G_{K_{\cyc}}$ of $\Frob_K^f$, splitting the
        short exact sequences~\eqref{eqn:K-cyc ramification s.e.s}
        and~\eqref{eqn:abelian K-cyc ramification s.e.s}.  If $A$ is
        any discrete Artinian quotient of $\cO[[I^{\ab}_{\Kcyc}]],$ then we
        extend the continuous morphism $\cO[[I^{\ab}_{\Kcyc}]] \to A$ to a
        continuous morphism $\cO[[G_{K_{\cyc}}^{\ab}]]\to A$ by
        mapping $\sigma_{\cyc}$ to $1 \in A$.  We may view this latter
        morphism as a rank~$1$ representation of $G_{K_{\cyc}}$ with
        coefficients in~$A$; it thus gives rise to a rank one
        projective \'etale $\varphi$-module $M_A$ over $\A_{K,A}$, and
        therefore to a morphism $\Spf \cO[[I^{\ab}_{\Kcyc}]] \to \cR_1.$ We
        extend this to the morphism~\eqref{eqn: covering R1} by
        unramified twisting; more precisely, we have an induced
        morphism $\Spf \cO[[I^{\ab}_{\Kcyc}]]\times\Gmhat \to
        \cR_1\times\Gmhat$, and we compose with the morphism
        $\cR_1\times\Gmhat\to\cR_1$ given by taking the tensor product
        with the universal unramified rank one \'etale $\varphi$-module.%  where~$\varphi$
        % acts by multiplication a morphism $\Gmhat\to
        % A^\times$ corresponds to an element $a\in A^\times$, and we
        % have the corresponding unramified character sending $\Frob_K$ to~$a$.



%\mar{ME: Use formal deformations of the trivial char.\ plus unramified twists}


By construction (and the usual explicit description of the universal
Galois deformation ring of a character), % \mar{TG: we might want to
  % develop something here, as it's over~$\Kcyc$}
~\eqref{eqn: covering R1} is versal at finite type
points, as well as surjective on finite type points.   % \mar{ME: This is
% because it is constructed via deformation theory!} 
%We claim that the morphism
We claim that the canonical morphism
\nummultline
\label{eqn:morphism of diagonals}
\bigl( \Spf \cO[[I^{\ab}_{\Kcyc}]] \times \widehat{\mathbb G}_m\bigr)
\times_{\cR_1} 
\bigl( \Spf \cO[[I^{\ab}_{\Kcyc}]] \times \widehat{\mathbb G}_m\bigr)
\\
\to
\bigl( \Spf \cO[[I^{\ab}_{\Kcyc}]] \times \widehat{\mathbb G}_m\bigr)
\times_{\cO} 
\bigl( \Spf \cO[[I^{\ab}_{\Kcyc}]] \times \widehat{\mathbb G}_m\bigr)
\end{multline}
factors through the diagonal copy of
$ \Spf \cO[[I^{\ab}_{\Kcyc}]] \times \widehat{\mathbb G}_m$
in the target; given this, it follows from Lemmas~\ref{lem:fibre product and inertia}
and~\ref{lem:rank one automorphisms} that there is an isomorphism of groupoids
\[\bigl ( \Spf \cO[[I^{\ab}_{\Kcyc}]] 
	\times \widehat{\mathbb G}_m
	\bigr) \times \Gmhat\isoto\bigl ( \Spf \cO[[I^{\ab}_{\Kcyc}]] 
	\times \widehat{\mathbb G}_m
	\bigr)\times_{\cR_1}\bigl ( \Spf \cO[[I^{\ab}_{\Kcyc}]] 
	\times \widehat{\mathbb G}_m
	\bigr)\]
with $\widehat{\mathbb G}_m$ acting trivially on $\Spf \cO[[I^{\ab}_{\Kcyc}]]\times
\widehat{\mathbb G}_m.$
%given by the composite of projection to $\Bigl ( \Spf \cO[[I^{\ab}_{\Kcyc}]] 
%	\times \widehat{\mathbb G}_m
%	\Bigr) $ and the diagonal morphism is an isomorphism.  
%\mar{ME: Together with a computation of $U
%        \times_{\cR}  U \dots $ TG: I guess we should write a map and
%        reduce to checking on Artinian points, using Lemma~\ref{lem: checking closed subsheaf inclusion on versal
%    rings}, and we should be able to do that with~\eqref{eqn: deformation ring character torsor}.}
The proposition will then follow from Lemma~\ref{lem:presentations}
provided we show that ~\eqref{eqn: covering R1} is representable by algebraic
spaces (or, equivalently, by algebraic stacks).

%\mar{ME to add argument here}

It remains to verify the claimed factorization of~\eqref{eqn:morphism of diagonals},
as well as the claimed representability by algebraic stacks of~\eqref{eqn: covering R1}.
We establish each of these claims in turn.

%\mar{ME: Have to verify factorization claim}
To show that~\eqref{eqn:morphism of diagonals} factors through
the diagonal diagonal copy of
$ \Spf \cO[[I^{\ab}_{\Kcyc}]] \times \widehat{\mathbb G}_m$
in the target, it suffices,
by Lemma~\ref{lem: checking closed subsheaf inclusion on versal rings},
to show that if
$$
\Spec A \to 
 \bigl(\Spf \cO[[I^{\ab}_{\Kcyc}]] \times \widehat{\mathbb G}_m\bigr)
\times
 \bigl(\Spf \cO[[I^{\ab}_{\Kcyc}]] \times \widehat{\mathbb G}_m\bigr)
$$
is a morphism whose source is a finite type Artinian local $\cO$-algebra,
then the pull back of \eqref{eqn:morphism of diagonals}
to $\Spec A$ factors through the pull-back to $\Spec A$ of the diagonal morphism
$$
 \Spf \cO[[I^{\ab}_{\Kcyc}]] \times \widehat{\mathbb G}_m
\to
\bigl( \Spf \cO[[I^{\ab}_{\Kcyc}]] \times \widehat{\mathbb G}_m\bigr)
\times_{\cO} 
\bigl( \Spf \cO[[I^{\ab}_{\Kcyc}]] \times \widehat{\mathbb G}_m\bigr)
.$$  Since morphisms $\Spec A \to
 \Spf \cO[[I^{\ab}_{\Kcyc}]] \times \widehat{\mathbb G}_m$
%and morphisms $\Spec A \to \cR_1$ both classify the same data, namely
correspond to characters $\chi: G_{\Kcyc}^{\ab} \to A^{\times}$,
this amounts to the evident fact that if $\chi,\chi'$ are two such characters,
then the locus in $\Spec A$ over which $\chi$ and $\chi'$ become 
isomorphic coincides with the locus over which they coincide.

We now turn to proving that~\eqref{eqn: covering R1} is representable by algebraic
stacks.
To see this, it suffices to study the corresponding question modulo
some power $\varpi^a$ of the uniformizer in~$\cO$, so we work with the
stack~$\cR_{K,1}^a$ from now on.%  In fact, we can and do suppose
% that~$a$ is a multiple of~$e$ (we do this in order to be able to cite
% Proposition~\ref{prop: ramification bound for non basic abelian case
%   any a} below). \mar{TG: maybe this will go away?}

% \mar{TG: I'm inclined to pull some of this material (basically the
%   stuff not mentioning the stacks), and the definition below,
% out before this proposition.}
We next recall some of the constructions we made in
Section~\ref{subsec: defn of Xd}. We have the subfield
$\Kbasic\subseteq K$ of Definition~\ref{defn: Kbasic}, and the natural
morphism $\cR_{K,1}^a\to\cR_{\Kbasic,[K:\Kbasic]}^a$ given by the
forgetful map which regards a rank one $\A_{K,A}$-module as an
$\A_{\Kbasic,A}$-module of rank~$[K:\Kbasic]$. By Lemma~\ref{lem: explicit Ind
  description for R pulled back from K basic}, we can write
$\cR_{K,1}^a=\varinjlim_h\cR_{K,1,\Kbasic,h}^a$, where the algebraic
stack $\cR_{K,1,\Kbasic,h}^a$ is the scheme-theoretic image of the
base-changed morphism
  \[\cC_{[K:\Kbasic],h}^a\times_{\cR_{\Kbasic,[K:\Kbasic]}}\cR_{K,1}\to\cR_{K,1},\]
  where $\cC_{[K:\Kbasic],h}^a$ is 
% \mar{TG: this notation is perhaps
%   now also established in Section~\ref{subsec:
%   canonical weak Wach}.} By~\cite[Thm.\
% 5.4.20]{EGstacktheoreticimages} (see also Theorem~\ref{thm: C is a $p$-adic formal algebraic stack and related
%     properties}), we can write $\cR_{\Kbasic,[K:\Kbasic]}^a=\varinjlim_{h\ge
%   1}
%   \cR_{\Kbasic,[K:\Kbasic],h}^a$, where $\cR_{\Kbasic,[K:\Kbasic],h}^a$ is an
%   algebraic stack of finite presentation over~$\Spec
%   \cO/\varpi^a$. More precisely, $\cR_{\Kbasic,[K:\Kbasic],h}^a$ is the
%   scheme-theoretic image of the morphism \[\cC^a_{\Kbasic,[K:\Kbasic],h}\to
%     \cR_{\Kbasic,[K:\Kbasic]}^a,\] whose source is
  the algebraic stack of
  $\varphi$-modules over~$\A^+_{\Kbasic,A}$ of rank~$[K:\Kbasic]$ and~$T$-height at
  most~$h$. % By Lemma~\ref{lem: relative representability for R over
  %   Kbasic}, the base
  % change \[\cR_{K,1,h}^a:=\cR_{K,1}^a\times_{\cR_{\Kbasic,[K:\Kbasic]}^a}\cR_{\Kbasic,[K:\Kbasic],h}^a \]is an
  % algebraic stack of finite presentation over~$\Spec
  % \cO/\varpi^a$, and we have $\cR_{K,1}^a=\varinjlim_{h\ge
  % 1}
  % \cR_{K,1,h}^a$.
  Accordingly, it suffices to show that  
 each base-change %\mar{TG: just making up some notation for this for now.}
	$$\cYha:=\Bigl( \Spf  \cO[[I^{\ab}_{\Kcyc}]] \times \widehat{\mathbb G}_m\Bigr )
\times_{\cR_{K,1}} \cR_{K,1,\Kbasic,h}^a$$ is an algebraic stack.% \mar{TG: have to decide whether we are
  % doing class field theory or Galois theory description.}

By construction, $\cYha$ is a closed subsheaf of $\Spf  (\cO/\varpi^a)[[I^{\ab}_{\Kcyc}]]
\times (\Gm)_{\cO/\varpi^a}$. We claim that it is in fact a closed
subsheaf of
%$\Spec  (\cO/\varpi^a)[I^{\ab}_{\Kcyc}/I^{K,h,a/e}_{\Kcyc}] \times
$\Spec  (\cO/\varpi^a)[I^{\ab}_{\Kcyc}/U] \times
(\Gm)_{\cO/\varpi^a}$,
for some open subgroup $U$ of $I^{\ab}_{\Kcyc}$.
%where~$I^{K,h,a/e}_{\Kcyc}$ is as in
%Proposition~\ref{prop: ramification bound in rank one case with coefficients
%    and any K} below; in particular, this implies that~$\cYha$ is a scheme, which
%will complete the proof.\mar{TG: sort this out now that ramification
%  bound completely changed - hopefully not too hard to get an open
%  subgroup of~$I^\ab_{\Kcyc}$ from having one for~$I^\ab_{\Kcyc,s}$?}


% \mar{ME: Finish proof.  Need to show that any character over a finite Artinian
% $\cO$-algebra $A$ which is height $\leq h$ has bounded
% ramification. TG: not yet clear to me how we reduce to these, but
% didn't think much about it - going to assume it for now. TG: also not
% clear to me whether we just need to bound ramification or whether we
% also need to use the canonical extension to some finite extension
% of~$\Qp$, i.e.\ I don't know if there is a discreteness issue (I would
% guess that there is? But maybe the idea is just to prove that some
% big~$p$-power acts trivially?)}

Since~$\cY_h^a$ and~$\Spec(\cO/\varpi^a)[I^{\ab}_{\Kcyc}/U] \times
(\Gm)_{\cO/\varpi^a}$ are both
closed subsheaves of $\Spf (\cO/\varpi^a)[[I^{\ab}_{\Kcyc}]] \times
(\Gm)_{\cO/\varpi^a}$, it follows from Lemma~\ref{lem: checking closed subsheaf inclusion on versal
  rings} that % \mar{TG: we need to check the hypotheses on this once it's
% written}
% \mar{TG: is it clear that this is what we want to check in order to
%   apply that lemma? It looks like instead we want to take $\Spec A\to\Spf (\cO/\varpi^a)[[I^{\ab}_{\Kcyc}]] \times
% (\Gm)_{\cO/\varpi^a}$, and pullback the morphism $\cYha\to\Spf (\cO/\varpi^a)[[I^{\ab}_{\Kcyc}]] \times
% (\Gm)_{\cO/\varpi^a}$ along this, and then assert that it factors? ME: But doesn't
% this just replaced $\Spec A$ by a closed subscheme $\Spec A'$ which now factors
% through $\cYha$?  So essentially I've just started with this $\Spec A'$ situation,
% and written $A$ rather than $A'$.  Feel free to integrate this explanation into
% the argument if you'd like to. TG: commenting out }
it is enough to show that if~$A$ is a finite local Artinian
$\cO/\varpi^a$-algebra, then given any morphism $\Spec A\to \cYha$,
the composite
$\Spec A\to\cYha\to\Spf (\cO/\varpi^a)[[I^{\ab}_{\Kcyc}]] \times
(\Gm)_{\cO/\varpi^a}$ factors through
%$\Spec (\cO/\varpi^a)[I^{\ab}_{\Kcyc}/I^{K,h,a/e}_{\Kcyc}] \times
$\Spec (\cO/\varpi^a)[I^{\ab}_{\Kcyc}/U] \times
(\Gm)_{\cO/\varpi^a}$ % Since~$A$ is a product of local rings, we can
% immediately reduce to the case that~$A$ itself is also local.
for the open subgroup $U:=I^{a,h}$ % obtained as the preimage of
% the group $I^{K,a,h,N,s}$ of \mar{TG: fix this citation.}
of Proposition~\ref{prop: ramification bound in rank one case with
  coefficients and any K} below.
% under the continuous morphism
% $I_{K_{\cyc}}^{\ab} \to
% I_{K_{\cyc,s}}^{\ab}$. 
%that is independent of~$A$ (i.e.\ depends only
%on $a$ and~$h$).
%\mar{ME: Have to figure out which $U$ to take! ME: Should just be the inverse image
%of some open set in $I_{K_{\cyc,s}}^{\ab}$ along the (continuous!)
%morphism $I_{K_{\cyc}}^{\ab} \to
%I_{K_{\cyc,s}}^{\ab}$.}
If~$A$
is in fact a field, then by~\cite[Lem.\
3.2.14]{EGstacktheoreticimages}, the \'etale $\varphi$-module~$M_A$
corresponding to the morphism $\Spec A\to \cR_{K,1}$ attains
$(\Kbasic,T)$-height at most~$h$
in the sense of Definition~\ref{defn: Kbasic T height at most h} below
after passing to a finite extension of~$A$,
and the required factorisation is immediate from
Proposition~\ref{prop: ramification bound in rank one case with coefficients
    and any K}. % \mar{TG: this is the first mention of
%   $(\Kbasic,T)$-height, should decide whether to use it more generally
% or not.}
% \mar{ME: Field of definition issues to sort out, I think. TG: I
%   haven't been able to understand this comment! TG: OK, so the point
%   is that there is only this structure after replacing $A$ by a finite
% extension (of finite fields); can we just point this out, and point
% out that this doesn't affect the existence of the claimed
% factorisation, and be done?} 

Having established the result in the case that $A$ is a finite type
field, in order to prove the general local Artinian case, it suffices to prove a factorisation on the level of versal rings. % , which we
% will do with the aid of Lemmas~\ref{alem: scheme theoretic image of versal is versal} and~\ref{lem: criterion for Artin to map to scheme theoretic
%     image}.  
More precisely, let $\F'/\F$ be a finite extension, and let
$x:\Spec\F'\to\cYha$ be a finite type point of~$\cYha$, corresponding
to a representation $\rhobar:G_{\Kcyc}\to(\F')^\times$. As above, we
write~$R_{\rhobar}^{\square,\cO'}$ for the corresponding universal
lifting $\cO'$-algebra. Write $\Spf R^{a}_h$ for the scheme-theoretic
image of the morphism
\[
  \cC_{[K:\Kbasic],h}^a\times_{\cR_{\Kbasic,[K:\Kbasic]}}\cR_{K,1}\times_{\cR_{K,1}}\Spf
  R_{\rhobar}^{\square,\cO'}\to\Spf R_{\rhobar}^{\square,\cO'}.\] By
Lemma~\ref{alem: scheme theoretic image of versal is versal}, the
induced morphism $\Spf R^a_h\to \cR_{K,1,\Kbasic,h}^a$ is a versal
morphism at~$x$. Write~$R_{\rhobar}^{\square,\cO',U}$ for the quotient
of~$R_{\rhobar}^{\square,\cO'}$ corresponding to liftings which are
trivial on~$U$; %~$I^{K,h,a/e}_{\Kcyc}$;
then it suffices to show that $\Spf
R^a_h$ is a closed formal subscheme of ~$\Spf
R_{\rhobar}^{\square,\cO',U}$.

We now employ Lemma~\ref{lem: criterion for Artin to map to scheme
  theoretic image}. Exactly as in the proof of Lemma~\ref{lem:Artinian points of scheme theoretic images}, it
is enough (after possibly increasing~$\F'$) to show that if~$A$ is a
finite type local Artinian $R_{\rhobar}^{\square,\cO'}$-algebra with
residue field~$\F'$ for which the induced morphism
$\cC_{[K:\Kbasic],h}^a\times_{\cR_{\Kbasic,[K:\Kbasic]}}\cR_{K,1}\times_{\cR_{K,1}}\Spec
A\to \Spec A$ admits a section, then the morphism $\Spec A\to \Spf
R_{\rhobar}^{\square,\cO'}$ factors through~$\Spf
R_{\rhobar}^{\square,\cO',U}$. Since the existence of the section
implies (by definition) that the \'etale $\varphi$-module $M_A$ has
$(T,\Kbasic)$-height at most~$h$
(in the sense of Definition~\ref{defn: Kbasic T height at most h}),
we are done by Proposition~\ref{prop: ramification bound in rank one case with coefficients
    and any K}.
\end{proof}%\mar{TG: should have everything we need now?}

In the previous argument, we used the following definition,
which will play an important role in the technicalities
of Section~\ref{subsec:ram bound} below.
(See in particular the statements of
Propositions~\ref{prop: ramification bound in rank one case with coefficients and any K} and~\ref{prop:
  ramification bound in rank one case with coefficients
    and any K with GK action}.)
%\mar{ME: Note --- this definition is now duplicated in \S 3.5}
\begin{defn}\label{defn: Kbasic T height at most h}
  Let~$A$ be a finite local Artinian $\cO/\varpi^a$-algebra for
  some~$a\ge 1$, and let~$M$ be a rank one projective \'etale $\varphi$-module
  over~$\A_{K,A}$. % As in Section~\ref{subsec: canonical weak Wach}, we
We regard~$M$ as a rank~$[K:\Kbasic]$ \'etale $\varphi$-module~$M'$
  over~$\A_{\Kbasic,A}$, and we say that~$M$ \emph{has
    $(\Kbasic,T)$-height at most~$h$} if there is a projective
  $\varphi$-module~$\gM$ over~$\A_{\Kbasic,A}^+$ of $T$-height at
  most~$h$ such that $M'=\gM[1/T]$.
\end{defn}


%\mar{TG: I should think through this next paragraph.}
We next turn to describing $\cR^{\Gamma_\disc}_1$ explicitly.
We first note that
since~$I_K\cap G_{\Kcyc}=I_{\Kcyc}$, we have an exact sequence
$1 \to I_{K_{\cyc}} \to G_K \to \Frob_K^{\Zhat} \times \Gamma.$
We let $H \subseteq
\Frob_K^{\Zhat} \times \Gamma$
denote the image of $G_K$;
it is an open subgroup of 
$\Frob_K^{\Zhat} \times \Gamma$.
We write $H_{\disc} := H \cap (\Frob_K^{\Z} \times \Gamma_{\disc});$
then $H_{\disc}$ is dense in $H$, and is isomorphic to $\Z\times \Z$.
Write $\cO[H_{\disc}]$ for the group ring of $H_{\disc}$ over $\cO$,
and $\widehat{\cO[H_{\disc}]}$ for its $\varpi$-adic completion.
Then $\Spec \cO[H_{\disc}] \cong \Gm\times_\cO \Gm,$
and $\Spf \widehat{\cO[H_{\disc}]} \cong \widehat{\mathbb G}_m \times_{\cO}
\widehat{\mathbb G}_m.$  Let $I'$ denote the image of $I_{\Kcyc}$ in $G_K^{\ab}$;
equivalently, $I'$ is the quotient of $I^{\ab}_{\Kcyc}$ by
the closure of its subgroup of commutators $[ I^{\ab}_{\Kcyc},\Gamma]$.


\begin{prop}
	\label{prop:rank 1 R-Gamma-disc}
       	There is an isomorphism
	$$
\Bigl[ \Bigl ( \Spf \cO[[I']] \times \Spf \widehat{\cO[H_{\disc}]}
	\Bigr) / \widehat{\mathbb G}_m \Bigr]
\iso
	\cR_1^{\Gamma_{\disc}} $$
%	\mar{ME: Have to introduce more notation to describe this I think}
{\em (}where, in the formation of the quotient stack, the $\widehat{\mathbb
G}_m$-action is taken to be trivial{\em )}.
\end{prop}
\begin{proof}
We leave the construction of this isomorphism to the reader,
by combining the isomorphism of Proposition~\ref{prop:rank 1 R}
with the definition of $\cR_1^{\Gamma_{\disc}}$.
\end{proof}

\begin{remark}
	Note that the case $d=1$ is not representative of the general
	case in so far as the structure of $\cR$, and so also of
	$\cR^{\Gamma_{\disc}}$, is concerned.  Namely, when $d = 1$,
	the stacks $\cR$ and $\cR^{\Gamma_{\disc}}$ are formal
	algebraic stacks.  This will not be the case when $d > 1$.
	(They will instead be Ind-algebraic stacks whose underlying
	reduced substacks are Ind-algebraic but not algebraic.)
	The reason for this is that the $1$-dimensional mod $p$
	representations of $G_{K_{\cyc}}$ may be described
	as Frobenius twists of a finite number of characters
	when $d = 1$,
	so that in this case $\cR_{\red}$ is indeed an algebraic stack,
	while the spaces $\Ext^1_{G_{K_{\cyc}}}(
	\chi,\psi)$ are infinite dimensional, for any
	two mod $p$ characters $\chi$ and $\psi$ of $G_{K_{\cyc}}$,
	so that already in the case $d = 2$,
	the stack $\cR_{\red}$ is merely Ind-algebraic.
\end{remark}

We now give an explicit description of the stack~$\cX_1$. % The argument
% is very similar to that used to establish Proposition~\ref{prop:rank 1
%   R}. 
% ; we could in fact give a proof that is even closer to that of Proposition~\ref{prop:rank 1
%   R}, but we have found it more instructive to make use of the
% material on canonical actions on weak Wach modules that we developed in Section~\ref{subsec:
%   canonical weak Wach}.\mar{TG: will have to relate those to the
%   canonical actions on dual Galois representations in the next
%   section. TG: probably I'm now going to use this above, too, so these
% comments can be adjusted.}
\begin{prop}
	\label{prop:rank 1 X}
       	There is an isomorphism
	$$ \Bigl[ \Bigl ( \Spf \bigl( \cO[[I_K^{\ab}]]\bigr) \times \widehat{\mathbb G}_m
	\Bigr) / \widehat{\mathbb G}_m \Bigr] \iso \cX_1$$%\mar{TG: proof needed.}
{\em (}where, in the formation of the quotient stack, the $\widehat{\mathbb
G}_m$-action is taken to be trivial{\em )}.
\end{prop}%\mar{TG: $k_{\infty}^{\times}$ should be $k^\times$?}
\begin{proof}
%\mar{TG: does this proof have enough detail? ME: Looks fine to me, at least to be getting on with.}  
We prove this in the same way as
  Proposition~\ref{prop:rank 1 R}; we leave the details to the reader,
  indicating only the key differences. We can construct a morphism
  $\Spf \bigl( \cO[[I_K^{\ab}]]\bigr) \times \widehat{\mathbb G}_m\to\cX_1$ in
  exactly the same way as in the proof of Proposition~\ref{prop:rank 1
    R} (by choosing a lift $\sigma_K \in G_{K}$ of $\Frob_K$,
  splitting the short exact sequence~\eqref{eqn:K ramification s.e.s.}),
  and we need to prove that this morphism is representable by
  algebraic spaces. This can be done by arguing exactly as in the
  proof of Proposition~\ref{prop:rank 1 R}, with the isomorphism
  $\cX_{K,1}^a=\varinjlim_{h,s}\cX_{K,1,\Kbasic,h,s}^a$ of
  Lemma~\ref{lem: explicit Ind description for X pulled back from K
    basic}  replacing the isomorphism
  $\cR_{K,1}^a=\varinjlim_h\cR_{K,1,\Kbasic,h}^a$. 

  Bearing in mind the definition of~$\cX_{K,1,\Kbasic,h,s}^a$, we find that 
  we have to show that
  there is an open subgroup~$I_K^{h,s,a}$ of~$I_K^{\ab}$ such that if~$A$ is a finite
  local Artinian $\cO/\varpi^a$-algebra, and $M$ is a rank one \'etale
  $\varphi$-module over~$\A_{K,A}$ with the property that there is a
  rank~$[K:\Kbasic]$ $\varphi$-module $\gM$ over $\A_{\Kbasic,A}^+$, such
  that:
  \begin{itemize}
  \item $\gM[1/T]=M$ (where we are regarding~$M$ as an \'etale
    $\varphi$-module over~$\A_{\Kbasic,A}$),
  \item $\gM$ is of $T$-height at most~$h$, and
  \item the action of~$\Gamma_{\Kbasic,\disc}=\Gamma_{K,\disc}$ on $\gM[1/T]$ extends the canonical action of
    $\Gamma_{\Kbasic_s,\disc}$ given by Corollary~\ref{cor: Caruso Liu Galois
      action phi Gamma discrete version},
  \end{itemize}
then the action of~$I_K^{h,s,a}$ on~$T_A(M)$ is trivial. This is 
immediate from Proposition~\ref{prop: ramification bound in rank one case with coefficients
    and any K with GK action}.%\mar{TG: again fix cited ramification bound.}
\end{proof}


\begin{rem}
  \label{rem: rank one X not closed substack}
  Note that Propositions~\ref{prop:rank 1 R-Gamma-disc} and
  \ref{prop:rank 1 X} describe $\cX$ as a certain formal completion
  of $\cR^{\Gamma_{\disc}}$ (in the case $d = 1$). %  \mar{TG: amplify
    % this slightly?}
  In particular, the monomorphism
  $\cX \hookrightarrow \cR^{\Gamma_{\disc}}$ is not a closed immersion
  (in the case $d =1$, and presumably not in the general case either).
This helps to explain why
somewhat elaborate arguments were required in
Chapter~\ref{sec: phi modules and phi gamma modules}
to deduce properties of $\cX$ from the corresponding properties of
$\cR^{\Gamma_{\disc}}$.
%	  In the case~$d=1$, one
%  can give a relatively explicit description of~$\cX_1$ as a substack
%  of~$\cR_1^{\Gamma_{\disc}}$.
 % We content ourselves with working on
 % the special fibre in the case~$K=\Qp$. Let $A$ be an $\Fp$-algebra;
 % then a rank one \'etale $\varphi$-module with $A$-coefficients is a
 % rank one projective $A[[T]]$-module with a semi-linear action
 % of~$\varphi$. \mar{TG: this is going to have an actual argument;
 %   idea is to think about $\Gm$ action on underlying reduced.}
\end{rem}

\begin{rem}\label{rem: rank 1 versus rank d}
	% There is one last simpification of the situation
	% that occurs in the case $d = 1$, which we note here.
	It follows easily from Proposition~\ref{prop:rank 1 X} that the stack $\cX_1$ may be described
	as a moduli stack of $1$-dimensional continuous
	representations of the Weil group~$W_K$; indeed, we see that if
        $A$ is a $p$-adically complete $\cO$-algebra, then $\cX_1(\Spf
        A)$ is the groupoid of %pairs consisting of a
        continuous characters~$I_K^{\ab}\times \Z \to
        A^\times$ (if we identify $\widehat{\mathbb G}_m$ with the $p$-adically
completed group ring of $\Z$). %and an element~$a\in A^\times$,
Mapping the generator $1 \in \Z$ to $\sigma_K$ (a lift of Frobenius, 
as in the proof of Proposition~\ref{prop:rank 1 X})
yields an isomorphism $I_K^{\ab} \times \Z \iso W_K^{\ab}$.

%which corresponds to
%        the representation $W_K\to A^\times$ sending~$\sigma_K\to a$,
%        where~$\sigma_K$ is
%as in the proof of Proposition~\ref{prop:rank 1 X}.
%\mar{TG: not sure what I'm doing here, please correct!}

% \mar{TG: probably we
%           should make this explicit as a corollary?}
	% (This is not obvious to us {\em a priori},
	% but follows from the explicit description given in
	% Proposition~\ref{prop:rank 1 X}.) 
        This phenomenon
	doesn't persist in the general case.  Indeed,
	as noted in Section~\ref{subsec: families of extensions}, % \mar{ME: Add reference back to appropriate 
		% para.\ in the introduction},
	already in the case $d = 2$, there seems to be no description
        of $\cX_2$ as the moduli space of representations of 
	a group that is compatible with the description
       of its closed points in terms of representations of~$G_K$. 

One can attempt
%Note in particular that if one attempts
to adapt the proof of
Proposition~\ref{prop:rank 1 X} to the case~$d>1$, by constructing a
morphism from the moduli stack of continuous $d$-dimensional Weil--Deligne representations
(suitably understood\footnote{Since we are considering representations
in rings that are $p$-power torsion,
we cannot use the usual formulation of the Weil--Deligne group.  Rather,
we choose a finitely generated dense subgroup
$\Z \ltimes \Z[1/p]$ 
of the tame Galois group of $G_K$ (by choosing a lift of Frobenius and a lift
of a topological generator of tame inertia), and form the topological
group $WD_K$ by taking its preimage in $G_K$ (and equipping 
$\Z \ltimes \Z[1/p]$ with its discrete topology).})
to the stack~$\cX_d$.
However, when $d > 1$,
%then the corresponding
this morphism of stacks will not be representable by
algebraic spaces. The precise point where the proof that we've given 
in the rank $1$ case fails to
generalise %from the rank one case
is that once we are not passing to abelianized Galois
groups, the upper numbered ramification
groups are not open in the inertia group.
% \mar{TG: add a remark that if you tried to extend the proofs above to
%   $d$-dimensional representations, you would find that the morphisms
%   are not representable by algebraic spaces, and an explicit problem
%   with trying to verify this in the way that we do for~$d=1$ is that
%   the upper ramification groups aren't open.}
\end{rem}

\section{A ramification bound}% \mar{TG: version 4610 is before we
  % messed with $n$.}
\label{subsec:ram bound}
We end this chapter by proving the bounds on the ramification of a character valued in a  finite
Artinian $\cO$-algebra %  in terms of the $T$-height of the corresponding
% $(\varphi,\Gamma)$-modules
that were used in the proofs of
Propositions~\ref{prop:rank 1 R} and~\ref{prop:rank 1 X}. % \mar{TG: used anywhere else?}
There are well-established techniques for
proving such a bound, going back to~\cite{MR807070}. We find it
convenient to follow the arguments of~\cite{MR2745530}, and indeed we
will follow some of their arguments very closely. The main differences
between our setting and that of~\cite{MR2745530} are that we are
working in the cyclotomic setting, rather than the Kummer setting;
that we are considering representations of~$G_{K_{\cyc}}$ of finite
height, rather than semistable representations of~$G_K$; and that to
define a representation of finite height, we need to pass between~$K$
and~$\Kbasic$. None of these changes make a fundamental difference
to the argument.

In fact, while~\cite{MR2745530} go to some effort to work modulo an
arbitrary power of~$p$, and to optimise the bounds that they obtain,
we are content to give the simplest proof that we can of the existence
of a bound of the kind that we need. In particular, we are able to
reduce to the case of mod~$p$ representations, % \mar{TG: reference where
  % this happens below},
which simplifies much of the discussion.
% \mar{TG: I'm imagining writing everything for~$\GL_d$,
%   maybe even the stacks arguments where we are constructing the map
%   from~$U$, and then just observing at the very end that if~$d=1$ we
%   actually get representability in algebraic spaces, and otherwise
%   it's not true.}

% We now prove two lemmas using techniques based on those
% of~\cite{MR2745530}, although adapted to the cyclotomic rather than
% Kummer setting.
% \mar{TG: actually I'm going to scrap this and just follow CL!}
% We begin with Lemma~\ref{lem: canonical extension
%   finite height cyclotomic}, 
% which has an essentially
% identical proof to Lemma~\ref{lem: Caruso Liu Galois action on Kisin}
% which is based in turn on \cite[\S2]{MR2745530}; for the convenience
% of the reader, we sketch the argument in our current setting.

% \begin{lem}\label{lem: canonical extension finite height
%     cyclotomic}Assume that~$K$ is basic.
%           For any fixed $a,h$, there is a constant~$C(a,h)$ such that
%           for any $N\ge C(a,h)$,  there is a
%   positive integer~$s(a,h,N)$ with the property that for any finite type
%   $\cO/\varpi^a$-algebra~$A$, any
%   projective   $\varphi$-module~$\gM$ over~$\A_{K,A}^+$  of $T$-height at most~$h$, and 
%   any~$s\ge s(a,h,N)$, there is a unique extension of the trivial
%   action of~$G_{K_{\cyc}}$ on~$\gM$ to a  continuous action of~
%   $G_{K_{\cyc,s}}$  on
%   $\gMt:=\AAinf{A}\otimes_{\A_{K,A}^+}\gM$, which commutes with~$\varphi$
%   and is semi-linear with respect to the natural action
%   of~$G_{K_{\cyc,s}}$
%   on~$\AAinf{A},$ and has the additional property that for
%   all~$g\in G_{K_{\cyc,s}}$ we have $(g-1)(\gM)\subset
%   T^N\gMt$.  \end{lem}

% \begin{proof}% We follow the proof of Lemma~\ref{lem: canonical extension
%   % finite height cyclotomic} closely.
% Since the~$u$-adic and $T$-adic
% topologies on~$\AAinf{A}$ coincide, there is some integer~$k\ge 1$
% (depending only on~$a$) such that $T^k\AAinf{A}\subseteq u\AAinf{A}$
% and $u^k\AAinf{A}\subseteq T\AAinf{A}$. We
% let~$C(a,h)$ be any integer greater than~$ek(a+h)/(p-1)$.

%  To give a semi-linear action of~$G_{K_{\cyc,s}}$  on~$\gMt$ is
%   to give for each~$g\in G_{K_{\cyc,s}}$ a morphism of Breuil--Kisin--Fargues
%   modules \[\beta_g:g^*\gMt\to\gMt\] with the property that for
%   all~$g,h\in G_{K_{\cyc,s}}$, we have $\beta_{gh}=\beta_h\circ
%   h^*\beta_g$. % \mar{TG: I think this is the right way round, but it
%     % would be sensible to check for nonsense.}
%   Now, the requirement that
%   $(g-1)(\gM)\subseteq T^N\gMt$ implies that if we have two such
%   morphisms~$\beta_g,\beta'_g$ then
%   $(\beta_g-\beta'_g)(g^*\gMt)\subseteq T^N\gMt\subseteq u^{\lfloor N/k\rfloor}\gMt$, so
%   that~$\beta'_g=\beta_g$ by the uniqueness assertion of
%   Lemma~\ref{lem: Frobenius amplification lemma over Ainf}.

%   % Since the homomorphism $\gS_A\into\AAinf{A}$ takes~$T$ to the
%   % Teichm\"uller lift of~$(\pi^{1/p^n})_n$,\mar{TG: fix} 
%   % \mar{TG: have we actually
% %     discussed this above? I suspect that we use the injectivity of
% %     various coefficient rings into $W(\C^\flat)_A$ throughout the
% %     paper - do we prove it? ME: The faithful flatness result
% % in \S 3 on Galois descent gives injectivy for perfect rings.  Hopefully
% % there is an argument that shows that the non-perfect rings embed into
% % (rather than just map to) the perfect ones.}
%   We can and do choose~$s(a,h,N)$ sufficiently large that for
%   all~$s\ge s(a,h,N)$ and~$g\in G_{K_{\cyc,s}}$, we
%   have~$g(T)-T\in T^{pN}\AAinf{A}$; thus for
%   all~$\lambda\in \A_{K,A}^+$, we have
%   $g(\lambda)-\lambda\in T^{pN}\AAinf{A}$. The existence of the
%   required action can then be shown exactly as in the proof of Lemma~\ref{lem: Caruso Liu Galois action on Kisin}, replacing the use
%   Lemma~\ref{lem: adding projective Kisin
%   to get free} by an obvious analogue for  $\varphi$-modules
% over~$\A^+_{K,A}$  (which can be proved in exactly the same way).
% \end{proof} % ,\mar{TG:
%   we could just say that the proof is identical to the BK version from here?}
% %we can and do assume from now on that~$\gM$ is free; indeed, in general
%   we may write~$\gF=\gM\oplus\gP$ where~$\gF,\gP$ are respectively free and
%   projective $\varphi$-modules over~$\A^+_{K,A}$ 

% We now use the existence
%   part of Lemma~\ref{lem: Frobenius amplification lemma over Ainf} to
%   construct the required action.
  % \mar{TG: need to choose~$N$
    % and of course need to justify that~$s$ exists, but I think this
    % should be just the continuity of the action of $G_K$ on~$\Ainf$.}

% We claim that for each~$g\in G_{K_{\cyc,s}}$, we may define
%   an~$\AAinf{A}$-linear map  \[\alpha_g:g^*\gMt\to
%     \gMt\] with the following two properties:
% \begin{itemize}
% 	\item
% $(\alpha_g\circ\Phi_{g^*\gMt}-\Phi_{\gMt}\circ\varphi^*\alpha_g)(\varphi^*g^*\gMt)\subseteq T^{pN}\gMt.$
% \item
% For each $m \in \gM,$ we have $\alpha_g(1\otimes m) - m \in T^N \gMt.$
% \end{itemize}
% It suffices to construct such maps in the case when~$\gM$ is free;
% indeed, by an obvious analogue of Lemma~\ref{lem: adding projective Kisin
%   to get free} (which can be proved in exactly the same way),\mar{TG:
%   we could just say that the proof is identical to the BK version from here?}
% %we can and do assume from now on that~$\gM$ is free; indeed, in general
%   we may write~$\gF=\gM\oplus\gP$ where~$\gF,\gP$ are respectively free and
%   projective $\varphi$-modules over~$\A^+_{K,A}$ of height at most~$h$, and given a
%   morphism~$\alpha_g:g^*\gF^{\inf}\to\gF^{\inf}$ satisfying our
%   requirements (where we write $\gF^{\inf}=\AAinf{A}\otimes_{\gS_{A}}\gF$),
%   we may obtain the desired
%   morphism~$\alpha_g:g^*\gMt\to\gMt$ by projection onto the
%   corresponding direct summands.


% If then~$\gF$ is a free $\varphi$-module over~$\A^+_{K,A}$ of  height at most~$h$, we choose a
% basis~$f_1,\dots,f_d$ for~$\gF$, and then, for each~$g\in G_{K_{\cyc,s}}$, we define
%   an~$\AAinf{A}$-linear map  \[\alpha_g:g^*\gFt\to
%     \gFt\] (which depends on our choice of basis) by

%   \[\alpha_g(\sum_i\lambda_i\otimes f_i)=\sum_i\lambda_if_i.\] We now check
%   that~$\alpha_g$ has the claimed properties.
%   If~$\Phi_{\gF}(1\otimes f_i)=\sum_j\theta_{i,j}f_j$ (so
%   that~$\theta_{i,j}\in\A_{K,A}^+$), then 
%   for any $\sum_i\lambda_i\otimes
%     f_i\in\varphi^*g^*\gFt, $ we compute that
%     \[(\alpha_g\circ\Phi_{g^*\gFt}-\Phi_{\gFt}\circ\varphi^*\alpha_g)(\sum_i\lambda_i\otimes
%     f_i)=\sum_{i,j}\lambda_i(g(\theta_{i,j})-\theta_{i,j})f_j\in
%     T^{pN}\gFt,\] 
%   while for any $\sum_i \lambda_i f_i \in \gF$
%   (note then that each $\lambda_i \in \A_{K,A}^+$),
%   we compute that
%     $$\alpha_g\bigl(1\otimes (\sum_i \lambda_i f_i)\bigr)
%     - \sum_i \lambda_i f_i
%     = 
% 				 \alpha_g(\sum_ig(\lambda_i)\otimes
%                                  g_i)-\sum_i\lambda_ig_i = 
% 				 \sum_i(g(\lambda_i)-\lambda_i)g_i
% 				 \in T^N \gFt;
% 				 $$
%   in both computations we have used the fact that, by our choice of~$s$,
%   if~$\lambda\in\A_{K,A}^+$, 
%   then~$g(\lambda)-\lambda\in T^{pN}\gFt \subseteq T^N\gFt$.

%  Returning to the general case of a projective Breuil--Kisin
%  module~$\gM$ of height at most~$h$, it follows from % our
%   % assumption on~$s$, and
%   Lemma~\ref{lem: Frobenius amplification lemma
%     over Ainf} that there is a unique % $\AAinf{A}$-linear,
%   % $\varphi$-semi-linear 
% morphism of Breuil--Kisin--Fargues modules
%   \[\beta_g:g^*\gMt\to
%  \gMt\]
%   with~$\im(\beta_g-\alpha_g)\subseteq
%   T^N\AAinf{A}\otimes_{\A_{K,A}^+}\gM$. Since we have
%   $\alpha_{gh}=\alpha_h\circ h^*\alpha_g$, it follows from this
%   uniqueness that~$\beta_{gh}=\beta_h\circ h^*\beta_g$, % \mar{TG: does
%     % this answer the previous marginal comment?}
%   so we have
%   constructed an action of~$G_{K_{\cyc,s}}$ on~$\gMt$. It remains to check
%   that~$(g-1)(\gM)\subseteq T^N\gMt$; for this, note that,
%   for any $m \in \gM$, we have
%     $$(g-1)m = \beta_g(1\otimes m)  - m = (\beta_g - \alpha_g)(1\otimes m)
%     + \alpha_g(1\otimes m ) - m \in T^N \gMt;
%     $$
% here we have used the fact $\im(\beta_g-\alpha_g)\subseteq
%   T^N\AAinf{A}\otimes_{\A_{K,A}^+}\gM$, together with the second of the
%   properties satisfied by the maps $\alpha_g$.


% In the setting of Lemma~\ref{lem: canonical extension finite height
%   cyclotomic}, the $W(\C^\flat)_A$-module
% $M^{\inf}:=W(\C^\flat)_A\otimes_{\A^+_{K,A}}\gM$ is naturally a
% $(G_{K_{\cyc,s}},\varphi)$-module with~$A$-coefficients, in the sense
% of Definition~\ref{defn: GK module over Ainf}. If we suppose that~$A$
% is a finite Artinian $\cO/\varpi^a$-algebra, then we have the  associated
% $G_{K_{\cyc,s}}$-representation~$V(M^{\inf}):=(M^{\inf})^{\varphi=1}$.\mar{TG:
% there is going to be a bunch of confusion between $\gM$, $\gMt$ and
% $M^{\inf}$ in what follows, which needs to be sorted out.} 


Recall that~$K_{\cyc}/K$ is a Galois extension with Galois
group~$\Gamma_K\cong\Z_p$. For each~$s\ge 0$, we write~$K_{\cyc,s}$
for the unique subfield of~$K_{\cyc}$ which is cyclic over~$K$ of
degree~$p^s$. We write~$e=e(K/K_0)$.

% Let~$A$ be a finite Artinian~$\cO$-algebra, so that in particular~$A$
% is an~$\cO/\varpi^a$-algebra for some~$a\ge 1$.
Suppose that~$K$ is 
basic, and that ~$\gM$ is a free $\varphi$-module over $\A^+_{K}/p$ of
$T$-height at most~$h$.
% \mar{ME: We have $n$ here where we have $a$ throughout the rest of the
% paper, fwiw. TG: added a sentence about this, comment out if OK.}
 Then~$M:=\A_{K}\otimes_{\A^+_{K}}\gM$ is an
\'etale $\varphi$-module over~$\A_{K}/p$, and we have a corresponding
$G_{\Kcyc}$-representation $T(M)$ as in Section~\ref{subsec:Galois
  reps}. (Note that while elsewhere in the book we have worked with
$\cO/\varpi^a$-coefficients, it is more convenient to use
$\Z/p\Z$-coefficients in most of this section.) % ; but since we will
% eventually return to $\cO/\varpi^a$-coefficients, we prefer to use the
% letter~$n$, rather than~$a$.)
% (Note that since
% ~$A$ is finite Artinian, it is a product of local finite Artinian
% $\cO$-algebras, so the results of Section~\ref{subsec:Galois reps}
% immediately extend to~$A$.)

% For the time being % \mar{TG: make it clear where we come back from this
%   % assumption}
% we will assume that~$a=1$, so that~$A$ is in fact a
% finite-dimensional $\F$-vector space,  and
% We will in fact forget the
% $A$-module structure on~$T_A(M)$ (which is irrelevant for proving
% a bound on the ramification), and simply regard it as a finite-dimensional
% free $\Z/p^n$-module for some~$n\ge 1$.\mar{TG: note that~$a$ now has two meanings!! I
%   think this is the one I want to work with for now though. TG: OK,
%   writing~$n$ instead of~$a$. There is some pain, though, as I guess
%   we needn't be free over any $\Z/p^n\Z$?}
In order to compare directly to the
arguments of~\cite{MR2745530}, it is more convenient to work
 with the contragredient $G_{K_{\cyc}}$-representation $T(M)^{\vee}$,
which can also be computed as
%\mar{TG: changed this a bit; I think this is more accurate, and is what we need? ME: Looks good, commenting out.}
$T(M^{\vee}) = (\C^\flat\otimes_{\A_K} M^{\vee})^{\varphi=1}= \Hom_{\A_K,\varphi}\bigl(M,{\C}^\flat\bigr)= \Hom_{\A^+_K,\varphi}\bigl(M,{\C}^\flat\bigr).$
By~\cite[A.1.2.7, B.1.8.3]{MR1106901}, this can also be computed
via the functor %\mar{TG:put in refs here}
\[T^{\vee}(\gM):=\Hom_{\A^+_K,\varphi}\bigl(\gM,\O_{\C}^\flat\bigr).\]
More precisely, there is a diagram
$$\xymatrixcolsep{-0.1in}
\xymatrix{
T^{\vee}(\gM):=
\Hom_{\A^+_K,\varphi}\bigl(\gM,\O_{\C}^\flat\bigr)
\ar[rd]&&
\Hom_{\A^+_K,\varphi}\bigl(M,{\C}^\flat\bigr) = T(M^{\vee})\ar[ld] \\
&\Hom_{\A^+_K,\varphi}\bigl(\gM,{\C}^\flat\bigr)&}$$
(the first arrow being induced by the inclusion of 
$\O_{\C}^\flat$ in ${\C}^\flat$, and the second by
restriction from $M$ to~$\gM$),  with both arrows in fact being
isomorphisms (by \emph{op.\ cit.}).

% \mar{ME: Should be $W_n(\O_{\C}^\flat)$ instead (or as well),
% so that the subsequent discussion makes sense? TG: corrected,
% commenting out.}
% \mar{TG:
  % check this is correct and write it in terms of~$M$ instead - if we
  % can (I'm confused about whether this makes sense with
  % the~$G_{\cyc,s}$-action, I'm guessing not?). Is there any twist
  % going on, I would have guessed that we wanted
  % $\C^\flat/\cO_{\C^\flat}$ coefficients?}\mar{TG: quite possibly
  % instead put a dual on~$M$, not worrying about this for now.}

% \mar{TG:
% now forget coefficients (and hopefully reduce to the case of being
% flat over $\Z/p^a$), take the dual, and follow CL as closely as we
% can. I am optimistic that it will be more or less just an issue of
% tweaking a few constants (that we don't care about anyway). Also note
% that the comparison between~$K$ and~$\Kbasic$ involves relating to the
% induction of Galois representations, but this can be done after
% inverting~$T$ so hopefully we can just reference something in the
% literature.}


%\mar{TG: actually might be better to abandon all coefficients altogether??}
We will now follow the arguments of~\cite[\S 4]{MR2745530}; note that
the ring denoted~$R$ in~\cite{MR2745530} is~$\cO_\C^\flat$. For the
time being we will assume that we are in the following situation; the
rather complicated hypotheses here will allow us to interpret the
canonical actions of Corollary~\ref{cor:
      Caruso Liu Galois action phi Gamma discrete version} in terms of the constructions of~\cite[\S 4]{MR2745530}.
  % (Note that
% we do not claim that the association of~$V(M)$ to~$M$ is an
% equivalence of categories.) \mar{TG: need to explain that this is in
%   agreement with the existing~$G_{\Kcyc}$-representation}


\begin{hyp}\label{hyp: K is basic and M has height h and s is big enough}
  Assume that~$K$ is basic, and fix some~$h\ge 1$. Let~$\gM$ be a free
  $\varphi$-module over $\A^+_{K}/p$ of $T$-height at most~$h$. Fix
  the following:
  \begin{itemize}
  \item an integer~$b\ge eh/(p-1)$. 
  \item An integer ~$N> pb/e$ such
    that $N\ge N(1,h)$, where $N(1,h)$ is as in Corollary~\ref{cor: Caruso
      Liu Galois action phi Gamma discrete version}.  
  \item An integer~$s$ such that~$p^{s-1}>b/e$, and $s\ge s(1,h,N)$,
    where $s(1,h,N)$ is as in Corollary~\ref{cor: Caruso
      Liu Galois action phi Gamma discrete version}. 
  \end{itemize}
\end{hyp}
%\mar{TG: check what's going on here}
Note in  particular that under Hypothesis~\ref{hyp: K is basic and M has height h and s is big enough} we have $p^s> ph/(p-1)>h$.% \mar{TG:
  % check this is OK. Also is $e=e(K/K_0)$ defined?}

% Fix
% integers~$N\ge C(1,h)$ and $s\ge s(1,h,N)$ as in Lemma~\ref{lem: canonical extension finite height
%     cyclotomic}. 
We normalize the valuation~$v$ on~$K$ so
% \mar{ME: This discussion of valuations should be moved up (see above comment
% on valuations).}
that~$v(K^\times)=\Z$, % \mar{TG: this may or may not be the optimal
  % choice}
and continue to write~$v$ for the unique compatible
valuation on~$\overline{K}$, and for the induced
valuation on~$\cO_\C^\flat=\varprojlim\cO_{\overline{K}}/p$. Note that
with this convention we have~$v(T)=v(p)=e$ (where we are
abusively continuing to denote the image of~$T$ in~$\cO_\C^\flat$
by~$T$; this follows
straightforwardly from the definition of~$T$ as the trace of
$T' := ([\varepsilon]-1)$, which shows (after a simple computation)
that $v(T) = (p-1)v(T')$, along with the formula
$v(\zeta_{p^s}-1)=v(p)/p^{s-1}(p-1)$, which shows that $v(T') = v(p)/(p-1)$).
% \mar{TG: need to do something
%   about this trace? but maybe it's actually much easier now that we're
% in characteristic $p$ -- maybe it's literally just $(T')^{p-1}$? ME: Not sure
% that it's literally that, but I  just  computed that is congruent to that
% mod $(T')^p$, which is good enough. I added a slight elaboration to
% the text. TG: looks good, commenting out.}
We write $\mf{a}_{\cO_{\C}^\flat}^{>
  c}$ for the set of elements of~$\cO^\flat_{\C}$
of valuation
% \mar{ME: Should note that valuations are defined so that $v(p) =
%   e$. TG: done, commenting out}
 greater than~$c$, and define $\mf{a}_{\cO_{\C}^\flat}^{\ge
  c}$ in the same way. 
%  We write~$[\mf{a}_{\cO_{\C}^\flat}^{>
%   c}]$ for the
% ideal of~$W_n(\cO^\flat_{\C})$ generated by the elements~$[x]$ for~$x\in \mf{a}_{\cO_{\C}^\flat}^{>
%   c}$, and similarly~$[\mf{a}_{\cO_{\C}^\flat}^{\ge
%   c}]$. 
For any~$c\ge 0$,
% \mar{ME: It seemed  to me that maybe this interval should actually be
% $[0, ep^{s-n+1}).$ Perhaps it would be good to explain why we
% but a bound on $c$ at all, since this would help to determine what
% the correct bound is! TG: there is no reason at all to have this
% bound, it turns out, so I'm getting rid of it! Commenting out.}
% \mar{TG: probably fix elsewhere; and surely
  % as~$s$ is fixed we can replace~$c/p^s$ with~$c$ as long as we point
  % out that we are deviating from~\cite{MR2745530} in doing so? TG:
  % actually I think I will stick with this? Although right now I'm
  % confused about putting~$c=pb$ being in the correct range}
we
write 
\[T_{\cO_\C^\flat,c}^\vee(\gM):=\Hom_{\A^+_K,\varphi}\bigl(\gM,\cO_\C^\flat/\mf{a}_{\cO_\C^\flat}^{>c}\bigr),\]
so that there is a natural map $T^\vee(\gM)\to T^\vee_{\cO_\C^\flat,c}(\gM)$. For any~$c'\ge
c$ % \mar{TG: I seem to be ending up echoing CL's notation that I found
  % confusing, so hopefully when the dust settles I can find something
  % cleaner}
we write 
\[T_{\cO_\C^\flat,c',c}^\vee(\gM)=\Im(T_{\cO_\C^\flat,c'}^\vee(\gM)\to T_{\cO_\C^\flat,c}^\vee(\gM)),\]
where the    morphism is induced by the natural map $\mf{a}_{\cO_\C^\flat}^{>c'}\to
    \mf{a}_{\cO_\C^\flat}^{>c}$. % Note in particular that since
    % $p^s>p^nb/e\ge pb/e$, we
    % can apply this definition with~$c=b$, $c'=pb$.

\begin{lem}% \mar{TG: need to define~$b$, need to figure out this
    % slightly confusing thing with~$a$, $b$.}
  \label{lem: truncation for Ainf mod p representation}Assume that we are in the setting of Hypothesis~{\em\ref{hyp: K
    is basic and M has height h and s is big enough}}. Then % we have a
%   natural identification
% $T^\vee(\gM)=\Hom_{\A_{K}^+,\varphi}(\gM,W_n(\cO_\C^\flat))$. Furthermore,
the induced morphism
$T^\vee(\gM)\to T^\vee_{\cO_\C^\flat,b}(\gM)$
is injective, with image~$T^\vee_{\cO_\C^\flat,pb,b}(\gM)$.
\end{lem}
\begin{proof}%\mar{TG: no more $\gMt$, and basically just cite earlier arguments/CL}
 %  As above, we can use our fixed basis $e_1,\dots,e_d$ of~$\gM$ to
% identify $\Hom_{\cO_\C^\flat,\varphi}(\gM^{\inf},\C^\flat)$ with
% \numequation\label{eqn: explicit V star mod p}\{(x_1,\dots,x_d)\in(\C^\flat)^d\mid
%   (x_1^p,\dots,x_d^p)=(x_1,\dots,x_d)A\}.\end{equation}If we
% let~$x=\min_iv(x_i)$, then we have~$px=\min_iv(x_i^p)$, and
% since~$A\in M_d(\cO_\C^\flat)$, we conclude that~$px\ge x$, so that~$x\ge
% 0$. Thus
% $\Hom_{\cO_\C^\flat,\varphi}(\gM^{\inf},\C^\flat)=\Hom_{\cO_\C^\flat,\varphi}(\gM^{\inf},\cO_\C^\flat)$,
% which establishes the first part.
% so that
% \[V^*(M^{\inf})=\Hom_{\C^\flat,\varphi}(M^{\inf},\C^\flat)=\Hom_{\cO_\C^\flat,\varphi}(\gM^{\inf},\C^\flat),\]so
% our first  task is to show that


% To see that the morphism $V^*(M^{\inf})\to\Hom_{\cO_\C^\flat,\varphi}(\gM^{\inf},\cO_\C^\flat/\mf{a}_{\cO_\C^\flat}^{>b})$
% is injective, we need to show that if $x>b$ for all~$i$ then
% the~$x_i$ are all zero. If some~$x_i\ne 0$, then rewriting~\ref{eqn: explicit V star mod p} as 
% \[\{(x_1,\dots,x_d)\in(\C^\flat)^d\mid
%   (x_1^p,\dots,x_d^p)B=T^h(x_1,\dots,x_d)\},\] we have
% $x+eh =x+hv(T)\ge px$, so that $b\ge x$, a contradiction.% \mar{TG: need to adjust constants so
%   % that this works.}

% The remainder of the proof 
  This can be proved by a standard Frobenius amplification argument
  exactly as in the proof % \mar{TG: maybe need to actually write out a
%     sketch argument, so as to explain where the hypotheses are used?
%     TG: OK this isn't too hard, and actually it's nothing to do with
%     $s$ at all, will add the correct explanation later! TG: now done,
%     comment out if it's OK. ME: Can we record somewhere above
% that $v(T) = e$, with a quick explanation as to why?}
  of~\cite[Prop.\ 2.3.3]{MR2745530}. For example, to see the
  injectivity, we note that since~$\gM$ is finitely generated, any
  element of the kernel corresponds to a $\varphi$-equivariant
  morphism $\gM\to \mf{a}_{\cO_\C^\flat}^{\ge c}$ for
  some~$c>b$. Since ~$\gM$ has $T$-height at most~$h$, and
  $v(T^h)=eh$, the $\varphi$-equivariance implies that the image of
  this morphism is
  contained in $\mf{a}_{\cO_\C^\flat}^{\ge pc-eh}$. Since $c>b\ge
  eh/(p-1)$ by hypothesis, we see that the sequence
  $c,pc-eh,p(pc-eh)-eh,\dots$ goes to infinity, and the injectivity follows.

% using that by Hypothesis~\ref{hyp: K
%     is basic and M has height h and s is big enough}, we have
%     $p^{s-n+1}> pb/e$.
    % , we
    % can apply this definition with~$c=b$, $c'=pb$.;
  We leave surjectivity to the reader; it is a simpler version of the
  proof of Lemma~\ref{lem: amplification from char p to char 0}
  below. (See also the proof of Lemma~\ref{lem: Frobenius
    amplification lemma over Ainf} for an almost identical
  argument.)\end{proof}

Suppose given $c \geq 0$ satisfying $p^s > c/e$.
(E.g.\ $b$ and $pb$ both satisfy this condition.)
The Frobenius $\varphi$  is an automorphism of the perfect $k$-algebra
$\cO_\C^\flat$ which multiplies valuations by~$p$,
and so $\varphi^{-s}$ 
%The usual morphism
%$\theta: \cO_\C^\flat\to \cO_\C/p = \cO_{\overline{K}}/p$
%(which is just projection onto the first term,
%when we write $\cO_\C^\flat = \varprojlim \cO_\C/p$),
%when precomposed with $\varphi^{-s}$ (the resulting map
%now being projection onto the $(s+1)$st term of the projective limit),
induces the first of the following sequence of  $G_K$-equivariant 
%
%
%Assume that we are in the setting of Hypothesis~\ref{hyp: K
%    is basic and M has height h and s is big enough}. By definition,
%  we have~$\cO_\C^\flat=\varprojlim_\varphi\cO_{\overline{K}}/p$; % , where the
%  % transition morphisms are given by the Frobenius morphism $\varphi:x\mapsto
%  % x^p$.
%  accordingly, we have a morphism of $k$-algebras
%  $\cO_\C^\flat\to k\otimes_{k,\varphi^s}\cO_{\overline{K}}/p$,
%  given by projection from $\varprojlim_\varphi\cO_{\overline{K}}/p$ to the
%  $s$th factor. By~\cite[Lem.\ 2.5.2]{MR2745530}, for any $c$ satisfying
%  that $p^{s-n+1}>c/e$, % (in fact, it would be enough to assume that $p^{s-n+1}>b/e$),
%  this induces a $G_K$-equivariant
% \mar{ME: Can we indicate where this comes from?  In particular, maybe indicate
% why we are working mod~$p^n$, how this relates to the bounds on~$b$,
% etc. TG: will do this, but not sure what you mean by ``why we are
% working mod~$p^n$''? Everything in sight is mod~$p^n$ by
% definition. TG: I corrected the citation to \cite{MR2745530}, and it's
% literally a direct citation to the identical statement, so hopefully
% it's OK? ME: Rewrote this to refer to $\theta$, which is known (and probably
% introduced back in \S 2), rather than reintroducing it {\em ab initio}.
% Also, what is the ``this'' that holds? TG: thanks; that sentence was
% left over from earlier edits, now removed. Commenting out.}
  isomorphisms of $k$-algebras \numequation\label{eqn: isomorphism from
    truncated Ainf to truncated Kbar}
  \cO_\C^\flat/\mf{a}_{\cO_\C^\flat}^{>c}\isoto
 k\otimes_{k,\varphi^s}\cO_{\C}^\flat/\mf{a}_{\C}^{>c/p^s} 
\iso
 k\otimes_{k,\varphi^s}\cO_{\overline{K}}/\mf{a}_{\overline{K}}^{>c/p^s} ,
\end{equation}
%\mar{TG: check if we actually need for $pb$, which is also true?}
where $\mf{a}_{\overline{K}}^{>c/p^s}$ has the %obvious
evident meaning, and the second isomorphism
is induced by projection onto the first factor in the projective
limit
$\O_\C^\flat  = \varprojlim_{\varphi} \O_\C/p =
\varprojlim_\varphi\cO_{\overline{K}}/p$.
(The upper bound on $c$ ensures that $p \in \mf{a}_{\overline{K}}^{>c/p^s}$,
so that this projection does indeed induce the indicated morphism; it is then
straightforward to verify that it  is an isomorphism.)
%\mar{ME: Since we are now mod $p$, I don't think it makes sense to refer
%to $\theta$ anymore, so I rewrote this for what I hope is maximal clarity.
%It also brings out something I think I was worried about earlier (in my now-commenteed-out question about working mod $p^n$): namely, for the second arrow to even
%exist (let alone be an isomorphism) we need $p \in \mf{a}_{\overline{K}}^{> c/p^s},$
%which means we neeed $p^s > c/e$. TG: sorry, that was there at some
%point and is now back again! ME: This all seems sorted now, so I'm commenting 
%out this comment.}
This isomorphism in turn induces an isomorphism
\numequation
\label{eqn:alternative description of Tc}
T^{\vee}_{\cO_C^\flat,c}(\gM) \iso
\Hom_{\A_K^+,\varphi}(\gM, \cO_{\overline{K}}/\mf{a}_{\overline{K}}^{>c/p^s}).
\end{equation}

\begin{lem}
\label{lem:G K s action}
If $p^s > c/e \geq 0,$ then the action of $G_{K_{\cyc,s}}$ on
$\cO_\C^\flat$
induces an action of $G_{K_{\cyc,s}}$
on $T^{\vee}_{\cO_C^\flat,c}(\gM).$ 
\end{lem}
\begin{proof}
The action of $\varphi$ on
$\cO_\C^\flat$
is $G_K$-equivariant, %(and so in particular $G_{K_{\cyc,s}}$-equivariant),
and so the $G_K$-action on
$\cO_\C^\flat$
induces an action on
$\Hom_{\varphi}\bigl(\gM,\cO_\C^\flat/\mf{a}_{\cO_\C^\flat}^{>c}\bigr)$. 
The claim of the lemma is that the restriction
of this action to $G_{K_{\cyc,s}}$ preserves
the subobject
$T_{\cO_\C^\flat,c}^\vee(\gM):=\Hom_{\A^+_K,\varphi}\bigl(\gM,\cO_\C^\flat/\mf{a}_{\cO_\C^\flat}^{>c}\bigr)$
of
$\Hom_{\varphi}\bigl(\gM,\cO_\C^\flat/\mf{a}_{\cO_\C^\flat}^{>c}\bigr)$.

Rather than proving this directly,
we use the
isomorphisms~\eqref{eqn: isomorphism from truncated Ainf to truncated Kbar}
and~\eqref{eqn:alternative description of Tc}; taking these into account,
it suffices to prove that 
the  action of $G_{K_{\cyc,s}}$ on
$ \Hom_{\varphi}(\gM, \cO_{\overline{K}}/\mf{a}_{\overline{K}}^{>c/p^s})$
preserves the subobject
$\Hom_{\A_K^+,\varphi}(\gM, \cO_{\overline{K}}/\mf{a}_{\overline{K}}^{>c/p^s})$.
In other words, we have to check that the $G_{K_{\cyc,s}}$-action
preserves the property of being $\A_K^+$-equivariant.
For this, it suffices to check that $G_{K_{\cyc,s}}$ acts
trivially on the image of $T$ in $\cO_{\overline{K}}/\mf{a}_{\overline{K}}^{>c/p^s}$
under the isomorphism~\eqref{eqn: isomorphism from truncated Ainf to truncated Kbar}.
Since $T$ can be expressed as a power series in $T'$, it in fact suffices
to check this for the image of $T'$.  But this image is equal to
$\zeta_{p^{s+1}} - 1 \bmod \mf{a}_{\overline{K}}^{> c/p^s}$,
which  {\em is} fixed by $G_{K_{\cyc,s}}$ (since $\zeta_{p^{s+1}} \in K_{\cyc,s}$).
\end{proof}

%If we combine~\eqref{eqn:alternative description of Tc}
%with Lemma~\ref{lem: truncation for Ainf mod p representation}, 
%we obtain an isomorphism
%\numequation
%\label{eqn:alternative description of T}
%T^{\vee}(\gM) \iso \im\bigl(
%\Hom_{\A_K^+,\varphi}(\gM, \cO_{\overline{K}}/\mf{a}_{\overline{K}}^{>pb})
%\to
%\Hom_{\A_K^+,\varphi}(\gM, \cO_{\overline{K}}/\mf{a}_{\overline{K}}^{>b})
%\bigr)
%\end{equation}


%definition). % In particular,
% in the setting of Hypothesis~\ref{hyp: K
%     is basic and M has height h and s is big enough}, 
% this holds if~$c=pb$ or~$c=b$, by our
% hypothesis that $p^{s-n}>b/e$.
% \mar{ME: This is an ideal in $\cO_{\Kbar}$, not in a Witt vector ring,
% and so I don't think it should have any square brackets. TG: agreed,
% fixed, commented out.}
\begin{cor}\label{cor: we get a canonical action of G K s}Assume that we are in the setting of Hypothesis~{\em\ref{hyp: K
    is basic and M has height h and s is big enough}}. Then the natural
  action of~$G_{K_{\cyc,s}}$
  on~$\cO_\C^\flat$ %$W_n(\cO_\C^\flat)/[\mf{a}_{\cO_\C^\flat}^{>b}]$
  %and the isomorphism $T^\vee(\gM)\isoto T^\vee_{\cO_\C^\flat,pb,b}(\gM)$ of
  %Lemma~{\em\ref{lem: truncation for Ainf mod p representation}}
induces an
  extension of the action of~$G_{\Kcyc}$ on $T^\vee(\gM)$ to an action
  of~$G_{K_{\cyc,s}}$. % \mar{TG: this bit is a bit of a mess because we
    % don't have the definition of the $\overline{K}$, but maybe this
    % just requires a reference below now, as the notion still makes sense?}
\end{cor}
\begin{proof}
This follows directly from
Lemmas~\ref{lem: truncation for Ainf mod p representation}
and~\ref{lem:G K s action},
once we note that, by hypothesis, $p^s > pb/e > b/e.$
%
%
%
%  % \mar{ME: Can we actually define the action? TG: hopefully now done.}
%  % \mar{TG: add to statement and proof that this agrees with the
%  %   existing canonical action.}
%By the definition of $T^\vee_{\cO_\C^\flat,pb,b}(\gM)$, it is enough to
%show that for any ~$p^{s}e>c'\ge
%c$, the action of $G_{K_{\cyc,s}}$
%  on~$\cO_\C^\flat$ induces an action of~$G_{K_{\cyc,s}}$ on
%  $T_{\cO_\C^\flat,c'}^\vee(\gM)$ and $T_{\cO_\C^\flat,c}^\vee(\gM)$, which
%  is equivariant for the natural map % \mar{TG: I seem to be ending up echoing CL's notation that I found
%  % confusing, so hopefully when the dust settles I can find something
%  % cleaner}
%  $T_{\cO_\C^\flat,c'}^\vee(\gM)\to T_{\cO_\C^\flat,c}^\vee(\gM)$. To see
%  this, note that the action of $G_{K_{\cyc,s}}$
%  on~$\cO_\C^\flat$ induces an action on
%  $\cO_\C^\flat/\mf{a}_{\cO_\C^\flat}^{>c}$, and we claim that we can define
%  an action of~$G_{K_{\cyc,s}}$ on $T_{\cO_\C^\flat,c}^\vee(\gM)$ by its
%  action on $\cO_\C^\flat/\mf{a}_{\cO_\C^\flat}^{>c}$ (and the
%  trivial action on~$\gM$). It is immediate from the definitions that this action extends the action of~$G_{K_{\cyc}}$
%  on~$T^\vee(\gM)$, and that the actions for~$c',c$ are compatible.
%
%  In order for this definition to make sense (i.e.\ in order for it to
%  preserve $\A^+$-linearity of elements of $T_{\cO_\C^\flat,c}^\vee(\gM)$), we need to check that
%  the action of ~$G_{K_{\cyc,s}}$
%  on~$\cO_\C^\flat/\mf{a}_{\cO_\C^\flat}^{>c}\isoto
%  k\otimes_{k,\varphi^s}\cO_{\overline{K}}/\mf{a}_{\overline{K}}^{>c/p^s}$
%  induces the trivial action on the image of the element~$T$. % ,
%  % i.e.\ on the image
%  % of~$T_s\in\cO_{\overline{K}}$
%  % in~$ \cO_{\overline{K}}/\mf{a}_{\overline{K}}^{>b/p^s}$.
%  It is in turn enough to check that~$G_{K(\zeta_{p^{s+1}})}$ acts
%  trivially on the image of~$\zeta_{p^{s+1}}$, which is obvious.
\end{proof}
%\mar{TG: following setup and corollary should be read.}
By Corollary~\ref{cor: Caruso Liu Galois action phi Gamma discrete
  version} (and our hypothesized bound on~$s$), % there is a constant $N(1,h)$ such that for any
% $N\ge N(1,h)$, there is a constant $s(1,h,N)$ with the property that
% (uniformly for~$\gM$ satisfying our running assumptions) 
we can give $\gM[1/T]$ the structure of a
$(\varphi,\Gamma_{K_{\cyc,s}})$-module; this in particular induces an action
of~$G_{K_{\cyc,s}}$ on $T(M)$. % in
% such a way that if~$\gamma_s$ is a topological generator of
% $\Gamma_{K_{\cyc,s}}$, then $(\gamma_s-1)(\gM)\subseteq T^N\gM$. In
% particular, this means that $T(M)^\vee$ has an action of
% $G_{K_{\cyc,s}}$ for all $s\ge s(1,h,N)$.
The following corollary
expresses the compatibility of this action with the action constructed
in Corollary~\ref{cor: we get a canonical action of G K s}.

 
\begin{cor}% \mar{TG: this proof needs reading through. TG: I'm now
    % happy with it if you are.}
  \label{cor: the various canonical actions agree}% Suppose that
  % $N\ge N(1,h)$ and $s\ge s(1,h,N)$ are as in Corollary~{\em\ref{cor:
  %     Caruso Liu Galois action phi Gamma discrete version}}. Assume
  % further that $N\ge $. \mar{TG: may also need a bound on~$N$.} Then for
  % any $s\ge s(1,h,N)$ which satisfies Hypothesis~\ref{hyp: K is basic
  %   and M has height h and s is big enough}, 
Assume that we are in the setting of Hypothesis~{\em\ref{hyp: K
    is basic and M has height h and s is big enough}}. Then the two actions of
  $G_{K_{\cyc,s}}$ on $T(M)^\vee$ \emph{(}from Corollary~\emph{\ref{cor: we get a
    canonical action of G K s}}, and from Corollary~\emph{\ref{cor:
      Caruso Liu Galois action phi Gamma discrete version}}\emph{)} agree.
\end{cor}
\begin{proof}By the proof of Corollary~\ref{cor: Caruso Liu Galois action phi Gamma discrete
  version}, the semilinear action of $\Gamma_{K_{\cyc,s}}$ on~$\gM$
extends to a semilinear
action of $G_{K_{\cyc,s}}$ on $\cO_{\C}^\flat\otimes_{\A^+_K}\gM$,
which is uniquely determined by the properties that it commutes
with~$\varphi$, and satisfies $(g-1)(\gM)\subset
T^N(\cO_{\C}^\flat\otimes_{\A^+_K}\gM)$ for all~$g\in
G_{K_{\cyc,s}}$. As explained in Section~\ref{subsubsec:
Galois reps for GK phi}, we have
$T(M)=(\C^\flat\otimes_{\A^+_K}\gM)^{\varphi=1}$, with the action of $G_{K_{\cyc,s}}$ on
$T(M)$ being that inherited from its action on
$\cO_{\C}^\flat\otimes_{\A^+_K}\gM$. % In fact (for example by
% \cite[A.1.2.7, B.1.8.3]{MR1106901}, recalling that~$\gM$ is projective
% by hypothesis) we even have
% $T(M)=(\cO_{\C}^\flat\otimes_{\A^+_K}\gM)^{\varphi=1}$. \mar{TG: I
%   think this is completely false! e.g.\ imagine replacing $\gM$ by
%   $T^n\gM$ for~$n$ large. Rather I think what should be true is that
%   it's equal to the $\varphi$-invariants after tensoring $\gM$ with
%   $(\mf{a}^{> c})^{-1}\cO_{\C}^\flat/\cO_{\C}^\flat$. This probably
%   explains a lot of the trouble I had with unwinding to get a proof,
%   although I'm not sure of the best way to fix it.}

% \mar{TG: need
%   to justify this claim; I presume it is easy to unwind although I
%   don't think I know a reference, I'll try to put this in when I have time.
% ME: Since $\gM$  is projective, isn't this just a rephrasing of
% the  Fontaine citation with the big diagram earlier in this section?
% TG: good point, citation added, comment out if this is OK.}
% \mar{ME: Rephrasing what follows to emphasize that this is the natural
%   perfect pairing.}
As noted above, we also have
\begin{multline*}
T(M)^{\vee} = T(M^{\vee}) = T^{\vee}(\gM)
\\
=\Hom_{\A^+_K,\varphi}\bigl(\gM,\O_{\C}^\flat\bigr)
=\Hom_{{\C}^\flat,\varphi}\bigl({\C}^\flat\otimes_{\A^+_K}\gM,{\C}^\flat\bigr).
\end{multline*}
Now evaluation induces a pairing 
$$\bigl({\C}^\flat\otimes_{\A^+_K}\gM \bigr)\times
\Hom_{{\C}^\flat}\bigl({\C}^\flat\otimes_{\A^+_K}\gM,{\C}^\flat\bigr)
\to 
\C^\flat,$$
which restricts to a pairing
\numequation
\label{eqn:perfect pairing}
({\C}^\flat\otimes_{\A^+_K}\gM)^{\varphi = 1} \times
\Hom_{{\C}^\flat,\varphi}\bigl({\C}^\flat\otimes_{\A^+_K}\gM,{\C}^\flat\bigr)
\to 
(\C^\flat)^{\varphi = 1} = \Z/p\Z,
\end{equation}
and this latter pairing {\em is} the natural pairing of $G_{K_{\cyc}}$-representations
\[T(M) \times T(M)^{\vee} \to \Z/p\Z.\]
Therefore, to prove the lemma, we have to show that~\eqref{eqn:perfect pairing}
is furthermore $G_{K_{\cyc,s}}$-equivariant (with the action on~$T(M)$
being the action recalled above, coming from Corollary~\ref{cor: Caruso Liu Galois action phi Gamma discrete
  version}, and the action on~$T(M)^\vee=T^{\vee}(\gM)$ given by Corollary~\ref{cor: we get a canonical action of G K s}).
% pairing
%\[(\cO_{\C}^\flat\otimes_{\A^+_K}\gM)^{\varphi=1}\times T^\vee(\gM)\to
%  \Z/p\Z. \] Since
%$T^\vee(\gM)=\Hom_{\A^+_K,\varphi}\bigl(\gM,\O_{\C}^\flat\bigr)=\Hom_{\cO_{\C}^\flat,\varphi}\bigl(\cO_{\C}^\flat\otimes_{\A^+_K}\gM,\O_{\C}^\flat\bigr)$
%by definition, and $(\cO_{\C}^\flat)^{\varphi=1}=\Z/p\Z$, the obvious
%evaluation map gives us a perfect pairing to~$\Z/p\Z$, and we only
%need to check $G_{K_{\cyc,s}}$-equivariance.

To see this, note firstly that by the same argument that we used in
the proof of Lemma~\ref{lem: truncation for Ainf mod p
  representation}, it follows from Hypothesis~\ref{hyp: K is basic and
  M has height h and s is big enough}
that \[(\C^\flat\otimes_{\A^+_K}\gM)^{\varphi=1}\subseteq (\mf{a}^{\ge b})^{-1}\otimes_{\A^+_K}\gM.\]
In addition, we
have \begin{multline*}\Hom_{{\C}^\flat,\varphi}\bigl({\C}^\flat\otimes_{\A^+_K}\gM,{\C}^\flat\bigr)=\Hom_{\cO_{\C^\flat},\varphi}\bigl(\cO_{\C^\flat}\otimes_{\A^+_K}\gM,\cO_{\C^\flat}\bigr)\\
  \subseteq
  \Hom_{\cO_{\C^\flat}}\bigl(\cO_{\C^\flat}\otimes_{\A^+_K}\gM,\cO_{\C^\flat}\bigr)
  \isofrom\Hom_{\cO_{\C^\flat}}\bigl((\mf{a}^{\ge
    b})^{-1}\otimes_{\A^+_K}\gM,(\mf{a}^{\ge
    b})^{-1}\bigr)\end{multline*} with the last
isomorphism being induced by restricting the domain of a homomorphism from $(\mf{a}^{\ge
    b})^{-1}\otimes_{\A^+_K}\gM$ to
  $\cO_{\C^\flat}\otimes_{\A^+_K}\gM$. Bearing in mind this
  isomorphism, we obtain an evaluation pairing \[\bigl((\mf{a}^{\ge
      b})^{-1}\otimes_{\A^+_K}\gM\bigr) \times \Hom_{\cO_{\C^\flat}}\bigl((\cO_{\C^\flat}\otimes_{\A^+_K}\gM,\cO_{\C^\flat}\bigr)\to (\mf{a}^{\ge
      b})^{-1} \](and taking $\varphi$-invariants again recovers the pairing~\eqref{eqn:perfect pairing}).

  Fix some $\alpha\in\C^\flat$ with $v(\alpha)=-b$, so that
  $(\alpha)=(\mf{a}^{\ge b})^{-1}$. Let~$x$ be an element of
  $({\C}^\flat\otimes_{\A^+_K}\gM)^{\varphi = 1}$, and write
  $x=\alpha\sum_{i=1}^n\lambda_i\otimes x_i$ where
  $\lambda_i\in\cO_{\C}^\flat$ and $x_i\in\gM$. Let
  $ f\in
  \Hom_{\cO_{\C}^\flat,\varphi}\bigl(\cO_{\C}^\flat\otimes_{\A^+_K}\gM,\cO_{\C}^\flat\bigr)$,
  so that $x$ and $f$ pair to
  $\alpha f(\sum_{i=1}^n\lambda_ix_i)\in (\mf{a}^{\ge b})^{-1}$; of
  course, they in fact pair to an element
  of~$\Fp\subset\cO_\C^\flat \subset (\mf{a}^{\ge b})^{-1}$.
  Let~$g\in G_{\Kcyc,s}$ be arbitrary. Then $gx$ and $gf$ pair to
  $g(\alpha) (gf)(\sum_{i=1}^ng(\lambda_i)g(x_i))$. To establish the
  claimed $G_{\Kcyc,s}$-equivariance, we need to show that
  \[g(\alpha)\cdot (gf)\bigl(\sum_{i=1}^ng(\lambda_i)g(x_i)\bigr)=g\bigl(\alpha
    f\bigl(\sum_{i=1}^n\lambda_ix_i\bigr)\bigr).\] This is an equality of
  elements of~$\Fp$, and it therefore suffices to show that it holds
  modulo~$\m_{\cO_{\C}}$ when each side is considered as an element
  of~$\cO_{\C}^\flat$. Equivalently, after multiplication
  by~$\alpha^{-1}$ it suffices to show
  that
  \[(gf)\bigl(\sum_{i=1}^ng(\lambda_i)g(x_i)\bigr)\equiv
    g\bigl(f\bigl(\sum_{i=1}^n\lambda_ix_i\bigr)\bigr) \pmod{\mf{a}^{>
        b}}.\] Recall that
  $(g-1)(\gM)\subset T^N(\cO_{\C}^\flat\otimes_{\A^+_K}\gM)$. By
  Hypothesis~\ref{hyp: K is basic and M has height h and s is big
    enough}, we have $v(T^N)=eN>pb$, so
  that in particular we have~$T^N\in \mf{a}_{\cO_{\C}^\flat}^{>b}$. Thus
  $g(x_i)\equiv x_i\pmod{\mf{a}^{> b}}$, so it suffices to show
  that
  \[(gf)\bigl(\sum_{i=1}^ng(\lambda_i)x_i\bigr)\equiv
    g\bigl(f\bigl(\sum_{i=1}^n\lambda_ix_i\bigr)\bigr) \pmod{\mf{a}^{>
        b}},\] or equivalently
  that
  \[\sum_{i=1}^ng(\lambda_i)(gf)(x_i)\equiv
    \sum_{i=1}^ng(\lambda_i)g(f(x_i)) \pmod{\mf{a}^{> b}},\] so in
  turn it is enough to show that if $x\in\gM$ then
  $(gf)(x)\equiv g(f(x))\pmod{\mf{a}^{>b}}$. But this is true by the
  very definition of~$gf$ (which was defined via the action
  of~$G_{\Kcyc,s}$ on~ $T_{\cO_\C^\flat,b}^\vee(\gM)$, which in turn
  is defined via the action of~$G_{\Kcyc,s}$ on
  $\cO_{\C}^\flat/\mf{a}^{>b}$), so we are done.
\end{proof}

% \begin{rem}\mar{TG: adjust this remark to reflect the above corollary}
%   \label{rem: we could have used canonical actions on the level of
%     Wach modules}Alternatively, a version of Corollary~\ref{cor: we get a canonical
%     action of G K s} follows from Corollary~\ref{cor: Caruso Liu
%     Galois action phi Gamma discrete version}; we have chosen to give
%   this more direct proof in order to facilitate the comparison with~\cite{MR2745530}.
% \end{rem} 


% Assume now that~$n=1$.
Recall that we write~$k$ for the residue field
of~$K$, and since~$K$ is assumed basic we have $k_\infty=k$, so that we can
and do regard~$\gM$ as a free $\A_{K}^+/p=k[[T]]$-module. % \mar{ME: Maybe recall
% that $k_{\infty} = k$ b/c we're in the basic case.}
Write~$T_s$ for
$\tr_{K(\zeta_{p^{s+1}})/K_{\cyc,s}}(\zeta_{p^{s+1}}-1)\in\cO_{K_{\cyc,s}}$,
and note  that  we have a
homomorphism \numequation\label{eqn: A to Ks}\A_{K}^+/p=k[[T]]\to
  \cO_{K_{\cyc,s}}/p\end{equation} which sends~$T\to T_s$, and whose
restriction to~$k$ is given by $\varphi^{-s}$.
% \mar{TG: discuss $s$ versus $N$ versus~$h$ to get something nontrivial
%   here, and actually define~$b$! I guess $v(T_s)=1/p^s$, so we at
%   least need $p^s>h$.}

Let~$L$ be any algebraic extension of~$K_{\cyc,s}$
inside~$\overline{K}$. We % normalise the valuation~$v$ on~$K$ so
% \mar{ME: This discussion of valuations should be moved up (see above comment
% on valuations).}
% that~$v(K^\times)=\Z$, % \mar{TG: this may or may not be the optimal
%   % choice}
% and
continue to write~$v$ for the unique 
valuation on~$L$ with $v(p)=e$.  Note that we have~$v(T_s^h)=eh/p^s<e=v(p)$. For
any~$c\ge 0$ we write $\mf{a}_L^{> c}$  for the set of elements of~$L$
of valuation greater than~$c$, and  define  $\mf{a}_L^{\ge c}$ in the same way.  % \mar{TG: check if we need this
  % for $\ge c$ too}

Choose an (ordered) $\A_{K}^+/p$-basis $e_1,\dots,e_d$ of~$\gM$,
and %\mar{TG: check that we actually have some~$\gM$ fixed at this point!}
write $\varphi(e_1,\dots,e_d)=(e_1,\dots,e_d)A$ for some~$A\in
M_d(\A_{K}^+/p)$. By our assumption that~$\gM$ has $T$-height at most~$h$,
we can write $AB=T^h\Id_d$ for some $B\in M_d(\A_{K}^+/p)$. We can
and do choose matrices $\widetilde{A},\tB\in M_d(\cO_{K_{\cyc,s}})$ which
respectively lift the images of~$A,B$ under the homomorphism
$\A_{K}^+/p\to\cO_{K_{\cyc,s}}/p$ defined above. Since the image
of~$T_s^h$ in $\cO_{K_{\cyc,s}}/p$ is nonzero, we see that after
possibly multiplying~$\tB$ by an invertible matrix which is trivial
mod~$p/T_s^h$,
%\mar{ME: Maybe actually only trivial $\bmod \bigl(p/\im(T_s^h)\bigr)$?
%TG: agreed - is this an OK way to write this? ME: Looks OK to me,
%commenting out.}
we can and do assume that \numequation\label{eqn: Atilde
  times Btilde}\widetilde{A}\tB=T_s^h\Id_d.\end{equation}



We
set~\[\tT_L^*(\gM):=\{(x_1,\dots,x_d)\in\cO_L^d\mid
  (x_1^p,\dots,x_d^p)=(x_1,\dots,x_d)\widetilde{A}\}.\]
Note that since~$x\mapsto x^p$ is not a ring homomorphism on~$\cO_L$,
this is just a set with a $G_{K_{\cyc,s}}$-action, rather than
an~$\cO_L$-module. 


For any~$c\in [0,ep^s)$, % \mar{TG: probably fix elsewhere; and surely
  % as~$s$ is fixed we can replace~$c/p^s$ with~$c$ as long as we point
  % out that we are deviating from~\cite{MR2745530} in doing so? TG:
  % actually I think I will stick with this? Although right now I'm
  % confused about putting~$c=pb$ being in the correct range}
we
write 
% \[V_{\cO_\C^\flat,c}^*(\gM):=\Hom_{\A^+_K/,\varphi}(\gM,\cO_\C^\flat/\mf{a}_{\cO_\C^\flat}^{>c}),\]
\[T_{L,c}^*(\gM):=\Hom_{\A^+_K/p,\varphi}(\gM,(\cO_L/p)/\mf{a}_L^{>c/p^s}),\]where
the $\A^+_K$-algebra structure on $\cO_L/p$ is that induced by the
composite $\A^+_K/p \buildrel \text{\eqref{eqn: A to Ks}}\over\longrightarrow \cO_{K_s}/p\to\cO_L/p$. (The use of~$c/p^s$
rather than~$c$ in the definition of~$T_{L,c}^*$ comes from the
definition of the homomorphism~\eqref{eqn: A to Ks}; see also~\eqref{eqn: isomorphism from
    truncated Ainf to truncated Kbar}.) For any~$c'\ge
c$ % \mar{TG: I seem to be ending up echoing CL's notation that I found
  % confusing, so hopefully when the dust settles I can find something
  % cleaner}
we write 
%\[V_{\cO_\C^\flat,c',c}^*(\gM)=\Im(V_{\cO_\C^\flat,c'}^*(\gM)\to V_{\cO_\C^\flat,c}^*(\gM)),\]
\[T_{L,c',c}^*(\gM)=\Im(T_{L,c'}^*(\gM)\to T_{L,c}^*(\gM)),\]where the
    morphism is induced by the natural map %  and $\mf{a}_{\cO_\C^\flat}^{>c'}\to
    % \mf{a}_{\cO_\C^\flat}^{>c}$ and
    $\mf{a}_L^{>c'/p^s}\to
    \mf{a}_L^{>c/p^s}$. % Note in particular that since $p^s>pb/e$, we
    % can apply these definitions with~$c=b$, $c'=pb$.  
  %Note in particular that by Corollary~\ref{cor: we get a canonical action of G K s}, the isomorphism~\eqref{eqn: isomorphism from
  %  truncated Ainf to truncated Kbar} induces isomorphisms of $G_{K_{\cyc,s}}$-representations
%\numequation\label{eqn: iso of Galois representations C to
%  Kbar}T^\vee(\gM)\isoto T^\vee_{\cO_\C^\flat,pb,b}(\gM)\isoto T^\vee_{\overline{K},pb,b}(\gM).\end{equation}


Evaluating on the members
of the basis~$e_1,\dots,e_d$ of~$\gM$ yields an
identification
\[T_{L,c}^\vee(\gM)\iso\{(x_1,\dots,x_d)\in((\cO_L/p)/\mf{a}_L^{>c/p^s})^d\mid
  (x_1^p,\dots,x_d^p)=(x_1,\dots,x_d)\overline{A}\} \]
(where~$\overline{A}$ denotes the image of~$A$
%\mar{ME: I think it's better to write ``image of $\tilde{A}$'' here,
%since $A$ has entries in $\A_K^+/p$, and then $\tilde{A}$ is the thing
%with entries in $\cO_{K_{\cyc,s}}$ whose reduction mod $p$ equals
%the image of $A$.  (So  really the ideal thing to put here would
%just be that image, but we don't have notation for it!) TG: it looks
%to me like you sorted this out after writing it?}
in~$M_d((\cO_L/p)/\mf{a}_L^{>c/p^s})$), so that there is a natural map
$\tT_L^\vee(\gM)\to T_{L,c}^\vee(\gM)$. 


%\mar{TG: now explain how that maps to the mod~$p$ Galois representation.}
The following lemma is the analogue of~\cite[Lem.\ 4.1.4]{MR2745530}
in our setting, and the proof is essentially identical (and in fact
simpler, since we are only working modulo~$p$).
\begin{lem}\label{lem: amplification from char p to char 0}Assume that we are in the setting of Hypothesis~{\em\ref{hyp: K
    is basic and M has height h and s is big enough}}. Then the morphism $\tT_L^\vee(\gM)\to T_{L,b}^\vee(\gM)$ is injective,
  and has image $T_{L,pb,b}^\vee(\gM)$. % \mar{TG: this will be exactly the same as the image in
    % the one below.}
\end{lem}%\mar{TG: here we should be taking~$c=b$ etc.}
\begin{proof}
  We begin with injectivity. Suppose that we have two distinct
  tuples~$(x_1,\dots,x_d)$ and $(y_1,\dots,y_d)\in \tT_L^*(\gM)$ whose
  images in $T_{L,b}^\vee(\gM)$ coincide; then if we
  write~$z_i:=y_i-x_i$, we have $v(z_i)> b/p^s$ for each~$i$. % \mar{TG: really some
    % function of~$b$.}
  By~\eqref{eqn: Atilde
  times Btilde}, we
have \[((x_1+z_1)^p-x_1^p,\dots,(x_d+z_d)^p-x_d^p)\tB=(z_1,\dots,z_d)T_s^h.\]Writing
$(x_i+z_i)^p-x_i^p=z_i^p+pz_i(x_i^{p-1}+\dots)$, we see that
if~$z:=\min_iv(z_i)$, then we have
$eh/p^s+z\ge\min(pz, e+z)$. Since~$p^s>h$, we see that necessarily
$e + z  > eh/p^s + z$.   Hence it must be that
$eh/p^s+z\ge pz$, and thus that $eh/p^s(p-1)\ge z$. (This final deduction
is where we use our assumption that the tuples $(x_1,\dots,x_d)$
and $(y_1,\ldots,y_d)$ are distinct; this assumption ensures that $z < \infty$,
so that it is legitimate to subtract if from the two sides of an inequality.) 
However this contradicts our assumption that~$z>b/p^s$.
%\mar{ME: I need to think through where we're using that the $z_i$
%are not identically zero! TG: in this case we would have~$z=+\infty$,
%so all of these manipulations of inequalities are meaningless.}

We now turn to determining the image. By definition, the morphism $\tT_L^\vee(\gM)\to T_{L,b}^\vee(\gM)$ factors through $T_{L,pb,b}^\vee(\gM)$. Choose some $(\overline{x}_1,\dots,\overline{x}_d)\in
T^\vee_{L,pb}(\gM)$. It suffices to construct $(\tx_1,\dots,\tx_d)\in \tT_L^\vee(\gM)$
with~$\tx_i$ lifting the~$\overline{x}_i$ modulo~$\mf{a}^{>b/p^s}$ by
successive approximation.

Begin
by choosing  arbitrary lifts~$x_i$ of~$\overline{x}_i$ to~$\cO_L$.

We have
$T_s^{-h}\mf{a}_L^{>b/p^{s-1}}=\mf{a}_L^{>(pb-eh)/p^s}\subseteq
\mf{a}_L^{>b/p^s}$, so by our
assumption that $(\overline{x}_1,\dots,\overline{x}_d)\in
T^\vee_{L,pb}(\gM)$, we
have \[T_s^{-h}(x_1^p,\dots,x_d^p)\tB-(x_1,\dots,x_d)\in
  (T_s^{-h}\mf{a}_L^{>b/p^{s-1}})^d\subseteq (\mf{a}_L^{>b/p^s})^d,\]
so in fact we have
%\mar{ME: I didn't think through exactly why we replace $>b$ by $\geq b'$
%here, but did want to note an unfortunate typesetting consequence:
%$\mf{a}_{L}^{\geq b'/p^s}$
%and
%$\mf{a}_{\overline{L}}^{> b'/p^s}$
%are hard to distinguish! TG: true, but we don't use the latter
%notation at all, I guess? (although I guess your point is that
%$\overline{L}=\overline{K}$?) We could add a sentence explaining this
%or just not worry about it. I think we really need the $\ge$ version,
%or at least it would be a big headache to rewrite not to use it.}
\[T_s^{-h}(x_1^p,\dots,x_d^p)\tB-(x_1,\dots,x_d)\in
  (\mf{a}_L^{\ge b'/p^{s}})^d\]for some~$b'> b$. We
will construct Cauchy sequences~$(z^{(j)}_i) \in \mf{a}_L^{\ge b'/p^s}$ whose limits~$z_i$ are
such that taking~$\tx_i:=x_i+z_i$ gives the required lift.

To this end, we set~$z_i^{(0)}=0$, and define 
\[(z_1^{(j+1)},\dots,z_d^{(j+1)}):=T_s^{-h}((x_1+z_1^{(j)})^p,\dots,(x_d+z_d^{(j)})^p)\tB-(x_1,\dots,x_d).\]
To see that $z_i^{(j+1)}\in \mf{a}_L^{\ge b'/p^s}$, we  write
\[
  \begin{split}
    (z_1^{(j+1)},\dots,z_d^{(j+1)}):=T_s^{-h}((x_1+z_1^{(j)})^p-x_1^p,\dots,(x_d+z_d^{(j)})^p-x_d^p)\tB
    \\+(T_s^{-h}(x_1^p,\dots,x_d^p)\tB-(x_1,\dots,x_d)).
  \end{split}
\]  so we need
only check that each $(x_i+z_i^{(j)})^p-x_i^p\in
T_s^{h}\mf{a}_L^{\ge b'/p^s}$. Since~$(x_i+z_i^{(j)})^p-x_i^p\in(p
z_i^{(j)},(z_i^{(j)})^p)$, and~$z_i^{(j)}\in\mf{a}_L^{\ge b'/p^s}$, this holds.  % \mar{TG: again this is equal to something}

It remains to show that each~$(z^{(j)}_i)$ is
Cauchy. Set~$\epsilon=\min(e-eh/p^s,((p-1)b'-eh)/p^s)$, so that in particular $\epsilon
>0$. It is enough to show that $\min_i
v(z_i^{(j+1)}-z_i^{(j)})\ge \min_i
v(z_i^{(j+1)}-z_i^{(j)})+\epsilon$. To this end, we have
  \begin{multline*}
    (z_1^{(j+1)}-z_1^{(j)},\dots,z_d^{(j+1)}-z_d^{(j)})\\=T_s^{-h}\bigl((x_1+z_1^{(j)})^p-(x_1+z_1^{(j-1)})^p,\dots,(x_d+z_d^{(j)})^p-(x_d+z_d^{(j-1)})^p\bigr)\tB,
  \end{multline*}
so that it is enough to show that for each~$i$, we have
$(x_i+z_i^{(j)})^p-(x_i+z_i^{(j-1)})^p\in
(z_i^{(j)}-z_i^{(j-1)})T_s^{h}\mf{a}_L^{\ge\epsilon}=(z_i^{(j)}-z_i^{(j-1)})\mf{a}_L^{\ge
  eh/p^s+\epsilon}$.

For each~$i$, we
write \[(x_i+z_i^{(j)})^p-(x_i+z_i^{(j-1)})^p=\sum_{k=1}^p\binom{p}{k}x_i^{p-k}((z_i^{(j)})^k-(z_i^{(j-1)})^k),\]so
it is enough to note that for $1\le k\le p-1$ we have
$v(\binom{p}{k})=e \ge eh/p^s+\epsilon$, while for~$k=p$, we
write \[(z_i^{(j)})^p-(z_i^{(j-1)})^p=(z_i^{(j)}-z_i^{(j-1)})((z_i^{(j)})^{p-1}+\dots+(z_i^{(j-1)})^{p-1}),\]
and note that since~$z_i^{(j)},z_i^{(j-1)}\in \mf{a}_L^{\ge b'/p^s}$, the
second term here is contained in $\mf{a}_L^{\ge (p-1)b'/p^s}$, and we have
$(p-1)b'/p^s\ge eh/p^s+\epsilon$.
\end{proof}

%   induced by 
% \begin{lem}
%   \label{lem: relating Ainf and OKbar coefficients for Galois repn} The
%   homomorphism $V^*_{\cO_\C^\flat,pb,b}(\gM)\to V^*_{\overline{K},pb,b}(\gM)$
%   induced by \mar{TG: might not even be a lemma in the end?}
% \end{lem}


% We now argue as in the proof of Lemma~\ref{lem: Frobenius
%   amplification lemma over Ainf}. Fix some
% $f\in\Hom_{\cO_\C^\flat,\varphi}(\gM^{\inf},\cO_\C^\flat/\mf{a}_{\cO_\C^\flat}^{>b/p^{s-1}})$,
% and lift it arbitrarily to $\widetilde{f}\in
% \Hom_{\cO_\C^\flat,\varphi}(\gM^{\inf},\cO_\C^\flat)$. For any
% $\cO_{\C^\flat}$-linear map $h:\gM^{\inf}\to\cO_\C^\flat$, we write
% $\delta(h):=(h\circ\varphi-\varphi\circ
% h):\gM^{\inf}\to\cO_\C^\flat$; then the assumption on~$f$ implies that \[
%   \im(\widetilde{f}\circ\varphi-\varphi\circ f)\subseteq
%   \mf{a}_{\cO_\C^\flat}^{>b/p^{s-1}},\]and we will be done if we can
% find~$t:\gM^{\inf}\to \mf{a}_{\cO_\C^\flat}^{>b/p^s}$
% with~$\delta(t)=\delta(\widetilde{f})$ (because~$\widetilde{f}-t$ gives us the sought-after
% lift of~$f$). This follows exactly as in the proof of  Lemma~\ref{lem: Frobenius
%   amplification lemma over Ainf}.\end{proof}

% We claim in fact that given any $s:\gM^{\inf}\to
% \mf{a}_{\cO_\C^\flat}^{>b/p^{s-1}}$, we can find $t:\gM^{\inf}\to \mf{a}_{\cO_\C^\flat}^{>b+1}$
% with~$\delta(t)=s$. In fact, it suffices to show that given~$c\ge
% b$ and $s:\gM^{\inf}\to\mf{a}_{\cO_\C^\flat}^{>pc}$, we can find
% $h:\gM^{\inf}\to \mf{a}_{\cO_\C^\flat}^{>c+1}$ with
% $\im(\delta(h)-s)\subseteq \mf{a}_{\cO_\C^\flat}^{>p(c+1)}$; indeed,
% we can apply the result to~$s$, and then to $(s-\delta(h))$, and so
% on, to find~$t$ by successive approximation.

% Finally, note that for any $h:\gM^{\inf}\to
% \mf{a}_{\cO_\C^\flat}^{>c+1}$, we have $\varphi\circ h:\gM^{\inf}\to
% \mf{a}_{\cO_\C^\flat}^{>p(c+1)}$, so we only need to show that we can
% find~$h$ with 
% $\im(h\circ\varphi-s)\subseteq
% \mf{a}_{\cO_\C^\flat}^{>p(c+1)}$.\mar{TG: now apply finite height
%   assumption; not sure if I should actually be just rewriting this
%   with~$\Phi$ or indeed if it would be better to just not write this
%   argument at all, just refer to Lemma~\ref{lem: Frobenius
%   amplification lemma over Ainf}, and write out the $p$-adic
% version. Probably I will do that in fact!}
% \end{proof}
The following corollary is the analogue of~\cite[Thm.\
4.1.1]{MR2745530} in our setting.
\begin{cor}
  \label{cor: we get everything if and only if we contain the
    splitting field}Assume that we are in the setting of Hypothesis~{\em \ref{hyp: K
    is basic and M has height h and s is big enough}}, and let~$L$ be
  an algebraic extension of~$K_{\cyc,s}$ inside~$\overline{K}$, so
  that~$G_L$ acts naturally on~$T^\vee(\gM)$ by Corollary~{\em \ref{cor: we get a
    canonical action of G K s}}. 
  Then the natural injection
  $T^\vee_{L,pb,b}(\gM)\subseteq T^\vee_{\overline{K},pb,b}(\gM)$ is an
  isomorphism if and only if the action of~$G_{L}$ on~$T^\vee(\gM)$ 
  % to~$G_{K_{\cyc,s}}$ 
  is trivial.
\end{cor}
\begin{proof}By Lemma~\ref{lem: amplification from char p to char 0}
  (applied to each of~$L$ and~$\overline{K}$), the inclusion
  $T^\vee_{L,pb,b}(\gM)\subseteq T^\vee_{\overline{K},pb,b}(\gM)$ is an
  isomorphism if and only if the same is true of the inclusion
  $\tT_L^\vee(\gM)\subseteq \tT_{\overline{K}}^\vee(\gM)$. Directly
  from the definition of these latter sets, we see that
  $\tT_L^\vee(\gM)=(\tT_{\overline{K}}^\vee(\gM))^{G_L}$, so this containment
  holds if and only if~$G_L$ acts trivially on
  $\tT_{\overline{K}}^\vee(\gM)$, or equivalently, if and only if~$G_L$ acts
  trivially on~$T^\vee_{\overline{K},pb,b}(\gM)$.

  Now Lemmas~\ref{lem: truncation for Ainf mod p representation}
and~\ref{lem:G K s action} show that 
$T^\vee(\gM) \iso T^\vee_{\overline{K},pb,b}(\gM),$
and by its construction,
the $G_{K_{\cyc,s}}$-action on
$T^\vee(\gM)$ arises, via this isomorphism,
from the $G_{K_{\cyc,s}}$-action on~$T^\vee_{\overline{K},pb,b}(\gM)$.
The lemma follows.
%
%By~\eqref{eqn: iso
%    of Galois representations C to Kbar}, this holds if and only
%  if~$G_L$ acts trivially on~$T^\vee(\gM)$.  
\end{proof}

\begin{cor}
  \label{cor: the slightly weird statement about lifting
    homomorphisms} Assume that we are in the setting of Hypothesis~{\em\ref{hyp: K
    is basic and M has height h and s is big enough}}, and let $\rho$
  denote the representation $T(M)^{\vee} = T^{\vee}(\gM)$
  of~$G_{K_{\cyc,s}}$ given by Corollary~{\em\ref{cor: we get a
    canonical action of G K s}}. Let~$L$ be
  an algebraic extension of~$K_{\cyc,s}$ inside~$\overline{K}$. If
  there exists an $\cO_{K_{\cyc,s}}$-algebra homomorphism
  $\cO_{\overline{K}^{\ker \rho}}\to \cO_L/\mathfrak{a}_L^{>b/p^{s-1}}$, % \mar{TG:
    % exponent maybe off by a factor of~$p$}
  then  $\overline{K}^{\ker \rho}\subseteq L$.
\end{cor}
\begin{proof}Suppose that we have an  $\cO_{K_{\cyc,s}}$-algebra homomorphism
  $\eta:\cO_{\overline{K}^{\ker \rho}}\to
  \cO_L/\mathfrak{a}^{>b/p^{s-1}}$. We claim that~$\eta$ induces an injective
% \mar{ME: I think that here and below $\m$ should be $\mf{a}$; is that
%   right? TG: yes, fixed, commenting out.}
  homomorphism \[\cO_{\overline{K}^{\ker
      \rho}}/\mathfrak{a}_{\overline{K}^{\ker \rho}}^{>b/p^{s-1}}\to
  \cO_L/\mathfrak{a}_L^{>b/p^{s-1}}.\]

To  begin with,
let $F/K_{\cyc,s}$ be the maximal unramified subextension of $\overline{K}^{\ker
\rho}/K_{\cyc,s}$, and let $k_F$ denote its residue field (which
is then also the residue field of~$\overline{K}^{\ker \rho}$).
If $k_L$ denotes the residue field of~$L$, 
then the morphism $\eta$  induces an embedding $k_F \hookrightarrow
k_L$, which in turn lifts to an embedding $W(k_F) \hookrightarrow \cO_L$.
This embedding then induces an embedding
\numequation
\label{eqn:unram embedding}
\cO_F = \cO_{K_{\cyc,s}}\otimes_{W(k)} W(k_L)\hookrightarrow \cO_L.
\end{equation}
In particular,
we find that $\cO_L$ contains
$\cO_F$, 
or, equivalently, that  $L$  contains~$F$,
although the embedding~\eqref{eqn:unram embedding}
may not be the inclusion,
but rather the composite of this inclusion with an element
of $\Gal(F/K_{\cyc,s})$.
If we regard $\cO_{\overline{K}^{\rhobar}}$ (resp.\ $\cO_L$)
as a $\cO_F$-algebra via the inclusion (resp.\ via the embedding~\eqref{eqn:unram
embedding}),
then we see (from the very construction of~\eqref{eqn:unram embedding})
that $\eta$ is a morphism of~$\cO_F$-algebras.

Now,
let~$\pi$ be a uniformizer
  of~$\cO_{\overline{K}^{\ker \rho}}$, with
  image~$\eta(\pi)\in \cO_L/\mathfrak{a}_L^{>b/p^{s-1}}$; we claim that
  $\eta(\pi)$ is non-zero, so that its valuation is well-defined,
  and that this valuation is equal to that of~$\pi$. It follows
  immediately from this that the kernel of~$\eta$ is generated by
  $\mathfrak{a}_{\overline{K}^{\ker \rho}}^{>b/p^{s-1}}$, as required.
  %the valuation of~$\pi$ is equal to the valuation of~$\eta(\pi)$.
  To verify the claim, let~$E_\pi$ be the (Eisenstein) minimal polynomial
  of~$\pi$ over~$\cO_F$. %$\cO_{K_{\cyc,s}}$.
%\mar{ME: I guess {\em a priori} $\Kbar^{\ker \rho}$ may not be totally
%ramified  over $K_{\cyc,s}$, and so this Eisenstein poly.\ will
%have coeffs.\ in an unram.\ extension of $K_{\cyc,s}$, not nec.\ $K_{\cyc,s}$
%itself.  Probably this isn't an issue, but in the abelian case (i.e.\ when
%$\rho$ has abelian image) we can also switch the order, to write
%$\Kbar^{\ker \rho}$ as an unram.\ ext.\ of a totally ram.\ extension,
%and maybe we can just do that here?}
 Then~$\eta(\pi)$ is a root
  of~$\eta(E_\pi)$, and the non-leading coefficients of~$\eta(E_\pi)$
  all have valuation at least~$v(T_s)=e/p^s$, with equality holding
  for the constant coefficient.
  Since~$e/p^s\le eh/p^s\le (p-1)b/p^s<pb/p^s$, it follows 
%\mar{ME:Should $v(\eta)$ be $v(\pi)$?}
  that~$v(\eta(\pi))=v(\pi)$. % \mar{ME: Can we indicate why there {\em is} an induced homomorphism
%(given that $\eta$ is only a map of $\cO_{K_{\cyc,s}}$-algebras)?  I guess
%it's implicit in the computations with valuations below, but it would
%help to be explicit, since I got lost when I first tried to think
%about it. TG: you're right! Hopefully this reordering/rewriting makes
%it all OK? ME: Looks good!}

We then have a
  composite of
  injections \[ T^\vee_{\overline{K}^{\ker\rho},pb,b}(\gM)\subseteq
    T^\vee_{L,pb,b}(\gM)\subseteq T^\vee_{\overline{K},pb,b}(\gM)\](the
  second inclusion being induced by the natural map
  $\cO_L\to \cO_{\overline{K}}$). By Corollary~\ref{cor: we get everything if and only if we contain the
    splitting field} the composite is an isomorphism % source and target of this
(being an injection of isomorphic finite sets),
  so all of these inclusions are isomorphisms. Applying Corollary~\ref{cor: we get everything if and only if we contain the
    splitting field} again, we conclude that $\overline{K}^{\ker
    \rho}\subseteq L$, as required. 
\end{proof}
Recall that for each finite extension~$L$ of~$\Qp$ and each real number~$v$, we have the upper numbered
ramification subgroups $G_L^v$, as defined in~\cite[Chapter IV]{SerreLF}; these are the groups denoted~$G_L^{(v+1)}$
in~\cite{MR807070} (see~\cite[Rem.\ 1.2]{MR807070}). We can now prove our
first bound on the ramification of~$T(M)$, from which our subsequent
bounds will be deduced.
\begin{cor}
  \label{cor: ramification bound in the basic mod p case}Assume that we are in the setting of Hypothesis~{\em\ref{hyp: K
    is basic and M has height h and s is big enough}}. % and let $\rho$
  %be the representation of~$G_{K_{\cyc,s}}$ given by Corollary~{\em\ref{cor: we get a
  %  canonical action of G K s}}.
%If~$v\ge pb-1$ then~$\rho(G_{K_{\cyc,s}}^v)=\{1\}$.
%\mar{ME: I looked at the Yoshida reference, and it seemed to me
%that one gets $v > pb - 1$ rather than $v \ge pb-1$. TG: agreed,
%fixed; I decided to just change to $\ge pb$ but if you'd rather do
%$v>pb$ I don't mind. ME: Made that change}
If~$v> pb-1$, then~$G_{K_{\cyc,s}}^v$ acts   trivially
on~$T(M)$.
%   In particular, if~$v\ge pb-1$ then the action
% \mar{ME: Why are we singling out this intersection? TG: I have no
%   idea, and this sentence doesn't seem to be doing anything helpful,
%   so I'm commenting it out.}
%   of~$G_{\Kcyc}\cap G_{K_{\cyc,s}}^v$ on~$V(M)$ is trivial.
\end{cor}%\mar{TG: probably these two corollaries can now be
 % merged. ME: Good idea TG: we should just be able to delete the
 % second one, I think, but I won't do it yet.}
\begin{proof}
Let $\rho$ denote $T(M)^{\vee}$. 
Recall that for a real number~$m$, the extension $K^{\ker \rho}/K_{\cyc,s}$
satisfies $(\mathrm{P}_m)$ if, for any subfield $L$  of $\overline{K}$
containing $K_{\cyc,s}$,
the existence of an $\cO_{K_{\cyc,s}}$-algebra homomorphism
$\cO_{K^{\ker \rho}} \to \cO_L/\mf{a}^{\ge m/p^s}$ implies that
$K^{\ker \rho} \subseteq L$.
(This property was introduced by Fontaine,
and is discussed extensively in~\cite{MR2814776},
to which we refer the reader for further explanations
and references.  Note that the denominator
$p^s$ appears in $\mf{a}^{\geq m/p^s}$ because we always compute valuations
with respect to~$K$, and $K_{\cyc,s}$ is a totally ramified
extension of $K$ of degree~$p^s$.)

Corollary~\ref{cor: the slightly
    weird statement about lifting homomorphisms}
shows that $\inf\{m \, | \, (\mathrm{P}_m) \text{ holds for  }
K^{\ker \rho}/K_{\cyc,s} \} \leq pb.$
The present corollary then follows from this,
together with~\cite[Thm.\ 1.1]{MR2814776}.
%This is immediate from Corollary~\ref{cor: the slightly
%    weird statement about lifting homomorphisms} and~\cite[Thm.\ 1.1]{MR2814776}.
\end{proof}%\mar{TG: again need to resolve~$a$ versus~$n$ now.}

% We now combine Corollary~\ref{cor: ramification bound in the basic mod
%   p case} with the canonical actions on $(\varphi,\Gamma)$-modules
% that we considered in Section~\ref{subsec:
%   canonical weak Wach}.

% \begin{cor}
%   \label{cor: ramification bound in the basic mod p case with
%     coefficients}Assume that~$K$ is basic. Let~$A$ be a finite local
%   Artinian $k$-algebra, and let $\gM$ be a projective
%   $\varphi$-module over~$\A_{K,A}^+$ of $T$-height at
%   most~$h$.
% \mar{ME: Just to note that these hypotheses on $N$ and $s$ are now
% part of Hyp.\ 7.3.1.}
% Suppose that $N\ge N(1,h)$ and $s\ge s(1,h,N)$ are as in
%   Corollary~{\em\ref{cor: Caruso Liu Galois action phi Gamma discrete
%     version}}, so that in particular we obtain a
%   $G_{K_{\cyc,s}}$-representation $T_A(\gM[1/T])$ from  the canonical
%   $(\varphi,\Gamma_{K_{\cyc,s}})$-structure on $\gM[1/T]$. 

%   Then if~$N,s$ are sufficiently large \emph{(}depending only on~$K$ and~$h$\emph{)}, there is some~$v$ depending only on~$K,h,N$ and~$s$ such that
%   the action of~$G_{\Kcyc,s}^{v}$ on  $T_A(\gM[1/T])$ is trivial.
% \end{cor}%\mar{TG: again need to resolve~$a$ versus~$n$ now.}
% \begin{proof}This follows from Corollary~\ref{cor: ramification
%     bound in the basic mod p case}, noting that by definition (for suitably
%   large~$N$ and~$s$) the two canonical actions of~$G_{K_{\cyc,s}}$ on
%   $T^\vee(\gM[1/T])$ (one coming from Corollary~\ref{cor: Caruso Liu Galois action phi Gamma discrete
%     version}, and one from Corollary~\ref{cor: we get a canonical action of G K
%     s}) coincide.  \mar{TG: this coincidence is hopefully explained above.}
% \end{proof}

% \mar{TG: my hope is that this can now be rewritten; we can't get rid
%   of the $N(a,h)$ etc so we may need to make that more complicated (!)
% to be compatible with the Hypothesis above, but hopefully we can work
% on $T(M)$ instead of $M$ throughout?}

We now return to the setting of a general~$K$ (i.e.\ we do not assume
that~$K$ is basic). % , and we allow~$A$ to be an arbitrary finite Artinian~$\cO/\varpi^a$-algebra. % ~$A$ be a finite Artinian
% $\cO/\varpi^a$-algebra for some $a\ge 1$.
Let~$A$ be a finite local Artinian $\cO/\varpi^a$-algebra for
some~$a\ge 1$, and let~$M$ be a rank one projective \'etale
$\varphi$-module over~$\A_{K,A}$, which we assume to be of
$(\Kbasic,T)$-height at most~$h$ in the sense of Definition~\ref{defn:
  Kbasic T height at most h}; that is, we regard~$M$ as a
rank~$[K:\Kbasic]$ \'etale $\varphi$-module~$M'$
over~$\A_{\Kbasic,A}$, and we assume there is a projective
$\varphi$-module~$\gM$ over~$\A_{\Kbasic,A}^+$ of $T$-height at
most~$h$ such that $M'=\gM[1/T]$.


%We now choose: %integers $b,s,N$ satisfying the following conditions.
%
%  \begin{itemize}
%  \item an integer~$b\ge eh/(p-1)$. 
%  \item An integer ~$N> pb/e$ such \mar{TG: obviously the  $a$ version
%      should imply the $1$ version, but it seems safest to make this
%      explicit, as we didn't fix choices of the $N(a,e)$ etc. Remove
%      this comment if you're happy!}
%    that $N\ge N(a,h)$ and $N\ge N(1,h)$, where $N(a,h)$ and $N(1,h)$ are as in Corollary~\ref{cor: Caruso
%      Liu Galois action phi Gamma discrete version}, with~$K$ there
%    being replaced by~$\Kbasic$.
%  \item An integer~$s$ such that~$p^{s-1}>b/e$,  $s\ge s(a,h,N)$, and
%    $s\ge s(1,h,N)$
%    where  $s(a,h,N)$ and $s(1,h,N)$ are as in Corollary~\ref{cor: Caruso
%      Liu Galois action phi Gamma discrete version}, with~$K$ there
%    being replaced by~$\Kbasic$.
%  \end{itemize}
%
%Note in particular that Hypothesis~\ref{hyp: K is basic and M has height h and s
%  is big enough} holds, with~$K$ in the hypothesis replaced
%by~$\Kbasic$. 

We have the $G_{\Kbasic_{\cyc}}$-representation $T_A(M')$, which we may
compute as \[T_A(M')=(W(\C^\flat)\otimes_{\A_{\Kbasic}}M')^{\varphi=1}=(W(\C^\flat)\otimes_{\A_{\Kbasic}}M)^{\varphi=1}.\]
It follows that there is a $G_{K_{\cyc}}$-equivariant surjection
\numequation\label{eqn: equivariant surjection from T(M') to T(M)}T_A(M')=(W(\C^\flat)\otimes_{\A_{\Kbasic}}M)^{\varphi=1}\onto
(W(\C^\flat)\otimes_{\A_K}M)^{\varphi=1}=T_A(M).\end{equation}
% \mar{TG: hopefully this is OK, so we don't have to get into induced
%   representations and so on. ME: I think it's fine; I added a paranthetical
% comment.}
In fact~\eqref{eqn:induction via phi mods}
shows that
$T_A(M')$ can be identified with $\Ind_{G_{K_{\cyc}}}^{G_{\Kbasic_{\cyc}}} T_A(M)$,
and the morphism~\eqref{eqn: equivariant surjection from T(M') to T(M)} then 
becomes the natural $G_{K_{\cyc}}$-equivariant surjection.  


%Note that our assumptions imply in particular that the hypotheses of Corollary~\ref{cor: Caruso
%      Liu Galois action phi Gamma discrete version} hold, so that the 
%    action of~$G_{\Kbasic_{\cyc}}$ on~$T(M')$ has a canonical
%    extension to an action of $G_{\Kbasic_{\cyc,s}}$.
%   Suppose that $N\ge N(a,h)$ and $s\ge s(a,h,N)$
% are as in Corollary~\ref{cor: Caruso Liu Galois action phi Gamma
%   discrete version} (with~$K $ replaced by~$\Kbasic$),
% % \mar{ME: Is $M'$ notation for the restriction of $M$ to $\Kbasic$? TG:
% % yes, fixed, commenting out.}
% so that in particular if as usual we regard~$M$ as a rank~$[K:\Kbasic]$ \'etale $\varphi$-module~$M'$
%   over~$\A_{\Kbasic,A}$, then~$M'$ obtains the structure
% of an \'etale $(\varphi,\Gamma_{\Kbasic_{\cyc,s}})$-module. Then the
% rank~$[K:\Kbasic]$ projective \'etale $(\varphi,\Gamma_{K_s})$-module
% $\A_{K,A}\otimes_{\A_{\Kbasic,A}}M'$ decomposes as a direct sum of \'etale
% $(\varphi,\Gamma_{K_s})$-modules, with the action
% of~$\Gamma_{K_{\cyc,s}}$ permuting the summands.
% One of these summands is our original module $M$ itself,
% and so after
% increasing~$s$ in a manner depending only on~$K$, we can assume that
% $M$ is preserved by the $\Gamma_{K_{\cyc,s}}$-action,
% so that~$M$ has the canonical structure of a
% $(\varphi,\Gamma_{\Kcyc,s})$-module. 

\begin{prop}
  \label{prop: ramification bound in rank one case with coefficients
    and any K}Suppose that~
  $A$ is an Artinian $\cO/\varpi^a$-algebra for some~$a\ge 1$, and $M$
  is a rank 1 projective \'etale $\varphi$-module over~$\A_{K,A}$ of
  $(\Kbasic,T)$-height at most~$h$.% , and assume further that
  % Hypothesis~\ref{hyp: K is basic and M has height h and s is big
  %   enough} holds \emph{(}with $K$ replaced by~$\Kbasic$\emph{)}.
  
  Then % there is some~$v$ depending only on~$K,a,h,N$ and~$s$ such that
  % the action of~$G_{\Kcyc,s}^{v}$ on $T_A(M)$ is trivial. In
  % particular, 
  there is an open subgroup~$I^{a,h}$
  of~$I_{\Kcyc}^{\ab}$ % \mar{TG: this is maybe not what we want! I
    % don't remember right now for the applications whether it's enough for this result to
    % just consider an open subgroup of $I_{\Kcyc}^{\ab}$. It's OK in
    % the first application but probably not in the second, where we
    % also have an $I_K$-action?} 
  % \mar{TG: is this defined?}
  depending only on
  $K,a$ and~$h$ \emph{(}and not on~$A$ or~$M$\emph{)} such that the action of ~$I^{a,h}$ on
  $T_A(M)$ is trivial.
\end{prop}%\mar{TG: again need to resolve~$a$ versus~$n$ now.}
\begin{proof}
%We begin by recalling the relationship between
%upper numbered ramification groups and the filtration of the unit
%group. Let~$L$ be a finite extension of~$\Qp$. Recall that the Artin
%map identifies~$\cO_L^\times$ with~$I_L^{\ab}$. Furthermore, for each integer~$v\ge 1$
%the Artin map identifies the subgroup~$1+(\m_L)^v$ of~$\cO_L^\times$
%with the image of~$G_L^v$ in~$G_L^{\ab}$ (see~\cite[Chapter XV, \S 2,
%Thm.\ 2]{SerreLF}). In particular, for each~$v$ the image of $G_L^v$
%in~$I_L^{\ab}$ is open (and in particular has finite index).
%
We choose integers $b,s,N$ satisfying the following conditions:

  \begin{itemize}
  \item $b\ge eh/(p-1)$. 
  \item $N> pb/e$ is such % \mar{TG: obviously the  $a$ version
      %should imply the $1$ version, but it seems safest to make this
      %explicit, as we didn't fix choices of the $N(a,e)$ etc. Remove
      %this comment if you're happy!}
    that $N\ge N(i,h)$ for each $i \in [1,a]$, % and $N\ge N(1,h)$,
where the $N(i,h)$ %and $N(1,h)$
are as in Corollary~\ref{cor: Caruso
      Liu Galois action phi Gamma discrete version}, with~$K$ there
    being replaced by~$\Kbasic$.
  \item $s$ is such that~$p^{s-1}>b/e$,  and such that
$s\ge s(i,h,N)$ for each $i \in [1,a]$,
%and $s\ge s(1,h,N)$
    where  the $s(i,h,N)$ %and $s(1,h,N)$
are as in Corollary~\ref{cor: Caruso
      Liu Galois action phi Gamma discrete version}, with~$K$ there
    being replaced by~$\Kbasic$.
  \end{itemize}
Note in particular that Hypothesis~\ref{hyp: K is basic and M has height h and s
  is big enough} holds, with~$K$ in the hypothesis replaced
by~$\Kbasic$. 
Note also that our assumptions imply that the hypotheses of Corollary~\ref{cor: Caruso
      Liu Galois action phi Gamma discrete version} hold (with the field
$K$ there taken to be~$\Kbasic$), so that the 
    action of~$G_{\Kbasic_{\cyc}}$ on~$T_A(M')$ has a canonical
    extension to an action of $G_{\Kbasic_{\cyc,s}}$.
This canonical action on $T_A(M')$ induces 
the canonical action on
$\varpi^i T_A(M') = T_{A}(\varpi^i M')$ %\mar{TG: isn't
 % $T_A(\varpi^i M')$ still OK here? Or if you want to change it, I
 % guess $T_{A/\varpi^{a-i}}(\varpi^i M')$? ME: I'll just leave it as $A$ }
for each $i \in [0,a],$
and hence on each of the quotients
$\varpi^{i-1} T_A(M') /\varpi^i T_A(M') = T_{A/\varpi}(\varpi^{i-1} M'/\varpi^iM'),$
for $i \in [1,a]$.

%Choose integers $b,N,s$ satisfying the hypotheses above.
To ease notation,
write~$M_i:=\varpi^{i-1}M/\varpi^iM$, for $1\le i\le a$,
and let $M_i'$ denote $M_i$ 
% Let~$M'$ (resp.\ $M'_i$) denote~$M$ (resp.\ $M_i$)
regarded as an \'etale
$\varphi$-module over~$\A_{\Kbasic,A/\varpi}$. By Corollary~\ref{cor:
  ramification bound in the basic mod p case}, %\mar{TG: say something
  %about compatibility with $\varpi$-adic filtration to justify this?}
if~$v$ is sufficiently large (depending only on~$K$, $a$ and~$h$, and
our subsequent choices of~$b$, $N$ and~$s$, and in particular not
on~$A$ or~$M$), then the canonical action of~$G_{\Kbasic_{\cyc,s}}^v$
on each~$T_{A/\varpi}(M'_i)$ is trivial. It follows that there is
some~$m$ depending only on~$a$ (and the degree $[K:\Kbasic]$) such that for
any~$g\in G_{\Kbasic_{\cyc,s}}^v$, the action of~$g^{p^m}$
on~$T_A(M')$ is trivial.
%\mar{TG: before this was written in a way
%  that seemed to use that the action was abelian, but I don't think it does?}

  It follows from the definitions that 
  $G_{\Kbasic_{\cyc,s}}^v\cap
  G_{K_{\cyc,s}}=G_{K_{\cyc,s}}^{\psi_{K_{\cyc,s}/\Kbasic_{\cyc,s}}(v)}$,
  where~$\psi_{K_{\cyc,s}/\Kbasic_{\cyc,s}}$ is as in~\cite[Chapter IV]{SerreLF}), so if
  we take any~$w\ge \psi_{K_{\cyc,s}/\Kbasic_{\cyc,s}}(v)$,
then~$(G_{K_{\cyc,s}}^w)^{p^m}$ 
($ := \{g^{p^m} \, | \,
g\in G_{K_{\cyc,s}}^w\}$; note that this is just a subset, not a subgroup,
of~$G_{K_{\cyc,s}}$)
acts trivially on $T_A(M')$. 

The product $G_{K_{\cyc}}
(G_{K_{\cyc,s}}^w)^{p^m}$ 
is {\em a priori} a subset of $G_K$, but we claim that it is in fact
%\mar{ME: Add justification here.}
an open subgroup of~$G_K$, and thus (since it contains
$G_{K_{\cyc}}$) is equal to $G_{K_{\cyc,s'}}$  for
some $s' \geq s$ which depends only on~$s$ and $m$
(and so ultimately only on  $K$, $a$ and $h$).

To prove the claim, 
we briefly recall the relationship between
upper numbered ramification groups and the filtration of the unit
group: If~$L$ is a finite extension of~$\Qp$, then the Artin
map identifies~$\cO_L^\times$ with~$I_L^{\ab}$, and for each integer~$w\ge 1,$
identifies the subgroup~$1+(\m_L)^v$ of~$\cO_L^\times$
with the image of~$G_L^v$ in~$G_L^{\ab}$ (see~\cite[Chapter XV, \S 2,
Thm.\ 2]{SerreLF}).   It then identifies the subgroup $(1+(\m_L)^v)^{p^m}$
with the image of~$(G_L^v)^{p^m}$ in~$G_L^{\ab}$.  This former
subgroup is an open subgroup
of $\cO_L^{\times}$,
and so we find that for each~$v$, the image of $(G_L^v)^{p^m}$
in~$I_L^{\ab}$ is open.

Now the extension $K_{\cyc} / K_{\cyc,s}$ is an infinitely ramified
$\Gamma_{K_{\cyc,s}} \cong \Z_p$-extension,
and so the inertia group $I(K_{\cyc}/K)$ is finite index (equivalently,
open) in~$\Gamma_{K_{\cyc,s}}$.   Since $I(K_{\cyc}/K)$ is a quotient
of~$I_K^{\ab}$, the discussion of the preceding paragraph then shows that 
the image of
$(G_{K_{\cyc,s}}^w)^{p^m}$ 
in $\Gamma_{K_{\cyc,s}}$  is also an open subgroup. 
Since $\Gamma_{K_{\cyc,s}}$ is itself open in~$G_K$,
this implies the claim.

Since~\eqref{eqn: equivariant surjection from T(M') to T(M)}
is $G_{K_{\cyc}}$-equivariant, and since 
$(G_{K_{\cyc,s}}^w)^{p^m}$ 
acts trivially on~$T_A(M')$,
we see that the kernel 
of~\eqref{eqn: equivariant surjection from T(M') to T(M)}
is in fact~$G_{K_{\cyc,s'}}$-invariant, and thus that
the $G_{K_{\cyc}}$-action on $T_A(M)$ extends to a $G_{K_{\cyc,s'}}$-action,
which is abelian (since the $G_{K_{\cyc}}$-action
is abelian, as $T_A(M)$ has rank one over~$A$, while
$(G_{K_{\cyc,s}}^w)^{p^m}$ 
acts trivially). 

Now another application of~\cite[Chapter IV]{SerreLF}) shows that
$G_{K_{\cyc,s}}^w \cap \, G_{K_{\cyc,s'}} =
G_{K_{\cyc,s'}}^{w'}$ 
for some  $w'$ depending only on~$w$, and our preceding discussion
of abelian ramification theory shows that
$(G_{K_{\cyc,s'}}^{w'})^{p^m}$ 
has open image in $G_{K_{\cyc,s'}}^{\ab}$.
If we let $I'$  denote the preimage of this open image
(i.e.\  \[I' = \overline{[G_{K_{\cyc,s'}},G_{K_{\cyc,s'}}]} 
(G_{K_{\cyc,s'}}^{w'})^{p^m},\] where the overline
denotes closure in~$G_{K_{\cyc,s'}}$),
then $I'$ is an open subgroup  of $G_{K_{\cyc,s'}}$ which
acts trivially on~$T_A(M)$.
Thus $I'' := I' \cap G_{K_{\cyc}}$ also acts trivially on~$T_A(M)$.
Finally, we let $I^{a,h}$ denote the image of $I''$ in $I_K^{\ab}$.
(As the notation indicates, $I^{a,h}$  depends only on $a$,  $h$, and~$K$,
but not on~$A$ or~$M$, since this is true of each of the groups
$I'$ and $I''$ from which it is constructed.)
%
%
%Since~$M$ has rank one, the action of~$G_{K_{\cyc}}$ on~$T_A(M)$ is
%  abelian.  Let~$G^{a,h}:= G_{\Kbasic_{\cyc,s}}^v\cap G_{K_{\cyc}}$.    It follows from~\eqref{eqn: equivariant surjection from
%    T(M') to T(M)} that if~$g\in G^{a,h}$, then~$g^{p^m}$ acts
%  trivially on $T_A(M)$. Let~$I^{a,h}$ be the image
%  in~$I_{K_{\cyc}}^{\ab}$ of the set of~$p^m$th powers of elements
%  of~$G^{a,h}$; by construction, this is a subgroup of
%  ~$I_{K_{\cyc}}^{\ab}$ which acts trivially on~$T_A(M)$, and it only
%  remains to show that it is in fact an open subgroup.\mar{TG:
%    hopefully this is easy to finish! I've left in some older
%    arguments below. Maybe we should be constructing this group as a
%    preimage of a subgroup of $I_K^{\ab}$, without having to
%    make~$I_K^{\ab}$ itself act?}
% %  each~$T_{A/\varpi}(M'_i)$ is also % \mar{TG: not sure if we want to go
%%   % to here yet?}
%% trivial.
%
%
%
% , so that % Since~$M$ is a direct summand of
%   $\A_{K}\otimes_{\A_{\Kbasic}}M'$, and since
%   $K\cdot\Kbasic_{\cyc,s}=K_{\cyc,s}$, we see that in particular the
%   representation~$V(M)$ admits an extension
% \mar{ME: I'm confused about this argument: it looks like we're claiming
% that if $W$ is a rep'n of a group $G$, and if $H$ is a subgroup of $G$,
% then any $H$-invariant direct summand of $W$ is automatically $G$-invariant.  
% But probably I'm misunderstanding! TG: this is the mistake that I
% explained to you at breakfast\dots}
% to a representation
%   of~$G_{K_{\cyc,s}}$, with
   % the action of ~$G_{\Kbasic_{\cyc,s}}^v\cap G_{K_{\cyc,s}}$ on
   % each~$T_{A/\varpi}(M_i)$ is trivial. \mar{TG: again we're using
   %   $G_{K_{\cyc,s}}$ and we probably don't want to.}
  % to~$G_{\Kbasic_{\cyc,s}}^v\cap G_{K_{\cyc,s}}$ is trivial. 
  % (See~\cite[Rem.\ A.3.4.4]{MR1106901} for the relationship between the actions of
  % $G_{\Kbasic_{\cyc}}$ on~$M'$ and $G_{\Kcyc}$ on
  % $\A_{K}\otimes_{\A_{\Kbasic}}M'$.) 
%
%
%  It follows from the definitions that we \mar{TG: this is an old argument}
%  have~$G_{\Kbasic_{\cyc,s}}^v\cap
%  G_{K_{\cyc,s}}=G_{K_{\cyc,s}}^{\psi_{K_{\cyc,s}/\Kbasic_{\cyc,s}}(v)}$,
%  where~$\psi_{K_{\cyc,s}/\Kbasic_{\cyc,s}}$ is as in~\cite[Chapter IV]{SerreLF}), so if
%  we take any~$w\ge \psi_{K_{\cyc,s}/\Kbasic_{\cyc,s}}(v)$, then~$G_{K_{\cyc,s}}^w$ acts
%  trivially on each~$T_{A/\varpi}(M_i)$.
%
%
% % We
% %  claim that there is some~$w$ depending only on~$K$ and~$h$ such that
% %  ~$G_{K_{\cyc,s}}^w$ acts trivially on each~$V(M_i)$. To see this, it is
% %  enough to show that we can choose~$w$ so
% %  that~$G_{\Kbasic_{\cyc,s}}^v\cap G_{K_{\cyc,s}}\supseteq
% %  G_{K_{\cyc,s}}^w$. It follows from the definitions that we
% %  have~$G_{\Kbasic_{\cyc,s}}^v\cap
% %  G_{\Kcyc}=G_{K_{\cyc,s}}^{\psi_{K/\Kbasic}(v)}$,
% %  where~$\psi_{K/\Kbasic}$ is as in~\cite[Chapter IV]{SerreLF}), so we
% %  can take any $w\ge\psi_{K/\Kbasic}(pb-1)$. 
%  % \mar{TG: probably some
%    % very inconsistent notation} 
%
%  % Since~$M$ has rank one, the action of~$G_{K_{\cyc,s}}$ on~$T_A(M)$ is
%  % abelian. 
%  % Since $G_{K_{\cyc,s}}^w$ acts trivially on each~$T_{A/\varpi}(M_i)$, there is
%  % some~$m$ depending only on~$a$ such that for
%  % any~$g\in G_{K_{\cyc,s}}^w$, the action of~$g^{p^m}$ on~$V(M)$ is
%  % trivial.
%  Since the $p^m$-th power of an open subgroup \mar{TG: old argument;
%    I don't think we can quite use this as written.}
%  of~$\cO_{K_{\cyc,s}}^\times$ is open, it follows from the compatibility of the
%  upper ramification numbering with the filtration on~$\cO_{K_{\cyc,s}}^\times$
%  that if we increase~$w$ by a quantity depending only on~$K,w$
%  and~$a$, we can assume that~$G_{K_{\cyc,s}}^w$ acts trivially
%  on~$T_A(M)$. Since the image of~$G_{K_{\cyc,s}}^w$ in~$I_{K_{\cyc,s}}^{\ab}$
%  is open, the intersection of this image with~$I_{K_{\cyc,s}}^\ab$ gives
%  us the required open subgroup~$I^{K,a,h,N,s}$.% \mar{TG: I think
%    % this is a bit bogus as written.}
\end{proof}

The same argument allows us to prove the following variant of Proposition~\ref{prop: ramification bound in rank one case with coefficients
  and any K}.


\begin{prop}
  \label{prop: ramification bound in rank one case with coefficients
    and any K with GK action}Suppose that
  $A$ is an Artinian $\cO/\varpi^a$-algebra for some~$a\ge 1$, and $M$
  is a rank 1 projective \'etale $(\varphi,\Gamma)$-module
  over~$\A_{K,A}$. Assume furthermore that~
  $M$ has $(\Kbasic,T)$-height at most~$h$, and that ~$s$ is some
  sufficiently large integer with the property that
%\mar{TG: hypothesis for canonical action ME: I think this comment
%can be removed, but I wanted to be sure I hand't accidentally omitted some
%additional hypothesis that was supposed to be added! TG: I think
%something is needed, in fact, because we need some hypotheses that
%guarantee that there is in fact a canonical action! i.e.\ $s$ needs to
%be sufficiently large. I added something, although strictly speaking
%one should introduce $\gM$, $N$ and so on, I guess; but I'm inclined
%not to do so, as I don't think it will really help anyone read
%this. Comment out if you agree\dots}
the action of~$\Gamma_{\Kbasic,\disc}=\Gamma_{K,\disc}$ on $M$ extends the canonical action of
    $\Gamma_{\Kbasic_s,\disc}$ given by Corollary~\emph{\ref{cor: Caruso Liu Galois
      action phi Gamma discrete version}}. {\em (}In particular, we are
  assuming that~$s$ has been chosen to be large enough that this
  canonical action is defined.{\em )}
  
  Then there is an open subgroup~$I_K^{a,h,s}$ of~$I_{K}^{\ab}$
  depending only on $K,a,h$ and~$s$ \emph{(}and not on~$A$
  or~$M$\emph{)} such that the action of ~$I_K^{a,h,s}$ on $T_A(M)$ is
  trivial.
\end{prop}%\mar{TG: again need to resolve~$a$ versus~$n$ now.}
\begin{proof}
%We follow the proof
The proof is very similar to that
of Proposition~\ref{prop: ramification bound in rank one
case with coefficients and any K},
and we content ourselves with indicating the key modifications.
The main difference between the setting of that proposition and
the present one is that, in our present setting,
the morphism~\eqref{eqn: equivariant surjection from T(M') to T(M)}
is a $G_K$-equivariant morphism from a $G_{\Kbasic}$-representation
to a $G_K$-representation.  Furthermore, by assumption,
the $G_{\Kbasic_{\cyc,s}}$-action on its source is the canonical action.  

We now follow the argument 
of Proposition~\ref{prop: ramification bound in rank one
case with coefficients and any K}, and find that,
after possibly enlarging $s$ in a manner depending only on $a$, $h$,
and~$K$, some ramification group $G_{K_{\cyc,s}}^w$ acts
trivially on~$T_A(M')$ (the index $w$ also depending
only on $a$, $h$, and~$s$).    
In the present argument, there is no need to pass
to an auxiliary subgroup $G_{K_{\cyc,s'}}$, since
the morphism~\eqref{eqn: equivariant surjection from T(M') to T(M)}
is in particular already a $G_{K_{\cyc,s}}$-equivariant morphism
of $G_{K_{\cyc,s}}$-representations;
and the action of $G_{K_{\cyc,s}}$ on $T(M)$ is already abelian.
The argument then proceeds in the same manner,
to produce an open subgroup $I'$
of the inertia subgroup of $G_{K_{\cyc,s}}$ 
(and thus of $I_K$) which acts trivially on~$T_A(M)$.
(And, by construction, the group $I'$ depends only on $a$, $h$,
and~$s$, but not on~$A$ or~$M$.)
We then let
$I_K^{a,h,s}$ 
denote the image of $I'$ in~$I_K^{\ab}$.
%
%
%In particular after possibly enlarging~$s$ (in a fashion
%  independent of~$A$ and~$M$), we again find that if~$v$, $m$ are
%  sufficiently large (again, independently of~$A$ and~$M$), then for
%  each $g\in G_{\Kbasic_{\cyc,s}}^v\cap G_K$, the action of ~$g^{p^m}$
%  on~$T_A(M)$ is trivial.\mar{TG: again, hopefully easy to finish from
%  here?}
  \end{proof}

 


% As recalled in the proof of
% Proposition~\ref{prop:rank 1 R}, given an \'etale $\varphi$-module~$M$
% over~$\A_{K,A}/p^n$ of some rank~$d$, we can regard it as an \'etale $\varphi$-module
%  over~$\A_{\Kbasic}/p^n$ of rank~$d[K:\Kbasic]$, and as
% such it makes sense to ask whether there is a $\varphi$-module~$\gM$
% over~$\A_{\Kbasic}^+/p^n$ of $T$-height at most~$h$ such that
% $M=\gM[1/T]$; if this holds, we say that~$M$ has $(\Kbasic,T)$-height
% at most~$h$.\mar{TG: I now use this terminology earlier, rearrange it
%   to be in the correct order.}

% We wish to prove a ramification bound for
% the~$G_{\Kcyc}$-representation ~$V(M)$ in terms of~$n$ and~$h$. We
% will do this under the further assumption that the action
% of~$G_{\Kcyc}$ on~$V(M)$ factors through the maximal abelian extension
% of~$\Kcyc$. Before doing this,

% \mar{TG:
  % I think I'm  making a bit of a mess of this and related statements!}
% \begin{cor}
%   \label{cor: mod p discriminant bound in the non basic
%     case}Assume that~$n=1$, but do not assume that~$K$ is basic.  \mar{TG:
%   still struggling a bit to get the right statements here, is ``open
%   subgroup'' what we want? In the end it should all be fine because we
% really have that all the characters actually come (up to unramified
% twist) from characters of a finite extension of finite extensions
% of~$\Qp$, but I'm not sure if we actually need that. If we do then
% these statements can be easily rewritten.} Fix some integer~$h\ge 1$. Then there is an open subgroup~$I^{K,h,\ab}_{\Kcyc}$ of~$I_{\Kcyc}^\ab$ depending
% only on~$K,h$ such that 
% if $M$ has $(\Kbasic,T)$-height
% at most~$h$ for some~$h\ge 1$, and the action of~$G_{\Kcyc}$
% on~$V(M)$ factors through~$G_{\Kcyc}^\ab$, then the action of~$I^{K,h,\ab}_{\Kcyc}$ on~$V(M)$ is trivial.
% \end{cor}

% \begin{proof}Choose
%   integers~$b$, $s$ as in Hypothesis~\ref{hyp: K is basic and M has
%     height h and s is big enough} with~$K$ replaced by~$\Kbasic$ (so
%   that in particular~$b$ and~$s$ depend only on~$K$ and~$h$).
% Let~$M'$ denote~$M$ regarded as an \'etale $\varphi$-module
% over~$\A_{\Kbasic}$. By Corollary~\ref{cor: ramification bound in
%   the basic mod p case}, if~$v\ge pb-1$ then the action
%   of~$G_{\Kbasic_{\cyc,s}}^v$ on~$V(M')$ is
%   trivial. Since~$M$ is a direct summand of
%   $\A_{K}\otimes_{\A_{\Kbasic}}M'$, and since
%   $K\cdot\Kbasic_{\cyc,s}=K_{\cyc,s}$, we see that in particular the
%   representation~$V(M)$ admits an extension to a representation~$\rho$
%   of~$G_{K_{\cyc,s}}$, with the property that the restriction of~$\rho$
%   to~$G_{\Kbasic_{\cyc,s}}^v\cap G_{K_{\cyc,s}}$ is trivial. 
%   (See~\cite[Rem.\ A.3.4.4]{MR1106901} for the relationship between the actions of
%   $G_{\Kbasic_{\cyc}}$ on~$M'$ and $G_{\Kcyc}$ on
%   $\A_{K}\otimes_{\A_{\Kbasic}}M'$.)

% Since~$\Kcyc$ is an abelian extension of~$K_{\cyc,s}$,
% and~$\rho|_{G_{\Kcyc}}$ has abelian image by assumption, we see
% that~$\rho$ also has abelian image. We claim that there is some~$w$
% depending only on~$K$ and~$h$ such that
% ~$\rho(G_{K_{\cyc,s}}^w)=\{1\}$. To see this, it is enough to show that
% we can choose~$w$ so that~$G_{\Kbasic_{\cyc,s}}^v\cap G_{K_{\cyc,s}}\supseteq
% G_{K_{\cyc,s}}^w$. It follows from the definitions that we
% have~$G_{\Kbasic_{\cyc,s}}^v\cap
% G_{\Kcyc}=G_{K_{\cyc,s}}^{\psi_{K/\Kbasic}(v)}$, where~$\psi_{K/\Kbasic}$ is as
% in~\cite[Chapter IV]{SerreLF}), so we can take any
% $w\ge\psi_{K/\Kbasic}(pb-1)$.

% Since the image of~$G_{K_{\cyc,s}}^w$ in~$I_{K_{\cyc,s}}^{\ab}$ is open, the
% intersection of this image with~$I_{\Kcyc}^\ab$ gives us the required
% open subgroup~$I^{K,h,\ab}_{\Kcyc}$.
% \end{proof}

%Finally, we deduce our ramification bound for general~$n$.
% \begin{prop}
%   \label{prop: ramification bound for non basic abelian case any a}
% Let~$K/\Qp$ be an arbitrary finite extension. Fix
%   some~$h,n\ge 1$. Then there is an open subgroup~$I^{K,h,n}_{\Kcyc}$ of~$I_{\Kcyc}^{\ab}$ depending
% only on~$K,h,$ and~$n$ such that 
% if $M$ is a free
% \'etale $\varphi$-module over~$\A_{K}/p^n$ which has $(\Kbasic,T)$-height
% at most~$h$ for some~$h\ge 1$, and the action of~$G_{\Kcyc}$
% on~$V(M)$ factors through~$G_{\Kcyc}^\ab$, then the action of~$I^{K,h,n}_{\Kcyc}$ on~$V(M)$ is trivial.
% \end{prop}
% \begin{proof}Choose
%   integers~$b$, $s$ as in Hypothesis~\ref{hyp: K is basic and M has
%     height h and s is big enough} with~$K$ replaced by~$\Kbasic$ (so
%   that in particular~$b$ and~$s$ depend only on~$K$, $h$
%   and~$n$). % Set~$M_i:=p^{i-1}M/p^iM$, $1\le i\le n$.
% Let~$M'$ (resp.\ $M'_i$) denote~$M$ (resp.\ $M_i$) regarded as an \'etale $\varphi$-module
% over~$\A_{\Kbasic}$. By Corollary~\ref{cor: ramification bound in
%   the basic mod p case}, if~$v\ge pb-1$ then the action
%   of~$G_{\Kbasic_{\cyc,s}}^v$ on each~$V(M'_i)$ is
%   trivial. Since~$M$ is a direct summand of
%   $\A_{K}\otimes_{\A_{\Kbasic}}M'$, and since
%   $K\cdot\Kbasic_{\cyc,s}=K_{\cyc,s}$, we see that in particular the
%   representation~$V(M)$ admits an extension
% \mar{ME: I'm confused about this argument: it looks like we're claiming
% that if $W$ is a rep'n of a group $G$, and if $H$ is a subgroup of $G$,
% then any $H$-invariant direct summand of $W$ is automatically $G$-invariant.  
% But probably I'm misunderstanding! TG: this is the mistake that I
% explained to you at breakfast\dots}
% to a representation
%   of~$G_{K_{\cyc,s}}$, with the property that the restriction of each~$V(M_i)$
%   to~$G_{\Kbasic_{\cyc,s}}^v\cap G_{K_{\cyc,s}}$ is trivial. 
%   (See~\cite[Rem.\ A.3.4.4]{MR1106901} for the relationship between the actions of
%   $G_{\Kbasic_{\cyc}}$ on~$M'$ and $G_{\Kcyc}$ on
%   $\A_{K}\otimes_{\A_{\Kbasic}}M'$.) It follows from the definitions that we
%   have~$G_{\Kbasic_{\cyc,s}}^v\cap
%   G_{\Kcyc}=G_{K_{\cyc,s}}^{\psi_{K/\Kbasic}(v)}$,
%   where~$\psi_{K/\Kbasic}$ is as in~\cite[Chapter IV]{SerreLF}), so if
%   we take any~$w\ge \psi_{K/\Kbasic}(pb-1)$, then~$G_{K_{\cyc,s}}^w$ acts
%   trivially on each~$V(M_i)$; note that~$w$ depends only on~$K$ and~$h$.


%  % We
%  %  claim that there is some~$w$ depending only on~$K$ and~$h$ such that
%  %  ~$G_{K_{\cyc,s}}^w$ acts trivially on each~$V(M_i)$. To see this, it is
%  %  enough to show that we can choose~$w$ so
%  %  that~$G_{\Kbasic_{\cyc,s}}^v\cap G_{K_{\cyc,s}}\supseteq
%  %  G_{K_{\cyc,s}}^w$. It follows from the definitions that we
%  %  have~$G_{\Kbasic_{\cyc,s}}^v\cap
%  %  G_{\Kcyc}=G_{K_{\cyc,s}}^{\psi_{K/\Kbasic}(v)}$,
%  %  where~$\psi_{K/\Kbasic}$ is as in~\cite[Chapter IV]{SerreLF}), so we
%  %  can take any $w\ge\psi_{K/\Kbasic}(pb-1)$. 
%   % \mar{TG: probably some
%     % very inconsistent notation} 

%   Since~$\Kcyc$ is an abelian extension of~$K_{\cyc,s}$, and the
%   action of~$G_{K_{\cyc}}$ on~$V(M)$ is abelian by assumption, we see
%   that the action of~$G_{K_{\cyc,s}}$ on~$V(M)$ is also
%   abelian. \mar{TG: this is nonsense as written! Is it definitely
%     true? At worst I would guess that this needs to be rewritten to be
%     $A$-linear rank one.}
%   Since $G_{K_{\cyc,s}}^w$ acts trivially on each~$V(M_i)$, there is
%   some~$m$ depending only on~$n$ such that for
%   any~$g\in G_{K_{\cyc,s}}^w$, the action of~$g^{p^m}$ on~$V(M)$ is
%   trivial. Since the $p^m$-th power of an open subgroup
%   of~$\cO_K^\times$ is open, it follows from the compatibility of the
%   upper ramification numbering with the filtration on~$\cO_K^\times$
%   that if we increase~$w$ by a quantity depending only on~$K,w$
%   and~$n$, we can assume that~$G_{K_{\cyc,s}}^w$ acts trivially
%   on~$V(M)$. Since the image of~$G_{K_{\cyc,s}}^w$ in~$I_{K_{\cyc,s}}^{\ab}$
%   is open, the intersection of this image with~$I_{\Kcyc}^\ab$ gives
%   us the required open subgroup~$I^{K,h,n}_{\Kcyc}$.\mar{TG: I think
%     this is a bit bogus as written.}
%\end{proof}% \mar{TG: not sure if we can make this argument with~$p$th
  % powers at infinite level? After all, aren't there infinitely many
  % degree~$p$ extensions of~$\Kcyc$, just by Kummer theory? I am a bit worried that after all this, we
  % have to go back and redo everything to work mod~$\varpi^a$ after
  % all in all the arguments above (which I think would be quite a lot
  % of extra work\dots) Hopefully there is a fix using that the
  % extension isn't arbitrary but still has bounded $T$-height in some sense?}



  % We argue by induction on~$n$, the case~$n=1$ being
  % Corollary~\ref{cor: ramification bound in the basic mod p
  %   case}. If~$n>1$ then we can write~$M$ as the extension
  % of~$M/pM$ by~$pM$, on each of which the action
  % of~$I^{K,h,n-1,\ab}_{\Kcyc}$ is trivial by the inductive hypothesis. 

% \begin{lem}\label{lem: kummer bounding character}Assume that we are in the setting of Hypothesis~\ref{hyp: K
%     is basic and M has height h and s is big enough}. 
%   Then there is a positive integer $t(a,h,N)$ such
%   that the action of~$G_{\Kcyc,t}$ on~$V(M^{\inf})$ is \mar{TG: figure out phrasing: do we want an actual Galois
%     representation, and some inertia to act trivially?}
% \end{lem}
% \begin{proof}
  
% \end{proof}
% \begin{proof} Write $\gMt:=\AAinf{A}\otimes_{\A^+_{K,A}}\gM$.
% %    , $M=W(\C^\flat)_A\otimes_{\AAinf{A}}\gM$.
%   We
%   claim that there is a positive integer~$l$, depending only
%   on~$a,h,N$ (but not on~$\gM$) such that the 
%   natural morphisms $T^{-l}\gMt\to M^{\inf}$ and~$M^{\inf}\to M^{\inf}/T^{l}\gMt$ induce an
%   isomorphism between~$(M^{\inf})^{\varphi=1}$ and
%   $(T^{-l}\gMt/T^l\gMt)^{\varphi=1}$.  \mar{TG: this can be upgraded
%     to $(T^{-l}\gMt/T^{>0}\gMt)^{\varphi=1}$, which is I think the
%     correct analogue of what they have in CL, because we are on the
%     dual side; they would get Homs to $W(\cO_\C^\flat)/T^l$. The point
%   is that we know that~$\varphi(T)\in T\A^+$ and is~$T^p$ mod~$p$,
%   which is better than what Lemma~\ref{lem: phi on powers of T without assuming A plus
%       stable} tells us. }

% To see this, we claim firstly that for~$l$ sufficiently large, the
% morphism~$(1-\varphi):M^{\inf}\to M^{\inf}$ induces an isomorphism
% $(1-\varphi):T^l\gMt\to T^l\gMt$; so that the natural morphism
% $(M^{\inf})^{\varphi=1}\to(M^{\inf}/T^l\gMt)^{\varphi=1}$ is an isomorphism. To see
% this, note that by Lemma~\ref{lem: phi on powers of T without assuming A plus
%       stable}, for all~$l$ sufficiently large, we have
%     $\varphi(T^l\A^+_{K,A})\subseteq T^{l+1}\A^+_{K,A}$, so that
%     $\varphi(T^l\gM)\subseteq T^{l+1}\gM$, and in particular
%     $(1-\varphi)(T^l\gMt)\subseteq T^l\gMt$. Furthermore if~$x\in
%     T^l\gMt$ satisfies~$(1-\varphi)(x)=0$, then writing
%     $x=\varphi(x)=\varphi^2(x)=\dots$, we see that~$x\in T^n\gM$ for
%     all~$n\ge l$, so $x=0$, and $(1-\varphi):T^l\gMt\to T^l\gMt$ is
%     injective. To see that it is surjective, given~$y\in T^l\gMt$,
%     set~$z=\sum_{i\ge 0}\varphi^i(y)$; then $(1-\varphi)(z)=y$.

% It remains to show that if~$l$ is sufficiently large (independently of~$\gM$), and~$x\in
% (M^{\inf})^{\varphi=1}$, then~$x\in T^{-l}\gMt$. Suppose that~$x\in T^{-m}\gMt$
% for some~$m\ge 0$, and that $\varphi(x)=x$. Since~$\gM$ has height at
% most~$h$, we have $T^{m+h}\varphi(x)=T^{m+h}x\in\varphi(\gMt)$ for
% some~$y\in\gMt$. By Lemma~\ref{lem: phi on powers of T without assuming A plus
%       stable}, there is a constant~$C$ (independent of~$m$) such that
%     $x\in T^{-(m+h)/p-C}\gM$. In particular, if we choose and $l\ge
%     (h+pC+p)/(p-1)$, then   if~$m> l$, it follows that $x\in
%     T^{1-m}\gM$, and so by induction on~$m$, we conclude that $x\in
%     T^{-l}\gMt$, as claimed.

% We now consider the action of~$G_{K_{\cyc,s}}$ on
% $(T^{-l}\gMt/T^l\gMt)^{\varphi=1}$.\mar{TG: got as far as here.}
% \end{proof}

\chapter{A geometric Breuil--M\'ezard conjecture}\label{sec: BM}
% \mar{TG: Remark somewhere that the conjecture is possibly also
%   interesting from the point of view of the representation theory
%   of~$\GL_n(k)$.} 
 In this chapter
we explain a (for the most part conjectural) relationship between the
geometry of our potentially semistable and crystalline moduli
stacks~$\cX_d^{\crys,\lambdau,\tau}$
and~$\cX_d^{\semis,\lambdau,\tau}$, and the representation theory
of~$\GL_n(k)$. Throughout the chapter we fix a sufficiently large
coefficient field~$E$ with ring of integers~$\cO$ and residue
field~$\F$, and we largely omit it from our notation.

The starting point for this proposed relationship is
the Breuil--M\'ezard conjecture~\cite{BreuilMezard},
which is
a conjectural formula for the Hilbert--Samuel multiplicities of the
special fibres of the lifting
rings~$R^{\crys,\underline{\lambda},\tau}_{\rhobar}$ and
~$R^{\semis,\underline{\lambda},\tau}_{\rhobar}$, and which has important
applications to proving modularity lifting theorems via the
Taylor--Wiles method~\cite{KisinFM}. The conjecture was geometrized
in~\cite{MR3248725} and~\cite{emertongeerefinedBM}, by refining the
conjectural formula for the Hilbert--Samuel multiplicity to a
conjectural formula for the underlying cycle of the special fibres of
the lifting rings ~$R^{\crys,\underline{\lambda},\tau}_{\rhobar}$ and
~$R^{\semis,\underline{\lambda},\tau}_{\rhobar}$, considered as (equidimensional)
closed subschemes of the special fibre of the universal lifting
ring~$R_{\rhobar}^{\square}$. We refer to this generalization as the
``refined Breuil--M\'ezard conjecture'', and to the original
conjecture (or rather, its generalizations to~$\GL_n$) as the
``numerical Breuil--M\'ezard conjecture''.

Our aim in this chapter, then, is to ``globalize'' the conjectures
of~\cite{emertongeerefinedBM} by formulating versions of them
for the 
stacks~$\cX_d^{\crys,\lambdau,\tau}$
and~$\cX_d^{\crys,\lambdau,\tau}$; the conjectures
of~\cite{emertongeerefinedBM} can be recovered from these conjectures
by passing to versal rings at finite type points. (We caution the
reader that there is another kind of ``globalization'' that could be
considered, namely realizing local Galois representations as the
restrictions to decomposition groups of global representations, as
used in the proofs of some of the results
of~\cite{emertongeerefinedBM} and~\cite{geekisin}.
Other than in the proof of
Theorem~\ref{thm: Gee Kisin extended to semistable} below,
we don't  consider this 
kind of globalization in the present work.) This
generalization, which we will call the ``geometric Breuil--M\'ezard
conjecture'', seems to us to be the natural setting in which to
consider the Breuil--M\'ezard conjecture, and we will use our
description of the irreducible components of~$\cX_{d,\red}$ to deduce
new results about both the refined and numerical versions 
of the Breuil--M\'ezard conjecture.


\section[The qualitative geometric Breuil--M\'ezard
  conjecture]{The qualitative geometric Breuil--M\'ezard
  conjecture\sectionmark{The qualitative geometric BM
  conjecture}}\sectionmark{The qualitative geometric BM
  conjecture}\label{subsec: qualitative BM}If~$\lambdau$ is a regular
Hodge type, and~$\tau$ is any inertial type, then by Theorems~\ref{thm: existence of ss stack} and~\ref{thm: dimension of ss stack}, the
stacks~$\cX^{\crys,\lambdau,\tau}_d$ and~$\cX^{\semis,\lambdau,\tau}_d$
are finite type $p$-adic formal algebraic stacks over~$\cO$, which are
$\cO$-flat and equidimensional of dimension~$1+[K:\Qp]d(d-1)/2$. It
follows that their special fibres $\cXbar^{\crys,\lambdau,\tau}_d$ and~$\cXbar^{\semis,\lambdau,\tau}_d$ are algebraic stacks over~$\F$ which
are equidimensional of dimension~$[K:\Qp]d(d-1)/2$. Since
~$\cX^{\crys,\lambdau,\tau}_d$ and~$\cX^{\semis,\lambdau,\tau}_d$ are
closed substacks of~$\cX_d$, $\cXbar^{\crys,\lambdau,\tau}_d$
and~$\cXbar^{\semis,\lambdau,\tau}_d$  are closed substacks of the
special fibre~$\cXbar_d$, and their irreducible components (with the
induced reduced substack structure) are therefore closed substacks of
the algebraic stack~$\cXbar_{d,\red}$ (see~\cite[\href{https://stacks.math.columbia.edu/tag/0DR4}{Tag 0DR4}]{stacks-project}
for the theory of irreducible components of algebraic stacks and their
multiplicities). % \mar{TG:
  % should replace this with stacks project citations, and maybe
  % anything we're citing from here that isn't there could be moved to
  % an appendix of this paper?!}
Since~$\cXbar_{d,\red}$ is
equidimensional of dimension~$[K:\Qp]d(d-1)/2$ by
Theorem~\ref{thm:reduced dimension}, it follows that the irreducible
components of  $\cXbar^{\crys,\lambda,\tau}_d$
and~$\cXbar^{\semis,\lambda,\tau}_d$ are irreducible components
of~$\cXbar_{d,\red}$, and are therefore of the
form~$\cXbar_{d,\red}^{\underline{k}}$ for some Serre
weight~$\underline{k}$ (again by Theorem~\ref{thm:reduced
  dimension}).

For each~$\underline{k}$, we
write~$\mu_{\underline{k}}(\cXbar^{\crys,\lambdau,\tau}_d)$ and
$\mu_{\underline{k}}(\cXbar^{\semis,\lambdau,\tau}_d)$ for the
multiplicity
of~$\cXbar_{d,\red}^{\underline{k}}$ as a component
of~$\cXbar^{\crys,\lambdau,\tau}_d$
and~$\cXbar^{\semis,\lambdau,\tau}_d$. % , in the sense
% of~\cite[\S1]{EGcomponents}.
We
write~$Z_{\crys,\lambdau,\tau}=Z(\cXbar^{\crys,\lambdau,\tau}_d)$ and
~$Z_{\semis,\lambdau,\tau}=Z(\cXbar^{\semis,\lambdau,\tau}_d)$ for the
corresponding cycles, i.e.\ for the formal sums
\numequation\label{eqn: cris HS multiplicity stack}Z_{\crys,\lambdau,\tau}=\sum_{\underline{k}}\mu_{\underline{k}}(\cXbar^{\crys,\lambdau,\tau}_d)\cdot\cXbar_d^{\underline{k}}, \end{equation}
\numequation\label{eqn: ss HS multiplicity stack}Z_{\semis,\lambdau,\tau}=\sum_{\underline{k}}\mu_{\underline{k}}(\cXbar^{\semis,\lambdau,\tau}_d)\cdot\cXbar_d^{\underline{k}}, \end{equation}which
we regard as elements of the finitely generated free abelian group~$\Z[\cX_{d,\red}]$
whose generators are the irreducible
components~$\cXbar_d^{\underline{k}}$.

Now fix some representation~$\rhobar:G_K\to\GL_d(\F)$, corresponding to a
point $x:\Spec\F\to\cX_d$. (More generally, we could consider
representations valued in~$\GL_d(\F')$ for some finite
extension~$\F'/\F$, but in keeping with our attempt to keep the
notation in this chapter as uncluttered as possible by omitting~$\cO$,
it is convenient to suppose that~$\F$ has been chosen sufficiently
large.)  For each regular Hodge type~$\lambdau$ and inertial
type~$\tau$, we have effective versal morphisms
$\Spec
R_{\rhobar}^{\crys,\underline{\lambda},\tau}/\varpi\to\cXbar^{\crys,\lambdau,\tau}$
and
$\Spec
R_{\rhobar}^{\semis,\underline{\lambda},\tau}/\varpi\to\cXbar^{\semis,\lambdau,\tau}$
(see Corollary~\ref{cor: crystalline deformation rings are effective
  versal}),
as well as a (non-effective; see Remark~\ref{rem:non-effective}) versal morphism $\Spf R_{\rhobar}^{\square} \to \cX_d$
(see Proposition~\ref{prop:versal rings}).  % \mar{ME: Have we actually proved non-effectivity?
% I forget! ME: Added a remark about this in \S 6.}

% \mar{TG: probably we should make the superscripts be in the
  % same order for the stacks and the versal rings!}

For each~$\underline{k}$,
%we then set 
%\[\cC_{\underline{k}}(\rhobar):=\
%Spf R_{\rhobar}^\square\times_{\cX_d}\cX_d^{\underline{k}}.\]
we may consider the fibre product
$\Spf R_{\rhobar}^\square\times_{\cX_d}\cX_d^{\underline{k}}.$
This is {\em a priori} a closed formal subscheme of $\Spf R_{\rhobar}^{\square}$,
but since $R_{\rhobar}^{\square}$ is a complete local ring, it may
equally well be regarded as a closed subscheme 
of $\Spec R_{\rhobar}^{\square}$
(see Lemma~\ref{lem: alem closed formal algebraic scheme adic*}).

\begin{lemma}
\label{lem:computing cycle dimension}
%Regarded as a closed subscheme of $\Spec R_{\rhobar}^{\square}$,
The fibre product
$\Spf R_{\rhobar}^\square\times_{\cX_d}\cX_d^{\underline{k}}$,
when we
regard it as a closed subscheme of $\Spec R_{\rhobar}^{\square}$,
is equidimensional of dimension~$d^2+[K:\Qp]d(d-1)/2$.
\end{lemma}
\begin{proof}
As in the proof of Theorem~\ref{thm:reduced dimension},
we may find a regular Hodge type $\underline{\lambda}$ such that $\cX_d^{\underline{k}}$
is an irreducible component  of $(\overline{\cX}^{\crys, \underline{\lambda}})_{\red}.$
Thus
$\Spf R_{\rhobar}^\square\times_{\cX_d}\cX_d^{\underline{k}}$
is versal to a union of irreducible components of
the spectrum of the versal ring
$(R_{\rhobar}^{\crys,\underline{\lambda}}/\varpi)_{\red}$
to $(\overline{\cX}^{\crys,\underline{\lambda}})_{\red}$
(those that correspond to the various formal branches of $\cX_d^{\underline{k}}$  
passing through~$x$, in the terminology of \cite[\href{https://stacks.math.columbia.edu/tag/0DRA}{Tag
  0DRA}]{stacks-project})
% \cite[Def.~1.7]{EGcomponents}) \mar{TG: if I'm correct elsewhere, this
% is the last citation of this preprint, and I think we could cite
% \cite[\href{https://stacks.math.columbia.edu/tag/0DRA}{Tag
%   0DRA}]{stacks-project} instead?}
and hence is of the stated dimension, since 
$R_{\rhobar}^{\crys,\underline{\lambda}}/\varpi$,
and so also
$(R_{\rhobar}^{\crys,\underline{\lambda}}/\varpi)_{\red}$,
is equidimensional of this dimension.
\end{proof}

We let $\cC_{\underline{k}}(\rhobar)$ denote the 
$d^2+[K:\Qp]d(d-1)/2$-dimensional
cycle in $\Spec R_{\rhobar}^{\square}/\varpi$ underlying the fibre product
of Lemma~\ref{lem:computing cycle dimension}.
 % codimension $[K:\Qp]d(d+1)/2$\mar{TG: I'm adopting the
%Viewed in this  light, it is equidimensional of 
%dimension~$d^2+[K:\Qp]d(d-1)/2$ 
%and so finally, we may (and do) regard 
%$\cC_{\underline{k}}(\rhobar)$ as a cycle of
%dimension~$d^2+[K:\Qp]d(d-1)/2$ % codimension $[K:\Qp]d(d+1)/2$\mar{TG: I'm adopting the
  % phrasing from the introduction, but it seems problematic to me,
  % because I don't think we actually know the dimension
  % of~$R_{\rhobar}$ at this point? So probably it's better to say that
  % it is a cycle of dimension~$[K:\Qp]d(d-1)/2$?}
%in~$\Spec R_{\rhobar}/\varpi$. % (note that since~$\cX_d^{\underline{k}}$
% is algebraic, it has effective versal rings, so we really do get a
% closed subscheme of $\Spec R_{\rhobar}/\varpi$, rather than of
% $\Spf R_{\rhobar}/\varpi$).
%\mar{ME: Aren't closed subschemes of $\Spf A$ and $\Spec A$  in bijection
%anyway, for a complete Noetherian local ring~$A$?  Given this, I'm not quite
%sure about this remark. TG: agreed, commenting it out but leaving
%marginal comment for now}
%\mar{ME: Have to find the earlier point in the paper, somewhere in one of the
%arguments in \S 6 I guess, where this dimension statement is proved.
%ME: Actually, might just be easiest to  deduce it here.}

The following theorem gives a qualitative
version of the refined Breuil--M\'ezard
conjecture~\cite[Conj.\ 4.2.1]{emertongeerefinedBM}. While its statement is purely local, we do not know
how to prove it without making use of the stack~$\cX_d$.
\begin{thm}
  \label{thm: qualitative BM}Let $\rhobar:G_K\to\GL_d(\F)$ be a
  continuous representation. Then there are finitely many cycles of
  dimension~$d^2+[K:\Qp]d(d-1)/2$ in~$\Spec R_{\rhobar}^\square/\varpi$ such that
  for any regular Hodge type~$\lambdau$ and any inertial type~$\tau$,
  each of the special fibres~$\Spec
  R_{\rhobar}^{\crys,\underline{\lambda},\tau}/\varpi$ and~$\Spec
  R_{\rhobar}^{\semis,\underline{\lambda},\tau}/\varpi$ is set-theoretically
  supported on some union of these cycles.
\end{thm}
\begin{proof}
	% \mar{ME: This argument is probably implicitly using $G$-ring facts,
	% 	to know that multiplicities don't change when passing
	% 	to versal rings. TG: how about just having this stacks
        %       project citation?}
  We have
  $\Spf R_{\rhobar}^{\crys,\underline{\lambda},\tau}/\varpi =\Spf
  R^{\square}_{\rhobar}\times_{\cX_d}\cXbar^{\crys,\lambdau,\tau}$ and
  $\Spf R_{\rhobar}^{\semis,\underline{\lambda},\tau}/\varpi =\Spf
  R^{\square}_{\rhobar}\times_{\cX_d}\cXbar^{\crys,\lambdau,\tau}$. It
  follows from~\eqref{eqn: cris HS multiplicity stack} and~\eqref{eqn:
    ss HS multiplicity stack}, together with the definition
  of~$\cC_{\underline{k}}(\rhobar)$, that we may write the underlying
  cycles as \numequation\label{eqn: cris HS multiplicity def
    ring}Z(\Spec
  R_{\rhobar}^{\crys,\underline{\lambda},\tau}/\varpi)=\sum_{\underline{k}}\mu_{\underline{k}}(\cXbar^{\crys,\lambdau,\tau}_d)\cdot\cC_{\underline{k}}(\rhobar), \end{equation}
\numequation\label{eqn: ss HS multiplicity def ring}Z(\Spec
R_{\rhobar}^{\semis,\underline{\lambda},\tau}/\varpi)=\sum_{\underline{k}}\mu_{\underline{k}}(\cXbar^{\semis,\lambdau,\tau}_d)\cdot\cC_{\underline{k}}(\rhobar). \end{equation}(Note
that by~\cite[\href{https://stacks.math.columbia.edu/tag/0DRD}{Tag
  0DRD}]{stacks-project}, the multiplicities do not change
when passing to versal rings.) The
theorem follows immediately (taking our finite set of cycles to be
the~$\cC_{\underline{k}}(\rhobar)$). 
\end{proof}
We can regard this theorem as isolating the ``refined'' part
of~\cite[Conj.\ 4.2.1]{emertongeerefinedBM}; that is, we have taken
the original numerical Breuil--M\'ezard conjecture, formulated a geometric
refinement of it, and then removed the numerical part of the
conjecture. The numerical part of the conjecture (in the optic of this
chapter) consists of relating the multiplicities
$\mu_{\underline{k}}(\cXbar^{\crys,\lambdau,\tau}_d)$ and
$\mu_{\underline{k}}(\cXbar^{\semis,\lambdau,\tau}_d)$ to the 
representation theory of~$\GL_n(k)$, as we recall in the next section.

\section{Semistable and crystalline inertial types}\label{subsec:
  type theory} We now
briefly recall the ``inertial local Langlands correspondence''
for~$\GL_d$. Let~$\rec_p$ denote the local Langlands correspondence
for~$\Qpbar$-representations of $\GL_d(K)$, normalized as
in~\cite[\S1.8]{Gpatch}; this is a bijection between the isomorphism
classes of irreducible smooth $\Qpbar$-representations of~$\GL_d(K)$
and the isomorphism classes of $d$-dimensional semisimple
Weil--Deligne $\Qpbar$-representations of the Weil group~$W_K$.  We
have the following result, which is essentially due to
Schneider--Zink \cite{MR1728541}.

\begin{thm}Let ~$\tau:I_K\to\GL_d(\Qpbar)$ be an inertial type. Then
  there are finite-dimensional smooth irreducible
  $\Qpbar$-representations $\sigma^{\crys}(\tau)$ and
  $\sigma^{\semis}(\tau)$ of~$\GL_d(\cO_K)$ with the properties that
  if ~$\pi$ is an irreducible smooth $\Qpbar$-representation
  of~$\GL_d(K)$, then the $\Qpbar$-vector space
  $\Hom_{\GL_d(\cO_K)}(\sigma^{\crys}(\tau),\pi)$ \emph{(}resp.\
the $\Qpbar$-vector space  $\Hom_{\GL_d(\cO_K)}(\sigma^{\semis}(\tau),\pi)$\emph{)} has dimension at
  most~$1$, and is nonzero precisely if~$\rec_p(\pi)|_{I_F}\cong
  \tau$, and $N=0$ on~$\rec_p(\pi)$ \emph{(}resp.\, if~$\rec_p(\pi)|_{I_F}\cong
  \tau$, and~$\pi$ is generic\emph{)}.
\end{thm}
\begin{proof}
  See~\cite[Thm.\ 3.7]{Gpatch} for~$\sigma^{\crys}(\tau)$, and~\cite[Thm.\
  3.7]{MR3769675} together with~\cite[Thm.\ 2.1, Lem.\ 2.2]{2018arXiv180302693P}  for~$\sigma^{\semis}(\tau)$.
\end{proof}Note that we do not claim that the
representations~$\sigma^{\crys}(\tau)$ and~$\sigma^{\semis}(\tau)$ are
unique; the possible non-uniqueness of these representations is
of no importance for us.

For each regular Hodge
type~$\underline{\lambda}$ we let~$L_{\lambdau}$ be the corresponding
representation of~$\GL_d(\cO_K)$, defined as follows: For
each~$\sigma:K\into\Qpbar$, we
write~$\xi_{\sigma,i}=\lambda_{\sigma,i}-(d-i)$, so that
$\xi_{\sigma,1}\ge\dots\ge\xi_{\sigma,d}$. We view each
$\xi_\sigma:=(\xi_{\sigma,1},\dots,\xi_{\sigma,d})$ as a dominant
weight of the algebraic group~$\GL_d$ (with respect to the upper
triangular Borel subgroup), and we write~$M_{\xi_\sigma}$ for the
algebraic $\cO_K$-representation of~$\GL_d(\cO_K)$ of highest
weight~$\xi_\sigma$. Then we define
$L_{\lambdau}:=\otimes_{\sigma}M_{\xi_\sigma}\otimes_{\cO_K,\sigma}\cO$.

For each~$\tau$ we let~$\sigma^{\crys,\circ}(\tau)$,
$\sigma^{\semis,\circ}(\tau)$ denote choices of $\GL_d(\cO_K)$-stable
$\cO$-lattices in~$\sigma^{\crys}(\tau)$, $\sigma^{\semis}(\tau)$
respectively (the precise choices being unimportant). Then we write
$\sigmabar^{\crys}(\lambda,\tau)$, (resp.\ $\sigmabar^{\semis}(\lambda,\tau)$)
for the semisimplification of the $\F$-representation of~$\GL_d(k)$
given by
$L_{\lambdau}\otimes_\cO\sigma^{\crys,\circ}(\tau)\otimes_\cO\F$
(resp.\
$L_{\lambdau}\otimes_\cO\sigma^{\semis,\circ}(\tau)\otimes_\cO\F$). For
each Serre weight~$\underline{k}$, we write~$F_{\underline{k}}$ for
the corresponding irreducible $\F$-representation of~$\GL_d(k)$ (see
for example the appendix to~\cite{MR2541127}). Then there are unique
integers $n_{\underline{k}}^\crys(\lambda,\tau)$ and
  $n_{\underline{k}}^\semis(\lambda,\tau)$ such
    that \[\sigmabar^{\crys}(\lambda,\tau)\cong\oplus_{\underline{k}}F_{\underline{k}}^{\oplus
        n_{\underline{k}}^\semis(\lambda,\tau)},\] \[\sigmabar^{\semis}(\lambda,\tau)\cong\oplus_{\underline{k}}F_{\underline{k}}^{\oplus n_{\underline{k}}^\semis(\lambda,\tau)}.\]
Our geometric Breuil--M\'ezard conjecture is as follows. \index{Breuil--M\'ezard conjecture}
\begin{conj}
  \label{conj: geometric BM}There are cycles~$Z_{\underline{k}}$ with
  the property that for each regular Hodge type~$\lambdau$ and each
  inertial type~$\tau$, we have
  $Z_{\crys,\lambdau,\tau}=\sum_{\underline{k}}n_{\underline{k}}^\crys(\lambda,\tau)\cdot
  Z_{\underline{k}}$,  $Z_{\semis,\lambdau,\tau}=\sum_{\underline{k}}n_{\underline{k}}^\semis(\lambda,\tau)\cdot
  Z_{\underline{k}}$.
\end{conj}

For some motivation for the conjecture (coming from the
Taylor--Wiles patching method), see for example~\cite[Thm.\
5.5.2]{emertongeerefinedBM}. Some evidence for the conjecture is given
in the following sections. 

The expressions~\eqref{eqn: cris HS multiplicity stack} and~\eqref{eqn: ss HS multiplicity
stack} describe $Z_{\crys,\lambdau,\tau}$ and $Z_{\semis,\lambdau,\tau}$
as linear combinations of (the cycles underlying) the irreducible
components~$\overline{\cX}_d^{\underline{k}'}$, and thus each cycle $Z_{\underline{k}}$
(assuming that such cycles exist) will itself be a linear combination of the
various~$\overline{\cX}_d^{\underline{k}'}.$ 
%For some motivation for the conjecture (coming from the
%Taylor--Wiles patching method), see for example~\cite[Thm.\
%5.5.2]{emertongeerefinedBM}. Some evidence for the conjecture is given
%in the following sections. 
We expect that the
cycles~$Z_{\underline{k}}$ are effective, i.e.\ that each of them is a
linear
combination of the~$\cXbar_d^{\underline{k}'}$ with non-negative
(integer) coefficients. Note that the finitely many (conjectural)
cycles~$Z_{\underline{k}}$ are completely determined by the infinitely
many equations in Conjecture~\ref{conj: geometric BM} (see~\cite[Lem.\
4.1.1, Rem.\ 4.1.7(1)]{emertongeerefinedBM}).

%Note that
While the original motivation for Conjecture~\ref{conj:
  geometric BM} was to understand potentially semistable deformation
rings in terms of the representation theory of~$\GL_d$, it can also be
thought of as giving a geometric interpretation of the
multiplicities~$n_{\underline{k}}^\crys(\lambda,\tau)$.

    
\section[{The relationship between the numerical, refined and
  geometric Breuil--M\'ezard conjectures}]{The relationship between the numerical, refined and
  geometric Breuil--M\'ezard conjectures \sectionmark{The relationship between the BM conjectures}}\sectionmark{The relationship between the BM conjectures}\label{subsec: relating
  different BM conjectures}We now explain the relationship between
Conjecture~\ref{conj: geometric BM} and the conjectures
of~\cite{emertongeerefinedBM}. Suppose firstly that
Conjecture~\ref{conj: geometric BM} holds, and fix some
$\rhobar:G_K\to\GL_d(\F)$ corresponding to a point $x:\Spec\F\to\cX_d$. For
each~$\underline{k}$, we
set \[Z_{\underline{k}}(\rhobar):=\Spf
  R_{\rhobar}^\square\times_{\cX_d}Z_{\underline{k}},\]which
%by definition is
we regard as a
$d^2+[K:\Qp]d(d-1)/2$-dimensional cycle
in $\Spec R_{\rhobar}^\square/\varpi$.
(As noted above, the $Z_{\underline{k}}$ will be linear combinations 
of the cycles underlying the various $\cXbar_d^{\underline{k}'}$,
and so the the $Z_{\underline{k}}(\rhobar)$ will be linear
combinations of the various cycles $\cC_{\underline{k}'}(\rhobar)$;
thus Lemma~\ref{lem:computing cycle
dimension} shows that they are indeed 
$d^2+[K:\Qp]d(d-1)/2$-dimensional cycles.)

%\mar{ME: Shouldn't the dimension (here and below) be $d^2 + [K:\Q_p]d(d-1)/2$?  Also,
%perhaps add a justification that these give cycles of the indicated type,
%along the lines of the justification/explanation I added for  the
%$\cC_{\underline{k}}(\rhobar)$? TG: I fixed the dimensions. I'm not
%sure if I'm being stupid, but isn't this just an immediate consequence
%of Lemma~\ref{lem:computing cycle dimension}, because these cycles are
%a linear combination of those fibre products? ME: I think you're right,
%and I added some explanation.}
\begin{remark}
\label{rem:cycle alternative}
Equivalently, we can define
$Z_{\underline{k}}(\rhobar)$ to be the image of~$Z_{\underline{k}}$
under the natural map from $\Z[\cX_{d,\red}]$ to the group $\Z_{d^2+[K:\Qp]d(d-1)/2}(\Spec
R_{\rhobar}^\square/\varpi)$ of $d^2+[K:\Qp]d(d-1)/2$-dimensional cycles in $\Spec
R_{\rhobar}^\square/\varpi$ which is defined as follows: We
let~$R_{\rhobar}^{\alg}$ be the quotient of~$R_{\rhobar}^{\square}/\varpi$
which is a versal ring to~$\cX_{d,\red}$ at~$x$, so that 
by~\cite[\href{https://stacks.math.columbia.edu/tag/0DRB}{Tag 0DRB},\href{https://stacks.math.columbia.edu/tag/0DRD}{Tag 0DRD}]{stacks-project} we have a 
multiplicity-preserving surjection from the set of irreducible
components of~$\Spec R_{\rhobar}^{\alg}$ to the set of irreducible
components of~$\cX_{d,\red}$ containing~$x$; we then send any
irreducible component of~$\cX_{d,\red}$ not containing~$x$ to zero,
and send each irreducible component containing~$x$ to the sum of the corresponding irreducible
components of~$\Spec R_{\rhobar}^{\alg}$  in its preimage.
\end{remark}

Exactly as in the proof of
Theorem~\ref{thm: qualitative BM}, it follows that for each regular
type~$\underline{\lambda}$ and inertial type~$\tau$, we have
\numequation\label{eqn: refined version of BM cris}Z(\Spec
  R_{\rhobar}^{\crys,\underline{\lambda},\tau}/\varpi)=\sum_{\underline{k}}n_{\underline{k}}^\crys(\lambda,\tau)\cdot
  Z_{\underline{k}}(\rhobar),\end{equation} \numequation\label{eqn: refined version of BM ss}Z(\Spec
  R_{\rhobar}^{\semis,\underline{\lambda},\tau}/\varpi)=\sum_{\underline{k}}n_{\underline{k}}^\semis(\lambda,\tau)\cdot
  Z_{\underline{k}}(\rhobar).\end{equation} The first of these statements is
\cite[Conj.\ 4.2.1]{emertongeerefinedBM} (with the cycles~$\cC_a$
there being our cycles~$Z_{\underline{k}}(\rhobar)$), and the second
is the corresponding statement for potentially semistable lifting
rings.

We now relate this conjecture to the numerical version of the
Breuil--M\'ezard conjecture. We have a homomorphism $\Z_{d^2+[K:\Qp]d(d-1)/2}(\Spec
R_{\rhobar}^\square/\varpi)\to\Z$ defined by sending each cycle to its
Hilbert--Samuel multiplicity in the sense
of~\cite[\S2.1]{emertongeerefinedBM}. % \mar{TG: this is a bit
  % inadequate!}
Let $\mu_{\underline{k}}(\rhobar)$ denote the Hilbert--Samuel
multiplicity of the cycle $Z_{\underline{k}}(\rhobar)$, and write $e(\Spec
  R_{\rhobar}^{\crys,\underline{\lambda},\tau}/\varpi)$, $e(\Spec
  R_{\rhobar}^{\semis,\underline{\lambda},\tau}/\varpi)$ for the Hilbert--Samuel
  multiplicities of the indicated rings. Then it follows
  from~\eqref{eqn: refined version of BM cris} and~\eqref{eqn: refined
    version of BM ss} that we have
  \numequation\label{eqn: numerical version of BM cris}e(\Spec
  R_{\rhobar}^{\crys,\underline{\lambda},\tau}/\varpi)=\sum_{\underline{k}}n_{\underline{k}}^\crys(\lambda,\tau)\mu_{\underline{k}}(\rhobar),\end{equation} \numequation\label{eqn:
  numerical version of BM ss}e(\Spec
  R_{\rhobar}^{\semis,\underline{\lambda},\tau}/\varpi)=\sum_{\underline{k}}n_{\underline{k}}^\semis(\lambda,\tau)\mu_{\underline{k}}(\rhobar).\end{equation}
Then~\eqref{eqn: numerical version of BM cris}
is~\cite[Conj.\ 4.1.6]{emertongeerefinedBM}, and ~\eqref{eqn:
  numerical version of BM ss} is the corresponding semistable version.

Suppose now that for each $\rhobar:G_K\to\GL_d(\F)$, all regular Hodge
types~$\lambdau$ and all inertial types~$\tau$ we have~\eqref{eqn:
  numerical version of BM cris} and~\eqref{eqn: numerical version of
  BM ss} for some integers~$\mu_{\underline{k}}(\rhobar)$, but do not
assume Conjecture~\ref{conj: geometric BM} (so in particular we do not
presuppose any geometric interpretation for the
integers~$\mu_{\underline{k}}(\rhobar)$). % \mar{TG: at the moment I'm
%   confused about what we can formally deduce in this case! For
%   example, I think things might be clearer if we could show that each
%   component $\cX_d^{\underline{k}}$ contained smooth (or at least
%   unibranch) points, so that the corresponding (special fibre of the)
%   versal ring is irreducible. Even then I'm a bit confused about
%   whether we can deduce what we want in this generality, or if we need
% to start making some addition assumptions on the~$\cZ_{\underline{k}}$
% and/or the~$\mu_{\underline{k}}(\rhobar)$ - or if indeed some
% restrictions are forced upon them!}

For each $\underline{k}$ we choose a point
$x_{\underline{k}}:\Spec\F\to\cXbar_{d,\red}$ which is contained
in~$\cXbar^{\underline{k}}$ and not in any~$\cXbar^{\underline{k}'}$
for $\underline{k}'\ne\underline{k}$. We furthermore demand that
$x_{\underline{k}}$ is a smooth point of~$\cXbar_{d,\red}$. (Since
$\cXbar_{d,\red}$ is reduced and of finite type over~$\F$, there is a dense set
of points of~$\cXbar^{\underline{k}}$ satisfying these conditions.)
% \mar{TG: this would all be cleanest if we could furthermore choose it
%   to have formally smooth versal ring -- can we? TG: yes, we're on
%   something reduced and varieties have smooth
%   points.}
% Let~$e_{\underline{k}}$ denote the multiplicity
% of~$\cXbar^{\underline{k}}$ at~$x_{\underline{k}}$.\mar{TG: needs to be
%   defined/explained, but I'm getting tired.} 
Write~$\rhobar_{\underline{k}}:G_K\to\GL_d(\F)$ for the
representation corresponding to~$X_{\underline{k}}$, and
set \numequation\label{eqn: cycle from multiplicities}Z_{\underline{k}}:=\sum_{\underline{k}'}
  \mu_{\underline{k}}(\rhobar_{\underline{k}'})\cdot\cXbar^{\underline{k}'}.\end{equation}
% (Note that this is a priori only a $\Q$-linear, rather than
% $\Z$-linear, combination of the~$\cXbar^{\underline{k}}$.) 
Then for each
regular Hodge type~$\lambdau$ and inertial type~$\tau$, it follows
from~\eqref{eqn: numerical version of BM cris}
that \begin{align*}\sum_{\underline{k}}n_{\underline{k}}^\crys(\lambda,\tau)\cdot
  Z_{\underline{k}}&=
                     \sum_{\underline{k}}e(R_{\rhobar_{\underline{k}}}^{\crys,\underline{\lambda},\tau}/\varpi)
                     \cdot\cXbar^{\underline{k}}\\
                   &=\sum_{\underline{k}}\mu_{\underline{k}}(\cXbar^{\crys,\lambdau,\tau}_d)\cdot\cXbar^{\underline{k}}\\
                   &=Z_{\crys,\lambdau,\tau}, \end{align*}where we
                 used that $x_{\underline{k}}$ is a smooth point
                 of~$\cXbar_{d,\red}$ and is only contained  in ~$\cXbar^{\underline{k}}$
to conclude that
 $\mu_{\underline{k}}(\cXbar^{\crys,\lambdau,\tau}_d)=e(R_{\rhobar_{\underline{k}}}^{\crys,\underline{\lambda},\tau}/\varpi)$. % \mar{TG: again a little to
                   % explain here about local/global multiplicities and
                   % where the assumption of being a smooth point comes in.}
Similarly we have  $Z_{\semis,\lambdau,\tau}=\sum_{\underline{k}}n_{\underline{k}}^\semis(\lambda,\tau)\cdot
  Z_{\underline{k}}$, and we conclude that % up to possibly working with
  % cycles with rational (rather than integral) multiplicities, \mar{TG:
  % why rational, not integral?}
the
  geometric Breuil--M\'ezard conjecture (Conjecture~\ref{conj:
    geometric BM}) is equivalent to the numerical conjecture.
  \begin{rem}\label{rem: geom BM lets us go from sufficiently generic
      rhobar to all rhobar}%\mar{TG: refer to this lower down.}
    The argument that we just made shows that the geometric conjecture
    follows from knowing the numerical conjecture for sufficiently
    generic~$\rhobar$ (indeed, it is enough to check it for a single
    sufficiently generic~$\rhobar$ for each irreducible component
    of~$\cX_{d,\red}$), while in turn the geometric conjecture implies
    the numerical conjecture for all~$\rhobar$. % This reduction was in
    % fact one of our original motivations for constructing the stacks~$\cX_d$.
  \end{rem}% \mar{TG: prove that the stacks, and the multiplicities, are
    % actually unique - same arguments as in ~\cite{geekisin} or \cite{emertongeerefinedBM}.}


\section{The weight part of Serre's conjecture}\label{subsec: Serre
  weights}We now briefly explain the relationship between
Conjecture~\ref{conj: geometric BM} and the weight part of Serre's
conjecture. For more details, see~\cite{2015arXiv150902527G}
(particularly Section~6).

We expect that the cycles~$Z_{\underline{k}}$ will be effective, in
the sense that they are combinations of the~$\cX_d^{\underline{k}}$
with non-negative coefficients. This expectation is borne out in all
known examples (see the following sections), and in any case would be
a consequence of standard conjectures about the Taylor--Wiles
method. Indeed, the local cycles~$Z_{\underline{k}}(\rhobar)$ of Section~\ref{subsec: relating
  different BM conjectures} are conjecturally the supports of certain
``patched modules'', and in particular are effective; see~\cite[Thm.\
5.5.2]{emertongeerefinedBM}. The effectivity of
the~$Z_{\underline{k}}(\rhobar)$ would immediately imply the
effectivity of the~$Z_{\underline{k}}$.

The ``weight part of Serre's conjecture'' is perhaps more of a
conjectural conjecture than an actual conjecture: it should assign to
each~$\rhobar:G_K\to\GL_d(\F)$ a set~$W(\rhobar)$ of Serre weights, \index{$W(\rhobar)$}
with the property that if~$\rhobar$ is the restriction to a
decomposition group of a suitable
global representation (for example, an irreducible representation
coming from an automorphic form on a unitary group), then
~$\underline{k}\in W(\rhobar)$ if and only if~$\underline{k}$ is a
weight for the global representation (for example, in the sense that
the global representation corresponds to some mod~$p$ cohomology class
for a coefficient system corresponding to~$F_{\underline{k}}$).

Many conjectural definitions of the sets~$W(\rhobar)$ have been
proposed. Following~\cite{geekisin}, one definition is to assume the
Breuil--M\'ezard conjecture, for example in the form~\eqref{eqn:
  refined version of BM cris}, and define $W(\rhobar)$ to be the set
of~$\underline{k}$ for which $\mu_{\underline{k}}(\rhobar)>0$. While
this is less explicit than other definitions, it has the merit that it
would follow from standard conjectures about modularity lifting
theorems that it gives the correct set of weights; see~\cite[\S3,
4]{2015arXiv150902527G}.

Assume Conjecture~\ref{conj: geometric BM}, and assume that the
cycles~$Z_{\underline{k}}$ are effective. As explained in Section~\ref{subsec: relating
  different BM conjectures}, it follows that the numerical
Breuil--M\'ezard holds for every~$\rhobar$, with
$\mu_{\underline{k}}(\rhobar)$ being the Hilbert--Samuel multiplicity
of the cycle~$Z_{\underline{k}}(\rhobar)$, which
(since~$Z_{\underline{k}}(\rhobar)$ is effective) is positive if and
only if~$Z_{\underline{k}}(\rhobar)$ is nonzero, i.e.\ if and only
if~$Z_{\underline{k}}$ is supported at~$\rhobar$. Thus we can rephrase
the Breuil--M\'ezard version of the weight part of Serre's conjecture
as saying that $W(\rhobar)$ is the set of~$\underline{k}$ such that
$Z_{\underline{k}}$ is supported at~$\rhobar$.

Alternatively, we can rephrase this conjecture in the following way:
to each irreducible component of~$\cX_{d,\red}$, we assign the set of
weights~$\underline{k}$ with the property that~$Z_{\underline{k}}$ is
supported on this component. Then~$W(\rhobar)$ is simply the union of
the sets of weights for the irreducible components of~$\cX_{d,\red}$
which contain~$\rhobar$. As we explain in Section~\ref{subsec: CEGS},
if~$d=2$ then this description agrees with the other definitions
of~$W(\rhobar)$ in the literature, and therefore gives a
geometrization of the weight part of Serre's conjecture.
\section{The case of $\GL_2(\Qp)$}\label{subsec: GL2 Qp}% \mar{TG: it seems that
  % this is still only known for the Steinberg weights if~$p>3$; I'll
  % check with Vytas. Unfortunately I don't think anything we do helps
  % with the Steinberg weights!}
The numerical Breuil--M\'ezard conjecture for~$K=\Qp$ and~$d=2$ is completely
known, thanks in large part to Kisin's paper~\cite{KisinFM} (which
gave a proof in many cases by a mixture of local and global
techniques), and Pa{\v{s}}k{\=u}nas' paper~\cite{paskunasBM} which
reproved these results by purely local means (the $p$-adic local
Langlands correspondence), relaxed the hypotheses on~$\rhobar$, and
also proved the refined version of the correspondence. The remaining
cases not handled by these papers are proved in the
papers~\cite{HuTan,sandermultiplicities,2018arXiv180307451T},
culminating in the proof of the final cases for~$p=2$ by Tung in~\cite{2019arXiv190806174T}.
% for~$p\ge 3$ the conjecture is completely known, and it is expected
% that these methods will also be able to complete the remaining cases
% with~$p=2$ (see the introduction to~\cite{2018arXiv180307451T})

Consequently, by the discussion of Section~\ref{subsec: relating
  different BM conjectures}, Conjecture~\ref{conj: geometric BM} holds
for~$K=\Qp$ and~$d=2$. It follows easily from the explicit
description of~$\mu_{\underline{k}}(\rhobar)$ in the papers cited
above (or alternatively from the description for~$\GL_2(K)$ in
Section~\ref{subsec: CEGS} below)
that~$Z_{\underline{k}}=\cX_2^{\underline{k}}$ unless
$\underline{k}=(a+p-1,a)$ for some~$a$, in which case
$Z_{(a+p-1,a)}=\cX_2^{(a+p-1,a)}+\cX_2^{(a,a)}$.

% When~$p=3$, the conjecture is also mostly known; more precisely, it
% is known in all cases except when the representation~$\rhobar$ is
% reducible and the characters on the diagonal of $\rhobar$ have
% ratio~$\varepsilonbar=1$. % ($= \varepsilonbar^{-1}$ when $p \le
% % 3$).
% Unfortunately, % \mar{TG: not currently clear to me if we can actually
%   % prove this? We know that the points will all have the appropriate
%   % twist of Steinberg as a Serre weight, but it's not clear to me that
%   % we know what the points on the intersection of the trivial and
%   % Steinberg components are. Maybe one way to figure that out would be
%   % to observe that the ordinary deformation rings are smooth at points
%   % where the ratio isn't cyclotomic? We should add that if it
%   % works. TG: I think it's easier than this.}
% by Lemma~\ref{lem: points of Steinberg components are 1 by cyclo} below, the finite type points of the
% components~$\cX_d^{(a+p-1,a)}$ are entirely of this form, so it does not
% seem possible to use our formalism to deduce any of these remaining
% cases.
% \mar{TG: actually it is done at least for
  % $p=3$ \url{https://arxiv.org/pdf/1803.07451.pdf} and hopefully soon
  % for $p=3$.}

\section{$\GL_2(K)$: potentially Barsotti--Tate types}\label{subsec: CEGS}% \mar{TG: refer back to weight part of
  % Serre's conjecture at some point.}
In this section we assume
that $p$ is odd, %~$p~>~2$,
and explain some consequences of the results
of~\cite{geekisin} (which proved the numerical Breuil--M\'ezard
conjecture for $2$-dimensional potentially Barsotti--Tate
representations) and~\cite{CEGSKisinwithdd} (which studied moduli stacks of rank~$2$
Breuil--Kisin modules with tame descent data). We will take advantage
of Remark~\ref{rem: geom BM lets us go from sufficiently generic
      rhobar to all rhobar}.

We begin with the following slight extension of one of the main
results of~\cite{geekisin}. Let~$\underline{\BT}$ denote the minimal regular Hodge type,
  i.e.\ we have~$\underline{\BT}_{\sigma,1}=1$,
  $\underline{\BT}_{\sigma,2}=0$ for all $\sigma:K\into\Qpbar$.
\begin{thm}
  \label{thm: Gee Kisin extended to semistable}Let $K/\Qp$ be a finite
  extension with $p>2$, and let $\rhobar:G_K\to\GL_2(\Qpbar)$ be
  arbitrary. Then the numerical Breuil--M\'ezard conjecture holds for
  potentially crystalline and potentially semistable lifts
  of~$\rhobar$ of Hodge type~$\underline{\BT}$ and arbitrary inertial type~$\tau$.

  More precisely, there are unique non-negative
  integers~$\mu_{\underline{k}}(\rhobar)$ such that \eqref{eqn:
    numerical version of BM cris} and \eqref{eqn: numerical version of
    BM ss} both hold for~$\underline{\lambda}=\underline{\BT}$ and
  ~$\tau$ arbitrary.
\end{thm}
\begin{proof}If we remove the potentially semistable case and consider
  only potentially crystalline representations, the theorem
  is~\cite[Thm.\ A]{geekisin}, which is proved as~\cite[Cor.\
  4.5.6]{geekisin}. We now briefly explain how to modify the proofs
  in~\cite{geekisin} to prove the more general result; as writing out
  a full argument would be a lengthy exercise, and the arguments are
  completely unrelated to those of this book, we only explain the key
  points. Examining the proof of~\cite[Cor.\ 4.5.6]{geekisin}, we see
  that we just need to verify that the assertion of the first sentence
  of the proof of~\cite[Thm.\ 4.5.5]{geekisin} holds in this setting,
  i.e.\ that the equivalent conditions of~\cite[Lem.\
  4.3.9]{geekisin} hold. Exactly as in the proof of~\cite[Cor.\
  4.4.3]{geekisin}, it is enough to show that every irreducible
  component of a product of local deformation rings is witnessed by an
  automorphic representation.

  By the usual Khare--Wintenberger argument, it suffices to prove this
  after making a solvable base change, and in particular we can
  suppose that all of the residual local Galois representations are
  trivial, the mod~$p$ cyclotomic character is trivial, the trivial
  mod~$p$ representation admits a non-ordinary crystalline lift,
  and the inertial types are all trivial. Now, any non-crystalline
  semistable representation of Hodge type~$\underline{\BT}$ is
  necessarily ordinary (indeed, it follows from a direct computation
  of the possible weakly admissible modules that all such
  representations are unramified twists of an extension of the inverse
  of the cyclotomic
  character by the trivial character), and by results of Kisin and Gee
  (see ~\cite[Prop.\ 4.3.1]{MR3778977}, \cite[Cor.\
  2.5.16]{KisinModularity}, and \cite[Prop.\ 2.3]{MR2280776}), as $p>2$ and~$\rhobar$ is
  trivial, the semistable ordinary deformation ring in question is a
  domain, as are the crystalline ordinary deformation ring, and the
  crystalline non-ordinary deformation ring (all in Hodge type~$\underline{\BT}$). It
  therefore suffices to show that in the situation of the proof
  of~\cite[Cor.\ 4.4.3]{geekisin}, given any decomposition of the set~$S$
  of places of~$F_1^+$ lying over~$p$ as~$S_{\textrm{ss}}\coprod S_{\textrm{crys-ord}}
  \coprod S_{\textrm{non-ord}}$, we can arrange to have a
  congruence to an automorphic representation~$\pi''$, having the properties
  that~$\pi''$ is unramified at the places lying over a place in~$S_{\textrm{crys-ord}}
  \coprod S_{\textrm{non-ord}}$,
  is an unramified twist of the Steinberg representation at the
  places  lying over a place in~$S_{\textrm{ss}}$, and is
  furthermore ordinary (resp.\ not ordinary) at the places lying over
  a place in ~$S_{\textrm{crys-ord}}$ (resp.\ $
  S_{\textrm{non-ord}}$). %\mar{TG: rewritten, should be checked.}

  The existence of such a representation if~$S_{\textrm{ss}}=S$
  follows from the construction of the global
  representation~$\overline{r}$ used in the proof of numerical
  Breuil--M\'ezard conjecture in~\cite{geekisin}; more precisely, it
  follows from~\cite[Prop.\ 8.2.1]{0905.4266}, which is applied in the
  proof of~\cite[Thm.\ A.2]{geekisin} (with the type function in the
  sense of~\cite{0905.4266} being~$C$ at all places above~$p$). Thus,
  to establish the general case, it suffices to establish the
  existence of suitable ``level lowering'' congruences. This is easily
  done by switching to a group which is ramified at the places
  lying over~$S_{\textrm{ss}}$, and then making
  congruences to automorphic representations which are either
  unramified and ordinary (in the case of places lying over
  $S_{\textrm{crys-ord}}$) or have cuspidal type (in the case of
  places lying over~$S_{\textrm{non-ord}}$), and then applying the
  Khare--Wintenberger argument again at the places lying
  over~$S_{\textrm{non-ord}}$. (See e.g.\ \cite[Lem.\
  3.5.3]{KisinModularity} for a similar argument for places away
  from~$p$.)
\end{proof}


% We begin by explaining the
% relation between the moduli stacks defined in~\cite{CEGSKisinwithdd}
% and the stacks defined in this paper.\mar{TG: I think we will in
%   particular want to make sure that before fixing types, when we
%   construct our stacks earlier in the paper we have one that is for
%   things becoming potentially crystalline over~$L$. Also there might
%   be a bit of pain about how the crystalline condition is being
%   imposed, because there is some weirdness about the existence of
%   extensions of the Galois action from $K_\infty$ to $K$ in the
%   crystalline case. I hope at least this map is in the right direction
% and that we can conclude that we have an isomorphism by virtue of
% flatness and having the same $\Zpbar$-points\dots}

% Let~$L/K$ be the maximal totally tamely
% ramified extension of the quadratic unramified extension of~$K$, and
% fix a choice of~$\piflat$.\mar{TG: need to check what we do; guessing
%   the point is to fix~$\piflat$ for~$K$ and then for~$L$ it's just
%   choosing a root of this which is the same as a root of~$\pi$ due to tameness} Then
% the stacks considered in~\cite{CEGSKisinwithdd} are the $p$-adic
% formal algebraic stack~$\cC^{\CEGS}$ of rank 2 Breuil--Kisin modules
% with descent data to~$G_K$.\mar{TG: need to understand how this
%   relates to BKF modules with descent data!}

% \begin{itemize}
% \item Hopefully we can phrase things in terms of a map from our stacks
%   to the CEGS stacks given by fixing~$\piflat$ and forgetting most of
%   the structure. One thing that might be tricky is the relationship of
%   Galois actions, e.g.\ of tame descent data versus BKF modules. I had
%   previously imagined that we might map the other way but this way
%   looks easier. Everything will be easy to compare on $\Zpbar$-points
%   so once we have a map I think we will be in good shape.
% \item Probably a better idea will actually be to compare them both as
%   closed substacks of some other stack, something like BKF modules
%   with a descent over $L$ for a single $\piflat$, and then a
%   $\Gal(L/K)$ action.
% \item Actually maybe better will be to use
%   $\Gal(L_{\infty}/K_{\infty})$? I'm possibly also coming around to
%   just working with Galois rather than Kisin, using $\Zp$-flatness
%   (and maybe (generic) reducedness? I guess we should think about
%   whether we're proving that, I guess we can get it from the
%   deformation rings?) and then thinking about the restriction from
%   ~$K$ to ~$K_\infty$.
 
% \item At one point I thought we would want to extend the geometric BM
%   results in \cite{CEGSKisinwithdd} to the case of general pot BT
%   types but actually I think we only need the numerical version so we
%   can just use~\cite{geekisin}.
% \item There is then the question of what we do about the Steinberg
%   component. I think in fact it's not going to be hard to just replace
%   ``potentially BT'' with ``potentially ss of Hodge type 0'' because
%   the missing things are just twists of Steinberg so just twists of
%   ordinary and we can handle these by hand? That feels a bit too good
%   to be true so maybe I'm being stupid.
% \end{itemize}

We say that a Serre weight~$\underline{k}$ for~$\GL_2$ is \index{Serre weight!Steinberg}
%\mar{ME: Something weird going on here --- is it possibly an indexing thing?}
\emph{Steinberg} if for each~$\sigmabar$ we have
$k_{\sigmabar,1}-k_{\sigmabar,2}=p-1$. If~$\underline{k}$ is Steinberg
then we define~$\underline{\tilde{k}}$ by
$\tilde{k}_{\sigmabar,1}=\tilde{k}_{\sigmabar,2}=k_{\sigmabar,2}$. 
\begin{thm}
  \label{thm: explicit BM cycles GL2 K}% \mar{TG: define Steinberg
    % weight. For everything except those we argue as in
    % \cite{CEGSKisinwithdd}: the weights can be seen in tame types and
    % those special fibres are generically reduced. For the Steinberg
    % weight it's enough to be computing the ``$N\ne 0$'' component
    % (i.e.\ ordinary which are twists of trivial by cyclo) for a
    % non-split $\rhobar$ and this is easily seen to be smooth (both
    % here and in the citation of Snowden above, note that since~$p>2$
    % we can reduce to the fixed determinant case.)}
 Continue to assume that~$d=2$ and~$p>2$. Then Conjecture~{\em \ref{conj: geometric BM}} holds
  for~$\underline{\lambda}=\underline{\BT}$ and~$\tau$ arbitrary, with
  the cycles~$Z_{\underline{k}}$ being as follows: if~$\underline{k}$
  is not Steinberg, then $Z_{\underline{k}}=\cX_2^{\underline{k}}$,
  while if $\underline{k}$ is Steinberg,
  $Z_{\underline{k}}=\cX_2^{\underline{k}}+\cX_2^{\underline{\tilde{k}}}$.

\end{thm}
\begin{proof}
We use the notation of Section~\ref{subsec: relating different BM
    conjectures}. By Theorem~\ref{thm: Gee Kisin extended to semistable} and the
  discussion of Section~\ref{subsec: relating different BM
    conjectures}, we need only show that the
  cycles~$Z_{\underline{k}}$ in the statement of the theorem are those
  determined by~\eqref{eqn: cycle from multiplicities}. Suppose
  firstly that~$\underline{k}$ is not Steinberg. Then by~\cite[Thm.\
  5.2.2~(2)]{CEGSKisinwithdd}, $\cX_2^{\underline{k}}$ has a dense set of
  finite type points with the property that their only non-Steinberg
  Serre weight is~$\underline{k}$. It follows from this and~\cite[Thm.\
  5.2.2~(3)]{CEGSKisinwithdd}  that if
  neither~$\underline{k}$ nor~$\underline{k}'$ is Steinberg, then
  $\mu_{\underline{k}}(\rhobar_{\underline{k}'})=\delta_{\underline{k},\underline{k}'}$.
  By construction, if~$\underline{k}'$ is Steinberg,
  then~$\rhobar_{\underline{k}'}$ is a twist of a tr\`es ramifi\'ee
  extension of the trivial character by the mod~$p$ cyclotomic
  character. By~\cite[Lem.\ B.5]{CEGSKisinwithdd}, this implies that
  $\mu_{\underline{k}}(\rhobar_{\underline{k}'})=0$.
  %unless~$\underline{k}=\underline{k}'$.
  The cycles
  $Z_{\underline{k}}$ are therefore as claimed if~$\underline{k}$ is
  non-Steinberg. 

  It remains to determine the values of
  $\mu_{\underline{k}}(\rhobar_{\underline{k}'})$ in the case
  that~$\underline{k}$ is Steinberg. By twisting, we can and do assume
  that~$k_{\sigmabar,2}=0$ for all~$\sigmabar$.  If we apply
  Theorem~\ref{thm: Gee Kisin extended to semistable} with $\tau$ being
  the trivial type, and recall  that $L_{\underline{\BT}}$  is the trivial representation
  and that $\sigma^{\semis}(\tau)$ is the Steinberg type, the reduction
  of which  is precisely the representation $F_{\underline{k}}$,
  we find that 
\numequation
\label{eqn:instance of BM}
Z_{\semis, \underline{\BT}, \tau} = Z_{\underline{k}}.
\end{equation}
  %being $\underline{\BT}$ in the case of liftings
  %of Hodge type~$\underline{\BT}$, we see that
  Thus
  $\mu_{\underline{k}}(\rhobar_{\underline{k}'})\ne 0$ if and only
  if~$\rhobar_{\underline{k}'}$ admits a semistable lift of Hodge
  type~$\underline{\BT}$. If this lift is in fact crystalline, then
  ~$\rhobar_{\underline{k}'}$ has~$\underline{\tilde{k}}$ as a Serre
  weight, so by another application of~\cite[Thm.\
  5.2.2~(2)]{CEGSKisinwithdd}, we see that either
  $\underline{k}'=\underline{\tilde{k}}$ or else that $\underline{k}'$ is
  also Steinberg.
  %$\underline{k}'=\underline{{k}}$.
  In the former case,
  $\rhobar_{\underline{k}'}$ is (by Remark~\ref{rem: open eigenvalue morphism implies ratio
          not cyclo}) an unramified twist of an extension of
  inverse of the mod~$p$ cyclotomic
  character by a
  non-trivial unramified character, so it does not admit a semistable non-crystalline lift of
  Hodge type~$\underline{\BT}$ (as all such lifts are unramified
  twists of an extension of the inverse of the cyclotomic  character by the trivial
  character); so we have
  $\mu_{\underline{k}}(\rhobar_{\underline{\tilde{k}}})=\mu_{\underline{\tilde{k}}}(\rhobar_{\underline{\tilde{k}}})=1$. 

  Finally, we are left with the task of
  computing $\mu_{\underline{k}}(\rhobar_{\underline{k}'})$
  when both $\underline{k}$ and $\underline{k'}$  are Steinberg.
%in the case
%  showing that
%  $\mu_{\underline{k}}(\rhobar_{\underline{k}})=1$ in the case
%  that~$\underline{k}$ is Steinberg.
Recall that by twisting, we are
  assuming that~$k_{\sigmabar,2}=0$ for all~$\sigmabar$.
The weight $\underline{k}'$ is then a twist of~$\underline{k}$.
We have to show that in this case we again
  have $\mu_{\underline{k}}(\rhobar_{\underline{k}'}) = \delta_{\underline{k},
\underline{k}'}.$
To see this, first note that
  since $\rhobar_{\underline{k}'}$  is  tr\`es ramifi\'ee, it does not admit a
  crystalline lift of Hodge type~$\underline{\BT}$, so all of its
  semistable lifts of Hodge type~$\underline{\BT}$ are given by
  unramified twists of  extensions of the inverse of the cyclotomic character by
  the trivial character.   Furthermore, the  reduction of any lattice in such
  an extension  is an  
  unramified twists of an extension of the inverse of the mod $\varpi$
  cyclotomic character by the trivial character, and hence  cannot equal
  $\rhobar_{\underline{k}'}$ unless $\underline{k}'$ equals $\underline{k}$
  (rather than being a non-trivial twist of it).
  If we take into account~\eqref{eqn:instance of BM}, we find
  that indeed 
  $\mu_{\underline{k}}(\rhobar_{\underline{k}'}) = 0$ when
  $\underline{k}'  \neq  \underline{k}$.
  
  Finally, any lift of $\rho_{\underline{k}}$ which is an 
  unramified twist of an
  extension of the inverse of the cyclotomic character by
  the trivial character is 
  automatically semistable of Hodge type~$\underline{\BT}$, so that
  $R^{\ss,\underline{\BT},\tau}_{\rhobar{\underline{k}}}$ is
  precisely the ordinary (framed) deformation ring parameterizing
  such lifts of $\rhobar_{\underline{k}}$.
  %which are unramified twists
  %of extensions of the inverse of the cyclotomic character by
  %the trivial character.
  A standard Galois cohomology
  calculation (that we leave to the reader)
  shows that this ring is formally smooth,
  and thus that
  $R^{\ss,\underline{\BT},\tau}_{\rhobar{\underline{k}}}/\varpi$ is
  also formally smooth.
  It follows that $\mu_{\underline{k}}(\rhobar_{\underline{k}})  = 1$,
  as claimed.
%
%
%
%Then by
%  construction~$\rhobar_{\underline{k}}$ is an unramified twist of a
%  tr\`es ramifi\'ee extension of the  inverse of the mod~$p$
%  cyclotomic character by the trivial character. Being  tr\`es ramifi\'ee, it does not admit a
%  crystalline lift of Hodge type~$\underline{\BT}$, so all of its
%  semistable lifts of Hodge type~$\underline{\BT}$ are given by
%  unramified twists of an extension of the inverse of the cyclotomic character by
%  the % \mar{TG: check this here and above/below; I guess with our
%    % conventions it's probably actually extensions of inverse
%    % cyclotomic by trivial?}
%  trivial character. Conversely, any lift of this form is
%  automatically semistable of Hodge type~$\underline{\BT}$, so it
%  suffices to show that the corresponding ordinary lifting ring is
%  formally smooth. This follows from a standard Galois cohomology
%  calculation that we leave to the reader.
\end{proof}
The following lemma makes precise which Galois representations occur
on Steinberg components. Note that the twist in the statement of the
lemma is determined by the values of~$k_{\underline{\sigma},2}$.

\begin{lem}
  \label{lem: points of Steinberg components are 1 by
    cyclo}If~$\underline{k}$ is Steinberg, then the $\Fpbar$-points of
  $\cX_2^{\underline{k}}$ are twists of an extension of the inverse of
  the mod~$p$ cyclotomic character by the trivial character.
\end{lem}
\begin{proof}Twisting, we may assume that~$k_{\underline{\sigma},2}=0$
  for all~$\sigmabar$. Let $\rhobar:G_K\to\GL_2(\F)$ correspond to a
  closed point of~$\cX_2^{\underline{k}}$. We claim that~$\rhobar$ is
  an unramified twist of an extension of~$\epsilonbar$ by~$1$. To
  see this, let~$R^{\underline{\BT},\operatorname{St}}_{\rhobar}$
  denote the $\Zp$-flat quotient
  of~$R^{\underline{\BT},\semis}_{\rhobar}$ determined by the
  irreducible components of the generic fibre which are not components
  of~$R^{\underline{\BT},\crys}_{\rhobar}$.

  It follows from Theorem~\ref{thm: explicit BM cycles GL2 K} that the
  cycle of
  $\Spec R^{\underline{\BT},\operatorname{St}}_{\rhobar}/\varpi$ is
  nonzero, so that in particular~$\rhobar$ must admit a semistable
  non-crystalline lift of Hodge type~$\BT$. As in the proof of
  Theorem~\ref{thm: Gee Kisin extended to semistable}, any such
  representation is an unramified twist of an extension
  of~$\epsilon^{-1}$ by the trivial character, so we are done. (We
  could presumably also phrase this argument on the level of the
  crystalline and semistable moduli stacks, but have chosen to present
  it in the more familiar setting of deformation rings.)
  \end{proof}

\begin{rem}
  \label{rem: we could compare to CEGS but we won't}The
  paper~\cite{CEGSKisinwithdd} constructs and studies moduli stacks of
  two-dimensional representations of~$G_{K_{\piflat,\infty}}$ (for a
  fixed choice of~$\piflat$) which are tamely potentially of height at
  most~$1$.  Since restriction from Barsotti--Tate representations
  of~$G_K$ to representations of~$G_{K_{\piflat,\infty}}$ of height at
  most~$1$ is an equivalence (by the results
  of~\cite{KisinModularity}), these stacks can be interpreted as
  stacks of~$G_K$-representations (and indeed are interpreted as such
  in~\cite{CEGSKisinwithdd}). It is presumably straightforward (using
  that all the stacks under consideration are of finite type, and are
  $\Zp$-flat and reduced, in order to reduce to a comparison of points
  over finite extensions of~$\Zp$) to identify them with the stacks
  considered in this book (for~$\underline{\lambda}=\underline{\BT}$
  and~$\tau$ a tame inertial type). (We were able to apply the results
  of~\cite{CEGSKisinwithdd} in the proof of Theorem~\ref{thm: explicit
    BM cycles GL2 K} without doing this because we only needed to use them on the level
  of versal rings, which are given by universal Galois lifting rings
  for both our stacks and those of~\cite{CEGSKisinwithdd}.) It may
  well be the case that the stacks of Kisin modules considered
  in~\cite{CEGSKisinwithdd} can be identified with certain of the
  moduli stacks of Breuil--Kisin--Fargues modules %~$\cC_{2,\crys,1}^{L/K}$ % \mar{TG: leaving this for now as versions
    % with descent data aren't yet in this paper}
  that we defined in Section~\ref{sec: semistable stack}. %  (for an
  % appropriate tame extension~$L/K$).
  Since we do not need to know
  this, we leave it as an exercise for the interested reader.
\end{rem}

\subsection{A lower bound}The patching arguments of~\cite{geekisin} easily
imply a general inequality, as we now record.

\begin{prop}
  \label{prop: GL2 geometric BM inequality}Assume that~$p>2$. Let~$d=2$, and let the~$Z_{\underline{k}}$ be
  as in the statement of Theorem~{\em \ref{thm: explicit BM cycles GL2
    K}}. Then for each regular Hodge type~$\underline{\lambda}$ and
  each inertial type~$\tau$, we have 
  $Z_{\crys,\lambdau,\tau}\ge\sum_{\underline{k}}n_{\underline{k}}^\crys(\lambda,\tau)\cdot
  Z_{\underline{k}}$,  and $Z_{\semis,\lambdau,\tau}\ge\sum_{\underline{k}}n_{\underline{k}}^\semis(\lambda,\tau)\cdot
  Z_{\underline{k}}$.
\end{prop}
\begin{rem}
  The meaning of the inequalities in the statement of
  Proposition~\ref{prop: GL2 geometric BM inequality} is the obvious
  one: each side is a linear combination of irreducible components,
  and the assertion is that for each irreducible component, the
  multiplicity on the left hand side is at least the multiplicity on
  the right hand side.
\end{rem}
\begin{proof}[Proof of Proposition~{\em \ref{prop: GL2 geometric BM
    inequality}}]As in Section~\ref{subsec: relating
  different BM conjectures}, it is enough to show that for
each~$\rhobar$, we have \[e(\Spec
  R_{\rhobar}^{\crys,\underline{\lambda},\tau}/\varpi)\ge\sum_{\underline{k}}n_{\underline{k}}^\crys(\lambda,\tau)\mu_{\underline{k}}(\rhobar),\] \[e(\Spec
  R_{\rhobar}^{\semis,\underline{\lambda},\tau}/\varpi)\ge\sum_{\underline{k}}n_{\underline{k}}^\semis(\lambda,\tau)\mu_{\underline{k}}(\rhobar),\]where
the~$\mu_{\underline{k}}(\rhobar)$ are the uniquely determined
integers from~\cite{geekisin}. (In fact, it is enough to show these
inequalities when $\rhobar=\rhobar_{\underline{k}}$ for
some~$\underline{k}$, but assuming this does not simplify our
arguments.) Since~\cite{geekisin} considers only potentially
crystalline representations, we only give the proof
for~$R_{\rhobar}^{\crys,\underline{\lambda},\tau}$; the proof in the potentially
semistable case is essentially identical, and we leave it to the reader.

We now examine the proof of~\cite[Thm.\ 4.5.5]{geekisin},
setting~$\rbar$ there to be our~$\rhobar$. As we've done throughout this book, we
write~$\underline{k}$ rather than~$\sigma$ for Serre weights. If we
ignore the assumption that every potentially crystalline lift of Hodge
type~$\lambdau$ and inertial type~$\tau$ is potentially
diagonalizable, then we do not know that the equivalent conditions
of~\cite[Lem.\ 4.3.9]{geekisin} hold, but examining the proof
of~\cite[Lem.\ 4.3.9]{geekisin}, we do know that the statement of
part~(4) of \emph{loc.\ cit.} can be replaced with an inequality (with
equality holding if and only if~$M_\infty$ is a faithful
$R_\infty$-module).

Returning to the proof of~\cite[Thm.\ 4.5.5]{geekisin}, we note that
the definition of~$\mu_{\underline{k}}(\rhobar)$ is such that the
quantity $\mu'_{\sigma_{\mathrm{gl}}}(\rbar)$ is simply the product of
the corresponding~$\mu_{\underline{k}}(\rhobar)$. In particular, if we
take $\lambda_v=\lambdau$ and~$\tau_v=\tau$ for each~$v$, then
(bearing in mind the discussion of the previous paragraph), the main
displayed equality becomes an inequality
\[e(\Spec
  R_{\rhobar}^{\crys,\underline{\lambda},\tau}/\varpi)^N\ge\left(\sum_{\underline{k}}n_{\underline{k}}^\crys(\lambda,\tau)\mu_{\underline{k}}(\rhobar)\right)^N, \]where
there are~$N$ places of~$F^+$ lying over~$p$. Since each side is
non-negative, the result follows.
\end{proof}
\begin{rem}\label{rem: BM inequality well known, doesn't obviously
    work for GLd though.}
  The argument used to prove Proposition~\ref{prop: GL2 geometric BM
    inequality} is well known to the experts, and goes back
  to~\cite[Lem.\ 2.2.11]{KisinFM}. Despite the relatively formal
  nature of the argument, it does not seem to be easy to prove an
  analogous statement for~$\GL_d$ with~$d>2$; the difficulty is in the
  step where we used~\cite[Thm.\ 4.5.5]{geekisin} to replace a global
  multiplicity with a product of local multiplicities. Without this
  argument (which crucially relies on the modularity lifting theorems
  of~\cite{KisinModularity}) one only obtains inequalities involving
  cycles in products of copies of~$\cX_d$.
\end{rem}
\section{Brief remarks on $\GL_{\lowercase{d}}$, $\lowercase{d}>2$}\label{subsec: GLn}We
expect the situation for~$\GL_d$, $d>2$, to be considerably more
complicated than that for~$\GL_2$. Experience to date suggests that
the weight part of Serre's conjecture in high dimension is
consistently more complicated than is anticipated, %expected,
and so it seems unwise to
engage in much speculation. In particular, the results of~\cite{le2020local} show that even for generic
weights~$\underline{k}$, we should not expect the
cycles~$Z_{\underline{k}}$ to have as simple a form as those
for~$\GL_2$.

In the light of Lemma~\ref{lem: maxnonsplit weight k implies crystalline lifts to
    shifts}, % \mar{TG: which is due to be upgraded to something with
  % a proof, and covering shifted weights}
  it seems reasonable to expect~$\cX_d^{\underline{k}}$ to contribute
  to~$Z_{\underline{k}}$, and it also seems reasonable to expect
  contributions from the ``shifted'' weights, as in
  Definition~\ref{defn: shifted weight} (see also~\cite[\S
  7.4]{2015arXiv150902527G}). The comparative weakness of the
  automorphy lifting theorems available to us in dimension greater
  than~$2$ prevents us from saying much more than this, although we
  refer the reader to~\cite[\S\S 3,4,6]{2015arXiv150902527G} for a
  discussion of the conjectural general relationship between the
  numerical Breuil--M\'ezard conjecture, the weight part of Serre's
  conjecture, automorphy lifting theorems, and the stacks~$\cX_d$.

\renewcommand{\theequation}{\Alph{chapter}.\arabic{section}}%.\arabic{subsubsection}}
\appendix

\chapter{Formal algebraic stacks}\label{app: formal algebraic
  stacks}The theory of formal
algebraic stacks is developed in~\cite{Emertonformalstacks}. In this
appendix we briefly summarise the parts of this theory that are used in
the body of the book, and also introduce some additional terminology
and establish some additional results which we will require.% \mar{TG: this appendix needs to be written. To
  % add: flat parts (move from earlier in the paper), nice obstruction
  % theory, underlying reduced substack, definition of (locally)
  % Noetherian.}

\subsection*{Algebraic stacks}We follow the terminology of~\cite{stacks-project}; in
particular, we write ``algebraic stack'' rather than ``Artin
stack''. More \index{algebraic stack}
precisely, an algebraic stack is a stack in groupoids in the \emph{fppf} topology,
whose diagonal is representable by algebraic spaces, which admits a smooth
surjection from a
scheme. See~\cite[\href{http://stacks.math.columbia.edu/tag/026N}{Tag
  026N}]{stacks-project} for a discussion of how this definition relates to
others in the literature. If~$S$ is a scheme, then by ``a stack
over~$S$'' we mean  a stack fibred in groupoids over the big
\emph{fppf} site of~$S$.

	We say that a morphism $\cX \to \cY$ of stacks over $S$
	is {\em representable by algebraic stacks} \index{representable by algebraic stacks}
	if for any morphism of stacks $\cZ \to \cY$ whose source is an algebraic
	stack, the fibre product $\cX \times_{\cY} \cZ$
	is again an algebraic stack.
(In the Stacks Project, the terminology {\em algebraic} is
used instead~\cite[\href{http://stacks.math.columbia.edu/tag/06CF}{Tag
  06CF}]{stacks-project}.) 
Note that a morphism from a sheaf to a
stack is representable by algebraic stacks if and only if it is representable
by algebraic spaces (this is easily verified,
or see e.g.\ \cite[Lem.~3.5]{Emertonformalstacks}). 

Following~\cite[\href{http://stacks.math.columbia.edu/tag/03YK}{Tag 03YK},
\href{http://stacks.math.columbia.edu/tag/04XB}{Tag 04XB}]{stacks-project},
we can define properties of morphisms representable by algebraic stacks
in the following way.

\begin{adefn}
	\label{df:properties of morphisms representable by alg stacks}
	If $P$ is a property of morphisms of algebraic stacks which is 
	{\em fppf} local on the target, and preserved by arbitrary
	base-change, then we say that a morphism $f:\cX \to \cY$ of
	stacks which is representable by algebraic stacks {\em has property
		$P$} if and only if for every algebraic stack $\cZ$
	and morphism $\cZ \to \cY$, the base-changed morphism of algebraic
	stacks $\cZ \times_{\cY} \cX \to \cZ$ has property $P$.
\end{adefn}

	When applying
	Definition~\ref{df:properties of morphisms representable by alg stacks},
	it suffices to consider the case when $\cZ$ is actually a scheme,
	or even an affine scheme.

Some properties $P$ to which we can apply
	Definition~\ref{df:properties of morphisms representable by alg stacks}
are being {\em locally of finite type}, {\em locally of finite presentation},
and {\em smooth}.   The following lemma provides alternative descriptions of
the latter two properties.

\begin{alemma}
\label{lem:characterizing properties}
Let $f:\cX \to \cY$ be a morphism of stacks which is representable by
algebraic stacks.
\begin{enumerate}
\item The morphism $f$ is locally of finite presentation if and only 
if it is {\em limit preserving on objects}, in the sense
of~\cite[\href{http://stacks.math.columbia.edu/tag/06CT}{Tag 06CT}]{stacks-project}.
\item The morphism $f$ is smooth if and only if it is locally of finite 
presentation and {\em formally smooth},
in the sense that it satisfies the usual infinitesimal lifting property:
  for every affine $\cY$-scheme $T$, and every closed
  subscheme $T_0\into T$ defined by a nilpotent ideal sheaf, the functor
  $\Hom_\cY(T,\cX)\to\Hom_\cY(T_0,\cX)$ is essentially
  surjective.
\end{enumerate}
\end{alemma}
\begin{proof}
We first recall that a morphism of algebraic stacks is locally of
finite presentation if and only if it is limit preserving on 
objects~\cite[Lem.~2.3.16]{EGstacktheoreticimages}.

Suppose now that $f$ is locally of finite presentation,
that $T$ is an affine scheme, written as a projective limit of affine schemes
$T \iso \varprojlim T_i$, and suppose given a commutative diagram
$$\xymatrix{T  \ar[r] \ar[d] & T_i \ar[d] \\ \cX \ar[r] & \cY}$$
for some value of~$i$.
We must show that we can factor the left-hand vertical arrow through~$T_{i'}$,
for some $i'\geq i$, so that the evident resulting diagram again commutes.
The original diagram induces a morphism $T \to T_i\times_{\cY} \cX,$
and obtaining a factorization of the desired type amounts to obtaining
a factorization of this latter morphism through some~$T_{i'}$.  
In other words, we must show that the morphism
$T_i\times_{\cY} \cX \to T_i$
%$T \to T_i\times_{\cY} \cX$
is limit preserving on objects.  But this is a morphism of algebraic
stacks which is locally of finite presentation,
and so it is indeed limit preserving on objects.

Conversely, suppose that $f$ is limit preserving on objects.
We must show that for any morphism $T \to \cY$, the base-changed morphism
$T\times_{\cY} \cX\to T$ is locally of finite presentation.
However, this base-change is again limit preserving on
objects~\cite[\href{https://stacks.math.columbia.edu/tag/06CV}{Tag 06CV}]{stacks-project},
and since it is a morphism of algebraic stacks, it is indeed locally
of finite presentation.% \mar{TG: I don't think we've given a proof of
  % (2)?!}

We now turn to~(2). By Definition~\ref{df:properties of morphisms
  representable by alg stacks} and~\cite[Cor.\ B.9]{MR2818725} (that
is, by the result at hand in the case that~$\cY$ is itself an
algebraic stack), it is enough to show that~$f$ satisfies the
infinitesimal lifting property if and only if the base changed
morphism $Z\times_\cY\cX\to Z$ satisfies the infinitesimal lifting
property for all morphisms $Z\to\cY$ whose source is a scheme.

Assume first that ~$f$ satisfies the
infinitesimal lifting property, and let $T_0\to T$ be a nilpotent closed
immersion whose target is an affine $Z$-scheme. Given a morphism
$T_0\to Z\times_\cY\cX$, the infinitesimal lifting property lets us lift the composite $T_0\to
Z\times_\cY\cX\to \cX$ to a morphism~$T\to\cX$, and since we also have
a morphism $T\to Z$, we obtain the required morphism $T\to
Z\times_\cY\cX$.

Conversely, assumed that the base changed
morphism $Z\times_\cY\cX\to Z$ satisfies the infinitesimal lifting
property for all~$Z\to\cY$ with $Z$ a scheme, and let~$T_0\to T$ be as in the statement
of the proposition. Taking~$Z=T$, we can lift the induced morphism
$T_0\to T\times_{\cY}\cX$ to a morphism $T\to T\times_{\cY}\cX$, and
the composite of this morphism with the second projection gives us the
required lifting.\end{proof}




\subsection*{Formal algebraic spaces}% \mar{TG: include~\cite[Lem.\
% 8.18]{Emertonformalstacks}}
\index{formal algebraic space}
Following~\cite[\href{http://stacks.math.columbia.edu/tag/0AHW}{Tag
    0AHW}]{stacks-project}, an
affine formal algebraic space over a base scheme~ $S$
is a sheaf $X$ on the \emph{fppf} site of $S$
which admits a description as
an Ind-scheme
$X \iso \varinjlim_i X_i,$ where the $X_i$ are affine schemes
and the transition morphisms 
are thickenings 
(in the sense 
of~\cite[\href{http://stacks.math.columbia.edu/tag/04EX}{Tag
  04EX}]{stacks-project}). Here we allow the indexing set in the
inductive limit to be arbitrary; if it can be chosen to be the natural
numbers, then we say that~$X$ is \emph{countably
  indexed}.
A countably indexed affine formal algebraic space can be written in
the form $X \cong \varinjlim_n \Spec A/I_n$, where $A$ is a complete
\index{countably indexed}
topological ring equipped with a decreasing sequence $\{I_n\}_n$ of open 
ideals which are weak ideals of definition, i.e.\ consist of topologically
nilpotent elements 
(\cite[\href{https://stacks.math.columbia.edu/tag/0AMV}{Tag 0AMV}]{stacks-project}),
and which
form a fundamental basis of $0$ in~$A$ (see the discussion
of~\cite[\href{https://stacks.math.columbia.edu/tag/0AIH}{Tag 0AIH}]{stacks-project}).
We then write $X := \Spf A$
(following~\cite[\href{https://stacks.math.columbia.edu/tag/0AIF}{Tag 0AIF}]{stacks-project}).



A particular example of a countably indexed affine formal algebraic
space is an adic affine formal algebraic space. By definition, this
is of the form~$\Spf A$, where~$A$ is adic: that is, $A$ is a topological ring which is
complete (and separated, by our convention throughout this book), and
which admits an ideal of definition (that is, an ideal~$I$ whose
powers form a basis of open neighbourhoods of zero). We say that $\Spf
A$ is Noetherian if $A$ is adic and Noetherian. % \mar{TG: comment out if
% OK.}
We say that~$A$
(or $\Spf A$) is adic* if~$I$ can be taken to be finitely
generated. \index{adic*} 

More generally, we say that a complete topological ring~$A$ is weakly
\index{weakly admissible}
admissible if~$A$ contains an open ideal~$I$ consisting of
topologically nilpotent elements, \cite[\href{http://stacks.math.columbia.edu/tag/0AMV}{Tag
  0AMV}]{stacks-project}). 

A formal algebraic space over $S$ is a sheaf $X$ on the \emph{fppf}
site of $S$ which receives a morphism $\coprod U_i \to X$ which is
representable by schemes, \'etale, and surjective, and whose source is
a disjoint union of affine formal algebraic spaces~$U_i$. We say
that~$X$ is \emph{locally countably indexed} if the~$U_i$ can be
\index{locally countably indexed}
chosen to be countably indexed. We say that~$X$ is \emph{locally
  Noetherian} if the~$U_i$ can be taken to be Noetherian. % \mar{TG: I
%   moved this paragraph down below the ones discussing the affine case;
% comment out if OK.}

We will find the
following lemmas useful.

%\mar{ME: Phrase the following in an appropriate level of generality}
\begin{alem}
  \label{lem: alem closed formal algebraic scheme adic*}Let~$X$
  be an affine formal algebraic space over~$S$,
which is either countably indexed {\em (}e.g.\ an adic or adic* affine
formal algebraic space{\em )},
and hence of the form $\Spf A$ for some weakly admissible topological ring~$A$,
by~\cite[\href{https://stacks.math.columbia.edu/tag/0AIK}{Tag 0AIK}]{stacks-project},
%\mar{ME: Put in correct SP citation here, and check if terminology is correct too!
% ME: Done by TG; thanks!}
or else is of the form $\Spf A$, for a pro-Artinian ring $A$.
  Then if~$Y\to X$ is a
  closed immersion of formal algebraic spaces over~$S$, we have that~$Y$ is of
  the form~$\Spf B$ for some quotient~$B$ of~$A$ by a closed ideal, 
  endowed with its quotient topology.  
\end{alem}
\begin{proof}
In the countably indexed case,
  this is immediate from ~\cite[\href{https://stacks.math.columbia.edu/tag/0AIK}{Tag 0AIK},\href{http://stacks.math.columbia.edu/tag/0ANQ}{Tag
  0ANQ},\href{http://stacks.math.columbia.edu/tag/0APT}{Tag
  0APT}]{stacks-project}.
If $A$ is pro-Artinian, say $A = \varprojlim A_i$ with the $A_i$ Artinian,
then $\Spec A_i \times_{\Spf A} Y$ is a closed subscheme of $\Spec A_i$ for 
each index~$i$, and hence is of the form $\Spec B_i$ for some quotient 
$B_i$ of $B$.
Thus $$Y \iso \varinjlim_i \Spec A_i\times_{\Spf A} Y \iso \varinjlim \Spec B_i
\iso \Spf B,$$
where $B = \varprojlim_i B_i$.  The general theory of pro-Artinian rings
shows that $B$ is a quotient of $A$, since each $B_i$ is a quotient
of the corresponding~$A_i$~\cite[Rem.~2.2.7]{EGstacktheoreticimages}.
\end{proof}

\begin{alem}\emph{(\cite[Lem.\ 8.18]{Emertonformalstacks})}\label{alem: Noetherian flat affine formal algebraic spaces}	A morphism 
% \mar{ME: Recall the notion of (locally) Noeth.\ above, since it is used
% here and below TG: did this, comment out if OK.}
  of Noetherian affine formal algebraic spaces $\Spf B \to \Spf A$
  which is representable by algebraic spaces is {\em (}faithfully{\em
    )} flat if and only if the corresponding morphism $A \to B$ is
  {\em (}faithfully{\em )} flat.
  
\end{alem}

\subsection*{Ind-algebraic and formal algebraic
  stacks}\index{Ind-algebraic stack}
\begin{adefn}\label{adefn: Ind algebraic stack}(\cite[Defn.\ 4.2]{Emertonformalstacks})
  An \emph{Ind-algebraic stack}
  over a scheme~$S$ is a stack~$\cX$ over~$S$ which can be written as
  $\cX\cong\varinjlim_{i\in I}\cX_i$, where we are taking the
  2-colimit in the 2-category of stacks of a 2-directed system
  $\{\cX_i\}_{i\in I}$ of algebraic stacks over~$S$.
\end{adefn}
By \cite[Rem.\
  4.9]{Emertonformalstacks}, if $\cX\to\cY$ is representable by
  algebraic spaces, and~$\cY$ is an Ind-algebraic stack, then~$\cX$ is
  also an Ind-algebraic stack.
\begin{adefn}\label{adefn: formal algebraic stack}(\cite[Defn.\ 5.3]{Emertonformalstacks})
  A \emph{formal algebraic stack} over~$S$ is a stack~$\cX$ which
  admits a \index{$p$-adic formal algebraic stack}
  morphism $U\to\cX$ which is representable by algebraic spaces,
  smooth, and surjective, and whose source is a formal algebraic
  space. If the source can be chosen to be locally countably indexed,
  then we say that~$\cX$ is locally countably indexed.
\end{adefn}

\begin{adefn}\label{adefn: p adic formal algebraic stack}(\cite[Defn.\ 7.6]{Emertonformalstacks})
  A \emph{$p$-adic formal algebraic stack} is a formal algebraic stack
  $\cX$ over~$\Spec\Zp$ which admits a morphism $\cX\to\Spf\Zp$ which
  is representable by algebraic stacks; so in particular, we may
  write~$\cX$ as an Ind-algebraic stack by writing
  $\cX\cong\varinjlim_n\cX\times_{\Spf\Zp}\Spec\Z/p^n\Z$.
\end{adefn}
The relationship between formal algebraic stacks and Ind-algebraic
stacks is discussed in detail
in~\cite[\S6]{Emertonformalstacks}. Before explaining some of this
material, we need to recall some  preliminaries on finiteness
properties and underlying reduced substacks.
\subsection*{Underlying reduced substacks}If~$\cX/S$ is any stack,
then we let $(\cX_{\red})'$ be the full subcategory of~$\cX$ whose set
of objects consists of those $T\to\cX$ for which~$T$ is a reduced
$S$-scheme.
\begin{adefn}\label{adefn: underlying reduced substack}(\cite[Lem.\ 3.27]{Emertonformalstacks})
  The \emph{underlying reduced substack} $\cX_{\red}$ of~$\cX$ is the
  intersection of all of the substacks of~$\cX$ which
  \index{underlying reduced substack}
  contain~$(\cX_{\red})'$. 
\end{adefn}

\begin{alem}\label{alem: underlying reduced of Ind algebraic}\emph{(}\cite[Lem.\ 4.16]{Emertonformalstacks}\emph{)}
  If~$\cX$ is an Ind-algebraic stack, and we write
  $\varinjlim_i \cX_i \iso \cX$, for some $2$-directed system
  $\{\cX_i\}_{i \in \cI}$ of algebraic stacks, then the induced
  morphism $\varinjlim_i (\cX_i)_{\red} \to \cX_{\red}$ is an
  isomorphism, and so in particular $\cX_{\red}$ is again an
  Ind-algebraic stack.
\end{alem}
As the following lemma records, in the case that~$\cX$ is a formal
algebraic stack, then~$\cX_{\red}$ is algebraic, and (just as in the
case of formal schemes, or formal algebraic spaces), $\cX$ is a
thickening of~$\cX_{\red}$.
\begin{alemma}\emph{(}\cite[Lem.\ 5.26]{Emertonformalstacks}\emph{)}
	\label{alem:underlying reduced}
	If $\cX$ is is a formal algebraic stack over $S$,
	then $\cX_{\red}$ is a closed and reduced algebraic substack of $\cX$, 
	and the inclusion $\cX_{\red} \hookrightarrow \cX$
	is a thickening {\em (}in the sense that its base-change 
	over any algebraic space induces a thickening of algebraic
	spaces{\em )}.   Furthermore,
	any morphism $\cY \to \cX$ with $\cY$ a reduced
	algebraic stack factors through~$\cX_{\red}.$
\end{alemma}


\subsection*{Finiteness properties}Let~$\cX$ be a formal algebraic
stack~$\cX$. We say that~$\cX$ is quasi-compact if the algebraic
stack~$\cX_{\red}$ is quasi-compact. We say that~$\cX$ is
quasi-separated if the diagonal morphism of~$\cX$ (which is
automatically representable by algebraic spaces) is quasi-compact and
quasi-separated.

We say that ~$\cX$ is locally Noetherian if and only if it admits a
morphism $U\to\cX$ which is representable by algebraic spaces, smooth,
and surjective, whose source is a locally Noetherian formal algebraic
space. Finally, we say that~$\cX$ is Noetherian if it is locally
Noetherian, and it is quasi-compact and quasi-separated.



\subsection*{The relationship between Ind-algebraic and formal algebraic stacks}
In
one direction, we have the following lemma.


\begin{alem}\label{alem: formal implies Ind}\emph{(}\cite[Lem.\ 6.2]{Emertonformalstacks}\emph{)}
  	If $\cX$ is a quasi-compact and quasi-separated
	%\mar{ME: Define quasi-separated above}
	formal algebraic stack, then $\cX \cong \varinjlim_i \cX_i$ for
	a $2$-directed system $\{\cX_i\}_{i \in I}$ of quasi-compact and quasi-separated
algebraic
	stacks in which the transition
	morphisms are thickenings. 
      \end{alem}

      A partial converse to Lemma~\ref{alem: formal implies Ind} is
      proved as~\cite[Lem.\ 6.3]{Emertonformalstacks}. In particular,
      we have the following useful criteria for being a 
      formal algebraic stack.% \mar{TG: add in ~\cite[Cor.\
      %   6.6]{Emertonformalstacks} and change the citation of it earlier in the
      % paper}


\begin{aprop}
  \label{prop: formal stack 6.6}\emph{(}\cite[Cor.\ 6.6]{Emertonformalstacks}\emph{)} Suppose that $\cX$ is an
  Ind-algebraic stack that can be written as the $2$-colimit
  $\cX \iso \varinjlim \cX_n$ of a directed sequence
  $(\cX_n)_{n \geq 1}$ in which the $\cX_n$ are algebraic stacks, and
  the transition morphisms are closed immersions.  If $\cX_{\red}$ is
  a quasi-compact algebraic stack, then $\cX$ is a locally countably
  indexed formal algebraic stack.
\end{aprop}

    
\begin{aprop}
  \label{prop: criterion for p-adic formal algebraic stack}Suppose
  that~$\{\cX_n\}_{n\ge 1}$ is an inductive system of algebraic stacks
  over~$\Spec\Zp$,
  %whose transition morphisms are thickenings,\mar{TG:
    %isn't this automatic from the following conditions? ME: Yes, corrected}
such that each~$\cX_n$ in fact lies over~$\Spec\Z/p^n\Z$, and each
of the induced morphisms
$\cX_n\to\cX_{n+1}\times_{\Spec\Z/p^{n+1}\Z}\Z/p^n\Z$
is an isomorphism. % , and each~$\cX_n$ has quasi-compact
% diagonal. % \mar{TG: hypothesis on diagonal has been removed from
  % \cite{Emertonformalstacks} ?}
Then the Ind-algebraic stack $\cX:=\varinjlim_n\cX_n$
  is a $p$-adic formal algebraic stack.
\end{aprop}
\begin{proof}
  This is a special case of~\cite[Lem.\ 6.3, Ex.\ 7.8]{Emertonformalstacks}.
\end{proof}


%We will find the following lemma useful.

% \begin{alemma}
%   \label{alem: closed immersion is lfp}Let $\cX\to\cY$ be a closed
%   immersion, where~$\cY$ is a formal algebraic stack which is locally
%   of finite presentation.\mar{TG: or locally Noetherian?} Then
%   $\cX\to\cY$ is locally of finite presentation.
% \end{alemma}
% \begin{proof}
%   \mar{TG: ME to fill in.}
% \end{proof}
\subsection*{Scheme-theoretic images}The notion of the
scheme-theoretic image of a (quasi-compact) morphism of algebraic stacks
is developed
in~\cite[\href{https://stacks.math.columbia.edu/tag/0CMH}{Tag
  0CMH}]{stacks-project}; equivalent alternative presentations are
given in~\cite[\S 3.1]{EGstacktheoreticimages} and in~\cite[Ex.\
9.9]{Emertonformalstacks}. In~\cite[\S 6]{Emertonformalstacks} this is
extended to a definition of the scheme-theoretic image
between formal algebraic stacks that are quasi-compact and quasi-separated.
In fact, the definition there uses 
        Lemma~\ref{alem: formal implies Ind} to write the formal stacks
in question as Ind-algebraic stacks with transition morphisms 
being closed immersions, and this is the level of generality
that is appropriate to us here.

A robust theory of scheme-theoretic images requires
a quasi-compactness assumption on the morphism whose scheme-theoretic image
is being formed, and we begin by introducing the notion of quasi-compactness
which seems appropriate to our context.

\begin{adefn}
\label{def:ind-q.c.}
Let $\cX$ be an Ind-algebraic stack
which can be written as a $2$-colimit of quasi-compact algebraic stacks
for which the transition morphisms are closed immersions, say
        $\cX \cong \varinjlim \cX_{\lambda}$.
We say that a morphism of stacks $\cX \to \cY$
is {\em Ind-representable by algebraic stacks} if each of the induced
morphisms
\anumequation
\label{eqn:restricted morphism}
\cX_{\lambda} \to \cY
\end{equation}
is representable by algebraic stacks.
We say that such a morphism is furthermore {\em Ind-quasi-compact}
if each of the morphisms $\cX_{\lambda} \to \cY$ (representable by
algebraic stacks by assumption) is quasi-compact.
\end{adefn}

Using the fact that any two descriptions of $\cX$ as an Ind-algebraic
stack as above (i.e.\ as the $2$-colimit of quasi-compact algebraic
stacks with respect to transition morphisms that are closed immersions)
are mutually cofinal, one sees that the property of a morphism
being Ind-representable by algebraic stacks (resp.\ Ind-representable
by algebraic stacks and Ind-quasi-compact) is independent of
the choice of such a description of~$\cX$.


Suppose now, in the context of the Definition~\ref{def:ind-q.c.},
that $\cY$ is also an Ind-algebraic stack,
which admits a description as the $2$-colimit of algebraic stacks
with respect to transition morphisms that are closed immersions,
say $\cY \cong \varinjlim \cY_{\mu}$.
Then the induced morphism~\eqref{eqn:restricted morphism}
factors through
$\cY_{\mu}$ for some $\mu$
(since the $\cX_{\lambda}$ are quasi-compact
	and the transition morphisms between the $\cY_{\mu}$ are monomorphisms).
Since any morphism of algebraic stacks is representable by algebraic stacks,
and since closed immersions are representable by algebraic stacks,
we find that~\eqref{eqn:restricted morphism} is necessarily representable
by algebraic stacks, and 
thus the morphism $\cX \to \cY$ is necessarily Ind-representable by algebraic
stacks.  
Furthermore, 
the morphism~\eqref{eqn:restricted morphism} is quasi-compact
if and only if the induced morphism $\cX_{\lambda} \to \cY_{\mu}$
is quasi-compact for some (or equivalently any) allowable choice
of~$\mu$.
For example, if the $\cY_{\mu}$ are all quasi-separated,
then these induced morphisms are necessarily quasi-compact
(via the usual graph argument, since each $X_{\lambda}$ 
is quasi-compact), and so in this case any morphism
$\cX\to \cY$ is necessarily Ind-quasi-compact.

Suppose now that the morphism $\cX \to \cY$ {\em is} Ind-quasi-compact,
or equivalently (as we have just explained), that the various induced morphisms 
$\cX_{\lambda} \to \cY_{\mu}$ are quasi-compact.
Then each of these induced morphisms
%
%
%       Thus we suppose that $\cX \to \cY$ is a morphism of Ind-algebraic stacks,
%       and that we may write 
%%
%%
%%	Namely, suppose that $\cX \to \cY$ is a morphism of formal algebraic
%%	stacks, each of which is quasi-compact and
%%        quasi-separated.\mar{TG: can we at least let~$\cY$ be
%%          Ind-algebraic, since that's what we want to use for~$\cX_d$?
%%        As ever, if this is a problem we'll have to swap sections
%%        around. It seems to me that the definition made in
%%        \cite{Emertonformalstacks} works for Ind-algebraic things, but
%%      since it is only done for formal algebraic there, I'm not sure\dots}
%% 
%	% The usual diagonal argument
%	% shows that $\cX \to \cY$ is then also quasi-compact and 
%	% quasi-separated. 
%%        Lemma~\ref{alem: formal implies Ind} allows us to write
%        $\cX \cong \varinjlim \cX_{\lambda}$ and
%        $\cY \cong \varinjlim \cY_{\mu}$, with the transition
%        morphisms in each case being closed immersions, and with the $\cX_{\lambda}$
%being  quasi-compact algebraic stacks,
%and the $\cY_{\mu}$ being quasi-compact and quasi-separated
%        algebraic stacks.  For each
%        index $\lambda,$ the composite morphism
%        $\cX_{\lambda} \to \cX \to \cY$ factors through one of the
%        $\cY_{\mu}$
%(since the $\cX_{\lambda}$ are quasi-compact
%	and the transition morphisms between the $\cY_{\mu}$ are monomorphisms).
%        The resulting morphism
%	$\cX_{\lambda} \to \cY_{\mu}$ 
%	is quasi-compact (by the standard graph argument, taking into account
%the fact that $X_{\lambda}$ is quasi-compact and that $Y_{\mu}$ is quasi-compact
%and quasi-separated),
%and so
has a scheme-theoretic image~$\cZ_\lambda$.  One easily checks that
	 $\cZ_{\lambda}$, thought of as a closed substack of $\cY$,
	 is independent of the particular choice of the index $\mu$
	 used in its definition.
	% and we may regard $\cY_\lambda$ as the scheme-theoretic
	% image of $\cX_{\lambda}$ in $\cY$. 
	Evidently $\cZ_{\lambda}$ is a closed substack
	of $\cZ_{\lambda'}$ if $\lambda \leq \lambda'$.
	%with the resulting morphism $\cZ_{\lambda}
	%\hookrightarrow \cZ_{\lambda'}$
	%being a thickening. 
        In particular, we may form the $2$-colimit
	$\varinjlim \cZ_{\lambda}$, which is an Ind-algebraic stack.
%	a formal algebraic stack, by Lemma~\ref{lem:ind to formal}. 
	There is a natural morphism $\varinjlim \cZ_{\lambda} \to
        \cY$.

\begin{adefn}
\label{def:scheme-theoretic image}\index{scheme-theoretic image}
In the preceding context,
we define the scheme-theoretic image of the Ind-quasi-compact morphism
        $\cX\to\cY$ to be~$\cZ~:=~\varinjlim \cZ_{\lambda}.$
\end{adefn}

It is easily verified that the scheme-theoretic image, so defined,
is independent of the chosen descriptions of $\cX$ and~$\cY$.
%choices of~$\cX_\lambda$ and~$\cY_\mu$.
% \mar{TG: guessing we want some
          % more properties of this - e.g.\ at least in our applications
          % it's presumably a closed substack of~$\cY$, and a formal
          % algebraic stack - but leaving this until it's clear what we
          % need.}
        There is a canonical monomorphism $\cZ \hookrightarrow \cY$.
        (We don't claim that this monomorphism is necessarily a closed immersion
        in this level of generality.)

        A very special case of this definition is the case that~$\cX \to \cY$
        is representable by algebraic spaces and quasi-compact,
        and~$\cY=\Spf A$ for a pro-Artinian ring~$A$.
%is an adic* affine formal algebraic space.
In this case the definition of the
        scheme-theoretic image coincides with~\cite[Defn.\ 3.2.15]{EGstacktheoreticimages} (i.e.\ the scheme-theoretic image can be computed via the scheme-theoretic images of the
        pull-backs of~$\cX$ to the discrete Artinian quotients of~$A$).
        In this particular case the scheme-theoretic image {\em is} a
        closed formal subspace of $\cY$, i.e.\ of the form $\Spf B$
        for some topological quotient $B$ of~$A$.

        \begin{adefn}
          \label{adefn: scheme theoretically dominant}Let
          $\cX$ and $\cY$ be 
 Ind-algebraic stacks which satisfy
the hypotheses introduced above, i.e.\ which may each be written as the $2$-colimit with
respect to closed immersions of \index{scheme-theoretically dominant}
%quasi-compact and quasi-separated
algebraic stacks, which are furthermore quasi-compact in the case 
of~$\cX$.
%formal algebraic stacks, each of which is quasi-compact and quasi-separated.
Then we say that an Ind-quasi-compact morphism
$\cX\to\cY$ is \emph{scheme-theoretically dominant}
          if the induced map $\cZ\to\cY$ is an isomorphism, where
          $\cZ$ is the scheme-theoretic image of $\cX\to\cY$.
        \end{adefn}
\
\begin{aremark}
If $\cX$ is a quasi-compact and quasi-separated formal algebraic stack,
then Lemma~\ref{alem: formal implies Ind} shows
that $\cX \cong \varinjlim \cX_{\lambda}$ with the $\cX_{\lambda}$ 
being quasi-compact and quasi-separated algebraic stacks, and the morphisms
being thickenings, and so in particular closed immersions.
The preceding discussion thus applies to morphisms $\cX \to \cY$
of quasi-compact and quasi-separated formal algebraic stacks, and shows that such
morphisms are necessarily Ind-representable by algebraic stacks and Ind-quasi-compact.
In particular,
Definition~\ref{adefn: scheme theoretically dominant} applies to such
morphisms, and in this case recovers the definition of scheme-theoretic dominance 
given in~\cite[Def.~6.13]{Emertonformalstacks}.

A closely related context in which the preceding definition applies if the following:
if $A$ is an adic* topological ring, with finitely generated ideal of definition~$I$,
and $\cX,\cY \to \Spf A$ are morphisms
of quasi-compact formal algebraic stacks that are representable by algebraic stacks,
then we may write $\cX \cong \varinjlim_n \cX_n$
and $\cY \cong \varinjlim_n \cY_n$,
where $\cX_n := \cX \times_{\Spf A} \Spec A/I^n$
and $\cY_n := \cY\times_{\Spf A} \Spec A/I^n$
are quasi-compact algebraic stacks.  A morphism $\cX \to \cY$ of stacks over $\Spf A$
is then necessarily representable by algebraic stacks~\cite[Lem.~7.10]{Emertonformalstacks},
and so is Ind-quasi-compact if and only if it quasi-compact in the usual sense.
We may then define its scheme-theoretic image, following Definition~\ref{def:scheme-theoretic
image},
or speak of such a morphism being scheme-theoretically dominant.
\end{aremark}

% \mar{TG: maybe change this to say ``the scheme-theoretic image of a
%   $p$-adic formal algebraic stack is a $p$-adic formal algebraic stack''?}
We now show (in Proposition~\ref{aprop: scheme theoretic image is p adic finite
    type II}) that under certain hypotheses the
scheme-theoretic image of a $p$-adic formal algebraic stack is also a
$p$-adic formal algebraic stack. The deduction of this result from
those of~\cite{Emertonformalstacks} will involve the following definition.  % If $\cX\to\cY$ is scheme-theoretically dominant, and~$\cX$ is a
% $p$-adic formal algebraic stack, we would like to be able to deduce
% that~$\cY$ is also a $p$-adic formal algebraic stack.
% Proposition~\ref{aprop: scheme theoretic image is p adic finite
%     type II} is a result of this kind.
  % ; it relies on the following
%   definition.\mar{TG: might not ultimately need this definition here?}
\begin{adefn}
  \label{adefn: locally Ind finite type}(\cite[Defn.\ 8.26, Rem.\
  8.39]{Emertonformalstacks}) We say that a formal algebraic
  stack~$\cX$ over a scheme $S$ is \index{Ind-locally of finite type} \emph{Ind-locally of finite type
    over~$S$} %(resp.\ \emph{locally Ind-finite presentation over~$S$})
  if there exists an isomorphism
$\cX \iso \varinjlim \cX_i,$
%$\coprod U_i \to\cX$ which is representable 
%by algebraic spaces, smooth,
%  and surjective, with each $U_i$ being an affine formal algebraic space
%  which is {\em Ind-finite type}, in the sense
%  that it can be written as an inductive limit~$U_i \isoto\varinjlim_j U_{i,j}$
%  of affine schemes which are of finite type % (resp.\ of finite
%  % presentation)
%  over~$S$, with the transition maps being thickenings.
where each $\cX_i$ is an algebraic stack locally of finite type over~$S$,
and the transition morphisms are thickenings.
\end{adefn}%\mar{TG: change thickenings to closed immersion? ME: Now sorted out}

\begin{aremark}
\label{rem:Ind-loc. f.t. criterion}
If $\cX$ is a quasi-compact and quasi-separated formal algebraic stack,
then in order to verify that $\cX$ is Ind-locally of finite type,
it suffices to exhibit an isomorphism
$\cX \iso \varinjlim \cX_i,$
where each $\cX_i$ is an algebraic stack locally of finite type over~$S$,
and the transition morphisms are closed immersions (not necessarily
thickenings) \cite[Lem.~8.29]{Emertonformalstacks}.
\end{aremark}
%\begin{alem}
%  \label{rem: instance of Ind finite type}
%  If~$\cX$ is a $2$-colimit of algebraic stacks of finite type
%  over~$S$, with the transition maps
%  being closed immersions, % \mar{TG: versus thickenings?}
%  then~$\cX$ is
%  locally Ind-finite type over~$S$.%\mar{TG: need to verify this, if it's actually true. The idea is to have something that we can just apply for~$\cX$.} 
%\end{alem}
%\begin{proof}
%It suffices to show that if $U \to \cX$ is a morphism whose source is an affine
%formal algebraic space, then $U$ is Ind-finite type.  
%If we write $\cX \cong \varinjlim_i \cX_i,$ where each $\cX_i$ is locally of finite type
%over $S$, with the transition morphisms being closed immersions,
%then we find that
%\numequation
%\label{eqn:ind f.t.}
%U \cong \varinjlim U\times_{\cX} \cX_i. 
%\end{equation}
%Here each $U_i:= U\times_{\cX} \cX_i$ is both an algebraic stack (since the projection to
%$\cX_i$ is representable by algebraic spaces) and a formal algebraic space
%(being a closed subspace of $U$), and thus is an algebraic space. 
%Furthermore, each $U_i$ is of finite type over $S$ (since $\cX_i$ is locally
%of finite type and $U \to \cX_i$ is locally of finite presentation, being smooth)
%\end{proof}
 
%% \mar{TG: eventually remove the following result.}        
% \begin{aprop}\label{aprop: scheme theoretic image is p adic finite
%     type}\emph{(}\cite[Prop.\ 10.5]{Emertonformalstacks}\emph{)} Suppose that we have quasi-compact and quasi-separated
%   formal algebraic stacks $\cX$ and $\cY$, and
%   a commutative diagram $$\xymatrix{\cX \ar[r] \ar[rd] & \cY \ar[d] \\
%     & \Spf \Zp}$$ in which the diagonal arrow makes $\cX$ into a
%   $p$-adic formal algebraic stack of finite type. Suppose also that
%   the horizontal morphism $\cX \to \cY$ is proper and
%   scheme-theoretically dominant, and that $\cY$ is locally Ind-finite
%   type over $\Spec \Zp$.

%   Then $\cY$ is a $p$-adic
%   formal algebraic stack of finite type.
% \end{aprop}


\begin{aprop}\label{aprop: scheme theoretic image is p adic finite
    type II} Suppose that we have 
  a commutative diagram $$\xymatrix{\cX \ar[r] \ar[rd] & \cY \ar[d] \\
    & \Spf \Zp}$$ in which the diagonal arrow makes $\cX$ into a
  $p$-adic formal algebraic stack of finite presentation, and where~$\cY$ is 
  an Ind-algebraic stack which can be written as  the inductive limit of
        algebraic stacks, each of finite presentation
        over~$\Spec\Z/p^a$ for some~$a\geq 1$,
with the transition maps being closed immersions.
% \mar{ME: There is a conflict b/w the hypothesis, which writes
% $\cY = \varinjlim_{\mu} \cY_{\mu}$ with $\cY_{\mu}$ f.p.\ over $\Z/p^a$,
% and the  argument, which assumes $\cY$ q.c.q.s., and {\em deduces} this
% structure.  We should figure out which is preferable. TG: the
% hypotheses should stay the same, and the proof should change. ME: Done ---
% if you're happy with what I wrote, you can comment this out.}
% \mar{TG: should we just have some
      %     terminology for this? Or maybe have a different hypothesis
      %     which applies in our case? Note that I have now relaxed
      %     $\cY$ to be Ind rather than formal.}\mar{TG: in the proof
      %     below, ``closed immersions'' have become ``thickenings''. Do
      %     we really need either? Note that we only apply this result
      %     once in the paper; in that case we know that the source is a
      %   $p$-adic formal algebraic stack of finite presentation which
      %   is flat over~$\Zp$ and has affine diagonal, the morphism is
      %   representable by algebraic spaces, proper, and of finite
      %   presentation; but the target is only Ind-algebraic, the
      %   inductive limit under closed immersions of algebraic stacks of
      % finite presentation, and the diagonal of the target is
      % representable by algebraic spaces, affine, and of finite presentation.}
% is 
%  locally of Ind-finite type % Noetherian
%   %formal algebraic stack.
%   over $\Z_p$.
%  \mar{TG: need to sort out hypotheses.} 

Suppose also that
  the horizontal morphism $\cX \to \cY$ {\em (}which, by the usual graph
argument, is seen to be representable by algebraic stacks{\em )} is proper. Then the
  scheme-theoretic image~$\cZ$ of this morphism 
  is a $p$-adic
  formal algebraic stack of finite type. If~$\cX$ is flat
  over~$\Spf\Zp$, then so is~$\cZ$.%\mar{TG: add proof of this last fact.}
\end{aprop}
\begin{proof}
  We deduce this from~\cite[Prop.\ 10.5]{Emertonformalstacks}. % \mar{TG:
    % argument needed.}
  Note firstly that~$\cX$ is quasi-compact and
  quasi-separated, since it is of finite presentation over~$\Spf\Zp$. Examining
  the hypotheses of~\cite[Prop.\ 10.5]{Emertonformalstacks}, we need
  to show that~$\cZ$ is formal algebraic,
that it is quasi-compact and quasi-separated, that it %$\cZ$
  is Ind-locally of finite type over~$\Spec\Zp$, % in the sense
  % of~\cite[Defn.\ 8.26, Rem.\ 8.39]{Emertonformalstacks},
  and that the
  induced morphism~$\cX\to\cZ$ (which will be representable by
% \mar{ME: I think ``rep.\ by alg.\ spaces'' should be ``rep.\ by alg.\ stacks''
% here.}
  algebraic stacks, by the usual graph argument) is proper.



  By hypothesis, we can write~$\cY\cong\varinjlim_\mu\cY_\mu$, where
  each~$\cY_{\mu}$ is an algebraic stack
%locally of finite type (or
%equivalently, locally of finite presentation)  over
%  some~$\Spec\Z_p$,
of finite presentation over some~$\Z/p^a$ 
(with~$a$ depending on~$\mu$),
and the
  transition maps are closed immersions.
%Since $\cY$ is quasi-compact and quasi-separated, the same is true of
%its closed substacks $\cY_{\mu}$, and so in fact each of these algebraic
%stacks is of finite presentation over $\Spec \Z_p$.
  %thickenings.\mar{TG: I don't think we know/need this - can't we just
  %  replace it with closed immersions? ME: Done}
If we write~$\cX\cong\varinjlim_a\cX_a$, where $\cX_a := \Spec \Z/p^a\otimes_{
\Spf \Z_p} \cX$, then by assumption each $\cX_a$ is an algebraic stack.
By definition, then, we have that
\anumequation
\label{eqn:Z description}
  \cZ:=\varinjlim_a\cZ_a,
\end{equation}
where~$\cZ_a$ is the
  scheme-theoretic image of~$\cX_a\to\cY_\mu$, for~$\mu$
  sufficiently large. 
  It follows from Lemma~\ref{lem:thick images} below
  that each of the closed immersions $\cZ_a \hookrightarrow \cZ_{a+1}$
  is a finite order thickening.  We conclude
  from~\cite[Lem.\ 6.3]{Emertonformalstacks} that $\cZ$ is in
  fact a formal algebraic stack.

  Since~$\cY_\mu$ is of finite presentation over~$\Spec\Z_p$, each
  closed substack~$\cZ_a$ is of finite presentation over $\Z/p^a$, and in
  particular quasi-compact and quasi-separated; it follows that~$\cZ$
  is quasi-compact and quasi-separated.
%\mar{ME: Fix this locally Ind-finite type discussion.}
  The description~\eqref{eqn:Z description} of $\cZ$ furthermore exhibits 
  $\cZ$ as being Ind-locally of finite type over $\Spec \Z_p$.
  %We claim that it is also
  %locally Ind-finite type over~$\Spec\Zp$. Indeed, since~$\cY$ is a
  %formal algebraic stack, we can choose a smooth surjection $U\to\cY$
  %whose source is an affine formal algebraic
  %space. Setting~$V_\lambda:=U\times_\cY\cZ_\lambda$,
  %and~$V:=\varinjlim_\lambda V_\lambda$, we see that the natural
  %morphism $V\to\cZ$ is a smooth surjection, that each~$V_\lambda$
  %is of finite type over~$\Spec\Zp$, and that the transition maps are
  %thickenings; so~$\cZ$ is  locally Ind-finite type over~$\Spec\Zp$ by
  %definition.


It remains to show that~$\cX\to\cZ$ is %representable by algebraic spaces and
proper.
If $T \to \cZ$ is any morphism whose source is a scheme, then since $\cZ \to \cY$
is a monomorphism, we find that there is an isomorphism
$T\times_{\cZ} \cX \to T\times_{\cY} \cX$.   Since by 
assumption  the target of this isomorphism
is an algebraic stack, proper over $T$, 
the same is true of the source.


The claim regarding flatness follows from Lemma~\ref{lem:p-adic flatness} below.
\end{proof}
%\mar{TG: still needs a proof that if~$\cX$ is $\Zp$-flat then so is~$\cZ$}
We used the following lemmas in the proof of the Proposition~\ref{aprop: scheme theoretic image is p adic finite
    type II}.
\begin{alemma}
\label{lem:thick images}
Consider a commutative diagram
of morphisms of algebraic stacks
 $$\xymatrix{ \cX \ar[r]\ar[d] & \cZ \ar[d] \\
\cX' \ar[r] & \cZ'}$$
with $\cZ'$ being quasi-compact,
in which the left-hand vertical arrow is an $n$th order thickening
for some $n \geq 1$,
the right-hand vertical arrow is a closed immersion, and
the lower horizontal arrow is quasi-compact, surjective,
and scheme-theoretically dominant.
Then the right-hand vertical arrow is then also an $n$th order thickening. 
\end{alemma}
\begin{proof}
Since $\cX \to \cX'$ and $\cX' \to \cZ'$ are surjective,
by assumption, the same is true of their composite $\cX \to \cZ'$,
and hence of the closed immersion $\cZ \to \cZ'$; thus this closed
immersion is a thickening.  To see that it is of order~$n$, we first
note that since $\cZ'$ is quasi-compact, we find a smooth surjection
$U' \to \cZ'$ whose source is an affine scheme; since the property
of a thickening being of finite order may be checked {\em fppf}
locally, and since the property of being quasi-compact
and scheme-theoretically dominant
is also preserved by flat base-change, after pulling back
our diagram over~$U'$, we may assume that $\cZ' = U'$ is an
affine scheme, and that $\cZ = U$ is a closed  subscheme.
Let $\cI$ be the ideal sheaf on $U'$ cutting out $U$,
and let $a$ be a section of $\cI$.  If $a'$ denotes the pull-back of
$a$ to $\cX'$, then $a'_{| \cX} = 0,$ and so~$(a')^n~=~0$,  
by assumption.
Since $\cX' \to U'$ is scheme-theoretically dominant, we find that~$a^n~=~0$.
Thus $U\hookrightarrow U'$ is indeed an $n$th order thickening.
\end{proof}

\begin{alemma}
\label{lem:p-adic flatness}
If $\cX \to \cY$ is a quasi-compact scheme-theoretically dominant morphism 
of $p$-adic formal algebraic stacks which are locally of finite type, % \mar{ME: Check f.t.\ vs.\ f.p (and
% in the lemma above too).}
and if $\cX$ is flat over $\Z_p$, then the same is true of $\cY$.
\end{alemma}
\begin{proof}
Let $V \to \cY$ be a morphism which is representable by algebraic spaces and smooth,
whose source is an affine formal algebraic space; so $V = \Spf B$ 
for some $p$-adically complete $\Z_p$-algebra $B$ that is topologically
of finite type.  It suffices to show that $B$ is flat over $\Z_p$.
Since $V \to \cY$ is in particular flat (being smooth),
the base-change morphism $\cX\times_{\cY} V \to V$ is again scheme-theoretically dominant.
Since it is also quasi-compact, and since $V$ is formally affine (and
so  quasi-compact),
%\mar{TG: this is written a bit strangely: what does $V$ have to do
%  with $U\to\cX$? ME: To get $U$ we need $\cX\times_{\cY} V$ to
%be q.c., which is what we deduce here.}
we may find a formal algebraic space $U = \Spf A$ endowed with a 
morphism $U \to \cX\times_{\cY}V $ which is representable by algebraic spaces,
smooth, and surjective; and thus also scheme-theoretically dominant.
Thus we may replace our original situation with the composite morphism
$\Spf A \to \Spf B$.  But in this context, scheme-theoretic dominance
amounts to the morphism $B \to A$ being injective; thus $B$ is $\Z_p$-flat
if $A$ is.
\end{proof}

\subsection*{Versality and versal rings}
\label{subsubsec:versal}
We will sometimes find it useful to study scheme-theoretic images in
terms of versal rings, and so we will recall some notation and
results from~\cite[\S 2.2]{EGstacktheoreticimages} related to this topic. 
%specialised to the setting of Ind-algebraic stacks.

If~$\Lambda$ is a Noetherian ring,
equipped with a finite ring map $\Lambda\to k$ whose target is a
field, then we let~$\cC_{\Lambda}$ denote the category whose objects
are Artinian local $\Lambda$-algebras~$A$ equipped with an
isomorphism of~$\Lambda$-algebras $A/\m_A\isoto k$. We
let~$\pro\cC_{\Lambda}$ be the corresponding category of formal
pro-objects, which (via passage to projective limits) we identify with
the category of topological pro-(discrete Artinian) local
$\Lambda$-algebras~$A$ equipped with a $\Lambda$-algebra isomorphism
$A/\m_A\isoto k$. By~\cite[Rem.\ 2.27]{EGstacktheoreticimages}, any
morphism $A\to B$ in ~$\pro\cC_{\Lambda}$ has closed image, and
induces a topological quotient map from its source to its image, so
that in particular $A\to B$ is surjective if and only if it is induced
by a compatible system of surjective morphisms in~$\cC_\Lambda$.

%Let~$\cX$ be an Ind-algebraic stack over a locally Noetherian base
%scheme $S$,
Fix a locally Noetherian base scheme~$S$,
and let $k$ be a finite type $\cO_S$-field,
i.e.\ $k$ is a field equipped with a morphism $\Spec k\to S$ of
finite type.   
Choose, as we may,
an affine open subscheme $\Spec\Lambda\subseteq S$ for which $\Spec k\to S$
factors as $\Spec k\to\Spec\Lambda\to S$, with $\Lambda\to k$ being finite.
(In what follows we fix such a choice of $\Lambda$, although the notions
that we define in terms of it are independent of this choice.) We now
define a category fibred in groupoids $\widehat{\cF}_x$ in the
following way; this category is an example of a
%Then $\widehat{\cX}_{x}$ is a
deformation category in the sense
of~\cite[\href{https://stacks.math.columbia.edu/tag/06J9}{Tag
  06J9}]{stacks-project}. In the
notation of~\cite[\href{https://stacks.math.columbia.edu/tag/07T2}{Tag
  07T2}]{stacks-project}, our category~$\widehat{\cF}_x$ is
denoted~$\cF_{\cF,k,x}$ (with a slightly unfortunate clash of notation
in the two instances of~$\cF$).
% \mar{TG: rearranged slightly and made it a citable definition; also
%   added what I think is a correct reference to an equivalent
%   definition in the Stacks project, but leaving this marginal comment
%   in case I've done something stupid!}

\begin{adefn}\label{adefn: deformation category}
  If $\cF$ is a category fibred in groupoids over~$S$, and if
  $x:\Spec k \to \cF$ is a morphism,
  % Let $x$ be an object of $\cF$ lying over $\Spec k$.  and let
  % $\Spec\Lambda\subseteq S$ be an affine open so that $\Spec k\to S$
  % factors as $\Spec k\to\Spec\Lambda\to S$, where $\Lambda\to k$ is
  % finite.  Write $p:\cX\to S$ for the tautological morphism.
  then we define a category~$\widehat{\cF}_{x}$, cofibred in groupoids
  over~$\cC_{\Lambda}$, as
  follows: %be the category defined as follows:
  For any object $A$ of $\cC_{\Lambda}$, the objects of
  $\widehat{\cF}_x(A)$ consist of morphisms $y:\Spec A \to \cF$,
  together with an isomorphism $\alpha: x \iso \ybar$ compatible with
  the given identification of $A/\mathfrak m$ with $k$; here $\ybar$
  denotes the induced morphism $\Spec A/\mathfrak m \to \cF$.  The set
  of morphisms between two objects $(y,\alpha)$ and $(y',\alpha')$ of
  $\widehat{\cF}_x(A)$ consists of the subset of morphisms
  $\beta: y\to y'$ in $\cF(A)$ for which
  $\alpha'\circ \overline{\beta} = \alpha$; here $\overline{\beta}$
  denotes the morphism $\ybar\to \ybar'$ induced by~$\beta$.  If
  $A \to B$ is a morphism in $\cC_{\Lambda}$, then the corresponding
  pushforward $\widehat{\cF}_x(A) \to \widehat{\cF}_x(B)$ is defined
  by pulling back morphisms to $\cF$ along the corresponding morphism
  of schemes $\Spec B \to \Spec A$.
\end{adefn}
%
%\begin{enumerate}
%\item
%Objects of $\widehat{\cF}_x$ are morphisms
%$x\to x'$ in $\cF$ such that $p(x')=\Spec A$, with
%  $A$ an Artinian local ring, and the morphism $\Spec k\to\Spec A$
%  given by $p(x)\to p(x')$ corresponds to a ring homomorphism $A\to k$
%  identifying $k$ with the residue field of $A$, and
%\item Morphisms in $\widehat{\cF}_x$ 
%morphisms $(x\to x')\to(x\to x'')$ are commutative diagrams in $\cX$ of
%  the form \[\xymatrix{x'&& x''\ar[ll]\\& x\ar[ul]\ar[ur]&}\]
%\end{enumerate}


%Recall that
%the category $\widehat{\cX}_{x}$ naturally extends to its
%completion, which by definition is the pro-category $\pro\widehat{\cX}_{x}$, which is a category cofibred in groupoids over
%$\pro{\cC}_{\Lambda}$.
%We now define what it means for a morphism
%$\Spf A_x \to \widehat{\cX}_x$  to be versal. 

If $\cF$ is a category fibred in groupoids over the locally Noetherian scheme $S$,
and if $x: \Spec k \to \cF$ is a $k$-valued point of~$\cF$,
for some finite type $\cO_S$-field,
then the notion of a versal ring to $\cF$ at $x$ 
is defined in~\cite[Def.~2.2.9]{EGstacktheoreticimages}.
Rather than recalling that definition here,
we will give a definition of versality in a greater level of generality that is convenient
for us.  We will then explain how the notion of versal ring is obtained as a particular
case.

\begin{adefn}
\label{def:versal}\index{versal}
Let $f:\cF \to \cG$ be a morphism of categories fibred in groupoids over the
locally Noetherian scheme~$S$, let $k$ be a finite type $\cO_S$-field,
let $x: \Spec k \to \cF$ be a $k$-valued point of~$\cF$,
%and let $y:= f\circ x: \Spec k \to \cF$
%denote the induced $k$-valued point of~$\cG$.
and let $y: \Spec k \to \cG$
be a $k$-valued point of~$\cG$ equipped with an isomorphism of $k$-valued points
$\alpha: f\circ x \iso y$.
%lying over the identity automorphism of~$k$.
%and let $y := f \circ x: \Spec k \to \cG$ denote the induced $k$-valued
%point of~$\cG$.
The morphism $f$ induces in an evident way a morphism 
$\widehat{f}_x: \widehat{\cF}_x \to \widehat{\cG}_y$
of categories cofibred in groupoids.
We say that $f$ is {\em versal} at $x$ 
if the the morphism $\widehat{f}_x$ is {\em smooth}
in the sense 
of~\cite[\href{https://stacks.math.columbia.edu/tag/06HG}{Tag 06HG}]{stacks-project};
that is, given a commutative diagram
$$\xymatrix{\Spec B \ar[r] \ar[d] & \Spec A \ar[d] \ar@{-->}[ld]\\
\widehat{\cF}_x \ar^-{\widehat{f}_x}[r] & \widehat{\cG}_{x}}$$
in which the upper arrow is the closed immersion corresponding
to a surjection $A \to B$ in $\cC_{\Lambda}$, 
 we can fill in the
dotted arrow 
% \mar{ME: I don't understand the parenthetical comment, and will delete
% it unless I'm missing something. TG: me neither, removing}
% (with a morphism coming from a morphism in $\pro\cC_\Lambda$)
so that the diagram remains commutative.
(Clearly this notion is independent of the particular choice of $y$ and $\alpha$;
more precisely
it holds for any such choice if it holds for one such choice, such as $y = f\circ x$
and $\alpha = \id$.)
\end{adefn}

%\begin{adefn}
\begin{aexample}
Let $S$ be a locally Noetherian scheme,
let $k$ be a finite type $\cO_S$-field~$k$,
and let $A$ be an object of $\pro\cC_{\Lambda}$;
recall in particular then that $A$ comes equipped with a chosen isomorphism
$\Spec A/\mathfrak m \iso k.$
We let $x': \Spec k \iso \Spec A/\mathfrak m \hookrightarrow \Spf A$
denote the induced $k$-valued point of $\Spf A$.

Now let $\cF$ be a category fibred in groupoids over~$S$,
let $x:\Spec k \to \cF$ be a $k$-valued point of $\cF$,
and suppose that $f:\Spf A \to \cF$ is a morphism, for which we can
find an isomorphism $f\circ x' \iso x$ of $k$-valued points of~$\cF$.
%Then we say that $f$ realizes $A$ as a versal ring to $\cF$ at $x$
%if and only if $f$ is versal at the point~$x'$.
Then the morphism $f$ is versal at the point $x'$ if and only
if $A$ is a versal ring to the $k$-valued point $x$ of $\cF$
in the sense of~\cite[Def.~2.2.9]{EGstacktheoreticimages}
%\end{adefn}
\end{aexample}

\begin{aexample}
Let $S$ be a locally Noetherian scheme, let $U$ be a locally finite type $S$-scheme,
and let $f:U \to \cF$ be a  morphism to a category fibred in groupoids over~$S$.
If $u\in U$ is a finite type point, giving rise to a morphism
$\Spec \kappa(u) \to U$, then $f$ is versal at this $\kappa(u)$-valued
point of $U$ if and only if $f$ is versal at the point $u$
in the sense of~\cite[Def.~2.4.4~(1)]{EGstacktheoreticimages};
cf.~\cite[Rem.~2.4.5]{EGstacktheoreticimages}.
\end{aexample}

%\begin{adefn}
%\label{def:versal ring}
%Let $A_x$ be a topological local $\Lambda$-algebra corresponding 
%to an element of $\pro\cC_{\Lambda}$.  We say that a morphism
%$\Spf A_x \to \widehat{\cX}_x$ is {\em versal} if it satisfies the infinitesimal lifting property with
%  respect to morphisms in $\cC_{\Lambda}$. More precisely,
%given a commutative diagram
%$$\xymatrix{\Spec A \ar[r] \ar[d] & \Spec B \ar[d] \ar@{-->}[ld]\\
%\Spf A_{x} \ar[r] & \widehat{\cX}_{x}}$$
%in which the upper arrow is the closed immersion corresponding
%to a surjection $B \to A$ in $\cC_{\Lambda}$, 
%and the left hand
%vertical arrow corresponds to a morphism in
%$\pro\cC_\Lambda$ (equivalently, it is continuous when~$A_x$ is given
%its projective limit topology, and~$A$ is given the discrete topology)  we can fill in the
%dotted arrow (with a morphism coming from a morphism in $\pro\cC_\Lambda$) so that the diagram remains commutative.
%\end{adefn}

%We say that a morphism
%$\Spf A_x \to \cX$ is versal if the corresponding morphism
%$\Spf A_x \to \widehat{\cX}_x$ is versal.

\begin{arem}\label{arem: versal rings for formal Ind}
  If~$\cX$ is an algebraic stack which is locally of finite
  presentation over a locally Noetherian scheme~$S$, then it admits
  (effective, Noetherian) versal rings at all finite type points~\cite[\href{https://stacks.math.columbia.edu/tag/0DR1}{Tag 0DR1}]{stacks-project}. If
  $X$ is an Ind-locally finite type algebraic space over~$S$, then it
  admits (canonical) versal rings at all finite type points,
  by~\cite[Lem.\ 4.2.14]{EGstacktheoreticimages}.
%it follows easily\mar{TG: cite something from proper about versality and
%    smoothness? also want to relate this condition to the definition
%    of a formal algebraic space?} that if  $\cX$ is a formal algebraic
%  stack which is Ind-locally finite type over~$S$,\mar{TG: haven't
%    thought about whether this notion makes sense, not going to worry
%    about it for now as we don't use this remark anyway.} then it
%  admits (canonical) versal rings at all finite type points.
%
%  The existence of versal rings for
%  Ind-algebraic stacks is less clear (to us, at least),
  We don't prove a general statement about the existence of versal
  rings for Ind-algebraic stacks, since in all the cases we consider
  in the body of the book we are able to construct them ``explicitly''
   (for example, in terms of Galois lifting rings, or lifting rings for
  \'etale $\varphi$-modules). %\mar{TG: or maybe
   % you can formally take some limit, at least in a filtered case?
   % Haven't thought about this, or about the formal algebraic case, so
   % I'm hedging my bets - I think we don't need a general statement. ME: Agreed}
\end{arem}


%Let~$A$ be an object of~$\cC_{\Lambda}$, and let $X\to \Spf A$ be a
%finite type morphism of formal algebraic spaces. Then we can define
%the scheme-theoretic image of $X\to \Spf R$ as in~\cite[Defn.\
%3.2.15]{EGstacktheoreticimages}: if we let~$A_i$ be any discrete
%Artinian quotient of~$A$, then we let $\Spec B_i$ be the
%scheme-theoretic image of $X_{A_i}:= X\times_{\Spf A} \Spec A_i\to \Spec A_i$. We
%set~$B=\varprojlim_i B_i$, so that $\Spf B$ is a formal subscheme
%of~$\Spf A$; since the morphisms $A_i\to B_i$ are surjective, it
%follows from the discussion above that $A\to B$ is a quotient map of
%topological rings, so that $\Spf B$ is a closed formal subscheme
%of~$\Spf A$.

The following lemma and its proof are essentially \cite[Lem.\
3.2.16]{EGstacktheoreticimages}, but since the setup there is
different, we give the details here.  Before stating the lemma, 
we introduce the setup. Let~$S$ be a locally Noetherian scheme.  We suppose given a morphism
  $\cX \to \cY$ of stacks over~$S$ which is representable by algebraic
stacks and proper, % \mar{TG: could add that $\cX\to\cY$ is locally of
  % finite presentation, although I'm not sure we need it in order to
  % deduce what we need for~$\cZ_\lambda$? We have it in
  % our only application of this result, where the morphism is also
  % representable by algebraic spaces, not just stacks}
and that $\cY$ is an Ind-algebraic stack which may
be written as the $2$-colimit of %quasi-compact and quasi-separated 
algebraic stacks which are of finite presentation over $S$ (and so in particular
quasi-compact and quasi-separated) with respect 
to transition morphisms that are closed immersions,  
say $\cY \cong \varinjlim \cY_{\lambda}.$
Then $\cX_{\lambda} := \cX \times_{\cY} \cY_{\lambda}$ is an algebraic stack,
and the projection $\cX_{\lambda} \to \cY_{\lambda}$ is proper, so that 
$\cX_{\lambda}$ is finite type over $S$ (and in particular quasi-compact,
and also quasi-separated, although we won't
use this latter fact).   Furthermore, we have an induced isomorphism
$\cX \cong \varinjlim_{\lambda} \cX_{\lambda}$, and so $\cX \to \cY$ 
is a morphism of Ind-algebraic stacks whose scheme-theoretic image~ $\cZ$ may be
defined.  Of course, in this context, since the Ind-structures on $\cX$
and $\cY$ are compatible, if we let $\cZ_{\lambda}$ denote the scheme-theoretic
image of $\cX_{\lambda}$ in $\cY_{\lambda}$, then this coincides with the
scheme-theoretic image of $\cX_{\lambda}$ in $\cY_{\lambda'}$ for any $\lambda' \geq
\lambda$, and we may write $\cZ \cong \varinjlim_{\lambda}
\cZ_{\lambda}$. Note that~$\cZ_\lambda$ is also of finite
presentation over~$S$, by Lemma~\ref{lem:finite presentation from closed immersion}.
%\mar{TG: add a couple of words of explanation? ME: Would need
%$Y_{\lambda}$ to be l.f.p.\ for this to hold automatically, so I'll just make
%this assumption.}

\begin{alemma}%\% mar{TG: need some finiteness condition, at least that
%     the $\cZ_\lambda$ are locally of finite presentation over a
%     locally Noetherian base~$S$. Leaving this for now as it would
%     probably be best to check that there aren't any other issues with
%     this argument! TG: I didn't see any other issues, so I think this
%     is exactly what we need? This lemma is only used twice in the body
%   of the paper, in fact. ME: I don't think q.c.q.s.\ makes sense for Ind-alg.\ stacks
% in general (or at least, is not defined for them), but what we need is that
% each $\cX_{\lambda}$ is q.c.q.s., to get the claimed factorization. ME: Started to
% tidy this up, but haven't fully come to grips with these comments yet.  Also
% haven't rewritten argument to reflect my changes in the phrasing of the lemma
% and its preamble. ME: Fixed some conflicting notation in the proof, and now everything
% is done (I hope) except the finiteness stuff.}
  \label{alem: scheme theoretic image of versal is versal}
%Suppose that
%  $\cX \to \cY$ is a morphism of stacks which is representable by algebraic
%stacks and proper, and that $\cY$ is an Ind-algebraic stack which may
%be written as the $2$-colimit of quasi-compact and quasi-separated 
%algebraic stacks with respect
%to transition morphisms that are closed immersions.  
%Then $\cX$ has a similar structure as an Ind-algebraic stack,
%so that the scheme-theoretic image $\cZ$ of $\cX$ in $\cY$ is defined.
%%  which is quasi-compact and quasi-separated, and let~$\cZ$ denote the
%%  scheme-theoretic image of this morphism.
Suppose that we are in the preceding situation, so that   $\cX \to
\cY$ is a morphism of stacks over a locally Noetherian base~$S$ which is representable by algebraic
stacks and proper, where $\cY$ is an Ind-algebraic stack which may
be written as the $2$-colimit of 
algebraic stacks which are of finite presentation over~$S$, with respect 
to transition morphisms that are closed immersions, and we write~$\cZ$
for the scheme-theoretic image of $\cX\to\cY$.


Suppose that~$x:\Spec k\to
  \cZ$ is a finite type point, and that $\Spf A_x\to\cY$ is a versal
  morphism for the composite $x:\Spec k\to
  \cZ\into\cY$.
Let~$\Spf B_x$ be the scheme-theoretic image of $\cX_{\Spf A_x}\to\Spf
A_x$. Then the morphism $\Spf B_x\to\cY$ factors through a versal
morphism $\Spf B_x\to \cZ$.
\end{alemma}
\begin{proof}
  By definition, we may write~ $A_x=\varprojlim A_i$,
  $B_x=\varprojlim B_i$, where the~$A_i$ are objects of~$\cC_\Lambda$,
  and $\Spec B_i$ is the scheme-theoretic image of
  $\cX_{A_i}\to \Spec A_i$. As explained immediately above, we write
  $\cX \cong \varinjlim \cX_{\lambda}$, $\cY\cong\varinjlim\cY_{\lambda}$, and
  $\cZ\cong\varinjlim \cZ_{\lambda}$, where the transition morphisms are
  %thickenings
  closed immersions
  of algebraic stacks, and~$\cZ_\lambda$ is the
  scheme-theoretic image of the morphism
  $\cX_{\lambda}\to\cY_{\lambda}$; in particular, the morphism $\cX_\lambda\to\cZ_\lambda$ is proper and scheme-theoretically
  dominant. % If we write $\cY \cong \varinjlim \cY_{\mu}$, then for
  % each $\lambda,$ the morphism $\cX_{\lambda}\to \cY$ factors through
  % $\cY_{\mu}$ for any sufficiently large~$\mu$, and the
  % scheme-theoretic image of this morphism is~$\cZ_\lambda$.

  If follows that for each~$i$, the composite $\Spec B_i\to\Spec A_i\to\cY$
  factors through~$\cZ$. Indeed, for~$\lambda$ sufficiently large the
  morphism $(\cX_\lambda)_{B_i}\to\Spec B_i$ is scheme-theoretically
  surjective, and  the morphism $(\cX_\lambda)_{B_i}\to\cY$ factors
  through $\cZ_\lambda$, so the morphism $\Spec B_i\to\cY$ also
  factors through~$\cZ_\lambda$.

We now show that a morphism $\Spec A\to\Spf A_x$, with $A$ an object of
$\cC_{\Lambda}$, 
% \mar{ME: Here and below in this argument, I replaced $\Lambda_\F$ by
%   $\cC_{\Lambda}$, which I think is our standard notation in the general
%   setting. TG: thanks; commenting out.}
factors through $\Spf B_x$ if and only if the composite $\Spec
A\to\Spf A_x\to\cY$ factors through~$\cZ$. In
one direction, if
$\Spec A\to\cY$ factors through~$\Spf B_x$, then it factors
through $\Spec B_i$ for some $i$, and hence 
through $\cZ$, as we saw above.


Conversely, if the composite $\Spec
A\to\Spf A_x\to\cY$ factors through~$\cZ$, then we claim 
that there exists factorisation $\Spec A\to\Spec B\to\cY$, where $B$ is
an object of~$\cC_{\Lambda}$, the morphism $\Spec
A\to\Spec B$ is a closed immersion, and $\cX_B\to \Spec B$ is
scheme-theoretically dominant. To see this, % \mar{TG: I was originally going to cite 3.2.19 in
  % proper.tex, but that was a pain to verify. Hopefully this is a
  % correct argument!}
note firstly that the composite $\Spec A\to\cY$ factors
through~$\cZ_\lambda$ for some~$\lambda$. Since~$\cZ_\lambda$ is
an algebraic stack which is locally of finite presentation over a
locally Noetherian base, % \mar{TG: I believe this is the one thing I was
  % referring to above - do we want to make it explicit that this is
  % locally of finite presentation over a locally Noetherian base, and
  % give a reference for this fact?} 
it admits effective Noetherian versal rings by Remark~\ref{arem: versal rings for formal Ind}, so
that $\Spec A\to\cZ_\lambda$  factors through a versal morphism $\Spec C_x\to\cZ_\lambda$ at the
finite type point of~$\cZ_\lambda$ induced by~$x$.  By ~\cite[Lem.\
1.6.3]{EGstacktheoreticimages}, we can find a factorisation  $\Spec
A\to\Spec R_x\to\Spec C_x$, where~$R_x$ is complete local Noetherian,
$\Spec R_x\to\Spec C_x$ is faithfully flat, and $\Spec A\to\Spec R_x$ is a closed
immersion. Since $\Spec R_x\to\Spec C_x$ is faithfully flat, and $\Spec
C_x\to\cZ_\lambda$ is flat
by~\cite[\href{https://stacks.math.columbia.edu/tag/0DR2}{Tag
  0DR2}]{stacks-project}, we see that the base changed morphism
$(\cX_\lambda)_{R_x}\to\Spec R_x$ is scheme-theoretically
dominant.  By~\cite[Lem.\ 3.2.4]{EGstacktheoreticimages} (applied to the proper
morphism of algebraic stacks~$\cX_\lambda\to\cZ_\lambda$), $R_x$
admits a cofinal collection of Artinian quotients~$R_i$ for which
$(\cX_\lambda)_{R_i}\to\Spec R_i$ is scheme-theoretically
dominant. The closed immersion $\Spec A\to\Spec R_x$ factors through
~$\Spec R_i$ for some~$R_i$, so we may take~$B=R_i$.

Now, by the versality of $\Spf A_x\to\cY$, we may lift the morphism
$\Spec B\to\cY$ to a morphism $\Spec B\to\Spf A_x$, which furthermore
we may factor as $\Spec B \to \Spec A_i \to \Spf A_x$, for some value
of $i$.  Since $\cX_B \to \Spec B$ is scheme-theoretically dominant,
the morphism $\Spec B \to \Spec A_i$ then factors through $\Spec B_i$,
%We may then
%factor the morphism $\cX_B\to\Spec B$ as the
%composite \[\cX_{B}\to \Spf B_x\times_{\Spf A_x}\Spec B\to\Spec B.\] The second
%arrow is then a scheme-theoretically dominant closed immersion, so it is an
%isomorphism, and $\Spec B\to\Spf A_x$ factors
and thus through $\Spf B_x$, as required.


It follows in particular that the composite
$\Spf B_x\to\Spf A_x\to\cY$ factors through a morphism
$\Spf B_x\to\cZ$. It remains to check that this morphism is
versal. This is formal. Suppose we are given a commutative
diagram \[\xymatrix{\Spec A_0\ar[r]\ar[d]&\Spf B_x\ar[r]\ar[d]&\Spf A_x\ar[d]\\
    \Spec A\ar[r]&\cZ\ar[r]&\cY }\]where the left hand vertical arrow
is a closed immersion, and $A_0, A$ are objects
of~% \mar{TG: should we change to $\Spec A_0,\Spec A$? I think this would
% be more consistent with what we write elsewhere. ME: Yes, I think that would 
% be sensible. TG: done, commenting out.}
$\cC_\Lambda$. By the versality of $\Spf A_x\to\cY$, we may lift the
composite $\Spec A\to\cY$ to a morphism $\Spec A\to\Spf A_x$. Since
the composite $\Spec A\to\Spf A_x\to\cY$ factors through
$\cZ$, the morphism $\Spec A\to \Spf A_x$ then factors through $\Spf
B_x$, as required.
\end{proof}

The following lemma records a useful property of versal rings in the Noetherian
context.

\begin{alem}
\label{lem:versal flatness}
Let $\cX$ is a locally Noetherian formal algebraic stack, let $R$ be a complete
Noetherian local ring with residue field~$k$,
and let
$f:\Spf R \to \cX$ be a morphism for which the induced morphism
$x:\Spec k \to \cX_{\red}$ is a finite type point, and which
is versal to $\cX$ at $x$.   Then the morphism $f$ is flat,
in the sense of~{\em \cite[Def.~8.42]{Emertonformalstacks}}.  
\end{alem}
\begin{proof}
Since $\cX$ is locally Noetherian, by assumption,
we may find a morphism $\Spf A \to \cX$ whose source
is a Noetherian affine formal algebraic space, which is representable 
by algebraic spaces and smooth, and whose image contains the point~$x$.
By definition, we have to show that the base-changed morphism
$\Spf A \times_{\cX} \Spf R \to \Spf A$ is flat.  

The fibre product $\Spf A \times_{\cX} \Spf R$ is a formal algebraic space
(which is locally Noetherian, since $\Spf R$ is
so~\cite[\href{https://stacks.math.columbia.edu/tag/0AQ7}{Tag 0AQ7}]{stacks-project}),
and so admits a morphism $\coprod \Spf B_i \to \Spf A \times_{\cX} \Spf R$
whose source is the disjoint union of Noetherian affine formal algebraic spaces,
and which is representable by algebraic spaces, \'etale, and surjective.
Again, it suffices to show that each of the morphisms $\Spf B_i \to \Spf A$
is flat, which by definition is equivalent to each of the induced morphisms $A \to B_i$
being flat. 
% \mar{ME: We had a proof of this written out in some iteration of [CEGS] I think; happy
% to include that or not, or maybe it's in [SP] --- I forget now, sorry!
% TG: isn't it proved in the proof of~\cite[Lem.\ 8.18]{Emertonformalstacks}?}
If we let $I$ denote an ideal of definition of the topology on $A$,
then it suffices to show that each morphism $A/I^n \to B_i/I^n$ is
flat (see e.g.\ the proof of~\cite[Lem.\ 8.18]{Emertonformalstacks}).

For this, it suffices in turn to show, for each maximal ideal $\mathfrak n$
of $B_i$,
that the induced 
morphism $A/I^n \to
(B_i/I^n)\, \hat{}$
is flat (where $(\text{--})\, \hat{}$ denotes
$\mathfrak n$-adic completion; note that if $\mathfrak m$ denotes the maximal ideal of~$R$,
then the topology on $B_i$ is the $\mathfrak m$-adic topology,
so that any maximal ideal $\mathfrak n$ of $B_i$ contains~$\mathfrak m$,
and also contains~$IB$). 
Now the induced morphism $\Spf \widehat{B}_i \to \Spf A$ is versal 
to the induced morphism $\Spec B_i/\mathfrak n \to \Spf A$ 
(as one sees by chasing through
the constructions, beginning from the versality of $\Spf R \to \cX$,
and taking into account the infinitesimal lifting property for \'etale
morphisms),
and thus the induced morphism
$\Spf (B_i/I^n)\, \hat{} \to \Spec A/I^n$ is also versal
(to the induced morphism
$\Spec B_i/\mathfrak n \to \Spec A/I^n$).
%(as one sees by chasing through
%the constructions, beginning from the versality of $\Spf R \to \cX$),
%\mar{ME: Should think through this versality claim, just because of possible
%residue field issues. ME: Seems fine, commenting out.}
Thus this morphism {\em is} flat, e.g.\
by~\cite[\href{https://stacks.math.columbia.edu/tag/0DR2}{Tag 0DR2}]{stacks-project},
and the lemma is proved.
\end{proof}


We will frequently find the following lemma useful.
\begin{alem}\label{lem: criterion for Artin to map to scheme theoretic
    image}% }\mar{TG: should be checked that I haven't blundered in the
    % formulation and proof, and that it is actually general enough for
    % our applications. If it's actually correct I guess we could
    % consider just putting this formulation into proper.tex\dots}
	Let $R \to S$ be a continuous surjection of objects in
	$\pro\cC_{\Lambda}$,
	and let $X\to \Spf R$ be a finite type morphism
	%whose source is a formal algebraic space,
	of formal algebraic spaces. % ,
	% %\mar{ME: Using some formal algebraic space terminology here,
	% %	perhas cite/explain it.}
	% and if $A$ is any discrete Artinian quotient of $R$
	% {\em (}so that there is a morphism $\Spec A \to \Spf R${\em )},
	% write $X_A := X\times_{\Spf R} \Spec A \to \Spec A$.
        
	Make the following assumption: 
	if $A$ is any finite-type Artinian local $R$-algebra
	for which the canonical morphism $R\to A$ factors through 
	a discrete quotient of $R$,
	and for which the canonical morphism $X_A \to \Spec A$
	admits a section,
	then the canonical morphism $R \to A$ furthermore factors
        through $S$.

        
        Then the scheme-theoretic image of~$X\to\Spf R$ is a closed
        formal subscheme of~$\Spf S$.
\end{alem}
\begin{proof}Writing~$\Spf T$ for the scheme-theoretic image of $X\to
  \Spf R$, we need to show that the surjection $R\to T$ factors
  through~$S$. By definition, we can
write~$T$ as an inverse limit of discrete Artinian quotients~$B$
for which the canonical morphism $X_B \to \Spec B$ is
scheme-theoretically dominant. It follows from \cite[Lem.\
5.4.15]{EGstacktheoreticimages} that for any such~$B$, the surjection
~$R\to B$ % \mar{TG: $A$ here should be $B$?}
factors through~$S$; so the surjection $R\to T$ factors
through~$S$, as required.
\end{proof}

% \mar{TG: older version from a few weeks ago; leaving it here for now
%   as I don't know if I was doing something sensible!}
% \begin{alem}\label{lem: criterion to contain scheme theoretic image
%     versal ring}Let~$R$ be a pro-Artinian topological local ring,
% let~$S$ be a
%   quotient of~$R$ by a closed ideal, \mar{TG: check
%     that this isn't too general - I'm imagining making the argument in
%     $\pro\cC_{\Lambda}$ with~$\Lambda=\Z$?  }
% and let $X\to \Spf R$ be a proper morphism
% 	%whose source is a formal algebraic space,
% 	of formal algebraic spaces.
%       If $A$ is any discrete Artinian quotient of $R$
% 	{\em (}so that there is a morphism $\Spec A \to \Spf R${\em )},
% 	write $X_A := X\times_{\Spf R} \Spec A \to \Spec A$.

% 	Make the following assumption: 
% 	if $A$ is any finite-type Artinian local $R$-algebra
% 	for which the canonical morphism $R\to A$ factors through 
% 	a discrete quotient of $R$,
% 	and for which the canonical morphism $X_A \to \Spec A$
% 	admits a section,
% 	then the canonical morphism $R \to A$ furthermore factors
%         through $S$.

%         Then the scheme-theoretic image of~$X\to\Spf R$ is a closed
%         sub-formal scheme of~$\Spf S$.
% \end{alem}


 \subsection*{Immersed substacks} 
 %\mar{ME: Need to add some references to {\tt formal-stacks.tex},
 %and perhaps slightly more explanation. ME: Just read over this,
 %and it seems fine for now; removing this comment.}

The following lemma is often useful for studying morphisms from a
stack in terms of morphisms from a cover of the stack.
	\begin{alemma}\cite[Lem.\ 3.18, 3.19]{Emertonformalstacks}
		\label{lem:stacks as quotients}
		Suppose that $\cX$ is a stack over a base scheme~$S$,
		%whose diagonal is representable by algebraic spaces,
		and that $U\to \cX$ is a morphism over~$S$ which is
                representable by algebraic spaces and whose source is
                a sheaf. Write $R := U\times_{\cX} U$.

                \begin{enumerate}
                \item Suppose that the projections
                  $R := U\times_{\cX} U \rightrightarrows U$ {\em
                    (}which are again representable by algebraic
                  spaces, being the base-change of morphisms which are
                  so representable{\em )} are flat and locally of
                  finite presentation.  Then the morphism $U\to \cX$
                  induces a monomorphism $[U/R]\to \cX$.

                 
                \item  Suppose that $U\to \cX$ is flat, surjective, and
                  locally of finite presentation.  Then the morphism
                  $U\to \cX$ induces an isomorphism $[U/R]\iso \cX$.
                \end{enumerate}


	\end{alemma}

In particular, if $\cX$ is a formal algebraic stack (over some base scheme $S$),
 then by definition there exists a morphism $U\to \cX$ 
 which is representable by algebraic spaces, smooth, and surjective,
 and whose source~$U$ is a formal algebraic space.
 In this case $R := U\times_{\cX}U$ is a formal algebraic
 space which is endowed in a natural way with the structure 
 of a groupoid in formal algebraic spaces over $U$, and (since smooth
 morphisms are in particular flat), by Lemma~\ref{lem:stacks as quotients} there is an isomorphism of stacks $[U/R] \iso \cX$. %(\cite[Lem.\ 5.25]{Emertonformalstacks}).
 
 Suppose now that $\cZ \hookrightarrow \cX$ is an immersion (in
the sense that it is representable by algebraic spaces,
and pulls back to an immersion over any test morphism $T \to \cX$
whose source is a scheme). % \mar{TG:
   % I don't think the general notion of an immersion of formal
   % algebraic stacks is defined here, or indeed in
   % \cite{Emertonformalstacks}, so probably there should be a
   % definition inline here? ME: Is what I just put in ok? TG: yes,
   % commenting out.}
 The
 induced morphism
 $W := U\times_{\cX} \cZ \hookrightarrow U$ is then an immersion of formal
 algebraic spaces, and $W$ is $R$-invariant, in (an evident generalization
 of) the sense 
 of~\cite[\href{http://stacks.math.columbia.edu/tag/044F}{Tag
  044F}]{stacks-project}. 
 Conversely, if $W$ is an $R$-invariant locally closed formal 
 algebraic subspace of $U$, and if we write $R_W := R\times_U W =
 W \times_U R$ for the restriction of $R$ to $W$, then
 the induced morphism $[W/R_W] \to [U/R] \iso \cX$ is an immersion.
 (In the context of algebraic stacks, this 
 is~\cite[\href{http://stacks.math.columbia.edu/tag/04YN}{Tag
  04YN}]{stacks-project}.)


\subsection*{Flat parts}Let~$\cO$ be the ring of integers in a
finite extension of~$\Qp$, and let~$\cX\to\Spf\cO$ be a $p$-adic formal
algebraic stack which is locally of finite type over~$\Spf\cO$. Then
by~\cite[Ex.\ 9.11]{Emertonformalstacks}, there is a closed substack~
$\cX_{\fl}$ of~$\cX$, the $\emph{flat part}$ of~$\cX$, which is the maximal substack of~$\cX$ which
flat over~$\Spf\cO$.  \index{$p$-adic formal algebraic stack!flat part}
% \mar{ME: This can also be done using the previous subsubsection ---
% 	take a cover of $\cX$ by $U = $ a disjoint union of $\Spf A$'s,
% 	take the flat part of each one, and check $R$-invariance. 
% 	Details to be added.}

More precisely, $\cX_{\fl}\to\Spf\cO$ is flat, % (since this morphism is
% representable by algebraic stacks, this has the usual meaning),
and if
$\cY\to\cX$ is a morphism of locally Noetherian formal algebraic
stacks for which the composite $\cY\to\Spf\cO$ is flat, then
$\cY\to\cX$ factors through~$\cX_{\fl}$.

% By~\cite{Emertonformalstacks}, a morphism $T\to\cX$ whose source is a
% scheme factors through~$\cX_{\fl}$ if and only if for any morphism
% $U\to\cX$ which is representable by algebraic spaces and smooth, and
% whose source is an affine formal algebraic space $U=\Spf A$, the base
% changed morphism $T\times_{\cX}U\to U$ factors through the $\cO$-flat
% part~$\Spf A_{\fl}$ of~$U$, where~$A_{\fl}$ is the quotient of~$A$ by
% its ideal of $p$-power torsion elements.% \mar{TG: not sure if we need
%   % this paragraph}




% We let~$\cS$ be a formal algebraic stack,
% we let~$\cD$ be the 2-category of formal algebraic stacks over~$\cS$,
% and we let~$\cC$ be a strong 2-subcategory of~$\cD$ (``strong''
% meaning that for any two objects of~$\cC$, the $1$-category of
% morphisms between these objects in~$\cC$ coincides with the
% corresponding $1$-category of morphisms in~$\cD$).  We assume
% that~$\cC$ has the following property:
% \begin{itemize}
% \item if~$\cX\to\cS$ is an object of~$\cC$, and if~$\cU\to \cX$ is a
%   morphism of formal algebraic stacks which is representable by
%   algebraic spaces and smooth, then the composite $\cU\to\cX\to\cS$ is
%   an object of~$\cC$.
% \end{itemize}
% We let~$\Aff_{\cC}$ be the full subcategory of~$\cC$ whose objects are
% formal algebraic spaces, and write $i:\Aff_\cC\to\cD$ for the
% inclusion. We suppose that we are given a functor~$F\Aff_\cC\to\cD$
% together with a natural transformation $F\to i$
% with the properties that:
% \begin{itemize}
% \item $F\to i$ is a closed immersion (in the sense that for any
%   object~$U$ of~$\Aff_{\cC}$, $F(U)\to U$ is a closed immersion).
% \item If $U\to V$ is a morphism in~$\Aff_\cC$ which is representable
%   by algebraic spaces and smooth, then the natural morphism $F(U)\to
%   U\times_V F(V)$ is an isomorphism.
% \end{itemize}
% \begin{aprop}
%   \label{prop: summary of section 9 of formal stack}Under the above
%   hypotheses, for each object~$\cX$ of~$\cC$, there is a closed
%   substack~$\tF(\cX)$ of~$\cX$, which is characterised by the property
%   that a morphism $T\to \cX$ whose source is a scheme factors
%   through~$\tF(\cX)$ if and only if   for each morphism $U\to\cX$ which is representable by algebraic
%   spaces and smooth, and whose source is an affine formal algebraic
%   space, the morphism $T\times_{\cX}U\to U$ factors through~$F(U)$.  
% \end{aprop}
% \begin{proof}
%   See~\cite[Lem.\ 9.2]{Emertonformalstacks}.
% \end{proof}

% Let~$\cS=\cC_{d,\semis,h}$, and let~$\cC$ denote
% the category of locally Noetherian formal algebraic stacks
% over~$\cC_{d,\semis,h}$, so that an object of~$\Aff_\cC$ is a morphism
% $\Spf B\to \cC_{d,\semis,h}$, where~$\Spf B$ is a Noetherian affine
% formal algebraic space. In particular, $B$ is complete with respect to
% an ideal containing~$pB$. We write~$B_\fl$ for the quotient of~$B$ by
% its ideal of~$p$-power torsion elements. Then we set
% $F(\Spf B)=\Spf
% B_\fl\times_{\cC_{d,\semis,h}}\cC_{d,\semis,h}$.\mar{TG: pst notation
%   and so on. Note that this definition looks a bit stupid, but I'm
%   guessing the point is going to be that the HT weight locus (and
%   perhaps the type locus too, depending on how we define it) is not
%   closed unless we're in the flat setting.}
% \mar{TG: fix $h$ and~$L$ but say ultimately independent of
%   choice. Similarly deal with twisting somewhere. Put in a lemma for this!}

% To see that~$F$ satisfies the required properties, we must show
% that~$F(\Spf B)\to \Spf B$ is a closed immersion, and that if~$\Spf
% C\to\Spf B$ is representable by algebraic spaces and smooth, then
% $F(\Spf C)=\Spf C\times_{\Spf B}F(\Spf B)$. 

\subsection*{Obstruction theory}Let~$\cX$ be a limit preserving Ind-algebraic stack
over a locally Noetherian scheme $S$. If $x:\Spec A \to \cX$ is a
morphism for which the composite morphism $\Spec A \to S$ factors
through an affine open subscheme of~$S$, then
in~\cite[\href{http://stacks.math.columbia.edu/tag/07Y9}{Tag
  07Y9}]{stacks-project} there is defined a functor $T_x$ from the
category of $A$-modules to itself, whose formation is also functorial
in the pair
$(x,A)$~\cite[\href{http://stacks.math.columbia.edu/tag/07YA}{Tag
  07YA}]{stacks-project}, and such that for any finitely generated $A$-module $M$, there
is a natural identification of $T_x(M)$ with the set of lifts of $x$
to morphisms $x': \Spec A[M] \to \cX$ (where $A[M]$ denotes the square
zero extension of $A$ by $M$). (These definitions apply to~$\cX$
by~\cite[Lem.\ 4.22]{Emertonformalstacks}.)
%\mar{ME: Finish this by recalling definition of $T_x(M)$ and related stuff.
%ME: Done.}

We make the following definition, which is a special case
of~\cite[Defn.\ 11.6]{Emertonformalstacks}; that definition
incorporates an auxiliary module in its definition for technical
reasons, but in our applications this module is zero, so we have
suppressed it here. We have also incorporated~\cite[Rem.\
11.9]{Emertonformalstacks}, which allows us to restrict to the case of
finitely generated $A$-modules~$M$.% )\mar{TG: having done this, remove
  % the corresponding reference in the body of the paper.}



\begin{adefn}
	\label{df:obstruction}
%We say that a category fibred in groupoids $\cX$ over
% Let $\cX$ be a category fibred in groupoids over 
% a locally Noetherian scheme $S$,
% which satisfies~(RS*).
%\mar{ME: This could be tidied up/elaborated on, and maybe will
%	need to be added to as the argument develops}
We say that $\cX$
admits a {\em nice obstruction theory} if, \index{nice obstruction theory}
%\mar{ME: Maybe not the best terminology? TG: it's fine with me\dots ME: OK}
for each finite type $S$-algebra $A$
which lies over an affine open subscheme of $S$,
equipped with a morphism $x: \Spec A \to \cX$, there 
exists a complex of $A$-modules $K^{\bullet}_{(x,A)}$
such that the following conditions are satisfied:

\begin{enumerate}
	\item The complex $K^{\bullet}_{(x,A)}$ is bounded above and has 
		finitely generated cohomology modules (or, equivalently,
		is isomorphic in the derived category to a bounded above complex
		of finitely generated $A$-modules).
	\item  The formation of $K^{\bullet}_{(x,A)}$ is compatible (in
		the derived category) with pull-back. More precisely,
if $f:\Spec B \to \Spec A$, inducing the morphism $y: \Spec B
\to \Spec A \to \cX$,
then there is a natural isomorphism in the derived category
$f^*K^{\bullet}_{(x,A)} \iso K^{\bullet}_{(y,B)}$.
	\item
For any pair $(x,A)$ as above,
and for any finitely generated $A$-module $M$,
we have an isomorphism $T_x(M) \isoto H^1(K^{\bullet}_{(x,A)} \Lotimes M)$
%short exact sequence % \mar{TG: should say what natural
% means - perhaps just functorial and compatible with base change, as above?
% ME: OK, made this clearer (hopefully).}
% $$0 \to C_{(x,A)}\otimes M \to T_x(M) \to H^1(K^{\bullet}_{(x,A)} \Lotimes M)
% \to 0$$
whose formation is functorial in $M$ and in the pair $(x,A)$.
% in the sense that if we have a homomorphism $A \to B$, and let
% $y$ denote the composite
% $$\Spec B \to \Spec A \buildrel x \over \longrightarrow \cX,$$
% and if $N$ is a $B$-module (regarded also as an $A$-module via
% the given morphism $A\to B$),
% then the diagram
% $$\xymatrix{
% 0 \ar[r] & C_{(x,A)}\otimes N \ar[r]\ar[d] & T_x(N) \ar[r] \ar[d] & 
% H^1(K^{\bullet}_{(x,A)} \Lotimes N ) \ar[r]\ar[d] & 0 \\
% 0 \ar[r] & C_{(y,B)}\otimes N \ar[r] & T_y(N) \ar[r]  & 
% H^1(K^{\bullet}_{(y,B)} \Lotimes N ) \ar[r] & 0 
% }$$
% commutes; 
% here the left-most and right-most vertical arrows are induced
% by the pull-back isomorphisms of condition~(2) above,
% while the middle vertical arrow is the natural map arising
% from the functoriality of $T$,
% as discussed
% in~\cite[\href{http://stacks.math.columbia.edu/tag/07YA}{Tag 07YA}]{stacks-project}.


%\mar{ME: Actually, we already need $(RS*)$ for the functor $T_x(\text{--})$ to 
%	be defined; 
%                  see~\cite[\href{http://stacks.math.columbia.edu/tag/07Y9}{Tag 07Y9}]{stacks-project}.}
%surjection
%$T_x(M) \to H^1(K^{\bullet}_{(x,A)} \Lotimes M)$,
%whose kernel is a locally free $A$-module of finite rank.
\item
	For any pair $(x,A)$ as above,
	and for any finitely generated $A$-module $M$,
	the cohomology module
$H^2(K^{\bullet}_{(x,A)}\Lotimes M)$ serves as an obstruction
module. In particular, for any square zero extension
	$$0 \to I \to A' \to A \to 0$$
	for which $A'$ is of finite type over $S$ (or, equivalently,
	for which $I$ is a finite $A$-module), 
we have a functorial obstruction element $o_x(A')
\in H^2(K^{\bullet}_{(x,A)} \Lotimes I)$, which vanishes if and only
if~$x$ can be lifted to~$\Spec A'$.
\end{enumerate}
\end{adefn}
% We need the following slightly technical definition.
% \begin{adefn}\label{adefn: locally countably indexed}
%   We say that~$\cX$ is \emph{locally countably indexed} if there
% is a representable morphism $U\to\cX$ which is smooth and surjective,
% and whose source is a union of affine formal algebraic spaces which
% are countably indexed.% , in the sense that they can be written as an
% % inductive limit of a sequence of thickenings of affine
% % schemes. \mar{TG: maybe now defined above?}
% \end{adefn}
The following result generalises the familiar fact that a
pro-representing object for a formal deformation ring whose tangent
space is finite dimensional is necessarily a complete Noetherian local
ring. % (Note that since in the statement of the following theorem we assume that~$\cX$ is locally Ind-finite
      %   presentation over a locally Noetherian scheme
      %   $S$, it is in particular limit preserving, as assumed in the
      %   definitions above.)
        % \mar{TG: now needs some finiteness hypothesis? TG: we're OK if
  % these hypotheses guarantee limit preserving? And if not we should
  % add something to have that. TG: could just add that $\cX$ is limit
  % preserving, as we only apply this to $\cX_d$ which is limit
  % preserving by 3.2.15}
\begin{athm}
	\label{thm:noetherian criterion}\emph{(}\cite[Thm.\ 11.13]{Emertonformalstacks}\emph{)}
	If $\cX$ is a %limit preserving
locally countably indexed and Ind-locally of finite
        type formal algebraic stack over a locally Noetherian scheme
        $S$, and if~$\cX$ admits a nice obstruction theory, then $\cX$ is
        locally Noetherian.
\end{athm}

(When comparing this result with the statement of 
	\cite[Thm.\ 11.13]{Emertonformalstacks},
the reader should
bear in mind that an Ind-locally of finite type formal algebraic stack 
over~$S$ is in particular the $2$-colimit of algebraic stacks that are locally of 
finite type (or equivalently, locally of finite presentation) over $S$,
and hence is limit preserving.  Also, 
	\cite[Lem.\ 8.38]{Emertonformalstacks} ensures that an Ind-locally of
finite type formal algebraic stack over $S$ is also locally of Ind-locally finite type ---
which is the condition that appears in~\cite[Thm.\ 11.13]{Emertonformalstacks}.)


\chapter{Graded modules and rigid analysis}
\label{sec: rigid analytic perspective}
In this appendix
we establish some (mostly straightforward)
results that combine graded ring techniques 
with various results related to completions that are of a rigid analytic flavour.
%\mar{TG: maybe you could add a sentence or two to say
%what we do here?}
\section{Associated graded algebras and modules}\label{subsec:
  associated graded}% \mar{TG: some of this material is now in
  % \cite{EGstacktheoreticimages}, possibly even all of it.}
Let $R$ be an Artinian local ring, with maximal ideal $I$ and residue field~$k$.

\begin{adefn} If $M$ is an $R$-module, then we let $\Gr^{\bullet} M$
	denote the graded $k$-vector space associated to the $I$-adic filtration
	on $M$ (so $\Gr^i M := I^iM/I^{i+1} M$).  
\end{adefn}

The formation of $\Gr^{\bullet} M$ is functorial in $M$.
If $M$ and $N$ are two $R$-modules, then there is a natural surjection
(of graded $k$-vector spaces)
\anumequation
\label{eqn:graded tensor}
\Gr^{\bullet} M \otimes_k \Gr^{\bullet} N \to \Gr^{\bullet}(M\otimes_R N).
\end{equation}
In particular, it follows that if $A$ is an $R$-algebra, then $\Gr^{\bullet} A$
is naturally a graded $k$-algebra, and if $M$ is an $A$-module,
then $\Gr^{\bullet} M$ is naturally a $\Gr^{\bullet} A$-module.

We recall the following basic lemmas.

\begin{alemma}
	\label{lem:graded Noetherian}
	If $A$ is an $R$-algebra, then $A$ is Noetherian
	if and only if $\Gr^{\bullet} A$ is Noetherian.
\end{alemma}
\begin{proof}
	This is easy and standard; see e.g.~\cite[Prop.~1.1]{MR1990669} 
	for a statement of the ``if'' direction of this result
	(which is the less obvious of the two directions)
	in a significantly more general setting.
\end{proof}

\begin{alemma}
	\label{lem:graded surjective}
	If $A$ is an $R$-algebra and if $M$ is an $A$-module,
	then the natural morphism $\Gr^{\bullet} A \otimes_{\Gr^0 A}
	\Gr^0 M \to \Gr^{\bullet} M$ is surjective.
	Furthermore, $M$ is non-zero if and only if $\Gr^0 M$ is non-zero,
	if and only if $\Gr^{\bullet} M$ is non-zero.
\end{alemma}
\begin{proof}
	The first claim is immediate.
	This implies in turn that if $\Gr^0 M$ is zero then 
	$\Gr^{\bullet} M$ is zero (the converse being evident), 
	and thus that $M$ is zero (since $\Gr^{\bullet} M$ is
	the associated graded of $M$ with respect to a finite length
	filtration beginning at $M$ and ending at $0$).
	(Essentially equivalently, one can observe directly that since $I$
	is nilpotent, if $M/IM = 0$ then $M = 0$.)
\end{proof}

\begin{alemma}
	\label{lem:graded fg}
	If $A$ is an $R$-algebra and $M$ is an $A$-module,
	then $M$ is finitely generated over $A$ if and only 
	if $\Gr^0 M$ is finitely generated over $\Gr^0 A$,
	if and only if $\Gr^{\bullet} M$ is finitely generated
	over $\Gr^{\bullet} A$.
\end{alemma}
\begin{proof} 
Clearly if $M$ is finitely generated over $A$, then
$\Gr^0 M := M/IM$ is finitely generated over $A/I$; the converse
assertion follows from Nakayama's lemma with respect to
the nilpotent ideal $I$.  If $\Gr^0 M$ is finitely generated 
over $\Gr^0 A$, then Lemma~\ref{lem:graded surjective}
shows that $\Gr^{\bullet} M$ is finitely generated over $\Gr^{\bullet} A$;
the converse assertion follows from the fact that (again by
Lemma~\ref{lem:graded surjective}) $\Gr^0 M$ may be obtained as the
quotient of $\Gr^{\bullet} M$ by the ideal $\Gr^{\bullet > 0} A$
of $\Gr^{\bullet} A$.
\end{proof}
\begin{alemma}
  \label{lem: graded iso is iso}Let $f:M\to N$ be a morphism of $A$-modules. Then~$f$ is an isomorphism if and only if the
  induced morphism $\Gr^{\bullet}f:\Gr^{\bullet}M\to\Gr^{\bullet}N$ is an
  isomorphism.%\mar{TG: feel free to improve/correct this!}
\end{alemma}
\begin{proof}
  If~$f$ is an isomorphism then certainly~$\Gr^{\bullet}f$ is an
  isomorphism. Conversely, if~$\Gr^{\bullet}f$ is an isomorphism, then
  in particular $M/IM\to N/IN$ is surjective, so~$f$ is surjective by
  Nakayama's lemma with respect to
the nilpotent ideal $I$. In addition, for each~$i$ the
  morphism~$I^iM/I^{i+1}M\to I^iN/I^{i+1}N$ is injective, so by an
  easy induction on~$i$ we see that any element of~$\ker f$ is
  contained in~$I^iM$ for all~$i$, and is thus zero, as required.
\end{proof}
\begin{alemma}
	\label{lem:graded flat}
If $A$ is an $R$-algebra, then an $A$-module $M$
is {\em (}faithfully{\em )} flat over $A$
if and only if $\Gr^{\bullet} M$ 
is {\em (}faithfully{\em )} flat over $\Gr^{\bullet} A.$
Furthermore, if any of these conditions holds,
then the natural morphism
$\Gr^{\bullet} A \otimes_{\Gr^0 A} \Gr^0 M \to
\Gr^{\bullet} M$
is an isomorphism.
\end{alemma}
\begin{proof}
This is~\cite[Lem.\ 5.5.37]{EGstacktheoreticimages}.  
\end{proof}

%\begin{proof}
%	If $M$ is flat over $A$, then considering the
%        result of tensoring $M$
%	by the various short exact sequences
%	$$0 \to I^nA \to I^mA \to I^mA/I^nA \to 0$$
%	we find that the natural
%	morphism $\Gr^i A \otimes_{\Gr^0 A} \Gr^0 M \to \Gr^i M$
%	is an isomorphism, for each~$i$.
%	This proves the final assertion of the lemma.
%	Furthermore, since (faithful) flatness is preserved
%	under base-change, we see first that $\Gr^0 M$ is flat
%	over $\Gr^0 A$ (and faithfully flat if $M$ is faithfully
%	flat over $A$), and then (using the result already proved)
%	that $\Gr^{\bullet} M$ is flat over $\Gr^{\bullet} A$
%	(and faithfully flat if $M$ is).
%	
%	It is not quite as obvious that flatness of $\Gr^{\bullet} M$
%	over $\Gr^{\bullet} A$ implies the flatness of $M$ over $A$,
%	but this is a standard fact in commutative algebra;
%	e.g.~it follows
%	from~\cite[\href{http://stacks.math.columbia.edu/tag/0AS8}{Tag 0AS8}]{stacks-project}. 
%	(If $i \geq 0,$ then base-changing via the map
%	$\Gr^{\bullet} A \to
%	\Gr^{\leq i} A := A/IA \oplus I/I^2 \oplus \cdots \oplus I^i/I^{i+1},$
%	we find that $\Gr^{\leq i} M := M/IM \oplus \cdots \oplus I^iM/I^{i+1}M$
%	is flat over $\Gr^{\leq i} A$.  In particular, the
%	embedding $I^i/I^{i+1} A =: \Gr^i A \hookrightarrow \Gr^{\leq i} A$
%	induces an embedding
%	$$I^i/I^{i+1} \otimes_A M
%	\iso \Gr^i A \otimes_{\Gr^{\leq i} A}
%	\Gr^{\leq i} M;$$ 
%	concretely, this means that the morphism 
%	$$I^i/I^{i+1} \otimes_A M \to I^iM/I^{i+1}M$$
%	is an embedding.  Letting $i$ vary, and recalling that $I$
%	is nilpotent, we deduce 
%	from~\cite[\href{http://stacks.math.columbia.edu/tag/0AS8}{Tag
%          0AS8}]{stacks-project} that $M$ is flat over $A$.)
%
%	If $\Gr^{\bullet} M$ is furthermore faithfully flat over
%	$\Gr^{\bullet} A$, then (since faithful flatness is preserved
%	under base-change) we see that $\Gr^0 M := M/IM$ is faithfully flat
%	over $\Gr^0 A := A/IA$.  Since an $A$-module vanishes
%	if and only if its reduction mod $I$ does (see Lemma~\ref{lem:graded
%		surjective}),
%	we conclude that $M$ is faithfully flat over $A$.
%\end{proof}

\begin{alemma}
	\label{lem:graded tensor}
	If $A$ is an $R$-algebra, and if $M$ and $N$ are $A$-modules,
	then the natural surjection~{\em (\ref{eqn:graded tensor})}
	induces a natural surjection 
	$$\Gr^{\bullet} M \otimes_{\Gr^{\bullet} A} \Gr^{\bullet} N \to
	\Gr^{\bullet} (M\otimes_A N).$$
\end{alemma}
\begin{proof}
	The natural surjection $M\otimes_R N \to M\otimes_A N$ induces
	a surjection $\Gr^{\bullet}(M\otimes_R N) \to \Gr^{\bullet} (M\otimes_A N).$  It is immediate that its composite with the surjection~(\ref{eqn:graded
		tensor}) factors through $\Gr^{\bullet} M \otimes_{\Gr^{\bullet}
		A} \Gr^{\bullet} N.$
\end{proof}

\begin{alemma}
	\label{lem:graded flat bis}
	If $A$ is an $R$-algebra, and if $M$ and $N$ are $A$-modules 
	with either $M$ or $N$ flat over $A$,
	then the surjection of Lemma~{\em \ref{lem:graded tensor}}
       	is a natural isomorphism
	$$\Gr^{\bullet} M \otimes_{\Gr^{\bullet} A} \Gr^{\bullet} N \iso
	\Gr^{\bullet} (M\otimes_A N).$$
\end{alemma}
\begin{proof}
	Suppose (without loss of generality) that $N$ is $A$-flat.
	Taking into account Lemma~\ref{lem:graded flat},
	we see that we have to show that natural morphism
	$\Gr^{\bullet} M \otimes_{\Gr^0 A} \Gr^0 N \to \Gr^{\bullet}(M\otimes_A N)$ is an isomorphism, i.e.\ that for each~$i \geq 0,$ the natural morphism
	$$I^iM /I^{i+1} M \otimes_{A/I} N/I \to I^i(M\otimes N)/I^{i+1}(M\otimes N)$$
	is an isomorphism.  This follows easily by tensoring $N$ over $A$
	with the various short exact sequences
$$0 \to I^n M \to I^m M \to I^mM/I^nM\to 0,$$
taking into account the flatness of $N$ over $A$.
\end{proof}


\section{A general setting}
\label{subsec:setting}
We fix an Artinian local ring $R$ with maximal ideal $I$
and residue field $k$, % \mar{ME: I'm choosing $I$ rather than $\mathfrak m$ so that later, $\mathfrak m$ is available for 
% 	the maximal ideal of rings like $\cO_{\C_p}$ which might appear.
% I don't know if this is sensible or not, though. TG: I think this is
% fine and is what we're doing, commenting out.}
as well an $R$-algebra $C^+$, and an element $u \in C^+$, satisfying the
following properties:
\begin{enumerate}[label=(\Alph*)]
	\item $u$ is a regular element (i.e.\ a non-zero divisor) of $C^+$.
	\item $C^+/u$ is a flat $R$-algebra.
	\item $C^+/I$ is a rank one complete valuation ring, and
		the image of $u$ in $C^+/I$ (which
		is necessarily non-zero, by (A)) is of
		positive valuation (i.e.\ lies in the maximal
		ideal of $C^+/I$).
\end{enumerate}

\begin{aremark}
	\label{rem:independence of u}
	In the preceding context, if $v \in C^+$ has non-zero
	image in $C^+/I$, then
	$v^n$ divides $u^m$ in $C^+$, for some $m, n > 0$.
	(Indeed, since $C^+/I$ is a rank one valuation ring
	in which the image of $u$ has positive valuation,
	we may write $u^a = v w + x,$ for some $a>0$, some $w \in C^+$,
	and some $x \in I C^+$.  Raising both
	sides of this equation to a sufficiently large power,
	remembering that $I$ is nilpotent,
	gives the claim.)
	
	In particular, if $v \in C^+$ is another element
	satisfying conditions~(A), (B), and (C) above, then
	the powers of $u$ and of $v$ are mutually cofinal with respect
	to divisibility (i.e.\ the $u$-adic and $v$-adic topologies
        on $C^+$ coincide), and consequently $C^+[1/u] = C^+[1/v]$.
\end{aremark}

\begin{aremark}
	\label{rem:flat implies free}
	Recall  that since $R$ is an Artinian local ring,
	an $R$-module is flat if and only if it is free~\cite[\href{https://stacks.math.columbia.edu/tag/051G}{Tag 051G}]{stacks-project}.
	(We will use this fact in the proof of Lemma~\ref{lem:vanishing
		locus} below.  For a proof, note first that if $M$ 
	is any $R$-module, then any basis of $M/IM$ lifts to
	a generating set of $M$.  If $M$ is furthermore flat,
	then lifting a basis of $M/IM$, we obtain a surjection
	$F\to M$ whose source is free, with the additional
        property that, if $N$ denotes its kernel, then $N/IN = 0$.
	Thus $N = 0$, and so $F\iso M$, as required.)

\end{aremark}

We begin by establishing some simple consequences of our assumptions
on $C^+$ and~$u$.

\begin{alemma}
	\label{lem:basic properties}\leavevmode
        \begin{enumerate}[label=\normalfont(\arabic*)]
        \item  For each $n \geq 1,$ the quotient $C^+/u^n$ is flat
	over $R$.
      \item  $C^+$ itself is flat over $R$.
      \item  The ring $C^+$ is $u$-adically complete.
        \end{enumerate}
	% {\em (1)} For each $n \geq 1,$ the quotient $C^+/u^n$ is flat
	% over $R$.

	% {\em (2)} $C^+$ itself is flat over $R$.

	% {\em (3)} The ring $C^+$ is $u$-adically complete.
\end{alemma}
\begin{proof}
Claim~(1) follows by an evident induction from assumptions~(A) and~(B)
together with a consideration of the short exact sequences
%\anumequation
%\label{eqn:u truncation}
$$
0 \to C^+/u^n \buildrel u\cdot \over \longrightarrow C^+/u^{n+1}
\to C^+/u \to 0.
$$
%\end{equation}

Now choose an increasing
filtration $(I_m)$ of $R$ by ideals so that $I_0 = 0$, and such that
each quotient $I_{m+1}/I_{m}$ is of length one (i.e.\ is isomorphic 
to $k$).
Taking into account~(1), the short exact sequences
$$0 \to I_{m} \to I_{m+1} \to I_{m+1}/I_m (\cong k = R/I) \to 0$$
give rise to short exact sequences
$$0 \to I_m\otimes_R (C^+/u^n) \to I_{m+1}\otimes_R (C^+/u^n) \to C^+/(I,u^n) \to 0.$$
Suppose that for some~$m$ we know that the natural morphism
$$I_m \otimes_R C^+ \to \varprojlim_{n} I_m \otimes_R (C^+/u^n)$$
is an isomorphism.
Then passing to the inverse limit over $n$, and taking into account 
both assumption~(C) and the fact that
the transition maps $I_m\otimes_R (C^+/u^{n+1}) \to I_m\otimes_R (C^+/u^n)$
are surjective, so that the relevant $\varprojlim^1$ vanishes,
we obtain a short exact sequence
$$0 \to I_m\otimes_R C^+ \to \varprojlim_n I_{m+1}\otimes_R (C^+/u^n)
\to C^+/I \to 0.$$
The five lemma shows that the natural morphism to this sequence
from the exact sequence
$$I_m\otimes_R C^+ \to I_{m+1}\otimes_R C^+ \to C^+/I \to 0$$
is then an isomorphism.
Proceeding by induction (the case~$m=0$ being trivial), we find (once we reach the
top of our filtration, so that $I_m = R$) that $C^+$ is $u$-adically
complete, and we also find that each of the morphisms
$$I_m\otimes_R C^+ \to C^+$$ 
is injective.  Since any ideal $J$ of $R$ can be placed in such a
filtration $(I_m)$, we find that $C^+$ is flat over $R$.  Thus~(2)
and~(3) are proved.
\end{proof}

We write $C := C^+[1/u]$. Since $C$ is a localisation of $C^+$,
which is $R$-flat by Lemma~\ref{lem:basic properties}~(2),
it is flat over $R$.
Note also that $\Gr^0C \cong (C^+/I)[1/u]$ is a field that is
complete with respect to a non-archimedean absolute value.
In particular, we can do rigid geometry over $\Gr^0 C$;
this is a key point, which we will exploit below.

\begin{adefn}
If $M$ is an $R$-module,
then we write $M\cotimes_R C^+$ to denote the $u$-adic completion
of the tensor product $M\otimes_R C^+$.  We also write
$M\cotimes_R C := (M\cotimes_R C^+)[1/u].$
\end{adefn}

\begin{aremark}
	Remark~\ref{rem:independence of u} shows
	that the formation of $M\cotimes_R C^+$ and $M\cotimes_R C$
	is independent of the choice of $u$ satisfying~(A), (B), and~(C)
	above.
\end{aremark}

\begin{aremark}
	In applications, we will apply this construction primarily %,
	%although not exclusively,
	in the case when $M = A$ is an $R$-algebra,
	in which case $A\cotimes_R C^+$ and $A\cotimes_R C$ are again
	$R$-algebras.
\end{aremark}

\begin{aexample}
	If we take $C^+ := R[[u]]$ (for an indeterminate $u$),
	then $A\cotimes_R C^+ = A[[u]],$ and $A\cotimes_R C = A((u))$.
	% Thus the set-up of this section may be regarded
	% as a generalisation of that of Section~\ref{sec: modules over power Laurent series rings}.
\end{aexample}

We next state and prove some additional
properties of $u$-adic completions that we will need.

\begin{alemma}
	\label{lem:non-zero divisor}
	If $M$ is an $R$-module, then multiplication by $u$ 
	is injective on $M\cotimes_A C^+$.
	Consequently {\em (}indeed, equivalently{\em )}, the
	natural map $M\cotimes_R C^+ \to M\cotimes_R C$ is
	injective.

	Furthermore, for each $m \geq 0$, the natural map
	$M\cotimes_R C^+ \to M\otimes_R (C^+/u^m)$
	is surjective,
	and induces an isomorphism
	$$(M\cotimes_R C^+)/u^m (M\cotimes_R C^+) \iso M\otimes_R (C^+/u^m).$$
\end{alemma}
\begin{proof}
	Note that the claim of the second paragraph
	is a general property of completion with respect to finitely 
	generated ideals
(see
e.g.~the statement and proof of
\cite[\href{http://stacks.math.columbia.edu/tag/05GG}{Tag
  05GG}]{stacks-project}). 
We will also deduce it as a byproduct of the proof of the claims
in the first paragraph, to which we now turn.

	Given the definition of $M\cotimes_R C$ as the localization
	$M\cotimes_R C^+[1/u],$ the second claim of the first
	paragraph is clearly
	equivalent to the first.
	As for this first claim,
	we note that for each $m, n \geq 1,$
	Lemma~\ref{lem:basic properties}~(1)
	shows that the terms in the short exact sequence
$$
0 \to C^+/u^n \buildrel u^m \cdot \over \longrightarrow C^+/u^{n+m}
\to C^+/u^m \to 0
$$
	are flat $R$-modules. 
	Thus, tensoring this short exact sequence
	with $M$ over $R$ yields a short exact sequence
	$$0 \to M\otimes_R (C^+/u^n) \to M\otimes_R (C^+/u^{n+1}) \to
	M\otimes_R (C^+/u)\to 0.$$
	Passing to the inverse limit over $n$,
	we obtain the exact sequence
	$$0 \to M\cotimes_R C^+ \buildrel u^m \cdot \over \longrightarrow
	M\cotimes_R C^+ \to M\otimes_R C^+/u^m C^+ $$
	from which the claims of the first paragraph follow.
	If we use the surjectivity of the transition maps in the
	inverse limit to infer
	that the relevant $\varinjlim^1$ vanishes, then we see that
	this sequence is even exact on the right, so that we also 
	obtain a confirmation of the claim of the second paragraph.
\end{proof}

%We next state and prove some additional
%properties of $u$-adic completions that we will need.

\begin{alemma}
	\label{lem:completion properties}\leavevmode
        \begin{enumerate}[label=\normalfont(\arabic*)]
        \item The functors $M \mapsto M\cotimes_R C^+$ and
          $M\mapsto M\cotimes_R C$ are exact. %\mar{TG: any reason that
           % the second tensor product had no subscript? I think it
           % should, so I added them.}
        \item  If $A$ is an $R$-algebra, if $J$ is a finitely
          generated ideal in $A$, and if $M$ is an $A$-module, then
          the natural map
	$$J (M\cotimes_R C^+) \to JM \cotimes_R C^+$$
        is an isomorphism, and consequently there is a short exact
        sequence
	$$0 \to J(M\cotimes_R C^+) \to M\cotimes_R C^+
	\to (M/JM)\cotimes_R C^+ \to 0.$$
      \item If $J_1$ and $J_2$ are ideals in an $R$-algebra $A$,
        with $J_1 \supseteq J_2$, and if $M$ is an $A$-module, then
        there is a natural isomorphism
$$J_1 (M\cotimes_R C^+)/J_2 (M\cotimes_R C^+) \iso (J_1 M/J_2 M) 
\cotimes_R C^+.$$
\end{enumerate}
\end{alemma}
\begin{proof}
	If $0 \to M_1 \to M_2 \to M_3 \to 0$ is a short exact sequence
	of $R$-modules, then Lemma~\ref{lem:basic properties}~(1) 
	implies that 
	$$0 \to M_1\otimes_R (C^+/u^n) \to M_2\otimes_R (C^+/u^n) \to 
	M_3\otimes_R (C^+/u^n) \to  0$$
	is exact for each $n$.  If we pass to the inverse limit over~$n$,
	and note that the transition morphisms are evidently surjective,
	so that the relevant $\varprojlim^1$ vanishes,
	we obtain a short exact sequence
	$$0 \to M_1\cotimes_R C^+ \to M_2\otimes_R C^+ \to 
	M_3\otimes_R C^+ \to  0,$$
	proving the first exactness claim of~(1).
	The second exactness claim follows from the first,
	together with the fact that the localization
	$C := C^+[1/u]$ is flat over $C$, and that there is an
	isomorphism
	$M \cotimes_R C \iso (M\cotimes_R C^+)\otimes_{C^+}C.$

%	We may write $C = \varinjlim_n u^{-n} C^+$ (the transition
%	maps being the evident inclusions),
%	and thus (since tensor products are compatible with formation
%	of direct limits) obtain the sequence
%	$$0 \to M_1\cotimes_R C \to M_2\otimes_R C \to 
%	M_3\otimes_R C \to  0$$
%	by passing to the direct limit over $n$ of the sequences
%	$$0 \to M_1\cotimes_R u^{-n}C^+ \to M_2\otimes_R u^{-n}C^+ \to 
%	M_3\otimes_R u^{-n}C^+ \to  0,$$
%	which have already been shown to be short exact.
%	Since the formation of direct limits is exact, we obtain
%	the second exactness claim of~(1).

%	In order to prove~(2), we will show that
%	if $J$ is an ideal in $R$, 
%	then there is an isomorphism
%	\numequation
%	\label{eqn:complete iso}
%	J (M\cotimes_R C^+) \iso (J M)\cotimes_R C^+
%\end{equation}
%	(i.e.\ we will first prove~(2) in the case when $J_2 = 0$). 
%	The general case of~(2) then follows from this special
%	case and the first exactness result of~(1).

	In order to prove~(2), 
	we apply~(1) to the short exact sequence $0 \to JM \to M
	\to M/JM \to 0$, obtaining a short exact sequence
	\anumequation
	\label{eqn:J ses}
	0 \to (J M)\cotimes_R C^+ \to M\cotimes_R C^+ 
	\to (M/JM)\cotimes_R C^+ \to 0.
\end{equation}
        Thus we may (and do) 
	regard $(JM)\cotimes_R C^+$ as a submodule of $M\cotimes_R C^+.$
        There is a consequent inclusion
	\anumequation
	\label{eqn:J inclusion}
	J (M\cotimes_R C^+)
	\subseteq (JM)\cotimes_R C^+,
        \end{equation}
	whose image is $u$-adically
	dense in the target.
	Since $J$ is a finitely generated ideal,
	%(as $R$ is Artinian, and thus Noetherian),
	we see that $J(M\cotimes_R C^+)$ is furthermore $u$-adically complete,
	and thus that~(\ref{eqn:J inclusion}) is actually an equality,
	establishing the first claim of~(2).

	%Consequently, we find that the natural morphism
	%$(JM)\cotimes_R C^+  \to 
	%M\cotimes_R C^+$ is a closed embedding (when the source
	%is equipped with its $u$-adic topology), and so we may (and do)
	%regard $(JM)\cotimes_R C^+$ as a $u$-adically closed
	%submodule of $M\cotimes_R C^+.$
%
%        Clearly there is an inclusion
%	\anumequation
%	\label{eqn:J inclusion}
%	J (M\cotimes_R C^+)
%	\subseteq (JM)\cotimes_R C^+,
%        \end{equation}
%	whose image is $u$-adically
%	dense in the target.
%	Since $J$ is a finitely generated ideal,
%	%(as $R$ is Artinian, and thus Noetherian),
%	we see that $J(M\cotimes_R C^+)$ is furthermore $u$-adically complete,
%	and thus that~(\ref{eqn:J inclusion}) is actually an equality,
%	establishing the first claim of~(2).
	%This establishes the required isomorphism~(\ref{eqn:complete iso}).
	Replacing $(JM)\cotimes_R C^+$ by $J(M\cotimes_R C^+)$ in
	the short exact sequence~(\ref{eqn:J ses}) yields the short
	exact sequence whose existence is asserted in the second claim
	of~(2).

	Claim~(3) follows directly from~(2), and the exactness result
	of~(1).
\end{proof}

%We are interested in studying finitely generated projective
%modules over $A\cotimes_R C$.  In the present general context, we don't
%know whether an analogue of Drinfeld's descent result for Tate modules
%(see~\cite[Thm.\ 3.3, 3.11]{MR2181808} and~\cite[Thm.\ 5.1.18]{EGstacktheoreticimages}) holds
%for arbitrary $R$-algebras $A$.  However, we will see that such
%a descent result is valid in the more restricted setting of finite
%type $R$-algebras.  In the case of finite type $k$-algebras,
%we will prove this using results from rigid analysis.  For finite type
%$R$-algebras that are not necessarily annihilated by $I$,
%we will use graded techniques to reduce to the case of $k$-algebras.
%We describe this reduction step first.

%\section{Associated graded algebras and modules}
%We maintain the notation ($R$, $C^+$, etc.) and assumptions from the
%preceding subsection.  

%\section{Associated graded algebras and modules}
%
%\begin{adefn} If $M$ is an $R$-module, then we let $\Gr^{\bullet} M$
%	denote the graded $k$-vector space associated to the $I$-adic filtration
%	on $M$ (so $\Gr^i M := I^iM/I^{i+1} M$).  
%\end{adefn}
%
%The formation of $\Gr^{\bullet} M$ is functorial in $M$.
%If $M$ and $N$ are two $R$-modules, then there is a natural surjection
%(of graded $k$-vector spaces)
%\anumequation
%\label{eqn:graded tensor}
%\Gr^{\bullet} M \otimes_k \Gr^{\bullet} N \to \Gr^{\bullet}(M\otimes_R N).
%\end{equation}
%In particular, it follows that if $A$ is an $R$-algebra, then $\Gr^{\bullet} A$
%is naturally a graded $k$-algebra, and if $M$ is an $A$-module,
%then $\Gr^{\bullet} M$ is naturally a $\Gr^{\bullet} A$-module.
%
%We recall the following basic lemmas.
%
%\begin{lemma}
%	\label{lem:graded surjective}
%	If $A$ is an $R$-algebra and if $M$ is an $A$-module,
%	then the natural morphism $\Gr^{\bullet} A \otimes_{\Gr^0 A}
%	\Gr^0 M \to \Gr^{\bullet} M$ is surjective.
%	Furthermore, $M$ is non-zero if and only if $\Gr^0 M$ is non-zero,
%	if and only if $\Gr^{\bullet} M$ is non-zero.
%\end{lemma}
%\begin{proof}
	%The first claim is immediate.
	%This implies in turn that if $\Gr^0 M$ is zero then 
	%$\Gr^{\bullet} M$ is zero (the converse being evident), 
	%and thus that $M$ is zero (since $\Gr^{\bullet} M$ is
	%the associated graded of $M$ with respect to a finite length
	%filtration beginning at $M$ and ending at $0$).
	%(Essentially equivalently, one can observe directly that since $I$
%	is nilpotent, if $M/IM = 0$ then $M = 0$.)
%\end{proof}
%
%\begin{lemma}
%	\label{lem:graded fg}
%	If $A$ is an $R$-algebra and $M$ is an $A$-module,
%	then $M$ is finitely generated over $A$ if and only 
%	if $\Gr^0 M$ is finitely generated over $\Gr^0 A$,
%	if and only if $\Gr^{\bullet} M$ is finitely generated
%	over $\Gr^{\bullet} A$.
%\end{lemma}
%\begin{proof} 
%Clearly if $M$ is finitely generated over $A$, then
%$\Gr^0 M := M/IM$ is finitely generated over $A/I$; the converse
%assertion follows from Nakayama's lemma with respect to
%the nilpotent ideal $I$.  If $\Gr^0 M$ is finitely generated 
%over $\Gr^0 A$, then Lemma~\ref{lem:graded surjective}
%shows that $\Gr^{\bullet} M$ is finitely generated over $\Gr^{\bullet} A$;
%the converse assertion follows from the fact that (again by
%Lemma~\ref{lem:graded surjective}) $\Gr^0 M$ may be obtained as the
%quotient of $\Gr^{\bullet} M$ by the ideal $\Gr^{\bullet > 0} A$
%of $\Gr^{\bullet} A$.
%\end{proof}
%
%\begin{lemma}
%	\label{lem:graded flat}
%If $A$ is an $R$-algebra, then an $A$-module $M$
%is {\em (}faithfully{\em )} flat over $A$
%if and only if $\Gr^{\bullet} M$ 
%is {\em (}faithfully{\em )} flat over $\Gr^{\bullet} A.$
%Furthermore, if any of these conditions holds,
%then the natural morphism
%$\Gr^{\bullet} A \otimes_{\Gr^0 A} \Gr^0 M \to
%\Gr^{\bullet} M$
%is an isomorphism.
%\end{lemma}
%\begin{proof}
%	If $M$ is flat over $A$, then considering the
%        result of tensoring $M$
%	by the various short exact sequences
%	$$0 \to I^nA \to I^mA \to I^mA/I^nA \to 0$$
%	we find that the natural
%	morphism $\Gr^i A \otimes_{\Gr^0 A} \Gr^0 M \to \Gr^i M$
%	is an isomorphism, for each $i$.
%	This proves the final assertion of the lemma.
%	Furthermore, since (faithful) flatness is preserved
%	under base-change, we see first that $\Gr^0 M$ is flat
%	over $\Gr^0 A$ (and faithfully flat if $M$ is faithfully
%	flat over $A$), and then (using the result already proved)
%	that $\Gr^{\bullet} M$ is flat over $\Gr^{\bullet} A$
%	(and faithfully flat if $M$ is).
%	
%	It is not quite as obvious that flatness of $\Gr^{\bullet} M$
%	over $\Gr^{\bullet} A$ implies the flatness of $M$ over $A$,
%	but this is a special case of.
%	\mar{ME: Cite something}
%	If $\Gr^{\bullet} M$ is furthermore faithfully flat over
%	$\Gr^{\bullet} A$, then (since faithful flatness is preserved
%	under base-change) we see that $\Gr^0 M := M/IM$ is faithfully flat
%	over $\Gr^0 A := A/IA$.  Since an $A$-module vanishes
%	if and only if its reduction mod $I$ does (see Lemma~\ref{lem:graded
%		surjective}),
%	we conclude that $M$ is faithfully flat over $A$.
%\end{proof}
%
%\begin{lemma}
%	\label{lem:graded tensor}
%	If $A$ is an $R$-algebra, and if $M$ and $N$ are $A$-modules,
%	then the natural surjection~{\em (\ref{eqn:graded tensor})}
%	induces a natural surjection 
%	$$\Gr^{\bullet} M \otimes_{\Gr^{\bullet} A} \Gr^{\bullet} N \to
%	\Gr^{\bullet} (M\otimes_A N).$$
%\end{lemma}
%\begin{proof}
%	The natural surjection $M\otimes_R N \to M\otimes_A N$ induces
%	a surjection $\Gr^{\bullet}(M\otimes_R N) \to \Gr^{\bullet} (M\otimes_A N).$  It is immediate that its composite with the surjection~(\ref{eqn:graded
%		tensor}) factors through $\Gr^{\bullet} M \otimes_{\Gr^{\bullet}
%		A} \Gr^{\bullet} N.$
%\end{proof}
%
%\begin{lemma}
%	\label{lem:graded flat bis}
%	If $A$ is an $R$-algebra, and if $M$ and $N$ are $A$-modules 
%	with either $M$ or $N$ flat over $A$,
%	then the surjection of Lemma~{\em \ref{lem:graded tensor}}
%       	is a natural isomorphism
%	$$\Gr^{\bullet} M \otimes_{\Gr^{\bullet} A} \Gr^{\bullet} N \iso
%	\Gr^{\bullet} (M\otimes_A N).$$
%\end{lemma}
%\begin{proof}
%	Suppose (without loss of generality) that $N$ is $A$-flat.
%	Taking into account Lemma~\ref{lem:graded flat},
%	we see that we have to show that natural morphism
%	$\Gr^{\bullet} M \otimes_{\Gr^0 A} \Gr^0 N \to \Gr^{\bullet}(M\otimes_A N)$ is an isomorphism, i.e.\ that for each $i \geq 0,$ the natural morphism
%	$$I^iM /I^{i+1} M \otimes_{A/I} N/I \to I^i(M\otimes N)/I^{i+1}(M\otimes N)$$
%	is an isomorphism.  This follows easily by tensoring $N$ over $A$
%	with the various short exact sequences
%$$0 \to I^n M \to I^m M \to I^mM/I^nM\to 0,$$
%taking into account the flatness of $N$ over $A$.
%\end{proof}
 
%Now let $M$ be an $R$-module.  We wish to apply the preceding discussion
%to $M\cotimes_R C^+$ and $M\cotimes_R C$.

The following lemma will allow us to use grading techniques
to study $u$-adic completions.

\begin{alemma}
	\label{lem:passing to graded context}
	If $M$ is an $R$-module,
	then there are natural isomorphisms

	{\em (1)}
	$\Gr^{\bullet} M \otimes_{k} \Gr^0 (C^+/u^n) \iso
	\Gr^{\bullet} \bigl(M\otimes_R
	(C^+/u^n)\bigr)$
{\em (}for any $n  \geq 1${\em )},

{\em (2)}
$\Gr^{\bullet} (M\cotimes_R C^+) \iso
\Gr^{\bullet} M \cotimes_{k} \Gr^0 C^+ ,$
	and

	{\em (3)}
	$\Gr^{\bullet} (M\cotimes_R C)
	\iso
	\Gr^{\bullet} M \cotimes_{k} \Gr^0 C.$ 
\end{alemma}
\begin{proof}
	Lemma~\ref{lem:basic properties}~(1)
	shows that $C^+/u^n$ is flat over $R$,
	and~(1) follows from this,
	together with Lemmas~\ref{lem:graded flat} and~\ref{lem:graded flat bis}.

	Lemma~\ref{lem:completion
		properties}~(3) (taking the inclusion $J_2 \subseteq
	J_1$ to be the various inclusions $I^{i+1} \subseteq I^i$
	in turn) shows that each
        $\Gr^i(M\cotimes_R C^+) \iso \Gr^i M \cotimes_R  C^+.$
        Since $I$ annihilates $\Gr^i(M)$, the target of this
        isomorphism can be rewritten as $\Gr^i M \cotimes_k \Gr^0 \C^+$.
        This proves~(2).
        


%\mar{ME: I think this is inefficient and maybe not even quite on-point,
%so rewrite slightly}
%	Since $C^+/u^n$ is flat over $R$, we see that the natural map
%	$(\Gr^0 C^+)/u^n \to \Gr^0( C^+/ u^n)$ is an isomorphism,
%	and combining this with the natural isomorphism of~(1),
%	we obtain a natural isomorphism
%	$$\Gr^{\bullet} M \otimes_{k} (\Gr^0 C^+)/u^n 
%	\iso \Gr^{\bullet} \bigl(M\otimes_R (C^+/u^n)\bigr).$$
%	Taking this into account, we find that to prove~(2), it suffices 
%	to show that the natural morphism
%	$$
%	\Gr^{\bullet}( M\cotimes_R C^+) = 
%	\Gr^{\bullet}\bigl(\varprojlim_n M\otimes_R (C^+/u^n)\bigr)
%	\to \varprojlim_n \Gr^{\bullet} \bigl(M\otimes_R (C^+/u^n)\bigr)
%	$$
%	is an isomorphism.
%	This follows directly from Lemma~\ref{lem:completion
%		properties}~(3) (taking the inclusion $J_2 \subseteq
%	J_1$ to be the various inclusions $I^{i+1} \subseteq I^i$
%	in turn).

	%As already noted in the proof of Lemma~\ref{lem:completion
	%	properties},
	The $R$-algebra $C$ may be described as the direct limit
	$ C \iso \varinjlim_n u^{-n}C^+$ (the transition
	maps being the obvious inclusions).
	Since the formation of direct limits is compatible with
	tensor products, we find that
	$$\Gr^0 C \iso \varinjlim_n \Gr^0 u^{-n} C^+.$$
	The isomorphism of~(3) is thus obtained from
	the isomorphisms
	$$\Gr^{\bullet} ( M\cotimes_R u^{-n} C^+) \iso
	\Gr^{\bullet} M \cotimes_{k} \Gr^0 u^{-n} C^+$$
	of~(2) by passing to the
	direct limit over $n$.
\end{proof}

\section{Loci of vanishing and of equivariance}
We maintain the notation of Section~\ref{subsec:setting},
and prove some slightly technical results about the ``locus
of vanishing'' and ``locus of equivariance'' of morphisms 
between finitely generated projective $A\cotimes_R C$-modules.

\begin{alemma}
	\label{lem:vanishing locus}
	Let $A$ be an $R$-algebra,
	and let $M$ be a finitely generated projective module 
	over $A\cotimes_R C$.  If $m \in M$, then the ``locus
	of vanishing'' of $m$ is a Zariski closed subset of $\Spec A$.
	More precisely, there is an ideal $J$ of $A$
	such that for any morphism $\phi:A \to B$ of $R$-algebras,
	the image of $m$ in $M_B := (B\cotimes_R C) \otimes_{A\cotimes_R C}
	M$ vanishes if and only if $J$ is contained in the kernel
	of~$\phi$.
\end{alemma}
\begin{proof}
	Since $M$ is finitely generated projective over $A\cotimes_R C$,
	we may find another finitely generated projective $A\cotimes_R C$-module $N$
	such that $M \oplus N$ is free of finite rank over $A\cotimes_R C$.
	Replacing $m \in M$ by $m \oplus 0 \in M \oplus N$, we reduce
	to the case when $M$ is free of finite rank, 
	say $M = A\cotimes_R C^{\oplus r}.$  Writing 
	$m = (x_1,\ldots,x_r)$ with $x_i \in A\cotimes_R C$,
	we see that it suffices to construct a corresponding ideal
	$J_i$ for each $x_i$; the ideal $J := \sum_i J_i$ % \mar{TG:
          % should be $\sum$?}
	will then satisfy the claim of the lemma.
	Thus we reduce to the proving the lemma in the case of an element
	$x \in A\cotimes_R C = A\cotimes_R C^+[1/u].$
	Recall from Lemma~\ref{lem:non-zero divisor} that
	$A\cotimes_R C^+$ is a subring of $A\cotimes_R C$.
	Thus, since $u$ is a unit in $A\otimes_R C$,
	we may replace $x$ by $u^n x$ for some sufficiently large
	value of $n$, and assume that $x \in A\cotimes_R C^+$.

	Note that tensoring the inclusion 
	$A\cotimes_R C^+ \hookrightarrow A\cotimes_R C$
	with $B\cotimes_R C^+$ over $A\cotimes_R C^+$
	induces the inclusion
	$$B\cotimes_R C^+ \hookrightarrow B\cotimes_R C 
	= B\cotimes_R C^+[1/u] = (B\cotimes_R C^+)\otimes_{A\cotimes_R C^+}
	A\cotimes_R C.$$
	Thus it suffices to construct an ideal $J$ in $A$
	such that the image of $x$ in $B\cotimes_R C^+$  vanishes
	if and only if the morphism $\phi:A \to B$ factors through $A/J$.
	Let $x_n$ denote the image of $x$ in $A\otimes_R (C^+/u^n)$,
	so that
	$$x = (x_n) \in A\cotimes_R C^+ = \varprojlim_n 
	A\otimes_R (C^+/u^n).$$
	If we let $y_n$ denote the image of $x_n$
	in $B\otimes_R (C^+/u^n)$,
	then the image $y$ of $x$ in $B\cotimes_R C$
	is equal to
       	$$(y_n) \in B\cotimes_R C^+ = \varprojlim_n 
	B\otimes_R (C^+/u^n).$$
	Thus it suffices to construct an ideal $J_n$ in $A$
	such that $y_n$ vanishes if and only if the morphism
	$A\to B$ factors through $A/J_n$; indeed, we may then take
	$J := \sum_n J_n$. %\mar{TG: should be $\sum$?}
	
	Lemma~\ref{lem:basic properties} shows that 
	$C^+/u^n$ is flat over $R$, and thus free over $R$.
	(See Remark~\ref{rem:flat implies free}.) %\mar{ME: Guessing that this is true for arbitrary Artinian local rings; it's certainly true
	%	for the case $R = \Z/p^a$, which is the ring of interest
	%	in applications. ME: Checked this --- it's true!}
	If we choose an isomorphism $C^+/u^n \cong R^{\oplus S},$
	then we may write 
	$x_n = (x_{n,s}) \in A^{\oplus S}.$
	If we let $J_{n,s}$ denote the ideal in $A$ generated
	by $x_{n,s}$, then $J_n := \sum_{s \in S} J_{n,s}$ % \mar{TG:
        % should be $\sum$?}
	is the required ideal.
\end{proof}

\begin{acor}
	\label{cor:vanishing locus}
	Let $A$ be an $R$-algebra,
	and let $M$ and $N$ be finitely generated projective modules
	over $A\cotimes_R C$.  If $f: M \to N$ is an
        $A\cotimes_R C$-module homomorphism,
	then the ``locus of vanishing''
	of $f$ is a Zariski closed subset of $\Spec A$.
	More precisely, there is an ideal $J$ of $A$
	such that for any morphism $\phi:A \to B$ of $R$-algebras,
	the base-change
	$$f_B :
	M_B := (B\cotimes_R C) \otimes_{A\cotimes_R C} M\to
	N_B := (B\cotimes_R C) \otimes_{A\cotimes_R C} N
	$$
       	vanishes if and only if $J$ is contained in the kernel
	of~$\phi$.
\end{acor}
\begin{proof}
	We may regard $f \in \Hom_{A\cotimes_R C}(M,N)$
	as an element of the module
	$$M^{\vee}\otimes_{A\cotimes_R C} N.$$
	The corollary then follows from Lemma~\ref{lem:vanishing
		locus} applied to $f$ so regarded.
\end{proof}

\begin{acor}
	\label{cor:equivariance locus}
	Let $\sigma$ be a ring automorphism of $C$,
	let $A$ be an $R$-algebra,
	and let $M$ and $N$ be finitely generated projective modules
	over $A\cotimes_R C$ each endowed with a $\sigma$-semi-linear
	automorphism, denoted $\sigma_M$ and $\sigma_N$ respectively. 
	If $f: M \to N$ is an
        $A\cotimes_R C$-module homomorphism,
	then the ``locus of $\sigma$-equivariance''
	of $f$ is a Zariski closed subset of $\Spec A$.
	More precisely, there is an ideal $J$ of $A$
	such that for any morphism $\phi:A \to B$ of $R$-algebras,
	the base-change
	$$f_B :
	M_B := (B\cotimes_R C) \otimes_{A\cotimes_R C} M\to
	N_B := (B\cotimes_R C) \otimes_{A\cotimes_R C} N
	$$
       	satisfies $f_B\circ \sigma_M = \sigma_N \circ f_B$
	if and only if $J$ is contained in the kernel
	of~$\phi$.
\end{acor}
\begin{proof}
	This follows from Corollary~\ref{cor:vanishing locus},
	applied to the morphism $f  - \sigma_N \circ f \circ \sigma_M^{-1}.$
\end{proof}
We also have the following variants on these results for $A\cotimes_RC^+$-modules.
\begin{alem}%\mar{TG: should be checked by ME}
  \label{lem: inclusion lattices closed condition}Let~$A$ be an
  $R$-algebra, and let~$M$ and~$N$ be finitely generated projective
  modules over~$A\cotimes_RC^+$. Let~$f:M\to N[1/u]$ be an
  $A\cotimes_RC^+$-module homomorphism. Then the locus where
  $f(M)\subseteq N$ is a Zariski closed subset of~$\Spec A$. More
  precisely, there is an ideal  $J$ of $A$
	such that for any morphism $\phi:A \to B$ of $R$-algebras,
	the base-change
	$$f_B :
	M_B := (B\cotimes_R C) \otimes_{A\cotimes_R C} M\to
	N_B[1/u] := (B\cotimes_R C) \otimes_{A\cotimes_R C} N[1/u]
	$$
       	satisfies $f_B(M_B) \subseteq N_B$
	if and only if $J$ is contained in the kernel
	of~$\phi$.
\end{alem}
\begin{proof}We argue as in the proof of Lemma~\ref{lem:vanishing
    locus}. Since~$M,N$ are finitely generated projective
  over~$A\cotimes_RC^+$, we may find finitely generated projective
  modules~$P,Q$ over~$A\cotimes_RC^+$ such that~$M\oplus P$
  and~$N\oplus Q$ are free. Replacing~$M$ by~$M\oplus P$, $N$
  by~$N\oplus Q$, and~$f$ by~$(f,0)$, we reduce to the case that~$M$,
  $N$ are both finite free~$A\cotimes_RC^+$-modules. After choosing
  bases we may suppose that~$M=(A\cotimes_RC^+)^{\oplus r}$,
  $N=(A\cotimes_RC^+)^{\oplus s}$, and considering the matrix
  representing~$f$, we reduce to the showing that for any
  element~$x\in A\cotimes_RC$, there is an ideal~$J$ of~$A$ such that
  $(\phi\otimes 1)(x)\in B\cotimes_RC^+$ if and only if~$J$ is
  contained in the kernel of~$\phi$. This is a special case of
  Lemma~\ref{lem: element of lattice closed condition} below.
\end{proof}


\begin{alem}%\mar{TG: should be checked by ME}
  \label{lem: element of lattice closed condition}Let~$A$ be an
  $R$-algebra, and let~$M$ be a finitely generated projective
  $A\cotimes_RC^+$-module. Let~$x$ be an element of~$M[1/u]$. Then the
  locus where $x\in M$ is a Zariski closed subset of~$\Spec A$. More
  precisely, there is an ideal $J$ of $A$ such that for any morphism
  $\phi:A \to B$ of $R$-algebras, the image of~$x$ in~$M_B[1/u]$ lies
  in~$M_B$ if and only if $J$ is contained in the kernel of~$\phi$.
\end{alem}
\begin{proof}We begin by arguing  as in the proof of Lemma~\ref{lem: inclusion
    lattices closed condition}. Since~$M$ is finitely generated
  projective over~$A\cotimes_RC^+$, we may find a finitely generated
  projective module~$P,Q$ over~$A\cotimes_RC^+$ such that~$M\oplus P$
  is free. Replacing~$M$ by~$M\oplus P$ and ~$x$ by~$(x,0)$, we reduce to the case that~$M$
  is a finite free~$A\cotimes_RC^+$-module. After choosing
  a basis we may suppose that~$M=(A\cotimes_RC^+)^{\oplus r}$, so that
  we can reduce to the case that~$M=A\cotimes_RC^+$.
  
  Choose~$n$ sufficiently large so that $u^nx\in A\cotimes_RC^+$. Then
  $(\phi\otimes 1)(x)\in B\cotimes_RC^+$ if and only if the image
  of~$(\phi\otimes 1)(x)$ in
  \[u^{-n}(B\cotimes_RC^+)/(B\cotimes_RC^+)=B\otimes_R(u^{-n}C^+)/C^+)\]
  vanishes. As in the proof of Lemma~\ref{lem:vanishing locus},
  $u^{-n}C^+/C^+$ is a free~$R$-module, so we can and do choose an
  isomorphism~$u^{-n}C^+/C^+\cong R^{\oplus S}$. Then if we
  write~$x=(x_s)\in A^{\oplus S}\cong A\otimes_R (u^{-n}C^+/C^+)$, we
  can take~$J$ to be the ideal generated by the~$x_s$.
\end{proof}

%We are interested in studying finitely generated projective
%modules over $A\cotimes_R C$.  In the present general context, we don't
%know whether an analogue of Drinfeld's descent result for Tate modules
%(see~\cite[Thm.\ 3.3, 3.11]{MR2181808} and~\cite[Thm.\ 5.1.18]{EGstacktheoreticimages}) holds
%for arbitrary $R$-algebras $A$.  However, we will see that such
%a descent result is valid in the more restricted setting of finite
%type $R$-algebras.  In the case of finite type $k$-algebras,
%we will prove this using results from rigid analysis.  For finite type
%$R$-algebras that are not necessarily annihilated by $I$,
%we will use grading techniques to reduce to the case of $k$-algebras.


%We next
%state and prove some properties of $u$-adic completions that we will need.
%
%\begin{alemma}
%	\label{lem:completion properties}
%	{\em (1)} The functors $M \mapsto M\cotimes_R C^+$
%        and $M\mapsto M\cotimes C$ are exact.
%
%	{\em (2)} If $A$ is an $R$-algebra, if $J$ is a finitely
%	generated ideal in $A$, and if $M$ is an $A$-module,
%	then the natural map
%	$$J (M\cotimes_R C^+) \to JM \cotimes_R C^+$$
%        is an isomorphism, and consequently there is a short
%	exact sequence
%	$$0 \to J(M\cotimes_R C^+) \to M\cotimes_R C^+
%	\to (M/JM)\cotimes_R C^+ \to 0.$$
%
%	{\em (3)}
%	If $J_1$ and $J_2$ are ideals in an $R$-algebra $A$,
%	with $J_1 \supseteq J_2$, and if $M$ is an $A$-module,
%       then there is a natural isomorphism
%$$J_1 (M\cotimes_R C^+)/J_2 (M\cotimes_R C^+) \iso (J_1 M/J_2 M) 
%\cotimes_R C^+.$$
%\end{alemma}
%\begin{proof}
%	If $0 \to M_1 \to M_2 \to M_3 \to 0$ is a short exact sequence
%	of $R$-modules, then Lemma~\ref{lem:basic properties} 
%	shows that 
%	$$0 \to M_1\otimes_R (C^+/u^n) \to M_2\otimes_R (C^+/u^n) \to 
%	M_3\otimes_R (C^+/u^n) \to  0$$
%	is exact for each $n$.  If we pass to the inverse limit over~$n$,
%	and note that the transition morphisms are evidently surjective,
%	so that the relevant $\varprojlim^1$ vanishes,
%	we obtain a short exact sequence
%	$$0 \to M_1\cotimes_R C^+ \to M_2\otimes_R C^+ \to 
%	M_3\otimes_R C^+ \to  0,$$
%	proving the first exactness claim of~(1).
%	The second exactness claim follows from the first,
%	together with the fact that the localization
%	$C := C^+[1/u]$ is flat over $C$, and that there is an
%	isomorphism
%	$M \cotimes_R C \iso (M\cotimes_R C^+)\otimes_{C^+}C.$
%
%%	We may write $C = \varinjlim_n u^{-n} C^+$ (the transition
%%	maps being the evident inclusions),
%%	and thus (since tensor products are compatible with formation
%%	of direct limits) obtain the sequence
%%	$$0 \to M_1\cotimes_R C \to M_2\otimes_R C \to 
%%	M_3\otimes_R C \to  0$$
%%	by passing to the direct limit over $n$ of the sequences
%%	$$0 \to M_1\cotimes_R u^{-n}C^+ \to M_2\otimes_R u^{-n}C^+ \to 
%%	M_3\otimes_R u^{-n}C^+ \to  0,$$
%%	which have already been shown to be short exact.
%%	Since the formation of direct limits is exact, we obtain
%%	the second exactness claim of~(1).
%
%%	In order to prove~(2), we will show that
%%	if $J$ is an ideal in $R$, 
%%	then there is an isomorphism
%%	\numequation
%%	\label{eqn:complete iso}
%%	J (M\cotimes_R C^+) \iso (J M)\cotimes_R C^+
%%\end{equation}
%%	(i.e.\ we will first prove~(2) in the case when $J_2 = 0$). 
%%	The general case of~(2) then follows from this special
%%	case and the first exactness result of~(1).
%
%	In order to prove~(2), 
%	we apply~(1) to the short exact sequence $0 \to JM \to M
%	\to M/JM \to 0$, obtaining a short exact sequence
%	\anumequation
%	\label{eqn:J ses}
%	0 \to (J M)\cotimes_R C^+ \to M\cotimes_R C^+ 
%	\to (M/JM)\cotimes_R C^+ \to 0.
%\end{equation}
%        Thus we may (and do) 
%	regard $(JM)\cotimes_R C^+$ as a submodule of $M\cotimes_R C^+.$
%        There is a consequent inclusion
%	\anumequation
%	\label{eqn:J inclusion}
%	J (M\cotimes_R C^+)
%	\subseteq (JM)\cotimes_R C^+,
%        \end{equation}
%	whose image is $u$-adically
%	dense in the target.
%	Since $J$ is a finitely generated ideal,
%	%(as $R$ is Artinian, and thus Noetherian),
%	we see that $J(M\cotimes_R C^+)$ is furthermore $u$-adically complete,
%	and thus that~(\ref{eqn:J inclusion}) is actually an equality,
%	establishing the first claim of~(2).
%
%	%Consequently, we find that the natural morphism
%	%$(JM)\cotimes_R C^+  \to 
%	%M\cotimes_R C^+$ is a closed embedding (when the source
%	%is equipped with its $u$-adic topology), and so we may (and do)
%	%regard $(JM)\cotimes_R C^+$ as a $u$-adically closed
%	%submodule of $M\cotimes_R C^+.$
%%
%%        Clearly there is an inclusion
%%	\anumequation
%%	\label{eqn:J inclusion}
%%	J (M\cotimes_R C^+)
%%	\subseteq (JM)\cotimes_R C^+,
%%        \end{equation}
%%	whose image is $u$-adically
%%	dense in the target.
%%	Since $J$ is a finitely generated ideal,
%%	%(as $R$ is Artinian, and thus Noetherian),
%%	we see that $J(M\cotimes_R C^+)$ is furthermore $u$-adically complete,
%%	and thus that~(\ref{eqn:J inclusion}) is actually an equality,
%%	establishing the first claim of~(2).
%	%This establishes the required isomorphism~(\ref{eqn:complete iso}).
%	Replacing $(JM)\cotimes_R C^+$ by $J(M\cotimes_R C^+)$ in
%	the short exact sequence~(\ref{eqn:J ses}) yields the short
%	exact sequence whose existence is asserted in the second claim
%	of~(2).
%
%	Claim~(3) follows directly from~(2), and the exactness result
%	of~(1).
%\end{proof}
%
%%We are interested in studying finitely generated projective
%%modules over $A\cotimes_R C$.  In the present general context, we don't
%%know whether an analogue of Drinfeld's descent result for Tate modules
%%(see~\cite[Thm.\ 3.3, 3.11]{MR2181808} and~\cite[Thm.\ 5.1.18]{EGstacktheoreticimages}) holds
%%for arbitrary $R$-algebras $A$.  However, we will see that such
%%a descent result is valid in the more restricted setting of finite
%%type $R$-algebras.  In the case of finite type $k$-algebras,
%%we will prove this using results from rigid analysis.  For finite type
%%$R$-algebras that are not necessarily annihilated by $I$,
%%we will use graded techniques to reduce to the case of $k$-algebras.
%%We describe this reduction step first.
%
%%\section{Associated graded algebras and modules}
%%We maintain the notation ($R$, $C^+$, etc.) and assumptions from the
%%preceding subsection.  
%
%%\section{Associated graded algebras and modules}
%%
%%\begin{adefn} If $M$ is an $R$-module, then we let $\Gr^{\bullet} M$
%%	denote the graded $k$-vector space associated to the $I$-adic filtration
%%	on $M$ (so $\Gr^i M := I^iM/I^{i+1} M$).  
%%\end{adefn}
%%
%%The formation of $\Gr^{\bullet} M$ is functorial in $M$.
%%If $M$ and $N$ are two $R$-modules, then there is a natural surjection
%%(of graded $k$-vector spaces)
%%\anumequation
%%\label{eqn:graded tensor}
%%\Gr^{\bullet} M \otimes_k \Gr^{\bullet} N \to \Gr^{\bullet}(M\otimes_R N).
%%\end{equation}
%%In particular, it follows that if $A$ is an $R$-algebra, then $\Gr^{\bullet} A$
%%is naturally a graded $k$-algebra, and if $M$ is an $A$-module,
%%then $\Gr^{\bullet} M$ is naturally a $\Gr^{\bullet} A$-module.
%%
%%We recall the following basic lemmas.
%%
%%\begin{lemma}
%%	\label{lem:graded surjective}
%%	If $A$ is an $R$-algebra and if $M$ is an $A$-module,
%%	then the natural morphism $\Gr^{\bullet} A \otimes_{\Gr^0 A}
%%	\Gr^0 M \to \Gr^{\bullet} M$ is surjective.
%%	Furthermore, $M$ is non-zero if and only if $\Gr^0 M$ is non-zero,
%%	if and only if $\Gr^{\bullet} M$ is non-zero.
%%\end{lemma}
%%\begin{proof}
%	%The first claim is immediate.
%	%This implies in turn that if $\Gr^0 M$ is zero then 
%	%$\Gr^{\bullet} M$ is zero (the converse being evident), 
%	%and thus that $M$ is zero (since $\Gr^{\bullet} M$ is
%	%the associated graded of $M$ with respect to a finite length
%	%filtration beginning at $M$ and ending at $0$).
%	%(Essentially equivalently, one can observe directly that since $I$
%%	is nilpotent, if $M/IM = 0$ then $M = 0$.)
%%\end{proof}
%%
%%\begin{lemma}
%%	\label{lem:graded fg}
%%	If $A$ is an $R$-algebra and $M$ is an $A$-module,
%%	then $M$ is finitely generated over $A$ if and only 
%%	if $\Gr^0 M$ is finitely generated over $\Gr^0 A$,
%%	if and only if $\Gr^{\bullet} M$ is finitely generated
%%	over $\Gr^{\bullet} A$.
%%\end{lemma}
%%\begin{proof} 
%%Clearly if $M$ is finitely generated over $A$, then
%%$\Gr^0 M := M/IM$ is finitely generated over $A/I$; the converse
%%assertion follows from Nakayama's lemma with respect to
%%the nilpotent ideal $I$.  If $\Gr^0 M$ is finitely generated 
%%over $\Gr^0 A$, then Lemma~\ref{lem:graded surjective}
%%shows that $\Gr^{\bullet} M$ is finitely generated over $\Gr^{\bullet} A$;
%%the converse assertion follows from the fact that (again by
%%Lemma~\ref{lem:graded surjective}) $\Gr^0 M$ may be obtained as the
%%quotient of $\Gr^{\bullet} M$ by the ideal $\Gr^{\bullet > 0} A$
%%of $\Gr^{\bullet} A$.
%%\end{proof}
%%
%%\begin{lemma}
%%	\label{lem:graded flat}
%%If $A$ is an $R$-algebra, then an $A$-module $M$
%%is {\em (}faithfully{\em )} flat over $A$
%%if and only if $\Gr^{\bullet} M$ 
%%is {\em (}faithfully{\em )} flat over $\Gr^{\bullet} A.$
%%Furthermore, if any of these conditions holds,
%%then the natural morphism
%%$\Gr^{\bullet} A \otimes_{\Gr^0 A} \Gr^0 M \to
%%\Gr^{\bullet} M$
%%is an isomorphism.
%%\end{lemma}
%%\begin{proof}
%%	If $M$ is flat over $A$, then considering the
%%        result of tensoring $M$
%%	by the various short exact sequences
%%	$$0 \to I^nA \to I^mA \to I^mA/I^nA \to 0$$
%%	we find that the natural
%%	morphism $\Gr^i A \otimes_{\Gr^0 A} \Gr^0 M \to \Gr^i M$
%%	is an isomorphism, for each $i$.
%%	This proves the final assertion of the lemma.
%%	Furthermore, since (faithful) flatness is preserved
%%	under base-change, we see first that $\Gr^0 M$ is flat
%%	over $\Gr^0 A$ (and faithfully flat if $M$ is faithfully
%%	flat over $A$), and then (using the result already proved)
%%	that $\Gr^{\bullet} M$ is flat over $\Gr^{\bullet} A$
%%	(and faithfully flat if $M$ is).
%%	
%%	It is not quite as obvious that flatness of $\Gr^{\bullet} M$
%%	over $\Gr^{\bullet} A$ implies the flatness of $M$ over $A$,
%%	but this is a special case of.
%%	\mar{ME: Cite something}
%%	If $\Gr^{\bullet} M$ is furthermore faithfully flat over
%%	$\Gr^{\bullet} A$, then (since faithful flatness is preserved
%%	under base-change) we see that $\Gr^0 M := M/IM$ is faithfully flat
%%	over $\Gr^0 A := A/IA$.  Since an $A$-module vanishes
%%	if and only if its reduction mod $I$ does (see Lemma~\ref{lem:graded
%%		surjective}),
%%	we conclude that $M$ is faithfully flat over $A$.
%%\end{proof}
%%
%%\begin{lemma}
%%	\label{lem:graded tensor}
%%	If $A$ is an $R$-algebra, and if $M$ and $N$ are $A$-modules,
%%	then the natural surjection~{\em (\ref{eqn:graded tensor})}
%%	induces a natural surjection 
%%	$$\Gr^{\bullet} M \otimes_{\Gr^{\bullet} A} \Gr^{\bullet} N \to
%%	\Gr^{\bullet} (M\otimes_A N).$$
%%\end{lemma}
%%\begin{proof}
%%	The natural surjection $M\otimes_R N \to M\otimes_A N$ induces
%%	a surjection $\Gr^{\bullet}(M\otimes_R N) \to \Gr^{\bullet} (M\otimes_A N).$  It is immediate that its composite with the surjection~(\ref{eqn:graded
%%		tensor}) factors through $\Gr^{\bullet} M \otimes_{\Gr^{\bullet}
%%		A} \Gr^{\bullet} N.$
%%\end{proof}
%%
%%\begin{lemma}
%%	\label{lem:graded flat bis}
%%	If $A$ is an $R$-algebra, and if $M$ and $N$ are $A$-modules 
%%	with either $M$ or $N$ flat over $A$,
%%	then the surjection of Lemma~{\em \ref{lem:graded tensor}}
%%       	is a natural isomorphism
%%	$$\Gr^{\bullet} M \otimes_{\Gr^{\bullet} A} \Gr^{\bullet} N \iso
%%	\Gr^{\bullet} (M\otimes_A N).$$
%%\end{lemma}
%%\begin{proof}
%%	Suppose (without loss of generality) that $N$ is $A$-flat.
%%	Taking into account Lemma~\ref{lem:graded flat},
%%	we see that we have to show that natural morphism
%%	$\Gr^{\bullet} M \otimes_{\Gr^0 A} \Gr^0 N \to \Gr^{\bullet}(M\otimes_A N)$ is an isomorphism, i.e.\ that for each $i \geq 0,$ the natural morphism
%%	$$I^iM /I^{i+1} M \otimes_{A/I} N/I \to I^i(M\otimes N)/I^{i+1}(M\otimes N)$$
%%	is an isomorphism.  This follows easily by tensoring $N$ over $A$
%%	with the various short exact sequences
%%$$0 \to I^n M \to I^m M \to I^mM/I^nM\to 0,$$
%%taking into account the flatness of $N$ over $A$.
%%\end{proof}
% 
%%Now let $M$ be an $R$-module.  We wish to apply the preceding discussion
%%to $M\cotimes_R C^+$ and $M\cotimes_R C$.
%
%\begin{alemma}
%	\label{lem:passing to graded context}
%	If $M$ is an $R$-module,
%	then there are natural isomorphisms
%
%	{\em (1)}
%	$\Gr^{\bullet} M \otimes_{k} \Gr^0 (C^+/u^n) \iso
%	\Gr^{\bullet} \bigl(M\otimes_R
%	(C^+/u^n)\bigr)$
%{\em (}for any $n  \geq 1${\em )},
%
%{\em (2)}
%$\Gr^{\bullet} (M\cotimes_R C^+) \iso
%\Gr^{\bullet} M \cotimes_{k} \Gr^0 C^+ ,$
%	and
%
%	{\em (3)}
%	$\Gr^{\bullet} (M\cotimes_R C)
%	\iso
%	\Gr^{\bullet} M \cotimes_{k} \Gr^0 C.$ 
%\end{alemma}
%\begin{proof}
%	Lemma~\ref{lem:basic properties}~(1)
%	shows that $C^+/u^n$ is flat over $R$,
%	and~(1) follows from this,
%	together with Lemmas~\ref{lem:graded flat} and~\ref{lem:graded flat bis}.
%
%	Since $C^+/u^n$ is flat over $R$, we see that the natural map
%	$(\Gr^0 C^+)/u^n \to \Gr^0( C^+/ u^n)$ is an isomorphism,
%	and combining this with the natural isomorphism of~(1),
%	we obtain a natural isomorphism
%	$$\Gr^{\bullet} M \otimes_{k} (\Gr^0 C^+)/u^n 
%	\iso \Gr^{\bullet} \bigl(M\otimes_R (C^+/u^n)\bigr).$$
%	Taking this into account, we find that to prove~(2), it suffices 
%	to show that the natural morphism
%	$$
%	\Gr^{\bullet}( M\cotimes_R C^+) = 
%	\Gr^{\bullet}\bigl(\varprojlim_n M\otimes_R (C^+/u^n)\bigr)
%	\to \varprojlim_n \Gr^{\bullet} \bigl(M\otimes_R (C^+/u^n)\bigr)
%	$$
%	is an isomorphism.
%	This follows directly from Lemma~\ref{lem:completion
%		properties}~(3) (taking the inclusion $J_2 \subseteq
%	J_1$ to be the various inclusion $I^{i+1} \subseteq I^i$
%	in turn).
%
%	%As already noted in the proof of Lemma~\ref{lem:completion
%	%	properties},
%	The $R$-algebra $C$ may be described as the direct limit
%	$ C \iso \varinjlim_n u^{-n}C^+$ (the transition
%	maps being the obvious inclusions).
%	Since the formation of direct limits is compatible with
%	tensor products, we find that
%	$$\Gr^0 C \iso \varinjlim_n \Gr^0 u^{-n} C^+.$$
%	The isomorphism of~(3) is thus obtained from
%	the isomorphisms
%	$$\Gr^{\bullet} ( M\cotimes_R u^{-n} C^+) \iso
%	\Gr^{\bullet} M \cotimes_{k} \Gr^0 u^{-n} C^+$$
%	of~(2) by passing to the
%	direct limit over $n$.
%\end{proof}

\section{Extending actions}
We continue to remain in the context of Section~\ref{subsec:setting}.
We will show that various sorts of actions on $C^+$ or $C$
can be extended to the completed tensor products $A\cotimes_R C^+$
and $A\cotimes_R C$.

Throughout our discussion,
we endow $C^+$ with the $u$-adic topology, while we endow $C$
with the unique topology which makes it a topological group,
and in which $C^+$ is an open subgroup.  
More generally, if $A$ is any $R$-algebra,
we again endow
$A\cotimes_R C^+$ with the $u$-adic topology,
while we endow
$A\cotimes_R C$ 
with the unique topology which makes it a topological group,
and in which $A\cotimes_R C^+$ (endowed with its $u$-adic topology)
is an open subgroup.  

\begin{alemma}
	\label{lem:extending endomorphisms}
	Suppose that $C^+$ {\em (}resp.\ $C${\em )}
	is equipped with a continuous $R$-linear endomorphism~$\varphi$.
	Then there is a unique extension of $\varphi$
	to a continuous $A$-linear endomorphism of $A\cotimes_R C^+$
	{\em (}resp.\ $A\cotimes_R C${\em )}.
\end{alemma}
\begin{proof}
	Lemma~\ref{lem:non-zero divisor} shows
	that $u$-adic topology on $A\cotimes_R C^+$ coincides
	with its inverse limit topology,
	if we recall that
	$A\cotimes_R C^+ := \varprojlim_n A\otimes_R (C^+/u^n),$
	and we endow each of the objects appearing in the 
	inverse limit with its discrete topology.
	Given this, it is immediate that a $u$-adically continuous
	$R$-linear endomorphism
	$\varphi$ of $C^+$ extends uniquely to a $u$-adically
	continuous $A$-linear endomorphism of $A\cotimes_R C^+$.
	
	The case when $\varphi$ is a continuous $R$-linear endomorphism
	of $C$ is only slightly more involved.  Namely, to say
	that $\varphi: C \to C$ is continuous is to say that,
	for some $m \geq 0$, we have an inclusion $\varphi( u^m C^+)
	\subseteq C^+$, and that the induced morphism
	$\varphi: u^m C^+ \to C^+$ is continuous, when each of the
	source and the target are endowed with their $u$-adic topologies.
	Again taking into account
	Lemma~\ref{lem:non-zero divisor},
       we then obtain a unique $A$-linear and continuous morphism
       \begin{multline*}
A\cotimes_R C =
\bigl( \varprojlim_n A\otimes_R (u^m C^+/u^{m+n}) \bigr) [1/u] 
\\
\longrightarrow
\bigl( \varprojlim_n A\otimes_R (C^+/u^{n}) \bigr) [1/u]  
= A\cotimes_R C,
\end{multline*}
as required.
\end{proof}

If $G$ is a topological group acting continuously on~$C^+$
(i.e.\ acting in such a way that the action map $G\times C^+ \to C^+$
is continuous), then for any $g \in G$ and $m\geq 0$, we may find $n \geq 0$,
and an open neighbourhood $U$ of~$g$,
such that $U \times u^n C^+ \subseteq u^m C^+$.  (This just expresses
the continuity of the action map at the point $(g,0) \in G\times C^+$.)
If $G$ is furthermore compact, then $G$ is covered by finitely many
such open sets $U$, and thus
for any given choice of~$m$, we may in fact find $n$
such that $G \cdot u^n C^+ \subseteq u^m C^+$.
An identical remark applies in the case of a compact topological group
acting on~$C$.

We will need to introduce an additional condition on $G$-actions on~$C$,
which is related to, but doesn't seem to follow from,
the previous considerations. 

\begin{adefn}\label{adefn: bounded group action}
	If $G$ is a group,
	we say that a $G$-action on $C$
	is {\em bounded} if, for any $M\geq 0,$ there exists
	$N \geq 0$ such that $G\cdot u^{-M} C^+ \subseteq u^{-N} C^+.$
\end{adefn} 

\begin{arem}
  \label{arem: bounded doesn't depend on u}
  It follows from Remark~\ref{rem:independence of u}
  that the notion of a group action being bounded is independent
  of the choice of $u\in C^+$ satisfying conditions~(A), (B), and~(C)
  of Section~\ref{subsec:setting}.
 % Note that since the
 % $u$-adic topology on~$C^+$ is independent of the choice of~$u$, the
 % notion of the action of~$G$ being bounded is also independent of~$u$.
\end{arem}

We then have the following lemma.

% \begin{alemma}\label{alem: bounded mod I implies bounded}
% 	If the group $G$ acts $R$-linearly on $C$,
% 	and the induced action on $C/I$ is bounded,
% 	then the given action on $C$ is also bounded.
% \end{alemma}
% \begin{proof}
% 	\mar{ME: Needs proof.  Or maybe not true? TG: I don't see how
%           to prove this, in fact - but I guess in the application we
%           can just do the $W_a(\O_{\C}^\flat)$ case and take~$u=[v]$
%           for some~$v$ with positive valuation, and then it's immediate?}
% \end{proof}

\begin{alemma}
	\label{lem:extending group actions}
	Suppose that $G$ is a compact topological group,
	and that $C^+$ {\em (}resp.~$C${\em )}
	is equipped with a continuous {\em (}resp.\ continuous
	and bounded{\em )} $R$-linear $G$-action. 
	%i.e.\ an $R$-linear action for which the action map
	%$G\times C^+ \to C^+$ {\em (}resp. $G\times C \to C${\em )}
	%is jointly continuous.
	Then there is a unique extension of this action 
	to a continuous $A$-linear action of $G$
	on $A\cotimes_R C^+$
	{\em (}resp.\ $A\cotimes_R C${\em )}.
\end{alemma}
\begin{proof}
The proof of this lemma is very similar to that of Lemma~\ref{lem:extending
endomorphisms}.  Indeed, applying that lemma (and taking into account
its uniqueness statement) we obtain the desired action of $G$.
It remains to confirm that this action is again continuous.

Consider first the case of a $G$-action on $C^+$.
For each $m \geq 0,$ we saw above that $G \cdot u^n C^+ \subseteq
u^m C^+$ if $n$ is sufficiently large, 
so that we obtain a well-defined map
$G\times A\cotimes_R C^+ = G\times \varprojlim_n A\otimes_R C^+/u^n 
\to A\otimes_R C^+/u^m$ which is continuous (since it is just obtained
from the continuous $G$-action on $C^+$ by tensoring with~$A$).
Passing to the projective limit in~$m$ then yields the desired continuity.

In the case of a $G$-action on $C$, we argue similarly:
namely, the assumptions of continuity and boundedness
of the $G$-action yield continuous morphisms
$$ G\times u^{-M} C^+/u^n C^+ \to  u^{-N} C^+/ u^m C^+;$$
Passing to the projective limit in $n$, then in $m$, and then to the
inductive limit in $N$, and then in $M$, we deduce the desired 
continuity of the $G$-action on $A\cotimes_R C$.
\end{proof}%\mar{TG: just to note that this isn't finished yet!}

\section{A rigid analytic perspective}
\label{subsec:rigid analytic}
We keep ourselves in the setting
of Section~\ref{subsec:setting}. %the preceding subsections.
We are interested in studying finitely generated projective
modules over $A\cotimes_R C$, especially their 
descent properties.
In the present general context, we don't
know whether an analogue of Drinfeld's descent result for Tate modules
(see~\cite[Thm.\ 3.3, 3.11]{MR2181808} and~\cite[Thm.\ 5.1.18]{EGstacktheoreticimages}) holds
for arbitrary $R$-algebras $A$.  However, we will see that such
a descent result is valid in the more restricted setting of finite
type $R$-algebras.  In the case of finite type $k$-algebras,
we will prove this using results from rigid analysis.  For finite type
$R$-algebras that are not necessarily annihilated by $I$,
we will use grading techniques to reduce to the case of $k$-algebras.

If $A$ is a finite type $k$-algebra,
then $\Spf (A\cotimes_k \Gr^0 C^+)$ (where the $\Spf$ 
is taken with respect to the $u$-adic topology on $A \cotimes_k
\Gr^0 C^+$) 
is a formal scheme of finite type over $\Spf \Gr^0 C^+$,
whose rigid analytic generic fibre
$\Max\Spec (A\cotimes_k \Gr^0 C)$
is an affinoid rigid analytic space over the complete
non-archimedean field $\Gr^0 C$.

%If we choose an embedding $\Spec A \hookrightarrow \mathbb A^n_{/k},$
%then the closed unit ball $\BB^n$ in the base-changed affine space
%$\mathbb A^n_{/\Gr^0 C}$ is well-defined independent of the choice
%of coordinates on $\mathbb A^n_{/k}$, 
%and $\Max\Spec A\cotimes_k \Gr^0 C$ 
%is obtained as the intersection of the affine scheme
%$\Spec A\otimes_k C$ and closed unit ball $\BB^n$, the intersection
%taking place in $\mathbb A^n_{/C}$.
%\mar{ME: This seems intuitively clear to
%	me, but I'm not sure if it's obvious if one starts to unwind the
%	definitions in a literal way; if not, we should add a more careful
%	explanation/proof.
%	Lemma~\ref{lem:completion properties}~(2) will help with this.}

As an application of this rigid analytic point-of-view (together
with graded techniques), we first establish the following
proposition.
%\mar{TG: expect to use~\cite[Prop.\ 5.2.1]{MR2774689} for faithful
%  flatness at some point.  ME: Actually, should cite discussion
%  at beginning of~\cite[\S 5.2]{MR2774689}, which gives flatness. }
 
\begin{aprop}\label{prop: rigid analysis FGK}\leavevmode
  
  \begin{enumerate}[label=\normalfont(\arabic*)]
  \item  If $A$ is a finite type $R$-algebra, then $A\cotimes_R C$ is
    Noetherian.
  \item  If $A \to B$ is a {\em (}faithfully{\em )} flat morphism
    of finite type $R$-algebras, then the induced morphisms
    $A\cotimes_R C^+ \to B\cotimes_R C^+$ and
    $A\cotimes_R C \to B\cotimes_R C$ are again {\em (}faithfully{\em
      )} flat.
  \end{enumerate}
\end{aprop}
\begin{proof}
	In the case when $A$ is in fact a $k$-algebra,
	it follows from the preceding discussion that
        $A\cotimes_R C$ is an affinoid algebra over~$\Gr^0 C$, and is therefore Noetherian by~\cite[Thm.\ 1, \S5.2.6]{MR746961}.
	In the general case, 
	Lemma~\ref{lem:passing to graded context}~(3)
	gives an isomorphism
	$\Gr^{\bullet} (A\cotimes_R C)
	\iso
	\Gr^{\bullet} A \cotimes_{k} \Gr^0 C.$ 
	Since $\Gr^{\bullet} A$ is a finite type $k$-algebra,
	it follows from the case already proved that
	the target of this isomorphism is Noetherian, and thus
	so is its source.  Thus $A\cotimes_R C$ is itself Noetherian,
	by Lemma~\ref{lem:graded Noetherian}.
	%\mar{ME: This final passage from the associated graded to the original
	%	algebra is pretty obvious, but we could put it in 
	%	the associated graded discussion, just for completeness's
	%	sake.}


        We explain how our rigid analytic perspective
	proves a part of~(2).
	If we suppose first that $A\to B$ is a flat
        morphism of $k$-algebras,
	then so is the morphism $A\otimes_R C \to B\otimes_R C.$
	Now the preceding discussion shows that the morphisms
	$\Max\Spec (A\cotimes_R C) \to \Max\Spec (A\otimes_R C)$
	and 
	$\Max\Spec (B\cotimes_R C) \to \Max\Spec (B\otimes_R C)$
	are open immersions of rigid analytic spaces over $\Gr^0 C$,
	so that the induced morphism %of affinoid spaces
	$\Max\Spec (A\cotimes_R C) \to \Max\Spec (B\cotimes_R C)$
	is also flat.  This proves (the $k$-algebra case of) the
	second flatness claim of~(2).  
	If $A\to B$ is a flat morphism
	of $R$-algebras that are not necessarily $k$-algebras,
	then we deduce from Lemma~\ref{lem:graded flat}
	that $\Gr^{\bullet} A \to \Gr^{\bullet} B$ is flat,
	and thus, from what we've already shown, that 
	$$\Gr^{\bullet} A \cotimes_{k} \Gr^0 C
	\to \Gr^{\bullet} B \cotimes_{k} \Gr^0 C$$
	is flat.  
	Lemma~\ref{lem:passing to graded context}~(3) then
	shows that the morphism
	\[\Gr^{\bullet} (A\cotimes_R C ) \to \Gr^{\bullet} (B\cotimes_R C)\]
        is flat,
	and one more application of Lemma~\ref{lem:graded flat}
	completes the proof of the second flatness claim of~(2).


	It is not clear to us whether one can deduce the second
	faithful flatness claim of~(2) by this style of argument.
	In order to prove this result, 
	as well as to obtain the first set of claims of~(2),
	we appeal to the results
	of~\cite{MR2774689}.  Indeed, it is clear that the
	second set of claims in~(2) follows immediately from the first,
	and it the first set of claims that we will now prove.

        We first note that if the morphism
	$A \to B$ is flat (resp.\ faithfully
	flat), then so are each of the morphisms
	$A\otimes_R C^+/u^n A \to B\otimes_R ^+C/u^n.$
	In the terminology
	of~\cite[\S 5.2]{MR2774689},
	the morphism
	$A\cotimes_R C^+ \to B\cotimes_R C^+$
	is adically flat (resp.\ adically faithfully flat); here
	we regard the source and target as being endowed with their
	$u$-adic topologies.  The discussion at the beginning
	of~\cite[\S 5.2]{MR2774689}
(see also our Proposition~\ref{prop: FGK flat})
then shows that
	$A\cotimes_R C^+ \to B\cotimes_R C^+$
	is flat.  In the
	adically faithfully flat case, we deduce in addition 
	from~\cite[Prop.~5.2.1~(2)]{MR2774689}
(again, see also Proposition~\ref{prop: FGK flat}) 
that $A\cotimes_R C^+ \to B\cotimes_R C^+$
	is faithfully flat;  
	note that
	$A\cotimes_R C^+$ and $B\cotimes_R C^+$ are Noetherian outside 
	$u$ by~(1) of the present proposition.
	%The second
	%set of claims in~(2) follows immediately from the first.
\end{proof}

\section{Descent for $A\cotimes_R C$-modules}
We now establish the descent result alluded to above; 
it is analogous to 
% Theorem~\ref{thm: fpqc locality of projective and locally free}
Drinfeld's \cite[Thm.\ 5.1.18~(1)]{EGstacktheoreticimages}
but is restricted to the context of finite type $R$-algebras.
We use grading techniques to reduce to the case of finite type
$k$-algebras, where we can then apply known results from rigid analysis.

\begin{atheorem}
	Let $A \to B$ be a faithfully flat morphism of 
	finite type $R$-algebras.  Then the functor
	$M \mapsto
	M\otimes_{A\cotimes_R C} (B\cotimes_R C)$
	induces an equivalence of categories between
	the category of finite type $A\cotimes_R C$-modules
	and the category of finite type $B\cotimes_R C$-modules
	equipped with descent data.
\end{atheorem}

\begin{aremark}\label{rem: what is descent data}Recall (for example
  from~\cite[\S 6.1]{MR1045822}) that if~$S\to S'$ is a ring
  homomorphism, then a descent datum for an~$S'$-module~$M'$ is an
  isomorphism of~$S'\otimes_S S'$-modules $S'\otimes_SM'\isoto
  M'\otimes_SS'$, satisfying a certain cocycle condition. Given a
  descent datum, we have a pair of morphisms of~$S$-modules $M'\to M'\otimes_SS'$, namely the
  obvious morphism and the composite of the obvious morphism $M'\to
  S'\otimes_SM'$ with the descent datum isomorphism. We let~$K$ denote
  the $S$-module given by the kernel of the difference of these two
  maps.  We
  refer to the functor $M'\mapsto K$,  from the category of
  $S'$-modules with descent data to the category of $S$-modules, as
  the \emph{kernel functor}.  (If~$S\to S'$ is faithfully flat, then the theory of faithfully
  flat descent shows that~$M'$ is the extension of scalars from~$S$ to~$S'$
  of~$K$.)%\mar{TG: you can doubtless improve this!}
	% \mar{ME: Explain slighty more what we mean by descent data,
	% 	and explain that descent is achieved by computing
	% 	a certain kernel, perhaps with a citation to BC's paper
	% 	\cite{conrad-rigid-descent}.}
\end{aremark}
\begin{proof}
	We first treat the case when $A$ and $B$ are $k$-algebras.
	In this case, the category of
	finite type $A\cotimes_R C$-modules (resp.\ finite type $B\cotimes_R C$-modules)
	is equivalent to the category of coherent sheaves on
       the affinoid rigid analytic space $\Max\Spec (A\cotimes_R C)$ (resp.\
       $\Max\Spec (B\cotimes_R C)$) over the field $\Gr^0 C$, 
       and the statement follows from faithfully flat descent
       for rigid analytic coherent sheaves \cite[Thm.~3.1]{MR1603849}.

       We reduce the general case of $R$-algebras to the case of 
       $k$-algebras via the usual graded arguments.
       To ease notation, we let $F:\cC \to \cD$ denote the base-change functor 
	$M \mapsto
	M\otimes_{A\cotimes_R C} (B\cotimes_R C)$
	from
	the category $\cC$ of finite type $A\cotimes_R C$-modules
	to
	the category $\cD$ of finite type $B\cotimes_R C$-modules
	equipped with descent data,
	and let $G:\cD \to \cC$ denote the kernel functor.
	% \mar{ME: Add formula for kernel functor. TG: will
        %   Remark~\ref{rem: what is descent data} do?}
	%from the category of finite type $B\cotimes_R C$-modules
	%equipped with descent data
	%to the category of finite type $A\cotimes_R C$-modules.
	There are evident natural transformations
	\anumequation
	\label{eqn:adjunctions}
	G\circ F \to \id_{\cC} \quad \text{ and } \quad F \circ G \to \id_{\cD},
\end{equation}
	which we claim are natural isomorphisms (so that
	$G$ provides a quasi-inverse to $F$, proving in particular
	that $F$ induces an equivalence of categories).

	We let $\overline{\cC}$ denote the category of finite type
	$\Gr^{0} (A\cotimes_R C)$-modules, % \mar{TG: this notation is
          % confusing - I think it means $\Gr^{0}( A\cotimes_R C)$?}
	and let $\overline{\cD}$ denote the category of finite type
	$\Gr^{0} (B\cotimes_R C)$-modules
	equipped with descent data to $\Gr^{0} (A\cotimes_R C)$.
	We let $\overline{F}:\overline{\cC} \to \overline{\cD}$
	denote the base-change functor
	$\text{--}\otimes_{\Gr^{0} (A\cotimes_R C)}
	\Gr^{0} (B\cotimes_R C)$,
	and let $\overline{G}:\overline{D} \to \overline{C}$
	denote the kernel functor. % \mar{ME: Maybe explicate ``kernel fucntor''
		% a little. TG: again, OK now?}
        The case of $k$-algebras that we have already treated shows that
	there are natural isomorphisms
	\anumequation
	\label{eqn:overline adjunctions}
	\overline{G}\circ \overline{F} \iso \id_{\overline{\cC}}
	\quad \text{ and } \quad
	\overline{F}\circ \overline{G} \iso \id_{\overline{\cD}}.
\end{equation}

	We think of passage to the associated graded $\Gr^{\bullet}$
	as inducing functors $\cC \to \overline{\cC}$ and also 
	$\cD\to \overline{\cD}$.
	It follows from 
	Lemmas~\ref{lem:graded flat} and~\ref{lem:graded flat bis},
        together with Proposition~\ref{prop: rigid analysis FGK},
	that there is a natural isomorphism
	$$\Gr^{\bullet} \circ F \iso \overline{F} \circ \Gr^{\bullet};$$
	there is also an evident natural isomorphism
	$$\Gr^{\bullet} \circ G \iso \overline{G} \circ \Gr^{\bullet};$$
	furthermore, these natural isomorphisms are compatible with the natural transformations~(\ref{eqn:adjunctions}) and~(\ref{eqn:overline adjunctions}).
	% \mar{ME: Probably should cite some of the previous grading
        %   results for this compatibility. TG: not sure what to cite.}
	Since the natural transformations~(\ref{eqn:overline adjunctions})
	are isomorphisms, the same is true of the natural
        transformations~(\ref{eqn:adjunctions}), by Lemma~\ref{lem: graded iso is iso}. 
	% \mar{ME: Have to prove that if $\Gr^{\bullet} f$ is an iso., then
	% 	so is $f$.}
	This completes the proof that $F$ is an equivalence.
       \end{proof}



% \section{The classical theory}
% \label{subsec:classical}
% We suppose now that $R$ is a valuation ring, complete with respect to
% a non-discrete rank one valuation, that $I = \mathfrak m$ 
% is the maximal ideal of $R$, and that $u$ is an element
% of $R$ of positive valuation.

% We let $K$ denote the fraction field of $R$, and let $L$ denote


\chapter{Topological groups and modules}
\label{app:topological groups}
In this appendix we recall some more-or-less well-known facts regarding topological groups
and modules,
for which we haven't located a convenient reference.
We found the note \cite{Kaye-Polish} to be a useful
reference for the basic facts regarding Polish topological groups.

Recall the following definition.

\begin{adefn} \index{Polish group}
	A topological space is called {\em Polish} if it is separable (i.e.\
	contains a countable dense subset) and completely metrizable.
	A topological group is called {\em Polish} if its underlying 
	topological space is Polish.  Similarly,
	a topological ring is called {\em Polish} if its underlying 
	topological space is Polish.  
\end{adefn}

Our interest in Polish groups is due to the following lemma and
corollary; see
also~\cite[\href{https://stacks.math.columbia.edu/tag/0CQW}{Tag
  0CQW}]{stacks-project} for a more algebraic proof of a closely
related result.
% \mar{ME: See how this relates to
% \cite[\href{http://stacks.math.columbia.edu/tag/0CQW}{Tag
%   0CQW}]{stacks-project}. TG: just add a remark about this but don't
% change anything.} 
\begin{alemma}
	\label{lem:open mapping}
	If $\phi: G \to H$ is a continuous homomorphism
        between two Polish topological groups,
	then the following are equivalent:
	\begin{enumerate}
		\item $\phi$ is open.
		\item $\phi(G)$ is not meagre in $H$.
	\end{enumerate}
\end{alemma}
\begin{proof}
	If~(1) holds, then $\phi(G)$ is a non-empty open subset
        of the Polish, and hence Baire, space $H$,
	and so is not meagre.  Thus~(1) implies~(2).
	For the converse, see e.g.~\cite[Thm.~18]{Kaye-Polish}.
\end{proof}

%\begin{alemma}
%	\label{lem:open mapping}
%	If $\phi: G \to H$ is a continuous homomorphism
%        between two topological groups,
%	then the following are equivalent:
%	\begin{enumerate}
%		\item $\phi$ is open.
%		\item For each open neighbourhood $U$ of $1$ in $G$,
%			the image $\phi(U)$ contains an open neighbourhood
%			of $1$ in $H$.
%	\end{enumerate}
%	If furthermore the underlying topological
%	spaces of $G$ and $H$ are Polish,
%	then these conditions are equivalent to
%	the following condition:
%	\begin{enumerate}
%		\item[(3)] $\phi(G)$ is not meagre in $H$.
%	\end{enumerate}
%\end{alemma}
%\begin{proof}
%	Clearly~(1) implies~(2).  Converesly, suppose that (2) holds,
%	and let $U$ be an open subset of $G$.  If $g \in U$,
%	choose an open neighbourhood $V$ of $1$ in $G$
%	such that $gV\subseteq U$,
%	and choose an open neighbourhood $W$ of $1$ in $H$ such that
%	$W \subseteq \phi(U)$.  Then
%	$$\phi(g) W \subseteq \phi(gV) \subseteq \phi(U),$$
%	showing that $\phi(U)$ contains an open neighbourhood of 
%	each of its point, and thus is open.  Hence~(1) holds.
%
%	Suppose now that $G$ and $H$ have Polish underlying topological spaces.
%	If~(1) holds, then $\phi(G)$ is a non-empty open subset
%        of the Polish, and hence Baire, space $H$,
%	and so is not meagre.  Thus~(3) holds.
%	Suppose conversely that~(3) holds,
%	and choose a countable dense subset $\{g_n\}$ of $G$.
%	If $U$ is an open neighbourhood of $1$, choose an open
%	neighbourhood $V$ of $1$ such that $V^{-1} V \subseteq U$.
%	Since $\{g_n\}$ is dense, we have that
%	$G = \bigcup_n g_n V$, and so $\phi(G) = \bigcup_n
%	\phi(g_n) \phi(V).$  Since $\phi(G)$ is non-meagre, 
%	one (and hence all) of the sets $\phi(g_n) \phi(V)$ is non-meagre;
%	equivalently, the set $\phi(V)$ is non-meagre.
%
%	Since $\phi$ is a continuous map between Polish spaces,
%	the image $\phi(V)$ of the open subset $V$ of $G$
%	is a subset of $H$ having the {\em property of Baire},
%	and thus $\phi(V^{-1} V) = \phi(V)^{-1} \phi(V)$ contains an open neighbourhood
%	$W$ of $1$ in $H$.  Since $V^{-1} V \subseteq U$, 
%	we find that~(2) holds.
%\end{proof}

\begin{acor}
	\label{cor:open polish}
	A continuous surjective homomorphism
       	of Polish topological groups is necessarily
	open.
\end{acor}
\begin{proof}
	A completely metrizable space is Baire, and hence is not a
	meagre subset of itself.  The corollary thus follows from
	the implication~``(2) $\implies$ (1)'' of Lemma~\ref{lem:open mapping}.
\end{proof}

Any metrizable space is first countable (i.e.\ each point contains a countable
neighbourhood basis).  
Recall that, conversely, if $G$ is a topological group, then $G$
is first countable if and only if the identity element $1$ admits
a countable neighbourhood basis,
and in this case, if
(the underlying topological space of) $G$ is Hausdorff then it is in fact
metrizable (this is the Birkhoff--Kakutani theorem),
% \mar{TG: doesn't this need a Hausdorff hypothesis? (and then it's the
%   Birkhoff--Kakutani theorem). The same applies to the various
%   mentions of this below.}
and even admits a (left or right) translation invariant metric.
In particular, (the underlying topological space of)
a Hausdorff topological group is separable and metrizable if and only 
if it is {\em second countable} (i.e.\ admits a countable
basis for its topology).

Recall also that a topological group $G$ admits a canonical uniform structure,
so that it makes sense to speak of $G$ being complete.
For the sake of completeness, we note that the underlying
topological space of $G$ being Polish %if and only if
forces $G$ to be complete.
% \mar{TG: is this true? How does it relate to Remark~\ref{rem:Polish} below?
% ME: I think this was just a mis-statement, now corrected.}

\begin{alemma}
	\label{lem:complete}
	If $G$ is a Polish topological group,
	then $G$ is complete.
\end{alemma}
\begin{proof}
	Since $G$ is metrizable, it is Hausdorff, and so we may
	consider the canonical embedding of topological groups
	$G \hookrightarrow \widehat{G}$ of $G$ into its completion.
	Since $G$ is separable and dense in~$\widehat{G}$, we see
	that $\widehat{G}$ is separable.  Also, $\widehat{G}$ is a complete
	metric space (it may be identified with the metric space
	completion of $G$ with respect to any invariant
	metric inducing the topology on $G$).  Thus $\widehat{G}$ is Polish.
	Any Polish subset of a Polish space is~$G_{\delta}$ (i.e.\ a
        countable intersection of open sets),
	and so in particular $G$ is $G_{\delta}$ inside~$\widehat{G}$.
	Since $\widehat{G}$ is Polish, and thus Baire,
	we find that $G$ is not meagre in~$\widehat{G}$.
	It follows from Lemma~\ref{lem:open mapping}
	that the inclusion $G\into\widehat{G}$ is open, and thus that $G$ is an open subgroup
	of its completion~$\widehat{G}$.  Since an open subgroup
	of a topological group is also closed, we find that
	$G$ closed, as well as dense, in~$\widehat{G}$,
	and thus that $G = \widehat{G}$, which is to say, $G$
	is complete, as claimed.
\end{proof}

%\begin{adefn}
%	We say that a topological group is {\em Polish}
%	if its underlying topological space is Polish.
%\end{adefn}

\begin{aremark}
	\label{rem:Polish}
It follows from Lemma~\ref{lem:complete}
and the discussion preceding it that
a topological group is Polish if and only if it is complete
and second countable.
\end{aremark}

We next
recall a result about Noetherian and Polish topological rings, 
which is inspired by a result of Grauert and Remmert
in the theory of classical Banach algebras
\cite[App.\ to \S I.5]{MR0316742}.
% \mar{ME: Not sure if this is the correct reference.
% 	I got it from a secondary source.}
%in the theory of classical Banach algebras.
(See~\cite[Prop.\ 3, \S3.7.3]{MR746961}  
for the analogous result in the context of non-archimedean Banach algebras.)

\begin{aprop}
	\label{prop:module topologies}
	Let $A$ be a Polish {\em (}or, equivalently,
	a \emph{(}Hausdorff\emph{)} complete and second countable{\em )} topological ring 
	which is Noetherian {\em (}as an abstract ring{\em )},
	and which contains an open additive subgroup that
	is closed under multiplication, and consists of topologically
	nilpotent elements.
	Then any finitely generated $A$-module $M$ has a unique
	completely metrizable topology with respect to which it becomes
	a topological $A$-module.  Furthermore, $M$ is complete
	with respect to this topology, any submodule of $M$
	is closed, and any morphism of finitely generated
	$A$-modules is automatically continuous, has closed image,
	and induces an open mapping from its domain onto its image.
\end{aprop}
\begin{proof}
  To begin with, suppose that $M$ is a finitely generated $A$-module
  which is furthermore endowed with a completely metrizable topology
  which makes it a topological $A$-module.  We claim first that $M$ is
  in fact complete, and that any surjection of $A$-modules $A^n \to M$
  (for some $n \geq 1$) is continuous and open. It then follows that
  $M$ is endowed with the quotient topology via this map, and thus the
  completely metrizable topology on $M$ is uniquely determined (if it
  exists).

	To see the claim, note that since $M$ is a finitely generated $A$-module,
	we may choose a surjective homomorphism of $A$-modules
	$A^n \to M$ for some $n~\geq~1$.  Furthermore,
	since $M$ is a topological $A$-module, any such surjection
	is continuous.  Since $A$ is separable as a topological
	space by assumption, we see that $M$ is as well.  (The image
	of any countable dense subset of $A^n$ is a countable dense 
	subset of $M$.)  Thus $M$ is Polish, as is $A^n$,
	and it follows from Lemma~\ref{lem:complete} that
	$M$ is complete (as a topological module), while
        it follows from Corollary~\ref{cor:open polish}	that
	the given surjection $A^n \to M$ is open, as claimed.
	

	We next claim that any $A$-submodule of $M$ is closed.
	To see this, let $N$ be a submodule of $M$,
	and let $\overline{N}$ denote its closure.
	Since $A$ is Noetherian by assumption, so is its finitely 
	generated module $M$, and thus $\overline{N}$ is finitely 
	generated, say by the elements $x_1,\ldots,x_n$.
	Applying the results proved above for $M$ to
	its closed (and hence completely metrizable topological) submodule
	$\overline{N}$, we find that
	the surjection $A^n \to \overline{N}$ given by
	$(a_1,\ldots,a_n) \mapsto \sum_{i=1}^n a_i x_i$ 
	is an open mapping.
        Consequently,
	if we let $I$ be an open additive subgroup of $A$ which
	satisfies $I^2 \subseteq I,$ and which consists of topologically
	nilpotent elements (such an $I$ exists by assumption),
	then $I x_1 + \cdots + I x_n $ is an open subset of $\overline{N}$.
	Since $N$ is dense in $\overline{N}$,
	we conclude that
	$N + I x_1 + \cdots + I x_n  = \overline{N},$
        and consequently we may write
        $$x_i = \sum_{i=1}^n a_{ij} x_j + y_i$$ for each $1 \leq i \leq n$,
	for some $y_i \in N$ and $a_{ij} \in I$. 
	Rearranging, we find that
	$$ \begin{pmatrix} 1 - a_{11} & - a_{1 2} & \cdots & - a_{1 n} \\
	       -a_{21} & 1 - a_{22} & \cdots & - a_{2 n} \\
	       \vdots & \vdots & \ddots & \vdots \\
-a_{n1} & - a_{2n} & \cdots & 1 - a_{nn} \end{pmatrix}
\begin{pmatrix} x_1 \\ x_2 \\ \vdots \\ x_n\end{pmatrix}
= 
\begin{pmatrix} y_1 \\ y_2 \\ \vdots \\ y_n\end{pmatrix}
$$	       
The determinant of the matrix 
	$ \begin{pmatrix} 1 - a_{11} & - a_{1 2} & \cdots & - a_{1 n} \\
	       -a_{21} & 1 - a_{22} & \cdots & - a_{2 n} \\
	       \vdots & \vdots & \ddots & \vdots \\
-a_{n1} & - a_{2n} & \cdots & 1 - a_{nn} \end{pmatrix}$
lies in $1 + I$, and thus is a unit (since $A$ is complete
and $I$ consists of topologically nilpotent elements).
Thus we find that $x_1,\ldots,x_n$ lies in the $A$-span
of $y_1,\ldots,y_n$, so that in fact $\overline{N} \subseteq N$.
Thus $N$ is indeed closed, as claimed.

Applying the preceding result to $A^n$, we find that
any $A$-submodule of $A^n$ is closed, and hence that any
finitely generated $A$-module $M$ is isomorphic to a quotient
$A^n/N$, where $N$ is a closed submodule of $A^n$.  The quotient
topology on $A^n/N$ makes it into a complete and metrizable
topological $A$-module, and thus $M$ does indeed admit a complete
and metrizable topological $A$-module structure.   We have already
seen that this topology is unique, and that any $A$-submodule of $M$
is closed.    Finally, we see that
homomorphisms between finitely generated $A$-modules are necessarily
continuous and that their images
(being $A$-submodules of their targets)
are necessarily closed, and a final application
of Corollary~\ref{cor:open polish} shows that they induce open
mappings onto their images.
\end{proof}

\chapter{Tate modules and continuity}\label{app: Tate modules and continuity}% \mar{TG: sort out appendix numbering}\label{sec: Tate modules and
% continuity}
In this appendix we study continuity
conditions for group actions on modules over Laurent series rings. We
begin with some results on the topology of such modules, before
introducing group actions and related notions.
%\section{Continuity conditions}
% We use the notation of Section~\ref{subsec: EG stuff}, although we make
% no use of the action of~$\varphi$.
% \mar{ME: I don't think this reference to \S~\ref{subsec: EG stuff} is
% 	correct anymore.  Probably it should be \S~\ref{subsec: moduli
% 		stacks of phi modules}, but there is so little
% 	notation needed, that we can probably just recall it directly.}


\section{Topologies and lattices}
We fix a finite extension
$E/\Qp$ with ring of integers~$\cO$, uniformizer~$\varpi$, and residue
field~$\F$, and we also fix a finite extension~$k/\F_p$. If $A$ is a $p$-adically complete~$\cO$-algebra, we write
$\Aplus_A:=(W(k)\otimes_{\Zp}A)[[T]]$, and we let $\AAA_A$ be the $p$-adic completion of~
$\Aplus_A[1/T]$.

\begin{arem}
  \label{rem:topologies on M}%By~\cite[Ex.\ 3.2.2]{MR2181808},
  Any
  finitely generated projective $\AAA_A$-module $M$ has a natural topology.
  Indeed, we may write $M$ as a direct summand of $\AAA_A^n$ for some
  $n \geq~1$.  We then endow this latter module with its product topology,
  and endow $M$ with the subspace topology. 
  More intrinsically,
  since $A$ is $p$-adically complete (by assumption),
  we may write $M = \varprojlim_a M/p^a M$.   Each of the quotients
  $M/p^a M$ is then a projective $(A/p^a A)((T))$-module,
  and so has natural topology, making it a Tate $A/p^a A$-module in
  the sense \index{Tate module}
  of~\cite{MR2181808}.
  The topology on $M$ is then the projective limit of these Tate module
  topologies.
  % . In particular, an \emph{fpqc}-locally free
  % \'etale $\varphi$-module~$M$ has a natural topology
  We note that multiplication by~$T$ is topologically nilpotent
  on~$M$.

  Similarly, any Zariski locally finite free $\Aplus_A$-module $\gM$ has a natural topology.
  We may describe this topology in an analogous manner to the
  case considered in the preceding paragraph.
  Namely, such a module
  is a finitely generated and projective $\Aplus_A$-module,
  and thus is a direct summand of $(\Aplus_A)^n$ for some $n \geq 0$.
  If we endow this latter module with its product topology
  then the natural topology on $\gM$ is its corresponding subspace topology. 
  In this case, though, this topology admits a more intrinsic and succinct
  description: it is the $(p,T)$-adic topology on $\gM$.
  % \mar{ME: Maybe add citation to \S 2 of {\tt etalephimodules.tex} to elaborate
  %         on what the natural topology actually is. ME: Alternatively,
  % could just explain directly, which I've now done.}
\end{arem}

We remind the reader that a topological group $G$
is said to be {\em Polish} if its underlying
topological space is Polish, i.e.\ is separable and completely metrizable.
As explained in
%Appendix~\ref{app:topological groups},
Remark~\ref{rem:Polish},
this is equivalent
to $G$ being complete (as a topological group) and second countable
(as a topological space).   



\begin{alemma}
	\label{lem:polish}
	If $A$ is a $p$-adically complete $\cO$-algebra
	for which $A/p$ is countable, then $\AAA_A$ is
	Polish,
	and consequently any finitely generated projective
	$\AAA_A$-module is Polish, when endowed with its canonical topology.
\end{alemma}
\begin{proof}
	%\mar{ME: Write something here}
	We will use the fact that a countable 
	product of Polish spaces is Polish,
	as is a closed subspace of a Polish space.
	It follows from these facts,
	and from the description of the
	natural topology on a finitely generated projective
	$\AAA_A$-module given in Remark~\ref{rem:topologies on M},
	that the canonical topology on any such module is Polish,
	provided that that $\AAA_A$ itself is Polish.

	We may write $\AAA_A = \varprojlim_a \AAA_{A/p^aA}$,
	so that $\AAA_A$ is a closed subset of the (countable)
	product of the various spaces $\AAA_{A/p^aA}$.
	It suffices, then, to show that each of these 
	spaces is Polish; in other words,
	we reduce to the case when $A$ is a countable $\Z/p^a$-algebra
	for some $a \geq 1$.

	Since $A$ is countable, so is each of the quotients
	$\Aplus_A/T^n$.  Choose a subset $X_n \subseteq
	\Aplus_A$ which maps bijectively onto $\Aplus_A/T^n$,
	and write $X := \bigcup_{m,n} T^{-m} X_n.$  
	Then $X$ is a countable dense subset of $\AAA_A$,
	and thus $\AAA_A$ is separable.  
	
	The topology on $\Aplus_A$ is the $T$-adic topology,
	and thus is metrizable (as is any $I$-adic topology
	on a ring).  Since $\Aplus_A$ is $T$-adically complete,
	it is in fact completely metrizable.  The same is then
	evidently true for
	\[\AAA_A  := \Aplus_A[1/T] = \bigcup_m T^{-m} \Aplus_A.\qedhere\]
\end{proof}

\begin{arem}
	\label{rem:Polish comparision}
	Clearly $T\Aplus_A$ is an open subgroup of $\AAA_A$
that is closed under multiplication
	and consists of topologically nilpotent elements.
	Thus if $A/p$ is countable, then $\AAA_A$ satisfies
	the conditions of Proposition~\ref{prop:module topologies},
and thus the canonical topology on finitely generated projective $\AAA_A$-modules
constructed in Remark~\ref{rem:topologies on M} is a particular
case of the canonical topology constructed on any finitely generated
$\AAA_A$-module in Proposition~\ref{prop:module topologies}.
\end{arem}



% \begin{rem}
%   \label{rem: p-adic definition here but mod p^n in proper.tex}Note
%   that if~$A$ is a $\Z/p^a\Z$-algebra, then $\AAA_A=\Aplus_A[1/u]$, so that
%   Definition~\ref{defn: etale phi module} coincides with the
%   definition in~\cite[\S5]{EGstacktheoreticimages} (in the special
%   case that $q=p$ and~$\varphi(u)=u^p$, which is the only case
%   considered there).
% \end{rem}

% \begin{lem}\mar{Need a Noetherian hypothesis on~$A$ to be
%     discussing~$\Spf$, or in the actual proof (if that exists)?}
%   \label{lem: can check freeness of Kisin and etale phi modules modulo
%     p^n}A finitely generated projective $\Aplus_A$-module~$\gM$ is
%   Zariski locally free of rank~$d$ if and only if for each~$n\ge 1$,
%   $\gM/p^n\gM$ is Zariski locally free of rank~$d$ as a
%   $\Aplus_{A/p^nA}$-module. Similarly, a finitely generated projective
%   $\AAA_A$-module~$M$ is \'etale \emph{(}resp.\ \emph{fppf}, resp.\
%   \emph{fpqc}\emph{)} locally free of rank~$d$ if and only if for
%   each~$n\ge 1$, $M/p^nM$ is \emph{(}resp.\ \emph{fppf}, resp.\ \emph{fpqc}\emph{)}
%   locally free of rank~$d$ as a $\cO_{\cE,A/p^nA}$-module.\mar{TG:
%     actually equivalent to just mod~$p$?}
% \end{lem}
% \begin{proof}For both parts the ``only if'' direction is clear. We
%   begin with the ``if'' direction for~$\gM$.\mar{TG: need to finish
%     this, assuming the statements for~$M$ are correct.} Since $\gM$ is
%   finite generated and projective, it is Zariski locally free as an
%   $A[[u]]$-module by~\cite[Prop.\
%   5.1.5]{EGstacktheoreticimages}. \mar{TG: I think it's easy to finish
%   the $[[u]]$ case but I'm not convinced we're doing the right thing
%   in the $((u))$ world.}

% Want to argue: locally free mod~$p$, lift generators, get by
% topological Nakayama \[0\to K\to A((u))^d\to M\to 0,\] and~$K$ is
% finitely generated projective and vanishes mod~$p$, so~$K=0$.\mar{TG:
%   scribbled note says ``write as quotient of free and lift explicitly''.}

% In~$\gM$ case, $\gM/u\gM$ is enough
% by~\cite[5.1.4]{EGstacktheoreticimages}. 
%   \end{proof}


We now present some additional facts related to the preceding concepts
which will be needed in the sequel.

%\mar{ME: Check this lemma isn't already in {\tt etalephimodules.tex}}
\begin{alem}
  \label{lem: can check projectivity of Kisin and etale phi modules modulo
    p^n}A finitely generated $\Aplus_A$-module is
  projective of rank~$d$ if and only if for each~$a\ge~1$,
  the quotient
  $\gM/p^a\gM$ is projective of rank~$d$ as an
  $\Aplus_{A/p^aA}$-module. Similarly, a finitely generated 
  $\AAA_A$-module~$M$ is projective of rank~$d$ if and only if for
  each~$a\ge~1$,
  the quotient
  $M/p^aM$ is projective of rank~$d$ as an $\AAA_{A/p^aA}$-module.
\end{alem}
\begin{proof}
  This is immediate from~\cite[Prop.\
0.7.2.10(ii)]{MR3075000}, applied to the $p$-adically complete
  rings~$\Aplus_A$ and $\AAA_A$.%\mar{TG: this should be checked!}
\end{proof}


   \begin{adefn}% \mar{TG: check that this definition coincides with
    % Drinfeld's definition of a lattice (to the extent that makes sense
    % - his may not require $\gS_A$-stability?)}
    \label{defn: lattice}If~$M$ is a finitely generated $\AAA_A$-module, then a
    \emph{lattice} in~$M$ is a finitely generated $\Aplus_A$-submodule
    $\gM\subseteq M$  whose $\AAA_A$-span \index{Tate module}
    is~$M$.% \mar{TG: didn't think much about this definition! Probably
      % torsion free is automatic, in fact\dots}
  \end{adefn}

  Note that any finitely generated $\AAA_A$-module contains a lattice.
  (If $f:\AAA_A^r \to M$ is a surjection from a finitely generated 
  free $\AAA_A$-module onto the finitely generated $\AAA_A$-module $M$,
  then the image of the restriction of $f$ to $(\Aplus_A)^r$ 
  is a lattice in~$M$.)

  We record some additional lemmas and a remark which apply in the case
  when $A$ is an $\cO/\varpi^a$-algebra for some $a \geq 1$.

  \begin{alemma}
	  \label{lem:lattice properties}
	  Let $A$ be an $\cO/\varpi^a$-module for some $a \geq 1$,
	  and let $M$ be a finitely generated $\AAA_A$-module.
	  
	  \begin{enumerate}
		  \item
	  If $\gM$ and $\gN$ are two lattices
	  contained in a finitely generated $\AAA_A$-module~$M$,
	  then there exists $n \geq 0$
	  such that $T^n\gM \subseteq \gN \subseteq T^{-n} \gM$.
  \item   If in addition either $A$ is Noetherian or $M$ is projective,
	  then any lattice in $M$ is $T$-adically complete.
  \item 
	  If $A$ is Noetherian, if $\gM$ is
	  a lattice in~$M$, and if $\gN$ is an $\Aplus_A$-submodule
	  of $M$ 
	  such that $T^n\gM \subseteq \gN \subseteq T^{-n} \gM$
	  for some $n \geq 0$, 
	  then $\gN$ is a lattice in~$M$.
  \end{enumerate}
  \end{alemma}
  \begin{proof}
	  To prove~(1), it suffices to prove one of the inclusions;
	  the reverse inclusion may then be obtained (possibly after
	  increasing $n$) by switching the roles of $\gM$ and $\gN$.
          Since $\AAA_A = \Aplus_A[1/T]$, we find 
	  that $\gN \subseteq \gM[1/T] = \bigcup_{n = 0}^{\infty} T^{-n}\gM.$
	  Since $\gN$ is finitely generated as an $\A_{A}^+$-module, 
	  we obtain that $\gN \subseteq T^{-n}\gM$ for some sufficiently
	  large value of $n$, as required. 

	  To prove~(2), we note that if $A$ is Noetherian, 
	  than $\Aplus_A$ is also Noetherian, and thus any finitely generated
	  $\Aplus_A$-module is $T$-adically complete, since $\A^+_A$ itself is.
	  If $M$ is projective, then we write $M$ as a direct summand 
	  of a finitely generated free $\AAA_A$-module,
	  so that $\gM$ may then be embedded
	  into a finitely generated free $\Aplus_A$-module.   Thus $\gM$
	  is $T$-adically separated, and hence the kernel of any
	  surjection $(\Aplus_A)^r \to \gM$ is $T$-adically closed.  Since
	  $(\Aplus_A)^r$ is $T$-adically complete, we conclude that
	  the same is true of $\gM$.

	  Suppose now that $A$ is Noetherian, and that we are in
	  the situation of~(3).  The inclusion $T^n \gM \subseteq \gN$
	  shows that $M = \gM[1/T] \subseteq \gN[1/T],$
	  so that $\gN$ generates $M$ as an $\AAA_A$-module.
	  We must show that $\gN$ is furthermore finitely generated
	  over $\Aplus_A$. For this, we first note that, since 
	  $\gM$ is $T$-adically complete (by (2)),
	  the inclusion $T^n \gM \subseteq \gN \subseteq T^{-n}\gM$
	  shows that $\gN$ is also $T$-adically complete.
	  It also shows that $\gN/T\gN$ is a subquotient
	  of the finitely generated $A$-module $T^{-n}\gM/T^{n+1}\gM$,
	  and thus is a finitely generated $A$-module (as $A$ is Noetherian).
	  Since $\gN$ is $T$-adically complete, an application
	  of the topological Nakayama lemma shows that~$\gN$ is finitely
	  generated over $\Aplus_A$, as required.
          %
	  %Suppose now that $A$ is Noetherian, and that we are in
	  %the situation of~(3).  The inclusion $T^n \gM \subseteq \gN$
	  %show that $M = \gM[1/T] \subseteq \gN[1/T],$
	  %so that $\gN$ generates $M$ as an $\AAA_A$-module.
	  %We must show that $\gN$ is furthermore finitely generated
	  %over $\Aplus_A$.
	  %For this, we note that $\gN/T\gN$ is a subquotient
	  %of the finitely generated $A$-module $T^{-n}\gM/T^{n+1}\gM$,
	  %and thus is a finitely generated $A$-module (as $A$ is Noetherian).
	  %Since $\gM$ is $T$-adically complete (by (2)) an application
	  %of the topological Nakayama lemma shows that $\gM$ is finitely
	  %generated over $\Aplus_A$, as required.\mar{TG: something is
          %wrong here; we want to show that~$\gN$ is finitely
          %generated, not~$\gM$, but as written~(2) doesn't apply to~$\gN$.}
  \end{proof}



    \begin{arem}
    \label{rem: lattice and Drinfeld}By~\cite[Thm.\
    5.1.14]{EGstacktheoreticimages} (a theorem of Drinfeld),
    if $A$ is an $\cO/\varpi^a$-algebra for some $a \geq 1$,
    then we may
    think of a finitely generated projective~$\AAA_A$-module as a Tate $A$-module $M$
    together with a topologically nilpotent automorphism~$T$. In this
    optic a lattice in our sense is precisely
    a lattice in the Tate module~$M$
    which is also an $\Aplus_A$-submodule (by definition,
    a lattice~$L\subseteq M$
    is an open submodule with the property that for every open
    submodule~$U\subseteq L$, the $A$-module $L/U$ is finitely
    generated), 
    as the following lemma shows.
    % \mar{TG: maybe the converse is true, that every
    %   $T$-stable lattice is a lattice? I think we just need to deduce
    %   that it's finitely generated over~$\Aplus_A$, which might follow
    %   from the SGA results cited in \cite[Prop.\
    %   5.1.8]{EGstacktheoreticimages}, but I'm not going to worry about
    % this right now. ME: I think the converse {\em is} true, and
    % have added a proof.}
  \end{arem}

  \begin{alemma}\label{lem:lattice equivalences}% \mar{TG: might be best to have this kind of material
    %   before we've ever put a $\varphi$ on $\Aplus_A$? So first of all
    % do ring theory stuff, then \'etale $\varphi$-modules, then Kisin modules.}
    If $A$ is an $\cO/\varpi^a$-algebra for some $a \geq 1$,
    and if $\gM$ is an $\Aplus_A$-submodule of a finitely generated
    projective $\AAA_A$-module~$M$,
	  then the following conditions on $\gM$ are equivalent:
	  \begin{enumerate}
		  \item $\gM$ is a lattice in $M$	
	  {\em (}in the sense of Definition~{\em \ref{defn: lattice})}
	  \item $\gM$ is open in~$M$,
		  and for every $\Aplus_A$-submodule $U$ of $\gM$
		  which is open in~$M$, the quotient $\gM/U$
		  is finitely generated over~$A$.
	  \item $\gM$ is open in~$M$,
		  and for every $A$-submodule $U$ of $\gM$
		  which is open in~$M$, the quotient $\gM/U$
		  is finitely generated over $A$
		  {\em (}i.e.\ $\gM$ is a lattice in~$M$,
		  when $M$ is thought of as a Tate module
		  over $A$ in the sense of Remark~{\em \ref{rem: lattice
				  and Drinfeld}).}
  \end{enumerate}
  \end{alemma}
  \begin{proof}
	  Suppose that $\gM$ is a lattice in $M$,
	  in the sense of Definition~\ref{defn: lattice}.
	  If we choose a surjection $(\Aplus_A)^r \to \gM$ for some $n \geq 0,$
	  then the induced surjection $\AAA_A^r \to M$ is a continuous
	  and  open map,
	  and so $\gM$ is open in~$M$, since $(\Aplus_A)^r$ is open in~$\AAA_A^r$.
	  Furthermore, since the submodules $T^n (\Aplus_A)^r$ ($n \geq 0$)
	  form a neighbourhood basis of $0$ in~$(\Aplus_A)^r$,
	  we see that their images $T^n\gM$ form a neighbourhood
	  basis of $0$ in~$\gM$.  Thus if $U$ is any open
	  $A$-submodule of~$\gM$, then $\gM/U$ is a quotient
	  of $\gM/T^n\gM$ for some~$n~\geq~0$,
	  and thus is finitely generated over~$A$.
	  Thus~(1) implies~(3).

	  Clearly~(3) implies~(2), and so we suppose that~(2) holds.
	  Let $\gN$ be an $\Aplus_A$-submodule of $M$ that is a lattice
	  in the sense of Definition~\ref{defn: lattice},
		  so that $M = \gN[1/T] = \bigcup_{n \geq 0} T^{-n} \gN$.
		  Then $\gM =
		  \bigcup_{n \geq 0} \gM \cap T^{-n}\gN,$
		  and thus
		 $$ \gM/(\gM \cap \gN) = 
		  \bigcup_{n \geq 0} (\gM \cap T^{-n}\gN) /
		  (\gM\cap  \gN).$$
		  But $\gM/(\gM \cap \gN)$ is finitely generated over $A$
		  by assumption,
		  and thus $$\gM/\gM \cap \gN = (\gM\cap T^{-n}\gN) /
		  (\gM\cap \gN),$$
		  for some sufficiently large value of~$n$,
		  implying that $\gM\subseteq T^{-n} \gN.$
		  In particular $\gM$ is $T$-adically complete,
		  being open (and hence closed) in
		  the lattice~$T^{-n}\gN$,
		  which is $T$-adically complete
		  by Lemma~\ref{lem:lattice properties}~(2).
	  Since $T$ is an automorphism of $M$, 
	  we find that $T\gM$ is an open submodule of $\gM$,
	  so that $\gM/T\gM$ is finitely generated over~$A$.
	  As $\gM$ is $T$-adically complete, 
	  the topological Nakayama lemma implies that $\gM$
	  is finitely generated over $\Aplus_A$.

	  Since $T$ is topologically nilpotent and $\gN$ is finitely
	  generated over $\Aplus_A$, we also find that $T^n\gN \subseteq 
	  \gM$ for some sufficiently large value of $n$.  Thus
	  $\gM[1/T] =\gN[1/T] = M.$  This completes the proof
	  that $\gM$ is a lattice in $M$
	  in the sense of Definition~\ref{defn: lattice},
	  showing that~(2) implies~(1).
  \end{proof}
  
  \begin{arem}
	  \label{rem:lattice cofinality}
	  It follows from Lemmas~\ref{lem:lattice properties}
	  and~\ref{lem:lattice equivalences}
	  that if $A$ is a~$\cO/\varpi^a$-algebra,
	  and $M$ is a finitely generated projective~$\AAA_A$-module,
	  then the lattices in $M$ form a neighbourhood basis
	  of the origin in~$M$.
  \end{arem}

  We also note the following technical lemma.

  \begin{alemma}
	  \label{lem:intersecting with lattice}
	  If $A\subseteq B$ is an inclusion of $\cO/\varpi^a$-algebras
	  for some $a \geq 1$,
	  with $A$ Noetherian,
	  if $M$ is a projective $\AAA_A$-module~$M$
	  with extension of scalars $M_B:=\A_B\otimes_{\A_A}M$ to $B$,
	  and if $\gM_B$ is a lattice in $M_B$,
	  then $\gM:=M \cap \gM_B$
	  {\em (}the intersection takes place in $M_B${\em )}
	  is a lattice in $M$,
	  having the additional property that $M \cap T^n \gM_B = T^n\gM$
	  for any integer~$n$.
  \end{alemma}
  \begin{proof}
	  Choose a projective complement to $M$,
	  i.e.\ a finitely generated projective $\AAA_A$-module
	  $N$ such that $M\oplus N$ is free over $\AAA_A$,
	  and let $\gN_B$ be a lattice in~$N_B$.
	  Then it suffices to prove the lemma with
	  $M$ replaced by $M\oplus N$ and with $\gM_B$
	  replaced by $\gM_B \oplus \gN_B$;
	  thus we may and do
	  suppose that $M$ is free over $\AAA_A$.
	  Choose a lattice $\gM'$ in $M$ which
	  is free over $\Aplus_A$,
	  and note that $\gM'_B := \Aplus_B\gM'$ is a lattice in $M_B$,
	  and that $M \cap T^n\gM'_B  = T^n\gM'$ for any integer $n$.

	  Lemma~\ref{lem:lattice properties}~(1) shows that
	  $T^n \gM'_B \subseteq \gM_B \subseteq T^{-n} \gM'_B$
	  for some sufficiently large value of $n$.
	  Intersecting with $M$, we find that
	  $T^n \gM' \subseteq \gM \subseteq T^{-n} \gM',$
	  so that $\gM$ is a lattice in $M$,
	  by Lemma~\ref{lem:lattice properties}~(3).
	  Finally,
	  by the definition of $\gM$,
	  we find that $M\cap T^n\gM_B = T^n\gM$ for any integer $n$.
  \end{proof}

% \mar{ME: Maybe put a short segue here, since we've completely
% changed topic. TG: how about this?}

  %\begin{alem}\mar{needs fixing! ME: Maybe just fold into the place
%	  where it's used. TG: needs to be done.}
%    \label{lem: extending Z to Zp}Let $G$ be a Hausdorff topological group
%    acting on a complete topological abelian group~$M$, and let
%    $G\times M\to M$ be a continuous map, making~$M$ a topological
%    $G$-module. Then the action of~$G$ on~$M$ extends uniquely to a
%    continuous action $\widehat{G}\times M\to M$, making~$M$ a
%    topological $\widehat{G}$-module.
%  \end{alem}
%  \begin{proof}
%    It suffices to show that $G\times M\to M$  is uniformly continuous.
%  \end{proof}

\section{Group actions}
We now prove some lemmas which allow us to check whether the action of
a topological group on a Tate module is continuous, and to extend such
an action from a group to its completion.
\begin{alemma}\label{lem: continuity and lattices}  Let $G$ be a
  topological group acting on a Tate module $M$, and assume that~$G$
  admits a neighbourhood basis of the identity consisting of open
  subgroups. Then the following are equivalent:
	\begin{enumerate}
		\item The action $G\times M \to M$ is continuous.
		\item
                  \begin{enumerate}
                  \item For each~$m\in M$, the map $G\to M$, $g\mapsto
                    gm$ is continuous at the identity of~$G$,
                  \item for each~$g\in G$, the map $M\to M$, $m\mapsto
                    gm$ is continuous, and
		   % \mar{ME: Is this
		%automatically true, or do we have to assume it? TG:
                %not sure - e.g. could there be a weird discontinuous
                %involution on~$M$? I would have guessed that this
                %could happen but I have no intuition for this. ME: I think
	        %you're right, so I'll remove this.}
                  \item for any lattice $L$ in $M$,
	 there is an open subgroup $H$ of $G$ that preserves~$L$.
 \end{enumerate}
                  \end{enumerate}
Furthermore, if these equivalent conditions hold, then the following stronger 
form of condition~{\em (2)(a)} holds:
%\mar{ME: Have to fix label; I've referred to it as (2)(a') in the proof below.}
\begin{enumerate}
\item[(2)(a')]
For each~$m \in M$, the map $G\to M$, $g \mapsto gm$ is equicontinuous
at the identity of~$G$.
\end{enumerate}
  % \mar{TG: should one add a third condition, which is
  % that for any two lattices~$L,L'$, there is an~$H$ with $HL\subseteq
  % L'$? That seems to be the condition we are using in
  % Lemma~\ref{lem:testing continuity on M mod T}. TG: this was
  % nonsense, commenting out.}
\end{alemma}
\begin{arem}
  \label{rem: neighbourhood basis profinite}The assumption in Lemma~\ref{lem: continuity and lattices} that~$G$
  admits a neighbourhood basis of open subgroups holds in particular
  if~$G$ is a profinite group, or if~$G$ is a subgroup of a profinite
  group with the subspace topology. In particular, it holds for the
  groups~$\Zp$ and for $\Z\subset\Z_p$.
\end{arem}
\begin{proof}[Proof of Lemma~\ref{lem: continuity and lattices}]
	Suppose first that~(1) holds; then clearly conditions~(2)(a)
	and~(2)(b) hold.  
	Let $L$ be a lattice in $M$.
	Since the action morphism $G \times M \to M$ is continuous,
	we may find an open subgroup $H$ of $G$ and an open submodule
	$U$ of~$M$, such that $H  U \subseteq L.$  Replacing $G$ by~$H$,
	we may thus suppose that $G U \subseteq L.$
	Now consider the morphism
	$G \times L/U \to M/L$  induced by the action morphism.
	Since $L$ is a lattice, the $A$-module $L/U$ is finitely generated.
	Let $\{m_i\}$ be a finite generating set.
	Since $M/L$ is discrete, for each generator~$m_i$,
	there is an open subgroup $H_i$ of $G$ such that $H_i m_i$ is
	constant, and thus equal to zero, modulo $L$.
        If we write $H := \bigcap_i H_i,$ then $H$ is an open subgroup
	of $G$ such $H m_i \subseteq L$ for every~$i$.  Consequently,
	$H L \subseteq L,$ verifying that~(2)(c) holds as well.

	Suppose conversely that~(2) holds. 
	We first note that we may strengthen~(2)(a) as follows:
	if $m \in M$, then the orbit map $g \mapsto gm$ is a continuous
	map $G \to M$.  Indeed, if $g_0 \in G$, then we may
	write this map as the composite of the continuous automorphism
	$g \mapsto g g_0^{-1}$ of $G$ and the orbit map $g \mapsto 
	g g_0 m.$  The latter orbit map is continuous at the identity
	of~$G$, by~(2)(a), and so the orbit map of $m$ is continuous
	at~$g_0$.

	We now wish to prove
	that the action map $G\times M \to M$ is continuous.
	Let $(g,m) \in G\times M.$   
	Any neighbourhood of the image $g m \in M$ contains
	a neighbourhood of the form $g m + L,$ where $L$ is a
	lattice in $M$.
	Given a lattice $L$, then, 
	we must find a neighbourhood of $(g,m)$ whose image
	lies in $g m + L.$

	% Any neighbourhood of the identity
	% in the product $G\times M$ contains a neighbourhood of the form
	% $H\times L'$, where $H$ is an open subgroup of $G$ and
	% $L'$ is a lattice in $M$,
	% %and $H$ is an open
	% %subgroup of $G$ that preserves $L'$ (this is where we apply~(2)),
	% and thus any neighbourhood of $(g,m)$ contains one
	% of the form $g H \times (m +L').$\mar{TG: something weird is
        %   happening here - $L'$ seems to be fixed in this sentence and
        % then varying in the next sentence.}
	Since $g$ is a continuous automorphism of $M$, by~(2)(a),
	we may find a lattice $L'$ such that $g L' \subseteq L.$
	By~(2)(c), we may find an open subgroup $H'\subseteq G$ that
	preserves~$L'.$  Then $H'$ acts on~$M/L'$, and since this quotient
	is discrete, and since the orbit maps for the action are continuous
	(by what we proved above),
        it follows that the action of $H'$ on $M/L'$ is smooth, % \mar{TG: I
        % don't understand this implication!}
	and so we may find an open subgroup $H$ of $H'$ that
	fixes the image of $m$ in $M/L'$.   Then we find that
	$(gH)(m+L') \subseteq gm + L,$ and since~$gH\times (m+L')$ is an
        open neighbourhood of $(g,m)\in G\times M$, %we are done.% as required.
        we have proved the required continuity.

Finally, suppose that conditions~(1) and~(2) hold; we must show that
the continuity condition of~(2)(a) 
can be upgraded  to the equicontinuity condition of~(2)(a'). 
To do this, we need to show that for any two
 lattices $L,L'\subseteq M$, there is an open subgroup~$H$ of~$G$ such
 that for each $m\in L$, we have $Hm\subseteq m+L'$. 

 Replacing~$L'$ by~$L\cap L'$, we can assume that~$L'\subseteq
 L$. By~(2)(c), we can choose~$H$ such that $H L'\subseteq L'$ and
 $H L\subseteq L$, so that we have an induced continuous map
 $H\times (L/L')\to L/L'$. Since~$L/L'$ is a finitely generated
 $A$-module and is discrete, after replacing~$H$ by an open subgroup we can assume
 that $H$ acts trivially on~$L/L'$, as required.
\end{proof}

%\mar{ME: Have to modify the statement/proof, in light of the improvement to D.12}
\begin{alem}\label{lem: extending
  Z to Zp}
	Let $G$ be a Hausdorff topological group, 
  which admits a neighbourhood basis of the identity consisting of open
  subgroups, and suppose furthermore that for any open subgroup $H$ of~$G$,
		  the quotient $G/H$ is finite; equivalently,
		  suppose that the completion $\widehat{G}$ of $G$
		  is profinite.

		  If $M$ is a Tate module,
		  endowed with a continuous action $G \times M \to M$,
		  then % the following conditions are equivalent:
	% \begin{enumerate}
	% 	\item The
                  the action $G\times M \to M$ extends
		       to a continuous action $\widehat{G} \times M
		       \to M$.
	 %        \item
         %          \begin{enumerate}
         %          \item For each lattice $L \subseteq M$,
	 %        	  the maps $G \to M$ 
	 %        	  defined by $g \mapsto gm$,
	 %        	  for $m~\in~L$, are equicontinuous at the
	 %        	  identity;
         %          \item for each~$g\in G$, the map $M\to M$, $m\mapsto
         %            gm$ is continuous; and
	 %           % \mar{ME: Is this
	 %        %automatically true, or do we have to assume it? TG:
         %        %not sure - e.g. could there be a weird discontinuous
         %        %involution on~$M$? I would have guessed that this
         %        %could happen but I have no intuition for this. ME: I think
	 %        %you're right, so I'll remove this.}
         %          \item for any lattice $L$ in $M$,
	 % there is an open subgroup $H$ of $G$ that preserves~$L$.
 % \end{enumerate}
 % \end{enumerate}
 \end{alem}
 \begin{proof}
	 % \mar{ME: We don't need~(1) implies~(2); just included for 
	 %         aesthetic reasons.  Add later.}
 %         Suppose that~(2) holds.  
 %         It follows from
 % Lemma~\ref{lem: continuity and lattices} that
We have to show that  the continuous action
 $G\times M \to M$  extends
 (necessarily uniquely, since $G$ is dense in $\widehat{G}$)
 to a continuous map $\widehat{G}\times M \to M$.  (The fact
 that this map will induce a $\widehat{G}$-action on $M$ follows
 from the corresponding fact for $G$, and the density of $G$ in $\widehat{G}$.)
 Since $M$ is the union of its lattices, it suffices to show
 that the induced map $G\times L \to M$ extends to a continuous
 map $\widehat{G} \times L \to M$, for each lattice $L~\subseteq~M$.
 For this, it suffices to show that the map $G\times L \to M$
 is uniformly continuous, for each lattice~$L$.
 That is, for any lattice $L' \subseteq M$, we have to find
 an open subgroup $H \subseteq G$ and a sublattice $L'' \subseteq L$
 such that $gH m + H L'' \subseteq gm + L'$ for all $g \in G$ and all
 $m \in L$.  

By Lemma~\ref{lem: continuity and lattices}, the conditions (2)(a)-(c)
of that lemma hold. We begin by showing that we may find a sublattice $L''$ of $L$
 such $G L'' \subseteq L'$.  For this, we first note that, by~(2)(c),
 we may find an open subgroup $H \subseteq G$ such
 that $H L' = L'$.  If let $\{g_i\}$ denote a (finite!) set
 of coset representatives for $H\backslash G,$ then
 since each $g_i$ induces a continuous automorphism of $M$,
 we may find a lattice $L_i$ such that $g_i L_i \subseteq L'$.  
 Taking into account Remark~\ref{rem:lattice cofinality},
 we may then find a lattice $L'' \subseteq L' \cap \bigcap_i L_i$,
 and by construction $G L'' \subseteq L'$.

 Condition~(2)(a') allows us to choose an open subgroup $H \subseteq G$
 such that $H m \subseteq m + L''$ for all $m \in L$.
 We then find that
 $$gH m + H L'' \subseteq gm + G L'' \subseteq gm + L'$$
 for all $g \in G$, as required.  \end{proof}

%  Conversely, suppose that~(1) holds. Replacing~$G$ by~$\widehat{G}$,
%  we may assume that~$G$ is profinite. By Lemma~\ref{lem:
%    continuity and lattices}, conditions~(2)(b) and~(2)(c) hold, while
%  for each~$m\in M$, the map $g\mapsto gm$ is continuous at the
%  identity of $G$. We need to upgrade this continuity to
%  equicontinuity. To do this, we need to show that for any two
%  lattices $L,L'\subseteq M$, there is an open subgroup~$H$ of~$G$ such
%  that for each $m\in L$, we have $Hm\subseteq m+L'$. 

%  Replacing~$L'$ by~$L\cap L'$, we can assume that~$L'\subseteq
%  L$. By~(2)(c), we can choose~$H$ such that $H L'\subseteq L'$ and
%  $H L\subseteq L$, so that we have an induced continuous map
%  $H\times (L/L')\to L/L'$. Since~$L/L'$ is a finitely generated
%  $A$-module, after replacing~$H$ by an open subgroup we can assume
%  that $H$ acts trivially on~$L/L'$, as required.

% \mar{TG: didn't think about this yet. TG: this is an argument but I
%   didn't seem to actually use the compactness of~$\widehat{G}$??}
%  \end{proof}

%\mar{ME: Another segue? TG: how about this?}

\section{$T$-quasi-linear endomorphisms}
In order to apply the preceding results to the case of interest to us
(the semilinear action of~$\Gamma$ on $(\varphi,\Gamma)$-modules), we
now introduce and study the notion of a $T$-quasi-linear endomorphism
of a finite projective $\A_A$-module. The relevance of this notion to
$(\varphi,\Gamma)$-modules is explained in
Lemma~\ref{lem:gamma minus 1 T quasi linear} below.

%In order to do
%this, it is convenient to be able to take the lattice $\gM$ 
%to be projective, so that it has good base-change properties,  
%and so we introduce some notions that will allow us to reduce
%to this case.
% We now introduce a notion that will allow us to work effectively
% with the criteria of Lemma~\ref{lem:testing continuity on M mod T}.

% \mar{ME: Still not sure about optimal placement of this
% 	material.  Maybe it should be merged with the preceding
% 	discussion of continuity of $\Gamma$-actions? ME: I made some version
% of this merge, but have to finish working out the precise presentation 
% of this material, and add necessary proofs.}
\begin{adefn}
If $M$ is a finite projective $\A_A$-module, % \mar{TG: make this more general than $\varphi$-module; possibly
% make later material more general than $\Gamma$-module.}
then a {\em $T$-quasi-linear} endomorphism of $M$
is a morphism $f: M \to M$ 
which is $W(k)\otimes_{\Z_p} A$-linear, % and $\varphi$-linear, % \mar{TG:
% what does $\varphi$-linear mean? Does it just mean that it commutes
% with $\varphi$? Do we ever actually need this commutation? I haven't
% found an argument that uses it yet.}
% \mar{ME: Could work in a more general setting --- $\Z/p^a$-algebras, or even
% 	$p$-adically complete algebras --- but I think the $\F_p$-case
% 	suffices.}
and which furthermore satisfies the following ($T$-quasi-linearity)
condition:
there exist power series $a(T) \in (\A_A^+)^{\times}$ and $b(T) \in (p,T)\A_A^+$
such that % \mar{TG: do you definitely want these rings rather
  % than~$\gS_A^\times$, $\gS_A$? TG: changed to these, commenting out.}
$$f(T m) = a(T) T f(m) + b(T)T m$$
for every $m \in M$.
\end{adefn}




%We will only apply the notion of $T$-quasi-linearity to endomorphisms of projective
%\'etale $\varphi$-modules.   As noted in Remark~\ref{rem:topologies on M},
%any such \'etale $\varphi$-module $M$ has a natural topology,
%and so it makes sense to speak of an endomorphism of $M$ being
%continuous, or topologically nilpotent. 

\begin{alem}
  \label{lem: powers of T quasilinear}If $f$ is $T$-quasi-linear, then
  for all~$n\in\Z$, we may
  write \[f(T^nm)=a(T)^nT^nf(m)+b_n(T)T^{n}m\] for all~$m\in M$, where $a(T) \in (\A_A^+)^{\times}$ and $b_n(T) \in (p,T)\A^+_A$.
\end{alem}
\begin{proof}Note that the case~$n=0$ is trivial (taking $b_0(T)=0$),
  while if the claim holds for some~$n\ge 1$, then we may write
  \[f(T^{-n}m)=a(T)^{-n}T^{-n}f(m)-a(T)^{-n}b_n(T)T^{-n}m. \]
It therefore suffices to prove the result for ~$n\ge 1$.  This may be proved by induction on~$n$, the case $n=1$ being the
  definition of $T$-quasi-linearity. Indeed, if the claim holds
  for~$n$, then we have
  \begin{align*}
    f(T^{n+1}m)=f(T^n(Tm))&=a(T)^nT^nf(Tm)+b_n(T)T^{n}(Tm)\\
                          &=a(T)^nT^n(a(T) T f(m) + b(T)T
                            m)+b_n(T)T^{n+1}m\\  &=a(T)^{n+1}T^{n+1} f(m) +(b_n(T)+b(T)a(T)^n)T^{n+1}m,
  \end{align*}as required.\end{proof}
% \begin{cor}\mar{TG: not quite sure now what I'm adding this for! not
%     adding proof for now}
%   \label{cor: quasilinear on scaled lattices}If $f$ is
%   $T$-quasi-linear, and $\gM$ is a lattice such that
%   $f(\gM)\subset\gM$, then $f(T^a\gM)\subseteq T^a\gM$ for all~$a\in\Z$.
% \end{cor}
% \begin{proof}
%\end{proof}
% We will only apply the notion of $T$-quasi-linearity to endomorphisms of projective
% \'etale $\varphi$-modules.  



\begin{alem}\label{lem: how f behaves on lattices} % \mar{ME: Have to add lemma about how $f$ behaves
  % on lattices}
Let~$A$ be an $\cO/\varpi^a$-algebra for some~$a\ge 1$, and let $M$ be a finite projective $\A_A$-module.  Let $f$ be a $T$-quasi-linear endomorphism
  of~$M$, and let~$\gM$ be a lattice in~$M$. Then  there is an
  integer~$m\ge 0$ such that for each~$s\in\Z$ and~$n\ge 0$ we have
  $f^n(T^s\gM)\subseteq T^{s-mn}\gM$.
\end{alem}
\begin{proof}
  Choose a finite
  set~$\{m_i\}$ of generators for~$\gM$ as an
  $\A^+_A$-module. Choose~$m$ sufficiently large that we
  have~$T^mf(m_i)\in \gM$ for each~$i$. Then by Lemma~\ref{lem: powers
    of T quasilinear} and the $W(k)\otimes_{\Zp}A$-linearity of~$f$, we
  see that for each~$s\in \Z$ we have $f(T^{s+m}\gM)\subseteq T^s\gM$,
  which gives the result in the case~$n=1$. The general case follows
  by induction on~$n$.
\end{proof}

As noted in Remark~\ref{rem:topologies on M},
any finite projective % \'etale $\varphi$
$\A_A$-module $M$ has a natural topology,
so it makes sense to speak of an endomorphism of $M$ being
continuous, or topologically nilpotent. 
\begin{alemma}\label{lem: quasilinear implies continuous}
	If~$A$ is an $\cO/\varpi^a$-algebra for some~$a\ge 1$, and $M$ is a finite projective $\A_A$-module,
	 then any $T$-quasi-linear endomorphism
	of $M$ is necessarily continuous.
\end{alemma}
\begin{proof}This is immediate from Lemma~\ref{lem: how f behaves on
    lattices}.%\mar{TG: hopefully correct?!}
\end{proof}

\begin{alemma}% \mar{TG: the statement and proof of this lemma have
    % changed, and it should be checked carefully; the same applies to
    % the material following it.}
	\label{lem:topological nilpotence criteria}
	If~$A$ is an $\cO/\varpi^a$-algebra for some~$a\ge 1$, and $M$ is a finite projective $\A_A$-module, and if $f$ is a
	$T$-quasi-linear endomorphism of $M$,
	then the following are equivalent:
	\begin{enumerate}
		\item $f$ is topologically nilpotent.
		\item There exists a lattice $\gM$ in $M$
			and some $n \geq 1$
			such that $f^n(\gM) \subseteq T \gM.$
                        	\item There exists a lattice $\gM$ in $M$
			and some $n \geq 1$
			such that $f^n(\gM) \subseteq (p,T) \gM.$
                                              \item There exists a lattice $\gM$ in $M$ such
                        that for any $m \geq 1$, there exists~$n_0$
			such that  for any~$s\in\Z$ and any $n\ge
                        n_0$, we have $f^n(T^{s}\gM) \subseteq T^{s+m}
                        \gM$.
  \item For any lattice $\gM$ in $M$ and any $m \geq 1$, there exists~$n_0$
			such that  for any~$s\in\Z$ and any $n\ge
                        n_0$, we have $f^n(T^{s}\gM) \subseteq T^{s+m} \gM$.
		% \item For any lattice $\gM$ in $M$, any~$m\ge 1$, and
                %   any~$s\ge 0$,
		% 	there exists an $n_0 \geq 1$ such
		% 	that $f^n(T^{-s}\gM) \subseteq T^m \gM$ if $n \geq n_0$.
	\end{enumerate}
\end{alemma}
\begin{proof}By Lemma~\ref{lem:lattice properties}~(1), we see
  that~$(4)\implies (5)\implies (1)\implies (2)\implies (3)$, so we
  only need to that~$(3)\implies (4)$. To this end, note firstly that
  it follows from Lemma~\ref{lem: powers of T quasilinear} and the
  $\Zp$-linearity of~$f$ that for
  each $i,j\ge 0$ and $s\in\Z$ we have
  \[f((p,T)^jT^sf^i(\gM))\subseteq (p,T)^{j}T^sf^{i+1}(\gM)+(p,T)^{j+1}T^sf^i(\gM). \] It
  follows (by induction on~$m$)
  that for each~$m\ge 0$ and $s\in\Z$ we have
  \[ f^m(T^s\gM)\subseteq \sum_{i=0}^m(p,T)^{m-i}T^sf^i(\gM).\] In
  particular, if we take $m=n$ and recall that
  $f^n(\gM)\subseteq(p,T)\gM$ by hypothesis, we find that
  \anumequation\label{eqn: base case of quasilinear induction}f^n(T^s\gM)\subseteq (p,T)T^s\gM+\sum_{i=1}^{n-1}(p,T)^{n-i}T^sf^i(\gM). \end{equation}

The same argument shows that if~$\gN$ is an $\A_A^+$-submodule of~$M$ with
  the property that
  \[\gN\subseteq \sum_{i=0}^{n-1}(p,T)^{a_i}T^sf^i(\gM)\]
      for non-negative integers $a_0,\dots,a_{n-1}$, then % it follows from
  % Lemma~\ref{lem: powers of T quasilinear} and our assumption that
  % $f^n(\gM) \subseteq T \gM$ that
  \[f(\gN)\subseteq  \sum_{i=0}^{n-1}(p,T)^{b_i}T^sf^i(\gM)\]   
  where~$b_0=\min(a_0+1,a_{n-1}+1)$, and~$b_i=\min(a_i+1,a_{i-1})$
  if~$i>0$.  % \[(b_{n-1},\dots,b_1,b_0)=(\min(a_{n-1}+1,a_{n-2}),\min(a_{n-2}+1,a_{n-3}),\dots,\min(a_1+1,a_0),\min(a_0+1,a_{n-1}+1). \]
  It follows by an easy induction on~$N$ (with the base case being
  given by~\eqref{eqn: base case of quasilinear induction}) that for all~$N\ge n$, if we
  write~$N+1=(q+1)n+r$ with $0\le r<n$, then we
  have
  \[f^N(T^{s}\gM)\subseteq \sum_{i=0}^{n-1}(p,T)^{c_i}T^sf^i(\gM)\]where
  \[(c_{n-1},\dots,c_0)=(q,q,\dots,q)+(r,r+1,\dots,n-1,1,2,\dots,r). \]
In particular we see that for all $N\ge n$ we
have   \[f^N(T^{s}\gM)\subseteq (p,T)^{\lfloor (N+1)/n\rfloor-1}T^s\sum_{i=0}^{n-1}f^i(\gM)\]

  Now, by Lemma % s~\ref{lem: powers of T quasilinear} and
  ~\ref{lem: how f
    behaves on lattices}, for any sufficiently large~$t$ we have
 ~$f^i(\gM)\subseteq T^{-t}\gM$ for $0\le i\le n-1$, so it follows
  that for $N\ge n$ we
have   \[f^N(T^{s}\gM)\subseteq (p,T)^{\lfloor
    (N+1)/n\rfloor-1}T^{s-t}\gM.\] Since~$p^a=0$ in~$A$, we also have
$(p,T)^{n+a-1}\subseteq (T^n)$ for all~$n\ge 0$, so that if $N\ge
an-1$ 
% \mar{ME: Just to check, would the more precise estimate be $N \ge an -
%   1$? TG: yes, changed, commenting out}
then we have \[f^N(T^{s}\gM)\subseteq T^{s+\lfloor
    (N+1)/n\rfloor-a-t}\gM.\] Since $\lfloor
    (N+1)/n\rfloor-a-t\to\infty$  as $N\to\infty$, we have~(4), as required.
  % Indeed, we may take~$t=m(n-1)$ where~$m$ is
  % as in Lemma~\ref{lem: how f behaves on lattices}, as we then have
  % $f^{n-1}(\gM)\subseteq f^{n-2}(T^{-m}\gM)\subseteq \dots\subseteq
  % f(T^{-(n-2)m}\gM)\subseteq T^{-(n-1)m}\gM$. 
  % \mar{TG: maybe it's worth upgrading Lemma~\ref{lem: how f 
  %   behaves on lattices} to make this easier? Leaving for now as it's
  % quite possible this argument is more complicated than it needs to
  % be, anyway!}
\end{proof}





% \begin{lemma}\label{lem: cts gamma iff top nilpt} If $A$ is a $\cO/\varpi^a$-algebra,
% 	and if $\Gamma_{\disc} := \langle \gamma \rangle$
% 	acts on a Tate module $M$ over~$A$,
% 	then this action is continuous if and only if 
% 	$\gamma - 1$ is a topologically nilpotent endomorphism
% 	of~$M$.
% \end{lemma}
% \begin{proof}
% 	\mar{ME: This will be a reinterptation of condition~(2) 
% 		in the preceding lemma.}\mar{TG: do we actually need
%                 this or is it OK to just have something in the context
%                 of Lemma~\ref{lem:testing
%                   continuity on M mod T}? I don't have a strong
%                 opinion. }
% \end{proof}

\begin{acor}% \mar{TG: not totally sure if I will keep this in this form,
  % but want to add something}
\label{cor: variants on topological nilpotence mod p}
Suppose that $A$ is an $\cO/\varpi^a$-algebra for
        some~$a\ge 1$.
	If $M$ is a finite projective $\A_A$-module, and if $f$ is a
	$T$-quasi-linear endomorphism of $M$,
	then the following are equivalent:
	\begin{enumerate}
		\item $f$ is topologically nilpotent.
		\item The action of~$f$ on ~$M\otimes_{\cO/\varpi^a}\F$
          is topologically nilpotent.
	\end{enumerate}
      \end{acor}
      \begin{proof}Obviously $(1)\implies (2)$. Conversely, if (2)
        holds, then by the equivalence of conditions (1) and (2) of
        Lemma~\ref{lem:topological nilpotence criteria} for the action
        of~$f$ on ~$M\otimes_{\cO/\varpi^a}\F$, we see that condition
        (3) of 
        Lemma~\ref{lem:topological nilpotence criteria} holds (for the
        action of~$f$ on~$M$), and therefore condition
        (1) of 
        Lemma~\ref{lem:topological nilpotence criteria} holds, as required.        
      \end{proof}
% \begin{proof}
%   Obviously $(1)\implies (2)$. We prove the converse by induction
%   on~$a$. Suppose that~$f$ is topologically nilpotent
%   modulo~$\varpi^a$. By Lemma~\ref{lem:topological nilpotence
%     criteria}, this implies that there  is
%   a lattice~$\gM$ in~$M$ and
%   an integer~$n\ge 1$ such that for all~$s\in\Z$, we have
%   $f^n(T^s\gM)\subseteq T^{s+1}\gM+\varpi^{a-1}\gM$. It follows from Lemma~\ref{lem: how f behaves
%     on lattices} that for some integer~$l$, we have
%   \anumequation\label{fn inductive}f^n(T^s\gM)\subseteq T^{s+1}\gM+\varpi^{a-1}T^{s+l}\gM\end{equation} for
%   all~$s\in\Z$.

%   Then we have \[f^{2n}(T^s\gM)\subseteq
%     f^n(T^{s+1}\gM)+\varpi^{a-1}f^n(T^{s+l}\gM), \]and
%   employing~\eqref{fn inductive} with~$s$ replaced by~$s+1$ and
%   by~$s+l$, and using that $(\varpi^{a-1})^2=0$, we
%   obtain \[f^{2n}(T^s\gM)\subseteq
%     T^{s+2}\gM+\varpi^{a-1}T^{s+l+1}\subseteq
%     T^{s+1}\gM+\varpi^{a-1}T^{s+l+1}. \]Thus at the expense of
%   replacing~$n$ by~$2n$, we can replace~$l$ by~$l+1$; so by induction
%   on~$l$, we see that we can enlarge~$n$ and assume that~$l\ge
%   1$. Then~\eqref{fn inductive} becomes~$f^n(T^s\gM)\subseteq
%   T^{s+1}\gM$, so in particular we have~$f^n(\gM)\subseteq T\gM$,
%   and~$f$ is topologically nilpotent by Lemma~\ref{lem:topological nilpotence criteria}.
% \end{proof}
The following lemmas provide the key examples of $T$-quasi-linear
endomorphisms, and explain our interest in the concept.

We suppose that~$\A_A$ is endowed with a continuous action
%of the subgroup~$\langle\gamma\rangle$
of~$\Zp$ by $A$-algebra automorphisms which preserve~$\A^+_A$,
and suppose further that, for some 
topological generator $\gamma$  of~$\Zp$, that
\anumequation\label{eqn: condition on gamma semi-linear}\gamma(T)-T\in(p,T)T\A^+_A.\end{equation} % \mar{TG: might
% move this to phi gamma section shortly, then reference it here?} 
% ,
% and in the question of their topological nilpotency.

\begin{alemma}
  \label{lem: all powers of gamma preserve the ideal}If~\eqref{eqn:
    condition on gamma semi-linear} holds, then $\gamma$ preserves the
  ideals~$(T)$ and~$(p,T)$ of~$\A^+_A$. Furthermore for each
  integer~$n\ge 1$ we have
  \anumequation\label{eqn: condition on powers of gamma semi-linear}\gamma^n(T)-T\in(p,T)T\A^+_A.\end{equation}
\end{alemma}
\begin{proof}The first claim follows immediately from~\eqref{eqn:
    condition on gamma semi-linear}, which shows that~$\gamma(T)$ is a
  unit multiple of~$T$, together with the fact that  $\gamma$  
preserves $\A^+_A$. Then~\eqref{eqn: condition on powers of gamma semi-linear}
    follows by induction on~$n$ (the case~$n=1$
  being~\eqref{eqn: condition on gamma semi-linear}).
\end{proof}

\begin{alemma}
	\label{lem:gamma minus 1 T quasi linear}
%        If~$\A_A$ is endowed with an action of~$\langle \gamma\rangle$
%        satisfying~\eqref{eqn: condition on gamma semi-linear}, and
If
        $M$ is a finite projective $\A_A$-module which
	is endowed with an action of the subgroup  $\langle
\gamma \rangle$  of $\Zp$ which  is semi-linear with  respect
to the given action of this group on~$\A_A$ {\em(}obtained
by restricting the $\Z_p$-action{\em )}, 
% \mar{ME: Based on the way this  is cited, I think we just want
% to assume $\gamma$ acts, right, not nec.\ all of $\Z_p$.  This
% doesn't change anything in the conclusion or proof. TG: agreed!
% haven't changed yet as not quite sure how best to reword everything.}
	then for any integer~ $n\ge 1$, $f := \gamma^n -1$ is a $T$-quasi-linear endomorphism of~$M$.
\end{alemma}
\begin{proof}% \mar{TG: needs to be thought through in the current level
  % of abstraction for $\gamma$ action.}
  We have $f(Tm)=\gamma^n(T)f(m)+(\gamma^n(T)-T)m$, so it follows from~\ref{eqn: condition on powers of gamma semi-linear} that~$f$ is $T$-quasi-linear.%  The rest
  % of the lemma is immediate from Lemma~\ref{lem: cts gamma iff top
  %   nilpt}.\mar{TG: that lemma has no proof as yet.}
\end{proof}




\begin{alem}% \mar{TG: I think this needs some further discussion of
    % jointly continuous versus separately continuous for group actions?}
	\label{lem:testing continuity on M mod T}% \mar{TG: currently I
        %   don't think this lemma is proved except if $a=1$. Possibly
        %   what we need to do is to work with $T$-quasilinearity modulo
        % $\varpi^a$? That seems a bit painful though because we would
        % have to work with $a(T), b(T)$ being in $\A_K$ rather than $\A_K^+$.}
        Suppose that $A$ is an $\cO/\varpi^a$-algebra for
        some~$a\ge 1$, and that~$\A_A$ is endowed with an action of~$\Zp$
        satisfying~\eqref{eqn: condition on gamma semi-linear}. Let~$M$
        be a finite projective $\A_A$-module, equipped with a
        semi-linear action of~$\langle \gamma\rangle\subset \Zp$. Then
        the following are equivalent:
        \begin{enumerate}
        \item\label{item: gamma cts} The action of~$\langle \gamma\rangle$ extends to a continuous
          action of~$\Zp$.
        \item\label{item: gamma cts mod p} The action of~$\langle \gamma\rangle$ on~$M\otimes_{\cO/\varpi^a}\F$ extends to a continuous
          action of~$\Zp$.
        \item\label{item: gamma minus 1 squeezing M to p T n} For any lattice~$\gM\subseteq M$, and any~$n\ge 1$, there exists~$s\ge 0$
          such that $(\gamma^{p^s}-1)^i(\gM)\subseteq (p,T)^n\gM$ for
          all~$i\ge 1$.
             %   \item For any lattice~$\gM\subseteq M$, there exists~$s\ge 0$
          % such that $(\gamma^{p^s}-1)(\gM)\subseteq (p,T)\gM$.
          \item\label{item: gamma minus 1 squeezing M to T} For any lattice~$\gM\subseteq M$, there exists~$s\ge 0$
          such that $(\gamma^{p^s}-1)(\gM)\subseteq T\gM$.
        % \item For some lattice~$\gM\subseteq M$ and some~$s\ge 0$, we have
        %   $(\gamma^{p^s}-1)(\gM)\subseteq T\gM$.
           \item\label{item: gamma minus 1 squeezing some M to p T} For some lattice~$\gM\subseteq M$ and some~$s\ge 0$, we have
          $(\gamma^{p^s}-1)(\gM)\subseteq (p,T)\gM$.
        \item\label{item: gamma minus 1 top nilpt mod p}  The action of~$\gamma-1$ on~$M\otimes_{\cO/\varpi^a}\F$
          is topologically nilpotent.
          \item\label{item: gamma minus 1 top nilpt}  The action of~$\gamma-1$ on~$M$
          is topologically nilpotent.
        \end{enumerate}
  \end{alem}
  \begin{proof}% We trivially have % $(1)\implies (2)$ and
    % $(3)\implies
    % (4)$. Conversely if~(4) holds, then  writing
    %     \[\gamma^{p^{s+n+1}} - 1 = ((\gamma^{p^{s+n}}-1)+1)^p-1,\] it
    %     follows from an easy induction on~$n$ that for each~$n\ge 0$,
    %     we have
    %     \[(\gamma^{p^{s+n}}-1)(\gM)\subseteq (p,T)^{n+1}\gM,\] so~(3)
    %     holds.
        % \mar{TG: I think that for semi-linearity reasons this
          % ``easy induction'' might need more work. TG: indeed I don't
          % currently understand what this argument should be!}

        Noting that if~$n\ge a$ then $(p,T)^n\subseteq (T)$, we see that
        \eqref{item: gamma minus 1 squeezing M to p T
          n}$\implies$\eqref{item: gamma minus 1 squeezing M to T},
        and by Lemma~\ref{lem:topological nilpotence criteria}
        and Corollary~\ref{cor: variants on topological nilpotence mod p} we
        have \eqref{item: gamma minus 1 squeezing M to
          T}$\implies$\eqref{item: gamma minus 1 squeezing some M to p
          T}$\implies$\eqref{item: gamma minus 1 top nilpt mod
          p}$\implies$\eqref{item: gamma minus 1 top nilpt}. (For \eqref{item: gamma minus 1 squeezing some M to p
          T}$\implies$\eqref{item: gamma minus 1 top nilpt mod
          p}, we also use that
        $(\gamma-1)^{p^s}\equiv(\gamma^{p^s}-1)\pmod{\varpi}$.)
        
%\mar{ME: I got a bit lost checking  this,
%because it sometimes seemed that we were applying quasi-linearity
%results from above to $\gamma^{p^s}-1$, rather than to $\gamma - 1$.
%And I didn't think through whether this was valid (but also wonder if
%I misunderstood). TG: unclear to me whether ``this'' is above or below
%the comment! Happy to try to clarify, I don't intend to apply
%quasi-linearity to $\gamma^{p^s}-1$, and don't see where it might be
%happening, though, so I don't know what to be doing. ME: I think this is now all good.}
        We next show that \eqref{item: gamma minus 1 top
          nilpt}$\implies$\eqref{item: gamma minus 1 squeezing M to p
          T n}. Suppose that~\eqref{item: gamma minus 1 top
          nilpt} holds. Recalling again that for each~$s\ge 0$ we have
        $(\gamma-1)^{p^s}\equiv(\gamma^{p^s}-1)\pmod{\varpi}$, it
        follows from Corollary~\ref{cor: variants on topological
          nilpotence mod p} and Lemma~\ref{lem:gamma minus 1 T quasi linear} that the action of $(\gamma^{p^s}-1)$
        on~$M$ is topologically nilpotent for each~$s\ge 0$. % \mar{TG:
          % not sure if using this in what follows}
        We now
        argue by induction on~$a$, noting that for $a=1$, the
        implication \eqref{item: gamma minus 1 top
          nilpt}$\implies$\eqref{item: gamma minus 1 squeezing M to p
          T n} is immediate from Lemma~\ref{lem:topological nilpotence
          criteria}. We may therefore assume
        that \[(\gamma^{p^s}-1)^i(\gM)\subseteq (p,T)^n\gM
          +\varpi^{a-1}M \]for all~$i\ge 1$. It follows in particular
        that $p(\gamma^{p^s}-1)^i(\gM)\subseteq (p,T)^n\gM$ for
        all~$i\ge 1$, so that for any~$t\ge 0$ and any~$i\ge 1$, we
        have \[((\gamma^{p^{s+t}}-1)^i-(\gamma^{p^s}-1)^{ip^t})(\gM)\subseteq
          (p,T)^n\gM.\](To see this, write
        $(\gamma^{p^{s+t}}-1)=((\gamma^{p^s}-1)+1)^{p^t}-1$ and use the
        binomial theorem.) It therefore suffices to show that there is
        some~$t\ge 1$ for which~$(\gamma^{p^s}-1)^{ip^t}(\gM)\subseteq
        (p,T)^n\gM$ for all~$i\ge 1$; but we have already seen
        that~$(\gamma^{p^s}-1)$ acts topologically nilpotently on~$M$,
        so by Lemma~\ref{lem:topological nilpotence criteria} we can
        even arrange that $(\gamma^{p^s}-1)^{ip^t}(\gM)\subseteq
        T^n\gM$ for all~$i\ge 1$. 

  
        
%         denotes % \mar{TG: I
% %   think there might be a gap here, I think possibly this is implicitly
% % assuming that $\gamma$ preserves $\A^+$. I suspect that here and below
% % the solution is to use continuity of the $\Gamma$ action on~$\A$ to show that for $s$ large enough,
% % $\gamma^{p^s}$ does preserve $\A^+$.}
% $\gM\otimes_{\cO/\varpi^a}\F$, we have
% $(\gamma-1)^{p^s}(\overline{\gM})=(\gamma^{p^s}-1)(\overline{\gM})\subset
% T\overline{\gM}$, so it follows from Lemma~\ref{lem:topological
%   nilpotence criteria} that~(8) holds. If on the other hand we assume
% that~(8) holds, then the same argument, together with the equivalence
% of conditions~(1) and~(4) of
% Lemma~\ref{lem:topological nilpotence
%   criteria}, shows that for any lattice~$\gM\subseteq M$, there exists ~$s\ge 0$
%           such that $(\gamma^{p^s}-1)(\gM)\subseteq (p,T)\gM$, % \mar{TG:
%           % I think this is wrong - rather it shows that $(\gamma^{p^s}-1)(\gM)\subseteq T\gM+pM$}
%         so that~(4)
%         holds.  Thus conditions (3)--(8) are all equivalent.

        We have shown the equivalence of conditions \eqref{item:
          gamma minus 1 squeezing M to p T n}--\eqref{item: gamma
          minus 1 top nilpt}. Suppose now that~\eqref{item: gamma cts} holds. Then Lemma~\ref{lem: continuity and
          lattices} shows that, for each lattice $\gM\subseteq M$,
we have~$\gamma^{p^s}(\gM)\subseteq\gM$
for all sufficiently large~$s\ge 0$,
and that furthermore, for each
        $m\in\gM$, there is some~ $t(m)$ such that if~$t\ge t(m)$, then
        $\gamma^{p^{t}}(m)\in m+(p,T)\gM$. Letting~$m_1,\dots,m_n$
        be generators for~$\gM$ as an $\A_{A}^+$-module, we see that
        if~$s$ is sufficiently large, then $(\gamma^{p^s}-1)(m_i)\in
        (p,T)\gM$ for each~$i$. Then if~$\lambda_i\in \A_{A}^+$, we
        have \[(\gamma^{p^s}-1)(\sum_i
          \lambda_im_i)=\sum_i\gamma^{p^s}(\lambda_i)(\gamma^{p^s}-1)(m_i)+\sum_i(\gamma^{p^s}-1)(\lambda_i)m_i.\]It
        follows from~\eqref{eqn: condition on powers of gamma semi-linear}
%from~\eqref{eqn: condition on gamma semi-linear} % \mar{TG: again, $\A$
          % versus $\A^+$}
        that~$(\gamma^{p^s}-1)(\lambda_i)\in T\A_{A}^+$; hence $(\gamma^{p^s}-1)(\gM)\subseteq
        (p,T)\gM$, and so~\eqref{item: gamma minus 1 squeezing some M to p T} holds.

Suppose now that the equivalent conditions \eqref{item:
          gamma minus 1 squeezing M to p T n}--\eqref{item: gamma
          minus 1 top nilpt} hold. We will show 
that~\eqref{item: gamma cts} holds.  By Lemma~\ref{lem: extending Z to
          Zp} (taking the group $G$ there to be $\langle \gamma \rangle \cong \Z$
endowed with its $p$-adic topology, so that $\widehat{G}~\cong~\Z_p$),
it is enough to show that  the conditions of
Lemma~\ref{lem: continuity and lattices}~(2) hold. % and~\ref{lem: extending Z to Zp},
% taking the group $G$ there to be $\langle \gamma \rangle \cong \Z$
% endowed with its $p$-adic topology
% (so that $\widehat{G}~\cong~\Z_p$). 
%Lemma~\ref{lem:
%  continuity and lattices}~(2). Note that by
%Lemma~\ref{lem: extending Z to Zp},
%mar{TG: now need to actually check the hypotheses of that lemma} it suffices to show that the action of~$\langle
%\gamma\rangle$ on~$M$ is continuous.

We begin with Lemma~\ref{lem: continuity and lattices}~(2)(b), the
condition that any~$g\in \langle \gamma\rangle$ acts continuously
on~$M$. It is enough to show that
if~$\gM\subseteq M$ is a lattice, then there
is a lattice~$\gN$ with $g(\gN)\subseteq \gM$.  This is in fact a general
property of semilinear automorphisms of~$\A_A$ which
preserve~$\A^+_A$. % :
% indeed, by Lemma~\ref{lem:lattice properties}~(1), 
Indeed, let~$m_1,\dots,m_n$ be
generators of~$\gM$ as an~$\A_{A}^+$-module, and let~$\gN$ be the
$\A_{A}^+$-module generated by $g^{-1}(m_1),\dots,g^{-1}(m_n)$. This
is a lattice, because~$g$ is an automorphism of~$M$, and it
follows easily from the semi-linearity of the action of~$g$ on~$M$ that
$g(\gN)\subseteq\gM$,  as required.% \mar{TG: again it looks like there's
  % something dubious about $g$ preserving $\A^+$ here.}

We now check Lemma~\ref{lem: continuity and lattices}~(2)(c). Let~$\gM\subseteq M$ be some lattice, fix a choice of $n\ge 1$,
and then choose~$s$ as in~\eqref{item: gamma minus 1 squeezing M to p
  T n}. It suffices to show that the subgroup~$H=\langle \gamma^{p^s}\rangle$ 
of~$\langle \gamma\rangle$ preserves~$\gM$. Since $(\gamma^{p^s}-1)(\gM)\subseteq
(p,T)^n\gM\subseteq\gM$, we certainly
have~$\gamma^{p^s}(\gM)\subseteq\gM$, so it suffices to show that
$\gamma^{-p^s}(\gM)\subseteq\gM$. For this, note that
since (as recalled above) %at the beginning of the lemma)
it follows from the congruence 
$(\gamma-1)^{p^s}\equiv(\gamma^{p^s}-1)\pmod{\varpi}$ that
$\gamma^{p^s}-1$ acts topologically nilpotently on~$M$, and
preserves~$\gM$, %so  % ~$(\gamma^{p^s}-1)(\gM)\subseteq (p,T)^n\gM\subseteq (p,T)\gM$, 
we see that for
any $m\in \gM$ we have
%may write
\[\gamma^{-p^s}(m)=(1-(1-\gamma^{p^s}))^{-1}(m)=m+(1-\gamma^{p^s})(m)+(1-\gamma^{p^s})^2(m)+\dots\in\gM, \]as
required.  

To complete the verification of the conditions of Lemma~\ref{lem:
  continuity and lattices}~(2), we need to show that for any
lattice $\gM~\subseteq~M$, 
the orbit maps $\langle \gamma \rangle \to M$, for the various $m \in \gM$,
are (equi)continuous at the identity of~$\langle \gamma \rangle$.
%each~$m\in M$, the map~$g\mapsto gm$ is continuous. Choose a lattice
%~$\gM$ containing~$m$ (as we may, by choosing any lattice and scaling
%by an appropriate power of~$T$).
It is enough to show that for
each~$n\ge 1$, we can find~$s$ sufficiently large such that
$H=\langle \gamma^{p^s}\rangle$ satisfies~$Hm\subseteq m+(p,T)^n\gM$,
for each $m \in \gM$,
or equivalently,
that $(h-1)(\gM)\subseteq (p,T)^n\gM$, for all~$h\in H$. 
We accomplish this by choosing~$s$ as in~\eqref{item: gamma minus 1
  squeezing M to p T n}.

It then suffices to prove
the stronger claim that~$(\gamma^{rp^s}-1)(\gM)\subset(p,T)^n\gM$ for
all~$r\in\Z$. We have already seen that $\gamma^{rp^s}(\gM)\subseteq\gM$. % Note firstly that since~$(\gamma^{p^s}-1)(\gM)\subset
% (p,T)^n\gM\subset\gM$, we have~$\gamma^{p^s}(\gM)\subset\gM$, so
% that~$\gamma^{rp^s}(\gM)\subseteq \gM$ for all~$r\ge 1$.
If~$r\ge
1$ we may write
$(\gamma^{rp^s}-1)=(\gamma^{p^s}-1)(1+\gamma^{p^s}+\dots+\gamma^{(r-1)p^s})$,
and since $(1+\gamma^{p^s}+\dots+\gamma^{(r-1)p^s})(\gM)\subset\gM$,
we have~$(\gamma^{rp^s}-1)(\gM)\subseteq (\gamma^{p^s}-1)(\gM)\subset
(p,T)^n\gM$, as
required. If~$r\le 0$, then the result follows by writing~$(\gamma^{-rp^s}-1)=-(\gamma^{rp^s}-1)(\gamma^{-p^s})^r$.% ,
% we see that to prove the claim for~$r\le 0$, it is enough to show that
% $\gamma^{-p^s}(\gM)\subset\gM$. For this, note that
% since~$(\gamma^{p^s}-1)(\gM)\subseteq (p,T)^n\gM\subseteq (p,T)\gM$, for
% any $m\in \gM$ we may
% write \[\gamma^{-p^s}(m)=(1-(1-\gamma^{p^s}))^{-1}(m)=m+(1-\gamma^{p^s})(m)+(1-\gamma^{p^s})^2(m)+\dots\in\gM, \]as required.


%we find that indeed
%$(h-1)(\gM)\subseteq (p,T)^n\gM$ for all~$h\in H$,
%as required.
%so since~$m\in\gM$, we are done. % \mar{TG: need to finish this too -
  % probably this and the previous point should be done in a similar way?}
        % To see
        % this, note firstly that by Lemma~\ref{lem: quasilinear implies
        %   continuous}, $(\gamma-1)$ is a continuous endomorphism of
        % $M\otimes_{\cO/\varpi^a}\F$, so~$\gamma$ is a continuous automorphism of
        % $M\otimes_{\cO/\varpi^a}\F$.
        
        Finally, applying the equivalence of~\eqref{item: gamma cts} and~\eqref{item: gamma minus 1 top nilpt mod p} with~$M$
        replaced by $M\otimes_{\cO/\varpi^a}\F$, we see that~\eqref{item: gamma cts mod p}
        and~\eqref{item: gamma minus 1 top nilpt mod p} are equivalent, as required.\end{proof}
    % Since $\Gamma\cong\Zp$, and the canonical topology
    %     on~$M$ is induced by the $(p,T)$-adic topology on the
    %     any lattice~$\gM$, we see that the action of~$\Gamma$ is
    %     continuous if and only if for some (equivalently, for every)
    %     lattice~$\gM\subseteq M$, and every~$n\ge 0$, there exists~$s(n)$
    %     such that $(\gamma^{p^s}-1)(\gM)\subseteq (p,T)^n\gM$ for all
    %     $s\ge s(n)$. If~$n\ge a$ then $(p,T)^n\subseteq (T)$, so we see
    %     that $(1)\implies (2)\implies (3)\implies (4)$.

    %     Suppose now that~(4) holds. We have
    %     $(\gamma^{p^s}-1)(\gM)\subseteq % T\gM\subset
    %     (p,T)\gM$. 
    %   \end{proof}

      % \begin{cor}
      %   \label{cor: can check continuity modulo p}Suppose that $A$ is an $\cO/\varpi^a$-algebra for
      %   some~$a\ge 1$. Let~$M$ be a projective \'etale $\varphi$-module with
      %   $A$-coefficients, equipped with an action
      %   of~$\Gamma_\disc$. Then the following are equivalent:
      %   \begin{enumerate}
      %    \item The action of~$\Gamma_{\disc}$ on~$M$ extends to a continuous
      %     action of~$\Gamma$. \item The action of~$\Gamma_{\disc}$ on~$M\otimes_{\cO/\varpi^a}\F$ extends to a continuous
      %     action of~$\Gamma$.
      %   \end{enumerate}

%       \end{cor}
%       \begin{proof}
% This follows from the equivalence of conditions~(1) and~(4) in Lemma~\ref{lem:testing continuity on M mod T}.
%       \end{proof}\mar{TG: numbering probably changed.}


The following lemmas will allow us to reduce the problem of investigating
the topological nilpotency of a $T$-quasi-linear endomorphism
from the projective case to the free case.


\begin{alemma}
	\label{lem:quasi-linear direct sum}
	If $M_1$ and $M_2$ are Tate modules over a ring~$A$,
	endowed with continuous endomorphisms $f_1$ and $f_2$
	respectively, then $f_1$ and $f_2$ are both topologically
	nilpotent if and only if the direct sum $f := f_1 \oplus f_2$
	is a topologically nilpotent endomorphism.
\end{alemma}
\begin{proof}% \mar{TG: needs proof still but I presume it will be
    % easy. TG: in fact it looks completely trivial - what am I
    % missing?!}
This is immediate from the definitions.
\end{proof}


% \mar{TG: could presumably adapt the following to $\cO/\varpi^a$-algebras,
%   leaving it until we're reading through though in case I screw
%   anything up doing that.}
\begin{alemma}
	% \mar{ME: Not sure we actually need the $\varphi$-linear aspect
	% 	of this, which would make it even easier! TG: I don't
        %         think it's so much easier in terms of explanations to
        %         take it out, though, so I'm leaving this
        %         alone. Commenting out.}
	\label{lem:quasi-linear complement}
	If $M$ is a finite projective $\A_A$-module
       	% \'etale $\varphi$-module over an $\F_p$-algebra~$A$, 
	endowed with a $T$-quasi-linear endomorphism $f$,
	then we may find a finite projective $\A_A$-module
	$N$, endowed with a $T$-quasi-linear endomorphism
	$g$ which is furthermore topologically nilpotent,
	such that $M\oplus N$ is a free $\A_A$-module.
\end{alemma}
\begin{proof}
  By % ~ \cite[Lem.\ 5.2.14]{EGstacktheoreticimages}, we may find a
definition, there is a $\A_A$-module   $N$ such that $M \oplus N$ is
  free.  We may then take $g$ to be the zero endomorphism of~$N$; this
  is evidently both $T$-quasi-linear and topologically nilpotent.
\end{proof}
If $A$ is an $\cO/\varpi^a$-algebra for some~$a\ge 1$, and $M$ is a
finite projective $\A_A$-module, equipped with a $T$-quasi-linear endomorphism $f$,
and if $A \to B$ is a morphism of $\cO/\varpi^a$-algebras, then we have a
base-changed $W(k)\otimes_\Zp B$-linear endomorphism of~$B\otimes_AM$.
This base-changed endomorphism is % also $T$-quasi-linear and therefore
$T$-adically continuous (by Lemma~\ref{lem: how f behaves on lattices}) % (where we are slightly abusing terminology by extending the
% notation of $T$-quasi-linearity to endomorphisms of modules
% over~$B\otimes_A\A_{A}$),
%\mar{ME: I wonder if it's possible to slightly recast the discussion
%	of quasi-linearity so that this case is literally included,
%	just since it's pretty important that all this material is
%        correct! TG: rearranged a bit, check that you're happy. ME: Looks good}
and it therefore induces an
endomorphism $f_B$ of the base-changed projective $\A_B$-module $M_B$ (which by definition is  the completion of~$B\otimes_AM$ for
the $T$-adic topology). The endomorphism~$f_B$ is evidently also $T$-quasi-linear.
% \mar{ME: Don't we need $f_B$ to be $T$-adically continuous for
% 	this to be defined?  I guess we should assume $f$ is $T$-adically
% 	continuous, and then check that $f_B$ is. TG: good point. Perhaps we
%         should just restrict (as you did before) to the quasi-linear
%         case, which we've shown is continuous? TG: done, commenting out}
(Note that since~$f$ is not necessarily
$A[[T]]$-linear, we do not define~$f_B$ by viewing~$M_B$ as
$B[[T]]\otimes_{A[[T]]} M$.) % If~$f$ is $T$-quasi-linear, then so is $f_B$. 
% of
% the base-changed \'etale $\varphi$-module $M_B$, defined by viewing
% $M_B$ as the completion of$B\otimes_AM$ for the $T$-adic..  \mar{ME:
%   Maybe elaborate on this construction, and make it a lemma, just
%   because $f$ is not actually $T$-linear, so we have to think of $M_B$
%   as being a completion of $B\otimes_A M$, rather than
%   $B[[T]]\otimes_{A[[T]]} M.$} 
If $f$ is furthermore topologically
nilpotent, then so is $f_B$ (for example, by Lemma~\ref{lem:topological nilpotence criteria}).

\begin{alemma}
	\label{lem:quasi-linear limit preserving}
	Let $\{A_i\}_{i \in I}$ be a directed system of $\cO/\varpi^a$-algebras,
	with $A = \varinjlim_{i \in I} A_i$, and suppose that $I$ admits
	a least element $i_0$.  Let $M$ be a finite projective $\A_{A_{i_0}}$-module,
	 let $f$ be a $T$-quasi-linear endomorphism of $M$,
	and suppose that the base-changed endomorphism $f_A$ is topologically
	nilpotent.  Then for some $i \in I,$ the base-changed
	endomorphism $f_{A_i}$ is topologically nilpotent.
\end{alemma}
\begin{proof}
	By Lemma~\ref{lem:quasi-linear complement},
	we may choose a finite projective  $\A_{A_{i_0}}$-module $N$,
	endowed with a topologically nilpotent $T$-quasi-linear
	endomorphism $g$, such that $M \oplus N$ is free.
	By Lemma~\ref{lem:quasi-linear direct sum}, the base-changed
        endomorphism $f_A \oplus g_A = (f\oplus g)_A$ of $M_A\oplus N_A$
        is topologically nilpotent, and it suffices to show 
	that $f_{A_i} \oplus g_{A_i} = (f\oplus g)_{A_i}$
	is topologically nilpotent
	for some $i \in I$.   Thus we may reduce to the case
	when $M$ is free, in which case we may also choose
	a free lattice $\gM$ contained in $M$.   Taking into account
	Lemma~\ref{lem:topological nilpotence criteria}, % and replacing
	% $f$ by $f^n$ for some $n \geq 1$, we may assume that
%\mar{ME: I don't think there's any actual problem, but I wonder
%if we shouldn't make this replacement, just b/c $f^n$ may not be
%quasi-linear. TG: hopefully fixed? ME: Great}
for some sufficiently large~	$n$ we have $f_A^n(\gM_A) \subseteq T\gM_A$, and it suffices to show that
	\anumequation
	\label{eqn:needed inclusion}
	f_{A_i}^n(\gM_{A_i}) \subseteq T \gM_{A_i}
\end{equation}
       	for some $i \in I$.
	Lemma~\ref{lem: how f behaves on lattices} shows that
	$f^n(T^r \gM) \subseteq T\gM$ for some $r\geq 0,$
       	and that $f^n(\gM) \subseteq T^{-s} \gM$
	for some $s \geq 0$.  Thus $f^n$ induces a morphism
	$\gM/T^r \gM \to T^{-s}\gM/T{\gM}$ of finite rank free
	$A_{i_0}$-modules,
	which, by assumption, vanishes after base-change to $A$.  
	Thus this morphism in fact vanishes after base-change
	to some $A_i,$ and consequently~(\ref{eqn:needed inclusion})
	does indeed hold for this choice of $A_i$.
\end{proof}

%\mar{ME: Initial guess at formulation of ``reduction to Artinian
%	quotients'' lemma. ME: Don't think we need this here, actually;
%commenting out.}
%\begin{lemma}
%	Suppose that $A$ is a finite type $\F_p$-algebra, that $M$ 
%	is an \'etale $\varphi$-module over $A$, and that $f$
%	is a $T$-quasi-linear endomorphism of $M$.
%	Suppose further that there exists $h \geq 0$ such that,
%	for each Artinian quotient $B$ of $A$, 
%	the base-changed module $M_B$ contains a lattice $\gM_B$
%	of height $\leq h$, such that $f_B(\gM_B) \subseteq T \gM_B$.
%	Then $f$ is topologically nilpotent.
%\end{lemma}
%\begin{proof}
%\end{proof}

%\chapter{Banach algebra valued points of formal algebraic stacks}
%\label{app:p inverted}
%As in Section~\ref{sec: phi modules and phi gamma modules},
%we fix a finite extension $E$ of $\Q_p$, with ring of integers $\cO$.
%As usual, by a {\em Banach algebra} over $E$, or $E$-{\em Banach algebra},
%we mean a complete
%topological
%$E$-algebra $A$ whose topology may be defined by an ultrametric norm.
%(With this definition, a Banach algebra $A$ does {\em not} come
%equipped with any particular choice of norm defining its topology.)
%
%\mar{ME: Not the best terminology --- not sure if there's something
%	standard. TG: I couldn't find anything.}
%\begin{adefn}
%	\label{def:orders}
%	If $A$ is an $E$-Banach algebra, 
%	then by an {\em order} in $A$, we mean a bounded open $\cO$-subalgebra
%	of $A$.
%\end{adefn}
%
%%\mar{ME: Not sure if we need this.}
%\begin{alemma}
%	\label{lem:orders}
%	If $A^{\circ}$ is a $\cO$-subalgebra of an
%	$E$-Banach algebra~$A$, 
%	then the following are equivalent:
%	\begin{enumerate}
%		\item $A^{\circ}$ is an order in $A$.
%		\item $A^{\circ}$ is closed and bounded,
%			and spans $A$ over~$E$.
%		\item There is a norm $| \, |$ on $A$ defining its topology,
%			and satisfying the additional condition
%			$|ab| \leq |a| |b|$ for all $a,b \in A$,
%			with respect to which $A^{\circ}$ is
%			the unit ball.
%	\end{enumerate}
%\end{alemma}
%\begin{proof}
%	The equivalence of (1) and (2) follows from the fact that
%for any $\cO$-submodule $L$ of a Banach space $B$ over $E$,
%it is equivalent for $L$ to
%be open and bounded, or to be closed and bounded and span $B$ over $L$.
%(This is the non-archimedean version of the Banach--Steinhaus theorem,
%in its formulation as the claim that Banach spaces are {\em barelled}.)
%
%If $| \, |$ is a norm on $A$ as in~(3), then clearly the unit ball
%with respect to this norm is an order in $A$.  Conversely,
%if $A^{\circ}$ is an order, then the gauge of $A^{\circ}$ is a norm
%defining the topology on $A$,
%with respect to which $A^{\circ}$ is the unit ball.  The
%fact that $A^{\circ}$ is closed under multiplication shows that this norm
%satisfies the condition of~(3).   Thus~(1) and~(3) are equivalent.
%\end{proof}
%
%\begin{alemma}
%	\label{lem:directed}
%	If $A$ is an $E$-Banach algebra, then the set of orders
%	in $A$ is non-empty. If $A$ is commutative,
%	then it is furthermore directed by inclusion.
%\end{alemma}
%\begin{proof}
%	Choose any norm $| \, |$ defining the topology of $A$.
%	Then since $A\times A \to A$ is continuous, we find
%	that $|a b| \leq C |a| | b|$, for some constant $C > 0$. 
%        Replacing $| \, |$ by the equivalent norm $C | \, |$,
%	we find that we may assume that in fact
%	$| a b | \leq |a| |b|.$  The unit ball with respect
%	to this norm then forms an order in $A$, 
%	by Lemma~\ref{lem:orders},
%	showing that the set of orders is non-empty.  
%
%	If $A^{\circ}$ and $B^{\circ}$ are two orders in~$A$,
%	choose (as we may, by Lemma~\ref{lem:orders}) a norm
%	defining the topology of $A$, with respect to which
%	$A^{\circ}$ is the unit ball.  Then since $B^{\circ}$ is
%	bounded, we see that the elements in the product
%	$A^{\circ} B^{\circ}$ are of uniformly bounded
%	norm.  Thus, if we let $C^{\circ}$ denote
%	the $\cO$-subalgebra of $A$ generated by $A^{\circ}$
%	and $B^{\circ}$, we see that $C^{\circ}$ is bounded.
%	(This is where we use the hypothesis that $A$ is commutative.)
%	Of course it is also open (as it contains $A^{\circ}$),
%       and thus it again forms an order in $A$,
%containing both $A^{\circ}$ and $B^{\circ}$.  This shows
%that the set of orders in $A$ is directed under inclusion.
%\end{proof}
%
%\begin{alemma}
%	\label{lem:mapping orders}
%	If $f: A \to B$ is a continuous homomorphism of $E$-Banach algebras,
%	and $A^{\circ}$ is an order in $A$, then the image $f(A^{\circ})$ 
%	is contained in an order in $B$.
%\end{alemma}
%\begin{proof}
%	\mar{ME: Give proof TG: do we care about the non-commutative
%          case? If it's commutative, can we argue as follows:
%          $f^{-1}(B^\circ)$ is open so contains $\lambda A^\circ$ for
%          some~$\lambda\ne 0$, so $f(A^\circ)\subset
%          \lambda^{-1}B^\circ$ is bounded, and $f(A^\circ)\subset
%          B^{\circ}\cdot f(A^\circ)$ with the RHS open and bounded?}
%\end{proof}
%
%Suppose now that $\cX$ is a formal algebraic stack over $\Spf \cO$.
%
%\begin{adefn}
%	\label{def:evaluation on Banach algebras}
%	If $A$ is a commutative $E$-Banach algebra,
%	then we define 
%	$\cX(A) := \varinjlim \cX(A^{\circ}),$
%	where the direct limit is understood as a $2$-colimit
%	of groupoids (so that $\cX(A)$ is again a groupoid),
%	taken over the set of orders $A^{\circ}$ in $A$
%	(which Lemma~\ref{lem:directed} shows is a non-empty
%	directed set).
%\end{adefn}
%
%\mar{ME: Explain how to define pull-backs, using Lemma~\ref{lem:mapping
%		orders}. TG: not sure what this is about - is it about
%            pulling back to~$A$? Something else?}

\chapter{Points, residual gerbes, and isotrivial
  families}\chaptermark{Points and residual gerbes}
\label{app:residual gerbes}
Recall that if $\cX$ is an algebraic stack, then the underlying set 
$|\cX|$ of points of $\cX$ is defined as the set of equivalence classes
of morphisms $\Spec K \to \cX$, with $K$ being a field; two such morphisms
are regarded as equivalent if they may be dominated by a common morphism
$\Spec L \to \cX$.  

If $X$ is a scheme (regarded as an algebraic stack), then the
set of underlying points $|X|$ is naturally identified with the
underlying 
set of points of $X$ in the usual sense,
and the equivalence class of 
morphisms representing a given point $x \in |X|$
has a canonical representative, namely the morphism $\Spec k(x) \to X$,
where $k(x)$ is the residue field of $x$.  More abstractly,
this morphism is a monomorphism, and this 
property characterizes it uniquely,
up to unique isomorphism,
among the morphisms in the equivalence class corresponding to~$x$.

If $X$ is an algebraic space, then it is not the case that every
equivalence class in $|X|$ admits a representative which is a monomorphism
(see
e.g.~\cite[\href{http://stacks.math.columbia.edu/tag/02Z7}{Tag
  02Z7}]{stacks-project}). 
However, under mild assumptions on $X$,
the points of $|X|$ do admit such representatives (which
are then unique up to unique isomorphism); this in particular
is the case if $X$ is quasi-separated.   See e.g.\
the discussion at the beginning
of~\cite[\href{http://stacks.math.columbia.edu/tag/03I7}{Tag 03I7}]{stacks-project}.  

If $\cX$ is genuinely an algebraic stack,
then it is not reasonable to expect the equivalence classes in $|\cX|$ 
to admit monomorphism representatives in general (even 
if $\cX$ is quasi-separated), since points of $\cX$ 
typically admit non-trivial stabilizers.  The aim of the theory
of residual gerbes is to provide a replacement for such representatives.

We recall the relevant
definition~\cite[\href{http://stacks.math.columbia.edu/tag/06MU}{Tag
  06MU}]{stacks-project}): 

\begin{adefn}\index{residual gerbe}
	Let $\cX$ be an algebraic stack.
	We say that the {\em residual gerbe} at a point $x \in |\cX|$ 
	exists if we may find a monomorphism
	$\cZ_x \hookrightarrow \cX$ such that $\cZ_x$ is reduced and
	locally Noetherian, $|\cZ_x|$ is a singleton,
	and the image of $|\cZ_x|$ in $|\cX|$ is equal to $x$.
	If such a monomorphism exists, 
        then $\cZ_x$ is unique up to unique isomorphism,
        and we refer to it as the {\em residual gerbe} at $x$.
\end{adefn}

In the case that $X$ is an algebraic space,
the residual gerbe $\cZ_x$ exists
at every point $x \in
|X|$~\cite[\href{http://stacks.math.columbia.edu/tag/06QZ}{Tag
  06QZ}]{stacks-project}, 
and is itself an algebraic space, called the {\em residual space} of
\index{residual space}
$X$ at~$x$~\cite[\href{http://stacks.math.columbia.edu/tag/06R0}{Tag
  06R0}]{stacks-project}. 
However, as already noted, in this case, if $X$ is furthermore
quasi-separated,
then the residual space at a point is simply the spectrum of a field.

If $\cX$ is an algebraic stack with quasi-compact diagonal
(e.g.\ if $\cX$ is quasi-separated),
then the residual gerbe exists at every point of
$|\cX|$~\cite[\href{http://stacks.math.columbia.edu/tag/06RD}{Tag
  06RD}]{stacks-project}). 

If $X$ is a scheme, then a {\em finite type point} of $X$
is a point $x \in X$ that is locally closed.  Equivalently,
these are the points for which the morphism $\Spec k(x) \to X$
is a finite type morphism, and may be characterized more abstractly
as those equivalence classes of morphisms $\Spec K \to X$ which
admit a representative which is locally of finite type.
This latter notion makes sense for an arbitrary algebraic stack,
and allows us to define the notion of a {\em finite type point} 
of an algebraic stack (see e.g.\ \cite[1.5.3]{EGstacktheoreticimages}).

We then have the following results, which provide analogues
of the topological characterization of the finite type points of a scheme
as being those points that are locally closed.

\begin{alemma}
	\label{lem:residual space immersion}
	Suppose that $X$ is an algebraic space,
	and that $x \in |X|$ is a finite type point with the property 
	that for any \'etale morphism $U \to X$ whose
	source is an affine scheme,  the fibre over $x$ 
	is finite.
	Then, if $Z_x$ denotes the residual space at $x$ in $X$,
	the canonical monomorphism $Z_x \hookrightarrow X$
	is an immersion.
\end{alemma}
\begin{proof}
	Consider the construction of $Z_x$ given
	in~\cite[\href{http://stacks.math.columbia.edu/tag/06QZ}{Tag
  06QZ}]{stacks-project}: we choose a surjective \'etale morphism
$U \to X$, form the union $$U' = \coprod_{u \in U \text{ lying
		over } x} \Spec k(u),$$
and then realize $Z_x \hookrightarrow X$ as a descent of the monomorphism
$U' \hookrightarrow U$.   In making this construction, we may replace
$U$ by any open subscheme containing $x$ in its image, 
and thus we may assume that $U$ is affine.  By assumption,
there are then only
finitely many points $u$ in $U$ lying over $x$,
these points are all finite type points of $U$,
and furthermore, none of these points are specializations of 
any of the
others~\cite[\href{http://stacks.math.columbia.edu/tag/03IM}{Tag
  03IM}]{stacks-project}.
Thus $U'\hookrightarrow U$ is in fact an immersion, and thus the same
is true of $Z_x \hookrightarrow X.$
\end{proof}

\begin{alemma}
	\label{lem:residual gerbe immersion}
	If $\cX$ is an algebraic stack whose diagonal is quasi-compact,
	if $x\in |\cX|$ is a finite type point,
	and if $\cZ_x$ is the residual gerbe of $\cX$ at $x$,
	then the canonical monomorphism
	$\cZ_x \hookrightarrow \cX$ is in fact an immersion.
\end{alemma}
\begin{proof}
	We prove this by examining the construction of
	$\cZ_x$ carried out in the proof 
        of~\cite[\href{http://stacks.math.columbia.edu/tag/06RD}{Tag
        06RD}]{stacks-project}. 
        The first step of the proof is to replace $\cX$
	by the closure of $x$ in $|\cX|$, regarded as a closed
	substack of $\cX$ with its induced reduced structure.
	The assumption that $\cX$ has quasi-compact diagonal
	implies that the morphism $\cI_{\cX} \to \cX$
	is quasi-compact, and thus that we may find a dense open substack
	$\cU$ of $\cX$ over which this morphism is flat and locally of
	finite presentation.
	Since $x$ is dense in $|\cX|$,
	we find that $x \in |\cU|$, and so replacing $\cX$
	by $\cU$, we may assume that $\cI_{\cX} \to \cX$ 
	is flat and locally of finite presentation.
        By~\cite[\href{http://stacks.math.columbia.edu/tag/06QJ}{Tag
        06QJ}]{stacks-project}, 
        this implies that $\cX$ is a gerbe over some algebraic space~$X$.
	If we let $Z_x$ denote the residual space at (the image in $|X|$ of) $x$
	in $|X|$, then $\cZ_x$ is obtained as the base-change of $Z_x$
	over the morphism $\cX \to X$.

	Note that, in the various reduction steps undertaken in the
	preceding argument, we replaced $\cX$ by an open substack
	of a closed substack; it thus suffices to verify
	that $\cZ_x \to \cX$ is an immersion after making these
	reductions.  Note that since immersions are monomorphisms,
	the diagonal of the stack obtained after these reductions
	are made is the base-change of the diagonal of the original
	stack $\cX$, and thus continues to be quasi-compact.
	Consequently, we may assume that we are in the case
	where $\cX$ is a gerbe over $X$.  Since $\cZ_x \hookrightarrow \cX$
	is then obtained as the base-change of $Z_x \hookrightarrow X,$
	it suffices to show that this latter morphism is an immersion.
	For this, if we take into account Lemma~\ref{lem:residual space
		immersion}, it suffices to show that $X$ can be chosen
	to be quasi-separated (as every \'etale morphism from an 
	affine scheme to a quasi-separated algebraic 
	space has finite fibres over every point of~$|\cX|$;
see the discussion at the beginning
of~\cite[\href{http://stacks.math.columbia.edu/tag/03I7}{Tag
  03I7}]{stacks-project}). 

If we examine the proof
        of~\cite[\href{http://stacks.math.columbia.edu/tag/06QJ}{Tag
        06QJ}]{stacks-project}, 
we see that $X$ is constructed as follows: We choose a smooth
surjective morphism $U \to \cX$ whose source is a scheme,
and write $R = U\times_{\cX} U,$ so that 
$\cX = [U/R]$.
We then factor the morphism $R \to U\times_{\Spec \Z} U$ 
through a morphism $R' \to U\times_{\Spec \Z} U$,
where $R' \to U\times_{\Spec \Z} U$
is a flat and locally of finite presentation equivalence relation,
and set $X = U/R'.$  
We note one additional aspect of the situation,
namely that the morphism $R\to R'$ is surjective.
The assumption that $\cX$ has quasi-compact diagonal then implies
that $R \to U\times_{\Spec \Z} U$ is quasi-compact, and since $R$ surjects
onto $R'$, we find that $R'\to U\times_{\Spec \Z} U$ is again quasi-compact.
Thus $X = U/R'$ is quasi-separated, as required.
\end{proof}

\begin{aexample}
	We note that quasi-separatedness (or some such hypothesis) is
	necessary for the truth of Lemma~\ref{lem:residual gerbe
	immersion}.  To illustrate this, let $G$ be an algebraic group
of positive dimension over an algebraically closed field $k$,
and let $X := G/G(k).$  Then $X$ contains a unique finite type
point $x$ --- namely, the equivalence
class of the monomorphism $\Spec k = G(k)/G(k) \hookrightarrow G/G(k)$
--- and this monomorphism realizes $\Spec k$ as the residual space
$Z_x$.  This monomorphism is {\em not} an immersion, 
since its pull-back to $G$ induces the monomorphism $G(k) \hookrightarrow
G,$ which is not an immersion.
\end{aexample}

Traditionally, in the theory of moduli problems, a family of 
some objects parameterized by a base scheme $T$ is called {\em
  isotrivial} \index{isotrivial}
if the isomorphism class of the members of the family is constant
over $T$.   From the view-point of morphisms to a moduli stack $\cX$,
this corresponds to the image of $T$ in $|\cX|$ being a singleton,
say $x$.  
We would like to conclude that the morphism $T \to \cX$
factors through the residual gerbe $\cZ_x$.
In practice, we often verify the ``constancy'' of the morphism $T \to \cX$
only at finite type points.
The following result gives sufficient conditions, under such a constancy
hypothesis
on the finite type points, for a morphism to factor through the residual
gerbe.

\begin{alemma}
	\label{lem:isotrivial}
	Let $\cX$ be an algebraic stack,
	and let $x \in |\cX|$ be a point for which
	the residual gerbe $\cZ_x$ exists,
	and for which the canonical monomorphism
	$\cZ_x \hookrightarrow \cX$ is an immersion.
	If $f:T \to \cX$ is a morphism
	whose domain is a reduced scheme,
	and for which all the finite type points of $T$
	map to the given point $x$,
	then $f$ factors through $\cZ_x$.
\end{alemma}
\begin{proof}
	Consider the fibre product $T \times_{\cX} \cZ_x$. 
	By assumption,
	the projection from this fibre product to $T$
	is an immersion 
	whose image contains every finite type point of $T$.
	A locally closed subset of $T$ that contains every finite type
	point is necessarily equal to $T$, and thus, since $T$ is
	reduced, this immersion is in fact an isomorphism.  
	Consequently the morphism $f$ factors through $\cZ_x$,
	as claimed.
\end{proof}

We end this discussion by explaining how the preceding discussion
generalizes to certain Ind-algebraic stacks.  Let $\{\cX_i\}_{i\in I}$
be a $2$-directed system of algebraic stacks, and assume that the 
transition morphisms are monomorphisms.  Let $\cX := \varinjlim_i \cX_i$
be the Ind-algebraic stack obtained as the $2$-direct limit of the $\cX_i$.
We may define the underlying set of points $|\cX|$ as equivalence classes
of morphisms from spectra of fields in the usual way, and since any
such morphism factors through some $\cX_i$, we find that 
$|\cX| = \varinjlim_i |\cX_i|$.  

Now suppose that $\cX$ has quasi-compact diagonal, or, equivalently (since
the transition morphisms are monomorphisms) that each $\cX_i$ has 
quasi-compact diagonal.  If $x \in |\cX|$, then $x \in |\cX_i|$ for 
some $i$, and the residual gerbe $\cZ_x$ at $x$ in $\cX_i$ exists.
The composite monomorphism
$\cZ_x \hookrightarrow \cX_i \hookrightarrow \cX_{i'}$
realizes $\cZ_x$ as the residual gerbe at $x$ in $\cX_{i'}$,
for any $i' \geq i$, and so we may regard $\cZ_x$ as being the 
residual gerbe at $x$ in $\cX$.
Lemmas~\ref{lem:residual gerbe immersion} and~\ref{lem:isotrivial}
immediately extend to the context of such Ind-algebraic stacks $\cX$.


%\chapter{Residual gerbes}\label{app:residual gerbes}
%
%% \section{Residual gerbes --- generalities} 
%% \label{subsec:gerbes}
%%\mar{ME: Maybe this should be moved to an appendix ?! ME: Great, now done,
%%removing this comment}
%% If we consider an $\Fbar_p$-valued point 
%% $x: \Spec \Fbar_p \to \cX_{d,\red},$
%% then we may base-change the residual gerbe to $\cX_{d,\red}$ at $x$
%% (which is a gerbe over $\F_p$) 
%% over $\Fbar_p$ via the composite $\Spec \Fbar_p \buildrel x \over
%% \longrightarrow \cX_{d,\red} \to \Spec \F_p$; equivalently,
%% we may regard $x$ as a point of the base-change $(\cX_{d,\red})_{/\Fbar_p}$, 
%% and then consider the residual gerbe of $(\cX_{d,\red})_{\Fbar_p}$ at
%% this point.
%% This residual gerbe is then of the form $[\Spec \Fbar_p/G]$,
%% for a finite type affine group scheme $G$ over $\Fbar_p$.
%% (The group $G$ may be described as the fibre
%% product $x \times_{\cX_{d,\red}} x$,
%% and hence its claimed properties follows from the
%% fact that the diagonal of $\cX_{d,\red}$ is affine and of finite type.)
%% If $\rhobar:G_K \to \GL_d(\Fbar_p)$ is the Galois representation
%% associated to $x$, then $G = \Aut_{G_K}(\rhobar).$
%
%
%We refer 
%to~\cite[\href{http://stacks.math.columbia.edu/tag/06MU}{Tag
%  06MU}]{stacks-project} for the notion of the residual
%gerbe at a point of an algebraic stack.
%It follows 
%from~\cite[\href{http://stacks.math.columbia.edu/tag/06RD}{Tag
%  06RD}]{stacks-project}
%that any point of a quasi-separated stack
%admits a residual gerbe.
%It also follows 
%from~\cite[\href{http://stacks.math.columbia.edu/tag/06G3}{Tag
%  06G3}]{stacks-project} that at if $x$
%is a finite type point of an algebraic stack $\cX$,
%then not only does the residual gerbe $\cZ_x$ exist,
%but the defining morphism $\cZ_x \hookrightarrow \cX$
%is locally of finite type.
%We wish to strengthen this result,
%under various additional hypotheses.
%
%
%\begin{alemma}\label{lem: definitions of locally finite presentation point} If $\cX$ is an algebraic stack,
%	and $x \in |\cX|$, then the following two 
%		conditions are equivalent:
%		\begin{enumerate}
%			\item There exists a morphism $\Spec k
%				\to \cX$ which is locally
%				of finite presentation and represents
%				$x$.
%			\item There exists a scheme $U$,
%				a closed point $u \in U$ for which
%				the natural morphism
%				$\Spec \kappa(u) \to U$ is
%				of finite presentation,
%				and a smooth morphism $U \to \cX$
%				taking $u$ to $x$.
%		\end{enumerate}
%	\end{alemma}
%        \begin{proof}% \mar{TG: add a proof; should be a simple
%            % modification of~\cite[\href{http://stacks.math.columbia.edu/tag/06FX}{Tag 06FX}]{stacks-project}.}
%          We follow the proof
%          of~\cite[\href{http://stacks.math.columbia.edu/tag/06FX}{Tag
%            06FX}]{stacks-project}. Firstly, if the second condition
%          holds, then the composite $\Spec\kappa(u)\to U\to\cX$ is a
%          composite of maps which are locally of finite presentation,
%          and is therefore also locally of finite presentation. Conversely,
%          suppose that $\Spec k\to\cX$ is locally of finite
%          presentation and represents~$x$, and let $U\to\cX$ be a
%          smooth cover by a scheme. Then $U\times_{\cX}\Spec k\to\Spec k$
%          is a surjective smooth morphism of algebraic spaces.
%          Choose an \'etale cover
%	  $V \to U\times_{\cX}\Spec k$ whose source is a scheme;
%	  then $V$ is a non-empty scheme,
%	  locally of finite presentation over $\Spec k,$
%	  and so we may find a point $v \in V$ of finite type over~$k$. 
%	  The image $u$ of $v$ in $U$ is then a finite type point
%	  of $U$, and indeed the morphism $\Spec \kappa(u) \to U$ 
%	  is of finite presentation.  Replacing $U$ by a sufficiently
%	  small open neighbourhood of $u$, we may assume
%	  that $u$ is in fact closed in $U$. By construction,
%	  the image of $u$ in $\cX$ is equal to $x$.
%	  \mar{TG: needs finishing ME: Done, hopefully}
%        \end{proof}
%
%
%\begin{adefn} If $\cX$ is an algebraic stack,
%	we say that a point $x \in |\cX|$ is {\em a point of
%		finite presentation}
%		if it satisfies the equivalent
%		conditions of Lemma~\ref{lem: definitions of locally finite presentation point}.
%	\end{adefn}
%
%	We recall that if $\cX$ is quasi-separated, then the underlying
%	topological space $|\cX|$ is sober \cite[Cor. 5.7.2]{MR1771927}.
%
%	\begin{alemma}
%		\label{lem:points of fp}
%		If $\cX$ is a quasi-separated algebraic stack,
%		and if $x$ is a 
%		point $\cX$ of finite presentation, then any
%		morphism $\Spec k \to \cX$ representing $x$ that
%		is locally of finite presentation is in fact of
%		finite presentation, and the point $x$ is locally
%		\mar{ME: This is the result that will let us do things
%			like prove that irreducible points are closed
%			in $|\cX_d|$ --- because they are locally closed
%			and admit no specializations. TG: I don't
%                        understand this - isn't the definition of a
%                        point being closed in $|\cX_d|$ just that it
%                        admits no specializations? ME: Maybe you're right;
%		I was thinking something that this comment was supposed
%to capture, but I probably got it a bit muddled.  I'll leave this for now
%just to remind myself that I thought that this could be useful for something.	}
%closed in $|\cX|$.
%	\end{alemma}
%	\begin{proof} Since $\Spec k$ is quasi-compact and $\cX$ is
%		quasi-separated, a standard argument with graphs
%		shows that $\Spec k \to \cX$ is quasi-compact.  
%		%If $\cX$ lies over the scheme $S$, then we may
%		If we consider the morphisms of algebraic spaces
%		$$\Spec k \to \Spec k \times_{\cX} \Spec k \to \Spec k
%		\times_{\Spec \Z} \Spec k,$$
%		then
%		the composite is a closed immersion %(since $\Spec k \times_S
%		%\Spec k$ is equal to $\Spec k \times_U \Spec k$,
%		%for any affine open subset $U$ of $S$ containing 
%		%the image of $\Spec k$, and 
%		(morphisms between affine
%		schemes are separated), and so in particular
%		is quasi-compact, and the second morphism
%		is (quasi-compact and) quasi-separated and representable
%		by algebraic spaces,
%		since it a base-change of the diagonal of $\cX$,
%		and $\cX$ is quasi-separated.  Thus the diagonal
%		of the second morphism (which is a base-change of the
%		second diagonal of $\cX$) is a quasi-compact monomorphism,
%		and so, again, a standard
%		argument with graphs shows that the first morphism
%		is a quasi-compact monomorphism, and so in particular
%		is quasi-compact and quasi-separated.   Thus
%		$\Spec k \to \cX$ is quasi-compact and quasi-separated
%		in addition to being locally of finite presentation,
%		and hence is indeed of finite presentation.
%
%		The image of a constructible set under a morphism
%		of finite presentation is again constructible,
%		and thus $\{x\}$ is a constructible subset of $|\cX|$.
%		Since it is a singleton, it is in fact locally closed
%		in $|\cX|$.
%	\end{proof}
%
%	\begin{alemma}
%         	\label{lem:gerbes of fp}
%%		If $\cX$ is an algebraic stack,
%		and if $x \in |\cX|$ is a point of finite
%		presentation, and if $\cZ_x$ is the residual
%		gerbe of $\cX$ at $x$,
%		then the defining morphism $\cZ_x \to \cX$
%		is locally of finite presentation.  If $\cX$
%		is quasi-separated, then it is furthermore 
%		a morphism of finite presentation.
%		If $\cX$ is in addition locally Noetherian,
%		then it is even an immersion of finite 
%		presentation.
%	\end{alemma}
%	\begin{proof}
%		The property of a morphism being locally of finite
%                presentation %\mar{TG: surely lfp, not fp? ME: Yes, fixed.}
%		may be checked \emph{fppf} locally on the source.
%                If we choose a representative $\Spec k \to \cX$
%		for $x$ which is locally of finite presentation,
%		then this morphism factors through a morphism
%		$\Spec k \to \cZ_x$~\cite[\href{http://stacks.math.columbia.edu/tag/06MW}{Tag
%  06MW}]{stacks-project}.  Since $\cZ_x \to \cX$ is a monomorphism,
%a standard graph argument shows that $\Spec k \to \cZ_x$ is 
%locally of finite presentation.  It is evidently surjective
%(since it induces an isomorphism of singleton underlying topological
%spaces), and it is flat (since this may be checked after base-changing
%to a cover of $\cZ_x$ by the spectrum of a field, and any morphism
%to the spectrum of a field is flat).   Thus $\Spec k \to \cZ_x$
%is an {\em fppf} cover,
%and its composite with the monomorphism $\cZ_x \hookrightarrow
%\cX$ is locally of finite presentation, by the choice of $\Spec k$.
%Thus $\cZ_x \to \cX$ is itself locally of finite presentation.
%
%If $\cX$ is futhermore quasi-separated, then Lemma~\ref{lem:points
%	of fp} shows that the morphism
%$\Spec k \to \cX$ is in fact quasi-compact, and since $\Spec k 
%\to \cZ_x$ is surjective, we see that $\cZ_x \hookrightarrow \cX$
%is quasi-compact.  This morphism is futhermore separated (being
%a monomorphism), and so in particular quasi-separated.  Thus $\cZ_x
%\hookrightarrow \cX$ is in fact of finite presentation in this case.
%It remains to show that if $\cX$ is 
%locally Noetherian, then this morphism is in fact an immersion.
%
%Let $\cT$ denote the underlying reduced locally closed substack of $\cX$
%corresponding to the locally closed subset $\{x\}$ of $\cX$.
%The morphism $\cZ_x \hookrightarrow \cX$ factors through $\cT$,
%and so replacing $\cX$ by $\cT$, we may assume that $\cX = \cT$,
%and thus that $\cZ_x \hookrightarrow \cX$ is a monomorphism
%inducing a bijection on underlying topological spaces.
%Since $\cX$ is locally Noetherian (by assumption), 
%it follows
%from~\cite[\href{http://stacks.math.columbia.edu/tag/06MR}{Tag
%  06MR}]{stacks-project} that this monomorphism is actually
%an isomorphism, as required.
%	\end{proof}


%\begin{alemma}
%	\label{lem:factoring through gerbes}
%	Suppose that $\cX$ is a locally Noetherian and quasi-separated
%	algebraic stack, 
%	and that $x \in |\cX|$ is a point of finite type {\em (}and
%	hence of finite presentation, by local Noetherianness{\em )}.
%	If $f:T \to \cX$
%	is a morphism whose source is a reduced scheme, which maps
%	each finite type point of $T$ to $x$, then $f$
%	factors through a morphism $T \to \cZ_x$.
%\end{alemma}
%\begin{proof}
%	Consider the fibre product $T \times_{\cX} \cZ_x$. 
%	Lemma~\ref{lem:gerbes of fp} shows that
%	the projection from this fibre product to $T$
%	is an immersion 
%	whose image contains every finite type point of $T$.
%	A locally closed subset of $T$ that contains every finite type
%	point is necessarily equal to $T$, and thus, since $T$ is
%	reduced, this immersion is in fact an isomorphism.  
%	Consequently the morphism $f$ factors through $\cZ_x$,
%	as claimed.
%\end{proof}  

% \chapter{Groups acting on categories fibred in groupoids}
% In this appendix we discuss notions related to the action of a group
% on a category fibred in groupoids.  Our main goal is to define the
% fixed point category of such an action, and to provide a description
% of this category, in the case of an action by a finitely presented
% group, which is sufficiently geometric to be useful in applications.
% Everything that we do is certainly
% well-known to experts in the theory of stacks
% and higher categories, but to reassure ourselves that everything
% works as we expect, we have found it convenient to write out the 
% basic ideas concretely in the particular setting that we will encounter
% in our applications.

% Our discussion uses some notions and terminology from
% the theory of $2$-categories.  For this, we follow the very down-to-earth
% treatment of this topic 
% in the Stacks Project~\cite[\href{http://stacks.math.columbia.edu/tag/0011}{Tag
%   0011}]{stacks-project}.  If the authors were more comfortable 
% with the homotopy-theoretic point of view on higher categories,
% then our discussion would surely be more streamlined; however, we have
% found it easiest to take the more down-to-earth approach.
% For example, at various points in our discussion
% we speak of diagrams of functors commuting on the nose, something
% which is probably not in the spirit of a more sophisticated treatment
% of higher categories.  Indeed, through this appendix, we adopt the
% convention that a commutative diagram of functors means a strictly
% commutative diagram, rather than a $2$-commutative one.
% (From a more sophisticated view-point, our approach could be viewed
% as an exercise in strictifying everything in sight.)

% \section{Group actions}
% If $\cC$ is a category fibred in groupoids,
% then $\Aut(\cC)$ is a $(2,1)$-category. 
% (It is a sub-$(2,1)$-category of the $(2,1)$-category of stacks.) 
% In fact, it is a particular kind of $(2,1)$-category:
% it contains a single object, namely $\cC$,
% and not only are the $2$-morphisms all isomorphisms,
% but the $1$-morphisms (i.e.\ the automorphisms of $\cC$) are weakly invertible,
% in that (by definition) they are equivalences of categories from $\cC$ to
% itself, and hence for each automorphism $\alpha: \cC \to \cC,$ there is another
% automorphism $\beta: \cC \to \cC$ such that $\alpha \circ \beta $  and $\beta \circ \alpha$ each are naturally equivalent to the identity
% (i.e.\ admit a $2$-morphism to $\id_{\cC}$, in the $2$-category $\Aut(\cC)$).
% Thus $\Aut(\cC)$ is an example of a $2$-groupoid with one object, which is 
% to say that it is an example of a $2$-group.

% \begin{aremark} Although the inverses in $\Aut(\cC)$ are merely weak,
% 	it is not as weak as some $2$-groups considered in the literature,
% 	insofar as composition is strictly associative (being simply
% 	composition of functors), and $\id_{\cC}$ is a strict identity.
% 	This will simplify the ensuing discussion, in comparison to
% 	the general theory.
% \end{aremark}

% If $G$ is a group, we may also regard it as a $2$-group, by regarding it
% as a category (i.e.\ a $1$-category) with a single object and with the
% elements of $G$ being the morphisms, which we then regard as a $(2,1)$-category
% by declaring the $2$-morphisms to consist just of identity morphisms.
% \mar{ME: Edit what follows to use $\star$ notation for ``horizontal''
% 	composition of $2$-morphisms, where appropriate.}
% A homomorphism $G \to \Aut(\cC)$ is then, by definition,
% a functor between these $2$-categories
% (what is sometimes called a {\em weak functor} or {\em pseudofunctor}).
% We spell out the details of this, in our context, in the following
% definition.
% (See e.g.~\cite[\href{http://stacks.math.columbia.edu/tag/003N}{Tag
%   003N}]{stacks-project}, where the notion of pseudofunctor between
% a category and a $2$-category is defined.)

% \begin{adefn}
% 	A homomorphism $(F,u,\alpha): G \to \Aut(\cC)$ consists of,
% for each $g \in G$, an isomorphism $F_g: \cC \to \cC$, together
% with a natural isomorphism\mar{TG: can we spell out exactly what this means?
%   I got bogged down in 2-categorical confusion when I tried to make
%   precise sense of this and of the question below about whether we
%   need the diagram with $F,F'$ swapped :-()} $u:  F_1 \cong \id_{\cC}$, and for each pair
% $(g,h) \in G\times G$ a natural isomorphism $\alpha_{g,h}: F_g\circ F_h \cong
% F_{gh}$, satisfying the following conditions:
% \begin{enumerate}
% 	\item If $g \in G$, then $\alpha_{g,1}$ is equal to
% 		$\id_{F_g} \circ u : F_g\circ F_1 \cong F_g,$ while
% 		$\alpha_{1,g}$ is equal to $u\circ \id_{F_g}: F_1 \circ F_g \cong F_g.$
% 	\item If $(g,h,k) \in G\times G \times G,$ then
% 	$$	\alpha_{gh, k} (\alpha_{g,h}\circ \id_{F_k}) = \alpha_{g,hk}
% 		(\id_{F_g} \circ \alpha_{h,k}).$$
% \end{enumerate}
% \end{adefn}

% %Since $G$ and $\Aut(\cC)$ are $2$-categories, the morphisms between them
% %form the objects of a $2$-category (in fact a $2$-groupoid), but we only
% %consider certain $1$-categorical aspects of the situation,
% %namely the $1$-morphisms between them;
% %these are the {\em equivalences} of the following definition.

% The morphisms between $2$-groups themselves have a higher
% categorical aspect.  Here, we only have need for the following particular
% notion.

% \begin{adefn}
% An {\em equivalence} between two morphisms
% $(F,u,\alpha), (F', u',\alpha'): G \to \Aut(\cC)$ consists
% of a collection of natural isomorphisms $\beta_g: F_g \cong F'_g$, one for each
% $g \in G$, such that for each $g,h \in G$ the following diagram commutes:
% \mar{ME: Check if we also need the diagram with $F$ and $F'$ swapped,
% 	or if that follows from this one.}
% $$
% \xymatrix{
% 	F_g \circ F_h \ar_-{{\id_{F_g} \circ \beta_h}}^-{\sim}[r]
% 	\ar_-{\alpha_{g,h}}^-{\sim}[d] & 
% 	F_g \circ F'_h \ar_-{{\beta_g \circ \id_{F'_h}}}^-{\sim}[r] &
% 	F'_g \circ F'_h \ar_-{{\alpha'_{g,h}}}^-{\sim}[d] \\
% 	F_{gh} \ar_-{\beta_{gh}}^-{\sim}[rr] && F'_{gh}}$$
% \end{adefn}\mar{TG: can we just have one arrow on the top for
%   $\beta_g\circ\beta_h$, or am I not understanding the definitions?}

% \begin{aremark}
% 	Equivalences, as defined above, are particular examples of what are
% 	called {\em strong transformation} between functors on bicategories in [Leinster].
% (We require the natural transformation $\sigma_{\cC}$ in the definition of
% transformation in [Leinster, \S 1.2] to be the identity.  Allowing more general
% choices of $\sigma_{\cC}$, as in that definition, would essentially mean
% that we regard morphisms $G \to \Aut(\cC)$ that are ``conjugate'' by a
% morphism in $\Aut(\cC)$ to be equivalent, and we don't want to do this.)
% 	\mar{ME: Fake citation alert. TG: it's a bit too vague for me
%           to know what it is! ``A survey of definitions of n-category''?}
% \end{aremark}

% \begin{aremark} Any morphism $G \to \Aut(\cC)$ is equivalent to one
% 	\mar{ME: This seems less important to me now than when I first wrote it, and we should check it's true if it stays here.}
% 	for which $F_1  = \id_{\cC}$ and $u$ is the identity.  
% \end{aremark}

% We now make the following definition.

% \begin{adefn} An {\em action} of a group $G$ on $\cC$ consists
% of a homomorphism (in the above sense) $F: G \to \Aut(\cC)$.
% \end{adefn}

% \begin{aexample} If $G$ is any group and $\cC$ is any category
% 	fibred in groupoids, we have the trivial action of $G$
% 	on $\cC$, defined by the trivial homomorphism
% 	$G \to \Aut(\cC)$ for which $F_g = \id_{\cC}$ for every
% 	$g \in G$.
% \end{aexample}

% \begin{aexample}
% 	\label{ex:duality of bundles}
% 	Maybe $\cE \mapsto \cE^{\vee}$ gives
% 	an action of the group of order two on the stack of rank $d$ vector
% 	bundles?
% \end{aexample}

% \section{Fixed points}
% Through this section, we suppose $G$ is a group acting on a
% category fibred in groupoids $\cC$, via the homomorphism
% $F: G \to \Aut(\cC).$

% \begin{adefn}
% A {\em fixed point} of this action consists of an object
% $x \in \cC$ together with a collection of isomorphisms
% $\imath_g: F_g(x) \to x$ such that:
% \begin{itemize} \item $\imath_1$ is the isomorphism induced by $u$;
% 	\item if $g, h \in G$, then $\imath_g \circ F_g(\imath_h) =
% 		\imath_{g h} \circ \alpha_{g,h}.$
% \end{itemize}
% A morphism of fixed points consists of a morphism of objects
% in $\cC$ that is compatible with the collections of isomorphism $\imath_g$.
% In this way the fixed points of the $G$-action form a category fibred in 
% groupoids that we denote by $\cC^G$.
% \end{adefn}

% \begin{aexample}
% 	\label{ex:fixed points of trivial action}
% 		If $G$ acts trivially on $\cC$, then giving an object
% 		of $\cC^G$ is the same as giving an object $x \in \cC$
% 		together with a homomorphism $G \to \Aut(x)$ (and a morphism
% 		in $\cC^G$ is a morphism in $\cC$ compatible with
% 		the $G$-action on source and target).  
% \end{aexample}

% \begin{aexample}
% 	\label{ex:fixed points of duality}
% 	If $G$ is the group of order two acting on $\Bun_d$ via duality
% 	of bundles, as in example~(\ref{ex:duality of bundles}),
% 	then $\cC^G$ is equivalent to the category of rank $d$-bundles
% 	equipped with a non-degenerate bilinear pairing.
% \end{aexample}

% In this last example, we see that $\cC$ is an algebraic stack,
% and the same is true of $\cC^G$.  We would like to show that this
% is a general phenomenon.
% In the context of Example~\ref{ex:fixed points of trivial action},
% i.e.\ when the $G$-action is trivial,
% if $\cC$ is an algebraic stack then  we can describe $\cC^G$ geometrically
% in terms of the inertia stack.  From such a description, it's
% easily seen that if $G$ is
% finitely presented, then $\cC^G$ is again
% an algebraic stack.  Our goal in what follows is to draw an analogous
% conclusion in the case of a possibly non-trivial action, by giving
% a description of $\cC^G$ in terms of various fibre products.


% Before doing this, we note two general properties of this fixed point
% construction.

% \begin{alemma}
% 	If $F$ and $F'$ are equivalent homomorphisms $G \to \Aut(\cC)$,
% 	then the fixed point categories for each of
% 	the induced actions on $\cC$ are equivalent.
% \end{alemma}
% \begin{proof}
% 	\mar{ME: Give this; just write down an explicit equivalence using
% 		the $\beta$'s.}
% \end{proof}

% \begin{alemma}
% 	\label{lem:fixed points preserve stacks}
% 	If $\cC$ is a stack in groupoids {\em (}over some base site{\em )}, 
% 	then $\cC^G$ is again a stack in groupoids.
% \end{alemma}
% \begin{proof}
% 	\mar{ME: We should think this through just to avoid blundering.}
% 	This follows immediately from the fact that (by assumption)
% 	objects and morphisms in $\cC$ satisfy descent.
% \end{proof}

% %\bigskip 
% %
% %*****************
% %
% %\smallskip
% %
% %What follows: * expressing $\cC^G$ in terms of fibre products and so on * , and 
% %then the crucial point that * we only to construct the $\imath_g$ for a 
% %set of generators, and to then verify the condition $\imath_r = u$ whenever
% %$u$ is a relation * .  (Also check that equivalent actions give equivalent fixed point categories.)
% %
% %\smallskip
% %
% %*****************
% %
% %\bigskip
% %
% %With an eye to our subsequent applications, we now give a more geometric 
% %description of $\cC^G$.

% To this end, for any element
% $g\in G$, we define $\cC^g$ to be the category fibred in groupoids
% whose objects are pairs $(x,\imath)$ with $x$ an object of $\cC$
% and $\imath: F_g(x) \iso x$ an isomorphism.  The morphisms are morphisms
% in $\cC$ that are compatible with the isomorphisms $\imath$.
% We begin by giving a description of $\cC^g$ as a $2$-fibre product.


% Let $\Gamma_g: \cC \to \cC \times \cC$
% (the product being taken over the base of $\cC$)
% \mar{ME: Check Stacks Project for correct terminology, since this
% 	is not a two-fibre product of categories fibred in groupoids,
% 	but rather an ``absolute'' product over the (fixed) base.}
% denote the graph of the
% functor $F_g$, i.e. $\Gamma_g(x) = \bigl(x,F_g(x)\bigr).$
% The $2$-fibre product 
% \mar{ME: Maybe this is a $2$-fibre product, in which case it would
% 	be good to say so, since I am trying to be pedantically precise 
% 	at this point.}
% $$\cC\times_{\Gamma_g, \cC\times\cC, \Delta} \cC$$
% then
% consists\footnote{It might be more correct to say that this is one possible
%        model of the $2$-fibre product, since the $2$-fibre product
% is only defined up to an appropriate equivalence.}      
% of tuples $(x,y,\alpha,\beta)$, with $x$ and $y$ being
% objects of $\cC$,
% and $\alpha$ and $\beta$ being isomorphisms
% $\alpha: x \iso y$ and $\beta: F_g(x) \iso y.$
% We may now define
% an equivalence of categories fibred in groupoids
% \anumequation
% \label{eqn:equivalence}
% \cC^g \iso \cC\times_{\Delta, \cC\times\cC, \Gamma_g} \cC
% \end{equation}
% via the functor
% $$(x,\imath) \mapsto (x, x, \id_x, \imath).$$
% We regard $\cC^g$ as lying over $\cC$ via the functor $(x,\imath) \mapsto
% x$.  The following diagram then evidently commutes:
% $$\xymatrix{ \cC^g \ar^-{\sim}_-{\text{(\ref{eqn:equivalence})}}[r]\ar[d]
% 		& \cC \times_{\Delta, \cC\times\cC,\Gamma_g} \cC\ar^{\quad
% 			\text{ proj.\ to $1$st factor}}[dl] \\
% 		\cC & }
% 	$$
% In the case when $g = 1$, we have a canonical functor
% \anumequation
% \label{eqn:unit section}
% \cC\to \cC^1,
% \end{equation}
% defined by $x \mapsto (x,u_x).$
% Composing the equivalence~(\ref{eqn:equivalence}) with this functor
% yields the fuctor
% \anumequation
% \label{eqn:composite unit section}
% x \mapsto (x,x,\id_x,u)
% \end{equation}
% from $\cC$ to $\cC_{\Gamma,\cC\times\cC,\Delta} \cC$.

% \begin{alemma}
% 	\label{lem:unit section is a monomorphism}
% 	The functor~{\em (\ref{eqn:unit section})} is fully faithful.
% \end{alemma}
% \begin{proof}
% 	The functor $(x,y,\alpha,\beta) \mapsto (x,y,\alpha,\beta\circ u)$
% 	induces an equivalence (indeed an isomorphism) of categories
% 	$\cC_{\Delta,\cC\times \cC,\Delta} \cC \iso \cC_{\Gamma_g,
% 		\cC\times\cC, \Delta},$
% 	whose composite with the double diagonal functor $\cC \to \cC_{\Delta,
% 		\cC\times\cC,\Delta} \cC$ is equal to the
% 	functor~(\ref{eqn:composite unit section}).
% 	Since the double diagonal functor is fully faithful, \mar{ME: I tried
% 		to find this statement in this level of generality in
% 		the Categories chapter of the Stacks Project, but failed.}
% 	so is~(\ref{eqn:composite unit section}), and thus so
% 	is~(\ref{eqn:unit section}).
% \end{proof}

% It will be helpful to extend the preceding definitions and constructions as 
% follows: if $(g_i)_{i \in I}$ is an $I$-tuple of elements of $G$  (indexed
% by some set $I$),
% then we define $\cC^{(g_i)_{i \in I}}$ to be
% the category fibred in groupoids whose objects are tuples
% $\bigl(x, (\imath_{i})_{i \in I})\bigr)$ consisting of
% an object $x \in \cC$ and an isomorphism $\imath_{i}: F_{g_i}(x) \iso x$
% for each $i \in I$, and whose morphisms are morphisms in $\cC$ compatible
% with the isomorphisms $\imath_{i}$.
% There is an obvious forgetful functor $\cC^{(g_i)_{i \in I}} \to \cC$
% 	via which we regard $\cC^{(g_i)_{i \in I}}$ as a category
% 	over $\cC$.
% 	If $I$ is a singleton, then we just recover the previously
% 	defined categories $\cC^g$. 

% If $I = \coprod_{j \in J} I_j$ is a partition
% of $I$,
% then there is an evident functor
% \anumequation
% \label{eqn:tuple equivalence}
% \bigl(x,(\imath_i)_{i \in I})\bigr) \to
% \bigl((x,(\imath_i)_{i \in I_j})\bigr)_{j \in J}
% \end{equation}
% from $\cC^{(g_i)_{i \in I}}$ to the $2$-fibre product $\bigtimes_{j \in J}
% \cC^{(g_i)_{i \in I_j}}$ over $\cC$,
% which is an equivalence.


% Now suppose that $g_1,g_2 \in G$.  We define a functor
% $$\mu_{(g_1,g_2)}:\cC^{(g_1,g_2)} \to \cC^{g_1g_2}$$
% as follows: on objects, it is given by 
% \mar{ME: Read through at some point to try to adopt definitions/directions
% 	of arrows/etc.\ so as to minimize the number of inverses that appear
% 	in formulas.}
% $$\bigl(x,(\imath_1,\imath_2)\bigr) \mapsto (x , 
% \imath_1 \circ F_{g_1}(\imath_2)\circ \alpha_{g_1,g_2}^{-1} ),$$
% while it is the identity on morphisms.

% \begin{alemma}
% 	If $g_1,g_2,g_3 \in G$, then the diagram 
% 	$$\xymatrix{ \cC^{(g_1,g_2,g_3)} \ar^-{\sim}_-{\text{\em (\ref{eqn:tuple
% 				equivalence})}}[d] 
% 		\ar[r]^-{\sim}_-{\text{\em (\ref{eqn:tuple equivalence})}}
% 		& \cC^{(g_1,g_2)}\times_{\cC} \cC^{g_3} \ar^-{\mu_{g_1,g_2}
% 			\times \id}[d] \\
% 		\cC^{g_1}\times_{\cC}\cC^{(g_2,g_3)} 
% 		\ar^-{\id\times \mu_{g_2,g_3}}[d] & \cC^{g_1 g_2}\times_{\cC}
% 		\cC^{g_3} \ar[d]^-{\mu_{g_1g_2, g_3}} \\
% 		\cC^{g_1}\times_{\cC} \cC^{g_2 g_3} \ar^-{\mu_{g_1,g_2g_3}}[r]
% 		&
% 		\cC^{g_1 g_2 g_3}}$$
% 	commutes.
% \end{alemma}
% \begin{proof}
% 	This follows from the cocycle condition on the $\alpha$'s.
% \end{proof}

% The diagram in the preceding lemma yields a functor
% $$\mu_{g_1,g_2,g_3}: \cC^{g_1,g_2,g_3} \to \cC^{g_1 g_2 g_3}.$$
% More generally, if we have an $m$-tuple $(g_1,\ldots, g_m)$ of
% elements of $G$, we obtain 
% a functor
% $$\mu_{g_1,\ldots,g_m} : \cC^{(g_1,\ldots,g_m)} \to \cC^{g_1\cdots g_m},$$
% well defined independently of the ``bracketing'' of the $n$-tuple into
% pairs (as follows by successive applications of the preceding lemma).
% Composing $\mu_{g_1,\ldots,g_m}$ with the equivalence
% $$\cC^{g_1}\times_{\cC} \cC^{g_2} \times_{\cC} \cdots \times_{\cC} \cC^{g_m}
% \iso \cC^{(g_1,\ldots,g_m)}$$
% coming from~(\ref{eqn:tuple equivalence}),
% we otain a functor
% $$\cC^{g_1}\times_{\cC} \cC^{g_2} \times_{\cC} \cdots \times_{\cC} \cC^{g_m}
% \to \cC^{g_1\ldots g_m}.$$

% \begin{alemma} Each of the composites
% 	$$\cC^g \cong \cC^g\times_{\cC} \cC 
% 	\buildrel\id\times \text{\em (\ref{eqn:unit section})}\over
% 	\hookrightarrow \cC^g \times_{\cC} \cC^1 \buildrel
% 	\mu_{g,1} \over \longrightarrow \cC^g$$
% 	and
% 	$$\cC^g \cong \cC \times_{\cC} \cC^g 
% 	\buildrel\text{\em (\ref{eqn:unit section})}\times \id \over
% 	\hookrightarrow \cC^1 \times_{\cC} \cC^g \buildrel
% 	\mu_{1,g} \over \longrightarrow \cC^g$$
% 	is equal to the identity. 
% \end{alemma}
% \begin{proof} This follows from the assumed relationship between each 
% 	of $F_{g,1}$ and $F_{1,g}$ and $u$.
% \end{proof}

% If $g \in G$, then we define an equivalence (indeed an isomorphism)
% of categories
% \anumequation
% \label{eqn:inverses}
% \cC^g \to \cC^{g^{-1}}
% \end{equation}
% \mar{ME: Fix this, using $\alpha$ and $u$ appropriately}
% via $(x,\imath) \mapsto (x,\imath^{-1})$.

% \begin{alemma}
% 	The diagram
% 	$$\xymatrix{
% 		\cC^g \times_{\cC} \cC^{g^{-1}} \ar^{\sim}_-{\text{\em
% 				(\ref{eqn:tuple equivalence})}}[r] \ar[d] &
% 		\cC^{(g_1,g_2)} \ar^-{\mu_{g,g^{-1}}}[d] \\
% 			\cC \ar^-{\text{\em (\ref{eqn:unit section})}}[r] &
% 			\cC^1}$$
% 			commutes.
% \end{alemma}

% If $g \in G$, then
% we have a diagonal functor
% $$\cC^g \to \cC^g \times_{\cC} \cC^g \cong \cC^{(g,g)}$$
% (the equivalence coming from~(\ref{eqn:tuple equivalence})).

% Fix an $n$-tuple $(g_1,\ldots,g_m)$ of elements of $G$.  
% Combining all the above constructions, we see that for any word $w$ in the
% $g_i$ (i.e.\ any element of the free group on $m$ generators),
% there is a well-defind functor
% $$\mu_w: \cC^{(g_1,\ldots,g_m)} \to \cC^{\overline{w}},$$
% where we have written $\overline{w}$ to denote the value of word $w$
% in the group $G$.

% Suppose now that $r_1,\ldots,r_n$ are $n$ given words, whose values in $G$
% is trivial (i.e.\ each $\overline{r}_i = 1$,
% i.e.\ the $r_i$ give valid relations between the $g_i$ in $G$).
% Then we obtain a functor
% $$\mu_{r_1}\times\cdots \times \mu_{r_n} : \cC^{(g_1,\ldots,g_m)}
% \to \cC^1\times_{\cC} \cdots \times_{\cC} \cC^1$$
% (with $n$ copies of $\cC^1$ in the target).
% Now taking the $n$-fold $2$-fibre product over $\cC$
% of~(\ref{eqn:unit section}) yields a
% fully faithful functor 
% $$\cC \hookrightarrow \cC^1\times_{\cC} \cdots \times_{\cC} \cC^1.$$
% Pulling back $\mu_{r_1}\times\cdots \times \mu_{r_n}$ along
% this fully faithful functor, we obtain a category fibred in groupoids
% that we denote by
% $\cC^{(g_1,\ldots,g_m); (r_1,\ldots,r_n)},$
% equipped with a fully faithful functor
% $$\cC^{(g_1,\ldots,g_m); (r_1,\ldots,r_n)} \hookrightarrow
% \cC^{(g_1,\ldots,g_m)}.$$

% \begin{aprop}
% 	\label{prop:fixed points as a fibre product}
% 	If $G$ is generated by $g_1,\ldots,g_m,$ with
% 	relations $r_1,\ldots,r_n,$ then the is an equivalence
% 	of categories fibred in groupoids
% 	$$\cC^G \cong \cC^{(g_1,\ldots,g_m); (r_1,\ldots,r_n)}.$$
% \end{aprop}
% \begin{proof}
% 	\mar{ME: Has to be written, or at least sketched.}
% \end{proof}

% When $G$ is finitely presented,
% the preceding proposition thus gives a description of $\cC^G$ in
% terms of various $2$-fibre product constructions.  We may thus
% deduce the following corollary.

% \mar{ME: We have to strengthen the following results to show that
% 	fixed point stacks inherit finiteness properties of the
% 	original stack.}

% \begin{acor}
% 	\label{cor:fixed points preserve algebraicity}
% 	If $\cC$ is an algebraic stack {\em (}over some base-scheme
% 	$S${\em )}, and if $G$ is finitely presented, then $\cC^G$ is an
% 	algebraic stack.
% \end{acor}
% \begin{proof}
% Proposition~\ref{prop:fixed points as a fibre product} shows
% that $\cC^G$ is equivalent to a category fibred in groupoids
% which can be constructed from $\cC$ by various $2$-fibre products of
% various morphism.  Since the $2$-fibre product of a pair 
% of algebraic stacks with respect to morphisms to a third algebraic
% stack is again an algebraic stack, we see that $\cC^G$ is an algebraic
% stack, as claimed.
% \end{proof}

% Actually, we will need a straightforward generalization
% of the preceding corollary to the case when $\cC$ is a certain
% kind of Ind-algebraic stack.

% \begin{acor}
% 	\label{cor:fixed points preserve Ind-algebraicity}
% \end{acor}
% \begin{proof}
% \end{proof}

% \begin{aremark}
% 	Our utilization of the various categories $\cC^{(g_1,\ldots,g_n)}$,
% 	and the morphisms between them,
% 	in the above discussion can probably be viewed as a particular instance
% 	of a general approach to equivariant homotopy theory via the bar
% 	construction.  (The first author thanks Nick Rozenblyum for a helpful 
% 	discussion related to the material of this appendix, in which this
% 	observation was made.)
% \end{aremark}
\chapter{Breuil--Kisin--Fargues modules and potentially semistable
  representations (by Toby Gee and Tong Liu)}\chaptermark{BKF modules
  and $G_K$-representations}\label{app: BKF pst}In this appendix we briefly
discuss the relationship between Breuil--Kisin--Fargues modules with
semilinear Galois actions, and potentially semistable Galois
representations. As explained in Remark~\ref{rem: probably don't need all choices of pi} below, we do not expect our results to be
optimal, but they suffice for our applications in the body of the
book. The results of this appendix were originally inspired
by~\cite{MR3127808}; the recent paper~\cite{2019arXiv190508555G}
corrects a mistake in~\cite{MR3127808} and independently proves related (and in some
cases stronger) versions of some of our results, by different
methods. We do not make any use of the arguments of
either~\cite{MR3127808,2019arXiv190508555G}, but instead
combine~\cite{1302.1888} with a result of Fargues (\cite[Thm.\
4.28]{2016arXiv160203148B}). 

% In keeping with the
% rest of the paper (but in contrast to the literature, e.g.\
% \cite{2016arXiv160203148B}) we only consider Breuil--Kisin--Fargues
% modules which are effective, in the sense that they
% are~$\varphi$-stable. This is harmless, as any Breuil--Kisin--Fargues
% module can be made effective by taking a sufficiently large Tate twist
% (in the sense of~\cite[Ex.\ 4.2.4]{2016arXiv160203148B}).
We use the notation introduced in the body of the book, in particular
in Section~\ref{subsec:rings}. % There is a natural ring homomorphism $\theta:\Ainf\to\cO_\C$ which for
% any~$x\in\cO_{\C}^\flat$ satisfies~$\theta([x])=x^\sharp$. \mar{ME: The
% notation $x^\sharp$ is not defined or used anywhere else,
% and I think we can just remove it here.  After all,
% $\theta$ is already defined, or at least recalled, in the main
% body of the text.  So maybe just write something like
% ``Recall the usual ring homomorphism $\theta: \dots$''.}
The kernel
of the usual ring homomorphism~$\theta:\Ainf\to\cO_\C$ is a principal ideal~$(\xi)$; one possible choice of~$\xi$
is $\mu/\varphi^{-1}(\mu)$, where ~$\mu=[\varepsilon]-1$ (for some
compatible choice of roots of unity
$\varepsilon=(1,\zeta_p,\zeta_{p^{2}},\dots)\in\cO_\C^\flat$). Recall that $\BdR^+ $ is the $\ker(\theta)$-adic completion of $\Ainf [\frac 1 p]$ and that $t:=\log[\varepsilon]\in \BdR^+$ is a generator of $\ker(\theta)$ in $\BdR^+$. %\mar{TL: note $t$ is not in $\Ainf$}

As usual, we let~$K$ be a finite extension of~$\Qp$ with residue
field~$k$. Recall that for each choice of uniformizer~$\pi$ of~$K$,
and each choice~$\pi^\flat\in\cO_\C^\flat$ of $p$-power roots of~$\pi$, we
write~$\gS_{\piflat}$ for~$\gS=W(k)[[u]]$, regarded as a subring
of~$\Ainf$ via $u\mapsto[\piflat]$. Write~$E_\pi(u)$ for the
Eisenstein polynomial for~$\pi$, and~$E_{\piflat}$ for its image in~$\Ainf$;
then $E_{\piflat}\in(\xi)$, because $\theta(E_{\piflat})=E_\pi(\pi)=0$.
Indeed $(E_{\pi^\flat}) = (\xi)$ (this follows for example from the
criterion given in~\cite[Rem.\ 3.11]{2016arXiv160203148B} and the
definition of an Eisenstein polynomial). % \mar{ME: I know it's probably standard,
% 	but is there a readily available citation for this? TG: how
%         about this?  ME: It took me a while to disentangle, but I guess
% I believe it :) }

In contrast to the body of the book, we do not use coefficients in
most of this appendix (we briefly consider $\cO$-coefficients at the end). Accordingly, we have the following definitions. 
\begin{adefn}\label{adefn: GK phi module}
	%\mar{ME: Shouldn't one still maintain the adjective ``\'etale'' here?}
  An \'etale $(\varphi,G_K)$-module is a finite free $W(\C^\flat)$-module $M$
  equipped with a $\varphi$-semilinear map $\varphi:M\to M$ which
  induces an isomorphism $\varphi^*M\isoto M$, together with a continuous
  semilinear action of~$G_K$ which commutes with~$\varphi$.
\end{adefn}
There is an equivalence of categories between the category of \'etale
~$(\varphi,G_K)$-modules $M$ of rank~$d$ and the category of
free~$\Zp$-modules~$T$ of rank~$d$ which are equipped with a
continuous action of~$G_K$. The Galois representation corresponding
to~$M$ is given by $T(M)=M^{\varphi=1}$.
(See Section~\ref{subsubsec: Galois reps for GK phi}.)
%(This is immediate from
%Proposition~\ref{prop: equivalences of categories to Ainf} and the
%results of Section~\ref{subsec:Galois reps}, but it is also easy to
%prove directly, since~$\C^\flat$ is an algebraically closed field of
%characteristic~$p$.) We write~$V(M):=T(M)\otimes_{\Zp}\Qp$.



% See Remark~\ref{rem: all BK BKF are effective} for a comparison of
% this definition to the notion of a Breuil--Kisin--Fargues module as
% defined in~\cite{2016arXiv160203148B}.
  
\begin{adefn}\label{adefn: BK module}
  Fix a choice of~$\pi^\flat$.  We define a Breuil--Kisin module of
  height at most~$h$ % with $A$-coefficients
  to be a finite free
  %\mar{ME: Shouldn't have ``$A$-coefficients'' here, right?}
  $\gS_{\pi^\flat}$-module~$\gM$ equipped with a~$\varphi$-semi-linear
  morphism $\varphi:\gM\to\gM$, with the property that the
  corresponding morphism % of $\AAinf{A}$-modules
  $\Phi_{\gM}:\varphi^*\gM\to\gM$ is injective, with cokernel killed
  by~$E_{\piflat}^h$.
\end{adefn}

\begin{adefn}
  \label{adefn: BKF module}A Breuil--Kisin--Fargues module of height
  at most~$h$ is a finite free $\Ainf$-module~$\gMt$ equipped with
  a~$\varphi$-semi-linear morphism $\varphi:\gMt\to\gMt$, with the
  property that the corresponding morphism
  $\Phi_{\gMt}:\varphi^*\gMt\to\gMt$ is injective, with cokernel
  killed by~$\xi^h$.
\end{adefn}

\begin{arem}  
If~$\gM$ is a Breuil--Kisin module we write~$\gMt$ for the
Breuil--Kisin--Fargues module $\Ainf\otimes_{\gS_{\pi^\flat}}\gM$
(this is indeed a Breuil--Kisin--Fargues module, because
$\gS_{\pi^\flat}\to\Ainf$ is faithfully flat, and
$(E_{\piflat})=(\xi)$). As noted in Remark~\ref{rem: we don't twist our
    embeddings by phi}, we are not twisting the embedding
  $\gS\to\Ainf$ by~$\varphi$, and accordingly,
  in Definition~\ref{adefn: BKF module},
  we demand that the
  cokernel of~$\varphi$ is killed by a power of~$\xi$, rather than a
  power of~$\varphi(\xi)$ as in~\cite{2016arXiv160203148B}.% \mar{TG:
    % need to be careful not to blunder somewhere with this!}

  Leaving this difference aside,
  our definition of a
  Breuil--Kisin--Fargues module is less general than that of
  ~\cite{2016arXiv160203148B}, in that we require~$\varphi$ to
  take~$\gM$ to itself; this corresponds to only considering Galois
  representations with non-negative Hodge--Tate weights. This
  definition is convenient for us, as it allows us to make direct
  reference to the literature on Breuil--Kisin modules. The
  restriction to non-negative Hodge--Tate weights is harmless in our
  main results, as we can reduce to this case by twisting by a large
  enough power of the cyclotomic character (the interpretation of
  which on Breuil--Kisin--Fargues modules is explained in~\cite[Ex.\
  4.24]{2016arXiv160203148B}). (We are also only considering free
  Breuil--Kisin--Fargues modules, rather than the more general
  possibilities considered in~\cite{2016arXiv160203148B}.)
\end{arem}

\begin{adefn}\label{adefn: BKF GK module}
  A Breuil--Kisin--Fargues $G_K$-module of height at most~$h$ is a Breuil--Kisin--Fargues
  module of height at most~$h$ which is equipped with a semilinear $G_K$-action which commutes
  with~$\varphi$.
\end{adefn}

\begin{arem}
Note that if~$\gMt$ is a Breuil--Kisin--Fargues $G_K$-module, then
$W(\C^\flat)\otimes_{\Ainf}\gMt$ is naturally an \'etale
%\mar{ME: ``\dots naturally a{\em n \'etale} $(G_K,\varphi)$-module \dots''?}
$(\varphi,G_K)$-module
in the sense of Definition~\ref{adefn: GK phi module}.
\end{arem}


%\begin{arem}
%  \label{arem: all BK BKF are effective}Our definition of a
%  Breuil--Kisin--Fargues module is less general than that of
%  ~\cite{2016arXiv160203148B}, in that we require~$\varphi$ to
%  take~$\gM$ to itself; this corresponds to only considering Galois
%  representations with non-positive Hodge--Tate weights. This
%  definition is convenient for us, as it allows us to make direct
%  reference to the literature on Breuil--Kisin modules. The
%  restriction to non-positive Hodge--Tate weights is harmless in our
%  main results, as we can reduce to this case by twisting by a large
%  enough power of the cyclotomic character (the interpretation of
%  which on Breuil--Kisin--Fargues modules is explained in~\cite[Ex.\
%  4.24]{2016arXiv160203148B}). (We are also only considering free
%  Breuil--Kisin--Fargues modules, rather than the more general
%  possibilities considered in~\cite{2016arXiv160203148B}.)
%\end{arem}

% \begin{adefn}
%   \label{adefn: BKF module with GK action}A Breuil--Kisin--Fargues $G_K$-module is
%   a Breuil--Kisin--Fargues module equipped with a continuous
%   semilinear action of~$G_K$ which commutes with~$\varphi$.
% \end{adefn}



Recall that for each choice of~$\piflat$ and each~$s\ge 0$ we
write~$K_{\pi^\flat,s}$ for $K(\pi^{1/p^s})$,
and~$K_{\pi^\flat,\infty}$ for~$\cup_sK_{\pi^\flat,s}$. 
\begin{adefn}
\label{adefn: descending BKF to BK}Let $\gMt$ be a
  Breuil--Kisin--Fargues $G_K$-module of height at most~$h$ with a semilinear $G_K$-action. Then
  we say that~$\gMt$ \emph{admits all descents}  if the following
  conditions hold.
  \begin{enumerate}
  \item\label{item: existence of descent Zp version} For every choice of~$\pi$ and~$\piflat$, there is a
    Breuil--Kisin module~$\gM_{\pi^\flat}$ of height at most~$h$ with
    $\gM_{\pi^\flat}\subset(\gMt)^{G_{K_{\pi^\flat,\infty}}}$
    for which the induced morphism
    $\Ainf\otimes_{\gS_{\pi^\flat}}\gM_{\pi^\flat} \to \gMt$
    is an isomorphism.
    %and
    %$\gMt=\Ainf\otimes_{\gS_{\pi^\flat}}\gM_{\pi^\flat}$.
  \item\label{appendix item: M mod u descends} The $W(k)$-submodule $\gM_{\piflat}/[\piflat]\gM_{\piflat}$ of
  $W(\overline{k})\otimes_{\Ainf}\gMt$ is independent of the
    choice of~$\pi$ and~$\piflat$.
  \item\label{appendix item: M mod E descends} The $\cO_K$-submodule $\varphi ^*\gM_{\piflat}/E_{\piflat}\varphi^*\gM_{\piflat}$
    of  $\cO_\C\otimes_{\theta, \Ainf}\varphi^*\gMt$  is independent of the
    choice of~$\pi$ and~$\piflat$.
  \end{enumerate}
% \mar{TL: we will see (2) is redundant TG: should we explain this
%   somewhere below? It doesn't seem completely trivial to extract from
%   the proof to me.}
\end{adefn}

\begin{arem}
  \label{arem: second condition is redundant}In fact
  condition~\eqref{appendix item: M mod u descends} in Definition~\ref{adefn:
    descending BKF to BK} is redundant; see
  Remark~\ref{arem: explanation that we didn't need a condition in BKF
    but that we want it for coefficients} below. However, we include
  the condition as it is useful when considering versions of the
  theory with coefficients and descent data.
\end{arem}



\begin{adefn}
  \label{adefn: crystalline descending BKF module}Let~$\gMt$ be a
  Breuil--Kisin--Fargues $G_K$-module which
  admits all descents. We say that~$\gMt$ is furthermore
  \emph{crystalline} if for each choice of~$\pi$ and~$\piflat$, and each~$g\in G_K$, we
  have \anumequation\label{appendix eqn: condition on BK GK for
    crystalline}(g-1)(\gM_{\piflat})\subset
  \varphi ^{-1}(\mu)[\piflat]\gMt.  \end{equation}% \mar{TG: hopefully correct but as
% ever need to be careful with~$\varphi$ twist - might be~$[\piflat]^p$?}\mar{TL : fixed the condition}
\end{adefn}


Definitions~\ref{adefn: descending BKF to BK} and~\ref{adefn:
  crystalline descending BKF module} are motivated by the
following result, whose proof occupies most of the rest of this
appendix.

\begin{athm}
  \label{athm: admits all descents if and only if semistable} Let~$M$
  be an \'etale $(\varphi,G_K)$-module. Then~$V(M)$ is semistable with
  Hodge--Tate weights in~$[0,h]$ if and only if there is a
  \emph{(}necessarily unique\emph{)}
  Breuil--Kisin--Fargues $G_K$-module $\gMt$ which is of height at
  most~$h$, which admits all descents, and which satisfies
  $M=W(\C^\flat)\otimes_{\Ainf}\gMt$.

  Furthermore, $V(M)$ is
  crystalline if and only if~$\gMt$ is crystalline.
\end{athm}
\begin{arem}
  \label{rem: probably don't need all choices of pi}
  As already noted in Remark~\ref{arem: second condition is redundant},
  condition~\eqref{appendix item: M mod u descends}
  of Definition~\ref{adefn: descending BKF to BK} is
  redundant, and it is plausible that
  condition~\eqref{appendix item: M mod E descends} is redundant as well
(i.e.\ that both condition
  condition~\eqref{appendix item: M mod u descends}
  and condition~\eqref{appendix item: M mod E descends} 
 in Definition~\ref{adefn: descending BKF to BK} are
  consequences of condition~\eqref{item: existence of descent Zp version}); but
  this does not seem to be obvious. Note though that it is not
  sufficient to demand the existence of a descent for a single choice
  of~$\pi$, as there are representations of finite height which are
  potentially semistable but not semistable (see for
  example~\cite[Ex.\ 4.2.1]{MR2558890}).
\end{arem}




We begin with some preliminary results. The following proposition and
its proof are due to Heng Du, and we thank him for allowing us to
include them here.

\begin{aprop}
  \label{aprop: BKF is de Rham if and only if GK basis mod E}Let~$\gMt$
  be a Breuil--Kisin--Fargues $G_K$-module with the property that
  $\C\otimes_{\theta , \Ainf}\varphi^*\gMt$ has a basis consisting of $G_K$-fixed 
  vectors. Let~$M=W(\C^\flat)\otimes_{\Ainf}\gMt$; then the
  $G_K$-representation ~$V(M)$ is de Rham.
\end{aprop}
\begin{arem}% \mar{TG: do we need to be careful about
    % $\varphi$-twists here?}
  \label{arem: de Rham implies existence of lattice}Using the
  equivalence of categories of~\cite[Thm.\ 4.28]{2016arXiv160203148B} (a theorem of Fargues),
  one can easily check that Proposition~\ref{aprop: BKF is de Rham if
    and only if GK basis mod E} admits a converse: namely that if~$M$
  is an \'etale $(\varphi,G_K)$-module with the property that $T(M)$ is a
  $\Zp$-lattice in a de Rham representation of~$G_K$ with non-negative
  Hodge--Tate weights, then there is a Breuil--Kisin--Fargues
  module~$\gMt$ with the properties in the statement of
  Proposition~\ref{aprop: BKF is de Rham if and only if GK basis mod
    E}.
\end{arem}
\begin{proof}[Proof of Proposition~\ref{aprop: BKF is de Rham if and
    only if GK basis mod E}]% We have a natural
  % identification  \[W(\C^\flat)\otimes_{\Ainf}\gMt=W(\C^\flat)\otimes_{\Zp}T(M)\]
  % and thus an
By~\cite[Thm.\ 4.28]{2016arXiv160203148B} (and our assumption that $\varphi(\gMt)\subseteq\gMt$), we have  injections % \mar{TG: explain that this comes from
    % $\phi^*\gMt\into\gMt$}
%     \mar{ME: How are we passing from $W(\C^\flat)$ to $\BdR^+$?  Wouldn't 
% 	    this require us to have descended to $\Ainf$ first? I vaguely
%     remember this as being the kind of thing that is discussed in 
% Fontaine's Grothendieck Festschrift article --- namely, that finite height
% modules have to map into $\Ainf$ rather than just $W(\C^\flat)$. ME: I think
% that the relevant material is in \S B.1.8 of Fontaine's article, although,
% as always, I find the notation and statements there a bit tricky to navigate.
% (See also Lemma~2.1 of {\tt padichodge.pdf}.) TG: I though this might
% be the quickest way to sort things out. ME: OK  --- removed a superscript ``$+$'' in the final expression, to be compatible with [BMS]; hopefully that's
% ok.}
  \[\BdR^+\otimes_{\Ainf}\varphi^*\gMt\hookrightarrow \BdR^+\otimes_{\Ainf}\gMt
  \hookrightarrow \BdR\otimes_{\Zp}T(M).\] To show
  that~$V(M)$ is de Rham, we need to show that the
   $\BdR$-vector space $\BdR\otimes_{\Zp}T(M)$ has a 
  basis consisting of $G_K$-fixed vectors,
  % \mar{ME: For clarity, maybe say ``basis consisting of $G_K$-fixed vectors''
  %         (and also above and below)? (At least, I think this
  % is what is meant.)}
  so it suffices to show that the $\BdR^+$-lattice
  $\BdR^+\otimes_{\Ainf}\varphi^*\gMt$ has a 
  basis consisting of $G_K$-fixed vectors. % \mar{TG: it is enough to show.}
  Since~$\BdR^+$ is a complete discrete valuation ring with maximal
  ideal~$(\xi)$, it is enough to show that there are compatible
bases of
  $\BdR^+/(\xi^n)\otimes_{\Ainf}\varphi^*\gMt$ for all~$n\ge 1$,  consisting of   $G_K$-fixed vectors.

  In the case~$n=1$, since $\xi$ generates the kernel of~$\theta$,
%~$\BdR^+/(\xi^n)=\C$,
we have such a
  basis by hypothesis. Suppose that~$\{e_i^{(n)}\}_{i=1,\dots,d}$ is a
  basis of $G_K$-fixed vectors for $\BdR^+/(\xi^n)\otimes_{\Ainf}\varphi^*\gMt$, and
  let~$\{\te_i^{(n+1)}\}_{i=1,\dots,d}$ be any basis of $\BdR^+/(\xi^{n+1})\otimes_{\Ainf}\varphi^*\gMt$
  lifting~$\{e_i^{(n)}\}$. Recall that~$t=\log[\varepsilon]\in \BdR^+$ is a
  generator of~$(\xi)$, and for each $g\in G_K$, write \[g\cdot
    (\te_1^{(n+1)},\dots,\te_d^{(n+1)})=(\te_1^{(n+1)},\dots,\te_d^{(n+1)})(1_d+t^nB_g^{(n+1)})\]where
  we can view~$B_g^{(n+1)}$ as an element of~$M_d(B^+_{\dR}/(\xi))=M_d(\C)$.

  A simple calculation shows that $g\mapsto B_g^{(n+1)}$ is a
  continuous 1-cocycle valued in $M_d(\C(n))$. Since~$n\ge 1$, the
  corresponding cohomology group $H^1(G_K,M_d(\C(n)))$ vanishes,
  because $H^1(G_K,\C(n))=0$
  by a theorem of Tate--Sen. There is therefore some $A\in M_d(\C)$ such
  that for all~$g$ we have \[B_g^{(n+1)}=\chi^n(g)g(A)-A,\]
where $\chi$ is the $p$-adic cyclotomic character.   Then \[
    (e_1^{(n+1)},\dots,e_d^{(n+1)})=(\te_1^{(n+1)},\dots,\te_d^{(n+1)})(1_d-t^nA)\]is
  the required  basis of
  $\BdR^+/(\xi^{n+1})\otimes_{\Ainf}\varphi^*\gMt$ consisting of
  $G_K$-fixed vectors
 lifting~$\{e_i^{(n)}\}$.  
\end{proof}

\begin{alem}
  \label{alemma: G K infty generate G K}There is no proper closed
  subgroup of~$G_K$ containing all of the subgroups~$G_{K_{\piflat,\infty}}$.
\end{alem}
\begin{proof}Let~$s$ be the greatest integer with the property
  that~$K$ contains a primitive $p^s$th root of unity; equivalently,
  it is the greatest integer with the property that~$K_{\piflat,s}/K$
  is Galois over~$K$ (for one, or equivalently every, choice
  of $\pi$ and~$\piflat$). Then~$K_{\piflat,s'}$ depends only on~$\pi$
if $s' \leq s$, and so we write~$K_{\pi,s'}=K_{\piflat,s'}$ for such~$s'$.

  The subgroups $G_{K_{\piflat,\infty}}$
  for any fixed choice of~$\pi$ topologically generate~$G_{K_{\pi,s}}$. It therefore
  suffices to show that as~$\pi$ varies, the various normal
  subgroups~$G_{K_{\pi,s}}$ of $G_K$
  collectively generate~$G_K$. Let $H$ be the normal subgroup that they generate;
  then for every uniformizer~$\pi$, $G_K/H$ is a quotient of $G_K/G_{K_{\pi,s}}$,  a cyclic group  of order~$p^s$.
  If $H$ were a proper subgroup of $G_K$, then necessarily $s \geq 1$,
  and the subgroups $G_{K_{\pi,1}}$ would coincide for every~$\pi$ (since they
  would coincide with  the unique index~$p$ subgroup of~$H$).
  %If~$s=0$, there is nothing to prove. Otherwise, since
  %the extensions~$K_{\pi,s}/K$ are cyclic of degree~$p^s$, we see
  %that if the groups~$G_{K_{\pi,s}}$ generate a proper subgroup of~$G_K$, then the
  Thus the
  extensions~$K_{\pi,1}/K$ would have to all coincide. % \mar{TG: Tong asks
  % if there is a clearer way to write this sentence. ME: I rewrote the argument
  % a little so I can follow it more easily.  The Kummer theory sentence seemed
  % pretty clear to me, though. TG: slightly rearranged, commenting out.}
By Kummer theory,
  this would imply that the ratio of any two uniformizers of~$K$ is a
  $p$th power in~$K$, which is nonsense (for example, consider the
  uniformizers~$\pi$ and~$\pi+\pi^2$).
  \end{proof}

\begin{proof}[Proof of Theorem~\ref{athm: admits all descents if and
    only if semistable}]We begin with the semistable case. If~$V(M)$ is semistable then the existence
  of~$\gMt$ is a straightforward consequence of the results
  of~\cite{1302.1888}. In particular, the existence of a unique~$\gMt$
  satisfying condition~\eqref{item: existence of descent Zp version} of
  Definition~\ref{adefn: descending BKF to BK} follows from ~\cite[Thm.\
  2.2.1]{1302.1888}. % \mar{TG: why do conditions~\eqref{item: M mod u
    %   descends} and~\eqref{appendix item: M mod E descends} of
    % Definition~\ref{adefn: descending BKF to BK} hold? I couldn't see
    % how to read this off from~\cite{1302.1888}.}
    
  The proofs that conditions~\eqref{appendix item: M mod u descends} and
  \eqref{appendix item: M mod E descends} hold are implicit in the proof of 
  \cite[Prop. 4.2.1]{1302.1888}, as we now explain. % which proved that $M_{\st , K}(T)$ and
  % $M_{\dr}(T) $ are independent of the choice of $\pi$ and
  % $\pi ^\flat$. For the convenience of readers, here we sketch the
  % proof of \eqref{appendix item: M mod u descends} and \eqref{item: M mod E
  %   descends}.
  Write
  $\overline{\gM}_{\pi ^\flat}= \gM_{\pi ^\flat}/ [\pi ^\flat]
  \gM_{\pi^\flat} $. Since~$W(\overline{k})\otimes_{W(k)}\overline{\gM}_{\pi
    ^\flat}=W(\overline{k})\otimes_{\Ainf}\gMt$ is independent
  of~$\piflat$, it suffices to show that the $K_0$-vector space
  $D_{\pi^\flat} :=\overline{\gM}_{\pi ^\flat}[\frac 1 p] $ is
  independent of $\pi^\flat$.
  % \mar{ME: ``Since \dots, it suffices \dots'' --- I don't follow
  %         what's going on here.  Are there notational errors,
  %         or am I just completely missing the point? TG: corrected
  %         notational error, hopefully it makes sense now?}
  Let $S_{\pi^\flat} $ be the $p$-adic
  completion of
  $\gS_{\pi ^\flat} [\frac{E_{\pi^\flat}^i}{i!}, i \geq 1]$. By
  \cite[Prop.\ 6.2.1.1]{BreuilGriffith} and  \cite[\S 2]{liulattice2},
  % \mar{TG:
    % can we be more precise? Is it actually \cite[Prop.\ 6.2.1.1]{BreuilGriffith}? 
        % TL: not exactly Breuil's result because there is no Kisin module involved in his paper. }
  $D_{\pi^\flat}$ admits a section
  $s_{\pi ^\flat}: D_{\pi ^\flat} \to S_{\pi ^\flat}[\frac 1 p]
  \otimes_{\varphi, \gS_{\pi ^\flat}}\gM_{\pi ^\flat}$ so that
\anumequation
\label{eqn:strongly divisible iso}
   S_{\piflat}[\frac 1 p] \otimes _{K_0}s_{\piflat}(D_{\piflat}) =
  S_{\pi ^\flat}[\frac 1 p] \otimes_{\varphi, \gS_{\pi
      ^\flat}}\gM_{\piflat}.  
\end{equation}
In
  particular, we may regard $s_{\pi^\flat} (D_{\pi ^\flat})$ as a
  submodule of
  $\Bcris^+ \otimes_{\varphi, \gS_{\piflat}} \gM_{\piflat} \subseteq
   \Bcris^+ \otimes_{\Z_p}
   T(M)$.

 Write $\gu = \log([\pi^\flat]) \in \Bst^+= \Bcris^+ [\gu]$. Note that
 the \emph{set} $\Bst^+$ does not depend on the choice of $\piflat$,
 because, if ${\varpi^\flat}$  is another choice and we write $\gu' =
 \log ([\varpi^\flat])$, then $\Bcris^+ [\gu'] = \Bcris^+
 [\gu]$. Indeed, we have  $\log([{\varpi}^\flat]) = \log([\piflat]) +
 \lambda$ with $\lambda= \log([\frac{\varpi^\flat }{\piflat}]) \in
 \Bcris^+$. In particular, $D_\st (T(M)) : = ( \Bst^+ \otimes_{\Z_p} T(M) )^{G_K} $ is independent of~ $\piflat$.
   Furthermore,
  \cite[Prop. 2.6]{liulattice2} shows that we have a commutative diagram
% \mar{ME: I think the second horizontal arrows are also embeddings,
% and I've retypeset them to indicate this.}
   (where $i_{\piflat}$ is a $K_0$-linear isomorphism)
    $$\xymatrix{D_\st (T(M)) \ar[d]^-{i_{\piflat}}_\wr \ar@{^{(}->}[r]
      &   \Bst^+\otimes_{K_0}D_\st (T(M))\ar[d]^{\mod \gu} \ar@{^{(}->}[r] &
      \Bst^+  \otimes _{\Z_p} T(M) \ar[d]^{\mod \mathfrak u}
      \\s_{\piflat}(D_{\piflat}) \ar@{^{(}->}[r] & \Bcris^+
      \otimes_{\varphi, \gS_{\piflat}} \gM_{\piflat} \ar@{^{(}->}[r]&  \Bcris^+
      \otimes _{\Z_p} T(M)}$$    
 To show \eqref{appendix item: M mod u descends}, it therefore suffices to show that the
 image of the composite $$ \resizebox{\textwidth}{!}{\xymatrix {D_\st (T(M)) \to
   \Bst^+\otimes_{K_0} D_\st (T(M)) \ar[r]^-{\mod \gu\ \ } &
   \Bcris^+ \otimes_{\varphi, \gS_{\piflat}} \gM_{\piflat} \to W(\bar
   k)[1/p] \otimes_{\varphi, \Ainf} \gMt}} $$ is independent of
 $\piflat$. Here the last map is induced by $\nu : \Bcris^+ \to
 W(\bark)[\frac 1 p]$ which extends the natural projection $\Ainf \to
 W(\bark)$. It suffices in turn to show that the composite
% \mar{ME: ``It suffices in turn \ldots'' --- it  seemed to me that
% this involved applying the second horizontal arrows in the above
% diagram,  where $T(M)$  is factored out, which is why I wanted to
% indicate that they are embeddings.}
 $\xymatrix{\Bst^+ \ar[r]^-{\mod \gu} & \Bcris^+ \ar[r]^-\nu &
   W(\bark ) [\frac 1 p]}$ is independent of $\piflat$. This follows
 from the fact that $ \lambda= \log([\frac{\varpi^\flat }{\piflat}])$
 is in $\ker (\nu)$ (see the proof of~\cite[Lem.\ 2.10]{liulattice2}). 
 
 To prove \eqref{appendix item: M mod E descends}, it again suffices to prove
 the statement after inverting~$p$. For any subring $A \subset
 \BdR^+$, set $F^1 A = A \cap \xi \BdR^+$. It is easy to see that $F^1 \gS_{\piflat}= E_{\piflat}\gS_{\piflat}$ and also that $S_{\piflat}/ F^1 S_{\piflat} = \cO_K$.
(Note that
the inclusion $E_{\piflat} S_{\piflat} \subseteq F^1 S_{\piflat}$ is strict.) Now we again use the isomorphism~\eqref{eqn:strongly divisible iso}.
%$
% S_{\piflat}[\frac 1 p]\otimes _{K_0}s_{\piflat}(D_{\piflat})  =
% S_{\pi ^\flat}[\frac 1 p] \otimes_{\varphi, \gS_{\pi
%     ^\flat}}\gM_{\piflat}$. 
By reducing modulo $F^1 S_{\piflat}$ on the both sides of this identification, 
%and using this equality,
we conclude that $$K \otimes _{K_0} D_{\piflat} =  K \otimes_{K_0}(s_{\piflat}(D_{\piflat}) \mod [\piflat]) =  \varphi^* \gM_{\piflat} / E_{\piflat} \varphi^* \gM_{\piflat} [\frac 1 p].  $$
On the other hand, by tensoring $\Ainf$ via $\gS_{\piflat}$ to~\eqref{eqn:strongly
divisible iso},
%the isomorphism  $
% S_{\piflat}[\frac 1 p]\otimes _{K_0}s_{\piflat}(D_{\piflat})  =
% S_{\pi ^\flat}[\frac 1 p] \otimes_{\varphi, \gS_{\pi
%     ^\flat}}\gM_{\piflat}$,
we obtain the isomorphism  $\Bcris^+ \otimes _{K_0} s_{\piflat}(D_{\piflat})  =
      \Bcris^+ \otimes_{\Ainf }\varphi^*\gMt$.  Modulo $F^1 \Bcris^+$
      on both sides, a similar argument to the above shows that 
      $ \C \otimes_{K_0} D_{\piflat} = (\varphi^*\gMt / \xi \varphi^* \gMt) [\frac 1 p] =\C \otimes_{W{(\bark)}}\overline{\varphi^*\gMt}
       $, where we write $ \overline{\varphi^*\gMt}
 =W(\bark)  \otimes_{\Ainf}\varphi^* \gMt
 $. In summary we see that % \mar{TG: I
   % don't understand these equalities, can you elaborate?! Also should
   % $p$ be inverted on the rightmost term?} 
 \[\resizebox{\textwidth}{!}{$(\varphi^*\gM_{\piflat}/ E_{\piflat} \varphi^*\gM_{\piflat})[\frac 1 p ] = K \otimes_{K_0}D_{\piflat} \subseteq  \C\otimes _{W(\bar k)}\overline{\varphi^*\gMt}  = (\varphi^*\gMt/ \xi \varphi^*\gMt) [\frac 1 p].$}\]  
 Since we have shown that $D_{\piflat} \subseteq \overline{\varphi^*\gMt}$ is independent of
 $\piflat$, it follows that $\varphi^*\gM_{\piflat}/ E_{\piflat}
 \varphi^*\gM_{\piflat}$ is independent of $\piflat$, as required.
    
    

  Suppose conversely that we are given~$\gMt$ as in Definition~\ref{adefn: descending BKF to BK}. Write~$\overline{\gM}$
  for the $W(k)$-module of Definition~\ref{adefn: descending BKF to
    BK}~\eqref{appendix item: M mod u descends}, and $\overline{\gM}'$
  for the $\cO_K$-module of Definition~\ref{adefn: descending BKF to
    BK}~\eqref{appendix item: M mod E descends}.   Since for
  each~$\piflat$ the action of~$G_{K_{\piflat,\infty}}$
  on~$\gM_{\piflat}$ is trivial by assumption, it follows from
  Lemma~\ref{alemma: G K infty generate G K} that   $G_K$ acts
  trivially on $\overline{\gM}$ and~$\overline{\gM}'$.

  Since $\C\otimes_{\Ainf,\theta}\varphi^*\gMt=\C\otimes_{\cO_K}\varphi^*\overline{\gM}'$,
% \mar{ME: I added  what I think was a missing overline on the RHS of
%   this equality. TG: moved the overline inside the Frobenius pullback;
% commenting out.}
 we see
  that the hypotheses of Proposition~\ref{aprop: BKF is de Rham if and
    only if GK basis mod E} are satisfied, so that
  $V(M)$ is de Rham, and consequently potentially
  semistable. To show that $V(M)$ is semistable, we
  may replace~$K$ with an (infinite) unramified extension and assume that~$K$ is
  a complete discretely valued field with residue field $k=\bark$. In
  particular, we now have $\overline{\gM} = W(\bark) \otimes_{\Ainf}
  \gMt $, and we claim that $\overline{\gM}[\frac 1 p] $ with its $G_K$-action is isomorphic to $D_{\pst}(V(M))$.  
  If the claim holds,  then since $G_K$ acts trivially on~$\overline{\gM}$, $G_K$ acts
  trivially on~$D_{\pst}(V(M))$, so
  that~$V(M)$ is semistable, as required.
  
  We now prove the claim. Let $L /K$ be a finite Galois extension so that
  $V(M)|_{G_L}$ is semi-stable, and let $\gM_L$ be the Breuil--Kisin
  module attached to $T(M)|_{G_L}$ for some choice of~$\piflat_L$. By~ \cite[(2.12)]{liulattice2}, 
  $ \Ainf\otimes_{\gS_{\piflat_L}}\varphi^*\gM_L  $ injects into $\Ainf
  \otimes_{\Z_p}T(M)$. Furthermore, by~\cite[Lem.\ 2.9]{liulattice2}, $\varphi^*{\gM}^{\inf}_L :=
  \Ainf \otimes_{\gS_{\piflat_L}}\varphi^*\gM_L   $ is stable under the
  $G_K$-action on $ \Ainf\otimes_{\Z_p}T(M)$, and by~\cite[Cor.\ 2.12]{liulattice2}, $ W(k) \otimes_{\gS_{\piflat_L}}\varphi^*\gM_L = 
  W(k) \otimes_{\Ainf}\varphi^*{\gM}^{\inf}_L $ (recall that $k = \bark$) together with its $G_K$-action
  is isomorphic to a $W(k)$-lattice inside $D_{\st, L} (T(M))= (\Bst^+
  \otimes_{\Z_p} T(M)) ^{G_L} $. In summary, $W(k)\otimes_{\Ainf}\varphi^*
  {\gM}^{\inf}_L [1/p] $
%\mar{ME: I added what I think was a missing $[1/p]$ here.}
is isomorphic to $D_{\st, L} (T(M))$ as $G_K$-modules.

  It remains to show that $\varphi
  ^*{\gM}^{\inf}_L = \varphi ^*\gMt$. By
  \cite[Thm.\ 4.28]{2016arXiv160203148B}, it suffices to show that
  $ \BdR^+ \otimes_{\Ainf}\varphi^*{\gM}^{\inf}_L= \BdR^+\otimes_{\Ainf}  \varphi^* \gMt
  $.  By the proof of
    Proposition~\ref{aprop: BKF is de Rham if and only if GK basis mod
      E}, we see that
    \anumequation\label{eqn: first check on L Minf}
      \BdR^+\otimes_{\Ainf}\varphi^* \gMt = \BdR^+ \otimes_K
      D_\dR(T(M)).
    \end{equation}
Since $T(M)$ is semi-stable over
   $L$, it is well-known that \anumequation\label{eqn: second check on L Minf}
\Bst^+   \otimes_{\gS_{\piflat_L}} \varphi^*\gM_L = 
\Bst^+   \otimes_{W(k _L)} D_\st (T(M)|_{G_L}).\end{equation} (This can be easily seen from the proof
   of \cite[Cor.1.3.15]{KisinCrys}, or see \cite[\S 2]{liulattice2}
   for a more detailed discussion; the key point,  in terms of  the
   diagram  above, is that \cite[\S 2]{liulattice2} shows that $K_0 [\gu] \otimes_{K_0}D_\st (T(M)) =
   K_0 [\gu] \otimes_{K_0}s_{\piflat} (D_{\piflat})$.) % , hence $ \varphi^*\gM_{\piflat} \otimes_{ \gS_{\piflat} } \Bst^+ = D_\st (T(M)) \otimes_{W(k)} \Bst^+$.) 

From~\eqref{eqn: second check on L Minf} we obtain    % In summary, over $L$ we have $
 \anumequation\label{eqn: third check on L Minf} \BdR^+\otimes_{ \gS_{\piflat_L} }\varphi^*\gM_{L} = \BdR^+ \otimes_{L}D_\dR (T(M)|_{G_L}).  \end{equation}
Comparing~\eqref{eqn: first check on L Minf} and~\eqref{eqn: third
  check on L Minf} we have  $ \BdR^+ \otimes_{\Ainf}\varphi^*{\gM}^{\inf}_L= \BdR^+\otimes_{\Ainf}  \varphi^* \gMt
  $, as required.
% \Bst^+\otimes_{\gS_{\piflat_L}}\varphi^*\gM_L  = 
%  \Bst^+ \otimes_{W(k _L)}D_\st (T(M)|_{G_L})
%$. % and   $\BdR^+ \otimes_{\Ainf}  \varphi^*{\gM}^{\inf}_L = \BdR^+ \otimes_{L}  D_\dR(T(M)|_{G_L}) $ as required.  
 %  we can apply Theorem ~\ref{acor: admits all descents if and only if potentially semistable} and Proposition ~\ref{aprop: BKF is de Rham if and only if GK basis mod E} to $T(M)|_{G_L }$ and $\gM_L$. More precisely,  by necessary part of  Theorem ~\ref{acor: admits all descents if and only if potentially semistable}, we have  $\varphi^*\gM_L/ E_{\piflat_L} \varphi^*\gM_L \subseteq \varphi ^* {\gMt}_L / \xi \varphi ^* {\gMt}_L$ is independent of the choice of $\piflat_L $. Consequently Lemma \ref{alemma: G K infty generate G K} shows that $\varphi ^*\gM_L/ E_{\piflat_L} \varphi ^*\gM_L $ is a $G_L$-stable $\O_L$-lattice inside $\C \otimes_{\theta , \Ainf} \varphi^*{\gMt} _L$. Then the proof of Proposition \ref{aprop: BKF is de Rham if and only if GK basis mod E} shows that     $\BdR^+ \otimes_{\Ainf}  \varphi^*{\gM}^{\inf}_L = \BdR^+ \otimes_{L}  D_\dR(T(M)|_{G_L}) $. 
 % \mar{TG: I don't think we can
 %    cite this as it is assuming crystalline rather than
 %    semistable. On the other hand the result is presumably still fine
 %    with the same proof. Another possibility is to cite \cite[\S
 %    8.5]{2017arXiv171006145C}. But actually why can't we cite
 %    Proposition~\ref{aprop: BKF is de Rham if and only if GK basis mod
 %      E} twice, once over~$K$ and once over~$L$?  
    %        TL : fix it.}
    %        By the proof ofProposition~\ref{aprop: BKF is de Rham if and only if GK basis modE}, we se%e that $ \BdR^+\otimes_{\Ainf}\varphi^* \gMt  = \BdR^+   \otimes_K D_\dR(T(M)) $.
  % Now  $\BdR^+ \otimes_{\Ainf} \varphi^*{\gM}^{\inf}_L
  % = \BdR^+ \otimes_{\Ainf}  \varphi^* \gMt $
  % and hence $\varphi^* \gM^{\inf}_L= \varphi ^* {\gMt}$, as required.
  
  
Finally, we turn to the crystalline case. Given the above, the result
  is a consequence of~\cite[Thm.\ 3.8]{Ozeki2018}, as we
 now explain.  In our case, $f(u)$ in~\cite{Ozeki2018} is equal to
 $u^p$, so  the assumptions of~\cite[Thm.\ 3.8]{Ozeki2018} are
 automatically satisfied (as noted at the beginning of~\cite[\S
 3.2]{Ozeki2018}).  More precisely, fix some~$\piflat$, and let $\widehat K$ be the Galois
 closure of~$K _{\piflat}$. There is a subring
 $\widehat{\mathcal R} \subseteq \Ainf$ (constructed in~\cite{Ozeki2018}) such that $\gS_{\piflat} \subset
 \widehat{\mathcal R}$ and $G_K$ acts on $\widehat{\mathcal R}$ through
% \mar{ME: Is this really a correct description of $\widehat{G}$? TG:
%   no! Corrected, commenting out.}
 $\widehat G:=\Gal(\cup_{n\ge 1}K(\zeta_{p^n},\pi^{1/p^n})/K)$.
When $T(M)$ is crystalline, ~\cite[Thm.\
 3.8]{Ozeki2018} shows that $T(M)$ admits a $(\varphi, \widehat
 G)$-module, which by definition consists of the following data: 
  \begin{itemize}
  \item The Breuil-Kisin module $\gM_{\piflat}$ attached to $T(M)|_{G_{{\piflat},\infty}}$.
  \item A $\widehat G$-action on $\widehat{\gM}: = \widehat{\mathcal
      R} \otimes_{\gS_{\piflat}} \varphi ^* \gM_{\piflat}$ which
    commutes with the action of~$\varphi$ and satisfies~\eqref{appendix eqn: condition on BK GK for
    crystalline}.
  \item $(W(\C^\flat) \otimes_{\widehat{\mathcal R}} \widehat{\gM})^{\varphi=1} = T(M)$ as $G_K$-modules.  
  % \item The $\widehat G$-action on  $\widehat{\gM}$ satisfies condition \eqref{appendix eqn: condition on BK GK for
  %   crystalline}
  \end{itemize}(In fact \cite{Ozeki2018} uses contravariant functors,
  which can be easily translated to the covariant functors used here.)
 This proves that if $T(M)$ is crystalline then $\gMt$
 satisfies~\eqref{appendix eqn: condition on BK GK for
    crystalline}. Conversely, if $\gMt$ satisfies~\eqref{appendix eqn: condition on BK GK for
    crystalline} for \emph{one} fixed $\piflat$, then
 \cite[Lem.\ 3.15]{Ozeki2018} shows that $s_{\pi}(D_{\piflat}) \subset
 (\Bcris^+ \otimes_{\gS_{\piflat}} \varphi^*
 \gM_{\piflat})^{G_K}$. Hence $T(M)$ is crystalline, as required. \end{proof}
  %if $p>2$ then the necessity of the
  %condition is~\cite[Prop.\ 5.9]{GLSII}, while its sufficiency is~\cite[Thm.\
  %21]{1206.4751}.\mar{TG: need an argument in the case $p=2$.}
% \mar{TL: explain why (2) is redundant}
\begin{arem}\label{arem: explanation that we didn't need a condition
    in BKF but that we want it for coefficients}
  Note that condition~\eqref{appendix item: M mod u descends} in Definition \ref{adefn: descending BKF to
    BK} is redundant for proving Theorem \ref{athm: admits all
    descents if and only if semistable}. %  \mar{ME: Why is~(2) imposed
    % then, rather deduced as a consequence? TG: maybe I'm blundering,
    % but I think that while~(2) follows from the other conditions in
    % this case (of $\Zp$-coefficients), it isn't clear why it follows
    % for general coefficients. Maybe we can reduce to this case?! But
    % in general I wanted to include it in the definition in order to be
    % able to define the stacks of fixed inertial type. TG: added
    % material on this commenting out.}
  In fact, if we
  remove~\eqref{appendix item: M mod u descends} from Definition~\ref{adefn: descending BKF to BK} then
  the necessity part of Theorem~\ref{athm: admits all descents if and
    only if semistable} of course still holds. For the sufficiency,
  note that in the above proof, we extended $K$ so that the residue
  field is $\bar k$; so we only use that $G_K$ (the inertia subgroup
  in this situation) acts trivially on
% \mar{ME: Could we use $\overline{\gMt}$ rather than $M$ here,
% since $M$ has quite a different connotation? TG: it seems to me that
% we don't need any notation at all, since we're only using it
% once. Commenting out.}
  $W(\bar k) \otimes _{\Ainf} \gMt$, and we do not need that
  $\gM_{\piflat}/ [\piflat] \gM_{\piflat} \subseteq W(\bar k) \otimes _{\Ainf} \gMt $ is independent
  of $\piflat$.

However, it is not immediately clear that the analogous condition is
redundant in the version of the theory with coefficients that we
consider in the body of the book (see Definition~\ref{defn:
  descending BKF to BK coefficients}), and this condition is used in defining our moduli stacks of potentially semistable
representations of given inertial type, so we include it here.
\end{arem}

  
  
  

\section{Potentially semistable representations}
It will be convenient for us to have a slight refinement of these
results, allowing us to discuss potentially semistable (and
potentially crystalline) representations. To this end, fix a finite Galois
extension~$L/K$. % \mar{TG: now redo the above in that framework. Is it
  % literally just demanding that there is something $G_K$-stable for
  % the data over~$L$?}

\begin{adefn}
  \label{adefn: potentially admits descents}Let $\gMt$ be a
  Breuil--Kisin--Fargues $G_K$-module of height at most~$h$. Then we say that~$\gMt$ \emph{admits all
    descents over~$L$} if the corresponding Breuil--Kisin--Fargues
  $G_L$-module (obtained by restricting the $G_K$-action on~$\gMt$ to~$G_L$) admits all descents (in the sense of
  Definition~\ref{adefn: descending BKF to BK}).%  \mar{TG: slightly
    % confused as to whether this is the correct thing to do.}
\end{adefn}

\begin{acor}
  \label{acor: admits all descents if and only if potentially semistable} Let~$M$
  be an \'etale $(\varphi,G_K)$-module. Then~$V(M)|_{G_L}$ is semistable with
  Hodge--Tate weights in~$[0,h]$ if and only if there is a \emph{(}necessarily unique\emph{)}
  Breuil--Kisin--Fargues $G_K$-module $\gMt$ which is of height at
  most~$h$, which admits all descents over~$L$, and which satisfies
  $M=W(\C^\flat)\otimes_{\Ainf}\gM$. Furthermore~$T(M)|_{G_L}$ is
  crystalline if and only if~$\gMt$ is crystalline as a
  Breuil--Kisin--Fargues $G_L$-module.
\end{acor}
\begin{proof}
  The sufficiency of the condition is immediate from
  Theorem~\ref{athm: admits all descents if and only if
    semistable}. For the necessity, by Theorem~\ref{athm: admits all
    descents if and only if semistable} there is a
  Breuil--Kisin--Fargues $G_L$-module $\gMt_L$ which admits all
  descents and satisfies $M=W(\C^\flat)\otimes_{\Ainf}\gM_L$, so we
  need only show that~$\gMt_L$ is $G_K$-stable. This follows from
  \cite[Lem.\ 2.9]{liulattice2}, or one can argue as
  follows: % use the following argument via \cite[Thm. 4.28]{2016arXiv160203148B}
  note that the proof of Theorem~\ref{athm: admits all descents if and
    only if semistable} shows that the Breuil--Kisin--Fargues module
  $\gMt_L : =\Ainf\otimes_{\gS_{\piflat_L}}\gM_L$ corresponds via
  \cite[Thm. 4.28]{2016arXiv160203148B} to the pair
  $(T(M), \BdR^+ \otimes_{L}D_\dR(T(M)|_{G_L}))$. This pair has a
  $G_K$-action, because $V(M)$ is de Rham, hence $\gM_L^{\inf}$ is
  $G_K$-stable by \cite[Thm. 4.28]{2016arXiv160203148B}.\end{proof}
  
  
%   \mar{TG:
%     hopefully it's actually true in the general case?! 
    
    
%      TL : Yes and here is another proof}
%    \end{proof}
\section{Hodge and inertial types}\label{subsec: Hodge and inertial
types}
We finally recall how to interpret Hodge--Tate weights and
inertial types in terms of Breuil--Kisin modules. Fix a finite
extension $E/\Qp$ with ring of integers~$\cO$, which is sufficiently
large that~$E$ contains the image of every embedding
$\sigma:K\into\Qpbar$. The definitions above admit obvious extensions
to the case of Breuil--Kisin--Fargues modules with $\cO$-coefficients
(see Definition~\ref{defn: descending BKF to BK coefficients}),
and since $\cO$ is a finite free~$\Zp$-module, the proofs of
Theorem~\ref{athm: admits all descents if and only if semistable} and
Corollary~\ref{acor: admits all descents if and only if potentially
  semistable} go
over unchanged in this setting.
% \mar{TG: need to do
  % this! Nothing much new should be required, just pointing out that
  % $D_{pst}=\overline{\gM}[1/p]$ as an $I_K$-representation (which I
  % think we show in the course of the proof of Theorem~\ref{athm:
  %   admits all descents if and only if semistable}, although perhaps
  % it should be extracted from there?), and that
  % we can read off Hodge--Tate weights from the filtration on $\overline{\gM}'$. 
  
  % TL: Do you want to write a little more precise statement? }

Suppose that~$\gMt$ is as in the statement of Corollary~\ref{acor:
  admits all descents if and only if potentially semistable}; so it is
a Breuil--Kisin--Fargues $G_K$-module of height at most~$h$, which
admits all descents over~$L$. Write~$l$ for the residue field of~$L$,
write~$\overline{\gM}$ for the $W(l)$-module
$\gM_{\piflat}/[\piflat]\gM_{\piflat}$ of Definition~\ref{adefn:
  descending BKF to BK}~\eqref{appendix item: M mod u descends}, and
$\overline{\gM}'$ for the $\cO_L$-module
$\varphi^*\gM_{\piflat}/E_{\piflat}\varphi^*\gM_{\piflat}$ of
Definition~\ref{adefn: descending BKF to BK}~\eqref{item: M mod E
  descends}. These modules are independent of the choice of~$\pi$
(which now denotes a uniformizer of~$L$) and of~$\piflat$, by definition.

The semilinear $G_K$-action on $\gMt$ induces semilinear actions
on~$\overline{\gM}$ and on~$\overline{\gM}'$, and as noted in the
proof of Theorem~\ref{athm: admits all descents if and only if
  semistable}, the action of $G_L$ on both modules is
trivial. Furthermore, by~\cite[Cor.\ 2.12]{liulattice2}, the inertial
  type $D_{\pst}(V(M))|_{I_K}$ is given by~$\overline{\gM}[1/p]$ with
  its action of~$I_{L/K}$. More precisely,  \cite[\S 2.3]{liulattice2} shows that 
  $W(\bark)[1/p] \otimes_{\Ainf}\varphi ^*\gMt$ is isomorphic to $D_{\pst}
  (T(M))\simeq W(\bark) \otimes_{W(l)}D_\st(T(M)|_{G_L})$ as $W(\bark)[1/p][G_K]$-modules. Furthermore, this isomorphism is compatible with the isomorphism (from the proof of Theorem \ref{athm: admits all descents if and only if semistable})
  \[\iota _{\piflat}: \xymatrix{\overline \gM [\frac 1 p] = D_{\piflat}
    \simeq s_{\piflat} (D_{\piflat}) & \ar[l]_-{i_{\piflat}}^-{\sim}
    D_\st(T(M)|_{G_L})}.\] Hence $\overline \gM [\frac 1 p]  \subset
  W(\bark)[1/p] \otimes_{\Ainf}\varphi ^*\gMt$ is endowed with an action of $I_{L/K}$ and is isomorphic to $D_{\pst} (V(M))$. 
  
  
  % \mar{TG: do we need to extend scalars to the
  % maximal unramified extension to know this?}


We now turn to the Hodge--Tate weights. We have a filtration on the $L\otimes_{\Qp}E$-module
$D_L=\overline{\gM'}[1/p]$, which is defined as follows. % \mar{TG: put a filtration on this using a
  % $\piflat$ - need to see that it is independent of $\piflat$, and
  % also that the semilinear $G_K$-action lets us descend it to $K$
  % from~$L$, hopefully this is just formal.}
For each~$\piflat$ we write
$\Phi_{\gM_{\piflat}}:\varphi^*\gM_{\piflat}\to\gM_{\piflat}$ and $f_\pi: \varphi^*\gM_{\piflat}[\frac 1 p] \twoheadrightarrow D_L $.  
For each~$ i\geq 0$ we define
$\Fil^i\varphi^*\gM_{\piflat}=\Phi_{\gM_{\piflat}}^{-1}(E_{\piflat}^i\gM_{\piflat})$ 
%We set $\Fil^i\varphi^*\gM_{\piflat}=\varphi^*\gM_{\piflat}$
%for~$i<0$.  
%Then for each $i$ we set
and $\Fil^i D_L = f_\pi (\Fil ^i \varphi ^* \gM_{\piflat})$. 
%$\Fil^iD_L: =(\Fil^i\varphi^*\gM_{\piflat}/E_{\piflat}\Fil^{i-1}\varphi^*\gM_{\piflat})[1/p]=(\Fil^i\varphi^*\gM_{\piflat}/(E_{\piflat}\varphi^*\gM_{\piflat}\cap\Fil^i\varphi^*\gM_{\piflat}))[1/p]$. 
Considering the isomorphism $\iota_{\piflat}: D_\st (V(M)|_{G_L}) \simeq D_{\piflat}$,  %via the isomorphism $i_{\piflat}: D_\st (T(M)|_{G_L}) \simeq s_{\piflat}(D_{\piflat})$, 
we obtain an isomorphism $D_\dR(V(M)|_{G_L}) \simeq L
\otimes_{W(l)[\frac 1 p]} D_\st (V(M)|_{G_L}) \simeq D_L$. By
\cite[Cor.\ 3.2.3, Thm.\ 3.4.1]{MR2388556}, this isomorphism respects
the  filtrations on each side. In particular, $\Fil^i D_L$ is independent of the choice of~$\piflat$. 

% \mar{TG: is it easy to show from something above that this is
%   independent of the choice of~$\piflat$?}
By Hilbert 90, $D_L$ and
its filtration descend to a filtration on a
$K\otimes_{\Qp}E$-module~$D_K\simeq D_\dR(V(M))$. % \mar{TG: I didn't check this! I presume
  % it's true, though} \mar{TL: Yes, because $D_L=\simeq D_\dR(T(M)|_{G_L}) $}
Then for each $\sigma:K\into E$, and
each~$i\in [0,h]$, the multiplicity of~$i$ in the multiset
$\HT_{\sigma}(V(M))$ is the dimension of the $i$-th graded piece of
$e_{\sigma}D_K$, where $e_\sigma\in (K\otimes_{\Qp}E)$ is the
idempotent corresponding to~$\sigma$.

  We summarise the preceding discussion in the following corollary.
\begin{acor}\label{acor: reading off inertial type and HT weights from BKF}
  In the setting of Corollary~\ref{acor: admits all descents if and
    only if potentially semistable}, the inertial type and Hodge type
  of~$V(M)$ are determined as follows: the inertial
  type~$D_{\pst}(V(M))|_{I_K}$ is given by~$\overline{\gM}[1/p]$ with
  its action of~$I_{L/K}$, while the Hodge type of~$V(M)$ is given by the
  jumps in the filtration on~$D_K$ described
  above.%\mar{TG: hopefully there's a cleaner way to phrase this part.}
\end{acor}
% \begin{arem}
%   \label{arem: explaining how to see Dpst from BKF GK module}In the
%   situation of Corollary~\ref{acor: admits all descents if and only if
%     potentially semistable}, note that for each choice of~$\piflat$
%   (where now~$\pi$ is a uniformiser of~$L$),
%   $\gM_{\piflat}/[\piflat]\gM_{\piflat}$\mar{TG: want to say something
%   about how this corresponds to~$D_{\pst}$, although I'm already
%   confused because it will have $W(l)$-coefficients rather
%   than~$W(k)$. On the other hand we probably only care about the
%   action of inertia, anyway, so probably this is OK?}
% \end{arem}

\emergencystretch=3em
\bibliographystyle{amsalpha}
\bibliography{universalBM}

\printindex
\end{document}

%%% Local Variables:
%%% mode: latex
%%% TeX-master: t
%%% End:
